% @see : https://coopmaths.fr/alea/?uuid=4c8c7&id=1AN11-3&n=1&d=10&s=2&alea=nawO&cd=1&uuid=0e984&id=can1F15&alea=1Alo&uuid=6f32d&id=can1F16&alea=dqSw&uuid=a1ba2&id=can1F14&alea=rkIp&uuid=3c690&id=can1F13&alea=EUDu&uuid=ebd89&id=1AN14-1&n=1&d=10&s=3&alea=ocYr&cd=0&uuid=ebd89&id=1AN14-1&n=1&d=10&s=4&alea=uAlP&cd=0&uuid=ebd89&id=1AN14-1&n=1&d=10&s=7&alea=vjN7&cd=0&uuid=a83c0&id=1AN14-4&n=1&d=10&s=1&alea=Njr3&cd=0&uuid=a83c0&id=1AN14-4&n=1&d=10&s=4&alea=4Vpz&cd=1
\documentclass[a4paper,11pt,fleqn]{article}
\usepackage{ProfMaquette}
\setKVdefault[Boulot]{CorrigeFin=true}
\usepackage{etoolbox}
\newbool{dys}
\setbool{dys}{false}          
\ifbool{dys}{
% POLICE DYS
\newcommand{\choiceFontsDys}[1]{
\ifstrequal{#1}{Fira}{%
  % Fira Sans + Fira Math
  % Description : Fira Sans est une police moderne, et Fira Math est une version mathématique compatible qui maintient le style sans-serif.
  % Utilisation : Fonctionne bien avec XeLaTeX et LuaLaTeX.
  \usepackage{fontspec}
  \setmainfont{Fira Sans}
  \setsansfont{Fira Sans}
  \usepackage{unicode-math}
  \setmathfont{Fira Math}
}{}
\ifstrequal{#1}{lmodern}{%
  % Latin Modern Sans
  % Description : Une version sans-serif de la célèbre police Latin Modern, qui est compatible avec mathastext.
  % Utilisation : Fonctionne avec pdfLaTeX, XeLaTeX et LuaLaTeX.
  \usepackage{lmodern}
  \renewcommand{\familydefault}{\sfdefault}
  \usepackage{mathastext}
}{}
\ifstrequal{#1}{tgheros}{%
  % TeX Gyre Heros + mathastext
  % Description : TeX Gyre Heros est une alternative à Helvetica, et le package mathastext permet d'utiliser la même police pour les mathématiques.
  % Utilisation : Fonctionne avec pdfLaTeX, XeLaTeX, ou LuaLaTeX.
  \usepackage{tgheros}
  \renewcommand{\familydefault}{\sfdefault}
  \usepackage{mathastext}
}{}
}
% Fira ou lmodern ou tgheros
\choiceFontsDys{Fira}

%\usepackage{unicode-math}
%\usepackage{fontspec}
%\setmainfont{TeX Gyre Schola}
%\setsansfont{TeX Gyre Schola}
%\setmathfont{TeX Gyre Schola Math}
%\setmainfont{OpenDyslexic}[Scale=1.0]
%\setsansfont{OpenDyslexic}[Scale=1.0]
%\setmathfont{OpenDyslexic}[Scale=1.0,range=up/{Latin,latin,num}]
\usepackage[fontsize=14]{scrextend}
\usepackage{setspace}
\setstretch{1.7}
% Une valeur d'environ 1.2 à 1.7 est souvent recommandée. Cela permet d'augmenter l'espace entre les lignes tout en offrant un léger espacement entre les mots.
\setlength{\spaceskip}{1.2em}  
% Une valeur d'environ 1.2em à 1.5em est couramment conseillée. Cela crée un espace plus ample entre les mots, ce qui peut aider à réduire la fatigue visuelle et à améliorer la fluidité de la lecture.
}{
% POLICE STANDARD
\usepackage{fontenc}
\usepackage[scaled=1]{helvet}
\usepackage[fontsize=12]{scrextend}
}
\usepackage[left=1.5cm,right=1.5cm,top=2cm,bottom=2cm]{geometry}
\usepackage[luatex]{hyperref}
\usepackage{tikz}
\usetikzlibrary{calc}
\usepackage{fancyhdr}
\pagestyle{fancy}
\renewcommand\headrulewidth{0pt}
\setlength{\headheight}{18pt}
\fancyhead[R]{\href{https://coopmaths.fr/alea}{Mathaléa}}
\fancyfoot[C]{\thepage}
\fancyfoot[R]{%
\begin{tikzpicture}[remember picture,overlay]
  \node[anchor=south east] at ($(current page.south east)+(-2,0.25cm)$) {\scriptsize {\bfseries \href{https://coopmaths.fr/}{Coopmaths.fr} -- \href{http://creativecommons.fr/licences/}{CC-BY-SA}}};
\end{tikzpicture}
}
\fancyhead[L]{
\begin{tikzpicture}[y=0.8, x=0.8, yscale=-0.04, xscale=0.04,remember picture, overlay,fill=orange!50,transform canvas={xshift=-1cm,yshift=1cm}]
%%%% Arc supérieur gauche%%%%
\path[fill](523,1424)..controls(474,1413)and(404,1372)..(362,1333)..controls(322,1295)and(313,1272)..(331,1254)..controls(348,1236)and(369,1245)..(410,1283)..controls(458,1328)and(517,1356)..(575,1362)..controls(635,1368)and(646,1375)..(643,1404)..controls(641,1428)and(641,1428)..(596,1430)..controls(571,1431)and(538,1428)..(523,1424)--cycle;
%%%% Dé face supérieur%%%%
\path[fill](512,1272)..controls(490,1260)and(195,878)..(195,861)..controls(195,854)and(198,846)..(202,843)..controls(210,838)and(677,772)..(707,772)..controls(720,772)and(737,781)..(753,796)..controls(792,833)and(1057,1179)..(1057,1193)..controls(1057,1200)and(1053,1209)..(1048,1212)..controls(1038,1220)and(590,1283)..(551,1282)..controls(539,1282)and(521,1278)..(512,1272)--cycle;
%%%% Dé faces gauche et droite%%%%
\path[fill](1061,1167)..controls(1050,1158)and(978,1068)..(900,967)..controls(792,829)and(756,777)..(753,756)--(748,729)--(724,745)..controls(704,759)and(660,767)..(456,794)..controls(322,813)and(207,825)..(200,822)..controls(193,820)and(187,812)..(187,804)..controls(188,797)and(229,688)..(279,563)..controls(349,390)and(376,331)..(391,320)..controls(406,309)and(462,299)..(649,273)..controls(780,254)and(897,240)..(907,241)..controls(918,243)and(927,249)..(928,256)..controls(930,264)and(912,315)..(889,372)..controls(866,429)and(848,476)..(849,477)..controls(851,479)and(872,432)..(897,373)..controls(936,276)and(942,266)..(960,266)..controls(975,266)and(999,292)..(1089,408)..controls(1281,654)and(1290,666)..(1290,691)..controls(1290,720)and(1104,1175)..(1090,1180)..controls(1085,1182)and (1071,1176)..(1061,1167)--cycle;
%%%% Arc inférieur bas%%%%
\path[fill](1329,861)..controls(1316,848)and(1317,844)..(1339,788)..controls(1364,726)and(1367,654)..(1347,591)..controls(1330,539)and(1338,522)..(1375,526)..controls(1395,528)and(1400,533)..(1412,566)..controls(1432,624)and(1426,760)..(1401,821)..controls(1386,861)and(1380,866)..(1361,868)..controls(1348,870)and(1334,866)..(1329,861)--cycle;
%%%% Arc inférieur gauche%%%%
\path[fill](196,373)..controls(181,358)and(186,335)..(213,294)..controls(252,237)and(304,190)..(363,161)..controls(435,124)and(472,127)..(472,170)..controls(472,183)and(462,192)..(414,213)..controls(350,243)and(303,283)..(264,343)..controls(239,383)and(216,393)..(196,373)--cycle;
\end{tikzpicture}
}
%%%%%% Style Fiche
\tcbset{%
  userfiche/.style={%
    %move upwards=-1cm,colback=red!75%
    top=5pt, left=5pt, right=5pt, colback=red!5!white%
  }%
}%
\tcbset{%
  userfichecor/.style={%
    %spread upwards=-1cm,colback=gray!5%
    top=5pt, left=5pt, right=5pt, colback=red!5!white%
  }%
}%

% Parametrages
\hypersetup{
    colorlinks=true,% On active la couleur pour les liens. Couleur par défaut rouge
    linkcolor=blue,% On définit la couleur pour les liens internes
    % filecolor=magenta,% On définit la couleur pour les liens vers les fichiers locaux      
    urlcolor=blue,% On définit la couleur pour les liens vers des sites web
    % pdftitle={Puissance Quatre},% On définit un titre pour le document pdf
    % pdfpagemode=FullScreen,% On fixe l'affichage par défaut à plein écran
}
\usepackage{qrcode}
\usepackage{mathrsfs}
\usepackage{enumitem}
\usepackage[french]{babel}
\setlength{\parindent}{0cm}
% loadPackagesFromContent
\usepackage{multicol}
\usepackage[table,svgnames]{xcolor}
\usepackage{needspace}
\usepackage{amsfonts}
\newcommand\R{\mathbb{R}}
\newcommand{\e}{\text{e}}
\usepackage{amssymb}

\newcount\tmpnum
\def\storedata#1#2{\tmpnum=0 \edef\tmp{\string#1}\storedataA#2\end}
\def\storedataA#1{\advance\tmpnum by1
   \ifx\end#1\else
      \expandafter\def\csname data:\tmp:\the\tmpnum\endcsname{#1}%
      \expandafter\storedataA\fi
}
\def\getdata[#1]#2{\csname data:\string#2:#1\endcsname}

\storedata\mydata{
	{ABADA Trixie}
	{AJNABER Oumnia}
	{AKASSA Moctar}
	{ASAR Gamal}
	{BEGHACHI Manelle}
	{BENHADRIA Naëlle}
	{BENMECHICHE Rostom Amine}
	{BIDAOUI Selma}
	{BILE Cyrelle}
	{BILLA Fatima}
	{BOU Corentin}
	{BOUCHLAGHMI Maryam}
	{CAMARA Hawa}
	{DIAWARA Dieynaba}
	{DZIUBA Kornélia}
	{ELONDA Christéla}
	{FAZILLEAU--SOULAT Emma}
	{FEUILLARD-LE QUILLEUC Enora}
	{FOFANA Rokiatou}
	{FRAZIER Lina}
	{GHAZOUANI Assia}
	{HAMADI Rilès}
	{HANG Sarah}
	{HASSINI Chayma}
	{HYDARA Amynata}
	{KEFI Abdelkader}
	{KHODIYA Saïma}
	{LE DRAPPIER Clovis}
	{LEIJOUR Sarah}
	{LETTO Marina}
	{LIN Alex}
	{LIN Elisa}
	{LLOYD Isham}
	{LOLLIER Roman}
	{MOKRANE Naïm}
	{MOYEN Alexandre}
	{QUENAULT Yanis}
	{RAMAUGÉ Madison}
	{SANG DAI Alice}
	{SARAVANAN Ahkhas}
	{SHAIKH Sana}
	{SRIMANIVANNAN Tharany}
	{TAHIROU Ismaël}
	{TOUNSI Hillal}
	{YANG Léo}
	{ZHENG Elodie}
}


\begin{document}

\begin{itemize}
	\item ABADA Trixie: \pageref{eleve:1}
	\item AJNABER Oumnia: \pageref{eleve:2}
	\item AKASSA Moctar: \pageref{eleve:3}
	\item ASAR Gamal: \pageref{eleve:4}
	\item BEGHACHI Manelle: \pageref{eleve:5}
	\item BENHADRIA Naëlle: \pageref{eleve:6}
	\item BENMECHICHE Rostom Amine: \pageref{eleve:7}
	\item BIDAOUI Selma: \pageref{eleve:8}
	\item BILE Cyrelle: \pageref{eleve:9}
	\item BILLA Fatima: \pageref{eleve:10}
	\item BOU Corentin: \pageref{eleve:11}
	\item BOUCHLAGHMI Maryam: \pageref{eleve:12}
	\item CAMARA Hawa: \pageref{eleve:13}
	\item DIAWARA Dieynaba: \pageref{eleve:14}
	\item DZIUBA Kornélia: \pageref{eleve:15}
	\item ELONDA Christéla: \pageref{eleve:16}
	\item FAZILLEAU--SOULAT Emma: \pageref{eleve:17}
	\item FEUILLARD-LE QUILLEUC Enora: \pageref{eleve:18}
	\item FOFANA Rokiatou: \pageref{eleve:19}
	\item FRAZIER Lina: \pageref{eleve:20}
	\item GHAZOUANI Assia: \pageref{eleve:21}
	\item HAMADI Rilès: \pageref{eleve:22}
	\item HANG Sarah: \pageref{eleve:23}
	\item HASSINI Chayma: \pageref{eleve:24}
	\item HYDARA Amynata: \pageref{eleve:25}
	\item KEFI Abdelkader: \pageref{eleve:26}
	\item KHODIYA Saïma: \pageref{eleve:27}
	\item LE DRAPPIER Clovis: \pageref{eleve:28}
	\item LEIJOUR Sarah: \pageref{eleve:29}
	\item LETTO Marina: \pageref{eleve:30}
	\item LIN Alex: \pageref{eleve:31}
	\item LIN Elisa: \pageref{eleve:32}
	\item LLOYD Isham: \pageref{eleve:33}
	\item LOLLIER Roman: \pageref{eleve:34}
	\item MOKRANE Naïm: \pageref{eleve:35}
	\item MOYEN Alexandre: \pageref{eleve:36}
	\item QUENAULT Yanis: \pageref{eleve:37}
	\item RAMAUGÉ Madison: \pageref{eleve:38}
	\item SANG DAI Alice: \pageref{eleve:39}
	\item SARAVANAN Ahkhas: \pageref{eleve:40}
	\item SHAIKH Sana: \pageref{eleve:41}
	\item SRIMANIVANNAN Tharany: \pageref{eleve:42}
	\item TAHIROU Ismaël: \pageref{eleve:43}
	\item TOUNSI Hillal: \pageref{eleve:44}
	\item YANG Léo: \pageref{eleve:45}
	\item ZHENG Elodie: \pageref{eleve:46}
\end{itemize}

\newpage
\begin{Maquette}[DM]{Niveau= ,Classe=\getdata[1]\mydata\ (v1), Date = 06/01/25  ,Theme=Exercices}


% @see : https://coopmaths.fr/alea?uuid=4c8c7&id=1AN11-3&n=1&d=10&s=2&alea=nawO22323232323423456789101112131415161718192021222324252627282930313233343536373839404142434445464748&cd=1&cols=2
\needspace{10\baselineskip}
\begin{exercice}[BaremeTotal=false]
\label{eleve:1}


 \begin{minipage}[t]{\linewidth} Soit $f$ une fonction dérivable sur $[-5;5]$ et $\mathcal{C}_f$ sa courbe représentative.\\On sait que  $f(-3)=4~~$ et que $~~f'(-3)=-3$.\\Déterminer une équation de la tangente $(T)$ à la courbe $\mathcal{C}_f$ au point d'abscisse $-3$,\\sans utiliser la formule de cours de l'équation de tangente. \end{minipage}

\end{exercice}

\begin{Solution}
 \begin{minipage}[t]{\linewidth} On sait que la tangente n'est pas une droite verticale, puisque la fonction est dérivable sur l'intervalle.\\On en déduit que la tangente $(T)$ au point d'abscisse $-3$, admet une équation réduite de la forme :  $(T) : y=m x + p$. 

\medskip
$\bullet$ Détermination de $m$ :\\On sait que le nombre dérivé en $-3$ est par définition, le coefficient directeur de la tangente au point d'abscisse $-3$.\\Par conséquent, on a déjà : $m=f'(-3)=-3$.\\ On en déduit que  $(T) : y= -3 x + p$\\$\bullet$ Détermination de $p$ :\\ Pour cela, on utilise que si $f(-3)=4~~$, alors le point $A$ de coordonnées $(-3;4)$ appartient à $\mathcal{C}_f$ mais aussi à $(T)$.\\ On peut écrire $A(-3;4) = \mathcal{C}_f \cap (T)$.\\ On remplace alors les coordonnées de $A(-3;4)$ dans l'équation  $(T) : y= -3 x + p$.\\ $\begin{aligned} \phantom{\iff}&A(-3;4)\in (T)\\ \iff& 4= -3 \times (-3)  + p\\ \iff& p=4 +3 \times (-3)  \\ \iff& p=4 -9   \\ \iff& p=-5\\\end{aligned}$\\On peut conclure que : $(T) : {\color[HTML]{f15929}\boldsymbol{y=-3x-5}}$. \end{minipage}

\end{Solution}

% @see : https://coopmaths.fr/alea?uuid=0e984&id=can1F15&n=1&d=10&alea=1Alo22323232323423456789101112131415161718192021222324252627282930313233343536373839404142434445464748&cd=1&cols=1
\needspace{10\baselineskip}
\begin{exercice}[BaremeTotal=false]


 La courbe représente une fonction $f$ et la droite est la tangente au point d'abscisse $2$.\\
        Déterminer $f'(2)$.\begin{center}\begin{tikzpicture}[baseline,scale = 0.5]

    \tikzset{
      point/.style={
        thick,
        draw,
        cross out,
        inner sep=0pt,
        minimum width=5pt,
        minimum height=5pt,
      },
    }
    \clip (-2,-2) rectangle (14,12);
    	
	\draw[color ={black},line width = 2,>=latex,->] (-2,0)--(14,0);
	\draw[color ={black},line width = 2,>=latex,->] (0,-2)--(0,12);
	\draw[color ={black},opacity = 0.5] (-2,1)--(14,1);
	\draw[color ={black},opacity = 0.5] (-2,2)--(14,2);
	\draw[color ={black},opacity = 0.5] (-2,3)--(14,3);
	\draw[color ={black},opacity = 0.5] (-2,4)--(14,4);
	\draw[color ={black},opacity = 0.5] (-2,5)--(14,5);
	\draw[color ={black},opacity = 0.5] (-2,6)--(14,6);
	\draw[color ={black},opacity = 0.5] (-2,7)--(14,7);
	\draw[color ={black},opacity = 0.5] (-2,8)--(14,8);
	\draw[color ={black},opacity = 0.5] (-2,9)--(14,9);
	\draw[color ={black},opacity = 0.5] (-2,10)--(14,10);
	\draw[color ={black},opacity = 0.5] (-2,11)--(14,11);
	\draw[color ={black},opacity = 0.5] (-2,12)--(14,12);
	\draw[color ={black},opacity = 0.5] (-2,-1)--(14,-1);
	\draw[color ={black},opacity = 0.5] (-2,-2)--(14,-2);
	\draw[color ={black},opacity = 0.5] (1,-2)--(1,12);
	\draw[color ={black},opacity = 0.5] (2,-2)--(2,12);
	\draw[color ={black},opacity = 0.5] (3,-2)--(3,12);
	\draw[color ={black},opacity = 0.5] (4,-2)--(4,12);
	\draw[color ={black},opacity = 0.5] (5,-2)--(5,12);
	\draw[color ={black},opacity = 0.5] (6,-2)--(6,12);
	\draw[color ={black},opacity = 0.5] (7,-2)--(7,12);
	\draw[color ={black},opacity = 0.5] (8,-2)--(8,12);
	\draw[color ={black},opacity = 0.5] (9,-2)--(9,12);
	\draw[color ={black},opacity = 0.5] (10,-2)--(10,12);
	\draw[color ={black},opacity = 0.5] (11,-2)--(11,12);
	\draw[color ={black},opacity = 0.5] (12,-2)--(12,12);
	\draw[color ={black},opacity = 0.5] (13,-2)--(13,12);
	\draw[color ={black},opacity = 0.5] (14,-2)--(14,12);
	\draw[color ={black},opacity = 0.5] (-1,-2)--(-1,12);
	\draw[color ={black},opacity = 0.5] (-2,-2)--(-2,12);
	\draw[color ={gray},opacity = 0.3] (-2,0)--(14,0);
	\draw[color ={gray},opacity = 0.3] (-2,0.5)--(14,0.5);
	\draw[color ={gray},opacity = 0.3] (-2,1.5)--(14,1.5);
	\draw[color ={gray},opacity = 0.3] (-2,2.5)--(14,2.5);
	\draw[color ={gray},opacity = 0.3] (-2,3.5)--(14,3.5);
	\draw[color ={gray},opacity = 0.3] (-2,4.5)--(14,4.5);
	\draw[color ={gray},opacity = 0.3] (-2,5.5)--(14,5.5);
	\draw[color ={gray},opacity = 0.3] (-2,6.5)--(14,6.5);
	\draw[color ={gray},opacity = 0.3] (-2,7.5)--(14,7.5);
	\draw[color ={gray},opacity = 0.3] (-2,8.5)--(14,8.5);
	\draw[color ={gray},opacity = 0.3] (-2,9.5)--(14,9.5);
	\draw[color ={gray},opacity = 0.3] (-2,10.5)--(14,10.5);
	\draw[color ={gray},opacity = 0.3] (-2,11.5)--(14,11.5);
	\draw[color ={gray},opacity = 0.3] (-2,0)--(14,0);
	\draw[color ={gray},opacity = 0.3] (-2,-0.5)--(14,-0.5);
	\draw[color ={gray},opacity = 0.3] (-2,-1.5)--(14,-1.5);
	\draw[color ={gray},opacity = 0.3] (0,-2)--(0,12);
	\draw[color ={gray},opacity = 0.3] (0.1,-2)--(0.1,12);
	\draw[color ={gray},opacity = 0.3] (0.2,-2)--(0.2,12);
	\draw[color ={gray},opacity = 0.3] (0.30000000000000004,-2)--(0.30000000000000004,12);
	\draw[color ={gray},opacity = 0.3] (0.4,-2)--(0.4,12);
	\draw[color ={gray},opacity = 0.3] (0.5,-2)--(0.5,12);
	\draw[color ={gray},opacity = 0.3] (0.6,-2)--(0.6,12);
	\draw[color ={gray},opacity = 0.3] (0.7,-2)--(0.7,12);
	\draw[color ={gray},opacity = 0.3] (0.7999999999999999,-2)--(0.7999999999999999,12);
	\draw[color ={gray},opacity = 0.3] (0.8999999999999999,-2)--(0.8999999999999999,12);
	\draw[color ={gray},opacity = 0.3] (1.0999999999999999,-2)--(1.0999999999999999,12);
	\draw[color ={gray},opacity = 0.3] (1.2,-2)--(1.2,12);
	\draw[color ={gray},opacity = 0.3] (1.3,-2)--(1.3,12);
	\draw[color ={gray},opacity = 0.3] (1.4000000000000001,-2)--(1.4000000000000001,12);
	\draw[color ={gray},opacity = 0.3] (1.5000000000000002,-2)--(1.5000000000000002,12);
	\draw[color ={gray},opacity = 0.3] (1.6000000000000003,-2)--(1.6000000000000003,12);
	\draw[color ={gray},opacity = 0.3] (1.7000000000000004,-2)--(1.7000000000000004,12);
	\draw[color ={gray},opacity = 0.3] (1.8000000000000005,-2)--(1.8000000000000005,12);
	\draw[color ={gray},opacity = 0.3] (1.9000000000000006,-2)--(1.9000000000000006,12);
	\draw[color ={gray},opacity = 0.3] (2.1000000000000005,-2)--(2.1000000000000005,12);
	\draw[color ={gray},opacity = 0.3] (2.2000000000000006,-2)--(2.2000000000000006,12);
	\draw[color ={gray},opacity = 0.3] (2.3000000000000007,-2)--(2.3000000000000007,12);
	\draw[color ={gray},opacity = 0.3] (2.400000000000001,-2)--(2.400000000000001,12);
	\draw[color ={gray},opacity = 0.3] (2.500000000000001,-2)--(2.500000000000001,12);
	\draw[color ={gray},opacity = 0.3] (2.600000000000001,-2)--(2.600000000000001,12);
	\draw[color ={gray},opacity = 0.3] (2.700000000000001,-2)--(2.700000000000001,12);
	\draw[color ={gray},opacity = 0.3] (2.800000000000001,-2)--(2.800000000000001,12);
	\draw[color ={gray},opacity = 0.3] (2.9000000000000012,-2)--(2.9000000000000012,12);
	\draw[color ={gray},opacity = 0.3] (3.1000000000000014,-2)--(3.1000000000000014,12);
	\draw[color ={gray},opacity = 0.3] (3.2000000000000015,-2)--(3.2000000000000015,12);
	\draw[color ={gray},opacity = 0.3] (3.3000000000000016,-2)--(3.3000000000000016,12);
	\draw[color ={gray},opacity = 0.3] (3.4000000000000017,-2)--(3.4000000000000017,12);
	\draw[color ={gray},opacity = 0.3] (3.5000000000000018,-2)--(3.5000000000000018,12);
	\draw[color ={gray},opacity = 0.3] (3.600000000000002,-2)--(3.600000000000002,12);
	\draw[color ={gray},opacity = 0.3] (3.700000000000002,-2)--(3.700000000000002,12);
	\draw[color ={gray},opacity = 0.3] (3.800000000000002,-2)--(3.800000000000002,12);
	\draw[color ={gray},opacity = 0.3] (3.900000000000002,-2)--(3.900000000000002,12);
	\draw[color ={gray},opacity = 0.3] (4.100000000000001,-2)--(4.100000000000001,12);
	\draw[color ={gray},opacity = 0.3] (4.200000000000001,-2)--(4.200000000000001,12);
	\draw[color ={gray},opacity = 0.3] (4.300000000000001,-2)--(4.300000000000001,12);
	\draw[color ={gray},opacity = 0.3] (4.4,-2)--(4.4,12);
	\draw[color ={gray},opacity = 0.3] (4.5,-2)--(4.5,12);
	\draw[color ={gray},opacity = 0.3] (4.6,-2)--(4.6,12);
	\draw[color ={gray},opacity = 0.3] (4.699999999999999,-2)--(4.699999999999999,12);
	\draw[color ={gray},opacity = 0.3] (4.799999999999999,-2)--(4.799999999999999,12);
	\draw[color ={gray},opacity = 0.3] (4.899999999999999,-2)--(4.899999999999999,12);
	\draw[color ={gray},opacity = 0.3] (5.099999999999998,-2)--(5.099999999999998,12);
	\draw[color ={gray},opacity = 0.3] (5.1999999999999975,-2)--(5.1999999999999975,12);
	\draw[color ={gray},opacity = 0.3] (5.299999999999997,-2)--(5.299999999999997,12);
	\draw[color ={gray},opacity = 0.3] (5.399999999999997,-2)--(5.399999999999997,12);
	\draw[color ={gray},opacity = 0.3] (5.4999999999999964,-2)--(5.4999999999999964,12);
	\draw[color ={gray},opacity = 0.3] (5.599999999999996,-2)--(5.599999999999996,12);
	\draw[color ={gray},opacity = 0.3] (5.699999999999996,-2)--(5.699999999999996,12);
	\draw[color ={gray},opacity = 0.3] (5.799999999999995,-2)--(5.799999999999995,12);
	\draw[color ={gray},opacity = 0.3] (5.899999999999995,-2)--(5.899999999999995,12);
	\draw[color ={gray},opacity = 0.3] (6.099999999999994,-2)--(6.099999999999994,12);
	\draw[color ={gray},opacity = 0.3] (6.199999999999994,-2)--(6.199999999999994,12);
	\draw[color ={gray},opacity = 0.3] (6.299999999999994,-2)--(6.299999999999994,12);
	\draw[color ={gray},opacity = 0.3] (6.399999999999993,-2)--(6.399999999999993,12);
	\draw[color ={gray},opacity = 0.3] (6.499999999999993,-2)--(6.499999999999993,12);
	\draw[color ={gray},opacity = 0.3] (6.5999999999999925,-2)--(6.5999999999999925,12);
	\draw[color ={gray},opacity = 0.3] (6.699999999999992,-2)--(6.699999999999992,12);
	\draw[color ={gray},opacity = 0.3] (6.799999999999992,-2)--(6.799999999999992,12);
	\draw[color ={gray},opacity = 0.3] (6.8999999999999915,-2)--(6.8999999999999915,12);
	\draw[color ={gray},opacity = 0.3] (7.099999999999991,-2)--(7.099999999999991,12);
	\draw[color ={gray},opacity = 0.3] (7.19999999999999,-2)--(7.19999999999999,12);
	\draw[color ={gray},opacity = 0.3] (7.29999999999999,-2)--(7.29999999999999,12);
	\draw[color ={gray},opacity = 0.3] (7.39999999999999,-2)--(7.39999999999999,12);
	\draw[color ={gray},opacity = 0.3] (7.499999999999989,-2)--(7.499999999999989,12);
	\draw[color ={gray},opacity = 0.3] (7.599999999999989,-2)--(7.599999999999989,12);
	\draw[color ={gray},opacity = 0.3] (7.699999999999989,-2)--(7.699999999999989,12);
	\draw[color ={gray},opacity = 0.3] (7.799999999999988,-2)--(7.799999999999988,12);
	\draw[color ={gray},opacity = 0.3] (7.899999999999988,-2)--(7.899999999999988,12);
	\draw[color ={gray},opacity = 0.3] (8.099999999999987,-2)--(8.099999999999987,12);
	\draw[color ={gray},opacity = 0.3] (8.199999999999987,-2)--(8.199999999999987,12);
	\draw[color ={gray},opacity = 0.3] (8.299999999999986,-2)--(8.299999999999986,12);
	\draw[color ={gray},opacity = 0.3] (8.399999999999986,-2)--(8.399999999999986,12);
	\draw[color ={gray},opacity = 0.3] (8.499999999999986,-2)--(8.499999999999986,12);
	\draw[color ={gray},opacity = 0.3] (8.599999999999985,-2)--(8.599999999999985,12);
	\draw[color ={gray},opacity = 0.3] (8.699999999999985,-2)--(8.699999999999985,12);
	\draw[color ={gray},opacity = 0.3] (8.799999999999985,-2)--(8.799999999999985,12);
	\draw[color ={gray},opacity = 0.3] (8.899999999999984,-2)--(8.899999999999984,12);
	\draw[color ={gray},opacity = 0.3] (9.099999999999984,-2)--(9.099999999999984,12);
	\draw[color ={gray},opacity = 0.3] (9.199999999999983,-2)--(9.199999999999983,12);
	\draw[color ={gray},opacity = 0.3] (9.299999999999983,-2)--(9.299999999999983,12);
	\draw[color ={gray},opacity = 0.3] (9.399999999999983,-2)--(9.399999999999983,12);
	\draw[color ={gray},opacity = 0.3] (9.499999999999982,-2)--(9.499999999999982,12);
	\draw[color ={gray},opacity = 0.3] (9.599999999999982,-2)--(9.599999999999982,12);
	\draw[color ={gray},opacity = 0.3] (9.699999999999982,-2)--(9.699999999999982,12);
	\draw[color ={gray},opacity = 0.3] (9.799999999999981,-2)--(9.799999999999981,12);
	\draw[color ={gray},opacity = 0.3] (9.89999999999998,-2)--(9.89999999999998,12);
	\draw[color ={gray},opacity = 0.3] (10.09999999999998,-2)--(10.09999999999998,12);
	\draw[color ={gray},opacity = 0.3] (10.19999999999998,-2)--(10.19999999999998,12);
	\draw[color ={gray},opacity = 0.3] (10.29999999999998,-2)--(10.29999999999998,12);
	\draw[color ={gray},opacity = 0.3] (10.399999999999979,-2)--(10.399999999999979,12);
	\draw[color ={gray},opacity = 0.3] (10.499999999999979,-2)--(10.499999999999979,12);
	\draw[color ={gray},opacity = 0.3] (10.599999999999978,-2)--(10.599999999999978,12);
	\draw[color ={gray},opacity = 0.3] (10.699999999999978,-2)--(10.699999999999978,12);
	\draw[color ={gray},opacity = 0.3] (10.799999999999978,-2)--(10.799999999999978,12);
	\draw[color ={gray},opacity = 0.3] (10.899999999999977,-2)--(10.899999999999977,12);
	\draw[color ={gray},opacity = 0.3] (11.099999999999977,-2)--(11.099999999999977,12);
	\draw[color ={gray},opacity = 0.3] (11.199999999999976,-2)--(11.199999999999976,12);
	\draw[color ={gray},opacity = 0.3] (11.299999999999976,-2)--(11.299999999999976,12);
	\draw[color ={gray},opacity = 0.3] (11.399999999999975,-2)--(11.399999999999975,12);
	\draw[color ={gray},opacity = 0.3] (11.499999999999975,-2)--(11.499999999999975,12);
	\draw[color ={gray},opacity = 0.3] (11.599999999999975,-2)--(11.599999999999975,12);
	\draw[color ={gray},opacity = 0.3] (11.699999999999974,-2)--(11.699999999999974,12);
	\draw[color ={gray},opacity = 0.3] (11.799999999999974,-2)--(11.799999999999974,12);
	\draw[color ={gray},opacity = 0.3] (11.899999999999974,-2)--(11.899999999999974,12);
	\draw[color ={gray},opacity = 0.3] (12.099999999999973,-2)--(12.099999999999973,12);
	\draw[color ={gray},opacity = 0.3] (12.199999999999973,-2)--(12.199999999999973,12);
	\draw[color ={gray},opacity = 0.3] (12.299999999999972,-2)--(12.299999999999972,12);
	\draw[color ={gray},opacity = 0.3] (12.399999999999972,-2)--(12.399999999999972,12);
	\draw[color ={gray},opacity = 0.3] (12.499999999999972,-2)--(12.499999999999972,12);
	\draw[color ={gray},opacity = 0.3] (12.599999999999971,-2)--(12.599999999999971,12);
	\draw[color ={gray},opacity = 0.3] (12.69999999999997,-2)--(12.69999999999997,12);
	\draw[color ={gray},opacity = 0.3] (12.79999999999997,-2)--(12.79999999999997,12);
	\draw[color ={gray},opacity = 0.3] (12.89999999999997,-2)--(12.89999999999997,12);
	\draw[color ={gray},opacity = 0.3] (13.09999999999997,-2)--(13.09999999999997,12);
	\draw[color ={gray},opacity = 0.3] (13.199999999999969,-2)--(13.199999999999969,12);
	\draw[color ={gray},opacity = 0.3] (13.299999999999969,-2)--(13.299999999999969,12);
	\draw[color ={gray},opacity = 0.3] (13.399999999999968,-2)--(13.399999999999968,12);
	\draw[color ={gray},opacity = 0.3] (13.499999999999968,-2)--(13.499999999999968,12);
	\draw[color ={gray},opacity = 0.3] (13.599999999999968,-2)--(13.599999999999968,12);
	\draw[color ={gray},opacity = 0.3] (13.699999999999967,-2)--(13.699999999999967,12);
	\draw[color ={gray},opacity = 0.3] (13.799999999999967,-2)--(13.799999999999967,12);
	\draw[color ={gray},opacity = 0.3] (13.899999999999967,-2)--(13.899999999999967,12);
	\draw[color ={gray},opacity = 0.3] (0,-2)--(0,12);
	\draw[color ={gray},opacity = 0.3] (-0.1,-2)--(-0.1,12);
	\draw[color ={gray},opacity = 0.3] (-0.2,-2)--(-0.2,12);
	\draw[color ={gray},opacity = 0.3] (-0.30000000000000004,-2)--(-0.30000000000000004,12);
	\draw[color ={gray},opacity = 0.3] (-0.4,-2)--(-0.4,12);
	\draw[color ={gray},opacity = 0.3] (-0.5,-2)--(-0.5,12);
	\draw[color ={gray},opacity = 0.3] (-0.6,-2)--(-0.6,12);
	\draw[color ={gray},opacity = 0.3] (-0.7,-2)--(-0.7,12);
	\draw[color ={gray},opacity = 0.3] (-0.7999999999999999,-2)--(-0.7999999999999999,12);
	\draw[color ={gray},opacity = 0.3] (-0.8999999999999999,-2)--(-0.8999999999999999,12);
	\draw[color ={gray},opacity = 0.3] (-1.0999999999999999,-2)--(-1.0999999999999999,12);
	\draw[color ={gray},opacity = 0.3] (-1.2,-2)--(-1.2,12);
	\draw[color ={gray},opacity = 0.3] (-1.3,-2)--(-1.3,12);
	\draw[color ={gray},opacity = 0.3] (-1.4000000000000001,-2)--(-1.4000000000000001,12);
	\draw[color ={gray},opacity = 0.3] (-1.5000000000000002,-2)--(-1.5000000000000002,12);
	\draw[color ={gray},opacity = 0.3] (-1.6000000000000003,-2)--(-1.6000000000000003,12);
	\draw[color ={gray},opacity = 0.3] (-1.7000000000000004,-2)--(-1.7000000000000004,12);
	\draw[color ={gray},opacity = 0.3] (-1.8000000000000005,-2)--(-1.8000000000000005,12);
	\draw[color ={gray},opacity = 0.3] (-1.9000000000000006,-2)--(-1.9000000000000006,12);
	\draw[color ={black},line width = 2] (2,-0.2)--(2,0.2);
	\draw[color ={black},line width = 2] (4,-0.2)--(4,0.2);
	\draw[color ={black},line width = 2] (6,-0.2)--(6,0.2);
	\draw[color ={black},line width = 2] (8,-0.2)--(8,0.2);
	\draw[color ={black},line width = 2] (10,-0.2)--(10,0.2);
	\draw[color ={black},line width = 2] (12,-0.2)--(12,0.2);
	\draw[color ={black},line width = 2] (-0.2,2)--(0.2,2);
	\draw[color ={black},line width = 2] (-0.2,4)--(0.2,4);
	\draw[color ={black},line width = 2] (-0.2,6)--(0.2,6);
	\draw[color ={black},line width = 2] (-0.2,8)--(0.2,8);
	\draw[color ={black},line width = 2] (-0.2,10)--(0.2,10);
	\draw (2,-0.4) node[anchor = center] {\scriptsize \color{black}{$1$}};
	\draw (4,-0.4) node[anchor = center] {\scriptsize \color{black}{$2$}};
	\draw (6,-0.4) node[anchor = center] {\scriptsize \color{black}{$3$}};
	\draw (8,-0.4) node[anchor = center] {\scriptsize \color{black}{$4$}};
	\draw (10,-0.4) node[anchor = center] {\scriptsize \color{black}{$5$}};
	\draw (-0.5,2.1) node[anchor = center] {\footnotesize \color{black}{$1$}};
	\draw (-0.5,4.1) node[anchor = center] {\footnotesize \color{black}{$2$}};
	\draw (-0.5,6.1) node[anchor = center] {\footnotesize \color{black}{$3$}};
	\draw (-0.5,8.1) node[anchor = center] {\footnotesize \color{black}{$4$}};
	\draw (-0.5,10.1) node[anchor = center] {\footnotesize \color{black}{$5$}};
	\draw [color={black}] (-0.3,-0.3) node[anchor = center,scale=1, rotate = 0] {O};
	
	\draw[color={blue},line width = 2] (0.6,8.67)--(1,6)--(1.4,4.86)--(1.8,4.22)--(2.2,3.82)--(2.6,3.54)--(3,3.33)--(3.4,3.18)--(3.8,3.05)--(4.2,2.95)--(4.6,2.87)--(5,2.8)--(5.4,2.74)--(5.8,2.69)--(6.2,2.65)--(6.6,2.61)--(7,2.57)--(7.4,2.54)--(7.8,2.51)--(8.2,2.49)--(8.6,2.47)--(9,2.44)--(9.4,2.43)--(9.8,2.41)--(10.2,2.39)--(10.6,2.38)--(11,2.36)--(11.4,2.35)--(11.8,2.34)--(12.2,2.33)--(12.6,2.32)--(13,2.31)--(13.4,2.3)--(13.8,2.29);
	
	\draw[color={red},line width = 2] (-2,4.5)--(-1.6,4.4)--(-1.2,4.3)--(-0.8,4.2)--(-0.4,4.1)--(0,4)--(0.4,3.9)--(0.8,3.8)--(1.2,3.7)--(1.6,3.6)--(2,3.5)--(2.4,3.4)--(2.8,3.3)--(3.2,3.2)--(3.6,3.1)--(4,3)--(4.4,2.9)--(4.8,2.8)--(5.2,2.7)--(5.6,2.6)--(6,2.5)--(6.4,2.4)--(6.8,2.3)--(7.2,2.2)--(7.6,2.1)--(8,2)--(8.4,1.9)--(8.8,1.8)--(9.2,1.7)--(9.6,1.6)--(10,1.5)--(10.4,1.4)--(10.8,1.3)--(11.2,1.2)--(11.6,1.1)--(12,1)--(12.4,0.9)--(12.8,0.8)--(13.2,0.7)--(13.6,0.6)--(14,0.5);

\end{tikzpicture}\end{center}
\end{exercice}

\begin{Solution}
 $f'(2)$ est donné par le coefficient directeur de la tangente à la courbe au point d'abscisse $2$, soit $\dfrac{-1}{4}=-\dfrac{1}{4}$.
\end{Solution}

% @see : https://coopmaths.fr/alea?uuid=6f32d&id=can1F16&n=1&d=10&alea=dqSw22323232323423456789101112131415161718192021222324252627282930313233343536373839404142434445464748&cd=1&cols=1
\needspace{10\baselineskip}
\newpage
\begin{exercice}[BaremeTotal=false]


 On donne les représentations graphiques d'une fonction et de sa dérivée.\\
        Donner l'équation réduite de la tangente à la courbe de $f$ en $x=0$. \\ \begin{minipage}{0.5\linewidth}
    \begin{tikzpicture}[baseline, scale=0.8]

    \tikzset{
      point/.style={
        thick,
        draw,
        cross out,
        inner sep=0pt,
        minimum width=5pt,
        minimum height=5pt,
      },
    }
    \clip (-4.5,-4.5) rectangle (5.75,13.5);
    	
	\draw[color ={black},>=latex,->] (-3,0)--(3,0);
	\draw[color ={black},>=latex,->] (0,-3)--(0,12);
	\draw[color ={black},opacity = 0.4] (-3,1)--(3,1);
	\draw[color ={black},opacity = 0.4] (-3,2)--(3,2);
	\draw[color ={black},opacity = 0.4] (-3,3)--(3,3);
	\draw[color ={black},opacity = 0.4] (-3,4)--(3,4);
	\draw[color ={black},opacity = 0.4] (-3,5)--(3,5);
	\draw[color ={black},opacity = 0.4] (-3,6)--(3,6);
	\draw[color ={black},opacity = 0.4] (-3,7)--(3,7);
	\draw[color ={black},opacity = 0.4] (-3,8)--(3,8);
	\draw[color ={black},opacity = 0.4] (-3,9)--(3,9);
	\draw[color ={black},opacity = 0.4] (-3,10)--(3,10);
	\draw[color ={black},opacity = 0.4] (-3,11)--(3,11);
	\draw[color ={black},opacity = 0.4] (-3,12)--(3,12);
	\draw[color ={black},opacity = 0.4] (-3,-1)--(3,-1);
	\draw[color ={black},opacity = 0.4] (-3,-2)--(3,-2);
	\draw[color ={black},opacity = 0.4] (-3,-3)--(3,-3);
	\draw[color ={black},opacity = 0.4] (1,-3)--(1,12);
	\draw[color ={black},opacity = 0.4] (2,-3)--(2,12);
	\draw[color ={black},opacity = 0.4] (3,-3)--(3,12);
	\draw[color ={black},opacity = 0.4] (-1,-3)--(-1,12);
	\draw[color ={black},opacity = 0.4] (-2,-3)--(-2,12);
	\draw[color ={black},opacity = 0.4] (-3,-3)--(-3,12);
	\draw[color ={gray},opacity = 0.3] (-3,0)--(3,0);
	\draw[color ={gray},opacity = 0.3] (-3,0.5)--(3,0.5);
	\draw[color ={gray},opacity = 0.3] (-3,1.5)--(3,1.5);
	\draw[color ={gray},opacity = 0.3] (-3,2.5)--(3,2.5);
	\draw[color ={gray},opacity = 0.3] (-3,3.5)--(3,3.5);
	\draw[color ={gray},opacity = 0.3] (-3,4.5)--(3,4.5);
	\draw[color ={gray},opacity = 0.3] (-3,5.5)--(3,5.5);
	\draw[color ={gray},opacity = 0.3] (-3,6.5)--(3,6.5);
	\draw[color ={gray},opacity = 0.3] (-3,7.5)--(3,7.5);
	\draw[color ={gray},opacity = 0.3] (-3,8.5)--(3,8.5);
	\draw[color ={gray},opacity = 0.3] (-3,9.5)--(3,9.5);
	\draw[color ={gray},opacity = 0.3] (-3,10.5)--(3,10.5);
	\draw[color ={gray},opacity = 0.3] (-3,11.5)--(3,11.5);
	\draw[color ={gray},opacity = 0.3] (-3,0)--(3,0);
	\draw[color ={gray},opacity = 0.3] (-3,-0.5)--(3,-0.5);
	\draw[color ={gray},opacity = 0.3] (-3,-1.5)--(3,-1.5);
	\draw[color ={gray},opacity = 0.3] (-3,-2.5)--(3,-2.5);
	\draw[color ={gray},opacity = 0.3] (0,-3)--(0,12);
	\draw[color ={gray},opacity = 0.3] (0,-3)--(0,12);
	\draw[color ={black},line width = 1.2] (1,-0.13)--(1,0.13);
	\draw[color ={black},line width = 1.2] (2,-0.13)--(2,0.13);
	\draw[color ={black},line width = 1.2] (3,-0.13)--(3,0.13);
	\draw[color ={black},line width = 1.2] (-1,-0.13)--(-1,0.13);
	\draw[color ={black},line width = 1.2] (-2,-0.13)--(-2,0.13);
	\draw[color ={black},line width = 1.2] (-3,-0.13)--(-3,0.13);
	\draw[color ={black},line width = 1.2] (-0.13,1)--(0.13,1);
	\draw[color ={black},line width = 1.2] (-0.13,2)--(0.13,2);
	\draw[color ={black},line width = 1.2] (-0.13,3)--(0.13,3);
	\draw[color ={black},line width = 1.2] (-0.13,4)--(0.13,4);
	\draw[color ={black},line width = 1.2] (-0.13,5)--(0.13,5);
	\draw[color ={black},line width = 1.2] (-0.13,6)--(0.13,6);
	\draw[color ={black},line width = 1.2] (-0.13,7)--(0.13,7);
	\draw[color ={black},line width = 1.2] (-0.13,8)--(0.13,8);
	\draw[color ={black},line width = 1.2] (-0.13,9)--(0.13,9);
	\draw[color ={black},line width = 1.2] (-0.13,10)--(0.13,10);
	\draw[color ={black},line width = 1.2] (-0.13,11)--(0.13,11);
	\draw[color ={black},line width = 1.2] (-0.13,12)--(0.13,12);
	\draw[color ={black},line width = 1.2] (-0.13,-1)--(0.13,-1);
	\draw[color ={black},line width = 1.2] (-0.13,-2)--(0.13,-2);
	\draw[color ={black},line width = 1.2] (-0.13,-3)--(0.13,-3);
	\draw (3,-0.4) node[anchor = center] {\scriptsize \color{black}{$3$}};
	\draw (-3,-0.4) node[anchor = center] {\scriptsize \color{black}{$-3$}};
	\draw (-0.5,3.1) node[anchor = center] {\footnotesize \color{black}{$3$}};
	\draw (-0.5,6.1) node[anchor = center] {\footnotesize \color{black}{$6$}};
	\draw (-0.5,9.1) node[anchor = center] {\footnotesize \color{black}{$9$}};
	\draw (-0.5,-2.9) node[anchor = center] {\footnotesize \color{black}{$-3$}};
	\draw [color={black}] (-0.3,-0.3) node[anchor = center,scale=1, rotate = 0] {O};
	\draw (3,10) node[anchor = center] {\footnotesize \color{blue}{$\Large \mathcal{C}_f$}};
	
	\draw[color={blue},line width = 2] (-1.8,12.44)--(-1.6,10.96)--(-1.4,9.56)--(-1.2,8.24)--(-1,7)--(-0.8,5.84)--(-0.6,4.76)--(-0.4,3.76)--(-0.2,2.84)--(0,2)--(0.2,1.24)--(0.4,0.56)--(0.6,-0.04)--(0.8,-0.56)--(1,-1)--(1.2,-1.36)--(1.4,-1.64)--(1.6,-1.84)--(1.8,-1.96)--(2,-2)--(2.2,-1.96)--(2.4,-1.84)--(2.6,-1.64)--(2.8,-1.36)--(3,-1);

\end{tikzpicture}\\
    \end{minipage}
    \begin{minipage}{0.5\linewidth}
    \begin{tikzpicture}[baseline, scale=0.8]

    \tikzset{
      point/.style={
        thick,
        draw,
        cross out,
        inner sep=0pt,
        minimum width=5pt,
        minimum height=5pt,
      },
    }
    \clip (-4.5,-6.5) rectangle (6.5,9.5);
    	
	\draw[color ={black},>=latex,->] (-3,0)--(3,0);
	\draw[color ={black},>=latex,->] (0,-5)--(0,8);
	\draw[color ={black},opacity = 0.4] (-3,1)--(3,1);
	\draw[color ={black},opacity = 0.4] (-3,2)--(3,2);
	\draw[color ={black},opacity = 0.4] (-3,3)--(3,3);
	\draw[color ={black},opacity = 0.4] (-3,4)--(3,4);
	\draw[color ={black},opacity = 0.4] (-3,5)--(3,5);
	\draw[color ={black},opacity = 0.4] (-3,6)--(3,6);
	\draw[color ={black},opacity = 0.4] (-3,7)--(3,7);
	\draw[color ={black},opacity = 0.4] (-3,8)--(3,8);
	\draw[color ={black},opacity = 0.4] (-3,-1)--(3,-1);
	\draw[color ={black},opacity = 0.4] (-3,-2)--(3,-2);
	\draw[color ={black},opacity = 0.4] (-3,-3)--(3,-3);
	\draw[color ={black},opacity = 0.4] (-3,-4)--(3,-4);
	\draw[color ={black},opacity = 0.4] (-3,-5)--(3,-5);
	\draw[color ={black},opacity = 0.4] (1,-5)--(1,8);
	\draw[color ={black},opacity = 0.4] (2,-5)--(2,8);
	\draw[color ={black},opacity = 0.4] (3,-5)--(3,8);
	\draw[color ={black},opacity = 0.4] (-1,-5)--(-1,8);
	\draw[color ={black},opacity = 0.4] (-2,-5)--(-2,8);
	\draw[color ={black},opacity = 0.4] (-3,-5)--(-3,8);
	\draw[color ={gray},opacity = 0.3] (-3,0)--(3,0);
	\draw[color ={gray},opacity = 0.3] (-3,0.5)--(3,0.5);
	\draw[color ={gray},opacity = 0.3] (-3,1.5)--(3,1.5);
	\draw[color ={gray},opacity = 0.3] (-3,2.5)--(3,2.5);
	\draw[color ={gray},opacity = 0.3] (-3,3.5)--(3,3.5);
	\draw[color ={gray},opacity = 0.3] (-3,4.5)--(3,4.5);
	\draw[color ={gray},opacity = 0.3] (-3,5.5)--(3,5.5);
	\draw[color ={gray},opacity = 0.3] (-3,6.5)--(3,6.5);
	\draw[color ={gray},opacity = 0.3] (-3,7.5)--(3,7.5);
	\draw[color ={gray},opacity = 0.3] (-3,0)--(3,0);
	\draw[color ={gray},opacity = 0.3] (-3,-0.5)--(3,-0.5);
	\draw[color ={gray},opacity = 0.3] (-3,-1.5)--(3,-1.5);
	\draw[color ={gray},opacity = 0.3] (-3,-2.5)--(3,-2.5);
	\draw[color ={gray},opacity = 0.3] (-3,-3.5)--(3,-3.5);
	\draw[color ={gray},opacity = 0.3] (-3,-4.5)--(3,-4.5);
	\draw[color ={gray},opacity = 0.3] (0,-5)--(0,8);
	\draw[color ={gray},opacity = 0.3] (0,-5)--(0,8);
	\draw[color ={black},line width = 1.2] (1,-0.13)--(1,0.13);
	\draw[color ={black},line width = 1.2] (2,-0.13)--(2,0.13);
	\draw[color ={black},line width = 1.2] (3,-0.13)--(3,0.13);
	\draw[color ={black},line width = 1.2] (-1,-0.13)--(-1,0.13);
	\draw[color ={black},line width = 1.2] (-2,-0.13)--(-2,0.13);
	\draw[color ={black},line width = 1.2] (-3,-0.13)--(-3,0.13);
	\draw[color ={black},line width = 1.2] (-0.13,1)--(0.13,1);
	\draw[color ={black},line width = 1.2] (-0.13,2)--(0.13,2);
	\draw[color ={black},line width = 1.2] (-0.13,3)--(0.13,3);
	\draw[color ={black},line width = 1.2] (-0.13,4)--(0.13,4);
	\draw[color ={black},line width = 1.2] (-0.13,5)--(0.13,5);
	\draw[color ={black},line width = 1.2] (-0.13,6)--(0.13,6);
	\draw[color ={black},line width = 1.2] (-0.13,7)--(0.13,7);
	\draw[color ={black},line width = 1.2] (-0.13,8)--(0.13,8);
	\draw[color ={black},line width = 1.2] (-0.13,9)--(0.13,9);
	\draw[color ={black},line width = 1.2] (-0.13,10)--(0.13,10);
	\draw[color ={black},line width = 1.2] (-0.13,11)--(0.13,11);
	\draw[color ={black},line width = 1.2] (-0.13,12)--(0.13,12);
	\draw[color ={black},line width = 1.2] (-0.13,-1)--(0.13,-1);
	\draw[color ={black},line width = 1.2] (-0.13,-2)--(0.13,-2);
	\draw[color ={black},line width = 1.2] (-0.13,-3)--(0.13,-3);
	\draw (3,-0.4) node[anchor = center] {\scriptsize \color{black}{$3$}};
	\draw (-3,-0.4) node[anchor = center] {\scriptsize \color{black}{$-3$}};
	\draw (-0.5,3.1) node[anchor = center] {\footnotesize \color{black}{$3$}};
	\draw (-0.5,6.1) node[anchor = center] {\footnotesize \color{black}{$6$}};
	\draw (-0.5,-2.9) node[anchor = center] {\footnotesize \color{black}{$-3$}};
	\draw [color={black}] (-0.3,-0.3) node[anchor = center,scale=1, rotate = 0] {O};
	\draw (3,6) node[anchor = center] {\footnotesize \color{red}{$\Large\mathcal{C}_f\prime$}};
	
	\draw[color={red},line width = 2] (-1,-6)--(-0.8,-5.6)--(-0.6,-5.2)--(-0.4,-4.8)--(-0.2,-4.4)--(0,-4)--(0.2,-3.6)--(0.4,-3.2)--(0.6,-2.8)--(0.8,-2.4)--(1,-2)--(1.2,-1.6)--(1.4,-1.2)--(1.6,-0.8)--(1.8,-0.4)--(2,0)--(2.2,0.4)--(2.4,0.8)--(2.6,1.2)--(2.8,1.6)--(3,2);

\end{tikzpicture}\\
    \end{minipage}
    
\end{exercice}

\begin{Solution}
 L'équation réduite de la tangente au point d'abscisse $0$ est  : $y=f'(0)(x-0)+f(0)$.\\
        On lit graphiquement $f(0)=2$ et $f'(0)=-4$.\\
        L'équation réduite de la tangente est donc donnée par :
        $y=-4(x+0)+2$, soit $y=-4x+2$.
\end{Solution}

% @see : https://coopmaths.fr/alea?uuid=a1ba2&id=can1F14&n=1&d=10&alea=rkIp22323232323423456789101112131415161718192021222324252627282930313233343536373839404142434445464748&cd=1&cols=1
\needspace{10\baselineskip}
\begin{exercice}[BaremeTotal=false]


Soit $f$ la fonction définie sur $\mathbb{R}$ par : $f(x)= x^2-5x-8$.\\
En calculant la limite du taux d'accroissement, déterminer $f'(-1)$.
\end{exercice}

\begin{Solution}
 $f$ est une fonction polynôme du second degré de la forme $f(x)=ax^2+bx+c$.\\
    La fonction dérivée est donnée par la somme des dérivées des fonctions $u$ et $v$ définies par $u(x)=x^2$ et $v(x)=-5x-8$.\\
     Comme $u'(x)=2x$ et $v'(x)=-5$, on obtient  $f'(x)=2x-5$. \\
     Ainsi, $f'(-1)=2\times (-1)-5=-7$.
\end{Solution}

% @see : https://coopmaths.fr/alea?uuid=3c690&id=can1F13&n=1&d=10&alea=EUDu22323232323423456789101112131415161718192021222324252627282930313233343536373839404142434445464748&cd=1&cols=1
\needspace{10\baselineskip}
\begin{exercice}[BaremeTotal=false]


 Déterminer le coefficient directeur de la tangente à la courbe représentative de la fonction cube au point d'abscisse $1$.
                        
\end{exercice}

\begin{Solution}
 Le coefficient directeur de la tangente au point d'abscisse $a$ est donné par le nombre dérivé $f'(a)$.\\
        La fonction $f$ définie par $f(x)=x^3$ a pour fonction dérivée la fonction $f'$ définie par $f'(x)=3x^2$.\\
        Comme $f'(1)=3\times 1^2=3$, le coefficient directeur de la tangente au point d'abscisse $1$ est : $3$. 
\end{Solution}

% @see : https://coopmaths.fr/alea?uuid=ebd89&id=1AN14-1&n=1&d=10&s=3&alea=ocYr22323232323423456789101112131415161718192021222324252627282930313233343536373839404142434445464748&cd=0&cols=1
\needspace{10\baselineskip}
\begin{exercice}[BaremeTotal=false]


 Donner la dérivée de la fonction $f$, dérivable pour tout $x\in\R$, définie par  $f(x)=-\dfrac{1}{5}x+5$.
\end{exercice}

\begin{Solution}
 L'expression de la dérivée de la fonction $f$ définie par $f(x)=-\dfrac{1}{5}x+5$ est : ${\color[HTML]{f15929}\boldsymbol{f'(x)=-\dfrac{1}{5}}}$.
\end{Solution}

% @see : https://coopmaths.fr/alea?uuid=ebd89&id=1AN14-1&n=1&d=10&s=4&alea=uAlP22323232323423456789101112131415161718192021222324252627282930313233343536373839404142434445464748&cd=0&cols=1
\needspace{10\baselineskip}
\begin{exercice}[BaremeTotal=false]


 Donner la dérivée de la fonction $f$, dérivable pour tout $x\in\R$, définie par  $f(x)=4x^{9}$.
\end{exercice}

\begin{Solution}
 L'expression de la dérivée de la fonction $f$ définie par $f(x)=4x^{9}$ est : ${\color[HTML]{f15929}\boldsymbol{f'(x)=36x^{8}}}$.
\end{Solution}

% @see : https://coopmaths.fr/alea?uuid=ebd89&id=1AN14-1&n=1&d=10&s=7&alea=vjN722323232323423456789101112131415161718192021222324252627282930313233343536373839404142434445464748&cd=0&cols=1
\needspace{10\baselineskip}
\begin{exercice}[BaremeTotal=false]


 Donner la dérivée de la fonction $f$, dérivable pour tout $x\in\R^*$, définie par  $f(x)=\dfrac{9}{10x}$.
\end{exercice}

\begin{Solution}
 L'expression de la dérivée de la fonction $f$ définie par $f(x)=\dfrac{9}{10x}$ est : ${\color[HTML]{f15929}\boldsymbol{f'(x)=-\dfrac{9}{10x^2}}}$.
\end{Solution}

% @see : https://coopmaths.fr/alea?uuid=a83c0&id=1AN14-4&n=1&d=10&s=1&alea=Njr322323232323423456789101112131415161718192021222324252627282930313233343536373839404142434445464748&cd=0&cols=1
\needspace{10\baselineskip}
\begin{exercice}[BaremeTotal=false]


 Donner l'expression de la dérivée de la fonction $f$ définie sur $\R^{*}$ par $f(x)=\dfrac{1}{x}+2x^{2}$.\\
\end{exercice}

\begin{Solution}
 L'expression de la dérivée de la fonction $f$ définie par $f(x)=\dfrac{1}{x}+2x^{2}$ est :\\$f'(x)={\color[HTML]{f15929}\boldsymbol{-\dfrac{1}{x^2}+4x}}$.\\
\end{Solution}

% @see : https://coopmaths.fr/alea?uuid=a83c0&id=1AN14-4&n=1&d=10&s=4&alea=4Vpz22323232323423456789101112131415161718192021222324252627282930313233343536373839404142434445464748&cd=1&cols=1
\needspace{10\baselineskip}
\begin{exercice}[BaremeTotal=false]


 Donner l'expression de la dérivée de la fonction $f$ définie sur $\R_+^{*}$ par $f(x)=-2x^{2}+\sqrt{x}+\dfrac{2}{x}$.\\
\end{exercice}

\begin{Solution}
 $(-2x^{2})^\prime=-4x$\\$(\sqrt{x})^\prime=\dfrac{1}{2\sqrt{x}}$\\$(\dfrac{2}{x})^\prime=-\dfrac{2}{x^2}$\\L'expression de la dérivée de la fonction $f$ définie par $f(x)=-2x^{2}+\sqrt{x}+\dfrac{2}{x}$ est :\\$f'(x)={\color[HTML]{f15929}\boldsymbol{-4x+\dfrac{1}{2\sqrt{x}}-\dfrac{2}{x^2}}}$.\\
\end{Solution}

\end{Maquette}
\clearpage
\begin{Maquette}[DM]{Niveau= ,Classe=\getdata[2]\mydata\ (v2), Date = 06/01/25  ,Theme=Exercices}


% @see : https://coopmaths.fr/alea?uuid=4c8c7&id=1AN11-3&n=1&d=10&s=2&alea=nawO223232323234234567891011121314151617181920212223242526272829303132333435363738394041424344454647482&cd=1&cols=2
\needspace{10\baselineskip}
\begin{exercice}[BaremeTotal=false]
\label{eleve:2}


 \begin{minipage}[t]{\linewidth} Soit $f$ une fonction dérivable sur $[-5;5]$ et $\mathcal{C}_f$ sa courbe représentative.\\On sait que  $f(4)=1~~$ et que $~~f'(4)=-4$.\\Déterminer une équation de la tangente $(T)$ à la courbe $\mathcal{C}_f$ au point d'abscisse $4$,\\sans utiliser la formule de cours de l'équation de tangente. \end{minipage}

\end{exercice}

\begin{Solution}
 \begin{minipage}[t]{\linewidth} On sait que la tangente n'est pas une droite verticale, puisque la fonction est dérivable sur l'intervalle.\\On en déduit que la tangente $(T)$ au point d'abscisse $4$, admet une équation réduite de la forme :  $(T) : y=m x + p$. 

\medskip
$\bullet$ Détermination de $m$ :\\On sait que le nombre dérivé en $4$ est par définition, le coefficient directeur de la tangente au point d'abscisse $4$.\\Par conséquent, on a déjà : $m=f'(4)=-4$.\\ On en déduit que  $(T) : y= -4 x + p$\\$\bullet$ Détermination de $p$ :\\ Pour cela, on utilise que si $f(4)=1~~$, alors le point $A$ de coordonnées $(4;1)$ appartient à $\mathcal{C}_f$ mais aussi à $(T)$.\\ On peut écrire $A(4;1) = \mathcal{C}_f \cap (T)$.\\ On remplace alors les coordonnées de $A(4;1)$ dans l'équation  $(T) : y= -4 x + p$.\\ $\begin{aligned} \phantom{\iff}&A(4;1)\in (T)\\ \iff& 1= -4 \times 4  + p\\ \iff& p=1 +4 \times 4  \\ \iff& p=1 +16   \\ \iff& p=17\\\end{aligned}$\\On peut conclure que : $(T) : {\color[HTML]{f15929}\boldsymbol{y=-4x+17}}$. \end{minipage}

\end{Solution}

% @see : https://coopmaths.fr/alea?uuid=0e984&id=can1F15&n=1&d=10&alea=1Alo223232323234234567891011121314151617181920212223242526272829303132333435363738394041424344454647482&cd=1&cols=1
\needspace{10\baselineskip}
\begin{exercice}[BaremeTotal=false]


 La courbe représente une fonction $f$ et la droite est la tangente au point d'abscisse $2$.\\
        Déterminer $f'(2)$.\begin{center}\begin{tikzpicture}[baseline,scale = 0.5]

    \tikzset{
      point/.style={
        thick,
        draw,
        cross out,
        inner sep=0pt,
        minimum width=5pt,
        minimum height=5pt,
      },
    }
    \clip (-2,-2) rectangle (14,12);
    	
	\draw[color ={black},line width = 2,>=latex,->] (-2,0)--(14,0);
	\draw[color ={black},line width = 2,>=latex,->] (0,-2)--(0,12);
	\draw[color ={black},opacity = 0.5] (-2,1)--(14,1);
	\draw[color ={black},opacity = 0.5] (-2,2)--(14,2);
	\draw[color ={black},opacity = 0.5] (-2,3)--(14,3);
	\draw[color ={black},opacity = 0.5] (-2,4)--(14,4);
	\draw[color ={black},opacity = 0.5] (-2,5)--(14,5);
	\draw[color ={black},opacity = 0.5] (-2,6)--(14,6);
	\draw[color ={black},opacity = 0.5] (-2,7)--(14,7);
	\draw[color ={black},opacity = 0.5] (-2,8)--(14,8);
	\draw[color ={black},opacity = 0.5] (-2,9)--(14,9);
	\draw[color ={black},opacity = 0.5] (-2,10)--(14,10);
	\draw[color ={black},opacity = 0.5] (-2,11)--(14,11);
	\draw[color ={black},opacity = 0.5] (-2,12)--(14,12);
	\draw[color ={black},opacity = 0.5] (-2,-1)--(14,-1);
	\draw[color ={black},opacity = 0.5] (-2,-2)--(14,-2);
	\draw[color ={black},opacity = 0.5] (1,-2)--(1,12);
	\draw[color ={black},opacity = 0.5] (2,-2)--(2,12);
	\draw[color ={black},opacity = 0.5] (3,-2)--(3,12);
	\draw[color ={black},opacity = 0.5] (4,-2)--(4,12);
	\draw[color ={black},opacity = 0.5] (5,-2)--(5,12);
	\draw[color ={black},opacity = 0.5] (6,-2)--(6,12);
	\draw[color ={black},opacity = 0.5] (7,-2)--(7,12);
	\draw[color ={black},opacity = 0.5] (8,-2)--(8,12);
	\draw[color ={black},opacity = 0.5] (9,-2)--(9,12);
	\draw[color ={black},opacity = 0.5] (10,-2)--(10,12);
	\draw[color ={black},opacity = 0.5] (11,-2)--(11,12);
	\draw[color ={black},opacity = 0.5] (12,-2)--(12,12);
	\draw[color ={black},opacity = 0.5] (13,-2)--(13,12);
	\draw[color ={black},opacity = 0.5] (14,-2)--(14,12);
	\draw[color ={black},opacity = 0.5] (-1,-2)--(-1,12);
	\draw[color ={black},opacity = 0.5] (-2,-2)--(-2,12);
	\draw[color ={gray},opacity = 0.3] (-2,0)--(14,0);
	\draw[color ={gray},opacity = 0.3] (-2,0.5)--(14,0.5);
	\draw[color ={gray},opacity = 0.3] (-2,1.5)--(14,1.5);
	\draw[color ={gray},opacity = 0.3] (-2,2.5)--(14,2.5);
	\draw[color ={gray},opacity = 0.3] (-2,3.5)--(14,3.5);
	\draw[color ={gray},opacity = 0.3] (-2,4.5)--(14,4.5);
	\draw[color ={gray},opacity = 0.3] (-2,5.5)--(14,5.5);
	\draw[color ={gray},opacity = 0.3] (-2,6.5)--(14,6.5);
	\draw[color ={gray},opacity = 0.3] (-2,7.5)--(14,7.5);
	\draw[color ={gray},opacity = 0.3] (-2,8.5)--(14,8.5);
	\draw[color ={gray},opacity = 0.3] (-2,9.5)--(14,9.5);
	\draw[color ={gray},opacity = 0.3] (-2,10.5)--(14,10.5);
	\draw[color ={gray},opacity = 0.3] (-2,11.5)--(14,11.5);
	\draw[color ={gray},opacity = 0.3] (-2,0)--(14,0);
	\draw[color ={gray},opacity = 0.3] (-2,-0.5)--(14,-0.5);
	\draw[color ={gray},opacity = 0.3] (-2,-1.5)--(14,-1.5);
	\draw[color ={gray},opacity = 0.3] (0,-2)--(0,12);
	\draw[color ={gray},opacity = 0.3] (0.1,-2)--(0.1,12);
	\draw[color ={gray},opacity = 0.3] (0.2,-2)--(0.2,12);
	\draw[color ={gray},opacity = 0.3] (0.30000000000000004,-2)--(0.30000000000000004,12);
	\draw[color ={gray},opacity = 0.3] (0.4,-2)--(0.4,12);
	\draw[color ={gray},opacity = 0.3] (0.5,-2)--(0.5,12);
	\draw[color ={gray},opacity = 0.3] (0.6,-2)--(0.6,12);
	\draw[color ={gray},opacity = 0.3] (0.7,-2)--(0.7,12);
	\draw[color ={gray},opacity = 0.3] (0.7999999999999999,-2)--(0.7999999999999999,12);
	\draw[color ={gray},opacity = 0.3] (0.8999999999999999,-2)--(0.8999999999999999,12);
	\draw[color ={gray},opacity = 0.3] (1.0999999999999999,-2)--(1.0999999999999999,12);
	\draw[color ={gray},opacity = 0.3] (1.2,-2)--(1.2,12);
	\draw[color ={gray},opacity = 0.3] (1.3,-2)--(1.3,12);
	\draw[color ={gray},opacity = 0.3] (1.4000000000000001,-2)--(1.4000000000000001,12);
	\draw[color ={gray},opacity = 0.3] (1.5000000000000002,-2)--(1.5000000000000002,12);
	\draw[color ={gray},opacity = 0.3] (1.6000000000000003,-2)--(1.6000000000000003,12);
	\draw[color ={gray},opacity = 0.3] (1.7000000000000004,-2)--(1.7000000000000004,12);
	\draw[color ={gray},opacity = 0.3] (1.8000000000000005,-2)--(1.8000000000000005,12);
	\draw[color ={gray},opacity = 0.3] (1.9000000000000006,-2)--(1.9000000000000006,12);
	\draw[color ={gray},opacity = 0.3] (2.1000000000000005,-2)--(2.1000000000000005,12);
	\draw[color ={gray},opacity = 0.3] (2.2000000000000006,-2)--(2.2000000000000006,12);
	\draw[color ={gray},opacity = 0.3] (2.3000000000000007,-2)--(2.3000000000000007,12);
	\draw[color ={gray},opacity = 0.3] (2.400000000000001,-2)--(2.400000000000001,12);
	\draw[color ={gray},opacity = 0.3] (2.500000000000001,-2)--(2.500000000000001,12);
	\draw[color ={gray},opacity = 0.3] (2.600000000000001,-2)--(2.600000000000001,12);
	\draw[color ={gray},opacity = 0.3] (2.700000000000001,-2)--(2.700000000000001,12);
	\draw[color ={gray},opacity = 0.3] (2.800000000000001,-2)--(2.800000000000001,12);
	\draw[color ={gray},opacity = 0.3] (2.9000000000000012,-2)--(2.9000000000000012,12);
	\draw[color ={gray},opacity = 0.3] (3.1000000000000014,-2)--(3.1000000000000014,12);
	\draw[color ={gray},opacity = 0.3] (3.2000000000000015,-2)--(3.2000000000000015,12);
	\draw[color ={gray},opacity = 0.3] (3.3000000000000016,-2)--(3.3000000000000016,12);
	\draw[color ={gray},opacity = 0.3] (3.4000000000000017,-2)--(3.4000000000000017,12);
	\draw[color ={gray},opacity = 0.3] (3.5000000000000018,-2)--(3.5000000000000018,12);
	\draw[color ={gray},opacity = 0.3] (3.600000000000002,-2)--(3.600000000000002,12);
	\draw[color ={gray},opacity = 0.3] (3.700000000000002,-2)--(3.700000000000002,12);
	\draw[color ={gray},opacity = 0.3] (3.800000000000002,-2)--(3.800000000000002,12);
	\draw[color ={gray},opacity = 0.3] (3.900000000000002,-2)--(3.900000000000002,12);
	\draw[color ={gray},opacity = 0.3] (4.100000000000001,-2)--(4.100000000000001,12);
	\draw[color ={gray},opacity = 0.3] (4.200000000000001,-2)--(4.200000000000001,12);
	\draw[color ={gray},opacity = 0.3] (4.300000000000001,-2)--(4.300000000000001,12);
	\draw[color ={gray},opacity = 0.3] (4.4,-2)--(4.4,12);
	\draw[color ={gray},opacity = 0.3] (4.5,-2)--(4.5,12);
	\draw[color ={gray},opacity = 0.3] (4.6,-2)--(4.6,12);
	\draw[color ={gray},opacity = 0.3] (4.699999999999999,-2)--(4.699999999999999,12);
	\draw[color ={gray},opacity = 0.3] (4.799999999999999,-2)--(4.799999999999999,12);
	\draw[color ={gray},opacity = 0.3] (4.899999999999999,-2)--(4.899999999999999,12);
	\draw[color ={gray},opacity = 0.3] (5.099999999999998,-2)--(5.099999999999998,12);
	\draw[color ={gray},opacity = 0.3] (5.1999999999999975,-2)--(5.1999999999999975,12);
	\draw[color ={gray},opacity = 0.3] (5.299999999999997,-2)--(5.299999999999997,12);
	\draw[color ={gray},opacity = 0.3] (5.399999999999997,-2)--(5.399999999999997,12);
	\draw[color ={gray},opacity = 0.3] (5.4999999999999964,-2)--(5.4999999999999964,12);
	\draw[color ={gray},opacity = 0.3] (5.599999999999996,-2)--(5.599999999999996,12);
	\draw[color ={gray},opacity = 0.3] (5.699999999999996,-2)--(5.699999999999996,12);
	\draw[color ={gray},opacity = 0.3] (5.799999999999995,-2)--(5.799999999999995,12);
	\draw[color ={gray},opacity = 0.3] (5.899999999999995,-2)--(5.899999999999995,12);
	\draw[color ={gray},opacity = 0.3] (6.099999999999994,-2)--(6.099999999999994,12);
	\draw[color ={gray},opacity = 0.3] (6.199999999999994,-2)--(6.199999999999994,12);
	\draw[color ={gray},opacity = 0.3] (6.299999999999994,-2)--(6.299999999999994,12);
	\draw[color ={gray},opacity = 0.3] (6.399999999999993,-2)--(6.399999999999993,12);
	\draw[color ={gray},opacity = 0.3] (6.499999999999993,-2)--(6.499999999999993,12);
	\draw[color ={gray},opacity = 0.3] (6.5999999999999925,-2)--(6.5999999999999925,12);
	\draw[color ={gray},opacity = 0.3] (6.699999999999992,-2)--(6.699999999999992,12);
	\draw[color ={gray},opacity = 0.3] (6.799999999999992,-2)--(6.799999999999992,12);
	\draw[color ={gray},opacity = 0.3] (6.8999999999999915,-2)--(6.8999999999999915,12);
	\draw[color ={gray},opacity = 0.3] (7.099999999999991,-2)--(7.099999999999991,12);
	\draw[color ={gray},opacity = 0.3] (7.19999999999999,-2)--(7.19999999999999,12);
	\draw[color ={gray},opacity = 0.3] (7.29999999999999,-2)--(7.29999999999999,12);
	\draw[color ={gray},opacity = 0.3] (7.39999999999999,-2)--(7.39999999999999,12);
	\draw[color ={gray},opacity = 0.3] (7.499999999999989,-2)--(7.499999999999989,12);
	\draw[color ={gray},opacity = 0.3] (7.599999999999989,-2)--(7.599999999999989,12);
	\draw[color ={gray},opacity = 0.3] (7.699999999999989,-2)--(7.699999999999989,12);
	\draw[color ={gray},opacity = 0.3] (7.799999999999988,-2)--(7.799999999999988,12);
	\draw[color ={gray},opacity = 0.3] (7.899999999999988,-2)--(7.899999999999988,12);
	\draw[color ={gray},opacity = 0.3] (8.099999999999987,-2)--(8.099999999999987,12);
	\draw[color ={gray},opacity = 0.3] (8.199999999999987,-2)--(8.199999999999987,12);
	\draw[color ={gray},opacity = 0.3] (8.299999999999986,-2)--(8.299999999999986,12);
	\draw[color ={gray},opacity = 0.3] (8.399999999999986,-2)--(8.399999999999986,12);
	\draw[color ={gray},opacity = 0.3] (8.499999999999986,-2)--(8.499999999999986,12);
	\draw[color ={gray},opacity = 0.3] (8.599999999999985,-2)--(8.599999999999985,12);
	\draw[color ={gray},opacity = 0.3] (8.699999999999985,-2)--(8.699999999999985,12);
	\draw[color ={gray},opacity = 0.3] (8.799999999999985,-2)--(8.799999999999985,12);
	\draw[color ={gray},opacity = 0.3] (8.899999999999984,-2)--(8.899999999999984,12);
	\draw[color ={gray},opacity = 0.3] (9.099999999999984,-2)--(9.099999999999984,12);
	\draw[color ={gray},opacity = 0.3] (9.199999999999983,-2)--(9.199999999999983,12);
	\draw[color ={gray},opacity = 0.3] (9.299999999999983,-2)--(9.299999999999983,12);
	\draw[color ={gray},opacity = 0.3] (9.399999999999983,-2)--(9.399999999999983,12);
	\draw[color ={gray},opacity = 0.3] (9.499999999999982,-2)--(9.499999999999982,12);
	\draw[color ={gray},opacity = 0.3] (9.599999999999982,-2)--(9.599999999999982,12);
	\draw[color ={gray},opacity = 0.3] (9.699999999999982,-2)--(9.699999999999982,12);
	\draw[color ={gray},opacity = 0.3] (9.799999999999981,-2)--(9.799999999999981,12);
	\draw[color ={gray},opacity = 0.3] (9.89999999999998,-2)--(9.89999999999998,12);
	\draw[color ={gray},opacity = 0.3] (10.09999999999998,-2)--(10.09999999999998,12);
	\draw[color ={gray},opacity = 0.3] (10.19999999999998,-2)--(10.19999999999998,12);
	\draw[color ={gray},opacity = 0.3] (10.29999999999998,-2)--(10.29999999999998,12);
	\draw[color ={gray},opacity = 0.3] (10.399999999999979,-2)--(10.399999999999979,12);
	\draw[color ={gray},opacity = 0.3] (10.499999999999979,-2)--(10.499999999999979,12);
	\draw[color ={gray},opacity = 0.3] (10.599999999999978,-2)--(10.599999999999978,12);
	\draw[color ={gray},opacity = 0.3] (10.699999999999978,-2)--(10.699999999999978,12);
	\draw[color ={gray},opacity = 0.3] (10.799999999999978,-2)--(10.799999999999978,12);
	\draw[color ={gray},opacity = 0.3] (10.899999999999977,-2)--(10.899999999999977,12);
	\draw[color ={gray},opacity = 0.3] (11.099999999999977,-2)--(11.099999999999977,12);
	\draw[color ={gray},opacity = 0.3] (11.199999999999976,-2)--(11.199999999999976,12);
	\draw[color ={gray},opacity = 0.3] (11.299999999999976,-2)--(11.299999999999976,12);
	\draw[color ={gray},opacity = 0.3] (11.399999999999975,-2)--(11.399999999999975,12);
	\draw[color ={gray},opacity = 0.3] (11.499999999999975,-2)--(11.499999999999975,12);
	\draw[color ={gray},opacity = 0.3] (11.599999999999975,-2)--(11.599999999999975,12);
	\draw[color ={gray},opacity = 0.3] (11.699999999999974,-2)--(11.699999999999974,12);
	\draw[color ={gray},opacity = 0.3] (11.799999999999974,-2)--(11.799999999999974,12);
	\draw[color ={gray},opacity = 0.3] (11.899999999999974,-2)--(11.899999999999974,12);
	\draw[color ={gray},opacity = 0.3] (12.099999999999973,-2)--(12.099999999999973,12);
	\draw[color ={gray},opacity = 0.3] (12.199999999999973,-2)--(12.199999999999973,12);
	\draw[color ={gray},opacity = 0.3] (12.299999999999972,-2)--(12.299999999999972,12);
	\draw[color ={gray},opacity = 0.3] (12.399999999999972,-2)--(12.399999999999972,12);
	\draw[color ={gray},opacity = 0.3] (12.499999999999972,-2)--(12.499999999999972,12);
	\draw[color ={gray},opacity = 0.3] (12.599999999999971,-2)--(12.599999999999971,12);
	\draw[color ={gray},opacity = 0.3] (12.69999999999997,-2)--(12.69999999999997,12);
	\draw[color ={gray},opacity = 0.3] (12.79999999999997,-2)--(12.79999999999997,12);
	\draw[color ={gray},opacity = 0.3] (12.89999999999997,-2)--(12.89999999999997,12);
	\draw[color ={gray},opacity = 0.3] (13.09999999999997,-2)--(13.09999999999997,12);
	\draw[color ={gray},opacity = 0.3] (13.199999999999969,-2)--(13.199999999999969,12);
	\draw[color ={gray},opacity = 0.3] (13.299999999999969,-2)--(13.299999999999969,12);
	\draw[color ={gray},opacity = 0.3] (13.399999999999968,-2)--(13.399999999999968,12);
	\draw[color ={gray},opacity = 0.3] (13.499999999999968,-2)--(13.499999999999968,12);
	\draw[color ={gray},opacity = 0.3] (13.599999999999968,-2)--(13.599999999999968,12);
	\draw[color ={gray},opacity = 0.3] (13.699999999999967,-2)--(13.699999999999967,12);
	\draw[color ={gray},opacity = 0.3] (13.799999999999967,-2)--(13.799999999999967,12);
	\draw[color ={gray},opacity = 0.3] (13.899999999999967,-2)--(13.899999999999967,12);
	\draw[color ={gray},opacity = 0.3] (0,-2)--(0,12);
	\draw[color ={gray},opacity = 0.3] (-0.1,-2)--(-0.1,12);
	\draw[color ={gray},opacity = 0.3] (-0.2,-2)--(-0.2,12);
	\draw[color ={gray},opacity = 0.3] (-0.30000000000000004,-2)--(-0.30000000000000004,12);
	\draw[color ={gray},opacity = 0.3] (-0.4,-2)--(-0.4,12);
	\draw[color ={gray},opacity = 0.3] (-0.5,-2)--(-0.5,12);
	\draw[color ={gray},opacity = 0.3] (-0.6,-2)--(-0.6,12);
	\draw[color ={gray},opacity = 0.3] (-0.7,-2)--(-0.7,12);
	\draw[color ={gray},opacity = 0.3] (-0.7999999999999999,-2)--(-0.7999999999999999,12);
	\draw[color ={gray},opacity = 0.3] (-0.8999999999999999,-2)--(-0.8999999999999999,12);
	\draw[color ={gray},opacity = 0.3] (-1.0999999999999999,-2)--(-1.0999999999999999,12);
	\draw[color ={gray},opacity = 0.3] (-1.2,-2)--(-1.2,12);
	\draw[color ={gray},opacity = 0.3] (-1.3,-2)--(-1.3,12);
	\draw[color ={gray},opacity = 0.3] (-1.4000000000000001,-2)--(-1.4000000000000001,12);
	\draw[color ={gray},opacity = 0.3] (-1.5000000000000002,-2)--(-1.5000000000000002,12);
	\draw[color ={gray},opacity = 0.3] (-1.6000000000000003,-2)--(-1.6000000000000003,12);
	\draw[color ={gray},opacity = 0.3] (-1.7000000000000004,-2)--(-1.7000000000000004,12);
	\draw[color ={gray},opacity = 0.3] (-1.8000000000000005,-2)--(-1.8000000000000005,12);
	\draw[color ={gray},opacity = 0.3] (-1.9000000000000006,-2)--(-1.9000000000000006,12);
	\draw[color ={black},line width = 2] (2,-0.2)--(2,0.2);
	\draw[color ={black},line width = 2] (4,-0.2)--(4,0.2);
	\draw[color ={black},line width = 2] (6,-0.2)--(6,0.2);
	\draw[color ={black},line width = 2] (8,-0.2)--(8,0.2);
	\draw[color ={black},line width = 2] (10,-0.2)--(10,0.2);
	\draw[color ={black},line width = 2] (12,-0.2)--(12,0.2);
	\draw[color ={black},line width = 2] (-0.2,2)--(0.2,2);
	\draw[color ={black},line width = 2] (-0.2,4)--(0.2,4);
	\draw[color ={black},line width = 2] (-0.2,6)--(0.2,6);
	\draw[color ={black},line width = 2] (-0.2,8)--(0.2,8);
	\draw[color ={black},line width = 2] (-0.2,10)--(0.2,10);
	\draw (2,-0.4) node[anchor = center] {\scriptsize \color{black}{$1$}};
	\draw (4,-0.4) node[anchor = center] {\scriptsize \color{black}{$2$}};
	\draw (6,-0.4) node[anchor = center] {\scriptsize \color{black}{$3$}};
	\draw (8,-0.4) node[anchor = center] {\scriptsize \color{black}{$4$}};
	\draw (10,-0.4) node[anchor = center] {\scriptsize \color{black}{$5$}};
	\draw (-0.5,2.1) node[anchor = center] {\footnotesize \color{black}{$1$}};
	\draw (-0.5,4.1) node[anchor = center] {\footnotesize \color{black}{$2$}};
	\draw (-0.5,6.1) node[anchor = center] {\footnotesize \color{black}{$3$}};
	\draw (-0.5,8.1) node[anchor = center] {\footnotesize \color{black}{$4$}};
	\draw (-0.5,10.1) node[anchor = center] {\footnotesize \color{black}{$5$}};
	\draw [color={black}] (-0.3,-0.3) node[anchor = center,scale=1, rotate = 0] {O};
	
	\draw[color={blue},line width = 2] (0.6,12.67)--(1,10)--(1.4,8.86)--(1.8,8.22)--(2.2,7.82)--(2.6,7.54)--(3,7.33)--(3.4,7.18)--(3.8,7.05)--(4.2,6.95)--(4.6,6.87)--(5,6.8)--(5.4,6.74)--(5.8,6.69)--(6.2,6.65)--(6.6,6.61)--(7,6.57)--(7.4,6.54)--(7.8,6.51)--(8.2,6.49)--(8.6,6.47)--(9,6.44)--(9.4,6.43)--(9.8,6.41)--(10.2,6.39)--(10.6,6.38)--(11,6.36)--(11.4,6.35)--(11.8,6.34)--(12.2,6.33)--(12.6,6.32)--(13,6.31)--(13.4,6.3)--(13.8,6.29);
	
	\draw[color={red},line width = 2] (-2,8.5)--(-1.6,8.4)--(-1.2,8.3)--(-0.8,8.2)--(-0.4,8.1)--(0,8)--(0.4,7.9)--(0.8,7.8)--(1.2,7.7)--(1.6,7.6)--(2,7.5)--(2.4,7.4)--(2.8,7.3)--(3.2,7.2)--(3.6,7.1)--(4,7)--(4.4,6.9)--(4.8,6.8)--(5.2,6.7)--(5.6,6.6)--(6,6.5)--(6.4,6.4)--(6.8,6.3)--(7.2,6.2)--(7.6,6.1)--(8,6)--(8.4,5.9)--(8.8,5.8)--(9.2,5.7)--(9.6,5.6)--(10,5.5)--(10.4,5.4)--(10.8,5.3)--(11.2,5.2)--(11.6,5.1)--(12,5)--(12.4,4.9)--(12.8,4.8)--(13.2,4.7)--(13.6,4.6)--(14,4.5);

\end{tikzpicture}\end{center}
\end{exercice}

\begin{Solution}
 $f'(2)$ est donné par le coefficient directeur de la tangente à la courbe au point d'abscisse $2$, soit $\dfrac{-1}{4}=-\dfrac{1}{4}$.
\end{Solution}

% @see : https://coopmaths.fr/alea?uuid=6f32d&id=can1F16&n=1&d=10&alea=dqSw223232323234234567891011121314151617181920212223242526272829303132333435363738394041424344454647482&cd=1&cols=1
\needspace{10\baselineskip}
\newpage
\begin{exercice}[BaremeTotal=false]


 On donne les représentations graphiques d'une fonction et de sa dérivée.\\
        Donner l'équation réduite de la tangente à la courbe de $f$ en $x=3$. \\ \begin{minipage}{0.5\linewidth}
    \begin{tikzpicture}[baseline, scale=0.8]

    \tikzset{
      point/.style={
        thick,
        draw,
        cross out,
        inner sep=0pt,
        minimum width=5pt,
        minimum height=5pt,
      },
    }
    \clip (-4.5,-4.5) rectangle (5.75,13.5);
    	
	\draw[color ={black},>=latex,->] (-3,0)--(3,0);
	\draw[color ={black},>=latex,->] (0,-3)--(0,12);
	\draw[color ={black},opacity = 0.4] (-3,1)--(3,1);
	\draw[color ={black},opacity = 0.4] (-3,2)--(3,2);
	\draw[color ={black},opacity = 0.4] (-3,3)--(3,3);
	\draw[color ={black},opacity = 0.4] (-3,4)--(3,4);
	\draw[color ={black},opacity = 0.4] (-3,5)--(3,5);
	\draw[color ={black},opacity = 0.4] (-3,6)--(3,6);
	\draw[color ={black},opacity = 0.4] (-3,7)--(3,7);
	\draw[color ={black},opacity = 0.4] (-3,8)--(3,8);
	\draw[color ={black},opacity = 0.4] (-3,9)--(3,9);
	\draw[color ={black},opacity = 0.4] (-3,10)--(3,10);
	\draw[color ={black},opacity = 0.4] (-3,11)--(3,11);
	\draw[color ={black},opacity = 0.4] (-3,12)--(3,12);
	\draw[color ={black},opacity = 0.4] (-3,-1)--(3,-1);
	\draw[color ={black},opacity = 0.4] (-3,-2)--(3,-2);
	\draw[color ={black},opacity = 0.4] (-3,-3)--(3,-3);
	\draw[color ={black},opacity = 0.4] (1,-3)--(1,12);
	\draw[color ={black},opacity = 0.4] (2,-3)--(2,12);
	\draw[color ={black},opacity = 0.4] (3,-3)--(3,12);
	\draw[color ={black},opacity = 0.4] (-1,-3)--(-1,12);
	\draw[color ={black},opacity = 0.4] (-2,-3)--(-2,12);
	\draw[color ={black},opacity = 0.4] (-3,-3)--(-3,12);
	\draw[color ={gray},opacity = 0.3] (-3,0)--(3,0);
	\draw[color ={gray},opacity = 0.3] (-3,0.5)--(3,0.5);
	\draw[color ={gray},opacity = 0.3] (-3,1.5)--(3,1.5);
	\draw[color ={gray},opacity = 0.3] (-3,2.5)--(3,2.5);
	\draw[color ={gray},opacity = 0.3] (-3,3.5)--(3,3.5);
	\draw[color ={gray},opacity = 0.3] (-3,4.5)--(3,4.5);
	\draw[color ={gray},opacity = 0.3] (-3,5.5)--(3,5.5);
	\draw[color ={gray},opacity = 0.3] (-3,6.5)--(3,6.5);
	\draw[color ={gray},opacity = 0.3] (-3,7.5)--(3,7.5);
	\draw[color ={gray},opacity = 0.3] (-3,8.5)--(3,8.5);
	\draw[color ={gray},opacity = 0.3] (-3,9.5)--(3,9.5);
	\draw[color ={gray},opacity = 0.3] (-3,10.5)--(3,10.5);
	\draw[color ={gray},opacity = 0.3] (-3,11.5)--(3,11.5);
	\draw[color ={gray},opacity = 0.3] (-3,0)--(3,0);
	\draw[color ={gray},opacity = 0.3] (-3,-0.5)--(3,-0.5);
	\draw[color ={gray},opacity = 0.3] (-3,-1.5)--(3,-1.5);
	\draw[color ={gray},opacity = 0.3] (-3,-2.5)--(3,-2.5);
	\draw[color ={gray},opacity = 0.3] (0,-3)--(0,12);
	\draw[color ={gray},opacity = 0.3] (0,-3)--(0,12);
	\draw[color ={black},line width = 1.2] (1,-0.13)--(1,0.13);
	\draw[color ={black},line width = 1.2] (2,-0.13)--(2,0.13);
	\draw[color ={black},line width = 1.2] (3,-0.13)--(3,0.13);
	\draw[color ={black},line width = 1.2] (-1,-0.13)--(-1,0.13);
	\draw[color ={black},line width = 1.2] (-2,-0.13)--(-2,0.13);
	\draw[color ={black},line width = 1.2] (-3,-0.13)--(-3,0.13);
	\draw[color ={black},line width = 1.2] (-0.13,1)--(0.13,1);
	\draw[color ={black},line width = 1.2] (-0.13,2)--(0.13,2);
	\draw[color ={black},line width = 1.2] (-0.13,3)--(0.13,3);
	\draw[color ={black},line width = 1.2] (-0.13,4)--(0.13,4);
	\draw[color ={black},line width = 1.2] (-0.13,5)--(0.13,5);
	\draw[color ={black},line width = 1.2] (-0.13,6)--(0.13,6);
	\draw[color ={black},line width = 1.2] (-0.13,7)--(0.13,7);
	\draw[color ={black},line width = 1.2] (-0.13,8)--(0.13,8);
	\draw[color ={black},line width = 1.2] (-0.13,9)--(0.13,9);
	\draw[color ={black},line width = 1.2] (-0.13,10)--(0.13,10);
	\draw[color ={black},line width = 1.2] (-0.13,11)--(0.13,11);
	\draw[color ={black},line width = 1.2] (-0.13,12)--(0.13,12);
	\draw[color ={black},line width = 1.2] (-0.13,-1)--(0.13,-1);
	\draw[color ={black},line width = 1.2] (-0.13,-2)--(0.13,-2);
	\draw[color ={black},line width = 1.2] (-0.13,-3)--(0.13,-3);
	\draw (3,-0.4) node[anchor = center] {\scriptsize \color{black}{$3$}};
	\draw (-3,-0.4) node[anchor = center] {\scriptsize \color{black}{$-3$}};
	\draw (-0.5,3.1) node[anchor = center] {\footnotesize \color{black}{$3$}};
	\draw (-0.5,6.1) node[anchor = center] {\footnotesize \color{black}{$6$}};
	\draw (-0.5,9.1) node[anchor = center] {\footnotesize \color{black}{$9$}};
	\draw (-0.5,-2.9) node[anchor = center] {\footnotesize \color{black}{$-3$}};
	\draw [color={black}] (-0.3,-0.3) node[anchor = center,scale=1, rotate = 0] {O};
	\draw (3,10) node[anchor = center] {\footnotesize \color{blue}{$\Large \mathcal{C}_f$}};
	
	\draw[color={blue},line width = 2] (-3,10)--(-2.8,8.84)--(-2.6,7.76)--(-2.4,6.76)--(-2.2,5.84)--(-2,5)--(-1.8,4.24)--(-1.6,3.56)--(-1.4,2.96)--(-1.2,2.44)--(-1,2)--(-0.8,1.64)--(-0.6,1.36)--(-0.4,1.16)--(-0.2,1.04)--(0,1)--(0.2,1.04)--(0.4,1.16)--(0.6,1.36)--(0.8,1.64)--(1,2)--(1.2,2.44)--(1.4,2.96)--(1.6,3.56)--(1.8,4.24)--(2,5)--(2.2,5.84)--(2.4,6.76)--(2.6,7.76)--(2.8,8.84)--(3,10);

\end{tikzpicture}\\
    \end{minipage}
    \begin{minipage}{0.5\linewidth}
    \begin{tikzpicture}[baseline, scale=0.8]

    \tikzset{
      point/.style={
        thick,
        draw,
        cross out,
        inner sep=0pt,
        minimum width=5pt,
        minimum height=5pt,
      },
    }
    \clip (-4.5,-6.5) rectangle (6.5,9.5);
    	
	\draw[color ={black},>=latex,->] (-3,0)--(3,0);
	\draw[color ={black},>=latex,->] (0,-5)--(0,8);
	\draw[color ={black},opacity = 0.4] (-3,1)--(3,1);
	\draw[color ={black},opacity = 0.4] (-3,2)--(3,2);
	\draw[color ={black},opacity = 0.4] (-3,3)--(3,3);
	\draw[color ={black},opacity = 0.4] (-3,4)--(3,4);
	\draw[color ={black},opacity = 0.4] (-3,5)--(3,5);
	\draw[color ={black},opacity = 0.4] (-3,6)--(3,6);
	\draw[color ={black},opacity = 0.4] (-3,7)--(3,7);
	\draw[color ={black},opacity = 0.4] (-3,8)--(3,8);
	\draw[color ={black},opacity = 0.4] (-3,-1)--(3,-1);
	\draw[color ={black},opacity = 0.4] (-3,-2)--(3,-2);
	\draw[color ={black},opacity = 0.4] (-3,-3)--(3,-3);
	\draw[color ={black},opacity = 0.4] (-3,-4)--(3,-4);
	\draw[color ={black},opacity = 0.4] (-3,-5)--(3,-5);
	\draw[color ={black},opacity = 0.4] (1,-5)--(1,8);
	\draw[color ={black},opacity = 0.4] (2,-5)--(2,8);
	\draw[color ={black},opacity = 0.4] (3,-5)--(3,8);
	\draw[color ={black},opacity = 0.4] (-1,-5)--(-1,8);
	\draw[color ={black},opacity = 0.4] (-2,-5)--(-2,8);
	\draw[color ={black},opacity = 0.4] (-3,-5)--(-3,8);
	\draw[color ={gray},opacity = 0.3] (-3,0)--(3,0);
	\draw[color ={gray},opacity = 0.3] (-3,0.5)--(3,0.5);
	\draw[color ={gray},opacity = 0.3] (-3,1.5)--(3,1.5);
	\draw[color ={gray},opacity = 0.3] (-3,2.5)--(3,2.5);
	\draw[color ={gray},opacity = 0.3] (-3,3.5)--(3,3.5);
	\draw[color ={gray},opacity = 0.3] (-3,4.5)--(3,4.5);
	\draw[color ={gray},opacity = 0.3] (-3,5.5)--(3,5.5);
	\draw[color ={gray},opacity = 0.3] (-3,6.5)--(3,6.5);
	\draw[color ={gray},opacity = 0.3] (-3,7.5)--(3,7.5);
	\draw[color ={gray},opacity = 0.3] (-3,0)--(3,0);
	\draw[color ={gray},opacity = 0.3] (-3,-0.5)--(3,-0.5);
	\draw[color ={gray},opacity = 0.3] (-3,-1.5)--(3,-1.5);
	\draw[color ={gray},opacity = 0.3] (-3,-2.5)--(3,-2.5);
	\draw[color ={gray},opacity = 0.3] (-3,-3.5)--(3,-3.5);
	\draw[color ={gray},opacity = 0.3] (-3,-4.5)--(3,-4.5);
	\draw[color ={gray},opacity = 0.3] (0,-5)--(0,8);
	\draw[color ={gray},opacity = 0.3] (0,-5)--(0,8);
	\draw[color ={black},line width = 1.2] (1,-0.13)--(1,0.13);
	\draw[color ={black},line width = 1.2] (2,-0.13)--(2,0.13);
	\draw[color ={black},line width = 1.2] (3,-0.13)--(3,0.13);
	\draw[color ={black},line width = 1.2] (-1,-0.13)--(-1,0.13);
	\draw[color ={black},line width = 1.2] (-2,-0.13)--(-2,0.13);
	\draw[color ={black},line width = 1.2] (-3,-0.13)--(-3,0.13);
	\draw[color ={black},line width = 1.2] (-0.13,1)--(0.13,1);
	\draw[color ={black},line width = 1.2] (-0.13,2)--(0.13,2);
	\draw[color ={black},line width = 1.2] (-0.13,3)--(0.13,3);
	\draw[color ={black},line width = 1.2] (-0.13,4)--(0.13,4);
	\draw[color ={black},line width = 1.2] (-0.13,5)--(0.13,5);
	\draw[color ={black},line width = 1.2] (-0.13,6)--(0.13,6);
	\draw[color ={black},line width = 1.2] (-0.13,7)--(0.13,7);
	\draw[color ={black},line width = 1.2] (-0.13,8)--(0.13,8);
	\draw[color ={black},line width = 1.2] (-0.13,9)--(0.13,9);
	\draw[color ={black},line width = 1.2] (-0.13,10)--(0.13,10);
	\draw[color ={black},line width = 1.2] (-0.13,11)--(0.13,11);
	\draw[color ={black},line width = 1.2] (-0.13,12)--(0.13,12);
	\draw[color ={black},line width = 1.2] (-0.13,-1)--(0.13,-1);
	\draw[color ={black},line width = 1.2] (-0.13,-2)--(0.13,-2);
	\draw[color ={black},line width = 1.2] (-0.13,-3)--(0.13,-3);
	\draw (3,-0.4) node[anchor = center] {\scriptsize \color{black}{$3$}};
	\draw (-3,-0.4) node[anchor = center] {\scriptsize \color{black}{$-3$}};
	\draw (-0.5,3.1) node[anchor = center] {\footnotesize \color{black}{$3$}};
	\draw (-0.5,6.1) node[anchor = center] {\footnotesize \color{black}{$6$}};
	\draw (-0.5,-2.9) node[anchor = center] {\footnotesize \color{black}{$-3$}};
	\draw [color={black}] (-0.3,-0.3) node[anchor = center,scale=1, rotate = 0] {O};
	\draw (3,6) node[anchor = center] {\footnotesize \color{red}{$\Large\mathcal{C}_f\prime$}};
	
	\draw[color={red},line width = 2] (-2.8,-5.6)--(-2.6,-5.2)--(-2.4,-4.8)--(-2.2,-4.4)--(-2,-4)--(-1.8,-3.6)--(-1.6,-3.2)--(-1.4,-2.8)--(-1.2,-2.4)--(-1,-2)--(-0.8,-1.6)--(-0.6,-1.2)--(-0.4,-0.8)--(-0.2,-0.4)--(0,0)--(0.2,0.4)--(0.4,0.8)--(0.6,1.2)--(0.8,1.6)--(1,2)--(1.2,2.4)--(1.4,2.8)--(1.6,3.2)--(1.8,3.6)--(2,4)--(2.2,4.4)--(2.4,4.8)--(2.6,5.2)--(2.8,5.6)--(3,6);

\end{tikzpicture}\\
    \end{minipage}
    
\end{exercice}

\begin{Solution}
 L'équation réduite de la tangente au point d'abscisse $3$ est  : $y=f'(3)(x-3)+f(3)$.\\
        On lit graphiquement $f(3)=10$ et $f'(3)=6$.\\
        L'équation réduite de la tangente est donc donnée par :
        $y=6(x-3)+10$, soit $y=6x-8$.
\end{Solution}

% @see : https://coopmaths.fr/alea?uuid=a1ba2&id=can1F14&n=1&d=10&alea=rkIp223232323234234567891011121314151617181920212223242526272829303132333435363738394041424344454647482&cd=1&cols=1
\needspace{10\baselineskip}
\begin{exercice}[BaremeTotal=false]


Soit $f$ la fonction définie sur $\mathbb{R}$ par : $f(x)= -2x^2+x+4$.\\
En calculant la limite du taux d'accroissement, déterminer $f'(-2)$.
\end{exercice}

\begin{Solution}
 $f$ est une fonction polynôme du second degré de la forme $f(x)=ax^2+bx+c$.\\
    La fonction dérivée est donnée par la somme des dérivées des fonctions $u$ et $v$ définies par $u(x)=-2x^2$ et $v(x)=x+4$.\\
     Comme $u'(x)=-4x$ et $v'(x)=1$, on obtient  $f'(x)=-4x+1$. \\
     Ainsi, $f'(-2)=-4\times (-2)+1=9$.
\end{Solution}

% @see : https://coopmaths.fr/alea?uuid=3c690&id=can1F13&n=1&d=10&alea=EUDu223232323234234567891011121314151617181920212223242526272829303132333435363738394041424344454647482&cd=1&cols=1
\needspace{10\baselineskip}
\begin{exercice}[BaremeTotal=false]


 Déterminer le coefficient directeur de la tangente à la courbe représentative de la fonction racine carrée au point d'abscisse $8$.
         
\end{exercice}

\begin{Solution}
 Le coefficient directeur de la tangente au point d'abscisse $a$ est donné par le nombre dérivé $f'(a)$.\\
        La fonction $f$ définie par $f(x)=\sqrt{x}$ a pour fonction dérivée la fonction $f'$ définie par $f'(x)=\dfrac{1}{2\sqrt{x}}$.\\Comme $f'(8)=\dfrac{1}{2\sqrt{8}}$, le coefficient directeur de la tangente au point d'abscisse $8$ est : $\dfrac{1}{2\sqrt{8}}$.
\end{Solution}

% @see : https://coopmaths.fr/alea?uuid=ebd89&id=1AN14-1&n=1&d=10&s=3&alea=ocYr223232323234234567891011121314151617181920212223242526272829303132333435363738394041424344454647482&cd=0&cols=1
\needspace{10\baselineskip}
\begin{exercice}[BaremeTotal=false]


 Donner la dérivée de la fonction $f$, dérivable pour tout $x\in\R$, définie par  $f(x)=\dfrac{-9x+8}{10}$.
\end{exercice}

\begin{Solution}
 L'expression de la dérivée de la fonction $f$ définie par $f(x)=\dfrac{-9x+8}{10}$ est : ${\color[HTML]{f15929}\boldsymbol{f'(x)=-\dfrac{9}{10}}}$.
\end{Solution}

% @see : https://coopmaths.fr/alea?uuid=ebd89&id=1AN14-1&n=1&d=10&s=4&alea=uAlP223232323234234567891011121314151617181920212223242526272829303132333435363738394041424344454647482&cd=0&cols=1
\needspace{10\baselineskip}
\begin{exercice}[BaremeTotal=false]


 Donner la dérivée de la fonction $f$, dérivable pour tout $x\in\R$, définie par  $f(x)=9x^{14}$.
\end{exercice}

\begin{Solution}
 L'expression de la dérivée de la fonction $f$ définie par $f(x)=9x^{14}$ est : ${\color[HTML]{f15929}\boldsymbol{f'(x)=126x^{13}}}$.
\end{Solution}

% @see : https://coopmaths.fr/alea?uuid=ebd89&id=1AN14-1&n=1&d=10&s=7&alea=vjN7223232323234234567891011121314151617181920212223242526272829303132333435363738394041424344454647482&cd=0&cols=1
\needspace{10\baselineskip}
\begin{exercice}[BaremeTotal=false]


 Donner la dérivée de la fonction $f$, dérivable pour tout $x\in\R^*$, définie par  $f(x)=\dfrac{-5}{9x}$.
\end{exercice}

\begin{Solution}
 L'expression de la dérivée de la fonction $f$ définie par $f(x)=\dfrac{-5}{9x}$ est : ${\color[HTML]{f15929}\boldsymbol{f'(x)=\dfrac{5}{9x^2}}}$.
\end{Solution}

% @see : https://coopmaths.fr/alea?uuid=a83c0&id=1AN14-4&n=1&d=10&s=1&alea=Njr3223232323234234567891011121314151617181920212223242526272829303132333435363738394041424344454647482&cd=0&cols=1
\needspace{10\baselineskip}
\begin{exercice}[BaremeTotal=false]


 Donner l'expression de la dérivée de la fonction $f$ définie sur $\R^{*}$ par $f(x)=\dfrac{9}{x}+2x$.\\
\end{exercice}

\begin{Solution}
 L'expression de la dérivée de la fonction $f$ définie par $f(x)=\dfrac{9}{x}+2x$ est :\\$f'(x)={\color[HTML]{f15929}\boldsymbol{-\dfrac{9}{x^2}+2}}$.\\
\end{Solution}

% @see : https://coopmaths.fr/alea?uuid=a83c0&id=1AN14-4&n=1&d=10&s=4&alea=4Vpz223232323234234567891011121314151617181920212223242526272829303132333435363738394041424344454647482&cd=1&cols=1
\needspace{10\baselineskip}
\begin{exercice}[BaremeTotal=false]


 Donner l'expression de la dérivée de la fonction $f$ définie sur $\R_+^{*}$ par $f(x)=\dfrac{4}{x}-3x^{2}-2\sqrt{x}$.\\
\end{exercice}

\begin{Solution}
 $(\dfrac{4}{x})^\prime=-\dfrac{4}{x^2}$\\$(-3x^{2})^\prime=-6x$\\$(-2\sqrt{x})^\prime=-\dfrac{2}{2\sqrt{x}}$\\L'expression de la dérivée de la fonction $f$ définie par $f(x)=\dfrac{4}{x}-3x^{2}-2\sqrt{x}$ est :\\$f'(x)={\color[HTML]{f15929}\boldsymbol{-\dfrac{4}{x^2}-6x-\dfrac{2}{2\sqrt{x}}}}$.\\
\end{Solution}

\end{Maquette}
\clearpage
\begin{Maquette}[DM]{Niveau= ,Classe=\getdata[3]\mydata\ (v3), Date = 06/01/25  ,Theme=Exercices}


% @see : https://coopmaths.fr/alea?uuid=4c8c7&id=1AN11-3&n=1&d=10&s=2&alea=nawO2232323232342345678910111213141516171819202122232425262728293031323334353637383940414243444546474823&cd=1&cols=2
\needspace{10\baselineskip}
\begin{exercice}[BaremeTotal=false]
\label{eleve:3}


 \begin{minipage}[t]{\linewidth} Soit $f$ une fonction dérivable sur $[-5;5]$ et $\mathcal{C}_f$ sa courbe représentative.\\On sait que  $f(4)=3~~$ et que $~~f'(4)=0$.\\Déterminer une équation de la tangente $(T)$ à la courbe $\mathcal{C}_f$ au point d'abscisse $4$,\\sans utiliser la formule de cours de l'équation de tangente. \end{minipage}

\end{exercice}

\begin{Solution}
 \begin{minipage}[t]{\linewidth} On sait que la tangente n'est pas une droite verticale, puisque la fonction est dérivable sur l'intervalle.\\On en déduit que la tangente $(T)$ au point d'abscisse $4$, admet une équation réduite de la forme :  $(T) : y=m x + p$. 

\medskip
$\bullet$ Détermination de $m$ :\\On sait que le nombre dérivé en $4$ est par définition, le coefficient directeur de la tangente au point d'abscisse $4$.\\Par conséquent, on a déjà : $m=f'(4)=0$.\\ On en déduit que  $(T) : y= 0 x + p$\\$\bullet$ Détermination de $p$ :\\ Pour cela, on utilise que si $f(4)=3~~$, alors le point $A$ de coordonnées $(4;3)$ appartient à $\mathcal{C}_f$ mais aussi à $(T)$.\\ On peut écrire $A(4;3) = \mathcal{C}_f \cap (T)$.\\ On remplace alors les coordonnées de $A(4;3)$ dans l'équation  $(T) : y= 0 x + p$.\\ $\begin{aligned} \phantom{\iff}&A(4;3)\in (T)\\ \iff& 3= 0 \times 4  + p\\ \iff& p=3 +0 \times 4  \\ \iff& p=3 +0   \\ \iff& p=3\\\end{aligned}$\\On peut conclure que : $(T) : {\color[HTML]{f15929}\boldsymbol{y=3}}$. \end{minipage}

\end{Solution}

% @see : https://coopmaths.fr/alea?uuid=0e984&id=can1F15&n=1&d=10&alea=1Alo2232323232342345678910111213141516171819202122232425262728293031323334353637383940414243444546474823&cd=1&cols=1
\needspace{10\baselineskip}
\begin{exercice}[BaremeTotal=false]


 La courbe représente une fonction $f$ et la droite est la tangente au point d'abscisse $0$.\\
        
        Déterminer $f'(0)$.\begin{center}\begin{tikzpicture}[baseline,scale = 0.5]

    \tikzset{
      point/.style={
        thick,
        draw,
        cross out,
        inner sep=0pt,
        minimum width=5pt,
        minimum height=5pt,
      },
    }
    \clip (-8,-3) rectangle (8,10);
    	
	\draw[color ={black},line width = 2,>=latex,->] (-8,0)--(8,0);
	\draw[color ={black},line width = 2,>=latex,->] (0,-3)--(0,10);
	\draw[color ={black},opacity = 0.5] (-8,1)--(8,1);
	\draw[color ={black},opacity = 0.5] (-8,2)--(8,2);
	\draw[color ={black},opacity = 0.5] (-8,3)--(8,3);
	\draw[color ={black},opacity = 0.5] (-8,4)--(8,4);
	\draw[color ={black},opacity = 0.5] (-8,5)--(8,5);
	\draw[color ={black},opacity = 0.5] (-8,6)--(8,6);
	\draw[color ={black},opacity = 0.5] (-8,7)--(8,7);
	\draw[color ={black},opacity = 0.5] (-8,8)--(8,8);
	\draw[color ={black},opacity = 0.5] (-8,9)--(8,9);
	\draw[color ={black},opacity = 0.5] (-8,10)--(8,10);
	\draw[color ={black},opacity = 0.5] (-8,-1)--(8,-1);
	\draw[color ={black},opacity = 0.5] (-8,-2)--(8,-2);
	\draw[color ={black},opacity = 0.5] (-8,-3)--(8,-3);
	\draw[color ={black},opacity = 0.5] (1,-3)--(1,10);
	\draw[color ={black},opacity = 0.5] (2,-3)--(2,10);
	\draw[color ={black},opacity = 0.5] (3,-3)--(3,10);
	\draw[color ={black},opacity = 0.5] (4,-3)--(4,10);
	\draw[color ={black},opacity = 0.5] (5,-3)--(5,10);
	\draw[color ={black},opacity = 0.5] (6,-3)--(6,10);
	\draw[color ={black},opacity = 0.5] (7,-3)--(7,10);
	\draw[color ={black},opacity = 0.5] (8,-3)--(8,10);
	\draw[color ={black},opacity = 0.5] (-1,-3)--(-1,10);
	\draw[color ={black},opacity = 0.5] (-2,-3)--(-2,10);
	\draw[color ={black},opacity = 0.5] (-3,-3)--(-3,10);
	\draw[color ={black},opacity = 0.5] (-4,-3)--(-4,10);
	\draw[color ={black},opacity = 0.5] (-5,-3)--(-5,10);
	\draw[color ={black},opacity = 0.5] (-6,-3)--(-6,10);
	\draw[color ={black},opacity = 0.5] (-7,-3)--(-7,10);
	\draw[color ={black},opacity = 0.5] (-8,-3)--(-8,10);
	\draw[color ={gray},opacity = 0.3] (-8,0)--(8,0);
	\draw[color ={gray},opacity = 0.3] (-8,0)--(8,0);
	\draw[color ={gray},opacity = 0.3] (0,-3)--(0,10);
	\draw[color ={gray},opacity = 0.3] (0,-3)--(0,10);
	\draw[color ={black},line width = 2] (2,-0.2)--(2,0.2);
	\draw[color ={black},line width = 2] (4,-0.2)--(4,0.2);
	\draw[color ={black},line width = 2] (6,-0.2)--(6,0.2);
	\draw[color ={black},line width = 2] (-2,-0.2)--(-2,0.2);
	\draw[color ={black},line width = 2] (-4,-0.2)--(-4,0.2);
	\draw[color ={black},line width = 2] (-6,-0.2)--(-6,0.2);
	\draw[color ={black},line width = 2] (-0.2,1)--(0.2,1);
	\draw[color ={black},line width = 2] (-0.2,2)--(0.2,2);
	\draw[color ={black},line width = 2] (-0.2,3)--(0.2,3);
	\draw[color ={black},line width = 2] (-0.2,4)--(0.2,4);
	\draw[color ={black},line width = 2] (-0.2,5)--(0.2,5);
	\draw[color ={black},line width = 2] (-0.2,6)--(0.2,6);
	\draw[color ={black},line width = 2] (-0.2,7)--(0.2,7);
	\draw[color ={black},line width = 2] (-0.2,8)--(0.2,8);
	\draw[color ={black},line width = 2] (-0.2,9)--(0.2,9);
	\draw[color ={black},line width = 2] (-0.2,-1)--(0.2,-1);
	\draw[color ={black},line width = 2] (-0.2,-2)--(0.2,-2);
	\draw (2,-0.4) node[anchor = center] {\scriptsize \color{black}{$1$}};
	\draw (4,-0.4) node[anchor = center] {\scriptsize \color{black}{$2$}};
	\draw (6,-0.4) node[anchor = center] {\scriptsize \color{black}{$3$}};
	\draw (-2,-0.4) node[anchor = center] {\scriptsize \color{black}{$-1$}};
	\draw (-4,-0.4) node[anchor = center] {\scriptsize \color{black}{$-2$}};
	\draw (-6,-0.4) node[anchor = center] {\scriptsize \color{black}{$-3$}};
	\draw (-0.8,2.1) node[anchor = center] {\footnotesize \color{black}{$2$}};
	\draw (-0.8,4.1) node[anchor = center] {\footnotesize \color{black}{$4$}};
	\draw (-0.8,6.1) node[anchor = center] {\footnotesize \color{black}{$6$}};
	\draw (-0.8,8.1) node[anchor = center] {\footnotesize \color{black}{$8$}};
	\draw [color={black}] (-0.3,-0.3) node[anchor = center,scale=1, rotate = 0] {O};
	
	\draw[color={blue},line width = 2] (-8,9)--(-7.6,7.84)--(-7.2,6.76)--(-6.8,5.76)--(-6.4,4.84)--(-6,4)--(-5.6,3.24)--(-5.2,2.56)--(-4.8,1.96)--(-4.4,1.44)--(-4,1)--(-3.6,0.64)--(-3.2,0.36)--(-2.8,0.16)--(-2.4,0.04)--(-2,0)--(-1.6,0.04)--(-1.2,0.16)--(-0.8,0.36)--(-0.4,0.64)--(0,1)--(0.4,1.44)--(0.8,1.96)--(1.2,2.56)--(1.6,3.24)--(2,4)--(2.4,4.84)--(2.8,5.76)--(3.2,6.76)--(3.6,7.84)--(4,9)--(4.4,10.24);
	
	\draw[color={red},line width = 2] (-4.8,-3.8)--(-4.4,-3.4)--(-4,-3)--(-3.6,-2.6)--(-3.2,-2.2)--(-2.8,-1.8)--(-2.4,-1.4)--(-2,-1)--(-1.6,-0.6)--(-1.2,-0.2)--(-0.8,0.2)--(-0.4,0.6)--(0,1)--(0.4,1.4)--(0.8,1.8)--(1.2,2.2)--(1.6,2.6)--(2,3)--(2.4,3.4)--(2.8,3.8)--(3.2,4.2)--(3.6,4.6)--(4,5)--(4.4,5.4)--(4.8,5.8)--(5.2,6.2)--(5.6,6.6)--(6,7)--(6.4,7.4)--(6.8,7.8)--(7.2,8.2)--(7.6,8.6)--(8,9);

\end{tikzpicture}\end{center}
\end{exercice}

\begin{Solution}
 $f'(0)$ est donné par le coefficient directeur de la tangente à la courbe au point d'abscisse $0$, soit $2$.
\end{Solution}

% @see : https://coopmaths.fr/alea?uuid=6f32d&id=can1F16&n=1&d=10&alea=dqSw2232323232342345678910111213141516171819202122232425262728293031323334353637383940414243444546474823&cd=1&cols=1
\needspace{10\baselineskip}
\newpage
\begin{exercice}[BaremeTotal=false]


 On donne les représentations graphiques d'une fonction et de sa dérivée.\\
        Donner l'équation réduite de la tangente à la courbe de $f$ en $x=2$. \\ \begin{minipage}{0.5\linewidth}
    \begin{tikzpicture}[baseline, scale=0.8]

    \tikzset{
      point/.style={
        thick,
        draw,
        cross out,
        inner sep=0pt,
        minimum width=5pt,
        minimum height=5pt,
      },
    }
    \clip (-4.5,-4.5) rectangle (5.75,13.5);
    	
	\draw[color ={black},>=latex,->] (-3,0)--(3,0);
	\draw[color ={black},>=latex,->] (0,-3)--(0,12);
	\draw[color ={black},opacity = 0.4] (-3,1)--(3,1);
	\draw[color ={black},opacity = 0.4] (-3,2)--(3,2);
	\draw[color ={black},opacity = 0.4] (-3,3)--(3,3);
	\draw[color ={black},opacity = 0.4] (-3,4)--(3,4);
	\draw[color ={black},opacity = 0.4] (-3,5)--(3,5);
	\draw[color ={black},opacity = 0.4] (-3,6)--(3,6);
	\draw[color ={black},opacity = 0.4] (-3,7)--(3,7);
	\draw[color ={black},opacity = 0.4] (-3,8)--(3,8);
	\draw[color ={black},opacity = 0.4] (-3,9)--(3,9);
	\draw[color ={black},opacity = 0.4] (-3,10)--(3,10);
	\draw[color ={black},opacity = 0.4] (-3,11)--(3,11);
	\draw[color ={black},opacity = 0.4] (-3,12)--(3,12);
	\draw[color ={black},opacity = 0.4] (-3,-1)--(3,-1);
	\draw[color ={black},opacity = 0.4] (-3,-2)--(3,-2);
	\draw[color ={black},opacity = 0.4] (-3,-3)--(3,-3);
	\draw[color ={black},opacity = 0.4] (1,-3)--(1,12);
	\draw[color ={black},opacity = 0.4] (2,-3)--(2,12);
	\draw[color ={black},opacity = 0.4] (3,-3)--(3,12);
	\draw[color ={black},opacity = 0.4] (-1,-3)--(-1,12);
	\draw[color ={black},opacity = 0.4] (-2,-3)--(-2,12);
	\draw[color ={black},opacity = 0.4] (-3,-3)--(-3,12);
	\draw[color ={gray},opacity = 0.3] (-3,0)--(3,0);
	\draw[color ={gray},opacity = 0.3] (-3,0.5)--(3,0.5);
	\draw[color ={gray},opacity = 0.3] (-3,1.5)--(3,1.5);
	\draw[color ={gray},opacity = 0.3] (-3,2.5)--(3,2.5);
	\draw[color ={gray},opacity = 0.3] (-3,3.5)--(3,3.5);
	\draw[color ={gray},opacity = 0.3] (-3,4.5)--(3,4.5);
	\draw[color ={gray},opacity = 0.3] (-3,5.5)--(3,5.5);
	\draw[color ={gray},opacity = 0.3] (-3,6.5)--(3,6.5);
	\draw[color ={gray},opacity = 0.3] (-3,7.5)--(3,7.5);
	\draw[color ={gray},opacity = 0.3] (-3,8.5)--(3,8.5);
	\draw[color ={gray},opacity = 0.3] (-3,9.5)--(3,9.5);
	\draw[color ={gray},opacity = 0.3] (-3,10.5)--(3,10.5);
	\draw[color ={gray},opacity = 0.3] (-3,11.5)--(3,11.5);
	\draw[color ={gray},opacity = 0.3] (-3,0)--(3,0);
	\draw[color ={gray},opacity = 0.3] (-3,-0.5)--(3,-0.5);
	\draw[color ={gray},opacity = 0.3] (-3,-1.5)--(3,-1.5);
	\draw[color ={gray},opacity = 0.3] (-3,-2.5)--(3,-2.5);
	\draw[color ={gray},opacity = 0.3] (0,-3)--(0,12);
	\draw[color ={gray},opacity = 0.3] (0,-3)--(0,12);
	\draw[color ={black},line width = 1.2] (1,-0.13)--(1,0.13);
	\draw[color ={black},line width = 1.2] (2,-0.13)--(2,0.13);
	\draw[color ={black},line width = 1.2] (3,-0.13)--(3,0.13);
	\draw[color ={black},line width = 1.2] (-1,-0.13)--(-1,0.13);
	\draw[color ={black},line width = 1.2] (-2,-0.13)--(-2,0.13);
	\draw[color ={black},line width = 1.2] (-3,-0.13)--(-3,0.13);
	\draw[color ={black},line width = 1.2] (-0.13,1)--(0.13,1);
	\draw[color ={black},line width = 1.2] (-0.13,2)--(0.13,2);
	\draw[color ={black},line width = 1.2] (-0.13,3)--(0.13,3);
	\draw[color ={black},line width = 1.2] (-0.13,4)--(0.13,4);
	\draw[color ={black},line width = 1.2] (-0.13,5)--(0.13,5);
	\draw[color ={black},line width = 1.2] (-0.13,6)--(0.13,6);
	\draw[color ={black},line width = 1.2] (-0.13,7)--(0.13,7);
	\draw[color ={black},line width = 1.2] (-0.13,8)--(0.13,8);
	\draw[color ={black},line width = 1.2] (-0.13,9)--(0.13,9);
	\draw[color ={black},line width = 1.2] (-0.13,10)--(0.13,10);
	\draw[color ={black},line width = 1.2] (-0.13,11)--(0.13,11);
	\draw[color ={black},line width = 1.2] (-0.13,12)--(0.13,12);
	\draw[color ={black},line width = 1.2] (-0.13,-1)--(0.13,-1);
	\draw[color ={black},line width = 1.2] (-0.13,-2)--(0.13,-2);
	\draw[color ={black},line width = 1.2] (-0.13,-3)--(0.13,-3);
	\draw (3,-0.4) node[anchor = center] {\scriptsize \color{black}{$3$}};
	\draw (-3,-0.4) node[anchor = center] {\scriptsize \color{black}{$-3$}};
	\draw (-0.5,3.1) node[anchor = center] {\footnotesize \color{black}{$3$}};
	\draw (-0.5,6.1) node[anchor = center] {\footnotesize \color{black}{$6$}};
	\draw (-0.5,9.1) node[anchor = center] {\footnotesize \color{black}{$9$}};
	\draw (-0.5,-2.9) node[anchor = center] {\footnotesize \color{black}{$-3$}};
	\draw [color={black}] (-0.3,-0.3) node[anchor = center,scale=1, rotate = 0] {O};
	\draw (3,10) node[anchor = center] {\footnotesize \color{blue}{$\Large \mathcal{C}_f$}};
	
	\draw[color={blue},line width = 2] (-2.6,12.96)--(-2.4,11.56)--(-2.2,10.24)--(-2,9)--(-1.8,7.84)--(-1.6,6.76)--(-1.4,5.76)--(-1.2,4.84)--(-1,4)--(-0.8,3.24)--(-0.6,2.56)--(-0.4,1.96)--(-0.2,1.44)--(0,1)--(0.2,0.64)--(0.4,0.36)--(0.6,0.16)--(0.8,0.04)--(1,0)--(1.2,0.04)--(1.4,0.16)--(1.6,0.36)--(1.8,0.64)--(2,1)--(2.2,1.44)--(2.4,1.96)--(2.6,2.56)--(2.8,3.24)--(3,4);

\end{tikzpicture}\\
    \end{minipage}
    \begin{minipage}{0.5\linewidth}
    \begin{tikzpicture}[baseline, scale=0.8]

    \tikzset{
      point/.style={
        thick,
        draw,
        cross out,
        inner sep=0pt,
        minimum width=5pt,
        minimum height=5pt,
      },
    }
    \clip (-4.5,-6.5) rectangle (6.5,9.5);
    	
	\draw[color ={black},>=latex,->] (-3,0)--(3,0);
	\draw[color ={black},>=latex,->] (0,-5)--(0,8);
	\draw[color ={black},opacity = 0.4] (-3,1)--(3,1);
	\draw[color ={black},opacity = 0.4] (-3,2)--(3,2);
	\draw[color ={black},opacity = 0.4] (-3,3)--(3,3);
	\draw[color ={black},opacity = 0.4] (-3,4)--(3,4);
	\draw[color ={black},opacity = 0.4] (-3,5)--(3,5);
	\draw[color ={black},opacity = 0.4] (-3,6)--(3,6);
	\draw[color ={black},opacity = 0.4] (-3,7)--(3,7);
	\draw[color ={black},opacity = 0.4] (-3,8)--(3,8);
	\draw[color ={black},opacity = 0.4] (-3,-1)--(3,-1);
	\draw[color ={black},opacity = 0.4] (-3,-2)--(3,-2);
	\draw[color ={black},opacity = 0.4] (-3,-3)--(3,-3);
	\draw[color ={black},opacity = 0.4] (-3,-4)--(3,-4);
	\draw[color ={black},opacity = 0.4] (-3,-5)--(3,-5);
	\draw[color ={black},opacity = 0.4] (1,-5)--(1,8);
	\draw[color ={black},opacity = 0.4] (2,-5)--(2,8);
	\draw[color ={black},opacity = 0.4] (3,-5)--(3,8);
	\draw[color ={black},opacity = 0.4] (-1,-5)--(-1,8);
	\draw[color ={black},opacity = 0.4] (-2,-5)--(-2,8);
	\draw[color ={black},opacity = 0.4] (-3,-5)--(-3,8);
	\draw[color ={gray},opacity = 0.3] (-3,0)--(3,0);
	\draw[color ={gray},opacity = 0.3] (-3,0.5)--(3,0.5);
	\draw[color ={gray},opacity = 0.3] (-3,1.5)--(3,1.5);
	\draw[color ={gray},opacity = 0.3] (-3,2.5)--(3,2.5);
	\draw[color ={gray},opacity = 0.3] (-3,3.5)--(3,3.5);
	\draw[color ={gray},opacity = 0.3] (-3,4.5)--(3,4.5);
	\draw[color ={gray},opacity = 0.3] (-3,5.5)--(3,5.5);
	\draw[color ={gray},opacity = 0.3] (-3,6.5)--(3,6.5);
	\draw[color ={gray},opacity = 0.3] (-3,7.5)--(3,7.5);
	\draw[color ={gray},opacity = 0.3] (-3,0)--(3,0);
	\draw[color ={gray},opacity = 0.3] (-3,-0.5)--(3,-0.5);
	\draw[color ={gray},opacity = 0.3] (-3,-1.5)--(3,-1.5);
	\draw[color ={gray},opacity = 0.3] (-3,-2.5)--(3,-2.5);
	\draw[color ={gray},opacity = 0.3] (-3,-3.5)--(3,-3.5);
	\draw[color ={gray},opacity = 0.3] (-3,-4.5)--(3,-4.5);
	\draw[color ={gray},opacity = 0.3] (0,-5)--(0,8);
	\draw[color ={gray},opacity = 0.3] (0,-5)--(0,8);
	\draw[color ={black},line width = 1.2] (1,-0.13)--(1,0.13);
	\draw[color ={black},line width = 1.2] (2,-0.13)--(2,0.13);
	\draw[color ={black},line width = 1.2] (3,-0.13)--(3,0.13);
	\draw[color ={black},line width = 1.2] (-1,-0.13)--(-1,0.13);
	\draw[color ={black},line width = 1.2] (-2,-0.13)--(-2,0.13);
	\draw[color ={black},line width = 1.2] (-3,-0.13)--(-3,0.13);
	\draw[color ={black},line width = 1.2] (-0.13,1)--(0.13,1);
	\draw[color ={black},line width = 1.2] (-0.13,2)--(0.13,2);
	\draw[color ={black},line width = 1.2] (-0.13,3)--(0.13,3);
	\draw[color ={black},line width = 1.2] (-0.13,4)--(0.13,4);
	\draw[color ={black},line width = 1.2] (-0.13,5)--(0.13,5);
	\draw[color ={black},line width = 1.2] (-0.13,6)--(0.13,6);
	\draw[color ={black},line width = 1.2] (-0.13,7)--(0.13,7);
	\draw[color ={black},line width = 1.2] (-0.13,8)--(0.13,8);
	\draw[color ={black},line width = 1.2] (-0.13,9)--(0.13,9);
	\draw[color ={black},line width = 1.2] (-0.13,10)--(0.13,10);
	\draw[color ={black},line width = 1.2] (-0.13,11)--(0.13,11);
	\draw[color ={black},line width = 1.2] (-0.13,12)--(0.13,12);
	\draw[color ={black},line width = 1.2] (-0.13,-1)--(0.13,-1);
	\draw[color ={black},line width = 1.2] (-0.13,-2)--(0.13,-2);
	\draw[color ={black},line width = 1.2] (-0.13,-3)--(0.13,-3);
	\draw (3,-0.4) node[anchor = center] {\scriptsize \color{black}{$3$}};
	\draw (-3,-0.4) node[anchor = center] {\scriptsize \color{black}{$-3$}};
	\draw (-0.5,3.1) node[anchor = center] {\footnotesize \color{black}{$3$}};
	\draw (-0.5,6.1) node[anchor = center] {\footnotesize \color{black}{$6$}};
	\draw (-0.5,-2.9) node[anchor = center] {\footnotesize \color{black}{$-3$}};
	\draw [color={black}] (-0.3,-0.3) node[anchor = center,scale=1, rotate = 0] {O};
	\draw (3,6) node[anchor = center] {\footnotesize \color{red}{$\Large\mathcal{C}_f\prime$}};
	
	\draw[color={red},line width = 2] (-2,-6)--(-1.8,-5.6)--(-1.6,-5.2)--(-1.4,-4.8)--(-1.2,-4.4)--(-1,-4)--(-0.8,-3.6)--(-0.6,-3.2)--(-0.4,-2.8)--(-0.2,-2.4)--(0,-2)--(0.2,-1.6)--(0.4,-1.2)--(0.6,-0.8)--(0.8,-0.4)--(1,0)--(1.2,0.4)--(1.4,0.8)--(1.6,1.2)--(1.8,1.6)--(2,2)--(2.2,2.4)--(2.4,2.8)--(2.6,3.2)--(2.8,3.6)--(3,4);

\end{tikzpicture}\\
    \end{minipage}
    
\end{exercice}

\begin{Solution}
 L'équation réduite de la tangente au point d'abscisse $2$ est  : $y=f'(2)(x-2)+f(2)$.\\
        On lit graphiquement $f(2)=1$ et $f'(2)=2$.\\
        L'équation réduite de la tangente est donc donnée par :
        $y=2(x-2)+1$, soit $y=2x-3$.
\end{Solution}

% @see : https://coopmaths.fr/alea?uuid=a1ba2&id=can1F14&n=1&d=10&alea=rkIp2232323232342345678910111213141516171819202122232425262728293031323334353637383940414243444546474823&cd=1&cols=1
\needspace{10\baselineskip}
\begin{exercice}[BaremeTotal=false]


Soit $f$ la fonction définie sur $\mathbb{R}$ par :
          $f(x)= 5-5x^2$.\\
En calculant la limite du taux d'accroissement, déterminer $f'(2)$.
\end{exercice}

\begin{Solution}
 $f$ est une fonction polynôme du second degré de la forme $f(x)=ax^2+b$.\\
    La fonction dérivée est donnée par la somme des dérivées des fonctions $u$ et $v$ définies par $u(x)=-5x^2$ et $v(x)=5$.\\
     Comme $u'(x)=-10x$ et $v'(x)=0$, on obtient  $f'(x)=-10x$. \\
     Ainsi, $f'(2)=-10\times 2=-20$.
\end{Solution}

% @see : https://coopmaths.fr/alea?uuid=3c690&id=can1F13&n=1&d=10&alea=EUDu2232323232342345678910111213141516171819202122232425262728293031323334353637383940414243444546474823&cd=1&cols=1
\needspace{10\baselineskip}
\begin{exercice}[BaremeTotal=false]


 Déterminer le coefficient directeur de la tangente à la courbe représentative de la fonction cube au point d'abscisse $3$.
                        
\end{exercice}

\begin{Solution}
 Le coefficient directeur de la tangente au point d'abscisse $a$ est donné par le nombre dérivé $f'(a)$.\\
        La fonction $f$ définie par $f(x)=x^3$ a pour fonction dérivée la fonction $f'$ définie par $f'(x)=3x^2$.\\
        Comme $f'(3)=3\times 3^2=27$, le coefficient directeur de la tangente au point d'abscisse $3$ est : $27$. 
\end{Solution}

% @see : https://coopmaths.fr/alea?uuid=ebd89&id=1AN14-1&n=1&d=10&s=3&alea=ocYr2232323232342345678910111213141516171819202122232425262728293031323334353637383940414243444546474823&cd=0&cols=1
\needspace{10\baselineskip}
\begin{exercice}[BaremeTotal=false]


 Donner la dérivée de la fonction $f$, dérivable pour tout $x\in\R$, définie par  $f(x)=\dfrac{x+6}{5}$.
\end{exercice}

\begin{Solution}
 L'expression de la dérivée de la fonction $f$ définie par $f(x)=\dfrac{x+6}{5}$ est : ${\color[HTML]{f15929}\boldsymbol{f'(x)=\dfrac{1}{5}}}$.
\end{Solution}

% @see : https://coopmaths.fr/alea?uuid=ebd89&id=1AN14-1&n=1&d=10&s=4&alea=uAlP2232323232342345678910111213141516171819202122232425262728293031323334353637383940414243444546474823&cd=0&cols=1
\needspace{10\baselineskip}
\begin{exercice}[BaremeTotal=false]


 Donner la dérivée de la fonction $f$, dérivable pour tout $x\in\R$, définie par  $f(x)=10x^{6}$.
\end{exercice}

\begin{Solution}
 L'expression de la dérivée de la fonction $f$ définie par $f(x)=10x^{6}$ est : ${\color[HTML]{f15929}\boldsymbol{f'(x)=60x^{5}}}$.
\end{Solution}

% @see : https://coopmaths.fr/alea?uuid=ebd89&id=1AN14-1&n=1&d=10&s=7&alea=vjN72232323232342345678910111213141516171819202122232425262728293031323334353637383940414243444546474823&cd=0&cols=1
\needspace{10\baselineskip}
\begin{exercice}[BaremeTotal=false]


 Donner la dérivée de la fonction $f$, dérivable pour tout $x\in\R^*$, définie par  $f(x)=\dfrac{2}{9x}$.
\end{exercice}

\begin{Solution}
 L'expression de la dérivée de la fonction $f$ définie par $f(x)=\dfrac{2}{9x}$ est : ${\color[HTML]{f15929}\boldsymbol{f'(x)=-\dfrac{2}{9x^2}}}$.
\end{Solution}

% @see : https://coopmaths.fr/alea?uuid=a83c0&id=1AN14-4&n=1&d=10&s=1&alea=Njr32232323232342345678910111213141516171819202122232425262728293031323334353637383940414243444546474823&cd=0&cols=1
\needspace{10\baselineskip}
\begin{exercice}[BaremeTotal=false]


 Donner l'expression de la dérivée de la fonction $f$ définie sur $\R^{*}$ par $f(x)=3x^{2}-5x-\dfrac{4}{x}$.\\
\end{exercice}

\begin{Solution}
 L'expression de la dérivée de la fonction $f$ définie par $f(x)=3x^{2}-5x-\dfrac{4}{x}$ est :\\$f'(x)={\color[HTML]{f15929}\boldsymbol{6x-5+\dfrac{4}{x^2}}}$.\\
\end{Solution}

% @see : https://coopmaths.fr/alea?uuid=a83c0&id=1AN14-4&n=1&d=10&s=4&alea=4Vpz2232323232342345678910111213141516171819202122232425262728293031323334353637383940414243444546474823&cd=1&cols=1
\needspace{10\baselineskip}
\begin{exercice}[BaremeTotal=false]


 Donner l'expression de la dérivée de la fonction $f$ définie sur $\R_+^{*}$ par $f(x)=3x^{2}-2\sqrt{x}+\dfrac{5}{x}$.\\
\end{exercice}

\begin{Solution}
 $(3x^{2})^\prime=6x$\\$(-2\sqrt{x})^\prime=-\dfrac{2}{2\sqrt{x}}$\\$(\dfrac{5}{x})^\prime=-\dfrac{5}{x^2}$\\L'expression de la dérivée de la fonction $f$ définie par $f(x)=3x^{2}-2\sqrt{x}+\dfrac{5}{x}$ est :\\$f'(x)={\color[HTML]{f15929}\boldsymbol{6x-\dfrac{2}{2\sqrt{x}}-\dfrac{5}{x^2}}}$.\\
\end{Solution}

\end{Maquette}
\clearpage
\begin{Maquette}[DM]{Niveau= ,Classe=\getdata[4]\mydata\ (v4), Date = 06/01/25  ,Theme=Exercices}


% @see : https://coopmaths.fr/alea?uuid=4c8c7&id=1AN11-3&n=1&d=10&s=2&alea=nawO22323232323423456789101112131415161718192021222324252627282930313233343536373839404142434445464748234&cd=1&cols=2
\needspace{10\baselineskip}
\begin{exercice}[BaremeTotal=false]
\label{eleve:4}


 \begin{minipage}[t]{\linewidth} Soit $f$ une fonction dérivable sur $[-5;5]$ et $\mathcal{C}_f$ sa courbe représentative.\\On sait que  $f(4)=1~~$ et que $~~f'(4)=-1$.\\Déterminer une équation de la tangente $(T)$ à la courbe $\mathcal{C}_f$ au point d'abscisse $4$,\\sans utiliser la formule de cours de l'équation de tangente. \end{minipage}

\end{exercice}

\begin{Solution}
 \begin{minipage}[t]{\linewidth} On sait que la tangente n'est pas une droite verticale, puisque la fonction est dérivable sur l'intervalle.\\On en déduit que la tangente $(T)$ au point d'abscisse $4$, admet une équation réduite de la forme :  $(T) : y=m x + p$. 

\medskip
$\bullet$ Détermination de $m$ :\\On sait que le nombre dérivé en $4$ est par définition, le coefficient directeur de la tangente au point d'abscisse $4$.\\Par conséquent, on a déjà : $m=f'(4)=-1$.\\ On en déduit que  $(T) : y= -1 x + p$\\$\bullet$ Détermination de $p$ :\\ Pour cela, on utilise que si $f(4)=1~~$, alors le point $A$ de coordonnées $(4;1)$ appartient à $\mathcal{C}_f$ mais aussi à $(T)$.\\ On peut écrire $A(4;1) = \mathcal{C}_f \cap (T)$.\\ On remplace alors les coordonnées de $A(4;1)$ dans l'équation  $(T) : y= -1 x + p$.\\ $\begin{aligned} \phantom{\iff}&A(4;1)\in (T)\\ \iff& 1= -1 \times 4  + p\\ \iff& p=1 + \times 4  \\ \iff& p=1 +4   \\ \iff& p=5\\\end{aligned}$\\On peut conclure que : $(T) : {\color[HTML]{f15929}\boldsymbol{y=-x+5}}$. \end{minipage}

\end{Solution}

% @see : https://coopmaths.fr/alea?uuid=0e984&id=can1F15&n=1&d=10&alea=1Alo22323232323423456789101112131415161718192021222324252627282930313233343536373839404142434445464748234&cd=1&cols=1
\needspace{10\baselineskip}
\begin{exercice}[BaremeTotal=false]


 La courbe représente une fonction $f$ et la droite est la tangente au point d'abscisse $-1$.\\
        Déterminer $f'(-1)$.\begin{center}\begin{tikzpicture}[baseline,scale = 0.5]

    \tikzset{
      point/.style={
        thick,
        draw,
        cross out,
        inner sep=0pt,
        minimum width=5pt,
        minimum height=5pt,
      },
    }
    \clip (-8,-10) rectangle (8,3);
    	
	\draw[color ={black},line width = 2,>=latex,->] (-8,0)--(8,0);
	\draw[color ={black},line width = 2,>=latex,->] (0,-10)--(0,3);
	\draw[color ={black},opacity = 0.5] (-8,1)--(8,1);
	\draw[color ={black},opacity = 0.5] (-8,2)--(8,2);
	\draw[color ={black},opacity = 0.5] (-8,3)--(8,3);
	\draw[color ={black},opacity = 0.5] (-8,-1)--(8,-1);
	\draw[color ={black},opacity = 0.5] (-8,-2)--(8,-2);
	\draw[color ={black},opacity = 0.5] (-8,-3)--(8,-3);
	\draw[color ={black},opacity = 0.5] (-8,-4)--(8,-4);
	\draw[color ={black},opacity = 0.5] (-8,-5)--(8,-5);
	\draw[color ={black},opacity = 0.5] (-8,-6)--(8,-6);
	\draw[color ={black},opacity = 0.5] (-8,-7)--(8,-7);
	\draw[color ={black},opacity = 0.5] (-8,-8)--(8,-8);
	\draw[color ={black},opacity = 0.5] (-8,-9)--(8,-9);
	\draw[color ={black},opacity = 0.5] (-8,-10)--(8,-10);
	\draw[color ={black},opacity = 0.5] (1,-10)--(1,3);
	\draw[color ={black},opacity = 0.5] (2,-10)--(2,3);
	\draw[color ={black},opacity = 0.5] (3,-10)--(3,3);
	\draw[color ={black},opacity = 0.5] (4,-10)--(4,3);
	\draw[color ={black},opacity = 0.5] (5,-10)--(5,3);
	\draw[color ={black},opacity = 0.5] (6,-10)--(6,3);
	\draw[color ={black},opacity = 0.5] (7,-10)--(7,3);
	\draw[color ={black},opacity = 0.5] (8,-10)--(8,3);
	\draw[color ={black},opacity = 0.5] (-1,-10)--(-1,3);
	\draw[color ={black},opacity = 0.5] (-2,-10)--(-2,3);
	\draw[color ={black},opacity = 0.5] (-3,-10)--(-3,3);
	\draw[color ={black},opacity = 0.5] (-4,-10)--(-4,3);
	\draw[color ={black},opacity = 0.5] (-5,-10)--(-5,3);
	\draw[color ={black},opacity = 0.5] (-6,-10)--(-6,3);
	\draw[color ={black},opacity = 0.5] (-7,-10)--(-7,3);
	\draw[color ={black},opacity = 0.5] (-8,-10)--(-8,3);
	\draw[color ={gray},opacity = 0.3] (-8,0)--(8,0);
	\draw[color ={gray},opacity = 0.3] (-8,0)--(8,0);
	\draw[color ={gray},opacity = 0.3] (-8,-4)--(8,-4);
	\draw[color ={gray},opacity = 0.3] (-8,-5)--(8,-5);
	\draw[color ={gray},opacity = 0.3] (-8,-6)--(8,-6);
	\draw[color ={gray},opacity = 0.3] (-8,-7)--(8,-7);
	\draw[color ={gray},opacity = 0.3] (-8,-8)--(8,-8);
	\draw[color ={gray},opacity = 0.3] (-8,-9)--(8,-9);
	\draw[color ={gray},opacity = 0.3] (-8,-10)--(8,-10);
	\draw[color ={gray},opacity = 0.3] (0,-10)--(0,3);
	\draw[color ={gray},opacity = 0.3] (0,-10)--(0,3);
	\draw[color ={black},line width = 2] (2,-0.2)--(2,0.2);
	\draw[color ={black},line width = 2] (4,-0.2)--(4,0.2);
	\draw[color ={black},line width = 2] (6,-0.2)--(6,0.2);
	\draw[color ={black},line width = 2] (-2,-0.2)--(-2,0.2);
	\draw[color ={black},line width = 2] (-4,-0.2)--(-4,0.2);
	\draw[color ={black},line width = 2] (-6,-0.2)--(-6,0.2);
	\draw[color ={black},line width = 2] (-0.2,1)--(0.2,1);
	\draw[color ={black},line width = 2] (-0.2,2)--(0.2,2);
	\draw[color ={black},line width = 2] (-0.2,-1)--(0.2,-1);
	\draw[color ={black},line width = 2] (-0.2,-2)--(0.2,-2);
	\draw[color ={black},line width = 2] (-0.2,-3)--(0.2,-3);
	\draw[color ={black},line width = 2] (-0.2,-4)--(0.2,-4);
	\draw[color ={black},line width = 2] (-0.2,-5)--(0.2,-5);
	\draw[color ={black},line width = 2] (-0.2,-6)--(0.2,-6);
	\draw[color ={black},line width = 2] (-0.2,-7)--(0.2,-7);
	\draw[color ={black},line width = 2] (-0.2,-8)--(0.2,-8);
	\draw[color ={black},line width = 2] (-0.2,-9)--(0.2,-9);
	\draw (2,-0.4) node[anchor = center] {\scriptsize \color{black}{$1$}};
	\draw (4,-0.4) node[anchor = center] {\scriptsize \color{black}{$2$}};
	\draw (6,-0.4) node[anchor = center] {\scriptsize \color{black}{$3$}};
	\draw (-2,-0.4) node[anchor = center] {\scriptsize \color{black}{$-1$}};
	\draw (-4,-0.4) node[anchor = center] {\scriptsize \color{black}{$-2$}};
	\draw (-6,-0.4) node[anchor = center] {\scriptsize \color{black}{$-3$}};
	\draw (-0.8,2.1) node[anchor = center] {\footnotesize \color{black}{$2$}};
	\draw (-0.8,-1.9) node[anchor = center] {\footnotesize \color{black}{$-2$}};
	\draw (-0.8,-3.9) node[anchor = center] {\footnotesize \color{black}{$-4$}};
	\draw (-0.8,-5.9) node[anchor = center] {\footnotesize \color{black}{$-6$}};
	\draw (-0.8,-7.9) node[anchor = center] {\footnotesize \color{black}{$-8$}};
	\draw [color={black}] (-0.3,-0.3) node[anchor = center,scale=1, rotate = 0] {O};
	
	\draw[color={blue},line width = 2] (-4.8,-10.52)--(-4.4,-8.68)--(-4,-7)--(-3.6,-5.48)--(-3.2,-4.12)--(-2.8,-2.92)--(-2.4,-1.88)--(-2,-1)--(-1.6,-0.28)--(-1.2,0.28)--(-0.8,0.68)--(-0.4,0.92)--(0,1)--(0.4,0.92)--(0.8,0.68)--(1.2,0.28)--(1.6,-0.28)--(2,-1)--(2.4,-1.88)--(2.8,-2.92)--(3.2,-4.12)--(3.6,-5.48)--(4,-7)--(4.4,-8.68)--(4.8,-10.52);
	
	\draw[color={red},line width = 2] (-6.8,-10.6)--(-6.4,-9.8)--(-6,-9)--(-5.6,-8.2)--(-5.2,-7.4)--(-4.8,-6.6)--(-4.4,-5.8)--(-4,-5)--(-3.6,-4.2)--(-3.2,-3.4)--(-2.8,-2.6)--(-2.4,-1.8)--(-2,-1)--(-1.6,-0.2)--(-1.2,0.6)--(-0.8,1.4)--(-0.4,2.2)--(0,3)--(0.4,3.8);

\end{tikzpicture}\end{center}
\end{exercice}

\begin{Solution}
 $f'(-1)$ est donné par le coefficient directeur de la tangente à la courbe au point d'abscisse $-1$, soit $4$.
\end{Solution}

% @see : https://coopmaths.fr/alea?uuid=6f32d&id=can1F16&n=1&d=10&alea=dqSw22323232323423456789101112131415161718192021222324252627282930313233343536373839404142434445464748234&cd=1&cols=1
\needspace{10\baselineskip}
\newpage
\begin{exercice}[BaremeTotal=false]


 On donne les représentations graphiques d'une fonction et de sa dérivée.\\
        Donner l'équation réduite de la tangente à la courbe de $f$ en $x=1$. \\ \begin{minipage}{0.5\linewidth}
    \begin{tikzpicture}[baseline, scale=0.8]

    \tikzset{
      point/.style={
        thick,
        draw,
        cross out,
        inner sep=0pt,
        minimum width=5pt,
        minimum height=5pt,
      },
    }
    \clip (-4.5,-4.5) rectangle (5.75,13.5);
    	
	\draw[color ={black},>=latex,->] (-3,0)--(3,0);
	\draw[color ={black},>=latex,->] (0,-3)--(0,12);
	\draw[color ={black},opacity = 0.4] (-3,1)--(3,1);
	\draw[color ={black},opacity = 0.4] (-3,2)--(3,2);
	\draw[color ={black},opacity = 0.4] (-3,3)--(3,3);
	\draw[color ={black},opacity = 0.4] (-3,4)--(3,4);
	\draw[color ={black},opacity = 0.4] (-3,5)--(3,5);
	\draw[color ={black},opacity = 0.4] (-3,6)--(3,6);
	\draw[color ={black},opacity = 0.4] (-3,7)--(3,7);
	\draw[color ={black},opacity = 0.4] (-3,8)--(3,8);
	\draw[color ={black},opacity = 0.4] (-3,9)--(3,9);
	\draw[color ={black},opacity = 0.4] (-3,10)--(3,10);
	\draw[color ={black},opacity = 0.4] (-3,11)--(3,11);
	\draw[color ={black},opacity = 0.4] (-3,12)--(3,12);
	\draw[color ={black},opacity = 0.4] (-3,-1)--(3,-1);
	\draw[color ={black},opacity = 0.4] (-3,-2)--(3,-2);
	\draw[color ={black},opacity = 0.4] (-3,-3)--(3,-3);
	\draw[color ={black},opacity = 0.4] (1,-3)--(1,12);
	\draw[color ={black},opacity = 0.4] (2,-3)--(2,12);
	\draw[color ={black},opacity = 0.4] (3,-3)--(3,12);
	\draw[color ={black},opacity = 0.4] (-1,-3)--(-1,12);
	\draw[color ={black},opacity = 0.4] (-2,-3)--(-2,12);
	\draw[color ={black},opacity = 0.4] (-3,-3)--(-3,12);
	\draw[color ={gray},opacity = 0.3] (-3,0)--(3,0);
	\draw[color ={gray},opacity = 0.3] (-3,0.5)--(3,0.5);
	\draw[color ={gray},opacity = 0.3] (-3,1.5)--(3,1.5);
	\draw[color ={gray},opacity = 0.3] (-3,2.5)--(3,2.5);
	\draw[color ={gray},opacity = 0.3] (-3,3.5)--(3,3.5);
	\draw[color ={gray},opacity = 0.3] (-3,4.5)--(3,4.5);
	\draw[color ={gray},opacity = 0.3] (-3,5.5)--(3,5.5);
	\draw[color ={gray},opacity = 0.3] (-3,6.5)--(3,6.5);
	\draw[color ={gray},opacity = 0.3] (-3,7.5)--(3,7.5);
	\draw[color ={gray},opacity = 0.3] (-3,8.5)--(3,8.5);
	\draw[color ={gray},opacity = 0.3] (-3,9.5)--(3,9.5);
	\draw[color ={gray},opacity = 0.3] (-3,10.5)--(3,10.5);
	\draw[color ={gray},opacity = 0.3] (-3,11.5)--(3,11.5);
	\draw[color ={gray},opacity = 0.3] (-3,0)--(3,0);
	\draw[color ={gray},opacity = 0.3] (-3,-0.5)--(3,-0.5);
	\draw[color ={gray},opacity = 0.3] (-3,-1.5)--(3,-1.5);
	\draw[color ={gray},opacity = 0.3] (-3,-2.5)--(3,-2.5);
	\draw[color ={gray},opacity = 0.3] (0,-3)--(0,12);
	\draw[color ={gray},opacity = 0.3] (0,-3)--(0,12);
	\draw[color ={black},line width = 1.2] (1,-0.13)--(1,0.13);
	\draw[color ={black},line width = 1.2] (2,-0.13)--(2,0.13);
	\draw[color ={black},line width = 1.2] (3,-0.13)--(3,0.13);
	\draw[color ={black},line width = 1.2] (-1,-0.13)--(-1,0.13);
	\draw[color ={black},line width = 1.2] (-2,-0.13)--(-2,0.13);
	\draw[color ={black},line width = 1.2] (-3,-0.13)--(-3,0.13);
	\draw[color ={black},line width = 1.2] (-0.13,1)--(0.13,1);
	\draw[color ={black},line width = 1.2] (-0.13,2)--(0.13,2);
	\draw[color ={black},line width = 1.2] (-0.13,3)--(0.13,3);
	\draw[color ={black},line width = 1.2] (-0.13,4)--(0.13,4);
	\draw[color ={black},line width = 1.2] (-0.13,5)--(0.13,5);
	\draw[color ={black},line width = 1.2] (-0.13,6)--(0.13,6);
	\draw[color ={black},line width = 1.2] (-0.13,7)--(0.13,7);
	\draw[color ={black},line width = 1.2] (-0.13,8)--(0.13,8);
	\draw[color ={black},line width = 1.2] (-0.13,9)--(0.13,9);
	\draw[color ={black},line width = 1.2] (-0.13,10)--(0.13,10);
	\draw[color ={black},line width = 1.2] (-0.13,11)--(0.13,11);
	\draw[color ={black},line width = 1.2] (-0.13,12)--(0.13,12);
	\draw[color ={black},line width = 1.2] (-0.13,-1)--(0.13,-1);
	\draw[color ={black},line width = 1.2] (-0.13,-2)--(0.13,-2);
	\draw[color ={black},line width = 1.2] (-0.13,-3)--(0.13,-3);
	\draw (3,-0.4) node[anchor = center] {\scriptsize \color{black}{$3$}};
	\draw (-3,-0.4) node[anchor = center] {\scriptsize \color{black}{$-3$}};
	\draw (-0.5,3.1) node[anchor = center] {\footnotesize \color{black}{$3$}};
	\draw (-0.5,6.1) node[anchor = center] {\footnotesize \color{black}{$6$}};
	\draw (-0.5,9.1) node[anchor = center] {\footnotesize \color{black}{$9$}};
	\draw (-0.5,-2.9) node[anchor = center] {\footnotesize \color{black}{$-3$}};
	\draw [color={black}] (-0.3,-0.3) node[anchor = center,scale=1, rotate = 0] {O};
	\draw (3,10) node[anchor = center] {\footnotesize \color{blue}{$\Large \mathcal{C}_f$}};
	
	\draw[color={blue},line width = 2] (-3,-1)--(-2.8,-1.36)--(-2.6,-1.64)--(-2.4,-1.84)--(-2.2,-1.96)--(-2,-2)--(-1.8,-1.96)--(-1.6,-1.84)--(-1.4,-1.64)--(-1.2,-1.36)--(-1,-1)--(-0.8,-0.56)--(-0.6,-0.04)--(-0.4,0.56)--(-0.2,1.24)--(0,2)--(0.2,2.84)--(0.4,3.76)--(0.6,4.76)--(0.8,5.84)--(1,7)--(1.2,8.24)--(1.4,9.56)--(1.6,10.96)--(1.8,12.44);

\end{tikzpicture}\\
    \end{minipage}
    \begin{minipage}{0.5\linewidth}
    \begin{tikzpicture}[baseline, scale=0.8]

    \tikzset{
      point/.style={
        thick,
        draw,
        cross out,
        inner sep=0pt,
        minimum width=5pt,
        minimum height=5pt,
      },
    }
    \clip (-4.5,-6.5) rectangle (6.5,9.5);
    	
	\draw[color ={black},>=latex,->] (-3,0)--(3,0);
	\draw[color ={black},>=latex,->] (0,-5)--(0,8);
	\draw[color ={black},opacity = 0.4] (-3,1)--(3,1);
	\draw[color ={black},opacity = 0.4] (-3,2)--(3,2);
	\draw[color ={black},opacity = 0.4] (-3,3)--(3,3);
	\draw[color ={black},opacity = 0.4] (-3,4)--(3,4);
	\draw[color ={black},opacity = 0.4] (-3,5)--(3,5);
	\draw[color ={black},opacity = 0.4] (-3,6)--(3,6);
	\draw[color ={black},opacity = 0.4] (-3,7)--(3,7);
	\draw[color ={black},opacity = 0.4] (-3,8)--(3,8);
	\draw[color ={black},opacity = 0.4] (-3,-1)--(3,-1);
	\draw[color ={black},opacity = 0.4] (-3,-2)--(3,-2);
	\draw[color ={black},opacity = 0.4] (-3,-3)--(3,-3);
	\draw[color ={black},opacity = 0.4] (-3,-4)--(3,-4);
	\draw[color ={black},opacity = 0.4] (-3,-5)--(3,-5);
	\draw[color ={black},opacity = 0.4] (1,-5)--(1,8);
	\draw[color ={black},opacity = 0.4] (2,-5)--(2,8);
	\draw[color ={black},opacity = 0.4] (3,-5)--(3,8);
	\draw[color ={black},opacity = 0.4] (-1,-5)--(-1,8);
	\draw[color ={black},opacity = 0.4] (-2,-5)--(-2,8);
	\draw[color ={black},opacity = 0.4] (-3,-5)--(-3,8);
	\draw[color ={gray},opacity = 0.3] (-3,0)--(3,0);
	\draw[color ={gray},opacity = 0.3] (-3,0.5)--(3,0.5);
	\draw[color ={gray},opacity = 0.3] (-3,1.5)--(3,1.5);
	\draw[color ={gray},opacity = 0.3] (-3,2.5)--(3,2.5);
	\draw[color ={gray},opacity = 0.3] (-3,3.5)--(3,3.5);
	\draw[color ={gray},opacity = 0.3] (-3,4.5)--(3,4.5);
	\draw[color ={gray},opacity = 0.3] (-3,5.5)--(3,5.5);
	\draw[color ={gray},opacity = 0.3] (-3,6.5)--(3,6.5);
	\draw[color ={gray},opacity = 0.3] (-3,7.5)--(3,7.5);
	\draw[color ={gray},opacity = 0.3] (-3,0)--(3,0);
	\draw[color ={gray},opacity = 0.3] (-3,-0.5)--(3,-0.5);
	\draw[color ={gray},opacity = 0.3] (-3,-1.5)--(3,-1.5);
	\draw[color ={gray},opacity = 0.3] (-3,-2.5)--(3,-2.5);
	\draw[color ={gray},opacity = 0.3] (-3,-3.5)--(3,-3.5);
	\draw[color ={gray},opacity = 0.3] (-3,-4.5)--(3,-4.5);
	\draw[color ={gray},opacity = 0.3] (0,-5)--(0,8);
	\draw[color ={gray},opacity = 0.3] (0,-5)--(0,8);
	\draw[color ={black},line width = 1.2] (1,-0.13)--(1,0.13);
	\draw[color ={black},line width = 1.2] (2,-0.13)--(2,0.13);
	\draw[color ={black},line width = 1.2] (3,-0.13)--(3,0.13);
	\draw[color ={black},line width = 1.2] (-1,-0.13)--(-1,0.13);
	\draw[color ={black},line width = 1.2] (-2,-0.13)--(-2,0.13);
	\draw[color ={black},line width = 1.2] (-3,-0.13)--(-3,0.13);
	\draw[color ={black},line width = 1.2] (-0.13,1)--(0.13,1);
	\draw[color ={black},line width = 1.2] (-0.13,2)--(0.13,2);
	\draw[color ={black},line width = 1.2] (-0.13,3)--(0.13,3);
	\draw[color ={black},line width = 1.2] (-0.13,4)--(0.13,4);
	\draw[color ={black},line width = 1.2] (-0.13,5)--(0.13,5);
	\draw[color ={black},line width = 1.2] (-0.13,6)--(0.13,6);
	\draw[color ={black},line width = 1.2] (-0.13,7)--(0.13,7);
	\draw[color ={black},line width = 1.2] (-0.13,8)--(0.13,8);
	\draw[color ={black},line width = 1.2] (-0.13,9)--(0.13,9);
	\draw[color ={black},line width = 1.2] (-0.13,10)--(0.13,10);
	\draw[color ={black},line width = 1.2] (-0.13,11)--(0.13,11);
	\draw[color ={black},line width = 1.2] (-0.13,12)--(0.13,12);
	\draw[color ={black},line width = 1.2] (-0.13,-1)--(0.13,-1);
	\draw[color ={black},line width = 1.2] (-0.13,-2)--(0.13,-2);
	\draw[color ={black},line width = 1.2] (-0.13,-3)--(0.13,-3);
	\draw (3,-0.4) node[anchor = center] {\scriptsize \color{black}{$3$}};
	\draw (-3,-0.4) node[anchor = center] {\scriptsize \color{black}{$-3$}};
	\draw (-0.5,3.1) node[anchor = center] {\footnotesize \color{black}{$3$}};
	\draw (-0.5,6.1) node[anchor = center] {\footnotesize \color{black}{$6$}};
	\draw (-0.5,-2.9) node[anchor = center] {\footnotesize \color{black}{$-3$}};
	\draw [color={black}] (-0.3,-0.3) node[anchor = center,scale=1, rotate = 0] {O};
	\draw (3,6) node[anchor = center] {\footnotesize \color{red}{$\Large\mathcal{C}_f\prime$}};
	
	\draw[color={red},line width = 2] (-3,-2)--(-2.8,-1.6)--(-2.6,-1.2)--(-2.4,-0.8)--(-2.2,-0.4)--(-2,0)--(-1.8,0.4)--(-1.6,0.8)--(-1.4,1.2)--(-1.2,1.6)--(-1,2)--(-0.8,2.4)--(-0.6,2.8)--(-0.4,3.2)--(-0.2,3.6)--(0,4)--(0.2,4.4)--(0.4,4.8)--(0.6,5.2)--(0.8,5.6)--(1,6)--(1.2,6.4)--(1.4,6.8)--(1.6,7.2)--(1.8,7.6)--(2,8)--(2.2,8.4)--(2.4,8.8);

\end{tikzpicture}\\
    \end{minipage}
    
\end{exercice}

\begin{Solution}
 L'équation réduite de la tangente au point d'abscisse $1$ est  : $y=f'(1)(x-1)+f(1)$.\\
        On lit graphiquement $f(1)=7$ et $f'(1)=6$.\\
        L'équation réduite de la tangente est donc donnée par :
        $y=6(x-1)+7$, soit $y=6x+1$.
\end{Solution}

% @see : https://coopmaths.fr/alea?uuid=a1ba2&id=can1F14&n=1&d=10&alea=rkIp22323232323423456789101112131415161718192021222324252627282930313233343536373839404142434445464748234&cd=1&cols=1
\needspace{10\baselineskip}
\begin{exercice}[BaremeTotal=false]


Soit $f$ la fonction définie sur $\mathbb{R}$ par : $f(x)= -x^2+x+4$.\\
En calculant la limite du taux d'accroissement, déterminer $f'(2)$.
\end{exercice}

\begin{Solution}
 $f$ est une fonction polynôme du second degré de la forme $f(x)=ax^2+bx+c$.\\
    La fonction dérivée est donnée par la somme des dérivées des fonctions $u$ et $v$ définies par $u(x)=-x^2$ et $v(x)=x+4$.\\
     Comme $u'(x)=-2x$ et $v'(x)=1$, on obtient  $f'(x)=-2x+1$. \\
     Ainsi, $f'(2)=-2\times 2+1=-3$.
\end{Solution}

% @see : https://coopmaths.fr/alea?uuid=3c690&id=can1F13&n=1&d=10&alea=EUDu22323232323423456789101112131415161718192021222324252627282930313233343536373839404142434445464748234&cd=1&cols=1
\needspace{10\baselineskip}
\begin{exercice}[BaremeTotal=false]


 Déterminer le coefficient directeur de la tangente à la courbe représentative de la fonction inverse au point d'abscisse $-3$.
                 
\end{exercice}

\begin{Solution}
 Le coefficient directeur de la tangente au point d'abscisse $a$ est donné par le nombre dérivé $f'(a)$.\\
        La fonction $f$ définie par $f(x)=\dfrac{1}{x}$ a pour fonction dérivée la fonction $f'$ définie par $f'(x)=-\dfrac{1}{x^2}$.\\
Comme $f'(-3)=-\dfrac{1}{(-3)^2}=-\dfrac{1}{9}$, le coefficient directeur de la tangente au point d'abscisse $-3$ est : $-\dfrac{1}{9}$.
\end{Solution}

% @see : https://coopmaths.fr/alea?uuid=ebd89&id=1AN14-1&n=1&d=10&s=3&alea=ocYr22323232323423456789101112131415161718192021222324252627282930313233343536373839404142434445464748234&cd=0&cols=1
\needspace{10\baselineskip}
\begin{exercice}[BaremeTotal=false]


 Donner la dérivée de la fonction $f$, dérivable pour tout $x\in\R$, définie par  $f(x)=-\dfrac{8}{9}x+6$.
\end{exercice}

\begin{Solution}
 L'expression de la dérivée de la fonction $f$ définie par $f(x)=-\dfrac{8}{9}x+6$ est : ${\color[HTML]{f15929}\boldsymbol{f'(x)=-\dfrac{8}{9}}}$.
\end{Solution}

% @see : https://coopmaths.fr/alea?uuid=ebd89&id=1AN14-1&n=1&d=10&s=4&alea=uAlP22323232323423456789101112131415161718192021222324252627282930313233343536373839404142434445464748234&cd=0&cols=1
\needspace{10\baselineskip}
\begin{exercice}[BaremeTotal=false]


 Donner la dérivée de la fonction $f$, dérivable pour tout $x\in\R$, définie par  $f(x)=10x^{15}$.
\end{exercice}

\begin{Solution}
 L'expression de la dérivée de la fonction $f$ définie par $f(x)=10x^{15}$ est : ${\color[HTML]{f15929}\boldsymbol{f'(x)=150x^{14}}}$.
\end{Solution}

% @see : https://coopmaths.fr/alea?uuid=ebd89&id=1AN14-1&n=1&d=10&s=7&alea=vjN722323232323423456789101112131415161718192021222324252627282930313233343536373839404142434445464748234&cd=0&cols=1
\needspace{10\baselineskip}
\begin{exercice}[BaremeTotal=false]


 Donner la dérivée de la fonction $f$, dérivable pour tout $x\in\R^*$, définie par  $f(x)=\dfrac{-2}{7x}$.
\end{exercice}

\begin{Solution}
 L'expression de la dérivée de la fonction $f$ définie par $f(x)=\dfrac{-2}{7x}$ est : ${\color[HTML]{f15929}\boldsymbol{f'(x)=\dfrac{2}{7x^2}}}$.
\end{Solution}

% @see : https://coopmaths.fr/alea?uuid=a83c0&id=1AN14-4&n=1&d=10&s=1&alea=Njr322323232323423456789101112131415161718192021222324252627282930313233343536373839404142434445464748234&cd=0&cols=1
\needspace{10\baselineskip}
\begin{exercice}[BaremeTotal=false]


 Donner l'expression de la dérivée de la fonction $f$ définie sur $\R^{*}$ par $f(x)=-5x^{2}-5x+\dfrac{2}{x}$.\\
\end{exercice}

\begin{Solution}
 L'expression de la dérivée de la fonction $f$ définie par $f(x)=-5x^{2}-5x+\dfrac{2}{x}$ est :\\$f'(x)={\color[HTML]{f15929}\boldsymbol{-10x-5-\dfrac{2}{x^2}}}$.\\
\end{Solution}

% @see : https://coopmaths.fr/alea?uuid=a83c0&id=1AN14-4&n=1&d=10&s=4&alea=4Vpz22323232323423456789101112131415161718192021222324252627282930313233343536373839404142434445464748234&cd=1&cols=1
\needspace{10\baselineskip}
\begin{exercice}[BaremeTotal=false]


 Donner l'expression de la dérivée de la fonction $f$ définie sur $\R_+^{*}$ par $f(x)=\dfrac{4}{x}-5\sqrt{x}-2x^{2}$.\\
\end{exercice}

\begin{Solution}
 $(\dfrac{4}{x})^\prime=-\dfrac{4}{x^2}$\\$(-5\sqrt{x})^\prime=-\dfrac{5}{2\sqrt{x}}$\\$(-2x^{2})^\prime=-4x$\\L'expression de la dérivée de la fonction $f$ définie par $f(x)=\dfrac{4}{x}-5\sqrt{x}-2x^{2}$ est :\\$f'(x)={\color[HTML]{f15929}\boldsymbol{-\dfrac{4}{x^2}-\dfrac{5}{2\sqrt{x}}-4x}}$.\\
\end{Solution}

\end{Maquette}
\clearpage
\begin{Maquette}[DM]{Niveau= ,Classe=\getdata[5]\mydata\ (v5), Date = 06/01/25  ,Theme=Exercices}


% @see : https://coopmaths.fr/alea?uuid=4c8c7&id=1AN11-3&n=1&d=10&s=2&alea=nawO223232323234234567891011121314151617181920212223242526272829303132333435363738394041424344454647482345&cd=1&cols=2
\needspace{10\baselineskip}
\begin{exercice}[BaremeTotal=false]
\label{eleve:5}


 \begin{minipage}[t]{\linewidth} Soit $f$ une fonction dérivable sur $[-5;5]$ et $\mathcal{C}_f$ sa courbe représentative.\\On sait que  $f(4)=-5~~$ et que $~~f'(4)=5$.\\Déterminer une équation de la tangente $(T)$ à la courbe $\mathcal{C}_f$ au point d'abscisse $4$,\\sans utiliser la formule de cours de l'équation de tangente. \end{minipage}

\end{exercice}

\begin{Solution}
 \begin{minipage}[t]{\linewidth} On sait que la tangente n'est pas une droite verticale, puisque la fonction est dérivable sur l'intervalle.\\On en déduit que la tangente $(T)$ au point d'abscisse $4$, admet une équation réduite de la forme :  $(T) : y=m x + p$. 

\medskip
$\bullet$ Détermination de $m$ :\\On sait que le nombre dérivé en $4$ est par définition, le coefficient directeur de la tangente au point d'abscisse $4$.\\Par conséquent, on a déjà : $m=f'(4)=5$.\\ On en déduit que  $(T) : y= 5 x + p$\\$\bullet$ Détermination de $p$ :\\ Pour cela, on utilise que si $f(4)=-5~~$, alors le point $A$ de coordonnées $(4;-5)$ appartient à $\mathcal{C}_f$ mais aussi à $(T)$.\\ On peut écrire $A(4;-5) = \mathcal{C}_f \cap (T)$.\\ On remplace alors les coordonnées de $A(4;-5)$ dans l'équation  $(T) : y= 5 x + p$.\\ $\begin{aligned} \phantom{\iff}&A(4;-5)\in (T)\\ \iff& -5= 5 \times 4  + p\\ \iff& p=-5 -5 \times 4  \\ \iff& p=-5 -20   \\ \iff& p=-25\\\end{aligned}$\\On peut conclure que : $(T) : {\color[HTML]{f15929}\boldsymbol{y=5x-25}}$. \end{minipage}

\end{Solution}

% @see : https://coopmaths.fr/alea?uuid=0e984&id=can1F15&n=1&d=10&alea=1Alo223232323234234567891011121314151617181920212223242526272829303132333435363738394041424344454647482345&cd=1&cols=1
\needspace{10\baselineskip}
\begin{exercice}[BaremeTotal=false]


 La courbe représente une fonction $f$ et la droite est la tangente au point d'abscisse $-1$.\\
        Déterminer $f'(-1)$.\begin{center}\begin{tikzpicture}[baseline,scale = 0.5]

    \tikzset{
      point/.style={
        thick,
        draw,
        cross out,
        inner sep=0pt,
        minimum width=5pt,
        minimum height=5pt,
      },
    }
    \clip (-8,-10) rectangle (8,3);
    	
	\draw[color ={black},line width = 2,>=latex,->] (-8,0)--(8,0);
	\draw[color ={black},line width = 2,>=latex,->] (0,-10)--(0,3);
	\draw[color ={black},opacity = 0.5] (-8,1)--(8,1);
	\draw[color ={black},opacity = 0.5] (-8,2)--(8,2);
	\draw[color ={black},opacity = 0.5] (-8,3)--(8,3);
	\draw[color ={black},opacity = 0.5] (-8,-1)--(8,-1);
	\draw[color ={black},opacity = 0.5] (-8,-2)--(8,-2);
	\draw[color ={black},opacity = 0.5] (-8,-3)--(8,-3);
	\draw[color ={black},opacity = 0.5] (-8,-4)--(8,-4);
	\draw[color ={black},opacity = 0.5] (-8,-5)--(8,-5);
	\draw[color ={black},opacity = 0.5] (-8,-6)--(8,-6);
	\draw[color ={black},opacity = 0.5] (-8,-7)--(8,-7);
	\draw[color ={black},opacity = 0.5] (-8,-8)--(8,-8);
	\draw[color ={black},opacity = 0.5] (-8,-9)--(8,-9);
	\draw[color ={black},opacity = 0.5] (-8,-10)--(8,-10);
	\draw[color ={black},opacity = 0.5] (1,-10)--(1,3);
	\draw[color ={black},opacity = 0.5] (2,-10)--(2,3);
	\draw[color ={black},opacity = 0.5] (3,-10)--(3,3);
	\draw[color ={black},opacity = 0.5] (4,-10)--(4,3);
	\draw[color ={black},opacity = 0.5] (5,-10)--(5,3);
	\draw[color ={black},opacity = 0.5] (6,-10)--(6,3);
	\draw[color ={black},opacity = 0.5] (7,-10)--(7,3);
	\draw[color ={black},opacity = 0.5] (8,-10)--(8,3);
	\draw[color ={black},opacity = 0.5] (-1,-10)--(-1,3);
	\draw[color ={black},opacity = 0.5] (-2,-10)--(-2,3);
	\draw[color ={black},opacity = 0.5] (-3,-10)--(-3,3);
	\draw[color ={black},opacity = 0.5] (-4,-10)--(-4,3);
	\draw[color ={black},opacity = 0.5] (-5,-10)--(-5,3);
	\draw[color ={black},opacity = 0.5] (-6,-10)--(-6,3);
	\draw[color ={black},opacity = 0.5] (-7,-10)--(-7,3);
	\draw[color ={black},opacity = 0.5] (-8,-10)--(-8,3);
	\draw[color ={gray},opacity = 0.3] (-8,0)--(8,0);
	\draw[color ={gray},opacity = 0.3] (-8,0)--(8,0);
	\draw[color ={gray},opacity = 0.3] (-8,-4)--(8,-4);
	\draw[color ={gray},opacity = 0.3] (-8,-5)--(8,-5);
	\draw[color ={gray},opacity = 0.3] (-8,-6)--(8,-6);
	\draw[color ={gray},opacity = 0.3] (-8,-7)--(8,-7);
	\draw[color ={gray},opacity = 0.3] (-8,-8)--(8,-8);
	\draw[color ={gray},opacity = 0.3] (-8,-9)--(8,-9);
	\draw[color ={gray},opacity = 0.3] (-8,-10)--(8,-10);
	\draw[color ={gray},opacity = 0.3] (0,-10)--(0,3);
	\draw[color ={gray},opacity = 0.3] (0,-10)--(0,3);
	\draw[color ={black},line width = 2] (2,-0.2)--(2,0.2);
	\draw[color ={black},line width = 2] (4,-0.2)--(4,0.2);
	\draw[color ={black},line width = 2] (6,-0.2)--(6,0.2);
	\draw[color ={black},line width = 2] (-2,-0.2)--(-2,0.2);
	\draw[color ={black},line width = 2] (-4,-0.2)--(-4,0.2);
	\draw[color ={black},line width = 2] (-6,-0.2)--(-6,0.2);
	\draw[color ={black},line width = 2] (-0.2,1)--(0.2,1);
	\draw[color ={black},line width = 2] (-0.2,2)--(0.2,2);
	\draw[color ={black},line width = 2] (-0.2,-1)--(0.2,-1);
	\draw[color ={black},line width = 2] (-0.2,-2)--(0.2,-2);
	\draw[color ={black},line width = 2] (-0.2,-3)--(0.2,-3);
	\draw[color ={black},line width = 2] (-0.2,-4)--(0.2,-4);
	\draw[color ={black},line width = 2] (-0.2,-5)--(0.2,-5);
	\draw[color ={black},line width = 2] (-0.2,-6)--(0.2,-6);
	\draw[color ={black},line width = 2] (-0.2,-7)--(0.2,-7);
	\draw[color ={black},line width = 2] (-0.2,-8)--(0.2,-8);
	\draw[color ={black},line width = 2] (-0.2,-9)--(0.2,-9);
	\draw (2,-0.4) node[anchor = center] {\scriptsize \color{black}{$1$}};
	\draw (4,-0.4) node[anchor = center] {\scriptsize \color{black}{$2$}};
	\draw (6,-0.4) node[anchor = center] {\scriptsize \color{black}{$3$}};
	\draw (-2,-0.4) node[anchor = center] {\scriptsize \color{black}{$-1$}};
	\draw (-4,-0.4) node[anchor = center] {\scriptsize \color{black}{$-2$}};
	\draw (-6,-0.4) node[anchor = center] {\scriptsize \color{black}{$-3$}};
	\draw (-0.8,2.1) node[anchor = center] {\footnotesize \color{black}{$2$}};
	\draw (-0.8,-1.9) node[anchor = center] {\footnotesize \color{black}{$-2$}};
	\draw (-0.8,-3.9) node[anchor = center] {\footnotesize \color{black}{$-4$}};
	\draw (-0.8,-5.9) node[anchor = center] {\footnotesize \color{black}{$-6$}};
	\draw (-0.8,-7.9) node[anchor = center] {\footnotesize \color{black}{$-8$}};
	\draw [color={black}] (-0.3,-0.3) node[anchor = center,scale=1, rotate = 0] {O};
	
	\draw[color={blue},line width = 2] (-4.8,-10.52)--(-4.4,-8.68)--(-4,-7)--(-3.6,-5.48)--(-3.2,-4.12)--(-2.8,-2.92)--(-2.4,-1.88)--(-2,-1)--(-1.6,-0.28)--(-1.2,0.28)--(-0.8,0.68)--(-0.4,0.92)--(0,1)--(0.4,0.92)--(0.8,0.68)--(1.2,0.28)--(1.6,-0.28)--(2,-1)--(2.4,-1.88)--(2.8,-2.92)--(3.2,-4.12)--(3.6,-5.48)--(4,-7)--(4.4,-8.68)--(4.8,-10.52);
	
	\draw[color={red},line width = 2] (-6.8,-10.6)--(-6.4,-9.8)--(-6,-9)--(-5.6,-8.2)--(-5.2,-7.4)--(-4.8,-6.6)--(-4.4,-5.8)--(-4,-5)--(-3.6,-4.2)--(-3.2,-3.4)--(-2.8,-2.6)--(-2.4,-1.8)--(-2,-1)--(-1.6,-0.2)--(-1.2,0.6)--(-0.8,1.4)--(-0.4,2.2)--(0,3)--(0.4,3.8);

\end{tikzpicture}\end{center}
\end{exercice}

\begin{Solution}
 $f'(-1)$ est donné par le coefficient directeur de la tangente à la courbe au point d'abscisse $-1$, soit $4$.
\end{Solution}

% @see : https://coopmaths.fr/alea?uuid=6f32d&id=can1F16&n=1&d=10&alea=dqSw223232323234234567891011121314151617181920212223242526272829303132333435363738394041424344454647482345&cd=1&cols=1
\needspace{10\baselineskip}
\newpage
\begin{exercice}[BaremeTotal=false]


 On donne les représentations graphiques d'une fonction et de sa dérivée.\\
      Donner l'équation réduite de la tangente à la courbe de $f$ en $x=1$. \\ \begin{minipage}{0.5\linewidth}
    \begin{tikzpicture}[baseline, scale=0.8]

    \tikzset{
      point/.style={
        thick,
        draw,
        cross out,
        inner sep=0pt,
        minimum width=5pt,
        minimum height=5pt,
      },
    }
    \clip (-5.5,-9.5) rectangle (5.75,7.5);
    	
	\draw[color ={black},>=latex,->] (-4,0)--(4,0);
	\draw[color ={black},>=latex,->] (0,-8)--(0,6);
	\draw[color ={black},opacity = 0.4] (-4,1)--(4,1);
	\draw[color ={black},opacity = 0.4] (-4,2)--(4,2);
	\draw[color ={black},opacity = 0.4] (-4,3)--(4,3);
	\draw[color ={black},opacity = 0.4] (-4,4)--(4,4);
	\draw[color ={black},opacity = 0.4] (-4,5)--(4,5);
	\draw[color ={black},opacity = 0.4] (-4,6)--(4,6);
	\draw[color ={black},opacity = 0.4] (-4,-1)--(4,-1);
	\draw[color ={black},opacity = 0.4] (-4,-2)--(4,-2);
	\draw[color ={black},opacity = 0.4] (-4,-3)--(4,-3);
	\draw[color ={black},opacity = 0.4] (-4,-4)--(4,-4);
	\draw[color ={black},opacity = 0.4] (-4,-5)--(4,-5);
	\draw[color ={black},opacity = 0.4] (-4,-6)--(4,-6);
	\draw[color ={black},opacity = 0.4] (-4,-7)--(4,-7);
	\draw[color ={black},opacity = 0.4] (-4,-8)--(4,-8);
	\draw[color ={black},opacity = 0.4] (1,-8)--(1,6);
	\draw[color ={black},opacity = 0.4] (2,-8)--(2,6);
	\draw[color ={black},opacity = 0.4] (3,-8)--(3,6);
	\draw[color ={black},opacity = 0.4] (4,-8)--(4,6);
	\draw[color ={black},opacity = 0.4] (-1,-8)--(-1,6);
	\draw[color ={black},opacity = 0.4] (-2,-8)--(-2,6);
	\draw[color ={black},opacity = 0.4] (-3,-8)--(-3,6);
	\draw[color ={black},opacity = 0.4] (-4,-8)--(-4,6);
	\draw[color ={gray},opacity = 0.3] (-4,0)--(4,0);
	\draw[color ={gray},opacity = 0.3] (-4,0.5)--(4,0.5);
	\draw[color ={gray},opacity = 0.3] (-4,1.5)--(4,1.5);
	\draw[color ={gray},opacity = 0.3] (-4,2.5)--(4,2.5);
	\draw[color ={gray},opacity = 0.3] (-4,3.5)--(4,3.5);
	\draw[color ={gray},opacity = 0.3] (-4,4.5)--(4,4.5);
	\draw[color ={gray},opacity = 0.3] (-4,5.5)--(4,5.5);
	\draw[color ={gray},opacity = 0.3] (-4,0)--(4,0);
	\draw[color ={gray},opacity = 0.3] (-4,-0.5)--(4,-0.5);
	\draw[color ={gray},opacity = 0.3] (-4,-1.5)--(4,-1.5);
	\draw[color ={gray},opacity = 0.3] (-4,-2.5)--(4,-2.5);
	\draw[color ={gray},opacity = 0.3] (-4,-3.5)--(4,-3.5);
	\draw[color ={gray},opacity = 0.3] (-4,-4.5)--(4,-4.5);
	\draw[color ={gray},opacity = 0.3] (-4,-5.5)--(4,-5.5);
	\draw[color ={gray},opacity = 0.3] (-4,-6.5)--(4,-6.5);
	\draw[color ={gray},opacity = 0.3] (-4,-7)--(4,-7);
	\draw[color ={gray},opacity = 0.3] (-4,-7.5)--(4,-7.5);
	\draw[color ={gray},opacity = 0.3] (-4,-8)--(4,-8);
	\draw[color ={gray},opacity = 0.3] (0,-8)--(0,6);
	\draw[color ={gray},opacity = 0.3] (0,-8)--(0,6);
	\draw[color ={black},line width = 1.2] (1,-0.13)--(1,0.13);
	\draw[color ={black},line width = 1.2] (2,-0.13)--(2,0.13);
	\draw[color ={black},line width = 1.2] (3,-0.13)--(3,0.13);
	\draw[color ={black},line width = 1.2] (4,-0.13)--(4,0.13);
	\draw[color ={black},line width = 1.2] (-1,-0.13)--(-1,0.13);
	\draw[color ={black},line width = 1.2] (-2,-0.13)--(-2,0.13);
	\draw[color ={black},line width = 1.2] (-3,-0.13)--(-3,0.13);
	\draw[color ={black},line width = 1.2] (-4,-0.13)--(-4,0.13);
	\draw[color ={black},line width = 1.2] (-0.13,1)--(0.13,1);
	\draw[color ={black},line width = 1.2] (-0.13,2)--(0.13,2);
	\draw[color ={black},line width = 1.2] (-0.13,3)--(0.13,3);
	\draw[color ={black},line width = 1.2] (-0.13,4)--(0.13,4);
	\draw[color ={black},line width = 1.2] (-0.13,5)--(0.13,5);
	\draw[color ={black},line width = 1.2] (-0.13,6)--(0.13,6);
	\draw[color ={black},line width = 1.2] (-0.13,-1)--(0.13,-1);
	\draw[color ={black},line width = 1.2] (-0.13,-2)--(0.13,-2);
	\draw[color ={black},line width = 1.2] (-0.13,-3)--(0.13,-3);
	\draw[color ={black},line width = 1.2] (-0.13,-4)--(0.13,-4);
	\draw[color ={black},line width = 1.2] (-0.13,-5)--(0.13,-5);
	\draw[color ={black},line width = 1.2] (-0.13,-6)--(0.13,-6);
	\draw[color ={black},line width = 1.2] (-0.13,-7)--(0.13,-7);
	\draw[color ={black},line width = 1.2] (-0.13,-8)--(0.13,-8);
	\draw (2,-0.4) node[anchor = center] {\scriptsize \color{black}{$2$}};
	\draw (4,-0.4) node[anchor = center] {\scriptsize \color{black}{$4$}};
	\draw (-2,-0.4) node[anchor = center] {\scriptsize \color{black}{$-2$}};
	\draw (-4,-0.4) node[anchor = center] {\scriptsize \color{black}{$-4$}};
	\draw (-0.5,2.1) node[anchor = center] {\footnotesize \color{black}{$2$}};
	\draw (-0.5,4.1) node[anchor = center] {\footnotesize \color{black}{$4$}};
	\draw (-0.5,6.1) node[anchor = center] {\footnotesize \color{black}{$6$}};
	\draw (-0.5,-1.9) node[anchor = center] {\footnotesize \color{black}{$-2$}};
	\draw (-0.5,-3.9) node[anchor = center] {\footnotesize \color{black}{$-4$}};
	\draw (-0.5,-5.9) node[anchor = center] {\footnotesize \color{black}{$-6$}};
	\draw (-0.5,-7.9) node[anchor = center] {\footnotesize \color{black}{$-8$}};
	\draw [color={black}] (-0.3,-0.3) node[anchor = center,scale=1, rotate = 0] {O};
	\draw (3,4) node[anchor = center] {\footnotesize \color{blue}{$\Large \mathcal{C}_f$}};
	
	\draw[color={blue},line width = 2] (-3.2,-8.24)--(-3,-7)--(-2.8,-5.84)--(-2.6,-4.76)--(-2.4,-3.76)--(-2.2,-2.84)--(-2,-2)--(-1.8,-1.24)--(-1.6,-0.56)--(-1.4,0.04)--(-1.2,0.56)--(-1,1)--(-0.8,1.36)--(-0.6,1.64)--(-0.4,1.84)--(-0.2,1.96)--(0,2)--(0.2,1.96)--(0.4,1.84)--(0.6,1.64)--(0.8,1.36)--(1,1)--(1.2,0.56)--(1.4,0.04)--(1.6,-0.56)--(1.8,-1.24)--(2,-2)--(2.2,-2.84)--(2.4,-3.76)--(2.6,-4.76)--(2.8,-5.84)--(3,-7)--(3.2,-8.24);

\end{tikzpicture}\\
    \end{minipage}
    \begin{minipage}{0.5\linewidth}
    \begin{tikzpicture}[baseline, scale=0.8]

    \tikzset{
      point/.style={
        thick,
        draw,
        cross out,
        inner sep=0pt,
        minimum width=5pt,
        minimum height=5pt,
      },
    }
    \clip (-5.5,-9.5) rectangle (6.5,7.5);
    	
	\draw[color ={black},>=latex,->] (-4,0)--(4,0);
	\draw[color ={black},>=latex,->] (0,-8)--(0,6);
	\draw[color ={black},opacity = 0.4] (-4,1)--(4,1);
	\draw[color ={black},opacity = 0.4] (-4,2)--(4,2);
	\draw[color ={black},opacity = 0.4] (-4,3)--(4,3);
	\draw[color ={black},opacity = 0.4] (-4,4)--(4,4);
	\draw[color ={black},opacity = 0.4] (-4,5)--(4,5);
	\draw[color ={black},opacity = 0.4] (-4,6)--(4,6);
	\draw[color ={black},opacity = 0.4] (-4,-1)--(4,-1);
	\draw[color ={black},opacity = 0.4] (-4,-2)--(4,-2);
	\draw[color ={black},opacity = 0.4] (-4,-3)--(4,-3);
	\draw[color ={black},opacity = 0.4] (-4,-4)--(4,-4);
	\draw[color ={black},opacity = 0.4] (-4,-5)--(4,-5);
	\draw[color ={black},opacity = 0.4] (-4,-6)--(4,-6);
	\draw[color ={black},opacity = 0.4] (-4,-7)--(4,-7);
	\draw[color ={black},opacity = 0.4] (-4,-8)--(4,-8);
	\draw[color ={black},opacity = 0.4] (1,-8)--(1,6);
	\draw[color ={black},opacity = 0.4] (2,-8)--(2,6);
	\draw[color ={black},opacity = 0.4] (3,-8)--(3,6);
	\draw[color ={black},opacity = 0.4] (4,-8)--(4,6);
	\draw[color ={black},opacity = 0.4] (-1,-8)--(-1,6);
	\draw[color ={black},opacity = 0.4] (-2,-8)--(-2,6);
	\draw[color ={black},opacity = 0.4] (-3,-8)--(-3,6);
	\draw[color ={black},opacity = 0.4] (-4,-8)--(-4,6);
	\draw[color ={gray},opacity = 0.3] (-4,0)--(4,0);
	\draw[color ={gray},opacity = 0.3] (-4,0.5)--(4,0.5);
	\draw[color ={gray},opacity = 0.3] (-4,1.5)--(4,1.5);
	\draw[color ={gray},opacity = 0.3] (-4,2.5)--(4,2.5);
	\draw[color ={gray},opacity = 0.3] (-4,3.5)--(4,3.5);
	\draw[color ={gray},opacity = 0.3] (-4,4.5)--(4,4.5);
	\draw[color ={gray},opacity = 0.3] (-4,5.5)--(4,5.5);
	\draw[color ={gray},opacity = 0.3] (-4,0)--(4,0);
	\draw[color ={gray},opacity = 0.3] (-4,-0.5)--(4,-0.5);
	\draw[color ={gray},opacity = 0.3] (-4,-1.5)--(4,-1.5);
	\draw[color ={gray},opacity = 0.3] (-4,-2.5)--(4,-2.5);
	\draw[color ={gray},opacity = 0.3] (-4,-3.5)--(4,-3.5);
	\draw[color ={gray},opacity = 0.3] (-4,-4.5)--(4,-4.5);
	\draw[color ={gray},opacity = 0.3] (-4,-5.5)--(4,-5.5);
	\draw[color ={gray},opacity = 0.3] (-4,-6.5)--(4,-6.5);
	\draw[color ={gray},opacity = 0.3] (-4,-7)--(4,-7);
	\draw[color ={gray},opacity = 0.3] (-4,-7.5)--(4,-7.5);
	\draw[color ={gray},opacity = 0.3] (-4,-8)--(4,-8);
	\draw[color ={gray},opacity = 0.3] (0,-8)--(0,6);
	\draw[color ={gray},opacity = 0.3] (0,-8)--(0,6);
	\draw[color ={black},line width = 1.2] (1,-0.13)--(1,0.13);
	\draw[color ={black},line width = 1.2] (2,-0.13)--(2,0.13);
	\draw[color ={black},line width = 1.2] (3,-0.13)--(3,0.13);
	\draw[color ={black},line width = 1.2] (4,-0.13)--(4,0.13);
	\draw[color ={black},line width = 1.2] (-1,-0.13)--(-1,0.13);
	\draw[color ={black},line width = 1.2] (-2,-0.13)--(-2,0.13);
	\draw[color ={black},line width = 1.2] (-3,-0.13)--(-3,0.13);
	\draw[color ={black},line width = 1.2] (-4,-0.13)--(-4,0.13);
	\draw[color ={black},line width = 1.2] (-0.13,1)--(0.13,1);
	\draw[color ={black},line width = 1.2] (-0.13,2)--(0.13,2);
	\draw[color ={black},line width = 1.2] (-0.13,3)--(0.13,3);
	\draw[color ={black},line width = 1.2] (-0.13,4)--(0.13,4);
	\draw[color ={black},line width = 1.2] (-0.13,5)--(0.13,5);
	\draw[color ={black},line width = 1.2] (-0.13,6)--(0.13,6);
	\draw[color ={black},line width = 1.2] (-0.13,-1)--(0.13,-1);
	\draw[color ={black},line width = 1.2] (-0.13,-2)--(0.13,-2);
	\draw[color ={black},line width = 1.2] (-0.13,-3)--(0.13,-3);
	\draw[color ={black},line width = 1.2] (-0.13,-4)--(0.13,-4);
	\draw[color ={black},line width = 1.2] (-0.13,-5)--(0.13,-5);
	\draw[color ={black},line width = 1.2] (-0.13,-6)--(0.13,-6);
	\draw[color ={black},line width = 1.2] (-0.13,-7)--(0.13,-7);
	\draw[color ={black},line width = 1.2] (-0.13,-8)--(0.13,-8);
	\draw (2,-0.4) node[anchor = center] {\scriptsize \color{black}{$2$}};
	\draw (4,-0.4) node[anchor = center] {\scriptsize \color{black}{$4$}};
	\draw (-2,-0.4) node[anchor = center] {\scriptsize \color{black}{$-2$}};
	\draw (-4,-0.4) node[anchor = center] {\scriptsize \color{black}{$-4$}};
	\draw (-0.5,2.1) node[anchor = center] {\footnotesize \color{black}{$2$}};
	\draw (-0.5,4.1) node[anchor = center] {\footnotesize \color{black}{$4$}};
	\draw (-0.5,-1.9) node[anchor = center] {\footnotesize \color{black}{$-2$}};
	\draw (-0.5,-3.9) node[anchor = center] {\footnotesize \color{black}{$-4$}};
	\draw (-0.5,-5.9) node[anchor = center] {\footnotesize \color{black}{$-6$}};
	\draw (-0.5,-7.9) node[anchor = center] {\footnotesize \color{black}{$-8$}};
	\draw [color={black}] (-0.3,-0.3) node[anchor = center,scale=1, rotate = 0] {O};
	\draw (3,4) node[anchor = center] {\footnotesize \color{red}{$\Large\mathcal{C}_f\prime$}};
	
	\draw[color={red},line width = 2] (-3.4,6.8)--(-3.2,6.4)--(-3,6)--(-2.8,5.6)--(-2.6,5.2)--(-2.4,4.8)--(-2.2,4.4)--(-2,4)--(-1.8,3.6)--(-1.6,3.2)--(-1.4,2.8)--(-1.2,2.4)--(-1,2)--(-0.8,1.6)--(-0.6,1.2)--(-0.4,0.8)--(-0.2,0.4)--(0,0)--(0.2,-0.4)--(0.4,-0.8)--(0.6,-1.2)--(0.8,-1.6)--(1,-2)--(1.2,-2.4)--(1.4,-2.8)--(1.6,-3.2)--(1.8,-3.6)--(2,-4)--(2.2,-4.4)--(2.4,-4.8)--(2.6,-5.2)--(2.8,-5.6)--(3,-6)--(3.2,-6.4)--(3.4,-6.8)--(3.6,-7.2)--(3.8,-7.6)--(4,-8);

\end{tikzpicture}\\
    \end{minipage}
    
\end{exercice}

\begin{Solution}
 L'équation réduite de la tangente au point d'abscisse $1$ est  : $y=f'(1)(x-1)+f(1)$.\\
      On lit graphiquement $f(1)=1$ et $f'(1)=-2$.\\
      L'équation réduite de la tangente est donc donnée par :
      $y=-2(x-1)+1$, soit $y=-2x+3$.
\end{Solution}

% @see : https://coopmaths.fr/alea?uuid=a1ba2&id=can1F14&n=1&d=10&alea=rkIp223232323234234567891011121314151617181920212223242526272829303132333435363738394041424344454647482345&cd=1&cols=1
\needspace{10\baselineskip}
\begin{exercice}[BaremeTotal=false]


Soit $f$ la fonction définie sur $\mathbb{R}$ par : 
$f(x)= 8+2x^2+4x$.\\
En calculant la limite du taux d'accroissement, déterminer $f'(0)$.
\end{exercice}

\begin{Solution}
 $f$ est une fonction polynôme du second degré de la forme $f(x)=ax^2+bx+c$.\\
    La fonction dérivée est donnée par la somme des dérivées des fonctions $u$ et $v$ définies par $u(x)=2x^2$ et $v(x)=4x+8$.\\
     Comme $u'(x)=4x$ et $v'(x)=4$, on obtient  $f'(x)=4x+4$. \\
     Ainsi, $f'(0)=4\times 0+4=4$.
\end{Solution}

% @see : https://coopmaths.fr/alea?uuid=3c690&id=can1F13&n=1&d=10&alea=EUDu223232323234234567891011121314151617181920212223242526272829303132333435363738394041424344454647482345&cd=1&cols=1
\needspace{10\baselineskip}
\begin{exercice}[BaremeTotal=false]


 Déterminer le coefficient directeur de la tangente à la courbe représentative de la fonction cube au point d'abscisse $2$.
                        
\end{exercice}

\begin{Solution}
 Le coefficient directeur de la tangente au point d'abscisse $a$ est donné par le nombre dérivé $f'(a)$.\\
        La fonction $f$ définie par $f(x)=x^3$ a pour fonction dérivée la fonction $f'$ définie par $f'(x)=3x^2$.\\
        Comme $f'(2)=3\times 2^2=12$, le coefficient directeur de la tangente au point d'abscisse $2$ est : $12$. 
\end{Solution}

% @see : https://coopmaths.fr/alea?uuid=ebd89&id=1AN14-1&n=1&d=10&s=3&alea=ocYr223232323234234567891011121314151617181920212223242526272829303132333435363738394041424344454647482345&cd=0&cols=1
\needspace{10\baselineskip}
\begin{exercice}[BaremeTotal=false]


 Donner la dérivée de la fonction $f$, dérivable pour tout $x\in\R$, définie par  $f(x)=\dfrac{-x+4}{5}$.
\end{exercice}

\begin{Solution}
 L'expression de la dérivée de la fonction $f$ définie par $f(x)=\dfrac{-x+4}{5}$ est : ${\color[HTML]{f15929}\boldsymbol{f'(x)=-\dfrac{1}{5}}}$.
\end{Solution}

% @see : https://coopmaths.fr/alea?uuid=ebd89&id=1AN14-1&n=1&d=10&s=4&alea=uAlP223232323234234567891011121314151617181920212223242526272829303132333435363738394041424344454647482345&cd=0&cols=1
\needspace{10\baselineskip}
\begin{exercice}[BaremeTotal=false]


 Donner la dérivée de la fonction $f$, dérivable pour tout $x\in\R$, définie par  $f(x)=7x^{10}$.
\end{exercice}

\begin{Solution}
 L'expression de la dérivée de la fonction $f$ définie par $f(x)=7x^{10}$ est : ${\color[HTML]{f15929}\boldsymbol{f'(x)=70x^{9}}}$.
\end{Solution}

% @see : https://coopmaths.fr/alea?uuid=ebd89&id=1AN14-1&n=1&d=10&s=7&alea=vjN7223232323234234567891011121314151617181920212223242526272829303132333435363738394041424344454647482345&cd=0&cols=1
\needspace{10\baselineskip}
\begin{exercice}[BaremeTotal=false]


 Donner la dérivée de la fonction $f$, dérivable pour tout $x\in\R^*$, définie par  $f(x)=\dfrac{7}{8x}$.
\end{exercice}

\begin{Solution}
 L'expression de la dérivée de la fonction $f$ définie par $f(x)=\dfrac{7}{8x}$ est : ${\color[HTML]{f15929}\boldsymbol{f'(x)=-\dfrac{7}{8x^2}}}$.
\end{Solution}

% @see : https://coopmaths.fr/alea?uuid=a83c0&id=1AN14-4&n=1&d=10&s=1&alea=Njr3223232323234234567891011121314151617181920212223242526272829303132333435363738394041424344454647482345&cd=0&cols=1
\needspace{10\baselineskip}
\begin{exercice}[BaremeTotal=false]


 Donner l'expression de la dérivée de la fonction $f$ définie sur $\R^{*}$ par $f(x)=4x^{2}-2x-3-\dfrac{2}{x}$.\\
\end{exercice}

\begin{Solution}
 L'expression de la dérivée de la fonction $f$ définie par $f(x)=4x^{2}-2x-3-\dfrac{2}{x}$ est :\\$f'(x)={\color[HTML]{f15929}\boldsymbol{8x-2+\dfrac{2}{x^2}}}$.\\
\end{Solution}

% @see : https://coopmaths.fr/alea?uuid=a83c0&id=1AN14-4&n=1&d=10&s=4&alea=4Vpz223232323234234567891011121314151617181920212223242526272829303132333435363738394041424344454647482345&cd=1&cols=1
\needspace{10\baselineskip}
\begin{exercice}[BaremeTotal=false]


 Donner l'expression de la dérivée de la fonction $f$ définie sur $\R_+^{*}$ par $f(x)=2x^{2}+2x+\dfrac{4}{x}-3\sqrt{x}$.\\
\end{exercice}

\begin{Solution}
 $(2x^{2}+2x)^\prime=4x+2$\\$(\dfrac{4}{x})^\prime=-\dfrac{4}{x^2}$\\$(-3\sqrt{x})^\prime=-\dfrac{3}{2\sqrt{x}}$\\L'expression de la dérivée de la fonction $f$ définie par $f(x)=2x^{2}+2x+\dfrac{4}{x}-3\sqrt{x}$ est :\\$f'(x)={\color[HTML]{f15929}\boldsymbol{4x+2-\dfrac{4}{x^2}-\dfrac{3}{2\sqrt{x}}}}$.\\
\end{Solution}

\end{Maquette}
\clearpage
\begin{Maquette}[DM]{Niveau= ,Classe=\getdata[6]\mydata\ (v6), Date = 06/01/25  ,Theme=Exercices}


% @see : https://coopmaths.fr/alea?uuid=4c8c7&id=1AN11-3&n=1&d=10&s=2&alea=nawO2232323232342345678910111213141516171819202122232425262728293031323334353637383940414243444546474823456&cd=1&cols=2
\needspace{10\baselineskip}
\begin{exercice}[BaremeTotal=false]
\label{eleve:6}


 \begin{minipage}[t]{\linewidth} Soit $f$ une fonction dérivable sur $[-5;5]$ et $\mathcal{C}_f$ sa courbe représentative.\\On sait que  $f(0)=0~~$ et que $~~f'(0)=-5$.\\Déterminer une équation de la tangente $(T)$ à la courbe $\mathcal{C}_f$ au point d'abscisse $0$,\\sans utiliser la formule de cours de l'équation de tangente. \end{minipage}

\end{exercice}

\begin{Solution}
 \begin{minipage}[t]{\linewidth} On sait que la tangente n'est pas une droite verticale, puisque la fonction est dérivable sur l'intervalle.\\On en déduit que la tangente $(T)$ au point d'abscisse $0$, admet une équation réduite de la forme :  $(T) : y=m x + p$. 

\medskip
$\bullet$ Détermination de $m$ :\\On sait que le nombre dérivé en $0$ est par définition, le coefficient directeur de la tangente au point d'abscisse $0$.\\Par conséquent, on a déjà : $m=f'(0)=-5$.\\ On en déduit que  $(T) : y= -5 x + p$\\$\bullet$ Détermination de $p$ :\\ Pour cela, on utilise que si $f(0)=0~~$, alors le point $A$ de coordonnées $(0;0)$ appartient à $\mathcal{C}_f$ mais aussi à $(T)$.\\ On peut écrire $A(0;0) = \mathcal{C}_f \cap (T)$.\\ On remplace alors les coordonnées de $A(0;0)$ dans l'équation  $(T) : y= -5 x + p$.\\ $\begin{aligned} \phantom{\iff}&A(0;0)\in (T)\\ \iff& 0= -5 \times 0  + p\\ \iff& p=0 +5 \times 0  \\ \iff& p=0 +0   \\ \iff& p=0\\\end{aligned}$\\On peut conclure que : $(T) : {\color[HTML]{f15929}\boldsymbol{y=-5x}}$. \end{minipage}

\end{Solution}

% @see : https://coopmaths.fr/alea?uuid=0e984&id=can1F15&n=1&d=10&alea=1Alo2232323232342345678910111213141516171819202122232425262728293031323334353637383940414243444546474823456&cd=1&cols=1
\needspace{10\baselineskip}
\begin{exercice}[BaremeTotal=false]


 La courbe représente une fonction $f$ et la droite est la tangente au point d'abscisse $0$.\\
        Déterminer $f'(0)$.\begin{center}\begin{tikzpicture}[baseline,scale = 0.5]

    \tikzset{
      point/.style={
        thick,
        draw,
        cross out,
        inner sep=0pt,
        minimum width=5pt,
        minimum height=5pt,
      },
    }
    \clip (-8,-10) rectangle (8,3);
    	
	\draw[color ={black},line width = 2,>=latex,->] (-8,0)--(8,0);
	\draw[color ={black},line width = 2,>=latex,->] (0,-10)--(0,3);
	\draw[color ={black},opacity = 0.5] (-8,1)--(8,1);
	\draw[color ={black},opacity = 0.5] (-8,2)--(8,2);
	\draw[color ={black},opacity = 0.5] (-8,3)--(8,3);
	\draw[color ={black},opacity = 0.5] (-8,-1)--(8,-1);
	\draw[color ={black},opacity = 0.5] (-8,-2)--(8,-2);
	\draw[color ={black},opacity = 0.5] (-8,-3)--(8,-3);
	\draw[color ={black},opacity = 0.5] (-8,-4)--(8,-4);
	\draw[color ={black},opacity = 0.5] (-8,-5)--(8,-5);
	\draw[color ={black},opacity = 0.5] (-8,-6)--(8,-6);
	\draw[color ={black},opacity = 0.5] (-8,-7)--(8,-7);
	\draw[color ={black},opacity = 0.5] (-8,-8)--(8,-8);
	\draw[color ={black},opacity = 0.5] (-8,-9)--(8,-9);
	\draw[color ={black},opacity = 0.5] (-8,-10)--(8,-10);
	\draw[color ={black},opacity = 0.5] (1,-10)--(1,3);
	\draw[color ={black},opacity = 0.5] (2,-10)--(2,3);
	\draw[color ={black},opacity = 0.5] (3,-10)--(3,3);
	\draw[color ={black},opacity = 0.5] (4,-10)--(4,3);
	\draw[color ={black},opacity = 0.5] (5,-10)--(5,3);
	\draw[color ={black},opacity = 0.5] (6,-10)--(6,3);
	\draw[color ={black},opacity = 0.5] (7,-10)--(7,3);
	\draw[color ={black},opacity = 0.5] (8,-10)--(8,3);
	\draw[color ={black},opacity = 0.5] (-1,-10)--(-1,3);
	\draw[color ={black},opacity = 0.5] (-2,-10)--(-2,3);
	\draw[color ={black},opacity = 0.5] (-3,-10)--(-3,3);
	\draw[color ={black},opacity = 0.5] (-4,-10)--(-4,3);
	\draw[color ={black},opacity = 0.5] (-5,-10)--(-5,3);
	\draw[color ={black},opacity = 0.5] (-6,-10)--(-6,3);
	\draw[color ={black},opacity = 0.5] (-7,-10)--(-7,3);
	\draw[color ={black},opacity = 0.5] (-8,-10)--(-8,3);
	\draw[color ={gray},opacity = 0.3] (-8,0)--(8,0);
	\draw[color ={gray},opacity = 0.3] (-8,0)--(8,0);
	\draw[color ={gray},opacity = 0.3] (-8,-4)--(8,-4);
	\draw[color ={gray},opacity = 0.3] (-8,-5)--(8,-5);
	\draw[color ={gray},opacity = 0.3] (-8,-6)--(8,-6);
	\draw[color ={gray},opacity = 0.3] (-8,-7)--(8,-7);
	\draw[color ={gray},opacity = 0.3] (-8,-8)--(8,-8);
	\draw[color ={gray},opacity = 0.3] (-8,-9)--(8,-9);
	\draw[color ={gray},opacity = 0.3] (-8,-10)--(8,-10);
	\draw[color ={gray},opacity = 0.3] (0,-10)--(0,3);
	\draw[color ={gray},opacity = 0.3] (0,-10)--(0,3);
	\draw[color ={black},line width = 2] (2,-0.2)--(2,0.2);
	\draw[color ={black},line width = 2] (4,-0.2)--(4,0.2);
	\draw[color ={black},line width = 2] (6,-0.2)--(6,0.2);
	\draw[color ={black},line width = 2] (-2,-0.2)--(-2,0.2);
	\draw[color ={black},line width = 2] (-4,-0.2)--(-4,0.2);
	\draw[color ={black},line width = 2] (-6,-0.2)--(-6,0.2);
	\draw[color ={black},line width = 2] (-0.2,1)--(0.2,1);
	\draw[color ={black},line width = 2] (-0.2,2)--(0.2,2);
	\draw[color ={black},line width = 2] (-0.2,-1)--(0.2,-1);
	\draw[color ={black},line width = 2] (-0.2,-2)--(0.2,-2);
	\draw[color ={black},line width = 2] (-0.2,-3)--(0.2,-3);
	\draw[color ={black},line width = 2] (-0.2,-4)--(0.2,-4);
	\draw[color ={black},line width = 2] (-0.2,-5)--(0.2,-5);
	\draw[color ={black},line width = 2] (-0.2,-6)--(0.2,-6);
	\draw[color ={black},line width = 2] (-0.2,-7)--(0.2,-7);
	\draw[color ={black},line width = 2] (-0.2,-8)--(0.2,-8);
	\draw[color ={black},line width = 2] (-0.2,-9)--(0.2,-9);
	\draw (2,-0.4) node[anchor = center] {\scriptsize \color{black}{$1$}};
	\draw (4,-0.4) node[anchor = center] {\scriptsize \color{black}{$2$}};
	\draw (6,-0.4) node[anchor = center] {\scriptsize \color{black}{$3$}};
	\draw (-2,-0.4) node[anchor = center] {\scriptsize \color{black}{$-1$}};
	\draw (-4,-0.4) node[anchor = center] {\scriptsize \color{black}{$-2$}};
	\draw (-6,-0.4) node[anchor = center] {\scriptsize \color{black}{$-3$}};
	\draw (-0.8,2.1) node[anchor = center] {\footnotesize \color{black}{$2$}};
	\draw (-0.8,-1.9) node[anchor = center] {\footnotesize \color{black}{$-2$}};
	\draw (-0.8,-3.9) node[anchor = center] {\footnotesize \color{black}{$-4$}};
	\draw (-0.8,-5.9) node[anchor = center] {\footnotesize \color{black}{$-6$}};
	\draw (-0.8,-7.9) node[anchor = center] {\footnotesize \color{black}{$-8$}};
	\draw [color={black}] (-0.3,-0.3) node[anchor = center,scale=1, rotate = 0] {O};
	
	\draw[color={blue},line width = 2] (-8,-8)--(-7.6,-6.84)--(-7.2,-5.76)--(-6.8,-4.76)--(-6.4,-3.84)--(-6,-3)--(-5.6,-2.24)--(-5.2,-1.56)--(-4.8,-0.96)--(-4.4,-0.44)--(-4,0)--(-3.6,0.36)--(-3.2,0.64)--(-2.8,0.84)--(-2.4,0.96)--(-2,1)--(-1.6,0.96)--(-1.2,0.84)--(-0.8,0.64)--(-0.4,0.36)--(0,0)--(0.4,-0.44)--(0.8,-0.96)--(1.2,-1.56)--(1.6,-2.24)--(2,-3)--(2.4,-3.84)--(2.8,-4.76)--(3.2,-5.76)--(3.6,-6.84)--(4,-8)--(4.4,-9.24)--(4.8,-10.56);
	
	\draw[color={red},line width = 2] (-4,4)--(-3.6,3.6)--(-3.2,3.2)--(-2.8,2.8)--(-2.4,2.4)--(-2,2)--(-1.6,1.6)--(-1.2,1.2)--(-0.8,0.8)--(-0.4,0.4)--(0,0)--(0.4,-0.4)--(0.8,-0.8)--(1.2,-1.2)--(1.6,-1.6)--(2,-2)--(2.4,-2.4)--(2.8,-2.8)--(3.2,-3.2)--(3.6,-3.6)--(4,-4)--(4.4,-4.4)--(4.8,-4.8)--(5.2,-5.2)--(5.6,-5.6)--(6,-6)--(6.4,-6.4)--(6.8,-6.8)--(7.2,-7.2)--(7.6,-7.6)--(8,-8);

\end{tikzpicture}\end{center}
\end{exercice}

\begin{Solution}
 $f'(0)$ est donné par le coefficient directeur de la tangente à la courbe au point d'abscisse $0$, soit $-2$.
\end{Solution}

% @see : https://coopmaths.fr/alea?uuid=6f32d&id=can1F16&n=1&d=10&alea=dqSw2232323232342345678910111213141516171819202122232425262728293031323334353637383940414243444546474823456&cd=1&cols=1
\needspace{10\baselineskip}
\newpage
\begin{exercice}[BaremeTotal=false]


 On donne les représentations graphiques d'une fonction et de sa dérivée.\\
        Donner l'équation réduite de la tangente à la courbe de $f$ en $x=1$. \\ \begin{minipage}{0.5\linewidth}
    \begin{tikzpicture}[baseline, scale=0.8]

    \tikzset{
      point/.style={
        thick,
        draw,
        cross out,
        inner sep=0pt,
        minimum width=5pt,
        minimum height=5pt,
      },
    }
    \clip (-4.5,-4.5) rectangle (5.75,13.5);
    	
	\draw[color ={black},>=latex,->] (-3,0)--(3,0);
	\draw[color ={black},>=latex,->] (0,-3)--(0,12);
	\draw[color ={black},opacity = 0.4] (-3,1)--(3,1);
	\draw[color ={black},opacity = 0.4] (-3,2)--(3,2);
	\draw[color ={black},opacity = 0.4] (-3,3)--(3,3);
	\draw[color ={black},opacity = 0.4] (-3,4)--(3,4);
	\draw[color ={black},opacity = 0.4] (-3,5)--(3,5);
	\draw[color ={black},opacity = 0.4] (-3,6)--(3,6);
	\draw[color ={black},opacity = 0.4] (-3,7)--(3,7);
	\draw[color ={black},opacity = 0.4] (-3,8)--(3,8);
	\draw[color ={black},opacity = 0.4] (-3,9)--(3,9);
	\draw[color ={black},opacity = 0.4] (-3,10)--(3,10);
	\draw[color ={black},opacity = 0.4] (-3,11)--(3,11);
	\draw[color ={black},opacity = 0.4] (-3,12)--(3,12);
	\draw[color ={black},opacity = 0.4] (-3,-1)--(3,-1);
	\draw[color ={black},opacity = 0.4] (-3,-2)--(3,-2);
	\draw[color ={black},opacity = 0.4] (-3,-3)--(3,-3);
	\draw[color ={black},opacity = 0.4] (1,-3)--(1,12);
	\draw[color ={black},opacity = 0.4] (2,-3)--(2,12);
	\draw[color ={black},opacity = 0.4] (3,-3)--(3,12);
	\draw[color ={black},opacity = 0.4] (-1,-3)--(-1,12);
	\draw[color ={black},opacity = 0.4] (-2,-3)--(-2,12);
	\draw[color ={black},opacity = 0.4] (-3,-3)--(-3,12);
	\draw[color ={gray},opacity = 0.3] (-3,0)--(3,0);
	\draw[color ={gray},opacity = 0.3] (-3,0.5)--(3,0.5);
	\draw[color ={gray},opacity = 0.3] (-3,1.5)--(3,1.5);
	\draw[color ={gray},opacity = 0.3] (-3,2.5)--(3,2.5);
	\draw[color ={gray},opacity = 0.3] (-3,3.5)--(3,3.5);
	\draw[color ={gray},opacity = 0.3] (-3,4.5)--(3,4.5);
	\draw[color ={gray},opacity = 0.3] (-3,5.5)--(3,5.5);
	\draw[color ={gray},opacity = 0.3] (-3,6.5)--(3,6.5);
	\draw[color ={gray},opacity = 0.3] (-3,7.5)--(3,7.5);
	\draw[color ={gray},opacity = 0.3] (-3,8.5)--(3,8.5);
	\draw[color ={gray},opacity = 0.3] (-3,9.5)--(3,9.5);
	\draw[color ={gray},opacity = 0.3] (-3,10.5)--(3,10.5);
	\draw[color ={gray},opacity = 0.3] (-3,11.5)--(3,11.5);
	\draw[color ={gray},opacity = 0.3] (-3,0)--(3,0);
	\draw[color ={gray},opacity = 0.3] (-3,-0.5)--(3,-0.5);
	\draw[color ={gray},opacity = 0.3] (-3,-1.5)--(3,-1.5);
	\draw[color ={gray},opacity = 0.3] (-3,-2.5)--(3,-2.5);
	\draw[color ={gray},opacity = 0.3] (0,-3)--(0,12);
	\draw[color ={gray},opacity = 0.3] (0,-3)--(0,12);
	\draw[color ={black},line width = 1.2] (1,-0.13)--(1,0.13);
	\draw[color ={black},line width = 1.2] (2,-0.13)--(2,0.13);
	\draw[color ={black},line width = 1.2] (3,-0.13)--(3,0.13);
	\draw[color ={black},line width = 1.2] (-1,-0.13)--(-1,0.13);
	\draw[color ={black},line width = 1.2] (-2,-0.13)--(-2,0.13);
	\draw[color ={black},line width = 1.2] (-3,-0.13)--(-3,0.13);
	\draw[color ={black},line width = 1.2] (-0.13,1)--(0.13,1);
	\draw[color ={black},line width = 1.2] (-0.13,2)--(0.13,2);
	\draw[color ={black},line width = 1.2] (-0.13,3)--(0.13,3);
	\draw[color ={black},line width = 1.2] (-0.13,4)--(0.13,4);
	\draw[color ={black},line width = 1.2] (-0.13,5)--(0.13,5);
	\draw[color ={black},line width = 1.2] (-0.13,6)--(0.13,6);
	\draw[color ={black},line width = 1.2] (-0.13,7)--(0.13,7);
	\draw[color ={black},line width = 1.2] (-0.13,8)--(0.13,8);
	\draw[color ={black},line width = 1.2] (-0.13,9)--(0.13,9);
	\draw[color ={black},line width = 1.2] (-0.13,10)--(0.13,10);
	\draw[color ={black},line width = 1.2] (-0.13,11)--(0.13,11);
	\draw[color ={black},line width = 1.2] (-0.13,12)--(0.13,12);
	\draw[color ={black},line width = 1.2] (-0.13,-1)--(0.13,-1);
	\draw[color ={black},line width = 1.2] (-0.13,-2)--(0.13,-2);
	\draw[color ={black},line width = 1.2] (-0.13,-3)--(0.13,-3);
	\draw (3,-0.4) node[anchor = center] {\scriptsize \color{black}{$3$}};
	\draw (-3,-0.4) node[anchor = center] {\scriptsize \color{black}{$-3$}};
	\draw (-0.5,3.1) node[anchor = center] {\footnotesize \color{black}{$3$}};
	\draw (-0.5,6.1) node[anchor = center] {\footnotesize \color{black}{$6$}};
	\draw (-0.5,9.1) node[anchor = center] {\footnotesize \color{black}{$9$}};
	\draw (-0.5,-2.9) node[anchor = center] {\footnotesize \color{black}{$-3$}};
	\draw [color={black}] (-0.3,-0.3) node[anchor = center,scale=1, rotate = 0] {O};
	\draw (3,10) node[anchor = center] {\footnotesize \color{blue}{$\Large \mathcal{C}_f$}};
	
	\draw[color={blue},line width = 2] (-3,7)--(-2.8,5.84)--(-2.6,4.76)--(-2.4,3.76)--(-2.2,2.84)--(-2,2)--(-1.8,1.24)--(-1.6,0.56)--(-1.4,-0.04)--(-1.2,-0.56)--(-1,-1)--(-0.8,-1.36)--(-0.6,-1.64)--(-0.4,-1.84)--(-0.2,-1.96)--(0,-2)--(0.2,-1.96)--(0.4,-1.84)--(0.6,-1.64)--(0.8,-1.36)--(1,-1)--(1.2,-0.56)--(1.4,-0.04)--(1.6,0.56)--(1.8,1.24)--(2,2)--(2.2,2.84)--(2.4,3.76)--(2.6,4.76)--(2.8,5.84)--(3,7);

\end{tikzpicture}\\
    \end{minipage}
    \begin{minipage}{0.5\linewidth}
    \begin{tikzpicture}[baseline, scale=0.8]

    \tikzset{
      point/.style={
        thick,
        draw,
        cross out,
        inner sep=0pt,
        minimum width=5pt,
        minimum height=5pt,
      },
    }
    \clip (-4.5,-6.5) rectangle (6.5,9.5);
    	
	\draw[color ={black},>=latex,->] (-3,0)--(3,0);
	\draw[color ={black},>=latex,->] (0,-5)--(0,8);
	\draw[color ={black},opacity = 0.4] (-3,1)--(3,1);
	\draw[color ={black},opacity = 0.4] (-3,2)--(3,2);
	\draw[color ={black},opacity = 0.4] (-3,3)--(3,3);
	\draw[color ={black},opacity = 0.4] (-3,4)--(3,4);
	\draw[color ={black},opacity = 0.4] (-3,5)--(3,5);
	\draw[color ={black},opacity = 0.4] (-3,6)--(3,6);
	\draw[color ={black},opacity = 0.4] (-3,7)--(3,7);
	\draw[color ={black},opacity = 0.4] (-3,8)--(3,8);
	\draw[color ={black},opacity = 0.4] (-3,-1)--(3,-1);
	\draw[color ={black},opacity = 0.4] (-3,-2)--(3,-2);
	\draw[color ={black},opacity = 0.4] (-3,-3)--(3,-3);
	\draw[color ={black},opacity = 0.4] (-3,-4)--(3,-4);
	\draw[color ={black},opacity = 0.4] (-3,-5)--(3,-5);
	\draw[color ={black},opacity = 0.4] (1,-5)--(1,8);
	\draw[color ={black},opacity = 0.4] (2,-5)--(2,8);
	\draw[color ={black},opacity = 0.4] (3,-5)--(3,8);
	\draw[color ={black},opacity = 0.4] (-1,-5)--(-1,8);
	\draw[color ={black},opacity = 0.4] (-2,-5)--(-2,8);
	\draw[color ={black},opacity = 0.4] (-3,-5)--(-3,8);
	\draw[color ={gray},opacity = 0.3] (-3,0)--(3,0);
	\draw[color ={gray},opacity = 0.3] (-3,0.5)--(3,0.5);
	\draw[color ={gray},opacity = 0.3] (-3,1.5)--(3,1.5);
	\draw[color ={gray},opacity = 0.3] (-3,2.5)--(3,2.5);
	\draw[color ={gray},opacity = 0.3] (-3,3.5)--(3,3.5);
	\draw[color ={gray},opacity = 0.3] (-3,4.5)--(3,4.5);
	\draw[color ={gray},opacity = 0.3] (-3,5.5)--(3,5.5);
	\draw[color ={gray},opacity = 0.3] (-3,6.5)--(3,6.5);
	\draw[color ={gray},opacity = 0.3] (-3,7.5)--(3,7.5);
	\draw[color ={gray},opacity = 0.3] (-3,0)--(3,0);
	\draw[color ={gray},opacity = 0.3] (-3,-0.5)--(3,-0.5);
	\draw[color ={gray},opacity = 0.3] (-3,-1.5)--(3,-1.5);
	\draw[color ={gray},opacity = 0.3] (-3,-2.5)--(3,-2.5);
	\draw[color ={gray},opacity = 0.3] (-3,-3.5)--(3,-3.5);
	\draw[color ={gray},opacity = 0.3] (-3,-4.5)--(3,-4.5);
	\draw[color ={gray},opacity = 0.3] (0,-5)--(0,8);
	\draw[color ={gray},opacity = 0.3] (0,-5)--(0,8);
	\draw[color ={black},line width = 1.2] (1,-0.13)--(1,0.13);
	\draw[color ={black},line width = 1.2] (2,-0.13)--(2,0.13);
	\draw[color ={black},line width = 1.2] (3,-0.13)--(3,0.13);
	\draw[color ={black},line width = 1.2] (-1,-0.13)--(-1,0.13);
	\draw[color ={black},line width = 1.2] (-2,-0.13)--(-2,0.13);
	\draw[color ={black},line width = 1.2] (-3,-0.13)--(-3,0.13);
	\draw[color ={black},line width = 1.2] (-0.13,1)--(0.13,1);
	\draw[color ={black},line width = 1.2] (-0.13,2)--(0.13,2);
	\draw[color ={black},line width = 1.2] (-0.13,3)--(0.13,3);
	\draw[color ={black},line width = 1.2] (-0.13,4)--(0.13,4);
	\draw[color ={black},line width = 1.2] (-0.13,5)--(0.13,5);
	\draw[color ={black},line width = 1.2] (-0.13,6)--(0.13,6);
	\draw[color ={black},line width = 1.2] (-0.13,7)--(0.13,7);
	\draw[color ={black},line width = 1.2] (-0.13,8)--(0.13,8);
	\draw[color ={black},line width = 1.2] (-0.13,9)--(0.13,9);
	\draw[color ={black},line width = 1.2] (-0.13,10)--(0.13,10);
	\draw[color ={black},line width = 1.2] (-0.13,11)--(0.13,11);
	\draw[color ={black},line width = 1.2] (-0.13,12)--(0.13,12);
	\draw[color ={black},line width = 1.2] (-0.13,-1)--(0.13,-1);
	\draw[color ={black},line width = 1.2] (-0.13,-2)--(0.13,-2);
	\draw[color ={black},line width = 1.2] (-0.13,-3)--(0.13,-3);
	\draw (3,-0.4) node[anchor = center] {\scriptsize \color{black}{$3$}};
	\draw (-3,-0.4) node[anchor = center] {\scriptsize \color{black}{$-3$}};
	\draw (-0.5,3.1) node[anchor = center] {\footnotesize \color{black}{$3$}};
	\draw (-0.5,6.1) node[anchor = center] {\footnotesize \color{black}{$6$}};
	\draw (-0.5,-2.9) node[anchor = center] {\footnotesize \color{black}{$-3$}};
	\draw [color={black}] (-0.3,-0.3) node[anchor = center,scale=1, rotate = 0] {O};
	\draw (3,6) node[anchor = center] {\footnotesize \color{red}{$\Large\mathcal{C}_f\prime$}};
	
	\draw[color={red},line width = 2] (-2.8,-5.6)--(-2.6,-5.2)--(-2.4,-4.8)--(-2.2,-4.4)--(-2,-4)--(-1.8,-3.6)--(-1.6,-3.2)--(-1.4,-2.8)--(-1.2,-2.4)--(-1,-2)--(-0.8,-1.6)--(-0.6,-1.2)--(-0.4,-0.8)--(-0.2,-0.4)--(0,0)--(0.2,0.4)--(0.4,0.8)--(0.6,1.2)--(0.8,1.6)--(1,2)--(1.2,2.4)--(1.4,2.8)--(1.6,3.2)--(1.8,3.6)--(2,4)--(2.2,4.4)--(2.4,4.8)--(2.6,5.2)--(2.8,5.6)--(3,6);

\end{tikzpicture}\\
    \end{minipage}
    
\end{exercice}

\begin{Solution}
 L'équation réduite de la tangente au point d'abscisse $1$ est  : $y=f'(1)(x-1)+f(1)$.\\
        On lit graphiquement $f(1)=-1$ et $f'(1)=2$.\\
        L'équation réduite de la tangente est donc donnée par :
        $y=2(x-1)-1$, soit $y=2x-3$.
\end{Solution}

% @see : https://coopmaths.fr/alea?uuid=a1ba2&id=can1F14&n=1&d=10&alea=rkIp2232323232342345678910111213141516171819202122232425262728293031323334353637383940414243444546474823456&cd=1&cols=1
\needspace{10\baselineskip}
\begin{exercice}[BaremeTotal=false]


Soit $f$ la fonction définie sur $\mathbb{R}$ par :
$f(x)= -5x-2+4x^2$.\\
En calculant la limite du taux d'accroissement, déterminer $f'(1)$.
\end{exercice}

\begin{Solution}
 $f$ est une fonction polynôme du second degré de la forme $f(x)=ax^2+bx+c$.\\
    La fonction dérivée est donnée par la somme des dérivées des fonctions $u$ et $v$ définies par $u(x)=4x^2$ et $v(x)=-5x-2$.\\
     Comme $u'(x)=8x$ et $v'(x)=-5$, on obtient  $f'(x)=8x-5$. \\
     Ainsi, $f'(1)=8\times 1-5=3$.
\end{Solution}

% @see : https://coopmaths.fr/alea?uuid=3c690&id=can1F13&n=1&d=10&alea=EUDu2232323232342345678910111213141516171819202122232425262728293031323334353637383940414243444546474823456&cd=1&cols=1
\needspace{10\baselineskip}
\begin{exercice}[BaremeTotal=false]


 Déterminer le coefficient directeur de la tangente à la courbe représentative de la fonction carré au point d'abscisse $-5$.    
\end{exercice}

\begin{Solution}
 Le coefficient directeur de la tangente au point d'abscisse $a$ est donné par le nombre dérivé $f'(a)$.\\
        La fonction $f$ définie par $f(x)=x^2$ a pour fonction dérivée la fonction $f'$ définie par $f'(x)=2x$.\\
        Comme $f'(-5)=2\times (-5)=-10$, le coefficient directeur de la tangente au point d'abscisse $-5$ est : $-10$. 
\end{Solution}

% @see : https://coopmaths.fr/alea?uuid=ebd89&id=1AN14-1&n=1&d=10&s=3&alea=ocYr2232323232342345678910111213141516171819202122232425262728293031323334353637383940414243444546474823456&cd=0&cols=1
\needspace{10\baselineskip}
\begin{exercice}[BaremeTotal=false]


 Donner la dérivée de la fonction $f$, dérivable pour tout $x\in\R$, définie par  $f(x)=\dfrac{8x+8}{9}$.
\end{exercice}

\begin{Solution}
 L'expression de la dérivée de la fonction $f$ définie par $f(x)=\dfrac{8x+8}{9}$ est : ${\color[HTML]{f15929}\boldsymbol{f'(x)=\dfrac{8}{9}}}$.
\end{Solution}

% @see : https://coopmaths.fr/alea?uuid=ebd89&id=1AN14-1&n=1&d=10&s=4&alea=uAlP2232323232342345678910111213141516171819202122232425262728293031323334353637383940414243444546474823456&cd=0&cols=1
\needspace{10\baselineskip}
\begin{exercice}[BaremeTotal=false]


 Donner la dérivée de la fonction $f$, dérivable pour tout $x\in\R$, définie par  $f(x)=-3x^{14}$.
\end{exercice}

\begin{Solution}
 L'expression de la dérivée de la fonction $f$ définie par $f(x)=-3x^{14}$ est : ${\color[HTML]{f15929}\boldsymbol{f'(x)=-42x^{13}}}$.
\end{Solution}

% @see : https://coopmaths.fr/alea?uuid=ebd89&id=1AN14-1&n=1&d=10&s=7&alea=vjN72232323232342345678910111213141516171819202122232425262728293031323334353637383940414243444546474823456&cd=0&cols=1
\needspace{10\baselineskip}
\begin{exercice}[BaremeTotal=false]


 Donner la dérivée de la fonction $f$, dérivable pour tout $x\in\R^*$, définie par  $f(x)=\dfrac{1}{9x}$.
\end{exercice}

\begin{Solution}
 L'expression de la dérivée de la fonction $f$ définie par $f(x)=\dfrac{1}{9x}$ est : ${\color[HTML]{f15929}\boldsymbol{f'(x)=-\dfrac{1}{9x^2}}}$.
\end{Solution}

% @see : https://coopmaths.fr/alea?uuid=a83c0&id=1AN14-4&n=1&d=10&s=1&alea=Njr32232323232342345678910111213141516171819202122232425262728293031323334353637383940414243444546474823456&cd=0&cols=1
\needspace{10\baselineskip}
\begin{exercice}[BaremeTotal=false]


 Donner l'expression de la dérivée de la fonction $f$ définie sur $\R^{*}$ par $f(x)=3x^{2}+3x+\dfrac{5}{x}$.\\
\end{exercice}

\begin{Solution}
 L'expression de la dérivée de la fonction $f$ définie par $f(x)=3x^{2}+3x+\dfrac{5}{x}$ est :\\$f'(x)={\color[HTML]{f15929}\boldsymbol{6x+3-\dfrac{5}{x^2}}}$.\\
\end{Solution}

% @see : https://coopmaths.fr/alea?uuid=a83c0&id=1AN14-4&n=1&d=10&s=4&alea=4Vpz2232323232342345678910111213141516171819202122232425262728293031323334353637383940414243444546474823456&cd=1&cols=1
\needspace{10\baselineskip}
\begin{exercice}[BaremeTotal=false]


 Donner l'expression de la dérivée de la fonction $f$ définie sur $\R_+^{*}$ par $f(x)=\dfrac{4}{x}-3x^{2}-2x+5\sqrt{x}$.\\
\end{exercice}

\begin{Solution}
 $(\dfrac{4}{x})^\prime=-\dfrac{4}{x^2}$\\$(-3x^{2}-2x)^\prime=-6x-2$\\$(5\sqrt{x})^\prime=\dfrac{5}{2\sqrt{x}}$\\L'expression de la dérivée de la fonction $f$ définie par $f(x)=\dfrac{4}{x}-3x^{2}-2x+5\sqrt{x}$ est :\\$f'(x)={\color[HTML]{f15929}\boldsymbol{-\dfrac{4}{x^2}-6x-2+\dfrac{5}{2\sqrt{x}}}}$.\\
\end{Solution}

\end{Maquette}
\clearpage
\begin{Maquette}[DM]{Niveau= ,Classe=\getdata[7]\mydata\ (v7), Date = 06/01/25  ,Theme=Exercices}


% @see : https://coopmaths.fr/alea?uuid=4c8c7&id=1AN11-3&n=1&d=10&s=2&alea=nawO22323232323423456789101112131415161718192021222324252627282930313233343536373839404142434445464748234567&cd=1&cols=2
\needspace{10\baselineskip}
\begin{exercice}[BaremeTotal=false]
\label{eleve:7}


 \begin{minipage}[t]{\linewidth} Soit $f$ une fonction dérivable sur $[-5;5]$ et $\mathcal{C}_f$ sa courbe représentative.\\On sait que  $f(3)=0~~$ et que $~~f'(3)=3$.\\Déterminer une équation de la tangente $(T)$ à la courbe $\mathcal{C}_f$ au point d'abscisse $3$,\\sans utiliser la formule de cours de l'équation de tangente. \end{minipage}

\end{exercice}

\begin{Solution}
 \begin{minipage}[t]{\linewidth} On sait que la tangente n'est pas une droite verticale, puisque la fonction est dérivable sur l'intervalle.\\On en déduit que la tangente $(T)$ au point d'abscisse $3$, admet une équation réduite de la forme :  $(T) : y=m x + p$. 

\medskip
$\bullet$ Détermination de $m$ :\\On sait que le nombre dérivé en $3$ est par définition, le coefficient directeur de la tangente au point d'abscisse $3$.\\Par conséquent, on a déjà : $m=f'(3)=3$.\\ On en déduit que  $(T) : y= 3 x + p$\\$\bullet$ Détermination de $p$ :\\ Pour cela, on utilise que si $f(3)=0~~$, alors le point $A$ de coordonnées $(3;0)$ appartient à $\mathcal{C}_f$ mais aussi à $(T)$.\\ On peut écrire $A(3;0) = \mathcal{C}_f \cap (T)$.\\ On remplace alors les coordonnées de $A(3;0)$ dans l'équation  $(T) : y= 3 x + p$.\\ $\begin{aligned} \phantom{\iff}&A(3;0)\in (T)\\ \iff& 0= 3 \times 3  + p\\ \iff& p=0 -3 \times 3  \\ \iff& p=0 -9   \\ \iff& p=-9\\\end{aligned}$\\On peut conclure que : $(T) : {\color[HTML]{f15929}\boldsymbol{y=3x-9}}$. \end{minipage}

\end{Solution}

% @see : https://coopmaths.fr/alea?uuid=0e984&id=can1F15&n=1&d=10&alea=1Alo22323232323423456789101112131415161718192021222324252627282930313233343536373839404142434445464748234567&cd=1&cols=1
\needspace{10\baselineskip}
\begin{exercice}[BaremeTotal=false]


 La courbe représente une fonction $f$ et la droite est la tangente au point d'abscisse $0$.\\
        Déterminer $f'(0)$.\begin{center}\begin{tikzpicture}[baseline,scale = 0.5]

    \tikzset{
      point/.style={
        thick,
        draw,
        cross out,
        inner sep=0pt,
        minimum width=5pt,
        minimum height=5pt,
      },
    }
    \clip (-8,-10) rectangle (8,3);
    	
	\draw[color ={black},line width = 2,>=latex,->] (-8,0)--(8,0);
	\draw[color ={black},line width = 2,>=latex,->] (0,-10)--(0,3);
	\draw[color ={black},opacity = 0.5] (-8,1)--(8,1);
	\draw[color ={black},opacity = 0.5] (-8,2)--(8,2);
	\draw[color ={black},opacity = 0.5] (-8,3)--(8,3);
	\draw[color ={black},opacity = 0.5] (-8,-1)--(8,-1);
	\draw[color ={black},opacity = 0.5] (-8,-2)--(8,-2);
	\draw[color ={black},opacity = 0.5] (-8,-3)--(8,-3);
	\draw[color ={black},opacity = 0.5] (-8,-4)--(8,-4);
	\draw[color ={black},opacity = 0.5] (-8,-5)--(8,-5);
	\draw[color ={black},opacity = 0.5] (-8,-6)--(8,-6);
	\draw[color ={black},opacity = 0.5] (-8,-7)--(8,-7);
	\draw[color ={black},opacity = 0.5] (-8,-8)--(8,-8);
	\draw[color ={black},opacity = 0.5] (-8,-9)--(8,-9);
	\draw[color ={black},opacity = 0.5] (-8,-10)--(8,-10);
	\draw[color ={black},opacity = 0.5] (1,-10)--(1,3);
	\draw[color ={black},opacity = 0.5] (2,-10)--(2,3);
	\draw[color ={black},opacity = 0.5] (3,-10)--(3,3);
	\draw[color ={black},opacity = 0.5] (4,-10)--(4,3);
	\draw[color ={black},opacity = 0.5] (5,-10)--(5,3);
	\draw[color ={black},opacity = 0.5] (6,-10)--(6,3);
	\draw[color ={black},opacity = 0.5] (7,-10)--(7,3);
	\draw[color ={black},opacity = 0.5] (8,-10)--(8,3);
	\draw[color ={black},opacity = 0.5] (-1,-10)--(-1,3);
	\draw[color ={black},opacity = 0.5] (-2,-10)--(-2,3);
	\draw[color ={black},opacity = 0.5] (-3,-10)--(-3,3);
	\draw[color ={black},opacity = 0.5] (-4,-10)--(-4,3);
	\draw[color ={black},opacity = 0.5] (-5,-10)--(-5,3);
	\draw[color ={black},opacity = 0.5] (-6,-10)--(-6,3);
	\draw[color ={black},opacity = 0.5] (-7,-10)--(-7,3);
	\draw[color ={black},opacity = 0.5] (-8,-10)--(-8,3);
	\draw[color ={gray},opacity = 0.3] (-8,0)--(8,0);
	\draw[color ={gray},opacity = 0.3] (-8,0)--(8,0);
	\draw[color ={gray},opacity = 0.3] (-8,-4)--(8,-4);
	\draw[color ={gray},opacity = 0.3] (-8,-5)--(8,-5);
	\draw[color ={gray},opacity = 0.3] (-8,-6)--(8,-6);
	\draw[color ={gray},opacity = 0.3] (-8,-7)--(8,-7);
	\draw[color ={gray},opacity = 0.3] (-8,-8)--(8,-8);
	\draw[color ={gray},opacity = 0.3] (-8,-9)--(8,-9);
	\draw[color ={gray},opacity = 0.3] (-8,-10)--(8,-10);
	\draw[color ={gray},opacity = 0.3] (0,-10)--(0,3);
	\draw[color ={gray},opacity = 0.3] (0,-10)--(0,3);
	\draw[color ={black},line width = 2] (2,-0.2)--(2,0.2);
	\draw[color ={black},line width = 2] (4,-0.2)--(4,0.2);
	\draw[color ={black},line width = 2] (6,-0.2)--(6,0.2);
	\draw[color ={black},line width = 2] (-2,-0.2)--(-2,0.2);
	\draw[color ={black},line width = 2] (-4,-0.2)--(-4,0.2);
	\draw[color ={black},line width = 2] (-6,-0.2)--(-6,0.2);
	\draw[color ={black},line width = 2] (-0.2,1)--(0.2,1);
	\draw[color ={black},line width = 2] (-0.2,2)--(0.2,2);
	\draw[color ={black},line width = 2] (-0.2,-1)--(0.2,-1);
	\draw[color ={black},line width = 2] (-0.2,-2)--(0.2,-2);
	\draw[color ={black},line width = 2] (-0.2,-3)--(0.2,-3);
	\draw[color ={black},line width = 2] (-0.2,-4)--(0.2,-4);
	\draw[color ={black},line width = 2] (-0.2,-5)--(0.2,-5);
	\draw[color ={black},line width = 2] (-0.2,-6)--(0.2,-6);
	\draw[color ={black},line width = 2] (-0.2,-7)--(0.2,-7);
	\draw[color ={black},line width = 2] (-0.2,-8)--(0.2,-8);
	\draw[color ={black},line width = 2] (-0.2,-9)--(0.2,-9);
	\draw (2,-0.4) node[anchor = center] {\scriptsize \color{black}{$1$}};
	\draw (4,-0.4) node[anchor = center] {\scriptsize \color{black}{$2$}};
	\draw (6,-0.4) node[anchor = center] {\scriptsize \color{black}{$3$}};
	\draw (-2,-0.4) node[anchor = center] {\scriptsize \color{black}{$-1$}};
	\draw (-4,-0.4) node[anchor = center] {\scriptsize \color{black}{$-2$}};
	\draw (-6,-0.4) node[anchor = center] {\scriptsize \color{black}{$-3$}};
	\draw (-0.8,2.1) node[anchor = center] {\footnotesize \color{black}{$2$}};
	\draw (-0.8,-1.9) node[anchor = center] {\footnotesize \color{black}{$-2$}};
	\draw (-0.8,-3.9) node[anchor = center] {\footnotesize \color{black}{$-4$}};
	\draw (-0.8,-5.9) node[anchor = center] {\footnotesize \color{black}{$-6$}};
	\draw (-0.8,-7.9) node[anchor = center] {\footnotesize \color{black}{$-8$}};
	\draw [color={black}] (-0.3,-0.3) node[anchor = center,scale=1, rotate = 0] {O};
	
	\draw[color={blue},line width = 2] (-8,-7)--(-7.6,-5.84)--(-7.2,-4.76)--(-6.8,-3.76)--(-6.4,-2.84)--(-6,-2)--(-5.6,-1.24)--(-5.2,-0.56)--(-4.8,0.04)--(-4.4,0.56)--(-4,1)--(-3.6,1.36)--(-3.2,1.64)--(-2.8,1.84)--(-2.4,1.96)--(-2,2)--(-1.6,1.96)--(-1.2,1.84)--(-0.8,1.64)--(-0.4,1.36)--(0,1)--(0.4,0.56)--(0.8,0.04)--(1.2,-0.56)--(1.6,-1.24)--(2,-2)--(2.4,-2.84)--(2.8,-3.76)--(3.2,-4.76)--(3.6,-5.84)--(4,-7)--(4.4,-8.24)--(4.8,-9.56)--(5.2,-10.96);
	
	\draw[color={red},line width = 2] (-2.8,3.8)--(-2.4,3.4)--(-2,3)--(-1.6,2.6)--(-1.2,2.2)--(-0.8,1.8)--(-0.4,1.4)--(0,1)--(0.4,0.6)--(0.8,0.2)--(1.2,-0.2)--(1.6,-0.6)--(2,-1)--(2.4,-1.4)--(2.8,-1.8)--(3.2,-2.2)--(3.6,-2.6)--(4,-3)--(4.4,-3.4)--(4.8,-3.8)--(5.2,-4.2)--(5.6,-4.6)--(6,-5)--(6.4,-5.4)--(6.8,-5.8)--(7.2,-6.2)--(7.6,-6.6)--(8,-7);

\end{tikzpicture}\end{center}
\end{exercice}

\begin{Solution}
 $f'(0)$ est donné par le coefficient directeur de la tangente à la courbe au point d'abscisse $0$, soit $-2$.
\end{Solution}

% @see : https://coopmaths.fr/alea?uuid=6f32d&id=can1F16&n=1&d=10&alea=dqSw22323232323423456789101112131415161718192021222324252627282930313233343536373839404142434445464748234567&cd=1&cols=1
\needspace{10\baselineskip}
\newpage
\begin{exercice}[BaremeTotal=false]


 On donne les représentations graphiques d'une fonction et de sa dérivée.\\
        Donner l'équation réduite de la tangente à la courbe de $f$ en $x=3$. \\ \begin{minipage}{0.5\linewidth}
    \begin{tikzpicture}[baseline, scale=0.8]

    \tikzset{
      point/.style={
        thick,
        draw,
        cross out,
        inner sep=0pt,
        minimum width=5pt,
        minimum height=5pt,
      },
    }
    \clip (-4.5,-4.5) rectangle (5.75,13.5);
    	
	\draw[color ={black},>=latex,->] (-3,0)--(3,0);
	\draw[color ={black},>=latex,->] (0,-3)--(0,12);
	\draw[color ={black},opacity = 0.4] (-3,1)--(3,1);
	\draw[color ={black},opacity = 0.4] (-3,2)--(3,2);
	\draw[color ={black},opacity = 0.4] (-3,3)--(3,3);
	\draw[color ={black},opacity = 0.4] (-3,4)--(3,4);
	\draw[color ={black},opacity = 0.4] (-3,5)--(3,5);
	\draw[color ={black},opacity = 0.4] (-3,6)--(3,6);
	\draw[color ={black},opacity = 0.4] (-3,7)--(3,7);
	\draw[color ={black},opacity = 0.4] (-3,8)--(3,8);
	\draw[color ={black},opacity = 0.4] (-3,9)--(3,9);
	\draw[color ={black},opacity = 0.4] (-3,10)--(3,10);
	\draw[color ={black},opacity = 0.4] (-3,11)--(3,11);
	\draw[color ={black},opacity = 0.4] (-3,12)--(3,12);
	\draw[color ={black},opacity = 0.4] (-3,-1)--(3,-1);
	\draw[color ={black},opacity = 0.4] (-3,-2)--(3,-2);
	\draw[color ={black},opacity = 0.4] (-3,-3)--(3,-3);
	\draw[color ={black},opacity = 0.4] (1,-3)--(1,12);
	\draw[color ={black},opacity = 0.4] (2,-3)--(2,12);
	\draw[color ={black},opacity = 0.4] (3,-3)--(3,12);
	\draw[color ={black},opacity = 0.4] (-1,-3)--(-1,12);
	\draw[color ={black},opacity = 0.4] (-2,-3)--(-2,12);
	\draw[color ={black},opacity = 0.4] (-3,-3)--(-3,12);
	\draw[color ={gray},opacity = 0.3] (-3,0)--(3,0);
	\draw[color ={gray},opacity = 0.3] (-3,0.5)--(3,0.5);
	\draw[color ={gray},opacity = 0.3] (-3,1.5)--(3,1.5);
	\draw[color ={gray},opacity = 0.3] (-3,2.5)--(3,2.5);
	\draw[color ={gray},opacity = 0.3] (-3,3.5)--(3,3.5);
	\draw[color ={gray},opacity = 0.3] (-3,4.5)--(3,4.5);
	\draw[color ={gray},opacity = 0.3] (-3,5.5)--(3,5.5);
	\draw[color ={gray},opacity = 0.3] (-3,6.5)--(3,6.5);
	\draw[color ={gray},opacity = 0.3] (-3,7.5)--(3,7.5);
	\draw[color ={gray},opacity = 0.3] (-3,8.5)--(3,8.5);
	\draw[color ={gray},opacity = 0.3] (-3,9.5)--(3,9.5);
	\draw[color ={gray},opacity = 0.3] (-3,10.5)--(3,10.5);
	\draw[color ={gray},opacity = 0.3] (-3,11.5)--(3,11.5);
	\draw[color ={gray},opacity = 0.3] (-3,0)--(3,0);
	\draw[color ={gray},opacity = 0.3] (-3,-0.5)--(3,-0.5);
	\draw[color ={gray},opacity = 0.3] (-3,-1.5)--(3,-1.5);
	\draw[color ={gray},opacity = 0.3] (-3,-2.5)--(3,-2.5);
	\draw[color ={gray},opacity = 0.3] (0,-3)--(0,12);
	\draw[color ={gray},opacity = 0.3] (0,-3)--(0,12);
	\draw[color ={black},line width = 1.2] (1,-0.13)--(1,0.13);
	\draw[color ={black},line width = 1.2] (2,-0.13)--(2,0.13);
	\draw[color ={black},line width = 1.2] (3,-0.13)--(3,0.13);
	\draw[color ={black},line width = 1.2] (-1,-0.13)--(-1,0.13);
	\draw[color ={black},line width = 1.2] (-2,-0.13)--(-2,0.13);
	\draw[color ={black},line width = 1.2] (-3,-0.13)--(-3,0.13);
	\draw[color ={black},line width = 1.2] (-0.13,1)--(0.13,1);
	\draw[color ={black},line width = 1.2] (-0.13,2)--(0.13,2);
	\draw[color ={black},line width = 1.2] (-0.13,3)--(0.13,3);
	\draw[color ={black},line width = 1.2] (-0.13,4)--(0.13,4);
	\draw[color ={black},line width = 1.2] (-0.13,5)--(0.13,5);
	\draw[color ={black},line width = 1.2] (-0.13,6)--(0.13,6);
	\draw[color ={black},line width = 1.2] (-0.13,7)--(0.13,7);
	\draw[color ={black},line width = 1.2] (-0.13,8)--(0.13,8);
	\draw[color ={black},line width = 1.2] (-0.13,9)--(0.13,9);
	\draw[color ={black},line width = 1.2] (-0.13,10)--(0.13,10);
	\draw[color ={black},line width = 1.2] (-0.13,11)--(0.13,11);
	\draw[color ={black},line width = 1.2] (-0.13,12)--(0.13,12);
	\draw[color ={black},line width = 1.2] (-0.13,-1)--(0.13,-1);
	\draw[color ={black},line width = 1.2] (-0.13,-2)--(0.13,-2);
	\draw[color ={black},line width = 1.2] (-0.13,-3)--(0.13,-3);
	\draw (3,-0.4) node[anchor = center] {\scriptsize \color{black}{$3$}};
	\draw (-3,-0.4) node[anchor = center] {\scriptsize \color{black}{$-3$}};
	\draw (-0.5,3.1) node[anchor = center] {\footnotesize \color{black}{$3$}};
	\draw (-0.5,6.1) node[anchor = center] {\footnotesize \color{black}{$6$}};
	\draw (-0.5,9.1) node[anchor = center] {\footnotesize \color{black}{$9$}};
	\draw (-0.5,-2.9) node[anchor = center] {\footnotesize \color{black}{$-3$}};
	\draw [color={black}] (-0.3,-0.3) node[anchor = center,scale=1, rotate = 0] {O};
	\draw (3,10) node[anchor = center] {\footnotesize \color{blue}{$\Large \mathcal{C}_f$}};
	
	\draw[color={blue},line width = 2] (-2.6,11.96)--(-2.4,10.56)--(-2.2,9.24)--(-2,8)--(-1.8,6.84)--(-1.6,5.76)--(-1.4,4.76)--(-1.2,3.84)--(-1,3)--(-0.8,2.24)--(-0.6,1.56)--(-0.4,0.96)--(-0.2,0.44)--(0,0)--(0.2,-0.36)--(0.4,-0.64)--(0.6,-0.84)--(0.8,-0.96)--(1,-1)--(1.2,-0.96)--(1.4,-0.84)--(1.6,-0.64)--(1.8,-0.36)--(2,0)--(2.2,0.44)--(2.4,0.96)--(2.6,1.56)--(2.8,2.24)--(3,3);

\end{tikzpicture}\\
    \end{minipage}
    \begin{minipage}{0.5\linewidth}
    \begin{tikzpicture}[baseline, scale=0.8]

    \tikzset{
      point/.style={
        thick,
        draw,
        cross out,
        inner sep=0pt,
        minimum width=5pt,
        minimum height=5pt,
      },
    }
    \clip (-4.5,-6.5) rectangle (6.5,9.5);
    	
	\draw[color ={black},>=latex,->] (-3,0)--(3,0);
	\draw[color ={black},>=latex,->] (0,-5)--(0,8);
	\draw[color ={black},opacity = 0.4] (-3,1)--(3,1);
	\draw[color ={black},opacity = 0.4] (-3,2)--(3,2);
	\draw[color ={black},opacity = 0.4] (-3,3)--(3,3);
	\draw[color ={black},opacity = 0.4] (-3,4)--(3,4);
	\draw[color ={black},opacity = 0.4] (-3,5)--(3,5);
	\draw[color ={black},opacity = 0.4] (-3,6)--(3,6);
	\draw[color ={black},opacity = 0.4] (-3,7)--(3,7);
	\draw[color ={black},opacity = 0.4] (-3,8)--(3,8);
	\draw[color ={black},opacity = 0.4] (-3,-1)--(3,-1);
	\draw[color ={black},opacity = 0.4] (-3,-2)--(3,-2);
	\draw[color ={black},opacity = 0.4] (-3,-3)--(3,-3);
	\draw[color ={black},opacity = 0.4] (-3,-4)--(3,-4);
	\draw[color ={black},opacity = 0.4] (-3,-5)--(3,-5);
	\draw[color ={black},opacity = 0.4] (1,-5)--(1,8);
	\draw[color ={black},opacity = 0.4] (2,-5)--(2,8);
	\draw[color ={black},opacity = 0.4] (3,-5)--(3,8);
	\draw[color ={black},opacity = 0.4] (-1,-5)--(-1,8);
	\draw[color ={black},opacity = 0.4] (-2,-5)--(-2,8);
	\draw[color ={black},opacity = 0.4] (-3,-5)--(-3,8);
	\draw[color ={gray},opacity = 0.3] (-3,0)--(3,0);
	\draw[color ={gray},opacity = 0.3] (-3,0.5)--(3,0.5);
	\draw[color ={gray},opacity = 0.3] (-3,1.5)--(3,1.5);
	\draw[color ={gray},opacity = 0.3] (-3,2.5)--(3,2.5);
	\draw[color ={gray},opacity = 0.3] (-3,3.5)--(3,3.5);
	\draw[color ={gray},opacity = 0.3] (-3,4.5)--(3,4.5);
	\draw[color ={gray},opacity = 0.3] (-3,5.5)--(3,5.5);
	\draw[color ={gray},opacity = 0.3] (-3,6.5)--(3,6.5);
	\draw[color ={gray},opacity = 0.3] (-3,7.5)--(3,7.5);
	\draw[color ={gray},opacity = 0.3] (-3,0)--(3,0);
	\draw[color ={gray},opacity = 0.3] (-3,-0.5)--(3,-0.5);
	\draw[color ={gray},opacity = 0.3] (-3,-1.5)--(3,-1.5);
	\draw[color ={gray},opacity = 0.3] (-3,-2.5)--(3,-2.5);
	\draw[color ={gray},opacity = 0.3] (-3,-3.5)--(3,-3.5);
	\draw[color ={gray},opacity = 0.3] (-3,-4.5)--(3,-4.5);
	\draw[color ={gray},opacity = 0.3] (0,-5)--(0,8);
	\draw[color ={gray},opacity = 0.3] (0,-5)--(0,8);
	\draw[color ={black},line width = 1.2] (1,-0.13)--(1,0.13);
	\draw[color ={black},line width = 1.2] (2,-0.13)--(2,0.13);
	\draw[color ={black},line width = 1.2] (3,-0.13)--(3,0.13);
	\draw[color ={black},line width = 1.2] (-1,-0.13)--(-1,0.13);
	\draw[color ={black},line width = 1.2] (-2,-0.13)--(-2,0.13);
	\draw[color ={black},line width = 1.2] (-3,-0.13)--(-3,0.13);
	\draw[color ={black},line width = 1.2] (-0.13,1)--(0.13,1);
	\draw[color ={black},line width = 1.2] (-0.13,2)--(0.13,2);
	\draw[color ={black},line width = 1.2] (-0.13,3)--(0.13,3);
	\draw[color ={black},line width = 1.2] (-0.13,4)--(0.13,4);
	\draw[color ={black},line width = 1.2] (-0.13,5)--(0.13,5);
	\draw[color ={black},line width = 1.2] (-0.13,6)--(0.13,6);
	\draw[color ={black},line width = 1.2] (-0.13,7)--(0.13,7);
	\draw[color ={black},line width = 1.2] (-0.13,8)--(0.13,8);
	\draw[color ={black},line width = 1.2] (-0.13,9)--(0.13,9);
	\draw[color ={black},line width = 1.2] (-0.13,10)--(0.13,10);
	\draw[color ={black},line width = 1.2] (-0.13,11)--(0.13,11);
	\draw[color ={black},line width = 1.2] (-0.13,12)--(0.13,12);
	\draw[color ={black},line width = 1.2] (-0.13,-1)--(0.13,-1);
	\draw[color ={black},line width = 1.2] (-0.13,-2)--(0.13,-2);
	\draw[color ={black},line width = 1.2] (-0.13,-3)--(0.13,-3);
	\draw (3,-0.4) node[anchor = center] {\scriptsize \color{black}{$3$}};
	\draw (-3,-0.4) node[anchor = center] {\scriptsize \color{black}{$-3$}};
	\draw (-0.5,3.1) node[anchor = center] {\footnotesize \color{black}{$3$}};
	\draw (-0.5,6.1) node[anchor = center] {\footnotesize \color{black}{$6$}};
	\draw (-0.5,-2.9) node[anchor = center] {\footnotesize \color{black}{$-3$}};
	\draw [color={black}] (-0.3,-0.3) node[anchor = center,scale=1, rotate = 0] {O};
	\draw (3,6) node[anchor = center] {\footnotesize \color{red}{$\Large\mathcal{C}_f\prime$}};
	
	\draw[color={red},line width = 2] (-2,-6)--(-1.8,-5.6)--(-1.6,-5.2)--(-1.4,-4.8)--(-1.2,-4.4)--(-1,-4)--(-0.8,-3.6)--(-0.6,-3.2)--(-0.4,-2.8)--(-0.2,-2.4)--(0,-2)--(0.2,-1.6)--(0.4,-1.2)--(0.6,-0.8)--(0.8,-0.4)--(1,0)--(1.2,0.4)--(1.4,0.8)--(1.6,1.2)--(1.8,1.6)--(2,2)--(2.2,2.4)--(2.4,2.8)--(2.6,3.2)--(2.8,3.6)--(3,4);

\end{tikzpicture}\\
    \end{minipage}
    
\end{exercice}

\begin{Solution}
 L'équation réduite de la tangente au point d'abscisse $3$ est  : $y=f'(3)(x-3)+f(3)$.\\
        On lit graphiquement $f(3)=3$ et $f'(3)=4$.\\
        L'équation réduite de la tangente est donc donnée par :
        $y=4(x-3)+3$, soit $y=4x-9$.
\end{Solution}

% @see : https://coopmaths.fr/alea?uuid=a1ba2&id=can1F14&n=1&d=10&alea=rkIp22323232323423456789101112131415161718192021222324252627282930313233343536373839404142434445464748234567&cd=1&cols=1
\needspace{10\baselineskip}
\begin{exercice}[BaremeTotal=false]


Soit $f$ la fonction définie sur $\mathbb{R}$ par :
$f(x)= -5x+7+5x^2$.\\
En calculant la limite du taux d'accroissement, déterminer $f'(-1)$.
\end{exercice}

\begin{Solution}
 $f$ est une fonction polynôme du second degré de la forme $f(x)=ax^2+bx+c$.\\
    La fonction dérivée est donnée par la somme des dérivées des fonctions $u$ et $v$ définies par $u(x)=5x^2$ et $v(x)=-5x+7$.\\
     Comme $u'(x)=10x$ et $v'(x)=-5$, on obtient  $f'(x)=10x-5$. \\
     Ainsi, $f'(-1)=10\times (-1)-5=-15$.
\end{Solution}

% @see : https://coopmaths.fr/alea?uuid=3c690&id=can1F13&n=1&d=10&alea=EUDu22323232323423456789101112131415161718192021222324252627282930313233343536373839404142434445464748234567&cd=1&cols=1
\needspace{10\baselineskip}
\begin{exercice}[BaremeTotal=false]


 Déterminer le coefficient directeur de la tangente à la courbe représentative de la fonction racine carrée au point d'abscisse $4$.
         
\end{exercice}

\begin{Solution}
 Le coefficient directeur de la tangente au point d'abscisse $a$ est donné par le nombre dérivé $f'(a)$.\\
        La fonction $f$ définie par $f(x)=\sqrt{x}$ a pour fonction dérivée la fonction $f'$ définie par $f'(x)=\dfrac{1}{2\sqrt{x}}$.\\Comme $f'(4)=\dfrac{1}{2\sqrt{4}}=\dfrac{1}{4}$, le coefficient directeur de la tangente au point d'abscisse $4$ est : $\dfrac{1}{4}$.
\end{Solution}

% @see : https://coopmaths.fr/alea?uuid=ebd89&id=1AN14-1&n=1&d=10&s=3&alea=ocYr22323232323423456789101112131415161718192021222324252627282930313233343536373839404142434445464748234567&cd=0&cols=1
\needspace{10\baselineskip}
\begin{exercice}[BaremeTotal=false]


 Donner la dérivée de la fonction $f$, dérivable pour tout $x\in\R$, définie par  $f(x)=\dfrac{2x+3}{3}$.
\end{exercice}

\begin{Solution}
 L'expression de la dérivée de la fonction $f$ définie par $f(x)=\dfrac{2x+3}{3}$ est : ${\color[HTML]{f15929}\boldsymbol{f'(x)=\dfrac{2}{3}}}$.
\end{Solution}

% @see : https://coopmaths.fr/alea?uuid=ebd89&id=1AN14-1&n=1&d=10&s=4&alea=uAlP22323232323423456789101112131415161718192021222324252627282930313233343536373839404142434445464748234567&cd=0&cols=1
\needspace{10\baselineskip}
\begin{exercice}[BaremeTotal=false]


 Donner la dérivée de la fonction $f$, dérivable pour tout $x\in\R$, définie par  $f(x)=-9x^{6}$.
\end{exercice}

\begin{Solution}
 L'expression de la dérivée de la fonction $f$ définie par $f(x)=-9x^{6}$ est : ${\color[HTML]{f15929}\boldsymbol{f'(x)=-54x^{5}}}$.
\end{Solution}

% @see : https://coopmaths.fr/alea?uuid=ebd89&id=1AN14-1&n=1&d=10&s=7&alea=vjN722323232323423456789101112131415161718192021222324252627282930313233343536373839404142434445464748234567&cd=0&cols=1
\needspace{10\baselineskip}
\begin{exercice}[BaremeTotal=false]


 Donner la dérivée de la fonction $f$, dérivable pour tout $x\in\R^*$, définie par  $f(x)=\dfrac{-1}{9x}$.
\end{exercice}

\begin{Solution}
 L'expression de la dérivée de la fonction $f$ définie par $f(x)=\dfrac{-1}{9x}$ est : ${\color[HTML]{f15929}\boldsymbol{f'(x)=\dfrac{1}{9x^2}}}$.
\end{Solution}

% @see : https://coopmaths.fr/alea?uuid=a83c0&id=1AN14-4&n=1&d=10&s=1&alea=Njr322323232323423456789101112131415161718192021222324252627282930313233343536373839404142434445464748234567&cd=0&cols=1
\needspace{10\baselineskip}
\begin{exercice}[BaremeTotal=false]


 Donner l'expression de la dérivée de la fonction $f$ définie sur $\R^{*}$ par $f(x)=\dfrac{5}{x}-3x^{2}$.\\
\end{exercice}

\begin{Solution}
 L'expression de la dérivée de la fonction $f$ définie par $f(x)=\dfrac{5}{x}-3x^{2}$ est :\\$f'(x)={\color[HTML]{f15929}\boldsymbol{-\dfrac{5}{x^2}-6x}}$.\\
\end{Solution}

% @see : https://coopmaths.fr/alea?uuid=a83c0&id=1AN14-4&n=1&d=10&s=4&alea=4Vpz22323232323423456789101112131415161718192021222324252627282930313233343536373839404142434445464748234567&cd=1&cols=1
\needspace{10\baselineskip}
\begin{exercice}[BaremeTotal=false]


 Donner l'expression de la dérivée de la fonction $f$ définie sur $\R_+^{*}$ par $f(x)=-4\sqrt{x}-\dfrac{3}{x}+5x^{2}-4x+3$.\\
\end{exercice}

\begin{Solution}
 $(-4\sqrt{x})^\prime=-\dfrac{4}{2\sqrt{x}}$\\$(-\dfrac{3}{x})^\prime=\dfrac{3}{x^2}$\\$(5x^{2}-4x+3)^\prime=10x-4$\\L'expression de la dérivée de la fonction $f$ définie par $f(x)=-4\sqrt{x}-\dfrac{3}{x}+5x^{2}-4x+3$ est :\\$f'(x)={\color[HTML]{f15929}\boldsymbol{-\dfrac{4}{2\sqrt{x}}+\dfrac{3}{x^2}+10x-4}}$.\\
\end{Solution}

\end{Maquette}
\clearpage
\begin{Maquette}[DM]{Niveau= ,Classe=\getdata[8]\mydata\ (v8), Date = 06/01/25  ,Theme=Exercices}


% @see : https://coopmaths.fr/alea?uuid=4c8c7&id=1AN11-3&n=1&d=10&s=2&alea=nawO223232323234234567891011121314151617181920212223242526272829303132333435363738394041424344454647482345678&cd=1&cols=2
\needspace{10\baselineskip}
\begin{exercice}[BaremeTotal=false]
\label{eleve:8}


 \begin{minipage}[t]{\linewidth} Soit $f$ une fonction dérivable sur $[-5;5]$ et $\mathcal{C}_f$ sa courbe représentative.\\On sait que  $f(-4)=4~~$ et que $~~f'(-4)=0$.\\Déterminer une équation de la tangente $(T)$ à la courbe $\mathcal{C}_f$ au point d'abscisse $-4$,\\sans utiliser la formule de cours de l'équation de tangente. \end{minipage}

\end{exercice}

\begin{Solution}
 \begin{minipage}[t]{\linewidth} On sait que la tangente n'est pas une droite verticale, puisque la fonction est dérivable sur l'intervalle.\\On en déduit que la tangente $(T)$ au point d'abscisse $-4$, admet une équation réduite de la forme :  $(T) : y=m x + p$. 

\medskip
$\bullet$ Détermination de $m$ :\\On sait que le nombre dérivé en $-4$ est par définition, le coefficient directeur de la tangente au point d'abscisse $-4$.\\Par conséquent, on a déjà : $m=f'(-4)=0$.\\ On en déduit que  $(T) : y= 0 x + p$\\$\bullet$ Détermination de $p$ :\\ Pour cela, on utilise que si $f(-4)=4~~$, alors le point $A$ de coordonnées $(-4;4)$ appartient à $\mathcal{C}_f$ mais aussi à $(T)$.\\ On peut écrire $A(-4;4) = \mathcal{C}_f \cap (T)$.\\ On remplace alors les coordonnées de $A(-4;4)$ dans l'équation  $(T) : y= 0 x + p$.\\ $\begin{aligned} \phantom{\iff}&A(-4;4)\in (T)\\ \iff& 4= 0 \times (-4)  + p\\ \iff& p=4 +0 \times (-4)  \\ \iff& p=4 +0   \\ \iff& p=4\\\end{aligned}$\\On peut conclure que : $(T) : {\color[HTML]{f15929}\boldsymbol{y=4}}$. \end{minipage}

\end{Solution}

% @see : https://coopmaths.fr/alea?uuid=0e984&id=can1F15&n=1&d=10&alea=1Alo223232323234234567891011121314151617181920212223242526272829303132333435363738394041424344454647482345678&cd=1&cols=1
\needspace{10\baselineskip}
\begin{exercice}[BaremeTotal=false]


 La courbe représente une fonction $f$ et la droite est la tangente au point d'abscisse $0$.\\
        Déterminer $f'(0)$.\begin{center}\begin{tikzpicture}[baseline,scale = 0.5]

    \tikzset{
      point/.style={
        thick,
        draw,
        cross out,
        inner sep=0pt,
        minimum width=5pt,
        minimum height=5pt,
      },
    }
    \clip (-10,-2) rectangle (4,12);
    	
	\draw[color ={black},line width = 2,>=latex,->] (-10,0)--(4,0);
	\draw[color ={black},line width = 2,>=latex,->] (0,-2)--(0,12);
	\draw[color ={black},opacity = 0.5] (-10,1)--(4,1);
	\draw[color ={black},opacity = 0.5] (-10,2)--(4,2);
	\draw[color ={black},opacity = 0.5] (-10,3)--(4,3);
	\draw[color ={black},opacity = 0.5] (-10,4)--(4,4);
	\draw[color ={black},opacity = 0.5] (-10,5)--(4,5);
	\draw[color ={black},opacity = 0.5] (-10,6)--(4,6);
	\draw[color ={black},opacity = 0.5] (-10,7)--(4,7);
	\draw[color ={black},opacity = 0.5] (-10,8)--(4,8);
	\draw[color ={black},opacity = 0.5] (-10,9)--(4,9);
	\draw[color ={black},opacity = 0.5] (-10,10)--(4,10);
	\draw[color ={black},opacity = 0.5] (-10,11)--(4,11);
	\draw[color ={black},opacity = 0.5] (-10,12)--(4,12);
	\draw[color ={black},opacity = 0.5] (-10,-1)--(4,-1);
	\draw[color ={black},opacity = 0.5] (-10,-2)--(4,-2);
	\draw[color ={black},opacity = 0.5] (1,-2)--(1,12);
	\draw[color ={black},opacity = 0.5] (2,-2)--(2,12);
	\draw[color ={black},opacity = 0.5] (3,-2)--(3,12);
	\draw[color ={black},opacity = 0.5] (4,-2)--(4,12);
	\draw[color ={black},opacity = 0.5] (-1,-2)--(-1,12);
	\draw[color ={black},opacity = 0.5] (-2,-2)--(-2,12);
	\draw[color ={black},opacity = 0.5] (-3,-2)--(-3,12);
	\draw[color ={black},opacity = 0.5] (-4,-2)--(-4,12);
	\draw[color ={black},opacity = 0.5] (-5,-2)--(-5,12);
	\draw[color ={black},opacity = 0.5] (-6,-2)--(-6,12);
	\draw[color ={black},opacity = 0.5] (-7,-2)--(-7,12);
	\draw[color ={black},opacity = 0.5] (-8,-2)--(-8,12);
	\draw[color ={black},opacity = 0.5] (-9,-2)--(-9,12);
	\draw[color ={black},opacity = 0.5] (-10,-2)--(-10,12);
	\draw[color ={gray},opacity = 0.3] (-10,0)--(4,0);
	\draw[color ={gray},opacity = 0.3] (-10,0)--(4,0);
	\draw[color ={gray},opacity = 0.3] (0,-2)--(0,12);
	\draw[color ={gray},opacity = 0.3] (0,-2)--(0,12);
	\draw[color ={gray},opacity = 0.3] (-5,-2)--(-5,12);
	\draw[color ={gray},opacity = 0.3] (-6,-2)--(-6,12);
	\draw[color ={gray},opacity = 0.3] (-7,-2)--(-7,12);
	\draw[color ={gray},opacity = 0.3] (-8,-2)--(-8,12);
	\draw[color ={gray},opacity = 0.3] (-9,-2)--(-9,12);
	\draw[color ={gray},opacity = 0.3] (-10,-2)--(-10,12);
	\draw[color ={black},line width = 2] (2,-0.2)--(2,0.2);
	\draw[color ={black},line width = 2] (-2,-0.2)--(-2,0.2);
	\draw[color ={black},line width = 2] (-4,-0.2)--(-4,0.2);
	\draw[color ={black},line width = 2] (-6,-0.2)--(-6,0.2);
	\draw[color ={black},line width = 2] (-8,-0.2)--(-8,0.2);
	\draw[color ={black},line width = 2] (-0.2,2)--(0.2,2);
	\draw[color ={black},line width = 2] (-0.2,4)--(0.2,4);
	\draw[color ={black},line width = 2] (-0.2,6)--(0.2,6);
	\draw[color ={black},line width = 2] (-0.2,8)--(0.2,8);
	\draw[color ={black},line width = 2] (-0.2,10)--(0.2,10);
	\draw (2,-0.4) node[anchor = center] {\scriptsize \color{black}{$1$}};
	\draw (-2,-0.4) node[anchor = center] {\scriptsize \color{black}{$-1$}};
	\draw (-4,-0.4) node[anchor = center] {\scriptsize \color{black}{$-2$}};
	\draw (-6,-0.4) node[anchor = center] {\scriptsize \color{black}{$-3$}};
	\draw (-8,-0.4) node[anchor = center] {\scriptsize \color{black}{$-4$}};
	\draw (-0.5,2.1) node[anchor = center] {\footnotesize \color{black}{$1$}};
	\draw (-0.5,4.1) node[anchor = center] {\footnotesize \color{black}{$2$}};
	\draw (-0.5,6.1) node[anchor = center] {\footnotesize \color{black}{$3$}};
	\draw (-0.5,8.1) node[anchor = center] {\footnotesize \color{black}{$4$}};
	\draw (-0.5,10.1) node[anchor = center] {\footnotesize \color{black}{$5$}};
	\draw [color={black}] (-0.3,-0.3) node[anchor = center,scale=1, rotate = 0] {O};
	
	\draw[color={blue},line width = 2] (-2.4,12.1)--(-2,8.96)--(-1.6,6.64)--(-1.2,4.92)--(-0.8,3.64)--(-0.4,2.7)--(0,2)--(0.4,1.48)--(0.8,1.1)--(1.2,0.81)--(1.6,0.6)--(2,0.45)--(2.4,0.33)--(2.8,0.24)--(3.2,0.18)--(3.6,0.13)--(4,0.1);
	
	\draw[color={red},line width = 2] (-8,14)--(-7.6,13.4)--(-7.2,12.8)--(-6.8,12.2)--(-6.4,11.6)--(-6,11)--(-5.6,10.4)--(-5.2,9.8)--(-4.8,9.2)--(-4.4,8.6)--(-4,8)--(-3.6,7.4)--(-3.2,6.8)--(-2.8,6.2)--(-2.4,5.6)--(-2,5)--(-1.6,4.4)--(-1.2,3.8)--(-0.8,3.2)--(-0.4,2.6)--(0,2)--(0.4,1.4)--(0.8,0.8)--(1.2,0.2)--(1.6,-0.4)--(2,-1)--(2.4,-1.6)--(2.8,-2.2)--(3.2,-2.8)--(3.6,-3.4);

\end{tikzpicture}\end{center}
\end{exercice}

\begin{Solution}
 $f'(0)$ est donné par le coefficient directeur de la tangente à la courbe au point d'abscisse $0$, soit $\dfrac{-3}{2}=-\dfrac{3}{2}$.
\end{Solution}

% @see : https://coopmaths.fr/alea?uuid=6f32d&id=can1F16&n=1&d=10&alea=dqSw223232323234234567891011121314151617181920212223242526272829303132333435363738394041424344454647482345678&cd=1&cols=1
\needspace{10\baselineskip}
\newpage
\begin{exercice}[BaremeTotal=false]


 On donne les représentations graphiques d'une fonction et de sa dérivée.\\
      Donner l'équation réduite de la tangente à la courbe de $f$ en $x=-2$. \\ \begin{minipage}{0.5\linewidth}
    \begin{tikzpicture}[baseline, scale=0.8]

    \tikzset{
      point/.style={
        thick,
        draw,
        cross out,
        inner sep=0pt,
        minimum width=5pt,
        minimum height=5pt,
      },
    }
    \clip (-5.5,-9.5) rectangle (5.75,7.5);
    	
	\draw[color ={black},>=latex,->] (-4,0)--(4,0);
	\draw[color ={black},>=latex,->] (0,-8)--(0,6);
	\draw[color ={black},opacity = 0.4] (-4,1)--(4,1);
	\draw[color ={black},opacity = 0.4] (-4,2)--(4,2);
	\draw[color ={black},opacity = 0.4] (-4,3)--(4,3);
	\draw[color ={black},opacity = 0.4] (-4,4)--(4,4);
	\draw[color ={black},opacity = 0.4] (-4,5)--(4,5);
	\draw[color ={black},opacity = 0.4] (-4,6)--(4,6);
	\draw[color ={black},opacity = 0.4] (-4,-1)--(4,-1);
	\draw[color ={black},opacity = 0.4] (-4,-2)--(4,-2);
	\draw[color ={black},opacity = 0.4] (-4,-3)--(4,-3);
	\draw[color ={black},opacity = 0.4] (-4,-4)--(4,-4);
	\draw[color ={black},opacity = 0.4] (-4,-5)--(4,-5);
	\draw[color ={black},opacity = 0.4] (-4,-6)--(4,-6);
	\draw[color ={black},opacity = 0.4] (-4,-7)--(4,-7);
	\draw[color ={black},opacity = 0.4] (-4,-8)--(4,-8);
	\draw[color ={black},opacity = 0.4] (1,-8)--(1,6);
	\draw[color ={black},opacity = 0.4] (2,-8)--(2,6);
	\draw[color ={black},opacity = 0.4] (3,-8)--(3,6);
	\draw[color ={black},opacity = 0.4] (4,-8)--(4,6);
	\draw[color ={black},opacity = 0.4] (-1,-8)--(-1,6);
	\draw[color ={black},opacity = 0.4] (-2,-8)--(-2,6);
	\draw[color ={black},opacity = 0.4] (-3,-8)--(-3,6);
	\draw[color ={black},opacity = 0.4] (-4,-8)--(-4,6);
	\draw[color ={gray},opacity = 0.3] (-4,0)--(4,0);
	\draw[color ={gray},opacity = 0.3] (-4,0.5)--(4,0.5);
	\draw[color ={gray},opacity = 0.3] (-4,1.5)--(4,1.5);
	\draw[color ={gray},opacity = 0.3] (-4,2.5)--(4,2.5);
	\draw[color ={gray},opacity = 0.3] (-4,3.5)--(4,3.5);
	\draw[color ={gray},opacity = 0.3] (-4,4.5)--(4,4.5);
	\draw[color ={gray},opacity = 0.3] (-4,5.5)--(4,5.5);
	\draw[color ={gray},opacity = 0.3] (-4,0)--(4,0);
	\draw[color ={gray},opacity = 0.3] (-4,-0.5)--(4,-0.5);
	\draw[color ={gray},opacity = 0.3] (-4,-1.5)--(4,-1.5);
	\draw[color ={gray},opacity = 0.3] (-4,-2.5)--(4,-2.5);
	\draw[color ={gray},opacity = 0.3] (-4,-3.5)--(4,-3.5);
	\draw[color ={gray},opacity = 0.3] (-4,-4.5)--(4,-4.5);
	\draw[color ={gray},opacity = 0.3] (-4,-5.5)--(4,-5.5);
	\draw[color ={gray},opacity = 0.3] (-4,-6.5)--(4,-6.5);
	\draw[color ={gray},opacity = 0.3] (-4,-7)--(4,-7);
	\draw[color ={gray},opacity = 0.3] (-4,-7.5)--(4,-7.5);
	\draw[color ={gray},opacity = 0.3] (-4,-8)--(4,-8);
	\draw[color ={gray},opacity = 0.3] (0,-8)--(0,6);
	\draw[color ={gray},opacity = 0.3] (0,-8)--(0,6);
	\draw[color ={black},line width = 1.2] (1,-0.13)--(1,0.13);
	\draw[color ={black},line width = 1.2] (2,-0.13)--(2,0.13);
	\draw[color ={black},line width = 1.2] (3,-0.13)--(3,0.13);
	\draw[color ={black},line width = 1.2] (4,-0.13)--(4,0.13);
	\draw[color ={black},line width = 1.2] (-1,-0.13)--(-1,0.13);
	\draw[color ={black},line width = 1.2] (-2,-0.13)--(-2,0.13);
	\draw[color ={black},line width = 1.2] (-3,-0.13)--(-3,0.13);
	\draw[color ={black},line width = 1.2] (-4,-0.13)--(-4,0.13);
	\draw[color ={black},line width = 1.2] (-0.13,1)--(0.13,1);
	\draw[color ={black},line width = 1.2] (-0.13,2)--(0.13,2);
	\draw[color ={black},line width = 1.2] (-0.13,3)--(0.13,3);
	\draw[color ={black},line width = 1.2] (-0.13,4)--(0.13,4);
	\draw[color ={black},line width = 1.2] (-0.13,5)--(0.13,5);
	\draw[color ={black},line width = 1.2] (-0.13,6)--(0.13,6);
	\draw[color ={black},line width = 1.2] (-0.13,-1)--(0.13,-1);
	\draw[color ={black},line width = 1.2] (-0.13,-2)--(0.13,-2);
	\draw[color ={black},line width = 1.2] (-0.13,-3)--(0.13,-3);
	\draw[color ={black},line width = 1.2] (-0.13,-4)--(0.13,-4);
	\draw[color ={black},line width = 1.2] (-0.13,-5)--(0.13,-5);
	\draw[color ={black},line width = 1.2] (-0.13,-6)--(0.13,-6);
	\draw[color ={black},line width = 1.2] (-0.13,-7)--(0.13,-7);
	\draw[color ={black},line width = 1.2] (-0.13,-8)--(0.13,-8);
	\draw (2,-0.4) node[anchor = center] {\scriptsize \color{black}{$2$}};
	\draw (4,-0.4) node[anchor = center] {\scriptsize \color{black}{$4$}};
	\draw (-2,-0.4) node[anchor = center] {\scriptsize \color{black}{$-2$}};
	\draw (-4,-0.4) node[anchor = center] {\scriptsize \color{black}{$-4$}};
	\draw (-0.5,2.1) node[anchor = center] {\footnotesize \color{black}{$2$}};
	\draw (-0.5,4.1) node[anchor = center] {\footnotesize \color{black}{$4$}};
	\draw (-0.5,6.1) node[anchor = center] {\footnotesize \color{black}{$6$}};
	\draw (-0.5,-1.9) node[anchor = center] {\footnotesize \color{black}{$-2$}};
	\draw (-0.5,-3.9) node[anchor = center] {\footnotesize \color{black}{$-4$}};
	\draw (-0.5,-5.9) node[anchor = center] {\footnotesize \color{black}{$-6$}};
	\draw (-0.5,-7.9) node[anchor = center] {\footnotesize \color{black}{$-8$}};
	\draw [color={black}] (-0.3,-0.3) node[anchor = center,scale=1, rotate = 0] {O};
	\draw (3,4) node[anchor = center] {\footnotesize \color{blue}{$\Large \mathcal{C}_f$}};
	
	\draw[color={blue},line width = 2] (-4,-6)--(-3.8,-4.84)--(-3.6,-3.76)--(-3.4,-2.76)--(-3.2,-1.84)--(-3,-1)--(-2.8,-0.24)--(-2.6,0.44)--(-2.4,1.04)--(-2.2,1.56)--(-2,2)--(-1.8,2.36)--(-1.6,2.64)--(-1.4,2.84)--(-1.2,2.96)--(-1,3)--(-0.8,2.96)--(-0.6,2.84)--(-0.4,2.64)--(-0.2,2.36)--(0,2)--(0.2,1.56)--(0.4,1.04)--(0.6,0.44)--(0.8,-0.24)--(1,-1)--(1.2,-1.84)--(1.4,-2.76)--(1.6,-3.76)--(1.8,-4.84)--(2,-6)--(2.2,-7.24)--(2.4,-8.56);

\end{tikzpicture}\\
    \end{minipage}
    \begin{minipage}{0.5\linewidth}
    \begin{tikzpicture}[baseline, scale=0.8]

    \tikzset{
      point/.style={
        thick,
        draw,
        cross out,
        inner sep=0pt,
        minimum width=5pt,
        minimum height=5pt,
      },
    }
    \clip (-5.5,-9.5) rectangle (6.5,7.5);
    	
	\draw[color ={black},>=latex,->] (-4,0)--(4,0);
	\draw[color ={black},>=latex,->] (0,-8)--(0,6);
	\draw[color ={black},opacity = 0.4] (-4,1)--(4,1);
	\draw[color ={black},opacity = 0.4] (-4,2)--(4,2);
	\draw[color ={black},opacity = 0.4] (-4,3)--(4,3);
	\draw[color ={black},opacity = 0.4] (-4,4)--(4,4);
	\draw[color ={black},opacity = 0.4] (-4,5)--(4,5);
	\draw[color ={black},opacity = 0.4] (-4,6)--(4,6);
	\draw[color ={black},opacity = 0.4] (-4,-1)--(4,-1);
	\draw[color ={black},opacity = 0.4] (-4,-2)--(4,-2);
	\draw[color ={black},opacity = 0.4] (-4,-3)--(4,-3);
	\draw[color ={black},opacity = 0.4] (-4,-4)--(4,-4);
	\draw[color ={black},opacity = 0.4] (-4,-5)--(4,-5);
	\draw[color ={black},opacity = 0.4] (-4,-6)--(4,-6);
	\draw[color ={black},opacity = 0.4] (-4,-7)--(4,-7);
	\draw[color ={black},opacity = 0.4] (-4,-8)--(4,-8);
	\draw[color ={black},opacity = 0.4] (1,-8)--(1,6);
	\draw[color ={black},opacity = 0.4] (2,-8)--(2,6);
	\draw[color ={black},opacity = 0.4] (3,-8)--(3,6);
	\draw[color ={black},opacity = 0.4] (4,-8)--(4,6);
	\draw[color ={black},opacity = 0.4] (-1,-8)--(-1,6);
	\draw[color ={black},opacity = 0.4] (-2,-8)--(-2,6);
	\draw[color ={black},opacity = 0.4] (-3,-8)--(-3,6);
	\draw[color ={black},opacity = 0.4] (-4,-8)--(-4,6);
	\draw[color ={gray},opacity = 0.3] (-4,0)--(4,0);
	\draw[color ={gray},opacity = 0.3] (-4,0.5)--(4,0.5);
	\draw[color ={gray},opacity = 0.3] (-4,1.5)--(4,1.5);
	\draw[color ={gray},opacity = 0.3] (-4,2.5)--(4,2.5);
	\draw[color ={gray},opacity = 0.3] (-4,3.5)--(4,3.5);
	\draw[color ={gray},opacity = 0.3] (-4,4.5)--(4,4.5);
	\draw[color ={gray},opacity = 0.3] (-4,5.5)--(4,5.5);
	\draw[color ={gray},opacity = 0.3] (-4,0)--(4,0);
	\draw[color ={gray},opacity = 0.3] (-4,-0.5)--(4,-0.5);
	\draw[color ={gray},opacity = 0.3] (-4,-1.5)--(4,-1.5);
	\draw[color ={gray},opacity = 0.3] (-4,-2.5)--(4,-2.5);
	\draw[color ={gray},opacity = 0.3] (-4,-3.5)--(4,-3.5);
	\draw[color ={gray},opacity = 0.3] (-4,-4.5)--(4,-4.5);
	\draw[color ={gray},opacity = 0.3] (-4,-5.5)--(4,-5.5);
	\draw[color ={gray},opacity = 0.3] (-4,-6.5)--(4,-6.5);
	\draw[color ={gray},opacity = 0.3] (-4,-7)--(4,-7);
	\draw[color ={gray},opacity = 0.3] (-4,-7.5)--(4,-7.5);
	\draw[color ={gray},opacity = 0.3] (-4,-8)--(4,-8);
	\draw[color ={gray},opacity = 0.3] (0,-8)--(0,6);
	\draw[color ={gray},opacity = 0.3] (0,-8)--(0,6);
	\draw[color ={black},line width = 1.2] (1,-0.13)--(1,0.13);
	\draw[color ={black},line width = 1.2] (2,-0.13)--(2,0.13);
	\draw[color ={black},line width = 1.2] (3,-0.13)--(3,0.13);
	\draw[color ={black},line width = 1.2] (4,-0.13)--(4,0.13);
	\draw[color ={black},line width = 1.2] (-1,-0.13)--(-1,0.13);
	\draw[color ={black},line width = 1.2] (-2,-0.13)--(-2,0.13);
	\draw[color ={black},line width = 1.2] (-3,-0.13)--(-3,0.13);
	\draw[color ={black},line width = 1.2] (-4,-0.13)--(-4,0.13);
	\draw[color ={black},line width = 1.2] (-0.13,1)--(0.13,1);
	\draw[color ={black},line width = 1.2] (-0.13,2)--(0.13,2);
	\draw[color ={black},line width = 1.2] (-0.13,3)--(0.13,3);
	\draw[color ={black},line width = 1.2] (-0.13,4)--(0.13,4);
	\draw[color ={black},line width = 1.2] (-0.13,5)--(0.13,5);
	\draw[color ={black},line width = 1.2] (-0.13,6)--(0.13,6);
	\draw[color ={black},line width = 1.2] (-0.13,-1)--(0.13,-1);
	\draw[color ={black},line width = 1.2] (-0.13,-2)--(0.13,-2);
	\draw[color ={black},line width = 1.2] (-0.13,-3)--(0.13,-3);
	\draw[color ={black},line width = 1.2] (-0.13,-4)--(0.13,-4);
	\draw[color ={black},line width = 1.2] (-0.13,-5)--(0.13,-5);
	\draw[color ={black},line width = 1.2] (-0.13,-6)--(0.13,-6);
	\draw[color ={black},line width = 1.2] (-0.13,-7)--(0.13,-7);
	\draw[color ={black},line width = 1.2] (-0.13,-8)--(0.13,-8);
	\draw (2,-0.4) node[anchor = center] {\scriptsize \color{black}{$2$}};
	\draw (4,-0.4) node[anchor = center] {\scriptsize \color{black}{$4$}};
	\draw (-2,-0.4) node[anchor = center] {\scriptsize \color{black}{$-2$}};
	\draw (-4,-0.4) node[anchor = center] {\scriptsize \color{black}{$-4$}};
	\draw (-0.5,2.1) node[anchor = center] {\footnotesize \color{black}{$2$}};
	\draw (-0.5,4.1) node[anchor = center] {\footnotesize \color{black}{$4$}};
	\draw (-0.5,-1.9) node[anchor = center] {\footnotesize \color{black}{$-2$}};
	\draw (-0.5,-3.9) node[anchor = center] {\footnotesize \color{black}{$-4$}};
	\draw (-0.5,-5.9) node[anchor = center] {\footnotesize \color{black}{$-6$}};
	\draw (-0.5,-7.9) node[anchor = center] {\footnotesize \color{black}{$-8$}};
	\draw [color={black}] (-0.3,-0.3) node[anchor = center,scale=1, rotate = 0] {O};
	\draw (3,4) node[anchor = center] {\footnotesize \color{red}{$\Large\mathcal{C}_f\prime$}};
	
	\draw[color={red},line width = 2] (-4,6)--(-3.8,5.6)--(-3.6,5.2)--(-3.4,4.8)--(-3.2,4.4)--(-3,4)--(-2.8,3.6)--(-2.6,3.2)--(-2.4,2.8)--(-2.2,2.4)--(-2,2)--(-1.8,1.6)--(-1.6,1.2)--(-1.4,0.8)--(-1.2,0.4)--(-1,0)--(-0.8,-0.4)--(-0.6,-0.8)--(-0.4,-1.2)--(-0.2,-1.6)--(0,-2)--(0.2,-2.4)--(0.4,-2.8)--(0.6,-3.2)--(0.8,-3.6)--(1,-4)--(1.2,-4.4)--(1.4,-4.8)--(1.6,-5.2)--(1.8,-5.6)--(2,-6)--(2.2,-6.4)--(2.4,-6.8)--(2.6,-7.2)--(2.8,-7.6)--(3,-8)--(3.2,-8.4)--(3.4,-8.8);

\end{tikzpicture}\\
    \end{minipage}
    
\end{exercice}

\begin{Solution}
 L'équation réduite de la tangente au point d'abscisse $-2$ est  : $y=f'(-2)(x-(-2))+f(-2)$.\\
      On lit graphiquement $f(-2)=2$ et $f'(-2)=2$.\\
      L'équation réduite de la tangente est donc donnée par :
      $y=2(x+2)+2$, soit $y=2x+6$.
\end{Solution}

% @see : https://coopmaths.fr/alea?uuid=a1ba2&id=can1F14&n=1&d=10&alea=rkIp223232323234234567891011121314151617181920212223242526272829303132333435363738394041424344454647482345678&cd=1&cols=1
\needspace{10\baselineskip}
\begin{exercice}[BaremeTotal=false]


Soit $f$ la fonction définie sur $\mathbb{R}$ par :
$f(x)= -2x+8-4x^2$.\\
En calculant la limite du taux d'accroissement, déterminer $f'(-2)$.
\end{exercice}

\begin{Solution}
 $f$ est une fonction polynôme du second degré de la forme $f(x)=ax^2+bx+c$.\\
    La fonction dérivée est donnée par la somme des dérivées des fonctions $u$ et $v$ définies par $u(x)=-4x^2$ et $v(x)=-2x+8$.\\
     Comme $u'(x)=-8x$ et $v'(x)=-2$, on obtient  $f'(x)=-8x-2$. \\
     Ainsi, $f'(-2)=-8\times (-2)-2=14$.
\end{Solution}

% @see : https://coopmaths.fr/alea?uuid=3c690&id=can1F13&n=1&d=10&alea=EUDu223232323234234567891011121314151617181920212223242526272829303132333435363738394041424344454647482345678&cd=1&cols=1
\needspace{10\baselineskip}
\begin{exercice}[BaremeTotal=false]


 Déterminer le coefficient directeur de la tangente à la courbe représentative de la fonction inverse au point d'abscisse $8$.
                 
\end{exercice}

\begin{Solution}
 Le coefficient directeur de la tangente au point d'abscisse $a$ est donné par le nombre dérivé $f'(a)$.\\
        La fonction $f$ définie par $f(x)=\dfrac{1}{x}$ a pour fonction dérivée la fonction $f'$ définie par $f'(x)=-\dfrac{1}{x^2}$.\\
Comme $f'(8)=-\dfrac{1}{8^2}=-\dfrac{1}{64}$, le coefficient directeur de la tangente au point d'abscisse $8$ est : $-\dfrac{1}{64}$.
\end{Solution}

% @see : https://coopmaths.fr/alea?uuid=ebd89&id=1AN14-1&n=1&d=10&s=3&alea=ocYr223232323234234567891011121314151617181920212223242526272829303132333435363738394041424344454647482345678&cd=0&cols=1
\needspace{10\baselineskip}
\begin{exercice}[BaremeTotal=false]


 Donner la dérivée de la fonction $f$, dérivable pour tout $x\in\R$, définie par  $f(x)=\dfrac{5x+2}{8}$.
\end{exercice}

\begin{Solution}
 L'expression de la dérivée de la fonction $f$ définie par $f(x)=\dfrac{5x+2}{8}$ est : ${\color[HTML]{f15929}\boldsymbol{f'(x)=\dfrac{5}{8}}}$.
\end{Solution}

% @see : https://coopmaths.fr/alea?uuid=ebd89&id=1AN14-1&n=1&d=10&s=4&alea=uAlP223232323234234567891011121314151617181920212223242526272829303132333435363738394041424344454647482345678&cd=0&cols=1
\needspace{10\baselineskip}
\begin{exercice}[BaremeTotal=false]


 Donner la dérivée de la fonction $f$, dérivable pour tout $x\in\R$, définie par  $f(x)=4x^{9}$.
\end{exercice}

\begin{Solution}
 L'expression de la dérivée de la fonction $f$ définie par $f(x)=4x^{9}$ est : ${\color[HTML]{f15929}\boldsymbol{f'(x)=36x^{8}}}$.
\end{Solution}

% @see : https://coopmaths.fr/alea?uuid=ebd89&id=1AN14-1&n=1&d=10&s=7&alea=vjN7223232323234234567891011121314151617181920212223242526272829303132333435363738394041424344454647482345678&cd=0&cols=1
\needspace{10\baselineskip}
\begin{exercice}[BaremeTotal=false]


 Donner la dérivée de la fonction $f$, dérivable pour tout $x\in\R^*$, définie par  $f(x)=\dfrac{3}{7x}$.
\end{exercice}

\begin{Solution}
 L'expression de la dérivée de la fonction $f$ définie par $f(x)=\dfrac{3}{7x}$ est : ${\color[HTML]{f15929}\boldsymbol{f'(x)=-\dfrac{3}{7x^2}}}$.
\end{Solution}

% @see : https://coopmaths.fr/alea?uuid=a83c0&id=1AN14-4&n=1&d=10&s=1&alea=Njr3223232323234234567891011121314151617181920212223242526272829303132333435363738394041424344454647482345678&cd=0&cols=1
\needspace{10\baselineskip}
\begin{exercice}[BaremeTotal=false]


 Donner l'expression de la dérivée de la fonction $f$ définie sur $\R^{*}$ par $f(x)=-\dfrac{9}{x}+3x^{2}+5x$.\\
\end{exercice}

\begin{Solution}
 L'expression de la dérivée de la fonction $f$ définie par $f(x)=-\dfrac{9}{x}+3x^{2}+5x$ est :\\$f'(x)={\color[HTML]{f15929}\boldsymbol{\dfrac{9}{x^2}+6x+5}}$.\\
\end{Solution}

% @see : https://coopmaths.fr/alea?uuid=a83c0&id=1AN14-4&n=1&d=10&s=4&alea=4Vpz223232323234234567891011121314151617181920212223242526272829303132333435363738394041424344454647482345678&cd=1&cols=1
\needspace{10\baselineskip}
\begin{exercice}[BaremeTotal=false]


 Donner l'expression de la dérivée de la fonction $f$ définie sur $\R_+^{*}$ par $f(x)=-2\sqrt{x}+5x^{2}-\dfrac{1}{x}$.\\
\end{exercice}

\begin{Solution}
 $(-2\sqrt{x})^\prime=-\dfrac{2}{2\sqrt{x}}$\\$(5x^{2})^\prime=10x$\\$(-\dfrac{1}{x})^\prime=\dfrac{1}{x^2}$\\L'expression de la dérivée de la fonction $f$ définie par $f(x)=-2\sqrt{x}+5x^{2}-\dfrac{1}{x}$ est :\\$f'(x)={\color[HTML]{f15929}\boldsymbol{-\dfrac{2}{2\sqrt{x}}+10x+\dfrac{1}{x^2}}}$.\\
\end{Solution}

\end{Maquette}
\clearpage
\begin{Maquette}[DM]{Niveau= ,Classe=\getdata[9]\mydata\ (v9), Date = 06/01/25  ,Theme=Exercices}


% @see : https://coopmaths.fr/alea?uuid=4c8c7&id=1AN11-3&n=1&d=10&s=2&alea=nawO2232323232342345678910111213141516171819202122232425262728293031323334353637383940414243444546474823456789&cd=1&cols=2
\needspace{10\baselineskip}
\begin{exercice}[BaremeTotal=false]
\label{eleve:9}


 \begin{minipage}[t]{\linewidth} Soit $f$ une fonction dérivable sur $[-5;5]$ et $\mathcal{C}_f$ sa courbe représentative.\\On sait que  $f(-5)=0~~$ et que $~~f'(-5)=3$.\\Déterminer une équation de la tangente $(T)$ à la courbe $\mathcal{C}_f$ au point d'abscisse $-5$,\\sans utiliser la formule de cours de l'équation de tangente. \end{minipage}

\end{exercice}

\begin{Solution}
 \begin{minipage}[t]{\linewidth} On sait que la tangente n'est pas une droite verticale, puisque la fonction est dérivable sur l'intervalle.\\On en déduit que la tangente $(T)$ au point d'abscisse $-5$, admet une équation réduite de la forme :  $(T) : y=m x + p$. 

\medskip
$\bullet$ Détermination de $m$ :\\On sait que le nombre dérivé en $-5$ est par définition, le coefficient directeur de la tangente au point d'abscisse $-5$.\\Par conséquent, on a déjà : $m=f'(-5)=3$.\\ On en déduit que  $(T) : y= 3 x + p$\\$\bullet$ Détermination de $p$ :\\ Pour cela, on utilise que si $f(-5)=0~~$, alors le point $A$ de coordonnées $(-5;0)$ appartient à $\mathcal{C}_f$ mais aussi à $(T)$.\\ On peut écrire $A(-5;0) = \mathcal{C}_f \cap (T)$.\\ On remplace alors les coordonnées de $A(-5;0)$ dans l'équation  $(T) : y= 3 x + p$.\\ $\begin{aligned} \phantom{\iff}&A(-5;0)\in (T)\\ \iff& 0= 3 \times (-5)  + p\\ \iff& p=0 -3 \times (-5)  \\ \iff& p=0 +15   \\ \iff& p=15\\\end{aligned}$\\On peut conclure que : $(T) : {\color[HTML]{f15929}\boldsymbol{y=3x+15}}$. \end{minipage}

\end{Solution}

% @see : https://coopmaths.fr/alea?uuid=0e984&id=can1F15&n=1&d=10&alea=1Alo2232323232342345678910111213141516171819202122232425262728293031323334353637383940414243444546474823456789&cd=1&cols=1
\needspace{10\baselineskip}
\begin{exercice}[BaremeTotal=false]


 La courbe représente une fonction $f$ et la droite est la tangente au point d'abscisse $1$.\\
        Déterminer $f'(1)$.\begin{center}\begin{tikzpicture}[baseline,scale = 0.5]

    \tikzset{
      point/.style={
        thick,
        draw,
        cross out,
        inner sep=0pt,
        minimum width=5pt,
        minimum height=5pt,
      },
    }
    \clip (-8,-10) rectangle (8,3);
    	
	\draw[color ={black},line width = 2,>=latex,->] (-8,0)--(8,0);
	\draw[color ={black},line width = 2,>=latex,->] (0,-10)--(0,3);
	\draw[color ={black},opacity = 0.5] (-8,1)--(8,1);
	\draw[color ={black},opacity = 0.5] (-8,2)--(8,2);
	\draw[color ={black},opacity = 0.5] (-8,3)--(8,3);
	\draw[color ={black},opacity = 0.5] (-8,-1)--(8,-1);
	\draw[color ={black},opacity = 0.5] (-8,-2)--(8,-2);
	\draw[color ={black},opacity = 0.5] (-8,-3)--(8,-3);
	\draw[color ={black},opacity = 0.5] (-8,-4)--(8,-4);
	\draw[color ={black},opacity = 0.5] (-8,-5)--(8,-5);
	\draw[color ={black},opacity = 0.5] (-8,-6)--(8,-6);
	\draw[color ={black},opacity = 0.5] (-8,-7)--(8,-7);
	\draw[color ={black},opacity = 0.5] (-8,-8)--(8,-8);
	\draw[color ={black},opacity = 0.5] (-8,-9)--(8,-9);
	\draw[color ={black},opacity = 0.5] (-8,-10)--(8,-10);
	\draw[color ={black},opacity = 0.5] (1,-10)--(1,3);
	\draw[color ={black},opacity = 0.5] (2,-10)--(2,3);
	\draw[color ={black},opacity = 0.5] (3,-10)--(3,3);
	\draw[color ={black},opacity = 0.5] (4,-10)--(4,3);
	\draw[color ={black},opacity = 0.5] (5,-10)--(5,3);
	\draw[color ={black},opacity = 0.5] (6,-10)--(6,3);
	\draw[color ={black},opacity = 0.5] (7,-10)--(7,3);
	\draw[color ={black},opacity = 0.5] (8,-10)--(8,3);
	\draw[color ={black},opacity = 0.5] (-1,-10)--(-1,3);
	\draw[color ={black},opacity = 0.5] (-2,-10)--(-2,3);
	\draw[color ={black},opacity = 0.5] (-3,-10)--(-3,3);
	\draw[color ={black},opacity = 0.5] (-4,-10)--(-4,3);
	\draw[color ={black},opacity = 0.5] (-5,-10)--(-5,3);
	\draw[color ={black},opacity = 0.5] (-6,-10)--(-6,3);
	\draw[color ={black},opacity = 0.5] (-7,-10)--(-7,3);
	\draw[color ={black},opacity = 0.5] (-8,-10)--(-8,3);
	\draw[color ={gray},opacity = 0.3] (-8,0)--(8,0);
	\draw[color ={gray},opacity = 0.3] (-8,0)--(8,0);
	\draw[color ={gray},opacity = 0.3] (-8,-4)--(8,-4);
	\draw[color ={gray},opacity = 0.3] (-8,-5)--(8,-5);
	\draw[color ={gray},opacity = 0.3] (-8,-6)--(8,-6);
	\draw[color ={gray},opacity = 0.3] (-8,-7)--(8,-7);
	\draw[color ={gray},opacity = 0.3] (-8,-8)--(8,-8);
	\draw[color ={gray},opacity = 0.3] (-8,-9)--(8,-9);
	\draw[color ={gray},opacity = 0.3] (-8,-10)--(8,-10);
	\draw[color ={gray},opacity = 0.3] (0,-10)--(0,3);
	\draw[color ={gray},opacity = 0.3] (0,-10)--(0,3);
	\draw[color ={black},line width = 2] (2,-0.2)--(2,0.2);
	\draw[color ={black},line width = 2] (4,-0.2)--(4,0.2);
	\draw[color ={black},line width = 2] (6,-0.2)--(6,0.2);
	\draw[color ={black},line width = 2] (-2,-0.2)--(-2,0.2);
	\draw[color ={black},line width = 2] (-4,-0.2)--(-4,0.2);
	\draw[color ={black},line width = 2] (-6,-0.2)--(-6,0.2);
	\draw[color ={black},line width = 2] (-0.2,1)--(0.2,1);
	\draw[color ={black},line width = 2] (-0.2,2)--(0.2,2);
	\draw[color ={black},line width = 2] (-0.2,-1)--(0.2,-1);
	\draw[color ={black},line width = 2] (-0.2,-2)--(0.2,-2);
	\draw[color ={black},line width = 2] (-0.2,-3)--(0.2,-3);
	\draw[color ={black},line width = 2] (-0.2,-4)--(0.2,-4);
	\draw[color ={black},line width = 2] (-0.2,-5)--(0.2,-5);
	\draw[color ={black},line width = 2] (-0.2,-6)--(0.2,-6);
	\draw[color ={black},line width = 2] (-0.2,-7)--(0.2,-7);
	\draw[color ={black},line width = 2] (-0.2,-8)--(0.2,-8);
	\draw[color ={black},line width = 2] (-0.2,-9)--(0.2,-9);
	\draw (2,-0.4) node[anchor = center] {\scriptsize \color{black}{$1$}};
	\draw (4,-0.4) node[anchor = center] {\scriptsize \color{black}{$2$}};
	\draw (6,-0.4) node[anchor = center] {\scriptsize \color{black}{$3$}};
	\draw (-2,-0.4) node[anchor = center] {\scriptsize \color{black}{$-1$}};
	\draw (-4,-0.4) node[anchor = center] {\scriptsize \color{black}{$-2$}};
	\draw (-6,-0.4) node[anchor = center] {\scriptsize \color{black}{$-3$}};
	\draw (-0.8,2.1) node[anchor = center] {\footnotesize \color{black}{$2$}};
	\draw (-0.8,-1.9) node[anchor = center] {\footnotesize \color{black}{$-2$}};
	\draw (-0.8,-3.9) node[anchor = center] {\footnotesize \color{black}{$-4$}};
	\draw (-0.8,-5.9) node[anchor = center] {\footnotesize \color{black}{$-6$}};
	\draw (-0.8,-7.9) node[anchor = center] {\footnotesize \color{black}{$-8$}};
	\draw [color={black}] (-0.3,-0.3) node[anchor = center,scale=1, rotate = 0] {O};
	
	\draw[color={blue},line width = 2] (-4.4,-10.24)--(-4,-9)--(-3.6,-7.84)--(-3.2,-6.76)--(-2.8,-5.76)--(-2.4,-4.84)--(-2,-4)--(-1.6,-3.24)--(-1.2,-2.56)--(-0.8,-1.96)--(-0.4,-1.44)--(0,-1)--(0.4,-0.64)--(0.8,-0.36)--(1.2,-0.16)--(1.6,-0.04)--(2,0)--(2.4,-0.04)--(2.8,-0.16)--(3.2,-0.36)--(3.6,-0.64)--(4,-1)--(4.4,-1.44)--(4.8,-1.96)--(5.2,-2.56)--(5.6,-3.24)--(6,-4)--(6.4,-4.84)--(6.8,-5.76)--(7.2,-6.76)--(7.6,-7.84)--(8,-9);
	
	\draw[color={red},line width = 2] (-8,0)--(-7.6,0)--(-7.2,0)--(-6.8,0)--(-6.4,0)--(-6,0)--(-5.6,0)--(-5.2,0)--(-4.8,0)--(-4.4,0)--(-4,0)--(-3.6,0)--(-3.2,0)--(-2.8,0)--(-2.4,0)--(-2,0)--(-1.6,0)--(-1.2,0)--(-0.8,0)--(-0.4,0)--(0,0)--(0.4,0)--(0.8,0)--(1.2,0)--(1.6,0)--(2,0)--(2.4,0)--(2.8,0)--(3.2,0)--(3.6,0)--(4,0)--(4.4,0)--(4.8,0)--(5.2,0)--(5.6,0)--(6,0)--(6.4,0)--(6.8,0)--(7.2,0)--(7.6,0)--(8,0);

\end{tikzpicture}\end{center}
\end{exercice}

\begin{Solution}
 $f'(1)$ est donné par le coefficient directeur de la tangente à la courbe au point d'abscisse $1$, soit $0$.
\end{Solution}

% @see : https://coopmaths.fr/alea?uuid=6f32d&id=can1F16&n=1&d=10&alea=dqSw2232323232342345678910111213141516171819202122232425262728293031323334353637383940414243444546474823456789&cd=1&cols=1
\needspace{10\baselineskip}
\newpage
\begin{exercice}[BaremeTotal=false]


 On donne les représentations graphiques d'une fonction et de sa dérivée.\\
        Donner l'équation réduite de la tangente à la courbe de $f$ en $x=-1$. \\ \begin{minipage}{0.5\linewidth}
    \begin{tikzpicture}[baseline, scale=0.8]

    \tikzset{
      point/.style={
        thick,
        draw,
        cross out,
        inner sep=0pt,
        minimum width=5pt,
        minimum height=5pt,
      },
    }
    \clip (-4.5,-4.5) rectangle (5.75,13.5);
    	
	\draw[color ={black},>=latex,->] (-3,0)--(3,0);
	\draw[color ={black},>=latex,->] (0,-3)--(0,12);
	\draw[color ={black},opacity = 0.4] (-3,1)--(3,1);
	\draw[color ={black},opacity = 0.4] (-3,2)--(3,2);
	\draw[color ={black},opacity = 0.4] (-3,3)--(3,3);
	\draw[color ={black},opacity = 0.4] (-3,4)--(3,4);
	\draw[color ={black},opacity = 0.4] (-3,5)--(3,5);
	\draw[color ={black},opacity = 0.4] (-3,6)--(3,6);
	\draw[color ={black},opacity = 0.4] (-3,7)--(3,7);
	\draw[color ={black},opacity = 0.4] (-3,8)--(3,8);
	\draw[color ={black},opacity = 0.4] (-3,9)--(3,9);
	\draw[color ={black},opacity = 0.4] (-3,10)--(3,10);
	\draw[color ={black},opacity = 0.4] (-3,11)--(3,11);
	\draw[color ={black},opacity = 0.4] (-3,12)--(3,12);
	\draw[color ={black},opacity = 0.4] (-3,-1)--(3,-1);
	\draw[color ={black},opacity = 0.4] (-3,-2)--(3,-2);
	\draw[color ={black},opacity = 0.4] (-3,-3)--(3,-3);
	\draw[color ={black},opacity = 0.4] (1,-3)--(1,12);
	\draw[color ={black},opacity = 0.4] (2,-3)--(2,12);
	\draw[color ={black},opacity = 0.4] (3,-3)--(3,12);
	\draw[color ={black},opacity = 0.4] (-1,-3)--(-1,12);
	\draw[color ={black},opacity = 0.4] (-2,-3)--(-2,12);
	\draw[color ={black},opacity = 0.4] (-3,-3)--(-3,12);
	\draw[color ={gray},opacity = 0.3] (-3,0)--(3,0);
	\draw[color ={gray},opacity = 0.3] (-3,0.5)--(3,0.5);
	\draw[color ={gray},opacity = 0.3] (-3,1.5)--(3,1.5);
	\draw[color ={gray},opacity = 0.3] (-3,2.5)--(3,2.5);
	\draw[color ={gray},opacity = 0.3] (-3,3.5)--(3,3.5);
	\draw[color ={gray},opacity = 0.3] (-3,4.5)--(3,4.5);
	\draw[color ={gray},opacity = 0.3] (-3,5.5)--(3,5.5);
	\draw[color ={gray},opacity = 0.3] (-3,6.5)--(3,6.5);
	\draw[color ={gray},opacity = 0.3] (-3,7.5)--(3,7.5);
	\draw[color ={gray},opacity = 0.3] (-3,8.5)--(3,8.5);
	\draw[color ={gray},opacity = 0.3] (-3,9.5)--(3,9.5);
	\draw[color ={gray},opacity = 0.3] (-3,10.5)--(3,10.5);
	\draw[color ={gray},opacity = 0.3] (-3,11.5)--(3,11.5);
	\draw[color ={gray},opacity = 0.3] (-3,0)--(3,0);
	\draw[color ={gray},opacity = 0.3] (-3,-0.5)--(3,-0.5);
	\draw[color ={gray},opacity = 0.3] (-3,-1.5)--(3,-1.5);
	\draw[color ={gray},opacity = 0.3] (-3,-2.5)--(3,-2.5);
	\draw[color ={gray},opacity = 0.3] (0,-3)--(0,12);
	\draw[color ={gray},opacity = 0.3] (0,-3)--(0,12);
	\draw[color ={black},line width = 1.2] (1,-0.13)--(1,0.13);
	\draw[color ={black},line width = 1.2] (2,-0.13)--(2,0.13);
	\draw[color ={black},line width = 1.2] (3,-0.13)--(3,0.13);
	\draw[color ={black},line width = 1.2] (-1,-0.13)--(-1,0.13);
	\draw[color ={black},line width = 1.2] (-2,-0.13)--(-2,0.13);
	\draw[color ={black},line width = 1.2] (-3,-0.13)--(-3,0.13);
	\draw[color ={black},line width = 1.2] (-0.13,1)--(0.13,1);
	\draw[color ={black},line width = 1.2] (-0.13,2)--(0.13,2);
	\draw[color ={black},line width = 1.2] (-0.13,3)--(0.13,3);
	\draw[color ={black},line width = 1.2] (-0.13,4)--(0.13,4);
	\draw[color ={black},line width = 1.2] (-0.13,5)--(0.13,5);
	\draw[color ={black},line width = 1.2] (-0.13,6)--(0.13,6);
	\draw[color ={black},line width = 1.2] (-0.13,7)--(0.13,7);
	\draw[color ={black},line width = 1.2] (-0.13,8)--(0.13,8);
	\draw[color ={black},line width = 1.2] (-0.13,9)--(0.13,9);
	\draw[color ={black},line width = 1.2] (-0.13,10)--(0.13,10);
	\draw[color ={black},line width = 1.2] (-0.13,11)--(0.13,11);
	\draw[color ={black},line width = 1.2] (-0.13,12)--(0.13,12);
	\draw[color ={black},line width = 1.2] (-0.13,-1)--(0.13,-1);
	\draw[color ={black},line width = 1.2] (-0.13,-2)--(0.13,-2);
	\draw[color ={black},line width = 1.2] (-0.13,-3)--(0.13,-3);
	\draw (3,-0.4) node[anchor = center] {\scriptsize \color{black}{$3$}};
	\draw (-3,-0.4) node[anchor = center] {\scriptsize \color{black}{$-3$}};
	\draw (-0.5,3.1) node[anchor = center] {\footnotesize \color{black}{$3$}};
	\draw (-0.5,6.1) node[anchor = center] {\footnotesize \color{black}{$6$}};
	\draw (-0.5,9.1) node[anchor = center] {\footnotesize \color{black}{$9$}};
	\draw (-0.5,-2.9) node[anchor = center] {\footnotesize \color{black}{$-3$}};
	\draw [color={black}] (-0.3,-0.3) node[anchor = center,scale=1, rotate = 0] {O};
	\draw (3,10) node[anchor = center] {\footnotesize \color{blue}{$\Large \mathcal{C}_f$}};
	
	\draw[color={blue},line width = 2] (-3,10)--(-2.8,8.84)--(-2.6,7.76)--(-2.4,6.76)--(-2.2,5.84)--(-2,5)--(-1.8,4.24)--(-1.6,3.56)--(-1.4,2.96)--(-1.2,2.44)--(-1,2)--(-0.8,1.64)--(-0.6,1.36)--(-0.4,1.16)--(-0.2,1.04)--(0,1)--(0.2,1.04)--(0.4,1.16)--(0.6,1.36)--(0.8,1.64)--(1,2)--(1.2,2.44)--(1.4,2.96)--(1.6,3.56)--(1.8,4.24)--(2,5)--(2.2,5.84)--(2.4,6.76)--(2.6,7.76)--(2.8,8.84)--(3,10);

\end{tikzpicture}\\
    \end{minipage}
    \begin{minipage}{0.5\linewidth}
    \begin{tikzpicture}[baseline, scale=0.8]

    \tikzset{
      point/.style={
        thick,
        draw,
        cross out,
        inner sep=0pt,
        minimum width=5pt,
        minimum height=5pt,
      },
    }
    \clip (-4.5,-6.5) rectangle (6.5,9.5);
    	
	\draw[color ={black},>=latex,->] (-3,0)--(3,0);
	\draw[color ={black},>=latex,->] (0,-5)--(0,8);
	\draw[color ={black},opacity = 0.4] (-3,1)--(3,1);
	\draw[color ={black},opacity = 0.4] (-3,2)--(3,2);
	\draw[color ={black},opacity = 0.4] (-3,3)--(3,3);
	\draw[color ={black},opacity = 0.4] (-3,4)--(3,4);
	\draw[color ={black},opacity = 0.4] (-3,5)--(3,5);
	\draw[color ={black},opacity = 0.4] (-3,6)--(3,6);
	\draw[color ={black},opacity = 0.4] (-3,7)--(3,7);
	\draw[color ={black},opacity = 0.4] (-3,8)--(3,8);
	\draw[color ={black},opacity = 0.4] (-3,-1)--(3,-1);
	\draw[color ={black},opacity = 0.4] (-3,-2)--(3,-2);
	\draw[color ={black},opacity = 0.4] (-3,-3)--(3,-3);
	\draw[color ={black},opacity = 0.4] (-3,-4)--(3,-4);
	\draw[color ={black},opacity = 0.4] (-3,-5)--(3,-5);
	\draw[color ={black},opacity = 0.4] (1,-5)--(1,8);
	\draw[color ={black},opacity = 0.4] (2,-5)--(2,8);
	\draw[color ={black},opacity = 0.4] (3,-5)--(3,8);
	\draw[color ={black},opacity = 0.4] (-1,-5)--(-1,8);
	\draw[color ={black},opacity = 0.4] (-2,-5)--(-2,8);
	\draw[color ={black},opacity = 0.4] (-3,-5)--(-3,8);
	\draw[color ={gray},opacity = 0.3] (-3,0)--(3,0);
	\draw[color ={gray},opacity = 0.3] (-3,0.5)--(3,0.5);
	\draw[color ={gray},opacity = 0.3] (-3,1.5)--(3,1.5);
	\draw[color ={gray},opacity = 0.3] (-3,2.5)--(3,2.5);
	\draw[color ={gray},opacity = 0.3] (-3,3.5)--(3,3.5);
	\draw[color ={gray},opacity = 0.3] (-3,4.5)--(3,4.5);
	\draw[color ={gray},opacity = 0.3] (-3,5.5)--(3,5.5);
	\draw[color ={gray},opacity = 0.3] (-3,6.5)--(3,6.5);
	\draw[color ={gray},opacity = 0.3] (-3,7.5)--(3,7.5);
	\draw[color ={gray},opacity = 0.3] (-3,0)--(3,0);
	\draw[color ={gray},opacity = 0.3] (-3,-0.5)--(3,-0.5);
	\draw[color ={gray},opacity = 0.3] (-3,-1.5)--(3,-1.5);
	\draw[color ={gray},opacity = 0.3] (-3,-2.5)--(3,-2.5);
	\draw[color ={gray},opacity = 0.3] (-3,-3.5)--(3,-3.5);
	\draw[color ={gray},opacity = 0.3] (-3,-4.5)--(3,-4.5);
	\draw[color ={gray},opacity = 0.3] (0,-5)--(0,8);
	\draw[color ={gray},opacity = 0.3] (0,-5)--(0,8);
	\draw[color ={black},line width = 1.2] (1,-0.13)--(1,0.13);
	\draw[color ={black},line width = 1.2] (2,-0.13)--(2,0.13);
	\draw[color ={black},line width = 1.2] (3,-0.13)--(3,0.13);
	\draw[color ={black},line width = 1.2] (-1,-0.13)--(-1,0.13);
	\draw[color ={black},line width = 1.2] (-2,-0.13)--(-2,0.13);
	\draw[color ={black},line width = 1.2] (-3,-0.13)--(-3,0.13);
	\draw[color ={black},line width = 1.2] (-0.13,1)--(0.13,1);
	\draw[color ={black},line width = 1.2] (-0.13,2)--(0.13,2);
	\draw[color ={black},line width = 1.2] (-0.13,3)--(0.13,3);
	\draw[color ={black},line width = 1.2] (-0.13,4)--(0.13,4);
	\draw[color ={black},line width = 1.2] (-0.13,5)--(0.13,5);
	\draw[color ={black},line width = 1.2] (-0.13,6)--(0.13,6);
	\draw[color ={black},line width = 1.2] (-0.13,7)--(0.13,7);
	\draw[color ={black},line width = 1.2] (-0.13,8)--(0.13,8);
	\draw[color ={black},line width = 1.2] (-0.13,9)--(0.13,9);
	\draw[color ={black},line width = 1.2] (-0.13,10)--(0.13,10);
	\draw[color ={black},line width = 1.2] (-0.13,11)--(0.13,11);
	\draw[color ={black},line width = 1.2] (-0.13,12)--(0.13,12);
	\draw[color ={black},line width = 1.2] (-0.13,-1)--(0.13,-1);
	\draw[color ={black},line width = 1.2] (-0.13,-2)--(0.13,-2);
	\draw[color ={black},line width = 1.2] (-0.13,-3)--(0.13,-3);
	\draw (3,-0.4) node[anchor = center] {\scriptsize \color{black}{$3$}};
	\draw (-3,-0.4) node[anchor = center] {\scriptsize \color{black}{$-3$}};
	\draw (-0.5,3.1) node[anchor = center] {\footnotesize \color{black}{$3$}};
	\draw (-0.5,6.1) node[anchor = center] {\footnotesize \color{black}{$6$}};
	\draw (-0.5,-2.9) node[anchor = center] {\footnotesize \color{black}{$-3$}};
	\draw [color={black}] (-0.3,-0.3) node[anchor = center,scale=1, rotate = 0] {O};
	\draw (3,6) node[anchor = center] {\footnotesize \color{red}{$\Large\mathcal{C}_f\prime$}};
	
	\draw[color={red},line width = 2] (-2.8,-5.6)--(-2.6,-5.2)--(-2.4,-4.8)--(-2.2,-4.4)--(-2,-4)--(-1.8,-3.6)--(-1.6,-3.2)--(-1.4,-2.8)--(-1.2,-2.4)--(-1,-2)--(-0.8,-1.6)--(-0.6,-1.2)--(-0.4,-0.8)--(-0.2,-0.4)--(0,0)--(0.2,0.4)--(0.4,0.8)--(0.6,1.2)--(0.8,1.6)--(1,2)--(1.2,2.4)--(1.4,2.8)--(1.6,3.2)--(1.8,3.6)--(2,4)--(2.2,4.4)--(2.4,4.8)--(2.6,5.2)--(2.8,5.6)--(3,6);

\end{tikzpicture}\\
    \end{minipage}
    
\end{exercice}

\begin{Solution}
 L'équation réduite de la tangente au point d'abscisse $-1$ est  : $y=f'(-1)(x-(-1))+f(-1)$.\\
        On lit graphiquement $f(-1)=2$ et $f'(-1)=-2$.\\
        L'équation réduite de la tangente est donc donnée par :
        $y=-2(x+1)+2$, soit $y=-2x$.
\end{Solution}

% @see : https://coopmaths.fr/alea?uuid=a1ba2&id=can1F14&n=1&d=10&alea=rkIp2232323232342345678910111213141516171819202122232425262728293031323334353637383940414243444546474823456789&cd=1&cols=1
\needspace{10\baselineskip}
\begin{exercice}[BaremeTotal=false]


 Soit $f$ la fonction définie sur $\mathbb{R}$ par : $f(x)= -2x^2+5x-7$.\\
        En calculant la limite du taux d'accroissement, déterminer $f'(-3)$.
\end{exercice}

\begin{Solution}
 $f$ est une fonction polynôme du second degré de la forme $f(x)=ax^2+bx+c$.\\
    La fonction dérivée est donnée par la somme des dérivées des fonctions $u$ et $v$ définies par $u(x)=-2x^2$ et $v(x)=5x-7$.\\
     Comme $u'(x)=-4x$ et $v'(x)=5$, on obtient  $f'(x)=-4x+5$. \\
     Ainsi, $f'(-3)=-4\times (-3)+5=17$.
\end{Solution}

% @see : https://coopmaths.fr/alea?uuid=3c690&id=can1F13&n=1&d=10&alea=EUDu2232323232342345678910111213141516171819202122232425262728293031323334353637383940414243444546474823456789&cd=1&cols=1
\needspace{10\baselineskip}
\begin{exercice}[BaremeTotal=false]


 Déterminer le coefficient directeur de la tangente à la courbe représentative de la fonction cube au point d'abscisse $3$.
                        
\end{exercice}

\begin{Solution}
 Le coefficient directeur de la tangente au point d'abscisse $a$ est donné par le nombre dérivé $f'(a)$.\\
        La fonction $f$ définie par $f(x)=x^3$ a pour fonction dérivée la fonction $f'$ définie par $f'(x)=3x^2$.\\
        Comme $f'(3)=3\times 3^2=27$, le coefficient directeur de la tangente au point d'abscisse $3$ est : $27$. 
\end{Solution}

% @see : https://coopmaths.fr/alea?uuid=ebd89&id=1AN14-1&n=1&d=10&s=3&alea=ocYr2232323232342345678910111213141516171819202122232425262728293031323334353637383940414243444546474823456789&cd=0&cols=1
\needspace{10\baselineskip}
\begin{exercice}[BaremeTotal=false]


 Donner la dérivée de la fonction $f$, dérivable pour tout $x\in\R$, définie par  $f(x)=\dfrac{-4x+4}{5}$.
\end{exercice}

\begin{Solution}
 L'expression de la dérivée de la fonction $f$ définie par $f(x)=\dfrac{-4x+4}{5}$ est : ${\color[HTML]{f15929}\boldsymbol{f'(x)=-\dfrac{4}{5}}}$.
\end{Solution}

% @see : https://coopmaths.fr/alea?uuid=ebd89&id=1AN14-1&n=1&d=10&s=4&alea=uAlP2232323232342345678910111213141516171819202122232425262728293031323334353637383940414243444546474823456789&cd=0&cols=1
\needspace{10\baselineskip}
\begin{exercice}[BaremeTotal=false]


 Donner la dérivée de la fonction $f$, dérivable pour tout $x\in\R$, définie par  $f(x)=3x^{11}$.
\end{exercice}

\begin{Solution}
 L'expression de la dérivée de la fonction $f$ définie par $f(x)=3x^{11}$ est : ${\color[HTML]{f15929}\boldsymbol{f'(x)=33x^{10}}}$.
\end{Solution}

% @see : https://coopmaths.fr/alea?uuid=ebd89&id=1AN14-1&n=1&d=10&s=7&alea=vjN72232323232342345678910111213141516171819202122232425262728293031323334353637383940414243444546474823456789&cd=0&cols=1
\needspace{10\baselineskip}
\begin{exercice}[BaremeTotal=false]


 Donner la dérivée de la fonction $f$, dérivable pour tout $x\in\R^*$, définie par  $f(x)=\dfrac{1}{4x}$.
\end{exercice}

\begin{Solution}
 L'expression de la dérivée de la fonction $f$ définie par $f(x)=\dfrac{1}{4x}$ est : ${\color[HTML]{f15929}\boldsymbol{f'(x)=-\dfrac{1}{4x^2}}}$.
\end{Solution}

% @see : https://coopmaths.fr/alea?uuid=a83c0&id=1AN14-4&n=1&d=10&s=1&alea=Njr32232323232342345678910111213141516171819202122232425262728293031323334353637383940414243444546474823456789&cd=0&cols=1
\needspace{10\baselineskip}
\begin{exercice}[BaremeTotal=false]


 Donner l'expression de la dérivée de la fonction $f$ définie sur $\R^{*}$ par $f(x)=2x^{2}+2x+\dfrac{7}{x}$.\\
\end{exercice}

\begin{Solution}
 L'expression de la dérivée de la fonction $f$ définie par $f(x)=2x^{2}+2x+\dfrac{7}{x}$ est :\\$f'(x)={\color[HTML]{f15929}\boldsymbol{4x+2-\dfrac{7}{x^2}}}$.\\
\end{Solution}

% @see : https://coopmaths.fr/alea?uuid=a83c0&id=1AN14-4&n=1&d=10&s=4&alea=4Vpz2232323232342345678910111213141516171819202122232425262728293031323334353637383940414243444546474823456789&cd=1&cols=1
\needspace{10\baselineskip}
\begin{exercice}[BaremeTotal=false]


 Donner l'expression de la dérivée de la fonction $f$ définie sur $\R_+^{*}$ par $f(x)=-\dfrac{4}{x}-4x^{2}-3x+3\sqrt{x}$.\\
\end{exercice}

\begin{Solution}
 $(-\dfrac{4}{x})^\prime=\dfrac{4}{x^2}$\\$(-4x^{2}-3x)^\prime=-8x-3$\\$(3\sqrt{x})^\prime=\dfrac{3}{2\sqrt{x}}$\\L'expression de la dérivée de la fonction $f$ définie par $f(x)=-\dfrac{4}{x}-4x^{2}-3x+3\sqrt{x}$ est :\\$f'(x)={\color[HTML]{f15929}\boldsymbol{\dfrac{4}{x^2}-8x-3+\dfrac{3}{2\sqrt{x}}}}$.\\
\end{Solution}

\end{Maquette}
\clearpage
\begin{Maquette}[DM]{Niveau= ,Classe=\getdata[10]\mydata\ (v10), Date = 06/01/25  ,Theme=Exercices}


% @see : https://coopmaths.fr/alea?uuid=4c8c7&id=1AN11-3&n=1&d=10&s=2&alea=nawO223232323234234567891011121314151617181920212223242526272829303132333435363738394041424344454647482345678910&cd=1&cols=2
\needspace{10\baselineskip}
\begin{exercice}[BaremeTotal=false]
\label{eleve:10}


 \begin{minipage}[t]{\linewidth} Soit $f$ une fonction dérivable sur $[-5;5]$ et $\mathcal{C}_f$ sa courbe représentative.\\On sait que  $f(-5)=2~~$ et que $~~f'(-5)=-1$.\\Déterminer une équation de la tangente $(T)$ à la courbe $\mathcal{C}_f$ au point d'abscisse $-5$,\\sans utiliser la formule de cours de l'équation de tangente. \end{minipage}

\end{exercice}

\begin{Solution}
 \begin{minipage}[t]{\linewidth} On sait que la tangente n'est pas une droite verticale, puisque la fonction est dérivable sur l'intervalle.\\On en déduit que la tangente $(T)$ au point d'abscisse $-5$, admet une équation réduite de la forme :  $(T) : y=m x + p$. 

\medskip
$\bullet$ Détermination de $m$ :\\On sait que le nombre dérivé en $-5$ est par définition, le coefficient directeur de la tangente au point d'abscisse $-5$.\\Par conséquent, on a déjà : $m=f'(-5)=-1$.\\ On en déduit que  $(T) : y= -1 x + p$\\$\bullet$ Détermination de $p$ :\\ Pour cela, on utilise que si $f(-5)=2~~$, alors le point $A$ de coordonnées $(-5;2)$ appartient à $\mathcal{C}_f$ mais aussi à $(T)$.\\ On peut écrire $A(-5;2) = \mathcal{C}_f \cap (T)$.\\ On remplace alors les coordonnées de $A(-5;2)$ dans l'équation  $(T) : y= -1 x + p$.\\ $\begin{aligned} \phantom{\iff}&A(-5;2)\in (T)\\ \iff& 2= -1 \times (-5)  + p\\ \iff& p=2 + \times (-5)  \\ \iff& p=2 -5   \\ \iff& p=-3\\\end{aligned}$\\On peut conclure que : $(T) : {\color[HTML]{f15929}\boldsymbol{y=-x-3}}$. \end{minipage}

\end{Solution}

% @see : https://coopmaths.fr/alea?uuid=0e984&id=can1F15&n=1&d=10&alea=1Alo223232323234234567891011121314151617181920212223242526272829303132333435363738394041424344454647482345678910&cd=1&cols=1
\needspace{10\baselineskip}
\begin{exercice}[BaremeTotal=false]


 La courbe représente une fonction $f$ et la droite est la tangente au point d'abscisse $0$.\\
        Déterminer $f'(0)$.\begin{center}\begin{tikzpicture}[baseline,scale = 0.5]

    \tikzset{
      point/.style={
        thick,
        draw,
        cross out,
        inner sep=0pt,
        minimum width=5pt,
        minimum height=5pt,
      },
    }
    \clip (-4,-2) rectangle (12,12);
    	
	\draw[color ={black},line width = 2,>=latex,->] (-4,0)--(12,0);
	\draw[color ={black},line width = 2,>=latex,->] (0,-2)--(0,12);
	\draw[color ={black},opacity = 0.5] (-4,1)--(12,1);
	\draw[color ={black},opacity = 0.5] (-4,2)--(12,2);
	\draw[color ={black},opacity = 0.5] (-4,3)--(12,3);
	\draw[color ={black},opacity = 0.5] (-4,4)--(12,4);
	\draw[color ={black},opacity = 0.5] (-4,5)--(12,5);
	\draw[color ={black},opacity = 0.5] (-4,6)--(12,6);
	\draw[color ={black},opacity = 0.5] (-4,7)--(12,7);
	\draw[color ={black},opacity = 0.5] (-4,8)--(12,8);
	\draw[color ={black},opacity = 0.5] (-4,9)--(12,9);
	\draw[color ={black},opacity = 0.5] (-4,10)--(12,10);
	\draw[color ={black},opacity = 0.5] (-4,11)--(12,11);
	\draw[color ={black},opacity = 0.5] (-4,12)--(12,12);
	\draw[color ={black},opacity = 0.5] (-4,-1)--(12,-1);
	\draw[color ={black},opacity = 0.5] (-4,-2)--(12,-2);
	\draw[color ={black},opacity = 0.5] (1,-2)--(1,12);
	\draw[color ={black},opacity = 0.5] (2,-2)--(2,12);
	\draw[color ={black},opacity = 0.5] (3,-2)--(3,12);
	\draw[color ={black},opacity = 0.5] (4,-2)--(4,12);
	\draw[color ={black},opacity = 0.5] (5,-2)--(5,12);
	\draw[color ={black},opacity = 0.5] (6,-2)--(6,12);
	\draw[color ={black},opacity = 0.5] (7,-2)--(7,12);
	\draw[color ={black},opacity = 0.5] (8,-2)--(8,12);
	\draw[color ={black},opacity = 0.5] (9,-2)--(9,12);
	\draw[color ={black},opacity = 0.5] (10,-2)--(10,12);
	\draw[color ={black},opacity = 0.5] (11,-2)--(11,12);
	\draw[color ={black},opacity = 0.5] (12,-2)--(12,12);
	\draw[color ={black},opacity = 0.5] (-1,-2)--(-1,12);
	\draw[color ={black},opacity = 0.5] (-2,-2)--(-2,12);
	\draw[color ={black},opacity = 0.5] (-3,-2)--(-3,12);
	\draw[color ={black},opacity = 0.5] (-4,-2)--(-4,12);
	\draw[color ={gray},opacity = 0.3] (-4,0)--(12,0);
	\draw[color ={gray},opacity = 0.3] (-4,0)--(12,0);
	\draw[color ={gray},opacity = 0.3] (0,-2)--(0,12);
	\draw[color ={gray},opacity = 0.3] (0,-2)--(0,12);
	\draw[color ={black},line width = 2] (2,-0.2)--(2,0.2);
	\draw[color ={black},line width = 2] (4,-0.2)--(4,0.2);
	\draw[color ={black},line width = 2] (6,-0.2)--(6,0.2);
	\draw[color ={black},line width = 2] (8,-0.2)--(8,0.2);
	\draw[color ={black},line width = 2] (10,-0.2)--(10,0.2);
	\draw[color ={black},line width = 2] (-2,-0.2)--(-2,0.2);
	\draw[color ={black},line width = 2] (-0.2,2)--(0.2,2);
	\draw[color ={black},line width = 2] (-0.2,4)--(0.2,4);
	\draw[color ={black},line width = 2] (-0.2,6)--(0.2,6);
	\draw[color ={black},line width = 2] (-0.2,8)--(0.2,8);
	\draw[color ={black},line width = 2] (-0.2,10)--(0.2,10);
	\draw (2,-0.4) node[anchor = center] {\scriptsize \color{black}{$1$}};
	\draw (4,-0.4) node[anchor = center] {\scriptsize \color{black}{$2$}};
	\draw (6,-0.4) node[anchor = center] {\scriptsize \color{black}{$3$}};
	\draw (8,-0.4) node[anchor = center] {\scriptsize \color{black}{$4$}};
	\draw (10,-0.4) node[anchor = center] {\scriptsize \color{black}{$5$}};
	\draw (-0.5,2.1) node[anchor = center] {\footnotesize \color{black}{$1$}};
	\draw (-0.5,4.1) node[anchor = center] {\footnotesize \color{black}{$2$}};
	\draw (-0.5,6.1) node[anchor = center] {\footnotesize \color{black}{$3$}};
	\draw (-0.5,8.1) node[anchor = center] {\footnotesize \color{black}{$4$}};
	\draw (-0.5,10.1) node[anchor = center] {\footnotesize \color{black}{$5$}};
	\draw [color={black}] (-0.3,-0.3) node[anchor = center,scale=1, rotate = 0] {O};
	
	\draw[color={blue},line width = 2] (-4,0.74)--(-3.6,0.81)--(-3.2,0.9)--(-2.8,0.99)--(-2.4,1.1)--(-2,1.21)--(-1.6,1.34)--(-1.2,1.48)--(-0.8,1.64)--(-0.4,1.81)--(0,2)--(0.4,2.21)--(0.8,2.44)--(1.2,2.7)--(1.6,2.98)--(2,3.3)--(2.4,3.64)--(2.8,4.03)--(3.2,4.45)--(3.6,4.92)--(4,5.44)--(4.4,6.01)--(4.8,6.64)--(5.2,7.34)--(5.6,8.11)--(6,8.96)--(6.4,9.91)--(6.8,10.95)--(7.2,12.1)--(7.6,13.37);
	
	\draw[color={red},line width = 2] (-4,0)--(-3.6,0.2)--(-3.2,0.4)--(-2.8,0.6)--(-2.4,0.8)--(-2,1)--(-1.6,1.2)--(-1.2,1.4)--(-0.8,1.6)--(-0.4,1.8)--(0,2)--(0.4,2.2)--(0.8,2.4)--(1.2,2.6)--(1.6,2.8)--(2,3)--(2.4,3.2)--(2.8,3.4)--(3.2,3.6)--(3.6,3.8)--(4,4)--(4.4,4.2)--(4.8,4.4)--(5.2,4.6)--(5.6,4.8)--(6,5)--(6.4,5.2)--(6.8,5.4)--(7.2,5.6)--(7.6,5.8)--(8,6)--(8.4,6.2)--(8.8,6.4)--(9.2,6.6)--(9.6,6.8)--(10,7)--(10.4,7.2)--(10.8,7.4)--(11.2,7.6)--(11.6,7.8)--(12,8);

\end{tikzpicture}\end{center}
\end{exercice}

\begin{Solution}
 $f'(0)$ est donné par le coefficient directeur de la tangente à la courbe au point d'abscisse $0$, soit $\dfrac{1}{2}$.
\end{Solution}

% @see : https://coopmaths.fr/alea?uuid=6f32d&id=can1F16&n=1&d=10&alea=dqSw223232323234234567891011121314151617181920212223242526272829303132333435363738394041424344454647482345678910&cd=1&cols=1
\needspace{10\baselineskip}
\newpage
\begin{exercice}[BaremeTotal=false]


 On donne les représentations graphiques d'une fonction et de sa dérivée.\\
      Donner l'équation réduite de la tangente à la courbe de $f$ en $x=0$. \\ \begin{minipage}{0.5\linewidth}
    \begin{tikzpicture}[baseline, scale=0.8]

    \tikzset{
      point/.style={
        thick,
        draw,
        cross out,
        inner sep=0pt,
        minimum width=5pt,
        minimum height=5pt,
      },
    }
    \clip (-5.5,-9.5) rectangle (5.75,7.5);
    	
	\draw[color ={black},>=latex,->] (-4,0)--(4,0);
	\draw[color ={black},>=latex,->] (0,-8)--(0,6);
	\draw[color ={black},opacity = 0.4] (-4,1)--(4,1);
	\draw[color ={black},opacity = 0.4] (-4,2)--(4,2);
	\draw[color ={black},opacity = 0.4] (-4,3)--(4,3);
	\draw[color ={black},opacity = 0.4] (-4,4)--(4,4);
	\draw[color ={black},opacity = 0.4] (-4,5)--(4,5);
	\draw[color ={black},opacity = 0.4] (-4,6)--(4,6);
	\draw[color ={black},opacity = 0.4] (-4,-1)--(4,-1);
	\draw[color ={black},opacity = 0.4] (-4,-2)--(4,-2);
	\draw[color ={black},opacity = 0.4] (-4,-3)--(4,-3);
	\draw[color ={black},opacity = 0.4] (-4,-4)--(4,-4);
	\draw[color ={black},opacity = 0.4] (-4,-5)--(4,-5);
	\draw[color ={black},opacity = 0.4] (-4,-6)--(4,-6);
	\draw[color ={black},opacity = 0.4] (-4,-7)--(4,-7);
	\draw[color ={black},opacity = 0.4] (-4,-8)--(4,-8);
	\draw[color ={black},opacity = 0.4] (1,-8)--(1,6);
	\draw[color ={black},opacity = 0.4] (2,-8)--(2,6);
	\draw[color ={black},opacity = 0.4] (3,-8)--(3,6);
	\draw[color ={black},opacity = 0.4] (4,-8)--(4,6);
	\draw[color ={black},opacity = 0.4] (-1,-8)--(-1,6);
	\draw[color ={black},opacity = 0.4] (-2,-8)--(-2,6);
	\draw[color ={black},opacity = 0.4] (-3,-8)--(-3,6);
	\draw[color ={black},opacity = 0.4] (-4,-8)--(-4,6);
	\draw[color ={gray},opacity = 0.3] (-4,0)--(4,0);
	\draw[color ={gray},opacity = 0.3] (-4,0.5)--(4,0.5);
	\draw[color ={gray},opacity = 0.3] (-4,1.5)--(4,1.5);
	\draw[color ={gray},opacity = 0.3] (-4,2.5)--(4,2.5);
	\draw[color ={gray},opacity = 0.3] (-4,3.5)--(4,3.5);
	\draw[color ={gray},opacity = 0.3] (-4,4.5)--(4,4.5);
	\draw[color ={gray},opacity = 0.3] (-4,5.5)--(4,5.5);
	\draw[color ={gray},opacity = 0.3] (-4,0)--(4,0);
	\draw[color ={gray},opacity = 0.3] (-4,-0.5)--(4,-0.5);
	\draw[color ={gray},opacity = 0.3] (-4,-1.5)--(4,-1.5);
	\draw[color ={gray},opacity = 0.3] (-4,-2.5)--(4,-2.5);
	\draw[color ={gray},opacity = 0.3] (-4,-3.5)--(4,-3.5);
	\draw[color ={gray},opacity = 0.3] (-4,-4.5)--(4,-4.5);
	\draw[color ={gray},opacity = 0.3] (-4,-5.5)--(4,-5.5);
	\draw[color ={gray},opacity = 0.3] (-4,-6.5)--(4,-6.5);
	\draw[color ={gray},opacity = 0.3] (-4,-7)--(4,-7);
	\draw[color ={gray},opacity = 0.3] (-4,-7.5)--(4,-7.5);
	\draw[color ={gray},opacity = 0.3] (-4,-8)--(4,-8);
	\draw[color ={gray},opacity = 0.3] (0,-8)--(0,6);
	\draw[color ={gray},opacity = 0.3] (0,-8)--(0,6);
	\draw[color ={black},line width = 1.2] (1,-0.13)--(1,0.13);
	\draw[color ={black},line width = 1.2] (2,-0.13)--(2,0.13);
	\draw[color ={black},line width = 1.2] (3,-0.13)--(3,0.13);
	\draw[color ={black},line width = 1.2] (4,-0.13)--(4,0.13);
	\draw[color ={black},line width = 1.2] (-1,-0.13)--(-1,0.13);
	\draw[color ={black},line width = 1.2] (-2,-0.13)--(-2,0.13);
	\draw[color ={black},line width = 1.2] (-3,-0.13)--(-3,0.13);
	\draw[color ={black},line width = 1.2] (-4,-0.13)--(-4,0.13);
	\draw[color ={black},line width = 1.2] (-0.13,1)--(0.13,1);
	\draw[color ={black},line width = 1.2] (-0.13,2)--(0.13,2);
	\draw[color ={black},line width = 1.2] (-0.13,3)--(0.13,3);
	\draw[color ={black},line width = 1.2] (-0.13,4)--(0.13,4);
	\draw[color ={black},line width = 1.2] (-0.13,5)--(0.13,5);
	\draw[color ={black},line width = 1.2] (-0.13,6)--(0.13,6);
	\draw[color ={black},line width = 1.2] (-0.13,-1)--(0.13,-1);
	\draw[color ={black},line width = 1.2] (-0.13,-2)--(0.13,-2);
	\draw[color ={black},line width = 1.2] (-0.13,-3)--(0.13,-3);
	\draw[color ={black},line width = 1.2] (-0.13,-4)--(0.13,-4);
	\draw[color ={black},line width = 1.2] (-0.13,-5)--(0.13,-5);
	\draw[color ={black},line width = 1.2] (-0.13,-6)--(0.13,-6);
	\draw[color ={black},line width = 1.2] (-0.13,-7)--(0.13,-7);
	\draw[color ={black},line width = 1.2] (-0.13,-8)--(0.13,-8);
	\draw (2,-0.4) node[anchor = center] {\scriptsize \color{black}{$2$}};
	\draw (4,-0.4) node[anchor = center] {\scriptsize \color{black}{$4$}};
	\draw (-2,-0.4) node[anchor = center] {\scriptsize \color{black}{$-2$}};
	\draw (-4,-0.4) node[anchor = center] {\scriptsize \color{black}{$-4$}};
	\draw (-0.5,2.1) node[anchor = center] {\footnotesize \color{black}{$2$}};
	\draw (-0.5,4.1) node[anchor = center] {\footnotesize \color{black}{$4$}};
	\draw (-0.5,6.1) node[anchor = center] {\footnotesize \color{black}{$6$}};
	\draw (-0.5,-1.9) node[anchor = center] {\footnotesize \color{black}{$-2$}};
	\draw (-0.5,-3.9) node[anchor = center] {\footnotesize \color{black}{$-4$}};
	\draw (-0.5,-5.9) node[anchor = center] {\footnotesize \color{black}{$-6$}};
	\draw (-0.5,-7.9) node[anchor = center] {\footnotesize \color{black}{$-8$}};
	\draw [color={black}] (-0.3,-0.3) node[anchor = center,scale=1, rotate = 0] {O};
	\draw (3,4) node[anchor = center] {\footnotesize \color{blue}{$\Large \mathcal{C}_f$}};
	
	\draw[color={blue},line width = 2] (-4,-8)--(-3.8,-6.84)--(-3.6,-5.76)--(-3.4,-4.76)--(-3.2,-3.84)--(-3,-3)--(-2.8,-2.24)--(-2.6,-1.56)--(-2.4,-0.96)--(-2.2,-0.44)--(-2,0)--(-1.8,0.36)--(-1.6,0.64)--(-1.4,0.84)--(-1.2,0.96)--(-1,1)--(-0.8,0.96)--(-0.6,0.84)--(-0.4,0.64)--(-0.2,0.36)--(0,0)--(0.2,-0.44)--(0.4,-0.96)--(0.6,-1.56)--(0.8,-2.24)--(1,-3)--(1.2,-3.84)--(1.4,-4.76)--(1.6,-5.76)--(1.8,-6.84)--(2,-8);

\end{tikzpicture}\\
    \end{minipage}
    \begin{minipage}{0.5\linewidth}
    \begin{tikzpicture}[baseline, scale=0.8]

    \tikzset{
      point/.style={
        thick,
        draw,
        cross out,
        inner sep=0pt,
        minimum width=5pt,
        minimum height=5pt,
      },
    }
    \clip (-5.5,-9.5) rectangle (6.5,7.5);
    	
	\draw[color ={black},>=latex,->] (-4,0)--(4,0);
	\draw[color ={black},>=latex,->] (0,-8)--(0,6);
	\draw[color ={black},opacity = 0.4] (-4,1)--(4,1);
	\draw[color ={black},opacity = 0.4] (-4,2)--(4,2);
	\draw[color ={black},opacity = 0.4] (-4,3)--(4,3);
	\draw[color ={black},opacity = 0.4] (-4,4)--(4,4);
	\draw[color ={black},opacity = 0.4] (-4,5)--(4,5);
	\draw[color ={black},opacity = 0.4] (-4,6)--(4,6);
	\draw[color ={black},opacity = 0.4] (-4,-1)--(4,-1);
	\draw[color ={black},opacity = 0.4] (-4,-2)--(4,-2);
	\draw[color ={black},opacity = 0.4] (-4,-3)--(4,-3);
	\draw[color ={black},opacity = 0.4] (-4,-4)--(4,-4);
	\draw[color ={black},opacity = 0.4] (-4,-5)--(4,-5);
	\draw[color ={black},opacity = 0.4] (-4,-6)--(4,-6);
	\draw[color ={black},opacity = 0.4] (-4,-7)--(4,-7);
	\draw[color ={black},opacity = 0.4] (-4,-8)--(4,-8);
	\draw[color ={black},opacity = 0.4] (1,-8)--(1,6);
	\draw[color ={black},opacity = 0.4] (2,-8)--(2,6);
	\draw[color ={black},opacity = 0.4] (3,-8)--(3,6);
	\draw[color ={black},opacity = 0.4] (4,-8)--(4,6);
	\draw[color ={black},opacity = 0.4] (-1,-8)--(-1,6);
	\draw[color ={black},opacity = 0.4] (-2,-8)--(-2,6);
	\draw[color ={black},opacity = 0.4] (-3,-8)--(-3,6);
	\draw[color ={black},opacity = 0.4] (-4,-8)--(-4,6);
	\draw[color ={gray},opacity = 0.3] (-4,0)--(4,0);
	\draw[color ={gray},opacity = 0.3] (-4,0.5)--(4,0.5);
	\draw[color ={gray},opacity = 0.3] (-4,1.5)--(4,1.5);
	\draw[color ={gray},opacity = 0.3] (-4,2.5)--(4,2.5);
	\draw[color ={gray},opacity = 0.3] (-4,3.5)--(4,3.5);
	\draw[color ={gray},opacity = 0.3] (-4,4.5)--(4,4.5);
	\draw[color ={gray},opacity = 0.3] (-4,5.5)--(4,5.5);
	\draw[color ={gray},opacity = 0.3] (-4,0)--(4,0);
	\draw[color ={gray},opacity = 0.3] (-4,-0.5)--(4,-0.5);
	\draw[color ={gray},opacity = 0.3] (-4,-1.5)--(4,-1.5);
	\draw[color ={gray},opacity = 0.3] (-4,-2.5)--(4,-2.5);
	\draw[color ={gray},opacity = 0.3] (-4,-3.5)--(4,-3.5);
	\draw[color ={gray},opacity = 0.3] (-4,-4.5)--(4,-4.5);
	\draw[color ={gray},opacity = 0.3] (-4,-5.5)--(4,-5.5);
	\draw[color ={gray},opacity = 0.3] (-4,-6.5)--(4,-6.5);
	\draw[color ={gray},opacity = 0.3] (-4,-7)--(4,-7);
	\draw[color ={gray},opacity = 0.3] (-4,-7.5)--(4,-7.5);
	\draw[color ={gray},opacity = 0.3] (-4,-8)--(4,-8);
	\draw[color ={gray},opacity = 0.3] (0,-8)--(0,6);
	\draw[color ={gray},opacity = 0.3] (0,-8)--(0,6);
	\draw[color ={black},line width = 1.2] (1,-0.13)--(1,0.13);
	\draw[color ={black},line width = 1.2] (2,-0.13)--(2,0.13);
	\draw[color ={black},line width = 1.2] (3,-0.13)--(3,0.13);
	\draw[color ={black},line width = 1.2] (4,-0.13)--(4,0.13);
	\draw[color ={black},line width = 1.2] (-1,-0.13)--(-1,0.13);
	\draw[color ={black},line width = 1.2] (-2,-0.13)--(-2,0.13);
	\draw[color ={black},line width = 1.2] (-3,-0.13)--(-3,0.13);
	\draw[color ={black},line width = 1.2] (-4,-0.13)--(-4,0.13);
	\draw[color ={black},line width = 1.2] (-0.13,1)--(0.13,1);
	\draw[color ={black},line width = 1.2] (-0.13,2)--(0.13,2);
	\draw[color ={black},line width = 1.2] (-0.13,3)--(0.13,3);
	\draw[color ={black},line width = 1.2] (-0.13,4)--(0.13,4);
	\draw[color ={black},line width = 1.2] (-0.13,5)--(0.13,5);
	\draw[color ={black},line width = 1.2] (-0.13,6)--(0.13,6);
	\draw[color ={black},line width = 1.2] (-0.13,-1)--(0.13,-1);
	\draw[color ={black},line width = 1.2] (-0.13,-2)--(0.13,-2);
	\draw[color ={black},line width = 1.2] (-0.13,-3)--(0.13,-3);
	\draw[color ={black},line width = 1.2] (-0.13,-4)--(0.13,-4);
	\draw[color ={black},line width = 1.2] (-0.13,-5)--(0.13,-5);
	\draw[color ={black},line width = 1.2] (-0.13,-6)--(0.13,-6);
	\draw[color ={black},line width = 1.2] (-0.13,-7)--(0.13,-7);
	\draw[color ={black},line width = 1.2] (-0.13,-8)--(0.13,-8);
	\draw (2,-0.4) node[anchor = center] {\scriptsize \color{black}{$2$}};
	\draw (4,-0.4) node[anchor = center] {\scriptsize \color{black}{$4$}};
	\draw (-2,-0.4) node[anchor = center] {\scriptsize \color{black}{$-2$}};
	\draw (-4,-0.4) node[anchor = center] {\scriptsize \color{black}{$-4$}};
	\draw (-0.5,2.1) node[anchor = center] {\footnotesize \color{black}{$2$}};
	\draw (-0.5,4.1) node[anchor = center] {\footnotesize \color{black}{$4$}};
	\draw (-0.5,-1.9) node[anchor = center] {\footnotesize \color{black}{$-2$}};
	\draw (-0.5,-3.9) node[anchor = center] {\footnotesize \color{black}{$-4$}};
	\draw (-0.5,-5.9) node[anchor = center] {\footnotesize \color{black}{$-6$}};
	\draw (-0.5,-7.9) node[anchor = center] {\footnotesize \color{black}{$-8$}};
	\draw [color={black}] (-0.3,-0.3) node[anchor = center,scale=1, rotate = 0] {O};
	\draw (3,4) node[anchor = center] {\footnotesize \color{red}{$\Large\mathcal{C}_f\prime$}};
	
	\draw[color={red},line width = 2] (-4,6)--(-3.8,5.6)--(-3.6,5.2)--(-3.4,4.8)--(-3.2,4.4)--(-3,4)--(-2.8,3.6)--(-2.6,3.2)--(-2.4,2.8)--(-2.2,2.4)--(-2,2)--(-1.8,1.6)--(-1.6,1.2)--(-1.4,0.8)--(-1.2,0.4)--(-1,0)--(-0.8,-0.4)--(-0.6,-0.8)--(-0.4,-1.2)--(-0.2,-1.6)--(0,-2)--(0.2,-2.4)--(0.4,-2.8)--(0.6,-3.2)--(0.8,-3.6)--(1,-4)--(1.2,-4.4)--(1.4,-4.8)--(1.6,-5.2)--(1.8,-5.6)--(2,-6)--(2.2,-6.4)--(2.4,-6.8)--(2.6,-7.2)--(2.8,-7.6)--(3,-8)--(3.2,-8.4)--(3.4,-8.8);

\end{tikzpicture}\\
    \end{minipage}
    
\end{exercice}

\begin{Solution}
 L'équation réduite de la tangente au point d'abscisse $0$ est  : $y=f'(0)(x-0)+f(0)$.\\
      On lit graphiquement $f(0)=0$ et $f'(0)=-2$.\\
      L'équation réduite de la tangente est donc donnée par :
      $y=-2(x+0)+0$, soit $y=-2x$.
\end{Solution}

% @see : https://coopmaths.fr/alea?uuid=a1ba2&id=can1F14&n=1&d=10&alea=rkIp223232323234234567891011121314151617181920212223242526272829303132333435363738394041424344454647482345678910&cd=1&cols=1
\needspace{10\baselineskip}
\begin{exercice}[BaremeTotal=false]


Soit $f$ la fonction définie sur $\mathbb{R}$ par : 
$f(x)= 7+2x^2+x$.\\
En calculant la limite du taux d'accroissement, déterminer $f'(2)$.
\end{exercice}

\begin{Solution}
 $f$ est une fonction polynôme du second degré de la forme $f(x)=ax^2+bx+c$.\\
    La fonction dérivée est donnée par la somme des dérivées des fonctions $u$ et $v$ définies par $u(x)=2x^2$ et $v(x)=x+7$.\\
     Comme $u'(x)=4x$ et $v'(x)=1$, on obtient  $f'(x)=4x+1$. \\
     Ainsi, $f'(2)=4\times 2+1=9$.
\end{Solution}

% @see : https://coopmaths.fr/alea?uuid=3c690&id=can1F13&n=1&d=10&alea=EUDu223232323234234567891011121314151617181920212223242526272829303132333435363738394041424344454647482345678910&cd=1&cols=1
\needspace{10\baselineskip}
\begin{exercice}[BaremeTotal=false]


 Déterminer le coefficient directeur de la tangente à la courbe représentative de la fonction racine carrée au point d'abscisse $18$.
         
\end{exercice}

\begin{Solution}
 Le coefficient directeur de la tangente au point d'abscisse $a$ est donné par le nombre dérivé $f'(a)$.\\
        La fonction $f$ définie par $f(x)=\sqrt{x}$ a pour fonction dérivée la fonction $f'$ définie par $f'(x)=\dfrac{1}{2\sqrt{x}}$.\\Comme $f'(18)=\dfrac{1}{2\sqrt{18}}$, le coefficient directeur de la tangente au point d'abscisse $18$ est : $\dfrac{1}{2\sqrt{18}}$.
\end{Solution}

% @see : https://coopmaths.fr/alea?uuid=ebd89&id=1AN14-1&n=1&d=10&s=3&alea=ocYr223232323234234567891011121314151617181920212223242526272829303132333435363738394041424344454647482345678910&cd=0&cols=1
\needspace{10\baselineskip}
\begin{exercice}[BaremeTotal=false]


 Donner la dérivée de la fonction $f$, dérivable pour tout $x\in\R$, définie par  $f(x)=\dfrac{1}{2}x+4$.
\end{exercice}

\begin{Solution}
 L'expression de la dérivée de la fonction $f$ définie par $f(x)=\dfrac{1}{2}x+4$ est : ${\color[HTML]{f15929}\boldsymbol{f'(x)=\dfrac{1}{2}}}$.
\end{Solution}

% @see : https://coopmaths.fr/alea?uuid=ebd89&id=1AN14-1&n=1&d=10&s=4&alea=uAlP223232323234234567891011121314151617181920212223242526272829303132333435363738394041424344454647482345678910&cd=0&cols=1
\needspace{10\baselineskip}
\begin{exercice}[BaremeTotal=false]


 Donner la dérivée de la fonction $f$, dérivable pour tout $x\in\R$, définie par  $f(x)=-5x^{14}$.
\end{exercice}

\begin{Solution}
 L'expression de la dérivée de la fonction $f$ définie par $f(x)=-5x^{14}$ est : ${\color[HTML]{f15929}\boldsymbol{f'(x)=-70x^{13}}}$.
\end{Solution}

% @see : https://coopmaths.fr/alea?uuid=ebd89&id=1AN14-1&n=1&d=10&s=7&alea=vjN7223232323234234567891011121314151617181920212223242526272829303132333435363738394041424344454647482345678910&cd=0&cols=1
\needspace{10\baselineskip}
\begin{exercice}[BaremeTotal=false]


 Donner la dérivée de la fonction $f$, dérivable pour tout $x\in\R^*$, définie par  $f(x)=\dfrac{-1}{10x}$.
\end{exercice}

\begin{Solution}
 L'expression de la dérivée de la fonction $f$ définie par $f(x)=\dfrac{-1}{10x}$ est : ${\color[HTML]{f15929}\boldsymbol{f'(x)=\dfrac{1}{10x^2}}}$.
\end{Solution}

% @see : https://coopmaths.fr/alea?uuid=a83c0&id=1AN14-4&n=1&d=10&s=1&alea=Njr3223232323234234567891011121314151617181920212223242526272829303132333435363738394041424344454647482345678910&cd=0&cols=1
\needspace{10\baselineskip}
\begin{exercice}[BaremeTotal=false]


 Donner l'expression de la dérivée de la fonction $f$ définie sur $\R^{*}$ par $f(x)=-4x^{2}+3x-2+\dfrac{9}{x}$.\\
\end{exercice}

\begin{Solution}
 L'expression de la dérivée de la fonction $f$ définie par $f(x)=-4x^{2}+3x-2+\dfrac{9}{x}$ est :\\$f'(x)={\color[HTML]{f15929}\boldsymbol{-8x+3-\dfrac{9}{x^2}}}$.\\
\end{Solution}

% @see : https://coopmaths.fr/alea?uuid=a83c0&id=1AN14-4&n=1&d=10&s=4&alea=4Vpz223232323234234567891011121314151617181920212223242526272829303132333435363738394041424344454647482345678910&cd=1&cols=1
\needspace{10\baselineskip}
\begin{exercice}[BaremeTotal=false]


 Donner l'expression de la dérivée de la fonction $f$ définie sur $\R_+^{*}$ par $f(x)=\dfrac{4}{x}+5\sqrt{x}+3x$.\\
\end{exercice}

\begin{Solution}
 $(\dfrac{4}{x})^\prime=-\dfrac{4}{x^2}$\\$(5\sqrt{x})^\prime=\dfrac{5}{2\sqrt{x}}$\\$(3x)^\prime=3$\\L'expression de la dérivée de la fonction $f$ définie par $f(x)=\dfrac{4}{x}+5\sqrt{x}+3x$ est :\\$f'(x)={\color[HTML]{f15929}\boldsymbol{-\dfrac{4}{x^2}+\dfrac{5}{2\sqrt{x}}+3}}$.\\
\end{Solution}

\end{Maquette}
\clearpage
\begin{Maquette}[DM]{Niveau= ,Classe=\getdata[11]\mydata\ (v11), Date = 06/01/25  ,Theme=Exercices}


% @see : https://coopmaths.fr/alea?uuid=4c8c7&id=1AN11-3&n=1&d=10&s=2&alea=nawO22323232323423456789101112131415161718192021222324252627282930313233343536373839404142434445464748234567891011&cd=1&cols=2
\needspace{10\baselineskip}
\begin{exercice}[BaremeTotal=false]
\label{eleve:11}


 \begin{minipage}[t]{\linewidth} Soit $f$ une fonction dérivable sur $[-5;5]$ et $\mathcal{C}_f$ sa courbe représentative.\\On sait que  $f(3)=5~~$ et que $~~f'(3)=-5$.\\Déterminer une équation de la tangente $(T)$ à la courbe $\mathcal{C}_f$ au point d'abscisse $3$,\\sans utiliser la formule de cours de l'équation de tangente. \end{minipage}

\end{exercice}

\begin{Solution}
 \begin{minipage}[t]{\linewidth} On sait que la tangente n'est pas une droite verticale, puisque la fonction est dérivable sur l'intervalle.\\On en déduit que la tangente $(T)$ au point d'abscisse $3$, admet une équation réduite de la forme :  $(T) : y=m x + p$. 

\medskip
$\bullet$ Détermination de $m$ :\\On sait que le nombre dérivé en $3$ est par définition, le coefficient directeur de la tangente au point d'abscisse $3$.\\Par conséquent, on a déjà : $m=f'(3)=-5$.\\ On en déduit que  $(T) : y= -5 x + p$\\$\bullet$ Détermination de $p$ :\\ Pour cela, on utilise que si $f(3)=5~~$, alors le point $A$ de coordonnées $(3;5)$ appartient à $\mathcal{C}_f$ mais aussi à $(T)$.\\ On peut écrire $A(3;5) = \mathcal{C}_f \cap (T)$.\\ On remplace alors les coordonnées de $A(3;5)$ dans l'équation  $(T) : y= -5 x + p$.\\ $\begin{aligned} \phantom{\iff}&A(3;5)\in (T)\\ \iff& 5= -5 \times 3  + p\\ \iff& p=5 +5 \times 3  \\ \iff& p=5 +15   \\ \iff& p=20\\\end{aligned}$\\On peut conclure que : $(T) : {\color[HTML]{f15929}\boldsymbol{y=-5x+20}}$. \end{minipage}

\end{Solution}

% @see : https://coopmaths.fr/alea?uuid=0e984&id=can1F15&n=1&d=10&alea=1Alo22323232323423456789101112131415161718192021222324252627282930313233343536373839404142434445464748234567891011&cd=1&cols=1
\needspace{10\baselineskip}
\begin{exercice}[BaremeTotal=false]


 La courbe représente une fonction $f$ et la droite est la tangente au point d'abscisse $0$.\\
        Déterminer $f'(0)$.\begin{center}\begin{tikzpicture}[baseline,scale = 0.5]

    \tikzset{
      point/.style={
        thick,
        draw,
        cross out,
        inner sep=0pt,
        minimum width=5pt,
        minimum height=5pt,
      },
    }
    \clip (-4,-2) rectangle (12,12);
    	
	\draw[color ={black},line width = 2,>=latex,->] (-4,0)--(12,0);
	\draw[color ={black},line width = 2,>=latex,->] (0,-2)--(0,12);
	\draw[color ={black},opacity = 0.5] (-4,1)--(12,1);
	\draw[color ={black},opacity = 0.5] (-4,2)--(12,2);
	\draw[color ={black},opacity = 0.5] (-4,3)--(12,3);
	\draw[color ={black},opacity = 0.5] (-4,4)--(12,4);
	\draw[color ={black},opacity = 0.5] (-4,5)--(12,5);
	\draw[color ={black},opacity = 0.5] (-4,6)--(12,6);
	\draw[color ={black},opacity = 0.5] (-4,7)--(12,7);
	\draw[color ={black},opacity = 0.5] (-4,8)--(12,8);
	\draw[color ={black},opacity = 0.5] (-4,9)--(12,9);
	\draw[color ={black},opacity = 0.5] (-4,10)--(12,10);
	\draw[color ={black},opacity = 0.5] (-4,11)--(12,11);
	\draw[color ={black},opacity = 0.5] (-4,12)--(12,12);
	\draw[color ={black},opacity = 0.5] (-4,-1)--(12,-1);
	\draw[color ={black},opacity = 0.5] (-4,-2)--(12,-2);
	\draw[color ={black},opacity = 0.5] (1,-2)--(1,12);
	\draw[color ={black},opacity = 0.5] (2,-2)--(2,12);
	\draw[color ={black},opacity = 0.5] (3,-2)--(3,12);
	\draw[color ={black},opacity = 0.5] (4,-2)--(4,12);
	\draw[color ={black},opacity = 0.5] (5,-2)--(5,12);
	\draw[color ={black},opacity = 0.5] (6,-2)--(6,12);
	\draw[color ={black},opacity = 0.5] (7,-2)--(7,12);
	\draw[color ={black},opacity = 0.5] (8,-2)--(8,12);
	\draw[color ={black},opacity = 0.5] (9,-2)--(9,12);
	\draw[color ={black},opacity = 0.5] (10,-2)--(10,12);
	\draw[color ={black},opacity = 0.5] (11,-2)--(11,12);
	\draw[color ={black},opacity = 0.5] (12,-2)--(12,12);
	\draw[color ={black},opacity = 0.5] (-1,-2)--(-1,12);
	\draw[color ={black},opacity = 0.5] (-2,-2)--(-2,12);
	\draw[color ={black},opacity = 0.5] (-3,-2)--(-3,12);
	\draw[color ={black},opacity = 0.5] (-4,-2)--(-4,12);
	\draw[color ={gray},opacity = 0.3] (-4,0)--(12,0);
	\draw[color ={gray},opacity = 0.3] (-4,0)--(12,0);
	\draw[color ={gray},opacity = 0.3] (0,-2)--(0,12);
	\draw[color ={gray},opacity = 0.3] (0,-2)--(0,12);
	\draw[color ={black},line width = 2] (2,-0.2)--(2,0.2);
	\draw[color ={black},line width = 2] (4,-0.2)--(4,0.2);
	\draw[color ={black},line width = 2] (6,-0.2)--(6,0.2);
	\draw[color ={black},line width = 2] (8,-0.2)--(8,0.2);
	\draw[color ={black},line width = 2] (10,-0.2)--(10,0.2);
	\draw[color ={black},line width = 2] (-2,-0.2)--(-2,0.2);
	\draw[color ={black},line width = 2] (-0.2,2)--(0.2,2);
	\draw[color ={black},line width = 2] (-0.2,4)--(0.2,4);
	\draw[color ={black},line width = 2] (-0.2,6)--(0.2,6);
	\draw[color ={black},line width = 2] (-0.2,8)--(0.2,8);
	\draw[color ={black},line width = 2] (-0.2,10)--(0.2,10);
	\draw (2,-0.4) node[anchor = center] {\scriptsize \color{black}{$1$}};
	\draw (4,-0.4) node[anchor = center] {\scriptsize \color{black}{$2$}};
	\draw (6,-0.4) node[anchor = center] {\scriptsize \color{black}{$3$}};
	\draw (8,-0.4) node[anchor = center] {\scriptsize \color{black}{$4$}};
	\draw (10,-0.4) node[anchor = center] {\scriptsize \color{black}{$5$}};
	\draw (-0.5,2.1) node[anchor = center] {\footnotesize \color{black}{$1$}};
	\draw (-0.5,4.1) node[anchor = center] {\footnotesize \color{black}{$2$}};
	\draw (-0.5,6.1) node[anchor = center] {\footnotesize \color{black}{$3$}};
	\draw (-0.5,8.1) node[anchor = center] {\footnotesize \color{black}{$4$}};
	\draw (-0.5,10.1) node[anchor = center] {\footnotesize \color{black}{$5$}};
	\draw [color={black}] (-0.3,-0.3) node[anchor = center,scale=1, rotate = 0] {O};
	
	\draw[color={blue},line width = 2] (-4,1.21)--(-3.6,1.28)--(-3.2,1.34)--(-2.8,1.41)--(-2.4,1.48)--(-2,1.56)--(-1.6,1.64)--(-1.2,1.72)--(-0.8,1.81)--(-0.4,1.9)--(0,2)--(0.4,2.1)--(0.8,2.21)--(1.2,2.32)--(1.6,2.44)--(2,2.57)--(2.4,2.7)--(2.8,2.84)--(3.2,2.98)--(3.6,3.14)--(4,3.3)--(4.4,3.47)--(4.8,3.64)--(5.2,3.83)--(5.6,4.03)--(6,4.23)--(6.4,4.45)--(6.8,4.68)--(7.2,4.92)--(7.6,5.17)--(8,5.44)--(8.4,5.72)--(8.8,6.01)--(9.2,6.32)--(9.6,6.64)--(10,6.98);
	
	\draw[color={red},line width = 2] (-4,1)--(-3.6,1.1)--(-3.2,1.2)--(-2.8,1.3)--(-2.4,1.4)--(-2,1.5)--(-1.6,1.6)--(-1.2,1.7)--(-0.8,1.8)--(-0.4,1.9)--(0,2)--(0.4,2.1)--(0.8,2.2)--(1.2,2.3)--(1.6,2.4)--(2,2.5)--(2.4,2.6)--(2.8,2.7)--(3.2,2.8)--(3.6,2.9)--(4,3)--(4.4,3.1)--(4.8,3.2)--(5.2,3.3)--(5.6,3.4)--(6,3.5)--(6.4,3.6)--(6.8,3.7)--(7.2,3.8)--(7.6,3.9)--(8,4)--(8.4,4.1)--(8.8,4.2)--(9.2,4.3)--(9.6,4.4)--(10,4.5)--(10.4,4.6)--(10.8,4.7)--(11.2,4.8)--(11.6,4.9)--(12,5);

\end{tikzpicture}\end{center}
\end{exercice}

\begin{Solution}
 $f'(0)$ est donné par le coefficient directeur de la tangente à la courbe au point d'abscisse $0$, soit $\dfrac{1}{4}$.
\end{Solution}

% @see : https://coopmaths.fr/alea?uuid=6f32d&id=can1F16&n=1&d=10&alea=dqSw22323232323423456789101112131415161718192021222324252627282930313233343536373839404142434445464748234567891011&cd=1&cols=1
\needspace{10\baselineskip}
\newpage
\begin{exercice}[BaremeTotal=false]


 On donne les représentations graphiques d'une fonction et de sa dérivée.\\
        Donner l'équation réduite de la tangente à la courbe de $f$ en $x=-2$. \\ \begin{minipage}{0.5\linewidth}
    \begin{tikzpicture}[baseline, scale=0.8]

    \tikzset{
      point/.style={
        thick,
        draw,
        cross out,
        inner sep=0pt,
        minimum width=5pt,
        minimum height=5pt,
      },
    }
    \clip (-4.5,-4.5) rectangle (5.75,13.5);
    	
	\draw[color ={black},>=latex,->] (-3,0)--(3,0);
	\draw[color ={black},>=latex,->] (0,-3)--(0,12);
	\draw[color ={black},opacity = 0.4] (-3,1)--(3,1);
	\draw[color ={black},opacity = 0.4] (-3,2)--(3,2);
	\draw[color ={black},opacity = 0.4] (-3,3)--(3,3);
	\draw[color ={black},opacity = 0.4] (-3,4)--(3,4);
	\draw[color ={black},opacity = 0.4] (-3,5)--(3,5);
	\draw[color ={black},opacity = 0.4] (-3,6)--(3,6);
	\draw[color ={black},opacity = 0.4] (-3,7)--(3,7);
	\draw[color ={black},opacity = 0.4] (-3,8)--(3,8);
	\draw[color ={black},opacity = 0.4] (-3,9)--(3,9);
	\draw[color ={black},opacity = 0.4] (-3,10)--(3,10);
	\draw[color ={black},opacity = 0.4] (-3,11)--(3,11);
	\draw[color ={black},opacity = 0.4] (-3,12)--(3,12);
	\draw[color ={black},opacity = 0.4] (-3,-1)--(3,-1);
	\draw[color ={black},opacity = 0.4] (-3,-2)--(3,-2);
	\draw[color ={black},opacity = 0.4] (-3,-3)--(3,-3);
	\draw[color ={black},opacity = 0.4] (1,-3)--(1,12);
	\draw[color ={black},opacity = 0.4] (2,-3)--(2,12);
	\draw[color ={black},opacity = 0.4] (3,-3)--(3,12);
	\draw[color ={black},opacity = 0.4] (-1,-3)--(-1,12);
	\draw[color ={black},opacity = 0.4] (-2,-3)--(-2,12);
	\draw[color ={black},opacity = 0.4] (-3,-3)--(-3,12);
	\draw[color ={gray},opacity = 0.3] (-3,0)--(3,0);
	\draw[color ={gray},opacity = 0.3] (-3,0.5)--(3,0.5);
	\draw[color ={gray},opacity = 0.3] (-3,1.5)--(3,1.5);
	\draw[color ={gray},opacity = 0.3] (-3,2.5)--(3,2.5);
	\draw[color ={gray},opacity = 0.3] (-3,3.5)--(3,3.5);
	\draw[color ={gray},opacity = 0.3] (-3,4.5)--(3,4.5);
	\draw[color ={gray},opacity = 0.3] (-3,5.5)--(3,5.5);
	\draw[color ={gray},opacity = 0.3] (-3,6.5)--(3,6.5);
	\draw[color ={gray},opacity = 0.3] (-3,7.5)--(3,7.5);
	\draw[color ={gray},opacity = 0.3] (-3,8.5)--(3,8.5);
	\draw[color ={gray},opacity = 0.3] (-3,9.5)--(3,9.5);
	\draw[color ={gray},opacity = 0.3] (-3,10.5)--(3,10.5);
	\draw[color ={gray},opacity = 0.3] (-3,11.5)--(3,11.5);
	\draw[color ={gray},opacity = 0.3] (-3,0)--(3,0);
	\draw[color ={gray},opacity = 0.3] (-3,-0.5)--(3,-0.5);
	\draw[color ={gray},opacity = 0.3] (-3,-1.5)--(3,-1.5);
	\draw[color ={gray},opacity = 0.3] (-3,-2.5)--(3,-2.5);
	\draw[color ={gray},opacity = 0.3] (0,-3)--(0,12);
	\draw[color ={gray},opacity = 0.3] (0,-3)--(0,12);
	\draw[color ={black},line width = 1.2] (1,-0.13)--(1,0.13);
	\draw[color ={black},line width = 1.2] (2,-0.13)--(2,0.13);
	\draw[color ={black},line width = 1.2] (3,-0.13)--(3,0.13);
	\draw[color ={black},line width = 1.2] (-1,-0.13)--(-1,0.13);
	\draw[color ={black},line width = 1.2] (-2,-0.13)--(-2,0.13);
	\draw[color ={black},line width = 1.2] (-3,-0.13)--(-3,0.13);
	\draw[color ={black},line width = 1.2] (-0.13,1)--(0.13,1);
	\draw[color ={black},line width = 1.2] (-0.13,2)--(0.13,2);
	\draw[color ={black},line width = 1.2] (-0.13,3)--(0.13,3);
	\draw[color ={black},line width = 1.2] (-0.13,4)--(0.13,4);
	\draw[color ={black},line width = 1.2] (-0.13,5)--(0.13,5);
	\draw[color ={black},line width = 1.2] (-0.13,6)--(0.13,6);
	\draw[color ={black},line width = 1.2] (-0.13,7)--(0.13,7);
	\draw[color ={black},line width = 1.2] (-0.13,8)--(0.13,8);
	\draw[color ={black},line width = 1.2] (-0.13,9)--(0.13,9);
	\draw[color ={black},line width = 1.2] (-0.13,10)--(0.13,10);
	\draw[color ={black},line width = 1.2] (-0.13,11)--(0.13,11);
	\draw[color ={black},line width = 1.2] (-0.13,12)--(0.13,12);
	\draw[color ={black},line width = 1.2] (-0.13,-1)--(0.13,-1);
	\draw[color ={black},line width = 1.2] (-0.13,-2)--(0.13,-2);
	\draw[color ={black},line width = 1.2] (-0.13,-3)--(0.13,-3);
	\draw (3,-0.4) node[anchor = center] {\scriptsize \color{black}{$3$}};
	\draw (-3,-0.4) node[anchor = center] {\scriptsize \color{black}{$-3$}};
	\draw (-0.5,3.1) node[anchor = center] {\footnotesize \color{black}{$3$}};
	\draw (-0.5,6.1) node[anchor = center] {\footnotesize \color{black}{$6$}};
	\draw (-0.5,9.1) node[anchor = center] {\footnotesize \color{black}{$9$}};
	\draw (-0.5,-2.9) node[anchor = center] {\footnotesize \color{black}{$-3$}};
	\draw [color={black}] (-0.3,-0.3) node[anchor = center,scale=1, rotate = 0] {O};
	\draw (3,10) node[anchor = center] {\footnotesize \color{blue}{$\Large \mathcal{C}_f$}};
	
	\draw[color={blue},line width = 2] (-3,4)--(-2.8,3.24)--(-2.6,2.56)--(-2.4,1.96)--(-2.2,1.44)--(-2,1)--(-1.8,0.64)--(-1.6,0.36)--(-1.4,0.16)--(-1.2,0.04)--(-1,0)--(-0.8,0.04)--(-0.6,0.16)--(-0.4,0.36)--(-0.2,0.64)--(0,1)--(0.2,1.44)--(0.4,1.96)--(0.6,2.56)--(0.8,3.24)--(1,4)--(1.2,4.84)--(1.4,5.76)--(1.6,6.76)--(1.8,7.84)--(2,9)--(2.2,10.24)--(2.4,11.56)--(2.6,12.96);

\end{tikzpicture}\\
    \end{minipage}
    \begin{minipage}{0.5\linewidth}
    \begin{tikzpicture}[baseline, scale=0.8]

    \tikzset{
      point/.style={
        thick,
        draw,
        cross out,
        inner sep=0pt,
        minimum width=5pt,
        minimum height=5pt,
      },
    }
    \clip (-4.5,-6.5) rectangle (6.5,9.5);
    	
	\draw[color ={black},>=latex,->] (-3,0)--(3,0);
	\draw[color ={black},>=latex,->] (0,-5)--(0,8);
	\draw[color ={black},opacity = 0.4] (-3,1)--(3,1);
	\draw[color ={black},opacity = 0.4] (-3,2)--(3,2);
	\draw[color ={black},opacity = 0.4] (-3,3)--(3,3);
	\draw[color ={black},opacity = 0.4] (-3,4)--(3,4);
	\draw[color ={black},opacity = 0.4] (-3,5)--(3,5);
	\draw[color ={black},opacity = 0.4] (-3,6)--(3,6);
	\draw[color ={black},opacity = 0.4] (-3,7)--(3,7);
	\draw[color ={black},opacity = 0.4] (-3,8)--(3,8);
	\draw[color ={black},opacity = 0.4] (-3,-1)--(3,-1);
	\draw[color ={black},opacity = 0.4] (-3,-2)--(3,-2);
	\draw[color ={black},opacity = 0.4] (-3,-3)--(3,-3);
	\draw[color ={black},opacity = 0.4] (-3,-4)--(3,-4);
	\draw[color ={black},opacity = 0.4] (-3,-5)--(3,-5);
	\draw[color ={black},opacity = 0.4] (1,-5)--(1,8);
	\draw[color ={black},opacity = 0.4] (2,-5)--(2,8);
	\draw[color ={black},opacity = 0.4] (3,-5)--(3,8);
	\draw[color ={black},opacity = 0.4] (-1,-5)--(-1,8);
	\draw[color ={black},opacity = 0.4] (-2,-5)--(-2,8);
	\draw[color ={black},opacity = 0.4] (-3,-5)--(-3,8);
	\draw[color ={gray},opacity = 0.3] (-3,0)--(3,0);
	\draw[color ={gray},opacity = 0.3] (-3,0.5)--(3,0.5);
	\draw[color ={gray},opacity = 0.3] (-3,1.5)--(3,1.5);
	\draw[color ={gray},opacity = 0.3] (-3,2.5)--(3,2.5);
	\draw[color ={gray},opacity = 0.3] (-3,3.5)--(3,3.5);
	\draw[color ={gray},opacity = 0.3] (-3,4.5)--(3,4.5);
	\draw[color ={gray},opacity = 0.3] (-3,5.5)--(3,5.5);
	\draw[color ={gray},opacity = 0.3] (-3,6.5)--(3,6.5);
	\draw[color ={gray},opacity = 0.3] (-3,7.5)--(3,7.5);
	\draw[color ={gray},opacity = 0.3] (-3,0)--(3,0);
	\draw[color ={gray},opacity = 0.3] (-3,-0.5)--(3,-0.5);
	\draw[color ={gray},opacity = 0.3] (-3,-1.5)--(3,-1.5);
	\draw[color ={gray},opacity = 0.3] (-3,-2.5)--(3,-2.5);
	\draw[color ={gray},opacity = 0.3] (-3,-3.5)--(3,-3.5);
	\draw[color ={gray},opacity = 0.3] (-3,-4.5)--(3,-4.5);
	\draw[color ={gray},opacity = 0.3] (0,-5)--(0,8);
	\draw[color ={gray},opacity = 0.3] (0,-5)--(0,8);
	\draw[color ={black},line width = 1.2] (1,-0.13)--(1,0.13);
	\draw[color ={black},line width = 1.2] (2,-0.13)--(2,0.13);
	\draw[color ={black},line width = 1.2] (3,-0.13)--(3,0.13);
	\draw[color ={black},line width = 1.2] (-1,-0.13)--(-1,0.13);
	\draw[color ={black},line width = 1.2] (-2,-0.13)--(-2,0.13);
	\draw[color ={black},line width = 1.2] (-3,-0.13)--(-3,0.13);
	\draw[color ={black},line width = 1.2] (-0.13,1)--(0.13,1);
	\draw[color ={black},line width = 1.2] (-0.13,2)--(0.13,2);
	\draw[color ={black},line width = 1.2] (-0.13,3)--(0.13,3);
	\draw[color ={black},line width = 1.2] (-0.13,4)--(0.13,4);
	\draw[color ={black},line width = 1.2] (-0.13,5)--(0.13,5);
	\draw[color ={black},line width = 1.2] (-0.13,6)--(0.13,6);
	\draw[color ={black},line width = 1.2] (-0.13,7)--(0.13,7);
	\draw[color ={black},line width = 1.2] (-0.13,8)--(0.13,8);
	\draw[color ={black},line width = 1.2] (-0.13,9)--(0.13,9);
	\draw[color ={black},line width = 1.2] (-0.13,10)--(0.13,10);
	\draw[color ={black},line width = 1.2] (-0.13,11)--(0.13,11);
	\draw[color ={black},line width = 1.2] (-0.13,12)--(0.13,12);
	\draw[color ={black},line width = 1.2] (-0.13,-1)--(0.13,-1);
	\draw[color ={black},line width = 1.2] (-0.13,-2)--(0.13,-2);
	\draw[color ={black},line width = 1.2] (-0.13,-3)--(0.13,-3);
	\draw (3,-0.4) node[anchor = center] {\scriptsize \color{black}{$3$}};
	\draw (-3,-0.4) node[anchor = center] {\scriptsize \color{black}{$-3$}};
	\draw (-0.5,3.1) node[anchor = center] {\footnotesize \color{black}{$3$}};
	\draw (-0.5,6.1) node[anchor = center] {\footnotesize \color{black}{$6$}};
	\draw (-0.5,-2.9) node[anchor = center] {\footnotesize \color{black}{$-3$}};
	\draw [color={black}] (-0.3,-0.3) node[anchor = center,scale=1, rotate = 0] {O};
	\draw (3,6) node[anchor = center] {\footnotesize \color{red}{$\Large\mathcal{C}_f\prime$}};
	
	\draw[color={red},line width = 2] (-3,-4)--(-2.8,-3.6)--(-2.6,-3.2)--(-2.4,-2.8)--(-2.2,-2.4)--(-2,-2)--(-1.8,-1.6)--(-1.6,-1.2)--(-1.4,-0.8)--(-1.2,-0.4)--(-1,0)--(-0.8,0.4)--(-0.6,0.8)--(-0.4,1.2)--(-0.2,1.6)--(0,2)--(0.2,2.4)--(0.4,2.8)--(0.6,3.2)--(0.8,3.6)--(1,4)--(1.2,4.4)--(1.4,4.8)--(1.6,5.2)--(1.8,5.6)--(2,6)--(2.2,6.4)--(2.4,6.8)--(2.6,7.2)--(2.8,7.6)--(3,8);

\end{tikzpicture}\\
    \end{minipage}
    
\end{exercice}

\begin{Solution}
 L'équation réduite de la tangente au point d'abscisse $-2$ est  : $y=f'(-2)(x-(-2))+f(-2)$.\\
        On lit graphiquement $f(-2)=1$ et $f'(-2)=-2$.\\
        L'équation réduite de la tangente est donc donnée par :
        $y=-2(x+2)+1$, soit $y=-2x-3$.
\end{Solution}

% @see : https://coopmaths.fr/alea?uuid=a1ba2&id=can1F14&n=1&d=10&alea=rkIp22323232323423456789101112131415161718192021222324252627282930313233343536373839404142434445464748234567891011&cd=1&cols=1
\needspace{10\baselineskip}
\begin{exercice}[BaremeTotal=false]


Soit $f$ la fonction définie sur $\mathbb{R}$ par : 
$f(x)= -8-4x^2-5x$.\\
En calculant la limite du taux d'accroissement, déterminer $f'(-3)$.
\end{exercice}

\begin{Solution}
 $f$ est une fonction polynôme du second degré de la forme $f(x)=ax^2+bx+c$.\\
    La fonction dérivée est donnée par la somme des dérivées des fonctions $u$ et $v$ définies par $u(x)=-4x^2$ et $v(x)=-5x-8$.\\
     Comme $u'(x)=-8x$ et $v'(x)=-5$, on obtient  $f'(x)=-8x-5$. \\
     Ainsi, $f'(-3)=-8\times (-3)-5=19$.
\end{Solution}

% @see : https://coopmaths.fr/alea?uuid=3c690&id=can1F13&n=1&d=10&alea=EUDu22323232323423456789101112131415161718192021222324252627282930313233343536373839404142434445464748234567891011&cd=1&cols=1
\needspace{10\baselineskip}
\begin{exercice}[BaremeTotal=false]


 Déterminer le coefficient directeur de la tangente à la courbe représentative de la fonction carré au point d'abscisse $11$.    
\end{exercice}

\begin{Solution}
 Le coefficient directeur de la tangente au point d'abscisse $a$ est donné par le nombre dérivé $f'(a)$.\\
        La fonction $f$ définie par $f(x)=x^2$ a pour fonction dérivée la fonction $f'$ définie par $f'(x)=2x$.\\
        Comme $f'(11)=2\times 11=22$, le coefficient directeur de la tangente au point d'abscisse $11$ est : $22$. 
\end{Solution}

% @see : https://coopmaths.fr/alea?uuid=ebd89&id=1AN14-1&n=1&d=10&s=3&alea=ocYr22323232323423456789101112131415161718192021222324252627282930313233343536373839404142434445464748234567891011&cd=0&cols=1
\needspace{10\baselineskip}
\begin{exercice}[BaremeTotal=false]


 Donner la dérivée de la fonction $f$, dérivable pour tout $x\in\R$, définie par  $f(x)=-\dfrac{1}{9}x+2$.
\end{exercice}

\begin{Solution}
 L'expression de la dérivée de la fonction $f$ définie par $f(x)=-\dfrac{1}{9}x+2$ est : ${\color[HTML]{f15929}\boldsymbol{f'(x)=-\dfrac{1}{9}}}$.
\end{Solution}

% @see : https://coopmaths.fr/alea?uuid=ebd89&id=1AN14-1&n=1&d=10&s=4&alea=uAlP22323232323423456789101112131415161718192021222324252627282930313233343536373839404142434445464748234567891011&cd=0&cols=1
\needspace{10\baselineskip}
\begin{exercice}[BaremeTotal=false]


 Donner la dérivée de la fonction $f$, dérivable pour tout $x\in\R$, définie par  $f(x)=-7x^{5}$.
\end{exercice}

\begin{Solution}
 L'expression de la dérivée de la fonction $f$ définie par $f(x)=-7x^{5}$ est : ${\color[HTML]{f15929}\boldsymbol{f'(x)=-35x^{4}}}$.
\end{Solution}

% @see : https://coopmaths.fr/alea?uuid=ebd89&id=1AN14-1&n=1&d=10&s=7&alea=vjN722323232323423456789101112131415161718192021222324252627282930313233343536373839404142434445464748234567891011&cd=0&cols=1
\needspace{10\baselineskip}
\begin{exercice}[BaremeTotal=false]


 Donner la dérivée de la fonction $f$, dérivable pour tout $x\in\R^*$, définie par  $f(x)=\dfrac{-1}{7x}$.
\end{exercice}

\begin{Solution}
 L'expression de la dérivée de la fonction $f$ définie par $f(x)=\dfrac{-1}{7x}$ est : ${\color[HTML]{f15929}\boldsymbol{f'(x)=\dfrac{1}{7x^2}}}$.
\end{Solution}

% @see : https://coopmaths.fr/alea?uuid=a83c0&id=1AN14-4&n=1&d=10&s=1&alea=Njr322323232323423456789101112131415161718192021222324252627282930313233343536373839404142434445464748234567891011&cd=0&cols=1
\needspace{10\baselineskip}
\begin{exercice}[BaremeTotal=false]


 Donner l'expression de la dérivée de la fonction $f$ définie sur $\R^{*}$ par $f(x)=-\dfrac{9}{x}-5x$.\\
\end{exercice}

\begin{Solution}
 L'expression de la dérivée de la fonction $f$ définie par $f(x)=-\dfrac{9}{x}-5x$ est :\\$f'(x)={\color[HTML]{f15929}\boldsymbol{\dfrac{9}{x^2}-5}}$.\\
\end{Solution}

% @see : https://coopmaths.fr/alea?uuid=a83c0&id=1AN14-4&n=1&d=10&s=4&alea=4Vpz22323232323423456789101112131415161718192021222324252627282930313233343536373839404142434445464748234567891011&cd=1&cols=1
\needspace{10\baselineskip}
\begin{exercice}[BaremeTotal=false]


 Donner l'expression de la dérivée de la fonction $f$ définie sur $\R_+^{*}$ par $f(x)=4\sqrt{x}+\dfrac{3}{x}+4x^{2}+2x$.\\
\end{exercice}

\begin{Solution}
 $(4\sqrt{x})^\prime=\dfrac{4}{2\sqrt{x}}$\\$(\dfrac{3}{x})^\prime=-\dfrac{3}{x^2}$\\$(4x^{2}+2x)^\prime=8x+2$\\L'expression de la dérivée de la fonction $f$ définie par $f(x)=4\sqrt{x}+\dfrac{3}{x}+4x^{2}+2x$ est :\\$f'(x)={\color[HTML]{f15929}\boldsymbol{\dfrac{4}{2\sqrt{x}}-\dfrac{3}{x^2}+8x+2}}$.\\
\end{Solution}

\end{Maquette}
\clearpage
\begin{Maquette}[DM]{Niveau= ,Classe=\getdata[12]\mydata\ (v12), Date = 06/01/25  ,Theme=Exercices}


% @see : https://coopmaths.fr/alea?uuid=4c8c7&id=1AN11-3&n=1&d=10&s=2&alea=nawO2232323232342345678910111213141516171819202122232425262728293031323334353637383940414243444546474823456789101112&cd=1&cols=2
\needspace{10\baselineskip}
\begin{exercice}[BaremeTotal=false]
\label{eleve:12}


 \begin{minipage}[t]{\linewidth} Soit $f$ une fonction dérivable sur $[-5;5]$ et $\mathcal{C}_f$ sa courbe représentative.\\On sait que  $f(4)=-5~~$ et que $~~f'(4)=4$.\\Déterminer une équation de la tangente $(T)$ à la courbe $\mathcal{C}_f$ au point d'abscisse $4$,\\sans utiliser la formule de cours de l'équation de tangente. \end{minipage}

\end{exercice}

\begin{Solution}
 \begin{minipage}[t]{\linewidth} On sait que la tangente n'est pas une droite verticale, puisque la fonction est dérivable sur l'intervalle.\\On en déduit que la tangente $(T)$ au point d'abscisse $4$, admet une équation réduite de la forme :  $(T) : y=m x + p$. 

\medskip
$\bullet$ Détermination de $m$ :\\On sait que le nombre dérivé en $4$ est par définition, le coefficient directeur de la tangente au point d'abscisse $4$.\\Par conséquent, on a déjà : $m=f'(4)=4$.\\ On en déduit que  $(T) : y= 4 x + p$\\$\bullet$ Détermination de $p$ :\\ Pour cela, on utilise que si $f(4)=-5~~$, alors le point $A$ de coordonnées $(4;-5)$ appartient à $\mathcal{C}_f$ mais aussi à $(T)$.\\ On peut écrire $A(4;-5) = \mathcal{C}_f \cap (T)$.\\ On remplace alors les coordonnées de $A(4;-5)$ dans l'équation  $(T) : y= 4 x + p$.\\ $\begin{aligned} \phantom{\iff}&A(4;-5)\in (T)\\ \iff& -5= 4 \times 4  + p\\ \iff& p=-5 -4 \times 4  \\ \iff& p=-5 -16   \\ \iff& p=-21\\\end{aligned}$\\On peut conclure que : $(T) : {\color[HTML]{f15929}\boldsymbol{y=4x-21}}$. \end{minipage}

\end{Solution}

% @see : https://coopmaths.fr/alea?uuid=0e984&id=can1F15&n=1&d=10&alea=1Alo2232323232342345678910111213141516171819202122232425262728293031323334353637383940414243444546474823456789101112&cd=1&cols=1
\needspace{10\baselineskip}
\begin{exercice}[BaremeTotal=false]


 La courbe représente une fonction $f$ et la droite est la tangente au point d'abscisse $-1$.\\
        
        Déterminer $f'(-1)$.\begin{center}\begin{tikzpicture}[baseline,scale = 0.5]

    \tikzset{
      point/.style={
        thick,
        draw,
        cross out,
        inner sep=0pt,
        minimum width=5pt,
        minimum height=5pt,
      },
    }
    \clip (-8,-3) rectangle (8,10);
    	
	\draw[color ={black},line width = 2,>=latex,->] (-8,0)--(8,0);
	\draw[color ={black},line width = 2,>=latex,->] (0,-3)--(0,10);
	\draw[color ={black},opacity = 0.5] (-8,1)--(8,1);
	\draw[color ={black},opacity = 0.5] (-8,2)--(8,2);
	\draw[color ={black},opacity = 0.5] (-8,3)--(8,3);
	\draw[color ={black},opacity = 0.5] (-8,4)--(8,4);
	\draw[color ={black},opacity = 0.5] (-8,5)--(8,5);
	\draw[color ={black},opacity = 0.5] (-8,6)--(8,6);
	\draw[color ={black},opacity = 0.5] (-8,7)--(8,7);
	\draw[color ={black},opacity = 0.5] (-8,8)--(8,8);
	\draw[color ={black},opacity = 0.5] (-8,9)--(8,9);
	\draw[color ={black},opacity = 0.5] (-8,10)--(8,10);
	\draw[color ={black},opacity = 0.5] (-8,-1)--(8,-1);
	\draw[color ={black},opacity = 0.5] (-8,-2)--(8,-2);
	\draw[color ={black},opacity = 0.5] (-8,-3)--(8,-3);
	\draw[color ={black},opacity = 0.5] (1,-3)--(1,10);
	\draw[color ={black},opacity = 0.5] (2,-3)--(2,10);
	\draw[color ={black},opacity = 0.5] (3,-3)--(3,10);
	\draw[color ={black},opacity = 0.5] (4,-3)--(4,10);
	\draw[color ={black},opacity = 0.5] (5,-3)--(5,10);
	\draw[color ={black},opacity = 0.5] (6,-3)--(6,10);
	\draw[color ={black},opacity = 0.5] (7,-3)--(7,10);
	\draw[color ={black},opacity = 0.5] (8,-3)--(8,10);
	\draw[color ={black},opacity = 0.5] (-1,-3)--(-1,10);
	\draw[color ={black},opacity = 0.5] (-2,-3)--(-2,10);
	\draw[color ={black},opacity = 0.5] (-3,-3)--(-3,10);
	\draw[color ={black},opacity = 0.5] (-4,-3)--(-4,10);
	\draw[color ={black},opacity = 0.5] (-5,-3)--(-5,10);
	\draw[color ={black},opacity = 0.5] (-6,-3)--(-6,10);
	\draw[color ={black},opacity = 0.5] (-7,-3)--(-7,10);
	\draw[color ={black},opacity = 0.5] (-8,-3)--(-8,10);
	\draw[color ={gray},opacity = 0.3] (-8,0)--(8,0);
	\draw[color ={gray},opacity = 0.3] (-8,0)--(8,0);
	\draw[color ={gray},opacity = 0.3] (0,-3)--(0,10);
	\draw[color ={gray},opacity = 0.3] (0,-3)--(0,10);
	\draw[color ={black},line width = 2] (2,-0.2)--(2,0.2);
	\draw[color ={black},line width = 2] (4,-0.2)--(4,0.2);
	\draw[color ={black},line width = 2] (6,-0.2)--(6,0.2);
	\draw[color ={black},line width = 2] (-2,-0.2)--(-2,0.2);
	\draw[color ={black},line width = 2] (-4,-0.2)--(-4,0.2);
	\draw[color ={black},line width = 2] (-6,-0.2)--(-6,0.2);
	\draw[color ={black},line width = 2] (-0.2,1)--(0.2,1);
	\draw[color ={black},line width = 2] (-0.2,2)--(0.2,2);
	\draw[color ={black},line width = 2] (-0.2,3)--(0.2,3);
	\draw[color ={black},line width = 2] (-0.2,4)--(0.2,4);
	\draw[color ={black},line width = 2] (-0.2,5)--(0.2,5);
	\draw[color ={black},line width = 2] (-0.2,6)--(0.2,6);
	\draw[color ={black},line width = 2] (-0.2,7)--(0.2,7);
	\draw[color ={black},line width = 2] (-0.2,8)--(0.2,8);
	\draw[color ={black},line width = 2] (-0.2,9)--(0.2,9);
	\draw[color ={black},line width = 2] (-0.2,-1)--(0.2,-1);
	\draw[color ={black},line width = 2] (-0.2,-2)--(0.2,-2);
	\draw (2,-0.4) node[anchor = center] {\scriptsize \color{black}{$1$}};
	\draw (4,-0.4) node[anchor = center] {\scriptsize \color{black}{$2$}};
	\draw (6,-0.4) node[anchor = center] {\scriptsize \color{black}{$3$}};
	\draw (-2,-0.4) node[anchor = center] {\scriptsize \color{black}{$-1$}};
	\draw (-4,-0.4) node[anchor = center] {\scriptsize \color{black}{$-2$}};
	\draw (-6,-0.4) node[anchor = center] {\scriptsize \color{black}{$-3$}};
	\draw (-0.8,2.1) node[anchor = center] {\footnotesize \color{black}{$2$}};
	\draw (-0.8,4.1) node[anchor = center] {\footnotesize \color{black}{$4$}};
	\draw (-0.8,6.1) node[anchor = center] {\footnotesize \color{black}{$6$}};
	\draw (-0.8,8.1) node[anchor = center] {\footnotesize \color{black}{$8$}};
	\draw [color={black}] (-0.3,-0.3) node[anchor = center,scale=1, rotate = 0] {O};
	
	\draw[color={blue},line width = 2] (-6.4,9.68)--(-6,8)--(-5.6,6.48)--(-5.2,5.12)--(-4.8,3.92)--(-4.4,2.88)--(-4,2)--(-3.6,1.28)--(-3.2,0.72)--(-2.8,0.32)--(-2.4,0.08)--(-2,0)--(-1.6,0.08)--(-1.2,0.32)--(-0.8,0.72)--(-0.4,1.28)--(0,2)--(0.4,2.88)--(0.8,3.92)--(1.2,5.12)--(1.6,6.48)--(2,8)--(2.4,9.68);
	
	\draw[color={red},line width = 2] (-8,0)--(-7.6,0)--(-7.2,0)--(-6.8,0)--(-6.4,0)--(-6,0)--(-5.6,0)--(-5.2,0)--(-4.8,0)--(-4.4,0)--(-4,0)--(-3.6,0)--(-3.2,0)--(-2.8,0)--(-2.4,0)--(-2,0)--(-1.6,0)--(-1.2,0)--(-0.8,0)--(-0.4,0)--(0,0)--(0.4,0)--(0.8,0)--(1.2,0)--(1.6,0)--(2,0)--(2.4,0)--(2.8,0)--(3.2,0)--(3.6,0)--(4,0)--(4.4,0)--(4.8,0)--(5.2,0)--(5.6,0)--(6,0)--(6.4,0)--(6.8,0)--(7.2,0)--(7.6,0)--(8,0);

\end{tikzpicture}\end{center}
\end{exercice}

\begin{Solution}
 $f'(-1)$ est donné par le coefficient directeur de la tangente à la courbe au point d'abscisse $-1$, soit $0$.
\end{Solution}

% @see : https://coopmaths.fr/alea?uuid=6f32d&id=can1F16&n=1&d=10&alea=dqSw2232323232342345678910111213141516171819202122232425262728293031323334353637383940414243444546474823456789101112&cd=1&cols=1
\needspace{10\baselineskip}
\newpage
\begin{exercice}[BaremeTotal=false]


 On donne les représentations graphiques d'une fonction et de sa dérivée.\\
        Donner l'équation réduite de la tangente à la courbe de $f$ en $x=1$. \\ \begin{minipage}{0.5\linewidth}
    \begin{tikzpicture}[baseline, scale=0.8]

    \tikzset{
      point/.style={
        thick,
        draw,
        cross out,
        inner sep=0pt,
        minimum width=5pt,
        minimum height=5pt,
      },
    }
    \clip (-4.5,-4.5) rectangle (5.75,13.5);
    	
	\draw[color ={black},>=latex,->] (-3,0)--(3,0);
	\draw[color ={black},>=latex,->] (0,-3)--(0,12);
	\draw[color ={black},opacity = 0.4] (-3,1)--(3,1);
	\draw[color ={black},opacity = 0.4] (-3,2)--(3,2);
	\draw[color ={black},opacity = 0.4] (-3,3)--(3,3);
	\draw[color ={black},opacity = 0.4] (-3,4)--(3,4);
	\draw[color ={black},opacity = 0.4] (-3,5)--(3,5);
	\draw[color ={black},opacity = 0.4] (-3,6)--(3,6);
	\draw[color ={black},opacity = 0.4] (-3,7)--(3,7);
	\draw[color ={black},opacity = 0.4] (-3,8)--(3,8);
	\draw[color ={black},opacity = 0.4] (-3,9)--(3,9);
	\draw[color ={black},opacity = 0.4] (-3,10)--(3,10);
	\draw[color ={black},opacity = 0.4] (-3,11)--(3,11);
	\draw[color ={black},opacity = 0.4] (-3,12)--(3,12);
	\draw[color ={black},opacity = 0.4] (-3,-1)--(3,-1);
	\draw[color ={black},opacity = 0.4] (-3,-2)--(3,-2);
	\draw[color ={black},opacity = 0.4] (-3,-3)--(3,-3);
	\draw[color ={black},opacity = 0.4] (1,-3)--(1,12);
	\draw[color ={black},opacity = 0.4] (2,-3)--(2,12);
	\draw[color ={black},opacity = 0.4] (3,-3)--(3,12);
	\draw[color ={black},opacity = 0.4] (-1,-3)--(-1,12);
	\draw[color ={black},opacity = 0.4] (-2,-3)--(-2,12);
	\draw[color ={black},opacity = 0.4] (-3,-3)--(-3,12);
	\draw[color ={gray},opacity = 0.3] (-3,0)--(3,0);
	\draw[color ={gray},opacity = 0.3] (-3,0.5)--(3,0.5);
	\draw[color ={gray},opacity = 0.3] (-3,1.5)--(3,1.5);
	\draw[color ={gray},opacity = 0.3] (-3,2.5)--(3,2.5);
	\draw[color ={gray},opacity = 0.3] (-3,3.5)--(3,3.5);
	\draw[color ={gray},opacity = 0.3] (-3,4.5)--(3,4.5);
	\draw[color ={gray},opacity = 0.3] (-3,5.5)--(3,5.5);
	\draw[color ={gray},opacity = 0.3] (-3,6.5)--(3,6.5);
	\draw[color ={gray},opacity = 0.3] (-3,7.5)--(3,7.5);
	\draw[color ={gray},opacity = 0.3] (-3,8.5)--(3,8.5);
	\draw[color ={gray},opacity = 0.3] (-3,9.5)--(3,9.5);
	\draw[color ={gray},opacity = 0.3] (-3,10.5)--(3,10.5);
	\draw[color ={gray},opacity = 0.3] (-3,11.5)--(3,11.5);
	\draw[color ={gray},opacity = 0.3] (-3,0)--(3,0);
	\draw[color ={gray},opacity = 0.3] (-3,-0.5)--(3,-0.5);
	\draw[color ={gray},opacity = 0.3] (-3,-1.5)--(3,-1.5);
	\draw[color ={gray},opacity = 0.3] (-3,-2.5)--(3,-2.5);
	\draw[color ={gray},opacity = 0.3] (0,-3)--(0,12);
	\draw[color ={gray},opacity = 0.3] (0,-3)--(0,12);
	\draw[color ={black},line width = 1.2] (1,-0.13)--(1,0.13);
	\draw[color ={black},line width = 1.2] (2,-0.13)--(2,0.13);
	\draw[color ={black},line width = 1.2] (3,-0.13)--(3,0.13);
	\draw[color ={black},line width = 1.2] (-1,-0.13)--(-1,0.13);
	\draw[color ={black},line width = 1.2] (-2,-0.13)--(-2,0.13);
	\draw[color ={black},line width = 1.2] (-3,-0.13)--(-3,0.13);
	\draw[color ={black},line width = 1.2] (-0.13,1)--(0.13,1);
	\draw[color ={black},line width = 1.2] (-0.13,2)--(0.13,2);
	\draw[color ={black},line width = 1.2] (-0.13,3)--(0.13,3);
	\draw[color ={black},line width = 1.2] (-0.13,4)--(0.13,4);
	\draw[color ={black},line width = 1.2] (-0.13,5)--(0.13,5);
	\draw[color ={black},line width = 1.2] (-0.13,6)--(0.13,6);
	\draw[color ={black},line width = 1.2] (-0.13,7)--(0.13,7);
	\draw[color ={black},line width = 1.2] (-0.13,8)--(0.13,8);
	\draw[color ={black},line width = 1.2] (-0.13,9)--(0.13,9);
	\draw[color ={black},line width = 1.2] (-0.13,10)--(0.13,10);
	\draw[color ={black},line width = 1.2] (-0.13,11)--(0.13,11);
	\draw[color ={black},line width = 1.2] (-0.13,12)--(0.13,12);
	\draw[color ={black},line width = 1.2] (-0.13,-1)--(0.13,-1);
	\draw[color ={black},line width = 1.2] (-0.13,-2)--(0.13,-2);
	\draw[color ={black},line width = 1.2] (-0.13,-3)--(0.13,-3);
	\draw (3,-0.4) node[anchor = center] {\scriptsize \color{black}{$3$}};
	\draw (-3,-0.4) node[anchor = center] {\scriptsize \color{black}{$-3$}};
	\draw (-0.5,3.1) node[anchor = center] {\footnotesize \color{black}{$3$}};
	\draw (-0.5,6.1) node[anchor = center] {\footnotesize \color{black}{$6$}};
	\draw (-0.5,9.1) node[anchor = center] {\footnotesize \color{black}{$9$}};
	\draw (-0.5,-2.9) node[anchor = center] {\footnotesize \color{black}{$-3$}};
	\draw [color={black}] (-0.3,-0.3) node[anchor = center,scale=1, rotate = 0] {O};
	\draw (3,10) node[anchor = center] {\footnotesize \color{blue}{$\Large \mathcal{C}_f$}};
	
	\draw[color={blue},line width = 2] (-3,5)--(-2.8,4.24)--(-2.6,3.56)--(-2.4,2.96)--(-2.2,2.44)--(-2,2)--(-1.8,1.64)--(-1.6,1.36)--(-1.4,1.16)--(-1.2,1.04)--(-1,1)--(-0.8,1.04)--(-0.6,1.16)--(-0.4,1.36)--(-0.2,1.64)--(0,2)--(0.2,2.44)--(0.4,2.96)--(0.6,3.56)--(0.8,4.24)--(1,5)--(1.2,5.84)--(1.4,6.76)--(1.6,7.76)--(1.8,8.84)--(2,10)--(2.2,11.24)--(2.4,12.56);

\end{tikzpicture}\\
    \end{minipage}
    \begin{minipage}{0.5\linewidth}
    \begin{tikzpicture}[baseline, scale=0.8]

    \tikzset{
      point/.style={
        thick,
        draw,
        cross out,
        inner sep=0pt,
        minimum width=5pt,
        minimum height=5pt,
      },
    }
    \clip (-4.5,-6.5) rectangle (6.5,9.5);
    	
	\draw[color ={black},>=latex,->] (-3,0)--(3,0);
	\draw[color ={black},>=latex,->] (0,-5)--(0,8);
	\draw[color ={black},opacity = 0.4] (-3,1)--(3,1);
	\draw[color ={black},opacity = 0.4] (-3,2)--(3,2);
	\draw[color ={black},opacity = 0.4] (-3,3)--(3,3);
	\draw[color ={black},opacity = 0.4] (-3,4)--(3,4);
	\draw[color ={black},opacity = 0.4] (-3,5)--(3,5);
	\draw[color ={black},opacity = 0.4] (-3,6)--(3,6);
	\draw[color ={black},opacity = 0.4] (-3,7)--(3,7);
	\draw[color ={black},opacity = 0.4] (-3,8)--(3,8);
	\draw[color ={black},opacity = 0.4] (-3,-1)--(3,-1);
	\draw[color ={black},opacity = 0.4] (-3,-2)--(3,-2);
	\draw[color ={black},opacity = 0.4] (-3,-3)--(3,-3);
	\draw[color ={black},opacity = 0.4] (-3,-4)--(3,-4);
	\draw[color ={black},opacity = 0.4] (-3,-5)--(3,-5);
	\draw[color ={black},opacity = 0.4] (1,-5)--(1,8);
	\draw[color ={black},opacity = 0.4] (2,-5)--(2,8);
	\draw[color ={black},opacity = 0.4] (3,-5)--(3,8);
	\draw[color ={black},opacity = 0.4] (-1,-5)--(-1,8);
	\draw[color ={black},opacity = 0.4] (-2,-5)--(-2,8);
	\draw[color ={black},opacity = 0.4] (-3,-5)--(-3,8);
	\draw[color ={gray},opacity = 0.3] (-3,0)--(3,0);
	\draw[color ={gray},opacity = 0.3] (-3,0.5)--(3,0.5);
	\draw[color ={gray},opacity = 0.3] (-3,1.5)--(3,1.5);
	\draw[color ={gray},opacity = 0.3] (-3,2.5)--(3,2.5);
	\draw[color ={gray},opacity = 0.3] (-3,3.5)--(3,3.5);
	\draw[color ={gray},opacity = 0.3] (-3,4.5)--(3,4.5);
	\draw[color ={gray},opacity = 0.3] (-3,5.5)--(3,5.5);
	\draw[color ={gray},opacity = 0.3] (-3,6.5)--(3,6.5);
	\draw[color ={gray},opacity = 0.3] (-3,7.5)--(3,7.5);
	\draw[color ={gray},opacity = 0.3] (-3,0)--(3,0);
	\draw[color ={gray},opacity = 0.3] (-3,-0.5)--(3,-0.5);
	\draw[color ={gray},opacity = 0.3] (-3,-1.5)--(3,-1.5);
	\draw[color ={gray},opacity = 0.3] (-3,-2.5)--(3,-2.5);
	\draw[color ={gray},opacity = 0.3] (-3,-3.5)--(3,-3.5);
	\draw[color ={gray},opacity = 0.3] (-3,-4.5)--(3,-4.5);
	\draw[color ={gray},opacity = 0.3] (0,-5)--(0,8);
	\draw[color ={gray},opacity = 0.3] (0,-5)--(0,8);
	\draw[color ={black},line width = 1.2] (1,-0.13)--(1,0.13);
	\draw[color ={black},line width = 1.2] (2,-0.13)--(2,0.13);
	\draw[color ={black},line width = 1.2] (3,-0.13)--(3,0.13);
	\draw[color ={black},line width = 1.2] (-1,-0.13)--(-1,0.13);
	\draw[color ={black},line width = 1.2] (-2,-0.13)--(-2,0.13);
	\draw[color ={black},line width = 1.2] (-3,-0.13)--(-3,0.13);
	\draw[color ={black},line width = 1.2] (-0.13,1)--(0.13,1);
	\draw[color ={black},line width = 1.2] (-0.13,2)--(0.13,2);
	\draw[color ={black},line width = 1.2] (-0.13,3)--(0.13,3);
	\draw[color ={black},line width = 1.2] (-0.13,4)--(0.13,4);
	\draw[color ={black},line width = 1.2] (-0.13,5)--(0.13,5);
	\draw[color ={black},line width = 1.2] (-0.13,6)--(0.13,6);
	\draw[color ={black},line width = 1.2] (-0.13,7)--(0.13,7);
	\draw[color ={black},line width = 1.2] (-0.13,8)--(0.13,8);
	\draw[color ={black},line width = 1.2] (-0.13,9)--(0.13,9);
	\draw[color ={black},line width = 1.2] (-0.13,10)--(0.13,10);
	\draw[color ={black},line width = 1.2] (-0.13,11)--(0.13,11);
	\draw[color ={black},line width = 1.2] (-0.13,12)--(0.13,12);
	\draw[color ={black},line width = 1.2] (-0.13,-1)--(0.13,-1);
	\draw[color ={black},line width = 1.2] (-0.13,-2)--(0.13,-2);
	\draw[color ={black},line width = 1.2] (-0.13,-3)--(0.13,-3);
	\draw (3,-0.4) node[anchor = center] {\scriptsize \color{black}{$3$}};
	\draw (-3,-0.4) node[anchor = center] {\scriptsize \color{black}{$-3$}};
	\draw (-0.5,3.1) node[anchor = center] {\footnotesize \color{black}{$3$}};
	\draw (-0.5,6.1) node[anchor = center] {\footnotesize \color{black}{$6$}};
	\draw (-0.5,-2.9) node[anchor = center] {\footnotesize \color{black}{$-3$}};
	\draw [color={black}] (-0.3,-0.3) node[anchor = center,scale=1, rotate = 0] {O};
	\draw (3,6) node[anchor = center] {\footnotesize \color{red}{$\Large\mathcal{C}_f\prime$}};
	
	\draw[color={red},line width = 2] (-3,-4)--(-2.8,-3.6)--(-2.6,-3.2)--(-2.4,-2.8)--(-2.2,-2.4)--(-2,-2)--(-1.8,-1.6)--(-1.6,-1.2)--(-1.4,-0.8)--(-1.2,-0.4)--(-1,0)--(-0.8,0.4)--(-0.6,0.8)--(-0.4,1.2)--(-0.2,1.6)--(0,2)--(0.2,2.4)--(0.4,2.8)--(0.6,3.2)--(0.8,3.6)--(1,4)--(1.2,4.4)--(1.4,4.8)--(1.6,5.2)--(1.8,5.6)--(2,6)--(2.2,6.4)--(2.4,6.8)--(2.6,7.2)--(2.8,7.6)--(3,8);

\end{tikzpicture}\\
    \end{minipage}
    
\end{exercice}

\begin{Solution}
 L'équation réduite de la tangente au point d'abscisse $1$ est  : $y=f'(1)(x-1)+f(1)$.\\
        On lit graphiquement $f(1)=5$ et $f'(1)=4$.\\
        L'équation réduite de la tangente est donc donnée par :
        $y=4(x-1)+5$, soit $y=4x+1$.
\end{Solution}

% @see : https://coopmaths.fr/alea?uuid=a1ba2&id=can1F14&n=1&d=10&alea=rkIp2232323232342345678910111213141516171819202122232425262728293031323334353637383940414243444546474823456789101112&cd=1&cols=1
\needspace{10\baselineskip}
\begin{exercice}[BaremeTotal=false]


 Soit $f$ la fonction définie sur $\mathbb{R}$ par : $f(x)= -5x^2+2x-7$.\\
        En calculant la limite du taux d'accroissement, déterminer $f'(1)$.
\end{exercice}

\begin{Solution}
 $f$ est une fonction polynôme du second degré de la forme $f(x)=ax^2+bx+c$.\\
    La fonction dérivée est donnée par la somme des dérivées des fonctions $u$ et $v$ définies par $u(x)=-5x^2$ et $v(x)=2x-7$.\\
     Comme $u'(x)=-10x$ et $v'(x)=2$, on obtient  $f'(x)=-10x+2$. \\
     Ainsi, $f'(1)=-10\times 1+2=-8$.
\end{Solution}

% @see : https://coopmaths.fr/alea?uuid=3c690&id=can1F13&n=1&d=10&alea=EUDu2232323232342345678910111213141516171819202122232425262728293031323334353637383940414243444546474823456789101112&cd=1&cols=1
\needspace{10\baselineskip}
\begin{exercice}[BaremeTotal=false]


 Déterminer le coefficient directeur de la tangente à la courbe représentative de la fonction carré au point d'abscisse $-5$.    
\end{exercice}

\begin{Solution}
 Le coefficient directeur de la tangente au point d'abscisse $a$ est donné par le nombre dérivé $f'(a)$.\\
        La fonction $f$ définie par $f(x)=x^2$ a pour fonction dérivée la fonction $f'$ définie par $f'(x)=2x$.\\
        Comme $f'(-5)=2\times (-5)=-10$, le coefficient directeur de la tangente au point d'abscisse $-5$ est : $-10$. 
\end{Solution}

% @see : https://coopmaths.fr/alea?uuid=ebd89&id=1AN14-1&n=1&d=10&s=3&alea=ocYr2232323232342345678910111213141516171819202122232425262728293031323334353637383940414243444546474823456789101112&cd=0&cols=1
\needspace{10\baselineskip}
\begin{exercice}[BaremeTotal=false]


 Donner la dérivée de la fonction $f$, dérivable pour tout $x\in\R$, définie par  $f(x)=-\dfrac{3}{7}x+10$.
\end{exercice}

\begin{Solution}
 L'expression de la dérivée de la fonction $f$ définie par $f(x)=-\dfrac{3}{7}x+10$ est : ${\color[HTML]{f15929}\boldsymbol{f'(x)=-\dfrac{3}{7}}}$.
\end{Solution}

% @see : https://coopmaths.fr/alea?uuid=ebd89&id=1AN14-1&n=1&d=10&s=4&alea=uAlP2232323232342345678910111213141516171819202122232425262728293031323334353637383940414243444546474823456789101112&cd=0&cols=1
\needspace{10\baselineskip}
\begin{exercice}[BaremeTotal=false]


 Donner la dérivée de la fonction $f$, dérivable pour tout $x\in\R$, définie par  $f(x)=-10x^{14}$.
\end{exercice}

\begin{Solution}
 L'expression de la dérivée de la fonction $f$ définie par $f(x)=-10x^{14}$ est : ${\color[HTML]{f15929}\boldsymbol{f'(x)=-140x^{13}}}$.
\end{Solution}

% @see : https://coopmaths.fr/alea?uuid=ebd89&id=1AN14-1&n=1&d=10&s=7&alea=vjN72232323232342345678910111213141516171819202122232425262728293031323334353637383940414243444546474823456789101112&cd=0&cols=1
\needspace{10\baselineskip}
\begin{exercice}[BaremeTotal=false]


 Donner la dérivée de la fonction $f$, dérivable pour tout $x\in\R^*$, définie par  $f(x)=\dfrac{-1}{8x}$.
\end{exercice}

\begin{Solution}
 L'expression de la dérivée de la fonction $f$ définie par $f(x)=\dfrac{-1}{8x}$ est : ${\color[HTML]{f15929}\boldsymbol{f'(x)=\dfrac{1}{8x^2}}}$.
\end{Solution}

% @see : https://coopmaths.fr/alea?uuid=a83c0&id=1AN14-4&n=1&d=10&s=1&alea=Njr32232323232342345678910111213141516171819202122232425262728293031323334353637383940414243444546474823456789101112&cd=0&cols=1
\needspace{10\baselineskip}
\begin{exercice}[BaremeTotal=false]


 Donner l'expression de la dérivée de la fonction $f$ définie sur $\R^{*}$ par $f(x)=-3x^{2}-\dfrac{1}{x}$.\\
\end{exercice}

\begin{Solution}
 L'expression de la dérivée de la fonction $f$ définie par $f(x)=-3x^{2}-\dfrac{1}{x}$ est :\\$f'(x)={\color[HTML]{f15929}\boldsymbol{-6x+\dfrac{1}{x^2}}}$.\\
\end{Solution}

% @see : https://coopmaths.fr/alea?uuid=a83c0&id=1AN14-4&n=1&d=10&s=4&alea=4Vpz2232323232342345678910111213141516171819202122232425262728293031323334353637383940414243444546474823456789101112&cd=1&cols=1
\needspace{10\baselineskip}
\begin{exercice}[BaremeTotal=false]


 Donner l'expression de la dérivée de la fonction $f$ définie sur $\R_+^{*}$ par $f(x)=-3x^{2}-5x+5-4\sqrt{x}-\dfrac{3}{x}$.\\
\end{exercice}

\begin{Solution}
 $(-3x^{2}-5x+5)^\prime=-6x-5$\\$(-4\sqrt{x})^\prime=-\dfrac{4}{2\sqrt{x}}$\\$(-\dfrac{3}{x})^\prime=\dfrac{3}{x^2}$\\L'expression de la dérivée de la fonction $f$ définie par $f(x)=-3x^{2}-5x+5-4\sqrt{x}-\dfrac{3}{x}$ est :\\$f'(x)={\color[HTML]{f15929}\boldsymbol{-6x-5-\dfrac{4}{2\sqrt{x}}+\dfrac{3}{x^2}}}$.\\
\end{Solution}

\end{Maquette}
\clearpage
\begin{Maquette}[DM]{Niveau= ,Classe=\getdata[13]\mydata\ (v13), Date = 06/01/25  ,Theme=Exercices}


% @see : https://coopmaths.fr/alea?uuid=4c8c7&id=1AN11-3&n=1&d=10&s=2&alea=nawO223232323234234567891011121314151617181920212223242526272829303132333435363738394041424344454647482345678910111213&cd=1&cols=2
\needspace{10\baselineskip}
\begin{exercice}[BaremeTotal=false]
\label{eleve:13}


 \begin{minipage}[t]{\linewidth} Soit $f$ une fonction dérivable sur $[-5;5]$ et $\mathcal{C}_f$ sa courbe représentative.\\On sait que  $f(-3)=3~~$ et que $~~f'(-3)=4$.\\Déterminer une équation de la tangente $(T)$ à la courbe $\mathcal{C}_f$ au point d'abscisse $-3$,\\sans utiliser la formule de cours de l'équation de tangente. \end{minipage}

\end{exercice}

\begin{Solution}
 \begin{minipage}[t]{\linewidth} On sait que la tangente n'est pas une droite verticale, puisque la fonction est dérivable sur l'intervalle.\\On en déduit que la tangente $(T)$ au point d'abscisse $-3$, admet une équation réduite de la forme :  $(T) : y=m x + p$. 

\medskip
$\bullet$ Détermination de $m$ :\\On sait que le nombre dérivé en $-3$ est par définition, le coefficient directeur de la tangente au point d'abscisse $-3$.\\Par conséquent, on a déjà : $m=f'(-3)=4$.\\ On en déduit que  $(T) : y= 4 x + p$\\$\bullet$ Détermination de $p$ :\\ Pour cela, on utilise que si $f(-3)=3~~$, alors le point $A$ de coordonnées $(-3;3)$ appartient à $\mathcal{C}_f$ mais aussi à $(T)$.\\ On peut écrire $A(-3;3) = \mathcal{C}_f \cap (T)$.\\ On remplace alors les coordonnées de $A(-3;3)$ dans l'équation  $(T) : y= 4 x + p$.\\ $\begin{aligned} \phantom{\iff}&A(-3;3)\in (T)\\ \iff& 3= 4 \times (-3)  + p\\ \iff& p=3 -4 \times (-3)  \\ \iff& p=3 +12   \\ \iff& p=15\\\end{aligned}$\\On peut conclure que : $(T) : {\color[HTML]{f15929}\boldsymbol{y=4x+15}}$. \end{minipage}

\end{Solution}

% @see : https://coopmaths.fr/alea?uuid=0e984&id=can1F15&n=1&d=10&alea=1Alo223232323234234567891011121314151617181920212223242526272829303132333435363738394041424344454647482345678910111213&cd=1&cols=1
\needspace{10\baselineskip}
\begin{exercice}[BaremeTotal=false]


 La courbe représente une fonction $f$ et la droite est la tangente au point d'abscisse $2$.\\
        Déterminer $f'(2)$.\begin{center}\begin{tikzpicture}[baseline,scale = 0.5]

    \tikzset{
      point/.style={
        thick,
        draw,
        cross out,
        inner sep=0pt,
        minimum width=5pt,
        minimum height=5pt,
      },
    }
    \clip (-2,-2) rectangle (14,12);
    	
	\draw[color ={black},line width = 2,>=latex,->] (-2,0)--(14,0);
	\draw[color ={black},line width = 2,>=latex,->] (0,-2)--(0,12);
	\draw[color ={black},opacity = 0.5] (-2,1)--(14,1);
	\draw[color ={black},opacity = 0.5] (-2,2)--(14,2);
	\draw[color ={black},opacity = 0.5] (-2,3)--(14,3);
	\draw[color ={black},opacity = 0.5] (-2,4)--(14,4);
	\draw[color ={black},opacity = 0.5] (-2,5)--(14,5);
	\draw[color ={black},opacity = 0.5] (-2,6)--(14,6);
	\draw[color ={black},opacity = 0.5] (-2,7)--(14,7);
	\draw[color ={black},opacity = 0.5] (-2,8)--(14,8);
	\draw[color ={black},opacity = 0.5] (-2,9)--(14,9);
	\draw[color ={black},opacity = 0.5] (-2,10)--(14,10);
	\draw[color ={black},opacity = 0.5] (-2,11)--(14,11);
	\draw[color ={black},opacity = 0.5] (-2,12)--(14,12);
	\draw[color ={black},opacity = 0.5] (-2,-1)--(14,-1);
	\draw[color ={black},opacity = 0.5] (-2,-2)--(14,-2);
	\draw[color ={black},opacity = 0.5] (1,-2)--(1,12);
	\draw[color ={black},opacity = 0.5] (2,-2)--(2,12);
	\draw[color ={black},opacity = 0.5] (3,-2)--(3,12);
	\draw[color ={black},opacity = 0.5] (4,-2)--(4,12);
	\draw[color ={black},opacity = 0.5] (5,-2)--(5,12);
	\draw[color ={black},opacity = 0.5] (6,-2)--(6,12);
	\draw[color ={black},opacity = 0.5] (7,-2)--(7,12);
	\draw[color ={black},opacity = 0.5] (8,-2)--(8,12);
	\draw[color ={black},opacity = 0.5] (9,-2)--(9,12);
	\draw[color ={black},opacity = 0.5] (10,-2)--(10,12);
	\draw[color ={black},opacity = 0.5] (11,-2)--(11,12);
	\draw[color ={black},opacity = 0.5] (12,-2)--(12,12);
	\draw[color ={black},opacity = 0.5] (13,-2)--(13,12);
	\draw[color ={black},opacity = 0.5] (14,-2)--(14,12);
	\draw[color ={black},opacity = 0.5] (-1,-2)--(-1,12);
	\draw[color ={black},opacity = 0.5] (-2,-2)--(-2,12);
	\draw[color ={gray},opacity = 0.3] (-2,0)--(14,0);
	\draw[color ={gray},opacity = 0.3] (-2,0.5)--(14,0.5);
	\draw[color ={gray},opacity = 0.3] (-2,1.5)--(14,1.5);
	\draw[color ={gray},opacity = 0.3] (-2,2.5)--(14,2.5);
	\draw[color ={gray},opacity = 0.3] (-2,3.5)--(14,3.5);
	\draw[color ={gray},opacity = 0.3] (-2,4.5)--(14,4.5);
	\draw[color ={gray},opacity = 0.3] (-2,5.5)--(14,5.5);
	\draw[color ={gray},opacity = 0.3] (-2,6.5)--(14,6.5);
	\draw[color ={gray},opacity = 0.3] (-2,7.5)--(14,7.5);
	\draw[color ={gray},opacity = 0.3] (-2,8.5)--(14,8.5);
	\draw[color ={gray},opacity = 0.3] (-2,9.5)--(14,9.5);
	\draw[color ={gray},opacity = 0.3] (-2,10.5)--(14,10.5);
	\draw[color ={gray},opacity = 0.3] (-2,11.5)--(14,11.5);
	\draw[color ={gray},opacity = 0.3] (-2,0)--(14,0);
	\draw[color ={gray},opacity = 0.3] (-2,-0.5)--(14,-0.5);
	\draw[color ={gray},opacity = 0.3] (-2,-1.5)--(14,-1.5);
	\draw[color ={gray},opacity = 0.3] (0,-2)--(0,12);
	\draw[color ={gray},opacity = 0.3] (0.1,-2)--(0.1,12);
	\draw[color ={gray},opacity = 0.3] (0.2,-2)--(0.2,12);
	\draw[color ={gray},opacity = 0.3] (0.30000000000000004,-2)--(0.30000000000000004,12);
	\draw[color ={gray},opacity = 0.3] (0.4,-2)--(0.4,12);
	\draw[color ={gray},opacity = 0.3] (0.5,-2)--(0.5,12);
	\draw[color ={gray},opacity = 0.3] (0.6,-2)--(0.6,12);
	\draw[color ={gray},opacity = 0.3] (0.7,-2)--(0.7,12);
	\draw[color ={gray},opacity = 0.3] (0.7999999999999999,-2)--(0.7999999999999999,12);
	\draw[color ={gray},opacity = 0.3] (0.8999999999999999,-2)--(0.8999999999999999,12);
	\draw[color ={gray},opacity = 0.3] (1.0999999999999999,-2)--(1.0999999999999999,12);
	\draw[color ={gray},opacity = 0.3] (1.2,-2)--(1.2,12);
	\draw[color ={gray},opacity = 0.3] (1.3,-2)--(1.3,12);
	\draw[color ={gray},opacity = 0.3] (1.4000000000000001,-2)--(1.4000000000000001,12);
	\draw[color ={gray},opacity = 0.3] (1.5000000000000002,-2)--(1.5000000000000002,12);
	\draw[color ={gray},opacity = 0.3] (1.6000000000000003,-2)--(1.6000000000000003,12);
	\draw[color ={gray},opacity = 0.3] (1.7000000000000004,-2)--(1.7000000000000004,12);
	\draw[color ={gray},opacity = 0.3] (1.8000000000000005,-2)--(1.8000000000000005,12);
	\draw[color ={gray},opacity = 0.3] (1.9000000000000006,-2)--(1.9000000000000006,12);
	\draw[color ={gray},opacity = 0.3] (2.1000000000000005,-2)--(2.1000000000000005,12);
	\draw[color ={gray},opacity = 0.3] (2.2000000000000006,-2)--(2.2000000000000006,12);
	\draw[color ={gray},opacity = 0.3] (2.3000000000000007,-2)--(2.3000000000000007,12);
	\draw[color ={gray},opacity = 0.3] (2.400000000000001,-2)--(2.400000000000001,12);
	\draw[color ={gray},opacity = 0.3] (2.500000000000001,-2)--(2.500000000000001,12);
	\draw[color ={gray},opacity = 0.3] (2.600000000000001,-2)--(2.600000000000001,12);
	\draw[color ={gray},opacity = 0.3] (2.700000000000001,-2)--(2.700000000000001,12);
	\draw[color ={gray},opacity = 0.3] (2.800000000000001,-2)--(2.800000000000001,12);
	\draw[color ={gray},opacity = 0.3] (2.9000000000000012,-2)--(2.9000000000000012,12);
	\draw[color ={gray},opacity = 0.3] (3.1000000000000014,-2)--(3.1000000000000014,12);
	\draw[color ={gray},opacity = 0.3] (3.2000000000000015,-2)--(3.2000000000000015,12);
	\draw[color ={gray},opacity = 0.3] (3.3000000000000016,-2)--(3.3000000000000016,12);
	\draw[color ={gray},opacity = 0.3] (3.4000000000000017,-2)--(3.4000000000000017,12);
	\draw[color ={gray},opacity = 0.3] (3.5000000000000018,-2)--(3.5000000000000018,12);
	\draw[color ={gray},opacity = 0.3] (3.600000000000002,-2)--(3.600000000000002,12);
	\draw[color ={gray},opacity = 0.3] (3.700000000000002,-2)--(3.700000000000002,12);
	\draw[color ={gray},opacity = 0.3] (3.800000000000002,-2)--(3.800000000000002,12);
	\draw[color ={gray},opacity = 0.3] (3.900000000000002,-2)--(3.900000000000002,12);
	\draw[color ={gray},opacity = 0.3] (4.100000000000001,-2)--(4.100000000000001,12);
	\draw[color ={gray},opacity = 0.3] (4.200000000000001,-2)--(4.200000000000001,12);
	\draw[color ={gray},opacity = 0.3] (4.300000000000001,-2)--(4.300000000000001,12);
	\draw[color ={gray},opacity = 0.3] (4.4,-2)--(4.4,12);
	\draw[color ={gray},opacity = 0.3] (4.5,-2)--(4.5,12);
	\draw[color ={gray},opacity = 0.3] (4.6,-2)--(4.6,12);
	\draw[color ={gray},opacity = 0.3] (4.699999999999999,-2)--(4.699999999999999,12);
	\draw[color ={gray},opacity = 0.3] (4.799999999999999,-2)--(4.799999999999999,12);
	\draw[color ={gray},opacity = 0.3] (4.899999999999999,-2)--(4.899999999999999,12);
	\draw[color ={gray},opacity = 0.3] (5.099999999999998,-2)--(5.099999999999998,12);
	\draw[color ={gray},opacity = 0.3] (5.1999999999999975,-2)--(5.1999999999999975,12);
	\draw[color ={gray},opacity = 0.3] (5.299999999999997,-2)--(5.299999999999997,12);
	\draw[color ={gray},opacity = 0.3] (5.399999999999997,-2)--(5.399999999999997,12);
	\draw[color ={gray},opacity = 0.3] (5.4999999999999964,-2)--(5.4999999999999964,12);
	\draw[color ={gray},opacity = 0.3] (5.599999999999996,-2)--(5.599999999999996,12);
	\draw[color ={gray},opacity = 0.3] (5.699999999999996,-2)--(5.699999999999996,12);
	\draw[color ={gray},opacity = 0.3] (5.799999999999995,-2)--(5.799999999999995,12);
	\draw[color ={gray},opacity = 0.3] (5.899999999999995,-2)--(5.899999999999995,12);
	\draw[color ={gray},opacity = 0.3] (6.099999999999994,-2)--(6.099999999999994,12);
	\draw[color ={gray},opacity = 0.3] (6.199999999999994,-2)--(6.199999999999994,12);
	\draw[color ={gray},opacity = 0.3] (6.299999999999994,-2)--(6.299999999999994,12);
	\draw[color ={gray},opacity = 0.3] (6.399999999999993,-2)--(6.399999999999993,12);
	\draw[color ={gray},opacity = 0.3] (6.499999999999993,-2)--(6.499999999999993,12);
	\draw[color ={gray},opacity = 0.3] (6.5999999999999925,-2)--(6.5999999999999925,12);
	\draw[color ={gray},opacity = 0.3] (6.699999999999992,-2)--(6.699999999999992,12);
	\draw[color ={gray},opacity = 0.3] (6.799999999999992,-2)--(6.799999999999992,12);
	\draw[color ={gray},opacity = 0.3] (6.8999999999999915,-2)--(6.8999999999999915,12);
	\draw[color ={gray},opacity = 0.3] (7.099999999999991,-2)--(7.099999999999991,12);
	\draw[color ={gray},opacity = 0.3] (7.19999999999999,-2)--(7.19999999999999,12);
	\draw[color ={gray},opacity = 0.3] (7.29999999999999,-2)--(7.29999999999999,12);
	\draw[color ={gray},opacity = 0.3] (7.39999999999999,-2)--(7.39999999999999,12);
	\draw[color ={gray},opacity = 0.3] (7.499999999999989,-2)--(7.499999999999989,12);
	\draw[color ={gray},opacity = 0.3] (7.599999999999989,-2)--(7.599999999999989,12);
	\draw[color ={gray},opacity = 0.3] (7.699999999999989,-2)--(7.699999999999989,12);
	\draw[color ={gray},opacity = 0.3] (7.799999999999988,-2)--(7.799999999999988,12);
	\draw[color ={gray},opacity = 0.3] (7.899999999999988,-2)--(7.899999999999988,12);
	\draw[color ={gray},opacity = 0.3] (8.099999999999987,-2)--(8.099999999999987,12);
	\draw[color ={gray},opacity = 0.3] (8.199999999999987,-2)--(8.199999999999987,12);
	\draw[color ={gray},opacity = 0.3] (8.299999999999986,-2)--(8.299999999999986,12);
	\draw[color ={gray},opacity = 0.3] (8.399999999999986,-2)--(8.399999999999986,12);
	\draw[color ={gray},opacity = 0.3] (8.499999999999986,-2)--(8.499999999999986,12);
	\draw[color ={gray},opacity = 0.3] (8.599999999999985,-2)--(8.599999999999985,12);
	\draw[color ={gray},opacity = 0.3] (8.699999999999985,-2)--(8.699999999999985,12);
	\draw[color ={gray},opacity = 0.3] (8.799999999999985,-2)--(8.799999999999985,12);
	\draw[color ={gray},opacity = 0.3] (8.899999999999984,-2)--(8.899999999999984,12);
	\draw[color ={gray},opacity = 0.3] (9.099999999999984,-2)--(9.099999999999984,12);
	\draw[color ={gray},opacity = 0.3] (9.199999999999983,-2)--(9.199999999999983,12);
	\draw[color ={gray},opacity = 0.3] (9.299999999999983,-2)--(9.299999999999983,12);
	\draw[color ={gray},opacity = 0.3] (9.399999999999983,-2)--(9.399999999999983,12);
	\draw[color ={gray},opacity = 0.3] (9.499999999999982,-2)--(9.499999999999982,12);
	\draw[color ={gray},opacity = 0.3] (9.599999999999982,-2)--(9.599999999999982,12);
	\draw[color ={gray},opacity = 0.3] (9.699999999999982,-2)--(9.699999999999982,12);
	\draw[color ={gray},opacity = 0.3] (9.799999999999981,-2)--(9.799999999999981,12);
	\draw[color ={gray},opacity = 0.3] (9.89999999999998,-2)--(9.89999999999998,12);
	\draw[color ={gray},opacity = 0.3] (10.09999999999998,-2)--(10.09999999999998,12);
	\draw[color ={gray},opacity = 0.3] (10.19999999999998,-2)--(10.19999999999998,12);
	\draw[color ={gray},opacity = 0.3] (10.29999999999998,-2)--(10.29999999999998,12);
	\draw[color ={gray},opacity = 0.3] (10.399999999999979,-2)--(10.399999999999979,12);
	\draw[color ={gray},opacity = 0.3] (10.499999999999979,-2)--(10.499999999999979,12);
	\draw[color ={gray},opacity = 0.3] (10.599999999999978,-2)--(10.599999999999978,12);
	\draw[color ={gray},opacity = 0.3] (10.699999999999978,-2)--(10.699999999999978,12);
	\draw[color ={gray},opacity = 0.3] (10.799999999999978,-2)--(10.799999999999978,12);
	\draw[color ={gray},opacity = 0.3] (10.899999999999977,-2)--(10.899999999999977,12);
	\draw[color ={gray},opacity = 0.3] (11.099999999999977,-2)--(11.099999999999977,12);
	\draw[color ={gray},opacity = 0.3] (11.199999999999976,-2)--(11.199999999999976,12);
	\draw[color ={gray},opacity = 0.3] (11.299999999999976,-2)--(11.299999999999976,12);
	\draw[color ={gray},opacity = 0.3] (11.399999999999975,-2)--(11.399999999999975,12);
	\draw[color ={gray},opacity = 0.3] (11.499999999999975,-2)--(11.499999999999975,12);
	\draw[color ={gray},opacity = 0.3] (11.599999999999975,-2)--(11.599999999999975,12);
	\draw[color ={gray},opacity = 0.3] (11.699999999999974,-2)--(11.699999999999974,12);
	\draw[color ={gray},opacity = 0.3] (11.799999999999974,-2)--(11.799999999999974,12);
	\draw[color ={gray},opacity = 0.3] (11.899999999999974,-2)--(11.899999999999974,12);
	\draw[color ={gray},opacity = 0.3] (12.099999999999973,-2)--(12.099999999999973,12);
	\draw[color ={gray},opacity = 0.3] (12.199999999999973,-2)--(12.199999999999973,12);
	\draw[color ={gray},opacity = 0.3] (12.299999999999972,-2)--(12.299999999999972,12);
	\draw[color ={gray},opacity = 0.3] (12.399999999999972,-2)--(12.399999999999972,12);
	\draw[color ={gray},opacity = 0.3] (12.499999999999972,-2)--(12.499999999999972,12);
	\draw[color ={gray},opacity = 0.3] (12.599999999999971,-2)--(12.599999999999971,12);
	\draw[color ={gray},opacity = 0.3] (12.69999999999997,-2)--(12.69999999999997,12);
	\draw[color ={gray},opacity = 0.3] (12.79999999999997,-2)--(12.79999999999997,12);
	\draw[color ={gray},opacity = 0.3] (12.89999999999997,-2)--(12.89999999999997,12);
	\draw[color ={gray},opacity = 0.3] (13.09999999999997,-2)--(13.09999999999997,12);
	\draw[color ={gray},opacity = 0.3] (13.199999999999969,-2)--(13.199999999999969,12);
	\draw[color ={gray},opacity = 0.3] (13.299999999999969,-2)--(13.299999999999969,12);
	\draw[color ={gray},opacity = 0.3] (13.399999999999968,-2)--(13.399999999999968,12);
	\draw[color ={gray},opacity = 0.3] (13.499999999999968,-2)--(13.499999999999968,12);
	\draw[color ={gray},opacity = 0.3] (13.599999999999968,-2)--(13.599999999999968,12);
	\draw[color ={gray},opacity = 0.3] (13.699999999999967,-2)--(13.699999999999967,12);
	\draw[color ={gray},opacity = 0.3] (13.799999999999967,-2)--(13.799999999999967,12);
	\draw[color ={gray},opacity = 0.3] (13.899999999999967,-2)--(13.899999999999967,12);
	\draw[color ={gray},opacity = 0.3] (0,-2)--(0,12);
	\draw[color ={gray},opacity = 0.3] (-0.1,-2)--(-0.1,12);
	\draw[color ={gray},opacity = 0.3] (-0.2,-2)--(-0.2,12);
	\draw[color ={gray},opacity = 0.3] (-0.30000000000000004,-2)--(-0.30000000000000004,12);
	\draw[color ={gray},opacity = 0.3] (-0.4,-2)--(-0.4,12);
	\draw[color ={gray},opacity = 0.3] (-0.5,-2)--(-0.5,12);
	\draw[color ={gray},opacity = 0.3] (-0.6,-2)--(-0.6,12);
	\draw[color ={gray},opacity = 0.3] (-0.7,-2)--(-0.7,12);
	\draw[color ={gray},opacity = 0.3] (-0.7999999999999999,-2)--(-0.7999999999999999,12);
	\draw[color ={gray},opacity = 0.3] (-0.8999999999999999,-2)--(-0.8999999999999999,12);
	\draw[color ={gray},opacity = 0.3] (-1.0999999999999999,-2)--(-1.0999999999999999,12);
	\draw[color ={gray},opacity = 0.3] (-1.2,-2)--(-1.2,12);
	\draw[color ={gray},opacity = 0.3] (-1.3,-2)--(-1.3,12);
	\draw[color ={gray},opacity = 0.3] (-1.4000000000000001,-2)--(-1.4000000000000001,12);
	\draw[color ={gray},opacity = 0.3] (-1.5000000000000002,-2)--(-1.5000000000000002,12);
	\draw[color ={gray},opacity = 0.3] (-1.6000000000000003,-2)--(-1.6000000000000003,12);
	\draw[color ={gray},opacity = 0.3] (-1.7000000000000004,-2)--(-1.7000000000000004,12);
	\draw[color ={gray},opacity = 0.3] (-1.8000000000000005,-2)--(-1.8000000000000005,12);
	\draw[color ={gray},opacity = 0.3] (-1.9000000000000006,-2)--(-1.9000000000000006,12);
	\draw[color ={black},line width = 2] (2,-0.2)--(2,0.2);
	\draw[color ={black},line width = 2] (4,-0.2)--(4,0.2);
	\draw[color ={black},line width = 2] (6,-0.2)--(6,0.2);
	\draw[color ={black},line width = 2] (8,-0.2)--(8,0.2);
	\draw[color ={black},line width = 2] (10,-0.2)--(10,0.2);
	\draw[color ={black},line width = 2] (12,-0.2)--(12,0.2);
	\draw[color ={black},line width = 2] (-0.2,2)--(0.2,2);
	\draw[color ={black},line width = 2] (-0.2,4)--(0.2,4);
	\draw[color ={black},line width = 2] (-0.2,6)--(0.2,6);
	\draw[color ={black},line width = 2] (-0.2,8)--(0.2,8);
	\draw[color ={black},line width = 2] (-0.2,10)--(0.2,10);
	\draw (2,-0.4) node[anchor = center] {\scriptsize \color{black}{$1$}};
	\draw (4,-0.4) node[anchor = center] {\scriptsize \color{black}{$2$}};
	\draw (6,-0.4) node[anchor = center] {\scriptsize \color{black}{$3$}};
	\draw (8,-0.4) node[anchor = center] {\scriptsize \color{black}{$4$}};
	\draw (10,-0.4) node[anchor = center] {\scriptsize \color{black}{$5$}};
	\draw (-0.5,2.1) node[anchor = center] {\footnotesize \color{black}{$1$}};
	\draw (-0.5,4.1) node[anchor = center] {\footnotesize \color{black}{$2$}};
	\draw (-0.5,6.1) node[anchor = center] {\footnotesize \color{black}{$3$}};
	\draw (-0.5,8.1) node[anchor = center] {\footnotesize \color{black}{$4$}};
	\draw (-0.5,10.1) node[anchor = center] {\footnotesize \color{black}{$5$}};
	\draw [color={black}] (-0.3,-0.3) node[anchor = center,scale=1, rotate = 0] {O};
	
	\draw[color={blue},line width = 2] (1,12)--(1.4,9.71)--(1.8,8.44)--(2.2,7.64)--(2.6,7.08)--(3,6.67)--(3.4,6.35)--(3.8,6.11)--(4.2,5.9)--(4.6,5.74)--(5,5.6)--(5.4,5.48)--(5.8,5.38)--(6.2,5.29)--(6.6,5.21)--(7,5.14)--(7.4,5.08)--(7.8,5.03)--(8.2,4.98)--(8.6,4.93)--(9,4.89)--(9.4,4.85)--(9.8,4.82)--(10.2,4.78)--(10.6,4.75)--(11,4.73)--(11.4,4.7)--(11.8,4.68)--(12.2,4.66)--(12.6,4.63)--(13,4.62)--(13.4,4.6)--(13.8,4.58);
	
	\draw[color={red},line width = 2] (-2,9)--(-1.6,8.8)--(-1.2,8.6)--(-0.8,8.4)--(-0.4,8.2)--(0,8)--(0.4,7.8)--(0.8,7.6)--(1.2,7.4)--(1.6,7.2)--(2,7)--(2.4,6.8)--(2.8,6.6)--(3.2,6.4)--(3.6,6.2)--(4,6)--(4.4,5.8)--(4.8,5.6)--(5.2,5.4)--(5.6,5.2)--(6,5)--(6.4,4.8)--(6.8,4.6)--(7.2,4.4)--(7.6,4.2)--(8,4)--(8.4,3.8)--(8.8,3.6)--(9.2,3.4)--(9.6,3.2)--(10,3)--(10.4,2.8)--(10.8,2.6)--(11.2,2.4)--(11.6,2.2)--(12,2)--(12.4,1.8)--(12.8,1.6)--(13.2,1.4)--(13.6,1.2)--(14,1);

\end{tikzpicture}\end{center}
\end{exercice}

\begin{Solution}
 $f'(2)$ est donné par le coefficient directeur de la tangente à la courbe au point d'abscisse $2$, soit $\dfrac{-2}{4}=-\dfrac{1{\color[HTML]{2563a5}\boldsymbol{\times2}} }{2{\color[HTML]{2563a5}\boldsymbol{\times2}}}=-\dfrac{1}{2}$.
\end{Solution}

% @see : https://coopmaths.fr/alea?uuid=6f32d&id=can1F16&n=1&d=10&alea=dqSw223232323234234567891011121314151617181920212223242526272829303132333435363738394041424344454647482345678910111213&cd=1&cols=1
\needspace{10\baselineskip}
\newpage
\begin{exercice}[BaremeTotal=false]


 On donne les représentations graphiques d'une fonction et de sa dérivée.\\
      Donner l'équation réduite de la tangente à la courbe de $f$ en $x=0$. \\ \begin{minipage}{0.5\linewidth}
    \begin{tikzpicture}[baseline, scale=0.8]

    \tikzset{
      point/.style={
        thick,
        draw,
        cross out,
        inner sep=0pt,
        minimum width=5pt,
        minimum height=5pt,
      },
    }
    \clip (-5.5,-9.5) rectangle (5.75,7.5);
    	
	\draw[color ={black},>=latex,->] (-4,0)--(4,0);
	\draw[color ={black},>=latex,->] (0,-8)--(0,6);
	\draw[color ={black},opacity = 0.4] (-4,1)--(4,1);
	\draw[color ={black},opacity = 0.4] (-4,2)--(4,2);
	\draw[color ={black},opacity = 0.4] (-4,3)--(4,3);
	\draw[color ={black},opacity = 0.4] (-4,4)--(4,4);
	\draw[color ={black},opacity = 0.4] (-4,5)--(4,5);
	\draw[color ={black},opacity = 0.4] (-4,6)--(4,6);
	\draw[color ={black},opacity = 0.4] (-4,-1)--(4,-1);
	\draw[color ={black},opacity = 0.4] (-4,-2)--(4,-2);
	\draw[color ={black},opacity = 0.4] (-4,-3)--(4,-3);
	\draw[color ={black},opacity = 0.4] (-4,-4)--(4,-4);
	\draw[color ={black},opacity = 0.4] (-4,-5)--(4,-5);
	\draw[color ={black},opacity = 0.4] (-4,-6)--(4,-6);
	\draw[color ={black},opacity = 0.4] (-4,-7)--(4,-7);
	\draw[color ={black},opacity = 0.4] (-4,-8)--(4,-8);
	\draw[color ={black},opacity = 0.4] (1,-8)--(1,6);
	\draw[color ={black},opacity = 0.4] (2,-8)--(2,6);
	\draw[color ={black},opacity = 0.4] (3,-8)--(3,6);
	\draw[color ={black},opacity = 0.4] (4,-8)--(4,6);
	\draw[color ={black},opacity = 0.4] (-1,-8)--(-1,6);
	\draw[color ={black},opacity = 0.4] (-2,-8)--(-2,6);
	\draw[color ={black},opacity = 0.4] (-3,-8)--(-3,6);
	\draw[color ={black},opacity = 0.4] (-4,-8)--(-4,6);
	\draw[color ={gray},opacity = 0.3] (-4,0)--(4,0);
	\draw[color ={gray},opacity = 0.3] (-4,0.5)--(4,0.5);
	\draw[color ={gray},opacity = 0.3] (-4,1.5)--(4,1.5);
	\draw[color ={gray},opacity = 0.3] (-4,2.5)--(4,2.5);
	\draw[color ={gray},opacity = 0.3] (-4,3.5)--(4,3.5);
	\draw[color ={gray},opacity = 0.3] (-4,4.5)--(4,4.5);
	\draw[color ={gray},opacity = 0.3] (-4,5.5)--(4,5.5);
	\draw[color ={gray},opacity = 0.3] (-4,0)--(4,0);
	\draw[color ={gray},opacity = 0.3] (-4,-0.5)--(4,-0.5);
	\draw[color ={gray},opacity = 0.3] (-4,-1.5)--(4,-1.5);
	\draw[color ={gray},opacity = 0.3] (-4,-2.5)--(4,-2.5);
	\draw[color ={gray},opacity = 0.3] (-4,-3.5)--(4,-3.5);
	\draw[color ={gray},opacity = 0.3] (-4,-4.5)--(4,-4.5);
	\draw[color ={gray},opacity = 0.3] (-4,-5.5)--(4,-5.5);
	\draw[color ={gray},opacity = 0.3] (-4,-6.5)--(4,-6.5);
	\draw[color ={gray},opacity = 0.3] (-4,-7)--(4,-7);
	\draw[color ={gray},opacity = 0.3] (-4,-7.5)--(4,-7.5);
	\draw[color ={gray},opacity = 0.3] (-4,-8)--(4,-8);
	\draw[color ={gray},opacity = 0.3] (0,-8)--(0,6);
	\draw[color ={gray},opacity = 0.3] (0,-8)--(0,6);
	\draw[color ={black},line width = 1.2] (1,-0.13)--(1,0.13);
	\draw[color ={black},line width = 1.2] (2,-0.13)--(2,0.13);
	\draw[color ={black},line width = 1.2] (3,-0.13)--(3,0.13);
	\draw[color ={black},line width = 1.2] (4,-0.13)--(4,0.13);
	\draw[color ={black},line width = 1.2] (-1,-0.13)--(-1,0.13);
	\draw[color ={black},line width = 1.2] (-2,-0.13)--(-2,0.13);
	\draw[color ={black},line width = 1.2] (-3,-0.13)--(-3,0.13);
	\draw[color ={black},line width = 1.2] (-4,-0.13)--(-4,0.13);
	\draw[color ={black},line width = 1.2] (-0.13,1)--(0.13,1);
	\draw[color ={black},line width = 1.2] (-0.13,2)--(0.13,2);
	\draw[color ={black},line width = 1.2] (-0.13,3)--(0.13,3);
	\draw[color ={black},line width = 1.2] (-0.13,4)--(0.13,4);
	\draw[color ={black},line width = 1.2] (-0.13,5)--(0.13,5);
	\draw[color ={black},line width = 1.2] (-0.13,6)--(0.13,6);
	\draw[color ={black},line width = 1.2] (-0.13,-1)--(0.13,-1);
	\draw[color ={black},line width = 1.2] (-0.13,-2)--(0.13,-2);
	\draw[color ={black},line width = 1.2] (-0.13,-3)--(0.13,-3);
	\draw[color ={black},line width = 1.2] (-0.13,-4)--(0.13,-4);
	\draw[color ={black},line width = 1.2] (-0.13,-5)--(0.13,-5);
	\draw[color ={black},line width = 1.2] (-0.13,-6)--(0.13,-6);
	\draw[color ={black},line width = 1.2] (-0.13,-7)--(0.13,-7);
	\draw[color ={black},line width = 1.2] (-0.13,-8)--(0.13,-8);
	\draw (2,-0.4) node[anchor = center] {\scriptsize \color{black}{$2$}};
	\draw (4,-0.4) node[anchor = center] {\scriptsize \color{black}{$4$}};
	\draw (-2,-0.4) node[anchor = center] {\scriptsize \color{black}{$-2$}};
	\draw (-4,-0.4) node[anchor = center] {\scriptsize \color{black}{$-4$}};
	\draw (-0.5,2.1) node[anchor = center] {\footnotesize \color{black}{$2$}};
	\draw (-0.5,4.1) node[anchor = center] {\footnotesize \color{black}{$4$}};
	\draw (-0.5,6.1) node[anchor = center] {\footnotesize \color{black}{$6$}};
	\draw (-0.5,-1.9) node[anchor = center] {\footnotesize \color{black}{$-2$}};
	\draw (-0.5,-3.9) node[anchor = center] {\footnotesize \color{black}{$-4$}};
	\draw (-0.5,-5.9) node[anchor = center] {\footnotesize \color{black}{$-6$}};
	\draw (-0.5,-7.9) node[anchor = center] {\footnotesize \color{black}{$-8$}};
	\draw [color={black}] (-0.3,-0.3) node[anchor = center,scale=1, rotate = 0] {O};
	\draw (3,4) node[anchor = center] {\footnotesize \color{blue}{$\Large \mathcal{C}_f$}};
	
	\draw[color={blue},line width = 2] (-4,-4)--(-3.8,-3.24)--(-3.6,-2.56)--(-3.4,-1.96)--(-3.2,-1.44)--(-3,-1)--(-2.8,-0.64)--(-2.6,-0.36)--(-2.4,-0.16)--(-2.2,-0.04)--(-2,0)--(-1.8,-0.04)--(-1.6,-0.16)--(-1.4,-0.36)--(-1.2,-0.64)--(-1,-1)--(-0.8,-1.44)--(-0.6,-1.96)--(-0.4,-2.56)--(-0.2,-3.24)--(0,-4)--(0.2,-4.84)--(0.4,-5.76)--(0.6,-6.76)--(0.8,-7.84);

\end{tikzpicture}\\
    \end{minipage}
    \begin{minipage}{0.5\linewidth}
    \begin{tikzpicture}[baseline, scale=0.8]

    \tikzset{
      point/.style={
        thick,
        draw,
        cross out,
        inner sep=0pt,
        minimum width=5pt,
        minimum height=5pt,
      },
    }
    \clip (-5.5,-9.5) rectangle (6.5,7.5);
    	
	\draw[color ={black},>=latex,->] (-4,0)--(4,0);
	\draw[color ={black},>=latex,->] (0,-8)--(0,6);
	\draw[color ={black},opacity = 0.4] (-4,1)--(4,1);
	\draw[color ={black},opacity = 0.4] (-4,2)--(4,2);
	\draw[color ={black},opacity = 0.4] (-4,3)--(4,3);
	\draw[color ={black},opacity = 0.4] (-4,4)--(4,4);
	\draw[color ={black},opacity = 0.4] (-4,5)--(4,5);
	\draw[color ={black},opacity = 0.4] (-4,6)--(4,6);
	\draw[color ={black},opacity = 0.4] (-4,-1)--(4,-1);
	\draw[color ={black},opacity = 0.4] (-4,-2)--(4,-2);
	\draw[color ={black},opacity = 0.4] (-4,-3)--(4,-3);
	\draw[color ={black},opacity = 0.4] (-4,-4)--(4,-4);
	\draw[color ={black},opacity = 0.4] (-4,-5)--(4,-5);
	\draw[color ={black},opacity = 0.4] (-4,-6)--(4,-6);
	\draw[color ={black},opacity = 0.4] (-4,-7)--(4,-7);
	\draw[color ={black},opacity = 0.4] (-4,-8)--(4,-8);
	\draw[color ={black},opacity = 0.4] (1,-8)--(1,6);
	\draw[color ={black},opacity = 0.4] (2,-8)--(2,6);
	\draw[color ={black},opacity = 0.4] (3,-8)--(3,6);
	\draw[color ={black},opacity = 0.4] (4,-8)--(4,6);
	\draw[color ={black},opacity = 0.4] (-1,-8)--(-1,6);
	\draw[color ={black},opacity = 0.4] (-2,-8)--(-2,6);
	\draw[color ={black},opacity = 0.4] (-3,-8)--(-3,6);
	\draw[color ={black},opacity = 0.4] (-4,-8)--(-4,6);
	\draw[color ={gray},opacity = 0.3] (-4,0)--(4,0);
	\draw[color ={gray},opacity = 0.3] (-4,0.5)--(4,0.5);
	\draw[color ={gray},opacity = 0.3] (-4,1.5)--(4,1.5);
	\draw[color ={gray},opacity = 0.3] (-4,2.5)--(4,2.5);
	\draw[color ={gray},opacity = 0.3] (-4,3.5)--(4,3.5);
	\draw[color ={gray},opacity = 0.3] (-4,4.5)--(4,4.5);
	\draw[color ={gray},opacity = 0.3] (-4,5.5)--(4,5.5);
	\draw[color ={gray},opacity = 0.3] (-4,0)--(4,0);
	\draw[color ={gray},opacity = 0.3] (-4,-0.5)--(4,-0.5);
	\draw[color ={gray},opacity = 0.3] (-4,-1.5)--(4,-1.5);
	\draw[color ={gray},opacity = 0.3] (-4,-2.5)--(4,-2.5);
	\draw[color ={gray},opacity = 0.3] (-4,-3.5)--(4,-3.5);
	\draw[color ={gray},opacity = 0.3] (-4,-4.5)--(4,-4.5);
	\draw[color ={gray},opacity = 0.3] (-4,-5.5)--(4,-5.5);
	\draw[color ={gray},opacity = 0.3] (-4,-6.5)--(4,-6.5);
	\draw[color ={gray},opacity = 0.3] (-4,-7)--(4,-7);
	\draw[color ={gray},opacity = 0.3] (-4,-7.5)--(4,-7.5);
	\draw[color ={gray},opacity = 0.3] (-4,-8)--(4,-8);
	\draw[color ={gray},opacity = 0.3] (0,-8)--(0,6);
	\draw[color ={gray},opacity = 0.3] (0,-8)--(0,6);
	\draw[color ={black},line width = 1.2] (1,-0.13)--(1,0.13);
	\draw[color ={black},line width = 1.2] (2,-0.13)--(2,0.13);
	\draw[color ={black},line width = 1.2] (3,-0.13)--(3,0.13);
	\draw[color ={black},line width = 1.2] (4,-0.13)--(4,0.13);
	\draw[color ={black},line width = 1.2] (-1,-0.13)--(-1,0.13);
	\draw[color ={black},line width = 1.2] (-2,-0.13)--(-2,0.13);
	\draw[color ={black},line width = 1.2] (-3,-0.13)--(-3,0.13);
	\draw[color ={black},line width = 1.2] (-4,-0.13)--(-4,0.13);
	\draw[color ={black},line width = 1.2] (-0.13,1)--(0.13,1);
	\draw[color ={black},line width = 1.2] (-0.13,2)--(0.13,2);
	\draw[color ={black},line width = 1.2] (-0.13,3)--(0.13,3);
	\draw[color ={black},line width = 1.2] (-0.13,4)--(0.13,4);
	\draw[color ={black},line width = 1.2] (-0.13,5)--(0.13,5);
	\draw[color ={black},line width = 1.2] (-0.13,6)--(0.13,6);
	\draw[color ={black},line width = 1.2] (-0.13,-1)--(0.13,-1);
	\draw[color ={black},line width = 1.2] (-0.13,-2)--(0.13,-2);
	\draw[color ={black},line width = 1.2] (-0.13,-3)--(0.13,-3);
	\draw[color ={black},line width = 1.2] (-0.13,-4)--(0.13,-4);
	\draw[color ={black},line width = 1.2] (-0.13,-5)--(0.13,-5);
	\draw[color ={black},line width = 1.2] (-0.13,-6)--(0.13,-6);
	\draw[color ={black},line width = 1.2] (-0.13,-7)--(0.13,-7);
	\draw[color ={black},line width = 1.2] (-0.13,-8)--(0.13,-8);
	\draw (2,-0.4) node[anchor = center] {\scriptsize \color{black}{$2$}};
	\draw (4,-0.4) node[anchor = center] {\scriptsize \color{black}{$4$}};
	\draw (-2,-0.4) node[anchor = center] {\scriptsize \color{black}{$-2$}};
	\draw (-4,-0.4) node[anchor = center] {\scriptsize \color{black}{$-4$}};
	\draw (-0.5,2.1) node[anchor = center] {\footnotesize \color{black}{$2$}};
	\draw (-0.5,4.1) node[anchor = center] {\footnotesize \color{black}{$4$}};
	\draw (-0.5,-1.9) node[anchor = center] {\footnotesize \color{black}{$-2$}};
	\draw (-0.5,-3.9) node[anchor = center] {\footnotesize \color{black}{$-4$}};
	\draw (-0.5,-5.9) node[anchor = center] {\footnotesize \color{black}{$-6$}};
	\draw (-0.5,-7.9) node[anchor = center] {\footnotesize \color{black}{$-8$}};
	\draw [color={black}] (-0.3,-0.3) node[anchor = center,scale=1, rotate = 0] {O};
	\draw (3,4) node[anchor = center] {\footnotesize \color{red}{$\Large\mathcal{C}_f\prime$}};
	
	\draw[color={red},line width = 2] (-4,4)--(-3.8,3.6)--(-3.6,3.2)--(-3.4,2.8)--(-3.2,2.4)--(-3,2)--(-2.8,1.6)--(-2.6,1.2)--(-2.4,0.8)--(-2.2,0.4)--(-2,0)--(-1.8,-0.4)--(-1.6,-0.8)--(-1.4,-1.2)--(-1.2,-1.6)--(-1,-2)--(-0.8,-2.4)--(-0.6,-2.8)--(-0.4,-3.2)--(-0.2,-3.6)--(0,-4)--(0.2,-4.4)--(0.4,-4.8)--(0.6,-5.2)--(0.8,-5.6)--(1,-6)--(1.2,-6.4)--(1.4,-6.8)--(1.6,-7.2)--(1.8,-7.6)--(2,-8)--(2.2,-8.4)--(2.4,-8.8);

\end{tikzpicture}\\
    \end{minipage}
    
\end{exercice}

\begin{Solution}
 L'équation réduite de la tangente au point d'abscisse $0$ est  : $y=f'(0)(x-0)+f(0)$.\\
      On lit graphiquement $f(0)=-4$ et $f'(0)=-4$.\\
      L'équation réduite de la tangente est donc donnée par :
      $y=-4(x+0)-4$, soit $y=-4x-4$.
\end{Solution}

% @see : https://coopmaths.fr/alea?uuid=a1ba2&id=can1F14&n=1&d=10&alea=rkIp223232323234234567891011121314151617181920212223242526272829303132333435363738394041424344454647482345678910111213&cd=1&cols=1
\needspace{10\baselineskip}
\begin{exercice}[BaremeTotal=false]


Soit $f$ la fonction définie sur $\mathbb{R}$ par : 
$f(x)= -3x^2-9$.\\
En calculant la limite du taux d'accroissement, déterminer $f'(1)$.
\end{exercice}

\begin{Solution}
 $f$ est une fonction polynôme du second degré de la forme $f(x)=ax^2+b$.\\
    La fonction dérivée est donnée par la somme des dérivées des fonctions $u$ et $v$ définies par $u(x)=-3x^2$ et $v(x)=-9$.\\
     Comme $u'(x)=-6x$ et $v'(x)=0$, on obtient  $f'(x)=-6x$. \\
     Ainsi, $f'(1)=-6\times 1=-6$.
\end{Solution}

% @see : https://coopmaths.fr/alea?uuid=3c690&id=can1F13&n=1&d=10&alea=EUDu223232323234234567891011121314151617181920212223242526272829303132333435363738394041424344454647482345678910111213&cd=1&cols=1
\needspace{10\baselineskip}
\begin{exercice}[BaremeTotal=false]


 Déterminer le coefficient directeur de la tangente à la courbe représentative de la fonction carré au point d'abscisse $8$.    
\end{exercice}

\begin{Solution}
 Le coefficient directeur de la tangente au point d'abscisse $a$ est donné par le nombre dérivé $f'(a)$.\\
        La fonction $f$ définie par $f(x)=x^2$ a pour fonction dérivée la fonction $f'$ définie par $f'(x)=2x$.\\
        Comme $f'(8)=2\times 8=16$, le coefficient directeur de la tangente au point d'abscisse $8$ est : $16$. 
\end{Solution}

% @see : https://coopmaths.fr/alea?uuid=ebd89&id=1AN14-1&n=1&d=10&s=3&alea=ocYr223232323234234567891011121314151617181920212223242526272829303132333435363738394041424344454647482345678910111213&cd=0&cols=1
\needspace{10\baselineskip}
\begin{exercice}[BaremeTotal=false]


 Donner la dérivée de la fonction $f$, dérivable pour tout $x\in\R$, définie par  $f(x)=\dfrac{2x+5}{3}$.
\end{exercice}

\begin{Solution}
 L'expression de la dérivée de la fonction $f$ définie par $f(x)=\dfrac{2x+5}{3}$ est : ${\color[HTML]{f15929}\boldsymbol{f'(x)=\dfrac{2}{3}}}$.
\end{Solution}

% @see : https://coopmaths.fr/alea?uuid=ebd89&id=1AN14-1&n=1&d=10&s=4&alea=uAlP223232323234234567891011121314151617181920212223242526272829303132333435363738394041424344454647482345678910111213&cd=0&cols=1
\needspace{10\baselineskip}
\begin{exercice}[BaremeTotal=false]


 Donner la dérivée de la fonction $f$, dérivable pour tout $x\in\R$, définie par  $f(x)=4x^{9}$.
\end{exercice}

\begin{Solution}
 L'expression de la dérivée de la fonction $f$ définie par $f(x)=4x^{9}$ est : ${\color[HTML]{f15929}\boldsymbol{f'(x)=36x^{8}}}$.
\end{Solution}

% @see : https://coopmaths.fr/alea?uuid=ebd89&id=1AN14-1&n=1&d=10&s=7&alea=vjN7223232323234234567891011121314151617181920212223242526272829303132333435363738394041424344454647482345678910111213&cd=0&cols=1
\needspace{10\baselineskip}
\begin{exercice}[BaremeTotal=false]


 Donner la dérivée de la fonction $f$, dérivable pour tout $x\in\R^*$, définie par  $f(x)=\dfrac{-1}{7x}$.
\end{exercice}

\begin{Solution}
 L'expression de la dérivée de la fonction $f$ définie par $f(x)=\dfrac{-1}{7x}$ est : ${\color[HTML]{f15929}\boldsymbol{f'(x)=\dfrac{1}{7x^2}}}$.
\end{Solution}

% @see : https://coopmaths.fr/alea?uuid=a83c0&id=1AN14-4&n=1&d=10&s=1&alea=Njr3223232323234234567891011121314151617181920212223242526272829303132333435363738394041424344454647482345678910111213&cd=0&cols=1
\needspace{10\baselineskip}
\begin{exercice}[BaremeTotal=false]


 Donner l'expression de la dérivée de la fonction $f$ définie sur $\R^{*}$ par $f(x)=-4x^{2}-4x-\dfrac{6}{x}$.\\
\end{exercice}

\begin{Solution}
 L'expression de la dérivée de la fonction $f$ définie par $f(x)=-4x^{2}-4x-\dfrac{6}{x}$ est :\\$f'(x)={\color[HTML]{f15929}\boldsymbol{-8x-4+\dfrac{6}{x^2}}}$.\\
\end{Solution}

% @see : https://coopmaths.fr/alea?uuid=a83c0&id=1AN14-4&n=1&d=10&s=4&alea=4Vpz223232323234234567891011121314151617181920212223242526272829303132333435363738394041424344454647482345678910111213&cd=1&cols=1
\needspace{10\baselineskip}
\begin{exercice}[BaremeTotal=false]


 Donner l'expression de la dérivée de la fonction $f$ définie sur $\R_+^{*}$ par $f(x)=-\dfrac{3}{x}+3x+5\sqrt{x}$.\\
\end{exercice}

\begin{Solution}
 $(-\dfrac{3}{x})^\prime=\dfrac{3}{x^2}$\\$(3x)^\prime=3$\\$(5\sqrt{x})^\prime=\dfrac{5}{2\sqrt{x}}$\\L'expression de la dérivée de la fonction $f$ définie par $f(x)=-\dfrac{3}{x}+3x+5\sqrt{x}$ est :\\$f'(x)={\color[HTML]{f15929}\boldsymbol{\dfrac{3}{x^2}+3+\dfrac{5}{2\sqrt{x}}}}$.\\
\end{Solution}

\end{Maquette}
\clearpage
\begin{Maquette}[DM]{Niveau= ,Classe=\getdata[14]\mydata\ (v14), Date = 06/01/25  ,Theme=Exercices}


% @see : https://coopmaths.fr/alea?uuid=4c8c7&id=1AN11-3&n=1&d=10&s=2&alea=nawO22323232323423456789101112131415161718192021222324252627282930313233343536373839404142434445464748234567891011121314&cd=1&cols=2
\needspace{10\baselineskip}
\begin{exercice}[BaremeTotal=false]
\label{eleve:14}


 \begin{minipage}[t]{\linewidth} Soit $f$ une fonction dérivable sur $[-5;5]$ et $\mathcal{C}_f$ sa courbe représentative.\\On sait que  $f(-4)=-5~~$ et que $~~f'(-4)=2$.\\Déterminer une équation de la tangente $(T)$ à la courbe $\mathcal{C}_f$ au point d'abscisse $-4$,\\sans utiliser la formule de cours de l'équation de tangente. \end{minipage}

\end{exercice}

\begin{Solution}
 \begin{minipage}[t]{\linewidth} On sait que la tangente n'est pas une droite verticale, puisque la fonction est dérivable sur l'intervalle.\\On en déduit que la tangente $(T)$ au point d'abscisse $-4$, admet une équation réduite de la forme :  $(T) : y=m x + p$. 

\medskip
$\bullet$ Détermination de $m$ :\\On sait que le nombre dérivé en $-4$ est par définition, le coefficient directeur de la tangente au point d'abscisse $-4$.\\Par conséquent, on a déjà : $m=f'(-4)=2$.\\ On en déduit que  $(T) : y= 2 x + p$\\$\bullet$ Détermination de $p$ :\\ Pour cela, on utilise que si $f(-4)=-5~~$, alors le point $A$ de coordonnées $(-4;-5)$ appartient à $\mathcal{C}_f$ mais aussi à $(T)$.\\ On peut écrire $A(-4;-5) = \mathcal{C}_f \cap (T)$.\\ On remplace alors les coordonnées de $A(-4;-5)$ dans l'équation  $(T) : y= 2 x + p$.\\ $\begin{aligned} \phantom{\iff}&A(-4;-5)\in (T)\\ \iff& -5= 2 \times (-4)  + p\\ \iff& p=-5 -2 \times (-4)  \\ \iff& p=-5 +8   \\ \iff& p=3\\\end{aligned}$\\On peut conclure que : $(T) : {\color[HTML]{f15929}\boldsymbol{y=2x+3}}$. \end{minipage}

\end{Solution}

% @see : https://coopmaths.fr/alea?uuid=0e984&id=can1F15&n=1&d=10&alea=1Alo22323232323423456789101112131415161718192021222324252627282930313233343536373839404142434445464748234567891011121314&cd=1&cols=1
\needspace{10\baselineskip}
\begin{exercice}[BaremeTotal=false]


 La courbe représente une fonction $f$ et la droite est la tangente au point d'abscisse $0$.\\
        Déterminer $f'(0)$.\begin{center}\begin{tikzpicture}[baseline,scale = 0.5]

    \tikzset{
      point/.style={
        thick,
        draw,
        cross out,
        inner sep=0pt,
        minimum width=5pt,
        minimum height=5pt,
      },
    }
    \clip (-10,-2) rectangle (4,12);
    	
	\draw[color ={black},line width = 2,>=latex,->] (-10,0)--(4,0);
	\draw[color ={black},line width = 2,>=latex,->] (0,-2)--(0,12);
	\draw[color ={black},opacity = 0.5] (-10,1)--(4,1);
	\draw[color ={black},opacity = 0.5] (-10,2)--(4,2);
	\draw[color ={black},opacity = 0.5] (-10,3)--(4,3);
	\draw[color ={black},opacity = 0.5] (-10,4)--(4,4);
	\draw[color ={black},opacity = 0.5] (-10,5)--(4,5);
	\draw[color ={black},opacity = 0.5] (-10,6)--(4,6);
	\draw[color ={black},opacity = 0.5] (-10,7)--(4,7);
	\draw[color ={black},opacity = 0.5] (-10,8)--(4,8);
	\draw[color ={black},opacity = 0.5] (-10,9)--(4,9);
	\draw[color ={black},opacity = 0.5] (-10,10)--(4,10);
	\draw[color ={black},opacity = 0.5] (-10,11)--(4,11);
	\draw[color ={black},opacity = 0.5] (-10,12)--(4,12);
	\draw[color ={black},opacity = 0.5] (-10,-1)--(4,-1);
	\draw[color ={black},opacity = 0.5] (-10,-2)--(4,-2);
	\draw[color ={black},opacity = 0.5] (1,-2)--(1,12);
	\draw[color ={black},opacity = 0.5] (2,-2)--(2,12);
	\draw[color ={black},opacity = 0.5] (3,-2)--(3,12);
	\draw[color ={black},opacity = 0.5] (4,-2)--(4,12);
	\draw[color ={black},opacity = 0.5] (-1,-2)--(-1,12);
	\draw[color ={black},opacity = 0.5] (-2,-2)--(-2,12);
	\draw[color ={black},opacity = 0.5] (-3,-2)--(-3,12);
	\draw[color ={black},opacity = 0.5] (-4,-2)--(-4,12);
	\draw[color ={black},opacity = 0.5] (-5,-2)--(-5,12);
	\draw[color ={black},opacity = 0.5] (-6,-2)--(-6,12);
	\draw[color ={black},opacity = 0.5] (-7,-2)--(-7,12);
	\draw[color ={black},opacity = 0.5] (-8,-2)--(-8,12);
	\draw[color ={black},opacity = 0.5] (-9,-2)--(-9,12);
	\draw[color ={black},opacity = 0.5] (-10,-2)--(-10,12);
	\draw[color ={gray},opacity = 0.3] (-10,0)--(4,0);
	\draw[color ={gray},opacity = 0.3] (-10,0)--(4,0);
	\draw[color ={gray},opacity = 0.3] (0,-2)--(0,12);
	\draw[color ={gray},opacity = 0.3] (0,-2)--(0,12);
	\draw[color ={gray},opacity = 0.3] (-5,-2)--(-5,12);
	\draw[color ={gray},opacity = 0.3] (-6,-2)--(-6,12);
	\draw[color ={gray},opacity = 0.3] (-7,-2)--(-7,12);
	\draw[color ={gray},opacity = 0.3] (-8,-2)--(-8,12);
	\draw[color ={gray},opacity = 0.3] (-9,-2)--(-9,12);
	\draw[color ={gray},opacity = 0.3] (-10,-2)--(-10,12);
	\draw[color ={black},line width = 2] (2,-0.2)--(2,0.2);
	\draw[color ={black},line width = 2] (-2,-0.2)--(-2,0.2);
	\draw[color ={black},line width = 2] (-4,-0.2)--(-4,0.2);
	\draw[color ={black},line width = 2] (-6,-0.2)--(-6,0.2);
	\draw[color ={black},line width = 2] (-8,-0.2)--(-8,0.2);
	\draw[color ={black},line width = 2] (-0.2,2)--(0.2,2);
	\draw[color ={black},line width = 2] (-0.2,4)--(0.2,4);
	\draw[color ={black},line width = 2] (-0.2,6)--(0.2,6);
	\draw[color ={black},line width = 2] (-0.2,8)--(0.2,8);
	\draw[color ={black},line width = 2] (-0.2,10)--(0.2,10);
	\draw (2,-0.4) node[anchor = center] {\scriptsize \color{black}{$1$}};
	\draw (-2,-0.4) node[anchor = center] {\scriptsize \color{black}{$-1$}};
	\draw (-4,-0.4) node[anchor = center] {\scriptsize \color{black}{$-2$}};
	\draw (-6,-0.4) node[anchor = center] {\scriptsize \color{black}{$-3$}};
	\draw (-8,-0.4) node[anchor = center] {\scriptsize \color{black}{$-4$}};
	\draw (-0.5,2.1) node[anchor = center] {\footnotesize \color{black}{$1$}};
	\draw (-0.5,4.1) node[anchor = center] {\footnotesize \color{black}{$2$}};
	\draw (-0.5,6.1) node[anchor = center] {\footnotesize \color{black}{$3$}};
	\draw (-0.5,8.1) node[anchor = center] {\footnotesize \color{black}{$4$}};
	\draw (-0.5,10.1) node[anchor = center] {\footnotesize \color{black}{$5$}};
	\draw [color={black}] (-0.3,-0.3) node[anchor = center,scale=1, rotate = 0] {O};
	
	\draw[color={blue},line width = 2] (-2.8,11.51)--(-2.4,8.96)--(-2,6.98)--(-1.6,5.44)--(-1.2,4.23)--(-0.8,3.3)--(-0.4,2.57)--(0,2)--(0.4,1.56)--(0.8,1.21)--(1.2,0.94)--(1.6,0.74)--(2,0.57)--(2.4,0.45)--(2.8,0.35)--(3.2,0.27)--(3.6,0.21)--(4,0.16);
	
	\draw[color={red},line width = 2] (-9.2,13.5)--(-8.8,13)--(-8.4,12.5)--(-8,12)--(-7.6,11.5)--(-7.2,11)--(-6.8,10.5)--(-6.4,10)--(-6,9.5)--(-5.6,9)--(-5.2,8.5)--(-4.8,8)--(-4.4,7.5)--(-4,7)--(-3.6,6.5)--(-3.2,6)--(-2.8,5.5)--(-2.4,5)--(-2,4.5)--(-1.6,4)--(-1.2,3.5)--(-0.8,3)--(-0.4,2.5)--(0,2)--(0.4,1.5)--(0.8,1)--(1.2,0.5)--(1.6,0)--(2,-0.5)--(2.4,-1)--(2.8,-1.5)--(3.2,-2)--(3.6,-2.5)--(4,-3);

\end{tikzpicture}\end{center}
\end{exercice}

\begin{Solution}
 $f'(0)$ est donné par le coefficient directeur de la tangente à la courbe au point d'abscisse $0$, soit $\dfrac{-5}{4}=-\dfrac{5}{4}$.
\end{Solution}

% @see : https://coopmaths.fr/alea?uuid=6f32d&id=can1F16&n=1&d=10&alea=dqSw22323232323423456789101112131415161718192021222324252627282930313233343536373839404142434445464748234567891011121314&cd=1&cols=1
\needspace{10\baselineskip}
\newpage
\begin{exercice}[BaremeTotal=false]


 On donne les représentations graphiques d'une fonction et de sa dérivée.\\
        Donner l'équation réduite de la tangente à la courbe de $f$ en $x=0$. \\ \begin{minipage}{0.5\linewidth}
    \begin{tikzpicture}[baseline, scale=0.8]

    \tikzset{
      point/.style={
        thick,
        draw,
        cross out,
        inner sep=0pt,
        minimum width=5pt,
        minimum height=5pt,
      },
    }
    \clip (-4.5,-4.5) rectangle (5.75,13.5);
    	
	\draw[color ={black},>=latex,->] (-3,0)--(3,0);
	\draw[color ={black},>=latex,->] (0,-3)--(0,12);
	\draw[color ={black},opacity = 0.4] (-3,1)--(3,1);
	\draw[color ={black},opacity = 0.4] (-3,2)--(3,2);
	\draw[color ={black},opacity = 0.4] (-3,3)--(3,3);
	\draw[color ={black},opacity = 0.4] (-3,4)--(3,4);
	\draw[color ={black},opacity = 0.4] (-3,5)--(3,5);
	\draw[color ={black},opacity = 0.4] (-3,6)--(3,6);
	\draw[color ={black},opacity = 0.4] (-3,7)--(3,7);
	\draw[color ={black},opacity = 0.4] (-3,8)--(3,8);
	\draw[color ={black},opacity = 0.4] (-3,9)--(3,9);
	\draw[color ={black},opacity = 0.4] (-3,10)--(3,10);
	\draw[color ={black},opacity = 0.4] (-3,11)--(3,11);
	\draw[color ={black},opacity = 0.4] (-3,12)--(3,12);
	\draw[color ={black},opacity = 0.4] (-3,-1)--(3,-1);
	\draw[color ={black},opacity = 0.4] (-3,-2)--(3,-2);
	\draw[color ={black},opacity = 0.4] (-3,-3)--(3,-3);
	\draw[color ={black},opacity = 0.4] (1,-3)--(1,12);
	\draw[color ={black},opacity = 0.4] (2,-3)--(2,12);
	\draw[color ={black},opacity = 0.4] (3,-3)--(3,12);
	\draw[color ={black},opacity = 0.4] (-1,-3)--(-1,12);
	\draw[color ={black},opacity = 0.4] (-2,-3)--(-2,12);
	\draw[color ={black},opacity = 0.4] (-3,-3)--(-3,12);
	\draw[color ={gray},opacity = 0.3] (-3,0)--(3,0);
	\draw[color ={gray},opacity = 0.3] (-3,0.5)--(3,0.5);
	\draw[color ={gray},opacity = 0.3] (-3,1.5)--(3,1.5);
	\draw[color ={gray},opacity = 0.3] (-3,2.5)--(3,2.5);
	\draw[color ={gray},opacity = 0.3] (-3,3.5)--(3,3.5);
	\draw[color ={gray},opacity = 0.3] (-3,4.5)--(3,4.5);
	\draw[color ={gray},opacity = 0.3] (-3,5.5)--(3,5.5);
	\draw[color ={gray},opacity = 0.3] (-3,6.5)--(3,6.5);
	\draw[color ={gray},opacity = 0.3] (-3,7.5)--(3,7.5);
	\draw[color ={gray},opacity = 0.3] (-3,8.5)--(3,8.5);
	\draw[color ={gray},opacity = 0.3] (-3,9.5)--(3,9.5);
	\draw[color ={gray},opacity = 0.3] (-3,10.5)--(3,10.5);
	\draw[color ={gray},opacity = 0.3] (-3,11.5)--(3,11.5);
	\draw[color ={gray},opacity = 0.3] (-3,0)--(3,0);
	\draw[color ={gray},opacity = 0.3] (-3,-0.5)--(3,-0.5);
	\draw[color ={gray},opacity = 0.3] (-3,-1.5)--(3,-1.5);
	\draw[color ={gray},opacity = 0.3] (-3,-2.5)--(3,-2.5);
	\draw[color ={gray},opacity = 0.3] (0,-3)--(0,12);
	\draw[color ={gray},opacity = 0.3] (0,-3)--(0,12);
	\draw[color ={black},line width = 1.2] (1,-0.13)--(1,0.13);
	\draw[color ={black},line width = 1.2] (2,-0.13)--(2,0.13);
	\draw[color ={black},line width = 1.2] (3,-0.13)--(3,0.13);
	\draw[color ={black},line width = 1.2] (-1,-0.13)--(-1,0.13);
	\draw[color ={black},line width = 1.2] (-2,-0.13)--(-2,0.13);
	\draw[color ={black},line width = 1.2] (-3,-0.13)--(-3,0.13);
	\draw[color ={black},line width = 1.2] (-0.13,1)--(0.13,1);
	\draw[color ={black},line width = 1.2] (-0.13,2)--(0.13,2);
	\draw[color ={black},line width = 1.2] (-0.13,3)--(0.13,3);
	\draw[color ={black},line width = 1.2] (-0.13,4)--(0.13,4);
	\draw[color ={black},line width = 1.2] (-0.13,5)--(0.13,5);
	\draw[color ={black},line width = 1.2] (-0.13,6)--(0.13,6);
	\draw[color ={black},line width = 1.2] (-0.13,7)--(0.13,7);
	\draw[color ={black},line width = 1.2] (-0.13,8)--(0.13,8);
	\draw[color ={black},line width = 1.2] (-0.13,9)--(0.13,9);
	\draw[color ={black},line width = 1.2] (-0.13,10)--(0.13,10);
	\draw[color ={black},line width = 1.2] (-0.13,11)--(0.13,11);
	\draw[color ={black},line width = 1.2] (-0.13,12)--(0.13,12);
	\draw[color ={black},line width = 1.2] (-0.13,-1)--(0.13,-1);
	\draw[color ={black},line width = 1.2] (-0.13,-2)--(0.13,-2);
	\draw[color ={black},line width = 1.2] (-0.13,-3)--(0.13,-3);
	\draw (3,-0.4) node[anchor = center] {\scriptsize \color{black}{$3$}};
	\draw (-3,-0.4) node[anchor = center] {\scriptsize \color{black}{$-3$}};
	\draw (-0.5,3.1) node[anchor = center] {\footnotesize \color{black}{$3$}};
	\draw (-0.5,6.1) node[anchor = center] {\footnotesize \color{black}{$6$}};
	\draw (-0.5,9.1) node[anchor = center] {\footnotesize \color{black}{$9$}};
	\draw (-0.5,-2.9) node[anchor = center] {\footnotesize \color{black}{$-3$}};
	\draw [color={black}] (-0.3,-0.3) node[anchor = center,scale=1, rotate = 0] {O};
	\draw (3,10) node[anchor = center] {\footnotesize \color{blue}{$\Large \mathcal{C}_f$}};
	
	\draw[color={blue},line width = 2] (-1.8,12.44)--(-1.6,10.96)--(-1.4,9.56)--(-1.2,8.24)--(-1,7)--(-0.8,5.84)--(-0.6,4.76)--(-0.4,3.76)--(-0.2,2.84)--(0,2)--(0.2,1.24)--(0.4,0.56)--(0.6,-0.04)--(0.8,-0.56)--(1,-1)--(1.2,-1.36)--(1.4,-1.64)--(1.6,-1.84)--(1.8,-1.96)--(2,-2)--(2.2,-1.96)--(2.4,-1.84)--(2.6,-1.64)--(2.8,-1.36)--(3,-1);

\end{tikzpicture}\\
    \end{minipage}
    \begin{minipage}{0.5\linewidth}
    \begin{tikzpicture}[baseline, scale=0.8]

    \tikzset{
      point/.style={
        thick,
        draw,
        cross out,
        inner sep=0pt,
        minimum width=5pt,
        minimum height=5pt,
      },
    }
    \clip (-4.5,-6.5) rectangle (6.5,9.5);
    	
	\draw[color ={black},>=latex,->] (-3,0)--(3,0);
	\draw[color ={black},>=latex,->] (0,-5)--(0,8);
	\draw[color ={black},opacity = 0.4] (-3,1)--(3,1);
	\draw[color ={black},opacity = 0.4] (-3,2)--(3,2);
	\draw[color ={black},opacity = 0.4] (-3,3)--(3,3);
	\draw[color ={black},opacity = 0.4] (-3,4)--(3,4);
	\draw[color ={black},opacity = 0.4] (-3,5)--(3,5);
	\draw[color ={black},opacity = 0.4] (-3,6)--(3,6);
	\draw[color ={black},opacity = 0.4] (-3,7)--(3,7);
	\draw[color ={black},opacity = 0.4] (-3,8)--(3,8);
	\draw[color ={black},opacity = 0.4] (-3,-1)--(3,-1);
	\draw[color ={black},opacity = 0.4] (-3,-2)--(3,-2);
	\draw[color ={black},opacity = 0.4] (-3,-3)--(3,-3);
	\draw[color ={black},opacity = 0.4] (-3,-4)--(3,-4);
	\draw[color ={black},opacity = 0.4] (-3,-5)--(3,-5);
	\draw[color ={black},opacity = 0.4] (1,-5)--(1,8);
	\draw[color ={black},opacity = 0.4] (2,-5)--(2,8);
	\draw[color ={black},opacity = 0.4] (3,-5)--(3,8);
	\draw[color ={black},opacity = 0.4] (-1,-5)--(-1,8);
	\draw[color ={black},opacity = 0.4] (-2,-5)--(-2,8);
	\draw[color ={black},opacity = 0.4] (-3,-5)--(-3,8);
	\draw[color ={gray},opacity = 0.3] (-3,0)--(3,0);
	\draw[color ={gray},opacity = 0.3] (-3,0.5)--(3,0.5);
	\draw[color ={gray},opacity = 0.3] (-3,1.5)--(3,1.5);
	\draw[color ={gray},opacity = 0.3] (-3,2.5)--(3,2.5);
	\draw[color ={gray},opacity = 0.3] (-3,3.5)--(3,3.5);
	\draw[color ={gray},opacity = 0.3] (-3,4.5)--(3,4.5);
	\draw[color ={gray},opacity = 0.3] (-3,5.5)--(3,5.5);
	\draw[color ={gray},opacity = 0.3] (-3,6.5)--(3,6.5);
	\draw[color ={gray},opacity = 0.3] (-3,7.5)--(3,7.5);
	\draw[color ={gray},opacity = 0.3] (-3,0)--(3,0);
	\draw[color ={gray},opacity = 0.3] (-3,-0.5)--(3,-0.5);
	\draw[color ={gray},opacity = 0.3] (-3,-1.5)--(3,-1.5);
	\draw[color ={gray},opacity = 0.3] (-3,-2.5)--(3,-2.5);
	\draw[color ={gray},opacity = 0.3] (-3,-3.5)--(3,-3.5);
	\draw[color ={gray},opacity = 0.3] (-3,-4.5)--(3,-4.5);
	\draw[color ={gray},opacity = 0.3] (0,-5)--(0,8);
	\draw[color ={gray},opacity = 0.3] (0,-5)--(0,8);
	\draw[color ={black},line width = 1.2] (1,-0.13)--(1,0.13);
	\draw[color ={black},line width = 1.2] (2,-0.13)--(2,0.13);
	\draw[color ={black},line width = 1.2] (3,-0.13)--(3,0.13);
	\draw[color ={black},line width = 1.2] (-1,-0.13)--(-1,0.13);
	\draw[color ={black},line width = 1.2] (-2,-0.13)--(-2,0.13);
	\draw[color ={black},line width = 1.2] (-3,-0.13)--(-3,0.13);
	\draw[color ={black},line width = 1.2] (-0.13,1)--(0.13,1);
	\draw[color ={black},line width = 1.2] (-0.13,2)--(0.13,2);
	\draw[color ={black},line width = 1.2] (-0.13,3)--(0.13,3);
	\draw[color ={black},line width = 1.2] (-0.13,4)--(0.13,4);
	\draw[color ={black},line width = 1.2] (-0.13,5)--(0.13,5);
	\draw[color ={black},line width = 1.2] (-0.13,6)--(0.13,6);
	\draw[color ={black},line width = 1.2] (-0.13,7)--(0.13,7);
	\draw[color ={black},line width = 1.2] (-0.13,8)--(0.13,8);
	\draw[color ={black},line width = 1.2] (-0.13,9)--(0.13,9);
	\draw[color ={black},line width = 1.2] (-0.13,10)--(0.13,10);
	\draw[color ={black},line width = 1.2] (-0.13,11)--(0.13,11);
	\draw[color ={black},line width = 1.2] (-0.13,12)--(0.13,12);
	\draw[color ={black},line width = 1.2] (-0.13,-1)--(0.13,-1);
	\draw[color ={black},line width = 1.2] (-0.13,-2)--(0.13,-2);
	\draw[color ={black},line width = 1.2] (-0.13,-3)--(0.13,-3);
	\draw (3,-0.4) node[anchor = center] {\scriptsize \color{black}{$3$}};
	\draw (-3,-0.4) node[anchor = center] {\scriptsize \color{black}{$-3$}};
	\draw (-0.5,3.1) node[anchor = center] {\footnotesize \color{black}{$3$}};
	\draw (-0.5,6.1) node[anchor = center] {\footnotesize \color{black}{$6$}};
	\draw (-0.5,-2.9) node[anchor = center] {\footnotesize \color{black}{$-3$}};
	\draw [color={black}] (-0.3,-0.3) node[anchor = center,scale=1, rotate = 0] {O};
	\draw (3,6) node[anchor = center] {\footnotesize \color{red}{$\Large\mathcal{C}_f\prime$}};
	
	\draw[color={red},line width = 2] (-1,-6)--(-0.8,-5.6)--(-0.6,-5.2)--(-0.4,-4.8)--(-0.2,-4.4)--(0,-4)--(0.2,-3.6)--(0.4,-3.2)--(0.6,-2.8)--(0.8,-2.4)--(1,-2)--(1.2,-1.6)--(1.4,-1.2)--(1.6,-0.8)--(1.8,-0.4)--(2,0)--(2.2,0.4)--(2.4,0.8)--(2.6,1.2)--(2.8,1.6)--(3,2);

\end{tikzpicture}\\
    \end{minipage}
    
\end{exercice}

\begin{Solution}
 L'équation réduite de la tangente au point d'abscisse $0$ est  : $y=f'(0)(x-0)+f(0)$.\\
        On lit graphiquement $f(0)=2$ et $f'(0)=-4$.\\
        L'équation réduite de la tangente est donc donnée par :
        $y=-4(x+0)+2$, soit $y=-4x+2$.
\end{Solution}

% @see : https://coopmaths.fr/alea?uuid=a1ba2&id=can1F14&n=1&d=10&alea=rkIp22323232323423456789101112131415161718192021222324252627282930313233343536373839404142434445464748234567891011121314&cd=1&cols=1
\needspace{10\baselineskip}
\begin{exercice}[BaremeTotal=false]


Soit $f$ la fonction définie sur $\mathbb{R}$ par : 
$f(x)= -5x^2+5$.\\
En calculant la limite du taux d'accroissement, déterminer $f'(-1)$.
\end{exercice}

\begin{Solution}
 $f$ est une fonction polynôme du second degré de la forme $f(x)=ax^2+b$.\\
    La fonction dérivée est donnée par la somme des dérivées des fonctions $u$ et $v$ définies par $u(x)=-5x^2$ et $v(x)=5$.\\
     Comme $u'(x)=-10x$ et $v'(x)=0$, on obtient  $f'(x)=-10x$. \\
     Ainsi, $f'(-1)=-10\times (-1)=10$.
\end{Solution}

% @see : https://coopmaths.fr/alea?uuid=3c690&id=can1F13&n=1&d=10&alea=EUDu22323232323423456789101112131415161718192021222324252627282930313233343536373839404142434445464748234567891011121314&cd=1&cols=1
\needspace{10\baselineskip}
\begin{exercice}[BaremeTotal=false]


 Déterminer le coefficient directeur de la tangente à la courbe représentative de la fonction cube au point d'abscisse $-3$.
                        
\end{exercice}

\begin{Solution}
 Le coefficient directeur de la tangente au point d'abscisse $a$ est donné par le nombre dérivé $f'(a)$.\\
        La fonction $f$ définie par $f(x)=x^3$ a pour fonction dérivée la fonction $f'$ définie par $f'(x)=3x^2$.\\
        Comme $f'(-3)=3\times (-3)^2=27$, le coefficient directeur de la tangente au point d'abscisse $-3$ est : $27$. 
\end{Solution}

% @see : https://coopmaths.fr/alea?uuid=ebd89&id=1AN14-1&n=1&d=10&s=3&alea=ocYr22323232323423456789101112131415161718192021222324252627282930313233343536373839404142434445464748234567891011121314&cd=0&cols=1
\needspace{10\baselineskip}
\begin{exercice}[BaremeTotal=false]


 Donner la dérivée de la fonction $f$, dérivable pour tout $x\in\R$, définie par  $f(x)=-\dfrac{4}{7}x+4$.
\end{exercice}

\begin{Solution}
 L'expression de la dérivée de la fonction $f$ définie par $f(x)=-\dfrac{4}{7}x+4$ est : ${\color[HTML]{f15929}\boldsymbol{f'(x)=-\dfrac{4}{7}}}$.
\end{Solution}

% @see : https://coopmaths.fr/alea?uuid=ebd89&id=1AN14-1&n=1&d=10&s=4&alea=uAlP22323232323423456789101112131415161718192021222324252627282930313233343536373839404142434445464748234567891011121314&cd=0&cols=1
\needspace{10\baselineskip}
\begin{exercice}[BaremeTotal=false]


 Donner la dérivée de la fonction $f$, dérivable pour tout $x\in\R$, définie par  $f(x)=8x^{4}$.
\end{exercice}

\begin{Solution}
 L'expression de la dérivée de la fonction $f$ définie par $f(x)=8x^{4}$ est : ${\color[HTML]{f15929}\boldsymbol{f'(x)=32x^{3}}}$.
\end{Solution}

% @see : https://coopmaths.fr/alea?uuid=ebd89&id=1AN14-1&n=1&d=10&s=7&alea=vjN722323232323423456789101112131415161718192021222324252627282930313233343536373839404142434445464748234567891011121314&cd=0&cols=1
\needspace{10\baselineskip}
\begin{exercice}[BaremeTotal=false]


 Donner la dérivée de la fonction $f$, dérivable pour tout $x\in\R^*$, définie par  $f(x)=\dfrac{3}{8x}$.
\end{exercice}

\begin{Solution}
 L'expression de la dérivée de la fonction $f$ définie par $f(x)=\dfrac{3}{8x}$ est : ${\color[HTML]{f15929}\boldsymbol{f'(x)=-\dfrac{3}{8x^2}}}$.
\end{Solution}

% @see : https://coopmaths.fr/alea?uuid=a83c0&id=1AN14-4&n=1&d=10&s=1&alea=Njr322323232323423456789101112131415161718192021222324252627282930313233343536373839404142434445464748234567891011121314&cd=0&cols=1
\needspace{10\baselineskip}
\begin{exercice}[BaremeTotal=false]


 Donner l'expression de la dérivée de la fonction $f$ définie sur $\R^{*}$ par $f(x)=\dfrac{8}{x}-2x^{2}$.\\
\end{exercice}

\begin{Solution}
 L'expression de la dérivée de la fonction $f$ définie par $f(x)=\dfrac{8}{x}-2x^{2}$ est :\\$f'(x)={\color[HTML]{f15929}\boldsymbol{-\dfrac{8}{x^2}-4x}}$.\\
\end{Solution}

% @see : https://coopmaths.fr/alea?uuid=a83c0&id=1AN14-4&n=1&d=10&s=4&alea=4Vpz22323232323423456789101112131415161718192021222324252627282930313233343536373839404142434445464748234567891011121314&cd=1&cols=1
\needspace{10\baselineskip}
\begin{exercice}[BaremeTotal=false]


 Donner l'expression de la dérivée de la fonction $f$ définie sur $\R_+^{*}$ par $f(x)=\dfrac{5}{x}-3x^{2}+4\sqrt{x}$.\\
\end{exercice}

\begin{Solution}
 $(\dfrac{5}{x})^\prime=-\dfrac{5}{x^2}$\\$(-3x^{2})^\prime=-6x$\\$(4\sqrt{x})^\prime=\dfrac{4}{2\sqrt{x}}$\\L'expression de la dérivée de la fonction $f$ définie par $f(x)=\dfrac{5}{x}-3x^{2}+4\sqrt{x}$ est :\\$f'(x)={\color[HTML]{f15929}\boldsymbol{-\dfrac{5}{x^2}-6x+\dfrac{4}{2\sqrt{x}}}}$.\\
\end{Solution}

\end{Maquette}
\clearpage
\begin{Maquette}[DM]{Niveau= ,Classe=\getdata[15]\mydata\ (v15), Date = 06/01/25  ,Theme=Exercices}


% @see : https://coopmaths.fr/alea?uuid=4c8c7&id=1AN11-3&n=1&d=10&s=2&alea=nawO2232323232342345678910111213141516171819202122232425262728293031323334353637383940414243444546474823456789101112131415&cd=1&cols=2
\needspace{10\baselineskip}
\begin{exercice}[BaremeTotal=false]
\label{eleve:15}


 \begin{minipage}[t]{\linewidth} Soit $f$ une fonction dérivable sur $[-5;5]$ et $\mathcal{C}_f$ sa courbe représentative.\\On sait que  $f(-2)=5~~$ et que $~~f'(-2)=-5$.\\Déterminer une équation de la tangente $(T)$ à la courbe $\mathcal{C}_f$ au point d'abscisse $-2$,\\sans utiliser la formule de cours de l'équation de tangente. \end{minipage}

\end{exercice}

\begin{Solution}
 \begin{minipage}[t]{\linewidth} On sait que la tangente n'est pas une droite verticale, puisque la fonction est dérivable sur l'intervalle.\\On en déduit que la tangente $(T)$ au point d'abscisse $-2$, admet une équation réduite de la forme :  $(T) : y=m x + p$. 

\medskip
$\bullet$ Détermination de $m$ :\\On sait que le nombre dérivé en $-2$ est par définition, le coefficient directeur de la tangente au point d'abscisse $-2$.\\Par conséquent, on a déjà : $m=f'(-2)=-5$.\\ On en déduit que  $(T) : y= -5 x + p$\\$\bullet$ Détermination de $p$ :\\ Pour cela, on utilise que si $f(-2)=5~~$, alors le point $A$ de coordonnées $(-2;5)$ appartient à $\mathcal{C}_f$ mais aussi à $(T)$.\\ On peut écrire $A(-2;5) = \mathcal{C}_f \cap (T)$.\\ On remplace alors les coordonnées de $A(-2;5)$ dans l'équation  $(T) : y= -5 x + p$.\\ $\begin{aligned} \phantom{\iff}&A(-2;5)\in (T)\\ \iff& 5= -5 \times (-2)  + p\\ \iff& p=5 +5 \times (-2)  \\ \iff& p=5 -10   \\ \iff& p=-5\\\end{aligned}$\\On peut conclure que : $(T) : {\color[HTML]{f15929}\boldsymbol{y=-5x-5}}$. \end{minipage}

\end{Solution}

% @see : https://coopmaths.fr/alea?uuid=0e984&id=can1F15&n=1&d=10&alea=1Alo2232323232342345678910111213141516171819202122232425262728293031323334353637383940414243444546474823456789101112131415&cd=1&cols=1
\needspace{10\baselineskip}
\begin{exercice}[BaremeTotal=false]


 La courbe représente une fonction $f$ et la droite est la tangente au point d'abscisse $0$.\\
        Déterminer $f'(0)$.\begin{center}\begin{tikzpicture}[baseline,scale = 0.5]

    \tikzset{
      point/.style={
        thick,
        draw,
        cross out,
        inner sep=0pt,
        minimum width=5pt,
        minimum height=5pt,
      },
    }
    \clip (-4,-2) rectangle (12,12);
    	
	\draw[color ={black},line width = 2,>=latex,->] (-4,0)--(12,0);
	\draw[color ={black},line width = 2,>=latex,->] (0,-2)--(0,12);
	\draw[color ={black},opacity = 0.5] (-4,1)--(12,1);
	\draw[color ={black},opacity = 0.5] (-4,2)--(12,2);
	\draw[color ={black},opacity = 0.5] (-4,3)--(12,3);
	\draw[color ={black},opacity = 0.5] (-4,4)--(12,4);
	\draw[color ={black},opacity = 0.5] (-4,5)--(12,5);
	\draw[color ={black},opacity = 0.5] (-4,6)--(12,6);
	\draw[color ={black},opacity = 0.5] (-4,7)--(12,7);
	\draw[color ={black},opacity = 0.5] (-4,8)--(12,8);
	\draw[color ={black},opacity = 0.5] (-4,9)--(12,9);
	\draw[color ={black},opacity = 0.5] (-4,10)--(12,10);
	\draw[color ={black},opacity = 0.5] (-4,11)--(12,11);
	\draw[color ={black},opacity = 0.5] (-4,12)--(12,12);
	\draw[color ={black},opacity = 0.5] (-4,-1)--(12,-1);
	\draw[color ={black},opacity = 0.5] (-4,-2)--(12,-2);
	\draw[color ={black},opacity = 0.5] (1,-2)--(1,12);
	\draw[color ={black},opacity = 0.5] (2,-2)--(2,12);
	\draw[color ={black},opacity = 0.5] (3,-2)--(3,12);
	\draw[color ={black},opacity = 0.5] (4,-2)--(4,12);
	\draw[color ={black},opacity = 0.5] (5,-2)--(5,12);
	\draw[color ={black},opacity = 0.5] (6,-2)--(6,12);
	\draw[color ={black},opacity = 0.5] (7,-2)--(7,12);
	\draw[color ={black},opacity = 0.5] (8,-2)--(8,12);
	\draw[color ={black},opacity = 0.5] (9,-2)--(9,12);
	\draw[color ={black},opacity = 0.5] (10,-2)--(10,12);
	\draw[color ={black},opacity = 0.5] (11,-2)--(11,12);
	\draw[color ={black},opacity = 0.5] (12,-2)--(12,12);
	\draw[color ={black},opacity = 0.5] (-1,-2)--(-1,12);
	\draw[color ={black},opacity = 0.5] (-2,-2)--(-2,12);
	\draw[color ={black},opacity = 0.5] (-3,-2)--(-3,12);
	\draw[color ={black},opacity = 0.5] (-4,-2)--(-4,12);
	\draw[color ={gray},opacity = 0.3] (-4,0)--(12,0);
	\draw[color ={gray},opacity = 0.3] (-4,0)--(12,0);
	\draw[color ={gray},opacity = 0.3] (0,-2)--(0,12);
	\draw[color ={gray},opacity = 0.3] (0,-2)--(0,12);
	\draw[color ={black},line width = 2] (2,-0.2)--(2,0.2);
	\draw[color ={black},line width = 2] (4,-0.2)--(4,0.2);
	\draw[color ={black},line width = 2] (6,-0.2)--(6,0.2);
	\draw[color ={black},line width = 2] (8,-0.2)--(8,0.2);
	\draw[color ={black},line width = 2] (10,-0.2)--(10,0.2);
	\draw[color ={black},line width = 2] (-2,-0.2)--(-2,0.2);
	\draw[color ={black},line width = 2] (-0.2,2)--(0.2,2);
	\draw[color ={black},line width = 2] (-0.2,4)--(0.2,4);
	\draw[color ={black},line width = 2] (-0.2,6)--(0.2,6);
	\draw[color ={black},line width = 2] (-0.2,8)--(0.2,8);
	\draw[color ={black},line width = 2] (-0.2,10)--(0.2,10);
	\draw (2,-0.4) node[anchor = center] {\scriptsize \color{black}{$1$}};
	\draw (4,-0.4) node[anchor = center] {\scriptsize \color{black}{$2$}};
	\draw (6,-0.4) node[anchor = center] {\scriptsize \color{black}{$3$}};
	\draw (8,-0.4) node[anchor = center] {\scriptsize \color{black}{$4$}};
	\draw (10,-0.4) node[anchor = center] {\scriptsize \color{black}{$5$}};
	\draw (-0.5,2.1) node[anchor = center] {\footnotesize \color{black}{$1$}};
	\draw (-0.5,4.1) node[anchor = center] {\footnotesize \color{black}{$2$}};
	\draw (-0.5,6.1) node[anchor = center] {\footnotesize \color{black}{$3$}};
	\draw (-0.5,8.1) node[anchor = center] {\footnotesize \color{black}{$4$}};
	\draw (-0.5,10.1) node[anchor = center] {\footnotesize \color{black}{$5$}};
	\draw [color={black}] (-0.3,-0.3) node[anchor = center,scale=1, rotate = 0] {O};
	
	\draw[color={blue},line width = 2] (-4,0.04)--(-3.6,0.05)--(-3.2,0.08)--(-2.8,0.12)--(-2.4,0.18)--(-2,0.27)--(-1.6,0.4)--(-1.2,0.6)--(-0.8,0.9)--(-0.4,1.34)--(0,2)--(0.4,2.98)--(0.8,4.45)--(1.2,6.64)--(1.6,9.91);
	
	\draw[color={red},line width = 2] (-2.8,-3.6)--(-2.4,-2.8)--(-2,-2)--(-1.6,-1.2)--(-1.2,-0.4)--(-0.8,0.4)--(-0.4,1.2)--(0,2)--(0.4,2.8)--(0.8,3.6)--(1.2,4.4)--(1.6,5.2)--(2,6)--(2.4,6.8)--(2.8,7.6)--(3.2,8.4)--(3.6,9.2)--(4,10)--(4.4,10.8)--(4.8,11.6)--(5.2,12.4)--(5.6,13.2);

\end{tikzpicture}\end{center}
\end{exercice}

\begin{Solution}
 $f'(0)$ est donné par le coefficient directeur de la tangente à la courbe au point d'abscisse $0$, soit $2=2$.
\end{Solution}

% @see : https://coopmaths.fr/alea?uuid=6f32d&id=can1F16&n=1&d=10&alea=dqSw2232323232342345678910111213141516171819202122232425262728293031323334353637383940414243444546474823456789101112131415&cd=1&cols=1
\needspace{10\baselineskip}
\newpage
\begin{exercice}[BaremeTotal=false]


 On donne les représentations graphiques d'une fonction et de sa dérivée.\\
      Donner l'équation réduite de la tangente à la courbe de $f$ en $x=2$. \\ \begin{minipage}{0.5\linewidth}
    \begin{tikzpicture}[baseline, scale=0.8]

    \tikzset{
      point/.style={
        thick,
        draw,
        cross out,
        inner sep=0pt,
        minimum width=5pt,
        minimum height=5pt,
      },
    }
    \clip (-5.5,-9.5) rectangle (5.75,7.5);
    	
	\draw[color ={black},>=latex,->] (-4,0)--(4,0);
	\draw[color ={black},>=latex,->] (0,-8)--(0,6);
	\draw[color ={black},opacity = 0.4] (-4,1)--(4,1);
	\draw[color ={black},opacity = 0.4] (-4,2)--(4,2);
	\draw[color ={black},opacity = 0.4] (-4,3)--(4,3);
	\draw[color ={black},opacity = 0.4] (-4,4)--(4,4);
	\draw[color ={black},opacity = 0.4] (-4,5)--(4,5);
	\draw[color ={black},opacity = 0.4] (-4,6)--(4,6);
	\draw[color ={black},opacity = 0.4] (-4,-1)--(4,-1);
	\draw[color ={black},opacity = 0.4] (-4,-2)--(4,-2);
	\draw[color ={black},opacity = 0.4] (-4,-3)--(4,-3);
	\draw[color ={black},opacity = 0.4] (-4,-4)--(4,-4);
	\draw[color ={black},opacity = 0.4] (-4,-5)--(4,-5);
	\draw[color ={black},opacity = 0.4] (-4,-6)--(4,-6);
	\draw[color ={black},opacity = 0.4] (-4,-7)--(4,-7);
	\draw[color ={black},opacity = 0.4] (-4,-8)--(4,-8);
	\draw[color ={black},opacity = 0.4] (1,-8)--(1,6);
	\draw[color ={black},opacity = 0.4] (2,-8)--(2,6);
	\draw[color ={black},opacity = 0.4] (3,-8)--(3,6);
	\draw[color ={black},opacity = 0.4] (4,-8)--(4,6);
	\draw[color ={black},opacity = 0.4] (-1,-8)--(-1,6);
	\draw[color ={black},opacity = 0.4] (-2,-8)--(-2,6);
	\draw[color ={black},opacity = 0.4] (-3,-8)--(-3,6);
	\draw[color ={black},opacity = 0.4] (-4,-8)--(-4,6);
	\draw[color ={gray},opacity = 0.3] (-4,0)--(4,0);
	\draw[color ={gray},opacity = 0.3] (-4,0.5)--(4,0.5);
	\draw[color ={gray},opacity = 0.3] (-4,1.5)--(4,1.5);
	\draw[color ={gray},opacity = 0.3] (-4,2.5)--(4,2.5);
	\draw[color ={gray},opacity = 0.3] (-4,3.5)--(4,3.5);
	\draw[color ={gray},opacity = 0.3] (-4,4.5)--(4,4.5);
	\draw[color ={gray},opacity = 0.3] (-4,5.5)--(4,5.5);
	\draw[color ={gray},opacity = 0.3] (-4,0)--(4,0);
	\draw[color ={gray},opacity = 0.3] (-4,-0.5)--(4,-0.5);
	\draw[color ={gray},opacity = 0.3] (-4,-1.5)--(4,-1.5);
	\draw[color ={gray},opacity = 0.3] (-4,-2.5)--(4,-2.5);
	\draw[color ={gray},opacity = 0.3] (-4,-3.5)--(4,-3.5);
	\draw[color ={gray},opacity = 0.3] (-4,-4.5)--(4,-4.5);
	\draw[color ={gray},opacity = 0.3] (-4,-5.5)--(4,-5.5);
	\draw[color ={gray},opacity = 0.3] (-4,-6.5)--(4,-6.5);
	\draw[color ={gray},opacity = 0.3] (-4,-7)--(4,-7);
	\draw[color ={gray},opacity = 0.3] (-4,-7.5)--(4,-7.5);
	\draw[color ={gray},opacity = 0.3] (-4,-8)--(4,-8);
	\draw[color ={gray},opacity = 0.3] (0,-8)--(0,6);
	\draw[color ={gray},opacity = 0.3] (0,-8)--(0,6);
	\draw[color ={black},line width = 1.2] (1,-0.13)--(1,0.13);
	\draw[color ={black},line width = 1.2] (2,-0.13)--(2,0.13);
	\draw[color ={black},line width = 1.2] (3,-0.13)--(3,0.13);
	\draw[color ={black},line width = 1.2] (4,-0.13)--(4,0.13);
	\draw[color ={black},line width = 1.2] (-1,-0.13)--(-1,0.13);
	\draw[color ={black},line width = 1.2] (-2,-0.13)--(-2,0.13);
	\draw[color ={black},line width = 1.2] (-3,-0.13)--(-3,0.13);
	\draw[color ={black},line width = 1.2] (-4,-0.13)--(-4,0.13);
	\draw[color ={black},line width = 1.2] (-0.13,1)--(0.13,1);
	\draw[color ={black},line width = 1.2] (-0.13,2)--(0.13,2);
	\draw[color ={black},line width = 1.2] (-0.13,3)--(0.13,3);
	\draw[color ={black},line width = 1.2] (-0.13,4)--(0.13,4);
	\draw[color ={black},line width = 1.2] (-0.13,5)--(0.13,5);
	\draw[color ={black},line width = 1.2] (-0.13,6)--(0.13,6);
	\draw[color ={black},line width = 1.2] (-0.13,-1)--(0.13,-1);
	\draw[color ={black},line width = 1.2] (-0.13,-2)--(0.13,-2);
	\draw[color ={black},line width = 1.2] (-0.13,-3)--(0.13,-3);
	\draw[color ={black},line width = 1.2] (-0.13,-4)--(0.13,-4);
	\draw[color ={black},line width = 1.2] (-0.13,-5)--(0.13,-5);
	\draw[color ={black},line width = 1.2] (-0.13,-6)--(0.13,-6);
	\draw[color ={black},line width = 1.2] (-0.13,-7)--(0.13,-7);
	\draw[color ={black},line width = 1.2] (-0.13,-8)--(0.13,-8);
	\draw (2,-0.4) node[anchor = center] {\scriptsize \color{black}{$2$}};
	\draw (4,-0.4) node[anchor = center] {\scriptsize \color{black}{$4$}};
	\draw (-2,-0.4) node[anchor = center] {\scriptsize \color{black}{$-2$}};
	\draw (-4,-0.4) node[anchor = center] {\scriptsize \color{black}{$-4$}};
	\draw (-0.5,2.1) node[anchor = center] {\footnotesize \color{black}{$2$}};
	\draw (-0.5,4.1) node[anchor = center] {\footnotesize \color{black}{$4$}};
	\draw (-0.5,6.1) node[anchor = center] {\footnotesize \color{black}{$6$}};
	\draw (-0.5,-1.9) node[anchor = center] {\footnotesize \color{black}{$-2$}};
	\draw (-0.5,-3.9) node[anchor = center] {\footnotesize \color{black}{$-4$}};
	\draw (-0.5,-5.9) node[anchor = center] {\footnotesize \color{black}{$-6$}};
	\draw (-0.5,-7.9) node[anchor = center] {\footnotesize \color{black}{$-8$}};
	\draw [color={black}] (-0.3,-0.3) node[anchor = center,scale=1, rotate = 0] {O};
	\draw (3,4) node[anchor = center] {\footnotesize \color{blue}{$\Large \mathcal{C}_f$}};
	
	\draw[color={blue},line width = 2] (-3,-8)--(-2.8,-6.84)--(-2.6,-5.76)--(-2.4,-4.76)--(-2.2,-3.84)--(-2,-3)--(-1.8,-2.24)--(-1.6,-1.56)--(-1.4,-0.96)--(-1.2,-0.44)--(-1,0)--(-0.8,0.36)--(-0.6,0.64)--(-0.4,0.84)--(-0.2,0.96)--(0,1)--(0.2,0.96)--(0.4,0.84)--(0.6,0.64)--(0.8,0.36)--(1,0)--(1.2,-0.44)--(1.4,-0.96)--(1.6,-1.56)--(1.8,-2.24)--(2,-3)--(2.2,-3.84)--(2.4,-4.76)--(2.6,-5.76)--(2.8,-6.84)--(3,-8);

\end{tikzpicture}\\
    \end{minipage}
    \begin{minipage}{0.5\linewidth}
    \begin{tikzpicture}[baseline, scale=0.8]

    \tikzset{
      point/.style={
        thick,
        draw,
        cross out,
        inner sep=0pt,
        minimum width=5pt,
        minimum height=5pt,
      },
    }
    \clip (-5.5,-9.5) rectangle (6.5,7.5);
    	
	\draw[color ={black},>=latex,->] (-4,0)--(4,0);
	\draw[color ={black},>=latex,->] (0,-8)--(0,6);
	\draw[color ={black},opacity = 0.4] (-4,1)--(4,1);
	\draw[color ={black},opacity = 0.4] (-4,2)--(4,2);
	\draw[color ={black},opacity = 0.4] (-4,3)--(4,3);
	\draw[color ={black},opacity = 0.4] (-4,4)--(4,4);
	\draw[color ={black},opacity = 0.4] (-4,5)--(4,5);
	\draw[color ={black},opacity = 0.4] (-4,6)--(4,6);
	\draw[color ={black},opacity = 0.4] (-4,-1)--(4,-1);
	\draw[color ={black},opacity = 0.4] (-4,-2)--(4,-2);
	\draw[color ={black},opacity = 0.4] (-4,-3)--(4,-3);
	\draw[color ={black},opacity = 0.4] (-4,-4)--(4,-4);
	\draw[color ={black},opacity = 0.4] (-4,-5)--(4,-5);
	\draw[color ={black},opacity = 0.4] (-4,-6)--(4,-6);
	\draw[color ={black},opacity = 0.4] (-4,-7)--(4,-7);
	\draw[color ={black},opacity = 0.4] (-4,-8)--(4,-8);
	\draw[color ={black},opacity = 0.4] (1,-8)--(1,6);
	\draw[color ={black},opacity = 0.4] (2,-8)--(2,6);
	\draw[color ={black},opacity = 0.4] (3,-8)--(3,6);
	\draw[color ={black},opacity = 0.4] (4,-8)--(4,6);
	\draw[color ={black},opacity = 0.4] (-1,-8)--(-1,6);
	\draw[color ={black},opacity = 0.4] (-2,-8)--(-2,6);
	\draw[color ={black},opacity = 0.4] (-3,-8)--(-3,6);
	\draw[color ={black},opacity = 0.4] (-4,-8)--(-4,6);
	\draw[color ={gray},opacity = 0.3] (-4,0)--(4,0);
	\draw[color ={gray},opacity = 0.3] (-4,0.5)--(4,0.5);
	\draw[color ={gray},opacity = 0.3] (-4,1.5)--(4,1.5);
	\draw[color ={gray},opacity = 0.3] (-4,2.5)--(4,2.5);
	\draw[color ={gray},opacity = 0.3] (-4,3.5)--(4,3.5);
	\draw[color ={gray},opacity = 0.3] (-4,4.5)--(4,4.5);
	\draw[color ={gray},opacity = 0.3] (-4,5.5)--(4,5.5);
	\draw[color ={gray},opacity = 0.3] (-4,0)--(4,0);
	\draw[color ={gray},opacity = 0.3] (-4,-0.5)--(4,-0.5);
	\draw[color ={gray},opacity = 0.3] (-4,-1.5)--(4,-1.5);
	\draw[color ={gray},opacity = 0.3] (-4,-2.5)--(4,-2.5);
	\draw[color ={gray},opacity = 0.3] (-4,-3.5)--(4,-3.5);
	\draw[color ={gray},opacity = 0.3] (-4,-4.5)--(4,-4.5);
	\draw[color ={gray},opacity = 0.3] (-4,-5.5)--(4,-5.5);
	\draw[color ={gray},opacity = 0.3] (-4,-6.5)--(4,-6.5);
	\draw[color ={gray},opacity = 0.3] (-4,-7)--(4,-7);
	\draw[color ={gray},opacity = 0.3] (-4,-7.5)--(4,-7.5);
	\draw[color ={gray},opacity = 0.3] (-4,-8)--(4,-8);
	\draw[color ={gray},opacity = 0.3] (0,-8)--(0,6);
	\draw[color ={gray},opacity = 0.3] (0,-8)--(0,6);
	\draw[color ={black},line width = 1.2] (1,-0.13)--(1,0.13);
	\draw[color ={black},line width = 1.2] (2,-0.13)--(2,0.13);
	\draw[color ={black},line width = 1.2] (3,-0.13)--(3,0.13);
	\draw[color ={black},line width = 1.2] (4,-0.13)--(4,0.13);
	\draw[color ={black},line width = 1.2] (-1,-0.13)--(-1,0.13);
	\draw[color ={black},line width = 1.2] (-2,-0.13)--(-2,0.13);
	\draw[color ={black},line width = 1.2] (-3,-0.13)--(-3,0.13);
	\draw[color ={black},line width = 1.2] (-4,-0.13)--(-4,0.13);
	\draw[color ={black},line width = 1.2] (-0.13,1)--(0.13,1);
	\draw[color ={black},line width = 1.2] (-0.13,2)--(0.13,2);
	\draw[color ={black},line width = 1.2] (-0.13,3)--(0.13,3);
	\draw[color ={black},line width = 1.2] (-0.13,4)--(0.13,4);
	\draw[color ={black},line width = 1.2] (-0.13,5)--(0.13,5);
	\draw[color ={black},line width = 1.2] (-0.13,6)--(0.13,6);
	\draw[color ={black},line width = 1.2] (-0.13,-1)--(0.13,-1);
	\draw[color ={black},line width = 1.2] (-0.13,-2)--(0.13,-2);
	\draw[color ={black},line width = 1.2] (-0.13,-3)--(0.13,-3);
	\draw[color ={black},line width = 1.2] (-0.13,-4)--(0.13,-4);
	\draw[color ={black},line width = 1.2] (-0.13,-5)--(0.13,-5);
	\draw[color ={black},line width = 1.2] (-0.13,-6)--(0.13,-6);
	\draw[color ={black},line width = 1.2] (-0.13,-7)--(0.13,-7);
	\draw[color ={black},line width = 1.2] (-0.13,-8)--(0.13,-8);
	\draw (2,-0.4) node[anchor = center] {\scriptsize \color{black}{$2$}};
	\draw (4,-0.4) node[anchor = center] {\scriptsize \color{black}{$4$}};
	\draw (-2,-0.4) node[anchor = center] {\scriptsize \color{black}{$-2$}};
	\draw (-4,-0.4) node[anchor = center] {\scriptsize \color{black}{$-4$}};
	\draw (-0.5,2.1) node[anchor = center] {\footnotesize \color{black}{$2$}};
	\draw (-0.5,4.1) node[anchor = center] {\footnotesize \color{black}{$4$}};
	\draw (-0.5,-1.9) node[anchor = center] {\footnotesize \color{black}{$-2$}};
	\draw (-0.5,-3.9) node[anchor = center] {\footnotesize \color{black}{$-4$}};
	\draw (-0.5,-5.9) node[anchor = center] {\footnotesize \color{black}{$-6$}};
	\draw (-0.5,-7.9) node[anchor = center] {\footnotesize \color{black}{$-8$}};
	\draw [color={black}] (-0.3,-0.3) node[anchor = center,scale=1, rotate = 0] {O};
	\draw (3,4) node[anchor = center] {\footnotesize \color{red}{$\Large\mathcal{C}_f\prime$}};
	
	\draw[color={red},line width = 2] (-3.4,6.8)--(-3.2,6.4)--(-3,6)--(-2.8,5.6)--(-2.6,5.2)--(-2.4,4.8)--(-2.2,4.4)--(-2,4)--(-1.8,3.6)--(-1.6,3.2)--(-1.4,2.8)--(-1.2,2.4)--(-1,2)--(-0.8,1.6)--(-0.6,1.2)--(-0.4,0.8)--(-0.2,0.4)--(0,0)--(0.2,-0.4)--(0.4,-0.8)--(0.6,-1.2)--(0.8,-1.6)--(1,-2)--(1.2,-2.4)--(1.4,-2.8)--(1.6,-3.2)--(1.8,-3.6)--(2,-4)--(2.2,-4.4)--(2.4,-4.8)--(2.6,-5.2)--(2.8,-5.6)--(3,-6)--(3.2,-6.4)--(3.4,-6.8)--(3.6,-7.2)--(3.8,-7.6)--(4,-8);

\end{tikzpicture}\\
    \end{minipage}
    
\end{exercice}

\begin{Solution}
 L'équation réduite de la tangente au point d'abscisse $2$ est  : $y=f'(2)(x-2)+f(2)$.\\
      On lit graphiquement $f(2)=-3$ et $f'(2)=-4$.\\
      L'équation réduite de la tangente est donc donnée par :
      $y=-4(x-2)-3$, soit $y=-4x+5$.
\end{Solution}

% @see : https://coopmaths.fr/alea?uuid=a1ba2&id=can1F14&n=1&d=10&alea=rkIp2232323232342345678910111213141516171819202122232425262728293031323334353637383940414243444546474823456789101112131415&cd=1&cols=1
\needspace{10\baselineskip}
\begin{exercice}[BaremeTotal=false]


Soit $f$ la fonction définie sur $\mathbb{R}$ par : 
$f(x)= 4-2x^2-3x$.\\
En calculant la limite du taux d'accroissement, déterminer $f'(1)$.
\end{exercice}

\begin{Solution}
 $f$ est une fonction polynôme du second degré de la forme $f(x)=ax^2+bx+c$.\\
    La fonction dérivée est donnée par la somme des dérivées des fonctions $u$ et $v$ définies par $u(x)=-2x^2$ et $v(x)=-3x+4$.\\
     Comme $u'(x)=-4x$ et $v'(x)=-3$, on obtient  $f'(x)=-4x-3$. \\
     Ainsi, $f'(1)=-4\times 1-3=-7$.
\end{Solution}

% @see : https://coopmaths.fr/alea?uuid=3c690&id=can1F13&n=1&d=10&alea=EUDu2232323232342345678910111213141516171819202122232425262728293031323334353637383940414243444546474823456789101112131415&cd=1&cols=1
\needspace{10\baselineskip}
\begin{exercice}[BaremeTotal=false]


 Déterminer le coefficient directeur de la tangente à la courbe représentative de la fonction carré au point d'abscisse $15$.    
\end{exercice}

\begin{Solution}
 Le coefficient directeur de la tangente au point d'abscisse $a$ est donné par le nombre dérivé $f'(a)$.\\
        La fonction $f$ définie par $f(x)=x^2$ a pour fonction dérivée la fonction $f'$ définie par $f'(x)=2x$.\\
        Comme $f'(15)=2\times 15=30$, le coefficient directeur de la tangente au point d'abscisse $15$ est : $30$. 
\end{Solution}

% @see : https://coopmaths.fr/alea?uuid=ebd89&id=1AN14-1&n=1&d=10&s=3&alea=ocYr2232323232342345678910111213141516171819202122232425262728293031323334353637383940414243444546474823456789101112131415&cd=0&cols=1
\needspace{10\baselineskip}
\begin{exercice}[BaremeTotal=false]


 Donner la dérivée de la fonction $f$, dérivable pour tout $x\in\R$, définie par  $f(x)=-\dfrac{2}{9}x+9$.
\end{exercice}

\begin{Solution}
 L'expression de la dérivée de la fonction $f$ définie par $f(x)=-\dfrac{2}{9}x+9$ est : ${\color[HTML]{f15929}\boldsymbol{f'(x)=-\dfrac{2}{9}}}$.
\end{Solution}

% @see : https://coopmaths.fr/alea?uuid=ebd89&id=1AN14-1&n=1&d=10&s=4&alea=uAlP2232323232342345678910111213141516171819202122232425262728293031323334353637383940414243444546474823456789101112131415&cd=0&cols=1
\needspace{10\baselineskip}
\begin{exercice}[BaremeTotal=false]


 Donner la dérivée de la fonction $f$, dérivable pour tout $x\in\R$, définie par  $f(x)=-6x^{15}$.
\end{exercice}

\begin{Solution}
 L'expression de la dérivée de la fonction $f$ définie par $f(x)=-6x^{15}$ est : ${\color[HTML]{f15929}\boldsymbol{f'(x)=-90x^{14}}}$.
\end{Solution}

% @see : https://coopmaths.fr/alea?uuid=ebd89&id=1AN14-1&n=1&d=10&s=7&alea=vjN72232323232342345678910111213141516171819202122232425262728293031323334353637383940414243444546474823456789101112131415&cd=0&cols=1
\needspace{10\baselineskip}
\begin{exercice}[BaremeTotal=false]


 Donner la dérivée de la fonction $f$, dérivable pour tout $x\in\R^*$, définie par  $f(x)=\dfrac{2}{7x}$.
\end{exercice}

\begin{Solution}
 L'expression de la dérivée de la fonction $f$ définie par $f(x)=\dfrac{2}{7x}$ est : ${\color[HTML]{f15929}\boldsymbol{f'(x)=-\dfrac{2}{7x^2}}}$.
\end{Solution}

% @see : https://coopmaths.fr/alea?uuid=a83c0&id=1AN14-4&n=1&d=10&s=1&alea=Njr32232323232342345678910111213141516171819202122232425262728293031323334353637383940414243444546474823456789101112131415&cd=0&cols=1
\needspace{10\baselineskip}
\begin{exercice}[BaremeTotal=false]


 Donner l'expression de la dérivée de la fonction $f$ définie sur $\R^{*}$ par $f(x)=-2x^{2}-2x+5+\dfrac{4}{x}$.\\
\end{exercice}

\begin{Solution}
 L'expression de la dérivée de la fonction $f$ définie par $f(x)=-2x^{2}-2x+5+\dfrac{4}{x}$ est :\\$f'(x)={\color[HTML]{f15929}\boldsymbol{-4x-2-\dfrac{4}{x^2}}}$.\\
\end{Solution}

% @see : https://coopmaths.fr/alea?uuid=a83c0&id=1AN14-4&n=1&d=10&s=4&alea=4Vpz2232323232342345678910111213141516171819202122232425262728293031323334353637383940414243444546474823456789101112131415&cd=1&cols=1
\needspace{10\baselineskip}
\begin{exercice}[BaremeTotal=false]


 Donner l'expression de la dérivée de la fonction $f$ définie sur $\R_+^{*}$ par $f(x)=-\sqrt{x}+\dfrac{4}{x}-5x^{2}+2x$.\\
\end{exercice}

\begin{Solution}
 $(-\sqrt{x})^\prime=-\dfrac{1}{2\sqrt{x}}$\\$(\dfrac{4}{x})^\prime=-\dfrac{4}{x^2}$\\$(-5x^{2}+2x)^\prime=-10x+2$\\L'expression de la dérivée de la fonction $f$ définie par $f(x)=-\sqrt{x}+\dfrac{4}{x}-5x^{2}+2x$ est :\\$f'(x)={\color[HTML]{f15929}\boldsymbol{-\dfrac{1}{2\sqrt{x}}-\dfrac{4}{x^2}-10x+2}}$.\\
\end{Solution}

\end{Maquette}
\clearpage
\begin{Maquette}[DM]{Niveau= ,Classe=\getdata[16]\mydata\ (v16), Date = 06/01/25  ,Theme=Exercices}


% @see : https://coopmaths.fr/alea?uuid=4c8c7&id=1AN11-3&n=1&d=10&s=2&alea=nawO223232323234234567891011121314151617181920212223242526272829303132333435363738394041424344454647482345678910111213141516&cd=1&cols=2
\needspace{10\baselineskip}
\begin{exercice}[BaremeTotal=false]
\label{eleve:16}


 \begin{minipage}[t]{\linewidth} Soit $f$ une fonction dérivable sur $[-5;5]$ et $\mathcal{C}_f$ sa courbe représentative.\\On sait que  $f(3)=-3~~$ et que $~~f'(3)=5$.\\Déterminer une équation de la tangente $(T)$ à la courbe $\mathcal{C}_f$ au point d'abscisse $3$,\\sans utiliser la formule de cours de l'équation de tangente. \end{minipage}

\end{exercice}

\begin{Solution}
 \begin{minipage}[t]{\linewidth} On sait que la tangente n'est pas une droite verticale, puisque la fonction est dérivable sur l'intervalle.\\On en déduit que la tangente $(T)$ au point d'abscisse $3$, admet une équation réduite de la forme :  $(T) : y=m x + p$. 

\medskip
$\bullet$ Détermination de $m$ :\\On sait que le nombre dérivé en $3$ est par définition, le coefficient directeur de la tangente au point d'abscisse $3$.\\Par conséquent, on a déjà : $m=f'(3)=5$.\\ On en déduit que  $(T) : y= 5 x + p$\\$\bullet$ Détermination de $p$ :\\ Pour cela, on utilise que si $f(3)=-3~~$, alors le point $A$ de coordonnées $(3;-3)$ appartient à $\mathcal{C}_f$ mais aussi à $(T)$.\\ On peut écrire $A(3;-3) = \mathcal{C}_f \cap (T)$.\\ On remplace alors les coordonnées de $A(3;-3)$ dans l'équation  $(T) : y= 5 x + p$.\\ $\begin{aligned} \phantom{\iff}&A(3;-3)\in (T)\\ \iff& -3= 5 \times 3  + p\\ \iff& p=-3 -5 \times 3  \\ \iff& p=-3 -15   \\ \iff& p=-18\\\end{aligned}$\\On peut conclure que : $(T) : {\color[HTML]{f15929}\boldsymbol{y=5x-18}}$. \end{minipage}

\end{Solution}

% @see : https://coopmaths.fr/alea?uuid=0e984&id=can1F15&n=1&d=10&alea=1Alo223232323234234567891011121314151617181920212223242526272829303132333435363738394041424344454647482345678910111213141516&cd=1&cols=1
\needspace{10\baselineskip}
\begin{exercice}[BaremeTotal=false]


 La courbe représente une fonction $f$ et la droite est la tangente au point d'abscisse $-1$.\\
        
        Déterminer $f'(-1)$.\begin{center}\begin{tikzpicture}[baseline,scale = 0.5]

    \tikzset{
      point/.style={
        thick,
        draw,
        cross out,
        inner sep=0pt,
        minimum width=5pt,
        minimum height=5pt,
      },
    }
    \clip (-8,-3) rectangle (8,10);
    	
	\draw[color ={black},line width = 2,>=latex,->] (-8,0)--(8,0);
	\draw[color ={black},line width = 2,>=latex,->] (0,-3)--(0,10);
	\draw[color ={black},opacity = 0.5] (-8,1)--(8,1);
	\draw[color ={black},opacity = 0.5] (-8,2)--(8,2);
	\draw[color ={black},opacity = 0.5] (-8,3)--(8,3);
	\draw[color ={black},opacity = 0.5] (-8,4)--(8,4);
	\draw[color ={black},opacity = 0.5] (-8,5)--(8,5);
	\draw[color ={black},opacity = 0.5] (-8,6)--(8,6);
	\draw[color ={black},opacity = 0.5] (-8,7)--(8,7);
	\draw[color ={black},opacity = 0.5] (-8,8)--(8,8);
	\draw[color ={black},opacity = 0.5] (-8,9)--(8,9);
	\draw[color ={black},opacity = 0.5] (-8,10)--(8,10);
	\draw[color ={black},opacity = 0.5] (-8,-1)--(8,-1);
	\draw[color ={black},opacity = 0.5] (-8,-2)--(8,-2);
	\draw[color ={black},opacity = 0.5] (-8,-3)--(8,-3);
	\draw[color ={black},opacity = 0.5] (1,-3)--(1,10);
	\draw[color ={black},opacity = 0.5] (2,-3)--(2,10);
	\draw[color ={black},opacity = 0.5] (3,-3)--(3,10);
	\draw[color ={black},opacity = 0.5] (4,-3)--(4,10);
	\draw[color ={black},opacity = 0.5] (5,-3)--(5,10);
	\draw[color ={black},opacity = 0.5] (6,-3)--(6,10);
	\draw[color ={black},opacity = 0.5] (7,-3)--(7,10);
	\draw[color ={black},opacity = 0.5] (8,-3)--(8,10);
	\draw[color ={black},opacity = 0.5] (-1,-3)--(-1,10);
	\draw[color ={black},opacity = 0.5] (-2,-3)--(-2,10);
	\draw[color ={black},opacity = 0.5] (-3,-3)--(-3,10);
	\draw[color ={black},opacity = 0.5] (-4,-3)--(-4,10);
	\draw[color ={black},opacity = 0.5] (-5,-3)--(-5,10);
	\draw[color ={black},opacity = 0.5] (-6,-3)--(-6,10);
	\draw[color ={black},opacity = 0.5] (-7,-3)--(-7,10);
	\draw[color ={black},opacity = 0.5] (-8,-3)--(-8,10);
	\draw[color ={gray},opacity = 0.3] (-8,0)--(8,0);
	\draw[color ={gray},opacity = 0.3] (-8,0)--(8,0);
	\draw[color ={gray},opacity = 0.3] (0,-3)--(0,10);
	\draw[color ={gray},opacity = 0.3] (0,-3)--(0,10);
	\draw[color ={black},line width = 2] (2,-0.2)--(2,0.2);
	\draw[color ={black},line width = 2] (4,-0.2)--(4,0.2);
	\draw[color ={black},line width = 2] (6,-0.2)--(6,0.2);
	\draw[color ={black},line width = 2] (-2,-0.2)--(-2,0.2);
	\draw[color ={black},line width = 2] (-4,-0.2)--(-4,0.2);
	\draw[color ={black},line width = 2] (-6,-0.2)--(-6,0.2);
	\draw[color ={black},line width = 2] (-0.2,1)--(0.2,1);
	\draw[color ={black},line width = 2] (-0.2,2)--(0.2,2);
	\draw[color ={black},line width = 2] (-0.2,3)--(0.2,3);
	\draw[color ={black},line width = 2] (-0.2,4)--(0.2,4);
	\draw[color ={black},line width = 2] (-0.2,5)--(0.2,5);
	\draw[color ={black},line width = 2] (-0.2,6)--(0.2,6);
	\draw[color ={black},line width = 2] (-0.2,7)--(0.2,7);
	\draw[color ={black},line width = 2] (-0.2,8)--(0.2,8);
	\draw[color ={black},line width = 2] (-0.2,9)--(0.2,9);
	\draw[color ={black},line width = 2] (-0.2,-1)--(0.2,-1);
	\draw[color ={black},line width = 2] (-0.2,-2)--(0.2,-2);
	\draw (2,-0.4) node[anchor = center] {\scriptsize \color{black}{$1$}};
	\draw (4,-0.4) node[anchor = center] {\scriptsize \color{black}{$2$}};
	\draw (6,-0.4) node[anchor = center] {\scriptsize \color{black}{$3$}};
	\draw (-2,-0.4) node[anchor = center] {\scriptsize \color{black}{$-1$}};
	\draw (-4,-0.4) node[anchor = center] {\scriptsize \color{black}{$-2$}};
	\draw (-6,-0.4) node[anchor = center] {\scriptsize \color{black}{$-3$}};
	\draw (-0.8,2.1) node[anchor = center] {\footnotesize \color{black}{$2$}};
	\draw (-0.8,4.1) node[anchor = center] {\footnotesize \color{black}{$4$}};
	\draw (-0.8,6.1) node[anchor = center] {\footnotesize \color{black}{$6$}};
	\draw (-0.8,8.1) node[anchor = center] {\footnotesize \color{black}{$8$}};
	\draw [color={black}] (-0.3,-0.3) node[anchor = center,scale=1, rotate = 0] {O};
	
	\draw[color={blue},line width = 2] (-4,11)--(-3.6,9.84)--(-3.2,8.76)--(-2.8,7.76)--(-2.4,6.84)--(-2,6)--(-1.6,5.24)--(-1.2,4.56)--(-0.8,3.96)--(-0.4,3.44)--(0,3)--(0.4,2.64)--(0.8,2.36)--(1.2,2.16)--(1.6,2.04)--(2,2)--(2.4,2.04)--(2.8,2.16)--(3.2,2.36)--(3.6,2.64)--(4,3)--(4.4,3.44)--(4.8,3.96)--(5.2,4.56)--(5.6,5.24)--(6,6)--(6.4,6.84)--(6.8,7.76)--(7.2,8.76)--(7.6,9.84);
	
	\draw[color={red},line width = 2] (-4.4,10.8)--(-4,10)--(-3.6,9.2)--(-3.2,8.4)--(-2.8,7.6)--(-2.4,6.8)--(-2,6)--(-1.6,5.2)--(-1.2,4.4)--(-0.8,3.6)--(-0.4,2.8)--(0,2)--(0.4,1.2)--(0.8,0.4)--(1.2,-0.4)--(1.6,-1.2)--(2,-2)--(2.4,-2.8)--(2.8,-3.6);

\end{tikzpicture}\end{center}
\end{exercice}

\begin{Solution}
 $f'(-1)$ est donné par le coefficient directeur de la tangente à la courbe au point d'abscisse $-1$, soit $-4$.
\end{Solution}

% @see : https://coopmaths.fr/alea?uuid=6f32d&id=can1F16&n=1&d=10&alea=dqSw223232323234234567891011121314151617181920212223242526272829303132333435363738394041424344454647482345678910111213141516&cd=1&cols=1
\needspace{10\baselineskip}
\newpage
\begin{exercice}[BaremeTotal=false]


 On donne les représentations graphiques d'une fonction et de sa dérivée.\\
      Donner l'équation réduite de la tangente à la courbe de $f$ en $x=-2$. \\ \begin{minipage}{0.5\linewidth}
    \begin{tikzpicture}[baseline, scale=0.8]

    \tikzset{
      point/.style={
        thick,
        draw,
        cross out,
        inner sep=0pt,
        minimum width=5pt,
        minimum height=5pt,
      },
    }
    \clip (-5.5,-9.5) rectangle (5.75,7.5);
    	
	\draw[color ={black},>=latex,->] (-4,0)--(4,0);
	\draw[color ={black},>=latex,->] (0,-8)--(0,6);
	\draw[color ={black},opacity = 0.4] (-4,1)--(4,1);
	\draw[color ={black},opacity = 0.4] (-4,2)--(4,2);
	\draw[color ={black},opacity = 0.4] (-4,3)--(4,3);
	\draw[color ={black},opacity = 0.4] (-4,4)--(4,4);
	\draw[color ={black},opacity = 0.4] (-4,5)--(4,5);
	\draw[color ={black},opacity = 0.4] (-4,6)--(4,6);
	\draw[color ={black},opacity = 0.4] (-4,-1)--(4,-1);
	\draw[color ={black},opacity = 0.4] (-4,-2)--(4,-2);
	\draw[color ={black},opacity = 0.4] (-4,-3)--(4,-3);
	\draw[color ={black},opacity = 0.4] (-4,-4)--(4,-4);
	\draw[color ={black},opacity = 0.4] (-4,-5)--(4,-5);
	\draw[color ={black},opacity = 0.4] (-4,-6)--(4,-6);
	\draw[color ={black},opacity = 0.4] (-4,-7)--(4,-7);
	\draw[color ={black},opacity = 0.4] (-4,-8)--(4,-8);
	\draw[color ={black},opacity = 0.4] (1,-8)--(1,6);
	\draw[color ={black},opacity = 0.4] (2,-8)--(2,6);
	\draw[color ={black},opacity = 0.4] (3,-8)--(3,6);
	\draw[color ={black},opacity = 0.4] (4,-8)--(4,6);
	\draw[color ={black},opacity = 0.4] (-1,-8)--(-1,6);
	\draw[color ={black},opacity = 0.4] (-2,-8)--(-2,6);
	\draw[color ={black},opacity = 0.4] (-3,-8)--(-3,6);
	\draw[color ={black},opacity = 0.4] (-4,-8)--(-4,6);
	\draw[color ={gray},opacity = 0.3] (-4,0)--(4,0);
	\draw[color ={gray},opacity = 0.3] (-4,0.5)--(4,0.5);
	\draw[color ={gray},opacity = 0.3] (-4,1.5)--(4,1.5);
	\draw[color ={gray},opacity = 0.3] (-4,2.5)--(4,2.5);
	\draw[color ={gray},opacity = 0.3] (-4,3.5)--(4,3.5);
	\draw[color ={gray},opacity = 0.3] (-4,4.5)--(4,4.5);
	\draw[color ={gray},opacity = 0.3] (-4,5.5)--(4,5.5);
	\draw[color ={gray},opacity = 0.3] (-4,0)--(4,0);
	\draw[color ={gray},opacity = 0.3] (-4,-0.5)--(4,-0.5);
	\draw[color ={gray},opacity = 0.3] (-4,-1.5)--(4,-1.5);
	\draw[color ={gray},opacity = 0.3] (-4,-2.5)--(4,-2.5);
	\draw[color ={gray},opacity = 0.3] (-4,-3.5)--(4,-3.5);
	\draw[color ={gray},opacity = 0.3] (-4,-4.5)--(4,-4.5);
	\draw[color ={gray},opacity = 0.3] (-4,-5.5)--(4,-5.5);
	\draw[color ={gray},opacity = 0.3] (-4,-6.5)--(4,-6.5);
	\draw[color ={gray},opacity = 0.3] (-4,-7)--(4,-7);
	\draw[color ={gray},opacity = 0.3] (-4,-7.5)--(4,-7.5);
	\draw[color ={gray},opacity = 0.3] (-4,-8)--(4,-8);
	\draw[color ={gray},opacity = 0.3] (0,-8)--(0,6);
	\draw[color ={gray},opacity = 0.3] (0,-8)--(0,6);
	\draw[color ={black},line width = 1.2] (1,-0.13)--(1,0.13);
	\draw[color ={black},line width = 1.2] (2,-0.13)--(2,0.13);
	\draw[color ={black},line width = 1.2] (3,-0.13)--(3,0.13);
	\draw[color ={black},line width = 1.2] (4,-0.13)--(4,0.13);
	\draw[color ={black},line width = 1.2] (-1,-0.13)--(-1,0.13);
	\draw[color ={black},line width = 1.2] (-2,-0.13)--(-2,0.13);
	\draw[color ={black},line width = 1.2] (-3,-0.13)--(-3,0.13);
	\draw[color ={black},line width = 1.2] (-4,-0.13)--(-4,0.13);
	\draw[color ={black},line width = 1.2] (-0.13,1)--(0.13,1);
	\draw[color ={black},line width = 1.2] (-0.13,2)--(0.13,2);
	\draw[color ={black},line width = 1.2] (-0.13,3)--(0.13,3);
	\draw[color ={black},line width = 1.2] (-0.13,4)--(0.13,4);
	\draw[color ={black},line width = 1.2] (-0.13,5)--(0.13,5);
	\draw[color ={black},line width = 1.2] (-0.13,6)--(0.13,6);
	\draw[color ={black},line width = 1.2] (-0.13,-1)--(0.13,-1);
	\draw[color ={black},line width = 1.2] (-0.13,-2)--(0.13,-2);
	\draw[color ={black},line width = 1.2] (-0.13,-3)--(0.13,-3);
	\draw[color ={black},line width = 1.2] (-0.13,-4)--(0.13,-4);
	\draw[color ={black},line width = 1.2] (-0.13,-5)--(0.13,-5);
	\draw[color ={black},line width = 1.2] (-0.13,-6)--(0.13,-6);
	\draw[color ={black},line width = 1.2] (-0.13,-7)--(0.13,-7);
	\draw[color ={black},line width = 1.2] (-0.13,-8)--(0.13,-8);
	\draw (2,-0.4) node[anchor = center] {\scriptsize \color{black}{$2$}};
	\draw (4,-0.4) node[anchor = center] {\scriptsize \color{black}{$4$}};
	\draw (-2,-0.4) node[anchor = center] {\scriptsize \color{black}{$-2$}};
	\draw (-4,-0.4) node[anchor = center] {\scriptsize \color{black}{$-4$}};
	\draw (-0.5,2.1) node[anchor = center] {\footnotesize \color{black}{$2$}};
	\draw (-0.5,4.1) node[anchor = center] {\footnotesize \color{black}{$4$}};
	\draw (-0.5,6.1) node[anchor = center] {\footnotesize \color{black}{$6$}};
	\draw (-0.5,-1.9) node[anchor = center] {\footnotesize \color{black}{$-2$}};
	\draw (-0.5,-3.9) node[anchor = center] {\footnotesize \color{black}{$-4$}};
	\draw (-0.5,-5.9) node[anchor = center] {\footnotesize \color{black}{$-6$}};
	\draw (-0.5,-7.9) node[anchor = center] {\footnotesize \color{black}{$-8$}};
	\draw [color={black}] (-0.3,-0.3) node[anchor = center,scale=1, rotate = 0] {O};
	\draw (3,4) node[anchor = center] {\footnotesize \color{blue}{$\Large \mathcal{C}_f$}};
	
	\draw[color={blue},line width = 2] (-4,-6)--(-3.8,-4.84)--(-3.6,-3.76)--(-3.4,-2.76)--(-3.2,-1.84)--(-3,-1)--(-2.8,-0.24)--(-2.6,0.44)--(-2.4,1.04)--(-2.2,1.56)--(-2,2)--(-1.8,2.36)--(-1.6,2.64)--(-1.4,2.84)--(-1.2,2.96)--(-1,3)--(-0.8,2.96)--(-0.6,2.84)--(-0.4,2.64)--(-0.2,2.36)--(0,2)--(0.2,1.56)--(0.4,1.04)--(0.6,0.44)--(0.8,-0.24)--(1,-1)--(1.2,-1.84)--(1.4,-2.76)--(1.6,-3.76)--(1.8,-4.84)--(2,-6)--(2.2,-7.24)--(2.4,-8.56);

\end{tikzpicture}\\
    \end{minipage}
    \begin{minipage}{0.5\linewidth}
    \begin{tikzpicture}[baseline, scale=0.8]

    \tikzset{
      point/.style={
        thick,
        draw,
        cross out,
        inner sep=0pt,
        minimum width=5pt,
        minimum height=5pt,
      },
    }
    \clip (-5.5,-9.5) rectangle (6.5,7.5);
    	
	\draw[color ={black},>=latex,->] (-4,0)--(4,0);
	\draw[color ={black},>=latex,->] (0,-8)--(0,6);
	\draw[color ={black},opacity = 0.4] (-4,1)--(4,1);
	\draw[color ={black},opacity = 0.4] (-4,2)--(4,2);
	\draw[color ={black},opacity = 0.4] (-4,3)--(4,3);
	\draw[color ={black},opacity = 0.4] (-4,4)--(4,4);
	\draw[color ={black},opacity = 0.4] (-4,5)--(4,5);
	\draw[color ={black},opacity = 0.4] (-4,6)--(4,6);
	\draw[color ={black},opacity = 0.4] (-4,-1)--(4,-1);
	\draw[color ={black},opacity = 0.4] (-4,-2)--(4,-2);
	\draw[color ={black},opacity = 0.4] (-4,-3)--(4,-3);
	\draw[color ={black},opacity = 0.4] (-4,-4)--(4,-4);
	\draw[color ={black},opacity = 0.4] (-4,-5)--(4,-5);
	\draw[color ={black},opacity = 0.4] (-4,-6)--(4,-6);
	\draw[color ={black},opacity = 0.4] (-4,-7)--(4,-7);
	\draw[color ={black},opacity = 0.4] (-4,-8)--(4,-8);
	\draw[color ={black},opacity = 0.4] (1,-8)--(1,6);
	\draw[color ={black},opacity = 0.4] (2,-8)--(2,6);
	\draw[color ={black},opacity = 0.4] (3,-8)--(3,6);
	\draw[color ={black},opacity = 0.4] (4,-8)--(4,6);
	\draw[color ={black},opacity = 0.4] (-1,-8)--(-1,6);
	\draw[color ={black},opacity = 0.4] (-2,-8)--(-2,6);
	\draw[color ={black},opacity = 0.4] (-3,-8)--(-3,6);
	\draw[color ={black},opacity = 0.4] (-4,-8)--(-4,6);
	\draw[color ={gray},opacity = 0.3] (-4,0)--(4,0);
	\draw[color ={gray},opacity = 0.3] (-4,0.5)--(4,0.5);
	\draw[color ={gray},opacity = 0.3] (-4,1.5)--(4,1.5);
	\draw[color ={gray},opacity = 0.3] (-4,2.5)--(4,2.5);
	\draw[color ={gray},opacity = 0.3] (-4,3.5)--(4,3.5);
	\draw[color ={gray},opacity = 0.3] (-4,4.5)--(4,4.5);
	\draw[color ={gray},opacity = 0.3] (-4,5.5)--(4,5.5);
	\draw[color ={gray},opacity = 0.3] (-4,0)--(4,0);
	\draw[color ={gray},opacity = 0.3] (-4,-0.5)--(4,-0.5);
	\draw[color ={gray},opacity = 0.3] (-4,-1.5)--(4,-1.5);
	\draw[color ={gray},opacity = 0.3] (-4,-2.5)--(4,-2.5);
	\draw[color ={gray},opacity = 0.3] (-4,-3.5)--(4,-3.5);
	\draw[color ={gray},opacity = 0.3] (-4,-4.5)--(4,-4.5);
	\draw[color ={gray},opacity = 0.3] (-4,-5.5)--(4,-5.5);
	\draw[color ={gray},opacity = 0.3] (-4,-6.5)--(4,-6.5);
	\draw[color ={gray},opacity = 0.3] (-4,-7)--(4,-7);
	\draw[color ={gray},opacity = 0.3] (-4,-7.5)--(4,-7.5);
	\draw[color ={gray},opacity = 0.3] (-4,-8)--(4,-8);
	\draw[color ={gray},opacity = 0.3] (0,-8)--(0,6);
	\draw[color ={gray},opacity = 0.3] (0,-8)--(0,6);
	\draw[color ={black},line width = 1.2] (1,-0.13)--(1,0.13);
	\draw[color ={black},line width = 1.2] (2,-0.13)--(2,0.13);
	\draw[color ={black},line width = 1.2] (3,-0.13)--(3,0.13);
	\draw[color ={black},line width = 1.2] (4,-0.13)--(4,0.13);
	\draw[color ={black},line width = 1.2] (-1,-0.13)--(-1,0.13);
	\draw[color ={black},line width = 1.2] (-2,-0.13)--(-2,0.13);
	\draw[color ={black},line width = 1.2] (-3,-0.13)--(-3,0.13);
	\draw[color ={black},line width = 1.2] (-4,-0.13)--(-4,0.13);
	\draw[color ={black},line width = 1.2] (-0.13,1)--(0.13,1);
	\draw[color ={black},line width = 1.2] (-0.13,2)--(0.13,2);
	\draw[color ={black},line width = 1.2] (-0.13,3)--(0.13,3);
	\draw[color ={black},line width = 1.2] (-0.13,4)--(0.13,4);
	\draw[color ={black},line width = 1.2] (-0.13,5)--(0.13,5);
	\draw[color ={black},line width = 1.2] (-0.13,6)--(0.13,6);
	\draw[color ={black},line width = 1.2] (-0.13,-1)--(0.13,-1);
	\draw[color ={black},line width = 1.2] (-0.13,-2)--(0.13,-2);
	\draw[color ={black},line width = 1.2] (-0.13,-3)--(0.13,-3);
	\draw[color ={black},line width = 1.2] (-0.13,-4)--(0.13,-4);
	\draw[color ={black},line width = 1.2] (-0.13,-5)--(0.13,-5);
	\draw[color ={black},line width = 1.2] (-0.13,-6)--(0.13,-6);
	\draw[color ={black},line width = 1.2] (-0.13,-7)--(0.13,-7);
	\draw[color ={black},line width = 1.2] (-0.13,-8)--(0.13,-8);
	\draw (2,-0.4) node[anchor = center] {\scriptsize \color{black}{$2$}};
	\draw (4,-0.4) node[anchor = center] {\scriptsize \color{black}{$4$}};
	\draw (-2,-0.4) node[anchor = center] {\scriptsize \color{black}{$-2$}};
	\draw (-4,-0.4) node[anchor = center] {\scriptsize \color{black}{$-4$}};
	\draw (-0.5,2.1) node[anchor = center] {\footnotesize \color{black}{$2$}};
	\draw (-0.5,4.1) node[anchor = center] {\footnotesize \color{black}{$4$}};
	\draw (-0.5,-1.9) node[anchor = center] {\footnotesize \color{black}{$-2$}};
	\draw (-0.5,-3.9) node[anchor = center] {\footnotesize \color{black}{$-4$}};
	\draw (-0.5,-5.9) node[anchor = center] {\footnotesize \color{black}{$-6$}};
	\draw (-0.5,-7.9) node[anchor = center] {\footnotesize \color{black}{$-8$}};
	\draw [color={black}] (-0.3,-0.3) node[anchor = center,scale=1, rotate = 0] {O};
	\draw (3,4) node[anchor = center] {\footnotesize \color{red}{$\Large\mathcal{C}_f\prime$}};
	
	\draw[color={red},line width = 2] (-4,6)--(-3.8,5.6)--(-3.6,5.2)--(-3.4,4.8)--(-3.2,4.4)--(-3,4)--(-2.8,3.6)--(-2.6,3.2)--(-2.4,2.8)--(-2.2,2.4)--(-2,2)--(-1.8,1.6)--(-1.6,1.2)--(-1.4,0.8)--(-1.2,0.4)--(-1,0)--(-0.8,-0.4)--(-0.6,-0.8)--(-0.4,-1.2)--(-0.2,-1.6)--(0,-2)--(0.2,-2.4)--(0.4,-2.8)--(0.6,-3.2)--(0.8,-3.6)--(1,-4)--(1.2,-4.4)--(1.4,-4.8)--(1.6,-5.2)--(1.8,-5.6)--(2,-6)--(2.2,-6.4)--(2.4,-6.8)--(2.6,-7.2)--(2.8,-7.6)--(3,-8)--(3.2,-8.4)--(3.4,-8.8);

\end{tikzpicture}\\
    \end{minipage}
    
\end{exercice}

\begin{Solution}
 L'équation réduite de la tangente au point d'abscisse $-2$ est  : $y=f'(-2)(x-(-2))+f(-2)$.\\
      On lit graphiquement $f(-2)=2$ et $f'(-2)=2$.\\
      L'équation réduite de la tangente est donc donnée par :
      $y=2(x+2)+2$, soit $y=2x+6$.
\end{Solution}

% @see : https://coopmaths.fr/alea?uuid=a1ba2&id=can1F14&n=1&d=10&alea=rkIp223232323234234567891011121314151617181920212223242526272829303132333435363738394041424344454647482345678910111213141516&cd=1&cols=1
\needspace{10\baselineskip}
\begin{exercice}[BaremeTotal=false]


Soit $f$ la fonction définie sur $\mathbb{R}$ par : 
$f(x)= 7x^2-1$.\\
En calculant la limite du taux d'accroissement, déterminer $f'(-1)$.
\end{exercice}

\begin{Solution}
 $f$ est une fonction polynôme du second degré de la forme $f(x)=ax^2+b$.\\
    La fonction dérivée est donnée par la somme des dérivées des fonctions $u$ et $v$ définies par $u(x)=7x^2$ et $v(x)=-1$.\\
     Comme $u'(x)=14x$ et $v'(x)=0$, on obtient  $f'(x)=14x$. \\
     Ainsi, $f'(-1)=14\times (-1)=-14$.
\end{Solution}

% @see : https://coopmaths.fr/alea?uuid=3c690&id=can1F13&n=1&d=10&alea=EUDu223232323234234567891011121314151617181920212223242526272829303132333435363738394041424344454647482345678910111213141516&cd=1&cols=1
\needspace{10\baselineskip}
\begin{exercice}[BaremeTotal=false]


 Déterminer le coefficient directeur de la tangente à la courbe représentative de la fonction racine carrée au point d'abscisse $2$.
         
\end{exercice}

\begin{Solution}
 Le coefficient directeur de la tangente au point d'abscisse $a$ est donné par le nombre dérivé $f'(a)$.\\
        La fonction $f$ définie par $f(x)=\sqrt{x}$ a pour fonction dérivée la fonction $f'$ définie par $f'(x)=\dfrac{1}{2\sqrt{x}}$.\\Comme $f'(2)=\dfrac{1}{2\sqrt{2}}$, le coefficient directeur de la tangente au point d'abscisse $2$ est : $\dfrac{1}{2\sqrt{2}}$.
\end{Solution}

% @see : https://coopmaths.fr/alea?uuid=ebd89&id=1AN14-1&n=1&d=10&s=3&alea=ocYr223232323234234567891011121314151617181920212223242526272829303132333435363738394041424344454647482345678910111213141516&cd=0&cols=1
\needspace{10\baselineskip}
\begin{exercice}[BaremeTotal=false]


 Donner la dérivée de la fonction $f$, dérivable pour tout $x\in\R$, définie par  $f(x)=\dfrac{3}{4}x+7$.
\end{exercice}

\begin{Solution}
 L'expression de la dérivée de la fonction $f$ définie par $f(x)=\dfrac{3}{4}x+7$ est : ${\color[HTML]{f15929}\boldsymbol{f'(x)=\dfrac{3}{4}}}$.
\end{Solution}

% @see : https://coopmaths.fr/alea?uuid=ebd89&id=1AN14-1&n=1&d=10&s=4&alea=uAlP223232323234234567891011121314151617181920212223242526272829303132333435363738394041424344454647482345678910111213141516&cd=0&cols=1
\needspace{10\baselineskip}
\begin{exercice}[BaremeTotal=false]


 Donner la dérivée de la fonction $f$, dérivable pour tout $x\in\R$, définie par  $f(x)=-8x^{5}$.
\end{exercice}

\begin{Solution}
 L'expression de la dérivée de la fonction $f$ définie par $f(x)=-8x^{5}$ est : ${\color[HTML]{f15929}\boldsymbol{f'(x)=-40x^{4}}}$.
\end{Solution}

% @see : https://coopmaths.fr/alea?uuid=ebd89&id=1AN14-1&n=1&d=10&s=7&alea=vjN7223232323234234567891011121314151617181920212223242526272829303132333435363738394041424344454647482345678910111213141516&cd=0&cols=1
\needspace{10\baselineskip}
\begin{exercice}[BaremeTotal=false]


 Donner la dérivée de la fonction $f$, dérivable pour tout $x\in\R^*$, définie par  $f(x)=\dfrac{-1}{4x}$.
\end{exercice}

\begin{Solution}
 L'expression de la dérivée de la fonction $f$ définie par $f(x)=\dfrac{-1}{4x}$ est : ${\color[HTML]{f15929}\boldsymbol{f'(x)=\dfrac{1}{4x^2}}}$.
\end{Solution}

% @see : https://coopmaths.fr/alea?uuid=a83c0&id=1AN14-4&n=1&d=10&s=1&alea=Njr3223232323234234567891011121314151617181920212223242526272829303132333435363738394041424344454647482345678910111213141516&cd=0&cols=1
\needspace{10\baselineskip}
\begin{exercice}[BaremeTotal=false]


 Donner l'expression de la dérivée de la fonction $f$ définie sur $\R^{*}$ par $f(x)=5x^{2}+3x-4+\dfrac{4}{x}$.\\
\end{exercice}

\begin{Solution}
 L'expression de la dérivée de la fonction $f$ définie par $f(x)=5x^{2}+3x-4+\dfrac{4}{x}$ est :\\$f'(x)={\color[HTML]{f15929}\boldsymbol{10x+3-\dfrac{4}{x^2}}}$.\\
\end{Solution}

% @see : https://coopmaths.fr/alea?uuid=a83c0&id=1AN14-4&n=1&d=10&s=4&alea=4Vpz223232323234234567891011121314151617181920212223242526272829303132333435363738394041424344454647482345678910111213141516&cd=1&cols=1
\needspace{10\baselineskip}
\begin{exercice}[BaremeTotal=false]


 Donner l'expression de la dérivée de la fonction $f$ définie sur $\R_+^{*}$ par $f(x)=-\sqrt{x}+2x^{2}-5x+\dfrac{5}{x}$.\\
\end{exercice}

\begin{Solution}
 $(-\sqrt{x})^\prime=-\dfrac{1}{2\sqrt{x}}$\\$(2x^{2}-5x)^\prime=4x-5$\\$(\dfrac{5}{x})^\prime=-\dfrac{5}{x^2}$\\L'expression de la dérivée de la fonction $f$ définie par $f(x)=-\sqrt{x}+2x^{2}-5x+\dfrac{5}{x}$ est :\\$f'(x)={\color[HTML]{f15929}\boldsymbol{-\dfrac{1}{2\sqrt{x}}+4x-5-\dfrac{5}{x^2}}}$.\\
\end{Solution}

\end{Maquette}
\clearpage
\begin{Maquette}[DM]{Niveau= ,Classe=\getdata[17]\mydata\ (v17), Date = 06/01/25  ,Theme=Exercices}


% @see : https://coopmaths.fr/alea?uuid=4c8c7&id=1AN11-3&n=1&d=10&s=2&alea=nawO22323232323423456789101112131415161718192021222324252627282930313233343536373839404142434445464748234567891011121314151617&cd=1&cols=2
\needspace{10\baselineskip}
\begin{exercice}[BaremeTotal=false]
\label{eleve:17}


 \begin{minipage}[t]{\linewidth} Soit $f$ une fonction dérivable sur $[-5;5]$ et $\mathcal{C}_f$ sa courbe représentative.\\On sait que  $f(-3)=-5~~$ et que $~~f'(-3)=-5$.\\Déterminer une équation de la tangente $(T)$ à la courbe $\mathcal{C}_f$ au point d'abscisse $-3$,\\sans utiliser la formule de cours de l'équation de tangente. \end{minipage}

\end{exercice}

\begin{Solution}
 \begin{minipage}[t]{\linewidth} On sait que la tangente n'est pas une droite verticale, puisque la fonction est dérivable sur l'intervalle.\\On en déduit que la tangente $(T)$ au point d'abscisse $-3$, admet une équation réduite de la forme :  $(T) : y=m x + p$. 

\medskip
$\bullet$ Détermination de $m$ :\\On sait que le nombre dérivé en $-3$ est par définition, le coefficient directeur de la tangente au point d'abscisse $-3$.\\Par conséquent, on a déjà : $m=f'(-3)=-5$.\\ On en déduit que  $(T) : y= -5 x + p$\\$\bullet$ Détermination de $p$ :\\ Pour cela, on utilise que si $f(-3)=-5~~$, alors le point $A$ de coordonnées $(-3;-5)$ appartient à $\mathcal{C}_f$ mais aussi à $(T)$.\\ On peut écrire $A(-3;-5) = \mathcal{C}_f \cap (T)$.\\ On remplace alors les coordonnées de $A(-3;-5)$ dans l'équation  $(T) : y= -5 x + p$.\\ $\begin{aligned} \phantom{\iff}&A(-3;-5)\in (T)\\ \iff& -5= -5 \times (-3)  + p\\ \iff& p=-5 +5 \times (-3)  \\ \iff& p=-5 -15   \\ \iff& p=-20\\\end{aligned}$\\On peut conclure que : $(T) : {\color[HTML]{f15929}\boldsymbol{y=-5x-20}}$. \end{minipage}

\end{Solution}

% @see : https://coopmaths.fr/alea?uuid=0e984&id=can1F15&n=1&d=10&alea=1Alo22323232323423456789101112131415161718192021222324252627282930313233343536373839404142434445464748234567891011121314151617&cd=1&cols=1
\needspace{10\baselineskip}
\begin{exercice}[BaremeTotal=false]


 La courbe représente une fonction $f$ et la droite est la tangente au point d'abscisse $0$.\\
        Déterminer $f'(0)$.\begin{center}\begin{tikzpicture}[baseline,scale = 0.5]

    \tikzset{
      point/.style={
        thick,
        draw,
        cross out,
        inner sep=0pt,
        minimum width=5pt,
        minimum height=5pt,
      },
    }
    \clip (-4,-2) rectangle (12,12);
    	
	\draw[color ={black},line width = 2,>=latex,->] (-4,0)--(12,0);
	\draw[color ={black},line width = 2,>=latex,->] (0,-2)--(0,12);
	\draw[color ={black},opacity = 0.5] (-4,1)--(12,1);
	\draw[color ={black},opacity = 0.5] (-4,2)--(12,2);
	\draw[color ={black},opacity = 0.5] (-4,3)--(12,3);
	\draw[color ={black},opacity = 0.5] (-4,4)--(12,4);
	\draw[color ={black},opacity = 0.5] (-4,5)--(12,5);
	\draw[color ={black},opacity = 0.5] (-4,6)--(12,6);
	\draw[color ={black},opacity = 0.5] (-4,7)--(12,7);
	\draw[color ={black},opacity = 0.5] (-4,8)--(12,8);
	\draw[color ={black},opacity = 0.5] (-4,9)--(12,9);
	\draw[color ={black},opacity = 0.5] (-4,10)--(12,10);
	\draw[color ={black},opacity = 0.5] (-4,11)--(12,11);
	\draw[color ={black},opacity = 0.5] (-4,12)--(12,12);
	\draw[color ={black},opacity = 0.5] (-4,-1)--(12,-1);
	\draw[color ={black},opacity = 0.5] (-4,-2)--(12,-2);
	\draw[color ={black},opacity = 0.5] (1,-2)--(1,12);
	\draw[color ={black},opacity = 0.5] (2,-2)--(2,12);
	\draw[color ={black},opacity = 0.5] (3,-2)--(3,12);
	\draw[color ={black},opacity = 0.5] (4,-2)--(4,12);
	\draw[color ={black},opacity = 0.5] (5,-2)--(5,12);
	\draw[color ={black},opacity = 0.5] (6,-2)--(6,12);
	\draw[color ={black},opacity = 0.5] (7,-2)--(7,12);
	\draw[color ={black},opacity = 0.5] (8,-2)--(8,12);
	\draw[color ={black},opacity = 0.5] (9,-2)--(9,12);
	\draw[color ={black},opacity = 0.5] (10,-2)--(10,12);
	\draw[color ={black},opacity = 0.5] (11,-2)--(11,12);
	\draw[color ={black},opacity = 0.5] (12,-2)--(12,12);
	\draw[color ={black},opacity = 0.5] (-1,-2)--(-1,12);
	\draw[color ={black},opacity = 0.5] (-2,-2)--(-2,12);
	\draw[color ={black},opacity = 0.5] (-3,-2)--(-3,12);
	\draw[color ={black},opacity = 0.5] (-4,-2)--(-4,12);
	\draw[color ={gray},opacity = 0.3] (-4,0)--(12,0);
	\draw[color ={gray},opacity = 0.3] (-4,0)--(12,0);
	\draw[color ={gray},opacity = 0.3] (0,-2)--(0,12);
	\draw[color ={gray},opacity = 0.3] (0,-2)--(0,12);
	\draw[color ={black},line width = 2] (2,-0.2)--(2,0.2);
	\draw[color ={black},line width = 2] (4,-0.2)--(4,0.2);
	\draw[color ={black},line width = 2] (6,-0.2)--(6,0.2);
	\draw[color ={black},line width = 2] (8,-0.2)--(8,0.2);
	\draw[color ={black},line width = 2] (10,-0.2)--(10,0.2);
	\draw[color ={black},line width = 2] (-2,-0.2)--(-2,0.2);
	\draw[color ={black},line width = 2] (-0.2,2)--(0.2,2);
	\draw[color ={black},line width = 2] (-0.2,4)--(0.2,4);
	\draw[color ={black},line width = 2] (-0.2,6)--(0.2,6);
	\draw[color ={black},line width = 2] (-0.2,8)--(0.2,8);
	\draw[color ={black},line width = 2] (-0.2,10)--(0.2,10);
	\draw (2,-0.4) node[anchor = center] {\scriptsize \color{black}{$1$}};
	\draw (4,-0.4) node[anchor = center] {\scriptsize \color{black}{$2$}};
	\draw (6,-0.4) node[anchor = center] {\scriptsize \color{black}{$3$}};
	\draw (8,-0.4) node[anchor = center] {\scriptsize \color{black}{$4$}};
	\draw (10,-0.4) node[anchor = center] {\scriptsize \color{black}{$5$}};
	\draw (-0.5,2.1) node[anchor = center] {\footnotesize \color{black}{$1$}};
	\draw (-0.5,4.1) node[anchor = center] {\footnotesize \color{black}{$2$}};
	\draw (-0.5,6.1) node[anchor = center] {\footnotesize \color{black}{$3$}};
	\draw (-0.5,8.1) node[anchor = center] {\footnotesize \color{black}{$4$}};
	\draw (-0.5,10.1) node[anchor = center] {\footnotesize \color{black}{$5$}};
	\draw [color={black}] (-0.3,-0.3) node[anchor = center,scale=1, rotate = 0] {O};
	
	\draw[color={blue},line width = 2] (-4,0.53)--(-3.6,0.6)--(-3.2,0.69)--(-2.8,0.79)--(-2.4,0.9)--(-2,1.03)--(-1.6,1.17)--(-1.2,1.34)--(-0.8,1.53)--(-0.4,1.75)--(0,2)--(0.4,2.29)--(0.8,2.61)--(1.2,2.98)--(1.6,3.41)--(2,3.9)--(2.4,4.45)--(2.8,5.09)--(3.2,5.81)--(3.6,6.64)--(4,7.59)--(4.4,8.67)--(4.8,9.91)--(5.2,11.32)--(5.6,12.93);
	
	\draw[color={red},line width = 2] (-4,-0.67)--(-3.6,-0.4)--(-3.2,-0.13)--(-2.8,0.13)--(-2.4,0.4)--(-2,0.67)--(-1.6,0.93)--(-1.2,1.2)--(-0.8,1.47)--(-0.4,1.73)--(0,2)--(0.4,2.27)--(0.8,2.53)--(1.2,2.8)--(1.6,3.07)--(2,3.33)--(2.4,3.6)--(2.8,3.87)--(3.2,4.13)--(3.6,4.4)--(4,4.67)--(4.4,4.93)--(4.8,5.2)--(5.2,5.47)--(5.6,5.73)--(6,6)--(6.4,6.27)--(6.8,6.53)--(7.2,6.8)--(7.6,7.07)--(8,7.33)--(8.4,7.6)--(8.8,7.87)--(9.2,8.13)--(9.6,8.4)--(10,8.67)--(10.4,8.93)--(10.8,9.2)--(11.2,9.47)--(11.6,9.73)--(12,10);

\end{tikzpicture}\end{center}
\end{exercice}

\begin{Solution}
 $f'(0)$ est donné par le coefficient directeur de la tangente à la courbe au point d'abscisse $0$, soit $\dfrac{2}{3}$.
\end{Solution}

% @see : https://coopmaths.fr/alea?uuid=6f32d&id=can1F16&n=1&d=10&alea=dqSw22323232323423456789101112131415161718192021222324252627282930313233343536373839404142434445464748234567891011121314151617&cd=1&cols=1
\needspace{10\baselineskip}
\newpage
\begin{exercice}[BaremeTotal=false]


 On donne les représentations graphiques d'une fonction et de sa dérivée.\\
        Donner l'équation réduite de la tangente à la courbe de $f$ en $x=1$. \\ \begin{minipage}{0.5\linewidth}
    \begin{tikzpicture}[baseline, scale=0.8]

    \tikzset{
      point/.style={
        thick,
        draw,
        cross out,
        inner sep=0pt,
        minimum width=5pt,
        minimum height=5pt,
      },
    }
    \clip (-4.5,-4.5) rectangle (5.75,13.5);
    	
	\draw[color ={black},>=latex,->] (-3,0)--(3,0);
	\draw[color ={black},>=latex,->] (0,-3)--(0,12);
	\draw[color ={black},opacity = 0.4] (-3,1)--(3,1);
	\draw[color ={black},opacity = 0.4] (-3,2)--(3,2);
	\draw[color ={black},opacity = 0.4] (-3,3)--(3,3);
	\draw[color ={black},opacity = 0.4] (-3,4)--(3,4);
	\draw[color ={black},opacity = 0.4] (-3,5)--(3,5);
	\draw[color ={black},opacity = 0.4] (-3,6)--(3,6);
	\draw[color ={black},opacity = 0.4] (-3,7)--(3,7);
	\draw[color ={black},opacity = 0.4] (-3,8)--(3,8);
	\draw[color ={black},opacity = 0.4] (-3,9)--(3,9);
	\draw[color ={black},opacity = 0.4] (-3,10)--(3,10);
	\draw[color ={black},opacity = 0.4] (-3,11)--(3,11);
	\draw[color ={black},opacity = 0.4] (-3,12)--(3,12);
	\draw[color ={black},opacity = 0.4] (-3,-1)--(3,-1);
	\draw[color ={black},opacity = 0.4] (-3,-2)--(3,-2);
	\draw[color ={black},opacity = 0.4] (-3,-3)--(3,-3);
	\draw[color ={black},opacity = 0.4] (1,-3)--(1,12);
	\draw[color ={black},opacity = 0.4] (2,-3)--(2,12);
	\draw[color ={black},opacity = 0.4] (3,-3)--(3,12);
	\draw[color ={black},opacity = 0.4] (-1,-3)--(-1,12);
	\draw[color ={black},opacity = 0.4] (-2,-3)--(-2,12);
	\draw[color ={black},opacity = 0.4] (-3,-3)--(-3,12);
	\draw[color ={gray},opacity = 0.3] (-3,0)--(3,0);
	\draw[color ={gray},opacity = 0.3] (-3,0.5)--(3,0.5);
	\draw[color ={gray},opacity = 0.3] (-3,1.5)--(3,1.5);
	\draw[color ={gray},opacity = 0.3] (-3,2.5)--(3,2.5);
	\draw[color ={gray},opacity = 0.3] (-3,3.5)--(3,3.5);
	\draw[color ={gray},opacity = 0.3] (-3,4.5)--(3,4.5);
	\draw[color ={gray},opacity = 0.3] (-3,5.5)--(3,5.5);
	\draw[color ={gray},opacity = 0.3] (-3,6.5)--(3,6.5);
	\draw[color ={gray},opacity = 0.3] (-3,7.5)--(3,7.5);
	\draw[color ={gray},opacity = 0.3] (-3,8.5)--(3,8.5);
	\draw[color ={gray},opacity = 0.3] (-3,9.5)--(3,9.5);
	\draw[color ={gray},opacity = 0.3] (-3,10.5)--(3,10.5);
	\draw[color ={gray},opacity = 0.3] (-3,11.5)--(3,11.5);
	\draw[color ={gray},opacity = 0.3] (-3,0)--(3,0);
	\draw[color ={gray},opacity = 0.3] (-3,-0.5)--(3,-0.5);
	\draw[color ={gray},opacity = 0.3] (-3,-1.5)--(3,-1.5);
	\draw[color ={gray},opacity = 0.3] (-3,-2.5)--(3,-2.5);
	\draw[color ={gray},opacity = 0.3] (0,-3)--(0,12);
	\draw[color ={gray},opacity = 0.3] (0,-3)--(0,12);
	\draw[color ={black},line width = 1.2] (1,-0.13)--(1,0.13);
	\draw[color ={black},line width = 1.2] (2,-0.13)--(2,0.13);
	\draw[color ={black},line width = 1.2] (3,-0.13)--(3,0.13);
	\draw[color ={black},line width = 1.2] (-1,-0.13)--(-1,0.13);
	\draw[color ={black},line width = 1.2] (-2,-0.13)--(-2,0.13);
	\draw[color ={black},line width = 1.2] (-3,-0.13)--(-3,0.13);
	\draw[color ={black},line width = 1.2] (-0.13,1)--(0.13,1);
	\draw[color ={black},line width = 1.2] (-0.13,2)--(0.13,2);
	\draw[color ={black},line width = 1.2] (-0.13,3)--(0.13,3);
	\draw[color ={black},line width = 1.2] (-0.13,4)--(0.13,4);
	\draw[color ={black},line width = 1.2] (-0.13,5)--(0.13,5);
	\draw[color ={black},line width = 1.2] (-0.13,6)--(0.13,6);
	\draw[color ={black},line width = 1.2] (-0.13,7)--(0.13,7);
	\draw[color ={black},line width = 1.2] (-0.13,8)--(0.13,8);
	\draw[color ={black},line width = 1.2] (-0.13,9)--(0.13,9);
	\draw[color ={black},line width = 1.2] (-0.13,10)--(0.13,10);
	\draw[color ={black},line width = 1.2] (-0.13,11)--(0.13,11);
	\draw[color ={black},line width = 1.2] (-0.13,12)--(0.13,12);
	\draw[color ={black},line width = 1.2] (-0.13,-1)--(0.13,-1);
	\draw[color ={black},line width = 1.2] (-0.13,-2)--(0.13,-2);
	\draw[color ={black},line width = 1.2] (-0.13,-3)--(0.13,-3);
	\draw (3,-0.4) node[anchor = center] {\scriptsize \color{black}{$3$}};
	\draw (-3,-0.4) node[anchor = center] {\scriptsize \color{black}{$-3$}};
	\draw (-0.5,3.1) node[anchor = center] {\footnotesize \color{black}{$3$}};
	\draw (-0.5,6.1) node[anchor = center] {\footnotesize \color{black}{$6$}};
	\draw (-0.5,9.1) node[anchor = center] {\footnotesize \color{black}{$9$}};
	\draw (-0.5,-2.9) node[anchor = center] {\footnotesize \color{black}{$-3$}};
	\draw [color={black}] (-0.3,-0.3) node[anchor = center,scale=1, rotate = 0] {O};
	\draw (3,10) node[anchor = center] {\footnotesize \color{blue}{$\Large \mathcal{C}_f$}};
	
	\draw[color={blue},line width = 2] (-3,1)--(-2.8,0.64)--(-2.6,0.36)--(-2.4,0.16)--(-2.2,0.04)--(-2,0)--(-1.8,0.04)--(-1.6,0.16)--(-1.4,0.36)--(-1.2,0.64)--(-1,1)--(-0.8,1.44)--(-0.6,1.96)--(-0.4,2.56)--(-0.2,3.24)--(0,4)--(0.2,4.84)--(0.4,5.76)--(0.6,6.76)--(0.8,7.84)--(1,9)--(1.2,10.24)--(1.4,11.56)--(1.6,12.96);

\end{tikzpicture}\\
    \end{minipage}
    \begin{minipage}{0.5\linewidth}
    \begin{tikzpicture}[baseline, scale=0.8]

    \tikzset{
      point/.style={
        thick,
        draw,
        cross out,
        inner sep=0pt,
        minimum width=5pt,
        minimum height=5pt,
      },
    }
    \clip (-4.5,-6.5) rectangle (6.5,9.5);
    	
	\draw[color ={black},>=latex,->] (-3,0)--(3,0);
	\draw[color ={black},>=latex,->] (0,-5)--(0,8);
	\draw[color ={black},opacity = 0.4] (-3,1)--(3,1);
	\draw[color ={black},opacity = 0.4] (-3,2)--(3,2);
	\draw[color ={black},opacity = 0.4] (-3,3)--(3,3);
	\draw[color ={black},opacity = 0.4] (-3,4)--(3,4);
	\draw[color ={black},opacity = 0.4] (-3,5)--(3,5);
	\draw[color ={black},opacity = 0.4] (-3,6)--(3,6);
	\draw[color ={black},opacity = 0.4] (-3,7)--(3,7);
	\draw[color ={black},opacity = 0.4] (-3,8)--(3,8);
	\draw[color ={black},opacity = 0.4] (-3,-1)--(3,-1);
	\draw[color ={black},opacity = 0.4] (-3,-2)--(3,-2);
	\draw[color ={black},opacity = 0.4] (-3,-3)--(3,-3);
	\draw[color ={black},opacity = 0.4] (-3,-4)--(3,-4);
	\draw[color ={black},opacity = 0.4] (-3,-5)--(3,-5);
	\draw[color ={black},opacity = 0.4] (1,-5)--(1,8);
	\draw[color ={black},opacity = 0.4] (2,-5)--(2,8);
	\draw[color ={black},opacity = 0.4] (3,-5)--(3,8);
	\draw[color ={black},opacity = 0.4] (-1,-5)--(-1,8);
	\draw[color ={black},opacity = 0.4] (-2,-5)--(-2,8);
	\draw[color ={black},opacity = 0.4] (-3,-5)--(-3,8);
	\draw[color ={gray},opacity = 0.3] (-3,0)--(3,0);
	\draw[color ={gray},opacity = 0.3] (-3,0.5)--(3,0.5);
	\draw[color ={gray},opacity = 0.3] (-3,1.5)--(3,1.5);
	\draw[color ={gray},opacity = 0.3] (-3,2.5)--(3,2.5);
	\draw[color ={gray},opacity = 0.3] (-3,3.5)--(3,3.5);
	\draw[color ={gray},opacity = 0.3] (-3,4.5)--(3,4.5);
	\draw[color ={gray},opacity = 0.3] (-3,5.5)--(3,5.5);
	\draw[color ={gray},opacity = 0.3] (-3,6.5)--(3,6.5);
	\draw[color ={gray},opacity = 0.3] (-3,7.5)--(3,7.5);
	\draw[color ={gray},opacity = 0.3] (-3,0)--(3,0);
	\draw[color ={gray},opacity = 0.3] (-3,-0.5)--(3,-0.5);
	\draw[color ={gray},opacity = 0.3] (-3,-1.5)--(3,-1.5);
	\draw[color ={gray},opacity = 0.3] (-3,-2.5)--(3,-2.5);
	\draw[color ={gray},opacity = 0.3] (-3,-3.5)--(3,-3.5);
	\draw[color ={gray},opacity = 0.3] (-3,-4.5)--(3,-4.5);
	\draw[color ={gray},opacity = 0.3] (0,-5)--(0,8);
	\draw[color ={gray},opacity = 0.3] (0,-5)--(0,8);
	\draw[color ={black},line width = 1.2] (1,-0.13)--(1,0.13);
	\draw[color ={black},line width = 1.2] (2,-0.13)--(2,0.13);
	\draw[color ={black},line width = 1.2] (3,-0.13)--(3,0.13);
	\draw[color ={black},line width = 1.2] (-1,-0.13)--(-1,0.13);
	\draw[color ={black},line width = 1.2] (-2,-0.13)--(-2,0.13);
	\draw[color ={black},line width = 1.2] (-3,-0.13)--(-3,0.13);
	\draw[color ={black},line width = 1.2] (-0.13,1)--(0.13,1);
	\draw[color ={black},line width = 1.2] (-0.13,2)--(0.13,2);
	\draw[color ={black},line width = 1.2] (-0.13,3)--(0.13,3);
	\draw[color ={black},line width = 1.2] (-0.13,4)--(0.13,4);
	\draw[color ={black},line width = 1.2] (-0.13,5)--(0.13,5);
	\draw[color ={black},line width = 1.2] (-0.13,6)--(0.13,6);
	\draw[color ={black},line width = 1.2] (-0.13,7)--(0.13,7);
	\draw[color ={black},line width = 1.2] (-0.13,8)--(0.13,8);
	\draw[color ={black},line width = 1.2] (-0.13,9)--(0.13,9);
	\draw[color ={black},line width = 1.2] (-0.13,10)--(0.13,10);
	\draw[color ={black},line width = 1.2] (-0.13,11)--(0.13,11);
	\draw[color ={black},line width = 1.2] (-0.13,12)--(0.13,12);
	\draw[color ={black},line width = 1.2] (-0.13,-1)--(0.13,-1);
	\draw[color ={black},line width = 1.2] (-0.13,-2)--(0.13,-2);
	\draw[color ={black},line width = 1.2] (-0.13,-3)--(0.13,-3);
	\draw (3,-0.4) node[anchor = center] {\scriptsize \color{black}{$3$}};
	\draw (-3,-0.4) node[anchor = center] {\scriptsize \color{black}{$-3$}};
	\draw (-0.5,3.1) node[anchor = center] {\footnotesize \color{black}{$3$}};
	\draw (-0.5,6.1) node[anchor = center] {\footnotesize \color{black}{$6$}};
	\draw (-0.5,-2.9) node[anchor = center] {\footnotesize \color{black}{$-3$}};
	\draw [color={black}] (-0.3,-0.3) node[anchor = center,scale=1, rotate = 0] {O};
	\draw (3,6) node[anchor = center] {\footnotesize \color{red}{$\Large\mathcal{C}_f\prime$}};
	
	\draw[color={red},line width = 2] (-3,-2)--(-2.8,-1.6)--(-2.6,-1.2)--(-2.4,-0.8)--(-2.2,-0.4)--(-2,0)--(-1.8,0.4)--(-1.6,0.8)--(-1.4,1.2)--(-1.2,1.6)--(-1,2)--(-0.8,2.4)--(-0.6,2.8)--(-0.4,3.2)--(-0.2,3.6)--(0,4)--(0.2,4.4)--(0.4,4.8)--(0.6,5.2)--(0.8,5.6)--(1,6)--(1.2,6.4)--(1.4,6.8)--(1.6,7.2)--(1.8,7.6)--(2,8)--(2.2,8.4)--(2.4,8.8);

\end{tikzpicture}\\
    \end{minipage}
    
\end{exercice}

\begin{Solution}
 L'équation réduite de la tangente au point d'abscisse $1$ est  : $y=f'(1)(x-1)+f(1)$.\\
        On lit graphiquement $f(1)=9$ et $f'(1)=6$.\\
        L'équation réduite de la tangente est donc donnée par :
        $y=6(x-1)+9$, soit $y=6x+3$.
\end{Solution}

% @see : https://coopmaths.fr/alea?uuid=a1ba2&id=can1F14&n=1&d=10&alea=rkIp22323232323423456789101112131415161718192021222324252627282930313233343536373839404142434445464748234567891011121314151617&cd=1&cols=1
\needspace{10\baselineskip}
\begin{exercice}[BaremeTotal=false]


Soit $f$ la fonction définie sur $\mathbb{R}$ par :
$f(x)= 5x+6-2x^2$.\\
En calculant la limite du taux d'accroissement, déterminer $f'(-1)$.
\end{exercice}

\begin{Solution}
 $f$ est une fonction polynôme du second degré de la forme $f(x)=ax^2+bx+c$.\\
    La fonction dérivée est donnée par la somme des dérivées des fonctions $u$ et $v$ définies par $u(x)=-2x^2$ et $v(x)=5x+6$.\\
     Comme $u'(x)=-4x$ et $v'(x)=5$, on obtient  $f'(x)=-4x+5$. \\
     Ainsi, $f'(-1)=-4\times (-1)+5=9$.
\end{Solution}

% @see : https://coopmaths.fr/alea?uuid=3c690&id=can1F13&n=1&d=10&alea=EUDu22323232323423456789101112131415161718192021222324252627282930313233343536373839404142434445464748234567891011121314151617&cd=1&cols=1
\needspace{10\baselineskip}
\begin{exercice}[BaremeTotal=false]


 Déterminer le coefficient directeur de la tangente à la courbe représentative de la fonction carré au point d'abscisse $10$.    
\end{exercice}

\begin{Solution}
 Le coefficient directeur de la tangente au point d'abscisse $a$ est donné par le nombre dérivé $f'(a)$.\\
        La fonction $f$ définie par $f(x)=x^2$ a pour fonction dérivée la fonction $f'$ définie par $f'(x)=2x$.\\
        Comme $f'(10)=2\times 10=20$, le coefficient directeur de la tangente au point d'abscisse $10$ est : $20$. 
\end{Solution}

% @see : https://coopmaths.fr/alea?uuid=ebd89&id=1AN14-1&n=1&d=10&s=3&alea=ocYr22323232323423456789101112131415161718192021222324252627282930313233343536373839404142434445464748234567891011121314151617&cd=0&cols=1
\needspace{10\baselineskip}
\begin{exercice}[BaremeTotal=false]


 Donner la dérivée de la fonction $f$, dérivable pour tout $x\in\R$, définie par  $f(x)=\dfrac{-x+9}{9}$.
\end{exercice}

\begin{Solution}
 L'expression de la dérivée de la fonction $f$ définie par $f(x)=\dfrac{-x+9}{9}$ est : ${\color[HTML]{f15929}\boldsymbol{f'(x)=-\dfrac{1}{9}}}$.
\end{Solution}

% @see : https://coopmaths.fr/alea?uuid=ebd89&id=1AN14-1&n=1&d=10&s=4&alea=uAlP22323232323423456789101112131415161718192021222324252627282930313233343536373839404142434445464748234567891011121314151617&cd=0&cols=1
\needspace{10\baselineskip}
\begin{exercice}[BaremeTotal=false]


 Donner la dérivée de la fonction $f$, dérivable pour tout $x\in\R$, définie par  $f(x)=-10x^{5}$.
\end{exercice}

\begin{Solution}
 L'expression de la dérivée de la fonction $f$ définie par $f(x)=-10x^{5}$ est : ${\color[HTML]{f15929}\boldsymbol{f'(x)=-50x^{4}}}$.
\end{Solution}

% @see : https://coopmaths.fr/alea?uuid=ebd89&id=1AN14-1&n=1&d=10&s=7&alea=vjN722323232323423456789101112131415161718192021222324252627282930313233343536373839404142434445464748234567891011121314151617&cd=0&cols=1
\needspace{10\baselineskip}
\begin{exercice}[BaremeTotal=false]


 Donner la dérivée de la fonction $f$, dérivable pour tout $x\in\R^*$, définie par  $f(x)=\dfrac{-2}{5x}$.
\end{exercice}

\begin{Solution}
 L'expression de la dérivée de la fonction $f$ définie par $f(x)=\dfrac{-2}{5x}$ est : ${\color[HTML]{f15929}\boldsymbol{f'(x)=\dfrac{2}{5x^2}}}$.
\end{Solution}

% @see : https://coopmaths.fr/alea?uuid=a83c0&id=1AN14-4&n=1&d=10&s=1&alea=Njr322323232323423456789101112131415161718192021222324252627282930313233343536373839404142434445464748234567891011121314151617&cd=0&cols=1
\needspace{10\baselineskip}
\begin{exercice}[BaremeTotal=false]


 Donner l'expression de la dérivée de la fonction $f$ définie sur $\R^{*}$ par $f(x)=-4x^{2}+\dfrac{8}{x}$.\\
\end{exercice}

\begin{Solution}
 L'expression de la dérivée de la fonction $f$ définie par $f(x)=-4x^{2}+\dfrac{8}{x}$ est :\\$f'(x)={\color[HTML]{f15929}\boldsymbol{-8x-\dfrac{8}{x^2}}}$.\\
\end{Solution}

% @see : https://coopmaths.fr/alea?uuid=a83c0&id=1AN14-4&n=1&d=10&s=4&alea=4Vpz22323232323423456789101112131415161718192021222324252627282930313233343536373839404142434445464748234567891011121314151617&cd=1&cols=1
\needspace{10\baselineskip}
\begin{exercice}[BaremeTotal=false]


 Donner l'expression de la dérivée de la fonction $f$ définie sur $\R_+^{*}$ par $f(x)=-\dfrac{2}{x}+5\sqrt{x}+5x^{2}+4x-3$.\\
\end{exercice}

\begin{Solution}
 $(-\dfrac{2}{x})^\prime=\dfrac{2}{x^2}$\\$(5\sqrt{x})^\prime=\dfrac{5}{2\sqrt{x}}$\\$(5x^{2}+4x-3)^\prime=10x+4$\\L'expression de la dérivée de la fonction $f$ définie par $f(x)=-\dfrac{2}{x}+5\sqrt{x}+5x^{2}+4x-3$ est :\\$f'(x)={\color[HTML]{f15929}\boldsymbol{\dfrac{2}{x^2}+\dfrac{5}{2\sqrt{x}}+10x+4}}$.\\
\end{Solution}

\end{Maquette}
\clearpage
\begin{Maquette}[DM]{Niveau= ,Classe=\getdata[18]\mydata\ (v18), Date = 06/01/25  ,Theme=Exercices}


% @see : https://coopmaths.fr/alea?uuid=4c8c7&id=1AN11-3&n=1&d=10&s=2&alea=nawO2232323232342345678910111213141516171819202122232425262728293031323334353637383940414243444546474823456789101112131415161718&cd=1&cols=2
\needspace{10\baselineskip}
\begin{exercice}[BaremeTotal=false]
\label{eleve:18}


 \begin{minipage}[t]{\linewidth} Soit $f$ une fonction dérivable sur $[-5;5]$ et $\mathcal{C}_f$ sa courbe représentative.\\On sait que  $f(-1)=5~~$ et que $~~f'(-1)=3$.\\Déterminer une équation de la tangente $(T)$ à la courbe $\mathcal{C}_f$ au point d'abscisse $-1$,\\sans utiliser la formule de cours de l'équation de tangente. \end{minipage}

\end{exercice}

\begin{Solution}
 \begin{minipage}[t]{\linewidth} On sait que la tangente n'est pas une droite verticale, puisque la fonction est dérivable sur l'intervalle.\\On en déduit que la tangente $(T)$ au point d'abscisse $-1$, admet une équation réduite de la forme :  $(T) : y=m x + p$. 

\medskip
$\bullet$ Détermination de $m$ :\\On sait que le nombre dérivé en $-1$ est par définition, le coefficient directeur de la tangente au point d'abscisse $-1$.\\Par conséquent, on a déjà : $m=f'(-1)=3$.\\ On en déduit que  $(T) : y= 3 x + p$\\$\bullet$ Détermination de $p$ :\\ Pour cela, on utilise que si $f(-1)=5~~$, alors le point $A$ de coordonnées $(-1;5)$ appartient à $\mathcal{C}_f$ mais aussi à $(T)$.\\ On peut écrire $A(-1;5) = \mathcal{C}_f \cap (T)$.\\ On remplace alors les coordonnées de $A(-1;5)$ dans l'équation  $(T) : y= 3 x + p$.\\ $\begin{aligned} \phantom{\iff}&A(-1;5)\in (T)\\ \iff& 5= 3 \times (-1)  + p\\ \iff& p=5 -3 \times (-1)  \\ \iff& p=5 +3   \\ \iff& p=8\\\end{aligned}$\\On peut conclure que : $(T) : {\color[HTML]{f15929}\boldsymbol{y=3x+8}}$. \end{minipage}

\end{Solution}

% @see : https://coopmaths.fr/alea?uuid=0e984&id=can1F15&n=1&d=10&alea=1Alo2232323232342345678910111213141516171819202122232425262728293031323334353637383940414243444546474823456789101112131415161718&cd=1&cols=1
\needspace{10\baselineskip}
\begin{exercice}[BaremeTotal=false]


 La courbe représente une fonction $f$ et la droite est la tangente au point d'abscisse $0$.\\
        Déterminer $f'(0)$.\begin{center}\begin{tikzpicture}[baseline,scale = 0.5]

    \tikzset{
      point/.style={
        thick,
        draw,
        cross out,
        inner sep=0pt,
        minimum width=5pt,
        minimum height=5pt,
      },
    }
    \clip (-10,-2) rectangle (4,12);
    	
	\draw[color ={black},line width = 2,>=latex,->] (-10,0)--(4,0);
	\draw[color ={black},line width = 2,>=latex,->] (0,-2)--(0,12);
	\draw[color ={black},opacity = 0.5] (-10,1)--(4,1);
	\draw[color ={black},opacity = 0.5] (-10,2)--(4,2);
	\draw[color ={black},opacity = 0.5] (-10,3)--(4,3);
	\draw[color ={black},opacity = 0.5] (-10,4)--(4,4);
	\draw[color ={black},opacity = 0.5] (-10,5)--(4,5);
	\draw[color ={black},opacity = 0.5] (-10,6)--(4,6);
	\draw[color ={black},opacity = 0.5] (-10,7)--(4,7);
	\draw[color ={black},opacity = 0.5] (-10,8)--(4,8);
	\draw[color ={black},opacity = 0.5] (-10,9)--(4,9);
	\draw[color ={black},opacity = 0.5] (-10,10)--(4,10);
	\draw[color ={black},opacity = 0.5] (-10,11)--(4,11);
	\draw[color ={black},opacity = 0.5] (-10,12)--(4,12);
	\draw[color ={black},opacity = 0.5] (-10,-1)--(4,-1);
	\draw[color ={black},opacity = 0.5] (-10,-2)--(4,-2);
	\draw[color ={black},opacity = 0.5] (1,-2)--(1,12);
	\draw[color ={black},opacity = 0.5] (2,-2)--(2,12);
	\draw[color ={black},opacity = 0.5] (3,-2)--(3,12);
	\draw[color ={black},opacity = 0.5] (4,-2)--(4,12);
	\draw[color ={black},opacity = 0.5] (-1,-2)--(-1,12);
	\draw[color ={black},opacity = 0.5] (-2,-2)--(-2,12);
	\draw[color ={black},opacity = 0.5] (-3,-2)--(-3,12);
	\draw[color ={black},opacity = 0.5] (-4,-2)--(-4,12);
	\draw[color ={black},opacity = 0.5] (-5,-2)--(-5,12);
	\draw[color ={black},opacity = 0.5] (-6,-2)--(-6,12);
	\draw[color ={black},opacity = 0.5] (-7,-2)--(-7,12);
	\draw[color ={black},opacity = 0.5] (-8,-2)--(-8,12);
	\draw[color ={black},opacity = 0.5] (-9,-2)--(-9,12);
	\draw[color ={black},opacity = 0.5] (-10,-2)--(-10,12);
	\draw[color ={gray},opacity = 0.3] (-10,0)--(4,0);
	\draw[color ={gray},opacity = 0.3] (-10,0)--(4,0);
	\draw[color ={gray},opacity = 0.3] (0,-2)--(0,12);
	\draw[color ={gray},opacity = 0.3] (0,-2)--(0,12);
	\draw[color ={gray},opacity = 0.3] (-5,-2)--(-5,12);
	\draw[color ={gray},opacity = 0.3] (-6,-2)--(-6,12);
	\draw[color ={gray},opacity = 0.3] (-7,-2)--(-7,12);
	\draw[color ={gray},opacity = 0.3] (-8,-2)--(-8,12);
	\draw[color ={gray},opacity = 0.3] (-9,-2)--(-9,12);
	\draw[color ={gray},opacity = 0.3] (-10,-2)--(-10,12);
	\draw[color ={black},line width = 2] (2,-0.2)--(2,0.2);
	\draw[color ={black},line width = 2] (-2,-0.2)--(-2,0.2);
	\draw[color ={black},line width = 2] (-4,-0.2)--(-4,0.2);
	\draw[color ={black},line width = 2] (-6,-0.2)--(-6,0.2);
	\draw[color ={black},line width = 2] (-8,-0.2)--(-8,0.2);
	\draw[color ={black},line width = 2] (-0.2,2)--(0.2,2);
	\draw[color ={black},line width = 2] (-0.2,4)--(0.2,4);
	\draw[color ={black},line width = 2] (-0.2,6)--(0.2,6);
	\draw[color ={black},line width = 2] (-0.2,8)--(0.2,8);
	\draw[color ={black},line width = 2] (-0.2,10)--(0.2,10);
	\draw (2,-0.4) node[anchor = center] {\scriptsize \color{black}{$1$}};
	\draw (-2,-0.4) node[anchor = center] {\scriptsize \color{black}{$-1$}};
	\draw (-4,-0.4) node[anchor = center] {\scriptsize \color{black}{$-2$}};
	\draw (-6,-0.4) node[anchor = center] {\scriptsize \color{black}{$-3$}};
	\draw (-8,-0.4) node[anchor = center] {\scriptsize \color{black}{$-4$}};
	\draw (-0.5,2.1) node[anchor = center] {\footnotesize \color{black}{$1$}};
	\draw (-0.5,4.1) node[anchor = center] {\footnotesize \color{black}{$2$}};
	\draw (-0.5,6.1) node[anchor = center] {\footnotesize \color{black}{$3$}};
	\draw (-0.5,8.1) node[anchor = center] {\footnotesize \color{black}{$4$}};
	\draw (-0.5,10.1) node[anchor = center] {\footnotesize \color{black}{$5$}};
	\draw [color={black}] (-0.3,-0.3) node[anchor = center,scale=1, rotate = 0] {O};
	
	\draw[color={blue},line width = 2] (-2.4,12.1)--(-2,8.96)--(-1.6,6.64)--(-1.2,4.92)--(-0.8,3.64)--(-0.4,2.7)--(0,2)--(0.4,1.48)--(0.8,1.1)--(1.2,0.81)--(1.6,0.6)--(2,0.45)--(2.4,0.33)--(2.8,0.24)--(3.2,0.18)--(3.6,0.13)--(4,0.1);
	
	\draw[color={red},line width = 2] (-8,14)--(-7.6,13.4)--(-7.2,12.8)--(-6.8,12.2)--(-6.4,11.6)--(-6,11)--(-5.6,10.4)--(-5.2,9.8)--(-4.8,9.2)--(-4.4,8.6)--(-4,8)--(-3.6,7.4)--(-3.2,6.8)--(-2.8,6.2)--(-2.4,5.6)--(-2,5)--(-1.6,4.4)--(-1.2,3.8)--(-0.8,3.2)--(-0.4,2.6)--(0,2)--(0.4,1.4)--(0.8,0.8)--(1.2,0.2)--(1.6,-0.4)--(2,-1)--(2.4,-1.6)--(2.8,-2.2)--(3.2,-2.8)--(3.6,-3.4);

\end{tikzpicture}\end{center}
\end{exercice}

\begin{Solution}
 $f'(0)$ est donné par le coefficient directeur de la tangente à la courbe au point d'abscisse $0$, soit $\dfrac{-3}{2}=-\dfrac{3}{2}$.
\end{Solution}

% @see : https://coopmaths.fr/alea?uuid=6f32d&id=can1F16&n=1&d=10&alea=dqSw2232323232342345678910111213141516171819202122232425262728293031323334353637383940414243444546474823456789101112131415161718&cd=1&cols=1
\needspace{10\baselineskip}
\newpage
\begin{exercice}[BaremeTotal=false]


 On donne les représentations graphiques d'une fonction et de sa dérivée.\\
        Donner l'équation réduite de la tangente à la courbe de $f$ en $x=-1$. \\ \begin{minipage}{0.5\linewidth}
    \begin{tikzpicture}[baseline, scale=0.8]

    \tikzset{
      point/.style={
        thick,
        draw,
        cross out,
        inner sep=0pt,
        minimum width=5pt,
        minimum height=5pt,
      },
    }
    \clip (-4.5,-4.5) rectangle (5.75,13.5);
    	
	\draw[color ={black},>=latex,->] (-3,0)--(3,0);
	\draw[color ={black},>=latex,->] (0,-3)--(0,12);
	\draw[color ={black},opacity = 0.4] (-3,1)--(3,1);
	\draw[color ={black},opacity = 0.4] (-3,2)--(3,2);
	\draw[color ={black},opacity = 0.4] (-3,3)--(3,3);
	\draw[color ={black},opacity = 0.4] (-3,4)--(3,4);
	\draw[color ={black},opacity = 0.4] (-3,5)--(3,5);
	\draw[color ={black},opacity = 0.4] (-3,6)--(3,6);
	\draw[color ={black},opacity = 0.4] (-3,7)--(3,7);
	\draw[color ={black},opacity = 0.4] (-3,8)--(3,8);
	\draw[color ={black},opacity = 0.4] (-3,9)--(3,9);
	\draw[color ={black},opacity = 0.4] (-3,10)--(3,10);
	\draw[color ={black},opacity = 0.4] (-3,11)--(3,11);
	\draw[color ={black},opacity = 0.4] (-3,12)--(3,12);
	\draw[color ={black},opacity = 0.4] (-3,-1)--(3,-1);
	\draw[color ={black},opacity = 0.4] (-3,-2)--(3,-2);
	\draw[color ={black},opacity = 0.4] (-3,-3)--(3,-3);
	\draw[color ={black},opacity = 0.4] (1,-3)--(1,12);
	\draw[color ={black},opacity = 0.4] (2,-3)--(2,12);
	\draw[color ={black},opacity = 0.4] (3,-3)--(3,12);
	\draw[color ={black},opacity = 0.4] (-1,-3)--(-1,12);
	\draw[color ={black},opacity = 0.4] (-2,-3)--(-2,12);
	\draw[color ={black},opacity = 0.4] (-3,-3)--(-3,12);
	\draw[color ={gray},opacity = 0.3] (-3,0)--(3,0);
	\draw[color ={gray},opacity = 0.3] (-3,0.5)--(3,0.5);
	\draw[color ={gray},opacity = 0.3] (-3,1.5)--(3,1.5);
	\draw[color ={gray},opacity = 0.3] (-3,2.5)--(3,2.5);
	\draw[color ={gray},opacity = 0.3] (-3,3.5)--(3,3.5);
	\draw[color ={gray},opacity = 0.3] (-3,4.5)--(3,4.5);
	\draw[color ={gray},opacity = 0.3] (-3,5.5)--(3,5.5);
	\draw[color ={gray},opacity = 0.3] (-3,6.5)--(3,6.5);
	\draw[color ={gray},opacity = 0.3] (-3,7.5)--(3,7.5);
	\draw[color ={gray},opacity = 0.3] (-3,8.5)--(3,8.5);
	\draw[color ={gray},opacity = 0.3] (-3,9.5)--(3,9.5);
	\draw[color ={gray},opacity = 0.3] (-3,10.5)--(3,10.5);
	\draw[color ={gray},opacity = 0.3] (-3,11.5)--(3,11.5);
	\draw[color ={gray},opacity = 0.3] (-3,0)--(3,0);
	\draw[color ={gray},opacity = 0.3] (-3,-0.5)--(3,-0.5);
	\draw[color ={gray},opacity = 0.3] (-3,-1.5)--(3,-1.5);
	\draw[color ={gray},opacity = 0.3] (-3,-2.5)--(3,-2.5);
	\draw[color ={gray},opacity = 0.3] (0,-3)--(0,12);
	\draw[color ={gray},opacity = 0.3] (0,-3)--(0,12);
	\draw[color ={black},line width = 1.2] (1,-0.13)--(1,0.13);
	\draw[color ={black},line width = 1.2] (2,-0.13)--(2,0.13);
	\draw[color ={black},line width = 1.2] (3,-0.13)--(3,0.13);
	\draw[color ={black},line width = 1.2] (-1,-0.13)--(-1,0.13);
	\draw[color ={black},line width = 1.2] (-2,-0.13)--(-2,0.13);
	\draw[color ={black},line width = 1.2] (-3,-0.13)--(-3,0.13);
	\draw[color ={black},line width = 1.2] (-0.13,1)--(0.13,1);
	\draw[color ={black},line width = 1.2] (-0.13,2)--(0.13,2);
	\draw[color ={black},line width = 1.2] (-0.13,3)--(0.13,3);
	\draw[color ={black},line width = 1.2] (-0.13,4)--(0.13,4);
	\draw[color ={black},line width = 1.2] (-0.13,5)--(0.13,5);
	\draw[color ={black},line width = 1.2] (-0.13,6)--(0.13,6);
	\draw[color ={black},line width = 1.2] (-0.13,7)--(0.13,7);
	\draw[color ={black},line width = 1.2] (-0.13,8)--(0.13,8);
	\draw[color ={black},line width = 1.2] (-0.13,9)--(0.13,9);
	\draw[color ={black},line width = 1.2] (-0.13,10)--(0.13,10);
	\draw[color ={black},line width = 1.2] (-0.13,11)--(0.13,11);
	\draw[color ={black},line width = 1.2] (-0.13,12)--(0.13,12);
	\draw[color ={black},line width = 1.2] (-0.13,-1)--(0.13,-1);
	\draw[color ={black},line width = 1.2] (-0.13,-2)--(0.13,-2);
	\draw[color ={black},line width = 1.2] (-0.13,-3)--(0.13,-3);
	\draw (3,-0.4) node[anchor = center] {\scriptsize \color{black}{$3$}};
	\draw (-3,-0.4) node[anchor = center] {\scriptsize \color{black}{$-3$}};
	\draw (-0.5,3.1) node[anchor = center] {\footnotesize \color{black}{$3$}};
	\draw (-0.5,6.1) node[anchor = center] {\footnotesize \color{black}{$6$}};
	\draw (-0.5,9.1) node[anchor = center] {\footnotesize \color{black}{$9$}};
	\draw (-0.5,-2.9) node[anchor = center] {\footnotesize \color{black}{$-3$}};
	\draw [color={black}] (-0.3,-0.3) node[anchor = center,scale=1, rotate = 0] {O};
	\draw (3,10) node[anchor = center] {\footnotesize \color{blue}{$\Large \mathcal{C}_f$}};
	
	\draw[color={blue},line width = 2] (-3,8)--(-2.8,6.84)--(-2.6,5.76)--(-2.4,4.76)--(-2.2,3.84)--(-2,3)--(-1.8,2.24)--(-1.6,1.56)--(-1.4,0.96)--(-1.2,0.44)--(-1,0)--(-0.8,-0.36)--(-0.6,-0.64)--(-0.4,-0.84)--(-0.2,-0.96)--(0,-1)--(0.2,-0.96)--(0.4,-0.84)--(0.6,-0.64)--(0.8,-0.36)--(1,0)--(1.2,0.44)--(1.4,0.96)--(1.6,1.56)--(1.8,2.24)--(2,3)--(2.2,3.84)--(2.4,4.76)--(2.6,5.76)--(2.8,6.84)--(3,8);

\end{tikzpicture}\\
    \end{minipage}
    \begin{minipage}{0.5\linewidth}
    \begin{tikzpicture}[baseline, scale=0.8]

    \tikzset{
      point/.style={
        thick,
        draw,
        cross out,
        inner sep=0pt,
        minimum width=5pt,
        minimum height=5pt,
      },
    }
    \clip (-4.5,-6.5) rectangle (6.5,9.5);
    	
	\draw[color ={black},>=latex,->] (-3,0)--(3,0);
	\draw[color ={black},>=latex,->] (0,-5)--(0,8);
	\draw[color ={black},opacity = 0.4] (-3,1)--(3,1);
	\draw[color ={black},opacity = 0.4] (-3,2)--(3,2);
	\draw[color ={black},opacity = 0.4] (-3,3)--(3,3);
	\draw[color ={black},opacity = 0.4] (-3,4)--(3,4);
	\draw[color ={black},opacity = 0.4] (-3,5)--(3,5);
	\draw[color ={black},opacity = 0.4] (-3,6)--(3,6);
	\draw[color ={black},opacity = 0.4] (-3,7)--(3,7);
	\draw[color ={black},opacity = 0.4] (-3,8)--(3,8);
	\draw[color ={black},opacity = 0.4] (-3,-1)--(3,-1);
	\draw[color ={black},opacity = 0.4] (-3,-2)--(3,-2);
	\draw[color ={black},opacity = 0.4] (-3,-3)--(3,-3);
	\draw[color ={black},opacity = 0.4] (-3,-4)--(3,-4);
	\draw[color ={black},opacity = 0.4] (-3,-5)--(3,-5);
	\draw[color ={black},opacity = 0.4] (1,-5)--(1,8);
	\draw[color ={black},opacity = 0.4] (2,-5)--(2,8);
	\draw[color ={black},opacity = 0.4] (3,-5)--(3,8);
	\draw[color ={black},opacity = 0.4] (-1,-5)--(-1,8);
	\draw[color ={black},opacity = 0.4] (-2,-5)--(-2,8);
	\draw[color ={black},opacity = 0.4] (-3,-5)--(-3,8);
	\draw[color ={gray},opacity = 0.3] (-3,0)--(3,0);
	\draw[color ={gray},opacity = 0.3] (-3,0.5)--(3,0.5);
	\draw[color ={gray},opacity = 0.3] (-3,1.5)--(3,1.5);
	\draw[color ={gray},opacity = 0.3] (-3,2.5)--(3,2.5);
	\draw[color ={gray},opacity = 0.3] (-3,3.5)--(3,3.5);
	\draw[color ={gray},opacity = 0.3] (-3,4.5)--(3,4.5);
	\draw[color ={gray},opacity = 0.3] (-3,5.5)--(3,5.5);
	\draw[color ={gray},opacity = 0.3] (-3,6.5)--(3,6.5);
	\draw[color ={gray},opacity = 0.3] (-3,7.5)--(3,7.5);
	\draw[color ={gray},opacity = 0.3] (-3,0)--(3,0);
	\draw[color ={gray},opacity = 0.3] (-3,-0.5)--(3,-0.5);
	\draw[color ={gray},opacity = 0.3] (-3,-1.5)--(3,-1.5);
	\draw[color ={gray},opacity = 0.3] (-3,-2.5)--(3,-2.5);
	\draw[color ={gray},opacity = 0.3] (-3,-3.5)--(3,-3.5);
	\draw[color ={gray},opacity = 0.3] (-3,-4.5)--(3,-4.5);
	\draw[color ={gray},opacity = 0.3] (0,-5)--(0,8);
	\draw[color ={gray},opacity = 0.3] (0,-5)--(0,8);
	\draw[color ={black},line width = 1.2] (1,-0.13)--(1,0.13);
	\draw[color ={black},line width = 1.2] (2,-0.13)--(2,0.13);
	\draw[color ={black},line width = 1.2] (3,-0.13)--(3,0.13);
	\draw[color ={black},line width = 1.2] (-1,-0.13)--(-1,0.13);
	\draw[color ={black},line width = 1.2] (-2,-0.13)--(-2,0.13);
	\draw[color ={black},line width = 1.2] (-3,-0.13)--(-3,0.13);
	\draw[color ={black},line width = 1.2] (-0.13,1)--(0.13,1);
	\draw[color ={black},line width = 1.2] (-0.13,2)--(0.13,2);
	\draw[color ={black},line width = 1.2] (-0.13,3)--(0.13,3);
	\draw[color ={black},line width = 1.2] (-0.13,4)--(0.13,4);
	\draw[color ={black},line width = 1.2] (-0.13,5)--(0.13,5);
	\draw[color ={black},line width = 1.2] (-0.13,6)--(0.13,6);
	\draw[color ={black},line width = 1.2] (-0.13,7)--(0.13,7);
	\draw[color ={black},line width = 1.2] (-0.13,8)--(0.13,8);
	\draw[color ={black},line width = 1.2] (-0.13,9)--(0.13,9);
	\draw[color ={black},line width = 1.2] (-0.13,10)--(0.13,10);
	\draw[color ={black},line width = 1.2] (-0.13,11)--(0.13,11);
	\draw[color ={black},line width = 1.2] (-0.13,12)--(0.13,12);
	\draw[color ={black},line width = 1.2] (-0.13,-1)--(0.13,-1);
	\draw[color ={black},line width = 1.2] (-0.13,-2)--(0.13,-2);
	\draw[color ={black},line width = 1.2] (-0.13,-3)--(0.13,-3);
	\draw (3,-0.4) node[anchor = center] {\scriptsize \color{black}{$3$}};
	\draw (-3,-0.4) node[anchor = center] {\scriptsize \color{black}{$-3$}};
	\draw (-0.5,3.1) node[anchor = center] {\footnotesize \color{black}{$3$}};
	\draw (-0.5,6.1) node[anchor = center] {\footnotesize \color{black}{$6$}};
	\draw (-0.5,-2.9) node[anchor = center] {\footnotesize \color{black}{$-3$}};
	\draw [color={black}] (-0.3,-0.3) node[anchor = center,scale=1, rotate = 0] {O};
	\draw (3,6) node[anchor = center] {\footnotesize \color{red}{$\Large\mathcal{C}_f\prime$}};
	
	\draw[color={red},line width = 2] (-2.8,-5.6)--(-2.6,-5.2)--(-2.4,-4.8)--(-2.2,-4.4)--(-2,-4)--(-1.8,-3.6)--(-1.6,-3.2)--(-1.4,-2.8)--(-1.2,-2.4)--(-1,-2)--(-0.8,-1.6)--(-0.6,-1.2)--(-0.4,-0.8)--(-0.2,-0.4)--(0,0)--(0.2,0.4)--(0.4,0.8)--(0.6,1.2)--(0.8,1.6)--(1,2)--(1.2,2.4)--(1.4,2.8)--(1.6,3.2)--(1.8,3.6)--(2,4)--(2.2,4.4)--(2.4,4.8)--(2.6,5.2)--(2.8,5.6)--(3,6);

\end{tikzpicture}\\
    \end{minipage}
    
\end{exercice}

\begin{Solution}
 L'équation réduite de la tangente au point d'abscisse $-1$ est  : $y=f'(-1)(x-(-1))+f(-1)$.\\
        On lit graphiquement $f(-1)=0$ et $f'(-1)=-2$.\\
        L'équation réduite de la tangente est donc donnée par :
        $y=-2(x+1)+0$, soit $y=-2x-2$.
\end{Solution}

% @see : https://coopmaths.fr/alea?uuid=a1ba2&id=can1F14&n=1&d=10&alea=rkIp2232323232342345678910111213141516171819202122232425262728293031323334353637383940414243444546474823456789101112131415161718&cd=1&cols=1
\needspace{10\baselineskip}
\begin{exercice}[BaremeTotal=false]


Soit $f$ la fonction définie sur $\mathbb{R}$ par : 
$f(x)= 7x^2+1$.\\
En calculant la limite du taux d'accroissement, déterminer $f'(1)$.
\end{exercice}

\begin{Solution}
 $f$ est une fonction polynôme du second degré de la forme $f(x)=ax^2+b$.\\
    La fonction dérivée est donnée par la somme des dérivées des fonctions $u$ et $v$ définies par $u(x)=7x^2$ et $v(x)=1$.\\
     Comme $u'(x)=14x$ et $v'(x)=0$, on obtient  $f'(x)=14x$. \\
     Ainsi, $f'(1)=14\times 1=14$.
\end{Solution}

% @see : https://coopmaths.fr/alea?uuid=3c690&id=can1F13&n=1&d=10&alea=EUDu2232323232342345678910111213141516171819202122232425262728293031323334353637383940414243444546474823456789101112131415161718&cd=1&cols=1
\needspace{10\baselineskip}
\begin{exercice}[BaremeTotal=false]


 Déterminer le coefficient directeur de la tangente à la courbe représentative de la fonction carré au point d'abscisse $13$.    
\end{exercice}

\begin{Solution}
 Le coefficient directeur de la tangente au point d'abscisse $a$ est donné par le nombre dérivé $f'(a)$.\\
        La fonction $f$ définie par $f(x)=x^2$ a pour fonction dérivée la fonction $f'$ définie par $f'(x)=2x$.\\
        Comme $f'(13)=2\times 13=26$, le coefficient directeur de la tangente au point d'abscisse $13$ est : $26$. 
\end{Solution}

% @see : https://coopmaths.fr/alea?uuid=ebd89&id=1AN14-1&n=1&d=10&s=3&alea=ocYr2232323232342345678910111213141516171819202122232425262728293031323334353637383940414243444546474823456789101112131415161718&cd=0&cols=1
\needspace{10\baselineskip}
\begin{exercice}[BaremeTotal=false]


 Donner la dérivée de la fonction $f$, dérivable pour tout $x\in\R$, définie par  $f(x)=\dfrac{-2x+6}{7}$.
\end{exercice}

\begin{Solution}
 L'expression de la dérivée de la fonction $f$ définie par $f(x)=\dfrac{-2x+6}{7}$ est : ${\color[HTML]{f15929}\boldsymbol{f'(x)=-\dfrac{2}{7}}}$.
\end{Solution}

% @see : https://coopmaths.fr/alea?uuid=ebd89&id=1AN14-1&n=1&d=10&s=4&alea=uAlP2232323232342345678910111213141516171819202122232425262728293031323334353637383940414243444546474823456789101112131415161718&cd=0&cols=1
\needspace{10\baselineskip}
\begin{exercice}[BaremeTotal=false]


 Donner la dérivée de la fonction $f$, dérivable pour tout $x\in\R$, définie par  $f(x)=-3x^{10}$.
\end{exercice}

\begin{Solution}
 L'expression de la dérivée de la fonction $f$ définie par $f(x)=-3x^{10}$ est : ${\color[HTML]{f15929}\boldsymbol{f'(x)=-30x^{9}}}$.
\end{Solution}

% @see : https://coopmaths.fr/alea?uuid=ebd89&id=1AN14-1&n=1&d=10&s=7&alea=vjN72232323232342345678910111213141516171819202122232425262728293031323334353637383940414243444546474823456789101112131415161718&cd=0&cols=1
\needspace{10\baselineskip}
\begin{exercice}[BaremeTotal=false]


 Donner la dérivée de la fonction $f$, dérivable pour tout $x\in\R^*$, définie par  $f(x)=\dfrac{3}{4x}$.
\end{exercice}

\begin{Solution}
 L'expression de la dérivée de la fonction $f$ définie par $f(x)=\dfrac{3}{4x}$ est : ${\color[HTML]{f15929}\boldsymbol{f'(x)=-\dfrac{3}{4x^2}}}$.
\end{Solution}

% @see : https://coopmaths.fr/alea?uuid=a83c0&id=1AN14-4&n=1&d=10&s=1&alea=Njr32232323232342345678910111213141516171819202122232425262728293031323334353637383940414243444546474823456789101112131415161718&cd=0&cols=1
\needspace{10\baselineskip}
\begin{exercice}[BaremeTotal=false]


 Donner l'expression de la dérivée de la fonction $f$ définie sur $\R^{*}$ par $f(x)=\dfrac{3}{x}-5x^{2}+5x-2$.\\
\end{exercice}

\begin{Solution}
 L'expression de la dérivée de la fonction $f$ définie par $f(x)=\dfrac{3}{x}-5x^{2}+5x-2$ est :\\$f'(x)={\color[HTML]{f15929}\boldsymbol{-\dfrac{3}{x^2}-10x+5}}$.\\
\end{Solution}

% @see : https://coopmaths.fr/alea?uuid=a83c0&id=1AN14-4&n=1&d=10&s=4&alea=4Vpz2232323232342345678910111213141516171819202122232425262728293031323334353637383940414243444546474823456789101112131415161718&cd=1&cols=1
\needspace{10\baselineskip}
\begin{exercice}[BaremeTotal=false]


 Donner l'expression de la dérivée de la fonction $f$ définie sur $\R_+^{*}$ par $f(x)=5x-\dfrac{3}{x}+4\sqrt{x}$.\\
\end{exercice}

\begin{Solution}
 $(5x)^\prime=5$\\$(-\dfrac{3}{x})^\prime=\dfrac{3}{x^2}$\\$(4\sqrt{x})^\prime=\dfrac{4}{2\sqrt{x}}$\\L'expression de la dérivée de la fonction $f$ définie par $f(x)=5x-\dfrac{3}{x}+4\sqrt{x}$ est :\\$f'(x)={\color[HTML]{f15929}\boldsymbol{5+\dfrac{3}{x^2}+\dfrac{4}{2\sqrt{x}}}}$.\\
\end{Solution}

\end{Maquette}
\clearpage
\begin{Maquette}[DM]{Niveau= ,Classe=\getdata[19]\mydata\ (v19), Date = 06/01/25  ,Theme=Exercices}


% @see : https://coopmaths.fr/alea?uuid=4c8c7&id=1AN11-3&n=1&d=10&s=2&alea=nawO223232323234234567891011121314151617181920212223242526272829303132333435363738394041424344454647482345678910111213141516171819&cd=1&cols=2
\needspace{10\baselineskip}
\begin{exercice}[BaremeTotal=false]
\label{eleve:19}


 \begin{minipage}[t]{\linewidth} Soit $f$ une fonction dérivable sur $[-5;5]$ et $\mathcal{C}_f$ sa courbe représentative.\\On sait que  $f(5)=5~~$ et que $~~f'(5)=-2$.\\Déterminer une équation de la tangente $(T)$ à la courbe $\mathcal{C}_f$ au point d'abscisse $5$,\\sans utiliser la formule de cours de l'équation de tangente. \end{minipage}

\end{exercice}

\begin{Solution}
 \begin{minipage}[t]{\linewidth} On sait que la tangente n'est pas une droite verticale, puisque la fonction est dérivable sur l'intervalle.\\On en déduit que la tangente $(T)$ au point d'abscisse $5$, admet une équation réduite de la forme :  $(T) : y=m x + p$. 

\medskip
$\bullet$ Détermination de $m$ :\\On sait que le nombre dérivé en $5$ est par définition, le coefficient directeur de la tangente au point d'abscisse $5$.\\Par conséquent, on a déjà : $m=f'(5)=-2$.\\ On en déduit que  $(T) : y= -2 x + p$\\$\bullet$ Détermination de $p$ :\\ Pour cela, on utilise que si $f(5)=5~~$, alors le point $A$ de coordonnées $(5;5)$ appartient à $\mathcal{C}_f$ mais aussi à $(T)$.\\ On peut écrire $A(5;5) = \mathcal{C}_f \cap (T)$.\\ On remplace alors les coordonnées de $A(5;5)$ dans l'équation  $(T) : y= -2 x + p$.\\ $\begin{aligned} \phantom{\iff}&A(5;5)\in (T)\\ \iff& 5= -2 \times 5  + p\\ \iff& p=5 +2 \times 5  \\ \iff& p=5 +10   \\ \iff& p=15\\\end{aligned}$\\On peut conclure que : $(T) : {\color[HTML]{f15929}\boldsymbol{y=-2x+15}}$. \end{minipage}

\end{Solution}

% @see : https://coopmaths.fr/alea?uuid=0e984&id=can1F15&n=1&d=10&alea=1Alo223232323234234567891011121314151617181920212223242526272829303132333435363738394041424344454647482345678910111213141516171819&cd=1&cols=1
\needspace{10\baselineskip}
\begin{exercice}[BaremeTotal=false]


 La courbe représente une fonction $f$ et la droite est la tangente au point d'abscisse $0$.\\
        
        Déterminer $f'(0)$.\begin{center}\begin{tikzpicture}[baseline,scale = 0.5]

    \tikzset{
      point/.style={
        thick,
        draw,
        cross out,
        inner sep=0pt,
        minimum width=5pt,
        minimum height=5pt,
      },
    }
    \clip (-8,-3) rectangle (8,10);
    	
	\draw[color ={black},line width = 2,>=latex,->] (-8,0)--(8,0);
	\draw[color ={black},line width = 2,>=latex,->] (0,-3)--(0,10);
	\draw[color ={black},opacity = 0.5] (-8,1)--(8,1);
	\draw[color ={black},opacity = 0.5] (-8,2)--(8,2);
	\draw[color ={black},opacity = 0.5] (-8,3)--(8,3);
	\draw[color ={black},opacity = 0.5] (-8,4)--(8,4);
	\draw[color ={black},opacity = 0.5] (-8,5)--(8,5);
	\draw[color ={black},opacity = 0.5] (-8,6)--(8,6);
	\draw[color ={black},opacity = 0.5] (-8,7)--(8,7);
	\draw[color ={black},opacity = 0.5] (-8,8)--(8,8);
	\draw[color ={black},opacity = 0.5] (-8,9)--(8,9);
	\draw[color ={black},opacity = 0.5] (-8,10)--(8,10);
	\draw[color ={black},opacity = 0.5] (-8,-1)--(8,-1);
	\draw[color ={black},opacity = 0.5] (-8,-2)--(8,-2);
	\draw[color ={black},opacity = 0.5] (-8,-3)--(8,-3);
	\draw[color ={black},opacity = 0.5] (1,-3)--(1,10);
	\draw[color ={black},opacity = 0.5] (2,-3)--(2,10);
	\draw[color ={black},opacity = 0.5] (3,-3)--(3,10);
	\draw[color ={black},opacity = 0.5] (4,-3)--(4,10);
	\draw[color ={black},opacity = 0.5] (5,-3)--(5,10);
	\draw[color ={black},opacity = 0.5] (6,-3)--(6,10);
	\draw[color ={black},opacity = 0.5] (7,-3)--(7,10);
	\draw[color ={black},opacity = 0.5] (8,-3)--(8,10);
	\draw[color ={black},opacity = 0.5] (-1,-3)--(-1,10);
	\draw[color ={black},opacity = 0.5] (-2,-3)--(-2,10);
	\draw[color ={black},opacity = 0.5] (-3,-3)--(-3,10);
	\draw[color ={black},opacity = 0.5] (-4,-3)--(-4,10);
	\draw[color ={black},opacity = 0.5] (-5,-3)--(-5,10);
	\draw[color ={black},opacity = 0.5] (-6,-3)--(-6,10);
	\draw[color ={black},opacity = 0.5] (-7,-3)--(-7,10);
	\draw[color ={black},opacity = 0.5] (-8,-3)--(-8,10);
	\draw[color ={gray},opacity = 0.3] (-8,0)--(8,0);
	\draw[color ={gray},opacity = 0.3] (-8,0)--(8,0);
	\draw[color ={gray},opacity = 0.3] (0,-3)--(0,10);
	\draw[color ={gray},opacity = 0.3] (0,-3)--(0,10);
	\draw[color ={black},line width = 2] (2,-0.2)--(2,0.2);
	\draw[color ={black},line width = 2] (4,-0.2)--(4,0.2);
	\draw[color ={black},line width = 2] (6,-0.2)--(6,0.2);
	\draw[color ={black},line width = 2] (-2,-0.2)--(-2,0.2);
	\draw[color ={black},line width = 2] (-4,-0.2)--(-4,0.2);
	\draw[color ={black},line width = 2] (-6,-0.2)--(-6,0.2);
	\draw[color ={black},line width = 2] (-0.2,1)--(0.2,1);
	\draw[color ={black},line width = 2] (-0.2,2)--(0.2,2);
	\draw[color ={black},line width = 2] (-0.2,3)--(0.2,3);
	\draw[color ={black},line width = 2] (-0.2,4)--(0.2,4);
	\draw[color ={black},line width = 2] (-0.2,5)--(0.2,5);
	\draw[color ={black},line width = 2] (-0.2,6)--(0.2,6);
	\draw[color ={black},line width = 2] (-0.2,7)--(0.2,7);
	\draw[color ={black},line width = 2] (-0.2,8)--(0.2,8);
	\draw[color ={black},line width = 2] (-0.2,9)--(0.2,9);
	\draw[color ={black},line width = 2] (-0.2,-1)--(0.2,-1);
	\draw[color ={black},line width = 2] (-0.2,-2)--(0.2,-2);
	\draw (2,-0.4) node[anchor = center] {\scriptsize \color{black}{$1$}};
	\draw (4,-0.4) node[anchor = center] {\scriptsize \color{black}{$2$}};
	\draw (6,-0.4) node[anchor = center] {\scriptsize \color{black}{$3$}};
	\draw (-2,-0.4) node[anchor = center] {\scriptsize \color{black}{$-1$}};
	\draw (-4,-0.4) node[anchor = center] {\scriptsize \color{black}{$-2$}};
	\draw (-6,-0.4) node[anchor = center] {\scriptsize \color{black}{$-3$}};
	\draw (-0.8,2.1) node[anchor = center] {\footnotesize \color{black}{$2$}};
	\draw (-0.8,4.1) node[anchor = center] {\footnotesize \color{black}{$4$}};
	\draw (-0.8,6.1) node[anchor = center] {\footnotesize \color{black}{$6$}};
	\draw (-0.8,8.1) node[anchor = center] {\footnotesize \color{black}{$8$}};
	\draw [color={black}] (-0.3,-0.3) node[anchor = center,scale=1, rotate = 0] {O};
	
	\draw[color={blue},line width = 2] (-2.8,10.52)--(-2.4,8.68)--(-2,7)--(-1.6,5.48)--(-1.2,4.12)--(-0.8,2.92)--(-0.4,1.88)--(0,1)--(0.4,0.28)--(0.8,-0.28)--(1.2,-0.68)--(1.6,-0.92)--(2,-1)--(2.4,-0.92)--(2.8,-0.68)--(3.2,-0.28)--(3.6,0.28)--(4,1)--(4.4,1.88)--(4.8,2.92)--(5.2,4.12)--(5.6,5.48)--(6,7)--(6.4,8.68)--(6.8,10.52);
	
	\draw[color={red},line width = 2] (-4.8,10.6)--(-4.4,9.8)--(-4,9)--(-3.6,8.2)--(-3.2,7.4)--(-2.8,6.6)--(-2.4,5.8)--(-2,5)--(-1.6,4.2)--(-1.2,3.4)--(-0.8,2.6)--(-0.4,1.8)--(0,1)--(0.4,0.2)--(0.8,-0.6)--(1.2,-1.4)--(1.6,-2.2)--(2,-3)--(2.4,-3.8);

\end{tikzpicture}\end{center}
\end{exercice}

\begin{Solution}
 $f'(0)$ est donné par le coefficient directeur de la tangente à la courbe au point d'abscisse $0$, soit $-4$.
\end{Solution}

% @see : https://coopmaths.fr/alea?uuid=6f32d&id=can1F16&n=1&d=10&alea=dqSw223232323234234567891011121314151617181920212223242526272829303132333435363738394041424344454647482345678910111213141516171819&cd=1&cols=1
\needspace{10\baselineskip}
\newpage
\begin{exercice}[BaremeTotal=false]


 On donne les représentations graphiques d'une fonction et de sa dérivée.\\
        Donner l'équation réduite de la tangente à la courbe de $f$ en $x=2$. \\ \begin{minipage}{0.5\linewidth}
    \begin{tikzpicture}[baseline, scale=0.8]

    \tikzset{
      point/.style={
        thick,
        draw,
        cross out,
        inner sep=0pt,
        minimum width=5pt,
        minimum height=5pt,
      },
    }
    \clip (-4.5,-4.5) rectangle (5.75,13.5);
    	
	\draw[color ={black},>=latex,->] (-3,0)--(3,0);
	\draw[color ={black},>=latex,->] (0,-3)--(0,12);
	\draw[color ={black},opacity = 0.4] (-3,1)--(3,1);
	\draw[color ={black},opacity = 0.4] (-3,2)--(3,2);
	\draw[color ={black},opacity = 0.4] (-3,3)--(3,3);
	\draw[color ={black},opacity = 0.4] (-3,4)--(3,4);
	\draw[color ={black},opacity = 0.4] (-3,5)--(3,5);
	\draw[color ={black},opacity = 0.4] (-3,6)--(3,6);
	\draw[color ={black},opacity = 0.4] (-3,7)--(3,7);
	\draw[color ={black},opacity = 0.4] (-3,8)--(3,8);
	\draw[color ={black},opacity = 0.4] (-3,9)--(3,9);
	\draw[color ={black},opacity = 0.4] (-3,10)--(3,10);
	\draw[color ={black},opacity = 0.4] (-3,11)--(3,11);
	\draw[color ={black},opacity = 0.4] (-3,12)--(3,12);
	\draw[color ={black},opacity = 0.4] (-3,-1)--(3,-1);
	\draw[color ={black},opacity = 0.4] (-3,-2)--(3,-2);
	\draw[color ={black},opacity = 0.4] (-3,-3)--(3,-3);
	\draw[color ={black},opacity = 0.4] (1,-3)--(1,12);
	\draw[color ={black},opacity = 0.4] (2,-3)--(2,12);
	\draw[color ={black},opacity = 0.4] (3,-3)--(3,12);
	\draw[color ={black},opacity = 0.4] (-1,-3)--(-1,12);
	\draw[color ={black},opacity = 0.4] (-2,-3)--(-2,12);
	\draw[color ={black},opacity = 0.4] (-3,-3)--(-3,12);
	\draw[color ={gray},opacity = 0.3] (-3,0)--(3,0);
	\draw[color ={gray},opacity = 0.3] (-3,0.5)--(3,0.5);
	\draw[color ={gray},opacity = 0.3] (-3,1.5)--(3,1.5);
	\draw[color ={gray},opacity = 0.3] (-3,2.5)--(3,2.5);
	\draw[color ={gray},opacity = 0.3] (-3,3.5)--(3,3.5);
	\draw[color ={gray},opacity = 0.3] (-3,4.5)--(3,4.5);
	\draw[color ={gray},opacity = 0.3] (-3,5.5)--(3,5.5);
	\draw[color ={gray},opacity = 0.3] (-3,6.5)--(3,6.5);
	\draw[color ={gray},opacity = 0.3] (-3,7.5)--(3,7.5);
	\draw[color ={gray},opacity = 0.3] (-3,8.5)--(3,8.5);
	\draw[color ={gray},opacity = 0.3] (-3,9.5)--(3,9.5);
	\draw[color ={gray},opacity = 0.3] (-3,10.5)--(3,10.5);
	\draw[color ={gray},opacity = 0.3] (-3,11.5)--(3,11.5);
	\draw[color ={gray},opacity = 0.3] (-3,0)--(3,0);
	\draw[color ={gray},opacity = 0.3] (-3,-0.5)--(3,-0.5);
	\draw[color ={gray},opacity = 0.3] (-3,-1.5)--(3,-1.5);
	\draw[color ={gray},opacity = 0.3] (-3,-2.5)--(3,-2.5);
	\draw[color ={gray},opacity = 0.3] (0,-3)--(0,12);
	\draw[color ={gray},opacity = 0.3] (0,-3)--(0,12);
	\draw[color ={black},line width = 1.2] (1,-0.13)--(1,0.13);
	\draw[color ={black},line width = 1.2] (2,-0.13)--(2,0.13);
	\draw[color ={black},line width = 1.2] (3,-0.13)--(3,0.13);
	\draw[color ={black},line width = 1.2] (-1,-0.13)--(-1,0.13);
	\draw[color ={black},line width = 1.2] (-2,-0.13)--(-2,0.13);
	\draw[color ={black},line width = 1.2] (-3,-0.13)--(-3,0.13);
	\draw[color ={black},line width = 1.2] (-0.13,1)--(0.13,1);
	\draw[color ={black},line width = 1.2] (-0.13,2)--(0.13,2);
	\draw[color ={black},line width = 1.2] (-0.13,3)--(0.13,3);
	\draw[color ={black},line width = 1.2] (-0.13,4)--(0.13,4);
	\draw[color ={black},line width = 1.2] (-0.13,5)--(0.13,5);
	\draw[color ={black},line width = 1.2] (-0.13,6)--(0.13,6);
	\draw[color ={black},line width = 1.2] (-0.13,7)--(0.13,7);
	\draw[color ={black},line width = 1.2] (-0.13,8)--(0.13,8);
	\draw[color ={black},line width = 1.2] (-0.13,9)--(0.13,9);
	\draw[color ={black},line width = 1.2] (-0.13,10)--(0.13,10);
	\draw[color ={black},line width = 1.2] (-0.13,11)--(0.13,11);
	\draw[color ={black},line width = 1.2] (-0.13,12)--(0.13,12);
	\draw[color ={black},line width = 1.2] (-0.13,-1)--(0.13,-1);
	\draw[color ={black},line width = 1.2] (-0.13,-2)--(0.13,-2);
	\draw[color ={black},line width = 1.2] (-0.13,-3)--(0.13,-3);
	\draw (3,-0.4) node[anchor = center] {\scriptsize \color{black}{$3$}};
	\draw (-3,-0.4) node[anchor = center] {\scriptsize \color{black}{$-3$}};
	\draw (-0.5,3.1) node[anchor = center] {\footnotesize \color{black}{$3$}};
	\draw (-0.5,6.1) node[anchor = center] {\footnotesize \color{black}{$6$}};
	\draw (-0.5,9.1) node[anchor = center] {\footnotesize \color{black}{$9$}};
	\draw (-0.5,-2.9) node[anchor = center] {\footnotesize \color{black}{$-3$}};
	\draw [color={black}] (-0.3,-0.3) node[anchor = center,scale=1, rotate = 0] {O};
	\draw (3,10) node[anchor = center] {\footnotesize \color{blue}{$\Large \mathcal{C}_f$}};
	
	\draw[color={blue},line width = 2] (-2.2,12.24)--(-2,11)--(-1.8,9.84)--(-1.6,8.76)--(-1.4,7.76)--(-1.2,6.84)--(-1,6)--(-0.8,5.24)--(-0.6,4.56)--(-0.4,3.96)--(-0.2,3.44)--(0,3)--(0.2,2.64)--(0.4,2.36)--(0.6,2.16)--(0.8,2.04)--(1,2)--(1.2,2.04)--(1.4,2.16)--(1.6,2.36)--(1.8,2.64)--(2,3)--(2.2,3.44)--(2.4,3.96)--(2.6,4.56)--(2.8,5.24)--(3,6);

\end{tikzpicture}\\
    \end{minipage}
    \begin{minipage}{0.5\linewidth}
    \begin{tikzpicture}[baseline, scale=0.8]

    \tikzset{
      point/.style={
        thick,
        draw,
        cross out,
        inner sep=0pt,
        minimum width=5pt,
        minimum height=5pt,
      },
    }
    \clip (-4.5,-6.5) rectangle (6.5,9.5);
    	
	\draw[color ={black},>=latex,->] (-3,0)--(3,0);
	\draw[color ={black},>=latex,->] (0,-5)--(0,8);
	\draw[color ={black},opacity = 0.4] (-3,1)--(3,1);
	\draw[color ={black},opacity = 0.4] (-3,2)--(3,2);
	\draw[color ={black},opacity = 0.4] (-3,3)--(3,3);
	\draw[color ={black},opacity = 0.4] (-3,4)--(3,4);
	\draw[color ={black},opacity = 0.4] (-3,5)--(3,5);
	\draw[color ={black},opacity = 0.4] (-3,6)--(3,6);
	\draw[color ={black},opacity = 0.4] (-3,7)--(3,7);
	\draw[color ={black},opacity = 0.4] (-3,8)--(3,8);
	\draw[color ={black},opacity = 0.4] (-3,-1)--(3,-1);
	\draw[color ={black},opacity = 0.4] (-3,-2)--(3,-2);
	\draw[color ={black},opacity = 0.4] (-3,-3)--(3,-3);
	\draw[color ={black},opacity = 0.4] (-3,-4)--(3,-4);
	\draw[color ={black},opacity = 0.4] (-3,-5)--(3,-5);
	\draw[color ={black},opacity = 0.4] (1,-5)--(1,8);
	\draw[color ={black},opacity = 0.4] (2,-5)--(2,8);
	\draw[color ={black},opacity = 0.4] (3,-5)--(3,8);
	\draw[color ={black},opacity = 0.4] (-1,-5)--(-1,8);
	\draw[color ={black},opacity = 0.4] (-2,-5)--(-2,8);
	\draw[color ={black},opacity = 0.4] (-3,-5)--(-3,8);
	\draw[color ={gray},opacity = 0.3] (-3,0)--(3,0);
	\draw[color ={gray},opacity = 0.3] (-3,0.5)--(3,0.5);
	\draw[color ={gray},opacity = 0.3] (-3,1.5)--(3,1.5);
	\draw[color ={gray},opacity = 0.3] (-3,2.5)--(3,2.5);
	\draw[color ={gray},opacity = 0.3] (-3,3.5)--(3,3.5);
	\draw[color ={gray},opacity = 0.3] (-3,4.5)--(3,4.5);
	\draw[color ={gray},opacity = 0.3] (-3,5.5)--(3,5.5);
	\draw[color ={gray},opacity = 0.3] (-3,6.5)--(3,6.5);
	\draw[color ={gray},opacity = 0.3] (-3,7.5)--(3,7.5);
	\draw[color ={gray},opacity = 0.3] (-3,0)--(3,0);
	\draw[color ={gray},opacity = 0.3] (-3,-0.5)--(3,-0.5);
	\draw[color ={gray},opacity = 0.3] (-3,-1.5)--(3,-1.5);
	\draw[color ={gray},opacity = 0.3] (-3,-2.5)--(3,-2.5);
	\draw[color ={gray},opacity = 0.3] (-3,-3.5)--(3,-3.5);
	\draw[color ={gray},opacity = 0.3] (-3,-4.5)--(3,-4.5);
	\draw[color ={gray},opacity = 0.3] (0,-5)--(0,8);
	\draw[color ={gray},opacity = 0.3] (0,-5)--(0,8);
	\draw[color ={black},line width = 1.2] (1,-0.13)--(1,0.13);
	\draw[color ={black},line width = 1.2] (2,-0.13)--(2,0.13);
	\draw[color ={black},line width = 1.2] (3,-0.13)--(3,0.13);
	\draw[color ={black},line width = 1.2] (-1,-0.13)--(-1,0.13);
	\draw[color ={black},line width = 1.2] (-2,-0.13)--(-2,0.13);
	\draw[color ={black},line width = 1.2] (-3,-0.13)--(-3,0.13);
	\draw[color ={black},line width = 1.2] (-0.13,1)--(0.13,1);
	\draw[color ={black},line width = 1.2] (-0.13,2)--(0.13,2);
	\draw[color ={black},line width = 1.2] (-0.13,3)--(0.13,3);
	\draw[color ={black},line width = 1.2] (-0.13,4)--(0.13,4);
	\draw[color ={black},line width = 1.2] (-0.13,5)--(0.13,5);
	\draw[color ={black},line width = 1.2] (-0.13,6)--(0.13,6);
	\draw[color ={black},line width = 1.2] (-0.13,7)--(0.13,7);
	\draw[color ={black},line width = 1.2] (-0.13,8)--(0.13,8);
	\draw[color ={black},line width = 1.2] (-0.13,9)--(0.13,9);
	\draw[color ={black},line width = 1.2] (-0.13,10)--(0.13,10);
	\draw[color ={black},line width = 1.2] (-0.13,11)--(0.13,11);
	\draw[color ={black},line width = 1.2] (-0.13,12)--(0.13,12);
	\draw[color ={black},line width = 1.2] (-0.13,-1)--(0.13,-1);
	\draw[color ={black},line width = 1.2] (-0.13,-2)--(0.13,-2);
	\draw[color ={black},line width = 1.2] (-0.13,-3)--(0.13,-3);
	\draw (3,-0.4) node[anchor = center] {\scriptsize \color{black}{$3$}};
	\draw (-3,-0.4) node[anchor = center] {\scriptsize \color{black}{$-3$}};
	\draw (-0.5,3.1) node[anchor = center] {\footnotesize \color{black}{$3$}};
	\draw (-0.5,6.1) node[anchor = center] {\footnotesize \color{black}{$6$}};
	\draw (-0.5,-2.9) node[anchor = center] {\footnotesize \color{black}{$-3$}};
	\draw [color={black}] (-0.3,-0.3) node[anchor = center,scale=1, rotate = 0] {O};
	\draw (3,6) node[anchor = center] {\footnotesize \color{red}{$\Large\mathcal{C}_f\prime$}};
	
	\draw[color={red},line width = 2] (-2,-6)--(-1.8,-5.6)--(-1.6,-5.2)--(-1.4,-4.8)--(-1.2,-4.4)--(-1,-4)--(-0.8,-3.6)--(-0.6,-3.2)--(-0.4,-2.8)--(-0.2,-2.4)--(0,-2)--(0.2,-1.6)--(0.4,-1.2)--(0.6,-0.8)--(0.8,-0.4)--(1,0)--(1.2,0.4)--(1.4,0.8)--(1.6,1.2)--(1.8,1.6)--(2,2)--(2.2,2.4)--(2.4,2.8)--(2.6,3.2)--(2.8,3.6)--(3,4);

\end{tikzpicture}\\
    \end{minipage}
    
\end{exercice}

\begin{Solution}
 L'équation réduite de la tangente au point d'abscisse $2$ est  : $y=f'(2)(x-2)+f(2)$.\\
        On lit graphiquement $f(2)=3$ et $f'(2)=2$.\\
        L'équation réduite de la tangente est donc donnée par :
        $y=2(x-2)+3$, soit $y=2x-1$.
\end{Solution}

% @see : https://coopmaths.fr/alea?uuid=a1ba2&id=can1F14&n=1&d=10&alea=rkIp223232323234234567891011121314151617181920212223242526272829303132333435363738394041424344454647482345678910111213141516171819&cd=1&cols=1
\needspace{10\baselineskip}
\begin{exercice}[BaremeTotal=false]


 Soit $f$ la fonction définie sur $\mathbb{R}$ par : $f(x)= 3x^2+5x+6$.\\
        En calculant la limite du taux d'accroissement, déterminer $f'(-1)$.
\end{exercice}

\begin{Solution}
 $f$ est une fonction polynôme du second degré de la forme $f(x)=ax^2+bx+c$.\\
    La fonction dérivée est donnée par la somme des dérivées des fonctions $u$ et $v$ définies par $u(x)=3x^2$ et $v(x)=5x+6$.\\
     Comme $u'(x)=6x$ et $v'(x)=5$, on obtient  $f'(x)=6x+5$. \\
     Ainsi, $f'(-1)=6\times (-1)+5=-1$.
\end{Solution}

% @see : https://coopmaths.fr/alea?uuid=3c690&id=can1F13&n=1&d=10&alea=EUDu223232323234234567891011121314151617181920212223242526272829303132333435363738394041424344454647482345678910111213141516171819&cd=1&cols=1
\needspace{10\baselineskip}
\begin{exercice}[BaremeTotal=false]


 Déterminer le coefficient directeur de la tangente à la courbe représentative de la fonction racine carrée au point d'abscisse $25$.
         
\end{exercice}

\begin{Solution}
 Le coefficient directeur de la tangente au point d'abscisse $a$ est donné par le nombre dérivé $f'(a)$.\\
        La fonction $f$ définie par $f(x)=\sqrt{x}$ a pour fonction dérivée la fonction $f'$ définie par $f'(x)=\dfrac{1}{2\sqrt{x}}$.\\Comme $f'(25)=\dfrac{1}{2\sqrt{25}}=\dfrac{1}{10}$, le coefficient directeur de la tangente au point d'abscisse $25$ est : $\dfrac{1}{10}$.
\end{Solution}

% @see : https://coopmaths.fr/alea?uuid=ebd89&id=1AN14-1&n=1&d=10&s=3&alea=ocYr223232323234234567891011121314151617181920212223242526272829303132333435363738394041424344454647482345678910111213141516171819&cd=0&cols=1
\needspace{10\baselineskip}
\begin{exercice}[BaremeTotal=false]


 Donner la dérivée de la fonction $f$, dérivable pour tout $x\in\R$, définie par  $f(x)=\dfrac{6x+3}{7}$.
\end{exercice}

\begin{Solution}
 L'expression de la dérivée de la fonction $f$ définie par $f(x)=\dfrac{6x+3}{7}$ est : ${\color[HTML]{f15929}\boldsymbol{f'(x)=\dfrac{6}{7}}}$.
\end{Solution}

% @see : https://coopmaths.fr/alea?uuid=ebd89&id=1AN14-1&n=1&d=10&s=4&alea=uAlP223232323234234567891011121314151617181920212223242526272829303132333435363738394041424344454647482345678910111213141516171819&cd=0&cols=1
\needspace{10\baselineskip}
\begin{exercice}[BaremeTotal=false]


 Donner la dérivée de la fonction $f$, dérivable pour tout $x\in\R$, définie par  $f(x)=9x^{8}$.
\end{exercice}

\begin{Solution}
 L'expression de la dérivée de la fonction $f$ définie par $f(x)=9x^{8}$ est : ${\color[HTML]{f15929}\boldsymbol{f'(x)=72x^{7}}}$.
\end{Solution}

% @see : https://coopmaths.fr/alea?uuid=ebd89&id=1AN14-1&n=1&d=10&s=7&alea=vjN7223232323234234567891011121314151617181920212223242526272829303132333435363738394041424344454647482345678910111213141516171819&cd=0&cols=1
\needspace{10\baselineskip}
\begin{exercice}[BaremeTotal=false]


 Donner la dérivée de la fonction $f$, dérivable pour tout $x\in\R^*$, définie par  $f(x)=\dfrac{3}{10x}$.
\end{exercice}

\begin{Solution}
 L'expression de la dérivée de la fonction $f$ définie par $f(x)=\dfrac{3}{10x}$ est : ${\color[HTML]{f15929}\boldsymbol{f'(x)=-\dfrac{3}{10x^2}}}$.
\end{Solution}

% @see : https://coopmaths.fr/alea?uuid=a83c0&id=1AN14-4&n=1&d=10&s=1&alea=Njr3223232323234234567891011121314151617181920212223242526272829303132333435363738394041424344454647482345678910111213141516171819&cd=0&cols=1
\needspace{10\baselineskip}
\begin{exercice}[BaremeTotal=false]


 Donner l'expression de la dérivée de la fonction $f$ définie sur $\R^{*}$ par $f(x)=5x^{2}-3x-\dfrac{2}{x}$.\\
\end{exercice}

\begin{Solution}
 L'expression de la dérivée de la fonction $f$ définie par $f(x)=5x^{2}-3x-\dfrac{2}{x}$ est :\\$f'(x)={\color[HTML]{f15929}\boldsymbol{10x-3+\dfrac{2}{x^2}}}$.\\
\end{Solution}

% @see : https://coopmaths.fr/alea?uuid=a83c0&id=1AN14-4&n=1&d=10&s=4&alea=4Vpz223232323234234567891011121314151617181920212223242526272829303132333435363738394041424344454647482345678910111213141516171819&cd=1&cols=1
\needspace{10\baselineskip}
\begin{exercice}[BaremeTotal=false]


 Donner l'expression de la dérivée de la fonction $f$ définie sur $\R_+^{*}$ par $f(x)=4\sqrt{x}+3x^{2}+2x-4+\dfrac{5}{x}$.\\
\end{exercice}

\begin{Solution}
 $(4\sqrt{x})^\prime=\dfrac{4}{2\sqrt{x}}$\\$(3x^{2}+2x-4)^\prime=6x+2$\\$(\dfrac{5}{x})^\prime=-\dfrac{5}{x^2}$\\L'expression de la dérivée de la fonction $f$ définie par $f(x)=4\sqrt{x}+3x^{2}+2x-4+\dfrac{5}{x}$ est :\\$f'(x)={\color[HTML]{f15929}\boldsymbol{\dfrac{4}{2\sqrt{x}}+6x+2-\dfrac{5}{x^2}}}$.\\
\end{Solution}

\end{Maquette}
\clearpage
\begin{Maquette}[DM]{Niveau= ,Classe=\getdata[20]\mydata\ (v20), Date = 06/01/25  ,Theme=Exercices}


% @see : https://coopmaths.fr/alea?uuid=4c8c7&id=1AN11-3&n=1&d=10&s=2&alea=nawO22323232323423456789101112131415161718192021222324252627282930313233343536373839404142434445464748234567891011121314151617181920&cd=1&cols=2
\needspace{10\baselineskip}
\begin{exercice}[BaremeTotal=false]
\label{eleve:20}


 \begin{minipage}[t]{\linewidth} Soit $f$ une fonction dérivable sur $[-5;5]$ et $\mathcal{C}_f$ sa courbe représentative.\\On sait que  $f(1)=-4~~$ et que $~~f'(1)=-5$.\\Déterminer une équation de la tangente $(T)$ à la courbe $\mathcal{C}_f$ au point d'abscisse $1$,\\sans utiliser la formule de cours de l'équation de tangente. \end{minipage}

\end{exercice}

\begin{Solution}
 \begin{minipage}[t]{\linewidth} On sait que la tangente n'est pas une droite verticale, puisque la fonction est dérivable sur l'intervalle.\\On en déduit que la tangente $(T)$ au point d'abscisse $1$, admet une équation réduite de la forme :  $(T) : y=m x + p$. 

\medskip
$\bullet$ Détermination de $m$ :\\On sait que le nombre dérivé en $1$ est par définition, le coefficient directeur de la tangente au point d'abscisse $1$.\\Par conséquent, on a déjà : $m=f'(1)=-5$.\\ On en déduit que  $(T) : y= -5 x + p$\\$\bullet$ Détermination de $p$ :\\ Pour cela, on utilise que si $f(1)=-4~~$, alors le point $A$ de coordonnées $(1;-4)$ appartient à $\mathcal{C}_f$ mais aussi à $(T)$.\\ On peut écrire $A(1;-4) = \mathcal{C}_f \cap (T)$.\\ On remplace alors les coordonnées de $A(1;-4)$ dans l'équation  $(T) : y= -5 x + p$.\\ $\begin{aligned} \phantom{\iff}&A(1;-4)\in (T)\\ \iff& -4= -5 \times 1  + p\\ \iff& p=-4 +5 \times 1  \\ \iff& p=-4 +5   \\ \iff& p=1\\\end{aligned}$\\On peut conclure que : $(T) : {\color[HTML]{f15929}\boldsymbol{y=-5x+1}}$. \end{minipage}

\end{Solution}

% @see : https://coopmaths.fr/alea?uuid=0e984&id=can1F15&n=1&d=10&alea=1Alo22323232323423456789101112131415161718192021222324252627282930313233343536373839404142434445464748234567891011121314151617181920&cd=1&cols=1
\needspace{10\baselineskip}
\begin{exercice}[BaremeTotal=false]


 La courbe représente une fonction $f$ et la droite est la tangente au point d'abscisse $0$.\\
        Déterminer $f'(0)$.\begin{center}\begin{tikzpicture}[baseline,scale = 0.5]

    \tikzset{
      point/.style={
        thick,
        draw,
        cross out,
        inner sep=0pt,
        minimum width=5pt,
        minimum height=5pt,
      },
    }
    \clip (-10,-2) rectangle (4,12);
    	
	\draw[color ={black},line width = 2,>=latex,->] (-10,0)--(4,0);
	\draw[color ={black},line width = 2,>=latex,->] (0,-2)--(0,12);
	\draw[color ={black},opacity = 0.5] (-10,1)--(4,1);
	\draw[color ={black},opacity = 0.5] (-10,2)--(4,2);
	\draw[color ={black},opacity = 0.5] (-10,3)--(4,3);
	\draw[color ={black},opacity = 0.5] (-10,4)--(4,4);
	\draw[color ={black},opacity = 0.5] (-10,5)--(4,5);
	\draw[color ={black},opacity = 0.5] (-10,6)--(4,6);
	\draw[color ={black},opacity = 0.5] (-10,7)--(4,7);
	\draw[color ={black},opacity = 0.5] (-10,8)--(4,8);
	\draw[color ={black},opacity = 0.5] (-10,9)--(4,9);
	\draw[color ={black},opacity = 0.5] (-10,10)--(4,10);
	\draw[color ={black},opacity = 0.5] (-10,11)--(4,11);
	\draw[color ={black},opacity = 0.5] (-10,12)--(4,12);
	\draw[color ={black},opacity = 0.5] (-10,-1)--(4,-1);
	\draw[color ={black},opacity = 0.5] (-10,-2)--(4,-2);
	\draw[color ={black},opacity = 0.5] (1,-2)--(1,12);
	\draw[color ={black},opacity = 0.5] (2,-2)--(2,12);
	\draw[color ={black},opacity = 0.5] (3,-2)--(3,12);
	\draw[color ={black},opacity = 0.5] (4,-2)--(4,12);
	\draw[color ={black},opacity = 0.5] (-1,-2)--(-1,12);
	\draw[color ={black},opacity = 0.5] (-2,-2)--(-2,12);
	\draw[color ={black},opacity = 0.5] (-3,-2)--(-3,12);
	\draw[color ={black},opacity = 0.5] (-4,-2)--(-4,12);
	\draw[color ={black},opacity = 0.5] (-5,-2)--(-5,12);
	\draw[color ={black},opacity = 0.5] (-6,-2)--(-6,12);
	\draw[color ={black},opacity = 0.5] (-7,-2)--(-7,12);
	\draw[color ={black},opacity = 0.5] (-8,-2)--(-8,12);
	\draw[color ={black},opacity = 0.5] (-9,-2)--(-9,12);
	\draw[color ={black},opacity = 0.5] (-10,-2)--(-10,12);
	\draw[color ={gray},opacity = 0.3] (-10,0)--(4,0);
	\draw[color ={gray},opacity = 0.3] (-10,0)--(4,0);
	\draw[color ={gray},opacity = 0.3] (0,-2)--(0,12);
	\draw[color ={gray},opacity = 0.3] (0,-2)--(0,12);
	\draw[color ={gray},opacity = 0.3] (-5,-2)--(-5,12);
	\draw[color ={gray},opacity = 0.3] (-6,-2)--(-6,12);
	\draw[color ={gray},opacity = 0.3] (-7,-2)--(-7,12);
	\draw[color ={gray},opacity = 0.3] (-8,-2)--(-8,12);
	\draw[color ={gray},opacity = 0.3] (-9,-2)--(-9,12);
	\draw[color ={gray},opacity = 0.3] (-10,-2)--(-10,12);
	\draw[color ={black},line width = 2] (2,-0.2)--(2,0.2);
	\draw[color ={black},line width = 2] (-2,-0.2)--(-2,0.2);
	\draw[color ={black},line width = 2] (-4,-0.2)--(-4,0.2);
	\draw[color ={black},line width = 2] (-6,-0.2)--(-6,0.2);
	\draw[color ={black},line width = 2] (-8,-0.2)--(-8,0.2);
	\draw[color ={black},line width = 2] (-0.2,2)--(0.2,2);
	\draw[color ={black},line width = 2] (-0.2,4)--(0.2,4);
	\draw[color ={black},line width = 2] (-0.2,6)--(0.2,6);
	\draw[color ={black},line width = 2] (-0.2,8)--(0.2,8);
	\draw[color ={black},line width = 2] (-0.2,10)--(0.2,10);
	\draw (2,-0.4) node[anchor = center] {\scriptsize \color{black}{$1$}};
	\draw (-2,-0.4) node[anchor = center] {\scriptsize \color{black}{$-1$}};
	\draw (-4,-0.4) node[anchor = center] {\scriptsize \color{black}{$-2$}};
	\draw (-6,-0.4) node[anchor = center] {\scriptsize \color{black}{$-3$}};
	\draw (-8,-0.4) node[anchor = center] {\scriptsize \color{black}{$-4$}};
	\draw (-0.5,2.1) node[anchor = center] {\footnotesize \color{black}{$1$}};
	\draw (-0.5,4.1) node[anchor = center] {\footnotesize \color{black}{$2$}};
	\draw (-0.5,6.1) node[anchor = center] {\footnotesize \color{black}{$3$}};
	\draw (-0.5,8.1) node[anchor = center] {\footnotesize \color{black}{$4$}};
	\draw (-0.5,10.1) node[anchor = center] {\footnotesize \color{black}{$5$}};
	\draw [color={black}] (-0.3,-0.3) node[anchor = center,scale=1, rotate = 0] {O};
	
	\draw[color={blue},line width = 2] (-10,6.98)--(-9.6,6.64)--(-9.2,6.32)--(-8.8,6.01)--(-8.4,5.72)--(-8,5.44)--(-7.6,5.17)--(-7.2,4.92)--(-6.8,4.68)--(-6.4,4.45)--(-6,4.23)--(-5.6,4.03)--(-5.2,3.83)--(-4.8,3.64)--(-4.4,3.47)--(-4,3.3)--(-3.6,3.14)--(-3.2,2.98)--(-2.8,2.84)--(-2.4,2.7)--(-2,2.57)--(-1.6,2.44)--(-1.2,2.32)--(-0.8,2.21)--(-0.4,2.1)--(0,2)--(0.4,1.9)--(0.8,1.81)--(1.2,1.72)--(1.6,1.64)--(2,1.56)--(2.4,1.48)--(2.8,1.41)--(3.2,1.34)--(3.6,1.28)--(4,1.21);
	
	\draw[color={red},line width = 2] (-10,4.5)--(-9.6,4.4)--(-9.2,4.3)--(-8.8,4.2)--(-8.4,4.1)--(-8,4)--(-7.6,3.9)--(-7.2,3.8)--(-6.8,3.7)--(-6.4,3.6)--(-6,3.5)--(-5.6,3.4)--(-5.2,3.3)--(-4.8,3.2)--(-4.4,3.1)--(-4,3)--(-3.6,2.9)--(-3.2,2.8)--(-2.8,2.7)--(-2.4,2.6)--(-2,2.5)--(-1.6,2.4)--(-1.2,2.3)--(-0.8,2.2)--(-0.4,2.1)--(0,2)--(0.4,1.9)--(0.8,1.8)--(1.2,1.7)--(1.6,1.6)--(2,1.5)--(2.4,1.4)--(2.8,1.3)--(3.2,1.2)--(3.6,1.1)--(4,1);

\end{tikzpicture}\end{center}
\end{exercice}

\begin{Solution}
 $f'(0)$ est donné par le coefficient directeur de la tangente à la courbe au point d'abscisse $0$, soit $\dfrac{-1}{4}=-\dfrac{1}{4}$.
\end{Solution}

% @see : https://coopmaths.fr/alea?uuid=6f32d&id=can1F16&n=1&d=10&alea=dqSw22323232323423456789101112131415161718192021222324252627282930313233343536373839404142434445464748234567891011121314151617181920&cd=1&cols=1
\needspace{10\baselineskip}
\newpage
\begin{exercice}[BaremeTotal=false]


 On donne les représentations graphiques d'une fonction et de sa dérivée.\\
      Donner l'équation réduite de la tangente à la courbe de $f$ en $x=1$. \\ \begin{minipage}{0.5\linewidth}
    \begin{tikzpicture}[baseline, scale=0.8]

    \tikzset{
      point/.style={
        thick,
        draw,
        cross out,
        inner sep=0pt,
        minimum width=5pt,
        minimum height=5pt,
      },
    }
    \clip (-5.5,-9.5) rectangle (5.75,7.5);
    	
	\draw[color ={black},>=latex,->] (-4,0)--(4,0);
	\draw[color ={black},>=latex,->] (0,-8)--(0,6);
	\draw[color ={black},opacity = 0.4] (-4,1)--(4,1);
	\draw[color ={black},opacity = 0.4] (-4,2)--(4,2);
	\draw[color ={black},opacity = 0.4] (-4,3)--(4,3);
	\draw[color ={black},opacity = 0.4] (-4,4)--(4,4);
	\draw[color ={black},opacity = 0.4] (-4,5)--(4,5);
	\draw[color ={black},opacity = 0.4] (-4,6)--(4,6);
	\draw[color ={black},opacity = 0.4] (-4,-1)--(4,-1);
	\draw[color ={black},opacity = 0.4] (-4,-2)--(4,-2);
	\draw[color ={black},opacity = 0.4] (-4,-3)--(4,-3);
	\draw[color ={black},opacity = 0.4] (-4,-4)--(4,-4);
	\draw[color ={black},opacity = 0.4] (-4,-5)--(4,-5);
	\draw[color ={black},opacity = 0.4] (-4,-6)--(4,-6);
	\draw[color ={black},opacity = 0.4] (-4,-7)--(4,-7);
	\draw[color ={black},opacity = 0.4] (-4,-8)--(4,-8);
	\draw[color ={black},opacity = 0.4] (1,-8)--(1,6);
	\draw[color ={black},opacity = 0.4] (2,-8)--(2,6);
	\draw[color ={black},opacity = 0.4] (3,-8)--(3,6);
	\draw[color ={black},opacity = 0.4] (4,-8)--(4,6);
	\draw[color ={black},opacity = 0.4] (-1,-8)--(-1,6);
	\draw[color ={black},opacity = 0.4] (-2,-8)--(-2,6);
	\draw[color ={black},opacity = 0.4] (-3,-8)--(-3,6);
	\draw[color ={black},opacity = 0.4] (-4,-8)--(-4,6);
	\draw[color ={gray},opacity = 0.3] (-4,0)--(4,0);
	\draw[color ={gray},opacity = 0.3] (-4,0.5)--(4,0.5);
	\draw[color ={gray},opacity = 0.3] (-4,1.5)--(4,1.5);
	\draw[color ={gray},opacity = 0.3] (-4,2.5)--(4,2.5);
	\draw[color ={gray},opacity = 0.3] (-4,3.5)--(4,3.5);
	\draw[color ={gray},opacity = 0.3] (-4,4.5)--(4,4.5);
	\draw[color ={gray},opacity = 0.3] (-4,5.5)--(4,5.5);
	\draw[color ={gray},opacity = 0.3] (-4,0)--(4,0);
	\draw[color ={gray},opacity = 0.3] (-4,-0.5)--(4,-0.5);
	\draw[color ={gray},opacity = 0.3] (-4,-1.5)--(4,-1.5);
	\draw[color ={gray},opacity = 0.3] (-4,-2.5)--(4,-2.5);
	\draw[color ={gray},opacity = 0.3] (-4,-3.5)--(4,-3.5);
	\draw[color ={gray},opacity = 0.3] (-4,-4.5)--(4,-4.5);
	\draw[color ={gray},opacity = 0.3] (-4,-5.5)--(4,-5.5);
	\draw[color ={gray},opacity = 0.3] (-4,-6.5)--(4,-6.5);
	\draw[color ={gray},opacity = 0.3] (-4,-7)--(4,-7);
	\draw[color ={gray},opacity = 0.3] (-4,-7.5)--(4,-7.5);
	\draw[color ={gray},opacity = 0.3] (-4,-8)--(4,-8);
	\draw[color ={gray},opacity = 0.3] (0,-8)--(0,6);
	\draw[color ={gray},opacity = 0.3] (0,-8)--(0,6);
	\draw[color ={black},line width = 1.2] (1,-0.13)--(1,0.13);
	\draw[color ={black},line width = 1.2] (2,-0.13)--(2,0.13);
	\draw[color ={black},line width = 1.2] (3,-0.13)--(3,0.13);
	\draw[color ={black},line width = 1.2] (4,-0.13)--(4,0.13);
	\draw[color ={black},line width = 1.2] (-1,-0.13)--(-1,0.13);
	\draw[color ={black},line width = 1.2] (-2,-0.13)--(-2,0.13);
	\draw[color ={black},line width = 1.2] (-3,-0.13)--(-3,0.13);
	\draw[color ={black},line width = 1.2] (-4,-0.13)--(-4,0.13);
	\draw[color ={black},line width = 1.2] (-0.13,1)--(0.13,1);
	\draw[color ={black},line width = 1.2] (-0.13,2)--(0.13,2);
	\draw[color ={black},line width = 1.2] (-0.13,3)--(0.13,3);
	\draw[color ={black},line width = 1.2] (-0.13,4)--(0.13,4);
	\draw[color ={black},line width = 1.2] (-0.13,5)--(0.13,5);
	\draw[color ={black},line width = 1.2] (-0.13,6)--(0.13,6);
	\draw[color ={black},line width = 1.2] (-0.13,-1)--(0.13,-1);
	\draw[color ={black},line width = 1.2] (-0.13,-2)--(0.13,-2);
	\draw[color ={black},line width = 1.2] (-0.13,-3)--(0.13,-3);
	\draw[color ={black},line width = 1.2] (-0.13,-4)--(0.13,-4);
	\draw[color ={black},line width = 1.2] (-0.13,-5)--(0.13,-5);
	\draw[color ={black},line width = 1.2] (-0.13,-6)--(0.13,-6);
	\draw[color ={black},line width = 1.2] (-0.13,-7)--(0.13,-7);
	\draw[color ={black},line width = 1.2] (-0.13,-8)--(0.13,-8);
	\draw (2,-0.4) node[anchor = center] {\scriptsize \color{black}{$2$}};
	\draw (4,-0.4) node[anchor = center] {\scriptsize \color{black}{$4$}};
	\draw (-2,-0.4) node[anchor = center] {\scriptsize \color{black}{$-2$}};
	\draw (-4,-0.4) node[anchor = center] {\scriptsize \color{black}{$-4$}};
	\draw (-0.5,2.1) node[anchor = center] {\footnotesize \color{black}{$2$}};
	\draw (-0.5,4.1) node[anchor = center] {\footnotesize \color{black}{$4$}};
	\draw (-0.5,6.1) node[anchor = center] {\footnotesize \color{black}{$6$}};
	\draw (-0.5,-1.9) node[anchor = center] {\footnotesize \color{black}{$-2$}};
	\draw (-0.5,-3.9) node[anchor = center] {\footnotesize \color{black}{$-4$}};
	\draw (-0.5,-5.9) node[anchor = center] {\footnotesize \color{black}{$-6$}};
	\draw (-0.5,-7.9) node[anchor = center] {\footnotesize \color{black}{$-8$}};
	\draw [color={black}] (-0.3,-0.3) node[anchor = center,scale=1, rotate = 0] {O};
	\draw (3,4) node[anchor = center] {\footnotesize \color{blue}{$\Large \mathcal{C}_f$}};
	
	\draw[color={blue},line width = 2] (-1.6,-8.96)--(-1.4,-7.56)--(-1.2,-6.24)--(-1,-5)--(-0.8,-3.84)--(-0.6,-2.76)--(-0.4,-1.76)--(-0.2,-0.84)--(0,0)--(0.2,0.76)--(0.4,1.44)--(0.6,2.04)--(0.8,2.56)--(1,3)--(1.2,3.36)--(1.4,3.64)--(1.6,3.84)--(1.8,3.96)--(2,4)--(2.2,3.96)--(2.4,3.84)--(2.6,3.64)--(2.8,3.36)--(3,3)--(3.2,2.56)--(3.4,2.04)--(3.6,1.44)--(3.8,0.76)--(4,0);

\end{tikzpicture}\\
    \end{minipage}
    \begin{minipage}{0.5\linewidth}
    \begin{tikzpicture}[baseline, scale=0.8]

    \tikzset{
      point/.style={
        thick,
        draw,
        cross out,
        inner sep=0pt,
        minimum width=5pt,
        minimum height=5pt,
      },
    }
    \clip (-5.5,-9.5) rectangle (6.5,7.5);
    	
	\draw[color ={black},>=latex,->] (-4,0)--(4,0);
	\draw[color ={black},>=latex,->] (0,-8)--(0,6);
	\draw[color ={black},opacity = 0.4] (-4,1)--(4,1);
	\draw[color ={black},opacity = 0.4] (-4,2)--(4,2);
	\draw[color ={black},opacity = 0.4] (-4,3)--(4,3);
	\draw[color ={black},opacity = 0.4] (-4,4)--(4,4);
	\draw[color ={black},opacity = 0.4] (-4,5)--(4,5);
	\draw[color ={black},opacity = 0.4] (-4,6)--(4,6);
	\draw[color ={black},opacity = 0.4] (-4,-1)--(4,-1);
	\draw[color ={black},opacity = 0.4] (-4,-2)--(4,-2);
	\draw[color ={black},opacity = 0.4] (-4,-3)--(4,-3);
	\draw[color ={black},opacity = 0.4] (-4,-4)--(4,-4);
	\draw[color ={black},opacity = 0.4] (-4,-5)--(4,-5);
	\draw[color ={black},opacity = 0.4] (-4,-6)--(4,-6);
	\draw[color ={black},opacity = 0.4] (-4,-7)--(4,-7);
	\draw[color ={black},opacity = 0.4] (-4,-8)--(4,-8);
	\draw[color ={black},opacity = 0.4] (1,-8)--(1,6);
	\draw[color ={black},opacity = 0.4] (2,-8)--(2,6);
	\draw[color ={black},opacity = 0.4] (3,-8)--(3,6);
	\draw[color ={black},opacity = 0.4] (4,-8)--(4,6);
	\draw[color ={black},opacity = 0.4] (-1,-8)--(-1,6);
	\draw[color ={black},opacity = 0.4] (-2,-8)--(-2,6);
	\draw[color ={black},opacity = 0.4] (-3,-8)--(-3,6);
	\draw[color ={black},opacity = 0.4] (-4,-8)--(-4,6);
	\draw[color ={gray},opacity = 0.3] (-4,0)--(4,0);
	\draw[color ={gray},opacity = 0.3] (-4,0.5)--(4,0.5);
	\draw[color ={gray},opacity = 0.3] (-4,1.5)--(4,1.5);
	\draw[color ={gray},opacity = 0.3] (-4,2.5)--(4,2.5);
	\draw[color ={gray},opacity = 0.3] (-4,3.5)--(4,3.5);
	\draw[color ={gray},opacity = 0.3] (-4,4.5)--(4,4.5);
	\draw[color ={gray},opacity = 0.3] (-4,5.5)--(4,5.5);
	\draw[color ={gray},opacity = 0.3] (-4,0)--(4,0);
	\draw[color ={gray},opacity = 0.3] (-4,-0.5)--(4,-0.5);
	\draw[color ={gray},opacity = 0.3] (-4,-1.5)--(4,-1.5);
	\draw[color ={gray},opacity = 0.3] (-4,-2.5)--(4,-2.5);
	\draw[color ={gray},opacity = 0.3] (-4,-3.5)--(4,-3.5);
	\draw[color ={gray},opacity = 0.3] (-4,-4.5)--(4,-4.5);
	\draw[color ={gray},opacity = 0.3] (-4,-5.5)--(4,-5.5);
	\draw[color ={gray},opacity = 0.3] (-4,-6.5)--(4,-6.5);
	\draw[color ={gray},opacity = 0.3] (-4,-7)--(4,-7);
	\draw[color ={gray},opacity = 0.3] (-4,-7.5)--(4,-7.5);
	\draw[color ={gray},opacity = 0.3] (-4,-8)--(4,-8);
	\draw[color ={gray},opacity = 0.3] (0,-8)--(0,6);
	\draw[color ={gray},opacity = 0.3] (0,-8)--(0,6);
	\draw[color ={black},line width = 1.2] (1,-0.13)--(1,0.13);
	\draw[color ={black},line width = 1.2] (2,-0.13)--(2,0.13);
	\draw[color ={black},line width = 1.2] (3,-0.13)--(3,0.13);
	\draw[color ={black},line width = 1.2] (4,-0.13)--(4,0.13);
	\draw[color ={black},line width = 1.2] (-1,-0.13)--(-1,0.13);
	\draw[color ={black},line width = 1.2] (-2,-0.13)--(-2,0.13);
	\draw[color ={black},line width = 1.2] (-3,-0.13)--(-3,0.13);
	\draw[color ={black},line width = 1.2] (-4,-0.13)--(-4,0.13);
	\draw[color ={black},line width = 1.2] (-0.13,1)--(0.13,1);
	\draw[color ={black},line width = 1.2] (-0.13,2)--(0.13,2);
	\draw[color ={black},line width = 1.2] (-0.13,3)--(0.13,3);
	\draw[color ={black},line width = 1.2] (-0.13,4)--(0.13,4);
	\draw[color ={black},line width = 1.2] (-0.13,5)--(0.13,5);
	\draw[color ={black},line width = 1.2] (-0.13,6)--(0.13,6);
	\draw[color ={black},line width = 1.2] (-0.13,-1)--(0.13,-1);
	\draw[color ={black},line width = 1.2] (-0.13,-2)--(0.13,-2);
	\draw[color ={black},line width = 1.2] (-0.13,-3)--(0.13,-3);
	\draw[color ={black},line width = 1.2] (-0.13,-4)--(0.13,-4);
	\draw[color ={black},line width = 1.2] (-0.13,-5)--(0.13,-5);
	\draw[color ={black},line width = 1.2] (-0.13,-6)--(0.13,-6);
	\draw[color ={black},line width = 1.2] (-0.13,-7)--(0.13,-7);
	\draw[color ={black},line width = 1.2] (-0.13,-8)--(0.13,-8);
	\draw (2,-0.4) node[anchor = center] {\scriptsize \color{black}{$2$}};
	\draw (4,-0.4) node[anchor = center] {\scriptsize \color{black}{$4$}};
	\draw (-2,-0.4) node[anchor = center] {\scriptsize \color{black}{$-2$}};
	\draw (-4,-0.4) node[anchor = center] {\scriptsize \color{black}{$-4$}};
	\draw (-0.5,2.1) node[anchor = center] {\footnotesize \color{black}{$2$}};
	\draw (-0.5,4.1) node[anchor = center] {\footnotesize \color{black}{$4$}};
	\draw (-0.5,-1.9) node[anchor = center] {\footnotesize \color{black}{$-2$}};
	\draw (-0.5,-3.9) node[anchor = center] {\footnotesize \color{black}{$-4$}};
	\draw (-0.5,-5.9) node[anchor = center] {\footnotesize \color{black}{$-6$}};
	\draw (-0.5,-7.9) node[anchor = center] {\footnotesize \color{black}{$-8$}};
	\draw [color={black}] (-0.3,-0.3) node[anchor = center,scale=1, rotate = 0] {O};
	\draw (3,4) node[anchor = center] {\footnotesize \color{red}{$\Large\mathcal{C}_f\prime$}};
	
	\draw[color={red},line width = 2] (-1.4,6.8)--(-1.2,6.4)--(-1,6)--(-0.8,5.6)--(-0.6,5.2)--(-0.4,4.8)--(-0.2,4.4)--(0,4)--(0.2,3.6)--(0.4,3.2)--(0.6,2.8)--(0.8,2.4)--(1,2)--(1.2,1.6)--(1.4,1.2)--(1.6,0.8)--(1.8,0.4)--(2,0)--(2.2,-0.4)--(2.4,-0.8)--(2.6,-1.2)--(2.8,-1.6)--(3,-2)--(3.2,-2.4)--(3.4,-2.8)--(3.6,-3.2)--(3.8,-3.6)--(4,-4);

\end{tikzpicture}\\
    \end{minipage}
    
\end{exercice}

\begin{Solution}
 L'équation réduite de la tangente au point d'abscisse $1$ est  : $y=f'(1)(x-1)+f(1)$.\\
      On lit graphiquement $f(1)=3$ et $f'(1)=2$.\\
      L'équation réduite de la tangente est donc donnée par :
      $y=2(x-1)+3$, soit $y=2x+1$.
\end{Solution}

% @see : https://coopmaths.fr/alea?uuid=a1ba2&id=can1F14&n=1&d=10&alea=rkIp22323232323423456789101112131415161718192021222324252627282930313233343536373839404142434445464748234567891011121314151617181920&cd=1&cols=1
\needspace{10\baselineskip}
\begin{exercice}[BaremeTotal=false]


Soit $f$ la fonction définie sur $\mathbb{R}$ par :
$f(x)= 5+2x^2$.\\
En calculant la limite du taux d'accroissement, déterminer $f'(3)$.
\end{exercice}

\begin{Solution}
 $f$ est une fonction polynôme du second degré de la forme $f(x)=ax^2+b$.\\
    La fonction dérivée est donnée par la somme des dérivées des fonctions $u$ et $v$ définies par $u(x)=2x^2$ et $v(x)=5$.\\
     Comme $u'(x)=4x$ et $v'(x)=0$, on obtient  $f'(x)=4x$. \\
     Ainsi, $f'(3)=4\times 3=12$.
\end{Solution}

% @see : https://coopmaths.fr/alea?uuid=3c690&id=can1F13&n=1&d=10&alea=EUDu22323232323423456789101112131415161718192021222324252627282930313233343536373839404142434445464748234567891011121314151617181920&cd=1&cols=1
\needspace{10\baselineskip}
\begin{exercice}[BaremeTotal=false]


 Déterminer le coefficient directeur de la tangente à la courbe représentative de la fonction cube au point d'abscisse $-4$.
                        
\end{exercice}

\begin{Solution}
 Le coefficient directeur de la tangente au point d'abscisse $a$ est donné par le nombre dérivé $f'(a)$.\\
        La fonction $f$ définie par $f(x)=x^3$ a pour fonction dérivée la fonction $f'$ définie par $f'(x)=3x^2$.\\
        Comme $f'(-4)=3\times (-4)^2=48$, le coefficient directeur de la tangente au point d'abscisse $-4$ est : $48$. 
\end{Solution}

% @see : https://coopmaths.fr/alea?uuid=ebd89&id=1AN14-1&n=1&d=10&s=3&alea=ocYr22323232323423456789101112131415161718192021222324252627282930313233343536373839404142434445464748234567891011121314151617181920&cd=0&cols=1
\needspace{10\baselineskip}
\begin{exercice}[BaremeTotal=false]


 Donner la dérivée de la fonction $f$, dérivable pour tout $x\in\R$, définie par  $f(x)=-\dfrac{1}{9}x+1$.
\end{exercice}

\begin{Solution}
 L'expression de la dérivée de la fonction $f$ définie par $f(x)=-\dfrac{1}{9}x+1$ est : ${\color[HTML]{f15929}\boldsymbol{f'(x)=-\dfrac{1}{9}}}$.
\end{Solution}

% @see : https://coopmaths.fr/alea?uuid=ebd89&id=1AN14-1&n=1&d=10&s=4&alea=uAlP22323232323423456789101112131415161718192021222324252627282930313233343536373839404142434445464748234567891011121314151617181920&cd=0&cols=1
\needspace{10\baselineskip}
\begin{exercice}[BaremeTotal=false]


 Donner la dérivée de la fonction $f$, dérivable pour tout $x\in\R$, définie par  $f(x)=-10x^{14}$.
\end{exercice}

\begin{Solution}
 L'expression de la dérivée de la fonction $f$ définie par $f(x)=-10x^{14}$ est : ${\color[HTML]{f15929}\boldsymbol{f'(x)=-140x^{13}}}$.
\end{Solution}

% @see : https://coopmaths.fr/alea?uuid=ebd89&id=1AN14-1&n=1&d=10&s=7&alea=vjN722323232323423456789101112131415161718192021222324252627282930313233343536373839404142434445464748234567891011121314151617181920&cd=0&cols=1
\needspace{10\baselineskip}
\begin{exercice}[BaremeTotal=false]


 Donner la dérivée de la fonction $f$, dérivable pour tout $x\in\R^*$, définie par  $f(x)=\dfrac{-1}{3x}$.
\end{exercice}

\begin{Solution}
 L'expression de la dérivée de la fonction $f$ définie par $f(x)=\dfrac{-1}{3x}$ est : ${\color[HTML]{f15929}\boldsymbol{f'(x)=\dfrac{1}{3x^2}}}$.
\end{Solution}

% @see : https://coopmaths.fr/alea?uuid=a83c0&id=1AN14-4&n=1&d=10&s=1&alea=Njr322323232323423456789101112131415161718192021222324252627282930313233343536373839404142434445464748234567891011121314151617181920&cd=0&cols=1
\needspace{10\baselineskip}
\begin{exercice}[BaremeTotal=false]


 Donner l'expression de la dérivée de la fonction $f$ définie sur $\R^{*}$ par $f(x)=-2x^{2}-5x-2-\dfrac{1}{x}$.\\
\end{exercice}

\begin{Solution}
 L'expression de la dérivée de la fonction $f$ définie par $f(x)=-2x^{2}-5x-2-\dfrac{1}{x}$ est :\\$f'(x)={\color[HTML]{f15929}\boldsymbol{-4x-5+\dfrac{1}{x^2}}}$.\\
\end{Solution}

% @see : https://coopmaths.fr/alea?uuid=a83c0&id=1AN14-4&n=1&d=10&s=4&alea=4Vpz22323232323423456789101112131415161718192021222324252627282930313233343536373839404142434445464748234567891011121314151617181920&cd=1&cols=1
\needspace{10\baselineskip}
\begin{exercice}[BaremeTotal=false]


 Donner l'expression de la dérivée de la fonction $f$ définie sur $\R_+^{*}$ par $f(x)=-\dfrac{4}{x}+2\sqrt{x}+2x^{2}$.\\
\end{exercice}

\begin{Solution}
 $(-\dfrac{4}{x})^\prime=\dfrac{4}{x^2}$\\$(2\sqrt{x})^\prime=\dfrac{2}{2\sqrt{x}}$\\$(2x^{2})^\prime=4x$\\L'expression de la dérivée de la fonction $f$ définie par $f(x)=-\dfrac{4}{x}+2\sqrt{x}+2x^{2}$ est :\\$f'(x)={\color[HTML]{f15929}\boldsymbol{\dfrac{4}{x^2}+\dfrac{2}{2\sqrt{x}}+4x}}$.\\
\end{Solution}

\end{Maquette}
\clearpage
\begin{Maquette}[DM]{Niveau= ,Classe=\getdata[21]\mydata\ (v21), Date = 06/01/25  ,Theme=Exercices}


% @see : https://coopmaths.fr/alea?uuid=4c8c7&id=1AN11-3&n=1&d=10&s=2&alea=nawO2232323232342345678910111213141516171819202122232425262728293031323334353637383940414243444546474823456789101112131415161718192021&cd=1&cols=2
\needspace{10\baselineskip}
\begin{exercice}[BaremeTotal=false]
\label{eleve:21}


 \begin{minipage}[t]{\linewidth} Soit $f$ une fonction dérivable sur $[-5;5]$ et $\mathcal{C}_f$ sa courbe représentative.\\On sait que  $f(2)=-4~~$ et que $~~f'(2)=-4$.\\Déterminer une équation de la tangente $(T)$ à la courbe $\mathcal{C}_f$ au point d'abscisse $2$,\\sans utiliser la formule de cours de l'équation de tangente. \end{minipage}

\end{exercice}

\begin{Solution}
 \begin{minipage}[t]{\linewidth} On sait que la tangente n'est pas une droite verticale, puisque la fonction est dérivable sur l'intervalle.\\On en déduit que la tangente $(T)$ au point d'abscisse $2$, admet une équation réduite de la forme :  $(T) : y=m x + p$. 

\medskip
$\bullet$ Détermination de $m$ :\\On sait que le nombre dérivé en $2$ est par définition, le coefficient directeur de la tangente au point d'abscisse $2$.\\Par conséquent, on a déjà : $m=f'(2)=-4$.\\ On en déduit que  $(T) : y= -4 x + p$\\$\bullet$ Détermination de $p$ :\\ Pour cela, on utilise que si $f(2)=-4~~$, alors le point $A$ de coordonnées $(2;-4)$ appartient à $\mathcal{C}_f$ mais aussi à $(T)$.\\ On peut écrire $A(2;-4) = \mathcal{C}_f \cap (T)$.\\ On remplace alors les coordonnées de $A(2;-4)$ dans l'équation  $(T) : y= -4 x + p$.\\ $\begin{aligned} \phantom{\iff}&A(2;-4)\in (T)\\ \iff& -4= -4 \times 2  + p\\ \iff& p=-4 +4 \times 2  \\ \iff& p=-4 +8   \\ \iff& p=4\\\end{aligned}$\\On peut conclure que : $(T) : {\color[HTML]{f15929}\boldsymbol{y=-4x+4}}$. \end{minipage}

\end{Solution}

% @see : https://coopmaths.fr/alea?uuid=0e984&id=can1F15&n=1&d=10&alea=1Alo2232323232342345678910111213141516171819202122232425262728293031323334353637383940414243444546474823456789101112131415161718192021&cd=1&cols=1
\needspace{10\baselineskip}
\begin{exercice}[BaremeTotal=false]


 La courbe représente une fonction $f$ et la droite est la tangente au point d'abscisse $0$.\\
        Déterminer $f'(0)$.\begin{center}\begin{tikzpicture}[baseline,scale = 0.5]

    \tikzset{
      point/.style={
        thick,
        draw,
        cross out,
        inner sep=0pt,
        minimum width=5pt,
        minimum height=5pt,
      },
    }
    \clip (-10,-2) rectangle (4,12);
    	
	\draw[color ={black},line width = 2,>=latex,->] (-10,0)--(4,0);
	\draw[color ={black},line width = 2,>=latex,->] (0,-2)--(0,12);
	\draw[color ={black},opacity = 0.5] (-10,1)--(4,1);
	\draw[color ={black},opacity = 0.5] (-10,2)--(4,2);
	\draw[color ={black},opacity = 0.5] (-10,3)--(4,3);
	\draw[color ={black},opacity = 0.5] (-10,4)--(4,4);
	\draw[color ={black},opacity = 0.5] (-10,5)--(4,5);
	\draw[color ={black},opacity = 0.5] (-10,6)--(4,6);
	\draw[color ={black},opacity = 0.5] (-10,7)--(4,7);
	\draw[color ={black},opacity = 0.5] (-10,8)--(4,8);
	\draw[color ={black},opacity = 0.5] (-10,9)--(4,9);
	\draw[color ={black},opacity = 0.5] (-10,10)--(4,10);
	\draw[color ={black},opacity = 0.5] (-10,11)--(4,11);
	\draw[color ={black},opacity = 0.5] (-10,12)--(4,12);
	\draw[color ={black},opacity = 0.5] (-10,-1)--(4,-1);
	\draw[color ={black},opacity = 0.5] (-10,-2)--(4,-2);
	\draw[color ={black},opacity = 0.5] (1,-2)--(1,12);
	\draw[color ={black},opacity = 0.5] (2,-2)--(2,12);
	\draw[color ={black},opacity = 0.5] (3,-2)--(3,12);
	\draw[color ={black},opacity = 0.5] (4,-2)--(4,12);
	\draw[color ={black},opacity = 0.5] (-1,-2)--(-1,12);
	\draw[color ={black},opacity = 0.5] (-2,-2)--(-2,12);
	\draw[color ={black},opacity = 0.5] (-3,-2)--(-3,12);
	\draw[color ={black},opacity = 0.5] (-4,-2)--(-4,12);
	\draw[color ={black},opacity = 0.5] (-5,-2)--(-5,12);
	\draw[color ={black},opacity = 0.5] (-6,-2)--(-6,12);
	\draw[color ={black},opacity = 0.5] (-7,-2)--(-7,12);
	\draw[color ={black},opacity = 0.5] (-8,-2)--(-8,12);
	\draw[color ={black},opacity = 0.5] (-9,-2)--(-9,12);
	\draw[color ={black},opacity = 0.5] (-10,-2)--(-10,12);
	\draw[color ={gray},opacity = 0.3] (-10,0)--(4,0);
	\draw[color ={gray},opacity = 0.3] (-10,0)--(4,0);
	\draw[color ={gray},opacity = 0.3] (0,-2)--(0,12);
	\draw[color ={gray},opacity = 0.3] (0,-2)--(0,12);
	\draw[color ={gray},opacity = 0.3] (-5,-2)--(-5,12);
	\draw[color ={gray},opacity = 0.3] (-6,-2)--(-6,12);
	\draw[color ={gray},opacity = 0.3] (-7,-2)--(-7,12);
	\draw[color ={gray},opacity = 0.3] (-8,-2)--(-8,12);
	\draw[color ={gray},opacity = 0.3] (-9,-2)--(-9,12);
	\draw[color ={gray},opacity = 0.3] (-10,-2)--(-10,12);
	\draw[color ={black},line width = 2] (2,-0.2)--(2,0.2);
	\draw[color ={black},line width = 2] (-2,-0.2)--(-2,0.2);
	\draw[color ={black},line width = 2] (-4,-0.2)--(-4,0.2);
	\draw[color ={black},line width = 2] (-6,-0.2)--(-6,0.2);
	\draw[color ={black},line width = 2] (-8,-0.2)--(-8,0.2);
	\draw[color ={black},line width = 2] (-0.2,2)--(0.2,2);
	\draw[color ={black},line width = 2] (-0.2,4)--(0.2,4);
	\draw[color ={black},line width = 2] (-0.2,6)--(0.2,6);
	\draw[color ={black},line width = 2] (-0.2,8)--(0.2,8);
	\draw[color ={black},line width = 2] (-0.2,10)--(0.2,10);
	\draw (2,-0.4) node[anchor = center] {\scriptsize \color{black}{$1$}};
	\draw (-2,-0.4) node[anchor = center] {\scriptsize \color{black}{$-1$}};
	\draw (-4,-0.4) node[anchor = center] {\scriptsize \color{black}{$-2$}};
	\draw (-6,-0.4) node[anchor = center] {\scriptsize \color{black}{$-3$}};
	\draw (-8,-0.4) node[anchor = center] {\scriptsize \color{black}{$-4$}};
	\draw (-0.5,2.1) node[anchor = center] {\footnotesize \color{black}{$1$}};
	\draw (-0.5,4.1) node[anchor = center] {\footnotesize \color{black}{$2$}};
	\draw (-0.5,6.1) node[anchor = center] {\footnotesize \color{black}{$3$}};
	\draw (-0.5,8.1) node[anchor = center] {\footnotesize \color{black}{$4$}};
	\draw (-0.5,10.1) node[anchor = center] {\footnotesize \color{black}{$5$}};
	\draw [color={black}] (-0.3,-0.3) node[anchor = center,scale=1, rotate = 0] {O};
	
	\draw[color={blue},line width = 2] (-10,10.59)--(-9.6,9.91)--(-9.2,9.27)--(-8.8,8.67)--(-8.4,8.11)--(-8,7.59)--(-7.6,7.1)--(-7.2,6.64)--(-6.8,6.21)--(-6.4,5.81)--(-6,5.44)--(-5.6,5.09)--(-5.2,4.76)--(-4.8,4.45)--(-4.4,4.16)--(-4,3.9)--(-3.6,3.64)--(-3.2,3.41)--(-2.8,3.19)--(-2.4,2.98)--(-2,2.79)--(-1.6,2.61)--(-1.2,2.44)--(-0.8,2.29)--(-0.4,2.14)--(0,2)--(0.4,1.87)--(0.8,1.75)--(1.2,1.64)--(1.6,1.53)--(2,1.43)--(2.4,1.34)--(2.8,1.25)--(3.2,1.17)--(3.6,1.1)--(4,1.03);
	
	\draw[color={red},line width = 2] (-10,5.33)--(-9.6,5.2)--(-9.2,5.07)--(-8.8,4.93)--(-8.4,4.8)--(-8,4.67)--(-7.6,4.53)--(-7.2,4.4)--(-6.8,4.27)--(-6.4,4.13)--(-6,4)--(-5.6,3.87)--(-5.2,3.73)--(-4.8,3.6)--(-4.4,3.47)--(-4,3.33)--(-3.6,3.2)--(-3.2,3.07)--(-2.8,2.93)--(-2.4,2.8)--(-2,2.67)--(-1.6,2.53)--(-1.2,2.4)--(-0.8,2.27)--(-0.4,2.13)--(0,2)--(0.4,1.87)--(0.8,1.73)--(1.2,1.6)--(1.6,1.47)--(2,1.33)--(2.4,1.2)--(2.8,1.07)--(3.2,0.93)--(3.6,0.8)--(4,0.67);

\end{tikzpicture}\end{center}
\end{exercice}

\begin{Solution}
 $f'(0)$ est donné par le coefficient directeur de la tangente à la courbe au point d'abscisse $0$, soit $\dfrac{-1}{3}=-\dfrac{1}{3}$.
\end{Solution}

% @see : https://coopmaths.fr/alea?uuid=6f32d&id=can1F16&n=1&d=10&alea=dqSw2232323232342345678910111213141516171819202122232425262728293031323334353637383940414243444546474823456789101112131415161718192021&cd=1&cols=1
\needspace{10\baselineskip}
\newpage
\begin{exercice}[BaremeTotal=false]


 On donne les représentations graphiques d'une fonction et de sa dérivée.\\
        Donner l'équation réduite de la tangente à la courbe de $f$ en $x=-2$. \\ \begin{minipage}{0.5\linewidth}
    \begin{tikzpicture}[baseline, scale=0.8]

    \tikzset{
      point/.style={
        thick,
        draw,
        cross out,
        inner sep=0pt,
        minimum width=5pt,
        minimum height=5pt,
      },
    }
    \clip (-4.5,-4.5) rectangle (5.75,13.5);
    	
	\draw[color ={black},>=latex,->] (-3,0)--(3,0);
	\draw[color ={black},>=latex,->] (0,-3)--(0,12);
	\draw[color ={black},opacity = 0.4] (-3,1)--(3,1);
	\draw[color ={black},opacity = 0.4] (-3,2)--(3,2);
	\draw[color ={black},opacity = 0.4] (-3,3)--(3,3);
	\draw[color ={black},opacity = 0.4] (-3,4)--(3,4);
	\draw[color ={black},opacity = 0.4] (-3,5)--(3,5);
	\draw[color ={black},opacity = 0.4] (-3,6)--(3,6);
	\draw[color ={black},opacity = 0.4] (-3,7)--(3,7);
	\draw[color ={black},opacity = 0.4] (-3,8)--(3,8);
	\draw[color ={black},opacity = 0.4] (-3,9)--(3,9);
	\draw[color ={black},opacity = 0.4] (-3,10)--(3,10);
	\draw[color ={black},opacity = 0.4] (-3,11)--(3,11);
	\draw[color ={black},opacity = 0.4] (-3,12)--(3,12);
	\draw[color ={black},opacity = 0.4] (-3,-1)--(3,-1);
	\draw[color ={black},opacity = 0.4] (-3,-2)--(3,-2);
	\draw[color ={black},opacity = 0.4] (-3,-3)--(3,-3);
	\draw[color ={black},opacity = 0.4] (1,-3)--(1,12);
	\draw[color ={black},opacity = 0.4] (2,-3)--(2,12);
	\draw[color ={black},opacity = 0.4] (3,-3)--(3,12);
	\draw[color ={black},opacity = 0.4] (-1,-3)--(-1,12);
	\draw[color ={black},opacity = 0.4] (-2,-3)--(-2,12);
	\draw[color ={black},opacity = 0.4] (-3,-3)--(-3,12);
	\draw[color ={gray},opacity = 0.3] (-3,0)--(3,0);
	\draw[color ={gray},opacity = 0.3] (-3,0.5)--(3,0.5);
	\draw[color ={gray},opacity = 0.3] (-3,1.5)--(3,1.5);
	\draw[color ={gray},opacity = 0.3] (-3,2.5)--(3,2.5);
	\draw[color ={gray},opacity = 0.3] (-3,3.5)--(3,3.5);
	\draw[color ={gray},opacity = 0.3] (-3,4.5)--(3,4.5);
	\draw[color ={gray},opacity = 0.3] (-3,5.5)--(3,5.5);
	\draw[color ={gray},opacity = 0.3] (-3,6.5)--(3,6.5);
	\draw[color ={gray},opacity = 0.3] (-3,7.5)--(3,7.5);
	\draw[color ={gray},opacity = 0.3] (-3,8.5)--(3,8.5);
	\draw[color ={gray},opacity = 0.3] (-3,9.5)--(3,9.5);
	\draw[color ={gray},opacity = 0.3] (-3,10.5)--(3,10.5);
	\draw[color ={gray},opacity = 0.3] (-3,11.5)--(3,11.5);
	\draw[color ={gray},opacity = 0.3] (-3,0)--(3,0);
	\draw[color ={gray},opacity = 0.3] (-3,-0.5)--(3,-0.5);
	\draw[color ={gray},opacity = 0.3] (-3,-1.5)--(3,-1.5);
	\draw[color ={gray},opacity = 0.3] (-3,-2.5)--(3,-2.5);
	\draw[color ={gray},opacity = 0.3] (0,-3)--(0,12);
	\draw[color ={gray},opacity = 0.3] (0,-3)--(0,12);
	\draw[color ={black},line width = 1.2] (1,-0.13)--(1,0.13);
	\draw[color ={black},line width = 1.2] (2,-0.13)--(2,0.13);
	\draw[color ={black},line width = 1.2] (3,-0.13)--(3,0.13);
	\draw[color ={black},line width = 1.2] (-1,-0.13)--(-1,0.13);
	\draw[color ={black},line width = 1.2] (-2,-0.13)--(-2,0.13);
	\draw[color ={black},line width = 1.2] (-3,-0.13)--(-3,0.13);
	\draw[color ={black},line width = 1.2] (-0.13,1)--(0.13,1);
	\draw[color ={black},line width = 1.2] (-0.13,2)--(0.13,2);
	\draw[color ={black},line width = 1.2] (-0.13,3)--(0.13,3);
	\draw[color ={black},line width = 1.2] (-0.13,4)--(0.13,4);
	\draw[color ={black},line width = 1.2] (-0.13,5)--(0.13,5);
	\draw[color ={black},line width = 1.2] (-0.13,6)--(0.13,6);
	\draw[color ={black},line width = 1.2] (-0.13,7)--(0.13,7);
	\draw[color ={black},line width = 1.2] (-0.13,8)--(0.13,8);
	\draw[color ={black},line width = 1.2] (-0.13,9)--(0.13,9);
	\draw[color ={black},line width = 1.2] (-0.13,10)--(0.13,10);
	\draw[color ={black},line width = 1.2] (-0.13,11)--(0.13,11);
	\draw[color ={black},line width = 1.2] (-0.13,12)--(0.13,12);
	\draw[color ={black},line width = 1.2] (-0.13,-1)--(0.13,-1);
	\draw[color ={black},line width = 1.2] (-0.13,-2)--(0.13,-2);
	\draw[color ={black},line width = 1.2] (-0.13,-3)--(0.13,-3);
	\draw (3,-0.4) node[anchor = center] {\scriptsize \color{black}{$3$}};
	\draw (-3,-0.4) node[anchor = center] {\scriptsize \color{black}{$-3$}};
	\draw (-0.5,3.1) node[anchor = center] {\footnotesize \color{black}{$3$}};
	\draw (-0.5,6.1) node[anchor = center] {\footnotesize \color{black}{$6$}};
	\draw (-0.5,9.1) node[anchor = center] {\footnotesize \color{black}{$9$}};
	\draw (-0.5,-2.9) node[anchor = center] {\footnotesize \color{black}{$-3$}};
	\draw [color={black}] (-0.3,-0.3) node[anchor = center,scale=1, rotate = 0] {O};
	\draw (3,10) node[anchor = center] {\footnotesize \color{blue}{$\Large \mathcal{C}_f$}};
	
	\draw[color={blue},line width = 2] (-3,7)--(-2.8,5.84)--(-2.6,4.76)--(-2.4,3.76)--(-2.2,2.84)--(-2,2)--(-1.8,1.24)--(-1.6,0.56)--(-1.4,-0.04)--(-1.2,-0.56)--(-1,-1)--(-0.8,-1.36)--(-0.6,-1.64)--(-0.4,-1.84)--(-0.2,-1.96)--(0,-2)--(0.2,-1.96)--(0.4,-1.84)--(0.6,-1.64)--(0.8,-1.36)--(1,-1)--(1.2,-0.56)--(1.4,-0.04)--(1.6,0.56)--(1.8,1.24)--(2,2)--(2.2,2.84)--(2.4,3.76)--(2.6,4.76)--(2.8,5.84)--(3,7);

\end{tikzpicture}\\
    \end{minipage}
    \begin{minipage}{0.5\linewidth}
    \begin{tikzpicture}[baseline, scale=0.8]

    \tikzset{
      point/.style={
        thick,
        draw,
        cross out,
        inner sep=0pt,
        minimum width=5pt,
        minimum height=5pt,
      },
    }
    \clip (-4.5,-6.5) rectangle (6.5,9.5);
    	
	\draw[color ={black},>=latex,->] (-3,0)--(3,0);
	\draw[color ={black},>=latex,->] (0,-5)--(0,8);
	\draw[color ={black},opacity = 0.4] (-3,1)--(3,1);
	\draw[color ={black},opacity = 0.4] (-3,2)--(3,2);
	\draw[color ={black},opacity = 0.4] (-3,3)--(3,3);
	\draw[color ={black},opacity = 0.4] (-3,4)--(3,4);
	\draw[color ={black},opacity = 0.4] (-3,5)--(3,5);
	\draw[color ={black},opacity = 0.4] (-3,6)--(3,6);
	\draw[color ={black},opacity = 0.4] (-3,7)--(3,7);
	\draw[color ={black},opacity = 0.4] (-3,8)--(3,8);
	\draw[color ={black},opacity = 0.4] (-3,-1)--(3,-1);
	\draw[color ={black},opacity = 0.4] (-3,-2)--(3,-2);
	\draw[color ={black},opacity = 0.4] (-3,-3)--(3,-3);
	\draw[color ={black},opacity = 0.4] (-3,-4)--(3,-4);
	\draw[color ={black},opacity = 0.4] (-3,-5)--(3,-5);
	\draw[color ={black},opacity = 0.4] (1,-5)--(1,8);
	\draw[color ={black},opacity = 0.4] (2,-5)--(2,8);
	\draw[color ={black},opacity = 0.4] (3,-5)--(3,8);
	\draw[color ={black},opacity = 0.4] (-1,-5)--(-1,8);
	\draw[color ={black},opacity = 0.4] (-2,-5)--(-2,8);
	\draw[color ={black},opacity = 0.4] (-3,-5)--(-3,8);
	\draw[color ={gray},opacity = 0.3] (-3,0)--(3,0);
	\draw[color ={gray},opacity = 0.3] (-3,0.5)--(3,0.5);
	\draw[color ={gray},opacity = 0.3] (-3,1.5)--(3,1.5);
	\draw[color ={gray},opacity = 0.3] (-3,2.5)--(3,2.5);
	\draw[color ={gray},opacity = 0.3] (-3,3.5)--(3,3.5);
	\draw[color ={gray},opacity = 0.3] (-3,4.5)--(3,4.5);
	\draw[color ={gray},opacity = 0.3] (-3,5.5)--(3,5.5);
	\draw[color ={gray},opacity = 0.3] (-3,6.5)--(3,6.5);
	\draw[color ={gray},opacity = 0.3] (-3,7.5)--(3,7.5);
	\draw[color ={gray},opacity = 0.3] (-3,0)--(3,0);
	\draw[color ={gray},opacity = 0.3] (-3,-0.5)--(3,-0.5);
	\draw[color ={gray},opacity = 0.3] (-3,-1.5)--(3,-1.5);
	\draw[color ={gray},opacity = 0.3] (-3,-2.5)--(3,-2.5);
	\draw[color ={gray},opacity = 0.3] (-3,-3.5)--(3,-3.5);
	\draw[color ={gray},opacity = 0.3] (-3,-4.5)--(3,-4.5);
	\draw[color ={gray},opacity = 0.3] (0,-5)--(0,8);
	\draw[color ={gray},opacity = 0.3] (0,-5)--(0,8);
	\draw[color ={black},line width = 1.2] (1,-0.13)--(1,0.13);
	\draw[color ={black},line width = 1.2] (2,-0.13)--(2,0.13);
	\draw[color ={black},line width = 1.2] (3,-0.13)--(3,0.13);
	\draw[color ={black},line width = 1.2] (-1,-0.13)--(-1,0.13);
	\draw[color ={black},line width = 1.2] (-2,-0.13)--(-2,0.13);
	\draw[color ={black},line width = 1.2] (-3,-0.13)--(-3,0.13);
	\draw[color ={black},line width = 1.2] (-0.13,1)--(0.13,1);
	\draw[color ={black},line width = 1.2] (-0.13,2)--(0.13,2);
	\draw[color ={black},line width = 1.2] (-0.13,3)--(0.13,3);
	\draw[color ={black},line width = 1.2] (-0.13,4)--(0.13,4);
	\draw[color ={black},line width = 1.2] (-0.13,5)--(0.13,5);
	\draw[color ={black},line width = 1.2] (-0.13,6)--(0.13,6);
	\draw[color ={black},line width = 1.2] (-0.13,7)--(0.13,7);
	\draw[color ={black},line width = 1.2] (-0.13,8)--(0.13,8);
	\draw[color ={black},line width = 1.2] (-0.13,9)--(0.13,9);
	\draw[color ={black},line width = 1.2] (-0.13,10)--(0.13,10);
	\draw[color ={black},line width = 1.2] (-0.13,11)--(0.13,11);
	\draw[color ={black},line width = 1.2] (-0.13,12)--(0.13,12);
	\draw[color ={black},line width = 1.2] (-0.13,-1)--(0.13,-1);
	\draw[color ={black},line width = 1.2] (-0.13,-2)--(0.13,-2);
	\draw[color ={black},line width = 1.2] (-0.13,-3)--(0.13,-3);
	\draw (3,-0.4) node[anchor = center] {\scriptsize \color{black}{$3$}};
	\draw (-3,-0.4) node[anchor = center] {\scriptsize \color{black}{$-3$}};
	\draw (-0.5,3.1) node[anchor = center] {\footnotesize \color{black}{$3$}};
	\draw (-0.5,6.1) node[anchor = center] {\footnotesize \color{black}{$6$}};
	\draw (-0.5,-2.9) node[anchor = center] {\footnotesize \color{black}{$-3$}};
	\draw [color={black}] (-0.3,-0.3) node[anchor = center,scale=1, rotate = 0] {O};
	\draw (3,6) node[anchor = center] {\footnotesize \color{red}{$\Large\mathcal{C}_f\prime$}};
	
	\draw[color={red},line width = 2] (-2.8,-5.6)--(-2.6,-5.2)--(-2.4,-4.8)--(-2.2,-4.4)--(-2,-4)--(-1.8,-3.6)--(-1.6,-3.2)--(-1.4,-2.8)--(-1.2,-2.4)--(-1,-2)--(-0.8,-1.6)--(-0.6,-1.2)--(-0.4,-0.8)--(-0.2,-0.4)--(0,0)--(0.2,0.4)--(0.4,0.8)--(0.6,1.2)--(0.8,1.6)--(1,2)--(1.2,2.4)--(1.4,2.8)--(1.6,3.2)--(1.8,3.6)--(2,4)--(2.2,4.4)--(2.4,4.8)--(2.6,5.2)--(2.8,5.6)--(3,6);

\end{tikzpicture}\\
    \end{minipage}
    
\end{exercice}

\begin{Solution}
 L'équation réduite de la tangente au point d'abscisse $-2$ est  : $y=f'(-2)(x-(-2))+f(-2)$.\\
        On lit graphiquement $f(-2)=2$ et $f'(-2)=-4$.\\
        L'équation réduite de la tangente est donc donnée par :
        $y=-4(x+2)+2$, soit $y=-4x-6$.
\end{Solution}

% @see : https://coopmaths.fr/alea?uuid=a1ba2&id=can1F14&n=1&d=10&alea=rkIp2232323232342345678910111213141516171819202122232425262728293031323334353637383940414243444546474823456789101112131415161718192021&cd=1&cols=1
\needspace{10\baselineskip}
\begin{exercice}[BaremeTotal=false]


 Soit $f$ la fonction définie sur $\mathbb{R}$ par : $f(x)= -2x^2+2x-10$.\\
        En calculant la limite du taux d'accroissement, déterminer $f'(2)$.
\end{exercice}

\begin{Solution}
 $f$ est une fonction polynôme du second degré de la forme $f(x)=ax^2+bx+c$.\\
    La fonction dérivée est donnée par la somme des dérivées des fonctions $u$ et $v$ définies par $u(x)=-2x^2$ et $v(x)=2x-10$.\\
     Comme $u'(x)=-4x$ et $v'(x)=2$, on obtient  $f'(x)=-4x+2$. \\
     Ainsi, $f'(2)=-4\times 2+2=-6$.
\end{Solution}

% @see : https://coopmaths.fr/alea?uuid=3c690&id=can1F13&n=1&d=10&alea=EUDu2232323232342345678910111213141516171819202122232425262728293031323334353637383940414243444546474823456789101112131415161718192021&cd=1&cols=1
\needspace{10\baselineskip}
\begin{exercice}[BaremeTotal=false]


 Déterminer le coefficient directeur de la tangente à la courbe représentative de la fonction racine carrée au point d'abscisse $18$.
         
\end{exercice}

\begin{Solution}
 Le coefficient directeur de la tangente au point d'abscisse $a$ est donné par le nombre dérivé $f'(a)$.\\
        La fonction $f$ définie par $f(x)=\sqrt{x}$ a pour fonction dérivée la fonction $f'$ définie par $f'(x)=\dfrac{1}{2\sqrt{x}}$.\\Comme $f'(18)=\dfrac{1}{2\sqrt{18}}$, le coefficient directeur de la tangente au point d'abscisse $18$ est : $\dfrac{1}{2\sqrt{18}}$.
\end{Solution}

% @see : https://coopmaths.fr/alea?uuid=ebd89&id=1AN14-1&n=1&d=10&s=3&alea=ocYr2232323232342345678910111213141516171819202122232425262728293031323334353637383940414243444546474823456789101112131415161718192021&cd=0&cols=1
\needspace{10\baselineskip}
\begin{exercice}[BaremeTotal=false]


 Donner la dérivée de la fonction $f$, dérivable pour tout $x\in\R$, définie par  $f(x)=\dfrac{-4x+10}{5}$.
\end{exercice}

\begin{Solution}
 L'expression de la dérivée de la fonction $f$ définie par $f(x)=\dfrac{-4x+10}{5}$ est : ${\color[HTML]{f15929}\boldsymbol{f'(x)=-\dfrac{4}{5}}}$.
\end{Solution}

% @see : https://coopmaths.fr/alea?uuid=ebd89&id=1AN14-1&n=1&d=10&s=4&alea=uAlP2232323232342345678910111213141516171819202122232425262728293031323334353637383940414243444546474823456789101112131415161718192021&cd=0&cols=1
\needspace{10\baselineskip}
\begin{exercice}[BaremeTotal=false]


 Donner la dérivée de la fonction $f$, dérivable pour tout $x\in\R$, définie par  $f(x)=9x^{8}$.
\end{exercice}

\begin{Solution}
 L'expression de la dérivée de la fonction $f$ définie par $f(x)=9x^{8}$ est : ${\color[HTML]{f15929}\boldsymbol{f'(x)=72x^{7}}}$.
\end{Solution}

% @see : https://coopmaths.fr/alea?uuid=ebd89&id=1AN14-1&n=1&d=10&s=7&alea=vjN72232323232342345678910111213141516171819202122232425262728293031323334353637383940414243444546474823456789101112131415161718192021&cd=0&cols=1
\needspace{10\baselineskip}
\begin{exercice}[BaremeTotal=false]


 Donner la dérivée de la fonction $f$, dérivable pour tout $x\in\R^*$, définie par  $f(x)=\dfrac{-4}{9x}$.
\end{exercice}

\begin{Solution}
 L'expression de la dérivée de la fonction $f$ définie par $f(x)=\dfrac{-4}{9x}$ est : ${\color[HTML]{f15929}\boldsymbol{f'(x)=\dfrac{4}{9x^2}}}$.
\end{Solution}

% @see : https://coopmaths.fr/alea?uuid=a83c0&id=1AN14-4&n=1&d=10&s=1&alea=Njr32232323232342345678910111213141516171819202122232425262728293031323334353637383940414243444546474823456789101112131415161718192021&cd=0&cols=1
\needspace{10\baselineskip}
\begin{exercice}[BaremeTotal=false]


 Donner l'expression de la dérivée de la fonction $f$ définie sur $\R^{*}$ par $f(x)=5x^{2}-2x+5-\dfrac{7}{x}$.\\
\end{exercice}

\begin{Solution}
 L'expression de la dérivée de la fonction $f$ définie par $f(x)=5x^{2}-2x+5-\dfrac{7}{x}$ est :\\$f'(x)={\color[HTML]{f15929}\boldsymbol{10x-2+\dfrac{7}{x^2}}}$.\\
\end{Solution}

% @see : https://coopmaths.fr/alea?uuid=a83c0&id=1AN14-4&n=1&d=10&s=4&alea=4Vpz2232323232342345678910111213141516171819202122232425262728293031323334353637383940414243444546474823456789101112131415161718192021&cd=1&cols=1
\needspace{10\baselineskip}
\begin{exercice}[BaremeTotal=false]


 Donner l'expression de la dérivée de la fonction $f$ définie sur $\R_+^{*}$ par $f(x)=\dfrac{2}{x}-2\sqrt{x}-4x$.\\
\end{exercice}

\begin{Solution}
 $(\dfrac{2}{x})^\prime=-\dfrac{2}{x^2}$\\$(-2\sqrt{x})^\prime=-\dfrac{2}{2\sqrt{x}}$\\$(-4x)^\prime=-4$\\L'expression de la dérivée de la fonction $f$ définie par $f(x)=\dfrac{2}{x}-2\sqrt{x}-4x$ est :\\$f'(x)={\color[HTML]{f15929}\boldsymbol{-\dfrac{2}{x^2}-\dfrac{2}{2\sqrt{x}}-4}}$.\\
\end{Solution}

\end{Maquette}
\clearpage
\begin{Maquette}[DM]{Niveau= ,Classe=\getdata[22]\mydata\ (v22), Date = 06/01/25  ,Theme=Exercices}


% @see : https://coopmaths.fr/alea?uuid=4c8c7&id=1AN11-3&n=1&d=10&s=2&alea=nawO223232323234234567891011121314151617181920212223242526272829303132333435363738394041424344454647482345678910111213141516171819202122&cd=1&cols=2
\needspace{10\baselineskip}
\begin{exercice}[BaremeTotal=false]
\label{eleve:22}


 \begin{minipage}[t]{\linewidth} Soit $f$ une fonction dérivable sur $[-5;5]$ et $\mathcal{C}_f$ sa courbe représentative.\\On sait que  $f(-5)=1~~$ et que $~~f'(-5)=2$.\\Déterminer une équation de la tangente $(T)$ à la courbe $\mathcal{C}_f$ au point d'abscisse $-5$,\\sans utiliser la formule de cours de l'équation de tangente. \end{minipage}

\end{exercice}

\begin{Solution}
 \begin{minipage}[t]{\linewidth} On sait que la tangente n'est pas une droite verticale, puisque la fonction est dérivable sur l'intervalle.\\On en déduit que la tangente $(T)$ au point d'abscisse $-5$, admet une équation réduite de la forme :  $(T) : y=m x + p$. 

\medskip
$\bullet$ Détermination de $m$ :\\On sait que le nombre dérivé en $-5$ est par définition, le coefficient directeur de la tangente au point d'abscisse $-5$.\\Par conséquent, on a déjà : $m=f'(-5)=2$.\\ On en déduit que  $(T) : y= 2 x + p$\\$\bullet$ Détermination de $p$ :\\ Pour cela, on utilise que si $f(-5)=1~~$, alors le point $A$ de coordonnées $(-5;1)$ appartient à $\mathcal{C}_f$ mais aussi à $(T)$.\\ On peut écrire $A(-5;1) = \mathcal{C}_f \cap (T)$.\\ On remplace alors les coordonnées de $A(-5;1)$ dans l'équation  $(T) : y= 2 x + p$.\\ $\begin{aligned} \phantom{\iff}&A(-5;1)\in (T)\\ \iff& 1= 2 \times (-5)  + p\\ \iff& p=1 -2 \times (-5)  \\ \iff& p=1 +10   \\ \iff& p=11\\\end{aligned}$\\On peut conclure que : $(T) : {\color[HTML]{f15929}\boldsymbol{y=2x+11}}$. \end{minipage}

\end{Solution}

% @see : https://coopmaths.fr/alea?uuid=0e984&id=can1F15&n=1&d=10&alea=1Alo223232323234234567891011121314151617181920212223242526272829303132333435363738394041424344454647482345678910111213141516171819202122&cd=1&cols=1
\needspace{10\baselineskip}
\begin{exercice}[BaremeTotal=false]


 La courbe représente une fonction $f$ et la droite est la tangente au point d'abscisse $2$.\\
        Déterminer $f'(2)$.\begin{center}\begin{tikzpicture}[baseline,scale = 0.5]

    \tikzset{
      point/.style={
        thick,
        draw,
        cross out,
        inner sep=0pt,
        minimum width=5pt,
        minimum height=5pt,
      },
    }
    \clip (-2,-2) rectangle (14,12);
    	
	\draw[color ={black},line width = 2,>=latex,->] (-2,0)--(14,0);
	\draw[color ={black},line width = 2,>=latex,->] (0,-2)--(0,12);
	\draw[color ={black},opacity = 0.5] (-2,1)--(14,1);
	\draw[color ={black},opacity = 0.5] (-2,2)--(14,2);
	\draw[color ={black},opacity = 0.5] (-2,3)--(14,3);
	\draw[color ={black},opacity = 0.5] (-2,4)--(14,4);
	\draw[color ={black},opacity = 0.5] (-2,5)--(14,5);
	\draw[color ={black},opacity = 0.5] (-2,6)--(14,6);
	\draw[color ={black},opacity = 0.5] (-2,7)--(14,7);
	\draw[color ={black},opacity = 0.5] (-2,8)--(14,8);
	\draw[color ={black},opacity = 0.5] (-2,9)--(14,9);
	\draw[color ={black},opacity = 0.5] (-2,10)--(14,10);
	\draw[color ={black},opacity = 0.5] (-2,11)--(14,11);
	\draw[color ={black},opacity = 0.5] (-2,12)--(14,12);
	\draw[color ={black},opacity = 0.5] (-2,-1)--(14,-1);
	\draw[color ={black},opacity = 0.5] (-2,-2)--(14,-2);
	\draw[color ={black},opacity = 0.5] (1,-2)--(1,12);
	\draw[color ={black},opacity = 0.5] (2,-2)--(2,12);
	\draw[color ={black},opacity = 0.5] (3,-2)--(3,12);
	\draw[color ={black},opacity = 0.5] (4,-2)--(4,12);
	\draw[color ={black},opacity = 0.5] (5,-2)--(5,12);
	\draw[color ={black},opacity = 0.5] (6,-2)--(6,12);
	\draw[color ={black},opacity = 0.5] (7,-2)--(7,12);
	\draw[color ={black},opacity = 0.5] (8,-2)--(8,12);
	\draw[color ={black},opacity = 0.5] (9,-2)--(9,12);
	\draw[color ={black},opacity = 0.5] (10,-2)--(10,12);
	\draw[color ={black},opacity = 0.5] (11,-2)--(11,12);
	\draw[color ={black},opacity = 0.5] (12,-2)--(12,12);
	\draw[color ={black},opacity = 0.5] (13,-2)--(13,12);
	\draw[color ={black},opacity = 0.5] (14,-2)--(14,12);
	\draw[color ={black},opacity = 0.5] (-1,-2)--(-1,12);
	\draw[color ={black},opacity = 0.5] (-2,-2)--(-2,12);
	\draw[color ={gray},opacity = 0.3] (-2,0)--(14,0);
	\draw[color ={gray},opacity = 0.3] (-2,0.5)--(14,0.5);
	\draw[color ={gray},opacity = 0.3] (-2,1.5)--(14,1.5);
	\draw[color ={gray},opacity = 0.3] (-2,2.5)--(14,2.5);
	\draw[color ={gray},opacity = 0.3] (-2,3.5)--(14,3.5);
	\draw[color ={gray},opacity = 0.3] (-2,4.5)--(14,4.5);
	\draw[color ={gray},opacity = 0.3] (-2,5.5)--(14,5.5);
	\draw[color ={gray},opacity = 0.3] (-2,6.5)--(14,6.5);
	\draw[color ={gray},opacity = 0.3] (-2,7.5)--(14,7.5);
	\draw[color ={gray},opacity = 0.3] (-2,8.5)--(14,8.5);
	\draw[color ={gray},opacity = 0.3] (-2,9.5)--(14,9.5);
	\draw[color ={gray},opacity = 0.3] (-2,10.5)--(14,10.5);
	\draw[color ={gray},opacity = 0.3] (-2,11.5)--(14,11.5);
	\draw[color ={gray},opacity = 0.3] (-2,0)--(14,0);
	\draw[color ={gray},opacity = 0.3] (-2,-0.5)--(14,-0.5);
	\draw[color ={gray},opacity = 0.3] (-2,-1.5)--(14,-1.5);
	\draw[color ={gray},opacity = 0.3] (0,-2)--(0,12);
	\draw[color ={gray},opacity = 0.3] (0.1,-2)--(0.1,12);
	\draw[color ={gray},opacity = 0.3] (0.2,-2)--(0.2,12);
	\draw[color ={gray},opacity = 0.3] (0.30000000000000004,-2)--(0.30000000000000004,12);
	\draw[color ={gray},opacity = 0.3] (0.4,-2)--(0.4,12);
	\draw[color ={gray},opacity = 0.3] (0.5,-2)--(0.5,12);
	\draw[color ={gray},opacity = 0.3] (0.6,-2)--(0.6,12);
	\draw[color ={gray},opacity = 0.3] (0.7,-2)--(0.7,12);
	\draw[color ={gray},opacity = 0.3] (0.7999999999999999,-2)--(0.7999999999999999,12);
	\draw[color ={gray},opacity = 0.3] (0.8999999999999999,-2)--(0.8999999999999999,12);
	\draw[color ={gray},opacity = 0.3] (1.0999999999999999,-2)--(1.0999999999999999,12);
	\draw[color ={gray},opacity = 0.3] (1.2,-2)--(1.2,12);
	\draw[color ={gray},opacity = 0.3] (1.3,-2)--(1.3,12);
	\draw[color ={gray},opacity = 0.3] (1.4000000000000001,-2)--(1.4000000000000001,12);
	\draw[color ={gray},opacity = 0.3] (1.5000000000000002,-2)--(1.5000000000000002,12);
	\draw[color ={gray},opacity = 0.3] (1.6000000000000003,-2)--(1.6000000000000003,12);
	\draw[color ={gray},opacity = 0.3] (1.7000000000000004,-2)--(1.7000000000000004,12);
	\draw[color ={gray},opacity = 0.3] (1.8000000000000005,-2)--(1.8000000000000005,12);
	\draw[color ={gray},opacity = 0.3] (1.9000000000000006,-2)--(1.9000000000000006,12);
	\draw[color ={gray},opacity = 0.3] (2.1000000000000005,-2)--(2.1000000000000005,12);
	\draw[color ={gray},opacity = 0.3] (2.2000000000000006,-2)--(2.2000000000000006,12);
	\draw[color ={gray},opacity = 0.3] (2.3000000000000007,-2)--(2.3000000000000007,12);
	\draw[color ={gray},opacity = 0.3] (2.400000000000001,-2)--(2.400000000000001,12);
	\draw[color ={gray},opacity = 0.3] (2.500000000000001,-2)--(2.500000000000001,12);
	\draw[color ={gray},opacity = 0.3] (2.600000000000001,-2)--(2.600000000000001,12);
	\draw[color ={gray},opacity = 0.3] (2.700000000000001,-2)--(2.700000000000001,12);
	\draw[color ={gray},opacity = 0.3] (2.800000000000001,-2)--(2.800000000000001,12);
	\draw[color ={gray},opacity = 0.3] (2.9000000000000012,-2)--(2.9000000000000012,12);
	\draw[color ={gray},opacity = 0.3] (3.1000000000000014,-2)--(3.1000000000000014,12);
	\draw[color ={gray},opacity = 0.3] (3.2000000000000015,-2)--(3.2000000000000015,12);
	\draw[color ={gray},opacity = 0.3] (3.3000000000000016,-2)--(3.3000000000000016,12);
	\draw[color ={gray},opacity = 0.3] (3.4000000000000017,-2)--(3.4000000000000017,12);
	\draw[color ={gray},opacity = 0.3] (3.5000000000000018,-2)--(3.5000000000000018,12);
	\draw[color ={gray},opacity = 0.3] (3.600000000000002,-2)--(3.600000000000002,12);
	\draw[color ={gray},opacity = 0.3] (3.700000000000002,-2)--(3.700000000000002,12);
	\draw[color ={gray},opacity = 0.3] (3.800000000000002,-2)--(3.800000000000002,12);
	\draw[color ={gray},opacity = 0.3] (3.900000000000002,-2)--(3.900000000000002,12);
	\draw[color ={gray},opacity = 0.3] (4.100000000000001,-2)--(4.100000000000001,12);
	\draw[color ={gray},opacity = 0.3] (4.200000000000001,-2)--(4.200000000000001,12);
	\draw[color ={gray},opacity = 0.3] (4.300000000000001,-2)--(4.300000000000001,12);
	\draw[color ={gray},opacity = 0.3] (4.4,-2)--(4.4,12);
	\draw[color ={gray},opacity = 0.3] (4.5,-2)--(4.5,12);
	\draw[color ={gray},opacity = 0.3] (4.6,-2)--(4.6,12);
	\draw[color ={gray},opacity = 0.3] (4.699999999999999,-2)--(4.699999999999999,12);
	\draw[color ={gray},opacity = 0.3] (4.799999999999999,-2)--(4.799999999999999,12);
	\draw[color ={gray},opacity = 0.3] (4.899999999999999,-2)--(4.899999999999999,12);
	\draw[color ={gray},opacity = 0.3] (5.099999999999998,-2)--(5.099999999999998,12);
	\draw[color ={gray},opacity = 0.3] (5.1999999999999975,-2)--(5.1999999999999975,12);
	\draw[color ={gray},opacity = 0.3] (5.299999999999997,-2)--(5.299999999999997,12);
	\draw[color ={gray},opacity = 0.3] (5.399999999999997,-2)--(5.399999999999997,12);
	\draw[color ={gray},opacity = 0.3] (5.4999999999999964,-2)--(5.4999999999999964,12);
	\draw[color ={gray},opacity = 0.3] (5.599999999999996,-2)--(5.599999999999996,12);
	\draw[color ={gray},opacity = 0.3] (5.699999999999996,-2)--(5.699999999999996,12);
	\draw[color ={gray},opacity = 0.3] (5.799999999999995,-2)--(5.799999999999995,12);
	\draw[color ={gray},opacity = 0.3] (5.899999999999995,-2)--(5.899999999999995,12);
	\draw[color ={gray},opacity = 0.3] (6.099999999999994,-2)--(6.099999999999994,12);
	\draw[color ={gray},opacity = 0.3] (6.199999999999994,-2)--(6.199999999999994,12);
	\draw[color ={gray},opacity = 0.3] (6.299999999999994,-2)--(6.299999999999994,12);
	\draw[color ={gray},opacity = 0.3] (6.399999999999993,-2)--(6.399999999999993,12);
	\draw[color ={gray},opacity = 0.3] (6.499999999999993,-2)--(6.499999999999993,12);
	\draw[color ={gray},opacity = 0.3] (6.5999999999999925,-2)--(6.5999999999999925,12);
	\draw[color ={gray},opacity = 0.3] (6.699999999999992,-2)--(6.699999999999992,12);
	\draw[color ={gray},opacity = 0.3] (6.799999999999992,-2)--(6.799999999999992,12);
	\draw[color ={gray},opacity = 0.3] (6.8999999999999915,-2)--(6.8999999999999915,12);
	\draw[color ={gray},opacity = 0.3] (7.099999999999991,-2)--(7.099999999999991,12);
	\draw[color ={gray},opacity = 0.3] (7.19999999999999,-2)--(7.19999999999999,12);
	\draw[color ={gray},opacity = 0.3] (7.29999999999999,-2)--(7.29999999999999,12);
	\draw[color ={gray},opacity = 0.3] (7.39999999999999,-2)--(7.39999999999999,12);
	\draw[color ={gray},opacity = 0.3] (7.499999999999989,-2)--(7.499999999999989,12);
	\draw[color ={gray},opacity = 0.3] (7.599999999999989,-2)--(7.599999999999989,12);
	\draw[color ={gray},opacity = 0.3] (7.699999999999989,-2)--(7.699999999999989,12);
	\draw[color ={gray},opacity = 0.3] (7.799999999999988,-2)--(7.799999999999988,12);
	\draw[color ={gray},opacity = 0.3] (7.899999999999988,-2)--(7.899999999999988,12);
	\draw[color ={gray},opacity = 0.3] (8.099999999999987,-2)--(8.099999999999987,12);
	\draw[color ={gray},opacity = 0.3] (8.199999999999987,-2)--(8.199999999999987,12);
	\draw[color ={gray},opacity = 0.3] (8.299999999999986,-2)--(8.299999999999986,12);
	\draw[color ={gray},opacity = 0.3] (8.399999999999986,-2)--(8.399999999999986,12);
	\draw[color ={gray},opacity = 0.3] (8.499999999999986,-2)--(8.499999999999986,12);
	\draw[color ={gray},opacity = 0.3] (8.599999999999985,-2)--(8.599999999999985,12);
	\draw[color ={gray},opacity = 0.3] (8.699999999999985,-2)--(8.699999999999985,12);
	\draw[color ={gray},opacity = 0.3] (8.799999999999985,-2)--(8.799999999999985,12);
	\draw[color ={gray},opacity = 0.3] (8.899999999999984,-2)--(8.899999999999984,12);
	\draw[color ={gray},opacity = 0.3] (9.099999999999984,-2)--(9.099999999999984,12);
	\draw[color ={gray},opacity = 0.3] (9.199999999999983,-2)--(9.199999999999983,12);
	\draw[color ={gray},opacity = 0.3] (9.299999999999983,-2)--(9.299999999999983,12);
	\draw[color ={gray},opacity = 0.3] (9.399999999999983,-2)--(9.399999999999983,12);
	\draw[color ={gray},opacity = 0.3] (9.499999999999982,-2)--(9.499999999999982,12);
	\draw[color ={gray},opacity = 0.3] (9.599999999999982,-2)--(9.599999999999982,12);
	\draw[color ={gray},opacity = 0.3] (9.699999999999982,-2)--(9.699999999999982,12);
	\draw[color ={gray},opacity = 0.3] (9.799999999999981,-2)--(9.799999999999981,12);
	\draw[color ={gray},opacity = 0.3] (9.89999999999998,-2)--(9.89999999999998,12);
	\draw[color ={gray},opacity = 0.3] (10.09999999999998,-2)--(10.09999999999998,12);
	\draw[color ={gray},opacity = 0.3] (10.19999999999998,-2)--(10.19999999999998,12);
	\draw[color ={gray},opacity = 0.3] (10.29999999999998,-2)--(10.29999999999998,12);
	\draw[color ={gray},opacity = 0.3] (10.399999999999979,-2)--(10.399999999999979,12);
	\draw[color ={gray},opacity = 0.3] (10.499999999999979,-2)--(10.499999999999979,12);
	\draw[color ={gray},opacity = 0.3] (10.599999999999978,-2)--(10.599999999999978,12);
	\draw[color ={gray},opacity = 0.3] (10.699999999999978,-2)--(10.699999999999978,12);
	\draw[color ={gray},opacity = 0.3] (10.799999999999978,-2)--(10.799999999999978,12);
	\draw[color ={gray},opacity = 0.3] (10.899999999999977,-2)--(10.899999999999977,12);
	\draw[color ={gray},opacity = 0.3] (11.099999999999977,-2)--(11.099999999999977,12);
	\draw[color ={gray},opacity = 0.3] (11.199999999999976,-2)--(11.199999999999976,12);
	\draw[color ={gray},opacity = 0.3] (11.299999999999976,-2)--(11.299999999999976,12);
	\draw[color ={gray},opacity = 0.3] (11.399999999999975,-2)--(11.399999999999975,12);
	\draw[color ={gray},opacity = 0.3] (11.499999999999975,-2)--(11.499999999999975,12);
	\draw[color ={gray},opacity = 0.3] (11.599999999999975,-2)--(11.599999999999975,12);
	\draw[color ={gray},opacity = 0.3] (11.699999999999974,-2)--(11.699999999999974,12);
	\draw[color ={gray},opacity = 0.3] (11.799999999999974,-2)--(11.799999999999974,12);
	\draw[color ={gray},opacity = 0.3] (11.899999999999974,-2)--(11.899999999999974,12);
	\draw[color ={gray},opacity = 0.3] (12.099999999999973,-2)--(12.099999999999973,12);
	\draw[color ={gray},opacity = 0.3] (12.199999999999973,-2)--(12.199999999999973,12);
	\draw[color ={gray},opacity = 0.3] (12.299999999999972,-2)--(12.299999999999972,12);
	\draw[color ={gray},opacity = 0.3] (12.399999999999972,-2)--(12.399999999999972,12);
	\draw[color ={gray},opacity = 0.3] (12.499999999999972,-2)--(12.499999999999972,12);
	\draw[color ={gray},opacity = 0.3] (12.599999999999971,-2)--(12.599999999999971,12);
	\draw[color ={gray},opacity = 0.3] (12.69999999999997,-2)--(12.69999999999997,12);
	\draw[color ={gray},opacity = 0.3] (12.79999999999997,-2)--(12.79999999999997,12);
	\draw[color ={gray},opacity = 0.3] (12.89999999999997,-2)--(12.89999999999997,12);
	\draw[color ={gray},opacity = 0.3] (13.09999999999997,-2)--(13.09999999999997,12);
	\draw[color ={gray},opacity = 0.3] (13.199999999999969,-2)--(13.199999999999969,12);
	\draw[color ={gray},opacity = 0.3] (13.299999999999969,-2)--(13.299999999999969,12);
	\draw[color ={gray},opacity = 0.3] (13.399999999999968,-2)--(13.399999999999968,12);
	\draw[color ={gray},opacity = 0.3] (13.499999999999968,-2)--(13.499999999999968,12);
	\draw[color ={gray},opacity = 0.3] (13.599999999999968,-2)--(13.599999999999968,12);
	\draw[color ={gray},opacity = 0.3] (13.699999999999967,-2)--(13.699999999999967,12);
	\draw[color ={gray},opacity = 0.3] (13.799999999999967,-2)--(13.799999999999967,12);
	\draw[color ={gray},opacity = 0.3] (13.899999999999967,-2)--(13.899999999999967,12);
	\draw[color ={gray},opacity = 0.3] (0,-2)--(0,12);
	\draw[color ={gray},opacity = 0.3] (-0.1,-2)--(-0.1,12);
	\draw[color ={gray},opacity = 0.3] (-0.2,-2)--(-0.2,12);
	\draw[color ={gray},opacity = 0.3] (-0.30000000000000004,-2)--(-0.30000000000000004,12);
	\draw[color ={gray},opacity = 0.3] (-0.4,-2)--(-0.4,12);
	\draw[color ={gray},opacity = 0.3] (-0.5,-2)--(-0.5,12);
	\draw[color ={gray},opacity = 0.3] (-0.6,-2)--(-0.6,12);
	\draw[color ={gray},opacity = 0.3] (-0.7,-2)--(-0.7,12);
	\draw[color ={gray},opacity = 0.3] (-0.7999999999999999,-2)--(-0.7999999999999999,12);
	\draw[color ={gray},opacity = 0.3] (-0.8999999999999999,-2)--(-0.8999999999999999,12);
	\draw[color ={gray},opacity = 0.3] (-1.0999999999999999,-2)--(-1.0999999999999999,12);
	\draw[color ={gray},opacity = 0.3] (-1.2,-2)--(-1.2,12);
	\draw[color ={gray},opacity = 0.3] (-1.3,-2)--(-1.3,12);
	\draw[color ={gray},opacity = 0.3] (-1.4000000000000001,-2)--(-1.4000000000000001,12);
	\draw[color ={gray},opacity = 0.3] (-1.5000000000000002,-2)--(-1.5000000000000002,12);
	\draw[color ={gray},opacity = 0.3] (-1.6000000000000003,-2)--(-1.6000000000000003,12);
	\draw[color ={gray},opacity = 0.3] (-1.7000000000000004,-2)--(-1.7000000000000004,12);
	\draw[color ={gray},opacity = 0.3] (-1.8000000000000005,-2)--(-1.8000000000000005,12);
	\draw[color ={gray},opacity = 0.3] (-1.9000000000000006,-2)--(-1.9000000000000006,12);
	\draw[color ={black},line width = 2] (2,-0.2)--(2,0.2);
	\draw[color ={black},line width = 2] (4,-0.2)--(4,0.2);
	\draw[color ={black},line width = 2] (6,-0.2)--(6,0.2);
	\draw[color ={black},line width = 2] (8,-0.2)--(8,0.2);
	\draw[color ={black},line width = 2] (10,-0.2)--(10,0.2);
	\draw[color ={black},line width = 2] (12,-0.2)--(12,0.2);
	\draw[color ={black},line width = 2] (-0.2,2)--(0.2,2);
	\draw[color ={black},line width = 2] (-0.2,4)--(0.2,4);
	\draw[color ={black},line width = 2] (-0.2,6)--(0.2,6);
	\draw[color ={black},line width = 2] (-0.2,8)--(0.2,8);
	\draw[color ={black},line width = 2] (-0.2,10)--(0.2,10);
	\draw (2,-0.4) node[anchor = center] {\scriptsize \color{black}{$1$}};
	\draw (4,-0.4) node[anchor = center] {\scriptsize \color{black}{$2$}};
	\draw (6,-0.4) node[anchor = center] {\scriptsize \color{black}{$3$}};
	\draw (8,-0.4) node[anchor = center] {\scriptsize \color{black}{$4$}};
	\draw (10,-0.4) node[anchor = center] {\scriptsize \color{black}{$5$}};
	\draw (-0.5,2.1) node[anchor = center] {\footnotesize \color{black}{$1$}};
	\draw (-0.5,4.1) node[anchor = center] {\footnotesize \color{black}{$2$}};
	\draw (-0.5,6.1) node[anchor = center] {\footnotesize \color{black}{$3$}};
	\draw (-0.5,8.1) node[anchor = center] {\footnotesize \color{black}{$4$}};
	\draw (-0.5,10.1) node[anchor = center] {\footnotesize \color{black}{$5$}};
	\draw [color={black}] (-0.3,-0.3) node[anchor = center,scale=1, rotate = 0] {O};
	
	\draw[color={blue},line width = 2] (0.6,10.67)--(1,8)--(1.4,6.86)--(1.8,6.22)--(2.2,5.82)--(2.6,5.54)--(3,5.33)--(3.4,5.18)--(3.8,5.05)--(4.2,4.95)--(4.6,4.87)--(5,4.8)--(5.4,4.74)--(5.8,4.69)--(6.2,4.65)--(6.6,4.61)--(7,4.57)--(7.4,4.54)--(7.8,4.51)--(8.2,4.49)--(8.6,4.47)--(9,4.44)--(9.4,4.43)--(9.8,4.41)--(10.2,4.39)--(10.6,4.38)--(11,4.36)--(11.4,4.35)--(11.8,4.34)--(12.2,4.33)--(12.6,4.32)--(13,4.31)--(13.4,4.3)--(13.8,4.29);
	
	\draw[color={red},line width = 2] (-2,6.5)--(-1.6,6.4)--(-1.2,6.3)--(-0.8,6.2)--(-0.4,6.1)--(0,6)--(0.4,5.9)--(0.8,5.8)--(1.2,5.7)--(1.6,5.6)--(2,5.5)--(2.4,5.4)--(2.8,5.3)--(3.2,5.2)--(3.6,5.1)--(4,5)--(4.4,4.9)--(4.8,4.8)--(5.2,4.7)--(5.6,4.6)--(6,4.5)--(6.4,4.4)--(6.8,4.3)--(7.2,4.2)--(7.6,4.1)--(8,4)--(8.4,3.9)--(8.8,3.8)--(9.2,3.7)--(9.6,3.6)--(10,3.5)--(10.4,3.4)--(10.8,3.3)--(11.2,3.2)--(11.6,3.1)--(12,3)--(12.4,2.9)--(12.8,2.8)--(13.2,2.7)--(13.6,2.6)--(14,2.5);

\end{tikzpicture}\end{center}
\end{exercice}

\begin{Solution}
 $f'(2)$ est donné par le coefficient directeur de la tangente à la courbe au point d'abscisse $2$, soit $\dfrac{-1}{4}=-\dfrac{1}{4}$.
\end{Solution}

% @see : https://coopmaths.fr/alea?uuid=6f32d&id=can1F16&n=1&d=10&alea=dqSw223232323234234567891011121314151617181920212223242526272829303132333435363738394041424344454647482345678910111213141516171819202122&cd=1&cols=1
\needspace{10\baselineskip}
\newpage
\begin{exercice}[BaremeTotal=false]


 On donne les représentations graphiques d'une fonction et de sa dérivée.\\
        Donner l'équation réduite de la tangente à la courbe de $f$ en $x=2$. \\ \begin{minipage}{0.5\linewidth}
    \begin{tikzpicture}[baseline, scale=0.8]

    \tikzset{
      point/.style={
        thick,
        draw,
        cross out,
        inner sep=0pt,
        minimum width=5pt,
        minimum height=5pt,
      },
    }
    \clip (-4.5,-4.5) rectangle (5.75,13.5);
    	
	\draw[color ={black},>=latex,->] (-3,0)--(3,0);
	\draw[color ={black},>=latex,->] (0,-3)--(0,12);
	\draw[color ={black},opacity = 0.4] (-3,1)--(3,1);
	\draw[color ={black},opacity = 0.4] (-3,2)--(3,2);
	\draw[color ={black},opacity = 0.4] (-3,3)--(3,3);
	\draw[color ={black},opacity = 0.4] (-3,4)--(3,4);
	\draw[color ={black},opacity = 0.4] (-3,5)--(3,5);
	\draw[color ={black},opacity = 0.4] (-3,6)--(3,6);
	\draw[color ={black},opacity = 0.4] (-3,7)--(3,7);
	\draw[color ={black},opacity = 0.4] (-3,8)--(3,8);
	\draw[color ={black},opacity = 0.4] (-3,9)--(3,9);
	\draw[color ={black},opacity = 0.4] (-3,10)--(3,10);
	\draw[color ={black},opacity = 0.4] (-3,11)--(3,11);
	\draw[color ={black},opacity = 0.4] (-3,12)--(3,12);
	\draw[color ={black},opacity = 0.4] (-3,-1)--(3,-1);
	\draw[color ={black},opacity = 0.4] (-3,-2)--(3,-2);
	\draw[color ={black},opacity = 0.4] (-3,-3)--(3,-3);
	\draw[color ={black},opacity = 0.4] (1,-3)--(1,12);
	\draw[color ={black},opacity = 0.4] (2,-3)--(2,12);
	\draw[color ={black},opacity = 0.4] (3,-3)--(3,12);
	\draw[color ={black},opacity = 0.4] (-1,-3)--(-1,12);
	\draw[color ={black},opacity = 0.4] (-2,-3)--(-2,12);
	\draw[color ={black},opacity = 0.4] (-3,-3)--(-3,12);
	\draw[color ={gray},opacity = 0.3] (-3,0)--(3,0);
	\draw[color ={gray},opacity = 0.3] (-3,0.5)--(3,0.5);
	\draw[color ={gray},opacity = 0.3] (-3,1.5)--(3,1.5);
	\draw[color ={gray},opacity = 0.3] (-3,2.5)--(3,2.5);
	\draw[color ={gray},opacity = 0.3] (-3,3.5)--(3,3.5);
	\draw[color ={gray},opacity = 0.3] (-3,4.5)--(3,4.5);
	\draw[color ={gray},opacity = 0.3] (-3,5.5)--(3,5.5);
	\draw[color ={gray},opacity = 0.3] (-3,6.5)--(3,6.5);
	\draw[color ={gray},opacity = 0.3] (-3,7.5)--(3,7.5);
	\draw[color ={gray},opacity = 0.3] (-3,8.5)--(3,8.5);
	\draw[color ={gray},opacity = 0.3] (-3,9.5)--(3,9.5);
	\draw[color ={gray},opacity = 0.3] (-3,10.5)--(3,10.5);
	\draw[color ={gray},opacity = 0.3] (-3,11.5)--(3,11.5);
	\draw[color ={gray},opacity = 0.3] (-3,0)--(3,0);
	\draw[color ={gray},opacity = 0.3] (-3,-0.5)--(3,-0.5);
	\draw[color ={gray},opacity = 0.3] (-3,-1.5)--(3,-1.5);
	\draw[color ={gray},opacity = 0.3] (-3,-2.5)--(3,-2.5);
	\draw[color ={gray},opacity = 0.3] (0,-3)--(0,12);
	\draw[color ={gray},opacity = 0.3] (0,-3)--(0,12);
	\draw[color ={black},line width = 1.2] (1,-0.13)--(1,0.13);
	\draw[color ={black},line width = 1.2] (2,-0.13)--(2,0.13);
	\draw[color ={black},line width = 1.2] (3,-0.13)--(3,0.13);
	\draw[color ={black},line width = 1.2] (-1,-0.13)--(-1,0.13);
	\draw[color ={black},line width = 1.2] (-2,-0.13)--(-2,0.13);
	\draw[color ={black},line width = 1.2] (-3,-0.13)--(-3,0.13);
	\draw[color ={black},line width = 1.2] (-0.13,1)--(0.13,1);
	\draw[color ={black},line width = 1.2] (-0.13,2)--(0.13,2);
	\draw[color ={black},line width = 1.2] (-0.13,3)--(0.13,3);
	\draw[color ={black},line width = 1.2] (-0.13,4)--(0.13,4);
	\draw[color ={black},line width = 1.2] (-0.13,5)--(0.13,5);
	\draw[color ={black},line width = 1.2] (-0.13,6)--(0.13,6);
	\draw[color ={black},line width = 1.2] (-0.13,7)--(0.13,7);
	\draw[color ={black},line width = 1.2] (-0.13,8)--(0.13,8);
	\draw[color ={black},line width = 1.2] (-0.13,9)--(0.13,9);
	\draw[color ={black},line width = 1.2] (-0.13,10)--(0.13,10);
	\draw[color ={black},line width = 1.2] (-0.13,11)--(0.13,11);
	\draw[color ={black},line width = 1.2] (-0.13,12)--(0.13,12);
	\draw[color ={black},line width = 1.2] (-0.13,-1)--(0.13,-1);
	\draw[color ={black},line width = 1.2] (-0.13,-2)--(0.13,-2);
	\draw[color ={black},line width = 1.2] (-0.13,-3)--(0.13,-3);
	\draw (3,-0.4) node[anchor = center] {\scriptsize \color{black}{$3$}};
	\draw (-3,-0.4) node[anchor = center] {\scriptsize \color{black}{$-3$}};
	\draw (-0.5,3.1) node[anchor = center] {\footnotesize \color{black}{$3$}};
	\draw (-0.5,6.1) node[anchor = center] {\footnotesize \color{black}{$6$}};
	\draw (-0.5,9.1) node[anchor = center] {\footnotesize \color{black}{$9$}};
	\draw (-0.5,-2.9) node[anchor = center] {\footnotesize \color{black}{$-3$}};
	\draw [color={black}] (-0.3,-0.3) node[anchor = center,scale=1, rotate = 0] {O};
	\draw (3,10) node[anchor = center] {\footnotesize \color{blue}{$\Large \mathcal{C}_f$}};
	
	\draw[color={blue},line width = 2] (-2.6,11.96)--(-2.4,10.56)--(-2.2,9.24)--(-2,8)--(-1.8,6.84)--(-1.6,5.76)--(-1.4,4.76)--(-1.2,3.84)--(-1,3)--(-0.8,2.24)--(-0.6,1.56)--(-0.4,0.96)--(-0.2,0.44)--(0,0)--(0.2,-0.36)--(0.4,-0.64)--(0.6,-0.84)--(0.8,-0.96)--(1,-1)--(1.2,-0.96)--(1.4,-0.84)--(1.6,-0.64)--(1.8,-0.36)--(2,0)--(2.2,0.44)--(2.4,0.96)--(2.6,1.56)--(2.8,2.24)--(3,3);

\end{tikzpicture}\\
    \end{minipage}
    \begin{minipage}{0.5\linewidth}
    \begin{tikzpicture}[baseline, scale=0.8]

    \tikzset{
      point/.style={
        thick,
        draw,
        cross out,
        inner sep=0pt,
        minimum width=5pt,
        minimum height=5pt,
      },
    }
    \clip (-4.5,-6.5) rectangle (6.5,9.5);
    	
	\draw[color ={black},>=latex,->] (-3,0)--(3,0);
	\draw[color ={black},>=latex,->] (0,-5)--(0,8);
	\draw[color ={black},opacity = 0.4] (-3,1)--(3,1);
	\draw[color ={black},opacity = 0.4] (-3,2)--(3,2);
	\draw[color ={black},opacity = 0.4] (-3,3)--(3,3);
	\draw[color ={black},opacity = 0.4] (-3,4)--(3,4);
	\draw[color ={black},opacity = 0.4] (-3,5)--(3,5);
	\draw[color ={black},opacity = 0.4] (-3,6)--(3,6);
	\draw[color ={black},opacity = 0.4] (-3,7)--(3,7);
	\draw[color ={black},opacity = 0.4] (-3,8)--(3,8);
	\draw[color ={black},opacity = 0.4] (-3,-1)--(3,-1);
	\draw[color ={black},opacity = 0.4] (-3,-2)--(3,-2);
	\draw[color ={black},opacity = 0.4] (-3,-3)--(3,-3);
	\draw[color ={black},opacity = 0.4] (-3,-4)--(3,-4);
	\draw[color ={black},opacity = 0.4] (-3,-5)--(3,-5);
	\draw[color ={black},opacity = 0.4] (1,-5)--(1,8);
	\draw[color ={black},opacity = 0.4] (2,-5)--(2,8);
	\draw[color ={black},opacity = 0.4] (3,-5)--(3,8);
	\draw[color ={black},opacity = 0.4] (-1,-5)--(-1,8);
	\draw[color ={black},opacity = 0.4] (-2,-5)--(-2,8);
	\draw[color ={black},opacity = 0.4] (-3,-5)--(-3,8);
	\draw[color ={gray},opacity = 0.3] (-3,0)--(3,0);
	\draw[color ={gray},opacity = 0.3] (-3,0.5)--(3,0.5);
	\draw[color ={gray},opacity = 0.3] (-3,1.5)--(3,1.5);
	\draw[color ={gray},opacity = 0.3] (-3,2.5)--(3,2.5);
	\draw[color ={gray},opacity = 0.3] (-3,3.5)--(3,3.5);
	\draw[color ={gray},opacity = 0.3] (-3,4.5)--(3,4.5);
	\draw[color ={gray},opacity = 0.3] (-3,5.5)--(3,5.5);
	\draw[color ={gray},opacity = 0.3] (-3,6.5)--(3,6.5);
	\draw[color ={gray},opacity = 0.3] (-3,7.5)--(3,7.5);
	\draw[color ={gray},opacity = 0.3] (-3,0)--(3,0);
	\draw[color ={gray},opacity = 0.3] (-3,-0.5)--(3,-0.5);
	\draw[color ={gray},opacity = 0.3] (-3,-1.5)--(3,-1.5);
	\draw[color ={gray},opacity = 0.3] (-3,-2.5)--(3,-2.5);
	\draw[color ={gray},opacity = 0.3] (-3,-3.5)--(3,-3.5);
	\draw[color ={gray},opacity = 0.3] (-3,-4.5)--(3,-4.5);
	\draw[color ={gray},opacity = 0.3] (0,-5)--(0,8);
	\draw[color ={gray},opacity = 0.3] (0,-5)--(0,8);
	\draw[color ={black},line width = 1.2] (1,-0.13)--(1,0.13);
	\draw[color ={black},line width = 1.2] (2,-0.13)--(2,0.13);
	\draw[color ={black},line width = 1.2] (3,-0.13)--(3,0.13);
	\draw[color ={black},line width = 1.2] (-1,-0.13)--(-1,0.13);
	\draw[color ={black},line width = 1.2] (-2,-0.13)--(-2,0.13);
	\draw[color ={black},line width = 1.2] (-3,-0.13)--(-3,0.13);
	\draw[color ={black},line width = 1.2] (-0.13,1)--(0.13,1);
	\draw[color ={black},line width = 1.2] (-0.13,2)--(0.13,2);
	\draw[color ={black},line width = 1.2] (-0.13,3)--(0.13,3);
	\draw[color ={black},line width = 1.2] (-0.13,4)--(0.13,4);
	\draw[color ={black},line width = 1.2] (-0.13,5)--(0.13,5);
	\draw[color ={black},line width = 1.2] (-0.13,6)--(0.13,6);
	\draw[color ={black},line width = 1.2] (-0.13,7)--(0.13,7);
	\draw[color ={black},line width = 1.2] (-0.13,8)--(0.13,8);
	\draw[color ={black},line width = 1.2] (-0.13,9)--(0.13,9);
	\draw[color ={black},line width = 1.2] (-0.13,10)--(0.13,10);
	\draw[color ={black},line width = 1.2] (-0.13,11)--(0.13,11);
	\draw[color ={black},line width = 1.2] (-0.13,12)--(0.13,12);
	\draw[color ={black},line width = 1.2] (-0.13,-1)--(0.13,-1);
	\draw[color ={black},line width = 1.2] (-0.13,-2)--(0.13,-2);
	\draw[color ={black},line width = 1.2] (-0.13,-3)--(0.13,-3);
	\draw (3,-0.4) node[anchor = center] {\scriptsize \color{black}{$3$}};
	\draw (-3,-0.4) node[anchor = center] {\scriptsize \color{black}{$-3$}};
	\draw (-0.5,3.1) node[anchor = center] {\footnotesize \color{black}{$3$}};
	\draw (-0.5,6.1) node[anchor = center] {\footnotesize \color{black}{$6$}};
	\draw (-0.5,-2.9) node[anchor = center] {\footnotesize \color{black}{$-3$}};
	\draw [color={black}] (-0.3,-0.3) node[anchor = center,scale=1, rotate = 0] {O};
	\draw (3,6) node[anchor = center] {\footnotesize \color{red}{$\Large\mathcal{C}_f\prime$}};
	
	\draw[color={red},line width = 2] (-2,-6)--(-1.8,-5.6)--(-1.6,-5.2)--(-1.4,-4.8)--(-1.2,-4.4)--(-1,-4)--(-0.8,-3.6)--(-0.6,-3.2)--(-0.4,-2.8)--(-0.2,-2.4)--(0,-2)--(0.2,-1.6)--(0.4,-1.2)--(0.6,-0.8)--(0.8,-0.4)--(1,0)--(1.2,0.4)--(1.4,0.8)--(1.6,1.2)--(1.8,1.6)--(2,2)--(2.2,2.4)--(2.4,2.8)--(2.6,3.2)--(2.8,3.6)--(3,4);

\end{tikzpicture}\\
    \end{minipage}
    
\end{exercice}

\begin{Solution}
 L'équation réduite de la tangente au point d'abscisse $2$ est  : $y=f'(2)(x-2)+f(2)$.\\
        On lit graphiquement $f(2)=0$ et $f'(2)=2$.\\
        L'équation réduite de la tangente est donc donnée par :
        $y=2(x-2)+0$, soit $y=2x-4$.
\end{Solution}

% @see : https://coopmaths.fr/alea?uuid=a1ba2&id=can1F14&n=1&d=10&alea=rkIp223232323234234567891011121314151617181920212223242526272829303132333435363738394041424344454647482345678910111213141516171819202122&cd=1&cols=1
\needspace{10\baselineskip}
\begin{exercice}[BaremeTotal=false]


 Soit $f$ la fonction définie sur $\mathbb{R}$ par : $f(x)= 2x^2+3x-5$.\\
        En calculant la limite du taux d'accroissement, déterminer $f'(3)$.
\end{exercice}

\begin{Solution}
 $f$ est une fonction polynôme du second degré de la forme $f(x)=ax^2+bx+c$.\\
    La fonction dérivée est donnée par la somme des dérivées des fonctions $u$ et $v$ définies par $u(x)=2x^2$ et $v(x)=3x-5$.\\
     Comme $u'(x)=4x$ et $v'(x)=3$, on obtient  $f'(x)=4x+3$. \\
     Ainsi, $f'(3)=4\times 3+3=15$.
\end{Solution}

% @see : https://coopmaths.fr/alea?uuid=3c690&id=can1F13&n=1&d=10&alea=EUDu223232323234234567891011121314151617181920212223242526272829303132333435363738394041424344454647482345678910111213141516171819202122&cd=1&cols=1
\needspace{10\baselineskip}
\begin{exercice}[BaremeTotal=false]


 Déterminer le coefficient directeur de la tangente à la courbe représentative de la fonction racine carrée au point d'abscisse $14$.
         
\end{exercice}

\begin{Solution}
 Le coefficient directeur de la tangente au point d'abscisse $a$ est donné par le nombre dérivé $f'(a)$.\\
        La fonction $f$ définie par $f(x)=\sqrt{x}$ a pour fonction dérivée la fonction $f'$ définie par $f'(x)=\dfrac{1}{2\sqrt{x}}$.\\Comme $f'(14)=\dfrac{1}{2\sqrt{14}}$, le coefficient directeur de la tangente au point d'abscisse $14$ est : $\dfrac{1}{2\sqrt{14}}$.
\end{Solution}

% @see : https://coopmaths.fr/alea?uuid=ebd89&id=1AN14-1&n=1&d=10&s=3&alea=ocYr223232323234234567891011121314151617181920212223242526272829303132333435363738394041424344454647482345678910111213141516171819202122&cd=0&cols=1
\needspace{10\baselineskip}
\begin{exercice}[BaremeTotal=false]


 Donner la dérivée de la fonction $f$, dérivable pour tout $x\in\R$, définie par  $f(x)=\dfrac{4}{9}x+6$.
\end{exercice}

\begin{Solution}
 L'expression de la dérivée de la fonction $f$ définie par $f(x)=\dfrac{4}{9}x+6$ est : ${\color[HTML]{f15929}\boldsymbol{f'(x)=\dfrac{4}{9}}}$.
\end{Solution}

% @see : https://coopmaths.fr/alea?uuid=ebd89&id=1AN14-1&n=1&d=10&s=4&alea=uAlP223232323234234567891011121314151617181920212223242526272829303132333435363738394041424344454647482345678910111213141516171819202122&cd=0&cols=1
\needspace{10\baselineskip}
\begin{exercice}[BaremeTotal=false]


 Donner la dérivée de la fonction $f$, dérivable pour tout $x\in\R$, définie par  $f(x)=2x^{9}$.
\end{exercice}

\begin{Solution}
 L'expression de la dérivée de la fonction $f$ définie par $f(x)=2x^{9}$ est : ${\color[HTML]{f15929}\boldsymbol{f'(x)=18x^{8}}}$.
\end{Solution}

% @see : https://coopmaths.fr/alea?uuid=ebd89&id=1AN14-1&n=1&d=10&s=7&alea=vjN7223232323234234567891011121314151617181920212223242526272829303132333435363738394041424344454647482345678910111213141516171819202122&cd=0&cols=1
\needspace{10\baselineskip}
\begin{exercice}[BaremeTotal=false]


 Donner la dérivée de la fonction $f$, dérivable pour tout $x\in\R^*$, définie par  $f(x)=\dfrac{-5}{6x}$.
\end{exercice}

\begin{Solution}
 L'expression de la dérivée de la fonction $f$ définie par $f(x)=\dfrac{-5}{6x}$ est : ${\color[HTML]{f15929}\boldsymbol{f'(x)=\dfrac{5}{6x^2}}}$.
\end{Solution}

% @see : https://coopmaths.fr/alea?uuid=a83c0&id=1AN14-4&n=1&d=10&s=1&alea=Njr3223232323234234567891011121314151617181920212223242526272829303132333435363738394041424344454647482345678910111213141516171819202122&cd=0&cols=1
\needspace{10\baselineskip}
\begin{exercice}[BaremeTotal=false]


 Donner l'expression de la dérivée de la fonction $f$ définie sur $\R^{*}$ par $f(x)=\dfrac{3}{x}-2x^{2}$.\\
\end{exercice}

\begin{Solution}
 L'expression de la dérivée de la fonction $f$ définie par $f(x)=\dfrac{3}{x}-2x^{2}$ est :\\$f'(x)={\color[HTML]{f15929}\boldsymbol{-\dfrac{3}{x^2}-4x}}$.\\
\end{Solution}

% @see : https://coopmaths.fr/alea?uuid=a83c0&id=1AN14-4&n=1&d=10&s=4&alea=4Vpz223232323234234567891011121314151617181920212223242526272829303132333435363738394041424344454647482345678910111213141516171819202122&cd=1&cols=1
\needspace{10\baselineskip}
\begin{exercice}[BaremeTotal=false]


 Donner l'expression de la dérivée de la fonction $f$ définie sur $\R_+^{*}$ par $f(x)=-2x^{2}-4x+4\sqrt{x}-\dfrac{3}{x}$.\\
\end{exercice}

\begin{Solution}
 $(-2x^{2}-4x)^\prime=-4x-4$\\$(4\sqrt{x})^\prime=\dfrac{4}{2\sqrt{x}}$\\$(-\dfrac{3}{x})^\prime=\dfrac{3}{x^2}$\\L'expression de la dérivée de la fonction $f$ définie par $f(x)=-2x^{2}-4x+4\sqrt{x}-\dfrac{3}{x}$ est :\\$f'(x)={\color[HTML]{f15929}\boldsymbol{-4x-4+\dfrac{4}{2\sqrt{x}}+\dfrac{3}{x^2}}}$.\\
\end{Solution}

\end{Maquette}
\clearpage
\begin{Maquette}[DM]{Niveau= ,Classe=\getdata[23]\mydata\ (v23), Date = 06/01/25  ,Theme=Exercices}


% @see : https://coopmaths.fr/alea?uuid=4c8c7&id=1AN11-3&n=1&d=10&s=2&alea=nawO22323232323423456789101112131415161718192021222324252627282930313233343536373839404142434445464748234567891011121314151617181920212223&cd=1&cols=2
\needspace{10\baselineskip}
\begin{exercice}[BaremeTotal=false]
\label{eleve:23}


 \begin{minipage}[t]{\linewidth} Soit $f$ une fonction dérivable sur $[-5;5]$ et $\mathcal{C}_f$ sa courbe représentative.\\On sait que  $f(1)=-1~~$ et que $~~f'(1)=-2$.\\Déterminer une équation de la tangente $(T)$ à la courbe $\mathcal{C}_f$ au point d'abscisse $1$,\\sans utiliser la formule de cours de l'équation de tangente. \end{minipage}

\end{exercice}

\begin{Solution}
 \begin{minipage}[t]{\linewidth} On sait que la tangente n'est pas une droite verticale, puisque la fonction est dérivable sur l'intervalle.\\On en déduit que la tangente $(T)$ au point d'abscisse $1$, admet une équation réduite de la forme :  $(T) : y=m x + p$. 

\medskip
$\bullet$ Détermination de $m$ :\\On sait que le nombre dérivé en $1$ est par définition, le coefficient directeur de la tangente au point d'abscisse $1$.\\Par conséquent, on a déjà : $m=f'(1)=-2$.\\ On en déduit que  $(T) : y= -2 x + p$\\$\bullet$ Détermination de $p$ :\\ Pour cela, on utilise que si $f(1)=-1~~$, alors le point $A$ de coordonnées $(1;-1)$ appartient à $\mathcal{C}_f$ mais aussi à $(T)$.\\ On peut écrire $A(1;-1) = \mathcal{C}_f \cap (T)$.\\ On remplace alors les coordonnées de $A(1;-1)$ dans l'équation  $(T) : y= -2 x + p$.\\ $\begin{aligned} \phantom{\iff}&A(1;-1)\in (T)\\ \iff& -1= -2 \times 1  + p\\ \iff& p=-1 +2 \times 1  \\ \iff& p=-1 +2   \\ \iff& p=1\\\end{aligned}$\\On peut conclure que : $(T) : {\color[HTML]{f15929}\boldsymbol{y=-2x+1}}$. \end{minipage}

\end{Solution}

% @see : https://coopmaths.fr/alea?uuid=0e984&id=can1F15&n=1&d=10&alea=1Alo22323232323423456789101112131415161718192021222324252627282930313233343536373839404142434445464748234567891011121314151617181920212223&cd=1&cols=1
\needspace{10\baselineskip}
\begin{exercice}[BaremeTotal=false]


 La courbe représente une fonction $f$ et la droite est la tangente au point d'abscisse $1$.\\
        Déterminer $f'(1)$.\begin{center}\begin{tikzpicture}[baseline,scale = 0.5]

    \tikzset{
      point/.style={
        thick,
        draw,
        cross out,
        inner sep=0pt,
        minimum width=5pt,
        minimum height=5pt,
      },
    }
    \clip (-2,-2) rectangle (14,12);
    	
	\draw[color ={black},line width = 2,>=latex,->] (-2,0)--(14,0);
	\draw[color ={black},line width = 2,>=latex,->] (0,-2)--(0,12);
	\draw[color ={black},opacity = 0.5] (-2,1)--(14,1);
	\draw[color ={black},opacity = 0.5] (-2,2)--(14,2);
	\draw[color ={black},opacity = 0.5] (-2,3)--(14,3);
	\draw[color ={black},opacity = 0.5] (-2,4)--(14,4);
	\draw[color ={black},opacity = 0.5] (-2,5)--(14,5);
	\draw[color ={black},opacity = 0.5] (-2,6)--(14,6);
	\draw[color ={black},opacity = 0.5] (-2,7)--(14,7);
	\draw[color ={black},opacity = 0.5] (-2,8)--(14,8);
	\draw[color ={black},opacity = 0.5] (-2,9)--(14,9);
	\draw[color ={black},opacity = 0.5] (-2,10)--(14,10);
	\draw[color ={black},opacity = 0.5] (-2,11)--(14,11);
	\draw[color ={black},opacity = 0.5] (-2,12)--(14,12);
	\draw[color ={black},opacity = 0.5] (-2,-1)--(14,-1);
	\draw[color ={black},opacity = 0.5] (-2,-2)--(14,-2);
	\draw[color ={black},opacity = 0.5] (1,-2)--(1,12);
	\draw[color ={black},opacity = 0.5] (2,-2)--(2,12);
	\draw[color ={black},opacity = 0.5] (3,-2)--(3,12);
	\draw[color ={black},opacity = 0.5] (4,-2)--(4,12);
	\draw[color ={black},opacity = 0.5] (5,-2)--(5,12);
	\draw[color ={black},opacity = 0.5] (6,-2)--(6,12);
	\draw[color ={black},opacity = 0.5] (7,-2)--(7,12);
	\draw[color ={black},opacity = 0.5] (8,-2)--(8,12);
	\draw[color ={black},opacity = 0.5] (9,-2)--(9,12);
	\draw[color ={black},opacity = 0.5] (10,-2)--(10,12);
	\draw[color ={black},opacity = 0.5] (11,-2)--(11,12);
	\draw[color ={black},opacity = 0.5] (12,-2)--(12,12);
	\draw[color ={black},opacity = 0.5] (13,-2)--(13,12);
	\draw[color ={black},opacity = 0.5] (14,-2)--(14,12);
	\draw[color ={black},opacity = 0.5] (-1,-2)--(-1,12);
	\draw[color ={black},opacity = 0.5] (-2,-2)--(-2,12);
	\draw[color ={gray},opacity = 0.3] (-2,0)--(14,0);
	\draw[color ={gray},opacity = 0.3] (-2,0.5)--(14,0.5);
	\draw[color ={gray},opacity = 0.3] (-2,1.5)--(14,1.5);
	\draw[color ={gray},opacity = 0.3] (-2,2.5)--(14,2.5);
	\draw[color ={gray},opacity = 0.3] (-2,3.5)--(14,3.5);
	\draw[color ={gray},opacity = 0.3] (-2,4.5)--(14,4.5);
	\draw[color ={gray},opacity = 0.3] (-2,5.5)--(14,5.5);
	\draw[color ={gray},opacity = 0.3] (-2,6.5)--(14,6.5);
	\draw[color ={gray},opacity = 0.3] (-2,7.5)--(14,7.5);
	\draw[color ={gray},opacity = 0.3] (-2,8.5)--(14,8.5);
	\draw[color ={gray},opacity = 0.3] (-2,9.5)--(14,9.5);
	\draw[color ={gray},opacity = 0.3] (-2,10.5)--(14,10.5);
	\draw[color ={gray},opacity = 0.3] (-2,11.5)--(14,11.5);
	\draw[color ={gray},opacity = 0.3] (-2,0)--(14,0);
	\draw[color ={gray},opacity = 0.3] (-2,-0.5)--(14,-0.5);
	\draw[color ={gray},opacity = 0.3] (-2,-1.5)--(14,-1.5);
	\draw[color ={gray},opacity = 0.3] (0,-2)--(0,12);
	\draw[color ={gray},opacity = 0.3] (0.1,-2)--(0.1,12);
	\draw[color ={gray},opacity = 0.3] (0.2,-2)--(0.2,12);
	\draw[color ={gray},opacity = 0.3] (0.30000000000000004,-2)--(0.30000000000000004,12);
	\draw[color ={gray},opacity = 0.3] (0.4,-2)--(0.4,12);
	\draw[color ={gray},opacity = 0.3] (0.5,-2)--(0.5,12);
	\draw[color ={gray},opacity = 0.3] (0.6,-2)--(0.6,12);
	\draw[color ={gray},opacity = 0.3] (0.7,-2)--(0.7,12);
	\draw[color ={gray},opacity = 0.3] (0.7999999999999999,-2)--(0.7999999999999999,12);
	\draw[color ={gray},opacity = 0.3] (0.8999999999999999,-2)--(0.8999999999999999,12);
	\draw[color ={gray},opacity = 0.3] (1.0999999999999999,-2)--(1.0999999999999999,12);
	\draw[color ={gray},opacity = 0.3] (1.2,-2)--(1.2,12);
	\draw[color ={gray},opacity = 0.3] (1.3,-2)--(1.3,12);
	\draw[color ={gray},opacity = 0.3] (1.4000000000000001,-2)--(1.4000000000000001,12);
	\draw[color ={gray},opacity = 0.3] (1.5000000000000002,-2)--(1.5000000000000002,12);
	\draw[color ={gray},opacity = 0.3] (1.6000000000000003,-2)--(1.6000000000000003,12);
	\draw[color ={gray},opacity = 0.3] (1.7000000000000004,-2)--(1.7000000000000004,12);
	\draw[color ={gray},opacity = 0.3] (1.8000000000000005,-2)--(1.8000000000000005,12);
	\draw[color ={gray},opacity = 0.3] (1.9000000000000006,-2)--(1.9000000000000006,12);
	\draw[color ={gray},opacity = 0.3] (2.1000000000000005,-2)--(2.1000000000000005,12);
	\draw[color ={gray},opacity = 0.3] (2.2000000000000006,-2)--(2.2000000000000006,12);
	\draw[color ={gray},opacity = 0.3] (2.3000000000000007,-2)--(2.3000000000000007,12);
	\draw[color ={gray},opacity = 0.3] (2.400000000000001,-2)--(2.400000000000001,12);
	\draw[color ={gray},opacity = 0.3] (2.500000000000001,-2)--(2.500000000000001,12);
	\draw[color ={gray},opacity = 0.3] (2.600000000000001,-2)--(2.600000000000001,12);
	\draw[color ={gray},opacity = 0.3] (2.700000000000001,-2)--(2.700000000000001,12);
	\draw[color ={gray},opacity = 0.3] (2.800000000000001,-2)--(2.800000000000001,12);
	\draw[color ={gray},opacity = 0.3] (2.9000000000000012,-2)--(2.9000000000000012,12);
	\draw[color ={gray},opacity = 0.3] (3.1000000000000014,-2)--(3.1000000000000014,12);
	\draw[color ={gray},opacity = 0.3] (3.2000000000000015,-2)--(3.2000000000000015,12);
	\draw[color ={gray},opacity = 0.3] (3.3000000000000016,-2)--(3.3000000000000016,12);
	\draw[color ={gray},opacity = 0.3] (3.4000000000000017,-2)--(3.4000000000000017,12);
	\draw[color ={gray},opacity = 0.3] (3.5000000000000018,-2)--(3.5000000000000018,12);
	\draw[color ={gray},opacity = 0.3] (3.600000000000002,-2)--(3.600000000000002,12);
	\draw[color ={gray},opacity = 0.3] (3.700000000000002,-2)--(3.700000000000002,12);
	\draw[color ={gray},opacity = 0.3] (3.800000000000002,-2)--(3.800000000000002,12);
	\draw[color ={gray},opacity = 0.3] (3.900000000000002,-2)--(3.900000000000002,12);
	\draw[color ={gray},opacity = 0.3] (4.100000000000001,-2)--(4.100000000000001,12);
	\draw[color ={gray},opacity = 0.3] (4.200000000000001,-2)--(4.200000000000001,12);
	\draw[color ={gray},opacity = 0.3] (4.300000000000001,-2)--(4.300000000000001,12);
	\draw[color ={gray},opacity = 0.3] (4.4,-2)--(4.4,12);
	\draw[color ={gray},opacity = 0.3] (4.5,-2)--(4.5,12);
	\draw[color ={gray},opacity = 0.3] (4.6,-2)--(4.6,12);
	\draw[color ={gray},opacity = 0.3] (4.699999999999999,-2)--(4.699999999999999,12);
	\draw[color ={gray},opacity = 0.3] (4.799999999999999,-2)--(4.799999999999999,12);
	\draw[color ={gray},opacity = 0.3] (4.899999999999999,-2)--(4.899999999999999,12);
	\draw[color ={gray},opacity = 0.3] (5.099999999999998,-2)--(5.099999999999998,12);
	\draw[color ={gray},opacity = 0.3] (5.1999999999999975,-2)--(5.1999999999999975,12);
	\draw[color ={gray},opacity = 0.3] (5.299999999999997,-2)--(5.299999999999997,12);
	\draw[color ={gray},opacity = 0.3] (5.399999999999997,-2)--(5.399999999999997,12);
	\draw[color ={gray},opacity = 0.3] (5.4999999999999964,-2)--(5.4999999999999964,12);
	\draw[color ={gray},opacity = 0.3] (5.599999999999996,-2)--(5.599999999999996,12);
	\draw[color ={gray},opacity = 0.3] (5.699999999999996,-2)--(5.699999999999996,12);
	\draw[color ={gray},opacity = 0.3] (5.799999999999995,-2)--(5.799999999999995,12);
	\draw[color ={gray},opacity = 0.3] (5.899999999999995,-2)--(5.899999999999995,12);
	\draw[color ={gray},opacity = 0.3] (6.099999999999994,-2)--(6.099999999999994,12);
	\draw[color ={gray},opacity = 0.3] (6.199999999999994,-2)--(6.199999999999994,12);
	\draw[color ={gray},opacity = 0.3] (6.299999999999994,-2)--(6.299999999999994,12);
	\draw[color ={gray},opacity = 0.3] (6.399999999999993,-2)--(6.399999999999993,12);
	\draw[color ={gray},opacity = 0.3] (6.499999999999993,-2)--(6.499999999999993,12);
	\draw[color ={gray},opacity = 0.3] (6.5999999999999925,-2)--(6.5999999999999925,12);
	\draw[color ={gray},opacity = 0.3] (6.699999999999992,-2)--(6.699999999999992,12);
	\draw[color ={gray},opacity = 0.3] (6.799999999999992,-2)--(6.799999999999992,12);
	\draw[color ={gray},opacity = 0.3] (6.8999999999999915,-2)--(6.8999999999999915,12);
	\draw[color ={gray},opacity = 0.3] (7.099999999999991,-2)--(7.099999999999991,12);
	\draw[color ={gray},opacity = 0.3] (7.19999999999999,-2)--(7.19999999999999,12);
	\draw[color ={gray},opacity = 0.3] (7.29999999999999,-2)--(7.29999999999999,12);
	\draw[color ={gray},opacity = 0.3] (7.39999999999999,-2)--(7.39999999999999,12);
	\draw[color ={gray},opacity = 0.3] (7.499999999999989,-2)--(7.499999999999989,12);
	\draw[color ={gray},opacity = 0.3] (7.599999999999989,-2)--(7.599999999999989,12);
	\draw[color ={gray},opacity = 0.3] (7.699999999999989,-2)--(7.699999999999989,12);
	\draw[color ={gray},opacity = 0.3] (7.799999999999988,-2)--(7.799999999999988,12);
	\draw[color ={gray},opacity = 0.3] (7.899999999999988,-2)--(7.899999999999988,12);
	\draw[color ={gray},opacity = 0.3] (8.099999999999987,-2)--(8.099999999999987,12);
	\draw[color ={gray},opacity = 0.3] (8.199999999999987,-2)--(8.199999999999987,12);
	\draw[color ={gray},opacity = 0.3] (8.299999999999986,-2)--(8.299999999999986,12);
	\draw[color ={gray},opacity = 0.3] (8.399999999999986,-2)--(8.399999999999986,12);
	\draw[color ={gray},opacity = 0.3] (8.499999999999986,-2)--(8.499999999999986,12);
	\draw[color ={gray},opacity = 0.3] (8.599999999999985,-2)--(8.599999999999985,12);
	\draw[color ={gray},opacity = 0.3] (8.699999999999985,-2)--(8.699999999999985,12);
	\draw[color ={gray},opacity = 0.3] (8.799999999999985,-2)--(8.799999999999985,12);
	\draw[color ={gray},opacity = 0.3] (8.899999999999984,-2)--(8.899999999999984,12);
	\draw[color ={gray},opacity = 0.3] (9.099999999999984,-2)--(9.099999999999984,12);
	\draw[color ={gray},opacity = 0.3] (9.199999999999983,-2)--(9.199999999999983,12);
	\draw[color ={gray},opacity = 0.3] (9.299999999999983,-2)--(9.299999999999983,12);
	\draw[color ={gray},opacity = 0.3] (9.399999999999983,-2)--(9.399999999999983,12);
	\draw[color ={gray},opacity = 0.3] (9.499999999999982,-2)--(9.499999999999982,12);
	\draw[color ={gray},opacity = 0.3] (9.599999999999982,-2)--(9.599999999999982,12);
	\draw[color ={gray},opacity = 0.3] (9.699999999999982,-2)--(9.699999999999982,12);
	\draw[color ={gray},opacity = 0.3] (9.799999999999981,-2)--(9.799999999999981,12);
	\draw[color ={gray},opacity = 0.3] (9.89999999999998,-2)--(9.89999999999998,12);
	\draw[color ={gray},opacity = 0.3] (10.09999999999998,-2)--(10.09999999999998,12);
	\draw[color ={gray},opacity = 0.3] (10.19999999999998,-2)--(10.19999999999998,12);
	\draw[color ={gray},opacity = 0.3] (10.29999999999998,-2)--(10.29999999999998,12);
	\draw[color ={gray},opacity = 0.3] (10.399999999999979,-2)--(10.399999999999979,12);
	\draw[color ={gray},opacity = 0.3] (10.499999999999979,-2)--(10.499999999999979,12);
	\draw[color ={gray},opacity = 0.3] (10.599999999999978,-2)--(10.599999999999978,12);
	\draw[color ={gray},opacity = 0.3] (10.699999999999978,-2)--(10.699999999999978,12);
	\draw[color ={gray},opacity = 0.3] (10.799999999999978,-2)--(10.799999999999978,12);
	\draw[color ={gray},opacity = 0.3] (10.899999999999977,-2)--(10.899999999999977,12);
	\draw[color ={gray},opacity = 0.3] (11.099999999999977,-2)--(11.099999999999977,12);
	\draw[color ={gray},opacity = 0.3] (11.199999999999976,-2)--(11.199999999999976,12);
	\draw[color ={gray},opacity = 0.3] (11.299999999999976,-2)--(11.299999999999976,12);
	\draw[color ={gray},opacity = 0.3] (11.399999999999975,-2)--(11.399999999999975,12);
	\draw[color ={gray},opacity = 0.3] (11.499999999999975,-2)--(11.499999999999975,12);
	\draw[color ={gray},opacity = 0.3] (11.599999999999975,-2)--(11.599999999999975,12);
	\draw[color ={gray},opacity = 0.3] (11.699999999999974,-2)--(11.699999999999974,12);
	\draw[color ={gray},opacity = 0.3] (11.799999999999974,-2)--(11.799999999999974,12);
	\draw[color ={gray},opacity = 0.3] (11.899999999999974,-2)--(11.899999999999974,12);
	\draw[color ={gray},opacity = 0.3] (12.099999999999973,-2)--(12.099999999999973,12);
	\draw[color ={gray},opacity = 0.3] (12.199999999999973,-2)--(12.199999999999973,12);
	\draw[color ={gray},opacity = 0.3] (12.299999999999972,-2)--(12.299999999999972,12);
	\draw[color ={gray},opacity = 0.3] (12.399999999999972,-2)--(12.399999999999972,12);
	\draw[color ={gray},opacity = 0.3] (12.499999999999972,-2)--(12.499999999999972,12);
	\draw[color ={gray},opacity = 0.3] (12.599999999999971,-2)--(12.599999999999971,12);
	\draw[color ={gray},opacity = 0.3] (12.69999999999997,-2)--(12.69999999999997,12);
	\draw[color ={gray},opacity = 0.3] (12.79999999999997,-2)--(12.79999999999997,12);
	\draw[color ={gray},opacity = 0.3] (12.89999999999997,-2)--(12.89999999999997,12);
	\draw[color ={gray},opacity = 0.3] (13.09999999999997,-2)--(13.09999999999997,12);
	\draw[color ={gray},opacity = 0.3] (13.199999999999969,-2)--(13.199999999999969,12);
	\draw[color ={gray},opacity = 0.3] (13.299999999999969,-2)--(13.299999999999969,12);
	\draw[color ={gray},opacity = 0.3] (13.399999999999968,-2)--(13.399999999999968,12);
	\draw[color ={gray},opacity = 0.3] (13.499999999999968,-2)--(13.499999999999968,12);
	\draw[color ={gray},opacity = 0.3] (13.599999999999968,-2)--(13.599999999999968,12);
	\draw[color ={gray},opacity = 0.3] (13.699999999999967,-2)--(13.699999999999967,12);
	\draw[color ={gray},opacity = 0.3] (13.799999999999967,-2)--(13.799999999999967,12);
	\draw[color ={gray},opacity = 0.3] (13.899999999999967,-2)--(13.899999999999967,12);
	\draw[color ={gray},opacity = 0.3] (0,-2)--(0,12);
	\draw[color ={gray},opacity = 0.3] (-0.1,-2)--(-0.1,12);
	\draw[color ={gray},opacity = 0.3] (-0.2,-2)--(-0.2,12);
	\draw[color ={gray},opacity = 0.3] (-0.30000000000000004,-2)--(-0.30000000000000004,12);
	\draw[color ={gray},opacity = 0.3] (-0.4,-2)--(-0.4,12);
	\draw[color ={gray},opacity = 0.3] (-0.5,-2)--(-0.5,12);
	\draw[color ={gray},opacity = 0.3] (-0.6,-2)--(-0.6,12);
	\draw[color ={gray},opacity = 0.3] (-0.7,-2)--(-0.7,12);
	\draw[color ={gray},opacity = 0.3] (-0.7999999999999999,-2)--(-0.7999999999999999,12);
	\draw[color ={gray},opacity = 0.3] (-0.8999999999999999,-2)--(-0.8999999999999999,12);
	\draw[color ={gray},opacity = 0.3] (-1.0999999999999999,-2)--(-1.0999999999999999,12);
	\draw[color ={gray},opacity = 0.3] (-1.2,-2)--(-1.2,12);
	\draw[color ={gray},opacity = 0.3] (-1.3,-2)--(-1.3,12);
	\draw[color ={gray},opacity = 0.3] (-1.4000000000000001,-2)--(-1.4000000000000001,12);
	\draw[color ={gray},opacity = 0.3] (-1.5000000000000002,-2)--(-1.5000000000000002,12);
	\draw[color ={gray},opacity = 0.3] (-1.6000000000000003,-2)--(-1.6000000000000003,12);
	\draw[color ={gray},opacity = 0.3] (-1.7000000000000004,-2)--(-1.7000000000000004,12);
	\draw[color ={gray},opacity = 0.3] (-1.8000000000000005,-2)--(-1.8000000000000005,12);
	\draw[color ={gray},opacity = 0.3] (-1.9000000000000006,-2)--(-1.9000000000000006,12);
	\draw[color ={black},line width = 2] (2,-0.2)--(2,0.2);
	\draw[color ={black},line width = 2] (4,-0.2)--(4,0.2);
	\draw[color ={black},line width = 2] (6,-0.2)--(6,0.2);
	\draw[color ={black},line width = 2] (8,-0.2)--(8,0.2);
	\draw[color ={black},line width = 2] (10,-0.2)--(10,0.2);
	\draw[color ={black},line width = 2] (12,-0.2)--(12,0.2);
	\draw[color ={black},line width = 2] (-0.2,2)--(0.2,2);
	\draw[color ={black},line width = 2] (-0.2,4)--(0.2,4);
	\draw[color ={black},line width = 2] (-0.2,6)--(0.2,6);
	\draw[color ={black},line width = 2] (-0.2,8)--(0.2,8);
	\draw[color ={black},line width = 2] (-0.2,10)--(0.2,10);
	\draw (2,-0.4) node[anchor = center] {\scriptsize \color{black}{$1$}};
	\draw (4,-0.4) node[anchor = center] {\scriptsize \color{black}{$2$}};
	\draw (6,-0.4) node[anchor = center] {\scriptsize \color{black}{$3$}};
	\draw (8,-0.4) node[anchor = center] {\scriptsize \color{black}{$4$}};
	\draw (10,-0.4) node[anchor = center] {\scriptsize \color{black}{$5$}};
	\draw (-0.5,2.1) node[anchor = center] {\footnotesize \color{black}{$1$}};
	\draw (-0.5,4.1) node[anchor = center] {\footnotesize \color{black}{$2$}};
	\draw (-0.5,6.1) node[anchor = center] {\footnotesize \color{black}{$3$}};
	\draw (-0.5,8.1) node[anchor = center] {\footnotesize \color{black}{$4$}};
	\draw (-0.5,10.1) node[anchor = center] {\footnotesize \color{black}{$5$}};
	\draw [color={black}] (-0.3,-0.3) node[anchor = center,scale=1, rotate = 0] {O};
	
	\draw[color={blue},line width = 2] (1,10)--(1.4,7.71)--(1.8,6.44)--(2.2,5.64)--(2.6,5.08)--(3,4.67)--(3.4,4.35)--(3.8,4.11)--(4.2,3.9)--(4.6,3.74)--(5,3.6)--(5.4,3.48)--(5.8,3.38)--(6.2,3.29)--(6.6,3.21)--(7,3.14)--(7.4,3.08)--(7.8,3.03)--(8.2,2.98)--(8.6,2.93)--(9,2.89)--(9.4,2.85)--(9.8,2.82)--(10.2,2.78)--(10.6,2.75)--(11,2.73)--(11.4,2.7)--(11.8,2.68)--(12.2,2.66)--(12.6,2.63)--(13,2.62)--(13.4,2.6)--(13.8,2.58);
	
	\draw[color={red},line width = 2] (-1.6,13.2)--(-1.2,12.4)--(-0.8,11.6)--(-0.4,10.8)--(0,10)--(0.4,9.2)--(0.8,8.4)--(1.2,7.6)--(1.6,6.8)--(2,6)--(2.4,5.2)--(2.8,4.4)--(3.2,3.6)--(3.6,2.8)--(4,2)--(4.4,1.2)--(4.8,0.4)--(5.2,-0.4)--(5.6,-1.2)--(6,-2)--(6.4,-2.8)--(6.8,-3.6);

\end{tikzpicture}\end{center}
\end{exercice}

\begin{Solution}
 $f'(1)$ est donné par le coefficient directeur de la tangente à la courbe au point d'abscisse $1$, soit $-2=-2$.
\end{Solution}

% @see : https://coopmaths.fr/alea?uuid=6f32d&id=can1F16&n=1&d=10&alea=dqSw22323232323423456789101112131415161718192021222324252627282930313233343536373839404142434445464748234567891011121314151617181920212223&cd=1&cols=1
\needspace{10\baselineskip}
\newpage
\begin{exercice}[BaremeTotal=false]


 On donne les représentations graphiques d'une fonction et de sa dérivée.\\
        Donner l'équation réduite de la tangente à la courbe de $f$ en $x=0$. \\ \begin{minipage}{0.5\linewidth}
    \begin{tikzpicture}[baseline, scale=0.8]

    \tikzset{
      point/.style={
        thick,
        draw,
        cross out,
        inner sep=0pt,
        minimum width=5pt,
        minimum height=5pt,
      },
    }
    \clip (-4.5,-4.5) rectangle (5.75,13.5);
    	
	\draw[color ={black},>=latex,->] (-3,0)--(3,0);
	\draw[color ={black},>=latex,->] (0,-3)--(0,12);
	\draw[color ={black},opacity = 0.4] (-3,1)--(3,1);
	\draw[color ={black},opacity = 0.4] (-3,2)--(3,2);
	\draw[color ={black},opacity = 0.4] (-3,3)--(3,3);
	\draw[color ={black},opacity = 0.4] (-3,4)--(3,4);
	\draw[color ={black},opacity = 0.4] (-3,5)--(3,5);
	\draw[color ={black},opacity = 0.4] (-3,6)--(3,6);
	\draw[color ={black},opacity = 0.4] (-3,7)--(3,7);
	\draw[color ={black},opacity = 0.4] (-3,8)--(3,8);
	\draw[color ={black},opacity = 0.4] (-3,9)--(3,9);
	\draw[color ={black},opacity = 0.4] (-3,10)--(3,10);
	\draw[color ={black},opacity = 0.4] (-3,11)--(3,11);
	\draw[color ={black},opacity = 0.4] (-3,12)--(3,12);
	\draw[color ={black},opacity = 0.4] (-3,-1)--(3,-1);
	\draw[color ={black},opacity = 0.4] (-3,-2)--(3,-2);
	\draw[color ={black},opacity = 0.4] (-3,-3)--(3,-3);
	\draw[color ={black},opacity = 0.4] (1,-3)--(1,12);
	\draw[color ={black},opacity = 0.4] (2,-3)--(2,12);
	\draw[color ={black},opacity = 0.4] (3,-3)--(3,12);
	\draw[color ={black},opacity = 0.4] (-1,-3)--(-1,12);
	\draw[color ={black},opacity = 0.4] (-2,-3)--(-2,12);
	\draw[color ={black},opacity = 0.4] (-3,-3)--(-3,12);
	\draw[color ={gray},opacity = 0.3] (-3,0)--(3,0);
	\draw[color ={gray},opacity = 0.3] (-3,0.5)--(3,0.5);
	\draw[color ={gray},opacity = 0.3] (-3,1.5)--(3,1.5);
	\draw[color ={gray},opacity = 0.3] (-3,2.5)--(3,2.5);
	\draw[color ={gray},opacity = 0.3] (-3,3.5)--(3,3.5);
	\draw[color ={gray},opacity = 0.3] (-3,4.5)--(3,4.5);
	\draw[color ={gray},opacity = 0.3] (-3,5.5)--(3,5.5);
	\draw[color ={gray},opacity = 0.3] (-3,6.5)--(3,6.5);
	\draw[color ={gray},opacity = 0.3] (-3,7.5)--(3,7.5);
	\draw[color ={gray},opacity = 0.3] (-3,8.5)--(3,8.5);
	\draw[color ={gray},opacity = 0.3] (-3,9.5)--(3,9.5);
	\draw[color ={gray},opacity = 0.3] (-3,10.5)--(3,10.5);
	\draw[color ={gray},opacity = 0.3] (-3,11.5)--(3,11.5);
	\draw[color ={gray},opacity = 0.3] (-3,0)--(3,0);
	\draw[color ={gray},opacity = 0.3] (-3,-0.5)--(3,-0.5);
	\draw[color ={gray},opacity = 0.3] (-3,-1.5)--(3,-1.5);
	\draw[color ={gray},opacity = 0.3] (-3,-2.5)--(3,-2.5);
	\draw[color ={gray},opacity = 0.3] (0,-3)--(0,12);
	\draw[color ={gray},opacity = 0.3] (0,-3)--(0,12);
	\draw[color ={black},line width = 1.2] (1,-0.13)--(1,0.13);
	\draw[color ={black},line width = 1.2] (2,-0.13)--(2,0.13);
	\draw[color ={black},line width = 1.2] (3,-0.13)--(3,0.13);
	\draw[color ={black},line width = 1.2] (-1,-0.13)--(-1,0.13);
	\draw[color ={black},line width = 1.2] (-2,-0.13)--(-2,0.13);
	\draw[color ={black},line width = 1.2] (-3,-0.13)--(-3,0.13);
	\draw[color ={black},line width = 1.2] (-0.13,1)--(0.13,1);
	\draw[color ={black},line width = 1.2] (-0.13,2)--(0.13,2);
	\draw[color ={black},line width = 1.2] (-0.13,3)--(0.13,3);
	\draw[color ={black},line width = 1.2] (-0.13,4)--(0.13,4);
	\draw[color ={black},line width = 1.2] (-0.13,5)--(0.13,5);
	\draw[color ={black},line width = 1.2] (-0.13,6)--(0.13,6);
	\draw[color ={black},line width = 1.2] (-0.13,7)--(0.13,7);
	\draw[color ={black},line width = 1.2] (-0.13,8)--(0.13,8);
	\draw[color ={black},line width = 1.2] (-0.13,9)--(0.13,9);
	\draw[color ={black},line width = 1.2] (-0.13,10)--(0.13,10);
	\draw[color ={black},line width = 1.2] (-0.13,11)--(0.13,11);
	\draw[color ={black},line width = 1.2] (-0.13,12)--(0.13,12);
	\draw[color ={black},line width = 1.2] (-0.13,-1)--(0.13,-1);
	\draw[color ={black},line width = 1.2] (-0.13,-2)--(0.13,-2);
	\draw[color ={black},line width = 1.2] (-0.13,-3)--(0.13,-3);
	\draw (3,-0.4) node[anchor = center] {\scriptsize \color{black}{$3$}};
	\draw (-3,-0.4) node[anchor = center] {\scriptsize \color{black}{$-3$}};
	\draw (-0.5,3.1) node[anchor = center] {\footnotesize \color{black}{$3$}};
	\draw (-0.5,6.1) node[anchor = center] {\footnotesize \color{black}{$6$}};
	\draw (-0.5,9.1) node[anchor = center] {\footnotesize \color{black}{$9$}};
	\draw (-0.5,-2.9) node[anchor = center] {\footnotesize \color{black}{$-3$}};
	\draw [color={black}] (-0.3,-0.3) node[anchor = center,scale=1, rotate = 0] {O};
	\draw (3,10) node[anchor = center] {\footnotesize \color{blue}{$\Large \mathcal{C}_f$}};
	
	\draw[color={blue},line width = 2] (-3,4)--(-2.8,3.24)--(-2.6,2.56)--(-2.4,1.96)--(-2.2,1.44)--(-2,1)--(-1.8,0.64)--(-1.6,0.36)--(-1.4,0.16)--(-1.2,0.04)--(-1,0)--(-0.8,0.04)--(-0.6,0.16)--(-0.4,0.36)--(-0.2,0.64)--(0,1)--(0.2,1.44)--(0.4,1.96)--(0.6,2.56)--(0.8,3.24)--(1,4)--(1.2,4.84)--(1.4,5.76)--(1.6,6.76)--(1.8,7.84)--(2,9)--(2.2,10.24)--(2.4,11.56)--(2.6,12.96);

\end{tikzpicture}\\
    \end{minipage}
    \begin{minipage}{0.5\linewidth}
    \begin{tikzpicture}[baseline, scale=0.8]

    \tikzset{
      point/.style={
        thick,
        draw,
        cross out,
        inner sep=0pt,
        minimum width=5pt,
        minimum height=5pt,
      },
    }
    \clip (-4.5,-6.5) rectangle (6.5,9.5);
    	
	\draw[color ={black},>=latex,->] (-3,0)--(3,0);
	\draw[color ={black},>=latex,->] (0,-5)--(0,8);
	\draw[color ={black},opacity = 0.4] (-3,1)--(3,1);
	\draw[color ={black},opacity = 0.4] (-3,2)--(3,2);
	\draw[color ={black},opacity = 0.4] (-3,3)--(3,3);
	\draw[color ={black},opacity = 0.4] (-3,4)--(3,4);
	\draw[color ={black},opacity = 0.4] (-3,5)--(3,5);
	\draw[color ={black},opacity = 0.4] (-3,6)--(3,6);
	\draw[color ={black},opacity = 0.4] (-3,7)--(3,7);
	\draw[color ={black},opacity = 0.4] (-3,8)--(3,8);
	\draw[color ={black},opacity = 0.4] (-3,-1)--(3,-1);
	\draw[color ={black},opacity = 0.4] (-3,-2)--(3,-2);
	\draw[color ={black},opacity = 0.4] (-3,-3)--(3,-3);
	\draw[color ={black},opacity = 0.4] (-3,-4)--(3,-4);
	\draw[color ={black},opacity = 0.4] (-3,-5)--(3,-5);
	\draw[color ={black},opacity = 0.4] (1,-5)--(1,8);
	\draw[color ={black},opacity = 0.4] (2,-5)--(2,8);
	\draw[color ={black},opacity = 0.4] (3,-5)--(3,8);
	\draw[color ={black},opacity = 0.4] (-1,-5)--(-1,8);
	\draw[color ={black},opacity = 0.4] (-2,-5)--(-2,8);
	\draw[color ={black},opacity = 0.4] (-3,-5)--(-3,8);
	\draw[color ={gray},opacity = 0.3] (-3,0)--(3,0);
	\draw[color ={gray},opacity = 0.3] (-3,0.5)--(3,0.5);
	\draw[color ={gray},opacity = 0.3] (-3,1.5)--(3,1.5);
	\draw[color ={gray},opacity = 0.3] (-3,2.5)--(3,2.5);
	\draw[color ={gray},opacity = 0.3] (-3,3.5)--(3,3.5);
	\draw[color ={gray},opacity = 0.3] (-3,4.5)--(3,4.5);
	\draw[color ={gray},opacity = 0.3] (-3,5.5)--(3,5.5);
	\draw[color ={gray},opacity = 0.3] (-3,6.5)--(3,6.5);
	\draw[color ={gray},opacity = 0.3] (-3,7.5)--(3,7.5);
	\draw[color ={gray},opacity = 0.3] (-3,0)--(3,0);
	\draw[color ={gray},opacity = 0.3] (-3,-0.5)--(3,-0.5);
	\draw[color ={gray},opacity = 0.3] (-3,-1.5)--(3,-1.5);
	\draw[color ={gray},opacity = 0.3] (-3,-2.5)--(3,-2.5);
	\draw[color ={gray},opacity = 0.3] (-3,-3.5)--(3,-3.5);
	\draw[color ={gray},opacity = 0.3] (-3,-4.5)--(3,-4.5);
	\draw[color ={gray},opacity = 0.3] (0,-5)--(0,8);
	\draw[color ={gray},opacity = 0.3] (0,-5)--(0,8);
	\draw[color ={black},line width = 1.2] (1,-0.13)--(1,0.13);
	\draw[color ={black},line width = 1.2] (2,-0.13)--(2,0.13);
	\draw[color ={black},line width = 1.2] (3,-0.13)--(3,0.13);
	\draw[color ={black},line width = 1.2] (-1,-0.13)--(-1,0.13);
	\draw[color ={black},line width = 1.2] (-2,-0.13)--(-2,0.13);
	\draw[color ={black},line width = 1.2] (-3,-0.13)--(-3,0.13);
	\draw[color ={black},line width = 1.2] (-0.13,1)--(0.13,1);
	\draw[color ={black},line width = 1.2] (-0.13,2)--(0.13,2);
	\draw[color ={black},line width = 1.2] (-0.13,3)--(0.13,3);
	\draw[color ={black},line width = 1.2] (-0.13,4)--(0.13,4);
	\draw[color ={black},line width = 1.2] (-0.13,5)--(0.13,5);
	\draw[color ={black},line width = 1.2] (-0.13,6)--(0.13,6);
	\draw[color ={black},line width = 1.2] (-0.13,7)--(0.13,7);
	\draw[color ={black},line width = 1.2] (-0.13,8)--(0.13,8);
	\draw[color ={black},line width = 1.2] (-0.13,9)--(0.13,9);
	\draw[color ={black},line width = 1.2] (-0.13,10)--(0.13,10);
	\draw[color ={black},line width = 1.2] (-0.13,11)--(0.13,11);
	\draw[color ={black},line width = 1.2] (-0.13,12)--(0.13,12);
	\draw[color ={black},line width = 1.2] (-0.13,-1)--(0.13,-1);
	\draw[color ={black},line width = 1.2] (-0.13,-2)--(0.13,-2);
	\draw[color ={black},line width = 1.2] (-0.13,-3)--(0.13,-3);
	\draw (3,-0.4) node[anchor = center] {\scriptsize \color{black}{$3$}};
	\draw (-3,-0.4) node[anchor = center] {\scriptsize \color{black}{$-3$}};
	\draw (-0.5,3.1) node[anchor = center] {\footnotesize \color{black}{$3$}};
	\draw (-0.5,6.1) node[anchor = center] {\footnotesize \color{black}{$6$}};
	\draw (-0.5,-2.9) node[anchor = center] {\footnotesize \color{black}{$-3$}};
	\draw [color={black}] (-0.3,-0.3) node[anchor = center,scale=1, rotate = 0] {O};
	\draw (3,6) node[anchor = center] {\footnotesize \color{red}{$\Large\mathcal{C}_f\prime$}};
	
	\draw[color={red},line width = 2] (-3,-4)--(-2.8,-3.6)--(-2.6,-3.2)--(-2.4,-2.8)--(-2.2,-2.4)--(-2,-2)--(-1.8,-1.6)--(-1.6,-1.2)--(-1.4,-0.8)--(-1.2,-0.4)--(-1,0)--(-0.8,0.4)--(-0.6,0.8)--(-0.4,1.2)--(-0.2,1.6)--(0,2)--(0.2,2.4)--(0.4,2.8)--(0.6,3.2)--(0.8,3.6)--(1,4)--(1.2,4.4)--(1.4,4.8)--(1.6,5.2)--(1.8,5.6)--(2,6)--(2.2,6.4)--(2.4,6.8)--(2.6,7.2)--(2.8,7.6)--(3,8);

\end{tikzpicture}\\
    \end{minipage}
    
\end{exercice}

\begin{Solution}
 L'équation réduite de la tangente au point d'abscisse $0$ est  : $y=f'(0)(x-0)+f(0)$.\\
        On lit graphiquement $f(0)=1$ et $f'(0)=2$.\\
        L'équation réduite de la tangente est donc donnée par :
        $y=2(x+0)+1$, soit $y=2x+1$.
\end{Solution}

% @see : https://coopmaths.fr/alea?uuid=a1ba2&id=can1F14&n=1&d=10&alea=rkIp22323232323423456789101112131415161718192021222324252627282930313233343536373839404142434445464748234567891011121314151617181920212223&cd=1&cols=1
\needspace{10\baselineskip}
\begin{exercice}[BaremeTotal=false]


Soit $f$ la fonction définie sur $\mathbb{R}$ par : 
$f(x)= 1+4x^2+2x$.\\
En calculant la limite du taux d'accroissement, déterminer $f'(1)$.
\end{exercice}

\begin{Solution}
 $f$ est une fonction polynôme du second degré de la forme $f(x)=ax^2+bx+c$.\\
    La fonction dérivée est donnée par la somme des dérivées des fonctions $u$ et $v$ définies par $u(x)=4x^2$ et $v(x)=2x+1$.\\
     Comme $u'(x)=8x$ et $v'(x)=2$, on obtient  $f'(x)=8x+2$. \\
     Ainsi, $f'(1)=8\times 1+2=10$.
\end{Solution}

% @see : https://coopmaths.fr/alea?uuid=3c690&id=can1F13&n=1&d=10&alea=EUDu22323232323423456789101112131415161718192021222324252627282930313233343536373839404142434445464748234567891011121314151617181920212223&cd=1&cols=1
\needspace{10\baselineskip}
\begin{exercice}[BaremeTotal=false]


 Déterminer le coefficient directeur de la tangente à la courbe représentative de la fonction inverse au point d'abscisse $6$.
                 
\end{exercice}

\begin{Solution}
 Le coefficient directeur de la tangente au point d'abscisse $a$ est donné par le nombre dérivé $f'(a)$.\\
        La fonction $f$ définie par $f(x)=\dfrac{1}{x}$ a pour fonction dérivée la fonction $f'$ définie par $f'(x)=-\dfrac{1}{x^2}$.\\
Comme $f'(6)=-\dfrac{1}{6^2}=-\dfrac{1}{36}$, le coefficient directeur de la tangente au point d'abscisse $6$ est : $-\dfrac{1}{36}$.
\end{Solution}

% @see : https://coopmaths.fr/alea?uuid=ebd89&id=1AN14-1&n=1&d=10&s=3&alea=ocYr22323232323423456789101112131415161718192021222324252627282930313233343536373839404142434445464748234567891011121314151617181920212223&cd=0&cols=1
\needspace{10\baselineskip}
\begin{exercice}[BaremeTotal=false]


 Donner la dérivée de la fonction $f$, dérivable pour tout $x\in\R$, définie par  $f(x)=\dfrac{9}{10}x+4$.
\end{exercice}

\begin{Solution}
 L'expression de la dérivée de la fonction $f$ définie par $f(x)=\dfrac{9}{10}x+4$ est : ${\color[HTML]{f15929}\boldsymbol{f'(x)=\dfrac{9}{10}}}$.
\end{Solution}

% @see : https://coopmaths.fr/alea?uuid=ebd89&id=1AN14-1&n=1&d=10&s=4&alea=uAlP22323232323423456789101112131415161718192021222324252627282930313233343536373839404142434445464748234567891011121314151617181920212223&cd=0&cols=1
\needspace{10\baselineskip}
\begin{exercice}[BaremeTotal=false]


 Donner la dérivée de la fonction $f$, dérivable pour tout $x\in\R$, définie par  $f(x)=-5x^{8}$.
\end{exercice}

\begin{Solution}
 L'expression de la dérivée de la fonction $f$ définie par $f(x)=-5x^{8}$ est : ${\color[HTML]{f15929}\boldsymbol{f'(x)=-40x^{7}}}$.
\end{Solution}

% @see : https://coopmaths.fr/alea?uuid=ebd89&id=1AN14-1&n=1&d=10&s=7&alea=vjN722323232323423456789101112131415161718192021222324252627282930313233343536373839404142434445464748234567891011121314151617181920212223&cd=0&cols=1
\needspace{10\baselineskip}
\begin{exercice}[BaremeTotal=false]


 Donner la dérivée de la fonction $f$, dérivable pour tout $x\in\R^*$, définie par  $f(x)=\dfrac{-3}{4x}$.
\end{exercice}

\begin{Solution}
 L'expression de la dérivée de la fonction $f$ définie par $f(x)=\dfrac{-3}{4x}$ est : ${\color[HTML]{f15929}\boldsymbol{f'(x)=\dfrac{3}{4x^2}}}$.
\end{Solution}

% @see : https://coopmaths.fr/alea?uuid=a83c0&id=1AN14-4&n=1&d=10&s=1&alea=Njr322323232323423456789101112131415161718192021222324252627282930313233343536373839404142434445464748234567891011121314151617181920212223&cd=0&cols=1
\needspace{10\baselineskip}
\begin{exercice}[BaremeTotal=false]


 Donner l'expression de la dérivée de la fonction $f$ définie sur $\R^{*}$ par $f(x)=5x^{2}+3x-\dfrac{8}{x}$.\\
\end{exercice}

\begin{Solution}
 L'expression de la dérivée de la fonction $f$ définie par $f(x)=5x^{2}+3x-\dfrac{8}{x}$ est :\\$f'(x)={\color[HTML]{f15929}\boldsymbol{10x+3+\dfrac{8}{x^2}}}$.\\
\end{Solution}

% @see : https://coopmaths.fr/alea?uuid=a83c0&id=1AN14-4&n=1&d=10&s=4&alea=4Vpz22323232323423456789101112131415161718192021222324252627282930313233343536373839404142434445464748234567891011121314151617181920212223&cd=1&cols=1
\needspace{10\baselineskip}
\begin{exercice}[BaremeTotal=false]


 Donner l'expression de la dérivée de la fonction $f$ définie sur $\R_+^{*}$ par $f(x)=-\dfrac{5}{x}-2x^{2}-2\sqrt{x}$.\\
\end{exercice}

\begin{Solution}
 $(-\dfrac{5}{x})^\prime=\dfrac{5}{x^2}$\\$(-2x^{2})^\prime=-4x$\\$(-2\sqrt{x})^\prime=-\dfrac{2}{2\sqrt{x}}$\\L'expression de la dérivée de la fonction $f$ définie par $f(x)=-\dfrac{5}{x}-2x^{2}-2\sqrt{x}$ est :\\$f'(x)={\color[HTML]{f15929}\boldsymbol{\dfrac{5}{x^2}-4x-\dfrac{2}{2\sqrt{x}}}}$.\\
\end{Solution}

\end{Maquette}
\clearpage
\begin{Maquette}[DM]{Niveau= ,Classe=\getdata[24]\mydata\ (v24), Date = 06/01/25  ,Theme=Exercices}


% @see : https://coopmaths.fr/alea?uuid=4c8c7&id=1AN11-3&n=1&d=10&s=2&alea=nawO2232323232342345678910111213141516171819202122232425262728293031323334353637383940414243444546474823456789101112131415161718192021222324&cd=1&cols=2
\needspace{10\baselineskip}
\begin{exercice}[BaremeTotal=false]
\label{eleve:24}


 \begin{minipage}[t]{\linewidth} Soit $f$ une fonction dérivable sur $[-5;5]$ et $\mathcal{C}_f$ sa courbe représentative.\\On sait que  $f(-3)=-5~~$ et que $~~f'(-3)=-5$.\\Déterminer une équation de la tangente $(T)$ à la courbe $\mathcal{C}_f$ au point d'abscisse $-3$,\\sans utiliser la formule de cours de l'équation de tangente. \end{minipage}

\end{exercice}

\begin{Solution}
 \begin{minipage}[t]{\linewidth} On sait que la tangente n'est pas une droite verticale, puisque la fonction est dérivable sur l'intervalle.\\On en déduit que la tangente $(T)$ au point d'abscisse $-3$, admet une équation réduite de la forme :  $(T) : y=m x + p$. 

\medskip
$\bullet$ Détermination de $m$ :\\On sait que le nombre dérivé en $-3$ est par définition, le coefficient directeur de la tangente au point d'abscisse $-3$.\\Par conséquent, on a déjà : $m=f'(-3)=-5$.\\ On en déduit que  $(T) : y= -5 x + p$\\$\bullet$ Détermination de $p$ :\\ Pour cela, on utilise que si $f(-3)=-5~~$, alors le point $A$ de coordonnées $(-3;-5)$ appartient à $\mathcal{C}_f$ mais aussi à $(T)$.\\ On peut écrire $A(-3;-5) = \mathcal{C}_f \cap (T)$.\\ On remplace alors les coordonnées de $A(-3;-5)$ dans l'équation  $(T) : y= -5 x + p$.\\ $\begin{aligned} \phantom{\iff}&A(-3;-5)\in (T)\\ \iff& -5= -5 \times (-3)  + p\\ \iff& p=-5 +5 \times (-3)  \\ \iff& p=-5 -15   \\ \iff& p=-20\\\end{aligned}$\\On peut conclure que : $(T) : {\color[HTML]{f15929}\boldsymbol{y=-5x-20}}$. \end{minipage}

\end{Solution}

% @see : https://coopmaths.fr/alea?uuid=0e984&id=can1F15&n=1&d=10&alea=1Alo2232323232342345678910111213141516171819202122232425262728293031323334353637383940414243444546474823456789101112131415161718192021222324&cd=1&cols=1
\needspace{10\baselineskip}
\begin{exercice}[BaremeTotal=false]


 La courbe représente une fonction $f$ et la droite est la tangente au point d'abscisse $2$.\\
        Déterminer $f'(2)$.\begin{center}\begin{tikzpicture}[baseline,scale = 0.5]

    \tikzset{
      point/.style={
        thick,
        draw,
        cross out,
        inner sep=0pt,
        minimum width=5pt,
        minimum height=5pt,
      },
    }
    \clip (-2,-2) rectangle (14,12);
    	
	\draw[color ={black},line width = 2,>=latex,->] (-2,0)--(14,0);
	\draw[color ={black},line width = 2,>=latex,->] (0,-2)--(0,12);
	\draw[color ={black},opacity = 0.5] (-2,1)--(14,1);
	\draw[color ={black},opacity = 0.5] (-2,2)--(14,2);
	\draw[color ={black},opacity = 0.5] (-2,3)--(14,3);
	\draw[color ={black},opacity = 0.5] (-2,4)--(14,4);
	\draw[color ={black},opacity = 0.5] (-2,5)--(14,5);
	\draw[color ={black},opacity = 0.5] (-2,6)--(14,6);
	\draw[color ={black},opacity = 0.5] (-2,7)--(14,7);
	\draw[color ={black},opacity = 0.5] (-2,8)--(14,8);
	\draw[color ={black},opacity = 0.5] (-2,9)--(14,9);
	\draw[color ={black},opacity = 0.5] (-2,10)--(14,10);
	\draw[color ={black},opacity = 0.5] (-2,11)--(14,11);
	\draw[color ={black},opacity = 0.5] (-2,12)--(14,12);
	\draw[color ={black},opacity = 0.5] (-2,-1)--(14,-1);
	\draw[color ={black},opacity = 0.5] (-2,-2)--(14,-2);
	\draw[color ={black},opacity = 0.5] (1,-2)--(1,12);
	\draw[color ={black},opacity = 0.5] (2,-2)--(2,12);
	\draw[color ={black},opacity = 0.5] (3,-2)--(3,12);
	\draw[color ={black},opacity = 0.5] (4,-2)--(4,12);
	\draw[color ={black},opacity = 0.5] (5,-2)--(5,12);
	\draw[color ={black},opacity = 0.5] (6,-2)--(6,12);
	\draw[color ={black},opacity = 0.5] (7,-2)--(7,12);
	\draw[color ={black},opacity = 0.5] (8,-2)--(8,12);
	\draw[color ={black},opacity = 0.5] (9,-2)--(9,12);
	\draw[color ={black},opacity = 0.5] (10,-2)--(10,12);
	\draw[color ={black},opacity = 0.5] (11,-2)--(11,12);
	\draw[color ={black},opacity = 0.5] (12,-2)--(12,12);
	\draw[color ={black},opacity = 0.5] (13,-2)--(13,12);
	\draw[color ={black},opacity = 0.5] (14,-2)--(14,12);
	\draw[color ={black},opacity = 0.5] (-1,-2)--(-1,12);
	\draw[color ={black},opacity = 0.5] (-2,-2)--(-2,12);
	\draw[color ={gray},opacity = 0.3] (-2,0)--(14,0);
	\draw[color ={gray},opacity = 0.3] (-2,0.5)--(14,0.5);
	\draw[color ={gray},opacity = 0.3] (-2,1.5)--(14,1.5);
	\draw[color ={gray},opacity = 0.3] (-2,2.5)--(14,2.5);
	\draw[color ={gray},opacity = 0.3] (-2,3.5)--(14,3.5);
	\draw[color ={gray},opacity = 0.3] (-2,4.5)--(14,4.5);
	\draw[color ={gray},opacity = 0.3] (-2,5.5)--(14,5.5);
	\draw[color ={gray},opacity = 0.3] (-2,6.5)--(14,6.5);
	\draw[color ={gray},opacity = 0.3] (-2,7.5)--(14,7.5);
	\draw[color ={gray},opacity = 0.3] (-2,8.5)--(14,8.5);
	\draw[color ={gray},opacity = 0.3] (-2,9.5)--(14,9.5);
	\draw[color ={gray},opacity = 0.3] (-2,10.5)--(14,10.5);
	\draw[color ={gray},opacity = 0.3] (-2,11.5)--(14,11.5);
	\draw[color ={gray},opacity = 0.3] (-2,0)--(14,0);
	\draw[color ={gray},opacity = 0.3] (-2,-0.5)--(14,-0.5);
	\draw[color ={gray},opacity = 0.3] (-2,-1.5)--(14,-1.5);
	\draw[color ={gray},opacity = 0.3] (0,-2)--(0,12);
	\draw[color ={gray},opacity = 0.3] (0.1,-2)--(0.1,12);
	\draw[color ={gray},opacity = 0.3] (0.2,-2)--(0.2,12);
	\draw[color ={gray},opacity = 0.3] (0.30000000000000004,-2)--(0.30000000000000004,12);
	\draw[color ={gray},opacity = 0.3] (0.4,-2)--(0.4,12);
	\draw[color ={gray},opacity = 0.3] (0.5,-2)--(0.5,12);
	\draw[color ={gray},opacity = 0.3] (0.6,-2)--(0.6,12);
	\draw[color ={gray},opacity = 0.3] (0.7,-2)--(0.7,12);
	\draw[color ={gray},opacity = 0.3] (0.7999999999999999,-2)--(0.7999999999999999,12);
	\draw[color ={gray},opacity = 0.3] (0.8999999999999999,-2)--(0.8999999999999999,12);
	\draw[color ={gray},opacity = 0.3] (1.0999999999999999,-2)--(1.0999999999999999,12);
	\draw[color ={gray},opacity = 0.3] (1.2,-2)--(1.2,12);
	\draw[color ={gray},opacity = 0.3] (1.3,-2)--(1.3,12);
	\draw[color ={gray},opacity = 0.3] (1.4000000000000001,-2)--(1.4000000000000001,12);
	\draw[color ={gray},opacity = 0.3] (1.5000000000000002,-2)--(1.5000000000000002,12);
	\draw[color ={gray},opacity = 0.3] (1.6000000000000003,-2)--(1.6000000000000003,12);
	\draw[color ={gray},opacity = 0.3] (1.7000000000000004,-2)--(1.7000000000000004,12);
	\draw[color ={gray},opacity = 0.3] (1.8000000000000005,-2)--(1.8000000000000005,12);
	\draw[color ={gray},opacity = 0.3] (1.9000000000000006,-2)--(1.9000000000000006,12);
	\draw[color ={gray},opacity = 0.3] (2.1000000000000005,-2)--(2.1000000000000005,12);
	\draw[color ={gray},opacity = 0.3] (2.2000000000000006,-2)--(2.2000000000000006,12);
	\draw[color ={gray},opacity = 0.3] (2.3000000000000007,-2)--(2.3000000000000007,12);
	\draw[color ={gray},opacity = 0.3] (2.400000000000001,-2)--(2.400000000000001,12);
	\draw[color ={gray},opacity = 0.3] (2.500000000000001,-2)--(2.500000000000001,12);
	\draw[color ={gray},opacity = 0.3] (2.600000000000001,-2)--(2.600000000000001,12);
	\draw[color ={gray},opacity = 0.3] (2.700000000000001,-2)--(2.700000000000001,12);
	\draw[color ={gray},opacity = 0.3] (2.800000000000001,-2)--(2.800000000000001,12);
	\draw[color ={gray},opacity = 0.3] (2.9000000000000012,-2)--(2.9000000000000012,12);
	\draw[color ={gray},opacity = 0.3] (3.1000000000000014,-2)--(3.1000000000000014,12);
	\draw[color ={gray},opacity = 0.3] (3.2000000000000015,-2)--(3.2000000000000015,12);
	\draw[color ={gray},opacity = 0.3] (3.3000000000000016,-2)--(3.3000000000000016,12);
	\draw[color ={gray},opacity = 0.3] (3.4000000000000017,-2)--(3.4000000000000017,12);
	\draw[color ={gray},opacity = 0.3] (3.5000000000000018,-2)--(3.5000000000000018,12);
	\draw[color ={gray},opacity = 0.3] (3.600000000000002,-2)--(3.600000000000002,12);
	\draw[color ={gray},opacity = 0.3] (3.700000000000002,-2)--(3.700000000000002,12);
	\draw[color ={gray},opacity = 0.3] (3.800000000000002,-2)--(3.800000000000002,12);
	\draw[color ={gray},opacity = 0.3] (3.900000000000002,-2)--(3.900000000000002,12);
	\draw[color ={gray},opacity = 0.3] (4.100000000000001,-2)--(4.100000000000001,12);
	\draw[color ={gray},opacity = 0.3] (4.200000000000001,-2)--(4.200000000000001,12);
	\draw[color ={gray},opacity = 0.3] (4.300000000000001,-2)--(4.300000000000001,12);
	\draw[color ={gray},opacity = 0.3] (4.4,-2)--(4.4,12);
	\draw[color ={gray},opacity = 0.3] (4.5,-2)--(4.5,12);
	\draw[color ={gray},opacity = 0.3] (4.6,-2)--(4.6,12);
	\draw[color ={gray},opacity = 0.3] (4.699999999999999,-2)--(4.699999999999999,12);
	\draw[color ={gray},opacity = 0.3] (4.799999999999999,-2)--(4.799999999999999,12);
	\draw[color ={gray},opacity = 0.3] (4.899999999999999,-2)--(4.899999999999999,12);
	\draw[color ={gray},opacity = 0.3] (5.099999999999998,-2)--(5.099999999999998,12);
	\draw[color ={gray},opacity = 0.3] (5.1999999999999975,-2)--(5.1999999999999975,12);
	\draw[color ={gray},opacity = 0.3] (5.299999999999997,-2)--(5.299999999999997,12);
	\draw[color ={gray},opacity = 0.3] (5.399999999999997,-2)--(5.399999999999997,12);
	\draw[color ={gray},opacity = 0.3] (5.4999999999999964,-2)--(5.4999999999999964,12);
	\draw[color ={gray},opacity = 0.3] (5.599999999999996,-2)--(5.599999999999996,12);
	\draw[color ={gray},opacity = 0.3] (5.699999999999996,-2)--(5.699999999999996,12);
	\draw[color ={gray},opacity = 0.3] (5.799999999999995,-2)--(5.799999999999995,12);
	\draw[color ={gray},opacity = 0.3] (5.899999999999995,-2)--(5.899999999999995,12);
	\draw[color ={gray},opacity = 0.3] (6.099999999999994,-2)--(6.099999999999994,12);
	\draw[color ={gray},opacity = 0.3] (6.199999999999994,-2)--(6.199999999999994,12);
	\draw[color ={gray},opacity = 0.3] (6.299999999999994,-2)--(6.299999999999994,12);
	\draw[color ={gray},opacity = 0.3] (6.399999999999993,-2)--(6.399999999999993,12);
	\draw[color ={gray},opacity = 0.3] (6.499999999999993,-2)--(6.499999999999993,12);
	\draw[color ={gray},opacity = 0.3] (6.5999999999999925,-2)--(6.5999999999999925,12);
	\draw[color ={gray},opacity = 0.3] (6.699999999999992,-2)--(6.699999999999992,12);
	\draw[color ={gray},opacity = 0.3] (6.799999999999992,-2)--(6.799999999999992,12);
	\draw[color ={gray},opacity = 0.3] (6.8999999999999915,-2)--(6.8999999999999915,12);
	\draw[color ={gray},opacity = 0.3] (7.099999999999991,-2)--(7.099999999999991,12);
	\draw[color ={gray},opacity = 0.3] (7.19999999999999,-2)--(7.19999999999999,12);
	\draw[color ={gray},opacity = 0.3] (7.29999999999999,-2)--(7.29999999999999,12);
	\draw[color ={gray},opacity = 0.3] (7.39999999999999,-2)--(7.39999999999999,12);
	\draw[color ={gray},opacity = 0.3] (7.499999999999989,-2)--(7.499999999999989,12);
	\draw[color ={gray},opacity = 0.3] (7.599999999999989,-2)--(7.599999999999989,12);
	\draw[color ={gray},opacity = 0.3] (7.699999999999989,-2)--(7.699999999999989,12);
	\draw[color ={gray},opacity = 0.3] (7.799999999999988,-2)--(7.799999999999988,12);
	\draw[color ={gray},opacity = 0.3] (7.899999999999988,-2)--(7.899999999999988,12);
	\draw[color ={gray},opacity = 0.3] (8.099999999999987,-2)--(8.099999999999987,12);
	\draw[color ={gray},opacity = 0.3] (8.199999999999987,-2)--(8.199999999999987,12);
	\draw[color ={gray},opacity = 0.3] (8.299999999999986,-2)--(8.299999999999986,12);
	\draw[color ={gray},opacity = 0.3] (8.399999999999986,-2)--(8.399999999999986,12);
	\draw[color ={gray},opacity = 0.3] (8.499999999999986,-2)--(8.499999999999986,12);
	\draw[color ={gray},opacity = 0.3] (8.599999999999985,-2)--(8.599999999999985,12);
	\draw[color ={gray},opacity = 0.3] (8.699999999999985,-2)--(8.699999999999985,12);
	\draw[color ={gray},opacity = 0.3] (8.799999999999985,-2)--(8.799999999999985,12);
	\draw[color ={gray},opacity = 0.3] (8.899999999999984,-2)--(8.899999999999984,12);
	\draw[color ={gray},opacity = 0.3] (9.099999999999984,-2)--(9.099999999999984,12);
	\draw[color ={gray},opacity = 0.3] (9.199999999999983,-2)--(9.199999999999983,12);
	\draw[color ={gray},opacity = 0.3] (9.299999999999983,-2)--(9.299999999999983,12);
	\draw[color ={gray},opacity = 0.3] (9.399999999999983,-2)--(9.399999999999983,12);
	\draw[color ={gray},opacity = 0.3] (9.499999999999982,-2)--(9.499999999999982,12);
	\draw[color ={gray},opacity = 0.3] (9.599999999999982,-2)--(9.599999999999982,12);
	\draw[color ={gray},opacity = 0.3] (9.699999999999982,-2)--(9.699999999999982,12);
	\draw[color ={gray},opacity = 0.3] (9.799999999999981,-2)--(9.799999999999981,12);
	\draw[color ={gray},opacity = 0.3] (9.89999999999998,-2)--(9.89999999999998,12);
	\draw[color ={gray},opacity = 0.3] (10.09999999999998,-2)--(10.09999999999998,12);
	\draw[color ={gray},opacity = 0.3] (10.19999999999998,-2)--(10.19999999999998,12);
	\draw[color ={gray},opacity = 0.3] (10.29999999999998,-2)--(10.29999999999998,12);
	\draw[color ={gray},opacity = 0.3] (10.399999999999979,-2)--(10.399999999999979,12);
	\draw[color ={gray},opacity = 0.3] (10.499999999999979,-2)--(10.499999999999979,12);
	\draw[color ={gray},opacity = 0.3] (10.599999999999978,-2)--(10.599999999999978,12);
	\draw[color ={gray},opacity = 0.3] (10.699999999999978,-2)--(10.699999999999978,12);
	\draw[color ={gray},opacity = 0.3] (10.799999999999978,-2)--(10.799999999999978,12);
	\draw[color ={gray},opacity = 0.3] (10.899999999999977,-2)--(10.899999999999977,12);
	\draw[color ={gray},opacity = 0.3] (11.099999999999977,-2)--(11.099999999999977,12);
	\draw[color ={gray},opacity = 0.3] (11.199999999999976,-2)--(11.199999999999976,12);
	\draw[color ={gray},opacity = 0.3] (11.299999999999976,-2)--(11.299999999999976,12);
	\draw[color ={gray},opacity = 0.3] (11.399999999999975,-2)--(11.399999999999975,12);
	\draw[color ={gray},opacity = 0.3] (11.499999999999975,-2)--(11.499999999999975,12);
	\draw[color ={gray},opacity = 0.3] (11.599999999999975,-2)--(11.599999999999975,12);
	\draw[color ={gray},opacity = 0.3] (11.699999999999974,-2)--(11.699999999999974,12);
	\draw[color ={gray},opacity = 0.3] (11.799999999999974,-2)--(11.799999999999974,12);
	\draw[color ={gray},opacity = 0.3] (11.899999999999974,-2)--(11.899999999999974,12);
	\draw[color ={gray},opacity = 0.3] (12.099999999999973,-2)--(12.099999999999973,12);
	\draw[color ={gray},opacity = 0.3] (12.199999999999973,-2)--(12.199999999999973,12);
	\draw[color ={gray},opacity = 0.3] (12.299999999999972,-2)--(12.299999999999972,12);
	\draw[color ={gray},opacity = 0.3] (12.399999999999972,-2)--(12.399999999999972,12);
	\draw[color ={gray},opacity = 0.3] (12.499999999999972,-2)--(12.499999999999972,12);
	\draw[color ={gray},opacity = 0.3] (12.599999999999971,-2)--(12.599999999999971,12);
	\draw[color ={gray},opacity = 0.3] (12.69999999999997,-2)--(12.69999999999997,12);
	\draw[color ={gray},opacity = 0.3] (12.79999999999997,-2)--(12.79999999999997,12);
	\draw[color ={gray},opacity = 0.3] (12.89999999999997,-2)--(12.89999999999997,12);
	\draw[color ={gray},opacity = 0.3] (13.09999999999997,-2)--(13.09999999999997,12);
	\draw[color ={gray},opacity = 0.3] (13.199999999999969,-2)--(13.199999999999969,12);
	\draw[color ={gray},opacity = 0.3] (13.299999999999969,-2)--(13.299999999999969,12);
	\draw[color ={gray},opacity = 0.3] (13.399999999999968,-2)--(13.399999999999968,12);
	\draw[color ={gray},opacity = 0.3] (13.499999999999968,-2)--(13.499999999999968,12);
	\draw[color ={gray},opacity = 0.3] (13.599999999999968,-2)--(13.599999999999968,12);
	\draw[color ={gray},opacity = 0.3] (13.699999999999967,-2)--(13.699999999999967,12);
	\draw[color ={gray},opacity = 0.3] (13.799999999999967,-2)--(13.799999999999967,12);
	\draw[color ={gray},opacity = 0.3] (13.899999999999967,-2)--(13.899999999999967,12);
	\draw[color ={gray},opacity = 0.3] (0,-2)--(0,12);
	\draw[color ={gray},opacity = 0.3] (-0.1,-2)--(-0.1,12);
	\draw[color ={gray},opacity = 0.3] (-0.2,-2)--(-0.2,12);
	\draw[color ={gray},opacity = 0.3] (-0.30000000000000004,-2)--(-0.30000000000000004,12);
	\draw[color ={gray},opacity = 0.3] (-0.4,-2)--(-0.4,12);
	\draw[color ={gray},opacity = 0.3] (-0.5,-2)--(-0.5,12);
	\draw[color ={gray},opacity = 0.3] (-0.6,-2)--(-0.6,12);
	\draw[color ={gray},opacity = 0.3] (-0.7,-2)--(-0.7,12);
	\draw[color ={gray},opacity = 0.3] (-0.7999999999999999,-2)--(-0.7999999999999999,12);
	\draw[color ={gray},opacity = 0.3] (-0.8999999999999999,-2)--(-0.8999999999999999,12);
	\draw[color ={gray},opacity = 0.3] (-1.0999999999999999,-2)--(-1.0999999999999999,12);
	\draw[color ={gray},opacity = 0.3] (-1.2,-2)--(-1.2,12);
	\draw[color ={gray},opacity = 0.3] (-1.3,-2)--(-1.3,12);
	\draw[color ={gray},opacity = 0.3] (-1.4000000000000001,-2)--(-1.4000000000000001,12);
	\draw[color ={gray},opacity = 0.3] (-1.5000000000000002,-2)--(-1.5000000000000002,12);
	\draw[color ={gray},opacity = 0.3] (-1.6000000000000003,-2)--(-1.6000000000000003,12);
	\draw[color ={gray},opacity = 0.3] (-1.7000000000000004,-2)--(-1.7000000000000004,12);
	\draw[color ={gray},opacity = 0.3] (-1.8000000000000005,-2)--(-1.8000000000000005,12);
	\draw[color ={gray},opacity = 0.3] (-1.9000000000000006,-2)--(-1.9000000000000006,12);
	\draw[color ={black},line width = 2] (2,-0.2)--(2,0.2);
	\draw[color ={black},line width = 2] (4,-0.2)--(4,0.2);
	\draw[color ={black},line width = 2] (6,-0.2)--(6,0.2);
	\draw[color ={black},line width = 2] (8,-0.2)--(8,0.2);
	\draw[color ={black},line width = 2] (10,-0.2)--(10,0.2);
	\draw[color ={black},line width = 2] (12,-0.2)--(12,0.2);
	\draw[color ={black},line width = 2] (-0.2,2)--(0.2,2);
	\draw[color ={black},line width = 2] (-0.2,4)--(0.2,4);
	\draw[color ={black},line width = 2] (-0.2,6)--(0.2,6);
	\draw[color ={black},line width = 2] (-0.2,8)--(0.2,8);
	\draw[color ={black},line width = 2] (-0.2,10)--(0.2,10);
	\draw (2,-0.4) node[anchor = center] {\scriptsize \color{black}{$1$}};
	\draw (4,-0.4) node[anchor = center] {\scriptsize \color{black}{$2$}};
	\draw (6,-0.4) node[anchor = center] {\scriptsize \color{black}{$3$}};
	\draw (8,-0.4) node[anchor = center] {\scriptsize \color{black}{$4$}};
	\draw (10,-0.4) node[anchor = center] {\scriptsize \color{black}{$5$}};
	\draw (-0.5,2.1) node[anchor = center] {\footnotesize \color{black}{$1$}};
	\draw (-0.5,4.1) node[anchor = center] {\footnotesize \color{black}{$2$}};
	\draw (-0.5,6.1) node[anchor = center] {\footnotesize \color{black}{$3$}};
	\draw (-0.5,8.1) node[anchor = center] {\footnotesize \color{black}{$4$}};
	\draw (-0.5,10.1) node[anchor = center] {\footnotesize \color{black}{$5$}};
	\draw [color={black}] (-0.3,-0.3) node[anchor = center,scale=1, rotate = 0] {O};
	
	\draw[color={blue},line width = 2] (1,10)--(1.4,7.71)--(1.8,6.44)--(2.2,5.64)--(2.6,5.08)--(3,4.67)--(3.4,4.35)--(3.8,4.11)--(4.2,3.9)--(4.6,3.74)--(5,3.6)--(5.4,3.48)--(5.8,3.38)--(6.2,3.29)--(6.6,3.21)--(7,3.14)--(7.4,3.08)--(7.8,3.03)--(8.2,2.98)--(8.6,2.93)--(9,2.89)--(9.4,2.85)--(9.8,2.82)--(10.2,2.78)--(10.6,2.75)--(11,2.73)--(11.4,2.7)--(11.8,2.68)--(12.2,2.66)--(12.6,2.63)--(13,2.62)--(13.4,2.6)--(13.8,2.58);
	
	\draw[color={red},line width = 2] (-2,7)--(-1.6,6.8)--(-1.2,6.6)--(-0.8,6.4)--(-0.4,6.2)--(0,6)--(0.4,5.8)--(0.8,5.6)--(1.2,5.4)--(1.6,5.2)--(2,5)--(2.4,4.8)--(2.8,4.6)--(3.2,4.4)--(3.6,4.2)--(4,4)--(4.4,3.8)--(4.8,3.6)--(5.2,3.4)--(5.6,3.2)--(6,3)--(6.4,2.8)--(6.8,2.6)--(7.2,2.4)--(7.6,2.2)--(8,2)--(8.4,1.8)--(8.8,1.6)--(9.2,1.4)--(9.6,1.2)--(10,1)--(10.4,0.8)--(10.8,0.6)--(11.2,0.4)--(11.6,0.2)--(12,0)--(12.4,-0.2)--(12.8,-0.4)--(13.2,-0.6)--(13.6,-0.8)--(14,-1);

\end{tikzpicture}\end{center}
\end{exercice}

\begin{Solution}
 $f'(2)$ est donné par le coefficient directeur de la tangente à la courbe au point d'abscisse $2$, soit $\dfrac{-2}{4}=-\dfrac{1{\color[HTML]{2563a5}\boldsymbol{\times2}} }{2{\color[HTML]{2563a5}\boldsymbol{\times2}}}=-\dfrac{1}{2}$.
\end{Solution}

% @see : https://coopmaths.fr/alea?uuid=6f32d&id=can1F16&n=1&d=10&alea=dqSw2232323232342345678910111213141516171819202122232425262728293031323334353637383940414243444546474823456789101112131415161718192021222324&cd=1&cols=1
\needspace{10\baselineskip}
\newpage
\begin{exercice}[BaremeTotal=false]


 On donne les représentations graphiques d'une fonction et de sa dérivée.\\
      Donner l'équation réduite de la tangente à la courbe de $f$ en $x=-1$. \\ \begin{minipage}{0.5\linewidth}
    \begin{tikzpicture}[baseline, scale=0.8]

    \tikzset{
      point/.style={
        thick,
        draw,
        cross out,
        inner sep=0pt,
        minimum width=5pt,
        minimum height=5pt,
      },
    }
    \clip (-5.5,-9.5) rectangle (5.75,7.5);
    	
	\draw[color ={black},>=latex,->] (-4,0)--(4,0);
	\draw[color ={black},>=latex,->] (0,-8)--(0,6);
	\draw[color ={black},opacity = 0.4] (-4,1)--(4,1);
	\draw[color ={black},opacity = 0.4] (-4,2)--(4,2);
	\draw[color ={black},opacity = 0.4] (-4,3)--(4,3);
	\draw[color ={black},opacity = 0.4] (-4,4)--(4,4);
	\draw[color ={black},opacity = 0.4] (-4,5)--(4,5);
	\draw[color ={black},opacity = 0.4] (-4,6)--(4,6);
	\draw[color ={black},opacity = 0.4] (-4,-1)--(4,-1);
	\draw[color ={black},opacity = 0.4] (-4,-2)--(4,-2);
	\draw[color ={black},opacity = 0.4] (-4,-3)--(4,-3);
	\draw[color ={black},opacity = 0.4] (-4,-4)--(4,-4);
	\draw[color ={black},opacity = 0.4] (-4,-5)--(4,-5);
	\draw[color ={black},opacity = 0.4] (-4,-6)--(4,-6);
	\draw[color ={black},opacity = 0.4] (-4,-7)--(4,-7);
	\draw[color ={black},opacity = 0.4] (-4,-8)--(4,-8);
	\draw[color ={black},opacity = 0.4] (1,-8)--(1,6);
	\draw[color ={black},opacity = 0.4] (2,-8)--(2,6);
	\draw[color ={black},opacity = 0.4] (3,-8)--(3,6);
	\draw[color ={black},opacity = 0.4] (4,-8)--(4,6);
	\draw[color ={black},opacity = 0.4] (-1,-8)--(-1,6);
	\draw[color ={black},opacity = 0.4] (-2,-8)--(-2,6);
	\draw[color ={black},opacity = 0.4] (-3,-8)--(-3,6);
	\draw[color ={black},opacity = 0.4] (-4,-8)--(-4,6);
	\draw[color ={gray},opacity = 0.3] (-4,0)--(4,0);
	\draw[color ={gray},opacity = 0.3] (-4,0.5)--(4,0.5);
	\draw[color ={gray},opacity = 0.3] (-4,1.5)--(4,1.5);
	\draw[color ={gray},opacity = 0.3] (-4,2.5)--(4,2.5);
	\draw[color ={gray},opacity = 0.3] (-4,3.5)--(4,3.5);
	\draw[color ={gray},opacity = 0.3] (-4,4.5)--(4,4.5);
	\draw[color ={gray},opacity = 0.3] (-4,5.5)--(4,5.5);
	\draw[color ={gray},opacity = 0.3] (-4,0)--(4,0);
	\draw[color ={gray},opacity = 0.3] (-4,-0.5)--(4,-0.5);
	\draw[color ={gray},opacity = 0.3] (-4,-1.5)--(4,-1.5);
	\draw[color ={gray},opacity = 0.3] (-4,-2.5)--(4,-2.5);
	\draw[color ={gray},opacity = 0.3] (-4,-3.5)--(4,-3.5);
	\draw[color ={gray},opacity = 0.3] (-4,-4.5)--(4,-4.5);
	\draw[color ={gray},opacity = 0.3] (-4,-5.5)--(4,-5.5);
	\draw[color ={gray},opacity = 0.3] (-4,-6.5)--(4,-6.5);
	\draw[color ={gray},opacity = 0.3] (-4,-7)--(4,-7);
	\draw[color ={gray},opacity = 0.3] (-4,-7.5)--(4,-7.5);
	\draw[color ={gray},opacity = 0.3] (-4,-8)--(4,-8);
	\draw[color ={gray},opacity = 0.3] (0,-8)--(0,6);
	\draw[color ={gray},opacity = 0.3] (0,-8)--(0,6);
	\draw[color ={black},line width = 1.2] (1,-0.13)--(1,0.13);
	\draw[color ={black},line width = 1.2] (2,-0.13)--(2,0.13);
	\draw[color ={black},line width = 1.2] (3,-0.13)--(3,0.13);
	\draw[color ={black},line width = 1.2] (4,-0.13)--(4,0.13);
	\draw[color ={black},line width = 1.2] (-1,-0.13)--(-1,0.13);
	\draw[color ={black},line width = 1.2] (-2,-0.13)--(-2,0.13);
	\draw[color ={black},line width = 1.2] (-3,-0.13)--(-3,0.13);
	\draw[color ={black},line width = 1.2] (-4,-0.13)--(-4,0.13);
	\draw[color ={black},line width = 1.2] (-0.13,1)--(0.13,1);
	\draw[color ={black},line width = 1.2] (-0.13,2)--(0.13,2);
	\draw[color ={black},line width = 1.2] (-0.13,3)--(0.13,3);
	\draw[color ={black},line width = 1.2] (-0.13,4)--(0.13,4);
	\draw[color ={black},line width = 1.2] (-0.13,5)--(0.13,5);
	\draw[color ={black},line width = 1.2] (-0.13,6)--(0.13,6);
	\draw[color ={black},line width = 1.2] (-0.13,-1)--(0.13,-1);
	\draw[color ={black},line width = 1.2] (-0.13,-2)--(0.13,-2);
	\draw[color ={black},line width = 1.2] (-0.13,-3)--(0.13,-3);
	\draw[color ={black},line width = 1.2] (-0.13,-4)--(0.13,-4);
	\draw[color ={black},line width = 1.2] (-0.13,-5)--(0.13,-5);
	\draw[color ={black},line width = 1.2] (-0.13,-6)--(0.13,-6);
	\draw[color ={black},line width = 1.2] (-0.13,-7)--(0.13,-7);
	\draw[color ={black},line width = 1.2] (-0.13,-8)--(0.13,-8);
	\draw (2,-0.4) node[anchor = center] {\scriptsize \color{black}{$2$}};
	\draw (4,-0.4) node[anchor = center] {\scriptsize \color{black}{$4$}};
	\draw (-2,-0.4) node[anchor = center] {\scriptsize \color{black}{$-2$}};
	\draw (-4,-0.4) node[anchor = center] {\scriptsize \color{black}{$-4$}};
	\draw (-0.5,2.1) node[anchor = center] {\footnotesize \color{black}{$2$}};
	\draw (-0.5,4.1) node[anchor = center] {\footnotesize \color{black}{$4$}};
	\draw (-0.5,6.1) node[anchor = center] {\footnotesize \color{black}{$6$}};
	\draw (-0.5,-1.9) node[anchor = center] {\footnotesize \color{black}{$-2$}};
	\draw (-0.5,-3.9) node[anchor = center] {\footnotesize \color{black}{$-4$}};
	\draw (-0.5,-5.9) node[anchor = center] {\footnotesize \color{black}{$-6$}};
	\draw (-0.5,-7.9) node[anchor = center] {\footnotesize \color{black}{$-8$}};
	\draw [color={black}] (-0.3,-0.3) node[anchor = center,scale=1, rotate = 0] {O};
	\draw (3,4) node[anchor = center] {\footnotesize \color{blue}{$\Large \mathcal{C}_f$}};
	
	\draw[color={blue},line width = 2] (-4,-2)--(-3.8,-1.24)--(-3.6,-0.56)--(-3.4,0.04)--(-3.2,0.56)--(-3,1)--(-2.8,1.36)--(-2.6,1.64)--(-2.4,1.84)--(-2.2,1.96)--(-2,2)--(-1.8,1.96)--(-1.6,1.84)--(-1.4,1.64)--(-1.2,1.36)--(-1,1)--(-0.8,0.56)--(-0.6,0.04)--(-0.4,-0.56)--(-0.2,-1.24)--(0,-2)--(0.2,-2.84)--(0.4,-3.76)--(0.6,-4.76)--(0.8,-5.84)--(1,-7)--(1.2,-8.24);

\end{tikzpicture}\\
    \end{minipage}
    \begin{minipage}{0.5\linewidth}
    \begin{tikzpicture}[baseline, scale=0.8]

    \tikzset{
      point/.style={
        thick,
        draw,
        cross out,
        inner sep=0pt,
        minimum width=5pt,
        minimum height=5pt,
      },
    }
    \clip (-5.5,-9.5) rectangle (6.5,7.5);
    	
	\draw[color ={black},>=latex,->] (-4,0)--(4,0);
	\draw[color ={black},>=latex,->] (0,-8)--(0,6);
	\draw[color ={black},opacity = 0.4] (-4,1)--(4,1);
	\draw[color ={black},opacity = 0.4] (-4,2)--(4,2);
	\draw[color ={black},opacity = 0.4] (-4,3)--(4,3);
	\draw[color ={black},opacity = 0.4] (-4,4)--(4,4);
	\draw[color ={black},opacity = 0.4] (-4,5)--(4,5);
	\draw[color ={black},opacity = 0.4] (-4,6)--(4,6);
	\draw[color ={black},opacity = 0.4] (-4,-1)--(4,-1);
	\draw[color ={black},opacity = 0.4] (-4,-2)--(4,-2);
	\draw[color ={black},opacity = 0.4] (-4,-3)--(4,-3);
	\draw[color ={black},opacity = 0.4] (-4,-4)--(4,-4);
	\draw[color ={black},opacity = 0.4] (-4,-5)--(4,-5);
	\draw[color ={black},opacity = 0.4] (-4,-6)--(4,-6);
	\draw[color ={black},opacity = 0.4] (-4,-7)--(4,-7);
	\draw[color ={black},opacity = 0.4] (-4,-8)--(4,-8);
	\draw[color ={black},opacity = 0.4] (1,-8)--(1,6);
	\draw[color ={black},opacity = 0.4] (2,-8)--(2,6);
	\draw[color ={black},opacity = 0.4] (3,-8)--(3,6);
	\draw[color ={black},opacity = 0.4] (4,-8)--(4,6);
	\draw[color ={black},opacity = 0.4] (-1,-8)--(-1,6);
	\draw[color ={black},opacity = 0.4] (-2,-8)--(-2,6);
	\draw[color ={black},opacity = 0.4] (-3,-8)--(-3,6);
	\draw[color ={black},opacity = 0.4] (-4,-8)--(-4,6);
	\draw[color ={gray},opacity = 0.3] (-4,0)--(4,0);
	\draw[color ={gray},opacity = 0.3] (-4,0.5)--(4,0.5);
	\draw[color ={gray},opacity = 0.3] (-4,1.5)--(4,1.5);
	\draw[color ={gray},opacity = 0.3] (-4,2.5)--(4,2.5);
	\draw[color ={gray},opacity = 0.3] (-4,3.5)--(4,3.5);
	\draw[color ={gray},opacity = 0.3] (-4,4.5)--(4,4.5);
	\draw[color ={gray},opacity = 0.3] (-4,5.5)--(4,5.5);
	\draw[color ={gray},opacity = 0.3] (-4,0)--(4,0);
	\draw[color ={gray},opacity = 0.3] (-4,-0.5)--(4,-0.5);
	\draw[color ={gray},opacity = 0.3] (-4,-1.5)--(4,-1.5);
	\draw[color ={gray},opacity = 0.3] (-4,-2.5)--(4,-2.5);
	\draw[color ={gray},opacity = 0.3] (-4,-3.5)--(4,-3.5);
	\draw[color ={gray},opacity = 0.3] (-4,-4.5)--(4,-4.5);
	\draw[color ={gray},opacity = 0.3] (-4,-5.5)--(4,-5.5);
	\draw[color ={gray},opacity = 0.3] (-4,-6.5)--(4,-6.5);
	\draw[color ={gray},opacity = 0.3] (-4,-7)--(4,-7);
	\draw[color ={gray},opacity = 0.3] (-4,-7.5)--(4,-7.5);
	\draw[color ={gray},opacity = 0.3] (-4,-8)--(4,-8);
	\draw[color ={gray},opacity = 0.3] (0,-8)--(0,6);
	\draw[color ={gray},opacity = 0.3] (0,-8)--(0,6);
	\draw[color ={black},line width = 1.2] (1,-0.13)--(1,0.13);
	\draw[color ={black},line width = 1.2] (2,-0.13)--(2,0.13);
	\draw[color ={black},line width = 1.2] (3,-0.13)--(3,0.13);
	\draw[color ={black},line width = 1.2] (4,-0.13)--(4,0.13);
	\draw[color ={black},line width = 1.2] (-1,-0.13)--(-1,0.13);
	\draw[color ={black},line width = 1.2] (-2,-0.13)--(-2,0.13);
	\draw[color ={black},line width = 1.2] (-3,-0.13)--(-3,0.13);
	\draw[color ={black},line width = 1.2] (-4,-0.13)--(-4,0.13);
	\draw[color ={black},line width = 1.2] (-0.13,1)--(0.13,1);
	\draw[color ={black},line width = 1.2] (-0.13,2)--(0.13,2);
	\draw[color ={black},line width = 1.2] (-0.13,3)--(0.13,3);
	\draw[color ={black},line width = 1.2] (-0.13,4)--(0.13,4);
	\draw[color ={black},line width = 1.2] (-0.13,5)--(0.13,5);
	\draw[color ={black},line width = 1.2] (-0.13,6)--(0.13,6);
	\draw[color ={black},line width = 1.2] (-0.13,-1)--(0.13,-1);
	\draw[color ={black},line width = 1.2] (-0.13,-2)--(0.13,-2);
	\draw[color ={black},line width = 1.2] (-0.13,-3)--(0.13,-3);
	\draw[color ={black},line width = 1.2] (-0.13,-4)--(0.13,-4);
	\draw[color ={black},line width = 1.2] (-0.13,-5)--(0.13,-5);
	\draw[color ={black},line width = 1.2] (-0.13,-6)--(0.13,-6);
	\draw[color ={black},line width = 1.2] (-0.13,-7)--(0.13,-7);
	\draw[color ={black},line width = 1.2] (-0.13,-8)--(0.13,-8);
	\draw (2,-0.4) node[anchor = center] {\scriptsize \color{black}{$2$}};
	\draw (4,-0.4) node[anchor = center] {\scriptsize \color{black}{$4$}};
	\draw (-2,-0.4) node[anchor = center] {\scriptsize \color{black}{$-2$}};
	\draw (-4,-0.4) node[anchor = center] {\scriptsize \color{black}{$-4$}};
	\draw (-0.5,2.1) node[anchor = center] {\footnotesize \color{black}{$2$}};
	\draw (-0.5,4.1) node[anchor = center] {\footnotesize \color{black}{$4$}};
	\draw (-0.5,-1.9) node[anchor = center] {\footnotesize \color{black}{$-2$}};
	\draw (-0.5,-3.9) node[anchor = center] {\footnotesize \color{black}{$-4$}};
	\draw (-0.5,-5.9) node[anchor = center] {\footnotesize \color{black}{$-6$}};
	\draw (-0.5,-7.9) node[anchor = center] {\footnotesize \color{black}{$-8$}};
	\draw [color={black}] (-0.3,-0.3) node[anchor = center,scale=1, rotate = 0] {O};
	\draw (3,4) node[anchor = center] {\footnotesize \color{red}{$\Large\mathcal{C}_f\prime$}};
	
	\draw[color={red},line width = 2] (-4,4)--(-3.8,3.6)--(-3.6,3.2)--(-3.4,2.8)--(-3.2,2.4)--(-3,2)--(-2.8,1.6)--(-2.6,1.2)--(-2.4,0.8)--(-2.2,0.4)--(-2,0)--(-1.8,-0.4)--(-1.6,-0.8)--(-1.4,-1.2)--(-1.2,-1.6)--(-1,-2)--(-0.8,-2.4)--(-0.6,-2.8)--(-0.4,-3.2)--(-0.2,-3.6)--(0,-4)--(0.2,-4.4)--(0.4,-4.8)--(0.6,-5.2)--(0.8,-5.6)--(1,-6)--(1.2,-6.4)--(1.4,-6.8)--(1.6,-7.2)--(1.8,-7.6)--(2,-8)--(2.2,-8.4)--(2.4,-8.8);

\end{tikzpicture}\\
    \end{minipage}
    
\end{exercice}

\begin{Solution}
 L'équation réduite de la tangente au point d'abscisse $-1$ est  : $y=f'(-1)(x-(-1))+f(-1)$.\\
      On lit graphiquement $f(-1)=1$ et $f'(-1)=-2$.\\
      L'équation réduite de la tangente est donc donnée par :
      $y=-2(x+1)+1$, soit $y=-2x-1$.
\end{Solution}

% @see : https://coopmaths.fr/alea?uuid=a1ba2&id=can1F14&n=1&d=10&alea=rkIp2232323232342345678910111213141516171819202122232425262728293031323334353637383940414243444546474823456789101112131415161718192021222324&cd=1&cols=1
\needspace{10\baselineskip}
\begin{exercice}[BaremeTotal=false]


 Soit $f$ la fonction définie sur $\mathbb{R}$ par : $f(x)= x^2-2x-2$.\\
        En calculant la limite du taux d'accroissement, déterminer $f'(3)$.
\end{exercice}

\begin{Solution}
 $f$ est une fonction polynôme du second degré de la forme $f(x)=ax^2+bx+c$.\\
    La fonction dérivée est donnée par la somme des dérivées des fonctions $u$ et $v$ définies par $u(x)=x^2$ et $v(x)=-2x-2$.\\
     Comme $u'(x)=2x$ et $v'(x)=-2$, on obtient  $f'(x)=2x-2$. \\
     Ainsi, $f'(3)=2\times 3-2=4$.
\end{Solution}

% @see : https://coopmaths.fr/alea?uuid=3c690&id=can1F13&n=1&d=10&alea=EUDu2232323232342345678910111213141516171819202122232425262728293031323334353637383940414243444546474823456789101112131415161718192021222324&cd=1&cols=1
\needspace{10\baselineskip}
\begin{exercice}[BaremeTotal=false]


 Déterminer le coefficient directeur de la tangente à la courbe représentative de la fonction inverse au point d'abscisse $-1$.
                 
\end{exercice}

\begin{Solution}
 Le coefficient directeur de la tangente au point d'abscisse $a$ est donné par le nombre dérivé $f'(a)$.\\
        La fonction $f$ définie par $f(x)=\dfrac{1}{x}$ a pour fonction dérivée la fonction $f'$ définie par $f'(x)=-\dfrac{1}{x^2}$.\\
Comme $f'(-1)=-\dfrac{1}{(-1)^2}=-\dfrac{1}{1}$, le coefficient directeur de la tangente au point d'abscisse $-1$ est : $-\dfrac{1}{1}$$=-1$.
\end{Solution}

% @see : https://coopmaths.fr/alea?uuid=ebd89&id=1AN14-1&n=1&d=10&s=3&alea=ocYr2232323232342345678910111213141516171819202122232425262728293031323334353637383940414243444546474823456789101112131415161718192021222324&cd=0&cols=1
\needspace{10\baselineskip}
\begin{exercice}[BaremeTotal=false]


 Donner la dérivée de la fonction $f$, dérivable pour tout $x\in\R$, définie par  $f(x)=\dfrac{-6x+5}{7}$.
\end{exercice}

\begin{Solution}
 L'expression de la dérivée de la fonction $f$ définie par $f(x)=\dfrac{-6x+5}{7}$ est : ${\color[HTML]{f15929}\boldsymbol{f'(x)=-\dfrac{6}{7}}}$.
\end{Solution}

% @see : https://coopmaths.fr/alea?uuid=ebd89&id=1AN14-1&n=1&d=10&s=4&alea=uAlP2232323232342345678910111213141516171819202122232425262728293031323334353637383940414243444546474823456789101112131415161718192021222324&cd=0&cols=1
\needspace{10\baselineskip}
\begin{exercice}[BaremeTotal=false]


 Donner la dérivée de la fonction $f$, dérivable pour tout $x\in\R$, définie par  $f(x)=-8x^{4}$.
\end{exercice}

\begin{Solution}
 L'expression de la dérivée de la fonction $f$ définie par $f(x)=-8x^{4}$ est : ${\color[HTML]{f15929}\boldsymbol{f'(x)=-32x^{3}}}$.
\end{Solution}

% @see : https://coopmaths.fr/alea?uuid=ebd89&id=1AN14-1&n=1&d=10&s=7&alea=vjN72232323232342345678910111213141516171819202122232425262728293031323334353637383940414243444546474823456789101112131415161718192021222324&cd=0&cols=1
\needspace{10\baselineskip}
\begin{exercice}[BaremeTotal=false]


 Donner la dérivée de la fonction $f$, dérivable pour tout $x\in\R^*$, définie par  $f(x)=\dfrac{-4}{9x}$.
\end{exercice}

\begin{Solution}
 L'expression de la dérivée de la fonction $f$ définie par $f(x)=\dfrac{-4}{9x}$ est : ${\color[HTML]{f15929}\boldsymbol{f'(x)=\dfrac{4}{9x^2}}}$.
\end{Solution}

% @see : https://coopmaths.fr/alea?uuid=a83c0&id=1AN14-4&n=1&d=10&s=1&alea=Njr32232323232342345678910111213141516171819202122232425262728293031323334353637383940414243444546474823456789101112131415161718192021222324&cd=0&cols=1
\needspace{10\baselineskip}
\begin{exercice}[BaremeTotal=false]


 Donner l'expression de la dérivée de la fonction $f$ définie sur $\R^{*}$ par $f(x)=-2x^{2}-\dfrac{9}{x}$.\\
\end{exercice}

\begin{Solution}
 L'expression de la dérivée de la fonction $f$ définie par $f(x)=-2x^{2}-\dfrac{9}{x}$ est :\\$f'(x)={\color[HTML]{f15929}\boldsymbol{-4x+\dfrac{9}{x^2}}}$.\\
\end{Solution}

% @see : https://coopmaths.fr/alea?uuid=a83c0&id=1AN14-4&n=1&d=10&s=4&alea=4Vpz2232323232342345678910111213141516171819202122232425262728293031323334353637383940414243444546474823456789101112131415161718192021222324&cd=1&cols=1
\needspace{10\baselineskip}
\begin{exercice}[BaremeTotal=false]


 Donner l'expression de la dérivée de la fonction $f$ définie sur $\R_+^{*}$ par $f(x)=-\sqrt{x}-\dfrac{5}{x}+4x^{2}+3x$.\\
\end{exercice}

\begin{Solution}
 $(-\sqrt{x})^\prime=-\dfrac{1}{2\sqrt{x}}$\\$(-\dfrac{5}{x})^\prime=\dfrac{5}{x^2}$\\$(4x^{2}+3x)^\prime=8x+3$\\L'expression de la dérivée de la fonction $f$ définie par $f(x)=-\sqrt{x}-\dfrac{5}{x}+4x^{2}+3x$ est :\\$f'(x)={\color[HTML]{f15929}\boldsymbol{-\dfrac{1}{2\sqrt{x}}+\dfrac{5}{x^2}+8x+3}}$.\\
\end{Solution}

\end{Maquette}
\clearpage
\begin{Maquette}[DM]{Niveau= ,Classe=\getdata[25]\mydata\ (v25), Date = 06/01/25  ,Theme=Exercices}


% @see : https://coopmaths.fr/alea?uuid=4c8c7&id=1AN11-3&n=1&d=10&s=2&alea=nawO223232323234234567891011121314151617181920212223242526272829303132333435363738394041424344454647482345678910111213141516171819202122232425&cd=1&cols=2
\needspace{10\baselineskip}
\begin{exercice}[BaremeTotal=false]
\label{eleve:25}


 \begin{minipage}[t]{\linewidth} Soit $f$ une fonction dérivable sur $[-5;5]$ et $\mathcal{C}_f$ sa courbe représentative.\\On sait que  $f(-5)=-1~~$ et que $~~f'(-5)=0$.\\Déterminer une équation de la tangente $(T)$ à la courbe $\mathcal{C}_f$ au point d'abscisse $-5$,\\sans utiliser la formule de cours de l'équation de tangente. \end{minipage}

\end{exercice}

\begin{Solution}
 \begin{minipage}[t]{\linewidth} On sait que la tangente n'est pas une droite verticale, puisque la fonction est dérivable sur l'intervalle.\\On en déduit que la tangente $(T)$ au point d'abscisse $-5$, admet une équation réduite de la forme :  $(T) : y=m x + p$. 

\medskip
$\bullet$ Détermination de $m$ :\\On sait que le nombre dérivé en $-5$ est par définition, le coefficient directeur de la tangente au point d'abscisse $-5$.\\Par conséquent, on a déjà : $m=f'(-5)=0$.\\ On en déduit que  $(T) : y= 0 x + p$\\$\bullet$ Détermination de $p$ :\\ Pour cela, on utilise que si $f(-5)=-1~~$, alors le point $A$ de coordonnées $(-5;-1)$ appartient à $\mathcal{C}_f$ mais aussi à $(T)$.\\ On peut écrire $A(-5;-1) = \mathcal{C}_f \cap (T)$.\\ On remplace alors les coordonnées de $A(-5;-1)$ dans l'équation  $(T) : y= 0 x + p$.\\ $\begin{aligned} \phantom{\iff}&A(-5;-1)\in (T)\\ \iff& -1= 0 \times (-5)  + p\\ \iff& p=-1 +0 \times (-5)  \\ \iff& p=-1 +0   \\ \iff& p=-1\\\end{aligned}$\\On peut conclure que : $(T) : {\color[HTML]{f15929}\boldsymbol{y=-1}}$. \end{minipage}

\end{Solution}

% @see : https://coopmaths.fr/alea?uuid=0e984&id=can1F15&n=1&d=10&alea=1Alo223232323234234567891011121314151617181920212223242526272829303132333435363738394041424344454647482345678910111213141516171819202122232425&cd=1&cols=1
\needspace{10\baselineskip}
\begin{exercice}[BaremeTotal=false]


 La courbe représente une fonction $f$ et la droite est la tangente au point d'abscisse $1$.\\
        Déterminer $f'(1)$.\begin{center}\begin{tikzpicture}[baseline,scale = 0.5]

    \tikzset{
      point/.style={
        thick,
        draw,
        cross out,
        inner sep=0pt,
        minimum width=5pt,
        minimum height=5pt,
      },
    }
    \clip (-2,-2) rectangle (14,12);
    	
	\draw[color ={black},line width = 2,>=latex,->] (-2,0)--(14,0);
	\draw[color ={black},line width = 2,>=latex,->] (0,-2)--(0,12);
	\draw[color ={black},opacity = 0.5] (-2,1)--(14,1);
	\draw[color ={black},opacity = 0.5] (-2,2)--(14,2);
	\draw[color ={black},opacity = 0.5] (-2,3)--(14,3);
	\draw[color ={black},opacity = 0.5] (-2,4)--(14,4);
	\draw[color ={black},opacity = 0.5] (-2,5)--(14,5);
	\draw[color ={black},opacity = 0.5] (-2,6)--(14,6);
	\draw[color ={black},opacity = 0.5] (-2,7)--(14,7);
	\draw[color ={black},opacity = 0.5] (-2,8)--(14,8);
	\draw[color ={black},opacity = 0.5] (-2,9)--(14,9);
	\draw[color ={black},opacity = 0.5] (-2,10)--(14,10);
	\draw[color ={black},opacity = 0.5] (-2,11)--(14,11);
	\draw[color ={black},opacity = 0.5] (-2,12)--(14,12);
	\draw[color ={black},opacity = 0.5] (-2,-1)--(14,-1);
	\draw[color ={black},opacity = 0.5] (-2,-2)--(14,-2);
	\draw[color ={black},opacity = 0.5] (1,-2)--(1,12);
	\draw[color ={black},opacity = 0.5] (2,-2)--(2,12);
	\draw[color ={black},opacity = 0.5] (3,-2)--(3,12);
	\draw[color ={black},opacity = 0.5] (4,-2)--(4,12);
	\draw[color ={black},opacity = 0.5] (5,-2)--(5,12);
	\draw[color ={black},opacity = 0.5] (6,-2)--(6,12);
	\draw[color ={black},opacity = 0.5] (7,-2)--(7,12);
	\draw[color ={black},opacity = 0.5] (8,-2)--(8,12);
	\draw[color ={black},opacity = 0.5] (9,-2)--(9,12);
	\draw[color ={black},opacity = 0.5] (10,-2)--(10,12);
	\draw[color ={black},opacity = 0.5] (11,-2)--(11,12);
	\draw[color ={black},opacity = 0.5] (12,-2)--(12,12);
	\draw[color ={black},opacity = 0.5] (13,-2)--(13,12);
	\draw[color ={black},opacity = 0.5] (14,-2)--(14,12);
	\draw[color ={black},opacity = 0.5] (-1,-2)--(-1,12);
	\draw[color ={black},opacity = 0.5] (-2,-2)--(-2,12);
	\draw[color ={gray},opacity = 0.3] (-2,0)--(14,0);
	\draw[color ={gray},opacity = 0.3] (-2,0.5)--(14,0.5);
	\draw[color ={gray},opacity = 0.3] (-2,1.5)--(14,1.5);
	\draw[color ={gray},opacity = 0.3] (-2,2.5)--(14,2.5);
	\draw[color ={gray},opacity = 0.3] (-2,3.5)--(14,3.5);
	\draw[color ={gray},opacity = 0.3] (-2,4.5)--(14,4.5);
	\draw[color ={gray},opacity = 0.3] (-2,5.5)--(14,5.5);
	\draw[color ={gray},opacity = 0.3] (-2,6.5)--(14,6.5);
	\draw[color ={gray},opacity = 0.3] (-2,7.5)--(14,7.5);
	\draw[color ={gray},opacity = 0.3] (-2,8.5)--(14,8.5);
	\draw[color ={gray},opacity = 0.3] (-2,9.5)--(14,9.5);
	\draw[color ={gray},opacity = 0.3] (-2,10.5)--(14,10.5);
	\draw[color ={gray},opacity = 0.3] (-2,11.5)--(14,11.5);
	\draw[color ={gray},opacity = 0.3] (-2,0)--(14,0);
	\draw[color ={gray},opacity = 0.3] (-2,-0.5)--(14,-0.5);
	\draw[color ={gray},opacity = 0.3] (-2,-1.5)--(14,-1.5);
	\draw[color ={gray},opacity = 0.3] (0,-2)--(0,12);
	\draw[color ={gray},opacity = 0.3] (0.1,-2)--(0.1,12);
	\draw[color ={gray},opacity = 0.3] (0.2,-2)--(0.2,12);
	\draw[color ={gray},opacity = 0.3] (0.30000000000000004,-2)--(0.30000000000000004,12);
	\draw[color ={gray},opacity = 0.3] (0.4,-2)--(0.4,12);
	\draw[color ={gray},opacity = 0.3] (0.5,-2)--(0.5,12);
	\draw[color ={gray},opacity = 0.3] (0.6,-2)--(0.6,12);
	\draw[color ={gray},opacity = 0.3] (0.7,-2)--(0.7,12);
	\draw[color ={gray},opacity = 0.3] (0.7999999999999999,-2)--(0.7999999999999999,12);
	\draw[color ={gray},opacity = 0.3] (0.8999999999999999,-2)--(0.8999999999999999,12);
	\draw[color ={gray},opacity = 0.3] (1.0999999999999999,-2)--(1.0999999999999999,12);
	\draw[color ={gray},opacity = 0.3] (1.2,-2)--(1.2,12);
	\draw[color ={gray},opacity = 0.3] (1.3,-2)--(1.3,12);
	\draw[color ={gray},opacity = 0.3] (1.4000000000000001,-2)--(1.4000000000000001,12);
	\draw[color ={gray},opacity = 0.3] (1.5000000000000002,-2)--(1.5000000000000002,12);
	\draw[color ={gray},opacity = 0.3] (1.6000000000000003,-2)--(1.6000000000000003,12);
	\draw[color ={gray},opacity = 0.3] (1.7000000000000004,-2)--(1.7000000000000004,12);
	\draw[color ={gray},opacity = 0.3] (1.8000000000000005,-2)--(1.8000000000000005,12);
	\draw[color ={gray},opacity = 0.3] (1.9000000000000006,-2)--(1.9000000000000006,12);
	\draw[color ={gray},opacity = 0.3] (2.1000000000000005,-2)--(2.1000000000000005,12);
	\draw[color ={gray},opacity = 0.3] (2.2000000000000006,-2)--(2.2000000000000006,12);
	\draw[color ={gray},opacity = 0.3] (2.3000000000000007,-2)--(2.3000000000000007,12);
	\draw[color ={gray},opacity = 0.3] (2.400000000000001,-2)--(2.400000000000001,12);
	\draw[color ={gray},opacity = 0.3] (2.500000000000001,-2)--(2.500000000000001,12);
	\draw[color ={gray},opacity = 0.3] (2.600000000000001,-2)--(2.600000000000001,12);
	\draw[color ={gray},opacity = 0.3] (2.700000000000001,-2)--(2.700000000000001,12);
	\draw[color ={gray},opacity = 0.3] (2.800000000000001,-2)--(2.800000000000001,12);
	\draw[color ={gray},opacity = 0.3] (2.9000000000000012,-2)--(2.9000000000000012,12);
	\draw[color ={gray},opacity = 0.3] (3.1000000000000014,-2)--(3.1000000000000014,12);
	\draw[color ={gray},opacity = 0.3] (3.2000000000000015,-2)--(3.2000000000000015,12);
	\draw[color ={gray},opacity = 0.3] (3.3000000000000016,-2)--(3.3000000000000016,12);
	\draw[color ={gray},opacity = 0.3] (3.4000000000000017,-2)--(3.4000000000000017,12);
	\draw[color ={gray},opacity = 0.3] (3.5000000000000018,-2)--(3.5000000000000018,12);
	\draw[color ={gray},opacity = 0.3] (3.600000000000002,-2)--(3.600000000000002,12);
	\draw[color ={gray},opacity = 0.3] (3.700000000000002,-2)--(3.700000000000002,12);
	\draw[color ={gray},opacity = 0.3] (3.800000000000002,-2)--(3.800000000000002,12);
	\draw[color ={gray},opacity = 0.3] (3.900000000000002,-2)--(3.900000000000002,12);
	\draw[color ={gray},opacity = 0.3] (4.100000000000001,-2)--(4.100000000000001,12);
	\draw[color ={gray},opacity = 0.3] (4.200000000000001,-2)--(4.200000000000001,12);
	\draw[color ={gray},opacity = 0.3] (4.300000000000001,-2)--(4.300000000000001,12);
	\draw[color ={gray},opacity = 0.3] (4.4,-2)--(4.4,12);
	\draw[color ={gray},opacity = 0.3] (4.5,-2)--(4.5,12);
	\draw[color ={gray},opacity = 0.3] (4.6,-2)--(4.6,12);
	\draw[color ={gray},opacity = 0.3] (4.699999999999999,-2)--(4.699999999999999,12);
	\draw[color ={gray},opacity = 0.3] (4.799999999999999,-2)--(4.799999999999999,12);
	\draw[color ={gray},opacity = 0.3] (4.899999999999999,-2)--(4.899999999999999,12);
	\draw[color ={gray},opacity = 0.3] (5.099999999999998,-2)--(5.099999999999998,12);
	\draw[color ={gray},opacity = 0.3] (5.1999999999999975,-2)--(5.1999999999999975,12);
	\draw[color ={gray},opacity = 0.3] (5.299999999999997,-2)--(5.299999999999997,12);
	\draw[color ={gray},opacity = 0.3] (5.399999999999997,-2)--(5.399999999999997,12);
	\draw[color ={gray},opacity = 0.3] (5.4999999999999964,-2)--(5.4999999999999964,12);
	\draw[color ={gray},opacity = 0.3] (5.599999999999996,-2)--(5.599999999999996,12);
	\draw[color ={gray},opacity = 0.3] (5.699999999999996,-2)--(5.699999999999996,12);
	\draw[color ={gray},opacity = 0.3] (5.799999999999995,-2)--(5.799999999999995,12);
	\draw[color ={gray},opacity = 0.3] (5.899999999999995,-2)--(5.899999999999995,12);
	\draw[color ={gray},opacity = 0.3] (6.099999999999994,-2)--(6.099999999999994,12);
	\draw[color ={gray},opacity = 0.3] (6.199999999999994,-2)--(6.199999999999994,12);
	\draw[color ={gray},opacity = 0.3] (6.299999999999994,-2)--(6.299999999999994,12);
	\draw[color ={gray},opacity = 0.3] (6.399999999999993,-2)--(6.399999999999993,12);
	\draw[color ={gray},opacity = 0.3] (6.499999999999993,-2)--(6.499999999999993,12);
	\draw[color ={gray},opacity = 0.3] (6.5999999999999925,-2)--(6.5999999999999925,12);
	\draw[color ={gray},opacity = 0.3] (6.699999999999992,-2)--(6.699999999999992,12);
	\draw[color ={gray},opacity = 0.3] (6.799999999999992,-2)--(6.799999999999992,12);
	\draw[color ={gray},opacity = 0.3] (6.8999999999999915,-2)--(6.8999999999999915,12);
	\draw[color ={gray},opacity = 0.3] (7.099999999999991,-2)--(7.099999999999991,12);
	\draw[color ={gray},opacity = 0.3] (7.19999999999999,-2)--(7.19999999999999,12);
	\draw[color ={gray},opacity = 0.3] (7.29999999999999,-2)--(7.29999999999999,12);
	\draw[color ={gray},opacity = 0.3] (7.39999999999999,-2)--(7.39999999999999,12);
	\draw[color ={gray},opacity = 0.3] (7.499999999999989,-2)--(7.499999999999989,12);
	\draw[color ={gray},opacity = 0.3] (7.599999999999989,-2)--(7.599999999999989,12);
	\draw[color ={gray},opacity = 0.3] (7.699999999999989,-2)--(7.699999999999989,12);
	\draw[color ={gray},opacity = 0.3] (7.799999999999988,-2)--(7.799999999999988,12);
	\draw[color ={gray},opacity = 0.3] (7.899999999999988,-2)--(7.899999999999988,12);
	\draw[color ={gray},opacity = 0.3] (8.099999999999987,-2)--(8.099999999999987,12);
	\draw[color ={gray},opacity = 0.3] (8.199999999999987,-2)--(8.199999999999987,12);
	\draw[color ={gray},opacity = 0.3] (8.299999999999986,-2)--(8.299999999999986,12);
	\draw[color ={gray},opacity = 0.3] (8.399999999999986,-2)--(8.399999999999986,12);
	\draw[color ={gray},opacity = 0.3] (8.499999999999986,-2)--(8.499999999999986,12);
	\draw[color ={gray},opacity = 0.3] (8.599999999999985,-2)--(8.599999999999985,12);
	\draw[color ={gray},opacity = 0.3] (8.699999999999985,-2)--(8.699999999999985,12);
	\draw[color ={gray},opacity = 0.3] (8.799999999999985,-2)--(8.799999999999985,12);
	\draw[color ={gray},opacity = 0.3] (8.899999999999984,-2)--(8.899999999999984,12);
	\draw[color ={gray},opacity = 0.3] (9.099999999999984,-2)--(9.099999999999984,12);
	\draw[color ={gray},opacity = 0.3] (9.199999999999983,-2)--(9.199999999999983,12);
	\draw[color ={gray},opacity = 0.3] (9.299999999999983,-2)--(9.299999999999983,12);
	\draw[color ={gray},opacity = 0.3] (9.399999999999983,-2)--(9.399999999999983,12);
	\draw[color ={gray},opacity = 0.3] (9.499999999999982,-2)--(9.499999999999982,12);
	\draw[color ={gray},opacity = 0.3] (9.599999999999982,-2)--(9.599999999999982,12);
	\draw[color ={gray},opacity = 0.3] (9.699999999999982,-2)--(9.699999999999982,12);
	\draw[color ={gray},opacity = 0.3] (9.799999999999981,-2)--(9.799999999999981,12);
	\draw[color ={gray},opacity = 0.3] (9.89999999999998,-2)--(9.89999999999998,12);
	\draw[color ={gray},opacity = 0.3] (10.09999999999998,-2)--(10.09999999999998,12);
	\draw[color ={gray},opacity = 0.3] (10.19999999999998,-2)--(10.19999999999998,12);
	\draw[color ={gray},opacity = 0.3] (10.29999999999998,-2)--(10.29999999999998,12);
	\draw[color ={gray},opacity = 0.3] (10.399999999999979,-2)--(10.399999999999979,12);
	\draw[color ={gray},opacity = 0.3] (10.499999999999979,-2)--(10.499999999999979,12);
	\draw[color ={gray},opacity = 0.3] (10.599999999999978,-2)--(10.599999999999978,12);
	\draw[color ={gray},opacity = 0.3] (10.699999999999978,-2)--(10.699999999999978,12);
	\draw[color ={gray},opacity = 0.3] (10.799999999999978,-2)--(10.799999999999978,12);
	\draw[color ={gray},opacity = 0.3] (10.899999999999977,-2)--(10.899999999999977,12);
	\draw[color ={gray},opacity = 0.3] (11.099999999999977,-2)--(11.099999999999977,12);
	\draw[color ={gray},opacity = 0.3] (11.199999999999976,-2)--(11.199999999999976,12);
	\draw[color ={gray},opacity = 0.3] (11.299999999999976,-2)--(11.299999999999976,12);
	\draw[color ={gray},opacity = 0.3] (11.399999999999975,-2)--(11.399999999999975,12);
	\draw[color ={gray},opacity = 0.3] (11.499999999999975,-2)--(11.499999999999975,12);
	\draw[color ={gray},opacity = 0.3] (11.599999999999975,-2)--(11.599999999999975,12);
	\draw[color ={gray},opacity = 0.3] (11.699999999999974,-2)--(11.699999999999974,12);
	\draw[color ={gray},opacity = 0.3] (11.799999999999974,-2)--(11.799999999999974,12);
	\draw[color ={gray},opacity = 0.3] (11.899999999999974,-2)--(11.899999999999974,12);
	\draw[color ={gray},opacity = 0.3] (12.099999999999973,-2)--(12.099999999999973,12);
	\draw[color ={gray},opacity = 0.3] (12.199999999999973,-2)--(12.199999999999973,12);
	\draw[color ={gray},opacity = 0.3] (12.299999999999972,-2)--(12.299999999999972,12);
	\draw[color ={gray},opacity = 0.3] (12.399999999999972,-2)--(12.399999999999972,12);
	\draw[color ={gray},opacity = 0.3] (12.499999999999972,-2)--(12.499999999999972,12);
	\draw[color ={gray},opacity = 0.3] (12.599999999999971,-2)--(12.599999999999971,12);
	\draw[color ={gray},opacity = 0.3] (12.69999999999997,-2)--(12.69999999999997,12);
	\draw[color ={gray},opacity = 0.3] (12.79999999999997,-2)--(12.79999999999997,12);
	\draw[color ={gray},opacity = 0.3] (12.89999999999997,-2)--(12.89999999999997,12);
	\draw[color ={gray},opacity = 0.3] (13.09999999999997,-2)--(13.09999999999997,12);
	\draw[color ={gray},opacity = 0.3] (13.199999999999969,-2)--(13.199999999999969,12);
	\draw[color ={gray},opacity = 0.3] (13.299999999999969,-2)--(13.299999999999969,12);
	\draw[color ={gray},opacity = 0.3] (13.399999999999968,-2)--(13.399999999999968,12);
	\draw[color ={gray},opacity = 0.3] (13.499999999999968,-2)--(13.499999999999968,12);
	\draw[color ={gray},opacity = 0.3] (13.599999999999968,-2)--(13.599999999999968,12);
	\draw[color ={gray},opacity = 0.3] (13.699999999999967,-2)--(13.699999999999967,12);
	\draw[color ={gray},opacity = 0.3] (13.799999999999967,-2)--(13.799999999999967,12);
	\draw[color ={gray},opacity = 0.3] (13.899999999999967,-2)--(13.899999999999967,12);
	\draw[color ={gray},opacity = 0.3] (0,-2)--(0,12);
	\draw[color ={gray},opacity = 0.3] (-0.1,-2)--(-0.1,12);
	\draw[color ={gray},opacity = 0.3] (-0.2,-2)--(-0.2,12);
	\draw[color ={gray},opacity = 0.3] (-0.30000000000000004,-2)--(-0.30000000000000004,12);
	\draw[color ={gray},opacity = 0.3] (-0.4,-2)--(-0.4,12);
	\draw[color ={gray},opacity = 0.3] (-0.5,-2)--(-0.5,12);
	\draw[color ={gray},opacity = 0.3] (-0.6,-2)--(-0.6,12);
	\draw[color ={gray},opacity = 0.3] (-0.7,-2)--(-0.7,12);
	\draw[color ={gray},opacity = 0.3] (-0.7999999999999999,-2)--(-0.7999999999999999,12);
	\draw[color ={gray},opacity = 0.3] (-0.8999999999999999,-2)--(-0.8999999999999999,12);
	\draw[color ={gray},opacity = 0.3] (-1.0999999999999999,-2)--(-1.0999999999999999,12);
	\draw[color ={gray},opacity = 0.3] (-1.2,-2)--(-1.2,12);
	\draw[color ={gray},opacity = 0.3] (-1.3,-2)--(-1.3,12);
	\draw[color ={gray},opacity = 0.3] (-1.4000000000000001,-2)--(-1.4000000000000001,12);
	\draw[color ={gray},opacity = 0.3] (-1.5000000000000002,-2)--(-1.5000000000000002,12);
	\draw[color ={gray},opacity = 0.3] (-1.6000000000000003,-2)--(-1.6000000000000003,12);
	\draw[color ={gray},opacity = 0.3] (-1.7000000000000004,-2)--(-1.7000000000000004,12);
	\draw[color ={gray},opacity = 0.3] (-1.8000000000000005,-2)--(-1.8000000000000005,12);
	\draw[color ={gray},opacity = 0.3] (-1.9000000000000006,-2)--(-1.9000000000000006,12);
	\draw[color ={black},line width = 2] (2,-0.2)--(2,0.2);
	\draw[color ={black},line width = 2] (4,-0.2)--(4,0.2);
	\draw[color ={black},line width = 2] (6,-0.2)--(6,0.2);
	\draw[color ={black},line width = 2] (8,-0.2)--(8,0.2);
	\draw[color ={black},line width = 2] (10,-0.2)--(10,0.2);
	\draw[color ={black},line width = 2] (12,-0.2)--(12,0.2);
	\draw[color ={black},line width = 2] (-0.2,2)--(0.2,2);
	\draw[color ={black},line width = 2] (-0.2,4)--(0.2,4);
	\draw[color ={black},line width = 2] (-0.2,6)--(0.2,6);
	\draw[color ={black},line width = 2] (-0.2,8)--(0.2,8);
	\draw[color ={black},line width = 2] (-0.2,10)--(0.2,10);
	\draw (2,-0.4) node[anchor = center] {\scriptsize \color{black}{$1$}};
	\draw (4,-0.4) node[anchor = center] {\scriptsize \color{black}{$2$}};
	\draw (6,-0.4) node[anchor = center] {\scriptsize \color{black}{$3$}};
	\draw (8,-0.4) node[anchor = center] {\scriptsize \color{black}{$4$}};
	\draw (10,-0.4) node[anchor = center] {\scriptsize \color{black}{$5$}};
	\draw (-0.5,2.1) node[anchor = center] {\footnotesize \color{black}{$1$}};
	\draw (-0.5,4.1) node[anchor = center] {\footnotesize \color{black}{$2$}};
	\draw (-0.5,6.1) node[anchor = center] {\footnotesize \color{black}{$3$}};
	\draw (-0.5,8.1) node[anchor = center] {\footnotesize \color{black}{$4$}};
	\draw (-0.5,10.1) node[anchor = center] {\footnotesize \color{black}{$5$}};
	\draw [color={black}] (-0.3,-0.3) node[anchor = center,scale=1, rotate = 0] {O};
	
	\draw[color={blue},line width = 2] (0.6,8.67)--(1,6)--(1.4,4.86)--(1.8,4.22)--(2.2,3.82)--(2.6,3.54)--(3,3.33)--(3.4,3.18)--(3.8,3.05)--(4.2,2.95)--(4.6,2.87)--(5,2.8)--(5.4,2.74)--(5.8,2.69)--(6.2,2.65)--(6.6,2.61)--(7,2.57)--(7.4,2.54)--(7.8,2.51)--(8.2,2.49)--(8.6,2.47)--(9,2.44)--(9.4,2.43)--(9.8,2.41)--(10.2,2.39)--(10.6,2.38)--(11,2.36)--(11.4,2.35)--(11.8,2.34)--(12.2,2.33)--(12.6,2.32)--(13,2.31)--(13.4,2.3)--(13.8,2.29);
	
	\draw[color={red},line width = 2] (-2,8)--(-1.6,7.6)--(-1.2,7.2)--(-0.8,6.8)--(-0.4,6.4)--(0,6)--(0.4,5.6)--(0.8,5.2)--(1.2,4.8)--(1.6,4.4)--(2,4)--(2.4,3.6)--(2.8,3.2)--(3.2,2.8)--(3.6,2.4)--(4,2)--(4.4,1.6)--(4.8,1.2)--(5.2,0.8)--(5.6,0.4)--(6,0)--(6.4,-0.4)--(6.8,-0.8)--(7.2,-1.2)--(7.6,-1.6)--(8,-2)--(8.4,-2.4)--(8.8,-2.8)--(9.2,-3.2)--(9.6,-3.6);

\end{tikzpicture}\end{center}
\end{exercice}

\begin{Solution}
 $f'(1)$ est donné par le coefficient directeur de la tangente à la courbe au point d'abscisse $1$, soit $-1=-1$.
\end{Solution}

% @see : https://coopmaths.fr/alea?uuid=6f32d&id=can1F16&n=1&d=10&alea=dqSw223232323234234567891011121314151617181920212223242526272829303132333435363738394041424344454647482345678910111213141516171819202122232425&cd=1&cols=1
\needspace{10\baselineskip}
\newpage
\begin{exercice}[BaremeTotal=false]


 On donne les représentations graphiques d'une fonction et de sa dérivée.\\
        Donner l'équation réduite de la tangente à la courbe de $f$ en $x=-2$. \\ \begin{minipage}{0.5\linewidth}
    \begin{tikzpicture}[baseline, scale=0.8]

    \tikzset{
      point/.style={
        thick,
        draw,
        cross out,
        inner sep=0pt,
        minimum width=5pt,
        minimum height=5pt,
      },
    }
    \clip (-4.5,-4.5) rectangle (5.75,13.5);
    	
	\draw[color ={black},>=latex,->] (-3,0)--(3,0);
	\draw[color ={black},>=latex,->] (0,-3)--(0,12);
	\draw[color ={black},opacity = 0.4] (-3,1)--(3,1);
	\draw[color ={black},opacity = 0.4] (-3,2)--(3,2);
	\draw[color ={black},opacity = 0.4] (-3,3)--(3,3);
	\draw[color ={black},opacity = 0.4] (-3,4)--(3,4);
	\draw[color ={black},opacity = 0.4] (-3,5)--(3,5);
	\draw[color ={black},opacity = 0.4] (-3,6)--(3,6);
	\draw[color ={black},opacity = 0.4] (-3,7)--(3,7);
	\draw[color ={black},opacity = 0.4] (-3,8)--(3,8);
	\draw[color ={black},opacity = 0.4] (-3,9)--(3,9);
	\draw[color ={black},opacity = 0.4] (-3,10)--(3,10);
	\draw[color ={black},opacity = 0.4] (-3,11)--(3,11);
	\draw[color ={black},opacity = 0.4] (-3,12)--(3,12);
	\draw[color ={black},opacity = 0.4] (-3,-1)--(3,-1);
	\draw[color ={black},opacity = 0.4] (-3,-2)--(3,-2);
	\draw[color ={black},opacity = 0.4] (-3,-3)--(3,-3);
	\draw[color ={black},opacity = 0.4] (1,-3)--(1,12);
	\draw[color ={black},opacity = 0.4] (2,-3)--(2,12);
	\draw[color ={black},opacity = 0.4] (3,-3)--(3,12);
	\draw[color ={black},opacity = 0.4] (-1,-3)--(-1,12);
	\draw[color ={black},opacity = 0.4] (-2,-3)--(-2,12);
	\draw[color ={black},opacity = 0.4] (-3,-3)--(-3,12);
	\draw[color ={gray},opacity = 0.3] (-3,0)--(3,0);
	\draw[color ={gray},opacity = 0.3] (-3,0.5)--(3,0.5);
	\draw[color ={gray},opacity = 0.3] (-3,1.5)--(3,1.5);
	\draw[color ={gray},opacity = 0.3] (-3,2.5)--(3,2.5);
	\draw[color ={gray},opacity = 0.3] (-3,3.5)--(3,3.5);
	\draw[color ={gray},opacity = 0.3] (-3,4.5)--(3,4.5);
	\draw[color ={gray},opacity = 0.3] (-3,5.5)--(3,5.5);
	\draw[color ={gray},opacity = 0.3] (-3,6.5)--(3,6.5);
	\draw[color ={gray},opacity = 0.3] (-3,7.5)--(3,7.5);
	\draw[color ={gray},opacity = 0.3] (-3,8.5)--(3,8.5);
	\draw[color ={gray},opacity = 0.3] (-3,9.5)--(3,9.5);
	\draw[color ={gray},opacity = 0.3] (-3,10.5)--(3,10.5);
	\draw[color ={gray},opacity = 0.3] (-3,11.5)--(3,11.5);
	\draw[color ={gray},opacity = 0.3] (-3,0)--(3,0);
	\draw[color ={gray},opacity = 0.3] (-3,-0.5)--(3,-0.5);
	\draw[color ={gray},opacity = 0.3] (-3,-1.5)--(3,-1.5);
	\draw[color ={gray},opacity = 0.3] (-3,-2.5)--(3,-2.5);
	\draw[color ={gray},opacity = 0.3] (0,-3)--(0,12);
	\draw[color ={gray},opacity = 0.3] (0,-3)--(0,12);
	\draw[color ={black},line width = 1.2] (1,-0.13)--(1,0.13);
	\draw[color ={black},line width = 1.2] (2,-0.13)--(2,0.13);
	\draw[color ={black},line width = 1.2] (3,-0.13)--(3,0.13);
	\draw[color ={black},line width = 1.2] (-1,-0.13)--(-1,0.13);
	\draw[color ={black},line width = 1.2] (-2,-0.13)--(-2,0.13);
	\draw[color ={black},line width = 1.2] (-3,-0.13)--(-3,0.13);
	\draw[color ={black},line width = 1.2] (-0.13,1)--(0.13,1);
	\draw[color ={black},line width = 1.2] (-0.13,2)--(0.13,2);
	\draw[color ={black},line width = 1.2] (-0.13,3)--(0.13,3);
	\draw[color ={black},line width = 1.2] (-0.13,4)--(0.13,4);
	\draw[color ={black},line width = 1.2] (-0.13,5)--(0.13,5);
	\draw[color ={black},line width = 1.2] (-0.13,6)--(0.13,6);
	\draw[color ={black},line width = 1.2] (-0.13,7)--(0.13,7);
	\draw[color ={black},line width = 1.2] (-0.13,8)--(0.13,8);
	\draw[color ={black},line width = 1.2] (-0.13,9)--(0.13,9);
	\draw[color ={black},line width = 1.2] (-0.13,10)--(0.13,10);
	\draw[color ={black},line width = 1.2] (-0.13,11)--(0.13,11);
	\draw[color ={black},line width = 1.2] (-0.13,12)--(0.13,12);
	\draw[color ={black},line width = 1.2] (-0.13,-1)--(0.13,-1);
	\draw[color ={black},line width = 1.2] (-0.13,-2)--(0.13,-2);
	\draw[color ={black},line width = 1.2] (-0.13,-3)--(0.13,-3);
	\draw (3,-0.4) node[anchor = center] {\scriptsize \color{black}{$3$}};
	\draw (-3,-0.4) node[anchor = center] {\scriptsize \color{black}{$-3$}};
	\draw (-0.5,3.1) node[anchor = center] {\footnotesize \color{black}{$3$}};
	\draw (-0.5,6.1) node[anchor = center] {\footnotesize \color{black}{$6$}};
	\draw (-0.5,9.1) node[anchor = center] {\footnotesize \color{black}{$9$}};
	\draw (-0.5,-2.9) node[anchor = center] {\footnotesize \color{black}{$-3$}};
	\draw [color={black}] (-0.3,-0.3) node[anchor = center,scale=1, rotate = 0] {O};
	\draw (3,10) node[anchor = center] {\footnotesize \color{blue}{$\Large \mathcal{C}_f$}};
	
	\draw[color={blue},line width = 2] (-3,2)--(-2.8,1.24)--(-2.6,0.56)--(-2.4,-0.04)--(-2.2,-0.56)--(-2,-1)--(-1.8,-1.36)--(-1.6,-1.64)--(-1.4,-1.84)--(-1.2,-1.96)--(-1,-2)--(-0.8,-1.96)--(-0.6,-1.84)--(-0.4,-1.64)--(-0.2,-1.36)--(0,-1)--(0.2,-0.56)--(0.4,-0.04)--(0.6,0.56)--(0.8,1.24)--(1,2)--(1.2,2.84)--(1.4,3.76)--(1.6,4.76)--(1.8,5.84)--(2,7)--(2.2,8.24)--(2.4,9.56)--(2.6,10.96)--(2.8,12.44);

\end{tikzpicture}\\
    \end{minipage}
    \begin{minipage}{0.5\linewidth}
    \begin{tikzpicture}[baseline, scale=0.8]

    \tikzset{
      point/.style={
        thick,
        draw,
        cross out,
        inner sep=0pt,
        minimum width=5pt,
        minimum height=5pt,
      },
    }
    \clip (-4.5,-6.5) rectangle (6.5,9.5);
    	
	\draw[color ={black},>=latex,->] (-3,0)--(3,0);
	\draw[color ={black},>=latex,->] (0,-5)--(0,8);
	\draw[color ={black},opacity = 0.4] (-3,1)--(3,1);
	\draw[color ={black},opacity = 0.4] (-3,2)--(3,2);
	\draw[color ={black},opacity = 0.4] (-3,3)--(3,3);
	\draw[color ={black},opacity = 0.4] (-3,4)--(3,4);
	\draw[color ={black},opacity = 0.4] (-3,5)--(3,5);
	\draw[color ={black},opacity = 0.4] (-3,6)--(3,6);
	\draw[color ={black},opacity = 0.4] (-3,7)--(3,7);
	\draw[color ={black},opacity = 0.4] (-3,8)--(3,8);
	\draw[color ={black},opacity = 0.4] (-3,-1)--(3,-1);
	\draw[color ={black},opacity = 0.4] (-3,-2)--(3,-2);
	\draw[color ={black},opacity = 0.4] (-3,-3)--(3,-3);
	\draw[color ={black},opacity = 0.4] (-3,-4)--(3,-4);
	\draw[color ={black},opacity = 0.4] (-3,-5)--(3,-5);
	\draw[color ={black},opacity = 0.4] (1,-5)--(1,8);
	\draw[color ={black},opacity = 0.4] (2,-5)--(2,8);
	\draw[color ={black},opacity = 0.4] (3,-5)--(3,8);
	\draw[color ={black},opacity = 0.4] (-1,-5)--(-1,8);
	\draw[color ={black},opacity = 0.4] (-2,-5)--(-2,8);
	\draw[color ={black},opacity = 0.4] (-3,-5)--(-3,8);
	\draw[color ={gray},opacity = 0.3] (-3,0)--(3,0);
	\draw[color ={gray},opacity = 0.3] (-3,0.5)--(3,0.5);
	\draw[color ={gray},opacity = 0.3] (-3,1.5)--(3,1.5);
	\draw[color ={gray},opacity = 0.3] (-3,2.5)--(3,2.5);
	\draw[color ={gray},opacity = 0.3] (-3,3.5)--(3,3.5);
	\draw[color ={gray},opacity = 0.3] (-3,4.5)--(3,4.5);
	\draw[color ={gray},opacity = 0.3] (-3,5.5)--(3,5.5);
	\draw[color ={gray},opacity = 0.3] (-3,6.5)--(3,6.5);
	\draw[color ={gray},opacity = 0.3] (-3,7.5)--(3,7.5);
	\draw[color ={gray},opacity = 0.3] (-3,0)--(3,0);
	\draw[color ={gray},opacity = 0.3] (-3,-0.5)--(3,-0.5);
	\draw[color ={gray},opacity = 0.3] (-3,-1.5)--(3,-1.5);
	\draw[color ={gray},opacity = 0.3] (-3,-2.5)--(3,-2.5);
	\draw[color ={gray},opacity = 0.3] (-3,-3.5)--(3,-3.5);
	\draw[color ={gray},opacity = 0.3] (-3,-4.5)--(3,-4.5);
	\draw[color ={gray},opacity = 0.3] (0,-5)--(0,8);
	\draw[color ={gray},opacity = 0.3] (0,-5)--(0,8);
	\draw[color ={black},line width = 1.2] (1,-0.13)--(1,0.13);
	\draw[color ={black},line width = 1.2] (2,-0.13)--(2,0.13);
	\draw[color ={black},line width = 1.2] (3,-0.13)--(3,0.13);
	\draw[color ={black},line width = 1.2] (-1,-0.13)--(-1,0.13);
	\draw[color ={black},line width = 1.2] (-2,-0.13)--(-2,0.13);
	\draw[color ={black},line width = 1.2] (-3,-0.13)--(-3,0.13);
	\draw[color ={black},line width = 1.2] (-0.13,1)--(0.13,1);
	\draw[color ={black},line width = 1.2] (-0.13,2)--(0.13,2);
	\draw[color ={black},line width = 1.2] (-0.13,3)--(0.13,3);
	\draw[color ={black},line width = 1.2] (-0.13,4)--(0.13,4);
	\draw[color ={black},line width = 1.2] (-0.13,5)--(0.13,5);
	\draw[color ={black},line width = 1.2] (-0.13,6)--(0.13,6);
	\draw[color ={black},line width = 1.2] (-0.13,7)--(0.13,7);
	\draw[color ={black},line width = 1.2] (-0.13,8)--(0.13,8);
	\draw[color ={black},line width = 1.2] (-0.13,9)--(0.13,9);
	\draw[color ={black},line width = 1.2] (-0.13,10)--(0.13,10);
	\draw[color ={black},line width = 1.2] (-0.13,11)--(0.13,11);
	\draw[color ={black},line width = 1.2] (-0.13,12)--(0.13,12);
	\draw[color ={black},line width = 1.2] (-0.13,-1)--(0.13,-1);
	\draw[color ={black},line width = 1.2] (-0.13,-2)--(0.13,-2);
	\draw[color ={black},line width = 1.2] (-0.13,-3)--(0.13,-3);
	\draw (3,-0.4) node[anchor = center] {\scriptsize \color{black}{$3$}};
	\draw (-3,-0.4) node[anchor = center] {\scriptsize \color{black}{$-3$}};
	\draw (-0.5,3.1) node[anchor = center] {\footnotesize \color{black}{$3$}};
	\draw (-0.5,6.1) node[anchor = center] {\footnotesize \color{black}{$6$}};
	\draw (-0.5,-2.9) node[anchor = center] {\footnotesize \color{black}{$-3$}};
	\draw [color={black}] (-0.3,-0.3) node[anchor = center,scale=1, rotate = 0] {O};
	\draw (3,6) node[anchor = center] {\footnotesize \color{red}{$\Large\mathcal{C}_f\prime$}};
	
	\draw[color={red},line width = 2] (-3,-4)--(-2.8,-3.6)--(-2.6,-3.2)--(-2.4,-2.8)--(-2.2,-2.4)--(-2,-2)--(-1.8,-1.6)--(-1.6,-1.2)--(-1.4,-0.8)--(-1.2,-0.4)--(-1,0)--(-0.8,0.4)--(-0.6,0.8)--(-0.4,1.2)--(-0.2,1.6)--(0,2)--(0.2,2.4)--(0.4,2.8)--(0.6,3.2)--(0.8,3.6)--(1,4)--(1.2,4.4)--(1.4,4.8)--(1.6,5.2)--(1.8,5.6)--(2,6)--(2.2,6.4)--(2.4,6.8)--(2.6,7.2)--(2.8,7.6)--(3,8);

\end{tikzpicture}\\
    \end{minipage}
    
\end{exercice}

\begin{Solution}
 L'équation réduite de la tangente au point d'abscisse $-2$ est  : $y=f'(-2)(x-(-2))+f(-2)$.\\
        On lit graphiquement $f(-2)=-1$ et $f'(-2)=-2$.\\
        L'équation réduite de la tangente est donc donnée par :
        $y=-2(x+2)-1$, soit $y=-2x-5$.
\end{Solution}

% @see : https://coopmaths.fr/alea?uuid=a1ba2&id=can1F14&n=1&d=10&alea=rkIp223232323234234567891011121314151617181920212223242526272829303132333435363738394041424344454647482345678910111213141516171819202122232425&cd=1&cols=1
\needspace{10\baselineskip}
\begin{exercice}[BaremeTotal=false]


Soit $f$ la fonction définie sur $\mathbb{R}$ par :
$f(x)= -x+9-x^2$.\\
En calculant la limite du taux d'accroissement, déterminer $f'(-3)$.
\end{exercice}

\begin{Solution}
 $f$ est une fonction polynôme du second degré de la forme $f(x)=ax^2+bx+c$.\\
    La fonction dérivée est donnée par la somme des dérivées des fonctions $u$ et $v$ définies par $u(x)=-x^2$ et $v(x)=-x+9$.\\
     Comme $u'(x)=-2x$ et $v'(x)=-1$, on obtient  $f'(x)=-2x-1$. \\
     Ainsi, $f'(-3)=-2\times (-3)-1=5$.
\end{Solution}

% @see : https://coopmaths.fr/alea?uuid=3c690&id=can1F13&n=1&d=10&alea=EUDu223232323234234567891011121314151617181920212223242526272829303132333435363738394041424344454647482345678910111213141516171819202122232425&cd=1&cols=1
\needspace{10\baselineskip}
\begin{exercice}[BaremeTotal=false]


 Déterminer le coefficient directeur de la tangente à la courbe représentative de la fonction racine carrée au point d'abscisse $7$.
         
\end{exercice}

\begin{Solution}
 Le coefficient directeur de la tangente au point d'abscisse $a$ est donné par le nombre dérivé $f'(a)$.\\
        La fonction $f$ définie par $f(x)=\sqrt{x}$ a pour fonction dérivée la fonction $f'$ définie par $f'(x)=\dfrac{1}{2\sqrt{x}}$.\\Comme $f'(7)=\dfrac{1}{2\sqrt{7}}$, le coefficient directeur de la tangente au point d'abscisse $7$ est : $\dfrac{1}{2\sqrt{7}}$.
\end{Solution}

% @see : https://coopmaths.fr/alea?uuid=ebd89&id=1AN14-1&n=1&d=10&s=3&alea=ocYr223232323234234567891011121314151617181920212223242526272829303132333435363738394041424344454647482345678910111213141516171819202122232425&cd=0&cols=1
\needspace{10\baselineskip}
\begin{exercice}[BaremeTotal=false]


 Donner la dérivée de la fonction $f$, dérivable pour tout $x\in\R$, définie par  $f(x)=-\dfrac{5}{7}x+6$.
\end{exercice}

\begin{Solution}
 L'expression de la dérivée de la fonction $f$ définie par $f(x)=-\dfrac{5}{7}x+6$ est : ${\color[HTML]{f15929}\boldsymbol{f'(x)=-\dfrac{5}{7}}}$.
\end{Solution}

% @see : https://coopmaths.fr/alea?uuid=ebd89&id=1AN14-1&n=1&d=10&s=4&alea=uAlP223232323234234567891011121314151617181920212223242526272829303132333435363738394041424344454647482345678910111213141516171819202122232425&cd=0&cols=1
\needspace{10\baselineskip}
\begin{exercice}[BaremeTotal=false]


 Donner la dérivée de la fonction $f$, dérivable pour tout $x\in\R$, définie par  $f(x)=10x^{14}$.
\end{exercice}

\begin{Solution}
 L'expression de la dérivée de la fonction $f$ définie par $f(x)=10x^{14}$ est : ${\color[HTML]{f15929}\boldsymbol{f'(x)=140x^{13}}}$.
\end{Solution}

% @see : https://coopmaths.fr/alea?uuid=ebd89&id=1AN14-1&n=1&d=10&s=7&alea=vjN7223232323234234567891011121314151617181920212223242526272829303132333435363738394041424344454647482345678910111213141516171819202122232425&cd=0&cols=1
\needspace{10\baselineskip}
\begin{exercice}[BaremeTotal=false]


 Donner la dérivée de la fonction $f$, dérivable pour tout $x\in\R^*$, définie par  $f(x)=\dfrac{-1}{10x}$.
\end{exercice}

\begin{Solution}
 L'expression de la dérivée de la fonction $f$ définie par $f(x)=\dfrac{-1}{10x}$ est : ${\color[HTML]{f15929}\boldsymbol{f'(x)=\dfrac{1}{10x^2}}}$.
\end{Solution}

% @see : https://coopmaths.fr/alea?uuid=a83c0&id=1AN14-4&n=1&d=10&s=1&alea=Njr3223232323234234567891011121314151617181920212223242526272829303132333435363738394041424344454647482345678910111213141516171819202122232425&cd=0&cols=1
\needspace{10\baselineskip}
\begin{exercice}[BaremeTotal=false]


 Donner l'expression de la dérivée de la fonction $f$ définie sur $\R^{*}$ par $f(x)=-\dfrac{5}{x}+2x$.\\
\end{exercice}

\begin{Solution}
 L'expression de la dérivée de la fonction $f$ définie par $f(x)=-\dfrac{5}{x}+2x$ est :\\$f'(x)={\color[HTML]{f15929}\boldsymbol{\dfrac{5}{x^2}+2}}$.\\
\end{Solution}

% @see : https://coopmaths.fr/alea?uuid=a83c0&id=1AN14-4&n=1&d=10&s=4&alea=4Vpz223232323234234567891011121314151617181920212223242526272829303132333435363738394041424344454647482345678910111213141516171819202122232425&cd=1&cols=1
\needspace{10\baselineskip}
\begin{exercice}[BaremeTotal=false]


 Donner l'expression de la dérivée de la fonction $f$ définie sur $\R_+^{*}$ par $f(x)=2\sqrt{x}-\dfrac{5}{x}+4x^{2}-3x$.\\
\end{exercice}

\begin{Solution}
 $(2\sqrt{x})^\prime=\dfrac{2}{2\sqrt{x}}$\\$(-\dfrac{5}{x})^\prime=\dfrac{5}{x^2}$\\$(4x^{2}-3x)^\prime=8x-3$\\L'expression de la dérivée de la fonction $f$ définie par $f(x)=2\sqrt{x}-\dfrac{5}{x}+4x^{2}-3x$ est :\\$f'(x)={\color[HTML]{f15929}\boldsymbol{\dfrac{2}{2\sqrt{x}}+\dfrac{5}{x^2}+8x-3}}$.\\
\end{Solution}

\end{Maquette}
\clearpage
\begin{Maquette}[DM]{Niveau= ,Classe=\getdata[26]\mydata\ (v26), Date = 06/01/25  ,Theme=Exercices}


% @see : https://coopmaths.fr/alea?uuid=4c8c7&id=1AN11-3&n=1&d=10&s=2&alea=nawO22323232323423456789101112131415161718192021222324252627282930313233343536373839404142434445464748234567891011121314151617181920212223242526&cd=1&cols=2
\needspace{10\baselineskip}
\begin{exercice}[BaremeTotal=false]
\label{eleve:26}


 \begin{minipage}[t]{\linewidth} Soit $f$ une fonction dérivable sur $[-5;5]$ et $\mathcal{C}_f$ sa courbe représentative.\\On sait que  $f(0)=-1~~$ et que $~~f'(0)=-4$.\\Déterminer une équation de la tangente $(T)$ à la courbe $\mathcal{C}_f$ au point d'abscisse $0$,\\sans utiliser la formule de cours de l'équation de tangente. \end{minipage}

\end{exercice}

\begin{Solution}
 \begin{minipage}[t]{\linewidth} On sait que la tangente n'est pas une droite verticale, puisque la fonction est dérivable sur l'intervalle.\\On en déduit que la tangente $(T)$ au point d'abscisse $0$, admet une équation réduite de la forme :  $(T) : y=m x + p$. 

\medskip
$\bullet$ Détermination de $m$ :\\On sait que le nombre dérivé en $0$ est par définition, le coefficient directeur de la tangente au point d'abscisse $0$.\\Par conséquent, on a déjà : $m=f'(0)=-4$.\\ On en déduit que  $(T) : y= -4 x + p$\\$\bullet$ Détermination de $p$ :\\ Pour cela, on utilise que si $f(0)=-1~~$, alors le point $A$ de coordonnées $(0;-1)$ appartient à $\mathcal{C}_f$ mais aussi à $(T)$.\\ On peut écrire $A(0;-1) = \mathcal{C}_f \cap (T)$.\\ On remplace alors les coordonnées de $A(0;-1)$ dans l'équation  $(T) : y= -4 x + p$.\\ $\begin{aligned} \phantom{\iff}&A(0;-1)\in (T)\\ \iff& -1= -4 \times 0  + p\\ \iff& p=-1 +4 \times 0  \\ \iff& p=-1 +0   \\ \iff& p=-1\\\end{aligned}$\\On peut conclure que : $(T) : {\color[HTML]{f15929}\boldsymbol{y=-4x-1}}$. \end{minipage}

\end{Solution}

% @see : https://coopmaths.fr/alea?uuid=0e984&id=can1F15&n=1&d=10&alea=1Alo22323232323423456789101112131415161718192021222324252627282930313233343536373839404142434445464748234567891011121314151617181920212223242526&cd=1&cols=1
\needspace{10\baselineskip}
\begin{exercice}[BaremeTotal=false]


 La courbe représente une fonction $f$ et la droite est la tangente au point d'abscisse $0$.\\
        Déterminer $f'(0)$.\begin{center}\begin{tikzpicture}[baseline,scale = 0.5]

    \tikzset{
      point/.style={
        thick,
        draw,
        cross out,
        inner sep=0pt,
        minimum width=5pt,
        minimum height=5pt,
      },
    }
    \clip (-10,-2) rectangle (4,12);
    	
	\draw[color ={black},line width = 2,>=latex,->] (-10,0)--(4,0);
	\draw[color ={black},line width = 2,>=latex,->] (0,-2)--(0,12);
	\draw[color ={black},opacity = 0.5] (-10,1)--(4,1);
	\draw[color ={black},opacity = 0.5] (-10,2)--(4,2);
	\draw[color ={black},opacity = 0.5] (-10,3)--(4,3);
	\draw[color ={black},opacity = 0.5] (-10,4)--(4,4);
	\draw[color ={black},opacity = 0.5] (-10,5)--(4,5);
	\draw[color ={black},opacity = 0.5] (-10,6)--(4,6);
	\draw[color ={black},opacity = 0.5] (-10,7)--(4,7);
	\draw[color ={black},opacity = 0.5] (-10,8)--(4,8);
	\draw[color ={black},opacity = 0.5] (-10,9)--(4,9);
	\draw[color ={black},opacity = 0.5] (-10,10)--(4,10);
	\draw[color ={black},opacity = 0.5] (-10,11)--(4,11);
	\draw[color ={black},opacity = 0.5] (-10,12)--(4,12);
	\draw[color ={black},opacity = 0.5] (-10,-1)--(4,-1);
	\draw[color ={black},opacity = 0.5] (-10,-2)--(4,-2);
	\draw[color ={black},opacity = 0.5] (1,-2)--(1,12);
	\draw[color ={black},opacity = 0.5] (2,-2)--(2,12);
	\draw[color ={black},opacity = 0.5] (3,-2)--(3,12);
	\draw[color ={black},opacity = 0.5] (4,-2)--(4,12);
	\draw[color ={black},opacity = 0.5] (-1,-2)--(-1,12);
	\draw[color ={black},opacity = 0.5] (-2,-2)--(-2,12);
	\draw[color ={black},opacity = 0.5] (-3,-2)--(-3,12);
	\draw[color ={black},opacity = 0.5] (-4,-2)--(-4,12);
	\draw[color ={black},opacity = 0.5] (-5,-2)--(-5,12);
	\draw[color ={black},opacity = 0.5] (-6,-2)--(-6,12);
	\draw[color ={black},opacity = 0.5] (-7,-2)--(-7,12);
	\draw[color ={black},opacity = 0.5] (-8,-2)--(-8,12);
	\draw[color ={black},opacity = 0.5] (-9,-2)--(-9,12);
	\draw[color ={black},opacity = 0.5] (-10,-2)--(-10,12);
	\draw[color ={gray},opacity = 0.3] (-10,0)--(4,0);
	\draw[color ={gray},opacity = 0.3] (-10,0)--(4,0);
	\draw[color ={gray},opacity = 0.3] (0,-2)--(0,12);
	\draw[color ={gray},opacity = 0.3] (0,-2)--(0,12);
	\draw[color ={gray},opacity = 0.3] (-5,-2)--(-5,12);
	\draw[color ={gray},opacity = 0.3] (-6,-2)--(-6,12);
	\draw[color ={gray},opacity = 0.3] (-7,-2)--(-7,12);
	\draw[color ={gray},opacity = 0.3] (-8,-2)--(-8,12);
	\draw[color ={gray},opacity = 0.3] (-9,-2)--(-9,12);
	\draw[color ={gray},opacity = 0.3] (-10,-2)--(-10,12);
	\draw[color ={black},line width = 2] (2,-0.2)--(2,0.2);
	\draw[color ={black},line width = 2] (-2,-0.2)--(-2,0.2);
	\draw[color ={black},line width = 2] (-4,-0.2)--(-4,0.2);
	\draw[color ={black},line width = 2] (-6,-0.2)--(-6,0.2);
	\draw[color ={black},line width = 2] (-8,-0.2)--(-8,0.2);
	\draw[color ={black},line width = 2] (-0.2,2)--(0.2,2);
	\draw[color ={black},line width = 2] (-0.2,4)--(0.2,4);
	\draw[color ={black},line width = 2] (-0.2,6)--(0.2,6);
	\draw[color ={black},line width = 2] (-0.2,8)--(0.2,8);
	\draw[color ={black},line width = 2] (-0.2,10)--(0.2,10);
	\draw (2,-0.4) node[anchor = center] {\scriptsize \color{black}{$1$}};
	\draw (-2,-0.4) node[anchor = center] {\scriptsize \color{black}{$-1$}};
	\draw (-4,-0.4) node[anchor = center] {\scriptsize \color{black}{$-2$}};
	\draw (-6,-0.4) node[anchor = center] {\scriptsize \color{black}{$-3$}};
	\draw (-8,-0.4) node[anchor = center] {\scriptsize \color{black}{$-4$}};
	\draw (-0.5,2.1) node[anchor = center] {\footnotesize \color{black}{$1$}};
	\draw (-0.5,4.1) node[anchor = center] {\footnotesize \color{black}{$2$}};
	\draw (-0.5,6.1) node[anchor = center] {\footnotesize \color{black}{$3$}};
	\draw (-0.5,8.1) node[anchor = center] {\footnotesize \color{black}{$4$}};
	\draw (-0.5,10.1) node[anchor = center] {\footnotesize \color{black}{$5$}};
	\draw [color={black}] (-0.3,-0.3) node[anchor = center,scale=1, rotate = 0] {O};
	
	\draw[color={blue},line width = 2] (-2.8,12.93)--(-2.4,9.91)--(-2,7.59)--(-1.6,5.81)--(-1.2,4.45)--(-0.8,3.41)--(-0.4,2.61)--(0,2)--(0.4,1.53)--(0.8,1.17)--(1.2,0.9)--(1.6,0.69)--(2,0.53)--(2.4,0.4)--(2.8,0.31)--(3.2,0.24)--(3.6,0.18)--(4,0.14);
	
	\draw[color={red},line width = 2] (-8.8,13.73)--(-8.4,13.2)--(-8,12.67)--(-7.6,12.13)--(-7.2,11.6)--(-6.8,11.07)--(-6.4,10.53)--(-6,10)--(-5.6,9.47)--(-5.2,8.93)--(-4.8,8.4)--(-4.4,7.87)--(-4,7.33)--(-3.6,6.8)--(-3.2,6.27)--(-2.8,5.73)--(-2.4,5.2)--(-2,4.67)--(-1.6,4.13)--(-1.2,3.6)--(-0.8,3.07)--(-0.4,2.53)--(0,2)--(0.4,1.47)--(0.8,0.93)--(1.2,0.4)--(1.6,-0.13)--(2,-0.67)--(2.4,-1.2)--(2.8,-1.73)--(3.2,-2.27)--(3.6,-2.8)--(4,-3.33);

\end{tikzpicture}\end{center}
\end{exercice}

\begin{Solution}
 $f'(0)$ est donné par le coefficient directeur de la tangente à la courbe au point d'abscisse $0$, soit $\dfrac{-4}{3}=-\dfrac{4}{3}$.
\end{Solution}

% @see : https://coopmaths.fr/alea?uuid=6f32d&id=can1F16&n=1&d=10&alea=dqSw22323232323423456789101112131415161718192021222324252627282930313233343536373839404142434445464748234567891011121314151617181920212223242526&cd=1&cols=1
\needspace{10\baselineskip}
\newpage
\begin{exercice}[BaremeTotal=false]


 On donne les représentations graphiques d'une fonction et de sa dérivée.\\
      Donner l'équation réduite de la tangente à la courbe de $f$ en $x=-1$. \\ \begin{minipage}{0.5\linewidth}
    \begin{tikzpicture}[baseline, scale=0.8]

    \tikzset{
      point/.style={
        thick,
        draw,
        cross out,
        inner sep=0pt,
        minimum width=5pt,
        minimum height=5pt,
      },
    }
    \clip (-5.5,-9.5) rectangle (5.75,7.5);
    	
	\draw[color ={black},>=latex,->] (-4,0)--(4,0);
	\draw[color ={black},>=latex,->] (0,-8)--(0,6);
	\draw[color ={black},opacity = 0.4] (-4,1)--(4,1);
	\draw[color ={black},opacity = 0.4] (-4,2)--(4,2);
	\draw[color ={black},opacity = 0.4] (-4,3)--(4,3);
	\draw[color ={black},opacity = 0.4] (-4,4)--(4,4);
	\draw[color ={black},opacity = 0.4] (-4,5)--(4,5);
	\draw[color ={black},opacity = 0.4] (-4,6)--(4,6);
	\draw[color ={black},opacity = 0.4] (-4,-1)--(4,-1);
	\draw[color ={black},opacity = 0.4] (-4,-2)--(4,-2);
	\draw[color ={black},opacity = 0.4] (-4,-3)--(4,-3);
	\draw[color ={black},opacity = 0.4] (-4,-4)--(4,-4);
	\draw[color ={black},opacity = 0.4] (-4,-5)--(4,-5);
	\draw[color ={black},opacity = 0.4] (-4,-6)--(4,-6);
	\draw[color ={black},opacity = 0.4] (-4,-7)--(4,-7);
	\draw[color ={black},opacity = 0.4] (-4,-8)--(4,-8);
	\draw[color ={black},opacity = 0.4] (1,-8)--(1,6);
	\draw[color ={black},opacity = 0.4] (2,-8)--(2,6);
	\draw[color ={black},opacity = 0.4] (3,-8)--(3,6);
	\draw[color ={black},opacity = 0.4] (4,-8)--(4,6);
	\draw[color ={black},opacity = 0.4] (-1,-8)--(-1,6);
	\draw[color ={black},opacity = 0.4] (-2,-8)--(-2,6);
	\draw[color ={black},opacity = 0.4] (-3,-8)--(-3,6);
	\draw[color ={black},opacity = 0.4] (-4,-8)--(-4,6);
	\draw[color ={gray},opacity = 0.3] (-4,0)--(4,0);
	\draw[color ={gray},opacity = 0.3] (-4,0.5)--(4,0.5);
	\draw[color ={gray},opacity = 0.3] (-4,1.5)--(4,1.5);
	\draw[color ={gray},opacity = 0.3] (-4,2.5)--(4,2.5);
	\draw[color ={gray},opacity = 0.3] (-4,3.5)--(4,3.5);
	\draw[color ={gray},opacity = 0.3] (-4,4.5)--(4,4.5);
	\draw[color ={gray},opacity = 0.3] (-4,5.5)--(4,5.5);
	\draw[color ={gray},opacity = 0.3] (-4,0)--(4,0);
	\draw[color ={gray},opacity = 0.3] (-4,-0.5)--(4,-0.5);
	\draw[color ={gray},opacity = 0.3] (-4,-1.5)--(4,-1.5);
	\draw[color ={gray},opacity = 0.3] (-4,-2.5)--(4,-2.5);
	\draw[color ={gray},opacity = 0.3] (-4,-3.5)--(4,-3.5);
	\draw[color ={gray},opacity = 0.3] (-4,-4.5)--(4,-4.5);
	\draw[color ={gray},opacity = 0.3] (-4,-5.5)--(4,-5.5);
	\draw[color ={gray},opacity = 0.3] (-4,-6.5)--(4,-6.5);
	\draw[color ={gray},opacity = 0.3] (-4,-7)--(4,-7);
	\draw[color ={gray},opacity = 0.3] (-4,-7.5)--(4,-7.5);
	\draw[color ={gray},opacity = 0.3] (-4,-8)--(4,-8);
	\draw[color ={gray},opacity = 0.3] (0,-8)--(0,6);
	\draw[color ={gray},opacity = 0.3] (0,-8)--(0,6);
	\draw[color ={black},line width = 1.2] (1,-0.13)--(1,0.13);
	\draw[color ={black},line width = 1.2] (2,-0.13)--(2,0.13);
	\draw[color ={black},line width = 1.2] (3,-0.13)--(3,0.13);
	\draw[color ={black},line width = 1.2] (4,-0.13)--(4,0.13);
	\draw[color ={black},line width = 1.2] (-1,-0.13)--(-1,0.13);
	\draw[color ={black},line width = 1.2] (-2,-0.13)--(-2,0.13);
	\draw[color ={black},line width = 1.2] (-3,-0.13)--(-3,0.13);
	\draw[color ={black},line width = 1.2] (-4,-0.13)--(-4,0.13);
	\draw[color ={black},line width = 1.2] (-0.13,1)--(0.13,1);
	\draw[color ={black},line width = 1.2] (-0.13,2)--(0.13,2);
	\draw[color ={black},line width = 1.2] (-0.13,3)--(0.13,3);
	\draw[color ={black},line width = 1.2] (-0.13,4)--(0.13,4);
	\draw[color ={black},line width = 1.2] (-0.13,5)--(0.13,5);
	\draw[color ={black},line width = 1.2] (-0.13,6)--(0.13,6);
	\draw[color ={black},line width = 1.2] (-0.13,-1)--(0.13,-1);
	\draw[color ={black},line width = 1.2] (-0.13,-2)--(0.13,-2);
	\draw[color ={black},line width = 1.2] (-0.13,-3)--(0.13,-3);
	\draw[color ={black},line width = 1.2] (-0.13,-4)--(0.13,-4);
	\draw[color ={black},line width = 1.2] (-0.13,-5)--(0.13,-5);
	\draw[color ={black},line width = 1.2] (-0.13,-6)--(0.13,-6);
	\draw[color ={black},line width = 1.2] (-0.13,-7)--(0.13,-7);
	\draw[color ={black},line width = 1.2] (-0.13,-8)--(0.13,-8);
	\draw (2,-0.4) node[anchor = center] {\scriptsize \color{black}{$2$}};
	\draw (4,-0.4) node[anchor = center] {\scriptsize \color{black}{$4$}};
	\draw (-2,-0.4) node[anchor = center] {\scriptsize \color{black}{$-2$}};
	\draw (-4,-0.4) node[anchor = center] {\scriptsize \color{black}{$-4$}};
	\draw (-0.5,2.1) node[anchor = center] {\footnotesize \color{black}{$2$}};
	\draw (-0.5,4.1) node[anchor = center] {\footnotesize \color{black}{$4$}};
	\draw (-0.5,6.1) node[anchor = center] {\footnotesize \color{black}{$6$}};
	\draw (-0.5,-1.9) node[anchor = center] {\footnotesize \color{black}{$-2$}};
	\draw (-0.5,-3.9) node[anchor = center] {\footnotesize \color{black}{$-4$}};
	\draw (-0.5,-5.9) node[anchor = center] {\footnotesize \color{black}{$-6$}};
	\draw (-0.5,-7.9) node[anchor = center] {\footnotesize \color{black}{$-8$}};
	\draw [color={black}] (-0.3,-0.3) node[anchor = center,scale=1, rotate = 0] {O};
	\draw (3,4) node[anchor = center] {\footnotesize \color{blue}{$\Large \mathcal{C}_f$}};
	
	\draw[color={blue},line width = 2] (-4,-2)--(-3.8,-1.24)--(-3.6,-0.56)--(-3.4,0.04)--(-3.2,0.56)--(-3,1)--(-2.8,1.36)--(-2.6,1.64)--(-2.4,1.84)--(-2.2,1.96)--(-2,2)--(-1.8,1.96)--(-1.6,1.84)--(-1.4,1.64)--(-1.2,1.36)--(-1,1)--(-0.8,0.56)--(-0.6,0.04)--(-0.4,-0.56)--(-0.2,-1.24)--(0,-2)--(0.2,-2.84)--(0.4,-3.76)--(0.6,-4.76)--(0.8,-5.84)--(1,-7)--(1.2,-8.24);

\end{tikzpicture}\\
    \end{minipage}
    \begin{minipage}{0.5\linewidth}
    \begin{tikzpicture}[baseline, scale=0.8]

    \tikzset{
      point/.style={
        thick,
        draw,
        cross out,
        inner sep=0pt,
        minimum width=5pt,
        minimum height=5pt,
      },
    }
    \clip (-5.5,-9.5) rectangle (6.5,7.5);
    	
	\draw[color ={black},>=latex,->] (-4,0)--(4,0);
	\draw[color ={black},>=latex,->] (0,-8)--(0,6);
	\draw[color ={black},opacity = 0.4] (-4,1)--(4,1);
	\draw[color ={black},opacity = 0.4] (-4,2)--(4,2);
	\draw[color ={black},opacity = 0.4] (-4,3)--(4,3);
	\draw[color ={black},opacity = 0.4] (-4,4)--(4,4);
	\draw[color ={black},opacity = 0.4] (-4,5)--(4,5);
	\draw[color ={black},opacity = 0.4] (-4,6)--(4,6);
	\draw[color ={black},opacity = 0.4] (-4,-1)--(4,-1);
	\draw[color ={black},opacity = 0.4] (-4,-2)--(4,-2);
	\draw[color ={black},opacity = 0.4] (-4,-3)--(4,-3);
	\draw[color ={black},opacity = 0.4] (-4,-4)--(4,-4);
	\draw[color ={black},opacity = 0.4] (-4,-5)--(4,-5);
	\draw[color ={black},opacity = 0.4] (-4,-6)--(4,-6);
	\draw[color ={black},opacity = 0.4] (-4,-7)--(4,-7);
	\draw[color ={black},opacity = 0.4] (-4,-8)--(4,-8);
	\draw[color ={black},opacity = 0.4] (1,-8)--(1,6);
	\draw[color ={black},opacity = 0.4] (2,-8)--(2,6);
	\draw[color ={black},opacity = 0.4] (3,-8)--(3,6);
	\draw[color ={black},opacity = 0.4] (4,-8)--(4,6);
	\draw[color ={black},opacity = 0.4] (-1,-8)--(-1,6);
	\draw[color ={black},opacity = 0.4] (-2,-8)--(-2,6);
	\draw[color ={black},opacity = 0.4] (-3,-8)--(-3,6);
	\draw[color ={black},opacity = 0.4] (-4,-8)--(-4,6);
	\draw[color ={gray},opacity = 0.3] (-4,0)--(4,0);
	\draw[color ={gray},opacity = 0.3] (-4,0.5)--(4,0.5);
	\draw[color ={gray},opacity = 0.3] (-4,1.5)--(4,1.5);
	\draw[color ={gray},opacity = 0.3] (-4,2.5)--(4,2.5);
	\draw[color ={gray},opacity = 0.3] (-4,3.5)--(4,3.5);
	\draw[color ={gray},opacity = 0.3] (-4,4.5)--(4,4.5);
	\draw[color ={gray},opacity = 0.3] (-4,5.5)--(4,5.5);
	\draw[color ={gray},opacity = 0.3] (-4,0)--(4,0);
	\draw[color ={gray},opacity = 0.3] (-4,-0.5)--(4,-0.5);
	\draw[color ={gray},opacity = 0.3] (-4,-1.5)--(4,-1.5);
	\draw[color ={gray},opacity = 0.3] (-4,-2.5)--(4,-2.5);
	\draw[color ={gray},opacity = 0.3] (-4,-3.5)--(4,-3.5);
	\draw[color ={gray},opacity = 0.3] (-4,-4.5)--(4,-4.5);
	\draw[color ={gray},opacity = 0.3] (-4,-5.5)--(4,-5.5);
	\draw[color ={gray},opacity = 0.3] (-4,-6.5)--(4,-6.5);
	\draw[color ={gray},opacity = 0.3] (-4,-7)--(4,-7);
	\draw[color ={gray},opacity = 0.3] (-4,-7.5)--(4,-7.5);
	\draw[color ={gray},opacity = 0.3] (-4,-8)--(4,-8);
	\draw[color ={gray},opacity = 0.3] (0,-8)--(0,6);
	\draw[color ={gray},opacity = 0.3] (0,-8)--(0,6);
	\draw[color ={black},line width = 1.2] (1,-0.13)--(1,0.13);
	\draw[color ={black},line width = 1.2] (2,-0.13)--(2,0.13);
	\draw[color ={black},line width = 1.2] (3,-0.13)--(3,0.13);
	\draw[color ={black},line width = 1.2] (4,-0.13)--(4,0.13);
	\draw[color ={black},line width = 1.2] (-1,-0.13)--(-1,0.13);
	\draw[color ={black},line width = 1.2] (-2,-0.13)--(-2,0.13);
	\draw[color ={black},line width = 1.2] (-3,-0.13)--(-3,0.13);
	\draw[color ={black},line width = 1.2] (-4,-0.13)--(-4,0.13);
	\draw[color ={black},line width = 1.2] (-0.13,1)--(0.13,1);
	\draw[color ={black},line width = 1.2] (-0.13,2)--(0.13,2);
	\draw[color ={black},line width = 1.2] (-0.13,3)--(0.13,3);
	\draw[color ={black},line width = 1.2] (-0.13,4)--(0.13,4);
	\draw[color ={black},line width = 1.2] (-0.13,5)--(0.13,5);
	\draw[color ={black},line width = 1.2] (-0.13,6)--(0.13,6);
	\draw[color ={black},line width = 1.2] (-0.13,-1)--(0.13,-1);
	\draw[color ={black},line width = 1.2] (-0.13,-2)--(0.13,-2);
	\draw[color ={black},line width = 1.2] (-0.13,-3)--(0.13,-3);
	\draw[color ={black},line width = 1.2] (-0.13,-4)--(0.13,-4);
	\draw[color ={black},line width = 1.2] (-0.13,-5)--(0.13,-5);
	\draw[color ={black},line width = 1.2] (-0.13,-6)--(0.13,-6);
	\draw[color ={black},line width = 1.2] (-0.13,-7)--(0.13,-7);
	\draw[color ={black},line width = 1.2] (-0.13,-8)--(0.13,-8);
	\draw (2,-0.4) node[anchor = center] {\scriptsize \color{black}{$2$}};
	\draw (4,-0.4) node[anchor = center] {\scriptsize \color{black}{$4$}};
	\draw (-2,-0.4) node[anchor = center] {\scriptsize \color{black}{$-2$}};
	\draw (-4,-0.4) node[anchor = center] {\scriptsize \color{black}{$-4$}};
	\draw (-0.5,2.1) node[anchor = center] {\footnotesize \color{black}{$2$}};
	\draw (-0.5,4.1) node[anchor = center] {\footnotesize \color{black}{$4$}};
	\draw (-0.5,-1.9) node[anchor = center] {\footnotesize \color{black}{$-2$}};
	\draw (-0.5,-3.9) node[anchor = center] {\footnotesize \color{black}{$-4$}};
	\draw (-0.5,-5.9) node[anchor = center] {\footnotesize \color{black}{$-6$}};
	\draw (-0.5,-7.9) node[anchor = center] {\footnotesize \color{black}{$-8$}};
	\draw [color={black}] (-0.3,-0.3) node[anchor = center,scale=1, rotate = 0] {O};
	\draw (3,4) node[anchor = center] {\footnotesize \color{red}{$\Large\mathcal{C}_f\prime$}};
	
	\draw[color={red},line width = 2] (-4,4)--(-3.8,3.6)--(-3.6,3.2)--(-3.4,2.8)--(-3.2,2.4)--(-3,2)--(-2.8,1.6)--(-2.6,1.2)--(-2.4,0.8)--(-2.2,0.4)--(-2,0)--(-1.8,-0.4)--(-1.6,-0.8)--(-1.4,-1.2)--(-1.2,-1.6)--(-1,-2)--(-0.8,-2.4)--(-0.6,-2.8)--(-0.4,-3.2)--(-0.2,-3.6)--(0,-4)--(0.2,-4.4)--(0.4,-4.8)--(0.6,-5.2)--(0.8,-5.6)--(1,-6)--(1.2,-6.4)--(1.4,-6.8)--(1.6,-7.2)--(1.8,-7.6)--(2,-8)--(2.2,-8.4)--(2.4,-8.8);

\end{tikzpicture}\\
    \end{minipage}
    
\end{exercice}

\begin{Solution}
 L'équation réduite de la tangente au point d'abscisse $-1$ est  : $y=f'(-1)(x-(-1))+f(-1)$.\\
      On lit graphiquement $f(-1)=1$ et $f'(-1)=-2$.\\
      L'équation réduite de la tangente est donc donnée par :
      $y=-2(x+1)+1$, soit $y=-2x-1$.
\end{Solution}

% @see : https://coopmaths.fr/alea?uuid=a1ba2&id=can1F14&n=1&d=10&alea=rkIp22323232323423456789101112131415161718192021222324252627282930313233343536373839404142434445464748234567891011121314151617181920212223242526&cd=1&cols=1
\needspace{10\baselineskip}
\begin{exercice}[BaremeTotal=false]


 Soit $f$ la fonction définie sur $\mathbb{R}$ par : $f(x)= -4x^2-x-4$.\\
        En calculant la limite du taux d'accroissement, déterminer $f'(-2)$.
\end{exercice}

\begin{Solution}
 $f$ est une fonction polynôme du second degré de la forme $f(x)=ax^2+bx+c$.\\
    La fonction dérivée est donnée par la somme des dérivées des fonctions $u$ et $v$ définies par $u(x)=-4x^2$ et $v(x)=-x-4$.\\
     Comme $u'(x)=-8x$ et $v'(x)=-1$, on obtient  $f'(x)=-8x-1$. \\
     Ainsi, $f'(-2)=-8\times (-2)-1=15$.
\end{Solution}

% @see : https://coopmaths.fr/alea?uuid=3c690&id=can1F13&n=1&d=10&alea=EUDu22323232323423456789101112131415161718192021222324252627282930313233343536373839404142434445464748234567891011121314151617181920212223242526&cd=1&cols=1
\needspace{10\baselineskip}
\begin{exercice}[BaremeTotal=false]


 Déterminer le coefficient directeur de la tangente à la courbe représentative de la fonction carré au point d'abscisse $5$.    
\end{exercice}

\begin{Solution}
 Le coefficient directeur de la tangente au point d'abscisse $a$ est donné par le nombre dérivé $f'(a)$.\\
        La fonction $f$ définie par $f(x)=x^2$ a pour fonction dérivée la fonction $f'$ définie par $f'(x)=2x$.\\
        Comme $f'(5)=2\times 5=10$, le coefficient directeur de la tangente au point d'abscisse $5$ est : $10$. 
\end{Solution}

% @see : https://coopmaths.fr/alea?uuid=ebd89&id=1AN14-1&n=1&d=10&s=3&alea=ocYr22323232323423456789101112131415161718192021222324252627282930313233343536373839404142434445464748234567891011121314151617181920212223242526&cd=0&cols=1
\needspace{10\baselineskip}
\begin{exercice}[BaremeTotal=false]


 Donner la dérivée de la fonction $f$, dérivable pour tout $x\in\R$, définie par  $f(x)=\dfrac{3x+7}{5}$.
\end{exercice}

\begin{Solution}
 L'expression de la dérivée de la fonction $f$ définie par $f(x)=\dfrac{3x+7}{5}$ est : ${\color[HTML]{f15929}\boldsymbol{f'(x)=\dfrac{3}{5}}}$.
\end{Solution}

% @see : https://coopmaths.fr/alea?uuid=ebd89&id=1AN14-1&n=1&d=10&s=4&alea=uAlP22323232323423456789101112131415161718192021222324252627282930313233343536373839404142434445464748234567891011121314151617181920212223242526&cd=0&cols=1
\needspace{10\baselineskip}
\begin{exercice}[BaremeTotal=false]


 Donner la dérivée de la fonction $f$, dérivable pour tout $x\in\R$, définie par  $f(x)=9x^{8}$.
\end{exercice}

\begin{Solution}
 L'expression de la dérivée de la fonction $f$ définie par $f(x)=9x^{8}$ est : ${\color[HTML]{f15929}\boldsymbol{f'(x)=72x^{7}}}$.
\end{Solution}

% @see : https://coopmaths.fr/alea?uuid=ebd89&id=1AN14-1&n=1&d=10&s=7&alea=vjN722323232323423456789101112131415161718192021222324252627282930313233343536373839404142434445464748234567891011121314151617181920212223242526&cd=0&cols=1
\needspace{10\baselineskip}
\begin{exercice}[BaremeTotal=false]


 Donner la dérivée de la fonction $f$, dérivable pour tout $x\in\R^*$, définie par  $f(x)=\dfrac{-4}{9x}$.
\end{exercice}

\begin{Solution}
 L'expression de la dérivée de la fonction $f$ définie par $f(x)=\dfrac{-4}{9x}$ est : ${\color[HTML]{f15929}\boldsymbol{f'(x)=\dfrac{4}{9x^2}}}$.
\end{Solution}

% @see : https://coopmaths.fr/alea?uuid=a83c0&id=1AN14-4&n=1&d=10&s=1&alea=Njr322323232323423456789101112131415161718192021222324252627282930313233343536373839404142434445464748234567891011121314151617181920212223242526&cd=0&cols=1
\needspace{10\baselineskip}
\begin{exercice}[BaremeTotal=false]


 Donner l'expression de la dérivée de la fonction $f$ définie sur $\R^{*}$ par $f(x)=\dfrac{9}{x}+4x^{2}$.\\
\end{exercice}

\begin{Solution}
 L'expression de la dérivée de la fonction $f$ définie par $f(x)=\dfrac{9}{x}+4x^{2}$ est :\\$f'(x)={\color[HTML]{f15929}\boldsymbol{-\dfrac{9}{x^2}+8x}}$.\\
\end{Solution}

% @see : https://coopmaths.fr/alea?uuid=a83c0&id=1AN14-4&n=1&d=10&s=4&alea=4Vpz22323232323423456789101112131415161718192021222324252627282930313233343536373839404142434445464748234567891011121314151617181920212223242526&cd=1&cols=1
\needspace{10\baselineskip}
\begin{exercice}[BaremeTotal=false]


 Donner l'expression de la dérivée de la fonction $f$ définie sur $\R_+^{*}$ par $f(x)=\dfrac{5}{x}+2\sqrt{x}-4x$.\\
\end{exercice}

\begin{Solution}
 $(\dfrac{5}{x})^\prime=-\dfrac{5}{x^2}$\\$(2\sqrt{x})^\prime=\dfrac{2}{2\sqrt{x}}$\\$(-4x)^\prime=-4$\\L'expression de la dérivée de la fonction $f$ définie par $f(x)=\dfrac{5}{x}+2\sqrt{x}-4x$ est :\\$f'(x)={\color[HTML]{f15929}\boldsymbol{-\dfrac{5}{x^2}+\dfrac{2}{2\sqrt{x}}-4}}$.\\
\end{Solution}

\end{Maquette}
\clearpage
\begin{Maquette}[DM]{Niveau= ,Classe=\getdata[27]\mydata\ (v27), Date = 06/01/25  ,Theme=Exercices}


% @see : https://coopmaths.fr/alea?uuid=4c8c7&id=1AN11-3&n=1&d=10&s=2&alea=nawO2232323232342345678910111213141516171819202122232425262728293031323334353637383940414243444546474823456789101112131415161718192021222324252627&cd=1&cols=2
\needspace{10\baselineskip}
\begin{exercice}[BaremeTotal=false]
\label{eleve:27}


 \begin{minipage}[t]{\linewidth} Soit $f$ une fonction dérivable sur $[-5;5]$ et $\mathcal{C}_f$ sa courbe représentative.\\On sait que  $f(4)=4~~$ et que $~~f'(4)=-2$.\\Déterminer une équation de la tangente $(T)$ à la courbe $\mathcal{C}_f$ au point d'abscisse $4$,\\sans utiliser la formule de cours de l'équation de tangente. \end{minipage}

\end{exercice}

\begin{Solution}
 \begin{minipage}[t]{\linewidth} On sait que la tangente n'est pas une droite verticale, puisque la fonction est dérivable sur l'intervalle.\\On en déduit que la tangente $(T)$ au point d'abscisse $4$, admet une équation réduite de la forme :  $(T) : y=m x + p$. 

\medskip
$\bullet$ Détermination de $m$ :\\On sait que le nombre dérivé en $4$ est par définition, le coefficient directeur de la tangente au point d'abscisse $4$.\\Par conséquent, on a déjà : $m=f'(4)=-2$.\\ On en déduit que  $(T) : y= -2 x + p$\\$\bullet$ Détermination de $p$ :\\ Pour cela, on utilise que si $f(4)=4~~$, alors le point $A$ de coordonnées $(4;4)$ appartient à $\mathcal{C}_f$ mais aussi à $(T)$.\\ On peut écrire $A(4;4) = \mathcal{C}_f \cap (T)$.\\ On remplace alors les coordonnées de $A(4;4)$ dans l'équation  $(T) : y= -2 x + p$.\\ $\begin{aligned} \phantom{\iff}&A(4;4)\in (T)\\ \iff& 4= -2 \times 4  + p\\ \iff& p=4 +2 \times 4  \\ \iff& p=4 +8   \\ \iff& p=12\\\end{aligned}$\\On peut conclure que : $(T) : {\color[HTML]{f15929}\boldsymbol{y=-2x+12}}$. \end{minipage}

\end{Solution}

% @see : https://coopmaths.fr/alea?uuid=0e984&id=can1F15&n=1&d=10&alea=1Alo2232323232342345678910111213141516171819202122232425262728293031323334353637383940414243444546474823456789101112131415161718192021222324252627&cd=1&cols=1
\needspace{10\baselineskip}
\begin{exercice}[BaremeTotal=false]


 La courbe représente une fonction $f$ et la droite est la tangente au point d'abscisse $-1$.\\
        Déterminer $f'(-1)$.\begin{center}\begin{tikzpicture}[baseline,scale = 0.5]

    \tikzset{
      point/.style={
        thick,
        draw,
        cross out,
        inner sep=0pt,
        minimum width=5pt,
        minimum height=5pt,
      },
    }
    \clip (-8,-10) rectangle (8,3);
    	
	\draw[color ={black},line width = 2,>=latex,->] (-8,0)--(8,0);
	\draw[color ={black},line width = 2,>=latex,->] (0,-10)--(0,3);
	\draw[color ={black},opacity = 0.5] (-8,1)--(8,1);
	\draw[color ={black},opacity = 0.5] (-8,2)--(8,2);
	\draw[color ={black},opacity = 0.5] (-8,3)--(8,3);
	\draw[color ={black},opacity = 0.5] (-8,-1)--(8,-1);
	\draw[color ={black},opacity = 0.5] (-8,-2)--(8,-2);
	\draw[color ={black},opacity = 0.5] (-8,-3)--(8,-3);
	\draw[color ={black},opacity = 0.5] (-8,-4)--(8,-4);
	\draw[color ={black},opacity = 0.5] (-8,-5)--(8,-5);
	\draw[color ={black},opacity = 0.5] (-8,-6)--(8,-6);
	\draw[color ={black},opacity = 0.5] (-8,-7)--(8,-7);
	\draw[color ={black},opacity = 0.5] (-8,-8)--(8,-8);
	\draw[color ={black},opacity = 0.5] (-8,-9)--(8,-9);
	\draw[color ={black},opacity = 0.5] (-8,-10)--(8,-10);
	\draw[color ={black},opacity = 0.5] (1,-10)--(1,3);
	\draw[color ={black},opacity = 0.5] (2,-10)--(2,3);
	\draw[color ={black},opacity = 0.5] (3,-10)--(3,3);
	\draw[color ={black},opacity = 0.5] (4,-10)--(4,3);
	\draw[color ={black},opacity = 0.5] (5,-10)--(5,3);
	\draw[color ={black},opacity = 0.5] (6,-10)--(6,3);
	\draw[color ={black},opacity = 0.5] (7,-10)--(7,3);
	\draw[color ={black},opacity = 0.5] (8,-10)--(8,3);
	\draw[color ={black},opacity = 0.5] (-1,-10)--(-1,3);
	\draw[color ={black},opacity = 0.5] (-2,-10)--(-2,3);
	\draw[color ={black},opacity = 0.5] (-3,-10)--(-3,3);
	\draw[color ={black},opacity = 0.5] (-4,-10)--(-4,3);
	\draw[color ={black},opacity = 0.5] (-5,-10)--(-5,3);
	\draw[color ={black},opacity = 0.5] (-6,-10)--(-6,3);
	\draw[color ={black},opacity = 0.5] (-7,-10)--(-7,3);
	\draw[color ={black},opacity = 0.5] (-8,-10)--(-8,3);
	\draw[color ={gray},opacity = 0.3] (-8,0)--(8,0);
	\draw[color ={gray},opacity = 0.3] (-8,0)--(8,0);
	\draw[color ={gray},opacity = 0.3] (-8,-4)--(8,-4);
	\draw[color ={gray},opacity = 0.3] (-8,-5)--(8,-5);
	\draw[color ={gray},opacity = 0.3] (-8,-6)--(8,-6);
	\draw[color ={gray},opacity = 0.3] (-8,-7)--(8,-7);
	\draw[color ={gray},opacity = 0.3] (-8,-8)--(8,-8);
	\draw[color ={gray},opacity = 0.3] (-8,-9)--(8,-9);
	\draw[color ={gray},opacity = 0.3] (-8,-10)--(8,-10);
	\draw[color ={gray},opacity = 0.3] (0,-10)--(0,3);
	\draw[color ={gray},opacity = 0.3] (0,-10)--(0,3);
	\draw[color ={black},line width = 2] (2,-0.2)--(2,0.2);
	\draw[color ={black},line width = 2] (4,-0.2)--(4,0.2);
	\draw[color ={black},line width = 2] (6,-0.2)--(6,0.2);
	\draw[color ={black},line width = 2] (-2,-0.2)--(-2,0.2);
	\draw[color ={black},line width = 2] (-4,-0.2)--(-4,0.2);
	\draw[color ={black},line width = 2] (-6,-0.2)--(-6,0.2);
	\draw[color ={black},line width = 2] (-0.2,1)--(0.2,1);
	\draw[color ={black},line width = 2] (-0.2,2)--(0.2,2);
	\draw[color ={black},line width = 2] (-0.2,-1)--(0.2,-1);
	\draw[color ={black},line width = 2] (-0.2,-2)--(0.2,-2);
	\draw[color ={black},line width = 2] (-0.2,-3)--(0.2,-3);
	\draw[color ={black},line width = 2] (-0.2,-4)--(0.2,-4);
	\draw[color ={black},line width = 2] (-0.2,-5)--(0.2,-5);
	\draw[color ={black},line width = 2] (-0.2,-6)--(0.2,-6);
	\draw[color ={black},line width = 2] (-0.2,-7)--(0.2,-7);
	\draw[color ={black},line width = 2] (-0.2,-8)--(0.2,-8);
	\draw[color ={black},line width = 2] (-0.2,-9)--(0.2,-9);
	\draw (2,-0.4) node[anchor = center] {\scriptsize \color{black}{$1$}};
	\draw (4,-0.4) node[anchor = center] {\scriptsize \color{black}{$2$}};
	\draw (6,-0.4) node[anchor = center] {\scriptsize \color{black}{$3$}};
	\draw (-2,-0.4) node[anchor = center] {\scriptsize \color{black}{$-1$}};
	\draw (-4,-0.4) node[anchor = center] {\scriptsize \color{black}{$-2$}};
	\draw (-6,-0.4) node[anchor = center] {\scriptsize \color{black}{$-3$}};
	\draw (-0.8,2.1) node[anchor = center] {\footnotesize \color{black}{$2$}};
	\draw (-0.8,-1.9) node[anchor = center] {\footnotesize \color{black}{$-2$}};
	\draw (-0.8,-3.9) node[anchor = center] {\footnotesize \color{black}{$-4$}};
	\draw (-0.8,-5.9) node[anchor = center] {\footnotesize \color{black}{$-6$}};
	\draw (-0.8,-7.9) node[anchor = center] {\footnotesize \color{black}{$-8$}};
	\draw [color={black}] (-0.3,-0.3) node[anchor = center,scale=1, rotate = 0] {O};
	
	\draw[color={blue},line width = 2] (-5.2,-10.96)--(-4.8,-9.56)--(-4.4,-8.24)--(-4,-7)--(-3.6,-5.84)--(-3.2,-4.76)--(-2.8,-3.76)--(-2.4,-2.84)--(-2,-2)--(-1.6,-1.24)--(-1.2,-0.56)--(-0.8,0.04)--(-0.4,0.56)--(0,1)--(0.4,1.36)--(0.8,1.64)--(1.2,1.84)--(1.6,1.96)--(2,2)--(2.4,1.96)--(2.8,1.84)--(3.2,1.64)--(3.6,1.36)--(4,1)--(4.4,0.56)--(4.8,0.04)--(5.2,-0.56)--(5.6,-1.24)--(6,-2)--(6.4,-2.84)--(6.8,-3.76)--(7.2,-4.76)--(7.6,-5.84)--(8,-7);
	
	\draw[color={red},line width = 2] (-6.4,-10.8)--(-6,-10)--(-5.6,-9.2)--(-5.2,-8.4)--(-4.8,-7.6)--(-4.4,-6.8)--(-4,-6)--(-3.6,-5.2)--(-3.2,-4.4)--(-2.8,-3.6)--(-2.4,-2.8)--(-2,-2)--(-1.6,-1.2)--(-1.2,-0.4)--(-0.8,0.4)--(-0.4,1.2)--(0,2)--(0.4,2.8)--(0.8,3.6);

\end{tikzpicture}\end{center}
\end{exercice}

\begin{Solution}
 $f'(-1)$ est donné par le coefficient directeur de la tangente à la courbe au point d'abscisse $-1$, soit $4$.
\end{Solution}

% @see : https://coopmaths.fr/alea?uuid=6f32d&id=can1F16&n=1&d=10&alea=dqSw2232323232342345678910111213141516171819202122232425262728293031323334353637383940414243444546474823456789101112131415161718192021222324252627&cd=1&cols=1
\needspace{10\baselineskip}
\newpage
\begin{exercice}[BaremeTotal=false]


 On donne les représentations graphiques d'une fonction et de sa dérivée.\\
        Donner l'équation réduite de la tangente à la courbe de $f$ en $x=1$. \\ \begin{minipage}{0.5\linewidth}
    \begin{tikzpicture}[baseline, scale=0.8]

    \tikzset{
      point/.style={
        thick,
        draw,
        cross out,
        inner sep=0pt,
        minimum width=5pt,
        minimum height=5pt,
      },
    }
    \clip (-4.5,-4.5) rectangle (5.75,13.5);
    	
	\draw[color ={black},>=latex,->] (-3,0)--(3,0);
	\draw[color ={black},>=latex,->] (0,-3)--(0,12);
	\draw[color ={black},opacity = 0.4] (-3,1)--(3,1);
	\draw[color ={black},opacity = 0.4] (-3,2)--(3,2);
	\draw[color ={black},opacity = 0.4] (-3,3)--(3,3);
	\draw[color ={black},opacity = 0.4] (-3,4)--(3,4);
	\draw[color ={black},opacity = 0.4] (-3,5)--(3,5);
	\draw[color ={black},opacity = 0.4] (-3,6)--(3,6);
	\draw[color ={black},opacity = 0.4] (-3,7)--(3,7);
	\draw[color ={black},opacity = 0.4] (-3,8)--(3,8);
	\draw[color ={black},opacity = 0.4] (-3,9)--(3,9);
	\draw[color ={black},opacity = 0.4] (-3,10)--(3,10);
	\draw[color ={black},opacity = 0.4] (-3,11)--(3,11);
	\draw[color ={black},opacity = 0.4] (-3,12)--(3,12);
	\draw[color ={black},opacity = 0.4] (-3,-1)--(3,-1);
	\draw[color ={black},opacity = 0.4] (-3,-2)--(3,-2);
	\draw[color ={black},opacity = 0.4] (-3,-3)--(3,-3);
	\draw[color ={black},opacity = 0.4] (1,-3)--(1,12);
	\draw[color ={black},opacity = 0.4] (2,-3)--(2,12);
	\draw[color ={black},opacity = 0.4] (3,-3)--(3,12);
	\draw[color ={black},opacity = 0.4] (-1,-3)--(-1,12);
	\draw[color ={black},opacity = 0.4] (-2,-3)--(-2,12);
	\draw[color ={black},opacity = 0.4] (-3,-3)--(-3,12);
	\draw[color ={gray},opacity = 0.3] (-3,0)--(3,0);
	\draw[color ={gray},opacity = 0.3] (-3,0.5)--(3,0.5);
	\draw[color ={gray},opacity = 0.3] (-3,1.5)--(3,1.5);
	\draw[color ={gray},opacity = 0.3] (-3,2.5)--(3,2.5);
	\draw[color ={gray},opacity = 0.3] (-3,3.5)--(3,3.5);
	\draw[color ={gray},opacity = 0.3] (-3,4.5)--(3,4.5);
	\draw[color ={gray},opacity = 0.3] (-3,5.5)--(3,5.5);
	\draw[color ={gray},opacity = 0.3] (-3,6.5)--(3,6.5);
	\draw[color ={gray},opacity = 0.3] (-3,7.5)--(3,7.5);
	\draw[color ={gray},opacity = 0.3] (-3,8.5)--(3,8.5);
	\draw[color ={gray},opacity = 0.3] (-3,9.5)--(3,9.5);
	\draw[color ={gray},opacity = 0.3] (-3,10.5)--(3,10.5);
	\draw[color ={gray},opacity = 0.3] (-3,11.5)--(3,11.5);
	\draw[color ={gray},opacity = 0.3] (-3,0)--(3,0);
	\draw[color ={gray},opacity = 0.3] (-3,-0.5)--(3,-0.5);
	\draw[color ={gray},opacity = 0.3] (-3,-1.5)--(3,-1.5);
	\draw[color ={gray},opacity = 0.3] (-3,-2.5)--(3,-2.5);
	\draw[color ={gray},opacity = 0.3] (0,-3)--(0,12);
	\draw[color ={gray},opacity = 0.3] (0,-3)--(0,12);
	\draw[color ={black},line width = 1.2] (1,-0.13)--(1,0.13);
	\draw[color ={black},line width = 1.2] (2,-0.13)--(2,0.13);
	\draw[color ={black},line width = 1.2] (3,-0.13)--(3,0.13);
	\draw[color ={black},line width = 1.2] (-1,-0.13)--(-1,0.13);
	\draw[color ={black},line width = 1.2] (-2,-0.13)--(-2,0.13);
	\draw[color ={black},line width = 1.2] (-3,-0.13)--(-3,0.13);
	\draw[color ={black},line width = 1.2] (-0.13,1)--(0.13,1);
	\draw[color ={black},line width = 1.2] (-0.13,2)--(0.13,2);
	\draw[color ={black},line width = 1.2] (-0.13,3)--(0.13,3);
	\draw[color ={black},line width = 1.2] (-0.13,4)--(0.13,4);
	\draw[color ={black},line width = 1.2] (-0.13,5)--(0.13,5);
	\draw[color ={black},line width = 1.2] (-0.13,6)--(0.13,6);
	\draw[color ={black},line width = 1.2] (-0.13,7)--(0.13,7);
	\draw[color ={black},line width = 1.2] (-0.13,8)--(0.13,8);
	\draw[color ={black},line width = 1.2] (-0.13,9)--(0.13,9);
	\draw[color ={black},line width = 1.2] (-0.13,10)--(0.13,10);
	\draw[color ={black},line width = 1.2] (-0.13,11)--(0.13,11);
	\draw[color ={black},line width = 1.2] (-0.13,12)--(0.13,12);
	\draw[color ={black},line width = 1.2] (-0.13,-1)--(0.13,-1);
	\draw[color ={black},line width = 1.2] (-0.13,-2)--(0.13,-2);
	\draw[color ={black},line width = 1.2] (-0.13,-3)--(0.13,-3);
	\draw (3,-0.4) node[anchor = center] {\scriptsize \color{black}{$3$}};
	\draw (-3,-0.4) node[anchor = center] {\scriptsize \color{black}{$-3$}};
	\draw (-0.5,3.1) node[anchor = center] {\footnotesize \color{black}{$3$}};
	\draw (-0.5,6.1) node[anchor = center] {\footnotesize \color{black}{$6$}};
	\draw (-0.5,9.1) node[anchor = center] {\footnotesize \color{black}{$9$}};
	\draw (-0.5,-2.9) node[anchor = center] {\footnotesize \color{black}{$-3$}};
	\draw [color={black}] (-0.3,-0.3) node[anchor = center,scale=1, rotate = 0] {O};
	\draw (3,10) node[anchor = center] {\footnotesize \color{blue}{$\Large \mathcal{C}_f$}};
	
	\draw[color={blue},line width = 2] (-3,9)--(-2.8,7.84)--(-2.6,6.76)--(-2.4,5.76)--(-2.2,4.84)--(-2,4)--(-1.8,3.24)--(-1.6,2.56)--(-1.4,1.96)--(-1.2,1.44)--(-1,1)--(-0.8,0.64)--(-0.6,0.36)--(-0.4,0.16)--(-0.2,0.04)--(0,0)--(0.2,0.04)--(0.4,0.16)--(0.6,0.36)--(0.8,0.64)--(1,1)--(1.2,1.44)--(1.4,1.96)--(1.6,2.56)--(1.8,3.24)--(2,4)--(2.2,4.84)--(2.4,5.76)--(2.6,6.76)--(2.8,7.84)--(3,9);

\end{tikzpicture}\\
    \end{minipage}
    \begin{minipage}{0.5\linewidth}
    \begin{tikzpicture}[baseline, scale=0.8]

    \tikzset{
      point/.style={
        thick,
        draw,
        cross out,
        inner sep=0pt,
        minimum width=5pt,
        minimum height=5pt,
      },
    }
    \clip (-4.5,-6.5) rectangle (6.5,9.5);
    	
	\draw[color ={black},>=latex,->] (-3,0)--(3,0);
	\draw[color ={black},>=latex,->] (0,-5)--(0,8);
	\draw[color ={black},opacity = 0.4] (-3,1)--(3,1);
	\draw[color ={black},opacity = 0.4] (-3,2)--(3,2);
	\draw[color ={black},opacity = 0.4] (-3,3)--(3,3);
	\draw[color ={black},opacity = 0.4] (-3,4)--(3,4);
	\draw[color ={black},opacity = 0.4] (-3,5)--(3,5);
	\draw[color ={black},opacity = 0.4] (-3,6)--(3,6);
	\draw[color ={black},opacity = 0.4] (-3,7)--(3,7);
	\draw[color ={black},opacity = 0.4] (-3,8)--(3,8);
	\draw[color ={black},opacity = 0.4] (-3,-1)--(3,-1);
	\draw[color ={black},opacity = 0.4] (-3,-2)--(3,-2);
	\draw[color ={black},opacity = 0.4] (-3,-3)--(3,-3);
	\draw[color ={black},opacity = 0.4] (-3,-4)--(3,-4);
	\draw[color ={black},opacity = 0.4] (-3,-5)--(3,-5);
	\draw[color ={black},opacity = 0.4] (1,-5)--(1,8);
	\draw[color ={black},opacity = 0.4] (2,-5)--(2,8);
	\draw[color ={black},opacity = 0.4] (3,-5)--(3,8);
	\draw[color ={black},opacity = 0.4] (-1,-5)--(-1,8);
	\draw[color ={black},opacity = 0.4] (-2,-5)--(-2,8);
	\draw[color ={black},opacity = 0.4] (-3,-5)--(-3,8);
	\draw[color ={gray},opacity = 0.3] (-3,0)--(3,0);
	\draw[color ={gray},opacity = 0.3] (-3,0.5)--(3,0.5);
	\draw[color ={gray},opacity = 0.3] (-3,1.5)--(3,1.5);
	\draw[color ={gray},opacity = 0.3] (-3,2.5)--(3,2.5);
	\draw[color ={gray},opacity = 0.3] (-3,3.5)--(3,3.5);
	\draw[color ={gray},opacity = 0.3] (-3,4.5)--(3,4.5);
	\draw[color ={gray},opacity = 0.3] (-3,5.5)--(3,5.5);
	\draw[color ={gray},opacity = 0.3] (-3,6.5)--(3,6.5);
	\draw[color ={gray},opacity = 0.3] (-3,7.5)--(3,7.5);
	\draw[color ={gray},opacity = 0.3] (-3,0)--(3,0);
	\draw[color ={gray},opacity = 0.3] (-3,-0.5)--(3,-0.5);
	\draw[color ={gray},opacity = 0.3] (-3,-1.5)--(3,-1.5);
	\draw[color ={gray},opacity = 0.3] (-3,-2.5)--(3,-2.5);
	\draw[color ={gray},opacity = 0.3] (-3,-3.5)--(3,-3.5);
	\draw[color ={gray},opacity = 0.3] (-3,-4.5)--(3,-4.5);
	\draw[color ={gray},opacity = 0.3] (0,-5)--(0,8);
	\draw[color ={gray},opacity = 0.3] (0,-5)--(0,8);
	\draw[color ={black},line width = 1.2] (1,-0.13)--(1,0.13);
	\draw[color ={black},line width = 1.2] (2,-0.13)--(2,0.13);
	\draw[color ={black},line width = 1.2] (3,-0.13)--(3,0.13);
	\draw[color ={black},line width = 1.2] (-1,-0.13)--(-1,0.13);
	\draw[color ={black},line width = 1.2] (-2,-0.13)--(-2,0.13);
	\draw[color ={black},line width = 1.2] (-3,-0.13)--(-3,0.13);
	\draw[color ={black},line width = 1.2] (-0.13,1)--(0.13,1);
	\draw[color ={black},line width = 1.2] (-0.13,2)--(0.13,2);
	\draw[color ={black},line width = 1.2] (-0.13,3)--(0.13,3);
	\draw[color ={black},line width = 1.2] (-0.13,4)--(0.13,4);
	\draw[color ={black},line width = 1.2] (-0.13,5)--(0.13,5);
	\draw[color ={black},line width = 1.2] (-0.13,6)--(0.13,6);
	\draw[color ={black},line width = 1.2] (-0.13,7)--(0.13,7);
	\draw[color ={black},line width = 1.2] (-0.13,8)--(0.13,8);
	\draw[color ={black},line width = 1.2] (-0.13,9)--(0.13,9);
	\draw[color ={black},line width = 1.2] (-0.13,10)--(0.13,10);
	\draw[color ={black},line width = 1.2] (-0.13,11)--(0.13,11);
	\draw[color ={black},line width = 1.2] (-0.13,12)--(0.13,12);
	\draw[color ={black},line width = 1.2] (-0.13,-1)--(0.13,-1);
	\draw[color ={black},line width = 1.2] (-0.13,-2)--(0.13,-2);
	\draw[color ={black},line width = 1.2] (-0.13,-3)--(0.13,-3);
	\draw (3,-0.4) node[anchor = center] {\scriptsize \color{black}{$3$}};
	\draw (-3,-0.4) node[anchor = center] {\scriptsize \color{black}{$-3$}};
	\draw (-0.5,3.1) node[anchor = center] {\footnotesize \color{black}{$3$}};
	\draw (-0.5,6.1) node[anchor = center] {\footnotesize \color{black}{$6$}};
	\draw (-0.5,-2.9) node[anchor = center] {\footnotesize \color{black}{$-3$}};
	\draw [color={black}] (-0.3,-0.3) node[anchor = center,scale=1, rotate = 0] {O};
	\draw (3,6) node[anchor = center] {\footnotesize \color{red}{$\Large\mathcal{C}_f\prime$}};
	
	\draw[color={red},line width = 2] (-2.8,-5.6)--(-2.6,-5.2)--(-2.4,-4.8)--(-2.2,-4.4)--(-2,-4)--(-1.8,-3.6)--(-1.6,-3.2)--(-1.4,-2.8)--(-1.2,-2.4)--(-1,-2)--(-0.8,-1.6)--(-0.6,-1.2)--(-0.4,-0.8)--(-0.2,-0.4)--(0,0)--(0.2,0.4)--(0.4,0.8)--(0.6,1.2)--(0.8,1.6)--(1,2)--(1.2,2.4)--(1.4,2.8)--(1.6,3.2)--(1.8,3.6)--(2,4)--(2.2,4.4)--(2.4,4.8)--(2.6,5.2)--(2.8,5.6)--(3,6);

\end{tikzpicture}\\
    \end{minipage}
    
\end{exercice}

\begin{Solution}
 L'équation réduite de la tangente au point d'abscisse $1$ est  : $y=f'(1)(x-1)+f(1)$.\\
        On lit graphiquement $f(1)=1$ et $f'(1)=2$.\\
        L'équation réduite de la tangente est donc donnée par :
        $y=2(x-1)+1$, soit $y=2x-1$.
\end{Solution}

% @see : https://coopmaths.fr/alea?uuid=a1ba2&id=can1F14&n=1&d=10&alea=rkIp2232323232342345678910111213141516171819202122232425262728293031323334353637383940414243444546474823456789101112131415161718192021222324252627&cd=1&cols=1
\needspace{10\baselineskip}
\begin{exercice}[BaremeTotal=false]


Soit $f$ la fonction définie sur $\mathbb{R}$ par : 
$f(x)= 6x^2-3$.\\
En calculant la limite du taux d'accroissement, déterminer $f'(5)$.
\end{exercice}

\begin{Solution}
 $f$ est une fonction polynôme du second degré de la forme $f(x)=ax^2+b$.\\
    La fonction dérivée est donnée par la somme des dérivées des fonctions $u$ et $v$ définies par $u(x)=6x^2$ et $v(x)=-3$.\\
     Comme $u'(x)=12x$ et $v'(x)=0$, on obtient  $f'(x)=12x$. \\
     Ainsi, $f'(5)=12\times 5=60$.
\end{Solution}

% @see : https://coopmaths.fr/alea?uuid=3c690&id=can1F13&n=1&d=10&alea=EUDu2232323232342345678910111213141516171819202122232425262728293031323334353637383940414243444546474823456789101112131415161718192021222324252627&cd=1&cols=1
\needspace{10\baselineskip}
\begin{exercice}[BaremeTotal=false]


 Déterminer le coefficient directeur de la tangente à la courbe représentative de la fonction inverse au point d'abscisse $6$.
                 
\end{exercice}

\begin{Solution}
 Le coefficient directeur de la tangente au point d'abscisse $a$ est donné par le nombre dérivé $f'(a)$.\\
        La fonction $f$ définie par $f(x)=\dfrac{1}{x}$ a pour fonction dérivée la fonction $f'$ définie par $f'(x)=-\dfrac{1}{x^2}$.\\
Comme $f'(6)=-\dfrac{1}{6^2}=-\dfrac{1}{36}$, le coefficient directeur de la tangente au point d'abscisse $6$ est : $-\dfrac{1}{36}$.
\end{Solution}

% @see : https://coopmaths.fr/alea?uuid=ebd89&id=1AN14-1&n=1&d=10&s=3&alea=ocYr2232323232342345678910111213141516171819202122232425262728293031323334353637383940414243444546474823456789101112131415161718192021222324252627&cd=0&cols=1
\needspace{10\baselineskip}
\begin{exercice}[BaremeTotal=false]


 Donner la dérivée de la fonction $f$, dérivable pour tout $x\in\R$, définie par  $f(x)=\dfrac{2}{9}x+5$.
\end{exercice}

\begin{Solution}
 L'expression de la dérivée de la fonction $f$ définie par $f(x)=\dfrac{2}{9}x+5$ est : ${\color[HTML]{f15929}\boldsymbol{f'(x)=\dfrac{2}{9}}}$.
\end{Solution}

% @see : https://coopmaths.fr/alea?uuid=ebd89&id=1AN14-1&n=1&d=10&s=4&alea=uAlP2232323232342345678910111213141516171819202122232425262728293031323334353637383940414243444546474823456789101112131415161718192021222324252627&cd=0&cols=1
\needspace{10\baselineskip}
\begin{exercice}[BaremeTotal=false]


 Donner la dérivée de la fonction $f$, dérivable pour tout $x\in\R$, définie par  $f(x)=-8x^{11}$.
\end{exercice}

\begin{Solution}
 L'expression de la dérivée de la fonction $f$ définie par $f(x)=-8x^{11}$ est : ${\color[HTML]{f15929}\boldsymbol{f'(x)=-88x^{10}}}$.
\end{Solution}

% @see : https://coopmaths.fr/alea?uuid=ebd89&id=1AN14-1&n=1&d=10&s=7&alea=vjN72232323232342345678910111213141516171819202122232425262728293031323334353637383940414243444546474823456789101112131415161718192021222324252627&cd=0&cols=1
\needspace{10\baselineskip}
\begin{exercice}[BaremeTotal=false]


 Donner la dérivée de la fonction $f$, dérivable pour tout $x\in\R^*$, définie par  $f(x)=\dfrac{-1}{9x}$.
\end{exercice}

\begin{Solution}
 L'expression de la dérivée de la fonction $f$ définie par $f(x)=\dfrac{-1}{9x}$ est : ${\color[HTML]{f15929}\boldsymbol{f'(x)=\dfrac{1}{9x^2}}}$.
\end{Solution}

% @see : https://coopmaths.fr/alea?uuid=a83c0&id=1AN14-4&n=1&d=10&s=1&alea=Njr32232323232342345678910111213141516171819202122232425262728293031323334353637383940414243444546474823456789101112131415161718192021222324252627&cd=0&cols=1
\needspace{10\baselineskip}
\begin{exercice}[BaremeTotal=false]


 Donner l'expression de la dérivée de la fonction $f$ définie sur $\R^{*}$ par $f(x)=5x^{2}+2x-\dfrac{8}{x}$.\\
\end{exercice}

\begin{Solution}
 L'expression de la dérivée de la fonction $f$ définie par $f(x)=5x^{2}+2x-\dfrac{8}{x}$ est :\\$f'(x)={\color[HTML]{f15929}\boldsymbol{10x+2+\dfrac{8}{x^2}}}$.\\
\end{Solution}

% @see : https://coopmaths.fr/alea?uuid=a83c0&id=1AN14-4&n=1&d=10&s=4&alea=4Vpz2232323232342345678910111213141516171819202122232425262728293031323334353637383940414243444546474823456789101112131415161718192021222324252627&cd=1&cols=1
\needspace{10\baselineskip}
\begin{exercice}[BaremeTotal=false]


 Donner l'expression de la dérivée de la fonction $f$ définie sur $\R_+^{*}$ par $f(x)=-2x^{2}-2x+3-5\sqrt{x}-\dfrac{2}{x}$.\\
\end{exercice}

\begin{Solution}
 $(-2x^{2}-2x+3)^\prime=-4x-2$\\$(-5\sqrt{x})^\prime=-\dfrac{5}{2\sqrt{x}}$\\$(-\dfrac{2}{x})^\prime=\dfrac{2}{x^2}$\\L'expression de la dérivée de la fonction $f$ définie par $f(x)=-2x^{2}-2x+3-5\sqrt{x}-\dfrac{2}{x}$ est :\\$f'(x)={\color[HTML]{f15929}\boldsymbol{-4x-2-\dfrac{5}{2\sqrt{x}}+\dfrac{2}{x^2}}}$.\\
\end{Solution}

\end{Maquette}
\clearpage
\begin{Maquette}[DM]{Niveau= ,Classe=\getdata[28]\mydata\ (v28), Date = 06/01/25  ,Theme=Exercices}


% @see : https://coopmaths.fr/alea?uuid=4c8c7&id=1AN11-3&n=1&d=10&s=2&alea=nawO223232323234234567891011121314151617181920212223242526272829303132333435363738394041424344454647482345678910111213141516171819202122232425262728&cd=1&cols=2
\needspace{10\baselineskip}
\begin{exercice}[BaremeTotal=false]
\label{eleve:28}


 \begin{minipage}[t]{\linewidth} Soit $f$ une fonction dérivable sur $[-5;5]$ et $\mathcal{C}_f$ sa courbe représentative.\\On sait que  $f(-2)=-2~~$ et que $~~f'(-2)=3$.\\Déterminer une équation de la tangente $(T)$ à la courbe $\mathcal{C}_f$ au point d'abscisse $-2$,\\sans utiliser la formule de cours de l'équation de tangente. \end{minipage}

\end{exercice}

\begin{Solution}
 \begin{minipage}[t]{\linewidth} On sait que la tangente n'est pas une droite verticale, puisque la fonction est dérivable sur l'intervalle.\\On en déduit que la tangente $(T)$ au point d'abscisse $-2$, admet une équation réduite de la forme :  $(T) : y=m x + p$. 

\medskip
$\bullet$ Détermination de $m$ :\\On sait que le nombre dérivé en $-2$ est par définition, le coefficient directeur de la tangente au point d'abscisse $-2$.\\Par conséquent, on a déjà : $m=f'(-2)=3$.\\ On en déduit que  $(T) : y= 3 x + p$\\$\bullet$ Détermination de $p$ :\\ Pour cela, on utilise que si $f(-2)=-2~~$, alors le point $A$ de coordonnées $(-2;-2)$ appartient à $\mathcal{C}_f$ mais aussi à $(T)$.\\ On peut écrire $A(-2;-2) = \mathcal{C}_f \cap (T)$.\\ On remplace alors les coordonnées de $A(-2;-2)$ dans l'équation  $(T) : y= 3 x + p$.\\ $\begin{aligned} \phantom{\iff}&A(-2;-2)\in (T)\\ \iff& -2= 3 \times (-2)  + p\\ \iff& p=-2 -3 \times (-2)  \\ \iff& p=-2 +6   \\ \iff& p=4\\\end{aligned}$\\On peut conclure que : $(T) : {\color[HTML]{f15929}\boldsymbol{y=3x+4}}$. \end{minipage}

\end{Solution}

% @see : https://coopmaths.fr/alea?uuid=0e984&id=can1F15&n=1&d=10&alea=1Alo223232323234234567891011121314151617181920212223242526272829303132333435363738394041424344454647482345678910111213141516171819202122232425262728&cd=1&cols=1
\needspace{10\baselineskip}
\begin{exercice}[BaremeTotal=false]


 La courbe représente une fonction $f$ et la droite est la tangente au point d'abscisse $0$.\\
        Déterminer $f'(0)$.\begin{center}\begin{tikzpicture}[baseline,scale = 0.5]

    \tikzset{
      point/.style={
        thick,
        draw,
        cross out,
        inner sep=0pt,
        minimum width=5pt,
        minimum height=5pt,
      },
    }
    \clip (-8,-10) rectangle (8,3);
    	
	\draw[color ={black},line width = 2,>=latex,->] (-8,0)--(8,0);
	\draw[color ={black},line width = 2,>=latex,->] (0,-10)--(0,3);
	\draw[color ={black},opacity = 0.5] (-8,1)--(8,1);
	\draw[color ={black},opacity = 0.5] (-8,2)--(8,2);
	\draw[color ={black},opacity = 0.5] (-8,3)--(8,3);
	\draw[color ={black},opacity = 0.5] (-8,-1)--(8,-1);
	\draw[color ={black},opacity = 0.5] (-8,-2)--(8,-2);
	\draw[color ={black},opacity = 0.5] (-8,-3)--(8,-3);
	\draw[color ={black},opacity = 0.5] (-8,-4)--(8,-4);
	\draw[color ={black},opacity = 0.5] (-8,-5)--(8,-5);
	\draw[color ={black},opacity = 0.5] (-8,-6)--(8,-6);
	\draw[color ={black},opacity = 0.5] (-8,-7)--(8,-7);
	\draw[color ={black},opacity = 0.5] (-8,-8)--(8,-8);
	\draw[color ={black},opacity = 0.5] (-8,-9)--(8,-9);
	\draw[color ={black},opacity = 0.5] (-8,-10)--(8,-10);
	\draw[color ={black},opacity = 0.5] (1,-10)--(1,3);
	\draw[color ={black},opacity = 0.5] (2,-10)--(2,3);
	\draw[color ={black},opacity = 0.5] (3,-10)--(3,3);
	\draw[color ={black},opacity = 0.5] (4,-10)--(4,3);
	\draw[color ={black},opacity = 0.5] (5,-10)--(5,3);
	\draw[color ={black},opacity = 0.5] (6,-10)--(6,3);
	\draw[color ={black},opacity = 0.5] (7,-10)--(7,3);
	\draw[color ={black},opacity = 0.5] (8,-10)--(8,3);
	\draw[color ={black},opacity = 0.5] (-1,-10)--(-1,3);
	\draw[color ={black},opacity = 0.5] (-2,-10)--(-2,3);
	\draw[color ={black},opacity = 0.5] (-3,-10)--(-3,3);
	\draw[color ={black},opacity = 0.5] (-4,-10)--(-4,3);
	\draw[color ={black},opacity = 0.5] (-5,-10)--(-5,3);
	\draw[color ={black},opacity = 0.5] (-6,-10)--(-6,3);
	\draw[color ={black},opacity = 0.5] (-7,-10)--(-7,3);
	\draw[color ={black},opacity = 0.5] (-8,-10)--(-8,3);
	\draw[color ={gray},opacity = 0.3] (-8,0)--(8,0);
	\draw[color ={gray},opacity = 0.3] (-8,0)--(8,0);
	\draw[color ={gray},opacity = 0.3] (-8,-4)--(8,-4);
	\draw[color ={gray},opacity = 0.3] (-8,-5)--(8,-5);
	\draw[color ={gray},opacity = 0.3] (-8,-6)--(8,-6);
	\draw[color ={gray},opacity = 0.3] (-8,-7)--(8,-7);
	\draw[color ={gray},opacity = 0.3] (-8,-8)--(8,-8);
	\draw[color ={gray},opacity = 0.3] (-8,-9)--(8,-9);
	\draw[color ={gray},opacity = 0.3] (-8,-10)--(8,-10);
	\draw[color ={gray},opacity = 0.3] (0,-10)--(0,3);
	\draw[color ={gray},opacity = 0.3] (0,-10)--(0,3);
	\draw[color ={black},line width = 2] (2,-0.2)--(2,0.2);
	\draw[color ={black},line width = 2] (4,-0.2)--(4,0.2);
	\draw[color ={black},line width = 2] (6,-0.2)--(6,0.2);
	\draw[color ={black},line width = 2] (-2,-0.2)--(-2,0.2);
	\draw[color ={black},line width = 2] (-4,-0.2)--(-4,0.2);
	\draw[color ={black},line width = 2] (-6,-0.2)--(-6,0.2);
	\draw[color ={black},line width = 2] (-0.2,1)--(0.2,1);
	\draw[color ={black},line width = 2] (-0.2,2)--(0.2,2);
	\draw[color ={black},line width = 2] (-0.2,-1)--(0.2,-1);
	\draw[color ={black},line width = 2] (-0.2,-2)--(0.2,-2);
	\draw[color ={black},line width = 2] (-0.2,-3)--(0.2,-3);
	\draw[color ={black},line width = 2] (-0.2,-4)--(0.2,-4);
	\draw[color ={black},line width = 2] (-0.2,-5)--(0.2,-5);
	\draw[color ={black},line width = 2] (-0.2,-6)--(0.2,-6);
	\draw[color ={black},line width = 2] (-0.2,-7)--(0.2,-7);
	\draw[color ={black},line width = 2] (-0.2,-8)--(0.2,-8);
	\draw[color ={black},line width = 2] (-0.2,-9)--(0.2,-9);
	\draw (2,-0.4) node[anchor = center] {\scriptsize \color{black}{$1$}};
	\draw (4,-0.4) node[anchor = center] {\scriptsize \color{black}{$2$}};
	\draw (6,-0.4) node[anchor = center] {\scriptsize \color{black}{$3$}};
	\draw (-2,-0.4) node[anchor = center] {\scriptsize \color{black}{$-1$}};
	\draw (-4,-0.4) node[anchor = center] {\scriptsize \color{black}{$-2$}};
	\draw (-6,-0.4) node[anchor = center] {\scriptsize \color{black}{$-3$}};
	\draw (-0.8,2.1) node[anchor = center] {\footnotesize \color{black}{$2$}};
	\draw (-0.8,-1.9) node[anchor = center] {\footnotesize \color{black}{$-2$}};
	\draw (-0.8,-3.9) node[anchor = center] {\footnotesize \color{black}{$-4$}};
	\draw (-0.8,-5.9) node[anchor = center] {\footnotesize \color{black}{$-6$}};
	\draw (-0.8,-7.9) node[anchor = center] {\footnotesize \color{black}{$-8$}};
	\draw [color={black}] (-0.3,-0.3) node[anchor = center,scale=1, rotate = 0] {O};
	
	\draw[color={blue},line width = 2] (-8,-7)--(-7.6,-5.84)--(-7.2,-4.76)--(-6.8,-3.76)--(-6.4,-2.84)--(-6,-2)--(-5.6,-1.24)--(-5.2,-0.56)--(-4.8,0.04)--(-4.4,0.56)--(-4,1)--(-3.6,1.36)--(-3.2,1.64)--(-2.8,1.84)--(-2.4,1.96)--(-2,2)--(-1.6,1.96)--(-1.2,1.84)--(-0.8,1.64)--(-0.4,1.36)--(0,1)--(0.4,0.56)--(0.8,0.04)--(1.2,-0.56)--(1.6,-1.24)--(2,-2)--(2.4,-2.84)--(2.8,-3.76)--(3.2,-4.76)--(3.6,-5.84)--(4,-7)--(4.4,-8.24)--(4.8,-9.56)--(5.2,-10.96);
	
	\draw[color={red},line width = 2] (-2.8,3.8)--(-2.4,3.4)--(-2,3)--(-1.6,2.6)--(-1.2,2.2)--(-0.8,1.8)--(-0.4,1.4)--(0,1)--(0.4,0.6)--(0.8,0.2)--(1.2,-0.2)--(1.6,-0.6)--(2,-1)--(2.4,-1.4)--(2.8,-1.8)--(3.2,-2.2)--(3.6,-2.6)--(4,-3)--(4.4,-3.4)--(4.8,-3.8)--(5.2,-4.2)--(5.6,-4.6)--(6,-5)--(6.4,-5.4)--(6.8,-5.8)--(7.2,-6.2)--(7.6,-6.6)--(8,-7);

\end{tikzpicture}\end{center}
\end{exercice}

\begin{Solution}
 $f'(0)$ est donné par le coefficient directeur de la tangente à la courbe au point d'abscisse $0$, soit $-2$.
\end{Solution}

% @see : https://coopmaths.fr/alea?uuid=6f32d&id=can1F16&n=1&d=10&alea=dqSw223232323234234567891011121314151617181920212223242526272829303132333435363738394041424344454647482345678910111213141516171819202122232425262728&cd=1&cols=1
\needspace{10\baselineskip}
\newpage
\begin{exercice}[BaremeTotal=false]


 On donne les représentations graphiques d'une fonction et de sa dérivée.\\
        Donner l'équation réduite de la tangente à la courbe de $f$ en $x=0$. \\ \begin{minipage}{0.5\linewidth}
    \begin{tikzpicture}[baseline, scale=0.8]

    \tikzset{
      point/.style={
        thick,
        draw,
        cross out,
        inner sep=0pt,
        minimum width=5pt,
        minimum height=5pt,
      },
    }
    \clip (-4.5,-4.5) rectangle (5.75,13.5);
    	
	\draw[color ={black},>=latex,->] (-3,0)--(3,0);
	\draw[color ={black},>=latex,->] (0,-3)--(0,12);
	\draw[color ={black},opacity = 0.4] (-3,1)--(3,1);
	\draw[color ={black},opacity = 0.4] (-3,2)--(3,2);
	\draw[color ={black},opacity = 0.4] (-3,3)--(3,3);
	\draw[color ={black},opacity = 0.4] (-3,4)--(3,4);
	\draw[color ={black},opacity = 0.4] (-3,5)--(3,5);
	\draw[color ={black},opacity = 0.4] (-3,6)--(3,6);
	\draw[color ={black},opacity = 0.4] (-3,7)--(3,7);
	\draw[color ={black},opacity = 0.4] (-3,8)--(3,8);
	\draw[color ={black},opacity = 0.4] (-3,9)--(3,9);
	\draw[color ={black},opacity = 0.4] (-3,10)--(3,10);
	\draw[color ={black},opacity = 0.4] (-3,11)--(3,11);
	\draw[color ={black},opacity = 0.4] (-3,12)--(3,12);
	\draw[color ={black},opacity = 0.4] (-3,-1)--(3,-1);
	\draw[color ={black},opacity = 0.4] (-3,-2)--(3,-2);
	\draw[color ={black},opacity = 0.4] (-3,-3)--(3,-3);
	\draw[color ={black},opacity = 0.4] (1,-3)--(1,12);
	\draw[color ={black},opacity = 0.4] (2,-3)--(2,12);
	\draw[color ={black},opacity = 0.4] (3,-3)--(3,12);
	\draw[color ={black},opacity = 0.4] (-1,-3)--(-1,12);
	\draw[color ={black},opacity = 0.4] (-2,-3)--(-2,12);
	\draw[color ={black},opacity = 0.4] (-3,-3)--(-3,12);
	\draw[color ={gray},opacity = 0.3] (-3,0)--(3,0);
	\draw[color ={gray},opacity = 0.3] (-3,0.5)--(3,0.5);
	\draw[color ={gray},opacity = 0.3] (-3,1.5)--(3,1.5);
	\draw[color ={gray},opacity = 0.3] (-3,2.5)--(3,2.5);
	\draw[color ={gray},opacity = 0.3] (-3,3.5)--(3,3.5);
	\draw[color ={gray},opacity = 0.3] (-3,4.5)--(3,4.5);
	\draw[color ={gray},opacity = 0.3] (-3,5.5)--(3,5.5);
	\draw[color ={gray},opacity = 0.3] (-3,6.5)--(3,6.5);
	\draw[color ={gray},opacity = 0.3] (-3,7.5)--(3,7.5);
	\draw[color ={gray},opacity = 0.3] (-3,8.5)--(3,8.5);
	\draw[color ={gray},opacity = 0.3] (-3,9.5)--(3,9.5);
	\draw[color ={gray},opacity = 0.3] (-3,10.5)--(3,10.5);
	\draw[color ={gray},opacity = 0.3] (-3,11.5)--(3,11.5);
	\draw[color ={gray},opacity = 0.3] (-3,0)--(3,0);
	\draw[color ={gray},opacity = 0.3] (-3,-0.5)--(3,-0.5);
	\draw[color ={gray},opacity = 0.3] (-3,-1.5)--(3,-1.5);
	\draw[color ={gray},opacity = 0.3] (-3,-2.5)--(3,-2.5);
	\draw[color ={gray},opacity = 0.3] (0,-3)--(0,12);
	\draw[color ={gray},opacity = 0.3] (0,-3)--(0,12);
	\draw[color ={black},line width = 1.2] (1,-0.13)--(1,0.13);
	\draw[color ={black},line width = 1.2] (2,-0.13)--(2,0.13);
	\draw[color ={black},line width = 1.2] (3,-0.13)--(3,0.13);
	\draw[color ={black},line width = 1.2] (-1,-0.13)--(-1,0.13);
	\draw[color ={black},line width = 1.2] (-2,-0.13)--(-2,0.13);
	\draw[color ={black},line width = 1.2] (-3,-0.13)--(-3,0.13);
	\draw[color ={black},line width = 1.2] (-0.13,1)--(0.13,1);
	\draw[color ={black},line width = 1.2] (-0.13,2)--(0.13,2);
	\draw[color ={black},line width = 1.2] (-0.13,3)--(0.13,3);
	\draw[color ={black},line width = 1.2] (-0.13,4)--(0.13,4);
	\draw[color ={black},line width = 1.2] (-0.13,5)--(0.13,5);
	\draw[color ={black},line width = 1.2] (-0.13,6)--(0.13,6);
	\draw[color ={black},line width = 1.2] (-0.13,7)--(0.13,7);
	\draw[color ={black},line width = 1.2] (-0.13,8)--(0.13,8);
	\draw[color ={black},line width = 1.2] (-0.13,9)--(0.13,9);
	\draw[color ={black},line width = 1.2] (-0.13,10)--(0.13,10);
	\draw[color ={black},line width = 1.2] (-0.13,11)--(0.13,11);
	\draw[color ={black},line width = 1.2] (-0.13,12)--(0.13,12);
	\draw[color ={black},line width = 1.2] (-0.13,-1)--(0.13,-1);
	\draw[color ={black},line width = 1.2] (-0.13,-2)--(0.13,-2);
	\draw[color ={black},line width = 1.2] (-0.13,-3)--(0.13,-3);
	\draw (3,-0.4) node[anchor = center] {\scriptsize \color{black}{$3$}};
	\draw (-3,-0.4) node[anchor = center] {\scriptsize \color{black}{$-3$}};
	\draw (-0.5,3.1) node[anchor = center] {\footnotesize \color{black}{$3$}};
	\draw (-0.5,6.1) node[anchor = center] {\footnotesize \color{black}{$6$}};
	\draw (-0.5,9.1) node[anchor = center] {\footnotesize \color{black}{$9$}};
	\draw (-0.5,-2.9) node[anchor = center] {\footnotesize \color{black}{$-3$}};
	\draw [color={black}] (-0.3,-0.3) node[anchor = center,scale=1, rotate = 0] {O};
	\draw (3,10) node[anchor = center] {\footnotesize \color{blue}{$\Large \mathcal{C}_f$}};
	
	\draw[color={blue},line width = 2] (-3,-1)--(-2.8,-1.36)--(-2.6,-1.64)--(-2.4,-1.84)--(-2.2,-1.96)--(-2,-2)--(-1.8,-1.96)--(-1.6,-1.84)--(-1.4,-1.64)--(-1.2,-1.36)--(-1,-1)--(-0.8,-0.56)--(-0.6,-0.04)--(-0.4,0.56)--(-0.2,1.24)--(0,2)--(0.2,2.84)--(0.4,3.76)--(0.6,4.76)--(0.8,5.84)--(1,7)--(1.2,8.24)--(1.4,9.56)--(1.6,10.96)--(1.8,12.44);

\end{tikzpicture}\\
    \end{minipage}
    \begin{minipage}{0.5\linewidth}
    \begin{tikzpicture}[baseline, scale=0.8]

    \tikzset{
      point/.style={
        thick,
        draw,
        cross out,
        inner sep=0pt,
        minimum width=5pt,
        minimum height=5pt,
      },
    }
    \clip (-4.5,-6.5) rectangle (6.5,9.5);
    	
	\draw[color ={black},>=latex,->] (-3,0)--(3,0);
	\draw[color ={black},>=latex,->] (0,-5)--(0,8);
	\draw[color ={black},opacity = 0.4] (-3,1)--(3,1);
	\draw[color ={black},opacity = 0.4] (-3,2)--(3,2);
	\draw[color ={black},opacity = 0.4] (-3,3)--(3,3);
	\draw[color ={black},opacity = 0.4] (-3,4)--(3,4);
	\draw[color ={black},opacity = 0.4] (-3,5)--(3,5);
	\draw[color ={black},opacity = 0.4] (-3,6)--(3,6);
	\draw[color ={black},opacity = 0.4] (-3,7)--(3,7);
	\draw[color ={black},opacity = 0.4] (-3,8)--(3,8);
	\draw[color ={black},opacity = 0.4] (-3,-1)--(3,-1);
	\draw[color ={black},opacity = 0.4] (-3,-2)--(3,-2);
	\draw[color ={black},opacity = 0.4] (-3,-3)--(3,-3);
	\draw[color ={black},opacity = 0.4] (-3,-4)--(3,-4);
	\draw[color ={black},opacity = 0.4] (-3,-5)--(3,-5);
	\draw[color ={black},opacity = 0.4] (1,-5)--(1,8);
	\draw[color ={black},opacity = 0.4] (2,-5)--(2,8);
	\draw[color ={black},opacity = 0.4] (3,-5)--(3,8);
	\draw[color ={black},opacity = 0.4] (-1,-5)--(-1,8);
	\draw[color ={black},opacity = 0.4] (-2,-5)--(-2,8);
	\draw[color ={black},opacity = 0.4] (-3,-5)--(-3,8);
	\draw[color ={gray},opacity = 0.3] (-3,0)--(3,0);
	\draw[color ={gray},opacity = 0.3] (-3,0.5)--(3,0.5);
	\draw[color ={gray},opacity = 0.3] (-3,1.5)--(3,1.5);
	\draw[color ={gray},opacity = 0.3] (-3,2.5)--(3,2.5);
	\draw[color ={gray},opacity = 0.3] (-3,3.5)--(3,3.5);
	\draw[color ={gray},opacity = 0.3] (-3,4.5)--(3,4.5);
	\draw[color ={gray},opacity = 0.3] (-3,5.5)--(3,5.5);
	\draw[color ={gray},opacity = 0.3] (-3,6.5)--(3,6.5);
	\draw[color ={gray},opacity = 0.3] (-3,7.5)--(3,7.5);
	\draw[color ={gray},opacity = 0.3] (-3,0)--(3,0);
	\draw[color ={gray},opacity = 0.3] (-3,-0.5)--(3,-0.5);
	\draw[color ={gray},opacity = 0.3] (-3,-1.5)--(3,-1.5);
	\draw[color ={gray},opacity = 0.3] (-3,-2.5)--(3,-2.5);
	\draw[color ={gray},opacity = 0.3] (-3,-3.5)--(3,-3.5);
	\draw[color ={gray},opacity = 0.3] (-3,-4.5)--(3,-4.5);
	\draw[color ={gray},opacity = 0.3] (0,-5)--(0,8);
	\draw[color ={gray},opacity = 0.3] (0,-5)--(0,8);
	\draw[color ={black},line width = 1.2] (1,-0.13)--(1,0.13);
	\draw[color ={black},line width = 1.2] (2,-0.13)--(2,0.13);
	\draw[color ={black},line width = 1.2] (3,-0.13)--(3,0.13);
	\draw[color ={black},line width = 1.2] (-1,-0.13)--(-1,0.13);
	\draw[color ={black},line width = 1.2] (-2,-0.13)--(-2,0.13);
	\draw[color ={black},line width = 1.2] (-3,-0.13)--(-3,0.13);
	\draw[color ={black},line width = 1.2] (-0.13,1)--(0.13,1);
	\draw[color ={black},line width = 1.2] (-0.13,2)--(0.13,2);
	\draw[color ={black},line width = 1.2] (-0.13,3)--(0.13,3);
	\draw[color ={black},line width = 1.2] (-0.13,4)--(0.13,4);
	\draw[color ={black},line width = 1.2] (-0.13,5)--(0.13,5);
	\draw[color ={black},line width = 1.2] (-0.13,6)--(0.13,6);
	\draw[color ={black},line width = 1.2] (-0.13,7)--(0.13,7);
	\draw[color ={black},line width = 1.2] (-0.13,8)--(0.13,8);
	\draw[color ={black},line width = 1.2] (-0.13,9)--(0.13,9);
	\draw[color ={black},line width = 1.2] (-0.13,10)--(0.13,10);
	\draw[color ={black},line width = 1.2] (-0.13,11)--(0.13,11);
	\draw[color ={black},line width = 1.2] (-0.13,12)--(0.13,12);
	\draw[color ={black},line width = 1.2] (-0.13,-1)--(0.13,-1);
	\draw[color ={black},line width = 1.2] (-0.13,-2)--(0.13,-2);
	\draw[color ={black},line width = 1.2] (-0.13,-3)--(0.13,-3);
	\draw (3,-0.4) node[anchor = center] {\scriptsize \color{black}{$3$}};
	\draw (-3,-0.4) node[anchor = center] {\scriptsize \color{black}{$-3$}};
	\draw (-0.5,3.1) node[anchor = center] {\footnotesize \color{black}{$3$}};
	\draw (-0.5,6.1) node[anchor = center] {\footnotesize \color{black}{$6$}};
	\draw (-0.5,-2.9) node[anchor = center] {\footnotesize \color{black}{$-3$}};
	\draw [color={black}] (-0.3,-0.3) node[anchor = center,scale=1, rotate = 0] {O};
	\draw (3,6) node[anchor = center] {\footnotesize \color{red}{$\Large\mathcal{C}_f\prime$}};
	
	\draw[color={red},line width = 2] (-3,-2)--(-2.8,-1.6)--(-2.6,-1.2)--(-2.4,-0.8)--(-2.2,-0.4)--(-2,0)--(-1.8,0.4)--(-1.6,0.8)--(-1.4,1.2)--(-1.2,1.6)--(-1,2)--(-0.8,2.4)--(-0.6,2.8)--(-0.4,3.2)--(-0.2,3.6)--(0,4)--(0.2,4.4)--(0.4,4.8)--(0.6,5.2)--(0.8,5.6)--(1,6)--(1.2,6.4)--(1.4,6.8)--(1.6,7.2)--(1.8,7.6)--(2,8)--(2.2,8.4)--(2.4,8.8);

\end{tikzpicture}\\
    \end{minipage}
    
\end{exercice}

\begin{Solution}
 L'équation réduite de la tangente au point d'abscisse $0$ est  : $y=f'(0)(x-0)+f(0)$.\\
        On lit graphiquement $f(0)=2$ et $f'(0)=4$.\\
        L'équation réduite de la tangente est donc donnée par :
        $y=4(x+0)+2$, soit $y=4x+2$.
\end{Solution}

% @see : https://coopmaths.fr/alea?uuid=a1ba2&id=can1F14&n=1&d=10&alea=rkIp223232323234234567891011121314151617181920212223242526272829303132333435363738394041424344454647482345678910111213141516171819202122232425262728&cd=1&cols=1
\needspace{10\baselineskip}
\begin{exercice}[BaremeTotal=false]


 Soit $f$ la fonction définie sur $\mathbb{R}$ par : $f(x)= x^2-2x+6$.\\
        En calculant la limite du taux d'accroissement, déterminer $f'(-3)$.
\end{exercice}

\begin{Solution}
 $f$ est une fonction polynôme du second degré de la forme $f(x)=ax^2+bx+c$.\\
    La fonction dérivée est donnée par la somme des dérivées des fonctions $u$ et $v$ définies par $u(x)=x^2$ et $v(x)=-2x+6$.\\
     Comme $u'(x)=2x$ et $v'(x)=-2$, on obtient  $f'(x)=2x-2$. \\
     Ainsi, $f'(-3)=2\times (-3)-2=-8$.
\end{Solution}

% @see : https://coopmaths.fr/alea?uuid=3c690&id=can1F13&n=1&d=10&alea=EUDu223232323234234567891011121314151617181920212223242526272829303132333435363738394041424344454647482345678910111213141516171819202122232425262728&cd=1&cols=1
\needspace{10\baselineskip}
\begin{exercice}[BaremeTotal=false]


 Déterminer le coefficient directeur de la tangente à la courbe représentative de la fonction cube au point d'abscisse $1$.
                        
\end{exercice}

\begin{Solution}
 Le coefficient directeur de la tangente au point d'abscisse $a$ est donné par le nombre dérivé $f'(a)$.\\
        La fonction $f$ définie par $f(x)=x^3$ a pour fonction dérivée la fonction $f'$ définie par $f'(x)=3x^2$.\\
        Comme $f'(1)=3\times 1^2=3$, le coefficient directeur de la tangente au point d'abscisse $1$ est : $3$. 
\end{Solution}

% @see : https://coopmaths.fr/alea?uuid=ebd89&id=1AN14-1&n=1&d=10&s=3&alea=ocYr223232323234234567891011121314151617181920212223242526272829303132333435363738394041424344454647482345678910111213141516171819202122232425262728&cd=0&cols=1
\needspace{10\baselineskip}
\begin{exercice}[BaremeTotal=false]


 Donner la dérivée de la fonction $f$, dérivable pour tout $x\in\R$, définie par  $f(x)=\dfrac{5x+2}{6}$.
\end{exercice}

\begin{Solution}
 L'expression de la dérivée de la fonction $f$ définie par $f(x)=\dfrac{5x+2}{6}$ est : ${\color[HTML]{f15929}\boldsymbol{f'(x)=\dfrac{5}{6}}}$.
\end{Solution}

% @see : https://coopmaths.fr/alea?uuid=ebd89&id=1AN14-1&n=1&d=10&s=4&alea=uAlP223232323234234567891011121314151617181920212223242526272829303132333435363738394041424344454647482345678910111213141516171819202122232425262728&cd=0&cols=1
\needspace{10\baselineskip}
\begin{exercice}[BaremeTotal=false]


 Donner la dérivée de la fonction $f$, dérivable pour tout $x\in\R$, définie par  $f(x)=4x^{4}$.
\end{exercice}

\begin{Solution}
 L'expression de la dérivée de la fonction $f$ définie par $f(x)=4x^{4}$ est : ${\color[HTML]{f15929}\boldsymbol{f'(x)=16x^{3}}}$.
\end{Solution}

% @see : https://coopmaths.fr/alea?uuid=ebd89&id=1AN14-1&n=1&d=10&s=7&alea=vjN7223232323234234567891011121314151617181920212223242526272829303132333435363738394041424344454647482345678910111213141516171819202122232425262728&cd=0&cols=1
\needspace{10\baselineskip}
\begin{exercice}[BaremeTotal=false]


 Donner la dérivée de la fonction $f$, dérivable pour tout $x\in\R^*$, définie par  $f(x)=\dfrac{-1}{5x}$.
\end{exercice}

\begin{Solution}
 L'expression de la dérivée de la fonction $f$ définie par $f(x)=\dfrac{-1}{5x}$ est : ${\color[HTML]{f15929}\boldsymbol{f'(x)=\dfrac{1}{5x^2}}}$.
\end{Solution}

% @see : https://coopmaths.fr/alea?uuid=a83c0&id=1AN14-4&n=1&d=10&s=1&alea=Njr3223232323234234567891011121314151617181920212223242526272829303132333435363738394041424344454647482345678910111213141516171819202122232425262728&cd=0&cols=1
\needspace{10\baselineskip}
\begin{exercice}[BaremeTotal=false]


 Donner l'expression de la dérivée de la fonction $f$ définie sur $\R^{*}$ par $f(x)=\dfrac{7}{x}-2x^{2}-2x-3$.\\
\end{exercice}

\begin{Solution}
 L'expression de la dérivée de la fonction $f$ définie par $f(x)=\dfrac{7}{x}-2x^{2}-2x-3$ est :\\$f'(x)={\color[HTML]{f15929}\boldsymbol{-\dfrac{7}{x^2}-4x-2}}$.\\
\end{Solution}

% @see : https://coopmaths.fr/alea?uuid=a83c0&id=1AN14-4&n=1&d=10&s=4&alea=4Vpz223232323234234567891011121314151617181920212223242526272829303132333435363738394041424344454647482345678910111213141516171819202122232425262728&cd=1&cols=1
\needspace{10\baselineskip}
\begin{exercice}[BaremeTotal=false]


 Donner l'expression de la dérivée de la fonction $f$ définie sur $\R_+^{*}$ par $f(x)=-3\sqrt{x}+5x^{2}+5x-4-\dfrac{2}{x}$.\\
\end{exercice}

\begin{Solution}
 $(-3\sqrt{x})^\prime=-\dfrac{3}{2\sqrt{x}}$\\$(5x^{2}+5x-4)^\prime=10x+5$\\$(-\dfrac{2}{x})^\prime=\dfrac{2}{x^2}$\\L'expression de la dérivée de la fonction $f$ définie par $f(x)=-3\sqrt{x}+5x^{2}+5x-4-\dfrac{2}{x}$ est :\\$f'(x)={\color[HTML]{f15929}\boldsymbol{-\dfrac{3}{2\sqrt{x}}+10x+5+\dfrac{2}{x^2}}}$.\\
\end{Solution}

\end{Maquette}
\clearpage
\begin{Maquette}[DM]{Niveau= ,Classe=\getdata[29]\mydata\ (v29), Date = 06/01/25  ,Theme=Exercices}


% @see : https://coopmaths.fr/alea?uuid=4c8c7&id=1AN11-3&n=1&d=10&s=2&alea=nawO22323232323423456789101112131415161718192021222324252627282930313233343536373839404142434445464748234567891011121314151617181920212223242526272829&cd=1&cols=2
\needspace{10\baselineskip}
\begin{exercice}[BaremeTotal=false]
\label{eleve:29}


 \begin{minipage}[t]{\linewidth} Soit $f$ une fonction dérivable sur $[-5;5]$ et $\mathcal{C}_f$ sa courbe représentative.\\On sait que  $f(3)=-3~~$ et que $~~f'(3)=-3$.\\Déterminer une équation de la tangente $(T)$ à la courbe $\mathcal{C}_f$ au point d'abscisse $3$,\\sans utiliser la formule de cours de l'équation de tangente. \end{minipage}

\end{exercice}

\begin{Solution}
 \begin{minipage}[t]{\linewidth} On sait que la tangente n'est pas une droite verticale, puisque la fonction est dérivable sur l'intervalle.\\On en déduit que la tangente $(T)$ au point d'abscisse $3$, admet une équation réduite de la forme :  $(T) : y=m x + p$. 

\medskip
$\bullet$ Détermination de $m$ :\\On sait que le nombre dérivé en $3$ est par définition, le coefficient directeur de la tangente au point d'abscisse $3$.\\Par conséquent, on a déjà : $m=f'(3)=-3$.\\ On en déduit que  $(T) : y= -3 x + p$\\$\bullet$ Détermination de $p$ :\\ Pour cela, on utilise que si $f(3)=-3~~$, alors le point $A$ de coordonnées $(3;-3)$ appartient à $\mathcal{C}_f$ mais aussi à $(T)$.\\ On peut écrire $A(3;-3) = \mathcal{C}_f \cap (T)$.\\ On remplace alors les coordonnées de $A(3;-3)$ dans l'équation  $(T) : y= -3 x + p$.\\ $\begin{aligned} \phantom{\iff}&A(3;-3)\in (T)\\ \iff& -3= -3 \times 3  + p\\ \iff& p=-3 +3 \times 3  \\ \iff& p=-3 +9   \\ \iff& p=6\\\end{aligned}$\\On peut conclure que : $(T) : {\color[HTML]{f15929}\boldsymbol{y=-3x+6}}$. \end{minipage}

\end{Solution}

% @see : https://coopmaths.fr/alea?uuid=0e984&id=can1F15&n=1&d=10&alea=1Alo22323232323423456789101112131415161718192021222324252627282930313233343536373839404142434445464748234567891011121314151617181920212223242526272829&cd=1&cols=1
\needspace{10\baselineskip}
\begin{exercice}[BaremeTotal=false]


 La courbe représente une fonction $f$ et la droite est la tangente au point d'abscisse $1$.\\
        
        Déterminer $f'(1)$.\begin{center}\begin{tikzpicture}[baseline,scale = 0.5]

    \tikzset{
      point/.style={
        thick,
        draw,
        cross out,
        inner sep=0pt,
        minimum width=5pt,
        minimum height=5pt,
      },
    }
    \clip (-8,-3) rectangle (8,10);
    	
	\draw[color ={black},line width = 2,>=latex,->] (-8,0)--(8,0);
	\draw[color ={black},line width = 2,>=latex,->] (0,-3)--(0,10);
	\draw[color ={black},opacity = 0.5] (-8,1)--(8,1);
	\draw[color ={black},opacity = 0.5] (-8,2)--(8,2);
	\draw[color ={black},opacity = 0.5] (-8,3)--(8,3);
	\draw[color ={black},opacity = 0.5] (-8,4)--(8,4);
	\draw[color ={black},opacity = 0.5] (-8,5)--(8,5);
	\draw[color ={black},opacity = 0.5] (-8,6)--(8,6);
	\draw[color ={black},opacity = 0.5] (-8,7)--(8,7);
	\draw[color ={black},opacity = 0.5] (-8,8)--(8,8);
	\draw[color ={black},opacity = 0.5] (-8,9)--(8,9);
	\draw[color ={black},opacity = 0.5] (-8,10)--(8,10);
	\draw[color ={black},opacity = 0.5] (-8,-1)--(8,-1);
	\draw[color ={black},opacity = 0.5] (-8,-2)--(8,-2);
	\draw[color ={black},opacity = 0.5] (-8,-3)--(8,-3);
	\draw[color ={black},opacity = 0.5] (1,-3)--(1,10);
	\draw[color ={black},opacity = 0.5] (2,-3)--(2,10);
	\draw[color ={black},opacity = 0.5] (3,-3)--(3,10);
	\draw[color ={black},opacity = 0.5] (4,-3)--(4,10);
	\draw[color ={black},opacity = 0.5] (5,-3)--(5,10);
	\draw[color ={black},opacity = 0.5] (6,-3)--(6,10);
	\draw[color ={black},opacity = 0.5] (7,-3)--(7,10);
	\draw[color ={black},opacity = 0.5] (8,-3)--(8,10);
	\draw[color ={black},opacity = 0.5] (-1,-3)--(-1,10);
	\draw[color ={black},opacity = 0.5] (-2,-3)--(-2,10);
	\draw[color ={black},opacity = 0.5] (-3,-3)--(-3,10);
	\draw[color ={black},opacity = 0.5] (-4,-3)--(-4,10);
	\draw[color ={black},opacity = 0.5] (-5,-3)--(-5,10);
	\draw[color ={black},opacity = 0.5] (-6,-3)--(-6,10);
	\draw[color ={black},opacity = 0.5] (-7,-3)--(-7,10);
	\draw[color ={black},opacity = 0.5] (-8,-3)--(-8,10);
	\draw[color ={gray},opacity = 0.3] (-8,0)--(8,0);
	\draw[color ={gray},opacity = 0.3] (-8,0)--(8,0);
	\draw[color ={gray},opacity = 0.3] (0,-3)--(0,10);
	\draw[color ={gray},opacity = 0.3] (0,-3)--(0,10);
	\draw[color ={black},line width = 2] (2,-0.2)--(2,0.2);
	\draw[color ={black},line width = 2] (4,-0.2)--(4,0.2);
	\draw[color ={black},line width = 2] (6,-0.2)--(6,0.2);
	\draw[color ={black},line width = 2] (-2,-0.2)--(-2,0.2);
	\draw[color ={black},line width = 2] (-4,-0.2)--(-4,0.2);
	\draw[color ={black},line width = 2] (-6,-0.2)--(-6,0.2);
	\draw[color ={black},line width = 2] (-0.2,1)--(0.2,1);
	\draw[color ={black},line width = 2] (-0.2,2)--(0.2,2);
	\draw[color ={black},line width = 2] (-0.2,3)--(0.2,3);
	\draw[color ={black},line width = 2] (-0.2,4)--(0.2,4);
	\draw[color ={black},line width = 2] (-0.2,5)--(0.2,5);
	\draw[color ={black},line width = 2] (-0.2,6)--(0.2,6);
	\draw[color ={black},line width = 2] (-0.2,7)--(0.2,7);
	\draw[color ={black},line width = 2] (-0.2,8)--(0.2,8);
	\draw[color ={black},line width = 2] (-0.2,9)--(0.2,9);
	\draw[color ={black},line width = 2] (-0.2,-1)--(0.2,-1);
	\draw[color ={black},line width = 2] (-0.2,-2)--(0.2,-2);
	\draw (2,-0.4) node[anchor = center] {\scriptsize \color{black}{$1$}};
	\draw (4,-0.4) node[anchor = center] {\scriptsize \color{black}{$2$}};
	\draw (6,-0.4) node[anchor = center] {\scriptsize \color{black}{$3$}};
	\draw (-2,-0.4) node[anchor = center] {\scriptsize \color{black}{$-1$}};
	\draw (-4,-0.4) node[anchor = center] {\scriptsize \color{black}{$-2$}};
	\draw (-6,-0.4) node[anchor = center] {\scriptsize \color{black}{$-3$}};
	\draw (-0.8,2.1) node[anchor = center] {\footnotesize \color{black}{$2$}};
	\draw (-0.8,4.1) node[anchor = center] {\footnotesize \color{black}{$4$}};
	\draw (-0.8,6.1) node[anchor = center] {\footnotesize \color{black}{$6$}};
	\draw (-0.8,8.1) node[anchor = center] {\footnotesize \color{black}{$8$}};
	\draw [color={black}] (-0.3,-0.3) node[anchor = center,scale=1, rotate = 0] {O};
	
	\draw[color={blue},line width = 2] (-4.4,10.24)--(-4,9)--(-3.6,7.84)--(-3.2,6.76)--(-2.8,5.76)--(-2.4,4.84)--(-2,4)--(-1.6,3.24)--(-1.2,2.56)--(-0.8,1.96)--(-0.4,1.44)--(0,1)--(0.4,0.64)--(0.8,0.36)--(1.2,0.16)--(1.6,0.04)--(2,0)--(2.4,0.04)--(2.8,0.16)--(3.2,0.36)--(3.6,0.64)--(4,1)--(4.4,1.44)--(4.8,1.96)--(5.2,2.56)--(5.6,3.24)--(6,4)--(6.4,4.84)--(6.8,5.76)--(7.2,6.76)--(7.6,7.84)--(8,9);
	
	\draw[color={red},line width = 2] (-8,0)--(-7.6,0)--(-7.2,0)--(-6.8,0)--(-6.4,0)--(-6,0)--(-5.6,0)--(-5.2,0)--(-4.8,0)--(-4.4,0)--(-4,0)--(-3.6,0)--(-3.2,0)--(-2.8,0)--(-2.4,0)--(-2,0)--(-1.6,0)--(-1.2,0)--(-0.8,0)--(-0.4,0)--(0,0)--(0.4,0)--(0.8,0)--(1.2,0)--(1.6,0)--(2,0)--(2.4,0)--(2.8,0)--(3.2,0)--(3.6,0)--(4,0)--(4.4,0)--(4.8,0)--(5.2,0)--(5.6,0)--(6,0)--(6.4,0)--(6.8,0)--(7.2,0)--(7.6,0)--(8,0);

\end{tikzpicture}\end{center}
\end{exercice}

\begin{Solution}
 $f'(1)$ est donné par le coefficient directeur de la tangente à la courbe au point d'abscisse $1$, soit $0$.
\end{Solution}

% @see : https://coopmaths.fr/alea?uuid=6f32d&id=can1F16&n=1&d=10&alea=dqSw22323232323423456789101112131415161718192021222324252627282930313233343536373839404142434445464748234567891011121314151617181920212223242526272829&cd=1&cols=1
\needspace{10\baselineskip}
\newpage
\begin{exercice}[BaremeTotal=false]


 On donne les représentations graphiques d'une fonction et de sa dérivée.\\
      Donner l'équation réduite de la tangente à la courbe de $f$ en $x=0$. \\ \begin{minipage}{0.5\linewidth}
    \begin{tikzpicture}[baseline, scale=0.8]

    \tikzset{
      point/.style={
        thick,
        draw,
        cross out,
        inner sep=0pt,
        minimum width=5pt,
        minimum height=5pt,
      },
    }
    \clip (-5.5,-9.5) rectangle (5.75,7.5);
    	
	\draw[color ={black},>=latex,->] (-4,0)--(4,0);
	\draw[color ={black},>=latex,->] (0,-8)--(0,6);
	\draw[color ={black},opacity = 0.4] (-4,1)--(4,1);
	\draw[color ={black},opacity = 0.4] (-4,2)--(4,2);
	\draw[color ={black},opacity = 0.4] (-4,3)--(4,3);
	\draw[color ={black},opacity = 0.4] (-4,4)--(4,4);
	\draw[color ={black},opacity = 0.4] (-4,5)--(4,5);
	\draw[color ={black},opacity = 0.4] (-4,6)--(4,6);
	\draw[color ={black},opacity = 0.4] (-4,-1)--(4,-1);
	\draw[color ={black},opacity = 0.4] (-4,-2)--(4,-2);
	\draw[color ={black},opacity = 0.4] (-4,-3)--(4,-3);
	\draw[color ={black},opacity = 0.4] (-4,-4)--(4,-4);
	\draw[color ={black},opacity = 0.4] (-4,-5)--(4,-5);
	\draw[color ={black},opacity = 0.4] (-4,-6)--(4,-6);
	\draw[color ={black},opacity = 0.4] (-4,-7)--(4,-7);
	\draw[color ={black},opacity = 0.4] (-4,-8)--(4,-8);
	\draw[color ={black},opacity = 0.4] (1,-8)--(1,6);
	\draw[color ={black},opacity = 0.4] (2,-8)--(2,6);
	\draw[color ={black},opacity = 0.4] (3,-8)--(3,6);
	\draw[color ={black},opacity = 0.4] (4,-8)--(4,6);
	\draw[color ={black},opacity = 0.4] (-1,-8)--(-1,6);
	\draw[color ={black},opacity = 0.4] (-2,-8)--(-2,6);
	\draw[color ={black},opacity = 0.4] (-3,-8)--(-3,6);
	\draw[color ={black},opacity = 0.4] (-4,-8)--(-4,6);
	\draw[color ={gray},opacity = 0.3] (-4,0)--(4,0);
	\draw[color ={gray},opacity = 0.3] (-4,0.5)--(4,0.5);
	\draw[color ={gray},opacity = 0.3] (-4,1.5)--(4,1.5);
	\draw[color ={gray},opacity = 0.3] (-4,2.5)--(4,2.5);
	\draw[color ={gray},opacity = 0.3] (-4,3.5)--(4,3.5);
	\draw[color ={gray},opacity = 0.3] (-4,4.5)--(4,4.5);
	\draw[color ={gray},opacity = 0.3] (-4,5.5)--(4,5.5);
	\draw[color ={gray},opacity = 0.3] (-4,0)--(4,0);
	\draw[color ={gray},opacity = 0.3] (-4,-0.5)--(4,-0.5);
	\draw[color ={gray},opacity = 0.3] (-4,-1.5)--(4,-1.5);
	\draw[color ={gray},opacity = 0.3] (-4,-2.5)--(4,-2.5);
	\draw[color ={gray},opacity = 0.3] (-4,-3.5)--(4,-3.5);
	\draw[color ={gray},opacity = 0.3] (-4,-4.5)--(4,-4.5);
	\draw[color ={gray},opacity = 0.3] (-4,-5.5)--(4,-5.5);
	\draw[color ={gray},opacity = 0.3] (-4,-6.5)--(4,-6.5);
	\draw[color ={gray},opacity = 0.3] (-4,-7)--(4,-7);
	\draw[color ={gray},opacity = 0.3] (-4,-7.5)--(4,-7.5);
	\draw[color ={gray},opacity = 0.3] (-4,-8)--(4,-8);
	\draw[color ={gray},opacity = 0.3] (0,-8)--(0,6);
	\draw[color ={gray},opacity = 0.3] (0,-8)--(0,6);
	\draw[color ={black},line width = 1.2] (1,-0.13)--(1,0.13);
	\draw[color ={black},line width = 1.2] (2,-0.13)--(2,0.13);
	\draw[color ={black},line width = 1.2] (3,-0.13)--(3,0.13);
	\draw[color ={black},line width = 1.2] (4,-0.13)--(4,0.13);
	\draw[color ={black},line width = 1.2] (-1,-0.13)--(-1,0.13);
	\draw[color ={black},line width = 1.2] (-2,-0.13)--(-2,0.13);
	\draw[color ={black},line width = 1.2] (-3,-0.13)--(-3,0.13);
	\draw[color ={black},line width = 1.2] (-4,-0.13)--(-4,0.13);
	\draw[color ={black},line width = 1.2] (-0.13,1)--(0.13,1);
	\draw[color ={black},line width = 1.2] (-0.13,2)--(0.13,2);
	\draw[color ={black},line width = 1.2] (-0.13,3)--(0.13,3);
	\draw[color ={black},line width = 1.2] (-0.13,4)--(0.13,4);
	\draw[color ={black},line width = 1.2] (-0.13,5)--(0.13,5);
	\draw[color ={black},line width = 1.2] (-0.13,6)--(0.13,6);
	\draw[color ={black},line width = 1.2] (-0.13,-1)--(0.13,-1);
	\draw[color ={black},line width = 1.2] (-0.13,-2)--(0.13,-2);
	\draw[color ={black},line width = 1.2] (-0.13,-3)--(0.13,-3);
	\draw[color ={black},line width = 1.2] (-0.13,-4)--(0.13,-4);
	\draw[color ={black},line width = 1.2] (-0.13,-5)--(0.13,-5);
	\draw[color ={black},line width = 1.2] (-0.13,-6)--(0.13,-6);
	\draw[color ={black},line width = 1.2] (-0.13,-7)--(0.13,-7);
	\draw[color ={black},line width = 1.2] (-0.13,-8)--(0.13,-8);
	\draw (2,-0.4) node[anchor = center] {\scriptsize \color{black}{$2$}};
	\draw (4,-0.4) node[anchor = center] {\scriptsize \color{black}{$4$}};
	\draw (-2,-0.4) node[anchor = center] {\scriptsize \color{black}{$-2$}};
	\draw (-4,-0.4) node[anchor = center] {\scriptsize \color{black}{$-4$}};
	\draw (-0.5,2.1) node[anchor = center] {\footnotesize \color{black}{$2$}};
	\draw (-0.5,4.1) node[anchor = center] {\footnotesize \color{black}{$4$}};
	\draw (-0.5,6.1) node[anchor = center] {\footnotesize \color{black}{$6$}};
	\draw (-0.5,-1.9) node[anchor = center] {\footnotesize \color{black}{$-2$}};
	\draw (-0.5,-3.9) node[anchor = center] {\footnotesize \color{black}{$-4$}};
	\draw (-0.5,-5.9) node[anchor = center] {\footnotesize \color{black}{$-6$}};
	\draw (-0.5,-7.9) node[anchor = center] {\footnotesize \color{black}{$-8$}};
	\draw [color={black}] (-0.3,-0.3) node[anchor = center,scale=1, rotate = 0] {O};
	\draw (3,4) node[anchor = center] {\footnotesize \color{blue}{$\Large \mathcal{C}_f$}};
	
	\draw[color={blue},line width = 2] (-2.4,-8.56)--(-2.2,-7.24)--(-2,-6)--(-1.8,-4.84)--(-1.6,-3.76)--(-1.4,-2.76)--(-1.2,-1.84)--(-1,-1)--(-0.8,-0.24)--(-0.6,0.44)--(-0.4,1.04)--(-0.2,1.56)--(0,2)--(0.2,2.36)--(0.4,2.64)--(0.6,2.84)--(0.8,2.96)--(1,3)--(1.2,2.96)--(1.4,2.84)--(1.6,2.64)--(1.8,2.36)--(2,2)--(2.2,1.56)--(2.4,1.04)--(2.6,0.44)--(2.8,-0.24)--(3,-1)--(3.2,-1.84)--(3.4,-2.76)--(3.6,-3.76)--(3.8,-4.84)--(4,-6);

\end{tikzpicture}\\
    \end{minipage}
    \begin{minipage}{0.5\linewidth}
    \begin{tikzpicture}[baseline, scale=0.8]

    \tikzset{
      point/.style={
        thick,
        draw,
        cross out,
        inner sep=0pt,
        minimum width=5pt,
        minimum height=5pt,
      },
    }
    \clip (-5.5,-9.5) rectangle (6.5,7.5);
    	
	\draw[color ={black},>=latex,->] (-4,0)--(4,0);
	\draw[color ={black},>=latex,->] (0,-8)--(0,6);
	\draw[color ={black},opacity = 0.4] (-4,1)--(4,1);
	\draw[color ={black},opacity = 0.4] (-4,2)--(4,2);
	\draw[color ={black},opacity = 0.4] (-4,3)--(4,3);
	\draw[color ={black},opacity = 0.4] (-4,4)--(4,4);
	\draw[color ={black},opacity = 0.4] (-4,5)--(4,5);
	\draw[color ={black},opacity = 0.4] (-4,6)--(4,6);
	\draw[color ={black},opacity = 0.4] (-4,-1)--(4,-1);
	\draw[color ={black},opacity = 0.4] (-4,-2)--(4,-2);
	\draw[color ={black},opacity = 0.4] (-4,-3)--(4,-3);
	\draw[color ={black},opacity = 0.4] (-4,-4)--(4,-4);
	\draw[color ={black},opacity = 0.4] (-4,-5)--(4,-5);
	\draw[color ={black},opacity = 0.4] (-4,-6)--(4,-6);
	\draw[color ={black},opacity = 0.4] (-4,-7)--(4,-7);
	\draw[color ={black},opacity = 0.4] (-4,-8)--(4,-8);
	\draw[color ={black},opacity = 0.4] (1,-8)--(1,6);
	\draw[color ={black},opacity = 0.4] (2,-8)--(2,6);
	\draw[color ={black},opacity = 0.4] (3,-8)--(3,6);
	\draw[color ={black},opacity = 0.4] (4,-8)--(4,6);
	\draw[color ={black},opacity = 0.4] (-1,-8)--(-1,6);
	\draw[color ={black},opacity = 0.4] (-2,-8)--(-2,6);
	\draw[color ={black},opacity = 0.4] (-3,-8)--(-3,6);
	\draw[color ={black},opacity = 0.4] (-4,-8)--(-4,6);
	\draw[color ={gray},opacity = 0.3] (-4,0)--(4,0);
	\draw[color ={gray},opacity = 0.3] (-4,0.5)--(4,0.5);
	\draw[color ={gray},opacity = 0.3] (-4,1.5)--(4,1.5);
	\draw[color ={gray},opacity = 0.3] (-4,2.5)--(4,2.5);
	\draw[color ={gray},opacity = 0.3] (-4,3.5)--(4,3.5);
	\draw[color ={gray},opacity = 0.3] (-4,4.5)--(4,4.5);
	\draw[color ={gray},opacity = 0.3] (-4,5.5)--(4,5.5);
	\draw[color ={gray},opacity = 0.3] (-4,0)--(4,0);
	\draw[color ={gray},opacity = 0.3] (-4,-0.5)--(4,-0.5);
	\draw[color ={gray},opacity = 0.3] (-4,-1.5)--(4,-1.5);
	\draw[color ={gray},opacity = 0.3] (-4,-2.5)--(4,-2.5);
	\draw[color ={gray},opacity = 0.3] (-4,-3.5)--(4,-3.5);
	\draw[color ={gray},opacity = 0.3] (-4,-4.5)--(4,-4.5);
	\draw[color ={gray},opacity = 0.3] (-4,-5.5)--(4,-5.5);
	\draw[color ={gray},opacity = 0.3] (-4,-6.5)--(4,-6.5);
	\draw[color ={gray},opacity = 0.3] (-4,-7)--(4,-7);
	\draw[color ={gray},opacity = 0.3] (-4,-7.5)--(4,-7.5);
	\draw[color ={gray},opacity = 0.3] (-4,-8)--(4,-8);
	\draw[color ={gray},opacity = 0.3] (0,-8)--(0,6);
	\draw[color ={gray},opacity = 0.3] (0,-8)--(0,6);
	\draw[color ={black},line width = 1.2] (1,-0.13)--(1,0.13);
	\draw[color ={black},line width = 1.2] (2,-0.13)--(2,0.13);
	\draw[color ={black},line width = 1.2] (3,-0.13)--(3,0.13);
	\draw[color ={black},line width = 1.2] (4,-0.13)--(4,0.13);
	\draw[color ={black},line width = 1.2] (-1,-0.13)--(-1,0.13);
	\draw[color ={black},line width = 1.2] (-2,-0.13)--(-2,0.13);
	\draw[color ={black},line width = 1.2] (-3,-0.13)--(-3,0.13);
	\draw[color ={black},line width = 1.2] (-4,-0.13)--(-4,0.13);
	\draw[color ={black},line width = 1.2] (-0.13,1)--(0.13,1);
	\draw[color ={black},line width = 1.2] (-0.13,2)--(0.13,2);
	\draw[color ={black},line width = 1.2] (-0.13,3)--(0.13,3);
	\draw[color ={black},line width = 1.2] (-0.13,4)--(0.13,4);
	\draw[color ={black},line width = 1.2] (-0.13,5)--(0.13,5);
	\draw[color ={black},line width = 1.2] (-0.13,6)--(0.13,6);
	\draw[color ={black},line width = 1.2] (-0.13,-1)--(0.13,-1);
	\draw[color ={black},line width = 1.2] (-0.13,-2)--(0.13,-2);
	\draw[color ={black},line width = 1.2] (-0.13,-3)--(0.13,-3);
	\draw[color ={black},line width = 1.2] (-0.13,-4)--(0.13,-4);
	\draw[color ={black},line width = 1.2] (-0.13,-5)--(0.13,-5);
	\draw[color ={black},line width = 1.2] (-0.13,-6)--(0.13,-6);
	\draw[color ={black},line width = 1.2] (-0.13,-7)--(0.13,-7);
	\draw[color ={black},line width = 1.2] (-0.13,-8)--(0.13,-8);
	\draw (2,-0.4) node[anchor = center] {\scriptsize \color{black}{$2$}};
	\draw (4,-0.4) node[anchor = center] {\scriptsize \color{black}{$4$}};
	\draw (-2,-0.4) node[anchor = center] {\scriptsize \color{black}{$-2$}};
	\draw (-4,-0.4) node[anchor = center] {\scriptsize \color{black}{$-4$}};
	\draw (-0.5,2.1) node[anchor = center] {\footnotesize \color{black}{$2$}};
	\draw (-0.5,4.1) node[anchor = center] {\footnotesize \color{black}{$4$}};
	\draw (-0.5,-1.9) node[anchor = center] {\footnotesize \color{black}{$-2$}};
	\draw (-0.5,-3.9) node[anchor = center] {\footnotesize \color{black}{$-4$}};
	\draw (-0.5,-5.9) node[anchor = center] {\footnotesize \color{black}{$-6$}};
	\draw (-0.5,-7.9) node[anchor = center] {\footnotesize \color{black}{$-8$}};
	\draw [color={black}] (-0.3,-0.3) node[anchor = center,scale=1, rotate = 0] {O};
	\draw (3,4) node[anchor = center] {\footnotesize \color{red}{$\Large\mathcal{C}_f\prime$}};
	
	\draw[color={red},line width = 2] (-2.4,6.8)--(-2.2,6.4)--(-2,6)--(-1.8,5.6)--(-1.6,5.2)--(-1.4,4.8)--(-1.2,4.4)--(-1,4)--(-0.8,3.6)--(-0.6,3.2)--(-0.4,2.8)--(-0.2,2.4)--(0,2)--(0.2,1.6)--(0.4,1.2)--(0.6,0.8)--(0.8,0.4)--(1,0)--(1.2,-0.4)--(1.4,-0.8)--(1.6,-1.2)--(1.8,-1.6)--(2,-2)--(2.2,-2.4)--(2.4,-2.8)--(2.6,-3.2)--(2.8,-3.6)--(3,-4)--(3.2,-4.4)--(3.4,-4.8)--(3.6,-5.2)--(3.8,-5.6)--(4,-6);

\end{tikzpicture}\\
    \end{minipage}
    
\end{exercice}

\begin{Solution}
 L'équation réduite de la tangente au point d'abscisse $0$ est  : $y=f'(0)(x-0)+f(0)$.\\
      On lit graphiquement $f(0)=2$ et $f'(0)=2$.\\
      L'équation réduite de la tangente est donc donnée par :
      $y=2(x+0)+2$, soit $y=2x+2$.
\end{Solution}

% @see : https://coopmaths.fr/alea?uuid=a1ba2&id=can1F14&n=1&d=10&alea=rkIp22323232323423456789101112131415161718192021222324252627282930313233343536373839404142434445464748234567891011121314151617181920212223242526272829&cd=1&cols=1
\needspace{10\baselineskip}
\begin{exercice}[BaremeTotal=false]


Soit $f$ la fonction définie sur $\mathbb{R}$ par : 
$f(x)= -10-5x^2-4x$.\\
En calculant la limite du taux d'accroissement, déterminer $f'(2)$.
\end{exercice}

\begin{Solution}
 $f$ est une fonction polynôme du second degré de la forme $f(x)=ax^2+bx+c$.\\
    La fonction dérivée est donnée par la somme des dérivées des fonctions $u$ et $v$ définies par $u(x)=-5x^2$ et $v(x)=-4x-10$.\\
     Comme $u'(x)=-10x$ et $v'(x)=-4$, on obtient  $f'(x)=-10x-4$. \\
     Ainsi, $f'(2)=-10\times 2-4=-24$.
\end{Solution}

% @see : https://coopmaths.fr/alea?uuid=3c690&id=can1F13&n=1&d=10&alea=EUDu22323232323423456789101112131415161718192021222324252627282930313233343536373839404142434445464748234567891011121314151617181920212223242526272829&cd=1&cols=1
\needspace{10\baselineskip}
\begin{exercice}[BaremeTotal=false]


 Déterminer le coefficient directeur de la tangente à la courbe représentative de la fonction racine carrée au point d'abscisse $12$.
         
\end{exercice}

\begin{Solution}
 Le coefficient directeur de la tangente au point d'abscisse $a$ est donné par le nombre dérivé $f'(a)$.\\
        La fonction $f$ définie par $f(x)=\sqrt{x}$ a pour fonction dérivée la fonction $f'$ définie par $f'(x)=\dfrac{1}{2\sqrt{x}}$.\\Comme $f'(12)=\dfrac{1}{2\sqrt{12}}$, le coefficient directeur de la tangente au point d'abscisse $12$ est : $\dfrac{1}{2\sqrt{12}}$.
\end{Solution}

% @see : https://coopmaths.fr/alea?uuid=ebd89&id=1AN14-1&n=1&d=10&s=3&alea=ocYr22323232323423456789101112131415161718192021222324252627282930313233343536373839404142434445464748234567891011121314151617181920212223242526272829&cd=0&cols=1
\needspace{10\baselineskip}
\begin{exercice}[BaremeTotal=false]


 Donner la dérivée de la fonction $f$, dérivable pour tout $x\in\R$, définie par  $f(x)=\dfrac{7x+6}{9}$.
\end{exercice}

\begin{Solution}
 L'expression de la dérivée de la fonction $f$ définie par $f(x)=\dfrac{7x+6}{9}$ est : ${\color[HTML]{f15929}\boldsymbol{f'(x)=\dfrac{7}{9}}}$.
\end{Solution}

% @see : https://coopmaths.fr/alea?uuid=ebd89&id=1AN14-1&n=1&d=10&s=4&alea=uAlP22323232323423456789101112131415161718192021222324252627282930313233343536373839404142434445464748234567891011121314151617181920212223242526272829&cd=0&cols=1
\needspace{10\baselineskip}
\begin{exercice}[BaremeTotal=false]


 Donner la dérivée de la fonction $f$, dérivable pour tout $x\in\R$, définie par  $f(x)=4x^{10}$.
\end{exercice}

\begin{Solution}
 L'expression de la dérivée de la fonction $f$ définie par $f(x)=4x^{10}$ est : ${\color[HTML]{f15929}\boldsymbol{f'(x)=40x^{9}}}$.
\end{Solution}

% @see : https://coopmaths.fr/alea?uuid=ebd89&id=1AN14-1&n=1&d=10&s=7&alea=vjN722323232323423456789101112131415161718192021222324252627282930313233343536373839404142434445464748234567891011121314151617181920212223242526272829&cd=0&cols=1
\needspace{10\baselineskip}
\begin{exercice}[BaremeTotal=false]


 Donner la dérivée de la fonction $f$, dérivable pour tout $x\in\R^*$, définie par  $f(x)=\dfrac{-3}{7x}$.
\end{exercice}

\begin{Solution}
 L'expression de la dérivée de la fonction $f$ définie par $f(x)=\dfrac{-3}{7x}$ est : ${\color[HTML]{f15929}\boldsymbol{f'(x)=\dfrac{3}{7x^2}}}$.
\end{Solution}

% @see : https://coopmaths.fr/alea?uuid=a83c0&id=1AN14-4&n=1&d=10&s=1&alea=Njr322323232323423456789101112131415161718192021222324252627282930313233343536373839404142434445464748234567891011121314151617181920212223242526272829&cd=0&cols=1
\needspace{10\baselineskip}
\begin{exercice}[BaremeTotal=false]


 Donner l'expression de la dérivée de la fonction $f$ définie sur $\R^{*}$ par $f(x)=-3x+\dfrac{1}{x}$.\\
\end{exercice}

\begin{Solution}
 L'expression de la dérivée de la fonction $f$ définie par $f(x)=-3x+\dfrac{1}{x}$ est :\\$f'(x)={\color[HTML]{f15929}\boldsymbol{-3-\dfrac{1}{x^2}}}$.\\
\end{Solution}

% @see : https://coopmaths.fr/alea?uuid=a83c0&id=1AN14-4&n=1&d=10&s=4&alea=4Vpz22323232323423456789101112131415161718192021222324252627282930313233343536373839404142434445464748234567891011121314151617181920212223242526272829&cd=1&cols=1
\needspace{10\baselineskip}
\begin{exercice}[BaremeTotal=false]


 Donner l'expression de la dérivée de la fonction $f$ définie sur $\R_+^{*}$ par $f(x)=2x-\dfrac{5}{x}+3\sqrt{x}$.\\
\end{exercice}

\begin{Solution}
 $(2x)^\prime=2$\\$(-\dfrac{5}{x})^\prime=\dfrac{5}{x^2}$\\$(3\sqrt{x})^\prime=\dfrac{3}{2\sqrt{x}}$\\L'expression de la dérivée de la fonction $f$ définie par $f(x)=2x-\dfrac{5}{x}+3\sqrt{x}$ est :\\$f'(x)={\color[HTML]{f15929}\boldsymbol{2+\dfrac{5}{x^2}+\dfrac{3}{2\sqrt{x}}}}$.\\
\end{Solution}

\end{Maquette}
\clearpage
\begin{Maquette}[DM]{Niveau= ,Classe=\getdata[30]\mydata\ (v30), Date = 06/01/25  ,Theme=Exercices}


% @see : https://coopmaths.fr/alea?uuid=4c8c7&id=1AN11-3&n=1&d=10&s=2&alea=nawO2232323232342345678910111213141516171819202122232425262728293031323334353637383940414243444546474823456789101112131415161718192021222324252627282930&cd=1&cols=2
\needspace{10\baselineskip}
\begin{exercice}[BaremeTotal=false]
\label{eleve:30}


 \begin{minipage}[t]{\linewidth} Soit $f$ une fonction dérivable sur $[-5;5]$ et $\mathcal{C}_f$ sa courbe représentative.\\On sait que  $f(-5)=-3~~$ et que $~~f'(-5)=-2$.\\Déterminer une équation de la tangente $(T)$ à la courbe $\mathcal{C}_f$ au point d'abscisse $-5$,\\sans utiliser la formule de cours de l'équation de tangente. \end{minipage}

\end{exercice}

\begin{Solution}
 \begin{minipage}[t]{\linewidth} On sait que la tangente n'est pas une droite verticale, puisque la fonction est dérivable sur l'intervalle.\\On en déduit que la tangente $(T)$ au point d'abscisse $-5$, admet une équation réduite de la forme :  $(T) : y=m x + p$. 

\medskip
$\bullet$ Détermination de $m$ :\\On sait que le nombre dérivé en $-5$ est par définition, le coefficient directeur de la tangente au point d'abscisse $-5$.\\Par conséquent, on a déjà : $m=f'(-5)=-2$.\\ On en déduit que  $(T) : y= -2 x + p$\\$\bullet$ Détermination de $p$ :\\ Pour cela, on utilise que si $f(-5)=-3~~$, alors le point $A$ de coordonnées $(-5;-3)$ appartient à $\mathcal{C}_f$ mais aussi à $(T)$.\\ On peut écrire $A(-5;-3) = \mathcal{C}_f \cap (T)$.\\ On remplace alors les coordonnées de $A(-5;-3)$ dans l'équation  $(T) : y= -2 x + p$.\\ $\begin{aligned} \phantom{\iff}&A(-5;-3)\in (T)\\ \iff& -3= -2 \times (-5)  + p\\ \iff& p=-3 +2 \times (-5)  \\ \iff& p=-3 -10   \\ \iff& p=-13\\\end{aligned}$\\On peut conclure que : $(T) : {\color[HTML]{f15929}\boldsymbol{y=-2x-13}}$. \end{minipage}

\end{Solution}

% @see : https://coopmaths.fr/alea?uuid=0e984&id=can1F15&n=1&d=10&alea=1Alo2232323232342345678910111213141516171819202122232425262728293031323334353637383940414243444546474823456789101112131415161718192021222324252627282930&cd=1&cols=1
\needspace{10\baselineskip}
\begin{exercice}[BaremeTotal=false]


 La courbe représente une fonction $f$ et la droite est la tangente au point d'abscisse $2$.\\
        Déterminer $f'(2)$.\begin{center}\begin{tikzpicture}[baseline,scale = 0.5]

    \tikzset{
      point/.style={
        thick,
        draw,
        cross out,
        inner sep=0pt,
        minimum width=5pt,
        minimum height=5pt,
      },
    }
    \clip (-2,-2) rectangle (14,12);
    	
	\draw[color ={black},line width = 2,>=latex,->] (-2,0)--(14,0);
	\draw[color ={black},line width = 2,>=latex,->] (0,-2)--(0,12);
	\draw[color ={black},opacity = 0.5] (-2,1)--(14,1);
	\draw[color ={black},opacity = 0.5] (-2,2)--(14,2);
	\draw[color ={black},opacity = 0.5] (-2,3)--(14,3);
	\draw[color ={black},opacity = 0.5] (-2,4)--(14,4);
	\draw[color ={black},opacity = 0.5] (-2,5)--(14,5);
	\draw[color ={black},opacity = 0.5] (-2,6)--(14,6);
	\draw[color ={black},opacity = 0.5] (-2,7)--(14,7);
	\draw[color ={black},opacity = 0.5] (-2,8)--(14,8);
	\draw[color ={black},opacity = 0.5] (-2,9)--(14,9);
	\draw[color ={black},opacity = 0.5] (-2,10)--(14,10);
	\draw[color ={black},opacity = 0.5] (-2,11)--(14,11);
	\draw[color ={black},opacity = 0.5] (-2,12)--(14,12);
	\draw[color ={black},opacity = 0.5] (-2,-1)--(14,-1);
	\draw[color ={black},opacity = 0.5] (-2,-2)--(14,-2);
	\draw[color ={black},opacity = 0.5] (1,-2)--(1,12);
	\draw[color ={black},opacity = 0.5] (2,-2)--(2,12);
	\draw[color ={black},opacity = 0.5] (3,-2)--(3,12);
	\draw[color ={black},opacity = 0.5] (4,-2)--(4,12);
	\draw[color ={black},opacity = 0.5] (5,-2)--(5,12);
	\draw[color ={black},opacity = 0.5] (6,-2)--(6,12);
	\draw[color ={black},opacity = 0.5] (7,-2)--(7,12);
	\draw[color ={black},opacity = 0.5] (8,-2)--(8,12);
	\draw[color ={black},opacity = 0.5] (9,-2)--(9,12);
	\draw[color ={black},opacity = 0.5] (10,-2)--(10,12);
	\draw[color ={black},opacity = 0.5] (11,-2)--(11,12);
	\draw[color ={black},opacity = 0.5] (12,-2)--(12,12);
	\draw[color ={black},opacity = 0.5] (13,-2)--(13,12);
	\draw[color ={black},opacity = 0.5] (14,-2)--(14,12);
	\draw[color ={black},opacity = 0.5] (-1,-2)--(-1,12);
	\draw[color ={black},opacity = 0.5] (-2,-2)--(-2,12);
	\draw[color ={gray},opacity = 0.3] (-2,0)--(14,0);
	\draw[color ={gray},opacity = 0.3] (-2,0.5)--(14,0.5);
	\draw[color ={gray},opacity = 0.3] (-2,1.5)--(14,1.5);
	\draw[color ={gray},opacity = 0.3] (-2,2.5)--(14,2.5);
	\draw[color ={gray},opacity = 0.3] (-2,3.5)--(14,3.5);
	\draw[color ={gray},opacity = 0.3] (-2,4.5)--(14,4.5);
	\draw[color ={gray},opacity = 0.3] (-2,5.5)--(14,5.5);
	\draw[color ={gray},opacity = 0.3] (-2,6.5)--(14,6.5);
	\draw[color ={gray},opacity = 0.3] (-2,7.5)--(14,7.5);
	\draw[color ={gray},opacity = 0.3] (-2,8.5)--(14,8.5);
	\draw[color ={gray},opacity = 0.3] (-2,9.5)--(14,9.5);
	\draw[color ={gray},opacity = 0.3] (-2,10.5)--(14,10.5);
	\draw[color ={gray},opacity = 0.3] (-2,11.5)--(14,11.5);
	\draw[color ={gray},opacity = 0.3] (-2,0)--(14,0);
	\draw[color ={gray},opacity = 0.3] (-2,-0.5)--(14,-0.5);
	\draw[color ={gray},opacity = 0.3] (-2,-1.5)--(14,-1.5);
	\draw[color ={gray},opacity = 0.3] (0,-2)--(0,12);
	\draw[color ={gray},opacity = 0.3] (0.1,-2)--(0.1,12);
	\draw[color ={gray},opacity = 0.3] (0.2,-2)--(0.2,12);
	\draw[color ={gray},opacity = 0.3] (0.30000000000000004,-2)--(0.30000000000000004,12);
	\draw[color ={gray},opacity = 0.3] (0.4,-2)--(0.4,12);
	\draw[color ={gray},opacity = 0.3] (0.5,-2)--(0.5,12);
	\draw[color ={gray},opacity = 0.3] (0.6,-2)--(0.6,12);
	\draw[color ={gray},opacity = 0.3] (0.7,-2)--(0.7,12);
	\draw[color ={gray},opacity = 0.3] (0.7999999999999999,-2)--(0.7999999999999999,12);
	\draw[color ={gray},opacity = 0.3] (0.8999999999999999,-2)--(0.8999999999999999,12);
	\draw[color ={gray},opacity = 0.3] (1.0999999999999999,-2)--(1.0999999999999999,12);
	\draw[color ={gray},opacity = 0.3] (1.2,-2)--(1.2,12);
	\draw[color ={gray},opacity = 0.3] (1.3,-2)--(1.3,12);
	\draw[color ={gray},opacity = 0.3] (1.4000000000000001,-2)--(1.4000000000000001,12);
	\draw[color ={gray},opacity = 0.3] (1.5000000000000002,-2)--(1.5000000000000002,12);
	\draw[color ={gray},opacity = 0.3] (1.6000000000000003,-2)--(1.6000000000000003,12);
	\draw[color ={gray},opacity = 0.3] (1.7000000000000004,-2)--(1.7000000000000004,12);
	\draw[color ={gray},opacity = 0.3] (1.8000000000000005,-2)--(1.8000000000000005,12);
	\draw[color ={gray},opacity = 0.3] (1.9000000000000006,-2)--(1.9000000000000006,12);
	\draw[color ={gray},opacity = 0.3] (2.1000000000000005,-2)--(2.1000000000000005,12);
	\draw[color ={gray},opacity = 0.3] (2.2000000000000006,-2)--(2.2000000000000006,12);
	\draw[color ={gray},opacity = 0.3] (2.3000000000000007,-2)--(2.3000000000000007,12);
	\draw[color ={gray},opacity = 0.3] (2.400000000000001,-2)--(2.400000000000001,12);
	\draw[color ={gray},opacity = 0.3] (2.500000000000001,-2)--(2.500000000000001,12);
	\draw[color ={gray},opacity = 0.3] (2.600000000000001,-2)--(2.600000000000001,12);
	\draw[color ={gray},opacity = 0.3] (2.700000000000001,-2)--(2.700000000000001,12);
	\draw[color ={gray},opacity = 0.3] (2.800000000000001,-2)--(2.800000000000001,12);
	\draw[color ={gray},opacity = 0.3] (2.9000000000000012,-2)--(2.9000000000000012,12);
	\draw[color ={gray},opacity = 0.3] (3.1000000000000014,-2)--(3.1000000000000014,12);
	\draw[color ={gray},opacity = 0.3] (3.2000000000000015,-2)--(3.2000000000000015,12);
	\draw[color ={gray},opacity = 0.3] (3.3000000000000016,-2)--(3.3000000000000016,12);
	\draw[color ={gray},opacity = 0.3] (3.4000000000000017,-2)--(3.4000000000000017,12);
	\draw[color ={gray},opacity = 0.3] (3.5000000000000018,-2)--(3.5000000000000018,12);
	\draw[color ={gray},opacity = 0.3] (3.600000000000002,-2)--(3.600000000000002,12);
	\draw[color ={gray},opacity = 0.3] (3.700000000000002,-2)--(3.700000000000002,12);
	\draw[color ={gray},opacity = 0.3] (3.800000000000002,-2)--(3.800000000000002,12);
	\draw[color ={gray},opacity = 0.3] (3.900000000000002,-2)--(3.900000000000002,12);
	\draw[color ={gray},opacity = 0.3] (4.100000000000001,-2)--(4.100000000000001,12);
	\draw[color ={gray},opacity = 0.3] (4.200000000000001,-2)--(4.200000000000001,12);
	\draw[color ={gray},opacity = 0.3] (4.300000000000001,-2)--(4.300000000000001,12);
	\draw[color ={gray},opacity = 0.3] (4.4,-2)--(4.4,12);
	\draw[color ={gray},opacity = 0.3] (4.5,-2)--(4.5,12);
	\draw[color ={gray},opacity = 0.3] (4.6,-2)--(4.6,12);
	\draw[color ={gray},opacity = 0.3] (4.699999999999999,-2)--(4.699999999999999,12);
	\draw[color ={gray},opacity = 0.3] (4.799999999999999,-2)--(4.799999999999999,12);
	\draw[color ={gray},opacity = 0.3] (4.899999999999999,-2)--(4.899999999999999,12);
	\draw[color ={gray},opacity = 0.3] (5.099999999999998,-2)--(5.099999999999998,12);
	\draw[color ={gray},opacity = 0.3] (5.1999999999999975,-2)--(5.1999999999999975,12);
	\draw[color ={gray},opacity = 0.3] (5.299999999999997,-2)--(5.299999999999997,12);
	\draw[color ={gray},opacity = 0.3] (5.399999999999997,-2)--(5.399999999999997,12);
	\draw[color ={gray},opacity = 0.3] (5.4999999999999964,-2)--(5.4999999999999964,12);
	\draw[color ={gray},opacity = 0.3] (5.599999999999996,-2)--(5.599999999999996,12);
	\draw[color ={gray},opacity = 0.3] (5.699999999999996,-2)--(5.699999999999996,12);
	\draw[color ={gray},opacity = 0.3] (5.799999999999995,-2)--(5.799999999999995,12);
	\draw[color ={gray},opacity = 0.3] (5.899999999999995,-2)--(5.899999999999995,12);
	\draw[color ={gray},opacity = 0.3] (6.099999999999994,-2)--(6.099999999999994,12);
	\draw[color ={gray},opacity = 0.3] (6.199999999999994,-2)--(6.199999999999994,12);
	\draw[color ={gray},opacity = 0.3] (6.299999999999994,-2)--(6.299999999999994,12);
	\draw[color ={gray},opacity = 0.3] (6.399999999999993,-2)--(6.399999999999993,12);
	\draw[color ={gray},opacity = 0.3] (6.499999999999993,-2)--(6.499999999999993,12);
	\draw[color ={gray},opacity = 0.3] (6.5999999999999925,-2)--(6.5999999999999925,12);
	\draw[color ={gray},opacity = 0.3] (6.699999999999992,-2)--(6.699999999999992,12);
	\draw[color ={gray},opacity = 0.3] (6.799999999999992,-2)--(6.799999999999992,12);
	\draw[color ={gray},opacity = 0.3] (6.8999999999999915,-2)--(6.8999999999999915,12);
	\draw[color ={gray},opacity = 0.3] (7.099999999999991,-2)--(7.099999999999991,12);
	\draw[color ={gray},opacity = 0.3] (7.19999999999999,-2)--(7.19999999999999,12);
	\draw[color ={gray},opacity = 0.3] (7.29999999999999,-2)--(7.29999999999999,12);
	\draw[color ={gray},opacity = 0.3] (7.39999999999999,-2)--(7.39999999999999,12);
	\draw[color ={gray},opacity = 0.3] (7.499999999999989,-2)--(7.499999999999989,12);
	\draw[color ={gray},opacity = 0.3] (7.599999999999989,-2)--(7.599999999999989,12);
	\draw[color ={gray},opacity = 0.3] (7.699999999999989,-2)--(7.699999999999989,12);
	\draw[color ={gray},opacity = 0.3] (7.799999999999988,-2)--(7.799999999999988,12);
	\draw[color ={gray},opacity = 0.3] (7.899999999999988,-2)--(7.899999999999988,12);
	\draw[color ={gray},opacity = 0.3] (8.099999999999987,-2)--(8.099999999999987,12);
	\draw[color ={gray},opacity = 0.3] (8.199999999999987,-2)--(8.199999999999987,12);
	\draw[color ={gray},opacity = 0.3] (8.299999999999986,-2)--(8.299999999999986,12);
	\draw[color ={gray},opacity = 0.3] (8.399999999999986,-2)--(8.399999999999986,12);
	\draw[color ={gray},opacity = 0.3] (8.499999999999986,-2)--(8.499999999999986,12);
	\draw[color ={gray},opacity = 0.3] (8.599999999999985,-2)--(8.599999999999985,12);
	\draw[color ={gray},opacity = 0.3] (8.699999999999985,-2)--(8.699999999999985,12);
	\draw[color ={gray},opacity = 0.3] (8.799999999999985,-2)--(8.799999999999985,12);
	\draw[color ={gray},opacity = 0.3] (8.899999999999984,-2)--(8.899999999999984,12);
	\draw[color ={gray},opacity = 0.3] (9.099999999999984,-2)--(9.099999999999984,12);
	\draw[color ={gray},opacity = 0.3] (9.199999999999983,-2)--(9.199999999999983,12);
	\draw[color ={gray},opacity = 0.3] (9.299999999999983,-2)--(9.299999999999983,12);
	\draw[color ={gray},opacity = 0.3] (9.399999999999983,-2)--(9.399999999999983,12);
	\draw[color ={gray},opacity = 0.3] (9.499999999999982,-2)--(9.499999999999982,12);
	\draw[color ={gray},opacity = 0.3] (9.599999999999982,-2)--(9.599999999999982,12);
	\draw[color ={gray},opacity = 0.3] (9.699999999999982,-2)--(9.699999999999982,12);
	\draw[color ={gray},opacity = 0.3] (9.799999999999981,-2)--(9.799999999999981,12);
	\draw[color ={gray},opacity = 0.3] (9.89999999999998,-2)--(9.89999999999998,12);
	\draw[color ={gray},opacity = 0.3] (10.09999999999998,-2)--(10.09999999999998,12);
	\draw[color ={gray},opacity = 0.3] (10.19999999999998,-2)--(10.19999999999998,12);
	\draw[color ={gray},opacity = 0.3] (10.29999999999998,-2)--(10.29999999999998,12);
	\draw[color ={gray},opacity = 0.3] (10.399999999999979,-2)--(10.399999999999979,12);
	\draw[color ={gray},opacity = 0.3] (10.499999999999979,-2)--(10.499999999999979,12);
	\draw[color ={gray},opacity = 0.3] (10.599999999999978,-2)--(10.599999999999978,12);
	\draw[color ={gray},opacity = 0.3] (10.699999999999978,-2)--(10.699999999999978,12);
	\draw[color ={gray},opacity = 0.3] (10.799999999999978,-2)--(10.799999999999978,12);
	\draw[color ={gray},opacity = 0.3] (10.899999999999977,-2)--(10.899999999999977,12);
	\draw[color ={gray},opacity = 0.3] (11.099999999999977,-2)--(11.099999999999977,12);
	\draw[color ={gray},opacity = 0.3] (11.199999999999976,-2)--(11.199999999999976,12);
	\draw[color ={gray},opacity = 0.3] (11.299999999999976,-2)--(11.299999999999976,12);
	\draw[color ={gray},opacity = 0.3] (11.399999999999975,-2)--(11.399999999999975,12);
	\draw[color ={gray},opacity = 0.3] (11.499999999999975,-2)--(11.499999999999975,12);
	\draw[color ={gray},opacity = 0.3] (11.599999999999975,-2)--(11.599999999999975,12);
	\draw[color ={gray},opacity = 0.3] (11.699999999999974,-2)--(11.699999999999974,12);
	\draw[color ={gray},opacity = 0.3] (11.799999999999974,-2)--(11.799999999999974,12);
	\draw[color ={gray},opacity = 0.3] (11.899999999999974,-2)--(11.899999999999974,12);
	\draw[color ={gray},opacity = 0.3] (12.099999999999973,-2)--(12.099999999999973,12);
	\draw[color ={gray},opacity = 0.3] (12.199999999999973,-2)--(12.199999999999973,12);
	\draw[color ={gray},opacity = 0.3] (12.299999999999972,-2)--(12.299999999999972,12);
	\draw[color ={gray},opacity = 0.3] (12.399999999999972,-2)--(12.399999999999972,12);
	\draw[color ={gray},opacity = 0.3] (12.499999999999972,-2)--(12.499999999999972,12);
	\draw[color ={gray},opacity = 0.3] (12.599999999999971,-2)--(12.599999999999971,12);
	\draw[color ={gray},opacity = 0.3] (12.69999999999997,-2)--(12.69999999999997,12);
	\draw[color ={gray},opacity = 0.3] (12.79999999999997,-2)--(12.79999999999997,12);
	\draw[color ={gray},opacity = 0.3] (12.89999999999997,-2)--(12.89999999999997,12);
	\draw[color ={gray},opacity = 0.3] (13.09999999999997,-2)--(13.09999999999997,12);
	\draw[color ={gray},opacity = 0.3] (13.199999999999969,-2)--(13.199999999999969,12);
	\draw[color ={gray},opacity = 0.3] (13.299999999999969,-2)--(13.299999999999969,12);
	\draw[color ={gray},opacity = 0.3] (13.399999999999968,-2)--(13.399999999999968,12);
	\draw[color ={gray},opacity = 0.3] (13.499999999999968,-2)--(13.499999999999968,12);
	\draw[color ={gray},opacity = 0.3] (13.599999999999968,-2)--(13.599999999999968,12);
	\draw[color ={gray},opacity = 0.3] (13.699999999999967,-2)--(13.699999999999967,12);
	\draw[color ={gray},opacity = 0.3] (13.799999999999967,-2)--(13.799999999999967,12);
	\draw[color ={gray},opacity = 0.3] (13.899999999999967,-2)--(13.899999999999967,12);
	\draw[color ={gray},opacity = 0.3] (0,-2)--(0,12);
	\draw[color ={gray},opacity = 0.3] (-0.1,-2)--(-0.1,12);
	\draw[color ={gray},opacity = 0.3] (-0.2,-2)--(-0.2,12);
	\draw[color ={gray},opacity = 0.3] (-0.30000000000000004,-2)--(-0.30000000000000004,12);
	\draw[color ={gray},opacity = 0.3] (-0.4,-2)--(-0.4,12);
	\draw[color ={gray},opacity = 0.3] (-0.5,-2)--(-0.5,12);
	\draw[color ={gray},opacity = 0.3] (-0.6,-2)--(-0.6,12);
	\draw[color ={gray},opacity = 0.3] (-0.7,-2)--(-0.7,12);
	\draw[color ={gray},opacity = 0.3] (-0.7999999999999999,-2)--(-0.7999999999999999,12);
	\draw[color ={gray},opacity = 0.3] (-0.8999999999999999,-2)--(-0.8999999999999999,12);
	\draw[color ={gray},opacity = 0.3] (-1.0999999999999999,-2)--(-1.0999999999999999,12);
	\draw[color ={gray},opacity = 0.3] (-1.2,-2)--(-1.2,12);
	\draw[color ={gray},opacity = 0.3] (-1.3,-2)--(-1.3,12);
	\draw[color ={gray},opacity = 0.3] (-1.4000000000000001,-2)--(-1.4000000000000001,12);
	\draw[color ={gray},opacity = 0.3] (-1.5000000000000002,-2)--(-1.5000000000000002,12);
	\draw[color ={gray},opacity = 0.3] (-1.6000000000000003,-2)--(-1.6000000000000003,12);
	\draw[color ={gray},opacity = 0.3] (-1.7000000000000004,-2)--(-1.7000000000000004,12);
	\draw[color ={gray},opacity = 0.3] (-1.8000000000000005,-2)--(-1.8000000000000005,12);
	\draw[color ={gray},opacity = 0.3] (-1.9000000000000006,-2)--(-1.9000000000000006,12);
	\draw[color ={black},line width = 2] (2,-0.2)--(2,0.2);
	\draw[color ={black},line width = 2] (4,-0.2)--(4,0.2);
	\draw[color ={black},line width = 2] (6,-0.2)--(6,0.2);
	\draw[color ={black},line width = 2] (8,-0.2)--(8,0.2);
	\draw[color ={black},line width = 2] (10,-0.2)--(10,0.2);
	\draw[color ={black},line width = 2] (12,-0.2)--(12,0.2);
	\draw[color ={black},line width = 2] (-0.2,2)--(0.2,2);
	\draw[color ={black},line width = 2] (-0.2,4)--(0.2,4);
	\draw[color ={black},line width = 2] (-0.2,6)--(0.2,6);
	\draw[color ={black},line width = 2] (-0.2,8)--(0.2,8);
	\draw[color ={black},line width = 2] (-0.2,10)--(0.2,10);
	\draw (2,-0.4) node[anchor = center] {\scriptsize \color{black}{$1$}};
	\draw (4,-0.4) node[anchor = center] {\scriptsize \color{black}{$2$}};
	\draw (6,-0.4) node[anchor = center] {\scriptsize \color{black}{$3$}};
	\draw (8,-0.4) node[anchor = center] {\scriptsize \color{black}{$4$}};
	\draw (10,-0.4) node[anchor = center] {\scriptsize \color{black}{$5$}};
	\draw (-0.5,2.1) node[anchor = center] {\footnotesize \color{black}{$1$}};
	\draw (-0.5,4.1) node[anchor = center] {\footnotesize \color{black}{$2$}};
	\draw (-0.5,6.1) node[anchor = center] {\footnotesize \color{black}{$3$}};
	\draw (-0.5,8.1) node[anchor = center] {\footnotesize \color{black}{$4$}};
	\draw (-0.5,10.1) node[anchor = center] {\footnotesize \color{black}{$5$}};
	\draw [color={black}] (-0.3,-0.3) node[anchor = center,scale=1, rotate = 0] {O};
	
	\draw[color={blue},line width = 2] (0.6,13.33)--(1,8)--(1.4,5.71)--(1.8,4.44)--(2.2,3.64)--(2.6,3.08)--(3,2.67)--(3.4,2.35)--(3.8,2.11)--(4.2,1.9)--(4.6,1.74)--(5,1.6)--(5.4,1.48)--(5.8,1.38)--(6.2,1.29)--(6.6,1.21)--(7,1.14)--(7.4,1.08)--(7.8,1.03)--(8.2,0.98)--(8.6,0.93)--(9,0.89)--(9.4,0.85)--(9.8,0.82)--(10.2,0.78)--(10.6,0.75)--(11,0.73)--(11.4,0.7)--(11.8,0.68)--(12.2,0.66)--(12.6,0.63)--(13,0.62)--(13.4,0.6)--(13.8,0.58);
	
	\draw[color={red},line width = 2] (-2,5)--(-1.6,4.8)--(-1.2,4.6)--(-0.8,4.4)--(-0.4,4.2)--(0,4)--(0.4,3.8)--(0.8,3.6)--(1.2,3.4)--(1.6,3.2)--(2,3)--(2.4,2.8)--(2.8,2.6)--(3.2,2.4)--(3.6,2.2)--(4,2)--(4.4,1.8)--(4.8,1.6)--(5.2,1.4)--(5.6,1.2)--(6,1)--(6.4,0.8)--(6.8,0.6)--(7.2,0.4)--(7.6,0.2)--(8,0)--(8.4,-0.2)--(8.8,-0.4)--(9.2,-0.6)--(9.6,-0.8)--(10,-1)--(10.4,-1.2)--(10.8,-1.4)--(11.2,-1.6)--(11.6,-1.8)--(12,-2)--(12.4,-2.2)--(12.8,-2.4)--(13.2,-2.6)--(13.6,-2.8)--(14,-3);

\end{tikzpicture}\end{center}
\end{exercice}

\begin{Solution}
 $f'(2)$ est donné par le coefficient directeur de la tangente à la courbe au point d'abscisse $2$, soit $\dfrac{-2}{4}=-\dfrac{1{\color[HTML]{2563a5}\boldsymbol{\times2}} }{2{\color[HTML]{2563a5}\boldsymbol{\times2}}}=-\dfrac{1}{2}$.
\end{Solution}

% @see : https://coopmaths.fr/alea?uuid=6f32d&id=can1F16&n=1&d=10&alea=dqSw2232323232342345678910111213141516171819202122232425262728293031323334353637383940414243444546474823456789101112131415161718192021222324252627282930&cd=1&cols=1
\needspace{10\baselineskip}
\newpage
\begin{exercice}[BaremeTotal=false]


 On donne les représentations graphiques d'une fonction et de sa dérivée.\\
        Donner l'équation réduite de la tangente à la courbe de $f$ en $x=3$. \\ \begin{minipage}{0.5\linewidth}
    \begin{tikzpicture}[baseline, scale=0.8]

    \tikzset{
      point/.style={
        thick,
        draw,
        cross out,
        inner sep=0pt,
        minimum width=5pt,
        minimum height=5pt,
      },
    }
    \clip (-4.5,-4.5) rectangle (5.75,13.5);
    	
	\draw[color ={black},>=latex,->] (-3,0)--(3,0);
	\draw[color ={black},>=latex,->] (0,-3)--(0,12);
	\draw[color ={black},opacity = 0.4] (-3,1)--(3,1);
	\draw[color ={black},opacity = 0.4] (-3,2)--(3,2);
	\draw[color ={black},opacity = 0.4] (-3,3)--(3,3);
	\draw[color ={black},opacity = 0.4] (-3,4)--(3,4);
	\draw[color ={black},opacity = 0.4] (-3,5)--(3,5);
	\draw[color ={black},opacity = 0.4] (-3,6)--(3,6);
	\draw[color ={black},opacity = 0.4] (-3,7)--(3,7);
	\draw[color ={black},opacity = 0.4] (-3,8)--(3,8);
	\draw[color ={black},opacity = 0.4] (-3,9)--(3,9);
	\draw[color ={black},opacity = 0.4] (-3,10)--(3,10);
	\draw[color ={black},opacity = 0.4] (-3,11)--(3,11);
	\draw[color ={black},opacity = 0.4] (-3,12)--(3,12);
	\draw[color ={black},opacity = 0.4] (-3,-1)--(3,-1);
	\draw[color ={black},opacity = 0.4] (-3,-2)--(3,-2);
	\draw[color ={black},opacity = 0.4] (-3,-3)--(3,-3);
	\draw[color ={black},opacity = 0.4] (1,-3)--(1,12);
	\draw[color ={black},opacity = 0.4] (2,-3)--(2,12);
	\draw[color ={black},opacity = 0.4] (3,-3)--(3,12);
	\draw[color ={black},opacity = 0.4] (-1,-3)--(-1,12);
	\draw[color ={black},opacity = 0.4] (-2,-3)--(-2,12);
	\draw[color ={black},opacity = 0.4] (-3,-3)--(-3,12);
	\draw[color ={gray},opacity = 0.3] (-3,0)--(3,0);
	\draw[color ={gray},opacity = 0.3] (-3,0.5)--(3,0.5);
	\draw[color ={gray},opacity = 0.3] (-3,1.5)--(3,1.5);
	\draw[color ={gray},opacity = 0.3] (-3,2.5)--(3,2.5);
	\draw[color ={gray},opacity = 0.3] (-3,3.5)--(3,3.5);
	\draw[color ={gray},opacity = 0.3] (-3,4.5)--(3,4.5);
	\draw[color ={gray},opacity = 0.3] (-3,5.5)--(3,5.5);
	\draw[color ={gray},opacity = 0.3] (-3,6.5)--(3,6.5);
	\draw[color ={gray},opacity = 0.3] (-3,7.5)--(3,7.5);
	\draw[color ={gray},opacity = 0.3] (-3,8.5)--(3,8.5);
	\draw[color ={gray},opacity = 0.3] (-3,9.5)--(3,9.5);
	\draw[color ={gray},opacity = 0.3] (-3,10.5)--(3,10.5);
	\draw[color ={gray},opacity = 0.3] (-3,11.5)--(3,11.5);
	\draw[color ={gray},opacity = 0.3] (-3,0)--(3,0);
	\draw[color ={gray},opacity = 0.3] (-3,-0.5)--(3,-0.5);
	\draw[color ={gray},opacity = 0.3] (-3,-1.5)--(3,-1.5);
	\draw[color ={gray},opacity = 0.3] (-3,-2.5)--(3,-2.5);
	\draw[color ={gray},opacity = 0.3] (0,-3)--(0,12);
	\draw[color ={gray},opacity = 0.3] (0,-3)--(0,12);
	\draw[color ={black},line width = 1.2] (1,-0.13)--(1,0.13);
	\draw[color ={black},line width = 1.2] (2,-0.13)--(2,0.13);
	\draw[color ={black},line width = 1.2] (3,-0.13)--(3,0.13);
	\draw[color ={black},line width = 1.2] (-1,-0.13)--(-1,0.13);
	\draw[color ={black},line width = 1.2] (-2,-0.13)--(-2,0.13);
	\draw[color ={black},line width = 1.2] (-3,-0.13)--(-3,0.13);
	\draw[color ={black},line width = 1.2] (-0.13,1)--(0.13,1);
	\draw[color ={black},line width = 1.2] (-0.13,2)--(0.13,2);
	\draw[color ={black},line width = 1.2] (-0.13,3)--(0.13,3);
	\draw[color ={black},line width = 1.2] (-0.13,4)--(0.13,4);
	\draw[color ={black},line width = 1.2] (-0.13,5)--(0.13,5);
	\draw[color ={black},line width = 1.2] (-0.13,6)--(0.13,6);
	\draw[color ={black},line width = 1.2] (-0.13,7)--(0.13,7);
	\draw[color ={black},line width = 1.2] (-0.13,8)--(0.13,8);
	\draw[color ={black},line width = 1.2] (-0.13,9)--(0.13,9);
	\draw[color ={black},line width = 1.2] (-0.13,10)--(0.13,10);
	\draw[color ={black},line width = 1.2] (-0.13,11)--(0.13,11);
	\draw[color ={black},line width = 1.2] (-0.13,12)--(0.13,12);
	\draw[color ={black},line width = 1.2] (-0.13,-1)--(0.13,-1);
	\draw[color ={black},line width = 1.2] (-0.13,-2)--(0.13,-2);
	\draw[color ={black},line width = 1.2] (-0.13,-3)--(0.13,-3);
	\draw (3,-0.4) node[anchor = center] {\scriptsize \color{black}{$3$}};
	\draw (-3,-0.4) node[anchor = center] {\scriptsize \color{black}{$-3$}};
	\draw (-0.5,3.1) node[anchor = center] {\footnotesize \color{black}{$3$}};
	\draw (-0.5,6.1) node[anchor = center] {\footnotesize \color{black}{$6$}};
	\draw (-0.5,9.1) node[anchor = center] {\footnotesize \color{black}{$9$}};
	\draw (-0.5,-2.9) node[anchor = center] {\footnotesize \color{black}{$-3$}};
	\draw [color={black}] (-0.3,-0.3) node[anchor = center,scale=1, rotate = 0] {O};
	\draw (3,10) node[anchor = center] {\footnotesize \color{blue}{$\Large \mathcal{C}_f$}};
	
	\draw[color={blue},line width = 2] (-1.4,12.56)--(-1.2,11.24)--(-1,10)--(-0.8,8.84)--(-0.6,7.76)--(-0.4,6.76)--(-0.2,5.84)--(0,5)--(0.2,4.24)--(0.4,3.56)--(0.6,2.96)--(0.8,2.44)--(1,2)--(1.2,1.64)--(1.4,1.36)--(1.6,1.16)--(1.8,1.04)--(2,1)--(2.2,1.04)--(2.4,1.16)--(2.6,1.36)--(2.8,1.64)--(3,2);

\end{tikzpicture}\\
    \end{minipage}
    \begin{minipage}{0.5\linewidth}
    \begin{tikzpicture}[baseline, scale=0.8]

    \tikzset{
      point/.style={
        thick,
        draw,
        cross out,
        inner sep=0pt,
        minimum width=5pt,
        minimum height=5pt,
      },
    }
    \clip (-4.5,-6.5) rectangle (6.5,9.5);
    	
	\draw[color ={black},>=latex,->] (-3,0)--(3,0);
	\draw[color ={black},>=latex,->] (0,-5)--(0,8);
	\draw[color ={black},opacity = 0.4] (-3,1)--(3,1);
	\draw[color ={black},opacity = 0.4] (-3,2)--(3,2);
	\draw[color ={black},opacity = 0.4] (-3,3)--(3,3);
	\draw[color ={black},opacity = 0.4] (-3,4)--(3,4);
	\draw[color ={black},opacity = 0.4] (-3,5)--(3,5);
	\draw[color ={black},opacity = 0.4] (-3,6)--(3,6);
	\draw[color ={black},opacity = 0.4] (-3,7)--(3,7);
	\draw[color ={black},opacity = 0.4] (-3,8)--(3,8);
	\draw[color ={black},opacity = 0.4] (-3,-1)--(3,-1);
	\draw[color ={black},opacity = 0.4] (-3,-2)--(3,-2);
	\draw[color ={black},opacity = 0.4] (-3,-3)--(3,-3);
	\draw[color ={black},opacity = 0.4] (-3,-4)--(3,-4);
	\draw[color ={black},opacity = 0.4] (-3,-5)--(3,-5);
	\draw[color ={black},opacity = 0.4] (1,-5)--(1,8);
	\draw[color ={black},opacity = 0.4] (2,-5)--(2,8);
	\draw[color ={black},opacity = 0.4] (3,-5)--(3,8);
	\draw[color ={black},opacity = 0.4] (-1,-5)--(-1,8);
	\draw[color ={black},opacity = 0.4] (-2,-5)--(-2,8);
	\draw[color ={black},opacity = 0.4] (-3,-5)--(-3,8);
	\draw[color ={gray},opacity = 0.3] (-3,0)--(3,0);
	\draw[color ={gray},opacity = 0.3] (-3,0.5)--(3,0.5);
	\draw[color ={gray},opacity = 0.3] (-3,1.5)--(3,1.5);
	\draw[color ={gray},opacity = 0.3] (-3,2.5)--(3,2.5);
	\draw[color ={gray},opacity = 0.3] (-3,3.5)--(3,3.5);
	\draw[color ={gray},opacity = 0.3] (-3,4.5)--(3,4.5);
	\draw[color ={gray},opacity = 0.3] (-3,5.5)--(3,5.5);
	\draw[color ={gray},opacity = 0.3] (-3,6.5)--(3,6.5);
	\draw[color ={gray},opacity = 0.3] (-3,7.5)--(3,7.5);
	\draw[color ={gray},opacity = 0.3] (-3,0)--(3,0);
	\draw[color ={gray},opacity = 0.3] (-3,-0.5)--(3,-0.5);
	\draw[color ={gray},opacity = 0.3] (-3,-1.5)--(3,-1.5);
	\draw[color ={gray},opacity = 0.3] (-3,-2.5)--(3,-2.5);
	\draw[color ={gray},opacity = 0.3] (-3,-3.5)--(3,-3.5);
	\draw[color ={gray},opacity = 0.3] (-3,-4.5)--(3,-4.5);
	\draw[color ={gray},opacity = 0.3] (0,-5)--(0,8);
	\draw[color ={gray},opacity = 0.3] (0,-5)--(0,8);
	\draw[color ={black},line width = 1.2] (1,-0.13)--(1,0.13);
	\draw[color ={black},line width = 1.2] (2,-0.13)--(2,0.13);
	\draw[color ={black},line width = 1.2] (3,-0.13)--(3,0.13);
	\draw[color ={black},line width = 1.2] (-1,-0.13)--(-1,0.13);
	\draw[color ={black},line width = 1.2] (-2,-0.13)--(-2,0.13);
	\draw[color ={black},line width = 1.2] (-3,-0.13)--(-3,0.13);
	\draw[color ={black},line width = 1.2] (-0.13,1)--(0.13,1);
	\draw[color ={black},line width = 1.2] (-0.13,2)--(0.13,2);
	\draw[color ={black},line width = 1.2] (-0.13,3)--(0.13,3);
	\draw[color ={black},line width = 1.2] (-0.13,4)--(0.13,4);
	\draw[color ={black},line width = 1.2] (-0.13,5)--(0.13,5);
	\draw[color ={black},line width = 1.2] (-0.13,6)--(0.13,6);
	\draw[color ={black},line width = 1.2] (-0.13,7)--(0.13,7);
	\draw[color ={black},line width = 1.2] (-0.13,8)--(0.13,8);
	\draw[color ={black},line width = 1.2] (-0.13,9)--(0.13,9);
	\draw[color ={black},line width = 1.2] (-0.13,10)--(0.13,10);
	\draw[color ={black},line width = 1.2] (-0.13,11)--(0.13,11);
	\draw[color ={black},line width = 1.2] (-0.13,12)--(0.13,12);
	\draw[color ={black},line width = 1.2] (-0.13,-1)--(0.13,-1);
	\draw[color ={black},line width = 1.2] (-0.13,-2)--(0.13,-2);
	\draw[color ={black},line width = 1.2] (-0.13,-3)--(0.13,-3);
	\draw (3,-0.4) node[anchor = center] {\scriptsize \color{black}{$3$}};
	\draw (-3,-0.4) node[anchor = center] {\scriptsize \color{black}{$-3$}};
	\draw (-0.5,3.1) node[anchor = center] {\footnotesize \color{black}{$3$}};
	\draw (-0.5,6.1) node[anchor = center] {\footnotesize \color{black}{$6$}};
	\draw (-0.5,-2.9) node[anchor = center] {\footnotesize \color{black}{$-3$}};
	\draw [color={black}] (-0.3,-0.3) node[anchor = center,scale=1, rotate = 0] {O};
	\draw (3,6) node[anchor = center] {\footnotesize \color{red}{$\Large\mathcal{C}_f\prime$}};
	
	\draw[color={red},line width = 2] (-1,-6)--(-0.8,-5.6)--(-0.6,-5.2)--(-0.4,-4.8)--(-0.2,-4.4)--(0,-4)--(0.2,-3.6)--(0.4,-3.2)--(0.6,-2.8)--(0.8,-2.4)--(1,-2)--(1.2,-1.6)--(1.4,-1.2)--(1.6,-0.8)--(1.8,-0.4)--(2,0)--(2.2,0.4)--(2.4,0.8)--(2.6,1.2)--(2.8,1.6)--(3,2);

\end{tikzpicture}\\
    \end{minipage}
    
\end{exercice}

\begin{Solution}
 L'équation réduite de la tangente au point d'abscisse $3$ est  : $y=f'(3)(x-3)+f(3)$.\\
        On lit graphiquement $f(3)=2$ et $f'(3)=2$.\\
        L'équation réduite de la tangente est donc donnée par :
        $y=2(x-3)+2$, soit $y=2x-4$.
\end{Solution}

% @see : https://coopmaths.fr/alea?uuid=a1ba2&id=can1F14&n=1&d=10&alea=rkIp2232323232342345678910111213141516171819202122232425262728293031323334353637383940414243444546474823456789101112131415161718192021222324252627282930&cd=1&cols=1
\needspace{10\baselineskip}
\begin{exercice}[BaremeTotal=false]


Soit $f$ la fonction définie sur $\mathbb{R}$ par : 
$f(x)= -7x^2+1$.\\
En calculant la limite du taux d'accroissement, déterminer $f'(-5)$.
\end{exercice}

\begin{Solution}
 $f$ est une fonction polynôme du second degré de la forme $f(x)=ax^2+b$.\\
    La fonction dérivée est donnée par la somme des dérivées des fonctions $u$ et $v$ définies par $u(x)=-7x^2$ et $v(x)=1$.\\
     Comme $u'(x)=-14x$ et $v'(x)=0$, on obtient  $f'(x)=-14x$. \\
     Ainsi, $f'(-5)=-14\times (-5)=70$.
\end{Solution}

% @see : https://coopmaths.fr/alea?uuid=3c690&id=can1F13&n=1&d=10&alea=EUDu2232323232342345678910111213141516171819202122232425262728293031323334353637383940414243444546474823456789101112131415161718192021222324252627282930&cd=1&cols=1
\needspace{10\baselineskip}
\begin{exercice}[BaremeTotal=false]


 Déterminer le coefficient directeur de la tangente à la courbe représentative de la fonction racine carrée au point d'abscisse $17$.
         
\end{exercice}

\begin{Solution}
 Le coefficient directeur de la tangente au point d'abscisse $a$ est donné par le nombre dérivé $f'(a)$.\\
        La fonction $f$ définie par $f(x)=\sqrt{x}$ a pour fonction dérivée la fonction $f'$ définie par $f'(x)=\dfrac{1}{2\sqrt{x}}$.\\Comme $f'(17)=\dfrac{1}{2\sqrt{17}}$, le coefficient directeur de la tangente au point d'abscisse $17$ est : $\dfrac{1}{2\sqrt{17}}$.
\end{Solution}

% @see : https://coopmaths.fr/alea?uuid=ebd89&id=1AN14-1&n=1&d=10&s=3&alea=ocYr2232323232342345678910111213141516171819202122232425262728293031323334353637383940414243444546474823456789101112131415161718192021222324252627282930&cd=0&cols=1
\needspace{10\baselineskip}
\begin{exercice}[BaremeTotal=false]


 Donner la dérivée de la fonction $f$, dérivable pour tout $x\in\R$, définie par  $f(x)=\dfrac{2}{3}x+7$.
\end{exercice}

\begin{Solution}
 L'expression de la dérivée de la fonction $f$ définie par $f(x)=\dfrac{2}{3}x+7$ est : ${\color[HTML]{f15929}\boldsymbol{f'(x)=\dfrac{2}{3}}}$.
\end{Solution}

% @see : https://coopmaths.fr/alea?uuid=ebd89&id=1AN14-1&n=1&d=10&s=4&alea=uAlP2232323232342345678910111213141516171819202122232425262728293031323334353637383940414243444546474823456789101112131415161718192021222324252627282930&cd=0&cols=1
\needspace{10\baselineskip}
\begin{exercice}[BaremeTotal=false]


 Donner la dérivée de la fonction $f$, dérivable pour tout $x\in\R$, définie par  $f(x)=-5x^{13}$.
\end{exercice}

\begin{Solution}
 L'expression de la dérivée de la fonction $f$ définie par $f(x)=-5x^{13}$ est : ${\color[HTML]{f15929}\boldsymbol{f'(x)=-65x^{12}}}$.
\end{Solution}

% @see : https://coopmaths.fr/alea?uuid=ebd89&id=1AN14-1&n=1&d=10&s=7&alea=vjN72232323232342345678910111213141516171819202122232425262728293031323334353637383940414243444546474823456789101112131415161718192021222324252627282930&cd=0&cols=1
\needspace{10\baselineskip}
\begin{exercice}[BaremeTotal=false]


 Donner la dérivée de la fonction $f$, dérivable pour tout $x\in\R^*$, définie par  $f(x)=\dfrac{-8}{9x}$.
\end{exercice}

\begin{Solution}
 L'expression de la dérivée de la fonction $f$ définie par $f(x)=\dfrac{-8}{9x}$ est : ${\color[HTML]{f15929}\boldsymbol{f'(x)=\dfrac{8}{9x^2}}}$.
\end{Solution}

% @see : https://coopmaths.fr/alea?uuid=a83c0&id=1AN14-4&n=1&d=10&s=1&alea=Njr32232323232342345678910111213141516171819202122232425262728293031323334353637383940414243444546474823456789101112131415161718192021222324252627282930&cd=0&cols=1
\needspace{10\baselineskip}
\begin{exercice}[BaremeTotal=false]


 Donner l'expression de la dérivée de la fonction $f$ définie sur $\R^{*}$ par $f(x)=\dfrac{7}{x}+4x^{2}$.\\
\end{exercice}

\begin{Solution}
 L'expression de la dérivée de la fonction $f$ définie par $f(x)=\dfrac{7}{x}+4x^{2}$ est :\\$f'(x)={\color[HTML]{f15929}\boldsymbol{-\dfrac{7}{x^2}+8x}}$.\\
\end{Solution}

% @see : https://coopmaths.fr/alea?uuid=a83c0&id=1AN14-4&n=1&d=10&s=4&alea=4Vpz2232323232342345678910111213141516171819202122232425262728293031323334353637383940414243444546474823456789101112131415161718192021222324252627282930&cd=1&cols=1
\needspace{10\baselineskip}
\begin{exercice}[BaremeTotal=false]


 Donner l'expression de la dérivée de la fonction $f$ définie sur $\R_+^{*}$ par $f(x)=\dfrac{5}{x}+3x^{2}-\sqrt{x}$.\\
\end{exercice}

\begin{Solution}
 $(\dfrac{5}{x})^\prime=-\dfrac{5}{x^2}$\\$(3x^{2})^\prime=6x$\\$(-\sqrt{x})^\prime=-\dfrac{1}{2\sqrt{x}}$\\L'expression de la dérivée de la fonction $f$ définie par $f(x)=\dfrac{5}{x}+3x^{2}-\sqrt{x}$ est :\\$f'(x)={\color[HTML]{f15929}\boldsymbol{-\dfrac{5}{x^2}+6x-\dfrac{1}{2\sqrt{x}}}}$.\\
\end{Solution}

\end{Maquette}
\clearpage
\begin{Maquette}[DM]{Niveau= ,Classe=\getdata[31]\mydata\ (v31), Date = 06/01/25  ,Theme=Exercices}


% @see : https://coopmaths.fr/alea?uuid=4c8c7&id=1AN11-3&n=1&d=10&s=2&alea=nawO223232323234234567891011121314151617181920212223242526272829303132333435363738394041424344454647482345678910111213141516171819202122232425262728293031&cd=1&cols=2
\needspace{10\baselineskip}
\begin{exercice}[BaremeTotal=false]
\label{eleve:31}


 \begin{minipage}[t]{\linewidth} Soit $f$ une fonction dérivable sur $[-5;5]$ et $\mathcal{C}_f$ sa courbe représentative.\\On sait que  $f(-2)=-5~~$ et que $~~f'(-2)=1$.\\Déterminer une équation de la tangente $(T)$ à la courbe $\mathcal{C}_f$ au point d'abscisse $-2$,\\sans utiliser la formule de cours de l'équation de tangente. \end{minipage}

\end{exercice}

\begin{Solution}
 \begin{minipage}[t]{\linewidth} On sait que la tangente n'est pas une droite verticale, puisque la fonction est dérivable sur l'intervalle.\\On en déduit que la tangente $(T)$ au point d'abscisse $-2$, admet une équation réduite de la forme :  $(T) : y=m x + p$. 

\medskip
$\bullet$ Détermination de $m$ :\\On sait que le nombre dérivé en $-2$ est par définition, le coefficient directeur de la tangente au point d'abscisse $-2$.\\Par conséquent, on a déjà : $m=f'(-2)=1$.\\ On en déduit que  $(T) : y= 1 x + p$\\$\bullet$ Détermination de $p$ :\\ Pour cela, on utilise que si $f(-2)=-5~~$, alors le point $A$ de coordonnées $(-2;-5)$ appartient à $\mathcal{C}_f$ mais aussi à $(T)$.\\ On peut écrire $A(-2;-5) = \mathcal{C}_f \cap (T)$.\\ On remplace alors les coordonnées de $A(-2;-5)$ dans l'équation  $(T) : y= 1 x + p$.\\ $\begin{aligned} \phantom{\iff}&A(-2;-5)\in (T)\\ \iff& -5= 1 \times (-2)  + p\\ \iff& p=-5 - \times (-2)  \\ \iff& p=-5 +2   \\ \iff& p=-3\\\end{aligned}$\\On peut conclure que : $(T) : {\color[HTML]{f15929}\boldsymbol{y=x-3}}$. \end{minipage}

\end{Solution}

% @see : https://coopmaths.fr/alea?uuid=0e984&id=can1F15&n=1&d=10&alea=1Alo223232323234234567891011121314151617181920212223242526272829303132333435363738394041424344454647482345678910111213141516171819202122232425262728293031&cd=1&cols=1
\needspace{10\baselineskip}
\begin{exercice}[BaremeTotal=false]


 La courbe représente une fonction $f$ et la droite est la tangente au point d'abscisse $1$.\\
        
        Déterminer $f'(1)$.\begin{center}\begin{tikzpicture}[baseline,scale = 0.5]

    \tikzset{
      point/.style={
        thick,
        draw,
        cross out,
        inner sep=0pt,
        minimum width=5pt,
        minimum height=5pt,
      },
    }
    \clip (-8,-3) rectangle (8,10);
    	
	\draw[color ={black},line width = 2,>=latex,->] (-8,0)--(8,0);
	\draw[color ={black},line width = 2,>=latex,->] (0,-3)--(0,10);
	\draw[color ={black},opacity = 0.5] (-8,1)--(8,1);
	\draw[color ={black},opacity = 0.5] (-8,2)--(8,2);
	\draw[color ={black},opacity = 0.5] (-8,3)--(8,3);
	\draw[color ={black},opacity = 0.5] (-8,4)--(8,4);
	\draw[color ={black},opacity = 0.5] (-8,5)--(8,5);
	\draw[color ={black},opacity = 0.5] (-8,6)--(8,6);
	\draw[color ={black},opacity = 0.5] (-8,7)--(8,7);
	\draw[color ={black},opacity = 0.5] (-8,8)--(8,8);
	\draw[color ={black},opacity = 0.5] (-8,9)--(8,9);
	\draw[color ={black},opacity = 0.5] (-8,10)--(8,10);
	\draw[color ={black},opacity = 0.5] (-8,-1)--(8,-1);
	\draw[color ={black},opacity = 0.5] (-8,-2)--(8,-2);
	\draw[color ={black},opacity = 0.5] (-8,-3)--(8,-3);
	\draw[color ={black},opacity = 0.5] (1,-3)--(1,10);
	\draw[color ={black},opacity = 0.5] (2,-3)--(2,10);
	\draw[color ={black},opacity = 0.5] (3,-3)--(3,10);
	\draw[color ={black},opacity = 0.5] (4,-3)--(4,10);
	\draw[color ={black},opacity = 0.5] (5,-3)--(5,10);
	\draw[color ={black},opacity = 0.5] (6,-3)--(6,10);
	\draw[color ={black},opacity = 0.5] (7,-3)--(7,10);
	\draw[color ={black},opacity = 0.5] (8,-3)--(8,10);
	\draw[color ={black},opacity = 0.5] (-1,-3)--(-1,10);
	\draw[color ={black},opacity = 0.5] (-2,-3)--(-2,10);
	\draw[color ={black},opacity = 0.5] (-3,-3)--(-3,10);
	\draw[color ={black},opacity = 0.5] (-4,-3)--(-4,10);
	\draw[color ={black},opacity = 0.5] (-5,-3)--(-5,10);
	\draw[color ={black},opacity = 0.5] (-6,-3)--(-6,10);
	\draw[color ={black},opacity = 0.5] (-7,-3)--(-7,10);
	\draw[color ={black},opacity = 0.5] (-8,-3)--(-8,10);
	\draw[color ={gray},opacity = 0.3] (-8,0)--(8,0);
	\draw[color ={gray},opacity = 0.3] (-8,0)--(8,0);
	\draw[color ={gray},opacity = 0.3] (0,-3)--(0,10);
	\draw[color ={gray},opacity = 0.3] (0,-3)--(0,10);
	\draw[color ={black},line width = 2] (2,-0.2)--(2,0.2);
	\draw[color ={black},line width = 2] (4,-0.2)--(4,0.2);
	\draw[color ={black},line width = 2] (6,-0.2)--(6,0.2);
	\draw[color ={black},line width = 2] (-2,-0.2)--(-2,0.2);
	\draw[color ={black},line width = 2] (-4,-0.2)--(-4,0.2);
	\draw[color ={black},line width = 2] (-6,-0.2)--(-6,0.2);
	\draw[color ={black},line width = 2] (-0.2,1)--(0.2,1);
	\draw[color ={black},line width = 2] (-0.2,2)--(0.2,2);
	\draw[color ={black},line width = 2] (-0.2,3)--(0.2,3);
	\draw[color ={black},line width = 2] (-0.2,4)--(0.2,4);
	\draw[color ={black},line width = 2] (-0.2,5)--(0.2,5);
	\draw[color ={black},line width = 2] (-0.2,6)--(0.2,6);
	\draw[color ={black},line width = 2] (-0.2,7)--(0.2,7);
	\draw[color ={black},line width = 2] (-0.2,8)--(0.2,8);
	\draw[color ={black},line width = 2] (-0.2,9)--(0.2,9);
	\draw[color ={black},line width = 2] (-0.2,-1)--(0.2,-1);
	\draw[color ={black},line width = 2] (-0.2,-2)--(0.2,-2);
	\draw (2,-0.4) node[anchor = center] {\scriptsize \color{black}{$1$}};
	\draw (4,-0.4) node[anchor = center] {\scriptsize \color{black}{$2$}};
	\draw (6,-0.4) node[anchor = center] {\scriptsize \color{black}{$3$}};
	\draw (-2,-0.4) node[anchor = center] {\scriptsize \color{black}{$-1$}};
	\draw (-4,-0.4) node[anchor = center] {\scriptsize \color{black}{$-2$}};
	\draw (-6,-0.4) node[anchor = center] {\scriptsize \color{black}{$-3$}};
	\draw (-0.8,2.1) node[anchor = center] {\footnotesize \color{black}{$2$}};
	\draw (-0.8,4.1) node[anchor = center] {\footnotesize \color{black}{$4$}};
	\draw (-0.8,6.1) node[anchor = center] {\footnotesize \color{black}{$6$}};
	\draw (-0.8,8.1) node[anchor = center] {\footnotesize \color{black}{$8$}};
	\draw [color={black}] (-0.3,-0.3) node[anchor = center,scale=1, rotate = 0] {O};
	
	\draw[color={blue},line width = 2] (-7.6,9.84)--(-7.2,8.76)--(-6.8,7.76)--(-6.4,6.84)--(-6,6)--(-5.6,5.24)--(-5.2,4.56)--(-4.8,3.96)--(-4.4,3.44)--(-4,3)--(-3.6,2.64)--(-3.2,2.36)--(-2.8,2.16)--(-2.4,2.04)--(-2,2)--(-1.6,2.04)--(-1.2,2.16)--(-0.8,2.36)--(-0.4,2.64)--(0,3)--(0.4,3.44)--(0.8,3.96)--(1.2,4.56)--(1.6,5.24)--(2,6)--(2.4,6.84)--(2.8,7.76)--(3.2,8.76)--(3.6,9.84);
	
	\draw[color={red},line width = 2] (-2.8,-3.6)--(-2.4,-2.8)--(-2,-2)--(-1.6,-1.2)--(-1.2,-0.4)--(-0.8,0.4)--(-0.4,1.2)--(0,2)--(0.4,2.8)--(0.8,3.6)--(1.2,4.4)--(1.6,5.2)--(2,6)--(2.4,6.8)--(2.8,7.6)--(3.2,8.4)--(3.6,9.2)--(4,10)--(4.4,10.8);

\end{tikzpicture}\end{center}
\end{exercice}

\begin{Solution}
 $f'(1)$ est donné par le coefficient directeur de la tangente à la courbe au point d'abscisse $1$, soit $4$.
\end{Solution}

% @see : https://coopmaths.fr/alea?uuid=6f32d&id=can1F16&n=1&d=10&alea=dqSw223232323234234567891011121314151617181920212223242526272829303132333435363738394041424344454647482345678910111213141516171819202122232425262728293031&cd=1&cols=1
\needspace{10\baselineskip}
\newpage
\begin{exercice}[BaremeTotal=false]


 On donne les représentations graphiques d'une fonction et de sa dérivée.\\
      Donner l'équation réduite de la tangente à la courbe de $f$ en $x=-1$. \\ \begin{minipage}{0.5\linewidth}
    \begin{tikzpicture}[baseline, scale=0.8]

    \tikzset{
      point/.style={
        thick,
        draw,
        cross out,
        inner sep=0pt,
        minimum width=5pt,
        minimum height=5pt,
      },
    }
    \clip (-5.5,-9.5) rectangle (5.75,7.5);
    	
	\draw[color ={black},>=latex,->] (-4,0)--(4,0);
	\draw[color ={black},>=latex,->] (0,-8)--(0,6);
	\draw[color ={black},opacity = 0.4] (-4,1)--(4,1);
	\draw[color ={black},opacity = 0.4] (-4,2)--(4,2);
	\draw[color ={black},opacity = 0.4] (-4,3)--(4,3);
	\draw[color ={black},opacity = 0.4] (-4,4)--(4,4);
	\draw[color ={black},opacity = 0.4] (-4,5)--(4,5);
	\draw[color ={black},opacity = 0.4] (-4,6)--(4,6);
	\draw[color ={black},opacity = 0.4] (-4,-1)--(4,-1);
	\draw[color ={black},opacity = 0.4] (-4,-2)--(4,-2);
	\draw[color ={black},opacity = 0.4] (-4,-3)--(4,-3);
	\draw[color ={black},opacity = 0.4] (-4,-4)--(4,-4);
	\draw[color ={black},opacity = 0.4] (-4,-5)--(4,-5);
	\draw[color ={black},opacity = 0.4] (-4,-6)--(4,-6);
	\draw[color ={black},opacity = 0.4] (-4,-7)--(4,-7);
	\draw[color ={black},opacity = 0.4] (-4,-8)--(4,-8);
	\draw[color ={black},opacity = 0.4] (1,-8)--(1,6);
	\draw[color ={black},opacity = 0.4] (2,-8)--(2,6);
	\draw[color ={black},opacity = 0.4] (3,-8)--(3,6);
	\draw[color ={black},opacity = 0.4] (4,-8)--(4,6);
	\draw[color ={black},opacity = 0.4] (-1,-8)--(-1,6);
	\draw[color ={black},opacity = 0.4] (-2,-8)--(-2,6);
	\draw[color ={black},opacity = 0.4] (-3,-8)--(-3,6);
	\draw[color ={black},opacity = 0.4] (-4,-8)--(-4,6);
	\draw[color ={gray},opacity = 0.3] (-4,0)--(4,0);
	\draw[color ={gray},opacity = 0.3] (-4,0.5)--(4,0.5);
	\draw[color ={gray},opacity = 0.3] (-4,1.5)--(4,1.5);
	\draw[color ={gray},opacity = 0.3] (-4,2.5)--(4,2.5);
	\draw[color ={gray},opacity = 0.3] (-4,3.5)--(4,3.5);
	\draw[color ={gray},opacity = 0.3] (-4,4.5)--(4,4.5);
	\draw[color ={gray},opacity = 0.3] (-4,5.5)--(4,5.5);
	\draw[color ={gray},opacity = 0.3] (-4,0)--(4,0);
	\draw[color ={gray},opacity = 0.3] (-4,-0.5)--(4,-0.5);
	\draw[color ={gray},opacity = 0.3] (-4,-1.5)--(4,-1.5);
	\draw[color ={gray},opacity = 0.3] (-4,-2.5)--(4,-2.5);
	\draw[color ={gray},opacity = 0.3] (-4,-3.5)--(4,-3.5);
	\draw[color ={gray},opacity = 0.3] (-4,-4.5)--(4,-4.5);
	\draw[color ={gray},opacity = 0.3] (-4,-5.5)--(4,-5.5);
	\draw[color ={gray},opacity = 0.3] (-4,-6.5)--(4,-6.5);
	\draw[color ={gray},opacity = 0.3] (-4,-7)--(4,-7);
	\draw[color ={gray},opacity = 0.3] (-4,-7.5)--(4,-7.5);
	\draw[color ={gray},opacity = 0.3] (-4,-8)--(4,-8);
	\draw[color ={gray},opacity = 0.3] (0,-8)--(0,6);
	\draw[color ={gray},opacity = 0.3] (0,-8)--(0,6);
	\draw[color ={black},line width = 1.2] (1,-0.13)--(1,0.13);
	\draw[color ={black},line width = 1.2] (2,-0.13)--(2,0.13);
	\draw[color ={black},line width = 1.2] (3,-0.13)--(3,0.13);
	\draw[color ={black},line width = 1.2] (4,-0.13)--(4,0.13);
	\draw[color ={black},line width = 1.2] (-1,-0.13)--(-1,0.13);
	\draw[color ={black},line width = 1.2] (-2,-0.13)--(-2,0.13);
	\draw[color ={black},line width = 1.2] (-3,-0.13)--(-3,0.13);
	\draw[color ={black},line width = 1.2] (-4,-0.13)--(-4,0.13);
	\draw[color ={black},line width = 1.2] (-0.13,1)--(0.13,1);
	\draw[color ={black},line width = 1.2] (-0.13,2)--(0.13,2);
	\draw[color ={black},line width = 1.2] (-0.13,3)--(0.13,3);
	\draw[color ={black},line width = 1.2] (-0.13,4)--(0.13,4);
	\draw[color ={black},line width = 1.2] (-0.13,5)--(0.13,5);
	\draw[color ={black},line width = 1.2] (-0.13,6)--(0.13,6);
	\draw[color ={black},line width = 1.2] (-0.13,-1)--(0.13,-1);
	\draw[color ={black},line width = 1.2] (-0.13,-2)--(0.13,-2);
	\draw[color ={black},line width = 1.2] (-0.13,-3)--(0.13,-3);
	\draw[color ={black},line width = 1.2] (-0.13,-4)--(0.13,-4);
	\draw[color ={black},line width = 1.2] (-0.13,-5)--(0.13,-5);
	\draw[color ={black},line width = 1.2] (-0.13,-6)--(0.13,-6);
	\draw[color ={black},line width = 1.2] (-0.13,-7)--(0.13,-7);
	\draw[color ={black},line width = 1.2] (-0.13,-8)--(0.13,-8);
	\draw (2,-0.4) node[anchor = center] {\scriptsize \color{black}{$2$}};
	\draw (4,-0.4) node[anchor = center] {\scriptsize \color{black}{$4$}};
	\draw (-2,-0.4) node[anchor = center] {\scriptsize \color{black}{$-2$}};
	\draw (-4,-0.4) node[anchor = center] {\scriptsize \color{black}{$-4$}};
	\draw (-0.5,2.1) node[anchor = center] {\footnotesize \color{black}{$2$}};
	\draw (-0.5,4.1) node[anchor = center] {\footnotesize \color{black}{$4$}};
	\draw (-0.5,6.1) node[anchor = center] {\footnotesize \color{black}{$6$}};
	\draw (-0.5,-1.9) node[anchor = center] {\footnotesize \color{black}{$-2$}};
	\draw (-0.5,-3.9) node[anchor = center] {\footnotesize \color{black}{$-4$}};
	\draw (-0.5,-5.9) node[anchor = center] {\footnotesize \color{black}{$-6$}};
	\draw (-0.5,-7.9) node[anchor = center] {\footnotesize \color{black}{$-8$}};
	\draw [color={black}] (-0.3,-0.3) node[anchor = center,scale=1, rotate = 0] {O};
	\draw (3,4) node[anchor = center] {\footnotesize \color{blue}{$\Large \mathcal{C}_f$}};
	
	\draw[color={blue},line width = 2] (-4,-2)--(-3.8,-1.24)--(-3.6,-0.56)--(-3.4,0.04)--(-3.2,0.56)--(-3,1)--(-2.8,1.36)--(-2.6,1.64)--(-2.4,1.84)--(-2.2,1.96)--(-2,2)--(-1.8,1.96)--(-1.6,1.84)--(-1.4,1.64)--(-1.2,1.36)--(-1,1)--(-0.8,0.56)--(-0.6,0.04)--(-0.4,-0.56)--(-0.2,-1.24)--(0,-2)--(0.2,-2.84)--(0.4,-3.76)--(0.6,-4.76)--(0.8,-5.84)--(1,-7)--(1.2,-8.24);

\end{tikzpicture}\\
    \end{minipage}
    \begin{minipage}{0.5\linewidth}
    \begin{tikzpicture}[baseline, scale=0.8]

    \tikzset{
      point/.style={
        thick,
        draw,
        cross out,
        inner sep=0pt,
        minimum width=5pt,
        minimum height=5pt,
      },
    }
    \clip (-5.5,-9.5) rectangle (6.5,7.5);
    	
	\draw[color ={black},>=latex,->] (-4,0)--(4,0);
	\draw[color ={black},>=latex,->] (0,-8)--(0,6);
	\draw[color ={black},opacity = 0.4] (-4,1)--(4,1);
	\draw[color ={black},opacity = 0.4] (-4,2)--(4,2);
	\draw[color ={black},opacity = 0.4] (-4,3)--(4,3);
	\draw[color ={black},opacity = 0.4] (-4,4)--(4,4);
	\draw[color ={black},opacity = 0.4] (-4,5)--(4,5);
	\draw[color ={black},opacity = 0.4] (-4,6)--(4,6);
	\draw[color ={black},opacity = 0.4] (-4,-1)--(4,-1);
	\draw[color ={black},opacity = 0.4] (-4,-2)--(4,-2);
	\draw[color ={black},opacity = 0.4] (-4,-3)--(4,-3);
	\draw[color ={black},opacity = 0.4] (-4,-4)--(4,-4);
	\draw[color ={black},opacity = 0.4] (-4,-5)--(4,-5);
	\draw[color ={black},opacity = 0.4] (-4,-6)--(4,-6);
	\draw[color ={black},opacity = 0.4] (-4,-7)--(4,-7);
	\draw[color ={black},opacity = 0.4] (-4,-8)--(4,-8);
	\draw[color ={black},opacity = 0.4] (1,-8)--(1,6);
	\draw[color ={black},opacity = 0.4] (2,-8)--(2,6);
	\draw[color ={black},opacity = 0.4] (3,-8)--(3,6);
	\draw[color ={black},opacity = 0.4] (4,-8)--(4,6);
	\draw[color ={black},opacity = 0.4] (-1,-8)--(-1,6);
	\draw[color ={black},opacity = 0.4] (-2,-8)--(-2,6);
	\draw[color ={black},opacity = 0.4] (-3,-8)--(-3,6);
	\draw[color ={black},opacity = 0.4] (-4,-8)--(-4,6);
	\draw[color ={gray},opacity = 0.3] (-4,0)--(4,0);
	\draw[color ={gray},opacity = 0.3] (-4,0.5)--(4,0.5);
	\draw[color ={gray},opacity = 0.3] (-4,1.5)--(4,1.5);
	\draw[color ={gray},opacity = 0.3] (-4,2.5)--(4,2.5);
	\draw[color ={gray},opacity = 0.3] (-4,3.5)--(4,3.5);
	\draw[color ={gray},opacity = 0.3] (-4,4.5)--(4,4.5);
	\draw[color ={gray},opacity = 0.3] (-4,5.5)--(4,5.5);
	\draw[color ={gray},opacity = 0.3] (-4,0)--(4,0);
	\draw[color ={gray},opacity = 0.3] (-4,-0.5)--(4,-0.5);
	\draw[color ={gray},opacity = 0.3] (-4,-1.5)--(4,-1.5);
	\draw[color ={gray},opacity = 0.3] (-4,-2.5)--(4,-2.5);
	\draw[color ={gray},opacity = 0.3] (-4,-3.5)--(4,-3.5);
	\draw[color ={gray},opacity = 0.3] (-4,-4.5)--(4,-4.5);
	\draw[color ={gray},opacity = 0.3] (-4,-5.5)--(4,-5.5);
	\draw[color ={gray},opacity = 0.3] (-4,-6.5)--(4,-6.5);
	\draw[color ={gray},opacity = 0.3] (-4,-7)--(4,-7);
	\draw[color ={gray},opacity = 0.3] (-4,-7.5)--(4,-7.5);
	\draw[color ={gray},opacity = 0.3] (-4,-8)--(4,-8);
	\draw[color ={gray},opacity = 0.3] (0,-8)--(0,6);
	\draw[color ={gray},opacity = 0.3] (0,-8)--(0,6);
	\draw[color ={black},line width = 1.2] (1,-0.13)--(1,0.13);
	\draw[color ={black},line width = 1.2] (2,-0.13)--(2,0.13);
	\draw[color ={black},line width = 1.2] (3,-0.13)--(3,0.13);
	\draw[color ={black},line width = 1.2] (4,-0.13)--(4,0.13);
	\draw[color ={black},line width = 1.2] (-1,-0.13)--(-1,0.13);
	\draw[color ={black},line width = 1.2] (-2,-0.13)--(-2,0.13);
	\draw[color ={black},line width = 1.2] (-3,-0.13)--(-3,0.13);
	\draw[color ={black},line width = 1.2] (-4,-0.13)--(-4,0.13);
	\draw[color ={black},line width = 1.2] (-0.13,1)--(0.13,1);
	\draw[color ={black},line width = 1.2] (-0.13,2)--(0.13,2);
	\draw[color ={black},line width = 1.2] (-0.13,3)--(0.13,3);
	\draw[color ={black},line width = 1.2] (-0.13,4)--(0.13,4);
	\draw[color ={black},line width = 1.2] (-0.13,5)--(0.13,5);
	\draw[color ={black},line width = 1.2] (-0.13,6)--(0.13,6);
	\draw[color ={black},line width = 1.2] (-0.13,-1)--(0.13,-1);
	\draw[color ={black},line width = 1.2] (-0.13,-2)--(0.13,-2);
	\draw[color ={black},line width = 1.2] (-0.13,-3)--(0.13,-3);
	\draw[color ={black},line width = 1.2] (-0.13,-4)--(0.13,-4);
	\draw[color ={black},line width = 1.2] (-0.13,-5)--(0.13,-5);
	\draw[color ={black},line width = 1.2] (-0.13,-6)--(0.13,-6);
	\draw[color ={black},line width = 1.2] (-0.13,-7)--(0.13,-7);
	\draw[color ={black},line width = 1.2] (-0.13,-8)--(0.13,-8);
	\draw (2,-0.4) node[anchor = center] {\scriptsize \color{black}{$2$}};
	\draw (4,-0.4) node[anchor = center] {\scriptsize \color{black}{$4$}};
	\draw (-2,-0.4) node[anchor = center] {\scriptsize \color{black}{$-2$}};
	\draw (-4,-0.4) node[anchor = center] {\scriptsize \color{black}{$-4$}};
	\draw (-0.5,2.1) node[anchor = center] {\footnotesize \color{black}{$2$}};
	\draw (-0.5,4.1) node[anchor = center] {\footnotesize \color{black}{$4$}};
	\draw (-0.5,-1.9) node[anchor = center] {\footnotesize \color{black}{$-2$}};
	\draw (-0.5,-3.9) node[anchor = center] {\footnotesize \color{black}{$-4$}};
	\draw (-0.5,-5.9) node[anchor = center] {\footnotesize \color{black}{$-6$}};
	\draw (-0.5,-7.9) node[anchor = center] {\footnotesize \color{black}{$-8$}};
	\draw [color={black}] (-0.3,-0.3) node[anchor = center,scale=1, rotate = 0] {O};
	\draw (3,4) node[anchor = center] {\footnotesize \color{red}{$\Large\mathcal{C}_f\prime$}};
	
	\draw[color={red},line width = 2] (-4,4)--(-3.8,3.6)--(-3.6,3.2)--(-3.4,2.8)--(-3.2,2.4)--(-3,2)--(-2.8,1.6)--(-2.6,1.2)--(-2.4,0.8)--(-2.2,0.4)--(-2,0)--(-1.8,-0.4)--(-1.6,-0.8)--(-1.4,-1.2)--(-1.2,-1.6)--(-1,-2)--(-0.8,-2.4)--(-0.6,-2.8)--(-0.4,-3.2)--(-0.2,-3.6)--(0,-4)--(0.2,-4.4)--(0.4,-4.8)--(0.6,-5.2)--(0.8,-5.6)--(1,-6)--(1.2,-6.4)--(1.4,-6.8)--(1.6,-7.2)--(1.8,-7.6)--(2,-8)--(2.2,-8.4)--(2.4,-8.8);

\end{tikzpicture}\\
    \end{minipage}
    
\end{exercice}

\begin{Solution}
 L'équation réduite de la tangente au point d'abscisse $-1$ est  : $y=f'(-1)(x-(-1))+f(-1)$.\\
      On lit graphiquement $f(-1)=1$ et $f'(-1)=-2$.\\
      L'équation réduite de la tangente est donc donnée par :
      $y=-2(x+1)+1$, soit $y=-2x-1$.
\end{Solution}

% @see : https://coopmaths.fr/alea?uuid=a1ba2&id=can1F14&n=1&d=10&alea=rkIp223232323234234567891011121314151617181920212223242526272829303132333435363738394041424344454647482345678910111213141516171819202122232425262728293031&cd=1&cols=1
\needspace{10\baselineskip}
\begin{exercice}[BaremeTotal=false]


Soit $f$ la fonction définie sur $\mathbb{R}$ par :
$f(x)= 8+4x^2$.\\
En calculant la limite du taux d'accroissement, déterminer $f'(0)$.
\end{exercice}

\begin{Solution}
 $f$ est une fonction polynôme du second degré de la forme $f(x)=ax^2+b$.\\
    La fonction dérivée est donnée par la somme des dérivées des fonctions $u$ et $v$ définies par $u(x)=4x^2$ et $v(x)=8$.\\
     Comme $u'(x)=8x$ et $v'(x)=0$, on obtient  $f'(x)=8x$. \\
     Ainsi, $f'(0)=8\times 0=0$.
\end{Solution}

% @see : https://coopmaths.fr/alea?uuid=3c690&id=can1F13&n=1&d=10&alea=EUDu223232323234234567891011121314151617181920212223242526272829303132333435363738394041424344454647482345678910111213141516171819202122232425262728293031&cd=1&cols=1
\needspace{10\baselineskip}
\begin{exercice}[BaremeTotal=false]


 Déterminer le coefficient directeur de la tangente à la courbe représentative de la fonction racine carrée au point d'abscisse $14$.
         
\end{exercice}

\begin{Solution}
 Le coefficient directeur de la tangente au point d'abscisse $a$ est donné par le nombre dérivé $f'(a)$.\\
        La fonction $f$ définie par $f(x)=\sqrt{x}$ a pour fonction dérivée la fonction $f'$ définie par $f'(x)=\dfrac{1}{2\sqrt{x}}$.\\Comme $f'(14)=\dfrac{1}{2\sqrt{14}}$, le coefficient directeur de la tangente au point d'abscisse $14$ est : $\dfrac{1}{2\sqrt{14}}$.
\end{Solution}

% @see : https://coopmaths.fr/alea?uuid=ebd89&id=1AN14-1&n=1&d=10&s=3&alea=ocYr223232323234234567891011121314151617181920212223242526272829303132333435363738394041424344454647482345678910111213141516171819202122232425262728293031&cd=0&cols=1
\needspace{10\baselineskip}
\begin{exercice}[BaremeTotal=false]


 Donner la dérivée de la fonction $f$, dérivable pour tout $x\in\R$, définie par  $f(x)=-\dfrac{3}{4}x+10$.
\end{exercice}

\begin{Solution}
 L'expression de la dérivée de la fonction $f$ définie par $f(x)=-\dfrac{3}{4}x+10$ est : ${\color[HTML]{f15929}\boldsymbol{f'(x)=-\dfrac{3}{4}}}$.
\end{Solution}

% @see : https://coopmaths.fr/alea?uuid=ebd89&id=1AN14-1&n=1&d=10&s=4&alea=uAlP223232323234234567891011121314151617181920212223242526272829303132333435363738394041424344454647482345678910111213141516171819202122232425262728293031&cd=0&cols=1
\needspace{10\baselineskip}
\begin{exercice}[BaremeTotal=false]


 Donner la dérivée de la fonction $f$, dérivable pour tout $x\in\R$, définie par  $f(x)=-3x^{15}$.
\end{exercice}

\begin{Solution}
 L'expression de la dérivée de la fonction $f$ définie par $f(x)=-3x^{15}$ est : ${\color[HTML]{f15929}\boldsymbol{f'(x)=-45x^{14}}}$.
\end{Solution}

% @see : https://coopmaths.fr/alea?uuid=ebd89&id=1AN14-1&n=1&d=10&s=7&alea=vjN7223232323234234567891011121314151617181920212223242526272829303132333435363738394041424344454647482345678910111213141516171819202122232425262728293031&cd=0&cols=1
\needspace{10\baselineskip}
\begin{exercice}[BaremeTotal=false]


 Donner la dérivée de la fonction $f$, dérivable pour tout $x\in\R^*$, définie par  $f(x)=\dfrac{-7}{8x}$.
\end{exercice}

\begin{Solution}
 L'expression de la dérivée de la fonction $f$ définie par $f(x)=\dfrac{-7}{8x}$ est : ${\color[HTML]{f15929}\boldsymbol{f'(x)=\dfrac{7}{8x^2}}}$.
\end{Solution}

% @see : https://coopmaths.fr/alea?uuid=a83c0&id=1AN14-4&n=1&d=10&s=1&alea=Njr3223232323234234567891011121314151617181920212223242526272829303132333435363738394041424344454647482345678910111213141516171819202122232425262728293031&cd=0&cols=1
\needspace{10\baselineskip}
\begin{exercice}[BaremeTotal=false]


 Donner l'expression de la dérivée de la fonction $f$ définie sur $\R^{*}$ par $f(x)=4x^{2}-2x-4+\dfrac{9}{x}$.\\
\end{exercice}

\begin{Solution}
 L'expression de la dérivée de la fonction $f$ définie par $f(x)=4x^{2}-2x-4+\dfrac{9}{x}$ est :\\$f'(x)={\color[HTML]{f15929}\boldsymbol{8x-2-\dfrac{9}{x^2}}}$.\\
\end{Solution}

% @see : https://coopmaths.fr/alea?uuid=a83c0&id=1AN14-4&n=1&d=10&s=4&alea=4Vpz223232323234234567891011121314151617181920212223242526272829303132333435363738394041424344454647482345678910111213141516171819202122232425262728293031&cd=1&cols=1
\needspace{10\baselineskip}
\begin{exercice}[BaremeTotal=false]


 Donner l'expression de la dérivée de la fonction $f$ définie sur $\R_+^{*}$ par $f(x)=2\sqrt{x}-2x^{2}+5x-5-\dfrac{1}{x}$.\\
\end{exercice}

\begin{Solution}
 $(2\sqrt{x})^\prime=\dfrac{2}{2\sqrt{x}}$\\$(-2x^{2}+5x-5)^\prime=-4x+5$\\$(-\dfrac{1}{x})^\prime=\dfrac{1}{x^2}$\\L'expression de la dérivée de la fonction $f$ définie par $f(x)=2\sqrt{x}-2x^{2}+5x-5-\dfrac{1}{x}$ est :\\$f'(x)={\color[HTML]{f15929}\boldsymbol{\dfrac{2}{2\sqrt{x}}-4x+5+\dfrac{1}{x^2}}}$.\\
\end{Solution}

\end{Maquette}
\clearpage
\begin{Maquette}[DM]{Niveau= ,Classe=\getdata[32]\mydata\ (v32), Date = 06/01/25  ,Theme=Exercices}


% @see : https://coopmaths.fr/alea?uuid=4c8c7&id=1AN11-3&n=1&d=10&s=2&alea=nawO22323232323423456789101112131415161718192021222324252627282930313233343536373839404142434445464748234567891011121314151617181920212223242526272829303132&cd=1&cols=2
\needspace{10\baselineskip}
\begin{exercice}[BaremeTotal=false]
\label{eleve:32}


 \begin{minipage}[t]{\linewidth} Soit $f$ une fonction dérivable sur $[-5;5]$ et $\mathcal{C}_f$ sa courbe représentative.\\On sait que  $f(4)=-5~~$ et que $~~f'(4)=1$.\\Déterminer une équation de la tangente $(T)$ à la courbe $\mathcal{C}_f$ au point d'abscisse $4$,\\sans utiliser la formule de cours de l'équation de tangente. \end{minipage}

\end{exercice}

\begin{Solution}
 \begin{minipage}[t]{\linewidth} On sait que la tangente n'est pas une droite verticale, puisque la fonction est dérivable sur l'intervalle.\\On en déduit que la tangente $(T)$ au point d'abscisse $4$, admet une équation réduite de la forme :  $(T) : y=m x + p$. 

\medskip
$\bullet$ Détermination de $m$ :\\On sait que le nombre dérivé en $4$ est par définition, le coefficient directeur de la tangente au point d'abscisse $4$.\\Par conséquent, on a déjà : $m=f'(4)=1$.\\ On en déduit que  $(T) : y= 1 x + p$\\$\bullet$ Détermination de $p$ :\\ Pour cela, on utilise que si $f(4)=-5~~$, alors le point $A$ de coordonnées $(4;-5)$ appartient à $\mathcal{C}_f$ mais aussi à $(T)$.\\ On peut écrire $A(4;-5) = \mathcal{C}_f \cap (T)$.\\ On remplace alors les coordonnées de $A(4;-5)$ dans l'équation  $(T) : y= 1 x + p$.\\ $\begin{aligned} \phantom{\iff}&A(4;-5)\in (T)\\ \iff& -5= 1 \times 4  + p\\ \iff& p=-5 - \times 4  \\ \iff& p=-5 -4   \\ \iff& p=-9\\\end{aligned}$\\On peut conclure que : $(T) : {\color[HTML]{f15929}\boldsymbol{y=x-9}}$. \end{minipage}

\end{Solution}

% @see : https://coopmaths.fr/alea?uuid=0e984&id=can1F15&n=1&d=10&alea=1Alo22323232323423456789101112131415161718192021222324252627282930313233343536373839404142434445464748234567891011121314151617181920212223242526272829303132&cd=1&cols=1
\needspace{10\baselineskip}
\begin{exercice}[BaremeTotal=false]


 La courbe représente une fonction $f$ et la droite est la tangente au point d'abscisse $1$.\\
        Déterminer $f'(1)$.\begin{center}\begin{tikzpicture}[baseline,scale = 0.5]

    \tikzset{
      point/.style={
        thick,
        draw,
        cross out,
        inner sep=0pt,
        minimum width=5pt,
        minimum height=5pt,
      },
    }
    \clip (-2,-2) rectangle (14,12);
    	
	\draw[color ={black},line width = 2,>=latex,->] (-2,0)--(14,0);
	\draw[color ={black},line width = 2,>=latex,->] (0,-2)--(0,12);
	\draw[color ={black},opacity = 0.5] (-2,1)--(14,1);
	\draw[color ={black},opacity = 0.5] (-2,2)--(14,2);
	\draw[color ={black},opacity = 0.5] (-2,3)--(14,3);
	\draw[color ={black},opacity = 0.5] (-2,4)--(14,4);
	\draw[color ={black},opacity = 0.5] (-2,5)--(14,5);
	\draw[color ={black},opacity = 0.5] (-2,6)--(14,6);
	\draw[color ={black},opacity = 0.5] (-2,7)--(14,7);
	\draw[color ={black},opacity = 0.5] (-2,8)--(14,8);
	\draw[color ={black},opacity = 0.5] (-2,9)--(14,9);
	\draw[color ={black},opacity = 0.5] (-2,10)--(14,10);
	\draw[color ={black},opacity = 0.5] (-2,11)--(14,11);
	\draw[color ={black},opacity = 0.5] (-2,12)--(14,12);
	\draw[color ={black},opacity = 0.5] (-2,-1)--(14,-1);
	\draw[color ={black},opacity = 0.5] (-2,-2)--(14,-2);
	\draw[color ={black},opacity = 0.5] (1,-2)--(1,12);
	\draw[color ={black},opacity = 0.5] (2,-2)--(2,12);
	\draw[color ={black},opacity = 0.5] (3,-2)--(3,12);
	\draw[color ={black},opacity = 0.5] (4,-2)--(4,12);
	\draw[color ={black},opacity = 0.5] (5,-2)--(5,12);
	\draw[color ={black},opacity = 0.5] (6,-2)--(6,12);
	\draw[color ={black},opacity = 0.5] (7,-2)--(7,12);
	\draw[color ={black},opacity = 0.5] (8,-2)--(8,12);
	\draw[color ={black},opacity = 0.5] (9,-2)--(9,12);
	\draw[color ={black},opacity = 0.5] (10,-2)--(10,12);
	\draw[color ={black},opacity = 0.5] (11,-2)--(11,12);
	\draw[color ={black},opacity = 0.5] (12,-2)--(12,12);
	\draw[color ={black},opacity = 0.5] (13,-2)--(13,12);
	\draw[color ={black},opacity = 0.5] (14,-2)--(14,12);
	\draw[color ={black},opacity = 0.5] (-1,-2)--(-1,12);
	\draw[color ={black},opacity = 0.5] (-2,-2)--(-2,12);
	\draw[color ={gray},opacity = 0.3] (-2,0)--(14,0);
	\draw[color ={gray},opacity = 0.3] (-2,0.5)--(14,0.5);
	\draw[color ={gray},opacity = 0.3] (-2,1.5)--(14,1.5);
	\draw[color ={gray},opacity = 0.3] (-2,2.5)--(14,2.5);
	\draw[color ={gray},opacity = 0.3] (-2,3.5)--(14,3.5);
	\draw[color ={gray},opacity = 0.3] (-2,4.5)--(14,4.5);
	\draw[color ={gray},opacity = 0.3] (-2,5.5)--(14,5.5);
	\draw[color ={gray},opacity = 0.3] (-2,6.5)--(14,6.5);
	\draw[color ={gray},opacity = 0.3] (-2,7.5)--(14,7.5);
	\draw[color ={gray},opacity = 0.3] (-2,8.5)--(14,8.5);
	\draw[color ={gray},opacity = 0.3] (-2,9.5)--(14,9.5);
	\draw[color ={gray},opacity = 0.3] (-2,10.5)--(14,10.5);
	\draw[color ={gray},opacity = 0.3] (-2,11.5)--(14,11.5);
	\draw[color ={gray},opacity = 0.3] (-2,0)--(14,0);
	\draw[color ={gray},opacity = 0.3] (-2,-0.5)--(14,-0.5);
	\draw[color ={gray},opacity = 0.3] (-2,-1.5)--(14,-1.5);
	\draw[color ={gray},opacity = 0.3] (0,-2)--(0,12);
	\draw[color ={gray},opacity = 0.3] (0.1,-2)--(0.1,12);
	\draw[color ={gray},opacity = 0.3] (0.2,-2)--(0.2,12);
	\draw[color ={gray},opacity = 0.3] (0.30000000000000004,-2)--(0.30000000000000004,12);
	\draw[color ={gray},opacity = 0.3] (0.4,-2)--(0.4,12);
	\draw[color ={gray},opacity = 0.3] (0.5,-2)--(0.5,12);
	\draw[color ={gray},opacity = 0.3] (0.6,-2)--(0.6,12);
	\draw[color ={gray},opacity = 0.3] (0.7,-2)--(0.7,12);
	\draw[color ={gray},opacity = 0.3] (0.7999999999999999,-2)--(0.7999999999999999,12);
	\draw[color ={gray},opacity = 0.3] (0.8999999999999999,-2)--(0.8999999999999999,12);
	\draw[color ={gray},opacity = 0.3] (1.0999999999999999,-2)--(1.0999999999999999,12);
	\draw[color ={gray},opacity = 0.3] (1.2,-2)--(1.2,12);
	\draw[color ={gray},opacity = 0.3] (1.3,-2)--(1.3,12);
	\draw[color ={gray},opacity = 0.3] (1.4000000000000001,-2)--(1.4000000000000001,12);
	\draw[color ={gray},opacity = 0.3] (1.5000000000000002,-2)--(1.5000000000000002,12);
	\draw[color ={gray},opacity = 0.3] (1.6000000000000003,-2)--(1.6000000000000003,12);
	\draw[color ={gray},opacity = 0.3] (1.7000000000000004,-2)--(1.7000000000000004,12);
	\draw[color ={gray},opacity = 0.3] (1.8000000000000005,-2)--(1.8000000000000005,12);
	\draw[color ={gray},opacity = 0.3] (1.9000000000000006,-2)--(1.9000000000000006,12);
	\draw[color ={gray},opacity = 0.3] (2.1000000000000005,-2)--(2.1000000000000005,12);
	\draw[color ={gray},opacity = 0.3] (2.2000000000000006,-2)--(2.2000000000000006,12);
	\draw[color ={gray},opacity = 0.3] (2.3000000000000007,-2)--(2.3000000000000007,12);
	\draw[color ={gray},opacity = 0.3] (2.400000000000001,-2)--(2.400000000000001,12);
	\draw[color ={gray},opacity = 0.3] (2.500000000000001,-2)--(2.500000000000001,12);
	\draw[color ={gray},opacity = 0.3] (2.600000000000001,-2)--(2.600000000000001,12);
	\draw[color ={gray},opacity = 0.3] (2.700000000000001,-2)--(2.700000000000001,12);
	\draw[color ={gray},opacity = 0.3] (2.800000000000001,-2)--(2.800000000000001,12);
	\draw[color ={gray},opacity = 0.3] (2.9000000000000012,-2)--(2.9000000000000012,12);
	\draw[color ={gray},opacity = 0.3] (3.1000000000000014,-2)--(3.1000000000000014,12);
	\draw[color ={gray},opacity = 0.3] (3.2000000000000015,-2)--(3.2000000000000015,12);
	\draw[color ={gray},opacity = 0.3] (3.3000000000000016,-2)--(3.3000000000000016,12);
	\draw[color ={gray},opacity = 0.3] (3.4000000000000017,-2)--(3.4000000000000017,12);
	\draw[color ={gray},opacity = 0.3] (3.5000000000000018,-2)--(3.5000000000000018,12);
	\draw[color ={gray},opacity = 0.3] (3.600000000000002,-2)--(3.600000000000002,12);
	\draw[color ={gray},opacity = 0.3] (3.700000000000002,-2)--(3.700000000000002,12);
	\draw[color ={gray},opacity = 0.3] (3.800000000000002,-2)--(3.800000000000002,12);
	\draw[color ={gray},opacity = 0.3] (3.900000000000002,-2)--(3.900000000000002,12);
	\draw[color ={gray},opacity = 0.3] (4.100000000000001,-2)--(4.100000000000001,12);
	\draw[color ={gray},opacity = 0.3] (4.200000000000001,-2)--(4.200000000000001,12);
	\draw[color ={gray},opacity = 0.3] (4.300000000000001,-2)--(4.300000000000001,12);
	\draw[color ={gray},opacity = 0.3] (4.4,-2)--(4.4,12);
	\draw[color ={gray},opacity = 0.3] (4.5,-2)--(4.5,12);
	\draw[color ={gray},opacity = 0.3] (4.6,-2)--(4.6,12);
	\draw[color ={gray},opacity = 0.3] (4.699999999999999,-2)--(4.699999999999999,12);
	\draw[color ={gray},opacity = 0.3] (4.799999999999999,-2)--(4.799999999999999,12);
	\draw[color ={gray},opacity = 0.3] (4.899999999999999,-2)--(4.899999999999999,12);
	\draw[color ={gray},opacity = 0.3] (5.099999999999998,-2)--(5.099999999999998,12);
	\draw[color ={gray},opacity = 0.3] (5.1999999999999975,-2)--(5.1999999999999975,12);
	\draw[color ={gray},opacity = 0.3] (5.299999999999997,-2)--(5.299999999999997,12);
	\draw[color ={gray},opacity = 0.3] (5.399999999999997,-2)--(5.399999999999997,12);
	\draw[color ={gray},opacity = 0.3] (5.4999999999999964,-2)--(5.4999999999999964,12);
	\draw[color ={gray},opacity = 0.3] (5.599999999999996,-2)--(5.599999999999996,12);
	\draw[color ={gray},opacity = 0.3] (5.699999999999996,-2)--(5.699999999999996,12);
	\draw[color ={gray},opacity = 0.3] (5.799999999999995,-2)--(5.799999999999995,12);
	\draw[color ={gray},opacity = 0.3] (5.899999999999995,-2)--(5.899999999999995,12);
	\draw[color ={gray},opacity = 0.3] (6.099999999999994,-2)--(6.099999999999994,12);
	\draw[color ={gray},opacity = 0.3] (6.199999999999994,-2)--(6.199999999999994,12);
	\draw[color ={gray},opacity = 0.3] (6.299999999999994,-2)--(6.299999999999994,12);
	\draw[color ={gray},opacity = 0.3] (6.399999999999993,-2)--(6.399999999999993,12);
	\draw[color ={gray},opacity = 0.3] (6.499999999999993,-2)--(6.499999999999993,12);
	\draw[color ={gray},opacity = 0.3] (6.5999999999999925,-2)--(6.5999999999999925,12);
	\draw[color ={gray},opacity = 0.3] (6.699999999999992,-2)--(6.699999999999992,12);
	\draw[color ={gray},opacity = 0.3] (6.799999999999992,-2)--(6.799999999999992,12);
	\draw[color ={gray},opacity = 0.3] (6.8999999999999915,-2)--(6.8999999999999915,12);
	\draw[color ={gray},opacity = 0.3] (7.099999999999991,-2)--(7.099999999999991,12);
	\draw[color ={gray},opacity = 0.3] (7.19999999999999,-2)--(7.19999999999999,12);
	\draw[color ={gray},opacity = 0.3] (7.29999999999999,-2)--(7.29999999999999,12);
	\draw[color ={gray},opacity = 0.3] (7.39999999999999,-2)--(7.39999999999999,12);
	\draw[color ={gray},opacity = 0.3] (7.499999999999989,-2)--(7.499999999999989,12);
	\draw[color ={gray},opacity = 0.3] (7.599999999999989,-2)--(7.599999999999989,12);
	\draw[color ={gray},opacity = 0.3] (7.699999999999989,-2)--(7.699999999999989,12);
	\draw[color ={gray},opacity = 0.3] (7.799999999999988,-2)--(7.799999999999988,12);
	\draw[color ={gray},opacity = 0.3] (7.899999999999988,-2)--(7.899999999999988,12);
	\draw[color ={gray},opacity = 0.3] (8.099999999999987,-2)--(8.099999999999987,12);
	\draw[color ={gray},opacity = 0.3] (8.199999999999987,-2)--(8.199999999999987,12);
	\draw[color ={gray},opacity = 0.3] (8.299999999999986,-2)--(8.299999999999986,12);
	\draw[color ={gray},opacity = 0.3] (8.399999999999986,-2)--(8.399999999999986,12);
	\draw[color ={gray},opacity = 0.3] (8.499999999999986,-2)--(8.499999999999986,12);
	\draw[color ={gray},opacity = 0.3] (8.599999999999985,-2)--(8.599999999999985,12);
	\draw[color ={gray},opacity = 0.3] (8.699999999999985,-2)--(8.699999999999985,12);
	\draw[color ={gray},opacity = 0.3] (8.799999999999985,-2)--(8.799999999999985,12);
	\draw[color ={gray},opacity = 0.3] (8.899999999999984,-2)--(8.899999999999984,12);
	\draw[color ={gray},opacity = 0.3] (9.099999999999984,-2)--(9.099999999999984,12);
	\draw[color ={gray},opacity = 0.3] (9.199999999999983,-2)--(9.199999999999983,12);
	\draw[color ={gray},opacity = 0.3] (9.299999999999983,-2)--(9.299999999999983,12);
	\draw[color ={gray},opacity = 0.3] (9.399999999999983,-2)--(9.399999999999983,12);
	\draw[color ={gray},opacity = 0.3] (9.499999999999982,-2)--(9.499999999999982,12);
	\draw[color ={gray},opacity = 0.3] (9.599999999999982,-2)--(9.599999999999982,12);
	\draw[color ={gray},opacity = 0.3] (9.699999999999982,-2)--(9.699999999999982,12);
	\draw[color ={gray},opacity = 0.3] (9.799999999999981,-2)--(9.799999999999981,12);
	\draw[color ={gray},opacity = 0.3] (9.89999999999998,-2)--(9.89999999999998,12);
	\draw[color ={gray},opacity = 0.3] (10.09999999999998,-2)--(10.09999999999998,12);
	\draw[color ={gray},opacity = 0.3] (10.19999999999998,-2)--(10.19999999999998,12);
	\draw[color ={gray},opacity = 0.3] (10.29999999999998,-2)--(10.29999999999998,12);
	\draw[color ={gray},opacity = 0.3] (10.399999999999979,-2)--(10.399999999999979,12);
	\draw[color ={gray},opacity = 0.3] (10.499999999999979,-2)--(10.499999999999979,12);
	\draw[color ={gray},opacity = 0.3] (10.599999999999978,-2)--(10.599999999999978,12);
	\draw[color ={gray},opacity = 0.3] (10.699999999999978,-2)--(10.699999999999978,12);
	\draw[color ={gray},opacity = 0.3] (10.799999999999978,-2)--(10.799999999999978,12);
	\draw[color ={gray},opacity = 0.3] (10.899999999999977,-2)--(10.899999999999977,12);
	\draw[color ={gray},opacity = 0.3] (11.099999999999977,-2)--(11.099999999999977,12);
	\draw[color ={gray},opacity = 0.3] (11.199999999999976,-2)--(11.199999999999976,12);
	\draw[color ={gray},opacity = 0.3] (11.299999999999976,-2)--(11.299999999999976,12);
	\draw[color ={gray},opacity = 0.3] (11.399999999999975,-2)--(11.399999999999975,12);
	\draw[color ={gray},opacity = 0.3] (11.499999999999975,-2)--(11.499999999999975,12);
	\draw[color ={gray},opacity = 0.3] (11.599999999999975,-2)--(11.599999999999975,12);
	\draw[color ={gray},opacity = 0.3] (11.699999999999974,-2)--(11.699999999999974,12);
	\draw[color ={gray},opacity = 0.3] (11.799999999999974,-2)--(11.799999999999974,12);
	\draw[color ={gray},opacity = 0.3] (11.899999999999974,-2)--(11.899999999999974,12);
	\draw[color ={gray},opacity = 0.3] (12.099999999999973,-2)--(12.099999999999973,12);
	\draw[color ={gray},opacity = 0.3] (12.199999999999973,-2)--(12.199999999999973,12);
	\draw[color ={gray},opacity = 0.3] (12.299999999999972,-2)--(12.299999999999972,12);
	\draw[color ={gray},opacity = 0.3] (12.399999999999972,-2)--(12.399999999999972,12);
	\draw[color ={gray},opacity = 0.3] (12.499999999999972,-2)--(12.499999999999972,12);
	\draw[color ={gray},opacity = 0.3] (12.599999999999971,-2)--(12.599999999999971,12);
	\draw[color ={gray},opacity = 0.3] (12.69999999999997,-2)--(12.69999999999997,12);
	\draw[color ={gray},opacity = 0.3] (12.79999999999997,-2)--(12.79999999999997,12);
	\draw[color ={gray},opacity = 0.3] (12.89999999999997,-2)--(12.89999999999997,12);
	\draw[color ={gray},opacity = 0.3] (13.09999999999997,-2)--(13.09999999999997,12);
	\draw[color ={gray},opacity = 0.3] (13.199999999999969,-2)--(13.199999999999969,12);
	\draw[color ={gray},opacity = 0.3] (13.299999999999969,-2)--(13.299999999999969,12);
	\draw[color ={gray},opacity = 0.3] (13.399999999999968,-2)--(13.399999999999968,12);
	\draw[color ={gray},opacity = 0.3] (13.499999999999968,-2)--(13.499999999999968,12);
	\draw[color ={gray},opacity = 0.3] (13.599999999999968,-2)--(13.599999999999968,12);
	\draw[color ={gray},opacity = 0.3] (13.699999999999967,-2)--(13.699999999999967,12);
	\draw[color ={gray},opacity = 0.3] (13.799999999999967,-2)--(13.799999999999967,12);
	\draw[color ={gray},opacity = 0.3] (13.899999999999967,-2)--(13.899999999999967,12);
	\draw[color ={gray},opacity = 0.3] (0,-2)--(0,12);
	\draw[color ={gray},opacity = 0.3] (-0.1,-2)--(-0.1,12);
	\draw[color ={gray},opacity = 0.3] (-0.2,-2)--(-0.2,12);
	\draw[color ={gray},opacity = 0.3] (-0.30000000000000004,-2)--(-0.30000000000000004,12);
	\draw[color ={gray},opacity = 0.3] (-0.4,-2)--(-0.4,12);
	\draw[color ={gray},opacity = 0.3] (-0.5,-2)--(-0.5,12);
	\draw[color ={gray},opacity = 0.3] (-0.6,-2)--(-0.6,12);
	\draw[color ={gray},opacity = 0.3] (-0.7,-2)--(-0.7,12);
	\draw[color ={gray},opacity = 0.3] (-0.7999999999999999,-2)--(-0.7999999999999999,12);
	\draw[color ={gray},opacity = 0.3] (-0.8999999999999999,-2)--(-0.8999999999999999,12);
	\draw[color ={gray},opacity = 0.3] (-1.0999999999999999,-2)--(-1.0999999999999999,12);
	\draw[color ={gray},opacity = 0.3] (-1.2,-2)--(-1.2,12);
	\draw[color ={gray},opacity = 0.3] (-1.3,-2)--(-1.3,12);
	\draw[color ={gray},opacity = 0.3] (-1.4000000000000001,-2)--(-1.4000000000000001,12);
	\draw[color ={gray},opacity = 0.3] (-1.5000000000000002,-2)--(-1.5000000000000002,12);
	\draw[color ={gray},opacity = 0.3] (-1.6000000000000003,-2)--(-1.6000000000000003,12);
	\draw[color ={gray},opacity = 0.3] (-1.7000000000000004,-2)--(-1.7000000000000004,12);
	\draw[color ={gray},opacity = 0.3] (-1.8000000000000005,-2)--(-1.8000000000000005,12);
	\draw[color ={gray},opacity = 0.3] (-1.9000000000000006,-2)--(-1.9000000000000006,12);
	\draw[color ={black},line width = 2] (2,-0.2)--(2,0.2);
	\draw[color ={black},line width = 2] (4,-0.2)--(4,0.2);
	\draw[color ={black},line width = 2] (6,-0.2)--(6,0.2);
	\draw[color ={black},line width = 2] (8,-0.2)--(8,0.2);
	\draw[color ={black},line width = 2] (10,-0.2)--(10,0.2);
	\draw[color ={black},line width = 2] (12,-0.2)--(12,0.2);
	\draw[color ={black},line width = 2] (-0.2,2)--(0.2,2);
	\draw[color ={black},line width = 2] (-0.2,4)--(0.2,4);
	\draw[color ={black},line width = 2] (-0.2,6)--(0.2,6);
	\draw[color ={black},line width = 2] (-0.2,8)--(0.2,8);
	\draw[color ={black},line width = 2] (-0.2,10)--(0.2,10);
	\draw (2,-0.4) node[anchor = center] {\scriptsize \color{black}{$1$}};
	\draw (4,-0.4) node[anchor = center] {\scriptsize \color{black}{$2$}};
	\draw (6,-0.4) node[anchor = center] {\scriptsize \color{black}{$3$}};
	\draw (8,-0.4) node[anchor = center] {\scriptsize \color{black}{$4$}};
	\draw (10,-0.4) node[anchor = center] {\scriptsize \color{black}{$5$}};
	\draw (-0.5,2.1) node[anchor = center] {\footnotesize \color{black}{$1$}};
	\draw (-0.5,4.1) node[anchor = center] {\footnotesize \color{black}{$2$}};
	\draw (-0.5,6.1) node[anchor = center] {\footnotesize \color{black}{$3$}};
	\draw (-0.5,8.1) node[anchor = center] {\footnotesize \color{black}{$4$}};
	\draw (-0.5,10.1) node[anchor = center] {\footnotesize \color{black}{$5$}};
	\draw [color={black}] (-0.3,-0.3) node[anchor = center,scale=1, rotate = 0] {O};
	
	\draw[color={blue},line width = 2] (0.6,10.67)--(1,8)--(1.4,6.86)--(1.8,6.22)--(2.2,5.82)--(2.6,5.54)--(3,5.33)--(3.4,5.18)--(3.8,5.05)--(4.2,4.95)--(4.6,4.87)--(5,4.8)--(5.4,4.74)--(5.8,4.69)--(6.2,4.65)--(6.6,4.61)--(7,4.57)--(7.4,4.54)--(7.8,4.51)--(8.2,4.49)--(8.6,4.47)--(9,4.44)--(9.4,4.43)--(9.8,4.41)--(10.2,4.39)--(10.6,4.38)--(11,4.36)--(11.4,4.35)--(11.8,4.34)--(12.2,4.33)--(12.6,4.32)--(13,4.31)--(13.4,4.3)--(13.8,4.29);
	
	\draw[color={red},line width = 2] (-2,10)--(-1.6,9.6)--(-1.2,9.2)--(-0.8,8.8)--(-0.4,8.4)--(0,8)--(0.4,7.6)--(0.8,7.2)--(1.2,6.8)--(1.6,6.4)--(2,6)--(2.4,5.6)--(2.8,5.2)--(3.2,4.8)--(3.6,4.4)--(4,4)--(4.4,3.6)--(4.8,3.2)--(5.2,2.8)--(5.6,2.4)--(6,2)--(6.4,1.6)--(6.8,1.2)--(7.2,0.8)--(7.6,0.4)--(8,0)--(8.4,-0.4)--(8.8,-0.8)--(9.2,-1.2)--(9.6,-1.6)--(10,-2)--(10.4,-2.4)--(10.8,-2.8)--(11.2,-3.2)--(11.6,-3.6);

\end{tikzpicture}\end{center}
\end{exercice}

\begin{Solution}
 $f'(1)$ est donné par le coefficient directeur de la tangente à la courbe au point d'abscisse $1$, soit $-1=-1$.
\end{Solution}

% @see : https://coopmaths.fr/alea?uuid=6f32d&id=can1F16&n=1&d=10&alea=dqSw22323232323423456789101112131415161718192021222324252627282930313233343536373839404142434445464748234567891011121314151617181920212223242526272829303132&cd=1&cols=1
\needspace{10\baselineskip}
\newpage
\begin{exercice}[BaremeTotal=false]


 On donne les représentations graphiques d'une fonction et de sa dérivée.\\
      Donner l'équation réduite de la tangente à la courbe de $f$ en $x=-2$. \\ \begin{minipage}{0.5\linewidth}
    \begin{tikzpicture}[baseline, scale=0.8]

    \tikzset{
      point/.style={
        thick,
        draw,
        cross out,
        inner sep=0pt,
        minimum width=5pt,
        minimum height=5pt,
      },
    }
    \clip (-5.5,-9.5) rectangle (5.75,7.5);
    	
	\draw[color ={black},>=latex,->] (-4,0)--(4,0);
	\draw[color ={black},>=latex,->] (0,-8)--(0,6);
	\draw[color ={black},opacity = 0.4] (-4,1)--(4,1);
	\draw[color ={black},opacity = 0.4] (-4,2)--(4,2);
	\draw[color ={black},opacity = 0.4] (-4,3)--(4,3);
	\draw[color ={black},opacity = 0.4] (-4,4)--(4,4);
	\draw[color ={black},opacity = 0.4] (-4,5)--(4,5);
	\draw[color ={black},opacity = 0.4] (-4,6)--(4,6);
	\draw[color ={black},opacity = 0.4] (-4,-1)--(4,-1);
	\draw[color ={black},opacity = 0.4] (-4,-2)--(4,-2);
	\draw[color ={black},opacity = 0.4] (-4,-3)--(4,-3);
	\draw[color ={black},opacity = 0.4] (-4,-4)--(4,-4);
	\draw[color ={black},opacity = 0.4] (-4,-5)--(4,-5);
	\draw[color ={black},opacity = 0.4] (-4,-6)--(4,-6);
	\draw[color ={black},opacity = 0.4] (-4,-7)--(4,-7);
	\draw[color ={black},opacity = 0.4] (-4,-8)--(4,-8);
	\draw[color ={black},opacity = 0.4] (1,-8)--(1,6);
	\draw[color ={black},opacity = 0.4] (2,-8)--(2,6);
	\draw[color ={black},opacity = 0.4] (3,-8)--(3,6);
	\draw[color ={black},opacity = 0.4] (4,-8)--(4,6);
	\draw[color ={black},opacity = 0.4] (-1,-8)--(-1,6);
	\draw[color ={black},opacity = 0.4] (-2,-8)--(-2,6);
	\draw[color ={black},opacity = 0.4] (-3,-8)--(-3,6);
	\draw[color ={black},opacity = 0.4] (-4,-8)--(-4,6);
	\draw[color ={gray},opacity = 0.3] (-4,0)--(4,0);
	\draw[color ={gray},opacity = 0.3] (-4,0.5)--(4,0.5);
	\draw[color ={gray},opacity = 0.3] (-4,1.5)--(4,1.5);
	\draw[color ={gray},opacity = 0.3] (-4,2.5)--(4,2.5);
	\draw[color ={gray},opacity = 0.3] (-4,3.5)--(4,3.5);
	\draw[color ={gray},opacity = 0.3] (-4,4.5)--(4,4.5);
	\draw[color ={gray},opacity = 0.3] (-4,5.5)--(4,5.5);
	\draw[color ={gray},opacity = 0.3] (-4,0)--(4,0);
	\draw[color ={gray},opacity = 0.3] (-4,-0.5)--(4,-0.5);
	\draw[color ={gray},opacity = 0.3] (-4,-1.5)--(4,-1.5);
	\draw[color ={gray},opacity = 0.3] (-4,-2.5)--(4,-2.5);
	\draw[color ={gray},opacity = 0.3] (-4,-3.5)--(4,-3.5);
	\draw[color ={gray},opacity = 0.3] (-4,-4.5)--(4,-4.5);
	\draw[color ={gray},opacity = 0.3] (-4,-5.5)--(4,-5.5);
	\draw[color ={gray},opacity = 0.3] (-4,-6.5)--(4,-6.5);
	\draw[color ={gray},opacity = 0.3] (-4,-7)--(4,-7);
	\draw[color ={gray},opacity = 0.3] (-4,-7.5)--(4,-7.5);
	\draw[color ={gray},opacity = 0.3] (-4,-8)--(4,-8);
	\draw[color ={gray},opacity = 0.3] (0,-8)--(0,6);
	\draw[color ={gray},opacity = 0.3] (0,-8)--(0,6);
	\draw[color ={black},line width = 1.2] (1,-0.13)--(1,0.13);
	\draw[color ={black},line width = 1.2] (2,-0.13)--(2,0.13);
	\draw[color ={black},line width = 1.2] (3,-0.13)--(3,0.13);
	\draw[color ={black},line width = 1.2] (4,-0.13)--(4,0.13);
	\draw[color ={black},line width = 1.2] (-1,-0.13)--(-1,0.13);
	\draw[color ={black},line width = 1.2] (-2,-0.13)--(-2,0.13);
	\draw[color ={black},line width = 1.2] (-3,-0.13)--(-3,0.13);
	\draw[color ={black},line width = 1.2] (-4,-0.13)--(-4,0.13);
	\draw[color ={black},line width = 1.2] (-0.13,1)--(0.13,1);
	\draw[color ={black},line width = 1.2] (-0.13,2)--(0.13,2);
	\draw[color ={black},line width = 1.2] (-0.13,3)--(0.13,3);
	\draw[color ={black},line width = 1.2] (-0.13,4)--(0.13,4);
	\draw[color ={black},line width = 1.2] (-0.13,5)--(0.13,5);
	\draw[color ={black},line width = 1.2] (-0.13,6)--(0.13,6);
	\draw[color ={black},line width = 1.2] (-0.13,-1)--(0.13,-1);
	\draw[color ={black},line width = 1.2] (-0.13,-2)--(0.13,-2);
	\draw[color ={black},line width = 1.2] (-0.13,-3)--(0.13,-3);
	\draw[color ={black},line width = 1.2] (-0.13,-4)--(0.13,-4);
	\draw[color ={black},line width = 1.2] (-0.13,-5)--(0.13,-5);
	\draw[color ={black},line width = 1.2] (-0.13,-6)--(0.13,-6);
	\draw[color ={black},line width = 1.2] (-0.13,-7)--(0.13,-7);
	\draw[color ={black},line width = 1.2] (-0.13,-8)--(0.13,-8);
	\draw (2,-0.4) node[anchor = center] {\scriptsize \color{black}{$2$}};
	\draw (4,-0.4) node[anchor = center] {\scriptsize \color{black}{$4$}};
	\draw (-2,-0.4) node[anchor = center] {\scriptsize \color{black}{$-2$}};
	\draw (-4,-0.4) node[anchor = center] {\scriptsize \color{black}{$-4$}};
	\draw (-0.5,2.1) node[anchor = center] {\footnotesize \color{black}{$2$}};
	\draw (-0.5,4.1) node[anchor = center] {\footnotesize \color{black}{$4$}};
	\draw (-0.5,6.1) node[anchor = center] {\footnotesize \color{black}{$6$}};
	\draw (-0.5,-1.9) node[anchor = center] {\footnotesize \color{black}{$-2$}};
	\draw (-0.5,-3.9) node[anchor = center] {\footnotesize \color{black}{$-4$}};
	\draw (-0.5,-5.9) node[anchor = center] {\footnotesize \color{black}{$-6$}};
	\draw (-0.5,-7.9) node[anchor = center] {\footnotesize \color{black}{$-8$}};
	\draw [color={black}] (-0.3,-0.3) node[anchor = center,scale=1, rotate = 0] {O};
	\draw (3,4) node[anchor = center] {\footnotesize \color{blue}{$\Large \mathcal{C}_f$}};
	
	\draw[color={blue},line width = 2] (-4,-7)--(-3.8,-5.84)--(-3.6,-4.76)--(-3.4,-3.76)--(-3.2,-2.84)--(-3,-2)--(-2.8,-1.24)--(-2.6,-0.56)--(-2.4,0.04)--(-2.2,0.56)--(-2,1)--(-1.8,1.36)--(-1.6,1.64)--(-1.4,1.84)--(-1.2,1.96)--(-1,2)--(-0.8,1.96)--(-0.6,1.84)--(-0.4,1.64)--(-0.2,1.36)--(0,1)--(0.2,0.56)--(0.4,0.04)--(0.6,-0.56)--(0.8,-1.24)--(1,-2)--(1.2,-2.84)--(1.4,-3.76)--(1.6,-4.76)--(1.8,-5.84)--(2,-7)--(2.2,-8.24);

\end{tikzpicture}\\
    \end{minipage}
    \begin{minipage}{0.5\linewidth}
    \begin{tikzpicture}[baseline, scale=0.8]

    \tikzset{
      point/.style={
        thick,
        draw,
        cross out,
        inner sep=0pt,
        minimum width=5pt,
        minimum height=5pt,
      },
    }
    \clip (-5.5,-9.5) rectangle (6.5,7.5);
    	
	\draw[color ={black},>=latex,->] (-4,0)--(4,0);
	\draw[color ={black},>=latex,->] (0,-8)--(0,6);
	\draw[color ={black},opacity = 0.4] (-4,1)--(4,1);
	\draw[color ={black},opacity = 0.4] (-4,2)--(4,2);
	\draw[color ={black},opacity = 0.4] (-4,3)--(4,3);
	\draw[color ={black},opacity = 0.4] (-4,4)--(4,4);
	\draw[color ={black},opacity = 0.4] (-4,5)--(4,5);
	\draw[color ={black},opacity = 0.4] (-4,6)--(4,6);
	\draw[color ={black},opacity = 0.4] (-4,-1)--(4,-1);
	\draw[color ={black},opacity = 0.4] (-4,-2)--(4,-2);
	\draw[color ={black},opacity = 0.4] (-4,-3)--(4,-3);
	\draw[color ={black},opacity = 0.4] (-4,-4)--(4,-4);
	\draw[color ={black},opacity = 0.4] (-4,-5)--(4,-5);
	\draw[color ={black},opacity = 0.4] (-4,-6)--(4,-6);
	\draw[color ={black},opacity = 0.4] (-4,-7)--(4,-7);
	\draw[color ={black},opacity = 0.4] (-4,-8)--(4,-8);
	\draw[color ={black},opacity = 0.4] (1,-8)--(1,6);
	\draw[color ={black},opacity = 0.4] (2,-8)--(2,6);
	\draw[color ={black},opacity = 0.4] (3,-8)--(3,6);
	\draw[color ={black},opacity = 0.4] (4,-8)--(4,6);
	\draw[color ={black},opacity = 0.4] (-1,-8)--(-1,6);
	\draw[color ={black},opacity = 0.4] (-2,-8)--(-2,6);
	\draw[color ={black},opacity = 0.4] (-3,-8)--(-3,6);
	\draw[color ={black},opacity = 0.4] (-4,-8)--(-4,6);
	\draw[color ={gray},opacity = 0.3] (-4,0)--(4,0);
	\draw[color ={gray},opacity = 0.3] (-4,0.5)--(4,0.5);
	\draw[color ={gray},opacity = 0.3] (-4,1.5)--(4,1.5);
	\draw[color ={gray},opacity = 0.3] (-4,2.5)--(4,2.5);
	\draw[color ={gray},opacity = 0.3] (-4,3.5)--(4,3.5);
	\draw[color ={gray},opacity = 0.3] (-4,4.5)--(4,4.5);
	\draw[color ={gray},opacity = 0.3] (-4,5.5)--(4,5.5);
	\draw[color ={gray},opacity = 0.3] (-4,0)--(4,0);
	\draw[color ={gray},opacity = 0.3] (-4,-0.5)--(4,-0.5);
	\draw[color ={gray},opacity = 0.3] (-4,-1.5)--(4,-1.5);
	\draw[color ={gray},opacity = 0.3] (-4,-2.5)--(4,-2.5);
	\draw[color ={gray},opacity = 0.3] (-4,-3.5)--(4,-3.5);
	\draw[color ={gray},opacity = 0.3] (-4,-4.5)--(4,-4.5);
	\draw[color ={gray},opacity = 0.3] (-4,-5.5)--(4,-5.5);
	\draw[color ={gray},opacity = 0.3] (-4,-6.5)--(4,-6.5);
	\draw[color ={gray},opacity = 0.3] (-4,-7)--(4,-7);
	\draw[color ={gray},opacity = 0.3] (-4,-7.5)--(4,-7.5);
	\draw[color ={gray},opacity = 0.3] (-4,-8)--(4,-8);
	\draw[color ={gray},opacity = 0.3] (0,-8)--(0,6);
	\draw[color ={gray},opacity = 0.3] (0,-8)--(0,6);
	\draw[color ={black},line width = 1.2] (1,-0.13)--(1,0.13);
	\draw[color ={black},line width = 1.2] (2,-0.13)--(2,0.13);
	\draw[color ={black},line width = 1.2] (3,-0.13)--(3,0.13);
	\draw[color ={black},line width = 1.2] (4,-0.13)--(4,0.13);
	\draw[color ={black},line width = 1.2] (-1,-0.13)--(-1,0.13);
	\draw[color ={black},line width = 1.2] (-2,-0.13)--(-2,0.13);
	\draw[color ={black},line width = 1.2] (-3,-0.13)--(-3,0.13);
	\draw[color ={black},line width = 1.2] (-4,-0.13)--(-4,0.13);
	\draw[color ={black},line width = 1.2] (-0.13,1)--(0.13,1);
	\draw[color ={black},line width = 1.2] (-0.13,2)--(0.13,2);
	\draw[color ={black},line width = 1.2] (-0.13,3)--(0.13,3);
	\draw[color ={black},line width = 1.2] (-0.13,4)--(0.13,4);
	\draw[color ={black},line width = 1.2] (-0.13,5)--(0.13,5);
	\draw[color ={black},line width = 1.2] (-0.13,6)--(0.13,6);
	\draw[color ={black},line width = 1.2] (-0.13,-1)--(0.13,-1);
	\draw[color ={black},line width = 1.2] (-0.13,-2)--(0.13,-2);
	\draw[color ={black},line width = 1.2] (-0.13,-3)--(0.13,-3);
	\draw[color ={black},line width = 1.2] (-0.13,-4)--(0.13,-4);
	\draw[color ={black},line width = 1.2] (-0.13,-5)--(0.13,-5);
	\draw[color ={black},line width = 1.2] (-0.13,-6)--(0.13,-6);
	\draw[color ={black},line width = 1.2] (-0.13,-7)--(0.13,-7);
	\draw[color ={black},line width = 1.2] (-0.13,-8)--(0.13,-8);
	\draw (2,-0.4) node[anchor = center] {\scriptsize \color{black}{$2$}};
	\draw (4,-0.4) node[anchor = center] {\scriptsize \color{black}{$4$}};
	\draw (-2,-0.4) node[anchor = center] {\scriptsize \color{black}{$-2$}};
	\draw (-4,-0.4) node[anchor = center] {\scriptsize \color{black}{$-4$}};
	\draw (-0.5,2.1) node[anchor = center] {\footnotesize \color{black}{$2$}};
	\draw (-0.5,4.1) node[anchor = center] {\footnotesize \color{black}{$4$}};
	\draw (-0.5,-1.9) node[anchor = center] {\footnotesize \color{black}{$-2$}};
	\draw (-0.5,-3.9) node[anchor = center] {\footnotesize \color{black}{$-4$}};
	\draw (-0.5,-5.9) node[anchor = center] {\footnotesize \color{black}{$-6$}};
	\draw (-0.5,-7.9) node[anchor = center] {\footnotesize \color{black}{$-8$}};
	\draw [color={black}] (-0.3,-0.3) node[anchor = center,scale=1, rotate = 0] {O};
	\draw (3,4) node[anchor = center] {\footnotesize \color{red}{$\Large\mathcal{C}_f\prime$}};
	
	\draw[color={red},line width = 2] (-4,6)--(-3.8,5.6)--(-3.6,5.2)--(-3.4,4.8)--(-3.2,4.4)--(-3,4)--(-2.8,3.6)--(-2.6,3.2)--(-2.4,2.8)--(-2.2,2.4)--(-2,2)--(-1.8,1.6)--(-1.6,1.2)--(-1.4,0.8)--(-1.2,0.4)--(-1,0)--(-0.8,-0.4)--(-0.6,-0.8)--(-0.4,-1.2)--(-0.2,-1.6)--(0,-2)--(0.2,-2.4)--(0.4,-2.8)--(0.6,-3.2)--(0.8,-3.6)--(1,-4)--(1.2,-4.4)--(1.4,-4.8)--(1.6,-5.2)--(1.8,-5.6)--(2,-6)--(2.2,-6.4)--(2.4,-6.8)--(2.6,-7.2)--(2.8,-7.6)--(3,-8)--(3.2,-8.4)--(3.4,-8.8);

\end{tikzpicture}\\
    \end{minipage}
    
\end{exercice}

\begin{Solution}
 L'équation réduite de la tangente au point d'abscisse $-2$ est  : $y=f'(-2)(x-(-2))+f(-2)$.\\
      On lit graphiquement $f(-2)=1$ et $f'(-2)=2$.\\
      L'équation réduite de la tangente est donc donnée par :
      $y=2(x+2)+1$, soit $y=2x+5$.
\end{Solution}

% @see : https://coopmaths.fr/alea?uuid=a1ba2&id=can1F14&n=1&d=10&alea=rkIp22323232323423456789101112131415161718192021222324252627282930313233343536373839404142434445464748234567891011121314151617181920212223242526272829303132&cd=1&cols=1
\needspace{10\baselineskip}
\begin{exercice}[BaremeTotal=false]


Soit $f$ la fonction définie sur $\mathbb{R}$ par : 
$f(x)= 3+5x^2-2x$.\\
En calculant la limite du taux d'accroissement, déterminer $f'(1)$.
\end{exercice}

\begin{Solution}
 $f$ est une fonction polynôme du second degré de la forme $f(x)=ax^2+bx+c$.\\
    La fonction dérivée est donnée par la somme des dérivées des fonctions $u$ et $v$ définies par $u(x)=5x^2$ et $v(x)=-2x+3$.\\
     Comme $u'(x)=10x$ et $v'(x)=-2$, on obtient  $f'(x)=10x-2$. \\
     Ainsi, $f'(1)=10\times 1-2=8$.
\end{Solution}

% @see : https://coopmaths.fr/alea?uuid=3c690&id=can1F13&n=1&d=10&alea=EUDu22323232323423456789101112131415161718192021222324252627282930313233343536373839404142434445464748234567891011121314151617181920212223242526272829303132&cd=1&cols=1
\needspace{10\baselineskip}
\begin{exercice}[BaremeTotal=false]


 Déterminer le coefficient directeur de la tangente à la courbe représentative de la fonction racine carrée au point d'abscisse $14$.
         
\end{exercice}

\begin{Solution}
 Le coefficient directeur de la tangente au point d'abscisse $a$ est donné par le nombre dérivé $f'(a)$.\\
        La fonction $f$ définie par $f(x)=\sqrt{x}$ a pour fonction dérivée la fonction $f'$ définie par $f'(x)=\dfrac{1}{2\sqrt{x}}$.\\Comme $f'(14)=\dfrac{1}{2\sqrt{14}}$, le coefficient directeur de la tangente au point d'abscisse $14$ est : $\dfrac{1}{2\sqrt{14}}$.
\end{Solution}

% @see : https://coopmaths.fr/alea?uuid=ebd89&id=1AN14-1&n=1&d=10&s=3&alea=ocYr22323232323423456789101112131415161718192021222324252627282930313233343536373839404142434445464748234567891011121314151617181920212223242526272829303132&cd=0&cols=1
\needspace{10\baselineskip}
\begin{exercice}[BaremeTotal=false]


 Donner la dérivée de la fonction $f$, dérivable pour tout $x\in\R$, définie par  $f(x)=-\dfrac{9}{10}x+8$.
\end{exercice}

\begin{Solution}
 L'expression de la dérivée de la fonction $f$ définie par $f(x)=-\dfrac{9}{10}x+8$ est : ${\color[HTML]{f15929}\boldsymbol{f'(x)=-\dfrac{9}{10}}}$.
\end{Solution}

% @see : https://coopmaths.fr/alea?uuid=ebd89&id=1AN14-1&n=1&d=10&s=4&alea=uAlP22323232323423456789101112131415161718192021222324252627282930313233343536373839404142434445464748234567891011121314151617181920212223242526272829303132&cd=0&cols=1
\needspace{10\baselineskip}
\begin{exercice}[BaremeTotal=false]


 Donner la dérivée de la fonction $f$, dérivable pour tout $x\in\R$, définie par  $f(x)=-6x^{8}$.
\end{exercice}

\begin{Solution}
 L'expression de la dérivée de la fonction $f$ définie par $f(x)=-6x^{8}$ est : ${\color[HTML]{f15929}\boldsymbol{f'(x)=-48x^{7}}}$.
\end{Solution}

% @see : https://coopmaths.fr/alea?uuid=ebd89&id=1AN14-1&n=1&d=10&s=7&alea=vjN722323232323423456789101112131415161718192021222324252627282930313233343536373839404142434445464748234567891011121314151617181920212223242526272829303132&cd=0&cols=1
\needspace{10\baselineskip}
\begin{exercice}[BaremeTotal=false]


 Donner la dérivée de la fonction $f$, dérivable pour tout $x\in\R^*$, définie par  $f(x)=\dfrac{1}{10x}$.
\end{exercice}

\begin{Solution}
 L'expression de la dérivée de la fonction $f$ définie par $f(x)=\dfrac{1}{10x}$ est : ${\color[HTML]{f15929}\boldsymbol{f'(x)=-\dfrac{1}{10x^2}}}$.
\end{Solution}

% @see : https://coopmaths.fr/alea?uuid=a83c0&id=1AN14-4&n=1&d=10&s=1&alea=Njr322323232323423456789101112131415161718192021222324252627282930313233343536373839404142434445464748234567891011121314151617181920212223242526272829303132&cd=0&cols=1
\needspace{10\baselineskip}
\begin{exercice}[BaremeTotal=false]


 Donner l'expression de la dérivée de la fonction $f$ définie sur $\R^{*}$ par $f(x)=\dfrac{8}{x}+5x^{2}+2x-3$.\\
\end{exercice}

\begin{Solution}
 L'expression de la dérivée de la fonction $f$ définie par $f(x)=\dfrac{8}{x}+5x^{2}+2x-3$ est :\\$f'(x)={\color[HTML]{f15929}\boldsymbol{-\dfrac{8}{x^2}+10x+2}}$.\\
\end{Solution}

% @see : https://coopmaths.fr/alea?uuid=a83c0&id=1AN14-4&n=1&d=10&s=4&alea=4Vpz22323232323423456789101112131415161718192021222324252627282930313233343536373839404142434445464748234567891011121314151617181920212223242526272829303132&cd=1&cols=1
\needspace{10\baselineskip}
\begin{exercice}[BaremeTotal=false]


 Donner l'expression de la dérivée de la fonction $f$ définie sur $\R_+^{*}$ par $f(x)=-5\sqrt{x}-\dfrac{3}{x}-2x^{2}+3x$.\\
\end{exercice}

\begin{Solution}
 $(-5\sqrt{x})^\prime=-\dfrac{5}{2\sqrt{x}}$\\$(-\dfrac{3}{x})^\prime=\dfrac{3}{x^2}$\\$(-2x^{2}+3x)^\prime=-4x+3$\\L'expression de la dérivée de la fonction $f$ définie par $f(x)=-5\sqrt{x}-\dfrac{3}{x}-2x^{2}+3x$ est :\\$f'(x)={\color[HTML]{f15929}\boldsymbol{-\dfrac{5}{2\sqrt{x}}+\dfrac{3}{x^2}-4x+3}}$.\\
\end{Solution}

\end{Maquette}
\clearpage
\begin{Maquette}[DM]{Niveau= ,Classe=\getdata[33]\mydata\ (v33), Date = 06/01/25  ,Theme=Exercices}


% @see : https://coopmaths.fr/alea?uuid=4c8c7&id=1AN11-3&n=1&d=10&s=2&alea=nawO2232323232342345678910111213141516171819202122232425262728293031323334353637383940414243444546474823456789101112131415161718192021222324252627282930313233&cd=1&cols=2
\needspace{10\baselineskip}
\begin{exercice}[BaremeTotal=false]
\label{eleve:33}


 \begin{minipage}[t]{\linewidth} Soit $f$ une fonction dérivable sur $[-5;5]$ et $\mathcal{C}_f$ sa courbe représentative.\\On sait que  $f(-5)=2~~$ et que $~~f'(-5)=5$.\\Déterminer une équation de la tangente $(T)$ à la courbe $\mathcal{C}_f$ au point d'abscisse $-5$,\\sans utiliser la formule de cours de l'équation de tangente. \end{minipage}

\end{exercice}

\begin{Solution}
 \begin{minipage}[t]{\linewidth} On sait que la tangente n'est pas une droite verticale, puisque la fonction est dérivable sur l'intervalle.\\On en déduit que la tangente $(T)$ au point d'abscisse $-5$, admet une équation réduite de la forme :  $(T) : y=m x + p$. 

\medskip
$\bullet$ Détermination de $m$ :\\On sait que le nombre dérivé en $-5$ est par définition, le coefficient directeur de la tangente au point d'abscisse $-5$.\\Par conséquent, on a déjà : $m=f'(-5)=5$.\\ On en déduit que  $(T) : y= 5 x + p$\\$\bullet$ Détermination de $p$ :\\ Pour cela, on utilise que si $f(-5)=2~~$, alors le point $A$ de coordonnées $(-5;2)$ appartient à $\mathcal{C}_f$ mais aussi à $(T)$.\\ On peut écrire $A(-5;2) = \mathcal{C}_f \cap (T)$.\\ On remplace alors les coordonnées de $A(-5;2)$ dans l'équation  $(T) : y= 5 x + p$.\\ $\begin{aligned} \phantom{\iff}&A(-5;2)\in (T)\\ \iff& 2= 5 \times (-5)  + p\\ \iff& p=2 -5 \times (-5)  \\ \iff& p=2 +25   \\ \iff& p=27\\\end{aligned}$\\On peut conclure que : $(T) : {\color[HTML]{f15929}\boldsymbol{y=5x+27}}$. \end{minipage}

\end{Solution}

% @see : https://coopmaths.fr/alea?uuid=0e984&id=can1F15&n=1&d=10&alea=1Alo2232323232342345678910111213141516171819202122232425262728293031323334353637383940414243444546474823456789101112131415161718192021222324252627282930313233&cd=1&cols=1
\needspace{10\baselineskip}
\begin{exercice}[BaremeTotal=false]


 La courbe représente une fonction $f$ et la droite est la tangente au point d'abscisse $1$.\\
        Déterminer $f'(1)$.\begin{center}\begin{tikzpicture}[baseline,scale = 0.5]

    \tikzset{
      point/.style={
        thick,
        draw,
        cross out,
        inner sep=0pt,
        minimum width=5pt,
        minimum height=5pt,
      },
    }
    \clip (-2,-2) rectangle (14,12);
    	
	\draw[color ={black},line width = 2,>=latex,->] (-2,0)--(14,0);
	\draw[color ={black},line width = 2,>=latex,->] (0,-2)--(0,12);
	\draw[color ={black},opacity = 0.5] (-2,1)--(14,1);
	\draw[color ={black},opacity = 0.5] (-2,2)--(14,2);
	\draw[color ={black},opacity = 0.5] (-2,3)--(14,3);
	\draw[color ={black},opacity = 0.5] (-2,4)--(14,4);
	\draw[color ={black},opacity = 0.5] (-2,5)--(14,5);
	\draw[color ={black},opacity = 0.5] (-2,6)--(14,6);
	\draw[color ={black},opacity = 0.5] (-2,7)--(14,7);
	\draw[color ={black},opacity = 0.5] (-2,8)--(14,8);
	\draw[color ={black},opacity = 0.5] (-2,9)--(14,9);
	\draw[color ={black},opacity = 0.5] (-2,10)--(14,10);
	\draw[color ={black},opacity = 0.5] (-2,11)--(14,11);
	\draw[color ={black},opacity = 0.5] (-2,12)--(14,12);
	\draw[color ={black},opacity = 0.5] (-2,-1)--(14,-1);
	\draw[color ={black},opacity = 0.5] (-2,-2)--(14,-2);
	\draw[color ={black},opacity = 0.5] (1,-2)--(1,12);
	\draw[color ={black},opacity = 0.5] (2,-2)--(2,12);
	\draw[color ={black},opacity = 0.5] (3,-2)--(3,12);
	\draw[color ={black},opacity = 0.5] (4,-2)--(4,12);
	\draw[color ={black},opacity = 0.5] (5,-2)--(5,12);
	\draw[color ={black},opacity = 0.5] (6,-2)--(6,12);
	\draw[color ={black},opacity = 0.5] (7,-2)--(7,12);
	\draw[color ={black},opacity = 0.5] (8,-2)--(8,12);
	\draw[color ={black},opacity = 0.5] (9,-2)--(9,12);
	\draw[color ={black},opacity = 0.5] (10,-2)--(10,12);
	\draw[color ={black},opacity = 0.5] (11,-2)--(11,12);
	\draw[color ={black},opacity = 0.5] (12,-2)--(12,12);
	\draw[color ={black},opacity = 0.5] (13,-2)--(13,12);
	\draw[color ={black},opacity = 0.5] (14,-2)--(14,12);
	\draw[color ={black},opacity = 0.5] (-1,-2)--(-1,12);
	\draw[color ={black},opacity = 0.5] (-2,-2)--(-2,12);
	\draw[color ={gray},opacity = 0.3] (-2,0)--(14,0);
	\draw[color ={gray},opacity = 0.3] (-2,0.5)--(14,0.5);
	\draw[color ={gray},opacity = 0.3] (-2,1.5)--(14,1.5);
	\draw[color ={gray},opacity = 0.3] (-2,2.5)--(14,2.5);
	\draw[color ={gray},opacity = 0.3] (-2,3.5)--(14,3.5);
	\draw[color ={gray},opacity = 0.3] (-2,4.5)--(14,4.5);
	\draw[color ={gray},opacity = 0.3] (-2,5.5)--(14,5.5);
	\draw[color ={gray},opacity = 0.3] (-2,6.5)--(14,6.5);
	\draw[color ={gray},opacity = 0.3] (-2,7.5)--(14,7.5);
	\draw[color ={gray},opacity = 0.3] (-2,8.5)--(14,8.5);
	\draw[color ={gray},opacity = 0.3] (-2,9.5)--(14,9.5);
	\draw[color ={gray},opacity = 0.3] (-2,10.5)--(14,10.5);
	\draw[color ={gray},opacity = 0.3] (-2,11.5)--(14,11.5);
	\draw[color ={gray},opacity = 0.3] (-2,0)--(14,0);
	\draw[color ={gray},opacity = 0.3] (-2,-0.5)--(14,-0.5);
	\draw[color ={gray},opacity = 0.3] (-2,-1.5)--(14,-1.5);
	\draw[color ={gray},opacity = 0.3] (0,-2)--(0,12);
	\draw[color ={gray},opacity = 0.3] (0.1,-2)--(0.1,12);
	\draw[color ={gray},opacity = 0.3] (0.2,-2)--(0.2,12);
	\draw[color ={gray},opacity = 0.3] (0.30000000000000004,-2)--(0.30000000000000004,12);
	\draw[color ={gray},opacity = 0.3] (0.4,-2)--(0.4,12);
	\draw[color ={gray},opacity = 0.3] (0.5,-2)--(0.5,12);
	\draw[color ={gray},opacity = 0.3] (0.6,-2)--(0.6,12);
	\draw[color ={gray},opacity = 0.3] (0.7,-2)--(0.7,12);
	\draw[color ={gray},opacity = 0.3] (0.7999999999999999,-2)--(0.7999999999999999,12);
	\draw[color ={gray},opacity = 0.3] (0.8999999999999999,-2)--(0.8999999999999999,12);
	\draw[color ={gray},opacity = 0.3] (1.0999999999999999,-2)--(1.0999999999999999,12);
	\draw[color ={gray},opacity = 0.3] (1.2,-2)--(1.2,12);
	\draw[color ={gray},opacity = 0.3] (1.3,-2)--(1.3,12);
	\draw[color ={gray},opacity = 0.3] (1.4000000000000001,-2)--(1.4000000000000001,12);
	\draw[color ={gray},opacity = 0.3] (1.5000000000000002,-2)--(1.5000000000000002,12);
	\draw[color ={gray},opacity = 0.3] (1.6000000000000003,-2)--(1.6000000000000003,12);
	\draw[color ={gray},opacity = 0.3] (1.7000000000000004,-2)--(1.7000000000000004,12);
	\draw[color ={gray},opacity = 0.3] (1.8000000000000005,-2)--(1.8000000000000005,12);
	\draw[color ={gray},opacity = 0.3] (1.9000000000000006,-2)--(1.9000000000000006,12);
	\draw[color ={gray},opacity = 0.3] (2.1000000000000005,-2)--(2.1000000000000005,12);
	\draw[color ={gray},opacity = 0.3] (2.2000000000000006,-2)--(2.2000000000000006,12);
	\draw[color ={gray},opacity = 0.3] (2.3000000000000007,-2)--(2.3000000000000007,12);
	\draw[color ={gray},opacity = 0.3] (2.400000000000001,-2)--(2.400000000000001,12);
	\draw[color ={gray},opacity = 0.3] (2.500000000000001,-2)--(2.500000000000001,12);
	\draw[color ={gray},opacity = 0.3] (2.600000000000001,-2)--(2.600000000000001,12);
	\draw[color ={gray},opacity = 0.3] (2.700000000000001,-2)--(2.700000000000001,12);
	\draw[color ={gray},opacity = 0.3] (2.800000000000001,-2)--(2.800000000000001,12);
	\draw[color ={gray},opacity = 0.3] (2.9000000000000012,-2)--(2.9000000000000012,12);
	\draw[color ={gray},opacity = 0.3] (3.1000000000000014,-2)--(3.1000000000000014,12);
	\draw[color ={gray},opacity = 0.3] (3.2000000000000015,-2)--(3.2000000000000015,12);
	\draw[color ={gray},opacity = 0.3] (3.3000000000000016,-2)--(3.3000000000000016,12);
	\draw[color ={gray},opacity = 0.3] (3.4000000000000017,-2)--(3.4000000000000017,12);
	\draw[color ={gray},opacity = 0.3] (3.5000000000000018,-2)--(3.5000000000000018,12);
	\draw[color ={gray},opacity = 0.3] (3.600000000000002,-2)--(3.600000000000002,12);
	\draw[color ={gray},opacity = 0.3] (3.700000000000002,-2)--(3.700000000000002,12);
	\draw[color ={gray},opacity = 0.3] (3.800000000000002,-2)--(3.800000000000002,12);
	\draw[color ={gray},opacity = 0.3] (3.900000000000002,-2)--(3.900000000000002,12);
	\draw[color ={gray},opacity = 0.3] (4.100000000000001,-2)--(4.100000000000001,12);
	\draw[color ={gray},opacity = 0.3] (4.200000000000001,-2)--(4.200000000000001,12);
	\draw[color ={gray},opacity = 0.3] (4.300000000000001,-2)--(4.300000000000001,12);
	\draw[color ={gray},opacity = 0.3] (4.4,-2)--(4.4,12);
	\draw[color ={gray},opacity = 0.3] (4.5,-2)--(4.5,12);
	\draw[color ={gray},opacity = 0.3] (4.6,-2)--(4.6,12);
	\draw[color ={gray},opacity = 0.3] (4.699999999999999,-2)--(4.699999999999999,12);
	\draw[color ={gray},opacity = 0.3] (4.799999999999999,-2)--(4.799999999999999,12);
	\draw[color ={gray},opacity = 0.3] (4.899999999999999,-2)--(4.899999999999999,12);
	\draw[color ={gray},opacity = 0.3] (5.099999999999998,-2)--(5.099999999999998,12);
	\draw[color ={gray},opacity = 0.3] (5.1999999999999975,-2)--(5.1999999999999975,12);
	\draw[color ={gray},opacity = 0.3] (5.299999999999997,-2)--(5.299999999999997,12);
	\draw[color ={gray},opacity = 0.3] (5.399999999999997,-2)--(5.399999999999997,12);
	\draw[color ={gray},opacity = 0.3] (5.4999999999999964,-2)--(5.4999999999999964,12);
	\draw[color ={gray},opacity = 0.3] (5.599999999999996,-2)--(5.599999999999996,12);
	\draw[color ={gray},opacity = 0.3] (5.699999999999996,-2)--(5.699999999999996,12);
	\draw[color ={gray},opacity = 0.3] (5.799999999999995,-2)--(5.799999999999995,12);
	\draw[color ={gray},opacity = 0.3] (5.899999999999995,-2)--(5.899999999999995,12);
	\draw[color ={gray},opacity = 0.3] (6.099999999999994,-2)--(6.099999999999994,12);
	\draw[color ={gray},opacity = 0.3] (6.199999999999994,-2)--(6.199999999999994,12);
	\draw[color ={gray},opacity = 0.3] (6.299999999999994,-2)--(6.299999999999994,12);
	\draw[color ={gray},opacity = 0.3] (6.399999999999993,-2)--(6.399999999999993,12);
	\draw[color ={gray},opacity = 0.3] (6.499999999999993,-2)--(6.499999999999993,12);
	\draw[color ={gray},opacity = 0.3] (6.5999999999999925,-2)--(6.5999999999999925,12);
	\draw[color ={gray},opacity = 0.3] (6.699999999999992,-2)--(6.699999999999992,12);
	\draw[color ={gray},opacity = 0.3] (6.799999999999992,-2)--(6.799999999999992,12);
	\draw[color ={gray},opacity = 0.3] (6.8999999999999915,-2)--(6.8999999999999915,12);
	\draw[color ={gray},opacity = 0.3] (7.099999999999991,-2)--(7.099999999999991,12);
	\draw[color ={gray},opacity = 0.3] (7.19999999999999,-2)--(7.19999999999999,12);
	\draw[color ={gray},opacity = 0.3] (7.29999999999999,-2)--(7.29999999999999,12);
	\draw[color ={gray},opacity = 0.3] (7.39999999999999,-2)--(7.39999999999999,12);
	\draw[color ={gray},opacity = 0.3] (7.499999999999989,-2)--(7.499999999999989,12);
	\draw[color ={gray},opacity = 0.3] (7.599999999999989,-2)--(7.599999999999989,12);
	\draw[color ={gray},opacity = 0.3] (7.699999999999989,-2)--(7.699999999999989,12);
	\draw[color ={gray},opacity = 0.3] (7.799999999999988,-2)--(7.799999999999988,12);
	\draw[color ={gray},opacity = 0.3] (7.899999999999988,-2)--(7.899999999999988,12);
	\draw[color ={gray},opacity = 0.3] (8.099999999999987,-2)--(8.099999999999987,12);
	\draw[color ={gray},opacity = 0.3] (8.199999999999987,-2)--(8.199999999999987,12);
	\draw[color ={gray},opacity = 0.3] (8.299999999999986,-2)--(8.299999999999986,12);
	\draw[color ={gray},opacity = 0.3] (8.399999999999986,-2)--(8.399999999999986,12);
	\draw[color ={gray},opacity = 0.3] (8.499999999999986,-2)--(8.499999999999986,12);
	\draw[color ={gray},opacity = 0.3] (8.599999999999985,-2)--(8.599999999999985,12);
	\draw[color ={gray},opacity = 0.3] (8.699999999999985,-2)--(8.699999999999985,12);
	\draw[color ={gray},opacity = 0.3] (8.799999999999985,-2)--(8.799999999999985,12);
	\draw[color ={gray},opacity = 0.3] (8.899999999999984,-2)--(8.899999999999984,12);
	\draw[color ={gray},opacity = 0.3] (9.099999999999984,-2)--(9.099999999999984,12);
	\draw[color ={gray},opacity = 0.3] (9.199999999999983,-2)--(9.199999999999983,12);
	\draw[color ={gray},opacity = 0.3] (9.299999999999983,-2)--(9.299999999999983,12);
	\draw[color ={gray},opacity = 0.3] (9.399999999999983,-2)--(9.399999999999983,12);
	\draw[color ={gray},opacity = 0.3] (9.499999999999982,-2)--(9.499999999999982,12);
	\draw[color ={gray},opacity = 0.3] (9.599999999999982,-2)--(9.599999999999982,12);
	\draw[color ={gray},opacity = 0.3] (9.699999999999982,-2)--(9.699999999999982,12);
	\draw[color ={gray},opacity = 0.3] (9.799999999999981,-2)--(9.799999999999981,12);
	\draw[color ={gray},opacity = 0.3] (9.89999999999998,-2)--(9.89999999999998,12);
	\draw[color ={gray},opacity = 0.3] (10.09999999999998,-2)--(10.09999999999998,12);
	\draw[color ={gray},opacity = 0.3] (10.19999999999998,-2)--(10.19999999999998,12);
	\draw[color ={gray},opacity = 0.3] (10.29999999999998,-2)--(10.29999999999998,12);
	\draw[color ={gray},opacity = 0.3] (10.399999999999979,-2)--(10.399999999999979,12);
	\draw[color ={gray},opacity = 0.3] (10.499999999999979,-2)--(10.499999999999979,12);
	\draw[color ={gray},opacity = 0.3] (10.599999999999978,-2)--(10.599999999999978,12);
	\draw[color ={gray},opacity = 0.3] (10.699999999999978,-2)--(10.699999999999978,12);
	\draw[color ={gray},opacity = 0.3] (10.799999999999978,-2)--(10.799999999999978,12);
	\draw[color ={gray},opacity = 0.3] (10.899999999999977,-2)--(10.899999999999977,12);
	\draw[color ={gray},opacity = 0.3] (11.099999999999977,-2)--(11.099999999999977,12);
	\draw[color ={gray},opacity = 0.3] (11.199999999999976,-2)--(11.199999999999976,12);
	\draw[color ={gray},opacity = 0.3] (11.299999999999976,-2)--(11.299999999999976,12);
	\draw[color ={gray},opacity = 0.3] (11.399999999999975,-2)--(11.399999999999975,12);
	\draw[color ={gray},opacity = 0.3] (11.499999999999975,-2)--(11.499999999999975,12);
	\draw[color ={gray},opacity = 0.3] (11.599999999999975,-2)--(11.599999999999975,12);
	\draw[color ={gray},opacity = 0.3] (11.699999999999974,-2)--(11.699999999999974,12);
	\draw[color ={gray},opacity = 0.3] (11.799999999999974,-2)--(11.799999999999974,12);
	\draw[color ={gray},opacity = 0.3] (11.899999999999974,-2)--(11.899999999999974,12);
	\draw[color ={gray},opacity = 0.3] (12.099999999999973,-2)--(12.099999999999973,12);
	\draw[color ={gray},opacity = 0.3] (12.199999999999973,-2)--(12.199999999999973,12);
	\draw[color ={gray},opacity = 0.3] (12.299999999999972,-2)--(12.299999999999972,12);
	\draw[color ={gray},opacity = 0.3] (12.399999999999972,-2)--(12.399999999999972,12);
	\draw[color ={gray},opacity = 0.3] (12.499999999999972,-2)--(12.499999999999972,12);
	\draw[color ={gray},opacity = 0.3] (12.599999999999971,-2)--(12.599999999999971,12);
	\draw[color ={gray},opacity = 0.3] (12.69999999999997,-2)--(12.69999999999997,12);
	\draw[color ={gray},opacity = 0.3] (12.79999999999997,-2)--(12.79999999999997,12);
	\draw[color ={gray},opacity = 0.3] (12.89999999999997,-2)--(12.89999999999997,12);
	\draw[color ={gray},opacity = 0.3] (13.09999999999997,-2)--(13.09999999999997,12);
	\draw[color ={gray},opacity = 0.3] (13.199999999999969,-2)--(13.199999999999969,12);
	\draw[color ={gray},opacity = 0.3] (13.299999999999969,-2)--(13.299999999999969,12);
	\draw[color ={gray},opacity = 0.3] (13.399999999999968,-2)--(13.399999999999968,12);
	\draw[color ={gray},opacity = 0.3] (13.499999999999968,-2)--(13.499999999999968,12);
	\draw[color ={gray},opacity = 0.3] (13.599999999999968,-2)--(13.599999999999968,12);
	\draw[color ={gray},opacity = 0.3] (13.699999999999967,-2)--(13.699999999999967,12);
	\draw[color ={gray},opacity = 0.3] (13.799999999999967,-2)--(13.799999999999967,12);
	\draw[color ={gray},opacity = 0.3] (13.899999999999967,-2)--(13.899999999999967,12);
	\draw[color ={gray},opacity = 0.3] (0,-2)--(0,12);
	\draw[color ={gray},opacity = 0.3] (-0.1,-2)--(-0.1,12);
	\draw[color ={gray},opacity = 0.3] (-0.2,-2)--(-0.2,12);
	\draw[color ={gray},opacity = 0.3] (-0.30000000000000004,-2)--(-0.30000000000000004,12);
	\draw[color ={gray},opacity = 0.3] (-0.4,-2)--(-0.4,12);
	\draw[color ={gray},opacity = 0.3] (-0.5,-2)--(-0.5,12);
	\draw[color ={gray},opacity = 0.3] (-0.6,-2)--(-0.6,12);
	\draw[color ={gray},opacity = 0.3] (-0.7,-2)--(-0.7,12);
	\draw[color ={gray},opacity = 0.3] (-0.7999999999999999,-2)--(-0.7999999999999999,12);
	\draw[color ={gray},opacity = 0.3] (-0.8999999999999999,-2)--(-0.8999999999999999,12);
	\draw[color ={gray},opacity = 0.3] (-1.0999999999999999,-2)--(-1.0999999999999999,12);
	\draw[color ={gray},opacity = 0.3] (-1.2,-2)--(-1.2,12);
	\draw[color ={gray},opacity = 0.3] (-1.3,-2)--(-1.3,12);
	\draw[color ={gray},opacity = 0.3] (-1.4000000000000001,-2)--(-1.4000000000000001,12);
	\draw[color ={gray},opacity = 0.3] (-1.5000000000000002,-2)--(-1.5000000000000002,12);
	\draw[color ={gray},opacity = 0.3] (-1.6000000000000003,-2)--(-1.6000000000000003,12);
	\draw[color ={gray},opacity = 0.3] (-1.7000000000000004,-2)--(-1.7000000000000004,12);
	\draw[color ={gray},opacity = 0.3] (-1.8000000000000005,-2)--(-1.8000000000000005,12);
	\draw[color ={gray},opacity = 0.3] (-1.9000000000000006,-2)--(-1.9000000000000006,12);
	\draw[color ={black},line width = 2] (2,-0.2)--(2,0.2);
	\draw[color ={black},line width = 2] (4,-0.2)--(4,0.2);
	\draw[color ={black},line width = 2] (6,-0.2)--(6,0.2);
	\draw[color ={black},line width = 2] (8,-0.2)--(8,0.2);
	\draw[color ={black},line width = 2] (10,-0.2)--(10,0.2);
	\draw[color ={black},line width = 2] (12,-0.2)--(12,0.2);
	\draw[color ={black},line width = 2] (-0.2,2)--(0.2,2);
	\draw[color ={black},line width = 2] (-0.2,4)--(0.2,4);
	\draw[color ={black},line width = 2] (-0.2,6)--(0.2,6);
	\draw[color ={black},line width = 2] (-0.2,8)--(0.2,8);
	\draw[color ={black},line width = 2] (-0.2,10)--(0.2,10);
	\draw (2,-0.4) node[anchor = center] {\scriptsize \color{black}{$1$}};
	\draw (4,-0.4) node[anchor = center] {\scriptsize \color{black}{$2$}};
	\draw (6,-0.4) node[anchor = center] {\scriptsize \color{black}{$3$}};
	\draw (8,-0.4) node[anchor = center] {\scriptsize \color{black}{$4$}};
	\draw (10,-0.4) node[anchor = center] {\scriptsize \color{black}{$5$}};
	\draw (-0.5,2.1) node[anchor = center] {\footnotesize \color{black}{$1$}};
	\draw (-0.5,4.1) node[anchor = center] {\footnotesize \color{black}{$2$}};
	\draw (-0.5,6.1) node[anchor = center] {\footnotesize \color{black}{$3$}};
	\draw (-0.5,8.1) node[anchor = center] {\footnotesize \color{black}{$4$}};
	\draw (-0.5,10.1) node[anchor = center] {\footnotesize \color{black}{$5$}};
	\draw [color={black}] (-0.3,-0.3) node[anchor = center,scale=1, rotate = 0] {O};
	
	\draw[color={blue},line width = 2] (1,10)--(1.4,7.71)--(1.8,6.44)--(2.2,5.64)--(2.6,5.08)--(3,4.67)--(3.4,4.35)--(3.8,4.11)--(4.2,3.9)--(4.6,3.74)--(5,3.6)--(5.4,3.48)--(5.8,3.38)--(6.2,3.29)--(6.6,3.21)--(7,3.14)--(7.4,3.08)--(7.8,3.03)--(8.2,2.98)--(8.6,2.93)--(9,2.89)--(9.4,2.85)--(9.8,2.82)--(10.2,2.78)--(10.6,2.75)--(11,2.73)--(11.4,2.7)--(11.8,2.68)--(12.2,2.66)--(12.6,2.63)--(13,2.62)--(13.4,2.6)--(13.8,2.58);
	
	\draw[color={red},line width = 2] (-1.6,13.2)--(-1.2,12.4)--(-0.8,11.6)--(-0.4,10.8)--(0,10)--(0.4,9.2)--(0.8,8.4)--(1.2,7.6)--(1.6,6.8)--(2,6)--(2.4,5.2)--(2.8,4.4)--(3.2,3.6)--(3.6,2.8)--(4,2)--(4.4,1.2)--(4.8,0.4)--(5.2,-0.4)--(5.6,-1.2)--(6,-2)--(6.4,-2.8)--(6.8,-3.6);

\end{tikzpicture}\end{center}
\end{exercice}

\begin{Solution}
 $f'(1)$ est donné par le coefficient directeur de la tangente à la courbe au point d'abscisse $1$, soit $-2=-2$.
\end{Solution}

% @see : https://coopmaths.fr/alea?uuid=6f32d&id=can1F16&n=1&d=10&alea=dqSw2232323232342345678910111213141516171819202122232425262728293031323334353637383940414243444546474823456789101112131415161718192021222324252627282930313233&cd=1&cols=1
\needspace{10\baselineskip}
\newpage
\begin{exercice}[BaremeTotal=false]


 On donne les représentations graphiques d'une fonction et de sa dérivée.\\
        Donner l'équation réduite de la tangente à la courbe de $f$ en $x=1$. \\ \begin{minipage}{0.5\linewidth}
    \begin{tikzpicture}[baseline, scale=0.8]

    \tikzset{
      point/.style={
        thick,
        draw,
        cross out,
        inner sep=0pt,
        minimum width=5pt,
        minimum height=5pt,
      },
    }
    \clip (-4.5,-4.5) rectangle (5.75,13.5);
    	
	\draw[color ={black},>=latex,->] (-3,0)--(3,0);
	\draw[color ={black},>=latex,->] (0,-3)--(0,12);
	\draw[color ={black},opacity = 0.4] (-3,1)--(3,1);
	\draw[color ={black},opacity = 0.4] (-3,2)--(3,2);
	\draw[color ={black},opacity = 0.4] (-3,3)--(3,3);
	\draw[color ={black},opacity = 0.4] (-3,4)--(3,4);
	\draw[color ={black},opacity = 0.4] (-3,5)--(3,5);
	\draw[color ={black},opacity = 0.4] (-3,6)--(3,6);
	\draw[color ={black},opacity = 0.4] (-3,7)--(3,7);
	\draw[color ={black},opacity = 0.4] (-3,8)--(3,8);
	\draw[color ={black},opacity = 0.4] (-3,9)--(3,9);
	\draw[color ={black},opacity = 0.4] (-3,10)--(3,10);
	\draw[color ={black},opacity = 0.4] (-3,11)--(3,11);
	\draw[color ={black},opacity = 0.4] (-3,12)--(3,12);
	\draw[color ={black},opacity = 0.4] (-3,-1)--(3,-1);
	\draw[color ={black},opacity = 0.4] (-3,-2)--(3,-2);
	\draw[color ={black},opacity = 0.4] (-3,-3)--(3,-3);
	\draw[color ={black},opacity = 0.4] (1,-3)--(1,12);
	\draw[color ={black},opacity = 0.4] (2,-3)--(2,12);
	\draw[color ={black},opacity = 0.4] (3,-3)--(3,12);
	\draw[color ={black},opacity = 0.4] (-1,-3)--(-1,12);
	\draw[color ={black},opacity = 0.4] (-2,-3)--(-2,12);
	\draw[color ={black},opacity = 0.4] (-3,-3)--(-3,12);
	\draw[color ={gray},opacity = 0.3] (-3,0)--(3,0);
	\draw[color ={gray},opacity = 0.3] (-3,0.5)--(3,0.5);
	\draw[color ={gray},opacity = 0.3] (-3,1.5)--(3,1.5);
	\draw[color ={gray},opacity = 0.3] (-3,2.5)--(3,2.5);
	\draw[color ={gray},opacity = 0.3] (-3,3.5)--(3,3.5);
	\draw[color ={gray},opacity = 0.3] (-3,4.5)--(3,4.5);
	\draw[color ={gray},opacity = 0.3] (-3,5.5)--(3,5.5);
	\draw[color ={gray},opacity = 0.3] (-3,6.5)--(3,6.5);
	\draw[color ={gray},opacity = 0.3] (-3,7.5)--(3,7.5);
	\draw[color ={gray},opacity = 0.3] (-3,8.5)--(3,8.5);
	\draw[color ={gray},opacity = 0.3] (-3,9.5)--(3,9.5);
	\draw[color ={gray},opacity = 0.3] (-3,10.5)--(3,10.5);
	\draw[color ={gray},opacity = 0.3] (-3,11.5)--(3,11.5);
	\draw[color ={gray},opacity = 0.3] (-3,0)--(3,0);
	\draw[color ={gray},opacity = 0.3] (-3,-0.5)--(3,-0.5);
	\draw[color ={gray},opacity = 0.3] (-3,-1.5)--(3,-1.5);
	\draw[color ={gray},opacity = 0.3] (-3,-2.5)--(3,-2.5);
	\draw[color ={gray},opacity = 0.3] (0,-3)--(0,12);
	\draw[color ={gray},opacity = 0.3] (0,-3)--(0,12);
	\draw[color ={black},line width = 1.2] (1,-0.13)--(1,0.13);
	\draw[color ={black},line width = 1.2] (2,-0.13)--(2,0.13);
	\draw[color ={black},line width = 1.2] (3,-0.13)--(3,0.13);
	\draw[color ={black},line width = 1.2] (-1,-0.13)--(-1,0.13);
	\draw[color ={black},line width = 1.2] (-2,-0.13)--(-2,0.13);
	\draw[color ={black},line width = 1.2] (-3,-0.13)--(-3,0.13);
	\draw[color ={black},line width = 1.2] (-0.13,1)--(0.13,1);
	\draw[color ={black},line width = 1.2] (-0.13,2)--(0.13,2);
	\draw[color ={black},line width = 1.2] (-0.13,3)--(0.13,3);
	\draw[color ={black},line width = 1.2] (-0.13,4)--(0.13,4);
	\draw[color ={black},line width = 1.2] (-0.13,5)--(0.13,5);
	\draw[color ={black},line width = 1.2] (-0.13,6)--(0.13,6);
	\draw[color ={black},line width = 1.2] (-0.13,7)--(0.13,7);
	\draw[color ={black},line width = 1.2] (-0.13,8)--(0.13,8);
	\draw[color ={black},line width = 1.2] (-0.13,9)--(0.13,9);
	\draw[color ={black},line width = 1.2] (-0.13,10)--(0.13,10);
	\draw[color ={black},line width = 1.2] (-0.13,11)--(0.13,11);
	\draw[color ={black},line width = 1.2] (-0.13,12)--(0.13,12);
	\draw[color ={black},line width = 1.2] (-0.13,-1)--(0.13,-1);
	\draw[color ={black},line width = 1.2] (-0.13,-2)--(0.13,-2);
	\draw[color ={black},line width = 1.2] (-0.13,-3)--(0.13,-3);
	\draw (3,-0.4) node[anchor = center] {\scriptsize \color{black}{$3$}};
	\draw (-3,-0.4) node[anchor = center] {\scriptsize \color{black}{$-3$}};
	\draw (-0.5,3.1) node[anchor = center] {\footnotesize \color{black}{$3$}};
	\draw (-0.5,6.1) node[anchor = center] {\footnotesize \color{black}{$6$}};
	\draw (-0.5,9.1) node[anchor = center] {\footnotesize \color{black}{$9$}};
	\draw (-0.5,-2.9) node[anchor = center] {\footnotesize \color{black}{$-3$}};
	\draw [color={black}] (-0.3,-0.3) node[anchor = center,scale=1, rotate = 0] {O};
	\draw (3,10) node[anchor = center] {\footnotesize \color{blue}{$\Large \mathcal{C}_f$}};
	
	\draw[color={blue},line width = 2] (-3,2)--(-2.8,1.24)--(-2.6,0.56)--(-2.4,-0.04)--(-2.2,-0.56)--(-2,-1)--(-1.8,-1.36)--(-1.6,-1.64)--(-1.4,-1.84)--(-1.2,-1.96)--(-1,-2)--(-0.8,-1.96)--(-0.6,-1.84)--(-0.4,-1.64)--(-0.2,-1.36)--(0,-1)--(0.2,-0.56)--(0.4,-0.04)--(0.6,0.56)--(0.8,1.24)--(1,2)--(1.2,2.84)--(1.4,3.76)--(1.6,4.76)--(1.8,5.84)--(2,7)--(2.2,8.24)--(2.4,9.56)--(2.6,10.96)--(2.8,12.44);

\end{tikzpicture}\\
    \end{minipage}
    \begin{minipage}{0.5\linewidth}
    \begin{tikzpicture}[baseline, scale=0.8]

    \tikzset{
      point/.style={
        thick,
        draw,
        cross out,
        inner sep=0pt,
        minimum width=5pt,
        minimum height=5pt,
      },
    }
    \clip (-4.5,-6.5) rectangle (6.5,9.5);
    	
	\draw[color ={black},>=latex,->] (-3,0)--(3,0);
	\draw[color ={black},>=latex,->] (0,-5)--(0,8);
	\draw[color ={black},opacity = 0.4] (-3,1)--(3,1);
	\draw[color ={black},opacity = 0.4] (-3,2)--(3,2);
	\draw[color ={black},opacity = 0.4] (-3,3)--(3,3);
	\draw[color ={black},opacity = 0.4] (-3,4)--(3,4);
	\draw[color ={black},opacity = 0.4] (-3,5)--(3,5);
	\draw[color ={black},opacity = 0.4] (-3,6)--(3,6);
	\draw[color ={black},opacity = 0.4] (-3,7)--(3,7);
	\draw[color ={black},opacity = 0.4] (-3,8)--(3,8);
	\draw[color ={black},opacity = 0.4] (-3,-1)--(3,-1);
	\draw[color ={black},opacity = 0.4] (-3,-2)--(3,-2);
	\draw[color ={black},opacity = 0.4] (-3,-3)--(3,-3);
	\draw[color ={black},opacity = 0.4] (-3,-4)--(3,-4);
	\draw[color ={black},opacity = 0.4] (-3,-5)--(3,-5);
	\draw[color ={black},opacity = 0.4] (1,-5)--(1,8);
	\draw[color ={black},opacity = 0.4] (2,-5)--(2,8);
	\draw[color ={black},opacity = 0.4] (3,-5)--(3,8);
	\draw[color ={black},opacity = 0.4] (-1,-5)--(-1,8);
	\draw[color ={black},opacity = 0.4] (-2,-5)--(-2,8);
	\draw[color ={black},opacity = 0.4] (-3,-5)--(-3,8);
	\draw[color ={gray},opacity = 0.3] (-3,0)--(3,0);
	\draw[color ={gray},opacity = 0.3] (-3,0.5)--(3,0.5);
	\draw[color ={gray},opacity = 0.3] (-3,1.5)--(3,1.5);
	\draw[color ={gray},opacity = 0.3] (-3,2.5)--(3,2.5);
	\draw[color ={gray},opacity = 0.3] (-3,3.5)--(3,3.5);
	\draw[color ={gray},opacity = 0.3] (-3,4.5)--(3,4.5);
	\draw[color ={gray},opacity = 0.3] (-3,5.5)--(3,5.5);
	\draw[color ={gray},opacity = 0.3] (-3,6.5)--(3,6.5);
	\draw[color ={gray},opacity = 0.3] (-3,7.5)--(3,7.5);
	\draw[color ={gray},opacity = 0.3] (-3,0)--(3,0);
	\draw[color ={gray},opacity = 0.3] (-3,-0.5)--(3,-0.5);
	\draw[color ={gray},opacity = 0.3] (-3,-1.5)--(3,-1.5);
	\draw[color ={gray},opacity = 0.3] (-3,-2.5)--(3,-2.5);
	\draw[color ={gray},opacity = 0.3] (-3,-3.5)--(3,-3.5);
	\draw[color ={gray},opacity = 0.3] (-3,-4.5)--(3,-4.5);
	\draw[color ={gray},opacity = 0.3] (0,-5)--(0,8);
	\draw[color ={gray},opacity = 0.3] (0,-5)--(0,8);
	\draw[color ={black},line width = 1.2] (1,-0.13)--(1,0.13);
	\draw[color ={black},line width = 1.2] (2,-0.13)--(2,0.13);
	\draw[color ={black},line width = 1.2] (3,-0.13)--(3,0.13);
	\draw[color ={black},line width = 1.2] (-1,-0.13)--(-1,0.13);
	\draw[color ={black},line width = 1.2] (-2,-0.13)--(-2,0.13);
	\draw[color ={black},line width = 1.2] (-3,-0.13)--(-3,0.13);
	\draw[color ={black},line width = 1.2] (-0.13,1)--(0.13,1);
	\draw[color ={black},line width = 1.2] (-0.13,2)--(0.13,2);
	\draw[color ={black},line width = 1.2] (-0.13,3)--(0.13,3);
	\draw[color ={black},line width = 1.2] (-0.13,4)--(0.13,4);
	\draw[color ={black},line width = 1.2] (-0.13,5)--(0.13,5);
	\draw[color ={black},line width = 1.2] (-0.13,6)--(0.13,6);
	\draw[color ={black},line width = 1.2] (-0.13,7)--(0.13,7);
	\draw[color ={black},line width = 1.2] (-0.13,8)--(0.13,8);
	\draw[color ={black},line width = 1.2] (-0.13,9)--(0.13,9);
	\draw[color ={black},line width = 1.2] (-0.13,10)--(0.13,10);
	\draw[color ={black},line width = 1.2] (-0.13,11)--(0.13,11);
	\draw[color ={black},line width = 1.2] (-0.13,12)--(0.13,12);
	\draw[color ={black},line width = 1.2] (-0.13,-1)--(0.13,-1);
	\draw[color ={black},line width = 1.2] (-0.13,-2)--(0.13,-2);
	\draw[color ={black},line width = 1.2] (-0.13,-3)--(0.13,-3);
	\draw (3,-0.4) node[anchor = center] {\scriptsize \color{black}{$3$}};
	\draw (-3,-0.4) node[anchor = center] {\scriptsize \color{black}{$-3$}};
	\draw (-0.5,3.1) node[anchor = center] {\footnotesize \color{black}{$3$}};
	\draw (-0.5,6.1) node[anchor = center] {\footnotesize \color{black}{$6$}};
	\draw (-0.5,-2.9) node[anchor = center] {\footnotesize \color{black}{$-3$}};
	\draw [color={black}] (-0.3,-0.3) node[anchor = center,scale=1, rotate = 0] {O};
	\draw (3,6) node[anchor = center] {\footnotesize \color{red}{$\Large\mathcal{C}_f\prime$}};
	
	\draw[color={red},line width = 2] (-3,-4)--(-2.8,-3.6)--(-2.6,-3.2)--(-2.4,-2.8)--(-2.2,-2.4)--(-2,-2)--(-1.8,-1.6)--(-1.6,-1.2)--(-1.4,-0.8)--(-1.2,-0.4)--(-1,0)--(-0.8,0.4)--(-0.6,0.8)--(-0.4,1.2)--(-0.2,1.6)--(0,2)--(0.2,2.4)--(0.4,2.8)--(0.6,3.2)--(0.8,3.6)--(1,4)--(1.2,4.4)--(1.4,4.8)--(1.6,5.2)--(1.8,5.6)--(2,6)--(2.2,6.4)--(2.4,6.8)--(2.6,7.2)--(2.8,7.6)--(3,8);

\end{tikzpicture}\\
    \end{minipage}
    
\end{exercice}

\begin{Solution}
 L'équation réduite de la tangente au point d'abscisse $1$ est  : $y=f'(1)(x-1)+f(1)$.\\
        On lit graphiquement $f(1)=2$ et $f'(1)=4$.\\
        L'équation réduite de la tangente est donc donnée par :
        $y=4(x-1)+2$, soit $y=4x-2$.
\end{Solution}

% @see : https://coopmaths.fr/alea?uuid=a1ba2&id=can1F14&n=1&d=10&alea=rkIp2232323232342345678910111213141516171819202122232425262728293031323334353637383940414243444546474823456789101112131415161718192021222324252627282930313233&cd=1&cols=1
\needspace{10\baselineskip}
\begin{exercice}[BaremeTotal=false]


 Soit $f$ la fonction définie sur $\mathbb{R}$ par : $f(x)= 2x^2-3x-5$.\\
        En calculant la limite du taux d'accroissement, déterminer $f'(2)$.
\end{exercice}

\begin{Solution}
 $f$ est une fonction polynôme du second degré de la forme $f(x)=ax^2+bx+c$.\\
    La fonction dérivée est donnée par la somme des dérivées des fonctions $u$ et $v$ définies par $u(x)=2x^2$ et $v(x)=-3x-5$.\\
     Comme $u'(x)=4x$ et $v'(x)=-3$, on obtient  $f'(x)=4x-3$. \\
     Ainsi, $f'(2)=4\times 2-3=5$.
\end{Solution}

% @see : https://coopmaths.fr/alea?uuid=3c690&id=can1F13&n=1&d=10&alea=EUDu2232323232342345678910111213141516171819202122232425262728293031323334353637383940414243444546474823456789101112131415161718192021222324252627282930313233&cd=1&cols=1
\needspace{10\baselineskip}
\begin{exercice}[BaremeTotal=false]


 Déterminer le coefficient directeur de la tangente à la courbe représentative de la fonction carré au point d'abscisse $9$.    
\end{exercice}

\begin{Solution}
 Le coefficient directeur de la tangente au point d'abscisse $a$ est donné par le nombre dérivé $f'(a)$.\\
        La fonction $f$ définie par $f(x)=x^2$ a pour fonction dérivée la fonction $f'$ définie par $f'(x)=2x$.\\
        Comme $f'(9)=2\times 9=18$, le coefficient directeur de la tangente au point d'abscisse $9$ est : $18$. 
\end{Solution}

% @see : https://coopmaths.fr/alea?uuid=ebd89&id=1AN14-1&n=1&d=10&s=3&alea=ocYr2232323232342345678910111213141516171819202122232425262728293031323334353637383940414243444546474823456789101112131415161718192021222324252627282930313233&cd=0&cols=1
\needspace{10\baselineskip}
\begin{exercice}[BaremeTotal=false]


 Donner la dérivée de la fonction $f$, dérivable pour tout $x\in\R$, définie par  $f(x)=-\dfrac{3}{10}x+9$.
\end{exercice}

\begin{Solution}
 L'expression de la dérivée de la fonction $f$ définie par $f(x)=-\dfrac{3}{10}x+9$ est : ${\color[HTML]{f15929}\boldsymbol{f'(x)=-\dfrac{3}{10}}}$.
\end{Solution}

% @see : https://coopmaths.fr/alea?uuid=ebd89&id=1AN14-1&n=1&d=10&s=4&alea=uAlP2232323232342345678910111213141516171819202122232425262728293031323334353637383940414243444546474823456789101112131415161718192021222324252627282930313233&cd=0&cols=1
\needspace{10\baselineskip}
\begin{exercice}[BaremeTotal=false]


 Donner la dérivée de la fonction $f$, dérivable pour tout $x\in\R$, définie par  $f(x)=-3x^{4}$.
\end{exercice}

\begin{Solution}
 L'expression de la dérivée de la fonction $f$ définie par $f(x)=-3x^{4}$ est : ${\color[HTML]{f15929}\boldsymbol{f'(x)=-12x^{3}}}$.
\end{Solution}

% @see : https://coopmaths.fr/alea?uuid=ebd89&id=1AN14-1&n=1&d=10&s=7&alea=vjN72232323232342345678910111213141516171819202122232425262728293031323334353637383940414243444546474823456789101112131415161718192021222324252627282930313233&cd=0&cols=1
\needspace{10\baselineskip}
\begin{exercice}[BaremeTotal=false]


 Donner la dérivée de la fonction $f$, dérivable pour tout $x\in\R^*$, définie par  $f(x)=\dfrac{2}{3x}$.
\end{exercice}

\begin{Solution}
 L'expression de la dérivée de la fonction $f$ définie par $f(x)=\dfrac{2}{3x}$ est : ${\color[HTML]{f15929}\boldsymbol{f'(x)=-\dfrac{2}{3x^2}}}$.
\end{Solution}

% @see : https://coopmaths.fr/alea?uuid=a83c0&id=1AN14-4&n=1&d=10&s=1&alea=Njr32232323232342345678910111213141516171819202122232425262728293031323334353637383940414243444546474823456789101112131415161718192021222324252627282930313233&cd=0&cols=1
\needspace{10\baselineskip}
\begin{exercice}[BaremeTotal=false]


 Donner l'expression de la dérivée de la fonction $f$ définie sur $\R^{*}$ par $f(x)=-2x-\dfrac{9}{x}$.\\
\end{exercice}

\begin{Solution}
 L'expression de la dérivée de la fonction $f$ définie par $f(x)=-2x-\dfrac{9}{x}$ est :\\$f'(x)={\color[HTML]{f15929}\boldsymbol{-2+\dfrac{9}{x^2}}}$.\\
\end{Solution}

% @see : https://coopmaths.fr/alea?uuid=a83c0&id=1AN14-4&n=1&d=10&s=4&alea=4Vpz2232323232342345678910111213141516171819202122232425262728293031323334353637383940414243444546474823456789101112131415161718192021222324252627282930313233&cd=1&cols=1
\needspace{10\baselineskip}
\begin{exercice}[BaremeTotal=false]


 Donner l'expression de la dérivée de la fonction $f$ définie sur $\R_+^{*}$ par $f(x)=-\dfrac{3}{x}+5\sqrt{x}-2x^{2}$.\\
\end{exercice}

\begin{Solution}
 $(-\dfrac{3}{x})^\prime=\dfrac{3}{x^2}$\\$(5\sqrt{x})^\prime=\dfrac{5}{2\sqrt{x}}$\\$(-2x^{2})^\prime=-4x$\\L'expression de la dérivée de la fonction $f$ définie par $f(x)=-\dfrac{3}{x}+5\sqrt{x}-2x^{2}$ est :\\$f'(x)={\color[HTML]{f15929}\boldsymbol{\dfrac{3}{x^2}+\dfrac{5}{2\sqrt{x}}-4x}}$.\\
\end{Solution}

\end{Maquette}
\clearpage
\begin{Maquette}[DM]{Niveau= ,Classe=\getdata[34]\mydata\ (v34), Date = 06/01/25  ,Theme=Exercices}


% @see : https://coopmaths.fr/alea?uuid=4c8c7&id=1AN11-3&n=1&d=10&s=2&alea=nawO223232323234234567891011121314151617181920212223242526272829303132333435363738394041424344454647482345678910111213141516171819202122232425262728293031323334&cd=1&cols=2
\needspace{10\baselineskip}
\begin{exercice}[BaremeTotal=false]
\label{eleve:34}


 \begin{minipage}[t]{\linewidth} Soit $f$ une fonction dérivable sur $[-5;5]$ et $\mathcal{C}_f$ sa courbe représentative.\\On sait que  $f(3)=5~~$ et que $~~f'(3)=-2$.\\Déterminer une équation de la tangente $(T)$ à la courbe $\mathcal{C}_f$ au point d'abscisse $3$,\\sans utiliser la formule de cours de l'équation de tangente. \end{minipage}

\end{exercice}

\begin{Solution}
 \begin{minipage}[t]{\linewidth} On sait que la tangente n'est pas une droite verticale, puisque la fonction est dérivable sur l'intervalle.\\On en déduit que la tangente $(T)$ au point d'abscisse $3$, admet une équation réduite de la forme :  $(T) : y=m x + p$. 

\medskip
$\bullet$ Détermination de $m$ :\\On sait que le nombre dérivé en $3$ est par définition, le coefficient directeur de la tangente au point d'abscisse $3$.\\Par conséquent, on a déjà : $m=f'(3)=-2$.\\ On en déduit que  $(T) : y= -2 x + p$\\$\bullet$ Détermination de $p$ :\\ Pour cela, on utilise que si $f(3)=5~~$, alors le point $A$ de coordonnées $(3;5)$ appartient à $\mathcal{C}_f$ mais aussi à $(T)$.\\ On peut écrire $A(3;5) = \mathcal{C}_f \cap (T)$.\\ On remplace alors les coordonnées de $A(3;5)$ dans l'équation  $(T) : y= -2 x + p$.\\ $\begin{aligned} \phantom{\iff}&A(3;5)\in (T)\\ \iff& 5= -2 \times 3  + p\\ \iff& p=5 +2 \times 3  \\ \iff& p=5 +6   \\ \iff& p=11\\\end{aligned}$\\On peut conclure que : $(T) : {\color[HTML]{f15929}\boldsymbol{y=-2x+11}}$. \end{minipage}

\end{Solution}

% @see : https://coopmaths.fr/alea?uuid=0e984&id=can1F15&n=1&d=10&alea=1Alo223232323234234567891011121314151617181920212223242526272829303132333435363738394041424344454647482345678910111213141516171819202122232425262728293031323334&cd=1&cols=1
\needspace{10\baselineskip}
\begin{exercice}[BaremeTotal=false]


 La courbe représente une fonction $f$ et la droite est la tangente au point d'abscisse $-1$.\\
        Déterminer $f'(-1)$.\begin{center}\begin{tikzpicture}[baseline,scale = 0.5]

    \tikzset{
      point/.style={
        thick,
        draw,
        cross out,
        inner sep=0pt,
        minimum width=5pt,
        minimum height=5pt,
      },
    }
    \clip (-8,-10) rectangle (8,3);
    	
	\draw[color ={black},line width = 2,>=latex,->] (-8,0)--(8,0);
	\draw[color ={black},line width = 2,>=latex,->] (0,-10)--(0,3);
	\draw[color ={black},opacity = 0.5] (-8,1)--(8,1);
	\draw[color ={black},opacity = 0.5] (-8,2)--(8,2);
	\draw[color ={black},opacity = 0.5] (-8,3)--(8,3);
	\draw[color ={black},opacity = 0.5] (-8,-1)--(8,-1);
	\draw[color ={black},opacity = 0.5] (-8,-2)--(8,-2);
	\draw[color ={black},opacity = 0.5] (-8,-3)--(8,-3);
	\draw[color ={black},opacity = 0.5] (-8,-4)--(8,-4);
	\draw[color ={black},opacity = 0.5] (-8,-5)--(8,-5);
	\draw[color ={black},opacity = 0.5] (-8,-6)--(8,-6);
	\draw[color ={black},opacity = 0.5] (-8,-7)--(8,-7);
	\draw[color ={black},opacity = 0.5] (-8,-8)--(8,-8);
	\draw[color ={black},opacity = 0.5] (-8,-9)--(8,-9);
	\draw[color ={black},opacity = 0.5] (-8,-10)--(8,-10);
	\draw[color ={black},opacity = 0.5] (1,-10)--(1,3);
	\draw[color ={black},opacity = 0.5] (2,-10)--(2,3);
	\draw[color ={black},opacity = 0.5] (3,-10)--(3,3);
	\draw[color ={black},opacity = 0.5] (4,-10)--(4,3);
	\draw[color ={black},opacity = 0.5] (5,-10)--(5,3);
	\draw[color ={black},opacity = 0.5] (6,-10)--(6,3);
	\draw[color ={black},opacity = 0.5] (7,-10)--(7,3);
	\draw[color ={black},opacity = 0.5] (8,-10)--(8,3);
	\draw[color ={black},opacity = 0.5] (-1,-10)--(-1,3);
	\draw[color ={black},opacity = 0.5] (-2,-10)--(-2,3);
	\draw[color ={black},opacity = 0.5] (-3,-10)--(-3,3);
	\draw[color ={black},opacity = 0.5] (-4,-10)--(-4,3);
	\draw[color ={black},opacity = 0.5] (-5,-10)--(-5,3);
	\draw[color ={black},opacity = 0.5] (-6,-10)--(-6,3);
	\draw[color ={black},opacity = 0.5] (-7,-10)--(-7,3);
	\draw[color ={black},opacity = 0.5] (-8,-10)--(-8,3);
	\draw[color ={gray},opacity = 0.3] (-8,0)--(8,0);
	\draw[color ={gray},opacity = 0.3] (-8,0)--(8,0);
	\draw[color ={gray},opacity = 0.3] (-8,-4)--(8,-4);
	\draw[color ={gray},opacity = 0.3] (-8,-5)--(8,-5);
	\draw[color ={gray},opacity = 0.3] (-8,-6)--(8,-6);
	\draw[color ={gray},opacity = 0.3] (-8,-7)--(8,-7);
	\draw[color ={gray},opacity = 0.3] (-8,-8)--(8,-8);
	\draw[color ={gray},opacity = 0.3] (-8,-9)--(8,-9);
	\draw[color ={gray},opacity = 0.3] (-8,-10)--(8,-10);
	\draw[color ={gray},opacity = 0.3] (0,-10)--(0,3);
	\draw[color ={gray},opacity = 0.3] (0,-10)--(0,3);
	\draw[color ={black},line width = 2] (2,-0.2)--(2,0.2);
	\draw[color ={black},line width = 2] (4,-0.2)--(4,0.2);
	\draw[color ={black},line width = 2] (6,-0.2)--(6,0.2);
	\draw[color ={black},line width = 2] (-2,-0.2)--(-2,0.2);
	\draw[color ={black},line width = 2] (-4,-0.2)--(-4,0.2);
	\draw[color ={black},line width = 2] (-6,-0.2)--(-6,0.2);
	\draw[color ={black},line width = 2] (-0.2,1)--(0.2,1);
	\draw[color ={black},line width = 2] (-0.2,2)--(0.2,2);
	\draw[color ={black},line width = 2] (-0.2,-1)--(0.2,-1);
	\draw[color ={black},line width = 2] (-0.2,-2)--(0.2,-2);
	\draw[color ={black},line width = 2] (-0.2,-3)--(0.2,-3);
	\draw[color ={black},line width = 2] (-0.2,-4)--(0.2,-4);
	\draw[color ={black},line width = 2] (-0.2,-5)--(0.2,-5);
	\draw[color ={black},line width = 2] (-0.2,-6)--(0.2,-6);
	\draw[color ={black},line width = 2] (-0.2,-7)--(0.2,-7);
	\draw[color ={black},line width = 2] (-0.2,-8)--(0.2,-8);
	\draw[color ={black},line width = 2] (-0.2,-9)--(0.2,-9);
	\draw (2,-0.4) node[anchor = center] {\scriptsize \color{black}{$1$}};
	\draw (4,-0.4) node[anchor = center] {\scriptsize \color{black}{$2$}};
	\draw (6,-0.4) node[anchor = center] {\scriptsize \color{black}{$3$}};
	\draw (-2,-0.4) node[anchor = center] {\scriptsize \color{black}{$-1$}};
	\draw (-4,-0.4) node[anchor = center] {\scriptsize \color{black}{$-2$}};
	\draw (-6,-0.4) node[anchor = center] {\scriptsize \color{black}{$-3$}};
	\draw (-0.8,2.1) node[anchor = center] {\footnotesize \color{black}{$2$}};
	\draw (-0.8,-1.9) node[anchor = center] {\footnotesize \color{black}{$-2$}};
	\draw (-0.8,-3.9) node[anchor = center] {\footnotesize \color{black}{$-4$}};
	\draw (-0.8,-5.9) node[anchor = center] {\footnotesize \color{black}{$-6$}};
	\draw (-0.8,-7.9) node[anchor = center] {\footnotesize \color{black}{$-8$}};
	\draw [color={black}] (-0.3,-0.3) node[anchor = center,scale=1, rotate = 0] {O};
	
	\draw[color={blue},line width = 2] (-6.4,-10.68)--(-6,-9)--(-5.6,-7.48)--(-5.2,-6.12)--(-4.8,-4.92)--(-4.4,-3.88)--(-4,-3)--(-3.6,-2.28)--(-3.2,-1.72)--(-2.8,-1.32)--(-2.4,-1.08)--(-2,-1)--(-1.6,-1.08)--(-1.2,-1.32)--(-0.8,-1.72)--(-0.4,-2.28)--(0,-3)--(0.4,-3.88)--(0.8,-4.92)--(1.2,-6.12)--(1.6,-7.48)--(2,-9)--(2.4,-10.68);
	
	\draw[color={red},line width = 2] (-8,-1)--(-7.6,-1)--(-7.2,-1)--(-6.8,-1)--(-6.4,-1)--(-6,-1)--(-5.6,-1)--(-5.2,-1)--(-4.8,-1)--(-4.4,-1)--(-4,-1)--(-3.6,-1)--(-3.2,-1)--(-2.8,-1)--(-2.4,-1)--(-2,-1)--(-1.6,-1)--(-1.2,-1)--(-0.8,-1)--(-0.4,-1)--(0,-1)--(0.4,-1)--(0.8,-1)--(1.2,-1)--(1.6,-1)--(2,-1)--(2.4,-1)--(2.8,-1)--(3.2,-1)--(3.6,-1)--(4,-1)--(4.4,-1)--(4.8,-1)--(5.2,-1)--(5.6,-1)--(6,-1)--(6.4,-1)--(6.8,-1)--(7.2,-1)--(7.6,-1)--(8,-1);

\end{tikzpicture}\end{center}
\end{exercice}

\begin{Solution}
 $f'(-1)$ est donné par le coefficient directeur de la tangente à la courbe au point d'abscisse $-1$, soit $0$.
\end{Solution}

% @see : https://coopmaths.fr/alea?uuid=6f32d&id=can1F16&n=1&d=10&alea=dqSw223232323234234567891011121314151617181920212223242526272829303132333435363738394041424344454647482345678910111213141516171819202122232425262728293031323334&cd=1&cols=1
\needspace{10\baselineskip}
\newpage
\begin{exercice}[BaremeTotal=false]


 On donne les représentations graphiques d'une fonction et de sa dérivée.\\
        Donner l'équation réduite de la tangente à la courbe de $f$ en $x=1$. \\ \begin{minipage}{0.5\linewidth}
    \begin{tikzpicture}[baseline, scale=0.8]

    \tikzset{
      point/.style={
        thick,
        draw,
        cross out,
        inner sep=0pt,
        minimum width=5pt,
        minimum height=5pt,
      },
    }
    \clip (-4.5,-4.5) rectangle (5.75,13.5);
    	
	\draw[color ={black},>=latex,->] (-3,0)--(3,0);
	\draw[color ={black},>=latex,->] (0,-3)--(0,12);
	\draw[color ={black},opacity = 0.4] (-3,1)--(3,1);
	\draw[color ={black},opacity = 0.4] (-3,2)--(3,2);
	\draw[color ={black},opacity = 0.4] (-3,3)--(3,3);
	\draw[color ={black},opacity = 0.4] (-3,4)--(3,4);
	\draw[color ={black},opacity = 0.4] (-3,5)--(3,5);
	\draw[color ={black},opacity = 0.4] (-3,6)--(3,6);
	\draw[color ={black},opacity = 0.4] (-3,7)--(3,7);
	\draw[color ={black},opacity = 0.4] (-3,8)--(3,8);
	\draw[color ={black},opacity = 0.4] (-3,9)--(3,9);
	\draw[color ={black},opacity = 0.4] (-3,10)--(3,10);
	\draw[color ={black},opacity = 0.4] (-3,11)--(3,11);
	\draw[color ={black},opacity = 0.4] (-3,12)--(3,12);
	\draw[color ={black},opacity = 0.4] (-3,-1)--(3,-1);
	\draw[color ={black},opacity = 0.4] (-3,-2)--(3,-2);
	\draw[color ={black},opacity = 0.4] (-3,-3)--(3,-3);
	\draw[color ={black},opacity = 0.4] (1,-3)--(1,12);
	\draw[color ={black},opacity = 0.4] (2,-3)--(2,12);
	\draw[color ={black},opacity = 0.4] (3,-3)--(3,12);
	\draw[color ={black},opacity = 0.4] (-1,-3)--(-1,12);
	\draw[color ={black},opacity = 0.4] (-2,-3)--(-2,12);
	\draw[color ={black},opacity = 0.4] (-3,-3)--(-3,12);
	\draw[color ={gray},opacity = 0.3] (-3,0)--(3,0);
	\draw[color ={gray},opacity = 0.3] (-3,0.5)--(3,0.5);
	\draw[color ={gray},opacity = 0.3] (-3,1.5)--(3,1.5);
	\draw[color ={gray},opacity = 0.3] (-3,2.5)--(3,2.5);
	\draw[color ={gray},opacity = 0.3] (-3,3.5)--(3,3.5);
	\draw[color ={gray},opacity = 0.3] (-3,4.5)--(3,4.5);
	\draw[color ={gray},opacity = 0.3] (-3,5.5)--(3,5.5);
	\draw[color ={gray},opacity = 0.3] (-3,6.5)--(3,6.5);
	\draw[color ={gray},opacity = 0.3] (-3,7.5)--(3,7.5);
	\draw[color ={gray},opacity = 0.3] (-3,8.5)--(3,8.5);
	\draw[color ={gray},opacity = 0.3] (-3,9.5)--(3,9.5);
	\draw[color ={gray},opacity = 0.3] (-3,10.5)--(3,10.5);
	\draw[color ={gray},opacity = 0.3] (-3,11.5)--(3,11.5);
	\draw[color ={gray},opacity = 0.3] (-3,0)--(3,0);
	\draw[color ={gray},opacity = 0.3] (-3,-0.5)--(3,-0.5);
	\draw[color ={gray},opacity = 0.3] (-3,-1.5)--(3,-1.5);
	\draw[color ={gray},opacity = 0.3] (-3,-2.5)--(3,-2.5);
	\draw[color ={gray},opacity = 0.3] (0,-3)--(0,12);
	\draw[color ={gray},opacity = 0.3] (0,-3)--(0,12);
	\draw[color ={black},line width = 1.2] (1,-0.13)--(1,0.13);
	\draw[color ={black},line width = 1.2] (2,-0.13)--(2,0.13);
	\draw[color ={black},line width = 1.2] (3,-0.13)--(3,0.13);
	\draw[color ={black},line width = 1.2] (-1,-0.13)--(-1,0.13);
	\draw[color ={black},line width = 1.2] (-2,-0.13)--(-2,0.13);
	\draw[color ={black},line width = 1.2] (-3,-0.13)--(-3,0.13);
	\draw[color ={black},line width = 1.2] (-0.13,1)--(0.13,1);
	\draw[color ={black},line width = 1.2] (-0.13,2)--(0.13,2);
	\draw[color ={black},line width = 1.2] (-0.13,3)--(0.13,3);
	\draw[color ={black},line width = 1.2] (-0.13,4)--(0.13,4);
	\draw[color ={black},line width = 1.2] (-0.13,5)--(0.13,5);
	\draw[color ={black},line width = 1.2] (-0.13,6)--(0.13,6);
	\draw[color ={black},line width = 1.2] (-0.13,7)--(0.13,7);
	\draw[color ={black},line width = 1.2] (-0.13,8)--(0.13,8);
	\draw[color ={black},line width = 1.2] (-0.13,9)--(0.13,9);
	\draw[color ={black},line width = 1.2] (-0.13,10)--(0.13,10);
	\draw[color ={black},line width = 1.2] (-0.13,11)--(0.13,11);
	\draw[color ={black},line width = 1.2] (-0.13,12)--(0.13,12);
	\draw[color ={black},line width = 1.2] (-0.13,-1)--(0.13,-1);
	\draw[color ={black},line width = 1.2] (-0.13,-2)--(0.13,-2);
	\draw[color ={black},line width = 1.2] (-0.13,-3)--(0.13,-3);
	\draw (3,-0.4) node[anchor = center] {\scriptsize \color{black}{$3$}};
	\draw (-3,-0.4) node[anchor = center] {\scriptsize \color{black}{$-3$}};
	\draw (-0.5,3.1) node[anchor = center] {\footnotesize \color{black}{$3$}};
	\draw (-0.5,6.1) node[anchor = center] {\footnotesize \color{black}{$6$}};
	\draw (-0.5,9.1) node[anchor = center] {\footnotesize \color{black}{$9$}};
	\draw (-0.5,-2.9) node[anchor = center] {\footnotesize \color{black}{$-3$}};
	\draw [color={black}] (-0.3,-0.3) node[anchor = center,scale=1, rotate = 0] {O};
	\draw (3,10) node[anchor = center] {\footnotesize \color{blue}{$\Large \mathcal{C}_f$}};
	
	\draw[color={blue},line width = 2] (-3,8)--(-2.8,6.84)--(-2.6,5.76)--(-2.4,4.76)--(-2.2,3.84)--(-2,3)--(-1.8,2.24)--(-1.6,1.56)--(-1.4,0.96)--(-1.2,0.44)--(-1,0)--(-0.8,-0.36)--(-0.6,-0.64)--(-0.4,-0.84)--(-0.2,-0.96)--(0,-1)--(0.2,-0.96)--(0.4,-0.84)--(0.6,-0.64)--(0.8,-0.36)--(1,0)--(1.2,0.44)--(1.4,0.96)--(1.6,1.56)--(1.8,2.24)--(2,3)--(2.2,3.84)--(2.4,4.76)--(2.6,5.76)--(2.8,6.84)--(3,8);

\end{tikzpicture}\\
    \end{minipage}
    \begin{minipage}{0.5\linewidth}
    \begin{tikzpicture}[baseline, scale=0.8]

    \tikzset{
      point/.style={
        thick,
        draw,
        cross out,
        inner sep=0pt,
        minimum width=5pt,
        minimum height=5pt,
      },
    }
    \clip (-4.5,-6.5) rectangle (6.5,9.5);
    	
	\draw[color ={black},>=latex,->] (-3,0)--(3,0);
	\draw[color ={black},>=latex,->] (0,-5)--(0,8);
	\draw[color ={black},opacity = 0.4] (-3,1)--(3,1);
	\draw[color ={black},opacity = 0.4] (-3,2)--(3,2);
	\draw[color ={black},opacity = 0.4] (-3,3)--(3,3);
	\draw[color ={black},opacity = 0.4] (-3,4)--(3,4);
	\draw[color ={black},opacity = 0.4] (-3,5)--(3,5);
	\draw[color ={black},opacity = 0.4] (-3,6)--(3,6);
	\draw[color ={black},opacity = 0.4] (-3,7)--(3,7);
	\draw[color ={black},opacity = 0.4] (-3,8)--(3,8);
	\draw[color ={black},opacity = 0.4] (-3,-1)--(3,-1);
	\draw[color ={black},opacity = 0.4] (-3,-2)--(3,-2);
	\draw[color ={black},opacity = 0.4] (-3,-3)--(3,-3);
	\draw[color ={black},opacity = 0.4] (-3,-4)--(3,-4);
	\draw[color ={black},opacity = 0.4] (-3,-5)--(3,-5);
	\draw[color ={black},opacity = 0.4] (1,-5)--(1,8);
	\draw[color ={black},opacity = 0.4] (2,-5)--(2,8);
	\draw[color ={black},opacity = 0.4] (3,-5)--(3,8);
	\draw[color ={black},opacity = 0.4] (-1,-5)--(-1,8);
	\draw[color ={black},opacity = 0.4] (-2,-5)--(-2,8);
	\draw[color ={black},opacity = 0.4] (-3,-5)--(-3,8);
	\draw[color ={gray},opacity = 0.3] (-3,0)--(3,0);
	\draw[color ={gray},opacity = 0.3] (-3,0.5)--(3,0.5);
	\draw[color ={gray},opacity = 0.3] (-3,1.5)--(3,1.5);
	\draw[color ={gray},opacity = 0.3] (-3,2.5)--(3,2.5);
	\draw[color ={gray},opacity = 0.3] (-3,3.5)--(3,3.5);
	\draw[color ={gray},opacity = 0.3] (-3,4.5)--(3,4.5);
	\draw[color ={gray},opacity = 0.3] (-3,5.5)--(3,5.5);
	\draw[color ={gray},opacity = 0.3] (-3,6.5)--(3,6.5);
	\draw[color ={gray},opacity = 0.3] (-3,7.5)--(3,7.5);
	\draw[color ={gray},opacity = 0.3] (-3,0)--(3,0);
	\draw[color ={gray},opacity = 0.3] (-3,-0.5)--(3,-0.5);
	\draw[color ={gray},opacity = 0.3] (-3,-1.5)--(3,-1.5);
	\draw[color ={gray},opacity = 0.3] (-3,-2.5)--(3,-2.5);
	\draw[color ={gray},opacity = 0.3] (-3,-3.5)--(3,-3.5);
	\draw[color ={gray},opacity = 0.3] (-3,-4.5)--(3,-4.5);
	\draw[color ={gray},opacity = 0.3] (0,-5)--(0,8);
	\draw[color ={gray},opacity = 0.3] (0,-5)--(0,8);
	\draw[color ={black},line width = 1.2] (1,-0.13)--(1,0.13);
	\draw[color ={black},line width = 1.2] (2,-0.13)--(2,0.13);
	\draw[color ={black},line width = 1.2] (3,-0.13)--(3,0.13);
	\draw[color ={black},line width = 1.2] (-1,-0.13)--(-1,0.13);
	\draw[color ={black},line width = 1.2] (-2,-0.13)--(-2,0.13);
	\draw[color ={black},line width = 1.2] (-3,-0.13)--(-3,0.13);
	\draw[color ={black},line width = 1.2] (-0.13,1)--(0.13,1);
	\draw[color ={black},line width = 1.2] (-0.13,2)--(0.13,2);
	\draw[color ={black},line width = 1.2] (-0.13,3)--(0.13,3);
	\draw[color ={black},line width = 1.2] (-0.13,4)--(0.13,4);
	\draw[color ={black},line width = 1.2] (-0.13,5)--(0.13,5);
	\draw[color ={black},line width = 1.2] (-0.13,6)--(0.13,6);
	\draw[color ={black},line width = 1.2] (-0.13,7)--(0.13,7);
	\draw[color ={black},line width = 1.2] (-0.13,8)--(0.13,8);
	\draw[color ={black},line width = 1.2] (-0.13,9)--(0.13,9);
	\draw[color ={black},line width = 1.2] (-0.13,10)--(0.13,10);
	\draw[color ={black},line width = 1.2] (-0.13,11)--(0.13,11);
	\draw[color ={black},line width = 1.2] (-0.13,12)--(0.13,12);
	\draw[color ={black},line width = 1.2] (-0.13,-1)--(0.13,-1);
	\draw[color ={black},line width = 1.2] (-0.13,-2)--(0.13,-2);
	\draw[color ={black},line width = 1.2] (-0.13,-3)--(0.13,-3);
	\draw (3,-0.4) node[anchor = center] {\scriptsize \color{black}{$3$}};
	\draw (-3,-0.4) node[anchor = center] {\scriptsize \color{black}{$-3$}};
	\draw (-0.5,3.1) node[anchor = center] {\footnotesize \color{black}{$3$}};
	\draw (-0.5,6.1) node[anchor = center] {\footnotesize \color{black}{$6$}};
	\draw (-0.5,-2.9) node[anchor = center] {\footnotesize \color{black}{$-3$}};
	\draw [color={black}] (-0.3,-0.3) node[anchor = center,scale=1, rotate = 0] {O};
	\draw (3,6) node[anchor = center] {\footnotesize \color{red}{$\Large\mathcal{C}_f\prime$}};
	
	\draw[color={red},line width = 2] (-2.8,-5.6)--(-2.6,-5.2)--(-2.4,-4.8)--(-2.2,-4.4)--(-2,-4)--(-1.8,-3.6)--(-1.6,-3.2)--(-1.4,-2.8)--(-1.2,-2.4)--(-1,-2)--(-0.8,-1.6)--(-0.6,-1.2)--(-0.4,-0.8)--(-0.2,-0.4)--(0,0)--(0.2,0.4)--(0.4,0.8)--(0.6,1.2)--(0.8,1.6)--(1,2)--(1.2,2.4)--(1.4,2.8)--(1.6,3.2)--(1.8,3.6)--(2,4)--(2.2,4.4)--(2.4,4.8)--(2.6,5.2)--(2.8,5.6)--(3,6);

\end{tikzpicture}\\
    \end{minipage}
    
\end{exercice}

\begin{Solution}
 L'équation réduite de la tangente au point d'abscisse $1$ est  : $y=f'(1)(x-1)+f(1)$.\\
        On lit graphiquement $f(1)=0$ et $f'(1)=2$.\\
        L'équation réduite de la tangente est donc donnée par :
        $y=2(x-1)+0$, soit $y=2x-2$.
\end{Solution}

% @see : https://coopmaths.fr/alea?uuid=a1ba2&id=can1F14&n=1&d=10&alea=rkIp223232323234234567891011121314151617181920212223242526272829303132333435363738394041424344454647482345678910111213141516171819202122232425262728293031323334&cd=1&cols=1
\needspace{10\baselineskip}
\begin{exercice}[BaremeTotal=false]


Soit $f$ la fonction définie sur $\mathbb{R}$ par : 
$f(x)= 7x^2+1$.\\
En calculant la limite du taux d'accroissement, déterminer $f'(4)$.
\end{exercice}

\begin{Solution}
 $f$ est une fonction polynôme du second degré de la forme $f(x)=ax^2+b$.\\
    La fonction dérivée est donnée par la somme des dérivées des fonctions $u$ et $v$ définies par $u(x)=7x^2$ et $v(x)=1$.\\
     Comme $u'(x)=14x$ et $v'(x)=0$, on obtient  $f'(x)=14x$. \\
     Ainsi, $f'(4)=14\times 4=56$.
\end{Solution}

% @see : https://coopmaths.fr/alea?uuid=3c690&id=can1F13&n=1&d=10&alea=EUDu223232323234234567891011121314151617181920212223242526272829303132333435363738394041424344454647482345678910111213141516171819202122232425262728293031323334&cd=1&cols=1
\needspace{10\baselineskip}
\begin{exercice}[BaremeTotal=false]


 Déterminer le coefficient directeur de la tangente à la courbe représentative de la fonction cube au point d'abscisse $2$.
                        
\end{exercice}

\begin{Solution}
 Le coefficient directeur de la tangente au point d'abscisse $a$ est donné par le nombre dérivé $f'(a)$.\\
        La fonction $f$ définie par $f(x)=x^3$ a pour fonction dérivée la fonction $f'$ définie par $f'(x)=3x^2$.\\
        Comme $f'(2)=3\times 2^2=12$, le coefficient directeur de la tangente au point d'abscisse $2$ est : $12$. 
\end{Solution}

% @see : https://coopmaths.fr/alea?uuid=ebd89&id=1AN14-1&n=1&d=10&s=3&alea=ocYr223232323234234567891011121314151617181920212223242526272829303132333435363738394041424344454647482345678910111213141516171819202122232425262728293031323334&cd=0&cols=1
\needspace{10\baselineskip}
\begin{exercice}[BaremeTotal=false]


 Donner la dérivée de la fonction $f$, dérivable pour tout $x\in\R$, définie par  $f(x)=\dfrac{4}{5}x+10$.
\end{exercice}

\begin{Solution}
 L'expression de la dérivée de la fonction $f$ définie par $f(x)=\dfrac{4}{5}x+10$ est : ${\color[HTML]{f15929}\boldsymbol{f'(x)=\dfrac{4}{5}}}$.
\end{Solution}

% @see : https://coopmaths.fr/alea?uuid=ebd89&id=1AN14-1&n=1&d=10&s=4&alea=uAlP223232323234234567891011121314151617181920212223242526272829303132333435363738394041424344454647482345678910111213141516171819202122232425262728293031323334&cd=0&cols=1
\needspace{10\baselineskip}
\begin{exercice}[BaremeTotal=false]


 Donner la dérivée de la fonction $f$, dérivable pour tout $x\in\R$, définie par  $f(x)=2x^{13}$.
\end{exercice}

\begin{Solution}
 L'expression de la dérivée de la fonction $f$ définie par $f(x)=2x^{13}$ est : ${\color[HTML]{f15929}\boldsymbol{f'(x)=26x^{12}}}$.
\end{Solution}

% @see : https://coopmaths.fr/alea?uuid=ebd89&id=1AN14-1&n=1&d=10&s=7&alea=vjN7223232323234234567891011121314151617181920212223242526272829303132333435363738394041424344454647482345678910111213141516171819202122232425262728293031323334&cd=0&cols=1
\needspace{10\baselineskip}
\begin{exercice}[BaremeTotal=false]


 Donner la dérivée de la fonction $f$, dérivable pour tout $x\in\R^*$, définie par  $f(x)=\dfrac{-1}{8x}$.
\end{exercice}

\begin{Solution}
 L'expression de la dérivée de la fonction $f$ définie par $f(x)=\dfrac{-1}{8x}$ est : ${\color[HTML]{f15929}\boldsymbol{f'(x)=\dfrac{1}{8x^2}}}$.
\end{Solution}

% @see : https://coopmaths.fr/alea?uuid=a83c0&id=1AN14-4&n=1&d=10&s=1&alea=Njr3223232323234234567891011121314151617181920212223242526272829303132333435363738394041424344454647482345678910111213141516171819202122232425262728293031323334&cd=0&cols=1
\needspace{10\baselineskip}
\begin{exercice}[BaremeTotal=false]


 Donner l'expression de la dérivée de la fonction $f$ définie sur $\R^{*}$ par $f(x)=-3x^{2}+5x-2-\dfrac{9}{x}$.\\
\end{exercice}

\begin{Solution}
 L'expression de la dérivée de la fonction $f$ définie par $f(x)=-3x^{2}+5x-2-\dfrac{9}{x}$ est :\\$f'(x)={\color[HTML]{f15929}\boldsymbol{-6x+5+\dfrac{9}{x^2}}}$.\\
\end{Solution}

% @see : https://coopmaths.fr/alea?uuid=a83c0&id=1AN14-4&n=1&d=10&s=4&alea=4Vpz223232323234234567891011121314151617181920212223242526272829303132333435363738394041424344454647482345678910111213141516171819202122232425262728293031323334&cd=1&cols=1
\needspace{10\baselineskip}
\begin{exercice}[BaremeTotal=false]


 Donner l'expression de la dérivée de la fonction $f$ définie sur $\R_+^{*}$ par $f(x)=\dfrac{4}{x}-3x^{2}-4x-4\sqrt{x}$.\\
\end{exercice}

\begin{Solution}
 $(\dfrac{4}{x})^\prime=-\dfrac{4}{x^2}$\\$(-3x^{2}-4x)^\prime=-6x-4$\\$(-4\sqrt{x})^\prime=-\dfrac{4}{2\sqrt{x}}$\\L'expression de la dérivée de la fonction $f$ définie par $f(x)=\dfrac{4}{x}-3x^{2}-4x-4\sqrt{x}$ est :\\$f'(x)={\color[HTML]{f15929}\boldsymbol{-\dfrac{4}{x^2}-6x-4-\dfrac{4}{2\sqrt{x}}}}$.\\
\end{Solution}

\end{Maquette}
\clearpage
\begin{Maquette}[DM]{Niveau= ,Classe=\getdata[35]\mydata\ (v35), Date = 06/01/25  ,Theme=Exercices}


% @see : https://coopmaths.fr/alea?uuid=4c8c7&id=1AN11-3&n=1&d=10&s=2&alea=nawO22323232323423456789101112131415161718192021222324252627282930313233343536373839404142434445464748234567891011121314151617181920212223242526272829303132333435&cd=1&cols=2
\needspace{10\baselineskip}
\begin{exercice}[BaremeTotal=false]
\label{eleve:35}


 \begin{minipage}[t]{\linewidth} Soit $f$ une fonction dérivable sur $[-5;5]$ et $\mathcal{C}_f$ sa courbe représentative.\\On sait que  $f(1)=4~~$ et que $~~f'(1)=3$.\\Déterminer une équation de la tangente $(T)$ à la courbe $\mathcal{C}_f$ au point d'abscisse $1$,\\sans utiliser la formule de cours de l'équation de tangente. \end{minipage}

\end{exercice}

\begin{Solution}
 \begin{minipage}[t]{\linewidth} On sait que la tangente n'est pas une droite verticale, puisque la fonction est dérivable sur l'intervalle.\\On en déduit que la tangente $(T)$ au point d'abscisse $1$, admet une équation réduite de la forme :  $(T) : y=m x + p$. 

\medskip
$\bullet$ Détermination de $m$ :\\On sait que le nombre dérivé en $1$ est par définition, le coefficient directeur de la tangente au point d'abscisse $1$.\\Par conséquent, on a déjà : $m=f'(1)=3$.\\ On en déduit que  $(T) : y= 3 x + p$\\$\bullet$ Détermination de $p$ :\\ Pour cela, on utilise que si $f(1)=4~~$, alors le point $A$ de coordonnées $(1;4)$ appartient à $\mathcal{C}_f$ mais aussi à $(T)$.\\ On peut écrire $A(1;4) = \mathcal{C}_f \cap (T)$.\\ On remplace alors les coordonnées de $A(1;4)$ dans l'équation  $(T) : y= 3 x + p$.\\ $\begin{aligned} \phantom{\iff}&A(1;4)\in (T)\\ \iff& 4= 3 \times 1  + p\\ \iff& p=4 -3 \times 1  \\ \iff& p=4 -3   \\ \iff& p=1\\\end{aligned}$\\On peut conclure que : $(T) : {\color[HTML]{f15929}\boldsymbol{y=3x+1}}$. \end{minipage}

\end{Solution}

% @see : https://coopmaths.fr/alea?uuid=0e984&id=can1F15&n=1&d=10&alea=1Alo22323232323423456789101112131415161718192021222324252627282930313233343536373839404142434445464748234567891011121314151617181920212223242526272829303132333435&cd=1&cols=1
\needspace{10\baselineskip}
\begin{exercice}[BaremeTotal=false]


 La courbe représente une fonction $f$ et la droite est la tangente au point d'abscisse $0$.\\
        
        Déterminer $f'(0)$.\begin{center}\begin{tikzpicture}[baseline,scale = 0.5]

    \tikzset{
      point/.style={
        thick,
        draw,
        cross out,
        inner sep=0pt,
        minimum width=5pt,
        minimum height=5pt,
      },
    }
    \clip (-8,-3) rectangle (8,10);
    	
	\draw[color ={black},line width = 2,>=latex,->] (-8,0)--(8,0);
	\draw[color ={black},line width = 2,>=latex,->] (0,-3)--(0,10);
	\draw[color ={black},opacity = 0.5] (-8,1)--(8,1);
	\draw[color ={black},opacity = 0.5] (-8,2)--(8,2);
	\draw[color ={black},opacity = 0.5] (-8,3)--(8,3);
	\draw[color ={black},opacity = 0.5] (-8,4)--(8,4);
	\draw[color ={black},opacity = 0.5] (-8,5)--(8,5);
	\draw[color ={black},opacity = 0.5] (-8,6)--(8,6);
	\draw[color ={black},opacity = 0.5] (-8,7)--(8,7);
	\draw[color ={black},opacity = 0.5] (-8,8)--(8,8);
	\draw[color ={black},opacity = 0.5] (-8,9)--(8,9);
	\draw[color ={black},opacity = 0.5] (-8,10)--(8,10);
	\draw[color ={black},opacity = 0.5] (-8,-1)--(8,-1);
	\draw[color ={black},opacity = 0.5] (-8,-2)--(8,-2);
	\draw[color ={black},opacity = 0.5] (-8,-3)--(8,-3);
	\draw[color ={black},opacity = 0.5] (1,-3)--(1,10);
	\draw[color ={black},opacity = 0.5] (2,-3)--(2,10);
	\draw[color ={black},opacity = 0.5] (3,-3)--(3,10);
	\draw[color ={black},opacity = 0.5] (4,-3)--(4,10);
	\draw[color ={black},opacity = 0.5] (5,-3)--(5,10);
	\draw[color ={black},opacity = 0.5] (6,-3)--(6,10);
	\draw[color ={black},opacity = 0.5] (7,-3)--(7,10);
	\draw[color ={black},opacity = 0.5] (8,-3)--(8,10);
	\draw[color ={black},opacity = 0.5] (-1,-3)--(-1,10);
	\draw[color ={black},opacity = 0.5] (-2,-3)--(-2,10);
	\draw[color ={black},opacity = 0.5] (-3,-3)--(-3,10);
	\draw[color ={black},opacity = 0.5] (-4,-3)--(-4,10);
	\draw[color ={black},opacity = 0.5] (-5,-3)--(-5,10);
	\draw[color ={black},opacity = 0.5] (-6,-3)--(-6,10);
	\draw[color ={black},opacity = 0.5] (-7,-3)--(-7,10);
	\draw[color ={black},opacity = 0.5] (-8,-3)--(-8,10);
	\draw[color ={gray},opacity = 0.3] (-8,0)--(8,0);
	\draw[color ={gray},opacity = 0.3] (-8,0)--(8,0);
	\draw[color ={gray},opacity = 0.3] (0,-3)--(0,10);
	\draw[color ={gray},opacity = 0.3] (0,-3)--(0,10);
	\draw[color ={black},line width = 2] (2,-0.2)--(2,0.2);
	\draw[color ={black},line width = 2] (4,-0.2)--(4,0.2);
	\draw[color ={black},line width = 2] (6,-0.2)--(6,0.2);
	\draw[color ={black},line width = 2] (-2,-0.2)--(-2,0.2);
	\draw[color ={black},line width = 2] (-4,-0.2)--(-4,0.2);
	\draw[color ={black},line width = 2] (-6,-0.2)--(-6,0.2);
	\draw[color ={black},line width = 2] (-0.2,1)--(0.2,1);
	\draw[color ={black},line width = 2] (-0.2,2)--(0.2,2);
	\draw[color ={black},line width = 2] (-0.2,3)--(0.2,3);
	\draw[color ={black},line width = 2] (-0.2,4)--(0.2,4);
	\draw[color ={black},line width = 2] (-0.2,5)--(0.2,5);
	\draw[color ={black},line width = 2] (-0.2,6)--(0.2,6);
	\draw[color ={black},line width = 2] (-0.2,7)--(0.2,7);
	\draw[color ={black},line width = 2] (-0.2,8)--(0.2,8);
	\draw[color ={black},line width = 2] (-0.2,9)--(0.2,9);
	\draw[color ={black},line width = 2] (-0.2,-1)--(0.2,-1);
	\draw[color ={black},line width = 2] (-0.2,-2)--(0.2,-2);
	\draw (2,-0.4) node[anchor = center] {\scriptsize \color{black}{$1$}};
	\draw (4,-0.4) node[anchor = center] {\scriptsize \color{black}{$2$}};
	\draw (6,-0.4) node[anchor = center] {\scriptsize \color{black}{$3$}};
	\draw (-2,-0.4) node[anchor = center] {\scriptsize \color{black}{$-1$}};
	\draw (-4,-0.4) node[anchor = center] {\scriptsize \color{black}{$-2$}};
	\draw (-6,-0.4) node[anchor = center] {\scriptsize \color{black}{$-3$}};
	\draw (-0.8,2.1) node[anchor = center] {\footnotesize \color{black}{$2$}};
	\draw (-0.8,4.1) node[anchor = center] {\footnotesize \color{black}{$4$}};
	\draw (-0.8,6.1) node[anchor = center] {\footnotesize \color{black}{$6$}};
	\draw (-0.8,8.1) node[anchor = center] {\footnotesize \color{black}{$8$}};
	\draw [color={black}] (-0.3,-0.3) node[anchor = center,scale=1, rotate = 0] {O};
	
	\draw[color={blue},line width = 2] (-2.4,10.68)--(-2,9)--(-1.6,7.48)--(-1.2,6.12)--(-0.8,4.92)--(-0.4,3.88)--(0,3)--(0.4,2.28)--(0.8,1.72)--(1.2,1.32)--(1.6,1.08)--(2,1)--(2.4,1.08)--(2.8,1.32)--(3.2,1.72)--(3.6,2.28)--(4,3)--(4.4,3.88)--(4.8,4.92)--(5.2,6.12)--(5.6,7.48)--(6,9)--(6.4,10.68);
	
	\draw[color={red},line width = 2] (-4,11)--(-3.6,10.2)--(-3.2,9.4)--(-2.8,8.6)--(-2.4,7.8)--(-2,7)--(-1.6,6.2)--(-1.2,5.4)--(-0.8,4.6)--(-0.4,3.8)--(0,3)--(0.4,2.2)--(0.8,1.4)--(1.2,0.6)--(1.6,-0.2)--(2,-1)--(2.4,-1.8)--(2.8,-2.6)--(3.2,-3.4);

\end{tikzpicture}\end{center}
\end{exercice}

\begin{Solution}
 $f'(0)$ est donné par le coefficient directeur de la tangente à la courbe au point d'abscisse $0$, soit $-4$.
\end{Solution}

% @see : https://coopmaths.fr/alea?uuid=6f32d&id=can1F16&n=1&d=10&alea=dqSw22323232323423456789101112131415161718192021222324252627282930313233343536373839404142434445464748234567891011121314151617181920212223242526272829303132333435&cd=1&cols=1
\needspace{10\baselineskip}
\newpage
\begin{exercice}[BaremeTotal=false]


 On donne les représentations graphiques d'une fonction et de sa dérivée.\\
      Donner l'équation réduite de la tangente à la courbe de $f$ en $x=0$. \\ \begin{minipage}{0.5\linewidth}
    \begin{tikzpicture}[baseline, scale=0.8]

    \tikzset{
      point/.style={
        thick,
        draw,
        cross out,
        inner sep=0pt,
        minimum width=5pt,
        minimum height=5pt,
      },
    }
    \clip (-5.5,-9.5) rectangle (5.75,7.5);
    	
	\draw[color ={black},>=latex,->] (-4,0)--(4,0);
	\draw[color ={black},>=latex,->] (0,-8)--(0,6);
	\draw[color ={black},opacity = 0.4] (-4,1)--(4,1);
	\draw[color ={black},opacity = 0.4] (-4,2)--(4,2);
	\draw[color ={black},opacity = 0.4] (-4,3)--(4,3);
	\draw[color ={black},opacity = 0.4] (-4,4)--(4,4);
	\draw[color ={black},opacity = 0.4] (-4,5)--(4,5);
	\draw[color ={black},opacity = 0.4] (-4,6)--(4,6);
	\draw[color ={black},opacity = 0.4] (-4,-1)--(4,-1);
	\draw[color ={black},opacity = 0.4] (-4,-2)--(4,-2);
	\draw[color ={black},opacity = 0.4] (-4,-3)--(4,-3);
	\draw[color ={black},opacity = 0.4] (-4,-4)--(4,-4);
	\draw[color ={black},opacity = 0.4] (-4,-5)--(4,-5);
	\draw[color ={black},opacity = 0.4] (-4,-6)--(4,-6);
	\draw[color ={black},opacity = 0.4] (-4,-7)--(4,-7);
	\draw[color ={black},opacity = 0.4] (-4,-8)--(4,-8);
	\draw[color ={black},opacity = 0.4] (1,-8)--(1,6);
	\draw[color ={black},opacity = 0.4] (2,-8)--(2,6);
	\draw[color ={black},opacity = 0.4] (3,-8)--(3,6);
	\draw[color ={black},opacity = 0.4] (4,-8)--(4,6);
	\draw[color ={black},opacity = 0.4] (-1,-8)--(-1,6);
	\draw[color ={black},opacity = 0.4] (-2,-8)--(-2,6);
	\draw[color ={black},opacity = 0.4] (-3,-8)--(-3,6);
	\draw[color ={black},opacity = 0.4] (-4,-8)--(-4,6);
	\draw[color ={gray},opacity = 0.3] (-4,0)--(4,0);
	\draw[color ={gray},opacity = 0.3] (-4,0.5)--(4,0.5);
	\draw[color ={gray},opacity = 0.3] (-4,1.5)--(4,1.5);
	\draw[color ={gray},opacity = 0.3] (-4,2.5)--(4,2.5);
	\draw[color ={gray},opacity = 0.3] (-4,3.5)--(4,3.5);
	\draw[color ={gray},opacity = 0.3] (-4,4.5)--(4,4.5);
	\draw[color ={gray},opacity = 0.3] (-4,5.5)--(4,5.5);
	\draw[color ={gray},opacity = 0.3] (-4,0)--(4,0);
	\draw[color ={gray},opacity = 0.3] (-4,-0.5)--(4,-0.5);
	\draw[color ={gray},opacity = 0.3] (-4,-1.5)--(4,-1.5);
	\draw[color ={gray},opacity = 0.3] (-4,-2.5)--(4,-2.5);
	\draw[color ={gray},opacity = 0.3] (-4,-3.5)--(4,-3.5);
	\draw[color ={gray},opacity = 0.3] (-4,-4.5)--(4,-4.5);
	\draw[color ={gray},opacity = 0.3] (-4,-5.5)--(4,-5.5);
	\draw[color ={gray},opacity = 0.3] (-4,-6.5)--(4,-6.5);
	\draw[color ={gray},opacity = 0.3] (-4,-7)--(4,-7);
	\draw[color ={gray},opacity = 0.3] (-4,-7.5)--(4,-7.5);
	\draw[color ={gray},opacity = 0.3] (-4,-8)--(4,-8);
	\draw[color ={gray},opacity = 0.3] (0,-8)--(0,6);
	\draw[color ={gray},opacity = 0.3] (0,-8)--(0,6);
	\draw[color ={black},line width = 1.2] (1,-0.13)--(1,0.13);
	\draw[color ={black},line width = 1.2] (2,-0.13)--(2,0.13);
	\draw[color ={black},line width = 1.2] (3,-0.13)--(3,0.13);
	\draw[color ={black},line width = 1.2] (4,-0.13)--(4,0.13);
	\draw[color ={black},line width = 1.2] (-1,-0.13)--(-1,0.13);
	\draw[color ={black},line width = 1.2] (-2,-0.13)--(-2,0.13);
	\draw[color ={black},line width = 1.2] (-3,-0.13)--(-3,0.13);
	\draw[color ={black},line width = 1.2] (-4,-0.13)--(-4,0.13);
	\draw[color ={black},line width = 1.2] (-0.13,1)--(0.13,1);
	\draw[color ={black},line width = 1.2] (-0.13,2)--(0.13,2);
	\draw[color ={black},line width = 1.2] (-0.13,3)--(0.13,3);
	\draw[color ={black},line width = 1.2] (-0.13,4)--(0.13,4);
	\draw[color ={black},line width = 1.2] (-0.13,5)--(0.13,5);
	\draw[color ={black},line width = 1.2] (-0.13,6)--(0.13,6);
	\draw[color ={black},line width = 1.2] (-0.13,-1)--(0.13,-1);
	\draw[color ={black},line width = 1.2] (-0.13,-2)--(0.13,-2);
	\draw[color ={black},line width = 1.2] (-0.13,-3)--(0.13,-3);
	\draw[color ={black},line width = 1.2] (-0.13,-4)--(0.13,-4);
	\draw[color ={black},line width = 1.2] (-0.13,-5)--(0.13,-5);
	\draw[color ={black},line width = 1.2] (-0.13,-6)--(0.13,-6);
	\draw[color ={black},line width = 1.2] (-0.13,-7)--(0.13,-7);
	\draw[color ={black},line width = 1.2] (-0.13,-8)--(0.13,-8);
	\draw (2,-0.4) node[anchor = center] {\scriptsize \color{black}{$2$}};
	\draw (4,-0.4) node[anchor = center] {\scriptsize \color{black}{$4$}};
	\draw (-2,-0.4) node[anchor = center] {\scriptsize \color{black}{$-2$}};
	\draw (-4,-0.4) node[anchor = center] {\scriptsize \color{black}{$-4$}};
	\draw (-0.5,2.1) node[anchor = center] {\footnotesize \color{black}{$2$}};
	\draw (-0.5,4.1) node[anchor = center] {\footnotesize \color{black}{$4$}};
	\draw (-0.5,6.1) node[anchor = center] {\footnotesize \color{black}{$6$}};
	\draw (-0.5,-1.9) node[anchor = center] {\footnotesize \color{black}{$-2$}};
	\draw (-0.5,-3.9) node[anchor = center] {\footnotesize \color{black}{$-4$}};
	\draw (-0.5,-5.9) node[anchor = center] {\footnotesize \color{black}{$-6$}};
	\draw (-0.5,-7.9) node[anchor = center] {\footnotesize \color{black}{$-8$}};
	\draw [color={black}] (-0.3,-0.3) node[anchor = center,scale=1, rotate = 0] {O};
	\draw (3,4) node[anchor = center] {\footnotesize \color{blue}{$\Large \mathcal{C}_f$}};
	
	\draw[color={blue},line width = 2] (-4,-1)--(-3.8,-0.24)--(-3.6,0.44)--(-3.4,1.04)--(-3.2,1.56)--(-3,2)--(-2.8,2.36)--(-2.6,2.64)--(-2.4,2.84)--(-2.2,2.96)--(-2,3)--(-1.8,2.96)--(-1.6,2.84)--(-1.4,2.64)--(-1.2,2.36)--(-1,2)--(-0.8,1.56)--(-0.6,1.04)--(-0.4,0.44)--(-0.2,-0.24)--(0,-1)--(0.2,-1.84)--(0.4,-2.76)--(0.6,-3.76)--(0.8,-4.84)--(1,-6)--(1.2,-7.24)--(1.4,-8.56);

\end{tikzpicture}\\
    \end{minipage}
    \begin{minipage}{0.5\linewidth}
    \begin{tikzpicture}[baseline, scale=0.8]

    \tikzset{
      point/.style={
        thick,
        draw,
        cross out,
        inner sep=0pt,
        minimum width=5pt,
        minimum height=5pt,
      },
    }
    \clip (-5.5,-9.5) rectangle (6.5,7.5);
    	
	\draw[color ={black},>=latex,->] (-4,0)--(4,0);
	\draw[color ={black},>=latex,->] (0,-8)--(0,6);
	\draw[color ={black},opacity = 0.4] (-4,1)--(4,1);
	\draw[color ={black},opacity = 0.4] (-4,2)--(4,2);
	\draw[color ={black},opacity = 0.4] (-4,3)--(4,3);
	\draw[color ={black},opacity = 0.4] (-4,4)--(4,4);
	\draw[color ={black},opacity = 0.4] (-4,5)--(4,5);
	\draw[color ={black},opacity = 0.4] (-4,6)--(4,6);
	\draw[color ={black},opacity = 0.4] (-4,-1)--(4,-1);
	\draw[color ={black},opacity = 0.4] (-4,-2)--(4,-2);
	\draw[color ={black},opacity = 0.4] (-4,-3)--(4,-3);
	\draw[color ={black},opacity = 0.4] (-4,-4)--(4,-4);
	\draw[color ={black},opacity = 0.4] (-4,-5)--(4,-5);
	\draw[color ={black},opacity = 0.4] (-4,-6)--(4,-6);
	\draw[color ={black},opacity = 0.4] (-4,-7)--(4,-7);
	\draw[color ={black},opacity = 0.4] (-4,-8)--(4,-8);
	\draw[color ={black},opacity = 0.4] (1,-8)--(1,6);
	\draw[color ={black},opacity = 0.4] (2,-8)--(2,6);
	\draw[color ={black},opacity = 0.4] (3,-8)--(3,6);
	\draw[color ={black},opacity = 0.4] (4,-8)--(4,6);
	\draw[color ={black},opacity = 0.4] (-1,-8)--(-1,6);
	\draw[color ={black},opacity = 0.4] (-2,-8)--(-2,6);
	\draw[color ={black},opacity = 0.4] (-3,-8)--(-3,6);
	\draw[color ={black},opacity = 0.4] (-4,-8)--(-4,6);
	\draw[color ={gray},opacity = 0.3] (-4,0)--(4,0);
	\draw[color ={gray},opacity = 0.3] (-4,0.5)--(4,0.5);
	\draw[color ={gray},opacity = 0.3] (-4,1.5)--(4,1.5);
	\draw[color ={gray},opacity = 0.3] (-4,2.5)--(4,2.5);
	\draw[color ={gray},opacity = 0.3] (-4,3.5)--(4,3.5);
	\draw[color ={gray},opacity = 0.3] (-4,4.5)--(4,4.5);
	\draw[color ={gray},opacity = 0.3] (-4,5.5)--(4,5.5);
	\draw[color ={gray},opacity = 0.3] (-4,0)--(4,0);
	\draw[color ={gray},opacity = 0.3] (-4,-0.5)--(4,-0.5);
	\draw[color ={gray},opacity = 0.3] (-4,-1.5)--(4,-1.5);
	\draw[color ={gray},opacity = 0.3] (-4,-2.5)--(4,-2.5);
	\draw[color ={gray},opacity = 0.3] (-4,-3.5)--(4,-3.5);
	\draw[color ={gray},opacity = 0.3] (-4,-4.5)--(4,-4.5);
	\draw[color ={gray},opacity = 0.3] (-4,-5.5)--(4,-5.5);
	\draw[color ={gray},opacity = 0.3] (-4,-6.5)--(4,-6.5);
	\draw[color ={gray},opacity = 0.3] (-4,-7)--(4,-7);
	\draw[color ={gray},opacity = 0.3] (-4,-7.5)--(4,-7.5);
	\draw[color ={gray},opacity = 0.3] (-4,-8)--(4,-8);
	\draw[color ={gray},opacity = 0.3] (0,-8)--(0,6);
	\draw[color ={gray},opacity = 0.3] (0,-8)--(0,6);
	\draw[color ={black},line width = 1.2] (1,-0.13)--(1,0.13);
	\draw[color ={black},line width = 1.2] (2,-0.13)--(2,0.13);
	\draw[color ={black},line width = 1.2] (3,-0.13)--(3,0.13);
	\draw[color ={black},line width = 1.2] (4,-0.13)--(4,0.13);
	\draw[color ={black},line width = 1.2] (-1,-0.13)--(-1,0.13);
	\draw[color ={black},line width = 1.2] (-2,-0.13)--(-2,0.13);
	\draw[color ={black},line width = 1.2] (-3,-0.13)--(-3,0.13);
	\draw[color ={black},line width = 1.2] (-4,-0.13)--(-4,0.13);
	\draw[color ={black},line width = 1.2] (-0.13,1)--(0.13,1);
	\draw[color ={black},line width = 1.2] (-0.13,2)--(0.13,2);
	\draw[color ={black},line width = 1.2] (-0.13,3)--(0.13,3);
	\draw[color ={black},line width = 1.2] (-0.13,4)--(0.13,4);
	\draw[color ={black},line width = 1.2] (-0.13,5)--(0.13,5);
	\draw[color ={black},line width = 1.2] (-0.13,6)--(0.13,6);
	\draw[color ={black},line width = 1.2] (-0.13,-1)--(0.13,-1);
	\draw[color ={black},line width = 1.2] (-0.13,-2)--(0.13,-2);
	\draw[color ={black},line width = 1.2] (-0.13,-3)--(0.13,-3);
	\draw[color ={black},line width = 1.2] (-0.13,-4)--(0.13,-4);
	\draw[color ={black},line width = 1.2] (-0.13,-5)--(0.13,-5);
	\draw[color ={black},line width = 1.2] (-0.13,-6)--(0.13,-6);
	\draw[color ={black},line width = 1.2] (-0.13,-7)--(0.13,-7);
	\draw[color ={black},line width = 1.2] (-0.13,-8)--(0.13,-8);
	\draw (2,-0.4) node[anchor = center] {\scriptsize \color{black}{$2$}};
	\draw (4,-0.4) node[anchor = center] {\scriptsize \color{black}{$4$}};
	\draw (-2,-0.4) node[anchor = center] {\scriptsize \color{black}{$-2$}};
	\draw (-4,-0.4) node[anchor = center] {\scriptsize \color{black}{$-4$}};
	\draw (-0.5,2.1) node[anchor = center] {\footnotesize \color{black}{$2$}};
	\draw (-0.5,4.1) node[anchor = center] {\footnotesize \color{black}{$4$}};
	\draw (-0.5,-1.9) node[anchor = center] {\footnotesize \color{black}{$-2$}};
	\draw (-0.5,-3.9) node[anchor = center] {\footnotesize \color{black}{$-4$}};
	\draw (-0.5,-5.9) node[anchor = center] {\footnotesize \color{black}{$-6$}};
	\draw (-0.5,-7.9) node[anchor = center] {\footnotesize \color{black}{$-8$}};
	\draw [color={black}] (-0.3,-0.3) node[anchor = center,scale=1, rotate = 0] {O};
	\draw (3,4) node[anchor = center] {\footnotesize \color{red}{$\Large\mathcal{C}_f\prime$}};
	
	\draw[color={red},line width = 2] (-4,4)--(-3.8,3.6)--(-3.6,3.2)--(-3.4,2.8)--(-3.2,2.4)--(-3,2)--(-2.8,1.6)--(-2.6,1.2)--(-2.4,0.8)--(-2.2,0.4)--(-2,0)--(-1.8,-0.4)--(-1.6,-0.8)--(-1.4,-1.2)--(-1.2,-1.6)--(-1,-2)--(-0.8,-2.4)--(-0.6,-2.8)--(-0.4,-3.2)--(-0.2,-3.6)--(0,-4)--(0.2,-4.4)--(0.4,-4.8)--(0.6,-5.2)--(0.8,-5.6)--(1,-6)--(1.2,-6.4)--(1.4,-6.8)--(1.6,-7.2)--(1.8,-7.6)--(2,-8)--(2.2,-8.4)--(2.4,-8.8);

\end{tikzpicture}\\
    \end{minipage}
    
\end{exercice}

\begin{Solution}
 L'équation réduite de la tangente au point d'abscisse $0$ est  : $y=f'(0)(x-0)+f(0)$.\\
      On lit graphiquement $f(0)=-1$ et $f'(0)=-4$.\\
      L'équation réduite de la tangente est donc donnée par :
      $y=-4(x+0)-1$, soit $y=-4x-1$.
\end{Solution}

% @see : https://coopmaths.fr/alea?uuid=a1ba2&id=can1F14&n=1&d=10&alea=rkIp22323232323423456789101112131415161718192021222324252627282930313233343536373839404142434445464748234567891011121314151617181920212223242526272829303132333435&cd=1&cols=1
\needspace{10\baselineskip}
\begin{exercice}[BaremeTotal=false]


Soit $f$ la fonction définie sur $\mathbb{R}$ par : 
$f(x)= -8x^2+10$.\\
En calculant la limite du taux d'accroissement, déterminer $f'(5)$.
\end{exercice}

\begin{Solution}
 $f$ est une fonction polynôme du second degré de la forme $f(x)=ax^2+b$.\\
    La fonction dérivée est donnée par la somme des dérivées des fonctions $u$ et $v$ définies par $u(x)=-8x^2$ et $v(x)=10$.\\
     Comme $u'(x)=-16x$ et $v'(x)=0$, on obtient  $f'(x)=-16x$. \\
     Ainsi, $f'(5)=-16\times 5=-80$.
\end{Solution}

% @see : https://coopmaths.fr/alea?uuid=3c690&id=can1F13&n=1&d=10&alea=EUDu22323232323423456789101112131415161718192021222324252627282930313233343536373839404142434445464748234567891011121314151617181920212223242526272829303132333435&cd=1&cols=1
\needspace{10\baselineskip}
\begin{exercice}[BaremeTotal=false]


 Déterminer le coefficient directeur de la tangente à la courbe représentative de la fonction racine carrée au point d'abscisse $12$.
         
\end{exercice}

\begin{Solution}
 Le coefficient directeur de la tangente au point d'abscisse $a$ est donné par le nombre dérivé $f'(a)$.\\
        La fonction $f$ définie par $f(x)=\sqrt{x}$ a pour fonction dérivée la fonction $f'$ définie par $f'(x)=\dfrac{1}{2\sqrt{x}}$.\\Comme $f'(12)=\dfrac{1}{2\sqrt{12}}$, le coefficient directeur de la tangente au point d'abscisse $12$ est : $\dfrac{1}{2\sqrt{12}}$.
\end{Solution}

% @see : https://coopmaths.fr/alea?uuid=ebd89&id=1AN14-1&n=1&d=10&s=3&alea=ocYr22323232323423456789101112131415161718192021222324252627282930313233343536373839404142434445464748234567891011121314151617181920212223242526272829303132333435&cd=0&cols=1
\needspace{10\baselineskip}
\begin{exercice}[BaremeTotal=false]


 Donner la dérivée de la fonction $f$, dérivable pour tout $x\in\R$, définie par  $f(x)=\dfrac{-4x+8}{7}$.
\end{exercice}

\begin{Solution}
 L'expression de la dérivée de la fonction $f$ définie par $f(x)=\dfrac{-4x+8}{7}$ est : ${\color[HTML]{f15929}\boldsymbol{f'(x)=-\dfrac{4}{7}}}$.
\end{Solution}

% @see : https://coopmaths.fr/alea?uuid=ebd89&id=1AN14-1&n=1&d=10&s=4&alea=uAlP22323232323423456789101112131415161718192021222324252627282930313233343536373839404142434445464748234567891011121314151617181920212223242526272829303132333435&cd=0&cols=1
\needspace{10\baselineskip}
\begin{exercice}[BaremeTotal=false]


 Donner la dérivée de la fonction $f$, dérivable pour tout $x\in\R$, définie par  $f(x)=-8x^{15}$.
\end{exercice}

\begin{Solution}
 L'expression de la dérivée de la fonction $f$ définie par $f(x)=-8x^{15}$ est : ${\color[HTML]{f15929}\boldsymbol{f'(x)=-120x^{14}}}$.
\end{Solution}

% @see : https://coopmaths.fr/alea?uuid=ebd89&id=1AN14-1&n=1&d=10&s=7&alea=vjN722323232323423456789101112131415161718192021222324252627282930313233343536373839404142434445464748234567891011121314151617181920212223242526272829303132333435&cd=0&cols=1
\needspace{10\baselineskip}
\begin{exercice}[BaremeTotal=false]


 Donner la dérivée de la fonction $f$, dérivable pour tout $x\in\R^*$, définie par  $f(x)=\dfrac{-4}{9x}$.
\end{exercice}

\begin{Solution}
 L'expression de la dérivée de la fonction $f$ définie par $f(x)=\dfrac{-4}{9x}$ est : ${\color[HTML]{f15929}\boldsymbol{f'(x)=\dfrac{4}{9x^2}}}$.
\end{Solution}

% @see : https://coopmaths.fr/alea?uuid=a83c0&id=1AN14-4&n=1&d=10&s=1&alea=Njr322323232323423456789101112131415161718192021222324252627282930313233343536373839404142434445464748234567891011121314151617181920212223242526272829303132333435&cd=0&cols=1
\needspace{10\baselineskip}
\begin{exercice}[BaremeTotal=false]


 Donner l'expression de la dérivée de la fonction $f$ définie sur $\R^{*}$ par $f(x)=\dfrac{3}{x}+4x^{2}$.\\
\end{exercice}

\begin{Solution}
 L'expression de la dérivée de la fonction $f$ définie par $f(x)=\dfrac{3}{x}+4x^{2}$ est :\\$f'(x)={\color[HTML]{f15929}\boldsymbol{-\dfrac{3}{x^2}+8x}}$.\\
\end{Solution}

% @see : https://coopmaths.fr/alea?uuid=a83c0&id=1AN14-4&n=1&d=10&s=4&alea=4Vpz22323232323423456789101112131415161718192021222324252627282930313233343536373839404142434445464748234567891011121314151617181920212223242526272829303132333435&cd=1&cols=1
\needspace{10\baselineskip}
\begin{exercice}[BaremeTotal=false]


 Donner l'expression de la dérivée de la fonction $f$ définie sur $\R_+^{*}$ par $f(x)=-2x^{2}-2x-2\sqrt{x}-\dfrac{3}{x}$.\\
\end{exercice}

\begin{Solution}
 $(-2x^{2}-2x)^\prime=-4x-2$\\$(-2\sqrt{x})^\prime=-\dfrac{2}{2\sqrt{x}}$\\$(-\dfrac{3}{x})^\prime=\dfrac{3}{x^2}$\\L'expression de la dérivée de la fonction $f$ définie par $f(x)=-2x^{2}-2x-2\sqrt{x}-\dfrac{3}{x}$ est :\\$f'(x)={\color[HTML]{f15929}\boldsymbol{-4x-2-\dfrac{2}{2\sqrt{x}}+\dfrac{3}{x^2}}}$.\\
\end{Solution}

\end{Maquette}
\clearpage
\begin{Maquette}[DM]{Niveau= ,Classe=\getdata[36]\mydata\ (v36), Date = 06/01/25  ,Theme=Exercices}


% @see : https://coopmaths.fr/alea?uuid=4c8c7&id=1AN11-3&n=1&d=10&s=2&alea=nawO2232323232342345678910111213141516171819202122232425262728293031323334353637383940414243444546474823456789101112131415161718192021222324252627282930313233343536&cd=1&cols=2
\needspace{10\baselineskip}
\begin{exercice}[BaremeTotal=false]
\label{eleve:36}


 \begin{minipage}[t]{\linewidth} Soit $f$ une fonction dérivable sur $[-5;5]$ et $\mathcal{C}_f$ sa courbe représentative.\\On sait que  $f(5)=-2~~$ et que $~~f'(5)=-5$.\\Déterminer une équation de la tangente $(T)$ à la courbe $\mathcal{C}_f$ au point d'abscisse $5$,\\sans utiliser la formule de cours de l'équation de tangente. \end{minipage}

\end{exercice}

\begin{Solution}
 \begin{minipage}[t]{\linewidth} On sait que la tangente n'est pas une droite verticale, puisque la fonction est dérivable sur l'intervalle.\\On en déduit que la tangente $(T)$ au point d'abscisse $5$, admet une équation réduite de la forme :  $(T) : y=m x + p$. 

\medskip
$\bullet$ Détermination de $m$ :\\On sait que le nombre dérivé en $5$ est par définition, le coefficient directeur de la tangente au point d'abscisse $5$.\\Par conséquent, on a déjà : $m=f'(5)=-5$.\\ On en déduit que  $(T) : y= -5 x + p$\\$\bullet$ Détermination de $p$ :\\ Pour cela, on utilise que si $f(5)=-2~~$, alors le point $A$ de coordonnées $(5;-2)$ appartient à $\mathcal{C}_f$ mais aussi à $(T)$.\\ On peut écrire $A(5;-2) = \mathcal{C}_f \cap (T)$.\\ On remplace alors les coordonnées de $A(5;-2)$ dans l'équation  $(T) : y= -5 x + p$.\\ $\begin{aligned} \phantom{\iff}&A(5;-2)\in (T)\\ \iff& -2= -5 \times 5  + p\\ \iff& p=-2 +5 \times 5  \\ \iff& p=-2 +25   \\ \iff& p=23\\\end{aligned}$\\On peut conclure que : $(T) : {\color[HTML]{f15929}\boldsymbol{y=-5x+23}}$. \end{minipage}

\end{Solution}

% @see : https://coopmaths.fr/alea?uuid=0e984&id=can1F15&n=1&d=10&alea=1Alo2232323232342345678910111213141516171819202122232425262728293031323334353637383940414243444546474823456789101112131415161718192021222324252627282930313233343536&cd=1&cols=1
\needspace{10\baselineskip}
\begin{exercice}[BaremeTotal=false]


 La courbe représente une fonction $f$ et la droite est la tangente au point d'abscisse $0$.\\
        Déterminer $f'(0)$.\begin{center}\begin{tikzpicture}[baseline,scale = 0.5]

    \tikzset{
      point/.style={
        thick,
        draw,
        cross out,
        inner sep=0pt,
        minimum width=5pt,
        minimum height=5pt,
      },
    }
    \clip (-10,-2) rectangle (4,12);
    	
	\draw[color ={black},line width = 2,>=latex,->] (-10,0)--(4,0);
	\draw[color ={black},line width = 2,>=latex,->] (0,-2)--(0,12);
	\draw[color ={black},opacity = 0.5] (-10,1)--(4,1);
	\draw[color ={black},opacity = 0.5] (-10,2)--(4,2);
	\draw[color ={black},opacity = 0.5] (-10,3)--(4,3);
	\draw[color ={black},opacity = 0.5] (-10,4)--(4,4);
	\draw[color ={black},opacity = 0.5] (-10,5)--(4,5);
	\draw[color ={black},opacity = 0.5] (-10,6)--(4,6);
	\draw[color ={black},opacity = 0.5] (-10,7)--(4,7);
	\draw[color ={black},opacity = 0.5] (-10,8)--(4,8);
	\draw[color ={black},opacity = 0.5] (-10,9)--(4,9);
	\draw[color ={black},opacity = 0.5] (-10,10)--(4,10);
	\draw[color ={black},opacity = 0.5] (-10,11)--(4,11);
	\draw[color ={black},opacity = 0.5] (-10,12)--(4,12);
	\draw[color ={black},opacity = 0.5] (-10,-1)--(4,-1);
	\draw[color ={black},opacity = 0.5] (-10,-2)--(4,-2);
	\draw[color ={black},opacity = 0.5] (1,-2)--(1,12);
	\draw[color ={black},opacity = 0.5] (2,-2)--(2,12);
	\draw[color ={black},opacity = 0.5] (3,-2)--(3,12);
	\draw[color ={black},opacity = 0.5] (4,-2)--(4,12);
	\draw[color ={black},opacity = 0.5] (-1,-2)--(-1,12);
	\draw[color ={black},opacity = 0.5] (-2,-2)--(-2,12);
	\draw[color ={black},opacity = 0.5] (-3,-2)--(-3,12);
	\draw[color ={black},opacity = 0.5] (-4,-2)--(-4,12);
	\draw[color ={black},opacity = 0.5] (-5,-2)--(-5,12);
	\draw[color ={black},opacity = 0.5] (-6,-2)--(-6,12);
	\draw[color ={black},opacity = 0.5] (-7,-2)--(-7,12);
	\draw[color ={black},opacity = 0.5] (-8,-2)--(-8,12);
	\draw[color ={black},opacity = 0.5] (-9,-2)--(-9,12);
	\draw[color ={black},opacity = 0.5] (-10,-2)--(-10,12);
	\draw[color ={gray},opacity = 0.3] (-10,0)--(4,0);
	\draw[color ={gray},opacity = 0.3] (-10,0)--(4,0);
	\draw[color ={gray},opacity = 0.3] (0,-2)--(0,12);
	\draw[color ={gray},opacity = 0.3] (0,-2)--(0,12);
	\draw[color ={gray},opacity = 0.3] (-5,-2)--(-5,12);
	\draw[color ={gray},opacity = 0.3] (-6,-2)--(-6,12);
	\draw[color ={gray},opacity = 0.3] (-7,-2)--(-7,12);
	\draw[color ={gray},opacity = 0.3] (-8,-2)--(-8,12);
	\draw[color ={gray},opacity = 0.3] (-9,-2)--(-9,12);
	\draw[color ={gray},opacity = 0.3] (-10,-2)--(-10,12);
	\draw[color ={black},line width = 2] (2,-0.2)--(2,0.2);
	\draw[color ={black},line width = 2] (-2,-0.2)--(-2,0.2);
	\draw[color ={black},line width = 2] (-4,-0.2)--(-4,0.2);
	\draw[color ={black},line width = 2] (-6,-0.2)--(-6,0.2);
	\draw[color ={black},line width = 2] (-8,-0.2)--(-8,0.2);
	\draw[color ={black},line width = 2] (-0.2,2)--(0.2,2);
	\draw[color ={black},line width = 2] (-0.2,4)--(0.2,4);
	\draw[color ={black},line width = 2] (-0.2,6)--(0.2,6);
	\draw[color ={black},line width = 2] (-0.2,8)--(0.2,8);
	\draw[color ={black},line width = 2] (-0.2,10)--(0.2,10);
	\draw (2,-0.4) node[anchor = center] {\scriptsize \color{black}{$1$}};
	\draw (-2,-0.4) node[anchor = center] {\scriptsize \color{black}{$-1$}};
	\draw (-4,-0.4) node[anchor = center] {\scriptsize \color{black}{$-2$}};
	\draw (-6,-0.4) node[anchor = center] {\scriptsize \color{black}{$-3$}};
	\draw (-8,-0.4) node[anchor = center] {\scriptsize \color{black}{$-4$}};
	\draw (-0.5,2.1) node[anchor = center] {\footnotesize \color{black}{$1$}};
	\draw (-0.5,4.1) node[anchor = center] {\footnotesize \color{black}{$2$}};
	\draw (-0.5,6.1) node[anchor = center] {\footnotesize \color{black}{$3$}};
	\draw (-0.5,8.1) node[anchor = center] {\footnotesize \color{black}{$4$}};
	\draw (-0.5,10.1) node[anchor = center] {\footnotesize \color{black}{$5$}};
	\draw [color={black}] (-0.3,-0.3) node[anchor = center,scale=1, rotate = 0] {O};
	
	\draw[color={blue},line width = 2] (-5.6,12.93)--(-5.2,11.32)--(-4.8,9.91)--(-4.4,8.67)--(-4,7.59)--(-3.6,6.64)--(-3.2,5.81)--(-2.8,5.09)--(-2.4,4.45)--(-2,3.9)--(-1.6,3.41)--(-1.2,2.98)--(-0.8,2.61)--(-0.4,2.29)--(0,2)--(0.4,1.75)--(0.8,1.53)--(1.2,1.34)--(1.6,1.17)--(2,1.03)--(2.4,0.9)--(2.8,0.79)--(3.2,0.69)--(3.6,0.6)--(4,0.53);
	
	\draw[color={red},line width = 2] (-10,8.67)--(-9.6,8.4)--(-9.2,8.13)--(-8.8,7.87)--(-8.4,7.6)--(-8,7.33)--(-7.6,7.07)--(-7.2,6.8)--(-6.8,6.53)--(-6.4,6.27)--(-6,6)--(-5.6,5.73)--(-5.2,5.47)--(-4.8,5.2)--(-4.4,4.93)--(-4,4.67)--(-3.6,4.4)--(-3.2,4.13)--(-2.8,3.87)--(-2.4,3.6)--(-2,3.33)--(-1.6,3.07)--(-1.2,2.8)--(-0.8,2.53)--(-0.4,2.27)--(0,2)--(0.4,1.73)--(0.8,1.47)--(1.2,1.2)--(1.6,0.93)--(2,0.67)--(2.4,0.4)--(2.8,0.13)--(3.2,-0.13)--(3.6,-0.4)--(4,-0.67);

\end{tikzpicture}\end{center}
\end{exercice}

\begin{Solution}
 $f'(0)$ est donné par le coefficient directeur de la tangente à la courbe au point d'abscisse $0$, soit $\dfrac{-2}{3}=-\dfrac{2}{3}$.
\end{Solution}

% @see : https://coopmaths.fr/alea?uuid=6f32d&id=can1F16&n=1&d=10&alea=dqSw2232323232342345678910111213141516171819202122232425262728293031323334353637383940414243444546474823456789101112131415161718192021222324252627282930313233343536&cd=1&cols=1
\needspace{10\baselineskip}
\newpage
\begin{exercice}[BaremeTotal=false]


 On donne les représentations graphiques d'une fonction et de sa dérivée.\\
        Donner l'équation réduite de la tangente à la courbe de $f$ en $x=2$. \\ \begin{minipage}{0.5\linewidth}
    \begin{tikzpicture}[baseline, scale=0.8]

    \tikzset{
      point/.style={
        thick,
        draw,
        cross out,
        inner sep=0pt,
        minimum width=5pt,
        minimum height=5pt,
      },
    }
    \clip (-4.5,-4.5) rectangle (5.75,13.5);
    	
	\draw[color ={black},>=latex,->] (-3,0)--(3,0);
	\draw[color ={black},>=latex,->] (0,-3)--(0,12);
	\draw[color ={black},opacity = 0.4] (-3,1)--(3,1);
	\draw[color ={black},opacity = 0.4] (-3,2)--(3,2);
	\draw[color ={black},opacity = 0.4] (-3,3)--(3,3);
	\draw[color ={black},opacity = 0.4] (-3,4)--(3,4);
	\draw[color ={black},opacity = 0.4] (-3,5)--(3,5);
	\draw[color ={black},opacity = 0.4] (-3,6)--(3,6);
	\draw[color ={black},opacity = 0.4] (-3,7)--(3,7);
	\draw[color ={black},opacity = 0.4] (-3,8)--(3,8);
	\draw[color ={black},opacity = 0.4] (-3,9)--(3,9);
	\draw[color ={black},opacity = 0.4] (-3,10)--(3,10);
	\draw[color ={black},opacity = 0.4] (-3,11)--(3,11);
	\draw[color ={black},opacity = 0.4] (-3,12)--(3,12);
	\draw[color ={black},opacity = 0.4] (-3,-1)--(3,-1);
	\draw[color ={black},opacity = 0.4] (-3,-2)--(3,-2);
	\draw[color ={black},opacity = 0.4] (-3,-3)--(3,-3);
	\draw[color ={black},opacity = 0.4] (1,-3)--(1,12);
	\draw[color ={black},opacity = 0.4] (2,-3)--(2,12);
	\draw[color ={black},opacity = 0.4] (3,-3)--(3,12);
	\draw[color ={black},opacity = 0.4] (-1,-3)--(-1,12);
	\draw[color ={black},opacity = 0.4] (-2,-3)--(-2,12);
	\draw[color ={black},opacity = 0.4] (-3,-3)--(-3,12);
	\draw[color ={gray},opacity = 0.3] (-3,0)--(3,0);
	\draw[color ={gray},opacity = 0.3] (-3,0.5)--(3,0.5);
	\draw[color ={gray},opacity = 0.3] (-3,1.5)--(3,1.5);
	\draw[color ={gray},opacity = 0.3] (-3,2.5)--(3,2.5);
	\draw[color ={gray},opacity = 0.3] (-3,3.5)--(3,3.5);
	\draw[color ={gray},opacity = 0.3] (-3,4.5)--(3,4.5);
	\draw[color ={gray},opacity = 0.3] (-3,5.5)--(3,5.5);
	\draw[color ={gray},opacity = 0.3] (-3,6.5)--(3,6.5);
	\draw[color ={gray},opacity = 0.3] (-3,7.5)--(3,7.5);
	\draw[color ={gray},opacity = 0.3] (-3,8.5)--(3,8.5);
	\draw[color ={gray},opacity = 0.3] (-3,9.5)--(3,9.5);
	\draw[color ={gray},opacity = 0.3] (-3,10.5)--(3,10.5);
	\draw[color ={gray},opacity = 0.3] (-3,11.5)--(3,11.5);
	\draw[color ={gray},opacity = 0.3] (-3,0)--(3,0);
	\draw[color ={gray},opacity = 0.3] (-3,-0.5)--(3,-0.5);
	\draw[color ={gray},opacity = 0.3] (-3,-1.5)--(3,-1.5);
	\draw[color ={gray},opacity = 0.3] (-3,-2.5)--(3,-2.5);
	\draw[color ={gray},opacity = 0.3] (0,-3)--(0,12);
	\draw[color ={gray},opacity = 0.3] (0,-3)--(0,12);
	\draw[color ={black},line width = 1.2] (1,-0.13)--(1,0.13);
	\draw[color ={black},line width = 1.2] (2,-0.13)--(2,0.13);
	\draw[color ={black},line width = 1.2] (3,-0.13)--(3,0.13);
	\draw[color ={black},line width = 1.2] (-1,-0.13)--(-1,0.13);
	\draw[color ={black},line width = 1.2] (-2,-0.13)--(-2,0.13);
	\draw[color ={black},line width = 1.2] (-3,-0.13)--(-3,0.13);
	\draw[color ={black},line width = 1.2] (-0.13,1)--(0.13,1);
	\draw[color ={black},line width = 1.2] (-0.13,2)--(0.13,2);
	\draw[color ={black},line width = 1.2] (-0.13,3)--(0.13,3);
	\draw[color ={black},line width = 1.2] (-0.13,4)--(0.13,4);
	\draw[color ={black},line width = 1.2] (-0.13,5)--(0.13,5);
	\draw[color ={black},line width = 1.2] (-0.13,6)--(0.13,6);
	\draw[color ={black},line width = 1.2] (-0.13,7)--(0.13,7);
	\draw[color ={black},line width = 1.2] (-0.13,8)--(0.13,8);
	\draw[color ={black},line width = 1.2] (-0.13,9)--(0.13,9);
	\draw[color ={black},line width = 1.2] (-0.13,10)--(0.13,10);
	\draw[color ={black},line width = 1.2] (-0.13,11)--(0.13,11);
	\draw[color ={black},line width = 1.2] (-0.13,12)--(0.13,12);
	\draw[color ={black},line width = 1.2] (-0.13,-1)--(0.13,-1);
	\draw[color ={black},line width = 1.2] (-0.13,-2)--(0.13,-2);
	\draw[color ={black},line width = 1.2] (-0.13,-3)--(0.13,-3);
	\draw (3,-0.4) node[anchor = center] {\scriptsize \color{black}{$3$}};
	\draw (-3,-0.4) node[anchor = center] {\scriptsize \color{black}{$-3$}};
	\draw (-0.5,3.1) node[anchor = center] {\footnotesize \color{black}{$3$}};
	\draw (-0.5,6.1) node[anchor = center] {\footnotesize \color{black}{$6$}};
	\draw (-0.5,9.1) node[anchor = center] {\footnotesize \color{black}{$9$}};
	\draw (-0.5,-2.9) node[anchor = center] {\footnotesize \color{black}{$-3$}};
	\draw [color={black}] (-0.3,-0.3) node[anchor = center,scale=1, rotate = 0] {O};
	\draw (3,10) node[anchor = center] {\footnotesize \color{blue}{$\Large \mathcal{C}_f$}};
	
	\draw[color={blue},line width = 2] (-2.6,12.96)--(-2.4,11.56)--(-2.2,10.24)--(-2,9)--(-1.8,7.84)--(-1.6,6.76)--(-1.4,5.76)--(-1.2,4.84)--(-1,4)--(-0.8,3.24)--(-0.6,2.56)--(-0.4,1.96)--(-0.2,1.44)--(0,1)--(0.2,0.64)--(0.4,0.36)--(0.6,0.16)--(0.8,0.04)--(1,0)--(1.2,0.04)--(1.4,0.16)--(1.6,0.36)--(1.8,0.64)--(2,1)--(2.2,1.44)--(2.4,1.96)--(2.6,2.56)--(2.8,3.24)--(3,4);

\end{tikzpicture}\\
    \end{minipage}
    \begin{minipage}{0.5\linewidth}
    \begin{tikzpicture}[baseline, scale=0.8]

    \tikzset{
      point/.style={
        thick,
        draw,
        cross out,
        inner sep=0pt,
        minimum width=5pt,
        minimum height=5pt,
      },
    }
    \clip (-4.5,-6.5) rectangle (6.5,9.5);
    	
	\draw[color ={black},>=latex,->] (-3,0)--(3,0);
	\draw[color ={black},>=latex,->] (0,-5)--(0,8);
	\draw[color ={black},opacity = 0.4] (-3,1)--(3,1);
	\draw[color ={black},opacity = 0.4] (-3,2)--(3,2);
	\draw[color ={black},opacity = 0.4] (-3,3)--(3,3);
	\draw[color ={black},opacity = 0.4] (-3,4)--(3,4);
	\draw[color ={black},opacity = 0.4] (-3,5)--(3,5);
	\draw[color ={black},opacity = 0.4] (-3,6)--(3,6);
	\draw[color ={black},opacity = 0.4] (-3,7)--(3,7);
	\draw[color ={black},opacity = 0.4] (-3,8)--(3,8);
	\draw[color ={black},opacity = 0.4] (-3,-1)--(3,-1);
	\draw[color ={black},opacity = 0.4] (-3,-2)--(3,-2);
	\draw[color ={black},opacity = 0.4] (-3,-3)--(3,-3);
	\draw[color ={black},opacity = 0.4] (-3,-4)--(3,-4);
	\draw[color ={black},opacity = 0.4] (-3,-5)--(3,-5);
	\draw[color ={black},opacity = 0.4] (1,-5)--(1,8);
	\draw[color ={black},opacity = 0.4] (2,-5)--(2,8);
	\draw[color ={black},opacity = 0.4] (3,-5)--(3,8);
	\draw[color ={black},opacity = 0.4] (-1,-5)--(-1,8);
	\draw[color ={black},opacity = 0.4] (-2,-5)--(-2,8);
	\draw[color ={black},opacity = 0.4] (-3,-5)--(-3,8);
	\draw[color ={gray},opacity = 0.3] (-3,0)--(3,0);
	\draw[color ={gray},opacity = 0.3] (-3,0.5)--(3,0.5);
	\draw[color ={gray},opacity = 0.3] (-3,1.5)--(3,1.5);
	\draw[color ={gray},opacity = 0.3] (-3,2.5)--(3,2.5);
	\draw[color ={gray},opacity = 0.3] (-3,3.5)--(3,3.5);
	\draw[color ={gray},opacity = 0.3] (-3,4.5)--(3,4.5);
	\draw[color ={gray},opacity = 0.3] (-3,5.5)--(3,5.5);
	\draw[color ={gray},opacity = 0.3] (-3,6.5)--(3,6.5);
	\draw[color ={gray},opacity = 0.3] (-3,7.5)--(3,7.5);
	\draw[color ={gray},opacity = 0.3] (-3,0)--(3,0);
	\draw[color ={gray},opacity = 0.3] (-3,-0.5)--(3,-0.5);
	\draw[color ={gray},opacity = 0.3] (-3,-1.5)--(3,-1.5);
	\draw[color ={gray},opacity = 0.3] (-3,-2.5)--(3,-2.5);
	\draw[color ={gray},opacity = 0.3] (-3,-3.5)--(3,-3.5);
	\draw[color ={gray},opacity = 0.3] (-3,-4.5)--(3,-4.5);
	\draw[color ={gray},opacity = 0.3] (0,-5)--(0,8);
	\draw[color ={gray},opacity = 0.3] (0,-5)--(0,8);
	\draw[color ={black},line width = 1.2] (1,-0.13)--(1,0.13);
	\draw[color ={black},line width = 1.2] (2,-0.13)--(2,0.13);
	\draw[color ={black},line width = 1.2] (3,-0.13)--(3,0.13);
	\draw[color ={black},line width = 1.2] (-1,-0.13)--(-1,0.13);
	\draw[color ={black},line width = 1.2] (-2,-0.13)--(-2,0.13);
	\draw[color ={black},line width = 1.2] (-3,-0.13)--(-3,0.13);
	\draw[color ={black},line width = 1.2] (-0.13,1)--(0.13,1);
	\draw[color ={black},line width = 1.2] (-0.13,2)--(0.13,2);
	\draw[color ={black},line width = 1.2] (-0.13,3)--(0.13,3);
	\draw[color ={black},line width = 1.2] (-0.13,4)--(0.13,4);
	\draw[color ={black},line width = 1.2] (-0.13,5)--(0.13,5);
	\draw[color ={black},line width = 1.2] (-0.13,6)--(0.13,6);
	\draw[color ={black},line width = 1.2] (-0.13,7)--(0.13,7);
	\draw[color ={black},line width = 1.2] (-0.13,8)--(0.13,8);
	\draw[color ={black},line width = 1.2] (-0.13,9)--(0.13,9);
	\draw[color ={black},line width = 1.2] (-0.13,10)--(0.13,10);
	\draw[color ={black},line width = 1.2] (-0.13,11)--(0.13,11);
	\draw[color ={black},line width = 1.2] (-0.13,12)--(0.13,12);
	\draw[color ={black},line width = 1.2] (-0.13,-1)--(0.13,-1);
	\draw[color ={black},line width = 1.2] (-0.13,-2)--(0.13,-2);
	\draw[color ={black},line width = 1.2] (-0.13,-3)--(0.13,-3);
	\draw (3,-0.4) node[anchor = center] {\scriptsize \color{black}{$3$}};
	\draw (-3,-0.4) node[anchor = center] {\scriptsize \color{black}{$-3$}};
	\draw (-0.5,3.1) node[anchor = center] {\footnotesize \color{black}{$3$}};
	\draw (-0.5,6.1) node[anchor = center] {\footnotesize \color{black}{$6$}};
	\draw (-0.5,-2.9) node[anchor = center] {\footnotesize \color{black}{$-3$}};
	\draw [color={black}] (-0.3,-0.3) node[anchor = center,scale=1, rotate = 0] {O};
	\draw (3,6) node[anchor = center] {\footnotesize \color{red}{$\Large\mathcal{C}_f\prime$}};
	
	\draw[color={red},line width = 2] (-2,-6)--(-1.8,-5.6)--(-1.6,-5.2)--(-1.4,-4.8)--(-1.2,-4.4)--(-1,-4)--(-0.8,-3.6)--(-0.6,-3.2)--(-0.4,-2.8)--(-0.2,-2.4)--(0,-2)--(0.2,-1.6)--(0.4,-1.2)--(0.6,-0.8)--(0.8,-0.4)--(1,0)--(1.2,0.4)--(1.4,0.8)--(1.6,1.2)--(1.8,1.6)--(2,2)--(2.2,2.4)--(2.4,2.8)--(2.6,3.2)--(2.8,3.6)--(3,4);

\end{tikzpicture}\\
    \end{minipage}
    
\end{exercice}

\begin{Solution}
 L'équation réduite de la tangente au point d'abscisse $2$ est  : $y=f'(2)(x-2)+f(2)$.\\
        On lit graphiquement $f(2)=1$ et $f'(2)=2$.\\
        L'équation réduite de la tangente est donc donnée par :
        $y=2(x-2)+1$, soit $y=2x-3$.
\end{Solution}

% @see : https://coopmaths.fr/alea?uuid=a1ba2&id=can1F14&n=1&d=10&alea=rkIp2232323232342345678910111213141516171819202122232425262728293031323334353637383940414243444546474823456789101112131415161718192021222324252627282930313233343536&cd=1&cols=1
\needspace{10\baselineskip}
\begin{exercice}[BaremeTotal=false]


Soit $f$ la fonction définie sur $\mathbb{R}$ par : 
$f(x)= -3x^2-4$.\\
En calculant la limite du taux d'accroissement, déterminer $f'(5)$.
\end{exercice}

\begin{Solution}
 $f$ est une fonction polynôme du second degré de la forme $f(x)=ax^2+b$.\\
    La fonction dérivée est donnée par la somme des dérivées des fonctions $u$ et $v$ définies par $u(x)=-3x^2$ et $v(x)=-4$.\\
     Comme $u'(x)=-6x$ et $v'(x)=0$, on obtient  $f'(x)=-6x$. \\
     Ainsi, $f'(5)=-6\times 5=-30$.
\end{Solution}

% @see : https://coopmaths.fr/alea?uuid=3c690&id=can1F13&n=1&d=10&alea=EUDu2232323232342345678910111213141516171819202122232425262728293031323334353637383940414243444546474823456789101112131415161718192021222324252627282930313233343536&cd=1&cols=1
\needspace{10\baselineskip}
\begin{exercice}[BaremeTotal=false]


 Déterminer le coefficient directeur de la tangente à la courbe représentative de la fonction cube au point d'abscisse $-2$.
                        
\end{exercice}

\begin{Solution}
 Le coefficient directeur de la tangente au point d'abscisse $a$ est donné par le nombre dérivé $f'(a)$.\\
        La fonction $f$ définie par $f(x)=x^3$ a pour fonction dérivée la fonction $f'$ définie par $f'(x)=3x^2$.\\
        Comme $f'(-2)=3\times (-2)^2=12$, le coefficient directeur de la tangente au point d'abscisse $-2$ est : $12$. 
\end{Solution}

% @see : https://coopmaths.fr/alea?uuid=ebd89&id=1AN14-1&n=1&d=10&s=3&alea=ocYr2232323232342345678910111213141516171819202122232425262728293031323334353637383940414243444546474823456789101112131415161718192021222324252627282930313233343536&cd=0&cols=1
\needspace{10\baselineskip}
\begin{exercice}[BaremeTotal=false]


 Donner la dérivée de la fonction $f$, dérivable pour tout $x\in\R$, définie par  $f(x)=\dfrac{3}{4}x+10$.
\end{exercice}

\begin{Solution}
 L'expression de la dérivée de la fonction $f$ définie par $f(x)=\dfrac{3}{4}x+10$ est : ${\color[HTML]{f15929}\boldsymbol{f'(x)=\dfrac{3}{4}}}$.
\end{Solution}

% @see : https://coopmaths.fr/alea?uuid=ebd89&id=1AN14-1&n=1&d=10&s=4&alea=uAlP2232323232342345678910111213141516171819202122232425262728293031323334353637383940414243444546474823456789101112131415161718192021222324252627282930313233343536&cd=0&cols=1
\needspace{10\baselineskip}
\begin{exercice}[BaremeTotal=false]


 Donner la dérivée de la fonction $f$, dérivable pour tout $x\in\R$, définie par  $f(x)=-4x^{6}$.
\end{exercice}

\begin{Solution}
 L'expression de la dérivée de la fonction $f$ définie par $f(x)=-4x^{6}$ est : ${\color[HTML]{f15929}\boldsymbol{f'(x)=-24x^{5}}}$.
\end{Solution}

% @see : https://coopmaths.fr/alea?uuid=ebd89&id=1AN14-1&n=1&d=10&s=7&alea=vjN72232323232342345678910111213141516171819202122232425262728293031323334353637383940414243444546474823456789101112131415161718192021222324252627282930313233343536&cd=0&cols=1
\needspace{10\baselineskip}
\begin{exercice}[BaremeTotal=false]


 Donner la dérivée de la fonction $f$, dérivable pour tout $x\in\R^*$, définie par  $f(x)=\dfrac{5}{7x}$.
\end{exercice}

\begin{Solution}
 L'expression de la dérivée de la fonction $f$ définie par $f(x)=\dfrac{5}{7x}$ est : ${\color[HTML]{f15929}\boldsymbol{f'(x)=-\dfrac{5}{7x^2}}}$.
\end{Solution}

% @see : https://coopmaths.fr/alea?uuid=a83c0&id=1AN14-4&n=1&d=10&s=1&alea=Njr32232323232342345678910111213141516171819202122232425262728293031323334353637383940414243444546474823456789101112131415161718192021222324252627282930313233343536&cd=0&cols=1
\needspace{10\baselineskip}
\begin{exercice}[BaremeTotal=false]


 Donner l'expression de la dérivée de la fonction $f$ définie sur $\R^{*}$ par $f(x)=-3x^{2}+\dfrac{1}{x}$.\\
\end{exercice}

\begin{Solution}
 L'expression de la dérivée de la fonction $f$ définie par $f(x)=-3x^{2}+\dfrac{1}{x}$ est :\\$f'(x)={\color[HTML]{f15929}\boldsymbol{-6x-\dfrac{1}{x^2}}}$.\\
\end{Solution}

% @see : https://coopmaths.fr/alea?uuid=a83c0&id=1AN14-4&n=1&d=10&s=4&alea=4Vpz2232323232342345678910111213141516171819202122232425262728293031323334353637383940414243444546474823456789101112131415161718192021222324252627282930313233343536&cd=1&cols=1
\needspace{10\baselineskip}
\begin{exercice}[BaremeTotal=false]


 Donner l'expression de la dérivée de la fonction $f$ définie sur $\R_+^{*}$ par $f(x)=-\dfrac{2}{x}-2x-4\sqrt{x}$.\\
\end{exercice}

\begin{Solution}
 $(-\dfrac{2}{x})^\prime=\dfrac{2}{x^2}$\\$(-2x)^\prime=-2$\\$(-4\sqrt{x})^\prime=-\dfrac{4}{2\sqrt{x}}$\\L'expression de la dérivée de la fonction $f$ définie par $f(x)=-\dfrac{2}{x}-2x-4\sqrt{x}$ est :\\$f'(x)={\color[HTML]{f15929}\boldsymbol{\dfrac{2}{x^2}-2-\dfrac{4}{2\sqrt{x}}}}$.\\
\end{Solution}

\end{Maquette}
\clearpage
\begin{Maquette}[DM]{Niveau= ,Classe=\getdata[37]\mydata\ (v37), Date = 06/01/25  ,Theme=Exercices}


% @see : https://coopmaths.fr/alea?uuid=4c8c7&id=1AN11-3&n=1&d=10&s=2&alea=nawO223232323234234567891011121314151617181920212223242526272829303132333435363738394041424344454647482345678910111213141516171819202122232425262728293031323334353637&cd=1&cols=2
\needspace{10\baselineskip}
\begin{exercice}[BaremeTotal=false]
\label{eleve:37}


 \begin{minipage}[t]{\linewidth} Soit $f$ une fonction dérivable sur $[-5;5]$ et $\mathcal{C}_f$ sa courbe représentative.\\On sait que  $f(3)=4~~$ et que $~~f'(3)=5$.\\Déterminer une équation de la tangente $(T)$ à la courbe $\mathcal{C}_f$ au point d'abscisse $3$,\\sans utiliser la formule de cours de l'équation de tangente. \end{minipage}

\end{exercice}

\begin{Solution}
 \begin{minipage}[t]{\linewidth} On sait que la tangente n'est pas une droite verticale, puisque la fonction est dérivable sur l'intervalle.\\On en déduit que la tangente $(T)$ au point d'abscisse $3$, admet une équation réduite de la forme :  $(T) : y=m x + p$. 

\medskip
$\bullet$ Détermination de $m$ :\\On sait que le nombre dérivé en $3$ est par définition, le coefficient directeur de la tangente au point d'abscisse $3$.\\Par conséquent, on a déjà : $m=f'(3)=5$.\\ On en déduit que  $(T) : y= 5 x + p$\\$\bullet$ Détermination de $p$ :\\ Pour cela, on utilise que si $f(3)=4~~$, alors le point $A$ de coordonnées $(3;4)$ appartient à $\mathcal{C}_f$ mais aussi à $(T)$.\\ On peut écrire $A(3;4) = \mathcal{C}_f \cap (T)$.\\ On remplace alors les coordonnées de $A(3;4)$ dans l'équation  $(T) : y= 5 x + p$.\\ $\begin{aligned} \phantom{\iff}&A(3;4)\in (T)\\ \iff& 4= 5 \times 3  + p\\ \iff& p=4 -5 \times 3  \\ \iff& p=4 -15   \\ \iff& p=-11\\\end{aligned}$\\On peut conclure que : $(T) : {\color[HTML]{f15929}\boldsymbol{y=5x-11}}$. \end{minipage}

\end{Solution}

% @see : https://coopmaths.fr/alea?uuid=0e984&id=can1F15&n=1&d=10&alea=1Alo223232323234234567891011121314151617181920212223242526272829303132333435363738394041424344454647482345678910111213141516171819202122232425262728293031323334353637&cd=1&cols=1
\needspace{10\baselineskip}
\begin{exercice}[BaremeTotal=false]


 La courbe représente une fonction $f$ et la droite est la tangente au point d'abscisse $2$.\\
        Déterminer $f'(2)$.\begin{center}\begin{tikzpicture}[baseline,scale = 0.5]

    \tikzset{
      point/.style={
        thick,
        draw,
        cross out,
        inner sep=0pt,
        minimum width=5pt,
        minimum height=5pt,
      },
    }
    \clip (-2,-2) rectangle (14,12);
    	
	\draw[color ={black},line width = 2,>=latex,->] (-2,0)--(14,0);
	\draw[color ={black},line width = 2,>=latex,->] (0,-2)--(0,12);
	\draw[color ={black},opacity = 0.5] (-2,1)--(14,1);
	\draw[color ={black},opacity = 0.5] (-2,2)--(14,2);
	\draw[color ={black},opacity = 0.5] (-2,3)--(14,3);
	\draw[color ={black},opacity = 0.5] (-2,4)--(14,4);
	\draw[color ={black},opacity = 0.5] (-2,5)--(14,5);
	\draw[color ={black},opacity = 0.5] (-2,6)--(14,6);
	\draw[color ={black},opacity = 0.5] (-2,7)--(14,7);
	\draw[color ={black},opacity = 0.5] (-2,8)--(14,8);
	\draw[color ={black},opacity = 0.5] (-2,9)--(14,9);
	\draw[color ={black},opacity = 0.5] (-2,10)--(14,10);
	\draw[color ={black},opacity = 0.5] (-2,11)--(14,11);
	\draw[color ={black},opacity = 0.5] (-2,12)--(14,12);
	\draw[color ={black},opacity = 0.5] (-2,-1)--(14,-1);
	\draw[color ={black},opacity = 0.5] (-2,-2)--(14,-2);
	\draw[color ={black},opacity = 0.5] (1,-2)--(1,12);
	\draw[color ={black},opacity = 0.5] (2,-2)--(2,12);
	\draw[color ={black},opacity = 0.5] (3,-2)--(3,12);
	\draw[color ={black},opacity = 0.5] (4,-2)--(4,12);
	\draw[color ={black},opacity = 0.5] (5,-2)--(5,12);
	\draw[color ={black},opacity = 0.5] (6,-2)--(6,12);
	\draw[color ={black},opacity = 0.5] (7,-2)--(7,12);
	\draw[color ={black},opacity = 0.5] (8,-2)--(8,12);
	\draw[color ={black},opacity = 0.5] (9,-2)--(9,12);
	\draw[color ={black},opacity = 0.5] (10,-2)--(10,12);
	\draw[color ={black},opacity = 0.5] (11,-2)--(11,12);
	\draw[color ={black},opacity = 0.5] (12,-2)--(12,12);
	\draw[color ={black},opacity = 0.5] (13,-2)--(13,12);
	\draw[color ={black},opacity = 0.5] (14,-2)--(14,12);
	\draw[color ={black},opacity = 0.5] (-1,-2)--(-1,12);
	\draw[color ={black},opacity = 0.5] (-2,-2)--(-2,12);
	\draw[color ={gray},opacity = 0.3] (-2,0)--(14,0);
	\draw[color ={gray},opacity = 0.3] (-2,0.5)--(14,0.5);
	\draw[color ={gray},opacity = 0.3] (-2,1.5)--(14,1.5);
	\draw[color ={gray},opacity = 0.3] (-2,2.5)--(14,2.5);
	\draw[color ={gray},opacity = 0.3] (-2,3.5)--(14,3.5);
	\draw[color ={gray},opacity = 0.3] (-2,4.5)--(14,4.5);
	\draw[color ={gray},opacity = 0.3] (-2,5.5)--(14,5.5);
	\draw[color ={gray},opacity = 0.3] (-2,6.5)--(14,6.5);
	\draw[color ={gray},opacity = 0.3] (-2,7.5)--(14,7.5);
	\draw[color ={gray},opacity = 0.3] (-2,8.5)--(14,8.5);
	\draw[color ={gray},opacity = 0.3] (-2,9.5)--(14,9.5);
	\draw[color ={gray},opacity = 0.3] (-2,10.5)--(14,10.5);
	\draw[color ={gray},opacity = 0.3] (-2,11.5)--(14,11.5);
	\draw[color ={gray},opacity = 0.3] (-2,0)--(14,0);
	\draw[color ={gray},opacity = 0.3] (-2,-0.5)--(14,-0.5);
	\draw[color ={gray},opacity = 0.3] (-2,-1.5)--(14,-1.5);
	\draw[color ={gray},opacity = 0.3] (0,-2)--(0,12);
	\draw[color ={gray},opacity = 0.3] (0.1,-2)--(0.1,12);
	\draw[color ={gray},opacity = 0.3] (0.2,-2)--(0.2,12);
	\draw[color ={gray},opacity = 0.3] (0.30000000000000004,-2)--(0.30000000000000004,12);
	\draw[color ={gray},opacity = 0.3] (0.4,-2)--(0.4,12);
	\draw[color ={gray},opacity = 0.3] (0.5,-2)--(0.5,12);
	\draw[color ={gray},opacity = 0.3] (0.6,-2)--(0.6,12);
	\draw[color ={gray},opacity = 0.3] (0.7,-2)--(0.7,12);
	\draw[color ={gray},opacity = 0.3] (0.7999999999999999,-2)--(0.7999999999999999,12);
	\draw[color ={gray},opacity = 0.3] (0.8999999999999999,-2)--(0.8999999999999999,12);
	\draw[color ={gray},opacity = 0.3] (1.0999999999999999,-2)--(1.0999999999999999,12);
	\draw[color ={gray},opacity = 0.3] (1.2,-2)--(1.2,12);
	\draw[color ={gray},opacity = 0.3] (1.3,-2)--(1.3,12);
	\draw[color ={gray},opacity = 0.3] (1.4000000000000001,-2)--(1.4000000000000001,12);
	\draw[color ={gray},opacity = 0.3] (1.5000000000000002,-2)--(1.5000000000000002,12);
	\draw[color ={gray},opacity = 0.3] (1.6000000000000003,-2)--(1.6000000000000003,12);
	\draw[color ={gray},opacity = 0.3] (1.7000000000000004,-2)--(1.7000000000000004,12);
	\draw[color ={gray},opacity = 0.3] (1.8000000000000005,-2)--(1.8000000000000005,12);
	\draw[color ={gray},opacity = 0.3] (1.9000000000000006,-2)--(1.9000000000000006,12);
	\draw[color ={gray},opacity = 0.3] (2.1000000000000005,-2)--(2.1000000000000005,12);
	\draw[color ={gray},opacity = 0.3] (2.2000000000000006,-2)--(2.2000000000000006,12);
	\draw[color ={gray},opacity = 0.3] (2.3000000000000007,-2)--(2.3000000000000007,12);
	\draw[color ={gray},opacity = 0.3] (2.400000000000001,-2)--(2.400000000000001,12);
	\draw[color ={gray},opacity = 0.3] (2.500000000000001,-2)--(2.500000000000001,12);
	\draw[color ={gray},opacity = 0.3] (2.600000000000001,-2)--(2.600000000000001,12);
	\draw[color ={gray},opacity = 0.3] (2.700000000000001,-2)--(2.700000000000001,12);
	\draw[color ={gray},opacity = 0.3] (2.800000000000001,-2)--(2.800000000000001,12);
	\draw[color ={gray},opacity = 0.3] (2.9000000000000012,-2)--(2.9000000000000012,12);
	\draw[color ={gray},opacity = 0.3] (3.1000000000000014,-2)--(3.1000000000000014,12);
	\draw[color ={gray},opacity = 0.3] (3.2000000000000015,-2)--(3.2000000000000015,12);
	\draw[color ={gray},opacity = 0.3] (3.3000000000000016,-2)--(3.3000000000000016,12);
	\draw[color ={gray},opacity = 0.3] (3.4000000000000017,-2)--(3.4000000000000017,12);
	\draw[color ={gray},opacity = 0.3] (3.5000000000000018,-2)--(3.5000000000000018,12);
	\draw[color ={gray},opacity = 0.3] (3.600000000000002,-2)--(3.600000000000002,12);
	\draw[color ={gray},opacity = 0.3] (3.700000000000002,-2)--(3.700000000000002,12);
	\draw[color ={gray},opacity = 0.3] (3.800000000000002,-2)--(3.800000000000002,12);
	\draw[color ={gray},opacity = 0.3] (3.900000000000002,-2)--(3.900000000000002,12);
	\draw[color ={gray},opacity = 0.3] (4.100000000000001,-2)--(4.100000000000001,12);
	\draw[color ={gray},opacity = 0.3] (4.200000000000001,-2)--(4.200000000000001,12);
	\draw[color ={gray},opacity = 0.3] (4.300000000000001,-2)--(4.300000000000001,12);
	\draw[color ={gray},opacity = 0.3] (4.4,-2)--(4.4,12);
	\draw[color ={gray},opacity = 0.3] (4.5,-2)--(4.5,12);
	\draw[color ={gray},opacity = 0.3] (4.6,-2)--(4.6,12);
	\draw[color ={gray},opacity = 0.3] (4.699999999999999,-2)--(4.699999999999999,12);
	\draw[color ={gray},opacity = 0.3] (4.799999999999999,-2)--(4.799999999999999,12);
	\draw[color ={gray},opacity = 0.3] (4.899999999999999,-2)--(4.899999999999999,12);
	\draw[color ={gray},opacity = 0.3] (5.099999999999998,-2)--(5.099999999999998,12);
	\draw[color ={gray},opacity = 0.3] (5.1999999999999975,-2)--(5.1999999999999975,12);
	\draw[color ={gray},opacity = 0.3] (5.299999999999997,-2)--(5.299999999999997,12);
	\draw[color ={gray},opacity = 0.3] (5.399999999999997,-2)--(5.399999999999997,12);
	\draw[color ={gray},opacity = 0.3] (5.4999999999999964,-2)--(5.4999999999999964,12);
	\draw[color ={gray},opacity = 0.3] (5.599999999999996,-2)--(5.599999999999996,12);
	\draw[color ={gray},opacity = 0.3] (5.699999999999996,-2)--(5.699999999999996,12);
	\draw[color ={gray},opacity = 0.3] (5.799999999999995,-2)--(5.799999999999995,12);
	\draw[color ={gray},opacity = 0.3] (5.899999999999995,-2)--(5.899999999999995,12);
	\draw[color ={gray},opacity = 0.3] (6.099999999999994,-2)--(6.099999999999994,12);
	\draw[color ={gray},opacity = 0.3] (6.199999999999994,-2)--(6.199999999999994,12);
	\draw[color ={gray},opacity = 0.3] (6.299999999999994,-2)--(6.299999999999994,12);
	\draw[color ={gray},opacity = 0.3] (6.399999999999993,-2)--(6.399999999999993,12);
	\draw[color ={gray},opacity = 0.3] (6.499999999999993,-2)--(6.499999999999993,12);
	\draw[color ={gray},opacity = 0.3] (6.5999999999999925,-2)--(6.5999999999999925,12);
	\draw[color ={gray},opacity = 0.3] (6.699999999999992,-2)--(6.699999999999992,12);
	\draw[color ={gray},opacity = 0.3] (6.799999999999992,-2)--(6.799999999999992,12);
	\draw[color ={gray},opacity = 0.3] (6.8999999999999915,-2)--(6.8999999999999915,12);
	\draw[color ={gray},opacity = 0.3] (7.099999999999991,-2)--(7.099999999999991,12);
	\draw[color ={gray},opacity = 0.3] (7.19999999999999,-2)--(7.19999999999999,12);
	\draw[color ={gray},opacity = 0.3] (7.29999999999999,-2)--(7.29999999999999,12);
	\draw[color ={gray},opacity = 0.3] (7.39999999999999,-2)--(7.39999999999999,12);
	\draw[color ={gray},opacity = 0.3] (7.499999999999989,-2)--(7.499999999999989,12);
	\draw[color ={gray},opacity = 0.3] (7.599999999999989,-2)--(7.599999999999989,12);
	\draw[color ={gray},opacity = 0.3] (7.699999999999989,-2)--(7.699999999999989,12);
	\draw[color ={gray},opacity = 0.3] (7.799999999999988,-2)--(7.799999999999988,12);
	\draw[color ={gray},opacity = 0.3] (7.899999999999988,-2)--(7.899999999999988,12);
	\draw[color ={gray},opacity = 0.3] (8.099999999999987,-2)--(8.099999999999987,12);
	\draw[color ={gray},opacity = 0.3] (8.199999999999987,-2)--(8.199999999999987,12);
	\draw[color ={gray},opacity = 0.3] (8.299999999999986,-2)--(8.299999999999986,12);
	\draw[color ={gray},opacity = 0.3] (8.399999999999986,-2)--(8.399999999999986,12);
	\draw[color ={gray},opacity = 0.3] (8.499999999999986,-2)--(8.499999999999986,12);
	\draw[color ={gray},opacity = 0.3] (8.599999999999985,-2)--(8.599999999999985,12);
	\draw[color ={gray},opacity = 0.3] (8.699999999999985,-2)--(8.699999999999985,12);
	\draw[color ={gray},opacity = 0.3] (8.799999999999985,-2)--(8.799999999999985,12);
	\draw[color ={gray},opacity = 0.3] (8.899999999999984,-2)--(8.899999999999984,12);
	\draw[color ={gray},opacity = 0.3] (9.099999999999984,-2)--(9.099999999999984,12);
	\draw[color ={gray},opacity = 0.3] (9.199999999999983,-2)--(9.199999999999983,12);
	\draw[color ={gray},opacity = 0.3] (9.299999999999983,-2)--(9.299999999999983,12);
	\draw[color ={gray},opacity = 0.3] (9.399999999999983,-2)--(9.399999999999983,12);
	\draw[color ={gray},opacity = 0.3] (9.499999999999982,-2)--(9.499999999999982,12);
	\draw[color ={gray},opacity = 0.3] (9.599999999999982,-2)--(9.599999999999982,12);
	\draw[color ={gray},opacity = 0.3] (9.699999999999982,-2)--(9.699999999999982,12);
	\draw[color ={gray},opacity = 0.3] (9.799999999999981,-2)--(9.799999999999981,12);
	\draw[color ={gray},opacity = 0.3] (9.89999999999998,-2)--(9.89999999999998,12);
	\draw[color ={gray},opacity = 0.3] (10.09999999999998,-2)--(10.09999999999998,12);
	\draw[color ={gray},opacity = 0.3] (10.19999999999998,-2)--(10.19999999999998,12);
	\draw[color ={gray},opacity = 0.3] (10.29999999999998,-2)--(10.29999999999998,12);
	\draw[color ={gray},opacity = 0.3] (10.399999999999979,-2)--(10.399999999999979,12);
	\draw[color ={gray},opacity = 0.3] (10.499999999999979,-2)--(10.499999999999979,12);
	\draw[color ={gray},opacity = 0.3] (10.599999999999978,-2)--(10.599999999999978,12);
	\draw[color ={gray},opacity = 0.3] (10.699999999999978,-2)--(10.699999999999978,12);
	\draw[color ={gray},opacity = 0.3] (10.799999999999978,-2)--(10.799999999999978,12);
	\draw[color ={gray},opacity = 0.3] (10.899999999999977,-2)--(10.899999999999977,12);
	\draw[color ={gray},opacity = 0.3] (11.099999999999977,-2)--(11.099999999999977,12);
	\draw[color ={gray},opacity = 0.3] (11.199999999999976,-2)--(11.199999999999976,12);
	\draw[color ={gray},opacity = 0.3] (11.299999999999976,-2)--(11.299999999999976,12);
	\draw[color ={gray},opacity = 0.3] (11.399999999999975,-2)--(11.399999999999975,12);
	\draw[color ={gray},opacity = 0.3] (11.499999999999975,-2)--(11.499999999999975,12);
	\draw[color ={gray},opacity = 0.3] (11.599999999999975,-2)--(11.599999999999975,12);
	\draw[color ={gray},opacity = 0.3] (11.699999999999974,-2)--(11.699999999999974,12);
	\draw[color ={gray},opacity = 0.3] (11.799999999999974,-2)--(11.799999999999974,12);
	\draw[color ={gray},opacity = 0.3] (11.899999999999974,-2)--(11.899999999999974,12);
	\draw[color ={gray},opacity = 0.3] (12.099999999999973,-2)--(12.099999999999973,12);
	\draw[color ={gray},opacity = 0.3] (12.199999999999973,-2)--(12.199999999999973,12);
	\draw[color ={gray},opacity = 0.3] (12.299999999999972,-2)--(12.299999999999972,12);
	\draw[color ={gray},opacity = 0.3] (12.399999999999972,-2)--(12.399999999999972,12);
	\draw[color ={gray},opacity = 0.3] (12.499999999999972,-2)--(12.499999999999972,12);
	\draw[color ={gray},opacity = 0.3] (12.599999999999971,-2)--(12.599999999999971,12);
	\draw[color ={gray},opacity = 0.3] (12.69999999999997,-2)--(12.69999999999997,12);
	\draw[color ={gray},opacity = 0.3] (12.79999999999997,-2)--(12.79999999999997,12);
	\draw[color ={gray},opacity = 0.3] (12.89999999999997,-2)--(12.89999999999997,12);
	\draw[color ={gray},opacity = 0.3] (13.09999999999997,-2)--(13.09999999999997,12);
	\draw[color ={gray},opacity = 0.3] (13.199999999999969,-2)--(13.199999999999969,12);
	\draw[color ={gray},opacity = 0.3] (13.299999999999969,-2)--(13.299999999999969,12);
	\draw[color ={gray},opacity = 0.3] (13.399999999999968,-2)--(13.399999999999968,12);
	\draw[color ={gray},opacity = 0.3] (13.499999999999968,-2)--(13.499999999999968,12);
	\draw[color ={gray},opacity = 0.3] (13.599999999999968,-2)--(13.599999999999968,12);
	\draw[color ={gray},opacity = 0.3] (13.699999999999967,-2)--(13.699999999999967,12);
	\draw[color ={gray},opacity = 0.3] (13.799999999999967,-2)--(13.799999999999967,12);
	\draw[color ={gray},opacity = 0.3] (13.899999999999967,-2)--(13.899999999999967,12);
	\draw[color ={gray},opacity = 0.3] (0,-2)--(0,12);
	\draw[color ={gray},opacity = 0.3] (-0.1,-2)--(-0.1,12);
	\draw[color ={gray},opacity = 0.3] (-0.2,-2)--(-0.2,12);
	\draw[color ={gray},opacity = 0.3] (-0.30000000000000004,-2)--(-0.30000000000000004,12);
	\draw[color ={gray},opacity = 0.3] (-0.4,-2)--(-0.4,12);
	\draw[color ={gray},opacity = 0.3] (-0.5,-2)--(-0.5,12);
	\draw[color ={gray},opacity = 0.3] (-0.6,-2)--(-0.6,12);
	\draw[color ={gray},opacity = 0.3] (-0.7,-2)--(-0.7,12);
	\draw[color ={gray},opacity = 0.3] (-0.7999999999999999,-2)--(-0.7999999999999999,12);
	\draw[color ={gray},opacity = 0.3] (-0.8999999999999999,-2)--(-0.8999999999999999,12);
	\draw[color ={gray},opacity = 0.3] (-1.0999999999999999,-2)--(-1.0999999999999999,12);
	\draw[color ={gray},opacity = 0.3] (-1.2,-2)--(-1.2,12);
	\draw[color ={gray},opacity = 0.3] (-1.3,-2)--(-1.3,12);
	\draw[color ={gray},opacity = 0.3] (-1.4000000000000001,-2)--(-1.4000000000000001,12);
	\draw[color ={gray},opacity = 0.3] (-1.5000000000000002,-2)--(-1.5000000000000002,12);
	\draw[color ={gray},opacity = 0.3] (-1.6000000000000003,-2)--(-1.6000000000000003,12);
	\draw[color ={gray},opacity = 0.3] (-1.7000000000000004,-2)--(-1.7000000000000004,12);
	\draw[color ={gray},opacity = 0.3] (-1.8000000000000005,-2)--(-1.8000000000000005,12);
	\draw[color ={gray},opacity = 0.3] (-1.9000000000000006,-2)--(-1.9000000000000006,12);
	\draw[color ={black},line width = 2] (2,-0.2)--(2,0.2);
	\draw[color ={black},line width = 2] (4,-0.2)--(4,0.2);
	\draw[color ={black},line width = 2] (6,-0.2)--(6,0.2);
	\draw[color ={black},line width = 2] (8,-0.2)--(8,0.2);
	\draw[color ={black},line width = 2] (10,-0.2)--(10,0.2);
	\draw[color ={black},line width = 2] (12,-0.2)--(12,0.2);
	\draw[color ={black},line width = 2] (-0.2,2)--(0.2,2);
	\draw[color ={black},line width = 2] (-0.2,4)--(0.2,4);
	\draw[color ={black},line width = 2] (-0.2,6)--(0.2,6);
	\draw[color ={black},line width = 2] (-0.2,8)--(0.2,8);
	\draw[color ={black},line width = 2] (-0.2,10)--(0.2,10);
	\draw (2,-0.4) node[anchor = center] {\scriptsize \color{black}{$1$}};
	\draw (4,-0.4) node[anchor = center] {\scriptsize \color{black}{$2$}};
	\draw (6,-0.4) node[anchor = center] {\scriptsize \color{black}{$3$}};
	\draw (8,-0.4) node[anchor = center] {\scriptsize \color{black}{$4$}};
	\draw (10,-0.4) node[anchor = center] {\scriptsize \color{black}{$5$}};
	\draw (-0.5,2.1) node[anchor = center] {\footnotesize \color{black}{$1$}};
	\draw (-0.5,4.1) node[anchor = center] {\footnotesize \color{black}{$2$}};
	\draw (-0.5,6.1) node[anchor = center] {\footnotesize \color{black}{$3$}};
	\draw (-0.5,8.1) node[anchor = center] {\footnotesize \color{black}{$4$}};
	\draw (-0.5,10.1) node[anchor = center] {\footnotesize \color{black}{$5$}};
	\draw [color={black}] (-0.3,-0.3) node[anchor = center,scale=1, rotate = 0] {O};
	
	\draw[color={blue},line width = 2] (1.4,11.71)--(1.8,10.44)--(2.2,9.64)--(2.6,9.08)--(3,8.67)--(3.4,8.35)--(3.8,8.11)--(4.2,7.9)--(4.6,7.74)--(5,7.6)--(5.4,7.48)--(5.8,7.38)--(6.2,7.29)--(6.6,7.21)--(7,7.14)--(7.4,7.08)--(7.8,7.03)--(8.2,6.98)--(8.6,6.93)--(9,6.89)--(9.4,6.85)--(9.8,6.82)--(10.2,6.78)--(10.6,6.75)--(11,6.73)--(11.4,6.7)--(11.8,6.68)--(12.2,6.66)--(12.6,6.63)--(13,6.62)--(13.4,6.6)--(13.8,6.58);
	
	\draw[color={red},line width = 2] (-2,11)--(-1.6,10.8)--(-1.2,10.6)--(-0.8,10.4)--(-0.4,10.2)--(0,10)--(0.4,9.8)--(0.8,9.6)--(1.2,9.4)--(1.6,9.2)--(2,9)--(2.4,8.8)--(2.8,8.6)--(3.2,8.4)--(3.6,8.2)--(4,8)--(4.4,7.8)--(4.8,7.6)--(5.2,7.4)--(5.6,7.2)--(6,7)--(6.4,6.8)--(6.8,6.6)--(7.2,6.4)--(7.6,6.2)--(8,6)--(8.4,5.8)--(8.8,5.6)--(9.2,5.4)--(9.6,5.2)--(10,5)--(10.4,4.8)--(10.8,4.6)--(11.2,4.4)--(11.6,4.2)--(12,4)--(12.4,3.8)--(12.8,3.6)--(13.2,3.4)--(13.6,3.2)--(14,3);

\end{tikzpicture}\end{center}
\end{exercice}

\begin{Solution}
 $f'(2)$ est donné par le coefficient directeur de la tangente à la courbe au point d'abscisse $2$, soit $\dfrac{-2}{4}=-\dfrac{1{\color[HTML]{2563a5}\boldsymbol{\times2}} }{2{\color[HTML]{2563a5}\boldsymbol{\times2}}}=-\dfrac{1}{2}$.
\end{Solution}

% @see : https://coopmaths.fr/alea?uuid=6f32d&id=can1F16&n=1&d=10&alea=dqSw223232323234234567891011121314151617181920212223242526272829303132333435363738394041424344454647482345678910111213141516171819202122232425262728293031323334353637&cd=1&cols=1
\needspace{10\baselineskip}
\newpage
\begin{exercice}[BaremeTotal=false]


 On donne les représentations graphiques d'une fonction et de sa dérivée.\\
        Donner l'équation réduite de la tangente à la courbe de $f$ en $x=3$. \\ \begin{minipage}{0.5\linewidth}
    \begin{tikzpicture}[baseline, scale=0.8]

    \tikzset{
      point/.style={
        thick,
        draw,
        cross out,
        inner sep=0pt,
        minimum width=5pt,
        minimum height=5pt,
      },
    }
    \clip (-4.5,-4.5) rectangle (5.75,13.5);
    	
	\draw[color ={black},>=latex,->] (-3,0)--(3,0);
	\draw[color ={black},>=latex,->] (0,-3)--(0,12);
	\draw[color ={black},opacity = 0.4] (-3,1)--(3,1);
	\draw[color ={black},opacity = 0.4] (-3,2)--(3,2);
	\draw[color ={black},opacity = 0.4] (-3,3)--(3,3);
	\draw[color ={black},opacity = 0.4] (-3,4)--(3,4);
	\draw[color ={black},opacity = 0.4] (-3,5)--(3,5);
	\draw[color ={black},opacity = 0.4] (-3,6)--(3,6);
	\draw[color ={black},opacity = 0.4] (-3,7)--(3,7);
	\draw[color ={black},opacity = 0.4] (-3,8)--(3,8);
	\draw[color ={black},opacity = 0.4] (-3,9)--(3,9);
	\draw[color ={black},opacity = 0.4] (-3,10)--(3,10);
	\draw[color ={black},opacity = 0.4] (-3,11)--(3,11);
	\draw[color ={black},opacity = 0.4] (-3,12)--(3,12);
	\draw[color ={black},opacity = 0.4] (-3,-1)--(3,-1);
	\draw[color ={black},opacity = 0.4] (-3,-2)--(3,-2);
	\draw[color ={black},opacity = 0.4] (-3,-3)--(3,-3);
	\draw[color ={black},opacity = 0.4] (1,-3)--(1,12);
	\draw[color ={black},opacity = 0.4] (2,-3)--(2,12);
	\draw[color ={black},opacity = 0.4] (3,-3)--(3,12);
	\draw[color ={black},opacity = 0.4] (-1,-3)--(-1,12);
	\draw[color ={black},opacity = 0.4] (-2,-3)--(-2,12);
	\draw[color ={black},opacity = 0.4] (-3,-3)--(-3,12);
	\draw[color ={gray},opacity = 0.3] (-3,0)--(3,0);
	\draw[color ={gray},opacity = 0.3] (-3,0.5)--(3,0.5);
	\draw[color ={gray},opacity = 0.3] (-3,1.5)--(3,1.5);
	\draw[color ={gray},opacity = 0.3] (-3,2.5)--(3,2.5);
	\draw[color ={gray},opacity = 0.3] (-3,3.5)--(3,3.5);
	\draw[color ={gray},opacity = 0.3] (-3,4.5)--(3,4.5);
	\draw[color ={gray},opacity = 0.3] (-3,5.5)--(3,5.5);
	\draw[color ={gray},opacity = 0.3] (-3,6.5)--(3,6.5);
	\draw[color ={gray},opacity = 0.3] (-3,7.5)--(3,7.5);
	\draw[color ={gray},opacity = 0.3] (-3,8.5)--(3,8.5);
	\draw[color ={gray},opacity = 0.3] (-3,9.5)--(3,9.5);
	\draw[color ={gray},opacity = 0.3] (-3,10.5)--(3,10.5);
	\draw[color ={gray},opacity = 0.3] (-3,11.5)--(3,11.5);
	\draw[color ={gray},opacity = 0.3] (-3,0)--(3,0);
	\draw[color ={gray},opacity = 0.3] (-3,-0.5)--(3,-0.5);
	\draw[color ={gray},opacity = 0.3] (-3,-1.5)--(3,-1.5);
	\draw[color ={gray},opacity = 0.3] (-3,-2.5)--(3,-2.5);
	\draw[color ={gray},opacity = 0.3] (0,-3)--(0,12);
	\draw[color ={gray},opacity = 0.3] (0,-3)--(0,12);
	\draw[color ={black},line width = 1.2] (1,-0.13)--(1,0.13);
	\draw[color ={black},line width = 1.2] (2,-0.13)--(2,0.13);
	\draw[color ={black},line width = 1.2] (3,-0.13)--(3,0.13);
	\draw[color ={black},line width = 1.2] (-1,-0.13)--(-1,0.13);
	\draw[color ={black},line width = 1.2] (-2,-0.13)--(-2,0.13);
	\draw[color ={black},line width = 1.2] (-3,-0.13)--(-3,0.13);
	\draw[color ={black},line width = 1.2] (-0.13,1)--(0.13,1);
	\draw[color ={black},line width = 1.2] (-0.13,2)--(0.13,2);
	\draw[color ={black},line width = 1.2] (-0.13,3)--(0.13,3);
	\draw[color ={black},line width = 1.2] (-0.13,4)--(0.13,4);
	\draw[color ={black},line width = 1.2] (-0.13,5)--(0.13,5);
	\draw[color ={black},line width = 1.2] (-0.13,6)--(0.13,6);
	\draw[color ={black},line width = 1.2] (-0.13,7)--(0.13,7);
	\draw[color ={black},line width = 1.2] (-0.13,8)--(0.13,8);
	\draw[color ={black},line width = 1.2] (-0.13,9)--(0.13,9);
	\draw[color ={black},line width = 1.2] (-0.13,10)--(0.13,10);
	\draw[color ={black},line width = 1.2] (-0.13,11)--(0.13,11);
	\draw[color ={black},line width = 1.2] (-0.13,12)--(0.13,12);
	\draw[color ={black},line width = 1.2] (-0.13,-1)--(0.13,-1);
	\draw[color ={black},line width = 1.2] (-0.13,-2)--(0.13,-2);
	\draw[color ={black},line width = 1.2] (-0.13,-3)--(0.13,-3);
	\draw (3,-0.4) node[anchor = center] {\scriptsize \color{black}{$3$}};
	\draw (-3,-0.4) node[anchor = center] {\scriptsize \color{black}{$-3$}};
	\draw (-0.5,3.1) node[anchor = center] {\footnotesize \color{black}{$3$}};
	\draw (-0.5,6.1) node[anchor = center] {\footnotesize \color{black}{$6$}};
	\draw (-0.5,9.1) node[anchor = center] {\footnotesize \color{black}{$9$}};
	\draw (-0.5,-2.9) node[anchor = center] {\footnotesize \color{black}{$-3$}};
	\draw [color={black}] (-0.3,-0.3) node[anchor = center,scale=1, rotate = 0] {O};
	\draw (3,10) node[anchor = center] {\footnotesize \color{blue}{$\Large \mathcal{C}_f$}};
	
	\draw[color={blue},line width = 2] (-1.4,12.56)--(-1.2,11.24)--(-1,10)--(-0.8,8.84)--(-0.6,7.76)--(-0.4,6.76)--(-0.2,5.84)--(0,5)--(0.2,4.24)--(0.4,3.56)--(0.6,2.96)--(0.8,2.44)--(1,2)--(1.2,1.64)--(1.4,1.36)--(1.6,1.16)--(1.8,1.04)--(2,1)--(2.2,1.04)--(2.4,1.16)--(2.6,1.36)--(2.8,1.64)--(3,2);

\end{tikzpicture}\\
    \end{minipage}
    \begin{minipage}{0.5\linewidth}
    \begin{tikzpicture}[baseline, scale=0.8]

    \tikzset{
      point/.style={
        thick,
        draw,
        cross out,
        inner sep=0pt,
        minimum width=5pt,
        minimum height=5pt,
      },
    }
    \clip (-4.5,-6.5) rectangle (6.5,9.5);
    	
	\draw[color ={black},>=latex,->] (-3,0)--(3,0);
	\draw[color ={black},>=latex,->] (0,-5)--(0,8);
	\draw[color ={black},opacity = 0.4] (-3,1)--(3,1);
	\draw[color ={black},opacity = 0.4] (-3,2)--(3,2);
	\draw[color ={black},opacity = 0.4] (-3,3)--(3,3);
	\draw[color ={black},opacity = 0.4] (-3,4)--(3,4);
	\draw[color ={black},opacity = 0.4] (-3,5)--(3,5);
	\draw[color ={black},opacity = 0.4] (-3,6)--(3,6);
	\draw[color ={black},opacity = 0.4] (-3,7)--(3,7);
	\draw[color ={black},opacity = 0.4] (-3,8)--(3,8);
	\draw[color ={black},opacity = 0.4] (-3,-1)--(3,-1);
	\draw[color ={black},opacity = 0.4] (-3,-2)--(3,-2);
	\draw[color ={black},opacity = 0.4] (-3,-3)--(3,-3);
	\draw[color ={black},opacity = 0.4] (-3,-4)--(3,-4);
	\draw[color ={black},opacity = 0.4] (-3,-5)--(3,-5);
	\draw[color ={black},opacity = 0.4] (1,-5)--(1,8);
	\draw[color ={black},opacity = 0.4] (2,-5)--(2,8);
	\draw[color ={black},opacity = 0.4] (3,-5)--(3,8);
	\draw[color ={black},opacity = 0.4] (-1,-5)--(-1,8);
	\draw[color ={black},opacity = 0.4] (-2,-5)--(-2,8);
	\draw[color ={black},opacity = 0.4] (-3,-5)--(-3,8);
	\draw[color ={gray},opacity = 0.3] (-3,0)--(3,0);
	\draw[color ={gray},opacity = 0.3] (-3,0.5)--(3,0.5);
	\draw[color ={gray},opacity = 0.3] (-3,1.5)--(3,1.5);
	\draw[color ={gray},opacity = 0.3] (-3,2.5)--(3,2.5);
	\draw[color ={gray},opacity = 0.3] (-3,3.5)--(3,3.5);
	\draw[color ={gray},opacity = 0.3] (-3,4.5)--(3,4.5);
	\draw[color ={gray},opacity = 0.3] (-3,5.5)--(3,5.5);
	\draw[color ={gray},opacity = 0.3] (-3,6.5)--(3,6.5);
	\draw[color ={gray},opacity = 0.3] (-3,7.5)--(3,7.5);
	\draw[color ={gray},opacity = 0.3] (-3,0)--(3,0);
	\draw[color ={gray},opacity = 0.3] (-3,-0.5)--(3,-0.5);
	\draw[color ={gray},opacity = 0.3] (-3,-1.5)--(3,-1.5);
	\draw[color ={gray},opacity = 0.3] (-3,-2.5)--(3,-2.5);
	\draw[color ={gray},opacity = 0.3] (-3,-3.5)--(3,-3.5);
	\draw[color ={gray},opacity = 0.3] (-3,-4.5)--(3,-4.5);
	\draw[color ={gray},opacity = 0.3] (0,-5)--(0,8);
	\draw[color ={gray},opacity = 0.3] (0,-5)--(0,8);
	\draw[color ={black},line width = 1.2] (1,-0.13)--(1,0.13);
	\draw[color ={black},line width = 1.2] (2,-0.13)--(2,0.13);
	\draw[color ={black},line width = 1.2] (3,-0.13)--(3,0.13);
	\draw[color ={black},line width = 1.2] (-1,-0.13)--(-1,0.13);
	\draw[color ={black},line width = 1.2] (-2,-0.13)--(-2,0.13);
	\draw[color ={black},line width = 1.2] (-3,-0.13)--(-3,0.13);
	\draw[color ={black},line width = 1.2] (-0.13,1)--(0.13,1);
	\draw[color ={black},line width = 1.2] (-0.13,2)--(0.13,2);
	\draw[color ={black},line width = 1.2] (-0.13,3)--(0.13,3);
	\draw[color ={black},line width = 1.2] (-0.13,4)--(0.13,4);
	\draw[color ={black},line width = 1.2] (-0.13,5)--(0.13,5);
	\draw[color ={black},line width = 1.2] (-0.13,6)--(0.13,6);
	\draw[color ={black},line width = 1.2] (-0.13,7)--(0.13,7);
	\draw[color ={black},line width = 1.2] (-0.13,8)--(0.13,8);
	\draw[color ={black},line width = 1.2] (-0.13,9)--(0.13,9);
	\draw[color ={black},line width = 1.2] (-0.13,10)--(0.13,10);
	\draw[color ={black},line width = 1.2] (-0.13,11)--(0.13,11);
	\draw[color ={black},line width = 1.2] (-0.13,12)--(0.13,12);
	\draw[color ={black},line width = 1.2] (-0.13,-1)--(0.13,-1);
	\draw[color ={black},line width = 1.2] (-0.13,-2)--(0.13,-2);
	\draw[color ={black},line width = 1.2] (-0.13,-3)--(0.13,-3);
	\draw (3,-0.4) node[anchor = center] {\scriptsize \color{black}{$3$}};
	\draw (-3,-0.4) node[anchor = center] {\scriptsize \color{black}{$-3$}};
	\draw (-0.5,3.1) node[anchor = center] {\footnotesize \color{black}{$3$}};
	\draw (-0.5,6.1) node[anchor = center] {\footnotesize \color{black}{$6$}};
	\draw (-0.5,-2.9) node[anchor = center] {\footnotesize \color{black}{$-3$}};
	\draw [color={black}] (-0.3,-0.3) node[anchor = center,scale=1, rotate = 0] {O};
	\draw (3,6) node[anchor = center] {\footnotesize \color{red}{$\Large\mathcal{C}_f\prime$}};
	
	\draw[color={red},line width = 2] (-1,-6)--(-0.8,-5.6)--(-0.6,-5.2)--(-0.4,-4.8)--(-0.2,-4.4)--(0,-4)--(0.2,-3.6)--(0.4,-3.2)--(0.6,-2.8)--(0.8,-2.4)--(1,-2)--(1.2,-1.6)--(1.4,-1.2)--(1.6,-0.8)--(1.8,-0.4)--(2,0)--(2.2,0.4)--(2.4,0.8)--(2.6,1.2)--(2.8,1.6)--(3,2);

\end{tikzpicture}\\
    \end{minipage}
    
\end{exercice}

\begin{Solution}
 L'équation réduite de la tangente au point d'abscisse $3$ est  : $y=f'(3)(x-3)+f(3)$.\\
        On lit graphiquement $f(3)=2$ et $f'(3)=2$.\\
        L'équation réduite de la tangente est donc donnée par :
        $y=2(x-3)+2$, soit $y=2x-4$.
\end{Solution}

% @see : https://coopmaths.fr/alea?uuid=a1ba2&id=can1F14&n=1&d=10&alea=rkIp223232323234234567891011121314151617181920212223242526272829303132333435363738394041424344454647482345678910111213141516171819202122232425262728293031323334353637&cd=1&cols=1
\needspace{10\baselineskip}
\begin{exercice}[BaremeTotal=false]


 Soit $f$ la fonction définie sur $\mathbb{R}$ par : $f(x)= -3x^2+4x+4$.\\
        En calculant la limite du taux d'accroissement, déterminer $f'(-2)$.
\end{exercice}

\begin{Solution}
 $f$ est une fonction polynôme du second degré de la forme $f(x)=ax^2+bx+c$.\\
    La fonction dérivée est donnée par la somme des dérivées des fonctions $u$ et $v$ définies par $u(x)=-3x^2$ et $v(x)=4x+4$.\\
     Comme $u'(x)=-6x$ et $v'(x)=4$, on obtient  $f'(x)=-6x+4$. \\
     Ainsi, $f'(-2)=-6\times (-2)+4=16$.
\end{Solution}

% @see : https://coopmaths.fr/alea?uuid=3c690&id=can1F13&n=1&d=10&alea=EUDu223232323234234567891011121314151617181920212223242526272829303132333435363738394041424344454647482345678910111213141516171819202122232425262728293031323334353637&cd=1&cols=1
\needspace{10\baselineskip}
\begin{exercice}[BaremeTotal=false]


 Déterminer le coefficient directeur de la tangente à la courbe représentative de la fonction inverse au point d'abscisse $-6$.
                 
\end{exercice}

\begin{Solution}
 Le coefficient directeur de la tangente au point d'abscisse $a$ est donné par le nombre dérivé $f'(a)$.\\
        La fonction $f$ définie par $f(x)=\dfrac{1}{x}$ a pour fonction dérivée la fonction $f'$ définie par $f'(x)=-\dfrac{1}{x^2}$.\\
Comme $f'(-6)=-\dfrac{1}{(-6)^2}=-\dfrac{1}{36}$, le coefficient directeur de la tangente au point d'abscisse $-6$ est : $-\dfrac{1}{36}$.
\end{Solution}

% @see : https://coopmaths.fr/alea?uuid=ebd89&id=1AN14-1&n=1&d=10&s=3&alea=ocYr223232323234234567891011121314151617181920212223242526272829303132333435363738394041424344454647482345678910111213141516171819202122232425262728293031323334353637&cd=0&cols=1
\needspace{10\baselineskip}
\begin{exercice}[BaremeTotal=false]


 Donner la dérivée de la fonction $f$, dérivable pour tout $x\in\R$, définie par  $f(x)=-\dfrac{3}{10}x+10$.
\end{exercice}

\begin{Solution}
 L'expression de la dérivée de la fonction $f$ définie par $f(x)=-\dfrac{3}{10}x+10$ est : ${\color[HTML]{f15929}\boldsymbol{f'(x)=-\dfrac{3}{10}}}$.
\end{Solution}

% @see : https://coopmaths.fr/alea?uuid=ebd89&id=1AN14-1&n=1&d=10&s=4&alea=uAlP223232323234234567891011121314151617181920212223242526272829303132333435363738394041424344454647482345678910111213141516171819202122232425262728293031323334353637&cd=0&cols=1
\needspace{10\baselineskip}
\begin{exercice}[BaremeTotal=false]


 Donner la dérivée de la fonction $f$, dérivable pour tout $x\in\R$, définie par  $f(x)=-2x^{8}$.
\end{exercice}

\begin{Solution}
 L'expression de la dérivée de la fonction $f$ définie par $f(x)=-2x^{8}$ est : ${\color[HTML]{f15929}\boldsymbol{f'(x)=-16x^{7}}}$.
\end{Solution}

% @see : https://coopmaths.fr/alea?uuid=ebd89&id=1AN14-1&n=1&d=10&s=7&alea=vjN7223232323234234567891011121314151617181920212223242526272829303132333435363738394041424344454647482345678910111213141516171819202122232425262728293031323334353637&cd=0&cols=1
\needspace{10\baselineskip}
\begin{exercice}[BaremeTotal=false]


 Donner la dérivée de la fonction $f$, dérivable pour tout $x\in\R^*$, définie par  $f(x)=\dfrac{-1}{9x}$.
\end{exercice}

\begin{Solution}
 L'expression de la dérivée de la fonction $f$ définie par $f(x)=\dfrac{-1}{9x}$ est : ${\color[HTML]{f15929}\boldsymbol{f'(x)=\dfrac{1}{9x^2}}}$.
\end{Solution}

% @see : https://coopmaths.fr/alea?uuid=a83c0&id=1AN14-4&n=1&d=10&s=1&alea=Njr3223232323234234567891011121314151617181920212223242526272829303132333435363738394041424344454647482345678910111213141516171819202122232425262728293031323334353637&cd=0&cols=1
\needspace{10\baselineskip}
\begin{exercice}[BaremeTotal=false]


 Donner l'expression de la dérivée de la fonction $f$ définie sur $\R^{*}$ par $f(x)=4x^{2}+5x+\dfrac{7}{x}$.\\
\end{exercice}

\begin{Solution}
 L'expression de la dérivée de la fonction $f$ définie par $f(x)=4x^{2}+5x+\dfrac{7}{x}$ est :\\$f'(x)={\color[HTML]{f15929}\boldsymbol{8x+5-\dfrac{7}{x^2}}}$.\\
\end{Solution}

% @see : https://coopmaths.fr/alea?uuid=a83c0&id=1AN14-4&n=1&d=10&s=4&alea=4Vpz223232323234234567891011121314151617181920212223242526272829303132333435363738394041424344454647482345678910111213141516171819202122232425262728293031323334353637&cd=1&cols=1
\needspace{10\baselineskip}
\begin{exercice}[BaremeTotal=false]


 Donner l'expression de la dérivée de la fonction $f$ définie sur $\R_+^{*}$ par $f(x)=-4x+\dfrac{3}{x}-5\sqrt{x}$.\\
\end{exercice}

\begin{Solution}
 $(-4x)^\prime=-4$\\$(\dfrac{3}{x})^\prime=-\dfrac{3}{x^2}$\\$(-5\sqrt{x})^\prime=-\dfrac{5}{2\sqrt{x}}$\\L'expression de la dérivée de la fonction $f$ définie par $f(x)=-4x+\dfrac{3}{x}-5\sqrt{x}$ est :\\$f'(x)={\color[HTML]{f15929}\boldsymbol{-4-\dfrac{3}{x^2}-\dfrac{5}{2\sqrt{x}}}}$.\\
\end{Solution}

\end{Maquette}
\clearpage
\begin{Maquette}[DM]{Niveau= ,Classe=\getdata[38]\mydata\ (v38), Date = 06/01/25  ,Theme=Exercices}


% @see : https://coopmaths.fr/alea?uuid=4c8c7&id=1AN11-3&n=1&d=10&s=2&alea=nawO22323232323423456789101112131415161718192021222324252627282930313233343536373839404142434445464748234567891011121314151617181920212223242526272829303132333435363738&cd=1&cols=2
\needspace{10\baselineskip}
\begin{exercice}[BaremeTotal=false]
\label{eleve:38}


 \begin{minipage}[t]{\linewidth} Soit $f$ une fonction dérivable sur $[-5;5]$ et $\mathcal{C}_f$ sa courbe représentative.\\On sait que  $f(-5)=-3~~$ et que $~~f'(-5)=-1$.\\Déterminer une équation de la tangente $(T)$ à la courbe $\mathcal{C}_f$ au point d'abscisse $-5$,\\sans utiliser la formule de cours de l'équation de tangente. \end{minipage}

\end{exercice}

\begin{Solution}
 \begin{minipage}[t]{\linewidth} On sait que la tangente n'est pas une droite verticale, puisque la fonction est dérivable sur l'intervalle.\\On en déduit que la tangente $(T)$ au point d'abscisse $-5$, admet une équation réduite de la forme :  $(T) : y=m x + p$. 

\medskip
$\bullet$ Détermination de $m$ :\\On sait que le nombre dérivé en $-5$ est par définition, le coefficient directeur de la tangente au point d'abscisse $-5$.\\Par conséquent, on a déjà : $m=f'(-5)=-1$.\\ On en déduit que  $(T) : y= -1 x + p$\\$\bullet$ Détermination de $p$ :\\ Pour cela, on utilise que si $f(-5)=-3~~$, alors le point $A$ de coordonnées $(-5;-3)$ appartient à $\mathcal{C}_f$ mais aussi à $(T)$.\\ On peut écrire $A(-5;-3) = \mathcal{C}_f \cap (T)$.\\ On remplace alors les coordonnées de $A(-5;-3)$ dans l'équation  $(T) : y= -1 x + p$.\\ $\begin{aligned} \phantom{\iff}&A(-5;-3)\in (T)\\ \iff& -3= -1 \times (-5)  + p\\ \iff& p=-3 + \times (-5)  \\ \iff& p=-3 -5   \\ \iff& p=-8\\\end{aligned}$\\On peut conclure que : $(T) : {\color[HTML]{f15929}\boldsymbol{y=-x-8}}$. \end{minipage}

\end{Solution}

% @see : https://coopmaths.fr/alea?uuid=0e984&id=can1F15&n=1&d=10&alea=1Alo22323232323423456789101112131415161718192021222324252627282930313233343536373839404142434445464748234567891011121314151617181920212223242526272829303132333435363738&cd=1&cols=1
\needspace{10\baselineskip}
\begin{exercice}[BaremeTotal=false]


 La courbe représente une fonction $f$ et la droite est la tangente au point d'abscisse $0$.\\
        Déterminer $f'(0)$.\begin{center}\begin{tikzpicture}[baseline,scale = 0.5]

    \tikzset{
      point/.style={
        thick,
        draw,
        cross out,
        inner sep=0pt,
        minimum width=5pt,
        minimum height=5pt,
      },
    }
    \clip (-4,-2) rectangle (12,12);
    	
	\draw[color ={black},line width = 2,>=latex,->] (-4,0)--(12,0);
	\draw[color ={black},line width = 2,>=latex,->] (0,-2)--(0,12);
	\draw[color ={black},opacity = 0.5] (-4,1)--(12,1);
	\draw[color ={black},opacity = 0.5] (-4,2)--(12,2);
	\draw[color ={black},opacity = 0.5] (-4,3)--(12,3);
	\draw[color ={black},opacity = 0.5] (-4,4)--(12,4);
	\draw[color ={black},opacity = 0.5] (-4,5)--(12,5);
	\draw[color ={black},opacity = 0.5] (-4,6)--(12,6);
	\draw[color ={black},opacity = 0.5] (-4,7)--(12,7);
	\draw[color ={black},opacity = 0.5] (-4,8)--(12,8);
	\draw[color ={black},opacity = 0.5] (-4,9)--(12,9);
	\draw[color ={black},opacity = 0.5] (-4,10)--(12,10);
	\draw[color ={black},opacity = 0.5] (-4,11)--(12,11);
	\draw[color ={black},opacity = 0.5] (-4,12)--(12,12);
	\draw[color ={black},opacity = 0.5] (-4,-1)--(12,-1);
	\draw[color ={black},opacity = 0.5] (-4,-2)--(12,-2);
	\draw[color ={black},opacity = 0.5] (1,-2)--(1,12);
	\draw[color ={black},opacity = 0.5] (2,-2)--(2,12);
	\draw[color ={black},opacity = 0.5] (3,-2)--(3,12);
	\draw[color ={black},opacity = 0.5] (4,-2)--(4,12);
	\draw[color ={black},opacity = 0.5] (5,-2)--(5,12);
	\draw[color ={black},opacity = 0.5] (6,-2)--(6,12);
	\draw[color ={black},opacity = 0.5] (7,-2)--(7,12);
	\draw[color ={black},opacity = 0.5] (8,-2)--(8,12);
	\draw[color ={black},opacity = 0.5] (9,-2)--(9,12);
	\draw[color ={black},opacity = 0.5] (10,-2)--(10,12);
	\draw[color ={black},opacity = 0.5] (11,-2)--(11,12);
	\draw[color ={black},opacity = 0.5] (12,-2)--(12,12);
	\draw[color ={black},opacity = 0.5] (-1,-2)--(-1,12);
	\draw[color ={black},opacity = 0.5] (-2,-2)--(-2,12);
	\draw[color ={black},opacity = 0.5] (-3,-2)--(-3,12);
	\draw[color ={black},opacity = 0.5] (-4,-2)--(-4,12);
	\draw[color ={gray},opacity = 0.3] (-4,0)--(12,0);
	\draw[color ={gray},opacity = 0.3] (-4,0)--(12,0);
	\draw[color ={gray},opacity = 0.3] (0,-2)--(0,12);
	\draw[color ={gray},opacity = 0.3] (0,-2)--(0,12);
	\draw[color ={black},line width = 2] (2,-0.2)--(2,0.2);
	\draw[color ={black},line width = 2] (4,-0.2)--(4,0.2);
	\draw[color ={black},line width = 2] (6,-0.2)--(6,0.2);
	\draw[color ={black},line width = 2] (8,-0.2)--(8,0.2);
	\draw[color ={black},line width = 2] (10,-0.2)--(10,0.2);
	\draw[color ={black},line width = 2] (-2,-0.2)--(-2,0.2);
	\draw[color ={black},line width = 2] (-0.2,2)--(0.2,2);
	\draw[color ={black},line width = 2] (-0.2,4)--(0.2,4);
	\draw[color ={black},line width = 2] (-0.2,6)--(0.2,6);
	\draw[color ={black},line width = 2] (-0.2,8)--(0.2,8);
	\draw[color ={black},line width = 2] (-0.2,10)--(0.2,10);
	\draw (2,-0.4) node[anchor = center] {\scriptsize \color{black}{$1$}};
	\draw (4,-0.4) node[anchor = center] {\scriptsize \color{black}{$2$}};
	\draw (6,-0.4) node[anchor = center] {\scriptsize \color{black}{$3$}};
	\draw (8,-0.4) node[anchor = center] {\scriptsize \color{black}{$4$}};
	\draw (10,-0.4) node[anchor = center] {\scriptsize \color{black}{$5$}};
	\draw (-0.5,2.1) node[anchor = center] {\footnotesize \color{black}{$1$}};
	\draw (-0.5,4.1) node[anchor = center] {\footnotesize \color{black}{$2$}};
	\draw (-0.5,6.1) node[anchor = center] {\footnotesize \color{black}{$3$}};
	\draw (-0.5,8.1) node[anchor = center] {\footnotesize \color{black}{$4$}};
	\draw (-0.5,10.1) node[anchor = center] {\footnotesize \color{black}{$5$}};
	\draw [color={black}] (-0.3,-0.3) node[anchor = center,scale=1, rotate = 0] {O};
	
	\draw[color={blue},line width = 2] (-4,0.07)--(-3.6,0.1)--(-3.2,0.14)--(-2.8,0.19)--(-2.4,0.27)--(-2,0.38)--(-1.6,0.53)--(-1.2,0.74)--(-0.8,1.03)--(-0.4,1.43)--(0,2)--(0.4,2.79)--(0.8,3.9)--(1.2,5.44)--(1.6,7.59)--(2,10.59);
	
	\draw[color={red},line width = 2] (-3.2,-3.33)--(-2.8,-2.67)--(-2.4,-2)--(-2,-1.33)--(-1.6,-0.67)--(-1.2,0)--(-0.8,0.67)--(-0.4,1.33)--(0,2)--(0.4,2.67)--(0.8,3.33)--(1.2,4)--(1.6,4.67)--(2,5.33)--(2.4,6)--(2.8,6.67)--(3.2,7.33)--(3.6,8)--(4,8.67)--(4.4,9.33)--(4.8,10)--(5.2,10.67)--(5.6,11.33)--(6,12)--(6.4,12.67)--(6.8,13.33);

\end{tikzpicture}\end{center}
\end{exercice}

\begin{Solution}
 $f'(0)$ est donné par le coefficient directeur de la tangente à la courbe au point d'abscisse $0$, soit $\dfrac{5}{3}$.
\end{Solution}

% @see : https://coopmaths.fr/alea?uuid=6f32d&id=can1F16&n=1&d=10&alea=dqSw22323232323423456789101112131415161718192021222324252627282930313233343536373839404142434445464748234567891011121314151617181920212223242526272829303132333435363738&cd=1&cols=1
\needspace{10\baselineskip}
\newpage
\begin{exercice}[BaremeTotal=false]


 On donne les représentations graphiques d'une fonction et de sa dérivée.\\
      Donner l'équation réduite de la tangente à la courbe de $f$ en $x=-2$. \\ \begin{minipage}{0.5\linewidth}
    \begin{tikzpicture}[baseline, scale=0.8]

    \tikzset{
      point/.style={
        thick,
        draw,
        cross out,
        inner sep=0pt,
        minimum width=5pt,
        minimum height=5pt,
      },
    }
    \clip (-5.5,-9.5) rectangle (5.75,7.5);
    	
	\draw[color ={black},>=latex,->] (-4,0)--(4,0);
	\draw[color ={black},>=latex,->] (0,-8)--(0,6);
	\draw[color ={black},opacity = 0.4] (-4,1)--(4,1);
	\draw[color ={black},opacity = 0.4] (-4,2)--(4,2);
	\draw[color ={black},opacity = 0.4] (-4,3)--(4,3);
	\draw[color ={black},opacity = 0.4] (-4,4)--(4,4);
	\draw[color ={black},opacity = 0.4] (-4,5)--(4,5);
	\draw[color ={black},opacity = 0.4] (-4,6)--(4,6);
	\draw[color ={black},opacity = 0.4] (-4,-1)--(4,-1);
	\draw[color ={black},opacity = 0.4] (-4,-2)--(4,-2);
	\draw[color ={black},opacity = 0.4] (-4,-3)--(4,-3);
	\draw[color ={black},opacity = 0.4] (-4,-4)--(4,-4);
	\draw[color ={black},opacity = 0.4] (-4,-5)--(4,-5);
	\draw[color ={black},opacity = 0.4] (-4,-6)--(4,-6);
	\draw[color ={black},opacity = 0.4] (-4,-7)--(4,-7);
	\draw[color ={black},opacity = 0.4] (-4,-8)--(4,-8);
	\draw[color ={black},opacity = 0.4] (1,-8)--(1,6);
	\draw[color ={black},opacity = 0.4] (2,-8)--(2,6);
	\draw[color ={black},opacity = 0.4] (3,-8)--(3,6);
	\draw[color ={black},opacity = 0.4] (4,-8)--(4,6);
	\draw[color ={black},opacity = 0.4] (-1,-8)--(-1,6);
	\draw[color ={black},opacity = 0.4] (-2,-8)--(-2,6);
	\draw[color ={black},opacity = 0.4] (-3,-8)--(-3,6);
	\draw[color ={black},opacity = 0.4] (-4,-8)--(-4,6);
	\draw[color ={gray},opacity = 0.3] (-4,0)--(4,0);
	\draw[color ={gray},opacity = 0.3] (-4,0.5)--(4,0.5);
	\draw[color ={gray},opacity = 0.3] (-4,1.5)--(4,1.5);
	\draw[color ={gray},opacity = 0.3] (-4,2.5)--(4,2.5);
	\draw[color ={gray},opacity = 0.3] (-4,3.5)--(4,3.5);
	\draw[color ={gray},opacity = 0.3] (-4,4.5)--(4,4.5);
	\draw[color ={gray},opacity = 0.3] (-4,5.5)--(4,5.5);
	\draw[color ={gray},opacity = 0.3] (-4,0)--(4,0);
	\draw[color ={gray},opacity = 0.3] (-4,-0.5)--(4,-0.5);
	\draw[color ={gray},opacity = 0.3] (-4,-1.5)--(4,-1.5);
	\draw[color ={gray},opacity = 0.3] (-4,-2.5)--(4,-2.5);
	\draw[color ={gray},opacity = 0.3] (-4,-3.5)--(4,-3.5);
	\draw[color ={gray},opacity = 0.3] (-4,-4.5)--(4,-4.5);
	\draw[color ={gray},opacity = 0.3] (-4,-5.5)--(4,-5.5);
	\draw[color ={gray},opacity = 0.3] (-4,-6.5)--(4,-6.5);
	\draw[color ={gray},opacity = 0.3] (-4,-7)--(4,-7);
	\draw[color ={gray},opacity = 0.3] (-4,-7.5)--(4,-7.5);
	\draw[color ={gray},opacity = 0.3] (-4,-8)--(4,-8);
	\draw[color ={gray},opacity = 0.3] (0,-8)--(0,6);
	\draw[color ={gray},opacity = 0.3] (0,-8)--(0,6);
	\draw[color ={black},line width = 1.2] (1,-0.13)--(1,0.13);
	\draw[color ={black},line width = 1.2] (2,-0.13)--(2,0.13);
	\draw[color ={black},line width = 1.2] (3,-0.13)--(3,0.13);
	\draw[color ={black},line width = 1.2] (4,-0.13)--(4,0.13);
	\draw[color ={black},line width = 1.2] (-1,-0.13)--(-1,0.13);
	\draw[color ={black},line width = 1.2] (-2,-0.13)--(-2,0.13);
	\draw[color ={black},line width = 1.2] (-3,-0.13)--(-3,0.13);
	\draw[color ={black},line width = 1.2] (-4,-0.13)--(-4,0.13);
	\draw[color ={black},line width = 1.2] (-0.13,1)--(0.13,1);
	\draw[color ={black},line width = 1.2] (-0.13,2)--(0.13,2);
	\draw[color ={black},line width = 1.2] (-0.13,3)--(0.13,3);
	\draw[color ={black},line width = 1.2] (-0.13,4)--(0.13,4);
	\draw[color ={black},line width = 1.2] (-0.13,5)--(0.13,5);
	\draw[color ={black},line width = 1.2] (-0.13,6)--(0.13,6);
	\draw[color ={black},line width = 1.2] (-0.13,-1)--(0.13,-1);
	\draw[color ={black},line width = 1.2] (-0.13,-2)--(0.13,-2);
	\draw[color ={black},line width = 1.2] (-0.13,-3)--(0.13,-3);
	\draw[color ={black},line width = 1.2] (-0.13,-4)--(0.13,-4);
	\draw[color ={black},line width = 1.2] (-0.13,-5)--(0.13,-5);
	\draw[color ={black},line width = 1.2] (-0.13,-6)--(0.13,-6);
	\draw[color ={black},line width = 1.2] (-0.13,-7)--(0.13,-7);
	\draw[color ={black},line width = 1.2] (-0.13,-8)--(0.13,-8);
	\draw (2,-0.4) node[anchor = center] {\scriptsize \color{black}{$2$}};
	\draw (4,-0.4) node[anchor = center] {\scriptsize \color{black}{$4$}};
	\draw (-2,-0.4) node[anchor = center] {\scriptsize \color{black}{$-2$}};
	\draw (-4,-0.4) node[anchor = center] {\scriptsize \color{black}{$-4$}};
	\draw (-0.5,2.1) node[anchor = center] {\footnotesize \color{black}{$2$}};
	\draw (-0.5,4.1) node[anchor = center] {\footnotesize \color{black}{$4$}};
	\draw (-0.5,6.1) node[anchor = center] {\footnotesize \color{black}{$6$}};
	\draw (-0.5,-1.9) node[anchor = center] {\footnotesize \color{black}{$-2$}};
	\draw (-0.5,-3.9) node[anchor = center] {\footnotesize \color{black}{$-4$}};
	\draw (-0.5,-5.9) node[anchor = center] {\footnotesize \color{black}{$-6$}};
	\draw (-0.5,-7.9) node[anchor = center] {\footnotesize \color{black}{$-8$}};
	\draw [color={black}] (-0.3,-0.3) node[anchor = center,scale=1, rotate = 0] {O};
	\draw (3,4) node[anchor = center] {\footnotesize \color{blue}{$\Large \mathcal{C}_f$}};
	
	\draw[color={blue},line width = 2] (-4,-8)--(-3.8,-6.84)--(-3.6,-5.76)--(-3.4,-4.76)--(-3.2,-3.84)--(-3,-3)--(-2.8,-2.24)--(-2.6,-1.56)--(-2.4,-0.96)--(-2.2,-0.44)--(-2,0)--(-1.8,0.36)--(-1.6,0.64)--(-1.4,0.84)--(-1.2,0.96)--(-1,1)--(-0.8,0.96)--(-0.6,0.84)--(-0.4,0.64)--(-0.2,0.36)--(0,0)--(0.2,-0.44)--(0.4,-0.96)--(0.6,-1.56)--(0.8,-2.24)--(1,-3)--(1.2,-3.84)--(1.4,-4.76)--(1.6,-5.76)--(1.8,-6.84)--(2,-8);

\end{tikzpicture}\\
    \end{minipage}
    \begin{minipage}{0.5\linewidth}
    \begin{tikzpicture}[baseline, scale=0.8]

    \tikzset{
      point/.style={
        thick,
        draw,
        cross out,
        inner sep=0pt,
        minimum width=5pt,
        minimum height=5pt,
      },
    }
    \clip (-5.5,-9.5) rectangle (6.5,7.5);
    	
	\draw[color ={black},>=latex,->] (-4,0)--(4,0);
	\draw[color ={black},>=latex,->] (0,-8)--(0,6);
	\draw[color ={black},opacity = 0.4] (-4,1)--(4,1);
	\draw[color ={black},opacity = 0.4] (-4,2)--(4,2);
	\draw[color ={black},opacity = 0.4] (-4,3)--(4,3);
	\draw[color ={black},opacity = 0.4] (-4,4)--(4,4);
	\draw[color ={black},opacity = 0.4] (-4,5)--(4,5);
	\draw[color ={black},opacity = 0.4] (-4,6)--(4,6);
	\draw[color ={black},opacity = 0.4] (-4,-1)--(4,-1);
	\draw[color ={black},opacity = 0.4] (-4,-2)--(4,-2);
	\draw[color ={black},opacity = 0.4] (-4,-3)--(4,-3);
	\draw[color ={black},opacity = 0.4] (-4,-4)--(4,-4);
	\draw[color ={black},opacity = 0.4] (-4,-5)--(4,-5);
	\draw[color ={black},opacity = 0.4] (-4,-6)--(4,-6);
	\draw[color ={black},opacity = 0.4] (-4,-7)--(4,-7);
	\draw[color ={black},opacity = 0.4] (-4,-8)--(4,-8);
	\draw[color ={black},opacity = 0.4] (1,-8)--(1,6);
	\draw[color ={black},opacity = 0.4] (2,-8)--(2,6);
	\draw[color ={black},opacity = 0.4] (3,-8)--(3,6);
	\draw[color ={black},opacity = 0.4] (4,-8)--(4,6);
	\draw[color ={black},opacity = 0.4] (-1,-8)--(-1,6);
	\draw[color ={black},opacity = 0.4] (-2,-8)--(-2,6);
	\draw[color ={black},opacity = 0.4] (-3,-8)--(-3,6);
	\draw[color ={black},opacity = 0.4] (-4,-8)--(-4,6);
	\draw[color ={gray},opacity = 0.3] (-4,0)--(4,0);
	\draw[color ={gray},opacity = 0.3] (-4,0.5)--(4,0.5);
	\draw[color ={gray},opacity = 0.3] (-4,1.5)--(4,1.5);
	\draw[color ={gray},opacity = 0.3] (-4,2.5)--(4,2.5);
	\draw[color ={gray},opacity = 0.3] (-4,3.5)--(4,3.5);
	\draw[color ={gray},opacity = 0.3] (-4,4.5)--(4,4.5);
	\draw[color ={gray},opacity = 0.3] (-4,5.5)--(4,5.5);
	\draw[color ={gray},opacity = 0.3] (-4,0)--(4,0);
	\draw[color ={gray},opacity = 0.3] (-4,-0.5)--(4,-0.5);
	\draw[color ={gray},opacity = 0.3] (-4,-1.5)--(4,-1.5);
	\draw[color ={gray},opacity = 0.3] (-4,-2.5)--(4,-2.5);
	\draw[color ={gray},opacity = 0.3] (-4,-3.5)--(4,-3.5);
	\draw[color ={gray},opacity = 0.3] (-4,-4.5)--(4,-4.5);
	\draw[color ={gray},opacity = 0.3] (-4,-5.5)--(4,-5.5);
	\draw[color ={gray},opacity = 0.3] (-4,-6.5)--(4,-6.5);
	\draw[color ={gray},opacity = 0.3] (-4,-7)--(4,-7);
	\draw[color ={gray},opacity = 0.3] (-4,-7.5)--(4,-7.5);
	\draw[color ={gray},opacity = 0.3] (-4,-8)--(4,-8);
	\draw[color ={gray},opacity = 0.3] (0,-8)--(0,6);
	\draw[color ={gray},opacity = 0.3] (0,-8)--(0,6);
	\draw[color ={black},line width = 1.2] (1,-0.13)--(1,0.13);
	\draw[color ={black},line width = 1.2] (2,-0.13)--(2,0.13);
	\draw[color ={black},line width = 1.2] (3,-0.13)--(3,0.13);
	\draw[color ={black},line width = 1.2] (4,-0.13)--(4,0.13);
	\draw[color ={black},line width = 1.2] (-1,-0.13)--(-1,0.13);
	\draw[color ={black},line width = 1.2] (-2,-0.13)--(-2,0.13);
	\draw[color ={black},line width = 1.2] (-3,-0.13)--(-3,0.13);
	\draw[color ={black},line width = 1.2] (-4,-0.13)--(-4,0.13);
	\draw[color ={black},line width = 1.2] (-0.13,1)--(0.13,1);
	\draw[color ={black},line width = 1.2] (-0.13,2)--(0.13,2);
	\draw[color ={black},line width = 1.2] (-0.13,3)--(0.13,3);
	\draw[color ={black},line width = 1.2] (-0.13,4)--(0.13,4);
	\draw[color ={black},line width = 1.2] (-0.13,5)--(0.13,5);
	\draw[color ={black},line width = 1.2] (-0.13,6)--(0.13,6);
	\draw[color ={black},line width = 1.2] (-0.13,-1)--(0.13,-1);
	\draw[color ={black},line width = 1.2] (-0.13,-2)--(0.13,-2);
	\draw[color ={black},line width = 1.2] (-0.13,-3)--(0.13,-3);
	\draw[color ={black},line width = 1.2] (-0.13,-4)--(0.13,-4);
	\draw[color ={black},line width = 1.2] (-0.13,-5)--(0.13,-5);
	\draw[color ={black},line width = 1.2] (-0.13,-6)--(0.13,-6);
	\draw[color ={black},line width = 1.2] (-0.13,-7)--(0.13,-7);
	\draw[color ={black},line width = 1.2] (-0.13,-8)--(0.13,-8);
	\draw (2,-0.4) node[anchor = center] {\scriptsize \color{black}{$2$}};
	\draw (4,-0.4) node[anchor = center] {\scriptsize \color{black}{$4$}};
	\draw (-2,-0.4) node[anchor = center] {\scriptsize \color{black}{$-2$}};
	\draw (-4,-0.4) node[anchor = center] {\scriptsize \color{black}{$-4$}};
	\draw (-0.5,2.1) node[anchor = center] {\footnotesize \color{black}{$2$}};
	\draw (-0.5,4.1) node[anchor = center] {\footnotesize \color{black}{$4$}};
	\draw (-0.5,-1.9) node[anchor = center] {\footnotesize \color{black}{$-2$}};
	\draw (-0.5,-3.9) node[anchor = center] {\footnotesize \color{black}{$-4$}};
	\draw (-0.5,-5.9) node[anchor = center] {\footnotesize \color{black}{$-6$}};
	\draw (-0.5,-7.9) node[anchor = center] {\footnotesize \color{black}{$-8$}};
	\draw [color={black}] (-0.3,-0.3) node[anchor = center,scale=1, rotate = 0] {O};
	\draw (3,4) node[anchor = center] {\footnotesize \color{red}{$\Large\mathcal{C}_f\prime$}};
	
	\draw[color={red},line width = 2] (-4,6)--(-3.8,5.6)--(-3.6,5.2)--(-3.4,4.8)--(-3.2,4.4)--(-3,4)--(-2.8,3.6)--(-2.6,3.2)--(-2.4,2.8)--(-2.2,2.4)--(-2,2)--(-1.8,1.6)--(-1.6,1.2)--(-1.4,0.8)--(-1.2,0.4)--(-1,0)--(-0.8,-0.4)--(-0.6,-0.8)--(-0.4,-1.2)--(-0.2,-1.6)--(0,-2)--(0.2,-2.4)--(0.4,-2.8)--(0.6,-3.2)--(0.8,-3.6)--(1,-4)--(1.2,-4.4)--(1.4,-4.8)--(1.6,-5.2)--(1.8,-5.6)--(2,-6)--(2.2,-6.4)--(2.4,-6.8)--(2.6,-7.2)--(2.8,-7.6)--(3,-8)--(3.2,-8.4)--(3.4,-8.8);

\end{tikzpicture}\\
    \end{minipage}
    
\end{exercice}

\begin{Solution}
 L'équation réduite de la tangente au point d'abscisse $-2$ est  : $y=f'(-2)(x-(-2))+f(-2)$.\\
      On lit graphiquement $f(-2)=0$ et $f'(-2)=2$.\\
      L'équation réduite de la tangente est donc donnée par :
      $y=2(x+2)+0$, soit $y=2x+4$.
\end{Solution}

% @see : https://coopmaths.fr/alea?uuid=a1ba2&id=can1F14&n=1&d=10&alea=rkIp22323232323423456789101112131415161718192021222324252627282930313233343536373839404142434445464748234567891011121314151617181920212223242526272829303132333435363738&cd=1&cols=1
\needspace{10\baselineskip}
\begin{exercice}[BaremeTotal=false]


 Soit $f$ la fonction définie sur $\mathbb{R}$ par : $f(x)= -3x^2+x-3$.\\
        En calculant la limite du taux d'accroissement, déterminer $f'(2)$.
\end{exercice}

\begin{Solution}
 $f$ est une fonction polynôme du second degré de la forme $f(x)=ax^2+bx+c$.\\
    La fonction dérivée est donnée par la somme des dérivées des fonctions $u$ et $v$ définies par $u(x)=-3x^2$ et $v(x)=x-3$.\\
     Comme $u'(x)=-6x$ et $v'(x)=1$, on obtient  $f'(x)=-6x+1$. \\
     Ainsi, $f'(2)=-6\times 2+1=-11$.
\end{Solution}

% @see : https://coopmaths.fr/alea?uuid=3c690&id=can1F13&n=1&d=10&alea=EUDu22323232323423456789101112131415161718192021222324252627282930313233343536373839404142434445464748234567891011121314151617181920212223242526272829303132333435363738&cd=1&cols=1
\needspace{10\baselineskip}
\begin{exercice}[BaremeTotal=false]


 Déterminer le coefficient directeur de la tangente à la courbe représentative de la fonction carré au point d'abscisse $-13$.    
\end{exercice}

\begin{Solution}
 Le coefficient directeur de la tangente au point d'abscisse $a$ est donné par le nombre dérivé $f'(a)$.\\
        La fonction $f$ définie par $f(x)=x^2$ a pour fonction dérivée la fonction $f'$ définie par $f'(x)=2x$.\\
        Comme $f'(-13)=2\times (-13)=-26$, le coefficient directeur de la tangente au point d'abscisse $-13$ est : $-26$. 
\end{Solution}

% @see : https://coopmaths.fr/alea?uuid=ebd89&id=1AN14-1&n=1&d=10&s=3&alea=ocYr22323232323423456789101112131415161718192021222324252627282930313233343536373839404142434445464748234567891011121314151617181920212223242526272829303132333435363738&cd=0&cols=1
\needspace{10\baselineskip}
\begin{exercice}[BaremeTotal=false]


 Donner la dérivée de la fonction $f$, dérivable pour tout $x\in\R$, définie par  $f(x)=\dfrac{8}{9}x+6$.
\end{exercice}

\begin{Solution}
 L'expression de la dérivée de la fonction $f$ définie par $f(x)=\dfrac{8}{9}x+6$ est : ${\color[HTML]{f15929}\boldsymbol{f'(x)=\dfrac{8}{9}}}$.
\end{Solution}

% @see : https://coopmaths.fr/alea?uuid=ebd89&id=1AN14-1&n=1&d=10&s=4&alea=uAlP22323232323423456789101112131415161718192021222324252627282930313233343536373839404142434445464748234567891011121314151617181920212223242526272829303132333435363738&cd=0&cols=1
\needspace{10\baselineskip}
\begin{exercice}[BaremeTotal=false]


 Donner la dérivée de la fonction $f$, dérivable pour tout $x\in\R$, définie par  $f(x)=6x^{7}$.
\end{exercice}

\begin{Solution}
 L'expression de la dérivée de la fonction $f$ définie par $f(x)=6x^{7}$ est : ${\color[HTML]{f15929}\boldsymbol{f'(x)=42x^{6}}}$.
\end{Solution}

% @see : https://coopmaths.fr/alea?uuid=ebd89&id=1AN14-1&n=1&d=10&s=7&alea=vjN722323232323423456789101112131415161718192021222324252627282930313233343536373839404142434445464748234567891011121314151617181920212223242526272829303132333435363738&cd=0&cols=1
\needspace{10\baselineskip}
\begin{exercice}[BaremeTotal=false]


 Donner la dérivée de la fonction $f$, dérivable pour tout $x\in\R^*$, définie par  $f(x)=\dfrac{1}{5x}$.
\end{exercice}

\begin{Solution}
 L'expression de la dérivée de la fonction $f$ définie par $f(x)=\dfrac{1}{5x}$ est : ${\color[HTML]{f15929}\boldsymbol{f'(x)=-\dfrac{1}{5x^2}}}$.
\end{Solution}

% @see : https://coopmaths.fr/alea?uuid=a83c0&id=1AN14-4&n=1&d=10&s=1&alea=Njr322323232323423456789101112131415161718192021222324252627282930313233343536373839404142434445464748234567891011121314151617181920212223242526272829303132333435363738&cd=0&cols=1
\needspace{10\baselineskip}
\begin{exercice}[BaremeTotal=false]


 Donner l'expression de la dérivée de la fonction $f$ définie sur $\R^{*}$ par $f(x)=\dfrac{4}{x}+5x^{2}$.\\
\end{exercice}

\begin{Solution}
 L'expression de la dérivée de la fonction $f$ définie par $f(x)=\dfrac{4}{x}+5x^{2}$ est :\\$f'(x)={\color[HTML]{f15929}\boldsymbol{-\dfrac{4}{x^2}+10x}}$.\\
\end{Solution}

% @see : https://coopmaths.fr/alea?uuid=a83c0&id=1AN14-4&n=1&d=10&s=4&alea=4Vpz22323232323423456789101112131415161718192021222324252627282930313233343536373839404142434445464748234567891011121314151617181920212223242526272829303132333435363738&cd=1&cols=1
\needspace{10\baselineskip}
\begin{exercice}[BaremeTotal=false]


 Donner l'expression de la dérivée de la fonction $f$ définie sur $\R_+^{*}$ par $f(x)=-5\sqrt{x}+\dfrac{3}{x}+3x$.\\
\end{exercice}

\begin{Solution}
 $(-5\sqrt{x})^\prime=-\dfrac{5}{2\sqrt{x}}$\\$(\dfrac{3}{x})^\prime=-\dfrac{3}{x^2}$\\$(3x)^\prime=3$\\L'expression de la dérivée de la fonction $f$ définie par $f(x)=-5\sqrt{x}+\dfrac{3}{x}+3x$ est :\\$f'(x)={\color[HTML]{f15929}\boldsymbol{-\dfrac{5}{2\sqrt{x}}-\dfrac{3}{x^2}+3}}$.\\
\end{Solution}

\end{Maquette}
\clearpage
\begin{Maquette}[DM]{Niveau= ,Classe=\getdata[39]\mydata\ (v39), Date = 06/01/25  ,Theme=Exercices}


% @see : https://coopmaths.fr/alea?uuid=4c8c7&id=1AN11-3&n=1&d=10&s=2&alea=nawO2232323232342345678910111213141516171819202122232425262728293031323334353637383940414243444546474823456789101112131415161718192021222324252627282930313233343536373839&cd=1&cols=2
\needspace{10\baselineskip}
\begin{exercice}[BaremeTotal=false]
\label{eleve:39}


 \begin{minipage}[t]{\linewidth} Soit $f$ une fonction dérivable sur $[-5;5]$ et $\mathcal{C}_f$ sa courbe représentative.\\On sait que  $f(-2)=5~~$ et que $~~f'(-2)=2$.\\Déterminer une équation de la tangente $(T)$ à la courbe $\mathcal{C}_f$ au point d'abscisse $-2$,\\sans utiliser la formule de cours de l'équation de tangente. \end{minipage}

\end{exercice}

\begin{Solution}
 \begin{minipage}[t]{\linewidth} On sait que la tangente n'est pas une droite verticale, puisque la fonction est dérivable sur l'intervalle.\\On en déduit que la tangente $(T)$ au point d'abscisse $-2$, admet une équation réduite de la forme :  $(T) : y=m x + p$. 

\medskip
$\bullet$ Détermination de $m$ :\\On sait que le nombre dérivé en $-2$ est par définition, le coefficient directeur de la tangente au point d'abscisse $-2$.\\Par conséquent, on a déjà : $m=f'(-2)=2$.\\ On en déduit que  $(T) : y= 2 x + p$\\$\bullet$ Détermination de $p$ :\\ Pour cela, on utilise que si $f(-2)=5~~$, alors le point $A$ de coordonnées $(-2;5)$ appartient à $\mathcal{C}_f$ mais aussi à $(T)$.\\ On peut écrire $A(-2;5) = \mathcal{C}_f \cap (T)$.\\ On remplace alors les coordonnées de $A(-2;5)$ dans l'équation  $(T) : y= 2 x + p$.\\ $\begin{aligned} \phantom{\iff}&A(-2;5)\in (T)\\ \iff& 5= 2 \times (-2)  + p\\ \iff& p=5 -2 \times (-2)  \\ \iff& p=5 +4   \\ \iff& p=9\\\end{aligned}$\\On peut conclure que : $(T) : {\color[HTML]{f15929}\boldsymbol{y=2x+9}}$. \end{minipage}

\end{Solution}

% @see : https://coopmaths.fr/alea?uuid=0e984&id=can1F15&n=1&d=10&alea=1Alo2232323232342345678910111213141516171819202122232425262728293031323334353637383940414243444546474823456789101112131415161718192021222324252627282930313233343536373839&cd=1&cols=1
\needspace{10\baselineskip}
\begin{exercice}[BaremeTotal=false]


 La courbe représente une fonction $f$ et la droite est la tangente au point d'abscisse $0$.\\
        Déterminer $f'(0)$.\begin{center}\begin{tikzpicture}[baseline,scale = 0.5]

    \tikzset{
      point/.style={
        thick,
        draw,
        cross out,
        inner sep=0pt,
        minimum width=5pt,
        minimum height=5pt,
      },
    }
    \clip (-10,-2) rectangle (4,12);
    	
	\draw[color ={black},line width = 2,>=latex,->] (-10,0)--(4,0);
	\draw[color ={black},line width = 2,>=latex,->] (0,-2)--(0,12);
	\draw[color ={black},opacity = 0.5] (-10,1)--(4,1);
	\draw[color ={black},opacity = 0.5] (-10,2)--(4,2);
	\draw[color ={black},opacity = 0.5] (-10,3)--(4,3);
	\draw[color ={black},opacity = 0.5] (-10,4)--(4,4);
	\draw[color ={black},opacity = 0.5] (-10,5)--(4,5);
	\draw[color ={black},opacity = 0.5] (-10,6)--(4,6);
	\draw[color ={black},opacity = 0.5] (-10,7)--(4,7);
	\draw[color ={black},opacity = 0.5] (-10,8)--(4,8);
	\draw[color ={black},opacity = 0.5] (-10,9)--(4,9);
	\draw[color ={black},opacity = 0.5] (-10,10)--(4,10);
	\draw[color ={black},opacity = 0.5] (-10,11)--(4,11);
	\draw[color ={black},opacity = 0.5] (-10,12)--(4,12);
	\draw[color ={black},opacity = 0.5] (-10,-1)--(4,-1);
	\draw[color ={black},opacity = 0.5] (-10,-2)--(4,-2);
	\draw[color ={black},opacity = 0.5] (1,-2)--(1,12);
	\draw[color ={black},opacity = 0.5] (2,-2)--(2,12);
	\draw[color ={black},opacity = 0.5] (3,-2)--(3,12);
	\draw[color ={black},opacity = 0.5] (4,-2)--(4,12);
	\draw[color ={black},opacity = 0.5] (-1,-2)--(-1,12);
	\draw[color ={black},opacity = 0.5] (-2,-2)--(-2,12);
	\draw[color ={black},opacity = 0.5] (-3,-2)--(-3,12);
	\draw[color ={black},opacity = 0.5] (-4,-2)--(-4,12);
	\draw[color ={black},opacity = 0.5] (-5,-2)--(-5,12);
	\draw[color ={black},opacity = 0.5] (-6,-2)--(-6,12);
	\draw[color ={black},opacity = 0.5] (-7,-2)--(-7,12);
	\draw[color ={black},opacity = 0.5] (-8,-2)--(-8,12);
	\draw[color ={black},opacity = 0.5] (-9,-2)--(-9,12);
	\draw[color ={black},opacity = 0.5] (-10,-2)--(-10,12);
	\draw[color ={gray},opacity = 0.3] (-10,0)--(4,0);
	\draw[color ={gray},opacity = 0.3] (-10,0)--(4,0);
	\draw[color ={gray},opacity = 0.3] (0,-2)--(0,12);
	\draw[color ={gray},opacity = 0.3] (0,-2)--(0,12);
	\draw[color ={gray},opacity = 0.3] (-5,-2)--(-5,12);
	\draw[color ={gray},opacity = 0.3] (-6,-2)--(-6,12);
	\draw[color ={gray},opacity = 0.3] (-7,-2)--(-7,12);
	\draw[color ={gray},opacity = 0.3] (-8,-2)--(-8,12);
	\draw[color ={gray},opacity = 0.3] (-9,-2)--(-9,12);
	\draw[color ={gray},opacity = 0.3] (-10,-2)--(-10,12);
	\draw[color ={black},line width = 2] (2,-0.2)--(2,0.2);
	\draw[color ={black},line width = 2] (-2,-0.2)--(-2,0.2);
	\draw[color ={black},line width = 2] (-4,-0.2)--(-4,0.2);
	\draw[color ={black},line width = 2] (-6,-0.2)--(-6,0.2);
	\draw[color ={black},line width = 2] (-8,-0.2)--(-8,0.2);
	\draw[color ={black},line width = 2] (-0.2,2)--(0.2,2);
	\draw[color ={black},line width = 2] (-0.2,4)--(0.2,4);
	\draw[color ={black},line width = 2] (-0.2,6)--(0.2,6);
	\draw[color ={black},line width = 2] (-0.2,8)--(0.2,8);
	\draw[color ={black},line width = 2] (-0.2,10)--(0.2,10);
	\draw (2,-0.4) node[anchor = center] {\scriptsize \color{black}{$1$}};
	\draw (-2,-0.4) node[anchor = center] {\scriptsize \color{black}{$-1$}};
	\draw (-4,-0.4) node[anchor = center] {\scriptsize \color{black}{$-2$}};
	\draw (-6,-0.4) node[anchor = center] {\scriptsize \color{black}{$-3$}};
	\draw (-8,-0.4) node[anchor = center] {\scriptsize \color{black}{$-4$}};
	\draw (-0.5,2.1) node[anchor = center] {\footnotesize \color{black}{$1$}};
	\draw (-0.5,4.1) node[anchor = center] {\footnotesize \color{black}{$2$}};
	\draw (-0.5,6.1) node[anchor = center] {\footnotesize \color{black}{$3$}};
	\draw (-0.5,8.1) node[anchor = center] {\footnotesize \color{black}{$4$}};
	\draw (-0.5,10.1) node[anchor = center] {\footnotesize \color{black}{$5$}};
	\draw [color={black}] (-0.3,-0.3) node[anchor = center,scale=1, rotate = 0] {O};
	
	\draw[color={blue},line width = 2] (-3.6,12.1)--(-3.2,9.91)--(-2.8,8.11)--(-2.4,6.64)--(-2,5.44)--(-1.6,4.45)--(-1.2,3.64)--(-0.8,2.98)--(-0.4,2.44)--(0,2)--(0.4,1.64)--(0.8,1.34)--(1.2,1.1)--(1.6,0.9)--(2,0.74)--(2.4,0.6)--(2.8,0.49)--(3.2,0.4)--(3.6,0.33)--(4,0.27);
	
	\draw[color={red},line width = 2] (-10,12)--(-9.6,11.6)--(-9.2,11.2)--(-8.8,10.8)--(-8.4,10.4)--(-8,10)--(-7.6,9.6)--(-7.2,9.2)--(-6.8,8.8)--(-6.4,8.4)--(-6,8)--(-5.6,7.6)--(-5.2,7.2)--(-4.8,6.8)--(-4.4,6.4)--(-4,6)--(-3.6,5.6)--(-3.2,5.2)--(-2.8,4.8)--(-2.4,4.4)--(-2,4)--(-1.6,3.6)--(-1.2,3.2)--(-0.8,2.8)--(-0.4,2.4)--(0,2)--(0.4,1.6)--(0.8,1.2)--(1.2,0.8)--(1.6,0.4)--(2,0)--(2.4,-0.4)--(2.8,-0.8)--(3.2,-1.2)--(3.6,-1.6)--(4,-2);

\end{tikzpicture}\end{center}
\end{exercice}

\begin{Solution}
 $f'(0)$ est donné par le coefficient directeur de la tangente à la courbe au point d'abscisse $0$, soit $-1=-1$.
\end{Solution}

% @see : https://coopmaths.fr/alea?uuid=6f32d&id=can1F16&n=1&d=10&alea=dqSw2232323232342345678910111213141516171819202122232425262728293031323334353637383940414243444546474823456789101112131415161718192021222324252627282930313233343536373839&cd=1&cols=1
\needspace{10\baselineskip}
\newpage
\begin{exercice}[BaremeTotal=false]


 On donne les représentations graphiques d'une fonction et de sa dérivée.\\
        Donner l'équation réduite de la tangente à la courbe de $f$ en $x=2$. \\ \begin{minipage}{0.5\linewidth}
    \begin{tikzpicture}[baseline, scale=0.8]

    \tikzset{
      point/.style={
        thick,
        draw,
        cross out,
        inner sep=0pt,
        minimum width=5pt,
        minimum height=5pt,
      },
    }
    \clip (-4.5,-4.5) rectangle (5.75,13.5);
    	
	\draw[color ={black},>=latex,->] (-3,0)--(3,0);
	\draw[color ={black},>=latex,->] (0,-3)--(0,12);
	\draw[color ={black},opacity = 0.4] (-3,1)--(3,1);
	\draw[color ={black},opacity = 0.4] (-3,2)--(3,2);
	\draw[color ={black},opacity = 0.4] (-3,3)--(3,3);
	\draw[color ={black},opacity = 0.4] (-3,4)--(3,4);
	\draw[color ={black},opacity = 0.4] (-3,5)--(3,5);
	\draw[color ={black},opacity = 0.4] (-3,6)--(3,6);
	\draw[color ={black},opacity = 0.4] (-3,7)--(3,7);
	\draw[color ={black},opacity = 0.4] (-3,8)--(3,8);
	\draw[color ={black},opacity = 0.4] (-3,9)--(3,9);
	\draw[color ={black},opacity = 0.4] (-3,10)--(3,10);
	\draw[color ={black},opacity = 0.4] (-3,11)--(3,11);
	\draw[color ={black},opacity = 0.4] (-3,12)--(3,12);
	\draw[color ={black},opacity = 0.4] (-3,-1)--(3,-1);
	\draw[color ={black},opacity = 0.4] (-3,-2)--(3,-2);
	\draw[color ={black},opacity = 0.4] (-3,-3)--(3,-3);
	\draw[color ={black},opacity = 0.4] (1,-3)--(1,12);
	\draw[color ={black},opacity = 0.4] (2,-3)--(2,12);
	\draw[color ={black},opacity = 0.4] (3,-3)--(3,12);
	\draw[color ={black},opacity = 0.4] (-1,-3)--(-1,12);
	\draw[color ={black},opacity = 0.4] (-2,-3)--(-2,12);
	\draw[color ={black},opacity = 0.4] (-3,-3)--(-3,12);
	\draw[color ={gray},opacity = 0.3] (-3,0)--(3,0);
	\draw[color ={gray},opacity = 0.3] (-3,0.5)--(3,0.5);
	\draw[color ={gray},opacity = 0.3] (-3,1.5)--(3,1.5);
	\draw[color ={gray},opacity = 0.3] (-3,2.5)--(3,2.5);
	\draw[color ={gray},opacity = 0.3] (-3,3.5)--(3,3.5);
	\draw[color ={gray},opacity = 0.3] (-3,4.5)--(3,4.5);
	\draw[color ={gray},opacity = 0.3] (-3,5.5)--(3,5.5);
	\draw[color ={gray},opacity = 0.3] (-3,6.5)--(3,6.5);
	\draw[color ={gray},opacity = 0.3] (-3,7.5)--(3,7.5);
	\draw[color ={gray},opacity = 0.3] (-3,8.5)--(3,8.5);
	\draw[color ={gray},opacity = 0.3] (-3,9.5)--(3,9.5);
	\draw[color ={gray},opacity = 0.3] (-3,10.5)--(3,10.5);
	\draw[color ={gray},opacity = 0.3] (-3,11.5)--(3,11.5);
	\draw[color ={gray},opacity = 0.3] (-3,0)--(3,0);
	\draw[color ={gray},opacity = 0.3] (-3,-0.5)--(3,-0.5);
	\draw[color ={gray},opacity = 0.3] (-3,-1.5)--(3,-1.5);
	\draw[color ={gray},opacity = 0.3] (-3,-2.5)--(3,-2.5);
	\draw[color ={gray},opacity = 0.3] (0,-3)--(0,12);
	\draw[color ={gray},opacity = 0.3] (0,-3)--(0,12);
	\draw[color ={black},line width = 1.2] (1,-0.13)--(1,0.13);
	\draw[color ={black},line width = 1.2] (2,-0.13)--(2,0.13);
	\draw[color ={black},line width = 1.2] (3,-0.13)--(3,0.13);
	\draw[color ={black},line width = 1.2] (-1,-0.13)--(-1,0.13);
	\draw[color ={black},line width = 1.2] (-2,-0.13)--(-2,0.13);
	\draw[color ={black},line width = 1.2] (-3,-0.13)--(-3,0.13);
	\draw[color ={black},line width = 1.2] (-0.13,1)--(0.13,1);
	\draw[color ={black},line width = 1.2] (-0.13,2)--(0.13,2);
	\draw[color ={black},line width = 1.2] (-0.13,3)--(0.13,3);
	\draw[color ={black},line width = 1.2] (-0.13,4)--(0.13,4);
	\draw[color ={black},line width = 1.2] (-0.13,5)--(0.13,5);
	\draw[color ={black},line width = 1.2] (-0.13,6)--(0.13,6);
	\draw[color ={black},line width = 1.2] (-0.13,7)--(0.13,7);
	\draw[color ={black},line width = 1.2] (-0.13,8)--(0.13,8);
	\draw[color ={black},line width = 1.2] (-0.13,9)--(0.13,9);
	\draw[color ={black},line width = 1.2] (-0.13,10)--(0.13,10);
	\draw[color ={black},line width = 1.2] (-0.13,11)--(0.13,11);
	\draw[color ={black},line width = 1.2] (-0.13,12)--(0.13,12);
	\draw[color ={black},line width = 1.2] (-0.13,-1)--(0.13,-1);
	\draw[color ={black},line width = 1.2] (-0.13,-2)--(0.13,-2);
	\draw[color ={black},line width = 1.2] (-0.13,-3)--(0.13,-3);
	\draw (3,-0.4) node[anchor = center] {\scriptsize \color{black}{$3$}};
	\draw (-3,-0.4) node[anchor = center] {\scriptsize \color{black}{$-3$}};
	\draw (-0.5,3.1) node[anchor = center] {\footnotesize \color{black}{$3$}};
	\draw (-0.5,6.1) node[anchor = center] {\footnotesize \color{black}{$6$}};
	\draw (-0.5,9.1) node[anchor = center] {\footnotesize \color{black}{$9$}};
	\draw (-0.5,-2.9) node[anchor = center] {\footnotesize \color{black}{$-3$}};
	\draw [color={black}] (-0.3,-0.3) node[anchor = center,scale=1, rotate = 0] {O};
	\draw (3,10) node[anchor = center] {\footnotesize \color{blue}{$\Large \mathcal{C}_f$}};
	
	\draw[color={blue},line width = 2] (-3,8)--(-2.8,6.84)--(-2.6,5.76)--(-2.4,4.76)--(-2.2,3.84)--(-2,3)--(-1.8,2.24)--(-1.6,1.56)--(-1.4,0.96)--(-1.2,0.44)--(-1,0)--(-0.8,-0.36)--(-0.6,-0.64)--(-0.4,-0.84)--(-0.2,-0.96)--(0,-1)--(0.2,-0.96)--(0.4,-0.84)--(0.6,-0.64)--(0.8,-0.36)--(1,0)--(1.2,0.44)--(1.4,0.96)--(1.6,1.56)--(1.8,2.24)--(2,3)--(2.2,3.84)--(2.4,4.76)--(2.6,5.76)--(2.8,6.84)--(3,8);

\end{tikzpicture}\\
    \end{minipage}
    \begin{minipage}{0.5\linewidth}
    \begin{tikzpicture}[baseline, scale=0.8]

    \tikzset{
      point/.style={
        thick,
        draw,
        cross out,
        inner sep=0pt,
        minimum width=5pt,
        minimum height=5pt,
      },
    }
    \clip (-4.5,-6.5) rectangle (6.5,9.5);
    	
	\draw[color ={black},>=latex,->] (-3,0)--(3,0);
	\draw[color ={black},>=latex,->] (0,-5)--(0,8);
	\draw[color ={black},opacity = 0.4] (-3,1)--(3,1);
	\draw[color ={black},opacity = 0.4] (-3,2)--(3,2);
	\draw[color ={black},opacity = 0.4] (-3,3)--(3,3);
	\draw[color ={black},opacity = 0.4] (-3,4)--(3,4);
	\draw[color ={black},opacity = 0.4] (-3,5)--(3,5);
	\draw[color ={black},opacity = 0.4] (-3,6)--(3,6);
	\draw[color ={black},opacity = 0.4] (-3,7)--(3,7);
	\draw[color ={black},opacity = 0.4] (-3,8)--(3,8);
	\draw[color ={black},opacity = 0.4] (-3,-1)--(3,-1);
	\draw[color ={black},opacity = 0.4] (-3,-2)--(3,-2);
	\draw[color ={black},opacity = 0.4] (-3,-3)--(3,-3);
	\draw[color ={black},opacity = 0.4] (-3,-4)--(3,-4);
	\draw[color ={black},opacity = 0.4] (-3,-5)--(3,-5);
	\draw[color ={black},opacity = 0.4] (1,-5)--(1,8);
	\draw[color ={black},opacity = 0.4] (2,-5)--(2,8);
	\draw[color ={black},opacity = 0.4] (3,-5)--(3,8);
	\draw[color ={black},opacity = 0.4] (-1,-5)--(-1,8);
	\draw[color ={black},opacity = 0.4] (-2,-5)--(-2,8);
	\draw[color ={black},opacity = 0.4] (-3,-5)--(-3,8);
	\draw[color ={gray},opacity = 0.3] (-3,0)--(3,0);
	\draw[color ={gray},opacity = 0.3] (-3,0.5)--(3,0.5);
	\draw[color ={gray},opacity = 0.3] (-3,1.5)--(3,1.5);
	\draw[color ={gray},opacity = 0.3] (-3,2.5)--(3,2.5);
	\draw[color ={gray},opacity = 0.3] (-3,3.5)--(3,3.5);
	\draw[color ={gray},opacity = 0.3] (-3,4.5)--(3,4.5);
	\draw[color ={gray},opacity = 0.3] (-3,5.5)--(3,5.5);
	\draw[color ={gray},opacity = 0.3] (-3,6.5)--(3,6.5);
	\draw[color ={gray},opacity = 0.3] (-3,7.5)--(3,7.5);
	\draw[color ={gray},opacity = 0.3] (-3,0)--(3,0);
	\draw[color ={gray},opacity = 0.3] (-3,-0.5)--(3,-0.5);
	\draw[color ={gray},opacity = 0.3] (-3,-1.5)--(3,-1.5);
	\draw[color ={gray},opacity = 0.3] (-3,-2.5)--(3,-2.5);
	\draw[color ={gray},opacity = 0.3] (-3,-3.5)--(3,-3.5);
	\draw[color ={gray},opacity = 0.3] (-3,-4.5)--(3,-4.5);
	\draw[color ={gray},opacity = 0.3] (0,-5)--(0,8);
	\draw[color ={gray},opacity = 0.3] (0,-5)--(0,8);
	\draw[color ={black},line width = 1.2] (1,-0.13)--(1,0.13);
	\draw[color ={black},line width = 1.2] (2,-0.13)--(2,0.13);
	\draw[color ={black},line width = 1.2] (3,-0.13)--(3,0.13);
	\draw[color ={black},line width = 1.2] (-1,-0.13)--(-1,0.13);
	\draw[color ={black},line width = 1.2] (-2,-0.13)--(-2,0.13);
	\draw[color ={black},line width = 1.2] (-3,-0.13)--(-3,0.13);
	\draw[color ={black},line width = 1.2] (-0.13,1)--(0.13,1);
	\draw[color ={black},line width = 1.2] (-0.13,2)--(0.13,2);
	\draw[color ={black},line width = 1.2] (-0.13,3)--(0.13,3);
	\draw[color ={black},line width = 1.2] (-0.13,4)--(0.13,4);
	\draw[color ={black},line width = 1.2] (-0.13,5)--(0.13,5);
	\draw[color ={black},line width = 1.2] (-0.13,6)--(0.13,6);
	\draw[color ={black},line width = 1.2] (-0.13,7)--(0.13,7);
	\draw[color ={black},line width = 1.2] (-0.13,8)--(0.13,8);
	\draw[color ={black},line width = 1.2] (-0.13,9)--(0.13,9);
	\draw[color ={black},line width = 1.2] (-0.13,10)--(0.13,10);
	\draw[color ={black},line width = 1.2] (-0.13,11)--(0.13,11);
	\draw[color ={black},line width = 1.2] (-0.13,12)--(0.13,12);
	\draw[color ={black},line width = 1.2] (-0.13,-1)--(0.13,-1);
	\draw[color ={black},line width = 1.2] (-0.13,-2)--(0.13,-2);
	\draw[color ={black},line width = 1.2] (-0.13,-3)--(0.13,-3);
	\draw (3,-0.4) node[anchor = center] {\scriptsize \color{black}{$3$}};
	\draw (-3,-0.4) node[anchor = center] {\scriptsize \color{black}{$-3$}};
	\draw (-0.5,3.1) node[anchor = center] {\footnotesize \color{black}{$3$}};
	\draw (-0.5,6.1) node[anchor = center] {\footnotesize \color{black}{$6$}};
	\draw (-0.5,-2.9) node[anchor = center] {\footnotesize \color{black}{$-3$}};
	\draw [color={black}] (-0.3,-0.3) node[anchor = center,scale=1, rotate = 0] {O};
	\draw (3,6) node[anchor = center] {\footnotesize \color{red}{$\Large\mathcal{C}_f\prime$}};
	
	\draw[color={red},line width = 2] (-2.8,-5.6)--(-2.6,-5.2)--(-2.4,-4.8)--(-2.2,-4.4)--(-2,-4)--(-1.8,-3.6)--(-1.6,-3.2)--(-1.4,-2.8)--(-1.2,-2.4)--(-1,-2)--(-0.8,-1.6)--(-0.6,-1.2)--(-0.4,-0.8)--(-0.2,-0.4)--(0,0)--(0.2,0.4)--(0.4,0.8)--(0.6,1.2)--(0.8,1.6)--(1,2)--(1.2,2.4)--(1.4,2.8)--(1.6,3.2)--(1.8,3.6)--(2,4)--(2.2,4.4)--(2.4,4.8)--(2.6,5.2)--(2.8,5.6)--(3,6);

\end{tikzpicture}\\
    \end{minipage}
    
\end{exercice}

\begin{Solution}
 L'équation réduite de la tangente au point d'abscisse $2$ est  : $y=f'(2)(x-2)+f(2)$.\\
        On lit graphiquement $f(2)=3$ et $f'(2)=4$.\\
        L'équation réduite de la tangente est donc donnée par :
        $y=4(x-2)+3$, soit $y=4x-5$.
\end{Solution}

% @see : https://coopmaths.fr/alea?uuid=a1ba2&id=can1F14&n=1&d=10&alea=rkIp2232323232342345678910111213141516171819202122232425262728293031323334353637383940414243444546474823456789101112131415161718192021222324252627282930313233343536373839&cd=1&cols=1
\needspace{10\baselineskip}
\begin{exercice}[BaremeTotal=false]


 Soit $f$ la fonction définie sur $\mathbb{R}$ par : $f(x)= -4x^2-x-7$.\\
        En calculant la limite du taux d'accroissement, déterminer $f'(2)$.
\end{exercice}

\begin{Solution}
 $f$ est une fonction polynôme du second degré de la forme $f(x)=ax^2+bx+c$.\\
    La fonction dérivée est donnée par la somme des dérivées des fonctions $u$ et $v$ définies par $u(x)=-4x^2$ et $v(x)=-x-7$.\\
     Comme $u'(x)=-8x$ et $v'(x)=-1$, on obtient  $f'(x)=-8x-1$. \\
     Ainsi, $f'(2)=-8\times 2-1=-17$.
\end{Solution}

% @see : https://coopmaths.fr/alea?uuid=3c690&id=can1F13&n=1&d=10&alea=EUDu2232323232342345678910111213141516171819202122232425262728293031323334353637383940414243444546474823456789101112131415161718192021222324252627282930313233343536373839&cd=1&cols=1
\needspace{10\baselineskip}
\begin{exercice}[BaremeTotal=false]


 Déterminer le coefficient directeur de la tangente à la courbe représentative de la fonction carré au point d'abscisse $-12$.    
\end{exercice}

\begin{Solution}
 Le coefficient directeur de la tangente au point d'abscisse $a$ est donné par le nombre dérivé $f'(a)$.\\
        La fonction $f$ définie par $f(x)=x^2$ a pour fonction dérivée la fonction $f'$ définie par $f'(x)=2x$.\\
        Comme $f'(-12)=2\times (-12)=-24$, le coefficient directeur de la tangente au point d'abscisse $-12$ est : $-24$. 
\end{Solution}

% @see : https://coopmaths.fr/alea?uuid=ebd89&id=1AN14-1&n=1&d=10&s=3&alea=ocYr2232323232342345678910111213141516171819202122232425262728293031323334353637383940414243444546474823456789101112131415161718192021222324252627282930313233343536373839&cd=0&cols=1
\needspace{10\baselineskip}
\begin{exercice}[BaremeTotal=false]


 Donner la dérivée de la fonction $f$, dérivable pour tout $x\in\R$, définie par  $f(x)=\dfrac{-x+2}{5}$.
\end{exercice}

\begin{Solution}
 L'expression de la dérivée de la fonction $f$ définie par $f(x)=\dfrac{-x+2}{5}$ est : ${\color[HTML]{f15929}\boldsymbol{f'(x)=-\dfrac{1}{5}}}$.
\end{Solution}

% @see : https://coopmaths.fr/alea?uuid=ebd89&id=1AN14-1&n=1&d=10&s=4&alea=uAlP2232323232342345678910111213141516171819202122232425262728293031323334353637383940414243444546474823456789101112131415161718192021222324252627282930313233343536373839&cd=0&cols=1
\needspace{10\baselineskip}
\begin{exercice}[BaremeTotal=false]


 Donner la dérivée de la fonction $f$, dérivable pour tout $x\in\R$, définie par  $f(x)=-6x^{14}$.
\end{exercice}

\begin{Solution}
 L'expression de la dérivée de la fonction $f$ définie par $f(x)=-6x^{14}$ est : ${\color[HTML]{f15929}\boldsymbol{f'(x)=-84x^{13}}}$.
\end{Solution}

% @see : https://coopmaths.fr/alea?uuid=ebd89&id=1AN14-1&n=1&d=10&s=7&alea=vjN72232323232342345678910111213141516171819202122232425262728293031323334353637383940414243444546474823456789101112131415161718192021222324252627282930313233343536373839&cd=0&cols=1
\needspace{10\baselineskip}
\begin{exercice}[BaremeTotal=false]


 Donner la dérivée de la fonction $f$, dérivable pour tout $x\in\R^*$, définie par  $f(x)=\dfrac{-3}{4x}$.
\end{exercice}

\begin{Solution}
 L'expression de la dérivée de la fonction $f$ définie par $f(x)=\dfrac{-3}{4x}$ est : ${\color[HTML]{f15929}\boldsymbol{f'(x)=\dfrac{3}{4x^2}}}$.
\end{Solution}

% @see : https://coopmaths.fr/alea?uuid=a83c0&id=1AN14-4&n=1&d=10&s=1&alea=Njr32232323232342345678910111213141516171819202122232425262728293031323334353637383940414243444546474823456789101112131415161718192021222324252627282930313233343536373839&cd=0&cols=1
\needspace{10\baselineskip}
\begin{exercice}[BaremeTotal=false]


 Donner l'expression de la dérivée de la fonction $f$ définie sur $\R^{*}$ par $f(x)=-5x-\dfrac{3}{x}$.\\
\end{exercice}

\begin{Solution}
 L'expression de la dérivée de la fonction $f$ définie par $f(x)=-5x-\dfrac{3}{x}$ est :\\$f'(x)={\color[HTML]{f15929}\boldsymbol{-5+\dfrac{3}{x^2}}}$.\\
\end{Solution}

% @see : https://coopmaths.fr/alea?uuid=a83c0&id=1AN14-4&n=1&d=10&s=4&alea=4Vpz2232323232342345678910111213141516171819202122232425262728293031323334353637383940414243444546474823456789101112131415161718192021222324252627282930313233343536373839&cd=1&cols=1
\needspace{10\baselineskip}
\begin{exercice}[BaremeTotal=false]


 Donner l'expression de la dérivée de la fonction $f$ définie sur $\R_+^{*}$ par $f(x)=4\sqrt{x}-\dfrac{4}{x}+3x^{2}-2x$.\\
\end{exercice}

\begin{Solution}
 $(4\sqrt{x})^\prime=\dfrac{4}{2\sqrt{x}}$\\$(-\dfrac{4}{x})^\prime=\dfrac{4}{x^2}$\\$(3x^{2}-2x)^\prime=6x-2$\\L'expression de la dérivée de la fonction $f$ définie par $f(x)=4\sqrt{x}-\dfrac{4}{x}+3x^{2}-2x$ est :\\$f'(x)={\color[HTML]{f15929}\boldsymbol{\dfrac{4}{2\sqrt{x}}+\dfrac{4}{x^2}+6x-2}}$.\\
\end{Solution}

\end{Maquette}
\clearpage
\begin{Maquette}[DM]{Niveau= ,Classe=\getdata[40]\mydata\ (v40), Date = 06/01/25  ,Theme=Exercices}


% @see : https://coopmaths.fr/alea?uuid=4c8c7&id=1AN11-3&n=1&d=10&s=2&alea=nawO223232323234234567891011121314151617181920212223242526272829303132333435363738394041424344454647482345678910111213141516171819202122232425262728293031323334353637383940&cd=1&cols=2
\needspace{10\baselineskip}
\begin{exercice}[BaremeTotal=false]
\label{eleve:40}


 \begin{minipage}[t]{\linewidth} Soit $f$ une fonction dérivable sur $[-5;5]$ et $\mathcal{C}_f$ sa courbe représentative.\\On sait que  $f(-4)=3~~$ et que $~~f'(-4)=3$.\\Déterminer une équation de la tangente $(T)$ à la courbe $\mathcal{C}_f$ au point d'abscisse $-4$,\\sans utiliser la formule de cours de l'équation de tangente. \end{minipage}

\end{exercice}

\begin{Solution}
 \begin{minipage}[t]{\linewidth} On sait que la tangente n'est pas une droite verticale, puisque la fonction est dérivable sur l'intervalle.\\On en déduit que la tangente $(T)$ au point d'abscisse $-4$, admet une équation réduite de la forme :  $(T) : y=m x + p$. 

\medskip
$\bullet$ Détermination de $m$ :\\On sait que le nombre dérivé en $-4$ est par définition, le coefficient directeur de la tangente au point d'abscisse $-4$.\\Par conséquent, on a déjà : $m=f'(-4)=3$.\\ On en déduit que  $(T) : y= 3 x + p$\\$\bullet$ Détermination de $p$ :\\ Pour cela, on utilise que si $f(-4)=3~~$, alors le point $A$ de coordonnées $(-4;3)$ appartient à $\mathcal{C}_f$ mais aussi à $(T)$.\\ On peut écrire $A(-4;3) = \mathcal{C}_f \cap (T)$.\\ On remplace alors les coordonnées de $A(-4;3)$ dans l'équation  $(T) : y= 3 x + p$.\\ $\begin{aligned} \phantom{\iff}&A(-4;3)\in (T)\\ \iff& 3= 3 \times (-4)  + p\\ \iff& p=3 -3 \times (-4)  \\ \iff& p=3 +12   \\ \iff& p=15\\\end{aligned}$\\On peut conclure que : $(T) : {\color[HTML]{f15929}\boldsymbol{y=3x+15}}$. \end{minipage}

\end{Solution}

% @see : https://coopmaths.fr/alea?uuid=0e984&id=can1F15&n=1&d=10&alea=1Alo223232323234234567891011121314151617181920212223242526272829303132333435363738394041424344454647482345678910111213141516171819202122232425262728293031323334353637383940&cd=1&cols=1
\needspace{10\baselineskip}
\begin{exercice}[BaremeTotal=false]


 La courbe représente une fonction $f$ et la droite est la tangente au point d'abscisse $2$.\\
        Déterminer $f'(2)$.\begin{center}\begin{tikzpicture}[baseline,scale = 0.5]

    \tikzset{
      point/.style={
        thick,
        draw,
        cross out,
        inner sep=0pt,
        minimum width=5pt,
        minimum height=5pt,
      },
    }
    \clip (-2,-2) rectangle (14,12);
    	
	\draw[color ={black},line width = 2,>=latex,->] (-2,0)--(14,0);
	\draw[color ={black},line width = 2,>=latex,->] (0,-2)--(0,12);
	\draw[color ={black},opacity = 0.5] (-2,1)--(14,1);
	\draw[color ={black},opacity = 0.5] (-2,2)--(14,2);
	\draw[color ={black},opacity = 0.5] (-2,3)--(14,3);
	\draw[color ={black},opacity = 0.5] (-2,4)--(14,4);
	\draw[color ={black},opacity = 0.5] (-2,5)--(14,5);
	\draw[color ={black},opacity = 0.5] (-2,6)--(14,6);
	\draw[color ={black},opacity = 0.5] (-2,7)--(14,7);
	\draw[color ={black},opacity = 0.5] (-2,8)--(14,8);
	\draw[color ={black},opacity = 0.5] (-2,9)--(14,9);
	\draw[color ={black},opacity = 0.5] (-2,10)--(14,10);
	\draw[color ={black},opacity = 0.5] (-2,11)--(14,11);
	\draw[color ={black},opacity = 0.5] (-2,12)--(14,12);
	\draw[color ={black},opacity = 0.5] (-2,-1)--(14,-1);
	\draw[color ={black},opacity = 0.5] (-2,-2)--(14,-2);
	\draw[color ={black},opacity = 0.5] (1,-2)--(1,12);
	\draw[color ={black},opacity = 0.5] (2,-2)--(2,12);
	\draw[color ={black},opacity = 0.5] (3,-2)--(3,12);
	\draw[color ={black},opacity = 0.5] (4,-2)--(4,12);
	\draw[color ={black},opacity = 0.5] (5,-2)--(5,12);
	\draw[color ={black},opacity = 0.5] (6,-2)--(6,12);
	\draw[color ={black},opacity = 0.5] (7,-2)--(7,12);
	\draw[color ={black},opacity = 0.5] (8,-2)--(8,12);
	\draw[color ={black},opacity = 0.5] (9,-2)--(9,12);
	\draw[color ={black},opacity = 0.5] (10,-2)--(10,12);
	\draw[color ={black},opacity = 0.5] (11,-2)--(11,12);
	\draw[color ={black},opacity = 0.5] (12,-2)--(12,12);
	\draw[color ={black},opacity = 0.5] (13,-2)--(13,12);
	\draw[color ={black},opacity = 0.5] (14,-2)--(14,12);
	\draw[color ={black},opacity = 0.5] (-1,-2)--(-1,12);
	\draw[color ={black},opacity = 0.5] (-2,-2)--(-2,12);
	\draw[color ={gray},opacity = 0.3] (-2,0)--(14,0);
	\draw[color ={gray},opacity = 0.3] (-2,0.5)--(14,0.5);
	\draw[color ={gray},opacity = 0.3] (-2,1.5)--(14,1.5);
	\draw[color ={gray},opacity = 0.3] (-2,2.5)--(14,2.5);
	\draw[color ={gray},opacity = 0.3] (-2,3.5)--(14,3.5);
	\draw[color ={gray},opacity = 0.3] (-2,4.5)--(14,4.5);
	\draw[color ={gray},opacity = 0.3] (-2,5.5)--(14,5.5);
	\draw[color ={gray},opacity = 0.3] (-2,6.5)--(14,6.5);
	\draw[color ={gray},opacity = 0.3] (-2,7.5)--(14,7.5);
	\draw[color ={gray},opacity = 0.3] (-2,8.5)--(14,8.5);
	\draw[color ={gray},opacity = 0.3] (-2,9.5)--(14,9.5);
	\draw[color ={gray},opacity = 0.3] (-2,10.5)--(14,10.5);
	\draw[color ={gray},opacity = 0.3] (-2,11.5)--(14,11.5);
	\draw[color ={gray},opacity = 0.3] (-2,0)--(14,0);
	\draw[color ={gray},opacity = 0.3] (-2,-0.5)--(14,-0.5);
	\draw[color ={gray},opacity = 0.3] (-2,-1.5)--(14,-1.5);
	\draw[color ={gray},opacity = 0.3] (0,-2)--(0,12);
	\draw[color ={gray},opacity = 0.3] (0.1,-2)--(0.1,12);
	\draw[color ={gray},opacity = 0.3] (0.2,-2)--(0.2,12);
	\draw[color ={gray},opacity = 0.3] (0.30000000000000004,-2)--(0.30000000000000004,12);
	\draw[color ={gray},opacity = 0.3] (0.4,-2)--(0.4,12);
	\draw[color ={gray},opacity = 0.3] (0.5,-2)--(0.5,12);
	\draw[color ={gray},opacity = 0.3] (0.6,-2)--(0.6,12);
	\draw[color ={gray},opacity = 0.3] (0.7,-2)--(0.7,12);
	\draw[color ={gray},opacity = 0.3] (0.7999999999999999,-2)--(0.7999999999999999,12);
	\draw[color ={gray},opacity = 0.3] (0.8999999999999999,-2)--(0.8999999999999999,12);
	\draw[color ={gray},opacity = 0.3] (1.0999999999999999,-2)--(1.0999999999999999,12);
	\draw[color ={gray},opacity = 0.3] (1.2,-2)--(1.2,12);
	\draw[color ={gray},opacity = 0.3] (1.3,-2)--(1.3,12);
	\draw[color ={gray},opacity = 0.3] (1.4000000000000001,-2)--(1.4000000000000001,12);
	\draw[color ={gray},opacity = 0.3] (1.5000000000000002,-2)--(1.5000000000000002,12);
	\draw[color ={gray},opacity = 0.3] (1.6000000000000003,-2)--(1.6000000000000003,12);
	\draw[color ={gray},opacity = 0.3] (1.7000000000000004,-2)--(1.7000000000000004,12);
	\draw[color ={gray},opacity = 0.3] (1.8000000000000005,-2)--(1.8000000000000005,12);
	\draw[color ={gray},opacity = 0.3] (1.9000000000000006,-2)--(1.9000000000000006,12);
	\draw[color ={gray},opacity = 0.3] (2.1000000000000005,-2)--(2.1000000000000005,12);
	\draw[color ={gray},opacity = 0.3] (2.2000000000000006,-2)--(2.2000000000000006,12);
	\draw[color ={gray},opacity = 0.3] (2.3000000000000007,-2)--(2.3000000000000007,12);
	\draw[color ={gray},opacity = 0.3] (2.400000000000001,-2)--(2.400000000000001,12);
	\draw[color ={gray},opacity = 0.3] (2.500000000000001,-2)--(2.500000000000001,12);
	\draw[color ={gray},opacity = 0.3] (2.600000000000001,-2)--(2.600000000000001,12);
	\draw[color ={gray},opacity = 0.3] (2.700000000000001,-2)--(2.700000000000001,12);
	\draw[color ={gray},opacity = 0.3] (2.800000000000001,-2)--(2.800000000000001,12);
	\draw[color ={gray},opacity = 0.3] (2.9000000000000012,-2)--(2.9000000000000012,12);
	\draw[color ={gray},opacity = 0.3] (3.1000000000000014,-2)--(3.1000000000000014,12);
	\draw[color ={gray},opacity = 0.3] (3.2000000000000015,-2)--(3.2000000000000015,12);
	\draw[color ={gray},opacity = 0.3] (3.3000000000000016,-2)--(3.3000000000000016,12);
	\draw[color ={gray},opacity = 0.3] (3.4000000000000017,-2)--(3.4000000000000017,12);
	\draw[color ={gray},opacity = 0.3] (3.5000000000000018,-2)--(3.5000000000000018,12);
	\draw[color ={gray},opacity = 0.3] (3.600000000000002,-2)--(3.600000000000002,12);
	\draw[color ={gray},opacity = 0.3] (3.700000000000002,-2)--(3.700000000000002,12);
	\draw[color ={gray},opacity = 0.3] (3.800000000000002,-2)--(3.800000000000002,12);
	\draw[color ={gray},opacity = 0.3] (3.900000000000002,-2)--(3.900000000000002,12);
	\draw[color ={gray},opacity = 0.3] (4.100000000000001,-2)--(4.100000000000001,12);
	\draw[color ={gray},opacity = 0.3] (4.200000000000001,-2)--(4.200000000000001,12);
	\draw[color ={gray},opacity = 0.3] (4.300000000000001,-2)--(4.300000000000001,12);
	\draw[color ={gray},opacity = 0.3] (4.4,-2)--(4.4,12);
	\draw[color ={gray},opacity = 0.3] (4.5,-2)--(4.5,12);
	\draw[color ={gray},opacity = 0.3] (4.6,-2)--(4.6,12);
	\draw[color ={gray},opacity = 0.3] (4.699999999999999,-2)--(4.699999999999999,12);
	\draw[color ={gray},opacity = 0.3] (4.799999999999999,-2)--(4.799999999999999,12);
	\draw[color ={gray},opacity = 0.3] (4.899999999999999,-2)--(4.899999999999999,12);
	\draw[color ={gray},opacity = 0.3] (5.099999999999998,-2)--(5.099999999999998,12);
	\draw[color ={gray},opacity = 0.3] (5.1999999999999975,-2)--(5.1999999999999975,12);
	\draw[color ={gray},opacity = 0.3] (5.299999999999997,-2)--(5.299999999999997,12);
	\draw[color ={gray},opacity = 0.3] (5.399999999999997,-2)--(5.399999999999997,12);
	\draw[color ={gray},opacity = 0.3] (5.4999999999999964,-2)--(5.4999999999999964,12);
	\draw[color ={gray},opacity = 0.3] (5.599999999999996,-2)--(5.599999999999996,12);
	\draw[color ={gray},opacity = 0.3] (5.699999999999996,-2)--(5.699999999999996,12);
	\draw[color ={gray},opacity = 0.3] (5.799999999999995,-2)--(5.799999999999995,12);
	\draw[color ={gray},opacity = 0.3] (5.899999999999995,-2)--(5.899999999999995,12);
	\draw[color ={gray},opacity = 0.3] (6.099999999999994,-2)--(6.099999999999994,12);
	\draw[color ={gray},opacity = 0.3] (6.199999999999994,-2)--(6.199999999999994,12);
	\draw[color ={gray},opacity = 0.3] (6.299999999999994,-2)--(6.299999999999994,12);
	\draw[color ={gray},opacity = 0.3] (6.399999999999993,-2)--(6.399999999999993,12);
	\draw[color ={gray},opacity = 0.3] (6.499999999999993,-2)--(6.499999999999993,12);
	\draw[color ={gray},opacity = 0.3] (6.5999999999999925,-2)--(6.5999999999999925,12);
	\draw[color ={gray},opacity = 0.3] (6.699999999999992,-2)--(6.699999999999992,12);
	\draw[color ={gray},opacity = 0.3] (6.799999999999992,-2)--(6.799999999999992,12);
	\draw[color ={gray},opacity = 0.3] (6.8999999999999915,-2)--(6.8999999999999915,12);
	\draw[color ={gray},opacity = 0.3] (7.099999999999991,-2)--(7.099999999999991,12);
	\draw[color ={gray},opacity = 0.3] (7.19999999999999,-2)--(7.19999999999999,12);
	\draw[color ={gray},opacity = 0.3] (7.29999999999999,-2)--(7.29999999999999,12);
	\draw[color ={gray},opacity = 0.3] (7.39999999999999,-2)--(7.39999999999999,12);
	\draw[color ={gray},opacity = 0.3] (7.499999999999989,-2)--(7.499999999999989,12);
	\draw[color ={gray},opacity = 0.3] (7.599999999999989,-2)--(7.599999999999989,12);
	\draw[color ={gray},opacity = 0.3] (7.699999999999989,-2)--(7.699999999999989,12);
	\draw[color ={gray},opacity = 0.3] (7.799999999999988,-2)--(7.799999999999988,12);
	\draw[color ={gray},opacity = 0.3] (7.899999999999988,-2)--(7.899999999999988,12);
	\draw[color ={gray},opacity = 0.3] (8.099999999999987,-2)--(8.099999999999987,12);
	\draw[color ={gray},opacity = 0.3] (8.199999999999987,-2)--(8.199999999999987,12);
	\draw[color ={gray},opacity = 0.3] (8.299999999999986,-2)--(8.299999999999986,12);
	\draw[color ={gray},opacity = 0.3] (8.399999999999986,-2)--(8.399999999999986,12);
	\draw[color ={gray},opacity = 0.3] (8.499999999999986,-2)--(8.499999999999986,12);
	\draw[color ={gray},opacity = 0.3] (8.599999999999985,-2)--(8.599999999999985,12);
	\draw[color ={gray},opacity = 0.3] (8.699999999999985,-2)--(8.699999999999985,12);
	\draw[color ={gray},opacity = 0.3] (8.799999999999985,-2)--(8.799999999999985,12);
	\draw[color ={gray},opacity = 0.3] (8.899999999999984,-2)--(8.899999999999984,12);
	\draw[color ={gray},opacity = 0.3] (9.099999999999984,-2)--(9.099999999999984,12);
	\draw[color ={gray},opacity = 0.3] (9.199999999999983,-2)--(9.199999999999983,12);
	\draw[color ={gray},opacity = 0.3] (9.299999999999983,-2)--(9.299999999999983,12);
	\draw[color ={gray},opacity = 0.3] (9.399999999999983,-2)--(9.399999999999983,12);
	\draw[color ={gray},opacity = 0.3] (9.499999999999982,-2)--(9.499999999999982,12);
	\draw[color ={gray},opacity = 0.3] (9.599999999999982,-2)--(9.599999999999982,12);
	\draw[color ={gray},opacity = 0.3] (9.699999999999982,-2)--(9.699999999999982,12);
	\draw[color ={gray},opacity = 0.3] (9.799999999999981,-2)--(9.799999999999981,12);
	\draw[color ={gray},opacity = 0.3] (9.89999999999998,-2)--(9.89999999999998,12);
	\draw[color ={gray},opacity = 0.3] (10.09999999999998,-2)--(10.09999999999998,12);
	\draw[color ={gray},opacity = 0.3] (10.19999999999998,-2)--(10.19999999999998,12);
	\draw[color ={gray},opacity = 0.3] (10.29999999999998,-2)--(10.29999999999998,12);
	\draw[color ={gray},opacity = 0.3] (10.399999999999979,-2)--(10.399999999999979,12);
	\draw[color ={gray},opacity = 0.3] (10.499999999999979,-2)--(10.499999999999979,12);
	\draw[color ={gray},opacity = 0.3] (10.599999999999978,-2)--(10.599999999999978,12);
	\draw[color ={gray},opacity = 0.3] (10.699999999999978,-2)--(10.699999999999978,12);
	\draw[color ={gray},opacity = 0.3] (10.799999999999978,-2)--(10.799999999999978,12);
	\draw[color ={gray},opacity = 0.3] (10.899999999999977,-2)--(10.899999999999977,12);
	\draw[color ={gray},opacity = 0.3] (11.099999999999977,-2)--(11.099999999999977,12);
	\draw[color ={gray},opacity = 0.3] (11.199999999999976,-2)--(11.199999999999976,12);
	\draw[color ={gray},opacity = 0.3] (11.299999999999976,-2)--(11.299999999999976,12);
	\draw[color ={gray},opacity = 0.3] (11.399999999999975,-2)--(11.399999999999975,12);
	\draw[color ={gray},opacity = 0.3] (11.499999999999975,-2)--(11.499999999999975,12);
	\draw[color ={gray},opacity = 0.3] (11.599999999999975,-2)--(11.599999999999975,12);
	\draw[color ={gray},opacity = 0.3] (11.699999999999974,-2)--(11.699999999999974,12);
	\draw[color ={gray},opacity = 0.3] (11.799999999999974,-2)--(11.799999999999974,12);
	\draw[color ={gray},opacity = 0.3] (11.899999999999974,-2)--(11.899999999999974,12);
	\draw[color ={gray},opacity = 0.3] (12.099999999999973,-2)--(12.099999999999973,12);
	\draw[color ={gray},opacity = 0.3] (12.199999999999973,-2)--(12.199999999999973,12);
	\draw[color ={gray},opacity = 0.3] (12.299999999999972,-2)--(12.299999999999972,12);
	\draw[color ={gray},opacity = 0.3] (12.399999999999972,-2)--(12.399999999999972,12);
	\draw[color ={gray},opacity = 0.3] (12.499999999999972,-2)--(12.499999999999972,12);
	\draw[color ={gray},opacity = 0.3] (12.599999999999971,-2)--(12.599999999999971,12);
	\draw[color ={gray},opacity = 0.3] (12.69999999999997,-2)--(12.69999999999997,12);
	\draw[color ={gray},opacity = 0.3] (12.79999999999997,-2)--(12.79999999999997,12);
	\draw[color ={gray},opacity = 0.3] (12.89999999999997,-2)--(12.89999999999997,12);
	\draw[color ={gray},opacity = 0.3] (13.09999999999997,-2)--(13.09999999999997,12);
	\draw[color ={gray},opacity = 0.3] (13.199999999999969,-2)--(13.199999999999969,12);
	\draw[color ={gray},opacity = 0.3] (13.299999999999969,-2)--(13.299999999999969,12);
	\draw[color ={gray},opacity = 0.3] (13.399999999999968,-2)--(13.399999999999968,12);
	\draw[color ={gray},opacity = 0.3] (13.499999999999968,-2)--(13.499999999999968,12);
	\draw[color ={gray},opacity = 0.3] (13.599999999999968,-2)--(13.599999999999968,12);
	\draw[color ={gray},opacity = 0.3] (13.699999999999967,-2)--(13.699999999999967,12);
	\draw[color ={gray},opacity = 0.3] (13.799999999999967,-2)--(13.799999999999967,12);
	\draw[color ={gray},opacity = 0.3] (13.899999999999967,-2)--(13.899999999999967,12);
	\draw[color ={gray},opacity = 0.3] (0,-2)--(0,12);
	\draw[color ={gray},opacity = 0.3] (-0.1,-2)--(-0.1,12);
	\draw[color ={gray},opacity = 0.3] (-0.2,-2)--(-0.2,12);
	\draw[color ={gray},opacity = 0.3] (-0.30000000000000004,-2)--(-0.30000000000000004,12);
	\draw[color ={gray},opacity = 0.3] (-0.4,-2)--(-0.4,12);
	\draw[color ={gray},opacity = 0.3] (-0.5,-2)--(-0.5,12);
	\draw[color ={gray},opacity = 0.3] (-0.6,-2)--(-0.6,12);
	\draw[color ={gray},opacity = 0.3] (-0.7,-2)--(-0.7,12);
	\draw[color ={gray},opacity = 0.3] (-0.7999999999999999,-2)--(-0.7999999999999999,12);
	\draw[color ={gray},opacity = 0.3] (-0.8999999999999999,-2)--(-0.8999999999999999,12);
	\draw[color ={gray},opacity = 0.3] (-1.0999999999999999,-2)--(-1.0999999999999999,12);
	\draw[color ={gray},opacity = 0.3] (-1.2,-2)--(-1.2,12);
	\draw[color ={gray},opacity = 0.3] (-1.3,-2)--(-1.3,12);
	\draw[color ={gray},opacity = 0.3] (-1.4000000000000001,-2)--(-1.4000000000000001,12);
	\draw[color ={gray},opacity = 0.3] (-1.5000000000000002,-2)--(-1.5000000000000002,12);
	\draw[color ={gray},opacity = 0.3] (-1.6000000000000003,-2)--(-1.6000000000000003,12);
	\draw[color ={gray},opacity = 0.3] (-1.7000000000000004,-2)--(-1.7000000000000004,12);
	\draw[color ={gray},opacity = 0.3] (-1.8000000000000005,-2)--(-1.8000000000000005,12);
	\draw[color ={gray},opacity = 0.3] (-1.9000000000000006,-2)--(-1.9000000000000006,12);
	\draw[color ={black},line width = 2] (2,-0.2)--(2,0.2);
	\draw[color ={black},line width = 2] (4,-0.2)--(4,0.2);
	\draw[color ={black},line width = 2] (6,-0.2)--(6,0.2);
	\draw[color ={black},line width = 2] (8,-0.2)--(8,0.2);
	\draw[color ={black},line width = 2] (10,-0.2)--(10,0.2);
	\draw[color ={black},line width = 2] (12,-0.2)--(12,0.2);
	\draw[color ={black},line width = 2] (-0.2,2)--(0.2,2);
	\draw[color ={black},line width = 2] (-0.2,4)--(0.2,4);
	\draw[color ={black},line width = 2] (-0.2,6)--(0.2,6);
	\draw[color ={black},line width = 2] (-0.2,8)--(0.2,8);
	\draw[color ={black},line width = 2] (-0.2,10)--(0.2,10);
	\draw (2,-0.4) node[anchor = center] {\scriptsize \color{black}{$1$}};
	\draw (4,-0.4) node[anchor = center] {\scriptsize \color{black}{$2$}};
	\draw (6,-0.4) node[anchor = center] {\scriptsize \color{black}{$3$}};
	\draw (8,-0.4) node[anchor = center] {\scriptsize \color{black}{$4$}};
	\draw (10,-0.4) node[anchor = center] {\scriptsize \color{black}{$5$}};
	\draw (-0.5,2.1) node[anchor = center] {\footnotesize \color{black}{$1$}};
	\draw (-0.5,4.1) node[anchor = center] {\footnotesize \color{black}{$2$}};
	\draw (-0.5,6.1) node[anchor = center] {\footnotesize \color{black}{$3$}};
	\draw (-0.5,8.1) node[anchor = center] {\footnotesize \color{black}{$4$}};
	\draw (-0.5,10.1) node[anchor = center] {\footnotesize \color{black}{$5$}};
	\draw [color={black}] (-0.3,-0.3) node[anchor = center,scale=1, rotate = 0] {O};
	
	\draw[color={blue},line width = 2] (0.6,10.67)--(1,8)--(1.4,6.86)--(1.8,6.22)--(2.2,5.82)--(2.6,5.54)--(3,5.33)--(3.4,5.18)--(3.8,5.05)--(4.2,4.95)--(4.6,4.87)--(5,4.8)--(5.4,4.74)--(5.8,4.69)--(6.2,4.65)--(6.6,4.61)--(7,4.57)--(7.4,4.54)--(7.8,4.51)--(8.2,4.49)--(8.6,4.47)--(9,4.44)--(9.4,4.43)--(9.8,4.41)--(10.2,4.39)--(10.6,4.38)--(11,4.36)--(11.4,4.35)--(11.8,4.34)--(12.2,4.33)--(12.6,4.32)--(13,4.31)--(13.4,4.3)--(13.8,4.29);
	
	\draw[color={red},line width = 2] (-2,6.5)--(-1.6,6.4)--(-1.2,6.3)--(-0.8,6.2)--(-0.4,6.1)--(0,6)--(0.4,5.9)--(0.8,5.8)--(1.2,5.7)--(1.6,5.6)--(2,5.5)--(2.4,5.4)--(2.8,5.3)--(3.2,5.2)--(3.6,5.1)--(4,5)--(4.4,4.9)--(4.8,4.8)--(5.2,4.7)--(5.6,4.6)--(6,4.5)--(6.4,4.4)--(6.8,4.3)--(7.2,4.2)--(7.6,4.1)--(8,4)--(8.4,3.9)--(8.8,3.8)--(9.2,3.7)--(9.6,3.6)--(10,3.5)--(10.4,3.4)--(10.8,3.3)--(11.2,3.2)--(11.6,3.1)--(12,3)--(12.4,2.9)--(12.8,2.8)--(13.2,2.7)--(13.6,2.6)--(14,2.5);

\end{tikzpicture}\end{center}
\end{exercice}

\begin{Solution}
 $f'(2)$ est donné par le coefficient directeur de la tangente à la courbe au point d'abscisse $2$, soit $\dfrac{-1}{4}=-\dfrac{1}{4}$.
\end{Solution}

% @see : https://coopmaths.fr/alea?uuid=6f32d&id=can1F16&n=1&d=10&alea=dqSw223232323234234567891011121314151617181920212223242526272829303132333435363738394041424344454647482345678910111213141516171819202122232425262728293031323334353637383940&cd=1&cols=1
\needspace{10\baselineskip}
\newpage
\begin{exercice}[BaremeTotal=false]


 On donne les représentations graphiques d'une fonction et de sa dérivée.\\
      Donner l'équation réduite de la tangente à la courbe de $f$ en $x=0$. \\ \begin{minipage}{0.5\linewidth}
    \begin{tikzpicture}[baseline, scale=0.8]

    \tikzset{
      point/.style={
        thick,
        draw,
        cross out,
        inner sep=0pt,
        minimum width=5pt,
        minimum height=5pt,
      },
    }
    \clip (-5.5,-9.5) rectangle (5.75,7.5);
    	
	\draw[color ={black},>=latex,->] (-4,0)--(4,0);
	\draw[color ={black},>=latex,->] (0,-8)--(0,6);
	\draw[color ={black},opacity = 0.4] (-4,1)--(4,1);
	\draw[color ={black},opacity = 0.4] (-4,2)--(4,2);
	\draw[color ={black},opacity = 0.4] (-4,3)--(4,3);
	\draw[color ={black},opacity = 0.4] (-4,4)--(4,4);
	\draw[color ={black},opacity = 0.4] (-4,5)--(4,5);
	\draw[color ={black},opacity = 0.4] (-4,6)--(4,6);
	\draw[color ={black},opacity = 0.4] (-4,-1)--(4,-1);
	\draw[color ={black},opacity = 0.4] (-4,-2)--(4,-2);
	\draw[color ={black},opacity = 0.4] (-4,-3)--(4,-3);
	\draw[color ={black},opacity = 0.4] (-4,-4)--(4,-4);
	\draw[color ={black},opacity = 0.4] (-4,-5)--(4,-5);
	\draw[color ={black},opacity = 0.4] (-4,-6)--(4,-6);
	\draw[color ={black},opacity = 0.4] (-4,-7)--(4,-7);
	\draw[color ={black},opacity = 0.4] (-4,-8)--(4,-8);
	\draw[color ={black},opacity = 0.4] (1,-8)--(1,6);
	\draw[color ={black},opacity = 0.4] (2,-8)--(2,6);
	\draw[color ={black},opacity = 0.4] (3,-8)--(3,6);
	\draw[color ={black},opacity = 0.4] (4,-8)--(4,6);
	\draw[color ={black},opacity = 0.4] (-1,-8)--(-1,6);
	\draw[color ={black},opacity = 0.4] (-2,-8)--(-2,6);
	\draw[color ={black},opacity = 0.4] (-3,-8)--(-3,6);
	\draw[color ={black},opacity = 0.4] (-4,-8)--(-4,6);
	\draw[color ={gray},opacity = 0.3] (-4,0)--(4,0);
	\draw[color ={gray},opacity = 0.3] (-4,0.5)--(4,0.5);
	\draw[color ={gray},opacity = 0.3] (-4,1.5)--(4,1.5);
	\draw[color ={gray},opacity = 0.3] (-4,2.5)--(4,2.5);
	\draw[color ={gray},opacity = 0.3] (-4,3.5)--(4,3.5);
	\draw[color ={gray},opacity = 0.3] (-4,4.5)--(4,4.5);
	\draw[color ={gray},opacity = 0.3] (-4,5.5)--(4,5.5);
	\draw[color ={gray},opacity = 0.3] (-4,0)--(4,0);
	\draw[color ={gray},opacity = 0.3] (-4,-0.5)--(4,-0.5);
	\draw[color ={gray},opacity = 0.3] (-4,-1.5)--(4,-1.5);
	\draw[color ={gray},opacity = 0.3] (-4,-2.5)--(4,-2.5);
	\draw[color ={gray},opacity = 0.3] (-4,-3.5)--(4,-3.5);
	\draw[color ={gray},opacity = 0.3] (-4,-4.5)--(4,-4.5);
	\draw[color ={gray},opacity = 0.3] (-4,-5.5)--(4,-5.5);
	\draw[color ={gray},opacity = 0.3] (-4,-6.5)--(4,-6.5);
	\draw[color ={gray},opacity = 0.3] (-4,-7)--(4,-7);
	\draw[color ={gray},opacity = 0.3] (-4,-7.5)--(4,-7.5);
	\draw[color ={gray},opacity = 0.3] (-4,-8)--(4,-8);
	\draw[color ={gray},opacity = 0.3] (0,-8)--(0,6);
	\draw[color ={gray},opacity = 0.3] (0,-8)--(0,6);
	\draw[color ={black},line width = 1.2] (1,-0.13)--(1,0.13);
	\draw[color ={black},line width = 1.2] (2,-0.13)--(2,0.13);
	\draw[color ={black},line width = 1.2] (3,-0.13)--(3,0.13);
	\draw[color ={black},line width = 1.2] (4,-0.13)--(4,0.13);
	\draw[color ={black},line width = 1.2] (-1,-0.13)--(-1,0.13);
	\draw[color ={black},line width = 1.2] (-2,-0.13)--(-2,0.13);
	\draw[color ={black},line width = 1.2] (-3,-0.13)--(-3,0.13);
	\draw[color ={black},line width = 1.2] (-4,-0.13)--(-4,0.13);
	\draw[color ={black},line width = 1.2] (-0.13,1)--(0.13,1);
	\draw[color ={black},line width = 1.2] (-0.13,2)--(0.13,2);
	\draw[color ={black},line width = 1.2] (-0.13,3)--(0.13,3);
	\draw[color ={black},line width = 1.2] (-0.13,4)--(0.13,4);
	\draw[color ={black},line width = 1.2] (-0.13,5)--(0.13,5);
	\draw[color ={black},line width = 1.2] (-0.13,6)--(0.13,6);
	\draw[color ={black},line width = 1.2] (-0.13,-1)--(0.13,-1);
	\draw[color ={black},line width = 1.2] (-0.13,-2)--(0.13,-2);
	\draw[color ={black},line width = 1.2] (-0.13,-3)--(0.13,-3);
	\draw[color ={black},line width = 1.2] (-0.13,-4)--(0.13,-4);
	\draw[color ={black},line width = 1.2] (-0.13,-5)--(0.13,-5);
	\draw[color ={black},line width = 1.2] (-0.13,-6)--(0.13,-6);
	\draw[color ={black},line width = 1.2] (-0.13,-7)--(0.13,-7);
	\draw[color ={black},line width = 1.2] (-0.13,-8)--(0.13,-8);
	\draw (2,-0.4) node[anchor = center] {\scriptsize \color{black}{$2$}};
	\draw (4,-0.4) node[anchor = center] {\scriptsize \color{black}{$4$}};
	\draw (-2,-0.4) node[anchor = center] {\scriptsize \color{black}{$-2$}};
	\draw (-4,-0.4) node[anchor = center] {\scriptsize \color{black}{$-4$}};
	\draw (-0.5,2.1) node[anchor = center] {\footnotesize \color{black}{$2$}};
	\draw (-0.5,4.1) node[anchor = center] {\footnotesize \color{black}{$4$}};
	\draw (-0.5,6.1) node[anchor = center] {\footnotesize \color{black}{$6$}};
	\draw (-0.5,-1.9) node[anchor = center] {\footnotesize \color{black}{$-2$}};
	\draw (-0.5,-3.9) node[anchor = center] {\footnotesize \color{black}{$-4$}};
	\draw (-0.5,-5.9) node[anchor = center] {\footnotesize \color{black}{$-6$}};
	\draw (-0.5,-7.9) node[anchor = center] {\footnotesize \color{black}{$-8$}};
	\draw [color={black}] (-0.3,-0.3) node[anchor = center,scale=1, rotate = 0] {O};
	\draw (3,4) node[anchor = center] {\footnotesize \color{blue}{$\Large \mathcal{C}_f$}};
	
	\draw[color={blue},line width = 2] (-4,-7)--(-3.8,-5.84)--(-3.6,-4.76)--(-3.4,-3.76)--(-3.2,-2.84)--(-3,-2)--(-2.8,-1.24)--(-2.6,-0.56)--(-2.4,0.04)--(-2.2,0.56)--(-2,1)--(-1.8,1.36)--(-1.6,1.64)--(-1.4,1.84)--(-1.2,1.96)--(-1,2)--(-0.8,1.96)--(-0.6,1.84)--(-0.4,1.64)--(-0.2,1.36)--(0,1)--(0.2,0.56)--(0.4,0.04)--(0.6,-0.56)--(0.8,-1.24)--(1,-2)--(1.2,-2.84)--(1.4,-3.76)--(1.6,-4.76)--(1.8,-5.84)--(2,-7)--(2.2,-8.24);

\end{tikzpicture}\\
    \end{minipage}
    \begin{minipage}{0.5\linewidth}
    \begin{tikzpicture}[baseline, scale=0.8]

    \tikzset{
      point/.style={
        thick,
        draw,
        cross out,
        inner sep=0pt,
        minimum width=5pt,
        minimum height=5pt,
      },
    }
    \clip (-5.5,-9.5) rectangle (6.5,7.5);
    	
	\draw[color ={black},>=latex,->] (-4,0)--(4,0);
	\draw[color ={black},>=latex,->] (0,-8)--(0,6);
	\draw[color ={black},opacity = 0.4] (-4,1)--(4,1);
	\draw[color ={black},opacity = 0.4] (-4,2)--(4,2);
	\draw[color ={black},opacity = 0.4] (-4,3)--(4,3);
	\draw[color ={black},opacity = 0.4] (-4,4)--(4,4);
	\draw[color ={black},opacity = 0.4] (-4,5)--(4,5);
	\draw[color ={black},opacity = 0.4] (-4,6)--(4,6);
	\draw[color ={black},opacity = 0.4] (-4,-1)--(4,-1);
	\draw[color ={black},opacity = 0.4] (-4,-2)--(4,-2);
	\draw[color ={black},opacity = 0.4] (-4,-3)--(4,-3);
	\draw[color ={black},opacity = 0.4] (-4,-4)--(4,-4);
	\draw[color ={black},opacity = 0.4] (-4,-5)--(4,-5);
	\draw[color ={black},opacity = 0.4] (-4,-6)--(4,-6);
	\draw[color ={black},opacity = 0.4] (-4,-7)--(4,-7);
	\draw[color ={black},opacity = 0.4] (-4,-8)--(4,-8);
	\draw[color ={black},opacity = 0.4] (1,-8)--(1,6);
	\draw[color ={black},opacity = 0.4] (2,-8)--(2,6);
	\draw[color ={black},opacity = 0.4] (3,-8)--(3,6);
	\draw[color ={black},opacity = 0.4] (4,-8)--(4,6);
	\draw[color ={black},opacity = 0.4] (-1,-8)--(-1,6);
	\draw[color ={black},opacity = 0.4] (-2,-8)--(-2,6);
	\draw[color ={black},opacity = 0.4] (-3,-8)--(-3,6);
	\draw[color ={black},opacity = 0.4] (-4,-8)--(-4,6);
	\draw[color ={gray},opacity = 0.3] (-4,0)--(4,0);
	\draw[color ={gray},opacity = 0.3] (-4,0.5)--(4,0.5);
	\draw[color ={gray},opacity = 0.3] (-4,1.5)--(4,1.5);
	\draw[color ={gray},opacity = 0.3] (-4,2.5)--(4,2.5);
	\draw[color ={gray},opacity = 0.3] (-4,3.5)--(4,3.5);
	\draw[color ={gray},opacity = 0.3] (-4,4.5)--(4,4.5);
	\draw[color ={gray},opacity = 0.3] (-4,5.5)--(4,5.5);
	\draw[color ={gray},opacity = 0.3] (-4,0)--(4,0);
	\draw[color ={gray},opacity = 0.3] (-4,-0.5)--(4,-0.5);
	\draw[color ={gray},opacity = 0.3] (-4,-1.5)--(4,-1.5);
	\draw[color ={gray},opacity = 0.3] (-4,-2.5)--(4,-2.5);
	\draw[color ={gray},opacity = 0.3] (-4,-3.5)--(4,-3.5);
	\draw[color ={gray},opacity = 0.3] (-4,-4.5)--(4,-4.5);
	\draw[color ={gray},opacity = 0.3] (-4,-5.5)--(4,-5.5);
	\draw[color ={gray},opacity = 0.3] (-4,-6.5)--(4,-6.5);
	\draw[color ={gray},opacity = 0.3] (-4,-7)--(4,-7);
	\draw[color ={gray},opacity = 0.3] (-4,-7.5)--(4,-7.5);
	\draw[color ={gray},opacity = 0.3] (-4,-8)--(4,-8);
	\draw[color ={gray},opacity = 0.3] (0,-8)--(0,6);
	\draw[color ={gray},opacity = 0.3] (0,-8)--(0,6);
	\draw[color ={black},line width = 1.2] (1,-0.13)--(1,0.13);
	\draw[color ={black},line width = 1.2] (2,-0.13)--(2,0.13);
	\draw[color ={black},line width = 1.2] (3,-0.13)--(3,0.13);
	\draw[color ={black},line width = 1.2] (4,-0.13)--(4,0.13);
	\draw[color ={black},line width = 1.2] (-1,-0.13)--(-1,0.13);
	\draw[color ={black},line width = 1.2] (-2,-0.13)--(-2,0.13);
	\draw[color ={black},line width = 1.2] (-3,-0.13)--(-3,0.13);
	\draw[color ={black},line width = 1.2] (-4,-0.13)--(-4,0.13);
	\draw[color ={black},line width = 1.2] (-0.13,1)--(0.13,1);
	\draw[color ={black},line width = 1.2] (-0.13,2)--(0.13,2);
	\draw[color ={black},line width = 1.2] (-0.13,3)--(0.13,3);
	\draw[color ={black},line width = 1.2] (-0.13,4)--(0.13,4);
	\draw[color ={black},line width = 1.2] (-0.13,5)--(0.13,5);
	\draw[color ={black},line width = 1.2] (-0.13,6)--(0.13,6);
	\draw[color ={black},line width = 1.2] (-0.13,-1)--(0.13,-1);
	\draw[color ={black},line width = 1.2] (-0.13,-2)--(0.13,-2);
	\draw[color ={black},line width = 1.2] (-0.13,-3)--(0.13,-3);
	\draw[color ={black},line width = 1.2] (-0.13,-4)--(0.13,-4);
	\draw[color ={black},line width = 1.2] (-0.13,-5)--(0.13,-5);
	\draw[color ={black},line width = 1.2] (-0.13,-6)--(0.13,-6);
	\draw[color ={black},line width = 1.2] (-0.13,-7)--(0.13,-7);
	\draw[color ={black},line width = 1.2] (-0.13,-8)--(0.13,-8);
	\draw (2,-0.4) node[anchor = center] {\scriptsize \color{black}{$2$}};
	\draw (4,-0.4) node[anchor = center] {\scriptsize \color{black}{$4$}};
	\draw (-2,-0.4) node[anchor = center] {\scriptsize \color{black}{$-2$}};
	\draw (-4,-0.4) node[anchor = center] {\scriptsize \color{black}{$-4$}};
	\draw (-0.5,2.1) node[anchor = center] {\footnotesize \color{black}{$2$}};
	\draw (-0.5,4.1) node[anchor = center] {\footnotesize \color{black}{$4$}};
	\draw (-0.5,-1.9) node[anchor = center] {\footnotesize \color{black}{$-2$}};
	\draw (-0.5,-3.9) node[anchor = center] {\footnotesize \color{black}{$-4$}};
	\draw (-0.5,-5.9) node[anchor = center] {\footnotesize \color{black}{$-6$}};
	\draw (-0.5,-7.9) node[anchor = center] {\footnotesize \color{black}{$-8$}};
	\draw [color={black}] (-0.3,-0.3) node[anchor = center,scale=1, rotate = 0] {O};
	\draw (3,4) node[anchor = center] {\footnotesize \color{red}{$\Large\mathcal{C}_f\prime$}};
	
	\draw[color={red},line width = 2] (-4,6)--(-3.8,5.6)--(-3.6,5.2)--(-3.4,4.8)--(-3.2,4.4)--(-3,4)--(-2.8,3.6)--(-2.6,3.2)--(-2.4,2.8)--(-2.2,2.4)--(-2,2)--(-1.8,1.6)--(-1.6,1.2)--(-1.4,0.8)--(-1.2,0.4)--(-1,0)--(-0.8,-0.4)--(-0.6,-0.8)--(-0.4,-1.2)--(-0.2,-1.6)--(0,-2)--(0.2,-2.4)--(0.4,-2.8)--(0.6,-3.2)--(0.8,-3.6)--(1,-4)--(1.2,-4.4)--(1.4,-4.8)--(1.6,-5.2)--(1.8,-5.6)--(2,-6)--(2.2,-6.4)--(2.4,-6.8)--(2.6,-7.2)--(2.8,-7.6)--(3,-8)--(3.2,-8.4)--(3.4,-8.8);

\end{tikzpicture}\\
    \end{minipage}
    
\end{exercice}

\begin{Solution}
 L'équation réduite de la tangente au point d'abscisse $0$ est  : $y=f'(0)(x-0)+f(0)$.\\
      On lit graphiquement $f(0)=1$ et $f'(0)=-2$.\\
      L'équation réduite de la tangente est donc donnée par :
      $y=-2(x+0)+1$, soit $y=-2x+1$.
\end{Solution}

% @see : https://coopmaths.fr/alea?uuid=a1ba2&id=can1F14&n=1&d=10&alea=rkIp223232323234234567891011121314151617181920212223242526272829303132333435363738394041424344454647482345678910111213141516171819202122232425262728293031323334353637383940&cd=1&cols=1
\needspace{10\baselineskip}
\begin{exercice}[BaremeTotal=false]


 Soit $f$ la fonction définie sur $\mathbb{R}$ par : $f(x)= 3x^2-x+1$.\\
        En calculant la limite du taux d'accroissement, déterminer $f'(-1)$.
\end{exercice}

\begin{Solution}
 $f$ est une fonction polynôme du second degré de la forme $f(x)=ax^2+bx+c$.\\
    La fonction dérivée est donnée par la somme des dérivées des fonctions $u$ et $v$ définies par $u(x)=3x^2$ et $v(x)=-x+1$.\\
     Comme $u'(x)=6x$ et $v'(x)=-1$, on obtient  $f'(x)=6x-1$. \\
     Ainsi, $f'(-1)=6\times (-1)-1=-7$.
\end{Solution}

% @see : https://coopmaths.fr/alea?uuid=3c690&id=can1F13&n=1&d=10&alea=EUDu223232323234234567891011121314151617181920212223242526272829303132333435363738394041424344454647482345678910111213141516171819202122232425262728293031323334353637383940&cd=1&cols=1
\needspace{10\baselineskip}
\begin{exercice}[BaremeTotal=false]


 Déterminer le coefficient directeur de la tangente à la courbe représentative de la fonction cube au point d'abscisse $-5$.
                        
\end{exercice}

\begin{Solution}
 Le coefficient directeur de la tangente au point d'abscisse $a$ est donné par le nombre dérivé $f'(a)$.\\
        La fonction $f$ définie par $f(x)=x^3$ a pour fonction dérivée la fonction $f'$ définie par $f'(x)=3x^2$.\\
        Comme $f'(-5)=3\times (-5)^2=75$, le coefficient directeur de la tangente au point d'abscisse $-5$ est : $75$. 
\end{Solution}

% @see : https://coopmaths.fr/alea?uuid=ebd89&id=1AN14-1&n=1&d=10&s=3&alea=ocYr223232323234234567891011121314151617181920212223242526272829303132333435363738394041424344454647482345678910111213141516171819202122232425262728293031323334353637383940&cd=0&cols=1
\needspace{10\baselineskip}
\begin{exercice}[BaremeTotal=false]


 Donner la dérivée de la fonction $f$, dérivable pour tout $x\in\R$, définie par  $f(x)=-\dfrac{3}{8}x+2$.
\end{exercice}

\begin{Solution}
 L'expression de la dérivée de la fonction $f$ définie par $f(x)=-\dfrac{3}{8}x+2$ est : ${\color[HTML]{f15929}\boldsymbol{f'(x)=-\dfrac{3}{8}}}$.
\end{Solution}

% @see : https://coopmaths.fr/alea?uuid=ebd89&id=1AN14-1&n=1&d=10&s=4&alea=uAlP223232323234234567891011121314151617181920212223242526272829303132333435363738394041424344454647482345678910111213141516171819202122232425262728293031323334353637383940&cd=0&cols=1
\needspace{10\baselineskip}
\begin{exercice}[BaremeTotal=false]


 Donner la dérivée de la fonction $f$, dérivable pour tout $x\in\R$, définie par  $f(x)=2x^{8}$.
\end{exercice}

\begin{Solution}
 L'expression de la dérivée de la fonction $f$ définie par $f(x)=2x^{8}$ est : ${\color[HTML]{f15929}\boldsymbol{f'(x)=16x^{7}}}$.
\end{Solution}

% @see : https://coopmaths.fr/alea?uuid=ebd89&id=1AN14-1&n=1&d=10&s=7&alea=vjN7223232323234234567891011121314151617181920212223242526272829303132333435363738394041424344454647482345678910111213141516171819202122232425262728293031323334353637383940&cd=0&cols=1
\needspace{10\baselineskip}
\begin{exercice}[BaremeTotal=false]


 Donner la dérivée de la fonction $f$, dérivable pour tout $x\in\R^*$, définie par  $f(x)=\dfrac{-3}{4x}$.
\end{exercice}

\begin{Solution}
 L'expression de la dérivée de la fonction $f$ définie par $f(x)=\dfrac{-3}{4x}$ est : ${\color[HTML]{f15929}\boldsymbol{f'(x)=\dfrac{3}{4x^2}}}$.
\end{Solution}

% @see : https://coopmaths.fr/alea?uuid=a83c0&id=1AN14-4&n=1&d=10&s=1&alea=Njr3223232323234234567891011121314151617181920212223242526272829303132333435363738394041424344454647482345678910111213141516171819202122232425262728293031323334353637383940&cd=0&cols=1
\needspace{10\baselineskip}
\begin{exercice}[BaremeTotal=false]


 Donner l'expression de la dérivée de la fonction $f$ définie sur $\R^{*}$ par $f(x)=-3x^{2}+3x-2-\dfrac{7}{x}$.\\
\end{exercice}

\begin{Solution}
 L'expression de la dérivée de la fonction $f$ définie par $f(x)=-3x^{2}+3x-2-\dfrac{7}{x}$ est :\\$f'(x)={\color[HTML]{f15929}\boldsymbol{-6x+3+\dfrac{7}{x^2}}}$.\\
\end{Solution}

% @see : https://coopmaths.fr/alea?uuid=a83c0&id=1AN14-4&n=1&d=10&s=4&alea=4Vpz223232323234234567891011121314151617181920212223242526272829303132333435363738394041424344454647482345678910111213141516171819202122232425262728293031323334353637383940&cd=1&cols=1
\needspace{10\baselineskip}
\begin{exercice}[BaremeTotal=false]


 Donner l'expression de la dérivée de la fonction $f$ définie sur $\R_+^{*}$ par $f(x)=2\sqrt{x}+3x^{2}+2x-\dfrac{4}{x}$.\\
\end{exercice}

\begin{Solution}
 $(2\sqrt{x})^\prime=\dfrac{2}{2\sqrt{x}}$\\$(3x^{2}+2x)^\prime=6x+2$\\$(-\dfrac{4}{x})^\prime=\dfrac{4}{x^2}$\\L'expression de la dérivée de la fonction $f$ définie par $f(x)=2\sqrt{x}+3x^{2}+2x-\dfrac{4}{x}$ est :\\$f'(x)={\color[HTML]{f15929}\boldsymbol{\dfrac{2}{2\sqrt{x}}+6x+2+\dfrac{4}{x^2}}}$.\\
\end{Solution}

\end{Maquette}
\clearpage
\begin{Maquette}[DM]{Niveau= ,Classe=\getdata[41]\mydata\ (v41), Date = 06/01/25  ,Theme=Exercices}


% @see : https://coopmaths.fr/alea?uuid=4c8c7&id=1AN11-3&n=1&d=10&s=2&alea=nawO22323232323423456789101112131415161718192021222324252627282930313233343536373839404142434445464748234567891011121314151617181920212223242526272829303132333435363738394041&cd=1&cols=2
\needspace{10\baselineskip}
\begin{exercice}[BaremeTotal=false]
\label{eleve:41}


 \begin{minipage}[t]{\linewidth} Soit $f$ une fonction dérivable sur $[-5;5]$ et $\mathcal{C}_f$ sa courbe représentative.\\On sait que  $f(0)=5~~$ et que $~~f'(0)=5$.\\Déterminer une équation de la tangente $(T)$ à la courbe $\mathcal{C}_f$ au point d'abscisse $0$,\\sans utiliser la formule de cours de l'équation de tangente. \end{minipage}

\end{exercice}

\begin{Solution}
 \begin{minipage}[t]{\linewidth} On sait que la tangente n'est pas une droite verticale, puisque la fonction est dérivable sur l'intervalle.\\On en déduit que la tangente $(T)$ au point d'abscisse $0$, admet une équation réduite de la forme :  $(T) : y=m x + p$. 

\medskip
$\bullet$ Détermination de $m$ :\\On sait que le nombre dérivé en $0$ est par définition, le coefficient directeur de la tangente au point d'abscisse $0$.\\Par conséquent, on a déjà : $m=f'(0)=5$.\\ On en déduit que  $(T) : y= 5 x + p$\\$\bullet$ Détermination de $p$ :\\ Pour cela, on utilise que si $f(0)=5~~$, alors le point $A$ de coordonnées $(0;5)$ appartient à $\mathcal{C}_f$ mais aussi à $(T)$.\\ On peut écrire $A(0;5) = \mathcal{C}_f \cap (T)$.\\ On remplace alors les coordonnées de $A(0;5)$ dans l'équation  $(T) : y= 5 x + p$.\\ $\begin{aligned} \phantom{\iff}&A(0;5)\in (T)\\ \iff& 5= 5 \times 0  + p\\ \iff& p=5 -5 \times 0  \\ \iff& p=5 +0   \\ \iff& p=5\\\end{aligned}$\\On peut conclure que : $(T) : {\color[HTML]{f15929}\boldsymbol{y=5x+5}}$. \end{minipage}

\end{Solution}

% @see : https://coopmaths.fr/alea?uuid=0e984&id=can1F15&n=1&d=10&alea=1Alo22323232323423456789101112131415161718192021222324252627282930313233343536373839404142434445464748234567891011121314151617181920212223242526272829303132333435363738394041&cd=1&cols=1
\needspace{10\baselineskip}
\begin{exercice}[BaremeTotal=false]


 La courbe représente une fonction $f$ et la droite est la tangente au point d'abscisse $0$.\\
        Déterminer $f'(0)$.\begin{center}\begin{tikzpicture}[baseline,scale = 0.5]

    \tikzset{
      point/.style={
        thick,
        draw,
        cross out,
        inner sep=0pt,
        minimum width=5pt,
        minimum height=5pt,
      },
    }
    \clip (-4,-2) rectangle (12,12);
    	
	\draw[color ={black},line width = 2,>=latex,->] (-4,0)--(12,0);
	\draw[color ={black},line width = 2,>=latex,->] (0,-2)--(0,12);
	\draw[color ={black},opacity = 0.5] (-4,1)--(12,1);
	\draw[color ={black},opacity = 0.5] (-4,2)--(12,2);
	\draw[color ={black},opacity = 0.5] (-4,3)--(12,3);
	\draw[color ={black},opacity = 0.5] (-4,4)--(12,4);
	\draw[color ={black},opacity = 0.5] (-4,5)--(12,5);
	\draw[color ={black},opacity = 0.5] (-4,6)--(12,6);
	\draw[color ={black},opacity = 0.5] (-4,7)--(12,7);
	\draw[color ={black},opacity = 0.5] (-4,8)--(12,8);
	\draw[color ={black},opacity = 0.5] (-4,9)--(12,9);
	\draw[color ={black},opacity = 0.5] (-4,10)--(12,10);
	\draw[color ={black},opacity = 0.5] (-4,11)--(12,11);
	\draw[color ={black},opacity = 0.5] (-4,12)--(12,12);
	\draw[color ={black},opacity = 0.5] (-4,-1)--(12,-1);
	\draw[color ={black},opacity = 0.5] (-4,-2)--(12,-2);
	\draw[color ={black},opacity = 0.5] (1,-2)--(1,12);
	\draw[color ={black},opacity = 0.5] (2,-2)--(2,12);
	\draw[color ={black},opacity = 0.5] (3,-2)--(3,12);
	\draw[color ={black},opacity = 0.5] (4,-2)--(4,12);
	\draw[color ={black},opacity = 0.5] (5,-2)--(5,12);
	\draw[color ={black},opacity = 0.5] (6,-2)--(6,12);
	\draw[color ={black},opacity = 0.5] (7,-2)--(7,12);
	\draw[color ={black},opacity = 0.5] (8,-2)--(8,12);
	\draw[color ={black},opacity = 0.5] (9,-2)--(9,12);
	\draw[color ={black},opacity = 0.5] (10,-2)--(10,12);
	\draw[color ={black},opacity = 0.5] (11,-2)--(11,12);
	\draw[color ={black},opacity = 0.5] (12,-2)--(12,12);
	\draw[color ={black},opacity = 0.5] (-1,-2)--(-1,12);
	\draw[color ={black},opacity = 0.5] (-2,-2)--(-2,12);
	\draw[color ={black},opacity = 0.5] (-3,-2)--(-3,12);
	\draw[color ={black},opacity = 0.5] (-4,-2)--(-4,12);
	\draw[color ={gray},opacity = 0.3] (-4,0)--(12,0);
	\draw[color ={gray},opacity = 0.3] (-4,0)--(12,0);
	\draw[color ={gray},opacity = 0.3] (0,-2)--(0,12);
	\draw[color ={gray},opacity = 0.3] (0,-2)--(0,12);
	\draw[color ={black},line width = 2] (2,-0.2)--(2,0.2);
	\draw[color ={black},line width = 2] (4,-0.2)--(4,0.2);
	\draw[color ={black},line width = 2] (6,-0.2)--(6,0.2);
	\draw[color ={black},line width = 2] (8,-0.2)--(8,0.2);
	\draw[color ={black},line width = 2] (10,-0.2)--(10,0.2);
	\draw[color ={black},line width = 2] (-2,-0.2)--(-2,0.2);
	\draw[color ={black},line width = 2] (-0.2,2)--(0.2,2);
	\draw[color ={black},line width = 2] (-0.2,4)--(0.2,4);
	\draw[color ={black},line width = 2] (-0.2,6)--(0.2,6);
	\draw[color ={black},line width = 2] (-0.2,8)--(0.2,8);
	\draw[color ={black},line width = 2] (-0.2,10)--(0.2,10);
	\draw (2,-0.4) node[anchor = center] {\scriptsize \color{black}{$1$}};
	\draw (4,-0.4) node[anchor = center] {\scriptsize \color{black}{$2$}};
	\draw (6,-0.4) node[anchor = center] {\scriptsize \color{black}{$3$}};
	\draw (8,-0.4) node[anchor = center] {\scriptsize \color{black}{$4$}};
	\draw (10,-0.4) node[anchor = center] {\scriptsize \color{black}{$5$}};
	\draw (-0.5,2.1) node[anchor = center] {\footnotesize \color{black}{$1$}};
	\draw (-0.5,4.1) node[anchor = center] {\footnotesize \color{black}{$2$}};
	\draw (-0.5,6.1) node[anchor = center] {\footnotesize \color{black}{$3$}};
	\draw (-0.5,8.1) node[anchor = center] {\footnotesize \color{black}{$4$}};
	\draw (-0.5,10.1) node[anchor = center] {\footnotesize \color{black}{$5$}};
	\draw [color={black}] (-0.3,-0.3) node[anchor = center,scale=1, rotate = 0] {O};
	
	\draw[color={blue},line width = 2] (-4,0.74)--(-3.6,0.81)--(-3.2,0.9)--(-2.8,0.99)--(-2.4,1.1)--(-2,1.21)--(-1.6,1.34)--(-1.2,1.48)--(-0.8,1.64)--(-0.4,1.81)--(0,2)--(0.4,2.21)--(0.8,2.44)--(1.2,2.7)--(1.6,2.98)--(2,3.3)--(2.4,3.64)--(2.8,4.03)--(3.2,4.45)--(3.6,4.92)--(4,5.44)--(4.4,6.01)--(4.8,6.64)--(5.2,7.34)--(5.6,8.11)--(6,8.96)--(6.4,9.91)--(6.8,10.95)--(7.2,12.1)--(7.6,13.37);
	
	\draw[color={red},line width = 2] (-4,0)--(-3.6,0.2)--(-3.2,0.4)--(-2.8,0.6)--(-2.4,0.8)--(-2,1)--(-1.6,1.2)--(-1.2,1.4)--(-0.8,1.6)--(-0.4,1.8)--(0,2)--(0.4,2.2)--(0.8,2.4)--(1.2,2.6)--(1.6,2.8)--(2,3)--(2.4,3.2)--(2.8,3.4)--(3.2,3.6)--(3.6,3.8)--(4,4)--(4.4,4.2)--(4.8,4.4)--(5.2,4.6)--(5.6,4.8)--(6,5)--(6.4,5.2)--(6.8,5.4)--(7.2,5.6)--(7.6,5.8)--(8,6)--(8.4,6.2)--(8.8,6.4)--(9.2,6.6)--(9.6,6.8)--(10,7)--(10.4,7.2)--(10.8,7.4)--(11.2,7.6)--(11.6,7.8)--(12,8);

\end{tikzpicture}\end{center}
\end{exercice}

\begin{Solution}
 $f'(0)$ est donné par le coefficient directeur de la tangente à la courbe au point d'abscisse $0$, soit $\dfrac{1}{2}$.
\end{Solution}

% @see : https://coopmaths.fr/alea?uuid=6f32d&id=can1F16&n=1&d=10&alea=dqSw22323232323423456789101112131415161718192021222324252627282930313233343536373839404142434445464748234567891011121314151617181920212223242526272829303132333435363738394041&cd=1&cols=1
\needspace{10\baselineskip}
\newpage
\begin{exercice}[BaremeTotal=false]


 On donne les représentations graphiques d'une fonction et de sa dérivée.\\
      Donner l'équation réduite de la tangente à la courbe de $f$ en $x=0$. \\ \begin{minipage}{0.5\linewidth}
    \begin{tikzpicture}[baseline, scale=0.8]

    \tikzset{
      point/.style={
        thick,
        draw,
        cross out,
        inner sep=0pt,
        minimum width=5pt,
        minimum height=5pt,
      },
    }
    \clip (-5.5,-9.5) rectangle (5.75,7.5);
    	
	\draw[color ={black},>=latex,->] (-4,0)--(4,0);
	\draw[color ={black},>=latex,->] (0,-8)--(0,6);
	\draw[color ={black},opacity = 0.4] (-4,1)--(4,1);
	\draw[color ={black},opacity = 0.4] (-4,2)--(4,2);
	\draw[color ={black},opacity = 0.4] (-4,3)--(4,3);
	\draw[color ={black},opacity = 0.4] (-4,4)--(4,4);
	\draw[color ={black},opacity = 0.4] (-4,5)--(4,5);
	\draw[color ={black},opacity = 0.4] (-4,6)--(4,6);
	\draw[color ={black},opacity = 0.4] (-4,-1)--(4,-1);
	\draw[color ={black},opacity = 0.4] (-4,-2)--(4,-2);
	\draw[color ={black},opacity = 0.4] (-4,-3)--(4,-3);
	\draw[color ={black},opacity = 0.4] (-4,-4)--(4,-4);
	\draw[color ={black},opacity = 0.4] (-4,-5)--(4,-5);
	\draw[color ={black},opacity = 0.4] (-4,-6)--(4,-6);
	\draw[color ={black},opacity = 0.4] (-4,-7)--(4,-7);
	\draw[color ={black},opacity = 0.4] (-4,-8)--(4,-8);
	\draw[color ={black},opacity = 0.4] (1,-8)--(1,6);
	\draw[color ={black},opacity = 0.4] (2,-8)--(2,6);
	\draw[color ={black},opacity = 0.4] (3,-8)--(3,6);
	\draw[color ={black},opacity = 0.4] (4,-8)--(4,6);
	\draw[color ={black},opacity = 0.4] (-1,-8)--(-1,6);
	\draw[color ={black},opacity = 0.4] (-2,-8)--(-2,6);
	\draw[color ={black},opacity = 0.4] (-3,-8)--(-3,6);
	\draw[color ={black},opacity = 0.4] (-4,-8)--(-4,6);
	\draw[color ={gray},opacity = 0.3] (-4,0)--(4,0);
	\draw[color ={gray},opacity = 0.3] (-4,0.5)--(4,0.5);
	\draw[color ={gray},opacity = 0.3] (-4,1.5)--(4,1.5);
	\draw[color ={gray},opacity = 0.3] (-4,2.5)--(4,2.5);
	\draw[color ={gray},opacity = 0.3] (-4,3.5)--(4,3.5);
	\draw[color ={gray},opacity = 0.3] (-4,4.5)--(4,4.5);
	\draw[color ={gray},opacity = 0.3] (-4,5.5)--(4,5.5);
	\draw[color ={gray},opacity = 0.3] (-4,0)--(4,0);
	\draw[color ={gray},opacity = 0.3] (-4,-0.5)--(4,-0.5);
	\draw[color ={gray},opacity = 0.3] (-4,-1.5)--(4,-1.5);
	\draw[color ={gray},opacity = 0.3] (-4,-2.5)--(4,-2.5);
	\draw[color ={gray},opacity = 0.3] (-4,-3.5)--(4,-3.5);
	\draw[color ={gray},opacity = 0.3] (-4,-4.5)--(4,-4.5);
	\draw[color ={gray},opacity = 0.3] (-4,-5.5)--(4,-5.5);
	\draw[color ={gray},opacity = 0.3] (-4,-6.5)--(4,-6.5);
	\draw[color ={gray},opacity = 0.3] (-4,-7)--(4,-7);
	\draw[color ={gray},opacity = 0.3] (-4,-7.5)--(4,-7.5);
	\draw[color ={gray},opacity = 0.3] (-4,-8)--(4,-8);
	\draw[color ={gray},opacity = 0.3] (0,-8)--(0,6);
	\draw[color ={gray},opacity = 0.3] (0,-8)--(0,6);
	\draw[color ={black},line width = 1.2] (1,-0.13)--(1,0.13);
	\draw[color ={black},line width = 1.2] (2,-0.13)--(2,0.13);
	\draw[color ={black},line width = 1.2] (3,-0.13)--(3,0.13);
	\draw[color ={black},line width = 1.2] (4,-0.13)--(4,0.13);
	\draw[color ={black},line width = 1.2] (-1,-0.13)--(-1,0.13);
	\draw[color ={black},line width = 1.2] (-2,-0.13)--(-2,0.13);
	\draw[color ={black},line width = 1.2] (-3,-0.13)--(-3,0.13);
	\draw[color ={black},line width = 1.2] (-4,-0.13)--(-4,0.13);
	\draw[color ={black},line width = 1.2] (-0.13,1)--(0.13,1);
	\draw[color ={black},line width = 1.2] (-0.13,2)--(0.13,2);
	\draw[color ={black},line width = 1.2] (-0.13,3)--(0.13,3);
	\draw[color ={black},line width = 1.2] (-0.13,4)--(0.13,4);
	\draw[color ={black},line width = 1.2] (-0.13,5)--(0.13,5);
	\draw[color ={black},line width = 1.2] (-0.13,6)--(0.13,6);
	\draw[color ={black},line width = 1.2] (-0.13,-1)--(0.13,-1);
	\draw[color ={black},line width = 1.2] (-0.13,-2)--(0.13,-2);
	\draw[color ={black},line width = 1.2] (-0.13,-3)--(0.13,-3);
	\draw[color ={black},line width = 1.2] (-0.13,-4)--(0.13,-4);
	\draw[color ={black},line width = 1.2] (-0.13,-5)--(0.13,-5);
	\draw[color ={black},line width = 1.2] (-0.13,-6)--(0.13,-6);
	\draw[color ={black},line width = 1.2] (-0.13,-7)--(0.13,-7);
	\draw[color ={black},line width = 1.2] (-0.13,-8)--(0.13,-8);
	\draw (2,-0.4) node[anchor = center] {\scriptsize \color{black}{$2$}};
	\draw (4,-0.4) node[anchor = center] {\scriptsize \color{black}{$4$}};
	\draw (-2,-0.4) node[anchor = center] {\scriptsize \color{black}{$-2$}};
	\draw (-4,-0.4) node[anchor = center] {\scriptsize \color{black}{$-4$}};
	\draw (-0.5,2.1) node[anchor = center] {\footnotesize \color{black}{$2$}};
	\draw (-0.5,4.1) node[anchor = center] {\footnotesize \color{black}{$4$}};
	\draw (-0.5,6.1) node[anchor = center] {\footnotesize \color{black}{$6$}};
	\draw (-0.5,-1.9) node[anchor = center] {\footnotesize \color{black}{$-2$}};
	\draw (-0.5,-3.9) node[anchor = center] {\footnotesize \color{black}{$-4$}};
	\draw (-0.5,-5.9) node[anchor = center] {\footnotesize \color{black}{$-6$}};
	\draw (-0.5,-7.9) node[anchor = center] {\footnotesize \color{black}{$-8$}};
	\draw [color={black}] (-0.3,-0.3) node[anchor = center,scale=1, rotate = 0] {O};
	\draw (3,4) node[anchor = center] {\footnotesize \color{blue}{$\Large \mathcal{C}_f$}};
	
	\draw[color={blue},line width = 2] (-4,-8)--(-3.8,-6.84)--(-3.6,-5.76)--(-3.4,-4.76)--(-3.2,-3.84)--(-3,-3)--(-2.8,-2.24)--(-2.6,-1.56)--(-2.4,-0.96)--(-2.2,-0.44)--(-2,0)--(-1.8,0.36)--(-1.6,0.64)--(-1.4,0.84)--(-1.2,0.96)--(-1,1)--(-0.8,0.96)--(-0.6,0.84)--(-0.4,0.64)--(-0.2,0.36)--(0,0)--(0.2,-0.44)--(0.4,-0.96)--(0.6,-1.56)--(0.8,-2.24)--(1,-3)--(1.2,-3.84)--(1.4,-4.76)--(1.6,-5.76)--(1.8,-6.84)--(2,-8);

\end{tikzpicture}\\
    \end{minipage}
    \begin{minipage}{0.5\linewidth}
    \begin{tikzpicture}[baseline, scale=0.8]

    \tikzset{
      point/.style={
        thick,
        draw,
        cross out,
        inner sep=0pt,
        minimum width=5pt,
        minimum height=5pt,
      },
    }
    \clip (-5.5,-9.5) rectangle (6.5,7.5);
    	
	\draw[color ={black},>=latex,->] (-4,0)--(4,0);
	\draw[color ={black},>=latex,->] (0,-8)--(0,6);
	\draw[color ={black},opacity = 0.4] (-4,1)--(4,1);
	\draw[color ={black},opacity = 0.4] (-4,2)--(4,2);
	\draw[color ={black},opacity = 0.4] (-4,3)--(4,3);
	\draw[color ={black},opacity = 0.4] (-4,4)--(4,4);
	\draw[color ={black},opacity = 0.4] (-4,5)--(4,5);
	\draw[color ={black},opacity = 0.4] (-4,6)--(4,6);
	\draw[color ={black},opacity = 0.4] (-4,-1)--(4,-1);
	\draw[color ={black},opacity = 0.4] (-4,-2)--(4,-2);
	\draw[color ={black},opacity = 0.4] (-4,-3)--(4,-3);
	\draw[color ={black},opacity = 0.4] (-4,-4)--(4,-4);
	\draw[color ={black},opacity = 0.4] (-4,-5)--(4,-5);
	\draw[color ={black},opacity = 0.4] (-4,-6)--(4,-6);
	\draw[color ={black},opacity = 0.4] (-4,-7)--(4,-7);
	\draw[color ={black},opacity = 0.4] (-4,-8)--(4,-8);
	\draw[color ={black},opacity = 0.4] (1,-8)--(1,6);
	\draw[color ={black},opacity = 0.4] (2,-8)--(2,6);
	\draw[color ={black},opacity = 0.4] (3,-8)--(3,6);
	\draw[color ={black},opacity = 0.4] (4,-8)--(4,6);
	\draw[color ={black},opacity = 0.4] (-1,-8)--(-1,6);
	\draw[color ={black},opacity = 0.4] (-2,-8)--(-2,6);
	\draw[color ={black},opacity = 0.4] (-3,-8)--(-3,6);
	\draw[color ={black},opacity = 0.4] (-4,-8)--(-4,6);
	\draw[color ={gray},opacity = 0.3] (-4,0)--(4,0);
	\draw[color ={gray},opacity = 0.3] (-4,0.5)--(4,0.5);
	\draw[color ={gray},opacity = 0.3] (-4,1.5)--(4,1.5);
	\draw[color ={gray},opacity = 0.3] (-4,2.5)--(4,2.5);
	\draw[color ={gray},opacity = 0.3] (-4,3.5)--(4,3.5);
	\draw[color ={gray},opacity = 0.3] (-4,4.5)--(4,4.5);
	\draw[color ={gray},opacity = 0.3] (-4,5.5)--(4,5.5);
	\draw[color ={gray},opacity = 0.3] (-4,0)--(4,0);
	\draw[color ={gray},opacity = 0.3] (-4,-0.5)--(4,-0.5);
	\draw[color ={gray},opacity = 0.3] (-4,-1.5)--(4,-1.5);
	\draw[color ={gray},opacity = 0.3] (-4,-2.5)--(4,-2.5);
	\draw[color ={gray},opacity = 0.3] (-4,-3.5)--(4,-3.5);
	\draw[color ={gray},opacity = 0.3] (-4,-4.5)--(4,-4.5);
	\draw[color ={gray},opacity = 0.3] (-4,-5.5)--(4,-5.5);
	\draw[color ={gray},opacity = 0.3] (-4,-6.5)--(4,-6.5);
	\draw[color ={gray},opacity = 0.3] (-4,-7)--(4,-7);
	\draw[color ={gray},opacity = 0.3] (-4,-7.5)--(4,-7.5);
	\draw[color ={gray},opacity = 0.3] (-4,-8)--(4,-8);
	\draw[color ={gray},opacity = 0.3] (0,-8)--(0,6);
	\draw[color ={gray},opacity = 0.3] (0,-8)--(0,6);
	\draw[color ={black},line width = 1.2] (1,-0.13)--(1,0.13);
	\draw[color ={black},line width = 1.2] (2,-0.13)--(2,0.13);
	\draw[color ={black},line width = 1.2] (3,-0.13)--(3,0.13);
	\draw[color ={black},line width = 1.2] (4,-0.13)--(4,0.13);
	\draw[color ={black},line width = 1.2] (-1,-0.13)--(-1,0.13);
	\draw[color ={black},line width = 1.2] (-2,-0.13)--(-2,0.13);
	\draw[color ={black},line width = 1.2] (-3,-0.13)--(-3,0.13);
	\draw[color ={black},line width = 1.2] (-4,-0.13)--(-4,0.13);
	\draw[color ={black},line width = 1.2] (-0.13,1)--(0.13,1);
	\draw[color ={black},line width = 1.2] (-0.13,2)--(0.13,2);
	\draw[color ={black},line width = 1.2] (-0.13,3)--(0.13,3);
	\draw[color ={black},line width = 1.2] (-0.13,4)--(0.13,4);
	\draw[color ={black},line width = 1.2] (-0.13,5)--(0.13,5);
	\draw[color ={black},line width = 1.2] (-0.13,6)--(0.13,6);
	\draw[color ={black},line width = 1.2] (-0.13,-1)--(0.13,-1);
	\draw[color ={black},line width = 1.2] (-0.13,-2)--(0.13,-2);
	\draw[color ={black},line width = 1.2] (-0.13,-3)--(0.13,-3);
	\draw[color ={black},line width = 1.2] (-0.13,-4)--(0.13,-4);
	\draw[color ={black},line width = 1.2] (-0.13,-5)--(0.13,-5);
	\draw[color ={black},line width = 1.2] (-0.13,-6)--(0.13,-6);
	\draw[color ={black},line width = 1.2] (-0.13,-7)--(0.13,-7);
	\draw[color ={black},line width = 1.2] (-0.13,-8)--(0.13,-8);
	\draw (2,-0.4) node[anchor = center] {\scriptsize \color{black}{$2$}};
	\draw (4,-0.4) node[anchor = center] {\scriptsize \color{black}{$4$}};
	\draw (-2,-0.4) node[anchor = center] {\scriptsize \color{black}{$-2$}};
	\draw (-4,-0.4) node[anchor = center] {\scriptsize \color{black}{$-4$}};
	\draw (-0.5,2.1) node[anchor = center] {\footnotesize \color{black}{$2$}};
	\draw (-0.5,4.1) node[anchor = center] {\footnotesize \color{black}{$4$}};
	\draw (-0.5,-1.9) node[anchor = center] {\footnotesize \color{black}{$-2$}};
	\draw (-0.5,-3.9) node[anchor = center] {\footnotesize \color{black}{$-4$}};
	\draw (-0.5,-5.9) node[anchor = center] {\footnotesize \color{black}{$-6$}};
	\draw (-0.5,-7.9) node[anchor = center] {\footnotesize \color{black}{$-8$}};
	\draw [color={black}] (-0.3,-0.3) node[anchor = center,scale=1, rotate = 0] {O};
	\draw (3,4) node[anchor = center] {\footnotesize \color{red}{$\Large\mathcal{C}_f\prime$}};
	
	\draw[color={red},line width = 2] (-4,6)--(-3.8,5.6)--(-3.6,5.2)--(-3.4,4.8)--(-3.2,4.4)--(-3,4)--(-2.8,3.6)--(-2.6,3.2)--(-2.4,2.8)--(-2.2,2.4)--(-2,2)--(-1.8,1.6)--(-1.6,1.2)--(-1.4,0.8)--(-1.2,0.4)--(-1,0)--(-0.8,-0.4)--(-0.6,-0.8)--(-0.4,-1.2)--(-0.2,-1.6)--(0,-2)--(0.2,-2.4)--(0.4,-2.8)--(0.6,-3.2)--(0.8,-3.6)--(1,-4)--(1.2,-4.4)--(1.4,-4.8)--(1.6,-5.2)--(1.8,-5.6)--(2,-6)--(2.2,-6.4)--(2.4,-6.8)--(2.6,-7.2)--(2.8,-7.6)--(3,-8)--(3.2,-8.4)--(3.4,-8.8);

\end{tikzpicture}\\
    \end{minipage}
    
\end{exercice}

\begin{Solution}
 L'équation réduite de la tangente au point d'abscisse $0$ est  : $y=f'(0)(x-0)+f(0)$.\\
      On lit graphiquement $f(0)=0$ et $f'(0)=-2$.\\
      L'équation réduite de la tangente est donc donnée par :
      $y=-2(x+0)+0$, soit $y=-2x$.
\end{Solution}

% @see : https://coopmaths.fr/alea?uuid=a1ba2&id=can1F14&n=1&d=10&alea=rkIp22323232323423456789101112131415161718192021222324252627282930313233343536373839404142434445464748234567891011121314151617181920212223242526272829303132333435363738394041&cd=1&cols=1
\needspace{10\baselineskip}
\begin{exercice}[BaremeTotal=false]


 Soit $f$ la fonction définie sur $\mathbb{R}$ par : $f(x)= 5x^2+3x-4$.\\
        En calculant la limite du taux d'accroissement, déterminer $f'(3)$.
\end{exercice}

\begin{Solution}
 $f$ est une fonction polynôme du second degré de la forme $f(x)=ax^2+bx+c$.\\
    La fonction dérivée est donnée par la somme des dérivées des fonctions $u$ et $v$ définies par $u(x)=5x^2$ et $v(x)=3x-4$.\\
     Comme $u'(x)=10x$ et $v'(x)=3$, on obtient  $f'(x)=10x+3$. \\
     Ainsi, $f'(3)=10\times 3+3=33$.
\end{Solution}

% @see : https://coopmaths.fr/alea?uuid=3c690&id=can1F13&n=1&d=10&alea=EUDu22323232323423456789101112131415161718192021222324252627282930313233343536373839404142434445464748234567891011121314151617181920212223242526272829303132333435363738394041&cd=1&cols=1
\needspace{10\baselineskip}
\begin{exercice}[BaremeTotal=false]


 Déterminer le coefficient directeur de la tangente à la courbe représentative de la fonction cube au point d'abscisse $-1$.
                        
\end{exercice}

\begin{Solution}
 Le coefficient directeur de la tangente au point d'abscisse $a$ est donné par le nombre dérivé $f'(a)$.\\
        La fonction $f$ définie par $f(x)=x^3$ a pour fonction dérivée la fonction $f'$ définie par $f'(x)=3x^2$.\\
        Comme $f'(-1)=3\times (-1)^2=3$, le coefficient directeur de la tangente au point d'abscisse $-1$ est : $3$. 
\end{Solution}

% @see : https://coopmaths.fr/alea?uuid=ebd89&id=1AN14-1&n=1&d=10&s=3&alea=ocYr22323232323423456789101112131415161718192021222324252627282930313233343536373839404142434445464748234567891011121314151617181920212223242526272829303132333435363738394041&cd=0&cols=1
\needspace{10\baselineskip}
\begin{exercice}[BaremeTotal=false]


 Donner la dérivée de la fonction $f$, dérivable pour tout $x\in\R$, définie par  $f(x)=\dfrac{5}{9}x+7$.
\end{exercice}

\begin{Solution}
 L'expression de la dérivée de la fonction $f$ définie par $f(x)=\dfrac{5}{9}x+7$ est : ${\color[HTML]{f15929}\boldsymbol{f'(x)=\dfrac{5}{9}}}$.
\end{Solution}

% @see : https://coopmaths.fr/alea?uuid=ebd89&id=1AN14-1&n=1&d=10&s=4&alea=uAlP22323232323423456789101112131415161718192021222324252627282930313233343536373839404142434445464748234567891011121314151617181920212223242526272829303132333435363738394041&cd=0&cols=1
\needspace{10\baselineskip}
\begin{exercice}[BaremeTotal=false]


 Donner la dérivée de la fonction $f$, dérivable pour tout $x\in\R$, définie par  $f(x)=-7x^{6}$.
\end{exercice}

\begin{Solution}
 L'expression de la dérivée de la fonction $f$ définie par $f(x)=-7x^{6}$ est : ${\color[HTML]{f15929}\boldsymbol{f'(x)=-42x^{5}}}$.
\end{Solution}

% @see : https://coopmaths.fr/alea?uuid=ebd89&id=1AN14-1&n=1&d=10&s=7&alea=vjN722323232323423456789101112131415161718192021222324252627282930313233343536373839404142434445464748234567891011121314151617181920212223242526272829303132333435363738394041&cd=0&cols=1
\needspace{10\baselineskip}
\begin{exercice}[BaremeTotal=false]


 Donner la dérivée de la fonction $f$, dérivable pour tout $x\in\R^*$, définie par  $f(x)=\dfrac{-4}{9x}$.
\end{exercice}

\begin{Solution}
 L'expression de la dérivée de la fonction $f$ définie par $f(x)=\dfrac{-4}{9x}$ est : ${\color[HTML]{f15929}\boldsymbol{f'(x)=\dfrac{4}{9x^2}}}$.
\end{Solution}

% @see : https://coopmaths.fr/alea?uuid=a83c0&id=1AN14-4&n=1&d=10&s=1&alea=Njr322323232323423456789101112131415161718192021222324252627282930313233343536373839404142434445464748234567891011121314151617181920212223242526272829303132333435363738394041&cd=0&cols=1
\needspace{10\baselineskip}
\begin{exercice}[BaremeTotal=false]


 Donner l'expression de la dérivée de la fonction $f$ définie sur $\R^{*}$ par $f(x)=-\dfrac{2}{x}-2x^{2}-3x+3$.\\
\end{exercice}

\begin{Solution}
 L'expression de la dérivée de la fonction $f$ définie par $f(x)=-\dfrac{2}{x}-2x^{2}-3x+3$ est :\\$f'(x)={\color[HTML]{f15929}\boldsymbol{\dfrac{2}{x^2}-4x-3}}$.\\
\end{Solution}

% @see : https://coopmaths.fr/alea?uuid=a83c0&id=1AN14-4&n=1&d=10&s=4&alea=4Vpz22323232323423456789101112131415161718192021222324252627282930313233343536373839404142434445464748234567891011121314151617181920212223242526272829303132333435363738394041&cd=1&cols=1
\needspace{10\baselineskip}
\begin{exercice}[BaremeTotal=false]


 Donner l'expression de la dérivée de la fonction $f$ définie sur $\R_+^{*}$ par $f(x)=-\dfrac{1}{x}-4\sqrt{x}+3x$.\\
\end{exercice}

\begin{Solution}
 $(-\dfrac{1}{x})^\prime=\dfrac{1}{x^2}$\\$(-4\sqrt{x})^\prime=-\dfrac{4}{2\sqrt{x}}$\\$(3x)^\prime=3$\\L'expression de la dérivée de la fonction $f$ définie par $f(x)=-\dfrac{1}{x}-4\sqrt{x}+3x$ est :\\$f'(x)={\color[HTML]{f15929}\boldsymbol{\dfrac{1}{x^2}-\dfrac{4}{2\sqrt{x}}+3}}$.\\
\end{Solution}

\end{Maquette}
\clearpage
\begin{Maquette}[DM]{Niveau= ,Classe=\getdata[42]\mydata\ (v42), Date = 06/01/25  ,Theme=Exercices}


% @see : https://coopmaths.fr/alea?uuid=4c8c7&id=1AN11-3&n=1&d=10&s=2&alea=nawO2232323232342345678910111213141516171819202122232425262728293031323334353637383940414243444546474823456789101112131415161718192021222324252627282930313233343536373839404142&cd=1&cols=2
\needspace{10\baselineskip}
\begin{exercice}[BaremeTotal=false]
\label{eleve:42}


 \begin{minipage}[t]{\linewidth} Soit $f$ une fonction dérivable sur $[-5;5]$ et $\mathcal{C}_f$ sa courbe représentative.\\On sait que  $f(1)=-5~~$ et que $~~f'(1)=-3$.\\Déterminer une équation de la tangente $(T)$ à la courbe $\mathcal{C}_f$ au point d'abscisse $1$,\\sans utiliser la formule de cours de l'équation de tangente. \end{minipage}

\end{exercice}

\begin{Solution}
 \begin{minipage}[t]{\linewidth} On sait que la tangente n'est pas une droite verticale, puisque la fonction est dérivable sur l'intervalle.\\On en déduit que la tangente $(T)$ au point d'abscisse $1$, admet une équation réduite de la forme :  $(T) : y=m x + p$. 

\medskip
$\bullet$ Détermination de $m$ :\\On sait que le nombre dérivé en $1$ est par définition, le coefficient directeur de la tangente au point d'abscisse $1$.\\Par conséquent, on a déjà : $m=f'(1)=-3$.\\ On en déduit que  $(T) : y= -3 x + p$\\$\bullet$ Détermination de $p$ :\\ Pour cela, on utilise que si $f(1)=-5~~$, alors le point $A$ de coordonnées $(1;-5)$ appartient à $\mathcal{C}_f$ mais aussi à $(T)$.\\ On peut écrire $A(1;-5) = \mathcal{C}_f \cap (T)$.\\ On remplace alors les coordonnées de $A(1;-5)$ dans l'équation  $(T) : y= -3 x + p$.\\ $\begin{aligned} \phantom{\iff}&A(1;-5)\in (T)\\ \iff& -5= -3 \times 1  + p\\ \iff& p=-5 +3 \times 1  \\ \iff& p=-5 +3   \\ \iff& p=-2\\\end{aligned}$\\On peut conclure que : $(T) : {\color[HTML]{f15929}\boldsymbol{y=-3x-2}}$. \end{minipage}

\end{Solution}

% @see : https://coopmaths.fr/alea?uuid=0e984&id=can1F15&n=1&d=10&alea=1Alo2232323232342345678910111213141516171819202122232425262728293031323334353637383940414243444546474823456789101112131415161718192021222324252627282930313233343536373839404142&cd=1&cols=1
\needspace{10\baselineskip}
\begin{exercice}[BaremeTotal=false]


 La courbe représente une fonction $f$ et la droite est la tangente au point d'abscisse $0$.\\
        Déterminer $f'(0)$.\begin{center}\begin{tikzpicture}[baseline,scale = 0.5]

    \tikzset{
      point/.style={
        thick,
        draw,
        cross out,
        inner sep=0pt,
        minimum width=5pt,
        minimum height=5pt,
      },
    }
    \clip (-4,-2) rectangle (12,12);
    	
	\draw[color ={black},line width = 2,>=latex,->] (-4,0)--(12,0);
	\draw[color ={black},line width = 2,>=latex,->] (0,-2)--(0,12);
	\draw[color ={black},opacity = 0.5] (-4,1)--(12,1);
	\draw[color ={black},opacity = 0.5] (-4,2)--(12,2);
	\draw[color ={black},opacity = 0.5] (-4,3)--(12,3);
	\draw[color ={black},opacity = 0.5] (-4,4)--(12,4);
	\draw[color ={black},opacity = 0.5] (-4,5)--(12,5);
	\draw[color ={black},opacity = 0.5] (-4,6)--(12,6);
	\draw[color ={black},opacity = 0.5] (-4,7)--(12,7);
	\draw[color ={black},opacity = 0.5] (-4,8)--(12,8);
	\draw[color ={black},opacity = 0.5] (-4,9)--(12,9);
	\draw[color ={black},opacity = 0.5] (-4,10)--(12,10);
	\draw[color ={black},opacity = 0.5] (-4,11)--(12,11);
	\draw[color ={black},opacity = 0.5] (-4,12)--(12,12);
	\draw[color ={black},opacity = 0.5] (-4,-1)--(12,-1);
	\draw[color ={black},opacity = 0.5] (-4,-2)--(12,-2);
	\draw[color ={black},opacity = 0.5] (1,-2)--(1,12);
	\draw[color ={black},opacity = 0.5] (2,-2)--(2,12);
	\draw[color ={black},opacity = 0.5] (3,-2)--(3,12);
	\draw[color ={black},opacity = 0.5] (4,-2)--(4,12);
	\draw[color ={black},opacity = 0.5] (5,-2)--(5,12);
	\draw[color ={black},opacity = 0.5] (6,-2)--(6,12);
	\draw[color ={black},opacity = 0.5] (7,-2)--(7,12);
	\draw[color ={black},opacity = 0.5] (8,-2)--(8,12);
	\draw[color ={black},opacity = 0.5] (9,-2)--(9,12);
	\draw[color ={black},opacity = 0.5] (10,-2)--(10,12);
	\draw[color ={black},opacity = 0.5] (11,-2)--(11,12);
	\draw[color ={black},opacity = 0.5] (12,-2)--(12,12);
	\draw[color ={black},opacity = 0.5] (-1,-2)--(-1,12);
	\draw[color ={black},opacity = 0.5] (-2,-2)--(-2,12);
	\draw[color ={black},opacity = 0.5] (-3,-2)--(-3,12);
	\draw[color ={black},opacity = 0.5] (-4,-2)--(-4,12);
	\draw[color ={gray},opacity = 0.3] (-4,0)--(12,0);
	\draw[color ={gray},opacity = 0.3] (-4,0)--(12,0);
	\draw[color ={gray},opacity = 0.3] (0,-2)--(0,12);
	\draw[color ={gray},opacity = 0.3] (0,-2)--(0,12);
	\draw[color ={black},line width = 2] (2,-0.2)--(2,0.2);
	\draw[color ={black},line width = 2] (4,-0.2)--(4,0.2);
	\draw[color ={black},line width = 2] (6,-0.2)--(6,0.2);
	\draw[color ={black},line width = 2] (8,-0.2)--(8,0.2);
	\draw[color ={black},line width = 2] (10,-0.2)--(10,0.2);
	\draw[color ={black},line width = 2] (-2,-0.2)--(-2,0.2);
	\draw[color ={black},line width = 2] (-0.2,2)--(0.2,2);
	\draw[color ={black},line width = 2] (-0.2,4)--(0.2,4);
	\draw[color ={black},line width = 2] (-0.2,6)--(0.2,6);
	\draw[color ={black},line width = 2] (-0.2,8)--(0.2,8);
	\draw[color ={black},line width = 2] (-0.2,10)--(0.2,10);
	\draw (2,-0.4) node[anchor = center] {\scriptsize \color{black}{$1$}};
	\draw (4,-0.4) node[anchor = center] {\scriptsize \color{black}{$2$}};
	\draw (6,-0.4) node[anchor = center] {\scriptsize \color{black}{$3$}};
	\draw (8,-0.4) node[anchor = center] {\scriptsize \color{black}{$4$}};
	\draw (10,-0.4) node[anchor = center] {\scriptsize \color{black}{$5$}};
	\draw (-0.5,2.1) node[anchor = center] {\footnotesize \color{black}{$1$}};
	\draw (-0.5,4.1) node[anchor = center] {\footnotesize \color{black}{$2$}};
	\draw (-0.5,6.1) node[anchor = center] {\footnotesize \color{black}{$3$}};
	\draw (-0.5,8.1) node[anchor = center] {\footnotesize \color{black}{$4$}};
	\draw (-0.5,10.1) node[anchor = center] {\footnotesize \color{black}{$5$}};
	\draw [color={black}] (-0.3,-0.3) node[anchor = center,scale=1, rotate = 0] {O};
	
	\draw[color={blue},line width = 2] (-4,1.03)--(-3.6,1.1)--(-3.2,1.17)--(-2.8,1.25)--(-2.4,1.34)--(-2,1.43)--(-1.6,1.53)--(-1.2,1.64)--(-0.8,1.75)--(-0.4,1.87)--(0,2)--(0.4,2.14)--(0.8,2.29)--(1.2,2.44)--(1.6,2.61)--(2,2.79)--(2.4,2.98)--(2.8,3.19)--(3.2,3.41)--(3.6,3.64)--(4,3.9)--(4.4,4.16)--(4.8,4.45)--(5.2,4.76)--(5.6,5.09)--(6,5.44)--(6.4,5.81)--(6.8,6.21)--(7.2,6.64)--(7.6,7.1)--(8,7.59)--(8.4,8.11)--(8.8,8.67)--(9.2,9.27)--(9.6,9.91)--(10,10.59);
	
	\draw[color={red},line width = 2] (-4,0.67)--(-3.6,0.8)--(-3.2,0.93)--(-2.8,1.07)--(-2.4,1.2)--(-2,1.33)--(-1.6,1.47)--(-1.2,1.6)--(-0.8,1.73)--(-0.4,1.87)--(0,2)--(0.4,2.13)--(0.8,2.27)--(1.2,2.4)--(1.6,2.53)--(2,2.67)--(2.4,2.8)--(2.8,2.93)--(3.2,3.07)--(3.6,3.2)--(4,3.33)--(4.4,3.47)--(4.8,3.6)--(5.2,3.73)--(5.6,3.87)--(6,4)--(6.4,4.13)--(6.8,4.27)--(7.2,4.4)--(7.6,4.53)--(8,4.67)--(8.4,4.8)--(8.8,4.93)--(9.2,5.07)--(9.6,5.2)--(10,5.33)--(10.4,5.47)--(10.8,5.6)--(11.2,5.73)--(11.6,5.87)--(12,6);

\end{tikzpicture}\end{center}
\end{exercice}

\begin{Solution}
 $f'(0)$ est donné par le coefficient directeur de la tangente à la courbe au point d'abscisse $0$, soit $\dfrac{1}{3}$.
\end{Solution}

% @see : https://coopmaths.fr/alea?uuid=6f32d&id=can1F16&n=1&d=10&alea=dqSw2232323232342345678910111213141516171819202122232425262728293031323334353637383940414243444546474823456789101112131415161718192021222324252627282930313233343536373839404142&cd=1&cols=1
\needspace{10\baselineskip}
\newpage
\begin{exercice}[BaremeTotal=false]


 On donne les représentations graphiques d'une fonction et de sa dérivée.\\
      Donner l'équation réduite de la tangente à la courbe de $f$ en $x=0$. \\ \begin{minipage}{0.5\linewidth}
    \begin{tikzpicture}[baseline, scale=0.8]

    \tikzset{
      point/.style={
        thick,
        draw,
        cross out,
        inner sep=0pt,
        minimum width=5pt,
        minimum height=5pt,
      },
    }
    \clip (-5.5,-9.5) rectangle (5.75,7.5);
    	
	\draw[color ={black},>=latex,->] (-4,0)--(4,0);
	\draw[color ={black},>=latex,->] (0,-8)--(0,6);
	\draw[color ={black},opacity = 0.4] (-4,1)--(4,1);
	\draw[color ={black},opacity = 0.4] (-4,2)--(4,2);
	\draw[color ={black},opacity = 0.4] (-4,3)--(4,3);
	\draw[color ={black},opacity = 0.4] (-4,4)--(4,4);
	\draw[color ={black},opacity = 0.4] (-4,5)--(4,5);
	\draw[color ={black},opacity = 0.4] (-4,6)--(4,6);
	\draw[color ={black},opacity = 0.4] (-4,-1)--(4,-1);
	\draw[color ={black},opacity = 0.4] (-4,-2)--(4,-2);
	\draw[color ={black},opacity = 0.4] (-4,-3)--(4,-3);
	\draw[color ={black},opacity = 0.4] (-4,-4)--(4,-4);
	\draw[color ={black},opacity = 0.4] (-4,-5)--(4,-5);
	\draw[color ={black},opacity = 0.4] (-4,-6)--(4,-6);
	\draw[color ={black},opacity = 0.4] (-4,-7)--(4,-7);
	\draw[color ={black},opacity = 0.4] (-4,-8)--(4,-8);
	\draw[color ={black},opacity = 0.4] (1,-8)--(1,6);
	\draw[color ={black},opacity = 0.4] (2,-8)--(2,6);
	\draw[color ={black},opacity = 0.4] (3,-8)--(3,6);
	\draw[color ={black},opacity = 0.4] (4,-8)--(4,6);
	\draw[color ={black},opacity = 0.4] (-1,-8)--(-1,6);
	\draw[color ={black},opacity = 0.4] (-2,-8)--(-2,6);
	\draw[color ={black},opacity = 0.4] (-3,-8)--(-3,6);
	\draw[color ={black},opacity = 0.4] (-4,-8)--(-4,6);
	\draw[color ={gray},opacity = 0.3] (-4,0)--(4,0);
	\draw[color ={gray},opacity = 0.3] (-4,0.5)--(4,0.5);
	\draw[color ={gray},opacity = 0.3] (-4,1.5)--(4,1.5);
	\draw[color ={gray},opacity = 0.3] (-4,2.5)--(4,2.5);
	\draw[color ={gray},opacity = 0.3] (-4,3.5)--(4,3.5);
	\draw[color ={gray},opacity = 0.3] (-4,4.5)--(4,4.5);
	\draw[color ={gray},opacity = 0.3] (-4,5.5)--(4,5.5);
	\draw[color ={gray},opacity = 0.3] (-4,0)--(4,0);
	\draw[color ={gray},opacity = 0.3] (-4,-0.5)--(4,-0.5);
	\draw[color ={gray},opacity = 0.3] (-4,-1.5)--(4,-1.5);
	\draw[color ={gray},opacity = 0.3] (-4,-2.5)--(4,-2.5);
	\draw[color ={gray},opacity = 0.3] (-4,-3.5)--(4,-3.5);
	\draw[color ={gray},opacity = 0.3] (-4,-4.5)--(4,-4.5);
	\draw[color ={gray},opacity = 0.3] (-4,-5.5)--(4,-5.5);
	\draw[color ={gray},opacity = 0.3] (-4,-6.5)--(4,-6.5);
	\draw[color ={gray},opacity = 0.3] (-4,-7)--(4,-7);
	\draw[color ={gray},opacity = 0.3] (-4,-7.5)--(4,-7.5);
	\draw[color ={gray},opacity = 0.3] (-4,-8)--(4,-8);
	\draw[color ={gray},opacity = 0.3] (0,-8)--(0,6);
	\draw[color ={gray},opacity = 0.3] (0,-8)--(0,6);
	\draw[color ={black},line width = 1.2] (1,-0.13)--(1,0.13);
	\draw[color ={black},line width = 1.2] (2,-0.13)--(2,0.13);
	\draw[color ={black},line width = 1.2] (3,-0.13)--(3,0.13);
	\draw[color ={black},line width = 1.2] (4,-0.13)--(4,0.13);
	\draw[color ={black},line width = 1.2] (-1,-0.13)--(-1,0.13);
	\draw[color ={black},line width = 1.2] (-2,-0.13)--(-2,0.13);
	\draw[color ={black},line width = 1.2] (-3,-0.13)--(-3,0.13);
	\draw[color ={black},line width = 1.2] (-4,-0.13)--(-4,0.13);
	\draw[color ={black},line width = 1.2] (-0.13,1)--(0.13,1);
	\draw[color ={black},line width = 1.2] (-0.13,2)--(0.13,2);
	\draw[color ={black},line width = 1.2] (-0.13,3)--(0.13,3);
	\draw[color ={black},line width = 1.2] (-0.13,4)--(0.13,4);
	\draw[color ={black},line width = 1.2] (-0.13,5)--(0.13,5);
	\draw[color ={black},line width = 1.2] (-0.13,6)--(0.13,6);
	\draw[color ={black},line width = 1.2] (-0.13,-1)--(0.13,-1);
	\draw[color ={black},line width = 1.2] (-0.13,-2)--(0.13,-2);
	\draw[color ={black},line width = 1.2] (-0.13,-3)--(0.13,-3);
	\draw[color ={black},line width = 1.2] (-0.13,-4)--(0.13,-4);
	\draw[color ={black},line width = 1.2] (-0.13,-5)--(0.13,-5);
	\draw[color ={black},line width = 1.2] (-0.13,-6)--(0.13,-6);
	\draw[color ={black},line width = 1.2] (-0.13,-7)--(0.13,-7);
	\draw[color ={black},line width = 1.2] (-0.13,-8)--(0.13,-8);
	\draw (2,-0.4) node[anchor = center] {\scriptsize \color{black}{$2$}};
	\draw (4,-0.4) node[anchor = center] {\scriptsize \color{black}{$4$}};
	\draw (-2,-0.4) node[anchor = center] {\scriptsize \color{black}{$-2$}};
	\draw (-4,-0.4) node[anchor = center] {\scriptsize \color{black}{$-4$}};
	\draw (-0.5,2.1) node[anchor = center] {\footnotesize \color{black}{$2$}};
	\draw (-0.5,4.1) node[anchor = center] {\footnotesize \color{black}{$4$}};
	\draw (-0.5,6.1) node[anchor = center] {\footnotesize \color{black}{$6$}};
	\draw (-0.5,-1.9) node[anchor = center] {\footnotesize \color{black}{$-2$}};
	\draw (-0.5,-3.9) node[anchor = center] {\footnotesize \color{black}{$-4$}};
	\draw (-0.5,-5.9) node[anchor = center] {\footnotesize \color{black}{$-6$}};
	\draw (-0.5,-7.9) node[anchor = center] {\footnotesize \color{black}{$-8$}};
	\draw [color={black}] (-0.3,-0.3) node[anchor = center,scale=1, rotate = 0] {O};
	\draw (3,4) node[anchor = center] {\footnotesize \color{blue}{$\Large \mathcal{C}_f$}};
	
	\draw[color={blue},line width = 2] (-2.6,-8.96)--(-2.4,-7.56)--(-2.2,-6.24)--(-2,-5)--(-1.8,-3.84)--(-1.6,-2.76)--(-1.4,-1.76)--(-1.2,-0.84)--(-1,0)--(-0.8,0.76)--(-0.6,1.44)--(-0.4,2.04)--(-0.2,2.56)--(0,3)--(0.2,3.36)--(0.4,3.64)--(0.6,3.84)--(0.8,3.96)--(1,4)--(1.2,3.96)--(1.4,3.84)--(1.6,3.64)--(1.8,3.36)--(2,3)--(2.2,2.56)--(2.4,2.04)--(2.6,1.44)--(2.8,0.76)--(3,0)--(3.2,-0.84)--(3.4,-1.76)--(3.6,-2.76)--(3.8,-3.84)--(4,-5);

\end{tikzpicture}\\
    \end{minipage}
    \begin{minipage}{0.5\linewidth}
    \begin{tikzpicture}[baseline, scale=0.8]

    \tikzset{
      point/.style={
        thick,
        draw,
        cross out,
        inner sep=0pt,
        minimum width=5pt,
        minimum height=5pt,
      },
    }
    \clip (-5.5,-9.5) rectangle (6.5,7.5);
    	
	\draw[color ={black},>=latex,->] (-4,0)--(4,0);
	\draw[color ={black},>=latex,->] (0,-8)--(0,6);
	\draw[color ={black},opacity = 0.4] (-4,1)--(4,1);
	\draw[color ={black},opacity = 0.4] (-4,2)--(4,2);
	\draw[color ={black},opacity = 0.4] (-4,3)--(4,3);
	\draw[color ={black},opacity = 0.4] (-4,4)--(4,4);
	\draw[color ={black},opacity = 0.4] (-4,5)--(4,5);
	\draw[color ={black},opacity = 0.4] (-4,6)--(4,6);
	\draw[color ={black},opacity = 0.4] (-4,-1)--(4,-1);
	\draw[color ={black},opacity = 0.4] (-4,-2)--(4,-2);
	\draw[color ={black},opacity = 0.4] (-4,-3)--(4,-3);
	\draw[color ={black},opacity = 0.4] (-4,-4)--(4,-4);
	\draw[color ={black},opacity = 0.4] (-4,-5)--(4,-5);
	\draw[color ={black},opacity = 0.4] (-4,-6)--(4,-6);
	\draw[color ={black},opacity = 0.4] (-4,-7)--(4,-7);
	\draw[color ={black},opacity = 0.4] (-4,-8)--(4,-8);
	\draw[color ={black},opacity = 0.4] (1,-8)--(1,6);
	\draw[color ={black},opacity = 0.4] (2,-8)--(2,6);
	\draw[color ={black},opacity = 0.4] (3,-8)--(3,6);
	\draw[color ={black},opacity = 0.4] (4,-8)--(4,6);
	\draw[color ={black},opacity = 0.4] (-1,-8)--(-1,6);
	\draw[color ={black},opacity = 0.4] (-2,-8)--(-2,6);
	\draw[color ={black},opacity = 0.4] (-3,-8)--(-3,6);
	\draw[color ={black},opacity = 0.4] (-4,-8)--(-4,6);
	\draw[color ={gray},opacity = 0.3] (-4,0)--(4,0);
	\draw[color ={gray},opacity = 0.3] (-4,0.5)--(4,0.5);
	\draw[color ={gray},opacity = 0.3] (-4,1.5)--(4,1.5);
	\draw[color ={gray},opacity = 0.3] (-4,2.5)--(4,2.5);
	\draw[color ={gray},opacity = 0.3] (-4,3.5)--(4,3.5);
	\draw[color ={gray},opacity = 0.3] (-4,4.5)--(4,4.5);
	\draw[color ={gray},opacity = 0.3] (-4,5.5)--(4,5.5);
	\draw[color ={gray},opacity = 0.3] (-4,0)--(4,0);
	\draw[color ={gray},opacity = 0.3] (-4,-0.5)--(4,-0.5);
	\draw[color ={gray},opacity = 0.3] (-4,-1.5)--(4,-1.5);
	\draw[color ={gray},opacity = 0.3] (-4,-2.5)--(4,-2.5);
	\draw[color ={gray},opacity = 0.3] (-4,-3.5)--(4,-3.5);
	\draw[color ={gray},opacity = 0.3] (-4,-4.5)--(4,-4.5);
	\draw[color ={gray},opacity = 0.3] (-4,-5.5)--(4,-5.5);
	\draw[color ={gray},opacity = 0.3] (-4,-6.5)--(4,-6.5);
	\draw[color ={gray},opacity = 0.3] (-4,-7)--(4,-7);
	\draw[color ={gray},opacity = 0.3] (-4,-7.5)--(4,-7.5);
	\draw[color ={gray},opacity = 0.3] (-4,-8)--(4,-8);
	\draw[color ={gray},opacity = 0.3] (0,-8)--(0,6);
	\draw[color ={gray},opacity = 0.3] (0,-8)--(0,6);
	\draw[color ={black},line width = 1.2] (1,-0.13)--(1,0.13);
	\draw[color ={black},line width = 1.2] (2,-0.13)--(2,0.13);
	\draw[color ={black},line width = 1.2] (3,-0.13)--(3,0.13);
	\draw[color ={black},line width = 1.2] (4,-0.13)--(4,0.13);
	\draw[color ={black},line width = 1.2] (-1,-0.13)--(-1,0.13);
	\draw[color ={black},line width = 1.2] (-2,-0.13)--(-2,0.13);
	\draw[color ={black},line width = 1.2] (-3,-0.13)--(-3,0.13);
	\draw[color ={black},line width = 1.2] (-4,-0.13)--(-4,0.13);
	\draw[color ={black},line width = 1.2] (-0.13,1)--(0.13,1);
	\draw[color ={black},line width = 1.2] (-0.13,2)--(0.13,2);
	\draw[color ={black},line width = 1.2] (-0.13,3)--(0.13,3);
	\draw[color ={black},line width = 1.2] (-0.13,4)--(0.13,4);
	\draw[color ={black},line width = 1.2] (-0.13,5)--(0.13,5);
	\draw[color ={black},line width = 1.2] (-0.13,6)--(0.13,6);
	\draw[color ={black},line width = 1.2] (-0.13,-1)--(0.13,-1);
	\draw[color ={black},line width = 1.2] (-0.13,-2)--(0.13,-2);
	\draw[color ={black},line width = 1.2] (-0.13,-3)--(0.13,-3);
	\draw[color ={black},line width = 1.2] (-0.13,-4)--(0.13,-4);
	\draw[color ={black},line width = 1.2] (-0.13,-5)--(0.13,-5);
	\draw[color ={black},line width = 1.2] (-0.13,-6)--(0.13,-6);
	\draw[color ={black},line width = 1.2] (-0.13,-7)--(0.13,-7);
	\draw[color ={black},line width = 1.2] (-0.13,-8)--(0.13,-8);
	\draw (2,-0.4) node[anchor = center] {\scriptsize \color{black}{$2$}};
	\draw (4,-0.4) node[anchor = center] {\scriptsize \color{black}{$4$}};
	\draw (-2,-0.4) node[anchor = center] {\scriptsize \color{black}{$-2$}};
	\draw (-4,-0.4) node[anchor = center] {\scriptsize \color{black}{$-4$}};
	\draw (-0.5,2.1) node[anchor = center] {\footnotesize \color{black}{$2$}};
	\draw (-0.5,4.1) node[anchor = center] {\footnotesize \color{black}{$4$}};
	\draw (-0.5,-1.9) node[anchor = center] {\footnotesize \color{black}{$-2$}};
	\draw (-0.5,-3.9) node[anchor = center] {\footnotesize \color{black}{$-4$}};
	\draw (-0.5,-5.9) node[anchor = center] {\footnotesize \color{black}{$-6$}};
	\draw (-0.5,-7.9) node[anchor = center] {\footnotesize \color{black}{$-8$}};
	\draw [color={black}] (-0.3,-0.3) node[anchor = center,scale=1, rotate = 0] {O};
	\draw (3,4) node[anchor = center] {\footnotesize \color{red}{$\Large\mathcal{C}_f\prime$}};
	
	\draw[color={red},line width = 2] (-2.4,6.8)--(-2.2,6.4)--(-2,6)--(-1.8,5.6)--(-1.6,5.2)--(-1.4,4.8)--(-1.2,4.4)--(-1,4)--(-0.8,3.6)--(-0.6,3.2)--(-0.4,2.8)--(-0.2,2.4)--(0,2)--(0.2,1.6)--(0.4,1.2)--(0.6,0.8)--(0.8,0.4)--(1,0)--(1.2,-0.4)--(1.4,-0.8)--(1.6,-1.2)--(1.8,-1.6)--(2,-2)--(2.2,-2.4)--(2.4,-2.8)--(2.6,-3.2)--(2.8,-3.6)--(3,-4)--(3.2,-4.4)--(3.4,-4.8)--(3.6,-5.2)--(3.8,-5.6)--(4,-6);

\end{tikzpicture}\\
    \end{minipage}
    
\end{exercice}

\begin{Solution}
 L'équation réduite de la tangente au point d'abscisse $0$ est  : $y=f'(0)(x-0)+f(0)$.\\
      On lit graphiquement $f(0)=3$ et $f'(0)=2$.\\
      L'équation réduite de la tangente est donc donnée par :
      $y=2(x+0)+3$, soit $y=2x+3$.
\end{Solution}

% @see : https://coopmaths.fr/alea?uuid=a1ba2&id=can1F14&n=1&d=10&alea=rkIp2232323232342345678910111213141516171819202122232425262728293031323334353637383940414243444546474823456789101112131415161718192021222324252627282930313233343536373839404142&cd=1&cols=1
\needspace{10\baselineskip}
\begin{exercice}[BaremeTotal=false]


 Soit $f$ la fonction définie sur $\mathbb{R}$ par : $f(x)= 5x^2+2x-4$.\\
        En calculant la limite du taux d'accroissement, déterminer $f'(0)$.
\end{exercice}

\begin{Solution}
 $f$ est une fonction polynôme du second degré de la forme $f(x)=ax^2+bx+c$.\\
    La fonction dérivée est donnée par la somme des dérivées des fonctions $u$ et $v$ définies par $u(x)=5x^2$ et $v(x)=2x-4$.\\
     Comme $u'(x)=10x$ et $v'(x)=2$, on obtient  $f'(x)=10x+2$. \\
     Ainsi, $f'(0)=10\times 0+2=2$.
\end{Solution}

% @see : https://coopmaths.fr/alea?uuid=3c690&id=can1F13&n=1&d=10&alea=EUDu2232323232342345678910111213141516171819202122232425262728293031323334353637383940414243444546474823456789101112131415161718192021222324252627282930313233343536373839404142&cd=1&cols=1
\needspace{10\baselineskip}
\begin{exercice}[BaremeTotal=false]


 Déterminer le coefficient directeur de la tangente à la courbe représentative de la fonction cube au point d'abscisse $1$.
                        
\end{exercice}

\begin{Solution}
 Le coefficient directeur de la tangente au point d'abscisse $a$ est donné par le nombre dérivé $f'(a)$.\\
        La fonction $f$ définie par $f(x)=x^3$ a pour fonction dérivée la fonction $f'$ définie par $f'(x)=3x^2$.\\
        Comme $f'(1)=3\times 1^2=3$, le coefficient directeur de la tangente au point d'abscisse $1$ est : $3$. 
\end{Solution}

% @see : https://coopmaths.fr/alea?uuid=ebd89&id=1AN14-1&n=1&d=10&s=3&alea=ocYr2232323232342345678910111213141516171819202122232425262728293031323334353637383940414243444546474823456789101112131415161718192021222324252627282930313233343536373839404142&cd=0&cols=1
\needspace{10\baselineskip}
\begin{exercice}[BaremeTotal=false]


 Donner la dérivée de la fonction $f$, dérivable pour tout $x\in\R$, définie par  $f(x)=-\dfrac{5}{8}x+9$.
\end{exercice}

\begin{Solution}
 L'expression de la dérivée de la fonction $f$ définie par $f(x)=-\dfrac{5}{8}x+9$ est : ${\color[HTML]{f15929}\boldsymbol{f'(x)=-\dfrac{5}{8}}}$.
\end{Solution}

% @see : https://coopmaths.fr/alea?uuid=ebd89&id=1AN14-1&n=1&d=10&s=4&alea=uAlP2232323232342345678910111213141516171819202122232425262728293031323334353637383940414243444546474823456789101112131415161718192021222324252627282930313233343536373839404142&cd=0&cols=1
\needspace{10\baselineskip}
\begin{exercice}[BaremeTotal=false]


 Donner la dérivée de la fonction $f$, dérivable pour tout $x\in\R$, définie par  $f(x)=-5x^{11}$.
\end{exercice}

\begin{Solution}
 L'expression de la dérivée de la fonction $f$ définie par $f(x)=-5x^{11}$ est : ${\color[HTML]{f15929}\boldsymbol{f'(x)=-55x^{10}}}$.
\end{Solution}

% @see : https://coopmaths.fr/alea?uuid=ebd89&id=1AN14-1&n=1&d=10&s=7&alea=vjN72232323232342345678910111213141516171819202122232425262728293031323334353637383940414243444546474823456789101112131415161718192021222324252627282930313233343536373839404142&cd=0&cols=1
\needspace{10\baselineskip}
\begin{exercice}[BaremeTotal=false]


 Donner la dérivée de la fonction $f$, dérivable pour tout $x\in\R^*$, définie par  $f(x)=\dfrac{-5}{7x}$.
\end{exercice}

\begin{Solution}
 L'expression de la dérivée de la fonction $f$ définie par $f(x)=\dfrac{-5}{7x}$ est : ${\color[HTML]{f15929}\boldsymbol{f'(x)=\dfrac{5}{7x^2}}}$.
\end{Solution}

% @see : https://coopmaths.fr/alea?uuid=a83c0&id=1AN14-4&n=1&d=10&s=1&alea=Njr32232323232342345678910111213141516171819202122232425262728293031323334353637383940414243444546474823456789101112131415161718192021222324252627282930313233343536373839404142&cd=0&cols=1
\needspace{10\baselineskip}
\begin{exercice}[BaremeTotal=false]


 Donner l'expression de la dérivée de la fonction $f$ définie sur $\R^{*}$ par $f(x)=-4x^{2}+\dfrac{8}{x}$.\\
\end{exercice}

\begin{Solution}
 L'expression de la dérivée de la fonction $f$ définie par $f(x)=-4x^{2}+\dfrac{8}{x}$ est :\\$f'(x)={\color[HTML]{f15929}\boldsymbol{-8x-\dfrac{8}{x^2}}}$.\\
\end{Solution}

% @see : https://coopmaths.fr/alea?uuid=a83c0&id=1AN14-4&n=1&d=10&s=4&alea=4Vpz2232323232342345678910111213141516171819202122232425262728293031323334353637383940414243444546474823456789101112131415161718192021222324252627282930313233343536373839404142&cd=1&cols=1
\needspace{10\baselineskip}
\begin{exercice}[BaremeTotal=false]


 Donner l'expression de la dérivée de la fonction $f$ définie sur $\R_+^{*}$ par $f(x)=-\dfrac{3}{x}+4x^{2}-4x-\sqrt{x}$.\\
\end{exercice}

\begin{Solution}
 $(-\dfrac{3}{x})^\prime=\dfrac{3}{x^2}$\\$(4x^{2}-4x)^\prime=8x-4$\\$(-\sqrt{x})^\prime=-\dfrac{1}{2\sqrt{x}}$\\L'expression de la dérivée de la fonction $f$ définie par $f(x)=-\dfrac{3}{x}+4x^{2}-4x-\sqrt{x}$ est :\\$f'(x)={\color[HTML]{f15929}\boldsymbol{\dfrac{3}{x^2}+8x-4-\dfrac{1}{2\sqrt{x}}}}$.\\
\end{Solution}

\end{Maquette}
\clearpage
\begin{Maquette}[DM]{Niveau= ,Classe=\getdata[43]\mydata\ (v43), Date = 06/01/25  ,Theme=Exercices}


% @see : https://coopmaths.fr/alea?uuid=4c8c7&id=1AN11-3&n=1&d=10&s=2&alea=nawO223232323234234567891011121314151617181920212223242526272829303132333435363738394041424344454647482345678910111213141516171819202122232425262728293031323334353637383940414243&cd=1&cols=2
\needspace{10\baselineskip}
\begin{exercice}[BaremeTotal=false]
\label{eleve:43}


 \begin{minipage}[t]{\linewidth} Soit $f$ une fonction dérivable sur $[-5;5]$ et $\mathcal{C}_f$ sa courbe représentative.\\On sait que  $f(-4)=-4~~$ et que $~~f'(-4)=-3$.\\Déterminer une équation de la tangente $(T)$ à la courbe $\mathcal{C}_f$ au point d'abscisse $-4$,\\sans utiliser la formule de cours de l'équation de tangente. \end{minipage}

\end{exercice}

\begin{Solution}
 \begin{minipage}[t]{\linewidth} On sait que la tangente n'est pas une droite verticale, puisque la fonction est dérivable sur l'intervalle.\\On en déduit que la tangente $(T)$ au point d'abscisse $-4$, admet une équation réduite de la forme :  $(T) : y=m x + p$. 

\medskip
$\bullet$ Détermination de $m$ :\\On sait que le nombre dérivé en $-4$ est par définition, le coefficient directeur de la tangente au point d'abscisse $-4$.\\Par conséquent, on a déjà : $m=f'(-4)=-3$.\\ On en déduit que  $(T) : y= -3 x + p$\\$\bullet$ Détermination de $p$ :\\ Pour cela, on utilise que si $f(-4)=-4~~$, alors le point $A$ de coordonnées $(-4;-4)$ appartient à $\mathcal{C}_f$ mais aussi à $(T)$.\\ On peut écrire $A(-4;-4) = \mathcal{C}_f \cap (T)$.\\ On remplace alors les coordonnées de $A(-4;-4)$ dans l'équation  $(T) : y= -3 x + p$.\\ $\begin{aligned} \phantom{\iff}&A(-4;-4)\in (T)\\ \iff& -4= -3 \times (-4)  + p\\ \iff& p=-4 +3 \times (-4)  \\ \iff& p=-4 -12   \\ \iff& p=-16\\\end{aligned}$\\On peut conclure que : $(T) : {\color[HTML]{f15929}\boldsymbol{y=-3x-16}}$. \end{minipage}

\end{Solution}

% @see : https://coopmaths.fr/alea?uuid=0e984&id=can1F15&n=1&d=10&alea=1Alo223232323234234567891011121314151617181920212223242526272829303132333435363738394041424344454647482345678910111213141516171819202122232425262728293031323334353637383940414243&cd=1&cols=1
\needspace{10\baselineskip}
\begin{exercice}[BaremeTotal=false]


 La courbe représente une fonction $f$ et la droite est la tangente au point d'abscisse $0$.\\
        Déterminer $f'(0)$.\begin{center}\begin{tikzpicture}[baseline,scale = 0.5]

    \tikzset{
      point/.style={
        thick,
        draw,
        cross out,
        inner sep=0pt,
        minimum width=5pt,
        minimum height=5pt,
      },
    }
    \clip (-10,-2) rectangle (4,12);
    	
	\draw[color ={black},line width = 2,>=latex,->] (-10,0)--(4,0);
	\draw[color ={black},line width = 2,>=latex,->] (0,-2)--(0,12);
	\draw[color ={black},opacity = 0.5] (-10,1)--(4,1);
	\draw[color ={black},opacity = 0.5] (-10,2)--(4,2);
	\draw[color ={black},opacity = 0.5] (-10,3)--(4,3);
	\draw[color ={black},opacity = 0.5] (-10,4)--(4,4);
	\draw[color ={black},opacity = 0.5] (-10,5)--(4,5);
	\draw[color ={black},opacity = 0.5] (-10,6)--(4,6);
	\draw[color ={black},opacity = 0.5] (-10,7)--(4,7);
	\draw[color ={black},opacity = 0.5] (-10,8)--(4,8);
	\draw[color ={black},opacity = 0.5] (-10,9)--(4,9);
	\draw[color ={black},opacity = 0.5] (-10,10)--(4,10);
	\draw[color ={black},opacity = 0.5] (-10,11)--(4,11);
	\draw[color ={black},opacity = 0.5] (-10,12)--(4,12);
	\draw[color ={black},opacity = 0.5] (-10,-1)--(4,-1);
	\draw[color ={black},opacity = 0.5] (-10,-2)--(4,-2);
	\draw[color ={black},opacity = 0.5] (1,-2)--(1,12);
	\draw[color ={black},opacity = 0.5] (2,-2)--(2,12);
	\draw[color ={black},opacity = 0.5] (3,-2)--(3,12);
	\draw[color ={black},opacity = 0.5] (4,-2)--(4,12);
	\draw[color ={black},opacity = 0.5] (-1,-2)--(-1,12);
	\draw[color ={black},opacity = 0.5] (-2,-2)--(-2,12);
	\draw[color ={black},opacity = 0.5] (-3,-2)--(-3,12);
	\draw[color ={black},opacity = 0.5] (-4,-2)--(-4,12);
	\draw[color ={black},opacity = 0.5] (-5,-2)--(-5,12);
	\draw[color ={black},opacity = 0.5] (-6,-2)--(-6,12);
	\draw[color ={black},opacity = 0.5] (-7,-2)--(-7,12);
	\draw[color ={black},opacity = 0.5] (-8,-2)--(-8,12);
	\draw[color ={black},opacity = 0.5] (-9,-2)--(-9,12);
	\draw[color ={black},opacity = 0.5] (-10,-2)--(-10,12);
	\draw[color ={gray},opacity = 0.3] (-10,0)--(4,0);
	\draw[color ={gray},opacity = 0.3] (-10,0)--(4,0);
	\draw[color ={gray},opacity = 0.3] (0,-2)--(0,12);
	\draw[color ={gray},opacity = 0.3] (0,-2)--(0,12);
	\draw[color ={gray},opacity = 0.3] (-5,-2)--(-5,12);
	\draw[color ={gray},opacity = 0.3] (-6,-2)--(-6,12);
	\draw[color ={gray},opacity = 0.3] (-7,-2)--(-7,12);
	\draw[color ={gray},opacity = 0.3] (-8,-2)--(-8,12);
	\draw[color ={gray},opacity = 0.3] (-9,-2)--(-9,12);
	\draw[color ={gray},opacity = 0.3] (-10,-2)--(-10,12);
	\draw[color ={black},line width = 2] (2,-0.2)--(2,0.2);
	\draw[color ={black},line width = 2] (-2,-0.2)--(-2,0.2);
	\draw[color ={black},line width = 2] (-4,-0.2)--(-4,0.2);
	\draw[color ={black},line width = 2] (-6,-0.2)--(-6,0.2);
	\draw[color ={black},line width = 2] (-8,-0.2)--(-8,0.2);
	\draw[color ={black},line width = 2] (-0.2,2)--(0.2,2);
	\draw[color ={black},line width = 2] (-0.2,4)--(0.2,4);
	\draw[color ={black},line width = 2] (-0.2,6)--(0.2,6);
	\draw[color ={black},line width = 2] (-0.2,8)--(0.2,8);
	\draw[color ={black},line width = 2] (-0.2,10)--(0.2,10);
	\draw (2,-0.4) node[anchor = center] {\scriptsize \color{black}{$1$}};
	\draw (-2,-0.4) node[anchor = center] {\scriptsize \color{black}{$-1$}};
	\draw (-4,-0.4) node[anchor = center] {\scriptsize \color{black}{$-2$}};
	\draw (-6,-0.4) node[anchor = center] {\scriptsize \color{black}{$-3$}};
	\draw (-8,-0.4) node[anchor = center] {\scriptsize \color{black}{$-4$}};
	\draw (-0.5,2.1) node[anchor = center] {\footnotesize \color{black}{$1$}};
	\draw (-0.5,4.1) node[anchor = center] {\footnotesize \color{black}{$2$}};
	\draw (-0.5,6.1) node[anchor = center] {\footnotesize \color{black}{$3$}};
	\draw (-0.5,8.1) node[anchor = center] {\footnotesize \color{black}{$4$}};
	\draw (-0.5,10.1) node[anchor = center] {\footnotesize \color{black}{$5$}};
	\draw [color={black}] (-0.3,-0.3) node[anchor = center,scale=1, rotate = 0] {O};
	
	\draw[color={blue},line width = 2] (-1.6,9.91)--(-1.2,6.64)--(-0.8,4.45)--(-0.4,2.98)--(0,2)--(0.4,1.34)--(0.8,0.9)--(1.2,0.6)--(1.6,0.4)--(2,0.27)--(2.4,0.18)--(2.8,0.12)--(3.2,0.08)--(3.6,0.05)--(4,0.04);
	
	\draw[color={red},line width = 2] (-6,14)--(-5.6,13.2)--(-5.2,12.4)--(-4.8,11.6)--(-4.4,10.8)--(-4,10)--(-3.6,9.2)--(-3.2,8.4)--(-2.8,7.6)--(-2.4,6.8)--(-2,6)--(-1.6,5.2)--(-1.2,4.4)--(-0.8,3.6)--(-0.4,2.8)--(0,2)--(0.4,1.2)--(0.8,0.4)--(1.2,-0.4)--(1.6,-1.2)--(2,-2)--(2.4,-2.8)--(2.8,-3.6);

\end{tikzpicture}\end{center}
\end{exercice}

\begin{Solution}
 $f'(0)$ est donné par le coefficient directeur de la tangente à la courbe au point d'abscisse $0$, soit $-2=-2$.
\end{Solution}

% @see : https://coopmaths.fr/alea?uuid=6f32d&id=can1F16&n=1&d=10&alea=dqSw223232323234234567891011121314151617181920212223242526272829303132333435363738394041424344454647482345678910111213141516171819202122232425262728293031323334353637383940414243&cd=1&cols=1
\needspace{10\baselineskip}
\newpage
\begin{exercice}[BaremeTotal=false]


 On donne les représentations graphiques d'une fonction et de sa dérivée.\\
        Donner l'équation réduite de la tangente à la courbe de $f$ en $x=1$. \\ \begin{minipage}{0.5\linewidth}
    \begin{tikzpicture}[baseline, scale=0.8]

    \tikzset{
      point/.style={
        thick,
        draw,
        cross out,
        inner sep=0pt,
        minimum width=5pt,
        minimum height=5pt,
      },
    }
    \clip (-4.5,-4.5) rectangle (5.75,13.5);
    	
	\draw[color ={black},>=latex,->] (-3,0)--(3,0);
	\draw[color ={black},>=latex,->] (0,-3)--(0,12);
	\draw[color ={black},opacity = 0.4] (-3,1)--(3,1);
	\draw[color ={black},opacity = 0.4] (-3,2)--(3,2);
	\draw[color ={black},opacity = 0.4] (-3,3)--(3,3);
	\draw[color ={black},opacity = 0.4] (-3,4)--(3,4);
	\draw[color ={black},opacity = 0.4] (-3,5)--(3,5);
	\draw[color ={black},opacity = 0.4] (-3,6)--(3,6);
	\draw[color ={black},opacity = 0.4] (-3,7)--(3,7);
	\draw[color ={black},opacity = 0.4] (-3,8)--(3,8);
	\draw[color ={black},opacity = 0.4] (-3,9)--(3,9);
	\draw[color ={black},opacity = 0.4] (-3,10)--(3,10);
	\draw[color ={black},opacity = 0.4] (-3,11)--(3,11);
	\draw[color ={black},opacity = 0.4] (-3,12)--(3,12);
	\draw[color ={black},opacity = 0.4] (-3,-1)--(3,-1);
	\draw[color ={black},opacity = 0.4] (-3,-2)--(3,-2);
	\draw[color ={black},opacity = 0.4] (-3,-3)--(3,-3);
	\draw[color ={black},opacity = 0.4] (1,-3)--(1,12);
	\draw[color ={black},opacity = 0.4] (2,-3)--(2,12);
	\draw[color ={black},opacity = 0.4] (3,-3)--(3,12);
	\draw[color ={black},opacity = 0.4] (-1,-3)--(-1,12);
	\draw[color ={black},opacity = 0.4] (-2,-3)--(-2,12);
	\draw[color ={black},opacity = 0.4] (-3,-3)--(-3,12);
	\draw[color ={gray},opacity = 0.3] (-3,0)--(3,0);
	\draw[color ={gray},opacity = 0.3] (-3,0.5)--(3,0.5);
	\draw[color ={gray},opacity = 0.3] (-3,1.5)--(3,1.5);
	\draw[color ={gray},opacity = 0.3] (-3,2.5)--(3,2.5);
	\draw[color ={gray},opacity = 0.3] (-3,3.5)--(3,3.5);
	\draw[color ={gray},opacity = 0.3] (-3,4.5)--(3,4.5);
	\draw[color ={gray},opacity = 0.3] (-3,5.5)--(3,5.5);
	\draw[color ={gray},opacity = 0.3] (-3,6.5)--(3,6.5);
	\draw[color ={gray},opacity = 0.3] (-3,7.5)--(3,7.5);
	\draw[color ={gray},opacity = 0.3] (-3,8.5)--(3,8.5);
	\draw[color ={gray},opacity = 0.3] (-3,9.5)--(3,9.5);
	\draw[color ={gray},opacity = 0.3] (-3,10.5)--(3,10.5);
	\draw[color ={gray},opacity = 0.3] (-3,11.5)--(3,11.5);
	\draw[color ={gray},opacity = 0.3] (-3,0)--(3,0);
	\draw[color ={gray},opacity = 0.3] (-3,-0.5)--(3,-0.5);
	\draw[color ={gray},opacity = 0.3] (-3,-1.5)--(3,-1.5);
	\draw[color ={gray},opacity = 0.3] (-3,-2.5)--(3,-2.5);
	\draw[color ={gray},opacity = 0.3] (0,-3)--(0,12);
	\draw[color ={gray},opacity = 0.3] (0,-3)--(0,12);
	\draw[color ={black},line width = 1.2] (1,-0.13)--(1,0.13);
	\draw[color ={black},line width = 1.2] (2,-0.13)--(2,0.13);
	\draw[color ={black},line width = 1.2] (3,-0.13)--(3,0.13);
	\draw[color ={black},line width = 1.2] (-1,-0.13)--(-1,0.13);
	\draw[color ={black},line width = 1.2] (-2,-0.13)--(-2,0.13);
	\draw[color ={black},line width = 1.2] (-3,-0.13)--(-3,0.13);
	\draw[color ={black},line width = 1.2] (-0.13,1)--(0.13,1);
	\draw[color ={black},line width = 1.2] (-0.13,2)--(0.13,2);
	\draw[color ={black},line width = 1.2] (-0.13,3)--(0.13,3);
	\draw[color ={black},line width = 1.2] (-0.13,4)--(0.13,4);
	\draw[color ={black},line width = 1.2] (-0.13,5)--(0.13,5);
	\draw[color ={black},line width = 1.2] (-0.13,6)--(0.13,6);
	\draw[color ={black},line width = 1.2] (-0.13,7)--(0.13,7);
	\draw[color ={black},line width = 1.2] (-0.13,8)--(0.13,8);
	\draw[color ={black},line width = 1.2] (-0.13,9)--(0.13,9);
	\draw[color ={black},line width = 1.2] (-0.13,10)--(0.13,10);
	\draw[color ={black},line width = 1.2] (-0.13,11)--(0.13,11);
	\draw[color ={black},line width = 1.2] (-0.13,12)--(0.13,12);
	\draw[color ={black},line width = 1.2] (-0.13,-1)--(0.13,-1);
	\draw[color ={black},line width = 1.2] (-0.13,-2)--(0.13,-2);
	\draw[color ={black},line width = 1.2] (-0.13,-3)--(0.13,-3);
	\draw (3,-0.4) node[anchor = center] {\scriptsize \color{black}{$3$}};
	\draw (-3,-0.4) node[anchor = center] {\scriptsize \color{black}{$-3$}};
	\draw (-0.5,3.1) node[anchor = center] {\footnotesize \color{black}{$3$}};
	\draw (-0.5,6.1) node[anchor = center] {\footnotesize \color{black}{$6$}};
	\draw (-0.5,9.1) node[anchor = center] {\footnotesize \color{black}{$9$}};
	\draw (-0.5,-2.9) node[anchor = center] {\footnotesize \color{black}{$-3$}};
	\draw [color={black}] (-0.3,-0.3) node[anchor = center,scale=1, rotate = 0] {O};
	\draw (3,10) node[anchor = center] {\footnotesize \color{blue}{$\Large \mathcal{C}_f$}};
	
	\draw[color={blue},line width = 2] (-3,10)--(-2.8,8.84)--(-2.6,7.76)--(-2.4,6.76)--(-2.2,5.84)--(-2,5)--(-1.8,4.24)--(-1.6,3.56)--(-1.4,2.96)--(-1.2,2.44)--(-1,2)--(-0.8,1.64)--(-0.6,1.36)--(-0.4,1.16)--(-0.2,1.04)--(0,1)--(0.2,1.04)--(0.4,1.16)--(0.6,1.36)--(0.8,1.64)--(1,2)--(1.2,2.44)--(1.4,2.96)--(1.6,3.56)--(1.8,4.24)--(2,5)--(2.2,5.84)--(2.4,6.76)--(2.6,7.76)--(2.8,8.84)--(3,10);

\end{tikzpicture}\\
    \end{minipage}
    \begin{minipage}{0.5\linewidth}
    \begin{tikzpicture}[baseline, scale=0.8]

    \tikzset{
      point/.style={
        thick,
        draw,
        cross out,
        inner sep=0pt,
        minimum width=5pt,
        minimum height=5pt,
      },
    }
    \clip (-4.5,-6.5) rectangle (6.5,9.5);
    	
	\draw[color ={black},>=latex,->] (-3,0)--(3,0);
	\draw[color ={black},>=latex,->] (0,-5)--(0,8);
	\draw[color ={black},opacity = 0.4] (-3,1)--(3,1);
	\draw[color ={black},opacity = 0.4] (-3,2)--(3,2);
	\draw[color ={black},opacity = 0.4] (-3,3)--(3,3);
	\draw[color ={black},opacity = 0.4] (-3,4)--(3,4);
	\draw[color ={black},opacity = 0.4] (-3,5)--(3,5);
	\draw[color ={black},opacity = 0.4] (-3,6)--(3,6);
	\draw[color ={black},opacity = 0.4] (-3,7)--(3,7);
	\draw[color ={black},opacity = 0.4] (-3,8)--(3,8);
	\draw[color ={black},opacity = 0.4] (-3,-1)--(3,-1);
	\draw[color ={black},opacity = 0.4] (-3,-2)--(3,-2);
	\draw[color ={black},opacity = 0.4] (-3,-3)--(3,-3);
	\draw[color ={black},opacity = 0.4] (-3,-4)--(3,-4);
	\draw[color ={black},opacity = 0.4] (-3,-5)--(3,-5);
	\draw[color ={black},opacity = 0.4] (1,-5)--(1,8);
	\draw[color ={black},opacity = 0.4] (2,-5)--(2,8);
	\draw[color ={black},opacity = 0.4] (3,-5)--(3,8);
	\draw[color ={black},opacity = 0.4] (-1,-5)--(-1,8);
	\draw[color ={black},opacity = 0.4] (-2,-5)--(-2,8);
	\draw[color ={black},opacity = 0.4] (-3,-5)--(-3,8);
	\draw[color ={gray},opacity = 0.3] (-3,0)--(3,0);
	\draw[color ={gray},opacity = 0.3] (-3,0.5)--(3,0.5);
	\draw[color ={gray},opacity = 0.3] (-3,1.5)--(3,1.5);
	\draw[color ={gray},opacity = 0.3] (-3,2.5)--(3,2.5);
	\draw[color ={gray},opacity = 0.3] (-3,3.5)--(3,3.5);
	\draw[color ={gray},opacity = 0.3] (-3,4.5)--(3,4.5);
	\draw[color ={gray},opacity = 0.3] (-3,5.5)--(3,5.5);
	\draw[color ={gray},opacity = 0.3] (-3,6.5)--(3,6.5);
	\draw[color ={gray},opacity = 0.3] (-3,7.5)--(3,7.5);
	\draw[color ={gray},opacity = 0.3] (-3,0)--(3,0);
	\draw[color ={gray},opacity = 0.3] (-3,-0.5)--(3,-0.5);
	\draw[color ={gray},opacity = 0.3] (-3,-1.5)--(3,-1.5);
	\draw[color ={gray},opacity = 0.3] (-3,-2.5)--(3,-2.5);
	\draw[color ={gray},opacity = 0.3] (-3,-3.5)--(3,-3.5);
	\draw[color ={gray},opacity = 0.3] (-3,-4.5)--(3,-4.5);
	\draw[color ={gray},opacity = 0.3] (0,-5)--(0,8);
	\draw[color ={gray},opacity = 0.3] (0,-5)--(0,8);
	\draw[color ={black},line width = 1.2] (1,-0.13)--(1,0.13);
	\draw[color ={black},line width = 1.2] (2,-0.13)--(2,0.13);
	\draw[color ={black},line width = 1.2] (3,-0.13)--(3,0.13);
	\draw[color ={black},line width = 1.2] (-1,-0.13)--(-1,0.13);
	\draw[color ={black},line width = 1.2] (-2,-0.13)--(-2,0.13);
	\draw[color ={black},line width = 1.2] (-3,-0.13)--(-3,0.13);
	\draw[color ={black},line width = 1.2] (-0.13,1)--(0.13,1);
	\draw[color ={black},line width = 1.2] (-0.13,2)--(0.13,2);
	\draw[color ={black},line width = 1.2] (-0.13,3)--(0.13,3);
	\draw[color ={black},line width = 1.2] (-0.13,4)--(0.13,4);
	\draw[color ={black},line width = 1.2] (-0.13,5)--(0.13,5);
	\draw[color ={black},line width = 1.2] (-0.13,6)--(0.13,6);
	\draw[color ={black},line width = 1.2] (-0.13,7)--(0.13,7);
	\draw[color ={black},line width = 1.2] (-0.13,8)--(0.13,8);
	\draw[color ={black},line width = 1.2] (-0.13,9)--(0.13,9);
	\draw[color ={black},line width = 1.2] (-0.13,10)--(0.13,10);
	\draw[color ={black},line width = 1.2] (-0.13,11)--(0.13,11);
	\draw[color ={black},line width = 1.2] (-0.13,12)--(0.13,12);
	\draw[color ={black},line width = 1.2] (-0.13,-1)--(0.13,-1);
	\draw[color ={black},line width = 1.2] (-0.13,-2)--(0.13,-2);
	\draw[color ={black},line width = 1.2] (-0.13,-3)--(0.13,-3);
	\draw (3,-0.4) node[anchor = center] {\scriptsize \color{black}{$3$}};
	\draw (-3,-0.4) node[anchor = center] {\scriptsize \color{black}{$-3$}};
	\draw (-0.5,3.1) node[anchor = center] {\footnotesize \color{black}{$3$}};
	\draw (-0.5,6.1) node[anchor = center] {\footnotesize \color{black}{$6$}};
	\draw (-0.5,-2.9) node[anchor = center] {\footnotesize \color{black}{$-3$}};
	\draw [color={black}] (-0.3,-0.3) node[anchor = center,scale=1, rotate = 0] {O};
	\draw (3,6) node[anchor = center] {\footnotesize \color{red}{$\Large\mathcal{C}_f\prime$}};
	
	\draw[color={red},line width = 2] (-2.8,-5.6)--(-2.6,-5.2)--(-2.4,-4.8)--(-2.2,-4.4)--(-2,-4)--(-1.8,-3.6)--(-1.6,-3.2)--(-1.4,-2.8)--(-1.2,-2.4)--(-1,-2)--(-0.8,-1.6)--(-0.6,-1.2)--(-0.4,-0.8)--(-0.2,-0.4)--(0,0)--(0.2,0.4)--(0.4,0.8)--(0.6,1.2)--(0.8,1.6)--(1,2)--(1.2,2.4)--(1.4,2.8)--(1.6,3.2)--(1.8,3.6)--(2,4)--(2.2,4.4)--(2.4,4.8)--(2.6,5.2)--(2.8,5.6)--(3,6);

\end{tikzpicture}\\
    \end{minipage}
    
\end{exercice}

\begin{Solution}
 L'équation réduite de la tangente au point d'abscisse $1$ est  : $y=f'(1)(x-1)+f(1)$.\\
        On lit graphiquement $f(1)=2$ et $f'(1)=2$.\\
        L'équation réduite de la tangente est donc donnée par :
        $y=2(x-1)+2$, soit $y=2x$.
\end{Solution}

% @see : https://coopmaths.fr/alea?uuid=a1ba2&id=can1F14&n=1&d=10&alea=rkIp223232323234234567891011121314151617181920212223242526272829303132333435363738394041424344454647482345678910111213141516171819202122232425262728293031323334353637383940414243&cd=1&cols=1
\needspace{10\baselineskip}
\begin{exercice}[BaremeTotal=false]


Soit $f$ la fonction définie sur $\mathbb{R}$ par :
$f(x)= -8+9x^2$.\\
En calculant la limite du taux d'accroissement, déterminer $f'(4)$.
\end{exercice}

\begin{Solution}
 $f$ est une fonction polynôme du second degré de la forme $f(x)=ax^2+b$.\\
    La fonction dérivée est donnée par la somme des dérivées des fonctions $u$ et $v$ définies par $u(x)=9x^2$ et $v(x)=-8$.\\
     Comme $u'(x)=18x$ et $v'(x)=0$, on obtient  $f'(x)=18x$. \\
     Ainsi, $f'(4)=18\times 4=72$.
\end{Solution}

% @see : https://coopmaths.fr/alea?uuid=3c690&id=can1F13&n=1&d=10&alea=EUDu223232323234234567891011121314151617181920212223242526272829303132333435363738394041424344454647482345678910111213141516171819202122232425262728293031323334353637383940414243&cd=1&cols=1
\needspace{10\baselineskip}
\begin{exercice}[BaremeTotal=false]


 Déterminer le coefficient directeur de la tangente à la courbe représentative de la fonction carré au point d'abscisse $-11$.    
\end{exercice}

\begin{Solution}
 Le coefficient directeur de la tangente au point d'abscisse $a$ est donné par le nombre dérivé $f'(a)$.\\
        La fonction $f$ définie par $f(x)=x^2$ a pour fonction dérivée la fonction $f'$ définie par $f'(x)=2x$.\\
        Comme $f'(-11)=2\times (-11)=-22$, le coefficient directeur de la tangente au point d'abscisse $-11$ est : $-22$. 
\end{Solution}

% @see : https://coopmaths.fr/alea?uuid=ebd89&id=1AN14-1&n=1&d=10&s=3&alea=ocYr223232323234234567891011121314151617181920212223242526272829303132333435363738394041424344454647482345678910111213141516171819202122232425262728293031323334353637383940414243&cd=0&cols=1
\needspace{10\baselineskip}
\begin{exercice}[BaremeTotal=false]


 Donner la dérivée de la fonction $f$, dérivable pour tout $x\in\R$, définie par  $f(x)=-\dfrac{3}{10}x+4$.
\end{exercice}

\begin{Solution}
 L'expression de la dérivée de la fonction $f$ définie par $f(x)=-\dfrac{3}{10}x+4$ est : ${\color[HTML]{f15929}\boldsymbol{f'(x)=-\dfrac{3}{10}}}$.
\end{Solution}

% @see : https://coopmaths.fr/alea?uuid=ebd89&id=1AN14-1&n=1&d=10&s=4&alea=uAlP223232323234234567891011121314151617181920212223242526272829303132333435363738394041424344454647482345678910111213141516171819202122232425262728293031323334353637383940414243&cd=0&cols=1
\needspace{10\baselineskip}
\begin{exercice}[BaremeTotal=false]


 Donner la dérivée de la fonction $f$, dérivable pour tout $x\in\R$, définie par  $f(x)=10x^{12}$.
\end{exercice}

\begin{Solution}
 L'expression de la dérivée de la fonction $f$ définie par $f(x)=10x^{12}$ est : ${\color[HTML]{f15929}\boldsymbol{f'(x)=120x^{11}}}$.
\end{Solution}

% @see : https://coopmaths.fr/alea?uuid=ebd89&id=1AN14-1&n=1&d=10&s=7&alea=vjN7223232323234234567891011121314151617181920212223242526272829303132333435363738394041424344454647482345678910111213141516171819202122232425262728293031323334353637383940414243&cd=0&cols=1
\needspace{10\baselineskip}
\begin{exercice}[BaremeTotal=false]


 Donner la dérivée de la fonction $f$, dérivable pour tout $x\in\R^*$, définie par  $f(x)=\dfrac{-2}{3x}$.
\end{exercice}

\begin{Solution}
 L'expression de la dérivée de la fonction $f$ définie par $f(x)=\dfrac{-2}{3x}$ est : ${\color[HTML]{f15929}\boldsymbol{f'(x)=\dfrac{2}{3x^2}}}$.
\end{Solution}

% @see : https://coopmaths.fr/alea?uuid=a83c0&id=1AN14-4&n=1&d=10&s=1&alea=Njr3223232323234234567891011121314151617181920212223242526272829303132333435363738394041424344454647482345678910111213141516171819202122232425262728293031323334353637383940414243&cd=0&cols=1
\needspace{10\baselineskip}
\begin{exercice}[BaremeTotal=false]


 Donner l'expression de la dérivée de la fonction $f$ définie sur $\R^{*}$ par $f(x)=-\dfrac{7}{x}-3x^{2}$.\\
\end{exercice}

\begin{Solution}
 L'expression de la dérivée de la fonction $f$ définie par $f(x)=-\dfrac{7}{x}-3x^{2}$ est :\\$f'(x)={\color[HTML]{f15929}\boldsymbol{\dfrac{7}{x^2}-6x}}$.\\
\end{Solution}

% @see : https://coopmaths.fr/alea?uuid=a83c0&id=1AN14-4&n=1&d=10&s=4&alea=4Vpz223232323234234567891011121314151617181920212223242526272829303132333435363738394041424344454647482345678910111213141516171819202122232425262728293031323334353637383940414243&cd=1&cols=1
\needspace{10\baselineskip}
\begin{exercice}[BaremeTotal=false]


 Donner l'expression de la dérivée de la fonction $f$ définie sur $\R_+^{*}$ par $f(x)=2x^{2}+5x+\dfrac{3}{x}+4\sqrt{x}$.\\
\end{exercice}

\begin{Solution}
 $(2x^{2}+5x)^\prime=4x+5$\\$(\dfrac{3}{x})^\prime=-\dfrac{3}{x^2}$\\$(4\sqrt{x})^\prime=\dfrac{4}{2\sqrt{x}}$\\L'expression de la dérivée de la fonction $f$ définie par $f(x)=2x^{2}+5x+\dfrac{3}{x}+4\sqrt{x}$ est :\\$f'(x)={\color[HTML]{f15929}\boldsymbol{4x+5-\dfrac{3}{x^2}+\dfrac{4}{2\sqrt{x}}}}$.\\
\end{Solution}

\end{Maquette}
\clearpage
\begin{Maquette}[DM]{Niveau= ,Classe=\getdata[44]\mydata\ (v44), Date = 06/01/25  ,Theme=Exercices}


% @see : https://coopmaths.fr/alea?uuid=4c8c7&id=1AN11-3&n=1&d=10&s=2&alea=nawO22323232323423456789101112131415161718192021222324252627282930313233343536373839404142434445464748234567891011121314151617181920212223242526272829303132333435363738394041424344&cd=1&cols=2
\needspace{10\baselineskip}
\begin{exercice}[BaremeTotal=false]
\label{eleve:44}


 \begin{minipage}[t]{\linewidth} Soit $f$ une fonction dérivable sur $[-5;5]$ et $\mathcal{C}_f$ sa courbe représentative.\\On sait que  $f(-3)=5~~$ et que $~~f'(-3)=3$.\\Déterminer une équation de la tangente $(T)$ à la courbe $\mathcal{C}_f$ au point d'abscisse $-3$,\\sans utiliser la formule de cours de l'équation de tangente. \end{minipage}

\end{exercice}

\begin{Solution}
 \begin{minipage}[t]{\linewidth} On sait que la tangente n'est pas une droite verticale, puisque la fonction est dérivable sur l'intervalle.\\On en déduit que la tangente $(T)$ au point d'abscisse $-3$, admet une équation réduite de la forme :  $(T) : y=m x + p$. 

\medskip
$\bullet$ Détermination de $m$ :\\On sait que le nombre dérivé en $-3$ est par définition, le coefficient directeur de la tangente au point d'abscisse $-3$.\\Par conséquent, on a déjà : $m=f'(-3)=3$.\\ On en déduit que  $(T) : y= 3 x + p$\\$\bullet$ Détermination de $p$ :\\ Pour cela, on utilise que si $f(-3)=5~~$, alors le point $A$ de coordonnées $(-3;5)$ appartient à $\mathcal{C}_f$ mais aussi à $(T)$.\\ On peut écrire $A(-3;5) = \mathcal{C}_f \cap (T)$.\\ On remplace alors les coordonnées de $A(-3;5)$ dans l'équation  $(T) : y= 3 x + p$.\\ $\begin{aligned} \phantom{\iff}&A(-3;5)\in (T)\\ \iff& 5= 3 \times (-3)  + p\\ \iff& p=5 -3 \times (-3)  \\ \iff& p=5 +9   \\ \iff& p=14\\\end{aligned}$\\On peut conclure que : $(T) : {\color[HTML]{f15929}\boldsymbol{y=3x+14}}$. \end{minipage}

\end{Solution}

% @see : https://coopmaths.fr/alea?uuid=0e984&id=can1F15&n=1&d=10&alea=1Alo22323232323423456789101112131415161718192021222324252627282930313233343536373839404142434445464748234567891011121314151617181920212223242526272829303132333435363738394041424344&cd=1&cols=1
\needspace{10\baselineskip}
\begin{exercice}[BaremeTotal=false]


 La courbe représente une fonction $f$ et la droite est la tangente au point d'abscisse $0$.\\
        Déterminer $f'(0)$.\begin{center}\begin{tikzpicture}[baseline,scale = 0.5]

    \tikzset{
      point/.style={
        thick,
        draw,
        cross out,
        inner sep=0pt,
        minimum width=5pt,
        minimum height=5pt,
      },
    }
    \clip (-10,-2) rectangle (4,12);
    	
	\draw[color ={black},line width = 2,>=latex,->] (-10,0)--(4,0);
	\draw[color ={black},line width = 2,>=latex,->] (0,-2)--(0,12);
	\draw[color ={black},opacity = 0.5] (-10,1)--(4,1);
	\draw[color ={black},opacity = 0.5] (-10,2)--(4,2);
	\draw[color ={black},opacity = 0.5] (-10,3)--(4,3);
	\draw[color ={black},opacity = 0.5] (-10,4)--(4,4);
	\draw[color ={black},opacity = 0.5] (-10,5)--(4,5);
	\draw[color ={black},opacity = 0.5] (-10,6)--(4,6);
	\draw[color ={black},opacity = 0.5] (-10,7)--(4,7);
	\draw[color ={black},opacity = 0.5] (-10,8)--(4,8);
	\draw[color ={black},opacity = 0.5] (-10,9)--(4,9);
	\draw[color ={black},opacity = 0.5] (-10,10)--(4,10);
	\draw[color ={black},opacity = 0.5] (-10,11)--(4,11);
	\draw[color ={black},opacity = 0.5] (-10,12)--(4,12);
	\draw[color ={black},opacity = 0.5] (-10,-1)--(4,-1);
	\draw[color ={black},opacity = 0.5] (-10,-2)--(4,-2);
	\draw[color ={black},opacity = 0.5] (1,-2)--(1,12);
	\draw[color ={black},opacity = 0.5] (2,-2)--(2,12);
	\draw[color ={black},opacity = 0.5] (3,-2)--(3,12);
	\draw[color ={black},opacity = 0.5] (4,-2)--(4,12);
	\draw[color ={black},opacity = 0.5] (-1,-2)--(-1,12);
	\draw[color ={black},opacity = 0.5] (-2,-2)--(-2,12);
	\draw[color ={black},opacity = 0.5] (-3,-2)--(-3,12);
	\draw[color ={black},opacity = 0.5] (-4,-2)--(-4,12);
	\draw[color ={black},opacity = 0.5] (-5,-2)--(-5,12);
	\draw[color ={black},opacity = 0.5] (-6,-2)--(-6,12);
	\draw[color ={black},opacity = 0.5] (-7,-2)--(-7,12);
	\draw[color ={black},opacity = 0.5] (-8,-2)--(-8,12);
	\draw[color ={black},opacity = 0.5] (-9,-2)--(-9,12);
	\draw[color ={black},opacity = 0.5] (-10,-2)--(-10,12);
	\draw[color ={gray},opacity = 0.3] (-10,0)--(4,0);
	\draw[color ={gray},opacity = 0.3] (-10,0)--(4,0);
	\draw[color ={gray},opacity = 0.3] (0,-2)--(0,12);
	\draw[color ={gray},opacity = 0.3] (0,-2)--(0,12);
	\draw[color ={gray},opacity = 0.3] (-5,-2)--(-5,12);
	\draw[color ={gray},opacity = 0.3] (-6,-2)--(-6,12);
	\draw[color ={gray},opacity = 0.3] (-7,-2)--(-7,12);
	\draw[color ={gray},opacity = 0.3] (-8,-2)--(-8,12);
	\draw[color ={gray},opacity = 0.3] (-9,-2)--(-9,12);
	\draw[color ={gray},opacity = 0.3] (-10,-2)--(-10,12);
	\draw[color ={black},line width = 2] (2,-0.2)--(2,0.2);
	\draw[color ={black},line width = 2] (-2,-0.2)--(-2,0.2);
	\draw[color ={black},line width = 2] (-4,-0.2)--(-4,0.2);
	\draw[color ={black},line width = 2] (-6,-0.2)--(-6,0.2);
	\draw[color ={black},line width = 2] (-8,-0.2)--(-8,0.2);
	\draw[color ={black},line width = 2] (-0.2,2)--(0.2,2);
	\draw[color ={black},line width = 2] (-0.2,4)--(0.2,4);
	\draw[color ={black},line width = 2] (-0.2,6)--(0.2,6);
	\draw[color ={black},line width = 2] (-0.2,8)--(0.2,8);
	\draw[color ={black},line width = 2] (-0.2,10)--(0.2,10);
	\draw (2,-0.4) node[anchor = center] {\scriptsize \color{black}{$1$}};
	\draw (-2,-0.4) node[anchor = center] {\scriptsize \color{black}{$-1$}};
	\draw (-4,-0.4) node[anchor = center] {\scriptsize \color{black}{$-2$}};
	\draw (-6,-0.4) node[anchor = center] {\scriptsize \color{black}{$-3$}};
	\draw (-8,-0.4) node[anchor = center] {\scriptsize \color{black}{$-4$}};
	\draw (-0.5,2.1) node[anchor = center] {\footnotesize \color{black}{$1$}};
	\draw (-0.5,4.1) node[anchor = center] {\footnotesize \color{black}{$2$}};
	\draw (-0.5,6.1) node[anchor = center] {\footnotesize \color{black}{$3$}};
	\draw (-0.5,8.1) node[anchor = center] {\footnotesize \color{black}{$4$}};
	\draw (-0.5,10.1) node[anchor = center] {\footnotesize \color{black}{$5$}};
	\draw [color={black}] (-0.3,-0.3) node[anchor = center,scale=1, rotate = 0] {O};
	
	\draw[color={blue},line width = 2] (-5.6,12.93)--(-5.2,11.32)--(-4.8,9.91)--(-4.4,8.67)--(-4,7.59)--(-3.6,6.64)--(-3.2,5.81)--(-2.8,5.09)--(-2.4,4.45)--(-2,3.9)--(-1.6,3.41)--(-1.2,2.98)--(-0.8,2.61)--(-0.4,2.29)--(0,2)--(0.4,1.75)--(0.8,1.53)--(1.2,1.34)--(1.6,1.17)--(2,1.03)--(2.4,0.9)--(2.8,0.79)--(3.2,0.69)--(3.6,0.6)--(4,0.53);
	
	\draw[color={red},line width = 2] (-10,8.67)--(-9.6,8.4)--(-9.2,8.13)--(-8.8,7.87)--(-8.4,7.6)--(-8,7.33)--(-7.6,7.07)--(-7.2,6.8)--(-6.8,6.53)--(-6.4,6.27)--(-6,6)--(-5.6,5.73)--(-5.2,5.47)--(-4.8,5.2)--(-4.4,4.93)--(-4,4.67)--(-3.6,4.4)--(-3.2,4.13)--(-2.8,3.87)--(-2.4,3.6)--(-2,3.33)--(-1.6,3.07)--(-1.2,2.8)--(-0.8,2.53)--(-0.4,2.27)--(0,2)--(0.4,1.73)--(0.8,1.47)--(1.2,1.2)--(1.6,0.93)--(2,0.67)--(2.4,0.4)--(2.8,0.13)--(3.2,-0.13)--(3.6,-0.4)--(4,-0.67);

\end{tikzpicture}\end{center}
\end{exercice}

\begin{Solution}
 $f'(0)$ est donné par le coefficient directeur de la tangente à la courbe au point d'abscisse $0$, soit $\dfrac{-2}{3}=-\dfrac{2}{3}$.
\end{Solution}

% @see : https://coopmaths.fr/alea?uuid=6f32d&id=can1F16&n=1&d=10&alea=dqSw22323232323423456789101112131415161718192021222324252627282930313233343536373839404142434445464748234567891011121314151617181920212223242526272829303132333435363738394041424344&cd=1&cols=1
\needspace{10\baselineskip}
\newpage
\begin{exercice}[BaremeTotal=false]


 On donne les représentations graphiques d'une fonction et de sa dérivée.\\
        Donner l'équation réduite de la tangente à la courbe de $f$ en $x=0$. \\ \begin{minipage}{0.5\linewidth}
    \begin{tikzpicture}[baseline, scale=0.8]

    \tikzset{
      point/.style={
        thick,
        draw,
        cross out,
        inner sep=0pt,
        minimum width=5pt,
        minimum height=5pt,
      },
    }
    \clip (-4.5,-4.5) rectangle (5.75,13.5);
    	
	\draw[color ={black},>=latex,->] (-3,0)--(3,0);
	\draw[color ={black},>=latex,->] (0,-3)--(0,12);
	\draw[color ={black},opacity = 0.4] (-3,1)--(3,1);
	\draw[color ={black},opacity = 0.4] (-3,2)--(3,2);
	\draw[color ={black},opacity = 0.4] (-3,3)--(3,3);
	\draw[color ={black},opacity = 0.4] (-3,4)--(3,4);
	\draw[color ={black},opacity = 0.4] (-3,5)--(3,5);
	\draw[color ={black},opacity = 0.4] (-3,6)--(3,6);
	\draw[color ={black},opacity = 0.4] (-3,7)--(3,7);
	\draw[color ={black},opacity = 0.4] (-3,8)--(3,8);
	\draw[color ={black},opacity = 0.4] (-3,9)--(3,9);
	\draw[color ={black},opacity = 0.4] (-3,10)--(3,10);
	\draw[color ={black},opacity = 0.4] (-3,11)--(3,11);
	\draw[color ={black},opacity = 0.4] (-3,12)--(3,12);
	\draw[color ={black},opacity = 0.4] (-3,-1)--(3,-1);
	\draw[color ={black},opacity = 0.4] (-3,-2)--(3,-2);
	\draw[color ={black},opacity = 0.4] (-3,-3)--(3,-3);
	\draw[color ={black},opacity = 0.4] (1,-3)--(1,12);
	\draw[color ={black},opacity = 0.4] (2,-3)--(2,12);
	\draw[color ={black},opacity = 0.4] (3,-3)--(3,12);
	\draw[color ={black},opacity = 0.4] (-1,-3)--(-1,12);
	\draw[color ={black},opacity = 0.4] (-2,-3)--(-2,12);
	\draw[color ={black},opacity = 0.4] (-3,-3)--(-3,12);
	\draw[color ={gray},opacity = 0.3] (-3,0)--(3,0);
	\draw[color ={gray},opacity = 0.3] (-3,0.5)--(3,0.5);
	\draw[color ={gray},opacity = 0.3] (-3,1.5)--(3,1.5);
	\draw[color ={gray},opacity = 0.3] (-3,2.5)--(3,2.5);
	\draw[color ={gray},opacity = 0.3] (-3,3.5)--(3,3.5);
	\draw[color ={gray},opacity = 0.3] (-3,4.5)--(3,4.5);
	\draw[color ={gray},opacity = 0.3] (-3,5.5)--(3,5.5);
	\draw[color ={gray},opacity = 0.3] (-3,6.5)--(3,6.5);
	\draw[color ={gray},opacity = 0.3] (-3,7.5)--(3,7.5);
	\draw[color ={gray},opacity = 0.3] (-3,8.5)--(3,8.5);
	\draw[color ={gray},opacity = 0.3] (-3,9.5)--(3,9.5);
	\draw[color ={gray},opacity = 0.3] (-3,10.5)--(3,10.5);
	\draw[color ={gray},opacity = 0.3] (-3,11.5)--(3,11.5);
	\draw[color ={gray},opacity = 0.3] (-3,0)--(3,0);
	\draw[color ={gray},opacity = 0.3] (-3,-0.5)--(3,-0.5);
	\draw[color ={gray},opacity = 0.3] (-3,-1.5)--(3,-1.5);
	\draw[color ={gray},opacity = 0.3] (-3,-2.5)--(3,-2.5);
	\draw[color ={gray},opacity = 0.3] (0,-3)--(0,12);
	\draw[color ={gray},opacity = 0.3] (0,-3)--(0,12);
	\draw[color ={black},line width = 1.2] (1,-0.13)--(1,0.13);
	\draw[color ={black},line width = 1.2] (2,-0.13)--(2,0.13);
	\draw[color ={black},line width = 1.2] (3,-0.13)--(3,0.13);
	\draw[color ={black},line width = 1.2] (-1,-0.13)--(-1,0.13);
	\draw[color ={black},line width = 1.2] (-2,-0.13)--(-2,0.13);
	\draw[color ={black},line width = 1.2] (-3,-0.13)--(-3,0.13);
	\draw[color ={black},line width = 1.2] (-0.13,1)--(0.13,1);
	\draw[color ={black},line width = 1.2] (-0.13,2)--(0.13,2);
	\draw[color ={black},line width = 1.2] (-0.13,3)--(0.13,3);
	\draw[color ={black},line width = 1.2] (-0.13,4)--(0.13,4);
	\draw[color ={black},line width = 1.2] (-0.13,5)--(0.13,5);
	\draw[color ={black},line width = 1.2] (-0.13,6)--(0.13,6);
	\draw[color ={black},line width = 1.2] (-0.13,7)--(0.13,7);
	\draw[color ={black},line width = 1.2] (-0.13,8)--(0.13,8);
	\draw[color ={black},line width = 1.2] (-0.13,9)--(0.13,9);
	\draw[color ={black},line width = 1.2] (-0.13,10)--(0.13,10);
	\draw[color ={black},line width = 1.2] (-0.13,11)--(0.13,11);
	\draw[color ={black},line width = 1.2] (-0.13,12)--(0.13,12);
	\draw[color ={black},line width = 1.2] (-0.13,-1)--(0.13,-1);
	\draw[color ={black},line width = 1.2] (-0.13,-2)--(0.13,-2);
	\draw[color ={black},line width = 1.2] (-0.13,-3)--(0.13,-3);
	\draw (3,-0.4) node[anchor = center] {\scriptsize \color{black}{$3$}};
	\draw (-3,-0.4) node[anchor = center] {\scriptsize \color{black}{$-3$}};
	\draw (-0.5,3.1) node[anchor = center] {\footnotesize \color{black}{$3$}};
	\draw (-0.5,6.1) node[anchor = center] {\footnotesize \color{black}{$6$}};
	\draw (-0.5,9.1) node[anchor = center] {\footnotesize \color{black}{$9$}};
	\draw (-0.5,-2.9) node[anchor = center] {\footnotesize \color{black}{$-3$}};
	\draw [color={black}] (-0.3,-0.3) node[anchor = center,scale=1, rotate = 0] {O};
	\draw (3,10) node[anchor = center] {\footnotesize \color{blue}{$\Large \mathcal{C}_f$}};
	
	\draw[color={blue},line width = 2] (-2.2,12.24)--(-2,11)--(-1.8,9.84)--(-1.6,8.76)--(-1.4,7.76)--(-1.2,6.84)--(-1,6)--(-0.8,5.24)--(-0.6,4.56)--(-0.4,3.96)--(-0.2,3.44)--(0,3)--(0.2,2.64)--(0.4,2.36)--(0.6,2.16)--(0.8,2.04)--(1,2)--(1.2,2.04)--(1.4,2.16)--(1.6,2.36)--(1.8,2.64)--(2,3)--(2.2,3.44)--(2.4,3.96)--(2.6,4.56)--(2.8,5.24)--(3,6);

\end{tikzpicture}\\
    \end{minipage}
    \begin{minipage}{0.5\linewidth}
    \begin{tikzpicture}[baseline, scale=0.8]

    \tikzset{
      point/.style={
        thick,
        draw,
        cross out,
        inner sep=0pt,
        minimum width=5pt,
        minimum height=5pt,
      },
    }
    \clip (-4.5,-6.5) rectangle (6.5,9.5);
    	
	\draw[color ={black},>=latex,->] (-3,0)--(3,0);
	\draw[color ={black},>=latex,->] (0,-5)--(0,8);
	\draw[color ={black},opacity = 0.4] (-3,1)--(3,1);
	\draw[color ={black},opacity = 0.4] (-3,2)--(3,2);
	\draw[color ={black},opacity = 0.4] (-3,3)--(3,3);
	\draw[color ={black},opacity = 0.4] (-3,4)--(3,4);
	\draw[color ={black},opacity = 0.4] (-3,5)--(3,5);
	\draw[color ={black},opacity = 0.4] (-3,6)--(3,6);
	\draw[color ={black},opacity = 0.4] (-3,7)--(3,7);
	\draw[color ={black},opacity = 0.4] (-3,8)--(3,8);
	\draw[color ={black},opacity = 0.4] (-3,-1)--(3,-1);
	\draw[color ={black},opacity = 0.4] (-3,-2)--(3,-2);
	\draw[color ={black},opacity = 0.4] (-3,-3)--(3,-3);
	\draw[color ={black},opacity = 0.4] (-3,-4)--(3,-4);
	\draw[color ={black},opacity = 0.4] (-3,-5)--(3,-5);
	\draw[color ={black},opacity = 0.4] (1,-5)--(1,8);
	\draw[color ={black},opacity = 0.4] (2,-5)--(2,8);
	\draw[color ={black},opacity = 0.4] (3,-5)--(3,8);
	\draw[color ={black},opacity = 0.4] (-1,-5)--(-1,8);
	\draw[color ={black},opacity = 0.4] (-2,-5)--(-2,8);
	\draw[color ={black},opacity = 0.4] (-3,-5)--(-3,8);
	\draw[color ={gray},opacity = 0.3] (-3,0)--(3,0);
	\draw[color ={gray},opacity = 0.3] (-3,0.5)--(3,0.5);
	\draw[color ={gray},opacity = 0.3] (-3,1.5)--(3,1.5);
	\draw[color ={gray},opacity = 0.3] (-3,2.5)--(3,2.5);
	\draw[color ={gray},opacity = 0.3] (-3,3.5)--(3,3.5);
	\draw[color ={gray},opacity = 0.3] (-3,4.5)--(3,4.5);
	\draw[color ={gray},opacity = 0.3] (-3,5.5)--(3,5.5);
	\draw[color ={gray},opacity = 0.3] (-3,6.5)--(3,6.5);
	\draw[color ={gray},opacity = 0.3] (-3,7.5)--(3,7.5);
	\draw[color ={gray},opacity = 0.3] (-3,0)--(3,0);
	\draw[color ={gray},opacity = 0.3] (-3,-0.5)--(3,-0.5);
	\draw[color ={gray},opacity = 0.3] (-3,-1.5)--(3,-1.5);
	\draw[color ={gray},opacity = 0.3] (-3,-2.5)--(3,-2.5);
	\draw[color ={gray},opacity = 0.3] (-3,-3.5)--(3,-3.5);
	\draw[color ={gray},opacity = 0.3] (-3,-4.5)--(3,-4.5);
	\draw[color ={gray},opacity = 0.3] (0,-5)--(0,8);
	\draw[color ={gray},opacity = 0.3] (0,-5)--(0,8);
	\draw[color ={black},line width = 1.2] (1,-0.13)--(1,0.13);
	\draw[color ={black},line width = 1.2] (2,-0.13)--(2,0.13);
	\draw[color ={black},line width = 1.2] (3,-0.13)--(3,0.13);
	\draw[color ={black},line width = 1.2] (-1,-0.13)--(-1,0.13);
	\draw[color ={black},line width = 1.2] (-2,-0.13)--(-2,0.13);
	\draw[color ={black},line width = 1.2] (-3,-0.13)--(-3,0.13);
	\draw[color ={black},line width = 1.2] (-0.13,1)--(0.13,1);
	\draw[color ={black},line width = 1.2] (-0.13,2)--(0.13,2);
	\draw[color ={black},line width = 1.2] (-0.13,3)--(0.13,3);
	\draw[color ={black},line width = 1.2] (-0.13,4)--(0.13,4);
	\draw[color ={black},line width = 1.2] (-0.13,5)--(0.13,5);
	\draw[color ={black},line width = 1.2] (-0.13,6)--(0.13,6);
	\draw[color ={black},line width = 1.2] (-0.13,7)--(0.13,7);
	\draw[color ={black},line width = 1.2] (-0.13,8)--(0.13,8);
	\draw[color ={black},line width = 1.2] (-0.13,9)--(0.13,9);
	\draw[color ={black},line width = 1.2] (-0.13,10)--(0.13,10);
	\draw[color ={black},line width = 1.2] (-0.13,11)--(0.13,11);
	\draw[color ={black},line width = 1.2] (-0.13,12)--(0.13,12);
	\draw[color ={black},line width = 1.2] (-0.13,-1)--(0.13,-1);
	\draw[color ={black},line width = 1.2] (-0.13,-2)--(0.13,-2);
	\draw[color ={black},line width = 1.2] (-0.13,-3)--(0.13,-3);
	\draw (3,-0.4) node[anchor = center] {\scriptsize \color{black}{$3$}};
	\draw (-3,-0.4) node[anchor = center] {\scriptsize \color{black}{$-3$}};
	\draw (-0.5,3.1) node[anchor = center] {\footnotesize \color{black}{$3$}};
	\draw (-0.5,6.1) node[anchor = center] {\footnotesize \color{black}{$6$}};
	\draw (-0.5,-2.9) node[anchor = center] {\footnotesize \color{black}{$-3$}};
	\draw [color={black}] (-0.3,-0.3) node[anchor = center,scale=1, rotate = 0] {O};
	\draw (3,6) node[anchor = center] {\footnotesize \color{red}{$\Large\mathcal{C}_f\prime$}};
	
	\draw[color={red},line width = 2] (-2,-6)--(-1.8,-5.6)--(-1.6,-5.2)--(-1.4,-4.8)--(-1.2,-4.4)--(-1,-4)--(-0.8,-3.6)--(-0.6,-3.2)--(-0.4,-2.8)--(-0.2,-2.4)--(0,-2)--(0.2,-1.6)--(0.4,-1.2)--(0.6,-0.8)--(0.8,-0.4)--(1,0)--(1.2,0.4)--(1.4,0.8)--(1.6,1.2)--(1.8,1.6)--(2,2)--(2.2,2.4)--(2.4,2.8)--(2.6,3.2)--(2.8,3.6)--(3,4);

\end{tikzpicture}\\
    \end{minipage}
    
\end{exercice}

\begin{Solution}
 L'équation réduite de la tangente au point d'abscisse $0$ est  : $y=f'(0)(x-0)+f(0)$.\\
        On lit graphiquement $f(0)=3$ et $f'(0)=-2$.\\
        L'équation réduite de la tangente est donc donnée par :
        $y=-2(x+0)+3$, soit $y=-2x+3$.
\end{Solution}

% @see : https://coopmaths.fr/alea?uuid=a1ba2&id=can1F14&n=1&d=10&alea=rkIp22323232323423456789101112131415161718192021222324252627282930313233343536373839404142434445464748234567891011121314151617181920212223242526272829303132333435363738394041424344&cd=1&cols=1
\needspace{10\baselineskip}
\begin{exercice}[BaremeTotal=false]


Soit $f$ la fonction définie sur $\mathbb{R}$ par : 
$f(x)= -3x^2-9$.\\
En calculant la limite du taux d'accroissement, déterminer $f'(-4)$.
\end{exercice}

\begin{Solution}
 $f$ est une fonction polynôme du second degré de la forme $f(x)=ax^2+b$.\\
    La fonction dérivée est donnée par la somme des dérivées des fonctions $u$ et $v$ définies par $u(x)=-3x^2$ et $v(x)=-9$.\\
     Comme $u'(x)=-6x$ et $v'(x)=0$, on obtient  $f'(x)=-6x$. \\
     Ainsi, $f'(-4)=-6\times (-4)=24$.
\end{Solution}

% @see : https://coopmaths.fr/alea?uuid=3c690&id=can1F13&n=1&d=10&alea=EUDu22323232323423456789101112131415161718192021222324252627282930313233343536373839404142434445464748234567891011121314151617181920212223242526272829303132333435363738394041424344&cd=1&cols=1
\needspace{10\baselineskip}
\begin{exercice}[BaremeTotal=false]


 Déterminer le coefficient directeur de la tangente à la courbe représentative de la fonction carré au point d'abscisse $-15$.    
\end{exercice}

\begin{Solution}
 Le coefficient directeur de la tangente au point d'abscisse $a$ est donné par le nombre dérivé $f'(a)$.\\
        La fonction $f$ définie par $f(x)=x^2$ a pour fonction dérivée la fonction $f'$ définie par $f'(x)=2x$.\\
        Comme $f'(-15)=2\times (-15)=-30$, le coefficient directeur de la tangente au point d'abscisse $-15$ est : $-30$. 
\end{Solution}

% @see : https://coopmaths.fr/alea?uuid=ebd89&id=1AN14-1&n=1&d=10&s=3&alea=ocYr22323232323423456789101112131415161718192021222324252627282930313233343536373839404142434445464748234567891011121314151617181920212223242526272829303132333435363738394041424344&cd=0&cols=1
\needspace{10\baselineskip}
\begin{exercice}[BaremeTotal=false]


 Donner la dérivée de la fonction $f$, dérivable pour tout $x\in\R$, définie par  $f(x)=\dfrac{5x+1}{8}$.
\end{exercice}

\begin{Solution}
 L'expression de la dérivée de la fonction $f$ définie par $f(x)=\dfrac{5x+1}{8}$ est : ${\color[HTML]{f15929}\boldsymbol{f'(x)=\dfrac{5}{8}}}$.
\end{Solution}

% @see : https://coopmaths.fr/alea?uuid=ebd89&id=1AN14-1&n=1&d=10&s=4&alea=uAlP22323232323423456789101112131415161718192021222324252627282930313233343536373839404142434445464748234567891011121314151617181920212223242526272829303132333435363738394041424344&cd=0&cols=1
\needspace{10\baselineskip}
\begin{exercice}[BaremeTotal=false]


 Donner la dérivée de la fonction $f$, dérivable pour tout $x\in\R$, définie par  $f(x)=5x^{12}$.
\end{exercice}

\begin{Solution}
 L'expression de la dérivée de la fonction $f$ définie par $f(x)=5x^{12}$ est : ${\color[HTML]{f15929}\boldsymbol{f'(x)=60x^{11}}}$.
\end{Solution}

% @see : https://coopmaths.fr/alea?uuid=ebd89&id=1AN14-1&n=1&d=10&s=7&alea=vjN722323232323423456789101112131415161718192021222324252627282930313233343536373839404142434445464748234567891011121314151617181920212223242526272829303132333435363738394041424344&cd=0&cols=1
\needspace{10\baselineskip}
\begin{exercice}[BaremeTotal=false]


 Donner la dérivée de la fonction $f$, dérivable pour tout $x\in\R^*$, définie par  $f(x)=\dfrac{-3}{7x}$.
\end{exercice}

\begin{Solution}
 L'expression de la dérivée de la fonction $f$ définie par $f(x)=\dfrac{-3}{7x}$ est : ${\color[HTML]{f15929}\boldsymbol{f'(x)=\dfrac{3}{7x^2}}}$.
\end{Solution}

% @see : https://coopmaths.fr/alea?uuid=a83c0&id=1AN14-4&n=1&d=10&s=1&alea=Njr322323232323423456789101112131415161718192021222324252627282930313233343536373839404142434445464748234567891011121314151617181920212223242526272829303132333435363738394041424344&cd=0&cols=1
\needspace{10\baselineskip}
\begin{exercice}[BaremeTotal=false]


 Donner l'expression de la dérivée de la fonction $f$ définie sur $\R^{*}$ par $f(x)=-\dfrac{5}{x}+5x$.\\
\end{exercice}

\begin{Solution}
 L'expression de la dérivée de la fonction $f$ définie par $f(x)=-\dfrac{5}{x}+5x$ est :\\$f'(x)={\color[HTML]{f15929}\boldsymbol{\dfrac{5}{x^2}+5}}$.\\
\end{Solution}

% @see : https://coopmaths.fr/alea?uuid=a83c0&id=1AN14-4&n=1&d=10&s=4&alea=4Vpz22323232323423456789101112131415161718192021222324252627282930313233343536373839404142434445464748234567891011121314151617181920212223242526272829303132333435363738394041424344&cd=1&cols=1
\needspace{10\baselineskip}
\begin{exercice}[BaremeTotal=false]


 Donner l'expression de la dérivée de la fonction $f$ définie sur $\R_+^{*}$ par $f(x)=-3x^{2}-4x-\dfrac{1}{x}-3\sqrt{x}$.\\
\end{exercice}

\begin{Solution}
 $(-3x^{2}-4x)^\prime=-6x-4$\\$(-\dfrac{1}{x})^\prime=\dfrac{1}{x^2}$\\$(-3\sqrt{x})^\prime=-\dfrac{3}{2\sqrt{x}}$\\L'expression de la dérivée de la fonction $f$ définie par $f(x)=-3x^{2}-4x-\dfrac{1}{x}-3\sqrt{x}$ est :\\$f'(x)={\color[HTML]{f15929}\boldsymbol{-6x-4+\dfrac{1}{x^2}-\dfrac{3}{2\sqrt{x}}}}$.\\
\end{Solution}

\end{Maquette}
\clearpage
\begin{Maquette}[DM]{Niveau= ,Classe=\getdata[45]\mydata\ (v45), Date = 06/01/25  ,Theme=Exercices}


% @see : https://coopmaths.fr/alea?uuid=4c8c7&id=1AN11-3&n=1&d=10&s=2&alea=nawO2232323232342345678910111213141516171819202122232425262728293031323334353637383940414243444546474823456789101112131415161718192021222324252627282930313233343536373839404142434445&cd=1&cols=2
\needspace{10\baselineskip}
\begin{exercice}[BaremeTotal=false]
\label{eleve:45}


 \begin{minipage}[t]{\linewidth} Soit $f$ une fonction dérivable sur $[-5;5]$ et $\mathcal{C}_f$ sa courbe représentative.\\On sait que  $f(1)=-4~~$ et que $~~f'(1)=1$.\\Déterminer une équation de la tangente $(T)$ à la courbe $\mathcal{C}_f$ au point d'abscisse $1$,\\sans utiliser la formule de cours de l'équation de tangente. \end{minipage}

\end{exercice}

\begin{Solution}
 \begin{minipage}[t]{\linewidth} On sait que la tangente n'est pas une droite verticale, puisque la fonction est dérivable sur l'intervalle.\\On en déduit que la tangente $(T)$ au point d'abscisse $1$, admet une équation réduite de la forme :  $(T) : y=m x + p$. 

\medskip
$\bullet$ Détermination de $m$ :\\On sait que le nombre dérivé en $1$ est par définition, le coefficient directeur de la tangente au point d'abscisse $1$.\\Par conséquent, on a déjà : $m=f'(1)=1$.\\ On en déduit que  $(T) : y= 1 x + p$\\$\bullet$ Détermination de $p$ :\\ Pour cela, on utilise que si $f(1)=-4~~$, alors le point $A$ de coordonnées $(1;-4)$ appartient à $\mathcal{C}_f$ mais aussi à $(T)$.\\ On peut écrire $A(1;-4) = \mathcal{C}_f \cap (T)$.\\ On remplace alors les coordonnées de $A(1;-4)$ dans l'équation  $(T) : y= 1 x + p$.\\ $\begin{aligned} \phantom{\iff}&A(1;-4)\in (T)\\ \iff& -4= 1 \times 1  + p\\ \iff& p=-4 - \times 1  \\ \iff& p=-4 -   \\ \iff& p=-5\\\end{aligned}$\\On peut conclure que : $(T) : {\color[HTML]{f15929}\boldsymbol{y=x-5}}$. \end{minipage}

\end{Solution}

% @see : https://coopmaths.fr/alea?uuid=0e984&id=can1F15&n=1&d=10&alea=1Alo2232323232342345678910111213141516171819202122232425262728293031323334353637383940414243444546474823456789101112131415161718192021222324252627282930313233343536373839404142434445&cd=1&cols=1
\needspace{10\baselineskip}
\begin{exercice}[BaremeTotal=false]


 La courbe représente une fonction $f$ et la droite est la tangente au point d'abscisse $0$.\\
        Déterminer $f'(0)$.\begin{center}\begin{tikzpicture}[baseline,scale = 0.5]

    \tikzset{
      point/.style={
        thick,
        draw,
        cross out,
        inner sep=0pt,
        minimum width=5pt,
        minimum height=5pt,
      },
    }
    \clip (-8,-10) rectangle (8,3);
    	
	\draw[color ={black},line width = 2,>=latex,->] (-8,0)--(8,0);
	\draw[color ={black},line width = 2,>=latex,->] (0,-10)--(0,3);
	\draw[color ={black},opacity = 0.5] (-8,1)--(8,1);
	\draw[color ={black},opacity = 0.5] (-8,2)--(8,2);
	\draw[color ={black},opacity = 0.5] (-8,3)--(8,3);
	\draw[color ={black},opacity = 0.5] (-8,-1)--(8,-1);
	\draw[color ={black},opacity = 0.5] (-8,-2)--(8,-2);
	\draw[color ={black},opacity = 0.5] (-8,-3)--(8,-3);
	\draw[color ={black},opacity = 0.5] (-8,-4)--(8,-4);
	\draw[color ={black},opacity = 0.5] (-8,-5)--(8,-5);
	\draw[color ={black},opacity = 0.5] (-8,-6)--(8,-6);
	\draw[color ={black},opacity = 0.5] (-8,-7)--(8,-7);
	\draw[color ={black},opacity = 0.5] (-8,-8)--(8,-8);
	\draw[color ={black},opacity = 0.5] (-8,-9)--(8,-9);
	\draw[color ={black},opacity = 0.5] (-8,-10)--(8,-10);
	\draw[color ={black},opacity = 0.5] (1,-10)--(1,3);
	\draw[color ={black},opacity = 0.5] (2,-10)--(2,3);
	\draw[color ={black},opacity = 0.5] (3,-10)--(3,3);
	\draw[color ={black},opacity = 0.5] (4,-10)--(4,3);
	\draw[color ={black},opacity = 0.5] (5,-10)--(5,3);
	\draw[color ={black},opacity = 0.5] (6,-10)--(6,3);
	\draw[color ={black},opacity = 0.5] (7,-10)--(7,3);
	\draw[color ={black},opacity = 0.5] (8,-10)--(8,3);
	\draw[color ={black},opacity = 0.5] (-1,-10)--(-1,3);
	\draw[color ={black},opacity = 0.5] (-2,-10)--(-2,3);
	\draw[color ={black},opacity = 0.5] (-3,-10)--(-3,3);
	\draw[color ={black},opacity = 0.5] (-4,-10)--(-4,3);
	\draw[color ={black},opacity = 0.5] (-5,-10)--(-5,3);
	\draw[color ={black},opacity = 0.5] (-6,-10)--(-6,3);
	\draw[color ={black},opacity = 0.5] (-7,-10)--(-7,3);
	\draw[color ={black},opacity = 0.5] (-8,-10)--(-8,3);
	\draw[color ={gray},opacity = 0.3] (-8,0)--(8,0);
	\draw[color ={gray},opacity = 0.3] (-8,0)--(8,0);
	\draw[color ={gray},opacity = 0.3] (-8,-4)--(8,-4);
	\draw[color ={gray},opacity = 0.3] (-8,-5)--(8,-5);
	\draw[color ={gray},opacity = 0.3] (-8,-6)--(8,-6);
	\draw[color ={gray},opacity = 0.3] (-8,-7)--(8,-7);
	\draw[color ={gray},opacity = 0.3] (-8,-8)--(8,-8);
	\draw[color ={gray},opacity = 0.3] (-8,-9)--(8,-9);
	\draw[color ={gray},opacity = 0.3] (-8,-10)--(8,-10);
	\draw[color ={gray},opacity = 0.3] (0,-10)--(0,3);
	\draw[color ={gray},opacity = 0.3] (0,-10)--(0,3);
	\draw[color ={black},line width = 2] (2,-0.2)--(2,0.2);
	\draw[color ={black},line width = 2] (4,-0.2)--(4,0.2);
	\draw[color ={black},line width = 2] (6,-0.2)--(6,0.2);
	\draw[color ={black},line width = 2] (-2,-0.2)--(-2,0.2);
	\draw[color ={black},line width = 2] (-4,-0.2)--(-4,0.2);
	\draw[color ={black},line width = 2] (-6,-0.2)--(-6,0.2);
	\draw[color ={black},line width = 2] (-0.2,1)--(0.2,1);
	\draw[color ={black},line width = 2] (-0.2,2)--(0.2,2);
	\draw[color ={black},line width = 2] (-0.2,-1)--(0.2,-1);
	\draw[color ={black},line width = 2] (-0.2,-2)--(0.2,-2);
	\draw[color ={black},line width = 2] (-0.2,-3)--(0.2,-3);
	\draw[color ={black},line width = 2] (-0.2,-4)--(0.2,-4);
	\draw[color ={black},line width = 2] (-0.2,-5)--(0.2,-5);
	\draw[color ={black},line width = 2] (-0.2,-6)--(0.2,-6);
	\draw[color ={black},line width = 2] (-0.2,-7)--(0.2,-7);
	\draw[color ={black},line width = 2] (-0.2,-8)--(0.2,-8);
	\draw[color ={black},line width = 2] (-0.2,-9)--(0.2,-9);
	\draw (2,-0.4) node[anchor = center] {\scriptsize \color{black}{$1$}};
	\draw (4,-0.4) node[anchor = center] {\scriptsize \color{black}{$2$}};
	\draw (6,-0.4) node[anchor = center] {\scriptsize \color{black}{$3$}};
	\draw (-2,-0.4) node[anchor = center] {\scriptsize \color{black}{$-1$}};
	\draw (-4,-0.4) node[anchor = center] {\scriptsize \color{black}{$-2$}};
	\draw (-6,-0.4) node[anchor = center] {\scriptsize \color{black}{$-3$}};
	\draw (-0.8,2.1) node[anchor = center] {\footnotesize \color{black}{$2$}};
	\draw (-0.8,-1.9) node[anchor = center] {\footnotesize \color{black}{$-2$}};
	\draw (-0.8,-3.9) node[anchor = center] {\footnotesize \color{black}{$-4$}};
	\draw (-0.8,-5.9) node[anchor = center] {\footnotesize \color{black}{$-6$}};
	\draw (-0.8,-7.9) node[anchor = center] {\footnotesize \color{black}{$-8$}};
	\draw [color={black}] (-0.3,-0.3) node[anchor = center,scale=1, rotate = 0] {O};
	
	\draw[color={blue},line width = 2] (-6,-11)--(-5.6,-9.84)--(-5.2,-8.76)--(-4.8,-7.76)--(-4.4,-6.84)--(-4,-6)--(-3.6,-5.24)--(-3.2,-4.56)--(-2.8,-3.96)--(-2.4,-3.44)--(-2,-3)--(-1.6,-2.64)--(-1.2,-2.36)--(-0.8,-2.16)--(-0.4,-2.04)--(0,-2)--(0.4,-2.04)--(0.8,-2.16)--(1.2,-2.36)--(1.6,-2.64)--(2,-3)--(2.4,-3.44)--(2.8,-3.96)--(3.2,-4.56)--(3.6,-5.24)--(4,-6)--(4.4,-6.84)--(4.8,-7.76)--(5.2,-8.76)--(5.6,-9.84);
	
	\draw[color={red},line width = 2] (-8,-2)--(-7.6,-2)--(-7.2,-2)--(-6.8,-2)--(-6.4,-2)--(-6,-2)--(-5.6,-2)--(-5.2,-2)--(-4.8,-2)--(-4.4,-2)--(-4,-2)--(-3.6,-2)--(-3.2,-2)--(-2.8,-2)--(-2.4,-2)--(-2,-2)--(-1.6,-2)--(-1.2,-2)--(-0.8,-2)--(-0.4,-2)--(0,-2)--(0.4,-2)--(0.8,-2)--(1.2,-2)--(1.6,-2)--(2,-2)--(2.4,-2)--(2.8,-2)--(3.2,-2)--(3.6,-2)--(4,-2)--(4.4,-2)--(4.8,-2)--(5.2,-2)--(5.6,-2)--(6,-2)--(6.4,-2)--(6.8,-2)--(7.2,-2)--(7.6,-2)--(8,-2);

\end{tikzpicture}\end{center}
\end{exercice}

\begin{Solution}
 $f'(0)$ est donné par le coefficient directeur de la tangente à la courbe au point d'abscisse $0$, soit $0$.
\end{Solution}

% @see : https://coopmaths.fr/alea?uuid=6f32d&id=can1F16&n=1&d=10&alea=dqSw2232323232342345678910111213141516171819202122232425262728293031323334353637383940414243444546474823456789101112131415161718192021222324252627282930313233343536373839404142434445&cd=1&cols=1
\needspace{10\baselineskip}
\newpage
\begin{exercice}[BaremeTotal=false]


 On donne les représentations graphiques d'une fonction et de sa dérivée.\\
      Donner l'équation réduite de la tangente à la courbe de $f$ en $x=-1$. \\ \begin{minipage}{0.5\linewidth}
    \begin{tikzpicture}[baseline, scale=0.8]

    \tikzset{
      point/.style={
        thick,
        draw,
        cross out,
        inner sep=0pt,
        minimum width=5pt,
        minimum height=5pt,
      },
    }
    \clip (-5.5,-9.5) rectangle (5.75,7.5);
    	
	\draw[color ={black},>=latex,->] (-4,0)--(4,0);
	\draw[color ={black},>=latex,->] (0,-8)--(0,6);
	\draw[color ={black},opacity = 0.4] (-4,1)--(4,1);
	\draw[color ={black},opacity = 0.4] (-4,2)--(4,2);
	\draw[color ={black},opacity = 0.4] (-4,3)--(4,3);
	\draw[color ={black},opacity = 0.4] (-4,4)--(4,4);
	\draw[color ={black},opacity = 0.4] (-4,5)--(4,5);
	\draw[color ={black},opacity = 0.4] (-4,6)--(4,6);
	\draw[color ={black},opacity = 0.4] (-4,-1)--(4,-1);
	\draw[color ={black},opacity = 0.4] (-4,-2)--(4,-2);
	\draw[color ={black},opacity = 0.4] (-4,-3)--(4,-3);
	\draw[color ={black},opacity = 0.4] (-4,-4)--(4,-4);
	\draw[color ={black},opacity = 0.4] (-4,-5)--(4,-5);
	\draw[color ={black},opacity = 0.4] (-4,-6)--(4,-6);
	\draw[color ={black},opacity = 0.4] (-4,-7)--(4,-7);
	\draw[color ={black},opacity = 0.4] (-4,-8)--(4,-8);
	\draw[color ={black},opacity = 0.4] (1,-8)--(1,6);
	\draw[color ={black},opacity = 0.4] (2,-8)--(2,6);
	\draw[color ={black},opacity = 0.4] (3,-8)--(3,6);
	\draw[color ={black},opacity = 0.4] (4,-8)--(4,6);
	\draw[color ={black},opacity = 0.4] (-1,-8)--(-1,6);
	\draw[color ={black},opacity = 0.4] (-2,-8)--(-2,6);
	\draw[color ={black},opacity = 0.4] (-3,-8)--(-3,6);
	\draw[color ={black},opacity = 0.4] (-4,-8)--(-4,6);
	\draw[color ={gray},opacity = 0.3] (-4,0)--(4,0);
	\draw[color ={gray},opacity = 0.3] (-4,0.5)--(4,0.5);
	\draw[color ={gray},opacity = 0.3] (-4,1.5)--(4,1.5);
	\draw[color ={gray},opacity = 0.3] (-4,2.5)--(4,2.5);
	\draw[color ={gray},opacity = 0.3] (-4,3.5)--(4,3.5);
	\draw[color ={gray},opacity = 0.3] (-4,4.5)--(4,4.5);
	\draw[color ={gray},opacity = 0.3] (-4,5.5)--(4,5.5);
	\draw[color ={gray},opacity = 0.3] (-4,0)--(4,0);
	\draw[color ={gray},opacity = 0.3] (-4,-0.5)--(4,-0.5);
	\draw[color ={gray},opacity = 0.3] (-4,-1.5)--(4,-1.5);
	\draw[color ={gray},opacity = 0.3] (-4,-2.5)--(4,-2.5);
	\draw[color ={gray},opacity = 0.3] (-4,-3.5)--(4,-3.5);
	\draw[color ={gray},opacity = 0.3] (-4,-4.5)--(4,-4.5);
	\draw[color ={gray},opacity = 0.3] (-4,-5.5)--(4,-5.5);
	\draw[color ={gray},opacity = 0.3] (-4,-6.5)--(4,-6.5);
	\draw[color ={gray},opacity = 0.3] (-4,-7)--(4,-7);
	\draw[color ={gray},opacity = 0.3] (-4,-7.5)--(4,-7.5);
	\draw[color ={gray},opacity = 0.3] (-4,-8)--(4,-8);
	\draw[color ={gray},opacity = 0.3] (0,-8)--(0,6);
	\draw[color ={gray},opacity = 0.3] (0,-8)--(0,6);
	\draw[color ={black},line width = 1.2] (1,-0.13)--(1,0.13);
	\draw[color ={black},line width = 1.2] (2,-0.13)--(2,0.13);
	\draw[color ={black},line width = 1.2] (3,-0.13)--(3,0.13);
	\draw[color ={black},line width = 1.2] (4,-0.13)--(4,0.13);
	\draw[color ={black},line width = 1.2] (-1,-0.13)--(-1,0.13);
	\draw[color ={black},line width = 1.2] (-2,-0.13)--(-2,0.13);
	\draw[color ={black},line width = 1.2] (-3,-0.13)--(-3,0.13);
	\draw[color ={black},line width = 1.2] (-4,-0.13)--(-4,0.13);
	\draw[color ={black},line width = 1.2] (-0.13,1)--(0.13,1);
	\draw[color ={black},line width = 1.2] (-0.13,2)--(0.13,2);
	\draw[color ={black},line width = 1.2] (-0.13,3)--(0.13,3);
	\draw[color ={black},line width = 1.2] (-0.13,4)--(0.13,4);
	\draw[color ={black},line width = 1.2] (-0.13,5)--(0.13,5);
	\draw[color ={black},line width = 1.2] (-0.13,6)--(0.13,6);
	\draw[color ={black},line width = 1.2] (-0.13,-1)--(0.13,-1);
	\draw[color ={black},line width = 1.2] (-0.13,-2)--(0.13,-2);
	\draw[color ={black},line width = 1.2] (-0.13,-3)--(0.13,-3);
	\draw[color ={black},line width = 1.2] (-0.13,-4)--(0.13,-4);
	\draw[color ={black},line width = 1.2] (-0.13,-5)--(0.13,-5);
	\draw[color ={black},line width = 1.2] (-0.13,-6)--(0.13,-6);
	\draw[color ={black},line width = 1.2] (-0.13,-7)--(0.13,-7);
	\draw[color ={black},line width = 1.2] (-0.13,-8)--(0.13,-8);
	\draw (2,-0.4) node[anchor = center] {\scriptsize \color{black}{$2$}};
	\draw (4,-0.4) node[anchor = center] {\scriptsize \color{black}{$4$}};
	\draw (-2,-0.4) node[anchor = center] {\scriptsize \color{black}{$-2$}};
	\draw (-4,-0.4) node[anchor = center] {\scriptsize \color{black}{$-4$}};
	\draw (-0.5,2.1) node[anchor = center] {\footnotesize \color{black}{$2$}};
	\draw (-0.5,4.1) node[anchor = center] {\footnotesize \color{black}{$4$}};
	\draw (-0.5,6.1) node[anchor = center] {\footnotesize \color{black}{$6$}};
	\draw (-0.5,-1.9) node[anchor = center] {\footnotesize \color{black}{$-2$}};
	\draw (-0.5,-3.9) node[anchor = center] {\footnotesize \color{black}{$-4$}};
	\draw (-0.5,-5.9) node[anchor = center] {\footnotesize \color{black}{$-6$}};
	\draw (-0.5,-7.9) node[anchor = center] {\footnotesize \color{black}{$-8$}};
	\draw [color={black}] (-0.3,-0.3) node[anchor = center,scale=1, rotate = 0] {O};
	\draw (3,4) node[anchor = center] {\footnotesize \color{blue}{$\Large \mathcal{C}_f$}};
	
	\draw[color={blue},line width = 2] (-3.4,-8.56)--(-3.2,-7.24)--(-3,-6)--(-2.8,-4.84)--(-2.6,-3.76)--(-2.4,-2.76)--(-2.2,-1.84)--(-2,-1)--(-1.8,-0.24)--(-1.6,0.44)--(-1.4,1.04)--(-1.2,1.56)--(-1,2)--(-0.8,2.36)--(-0.6,2.64)--(-0.4,2.84)--(-0.2,2.96)--(0,3)--(0.2,2.96)--(0.4,2.84)--(0.6,2.64)--(0.8,2.36)--(1,2)--(1.2,1.56)--(1.4,1.04)--(1.6,0.44)--(1.8,-0.24)--(2,-1)--(2.2,-1.84)--(2.4,-2.76)--(2.6,-3.76)--(2.8,-4.84)--(3,-6)--(3.2,-7.24)--(3.4,-8.56);

\end{tikzpicture}\\
    \end{minipage}
    \begin{minipage}{0.5\linewidth}
    \begin{tikzpicture}[baseline, scale=0.8]

    \tikzset{
      point/.style={
        thick,
        draw,
        cross out,
        inner sep=0pt,
        minimum width=5pt,
        minimum height=5pt,
      },
    }
    \clip (-5.5,-9.5) rectangle (6.5,7.5);
    	
	\draw[color ={black},>=latex,->] (-4,0)--(4,0);
	\draw[color ={black},>=latex,->] (0,-8)--(0,6);
	\draw[color ={black},opacity = 0.4] (-4,1)--(4,1);
	\draw[color ={black},opacity = 0.4] (-4,2)--(4,2);
	\draw[color ={black},opacity = 0.4] (-4,3)--(4,3);
	\draw[color ={black},opacity = 0.4] (-4,4)--(4,4);
	\draw[color ={black},opacity = 0.4] (-4,5)--(4,5);
	\draw[color ={black},opacity = 0.4] (-4,6)--(4,6);
	\draw[color ={black},opacity = 0.4] (-4,-1)--(4,-1);
	\draw[color ={black},opacity = 0.4] (-4,-2)--(4,-2);
	\draw[color ={black},opacity = 0.4] (-4,-3)--(4,-3);
	\draw[color ={black},opacity = 0.4] (-4,-4)--(4,-4);
	\draw[color ={black},opacity = 0.4] (-4,-5)--(4,-5);
	\draw[color ={black},opacity = 0.4] (-4,-6)--(4,-6);
	\draw[color ={black},opacity = 0.4] (-4,-7)--(4,-7);
	\draw[color ={black},opacity = 0.4] (-4,-8)--(4,-8);
	\draw[color ={black},opacity = 0.4] (1,-8)--(1,6);
	\draw[color ={black},opacity = 0.4] (2,-8)--(2,6);
	\draw[color ={black},opacity = 0.4] (3,-8)--(3,6);
	\draw[color ={black},opacity = 0.4] (4,-8)--(4,6);
	\draw[color ={black},opacity = 0.4] (-1,-8)--(-1,6);
	\draw[color ={black},opacity = 0.4] (-2,-8)--(-2,6);
	\draw[color ={black},opacity = 0.4] (-3,-8)--(-3,6);
	\draw[color ={black},opacity = 0.4] (-4,-8)--(-4,6);
	\draw[color ={gray},opacity = 0.3] (-4,0)--(4,0);
	\draw[color ={gray},opacity = 0.3] (-4,0.5)--(4,0.5);
	\draw[color ={gray},opacity = 0.3] (-4,1.5)--(4,1.5);
	\draw[color ={gray},opacity = 0.3] (-4,2.5)--(4,2.5);
	\draw[color ={gray},opacity = 0.3] (-4,3.5)--(4,3.5);
	\draw[color ={gray},opacity = 0.3] (-4,4.5)--(4,4.5);
	\draw[color ={gray},opacity = 0.3] (-4,5.5)--(4,5.5);
	\draw[color ={gray},opacity = 0.3] (-4,0)--(4,0);
	\draw[color ={gray},opacity = 0.3] (-4,-0.5)--(4,-0.5);
	\draw[color ={gray},opacity = 0.3] (-4,-1.5)--(4,-1.5);
	\draw[color ={gray},opacity = 0.3] (-4,-2.5)--(4,-2.5);
	\draw[color ={gray},opacity = 0.3] (-4,-3.5)--(4,-3.5);
	\draw[color ={gray},opacity = 0.3] (-4,-4.5)--(4,-4.5);
	\draw[color ={gray},opacity = 0.3] (-4,-5.5)--(4,-5.5);
	\draw[color ={gray},opacity = 0.3] (-4,-6.5)--(4,-6.5);
	\draw[color ={gray},opacity = 0.3] (-4,-7)--(4,-7);
	\draw[color ={gray},opacity = 0.3] (-4,-7.5)--(4,-7.5);
	\draw[color ={gray},opacity = 0.3] (-4,-8)--(4,-8);
	\draw[color ={gray},opacity = 0.3] (0,-8)--(0,6);
	\draw[color ={gray},opacity = 0.3] (0,-8)--(0,6);
	\draw[color ={black},line width = 1.2] (1,-0.13)--(1,0.13);
	\draw[color ={black},line width = 1.2] (2,-0.13)--(2,0.13);
	\draw[color ={black},line width = 1.2] (3,-0.13)--(3,0.13);
	\draw[color ={black},line width = 1.2] (4,-0.13)--(4,0.13);
	\draw[color ={black},line width = 1.2] (-1,-0.13)--(-1,0.13);
	\draw[color ={black},line width = 1.2] (-2,-0.13)--(-2,0.13);
	\draw[color ={black},line width = 1.2] (-3,-0.13)--(-3,0.13);
	\draw[color ={black},line width = 1.2] (-4,-0.13)--(-4,0.13);
	\draw[color ={black},line width = 1.2] (-0.13,1)--(0.13,1);
	\draw[color ={black},line width = 1.2] (-0.13,2)--(0.13,2);
	\draw[color ={black},line width = 1.2] (-0.13,3)--(0.13,3);
	\draw[color ={black},line width = 1.2] (-0.13,4)--(0.13,4);
	\draw[color ={black},line width = 1.2] (-0.13,5)--(0.13,5);
	\draw[color ={black},line width = 1.2] (-0.13,6)--(0.13,6);
	\draw[color ={black},line width = 1.2] (-0.13,-1)--(0.13,-1);
	\draw[color ={black},line width = 1.2] (-0.13,-2)--(0.13,-2);
	\draw[color ={black},line width = 1.2] (-0.13,-3)--(0.13,-3);
	\draw[color ={black},line width = 1.2] (-0.13,-4)--(0.13,-4);
	\draw[color ={black},line width = 1.2] (-0.13,-5)--(0.13,-5);
	\draw[color ={black},line width = 1.2] (-0.13,-6)--(0.13,-6);
	\draw[color ={black},line width = 1.2] (-0.13,-7)--(0.13,-7);
	\draw[color ={black},line width = 1.2] (-0.13,-8)--(0.13,-8);
	\draw (2,-0.4) node[anchor = center] {\scriptsize \color{black}{$2$}};
	\draw (4,-0.4) node[anchor = center] {\scriptsize \color{black}{$4$}};
	\draw (-2,-0.4) node[anchor = center] {\scriptsize \color{black}{$-2$}};
	\draw (-4,-0.4) node[anchor = center] {\scriptsize \color{black}{$-4$}};
	\draw (-0.5,2.1) node[anchor = center] {\footnotesize \color{black}{$2$}};
	\draw (-0.5,4.1) node[anchor = center] {\footnotesize \color{black}{$4$}};
	\draw (-0.5,-1.9) node[anchor = center] {\footnotesize \color{black}{$-2$}};
	\draw (-0.5,-3.9) node[anchor = center] {\footnotesize \color{black}{$-4$}};
	\draw (-0.5,-5.9) node[anchor = center] {\footnotesize \color{black}{$-6$}};
	\draw (-0.5,-7.9) node[anchor = center] {\footnotesize \color{black}{$-8$}};
	\draw [color={black}] (-0.3,-0.3) node[anchor = center,scale=1, rotate = 0] {O};
	\draw (3,4) node[anchor = center] {\footnotesize \color{red}{$\Large\mathcal{C}_f\prime$}};
	
	\draw[color={red},line width = 2] (-3.4,6.8)--(-3.2,6.4)--(-3,6)--(-2.8,5.6)--(-2.6,5.2)--(-2.4,4.8)--(-2.2,4.4)--(-2,4)--(-1.8,3.6)--(-1.6,3.2)--(-1.4,2.8)--(-1.2,2.4)--(-1,2)--(-0.8,1.6)--(-0.6,1.2)--(-0.4,0.8)--(-0.2,0.4)--(0,0)--(0.2,-0.4)--(0.4,-0.8)--(0.6,-1.2)--(0.8,-1.6)--(1,-2)--(1.2,-2.4)--(1.4,-2.8)--(1.6,-3.2)--(1.8,-3.6)--(2,-4)--(2.2,-4.4)--(2.4,-4.8)--(2.6,-5.2)--(2.8,-5.6)--(3,-6)--(3.2,-6.4)--(3.4,-6.8)--(3.6,-7.2)--(3.8,-7.6)--(4,-8);

\end{tikzpicture}\\
    \end{minipage}
    
\end{exercice}

\begin{Solution}
 L'équation réduite de la tangente au point d'abscisse $-1$ est  : $y=f'(-1)(x-(-1))+f(-1)$.\\
      On lit graphiquement $f(-1)=2$ et $f'(-1)=2$.\\
      L'équation réduite de la tangente est donc donnée par :
      $y=2(x+1)+2$, soit $y=2x+4$.
\end{Solution}

% @see : https://coopmaths.fr/alea?uuid=a1ba2&id=can1F14&n=1&d=10&alea=rkIp2232323232342345678910111213141516171819202122232425262728293031323334353637383940414243444546474823456789101112131415161718192021222324252627282930313233343536373839404142434445&cd=1&cols=1
\needspace{10\baselineskip}
\begin{exercice}[BaremeTotal=false]


Soit $f$ la fonction définie sur $\mathbb{R}$ par : 
$f(x)= 8x^2+9$.\\
En calculant la limite du taux d'accroissement, déterminer $f'(-2)$.
\end{exercice}

\begin{Solution}
 $f$ est une fonction polynôme du second degré de la forme $f(x)=ax^2+b$.\\
    La fonction dérivée est donnée par la somme des dérivées des fonctions $u$ et $v$ définies par $u(x)=8x^2$ et $v(x)=9$.\\
     Comme $u'(x)=16x$ et $v'(x)=0$, on obtient  $f'(x)=16x$. \\
     Ainsi, $f'(-2)=16\times (-2)=-32$.
\end{Solution}

% @see : https://coopmaths.fr/alea?uuid=3c690&id=can1F13&n=1&d=10&alea=EUDu2232323232342345678910111213141516171819202122232425262728293031323334353637383940414243444546474823456789101112131415161718192021222324252627282930313233343536373839404142434445&cd=1&cols=1
\needspace{10\baselineskip}
\begin{exercice}[BaremeTotal=false]


 Déterminer le coefficient directeur de la tangente à la courbe représentative de la fonction inverse au point d'abscisse $5$.
                 
\end{exercice}

\begin{Solution}
 Le coefficient directeur de la tangente au point d'abscisse $a$ est donné par le nombre dérivé $f'(a)$.\\
        La fonction $f$ définie par $f(x)=\dfrac{1}{x}$ a pour fonction dérivée la fonction $f'$ définie par $f'(x)=-\dfrac{1}{x^2}$.\\
Comme $f'(5)=-\dfrac{1}{5^2}=-\dfrac{1}{25}$, le coefficient directeur de la tangente au point d'abscisse $5$ est : $-\dfrac{1}{25}$.
\end{Solution}

% @see : https://coopmaths.fr/alea?uuid=ebd89&id=1AN14-1&n=1&d=10&s=3&alea=ocYr2232323232342345678910111213141516171819202122232425262728293031323334353637383940414243444546474823456789101112131415161718192021222324252627282930313233343536373839404142434445&cd=0&cols=1
\needspace{10\baselineskip}
\begin{exercice}[BaremeTotal=false]


 Donner la dérivée de la fonction $f$, dérivable pour tout $x\in\R$, définie par  $f(x)=\dfrac{3}{8}x+4$.
\end{exercice}

\begin{Solution}
 L'expression de la dérivée de la fonction $f$ définie par $f(x)=\dfrac{3}{8}x+4$ est : ${\color[HTML]{f15929}\boldsymbol{f'(x)=\dfrac{3}{8}}}$.
\end{Solution}

% @see : https://coopmaths.fr/alea?uuid=ebd89&id=1AN14-1&n=1&d=10&s=4&alea=uAlP2232323232342345678910111213141516171819202122232425262728293031323334353637383940414243444546474823456789101112131415161718192021222324252627282930313233343536373839404142434445&cd=0&cols=1
\needspace{10\baselineskip}
\begin{exercice}[BaremeTotal=false]


 Donner la dérivée de la fonction $f$, dérivable pour tout $x\in\R$, définie par  $f(x)=-2x^{8}$.
\end{exercice}

\begin{Solution}
 L'expression de la dérivée de la fonction $f$ définie par $f(x)=-2x^{8}$ est : ${\color[HTML]{f15929}\boldsymbol{f'(x)=-16x^{7}}}$.
\end{Solution}

% @see : https://coopmaths.fr/alea?uuid=ebd89&id=1AN14-1&n=1&d=10&s=7&alea=vjN72232323232342345678910111213141516171819202122232425262728293031323334353637383940414243444546474823456789101112131415161718192021222324252627282930313233343536373839404142434445&cd=0&cols=1
\needspace{10\baselineskip}
\begin{exercice}[BaremeTotal=false]


 Donner la dérivée de la fonction $f$, dérivable pour tout $x\in\R^*$, définie par  $f(x)=\dfrac{9}{10x}$.
\end{exercice}

\begin{Solution}
 L'expression de la dérivée de la fonction $f$ définie par $f(x)=\dfrac{9}{10x}$ est : ${\color[HTML]{f15929}\boldsymbol{f'(x)=-\dfrac{9}{10x^2}}}$.
\end{Solution}

% @see : https://coopmaths.fr/alea?uuid=a83c0&id=1AN14-4&n=1&d=10&s=1&alea=Njr32232323232342345678910111213141516171819202122232425262728293031323334353637383940414243444546474823456789101112131415161718192021222324252627282930313233343536373839404142434445&cd=0&cols=1
\needspace{10\baselineskip}
\begin{exercice}[BaremeTotal=false]


 Donner l'expression de la dérivée de la fonction $f$ définie sur $\R^{*}$ par $f(x)=\dfrac{4}{x}+5x^{2}-4x$.\\
\end{exercice}

\begin{Solution}
 L'expression de la dérivée de la fonction $f$ définie par $f(x)=\dfrac{4}{x}+5x^{2}-4x$ est :\\$f'(x)={\color[HTML]{f15929}\boldsymbol{-\dfrac{4}{x^2}+10x-4}}$.\\
\end{Solution}

% @see : https://coopmaths.fr/alea?uuid=a83c0&id=1AN14-4&n=1&d=10&s=4&alea=4Vpz2232323232342345678910111213141516171819202122232425262728293031323334353637383940414243444546474823456789101112131415161718192021222324252627282930313233343536373839404142434445&cd=1&cols=1
\needspace{10\baselineskip}
\begin{exercice}[BaremeTotal=false]


 Donner l'expression de la dérivée de la fonction $f$ définie sur $\R_+^{*}$ par $f(x)=-\sqrt{x}+3x+\dfrac{1}{x}$.\\
\end{exercice}

\begin{Solution}
 $(-\sqrt{x})^\prime=-\dfrac{1}{2\sqrt{x}}$\\$(3x)^\prime=3$\\$(\dfrac{1}{x})^\prime=-\dfrac{1}{x^2}$\\L'expression de la dérivée de la fonction $f$ définie par $f(x)=-\sqrt{x}+3x+\dfrac{1}{x}$ est :\\$f'(x)={\color[HTML]{f15929}\boldsymbol{-\dfrac{1}{2\sqrt{x}}+3-\dfrac{1}{x^2}}}$.\\
\end{Solution}

\end{Maquette}
\clearpage
\begin{Maquette}[DM]{Niveau= ,Classe=\getdata[46]\mydata\ (v46), Date = 06/01/25  ,Theme=Exercices}


% @see : https://coopmaths.fr/alea?uuid=4c8c7&id=1AN11-3&n=1&d=10&s=2&alea=nawO223232323234234567891011121314151617181920212223242526272829303132333435363738394041424344454647482345678910111213141516171819202122232425262728293031323334353637383940414243444546&cd=1&cols=2
\needspace{10\baselineskip}
\begin{exercice}[BaremeTotal=false]
\label{eleve:46}


 \begin{minipage}[t]{\linewidth} Soit $f$ une fonction dérivable sur $[-5;5]$ et $\mathcal{C}_f$ sa courbe représentative.\\On sait que  $f(-1)=-4~~$ et que $~~f'(-1)=-2$.\\Déterminer une équation de la tangente $(T)$ à la courbe $\mathcal{C}_f$ au point d'abscisse $-1$,\\sans utiliser la formule de cours de l'équation de tangente. \end{minipage}

\end{exercice}

\begin{Solution}
 \begin{minipage}[t]{\linewidth} On sait que la tangente n'est pas une droite verticale, puisque la fonction est dérivable sur l'intervalle.\\On en déduit que la tangente $(T)$ au point d'abscisse $-1$, admet une équation réduite de la forme :  $(T) : y=m x + p$. 

\medskip
$\bullet$ Détermination de $m$ :\\On sait que le nombre dérivé en $-1$ est par définition, le coefficient directeur de la tangente au point d'abscisse $-1$.\\Par conséquent, on a déjà : $m=f'(-1)=-2$.\\ On en déduit que  $(T) : y= -2 x + p$\\$\bullet$ Détermination de $p$ :\\ Pour cela, on utilise que si $f(-1)=-4~~$, alors le point $A$ de coordonnées $(-1;-4)$ appartient à $\mathcal{C}_f$ mais aussi à $(T)$.\\ On peut écrire $A(-1;-4) = \mathcal{C}_f \cap (T)$.\\ On remplace alors les coordonnées de $A(-1;-4)$ dans l'équation  $(T) : y= -2 x + p$.\\ $\begin{aligned} \phantom{\iff}&A(-1;-4)\in (T)\\ \iff& -4= -2 \times (-1)  + p\\ \iff& p=-4 +2 \times (-1)  \\ \iff& p=-4 -2   \\ \iff& p=-6\\\end{aligned}$\\On peut conclure que : $(T) : {\color[HTML]{f15929}\boldsymbol{y=-2x-6}}$. \end{minipage}

\end{Solution}

% @see : https://coopmaths.fr/alea?uuid=0e984&id=can1F15&n=1&d=10&alea=1Alo223232323234234567891011121314151617181920212223242526272829303132333435363738394041424344454647482345678910111213141516171819202122232425262728293031323334353637383940414243444546&cd=1&cols=1
\needspace{10\baselineskip}
\begin{exercice}[BaremeTotal=false]


 La courbe représente une fonction $f$ et la droite est la tangente au point d'abscisse $0$.\\
        Déterminer $f'(0)$.\begin{center}\begin{tikzpicture}[baseline,scale = 0.5]

    \tikzset{
      point/.style={
        thick,
        draw,
        cross out,
        inner sep=0pt,
        minimum width=5pt,
        minimum height=5pt,
      },
    }
    \clip (-4,-2) rectangle (12,12);
    	
	\draw[color ={black},line width = 2,>=latex,->] (-4,0)--(12,0);
	\draw[color ={black},line width = 2,>=latex,->] (0,-2)--(0,12);
	\draw[color ={black},opacity = 0.5] (-4,1)--(12,1);
	\draw[color ={black},opacity = 0.5] (-4,2)--(12,2);
	\draw[color ={black},opacity = 0.5] (-4,3)--(12,3);
	\draw[color ={black},opacity = 0.5] (-4,4)--(12,4);
	\draw[color ={black},opacity = 0.5] (-4,5)--(12,5);
	\draw[color ={black},opacity = 0.5] (-4,6)--(12,6);
	\draw[color ={black},opacity = 0.5] (-4,7)--(12,7);
	\draw[color ={black},opacity = 0.5] (-4,8)--(12,8);
	\draw[color ={black},opacity = 0.5] (-4,9)--(12,9);
	\draw[color ={black},opacity = 0.5] (-4,10)--(12,10);
	\draw[color ={black},opacity = 0.5] (-4,11)--(12,11);
	\draw[color ={black},opacity = 0.5] (-4,12)--(12,12);
	\draw[color ={black},opacity = 0.5] (-4,-1)--(12,-1);
	\draw[color ={black},opacity = 0.5] (-4,-2)--(12,-2);
	\draw[color ={black},opacity = 0.5] (1,-2)--(1,12);
	\draw[color ={black},opacity = 0.5] (2,-2)--(2,12);
	\draw[color ={black},opacity = 0.5] (3,-2)--(3,12);
	\draw[color ={black},opacity = 0.5] (4,-2)--(4,12);
	\draw[color ={black},opacity = 0.5] (5,-2)--(5,12);
	\draw[color ={black},opacity = 0.5] (6,-2)--(6,12);
	\draw[color ={black},opacity = 0.5] (7,-2)--(7,12);
	\draw[color ={black},opacity = 0.5] (8,-2)--(8,12);
	\draw[color ={black},opacity = 0.5] (9,-2)--(9,12);
	\draw[color ={black},opacity = 0.5] (10,-2)--(10,12);
	\draw[color ={black},opacity = 0.5] (11,-2)--(11,12);
	\draw[color ={black},opacity = 0.5] (12,-2)--(12,12);
	\draw[color ={black},opacity = 0.5] (-1,-2)--(-1,12);
	\draw[color ={black},opacity = 0.5] (-2,-2)--(-2,12);
	\draw[color ={black},opacity = 0.5] (-3,-2)--(-3,12);
	\draw[color ={black},opacity = 0.5] (-4,-2)--(-4,12);
	\draw[color ={gray},opacity = 0.3] (-4,0)--(12,0);
	\draw[color ={gray},opacity = 0.3] (-4,0)--(12,0);
	\draw[color ={gray},opacity = 0.3] (0,-2)--(0,12);
	\draw[color ={gray},opacity = 0.3] (0,-2)--(0,12);
	\draw[color ={black},line width = 2] (2,-0.2)--(2,0.2);
	\draw[color ={black},line width = 2] (4,-0.2)--(4,0.2);
	\draw[color ={black},line width = 2] (6,-0.2)--(6,0.2);
	\draw[color ={black},line width = 2] (8,-0.2)--(8,0.2);
	\draw[color ={black},line width = 2] (10,-0.2)--(10,0.2);
	\draw[color ={black},line width = 2] (-2,-0.2)--(-2,0.2);
	\draw[color ={black},line width = 2] (-0.2,2)--(0.2,2);
	\draw[color ={black},line width = 2] (-0.2,4)--(0.2,4);
	\draw[color ={black},line width = 2] (-0.2,6)--(0.2,6);
	\draw[color ={black},line width = 2] (-0.2,8)--(0.2,8);
	\draw[color ={black},line width = 2] (-0.2,10)--(0.2,10);
	\draw (2,-0.4) node[anchor = center] {\scriptsize \color{black}{$1$}};
	\draw (4,-0.4) node[anchor = center] {\scriptsize \color{black}{$2$}};
	\draw (6,-0.4) node[anchor = center] {\scriptsize \color{black}{$3$}};
	\draw (8,-0.4) node[anchor = center] {\scriptsize \color{black}{$4$}};
	\draw (10,-0.4) node[anchor = center] {\scriptsize \color{black}{$5$}};
	\draw (-0.5,2.1) node[anchor = center] {\footnotesize \color{black}{$1$}};
	\draw (-0.5,4.1) node[anchor = center] {\footnotesize \color{black}{$2$}};
	\draw (-0.5,6.1) node[anchor = center] {\footnotesize \color{black}{$3$}};
	\draw (-0.5,8.1) node[anchor = center] {\footnotesize \color{black}{$4$}};
	\draw (-0.5,10.1) node[anchor = center] {\footnotesize \color{black}{$5$}};
	\draw [color={black}] (-0.3,-0.3) node[anchor = center,scale=1, rotate = 0] {O};
	
	\draw[color={blue},line width = 2] (-4,0.1)--(-3.6,0.13)--(-3.2,0.18)--(-2.8,0.24)--(-2.4,0.33)--(-2,0.45)--(-1.6,0.6)--(-1.2,0.81)--(-0.8,1.1)--(-0.4,1.48)--(0,2)--(0.4,2.7)--(0.8,3.64)--(1.2,4.92)--(1.6,6.64)--(2,8.96)--(2.4,12.1);
	
	\draw[color={red},line width = 2] (-3.6,-3.4)--(-3.2,-2.8)--(-2.8,-2.2)--(-2.4,-1.6)--(-2,-1)--(-1.6,-0.4)--(-1.2,0.2)--(-0.8,0.8)--(-0.4,1.4)--(0,2)--(0.4,2.6)--(0.8,3.2)--(1.2,3.8)--(1.6,4.4)--(2,5)--(2.4,5.6)--(2.8,6.2)--(3.2,6.8)--(3.6,7.4)--(4,8)--(4.4,8.6)--(4.8,9.2)--(5.2,9.8)--(5.6,10.4)--(6,11)--(6.4,11.6)--(6.8,12.2)--(7.2,12.8)--(7.6,13.4);

\end{tikzpicture}\end{center}
\end{exercice}

\begin{Solution}
 $f'(0)$ est donné par le coefficient directeur de la tangente à la courbe au point d'abscisse $0$, soit $\dfrac{3}{2}$.
\end{Solution}

% @see : https://coopmaths.fr/alea?uuid=6f32d&id=can1F16&n=1&d=10&alea=dqSw223232323234234567891011121314151617181920212223242526272829303132333435363738394041424344454647482345678910111213141516171819202122232425262728293031323334353637383940414243444546&cd=1&cols=1
\needspace{10\baselineskip}
\newpage
\begin{exercice}[BaremeTotal=false]


 On donne les représentations graphiques d'une fonction et de sa dérivée.\\
      Donner l'équation réduite de la tangente à la courbe de $f$ en $x=-2$. \\ \begin{minipage}{0.5\linewidth}
    \begin{tikzpicture}[baseline, scale=0.8]

    \tikzset{
      point/.style={
        thick,
        draw,
        cross out,
        inner sep=0pt,
        minimum width=5pt,
        minimum height=5pt,
      },
    }
    \clip (-5.5,-9.5) rectangle (5.75,7.5);
    	
	\draw[color ={black},>=latex,->] (-4,0)--(4,0);
	\draw[color ={black},>=latex,->] (0,-8)--(0,6);
	\draw[color ={black},opacity = 0.4] (-4,1)--(4,1);
	\draw[color ={black},opacity = 0.4] (-4,2)--(4,2);
	\draw[color ={black},opacity = 0.4] (-4,3)--(4,3);
	\draw[color ={black},opacity = 0.4] (-4,4)--(4,4);
	\draw[color ={black},opacity = 0.4] (-4,5)--(4,5);
	\draw[color ={black},opacity = 0.4] (-4,6)--(4,6);
	\draw[color ={black},opacity = 0.4] (-4,-1)--(4,-1);
	\draw[color ={black},opacity = 0.4] (-4,-2)--(4,-2);
	\draw[color ={black},opacity = 0.4] (-4,-3)--(4,-3);
	\draw[color ={black},opacity = 0.4] (-4,-4)--(4,-4);
	\draw[color ={black},opacity = 0.4] (-4,-5)--(4,-5);
	\draw[color ={black},opacity = 0.4] (-4,-6)--(4,-6);
	\draw[color ={black},opacity = 0.4] (-4,-7)--(4,-7);
	\draw[color ={black},opacity = 0.4] (-4,-8)--(4,-8);
	\draw[color ={black},opacity = 0.4] (1,-8)--(1,6);
	\draw[color ={black},opacity = 0.4] (2,-8)--(2,6);
	\draw[color ={black},opacity = 0.4] (3,-8)--(3,6);
	\draw[color ={black},opacity = 0.4] (4,-8)--(4,6);
	\draw[color ={black},opacity = 0.4] (-1,-8)--(-1,6);
	\draw[color ={black},opacity = 0.4] (-2,-8)--(-2,6);
	\draw[color ={black},opacity = 0.4] (-3,-8)--(-3,6);
	\draw[color ={black},opacity = 0.4] (-4,-8)--(-4,6);
	\draw[color ={gray},opacity = 0.3] (-4,0)--(4,0);
	\draw[color ={gray},opacity = 0.3] (-4,0.5)--(4,0.5);
	\draw[color ={gray},opacity = 0.3] (-4,1.5)--(4,1.5);
	\draw[color ={gray},opacity = 0.3] (-4,2.5)--(4,2.5);
	\draw[color ={gray},opacity = 0.3] (-4,3.5)--(4,3.5);
	\draw[color ={gray},opacity = 0.3] (-4,4.5)--(4,4.5);
	\draw[color ={gray},opacity = 0.3] (-4,5.5)--(4,5.5);
	\draw[color ={gray},opacity = 0.3] (-4,0)--(4,0);
	\draw[color ={gray},opacity = 0.3] (-4,-0.5)--(4,-0.5);
	\draw[color ={gray},opacity = 0.3] (-4,-1.5)--(4,-1.5);
	\draw[color ={gray},opacity = 0.3] (-4,-2.5)--(4,-2.5);
	\draw[color ={gray},opacity = 0.3] (-4,-3.5)--(4,-3.5);
	\draw[color ={gray},opacity = 0.3] (-4,-4.5)--(4,-4.5);
	\draw[color ={gray},opacity = 0.3] (-4,-5.5)--(4,-5.5);
	\draw[color ={gray},opacity = 0.3] (-4,-6.5)--(4,-6.5);
	\draw[color ={gray},opacity = 0.3] (-4,-7)--(4,-7);
	\draw[color ={gray},opacity = 0.3] (-4,-7.5)--(4,-7.5);
	\draw[color ={gray},opacity = 0.3] (-4,-8)--(4,-8);
	\draw[color ={gray},opacity = 0.3] (0,-8)--(0,6);
	\draw[color ={gray},opacity = 0.3] (0,-8)--(0,6);
	\draw[color ={black},line width = 1.2] (1,-0.13)--(1,0.13);
	\draw[color ={black},line width = 1.2] (2,-0.13)--(2,0.13);
	\draw[color ={black},line width = 1.2] (3,-0.13)--(3,0.13);
	\draw[color ={black},line width = 1.2] (4,-0.13)--(4,0.13);
	\draw[color ={black},line width = 1.2] (-1,-0.13)--(-1,0.13);
	\draw[color ={black},line width = 1.2] (-2,-0.13)--(-2,0.13);
	\draw[color ={black},line width = 1.2] (-3,-0.13)--(-3,0.13);
	\draw[color ={black},line width = 1.2] (-4,-0.13)--(-4,0.13);
	\draw[color ={black},line width = 1.2] (-0.13,1)--(0.13,1);
	\draw[color ={black},line width = 1.2] (-0.13,2)--(0.13,2);
	\draw[color ={black},line width = 1.2] (-0.13,3)--(0.13,3);
	\draw[color ={black},line width = 1.2] (-0.13,4)--(0.13,4);
	\draw[color ={black},line width = 1.2] (-0.13,5)--(0.13,5);
	\draw[color ={black},line width = 1.2] (-0.13,6)--(0.13,6);
	\draw[color ={black},line width = 1.2] (-0.13,-1)--(0.13,-1);
	\draw[color ={black},line width = 1.2] (-0.13,-2)--(0.13,-2);
	\draw[color ={black},line width = 1.2] (-0.13,-3)--(0.13,-3);
	\draw[color ={black},line width = 1.2] (-0.13,-4)--(0.13,-4);
	\draw[color ={black},line width = 1.2] (-0.13,-5)--(0.13,-5);
	\draw[color ={black},line width = 1.2] (-0.13,-6)--(0.13,-6);
	\draw[color ={black},line width = 1.2] (-0.13,-7)--(0.13,-7);
	\draw[color ={black},line width = 1.2] (-0.13,-8)--(0.13,-8);
	\draw (2,-0.4) node[anchor = center] {\scriptsize \color{black}{$2$}};
	\draw (4,-0.4) node[anchor = center] {\scriptsize \color{black}{$4$}};
	\draw (-2,-0.4) node[anchor = center] {\scriptsize \color{black}{$-2$}};
	\draw (-4,-0.4) node[anchor = center] {\scriptsize \color{black}{$-4$}};
	\draw (-0.5,2.1) node[anchor = center] {\footnotesize \color{black}{$2$}};
	\draw (-0.5,4.1) node[anchor = center] {\footnotesize \color{black}{$4$}};
	\draw (-0.5,6.1) node[anchor = center] {\footnotesize \color{black}{$6$}};
	\draw (-0.5,-1.9) node[anchor = center] {\footnotesize \color{black}{$-2$}};
	\draw (-0.5,-3.9) node[anchor = center] {\footnotesize \color{black}{$-4$}};
	\draw (-0.5,-5.9) node[anchor = center] {\footnotesize \color{black}{$-6$}};
	\draw (-0.5,-7.9) node[anchor = center] {\footnotesize \color{black}{$-8$}};
	\draw [color={black}] (-0.3,-0.3) node[anchor = center,scale=1, rotate = 0] {O};
	\draw (3,4) node[anchor = center] {\footnotesize \color{blue}{$\Large \mathcal{C}_f$}};
	
	\draw[color={blue},line width = 2] (-4,-6)--(-3.8,-4.84)--(-3.6,-3.76)--(-3.4,-2.76)--(-3.2,-1.84)--(-3,-1)--(-2.8,-0.24)--(-2.6,0.44)--(-2.4,1.04)--(-2.2,1.56)--(-2,2)--(-1.8,2.36)--(-1.6,2.64)--(-1.4,2.84)--(-1.2,2.96)--(-1,3)--(-0.8,2.96)--(-0.6,2.84)--(-0.4,2.64)--(-0.2,2.36)--(0,2)--(0.2,1.56)--(0.4,1.04)--(0.6,0.44)--(0.8,-0.24)--(1,-1)--(1.2,-1.84)--(1.4,-2.76)--(1.6,-3.76)--(1.8,-4.84)--(2,-6)--(2.2,-7.24)--(2.4,-8.56);

\end{tikzpicture}\\
    \end{minipage}
    \begin{minipage}{0.5\linewidth}
    \begin{tikzpicture}[baseline, scale=0.8]

    \tikzset{
      point/.style={
        thick,
        draw,
        cross out,
        inner sep=0pt,
        minimum width=5pt,
        minimum height=5pt,
      },
    }
    \clip (-5.5,-9.5) rectangle (6.5,7.5);
    	
	\draw[color ={black},>=latex,->] (-4,0)--(4,0);
	\draw[color ={black},>=latex,->] (0,-8)--(0,6);
	\draw[color ={black},opacity = 0.4] (-4,1)--(4,1);
	\draw[color ={black},opacity = 0.4] (-4,2)--(4,2);
	\draw[color ={black},opacity = 0.4] (-4,3)--(4,3);
	\draw[color ={black},opacity = 0.4] (-4,4)--(4,4);
	\draw[color ={black},opacity = 0.4] (-4,5)--(4,5);
	\draw[color ={black},opacity = 0.4] (-4,6)--(4,6);
	\draw[color ={black},opacity = 0.4] (-4,-1)--(4,-1);
	\draw[color ={black},opacity = 0.4] (-4,-2)--(4,-2);
	\draw[color ={black},opacity = 0.4] (-4,-3)--(4,-3);
	\draw[color ={black},opacity = 0.4] (-4,-4)--(4,-4);
	\draw[color ={black},opacity = 0.4] (-4,-5)--(4,-5);
	\draw[color ={black},opacity = 0.4] (-4,-6)--(4,-6);
	\draw[color ={black},opacity = 0.4] (-4,-7)--(4,-7);
	\draw[color ={black},opacity = 0.4] (-4,-8)--(4,-8);
	\draw[color ={black},opacity = 0.4] (1,-8)--(1,6);
	\draw[color ={black},opacity = 0.4] (2,-8)--(2,6);
	\draw[color ={black},opacity = 0.4] (3,-8)--(3,6);
	\draw[color ={black},opacity = 0.4] (4,-8)--(4,6);
	\draw[color ={black},opacity = 0.4] (-1,-8)--(-1,6);
	\draw[color ={black},opacity = 0.4] (-2,-8)--(-2,6);
	\draw[color ={black},opacity = 0.4] (-3,-8)--(-3,6);
	\draw[color ={black},opacity = 0.4] (-4,-8)--(-4,6);
	\draw[color ={gray},opacity = 0.3] (-4,0)--(4,0);
	\draw[color ={gray},opacity = 0.3] (-4,0.5)--(4,0.5);
	\draw[color ={gray},opacity = 0.3] (-4,1.5)--(4,1.5);
	\draw[color ={gray},opacity = 0.3] (-4,2.5)--(4,2.5);
	\draw[color ={gray},opacity = 0.3] (-4,3.5)--(4,3.5);
	\draw[color ={gray},opacity = 0.3] (-4,4.5)--(4,4.5);
	\draw[color ={gray},opacity = 0.3] (-4,5.5)--(4,5.5);
	\draw[color ={gray},opacity = 0.3] (-4,0)--(4,0);
	\draw[color ={gray},opacity = 0.3] (-4,-0.5)--(4,-0.5);
	\draw[color ={gray},opacity = 0.3] (-4,-1.5)--(4,-1.5);
	\draw[color ={gray},opacity = 0.3] (-4,-2.5)--(4,-2.5);
	\draw[color ={gray},opacity = 0.3] (-4,-3.5)--(4,-3.5);
	\draw[color ={gray},opacity = 0.3] (-4,-4.5)--(4,-4.5);
	\draw[color ={gray},opacity = 0.3] (-4,-5.5)--(4,-5.5);
	\draw[color ={gray},opacity = 0.3] (-4,-6.5)--(4,-6.5);
	\draw[color ={gray},opacity = 0.3] (-4,-7)--(4,-7);
	\draw[color ={gray},opacity = 0.3] (-4,-7.5)--(4,-7.5);
	\draw[color ={gray},opacity = 0.3] (-4,-8)--(4,-8);
	\draw[color ={gray},opacity = 0.3] (0,-8)--(0,6);
	\draw[color ={gray},opacity = 0.3] (0,-8)--(0,6);
	\draw[color ={black},line width = 1.2] (1,-0.13)--(1,0.13);
	\draw[color ={black},line width = 1.2] (2,-0.13)--(2,0.13);
	\draw[color ={black},line width = 1.2] (3,-0.13)--(3,0.13);
	\draw[color ={black},line width = 1.2] (4,-0.13)--(4,0.13);
	\draw[color ={black},line width = 1.2] (-1,-0.13)--(-1,0.13);
	\draw[color ={black},line width = 1.2] (-2,-0.13)--(-2,0.13);
	\draw[color ={black},line width = 1.2] (-3,-0.13)--(-3,0.13);
	\draw[color ={black},line width = 1.2] (-4,-0.13)--(-4,0.13);
	\draw[color ={black},line width = 1.2] (-0.13,1)--(0.13,1);
	\draw[color ={black},line width = 1.2] (-0.13,2)--(0.13,2);
	\draw[color ={black},line width = 1.2] (-0.13,3)--(0.13,3);
	\draw[color ={black},line width = 1.2] (-0.13,4)--(0.13,4);
	\draw[color ={black},line width = 1.2] (-0.13,5)--(0.13,5);
	\draw[color ={black},line width = 1.2] (-0.13,6)--(0.13,6);
	\draw[color ={black},line width = 1.2] (-0.13,-1)--(0.13,-1);
	\draw[color ={black},line width = 1.2] (-0.13,-2)--(0.13,-2);
	\draw[color ={black},line width = 1.2] (-0.13,-3)--(0.13,-3);
	\draw[color ={black},line width = 1.2] (-0.13,-4)--(0.13,-4);
	\draw[color ={black},line width = 1.2] (-0.13,-5)--(0.13,-5);
	\draw[color ={black},line width = 1.2] (-0.13,-6)--(0.13,-6);
	\draw[color ={black},line width = 1.2] (-0.13,-7)--(0.13,-7);
	\draw[color ={black},line width = 1.2] (-0.13,-8)--(0.13,-8);
	\draw (2,-0.4) node[anchor = center] {\scriptsize \color{black}{$2$}};
	\draw (4,-0.4) node[anchor = center] {\scriptsize \color{black}{$4$}};
	\draw (-2,-0.4) node[anchor = center] {\scriptsize \color{black}{$-2$}};
	\draw (-4,-0.4) node[anchor = center] {\scriptsize \color{black}{$-4$}};
	\draw (-0.5,2.1) node[anchor = center] {\footnotesize \color{black}{$2$}};
	\draw (-0.5,4.1) node[anchor = center] {\footnotesize \color{black}{$4$}};
	\draw (-0.5,-1.9) node[anchor = center] {\footnotesize \color{black}{$-2$}};
	\draw (-0.5,-3.9) node[anchor = center] {\footnotesize \color{black}{$-4$}};
	\draw (-0.5,-5.9) node[anchor = center] {\footnotesize \color{black}{$-6$}};
	\draw (-0.5,-7.9) node[anchor = center] {\footnotesize \color{black}{$-8$}};
	\draw [color={black}] (-0.3,-0.3) node[anchor = center,scale=1, rotate = 0] {O};
	\draw (3,4) node[anchor = center] {\footnotesize \color{red}{$\Large\mathcal{C}_f\prime$}};
	
	\draw[color={red},line width = 2] (-4,6)--(-3.8,5.6)--(-3.6,5.2)--(-3.4,4.8)--(-3.2,4.4)--(-3,4)--(-2.8,3.6)--(-2.6,3.2)--(-2.4,2.8)--(-2.2,2.4)--(-2,2)--(-1.8,1.6)--(-1.6,1.2)--(-1.4,0.8)--(-1.2,0.4)--(-1,0)--(-0.8,-0.4)--(-0.6,-0.8)--(-0.4,-1.2)--(-0.2,-1.6)--(0,-2)--(0.2,-2.4)--(0.4,-2.8)--(0.6,-3.2)--(0.8,-3.6)--(1,-4)--(1.2,-4.4)--(1.4,-4.8)--(1.6,-5.2)--(1.8,-5.6)--(2,-6)--(2.2,-6.4)--(2.4,-6.8)--(2.6,-7.2)--(2.8,-7.6)--(3,-8)--(3.2,-8.4)--(3.4,-8.8);

\end{tikzpicture}\\
    \end{minipage}
    
\end{exercice}

\begin{Solution}
 L'équation réduite de la tangente au point d'abscisse $-2$ est  : $y=f'(-2)(x-(-2))+f(-2)$.\\
      On lit graphiquement $f(-2)=2$ et $f'(-2)=2$.\\
      L'équation réduite de la tangente est donc donnée par :
      $y=2(x+2)+2$, soit $y=2x+6$.
\end{Solution}

% @see : https://coopmaths.fr/alea?uuid=a1ba2&id=can1F14&n=1&d=10&alea=rkIp223232323234234567891011121314151617181920212223242526272829303132333435363738394041424344454647482345678910111213141516171819202122232425262728293031323334353637383940414243444546&cd=1&cols=1
\needspace{10\baselineskip}
\begin{exercice}[BaremeTotal=false]


 Soit $f$ la fonction définie sur $\mathbb{R}$ par : $f(x)= -4x^2+5x-3$.\\
        En calculant la limite du taux d'accroissement, déterminer $f'(2)$.
\end{exercice}

\begin{Solution}
 $f$ est une fonction polynôme du second degré de la forme $f(x)=ax^2+bx+c$.\\
    La fonction dérivée est donnée par la somme des dérivées des fonctions $u$ et $v$ définies par $u(x)=-4x^2$ et $v(x)=5x-3$.\\
     Comme $u'(x)=-8x$ et $v'(x)=5$, on obtient  $f'(x)=-8x+5$. \\
     Ainsi, $f'(2)=-8\times 2+5=-11$.
\end{Solution}

% @see : https://coopmaths.fr/alea?uuid=3c690&id=can1F13&n=1&d=10&alea=EUDu223232323234234567891011121314151617181920212223242526272829303132333435363738394041424344454647482345678910111213141516171819202122232425262728293031323334353637383940414243444546&cd=1&cols=1
\needspace{10\baselineskip}
\begin{exercice}[BaremeTotal=false]


 Déterminer le coefficient directeur de la tangente à la courbe représentative de la fonction racine carrée au point d'abscisse $20$.
         
\end{exercice}

\begin{Solution}
 Le coefficient directeur de la tangente au point d'abscisse $a$ est donné par le nombre dérivé $f'(a)$.\\
        La fonction $f$ définie par $f(x)=\sqrt{x}$ a pour fonction dérivée la fonction $f'$ définie par $f'(x)=\dfrac{1}{2\sqrt{x}}$.\\Comme $f'(20)=\dfrac{1}{2\sqrt{20}}$, le coefficient directeur de la tangente au point d'abscisse $20$ est : $\dfrac{1}{2\sqrt{20}}$.
\end{Solution}

% @see : https://coopmaths.fr/alea?uuid=ebd89&id=1AN14-1&n=1&d=10&s=3&alea=ocYr223232323234234567891011121314151617181920212223242526272829303132333435363738394041424344454647482345678910111213141516171819202122232425262728293031323334353637383940414243444546&cd=0&cols=1
\needspace{10\baselineskip}
\begin{exercice}[BaremeTotal=false]


 Donner la dérivée de la fonction $f$, dérivable pour tout $x\in\R$, définie par  $f(x)=-\dfrac{7}{10}x+7$.
\end{exercice}

\begin{Solution}
 L'expression de la dérivée de la fonction $f$ définie par $f(x)=-\dfrac{7}{10}x+7$ est : ${\color[HTML]{f15929}\boldsymbol{f'(x)=-\dfrac{7}{10}}}$.
\end{Solution}

% @see : https://coopmaths.fr/alea?uuid=ebd89&id=1AN14-1&n=1&d=10&s=4&alea=uAlP223232323234234567891011121314151617181920212223242526272829303132333435363738394041424344454647482345678910111213141516171819202122232425262728293031323334353637383940414243444546&cd=0&cols=1
\needspace{10\baselineskip}
\begin{exercice}[BaremeTotal=false]


 Donner la dérivée de la fonction $f$, dérivable pour tout $x\in\R$, définie par  $f(x)=4x^{10}$.
\end{exercice}

\begin{Solution}
 L'expression de la dérivée de la fonction $f$ définie par $f(x)=4x^{10}$ est : ${\color[HTML]{f15929}\boldsymbol{f'(x)=40x^{9}}}$.
\end{Solution}

% @see : https://coopmaths.fr/alea?uuid=ebd89&id=1AN14-1&n=1&d=10&s=7&alea=vjN7223232323234234567891011121314151617181920212223242526272829303132333435363738394041424344454647482345678910111213141516171819202122232425262728293031323334353637383940414243444546&cd=0&cols=1
\needspace{10\baselineskip}
\begin{exercice}[BaremeTotal=false]


 Donner la dérivée de la fonction $f$, dérivable pour tout $x\in\R^*$, définie par  $f(x)=\dfrac{-7}{8x}$.
\end{exercice}

\begin{Solution}
 L'expression de la dérivée de la fonction $f$ définie par $f(x)=\dfrac{-7}{8x}$ est : ${\color[HTML]{f15929}\boldsymbol{f'(x)=\dfrac{7}{8x^2}}}$.
\end{Solution}

% @see : https://coopmaths.fr/alea?uuid=a83c0&id=1AN14-4&n=1&d=10&s=1&alea=Njr3223232323234234567891011121314151617181920212223242526272829303132333435363738394041424344454647482345678910111213141516171819202122232425262728293031323334353637383940414243444546&cd=0&cols=1
\needspace{10\baselineskip}
\begin{exercice}[BaremeTotal=false]


 Donner l'expression de la dérivée de la fonction $f$ définie sur $\R^{*}$ par $f(x)=-\dfrac{6}{x}+4x^{2}-4x$.\\
\end{exercice}

\begin{Solution}
 L'expression de la dérivée de la fonction $f$ définie par $f(x)=-\dfrac{6}{x}+4x^{2}-4x$ est :\\$f'(x)={\color[HTML]{f15929}\boldsymbol{\dfrac{6}{x^2}+8x-4}}$.\\
\end{Solution}

% @see : https://coopmaths.fr/alea?uuid=a83c0&id=1AN14-4&n=1&d=10&s=4&alea=4Vpz223232323234234567891011121314151617181920212223242526272829303132333435363738394041424344454647482345678910111213141516171819202122232425262728293031323334353637383940414243444546&cd=1&cols=1
\needspace{10\baselineskip}
\begin{exercice}[BaremeTotal=false]


 Donner l'expression de la dérivée de la fonction $f$ définie sur $\R_+^{*}$ par $f(x)=-2\sqrt{x}-4x^{2}+4x-4+\dfrac{2}{x}$.\\
\end{exercice}

\begin{Solution}
 $(-2\sqrt{x})^\prime=-\dfrac{2}{2\sqrt{x}}$\\$(-4x^{2}+4x-4)^\prime=-8x+4$\\$(\dfrac{2}{x})^\prime=-\dfrac{2}{x^2}$\\L'expression de la dérivée de la fonction $f$ définie par $f(x)=-2\sqrt{x}-4x^{2}+4x-4+\dfrac{2}{x}$ est :\\$f'(x)={\color[HTML]{f15929}\boldsymbol{-\dfrac{2}{2\sqrt{x}}-8x+4-\dfrac{2}{x^2}}}$.\\
\end{Solution}

\end{Maquette}
\clearpage
\end{document}